\documentclass[12pt]{exam}
\usepackage[top=2cm, bottom=3cm, left=2cm, right=2cm]{geometry}
\usepackage{graphicx}
\usepackage{amsmath,amsfonts,amssymb,amsthm}
\usepackage{yhmath} % \wideparen
\usepackage[colorlinks=true, linkcolor=brown, urlcolor=brown]{hyperref}
\usepackage{xeCJK}
\usepackage{xcolor}
\usepackage{float}
\usepackage{enumitem}
\usepackage{tikz}
\usetikzlibrary{calc}
\usepackage{mathtools}
\usepackage{lipsum}
\usepackage{pgfornament}
\usepackage{multicol}
\usepackage{worldflags}
\usetikzlibrary{shapes.geometric, positioning}
\setCJKmainfont{KaiTi}

%\includeonly{微积分}
%\includeonly{代数,组合数学,微积分}
%\includeonly{几何}
%\includeonly{代数}
%\setCJKmainfont[BoldFont=Noto Sans CJK SC Bold]{AR PL KaitiM GB} 

%\begin{solution}
    %\textcolor{red}{(待验证)}
%\end{solution}

%ifprintanswers
%\begin{figure}[H]
    %\centering        
    %\includegraphics[width=0.5\textwidth]{images/image25.png}
%\end{figure}[]
%\fi

%简体字版本latex code的code snippet,注意一定要简体字,以\question开始然后才\begin{solution},尽量按照原文,用\[\]或\begin{align*}或$...$但不要用$$...$$,中文字不可以在\[\]里面,不要用\text{中文字},把\over换成\frac{}{},不要用\textbf{} ,\boxed及\ans,去掉\overline{},句号不要在\[...\]
%give copyable latex code for every screenshot i provide
%give copyable latex code in a chunk&
%write into single line of \[...\] with ,\quad in between
%put into single \[...\] with \Rightarrow in between
%write into begin align*

%unexplored/to expand later areas: 
%USSR 101,
%日本大学入学考试
%放缩法in不等式
%数形结合in不等式
%JEE
%Prime Newtons
%2d vector证明whatsapp 29 Jan
%两圆正交
%limit的泰勒展开
%u\cdot v^hat=u^hat\cdot v
%projection matrices are idempotent and symmetric
%\int_{b}^{a} y \, dx + \int_{f(a)}^{f(b)} x \, dy = b f(b) - a f(a)
%linear algebra线代启示录
%陈景润,华罗庚,指考,APMO,SEAMO,CTMC,CSMC,AMC,台湾,中国
%integral marathon, integration bee

%大排版/clean up issues:
%跟着solution排版
%由易至难/跟着课文走?跟着课文走+综合题/挑战题
%难度:眼光、麻烦、雾里看花、课外太空
%函数的递推题丢去数列
%三角不等式放哪?
%数列里有用到复数?
%圆锥曲线
%三角恒等式?
%参数化积分放哪？
%边长去掉|...|for |AB|,保留for|\vec{AB}|
%不等式要等号成立条件

%fixed issues:
%余弦定理求边时打横写
%\mathbb{R}显示不正常,但如果comment out \usepackage{amsmath, amssymb, amsfonts} 就不能放\square,
%\mathbb{R}显示不正常,\setmathfont{TeX Gyre Termes Math} for \unicodemath,but not compatible with xelatex? - \usepackage{yhmath} for \wideparen
%inverse trig,differential equations
%微分方程步骤中间用小写c, 最后答案用大C - C_1,C_2,...
%若指数为一积分,用exp()
%裂项求和-> 累加法/累乘法

\newcommand{\ans}[1]{{\color{blue}#1}}
\newcommand*{\perm}[2]{{}^{#1}\!P_{#2}}
\newcommand*{\comb}[2]{{}^{#1}C_{#2}}
\newcommand{\parallels}{\mathbin{\!/\mkern-5mu/\!}}
\newcommand{\congr}{\ \reflectbox{$\cong$} \ }
\renewcommand{\ge}{\geqslant}
\renewcommand{\geq}{\geqslant}
\renewcommand{\le}{\leqslant}
\renewcommand{\leq}{\leqslant}
\DeclarePairedDelimiter\ceil{\lceil}{\rceil}
\DeclarePairedDelimiter\floor{\lfloor}{\rfloor}
\newcommand{\separator}{% 
  {\par\noindent
  \begin{tikzpicture}
    \coordinate (C) at (0,0);
    \coordinate (D) at (\textwidth,0);
    \path (C) to[ornament=88,options/.append style={yshift=.52pt,scale=0.8}] (D);
  \end{tikzpicture}}
}
\printanswers
\shadedsolutions
\definecolor{SolutionColor}{rgb}{1.0,0.98,0.8} %https://latexcolor.com/
\definecolor{PageColor}{rgb}{1.0,0.98,0.8}
\renewcommand{\solutiontitle}{\noindent}
\solutionsreseteqcounter
%\SolutionEmphasis{\color{red}}
\linespread{1.3}
\flagsdefault[width=0.5cm]
\setlength{\footskip}{0.5cm} 
\setlength{\columnsep}{1cm}
\newcommand{\figureornament}[2][1]{%
  \raisebox{.5\fontcharht\font`0}{%
    \pgfornament[scale=#1,color=brown]{#2}%
  }%
}
\runningfooter{}{\figureornament[0.2]{11}\;\; \thepage \;\; \figureornament[0.2]{14}}{}

\begin{document}
\begin{titlepage}
    \pagecolor{PageColor}
    \begin{center}
        \vspace*{1cm}
        {\huge $2025$年循人中学高三期末考兼统考预试\\[0.5em] 
        高中组\\[0.5em] 
        高级数学(III)\\[0.5em]  
       (SC007)\\[3.5em]
       \fontsize{65pt}{36pt}\selectfont{训练集}\\[3em]}

       \vspace{0.2cm}
       {\Large 采题:李冬恒}

       \vspace{0.5cm}
       {\Large \today}
        \begin{tikzpicture}[remember picture,overlay,every node/.style={inner sep=0pt}]
        % \node [shift={(1cm,-1cm)},blue,scale=2,anchor=north west] (CNW)
        %at (current page.north west) {\pgfornament[height=1cm,width=1cm]{61}};
        %\node [shift={(-1cm,-1cm)},blue,scale=2,anchor=north east] (CNE)
        %at (current page.north east) {\pgfornament[height=1cm,width=1cm,symmetry=v]{61}};
        %\node [shift={(1cm,1cm)},blue,scale=2,anchor=south west] (CSW)
        %at (current page.south west) {\pgfornament[height=1cm,width=1cm,symmetry=h]{61}};
        %\node [shift={(-1cm,1cm)},blue,scale=2,anchor=south east] (CSE)
        %at (current page.south east) {\pgfornament[height=1cm,width=1cm,symmetry=c]{61}};
        %\pgfornamentline[color=brown]{CNW}{CNE}{2}{87}
        %\pgfornamentline{CSW}{CSE}{2}{87}
        %\pgfornamentline{CNW}{CSW}{3}{87}
        %\pgfornamentline{CNE}{CSE}{3}{87}
        \pgfornamentline[color=brown]{current page.north west}{current page.north east}{2}{87}
        \pgfornamentline{current page.south west}{current page.south east}{2}{87}
        \pgfornamentline{current page.north west}{current page.south west}{3}{87}
        \pgfornamentline{current page.north east}{current page.south east}{3}{87}
        \end{tikzpicture}%

    \end{center}


\end{titlepage}


\thispagestyle{empty} %clear away page 1
\begin{center}
  {\fontsize{30pt}{26pt}\selectfont 前言
  }
\end{center}
\vspace{1cm}

\begin{tikzpicture}[every node/.style={inner sep=0pt}]
\node[text width=11.5cm](Text){%
\qquad 此训练集的雏形要追溯到 2019 年,依稀记得为当时的备赛爬了大量的AMC历届考题,创了“AMC 10 精选”的文档,其中我认为是属于统考范围的难题,便收录于另一个文档。\\[10pt]

\qquad 尘封已久,直到今年5月,我在整理电脑资料时偶然翻出,看着当年的内容,心想:读了一个数学系,难免觉得发展空间实属不少。在这几个月内,我随性地在网络上寻找、采集、整理更有趣(恶心)、更值得深思(烧脑)的题目,于是诞生了——高级数学(III)训练集。\\[10pt]

\qquad 起初,我并没想过要为每一题整理出解析,但发现到不少当年的题目只附答案没附解法。可惜的是,我的脑袋算是停止运作了一年,发现有一些题目对于现在的我实在是难以下咽。好吧,此乃下下策,但我决定用LaTeX开始动手为每一题补上解法、重写排版、增加新的内容。也好,既是备忘录,也是为了防止未来的我看到这些题目时会再次怀疑人生:「咦？是不是百尺竿头,更蠢一步？」\\[10pt]

\qquad 依据统考范围,我把训练集大致分为代数、组合数学、几何及微积分。若你想自行动手尝试,可以在 \texttt{preamble} 中将 \verb|\printanswers| 这行注释掉。\\[10pt]

\qquad 绝大部分解法都由官方或我提供,只有某些是\texttt{ChatGpt}生成后由我再校对,之中一定会有纰漏,欢迎批评指正,也希望这份训练集能助你一路披荆斩棘、越战越勇！\\[10pt]
\hfill —— 冬恒} ;
\color{brown}
\node[shift={(-1.5cm,1.5cm)},anchor=north west](CNW)
at (Text.north west) {\pgfornament[width=1.75cm]{61}};
\node[shift={(1.5cm,1.5cm)},anchor=north east](CNE)
at (Text.north east) {\pgfornament[width=1.75cm,symmetry=v]{61}};
\node[shift={(-1.5cm,-1.5cm)},anchor=south west](CSW)
at (Text.south west) {\pgfornament[width=1.75cm,symmetry=h]{61}};
\node[shift={(1.5cm,-1.5cm)},anchor=south east](CSE)
at (Text.south east) {\pgfornament[width=1.75cm,symmetry=c]{61}};
\pgfornamenthline{CNW}{CNE}{north}{87}
\pgfornamenthline{CSW}{CSE}{south}{87}
\pgfornamentvline{CNW}{CSW}{west}{87}
\pgfornamentvline{CNE}{CSE}{east}{87}
\end{tikzpicture}

\pagebreak
\begin{center}
  {\fontsize{30pt}{26pt}\selectfont 目录
  }
\end{center}
\separator

\begin{multicols}{2}
\begin{center}
    {\Large 代数}
\end{center}
\hyperlink{一元二次方程、多项式}{一元二次方程、多项式} \dotfill \pageref{一元二次方程、多项式} \\
\hyperlink{因式定理、余式定理}{因式定理、余式定理} \dotfill \pageref{因式定理、余式定理} \\
\hyperlink{根式、绝对值}{根式、绝对值} \dotfill \pageref{根式、绝对值} \\
\hyperlink{指数与对数}{指数与对数} \dotfill \pageref{指数与对数} \\
\hyperlink{方程组}{方程组} \dotfill \pageref{方程组} \\
\hyperlink{取整}{取整} \dotfill \pageref{取整} \\
\hyperlink{函数}{函数} \dotfill \pageref{函数} \\
\hyperlink{不等式}{不等式} \dotfill \pageref{不等式} \\
\hyperlink{数列与级数}{数列与级数} \dotfill \pageref{数列与级数} \\
\hyperlink{二项展开式}{二项展开式} \dotfill \pageref{二项展开式} \\
\hyperlink{泰勒展开式}{泰勒展开式} \dotfill \pageref{泰勒展开式} \\
\hyperlink{矩阵}{矩阵} \dotfill \pageref{矩阵} \\
\hyperlink{行列式}{行列式} \dotfill \pageref{行列式} \\
\hyperlink{复数}{复数} \dotfill \pageref{复数} \\
\hyperlink{数学归纳法}{数学归纳法} \dotfill \pageref{数学归纳法} 
\begin{center}
    {\Large 组合数学}
\end{center}
\hyperlink{排列与组合}{排列与组合} \dotfill \pageref{排列与组合} \\
\hyperlink{概率、期望值}{概率、期望值} \dotfill \pageref{概率、期望值} \\
\hyperlink{统计}{统计} \dotfill \pageref{统计}
\columnbreak
\begin{center}
    {\Large 几何}
\end{center}
\hyperlink{解三角形}{解三角形} \dotfill \pageref{解三角形} \\
\hyperlink{三角函数}{三角函数} \dotfill \pageref{三角函数} \\
\hyperlink{反三角函数}{反三角函数} \dotfill \pageref{反三角函数} \\
\hyperlink{平面向量}{平面向量} \dotfill \pageref{平面向量} \\
\hyperlink{直角坐标}{直角坐标} \dotfill \pageref{直角坐标} \\
\hyperlink{圆锥曲线}{圆锥曲线} \dotfill \pageref{圆锥曲线} \\
\hyperlink{坐标变换}{坐标变换} \dotfill \pageref{坐标变换} \\
\hyperlink{轨迹方程式、参数方程式}{轨迹方程式、参数方程式} \dotfill \pageref{轨迹方程式、参数方程式} \\
\hyperlink{极坐标}{极坐标} \dotfill \pageref{极坐标} \\
\hyperlink{立体几何、空间向量}{立体几何、空间向量} \dotfill \pageref{立体几何、空间向量} 
\begin{center}
    {\Large 微积分}
\end{center}
\hyperlink{极限}{极限} \dotfill \pageref{极限} \\
\hyperlink{微分}{微分} \dotfill \pageref{微分} \\
\hyperlink{积分}{积分} \dotfill \pageref{积分} \\
\hyperlink{积分技巧}{积分技巧} \dotfill \pageref{积分技巧} \\
\hyperlink{微分方程}{微分方程} \dotfill \pageref{微分方程} 
\end{multicols}
\vfill
\noindent 总题数: \numquestions{}
\pagecolor{PageColor}
\
\vfil
\hfil  {\fontsize{50pt}{36pt}\selectfont{代数}} \hfil 
\vfil
\begin{tikzpicture}[remember picture,overlay,every node/.style={inner sep=0pt}]
        \node [shift={(1cm,-1cm)},brown,scale=2,anchor=north west] (CNW)
        at (current page.north west) {\pgfornament[height=1cm,width=1cm]{61}};
        \node [shift={(-1cm,-1cm)},brown,scale=2,anchor=north east] (CNE)
        at (current page.north east) {\pgfornament[height=1cm,width=1cm,symmetry=v]{61}};
        \node [shift={(1cm,1cm)},brown,scale=2,anchor=south west] (CSW)
        at (current page.south west) {\pgfornament[height=1cm,width=1cm,symmetry=h]{61}};
        \node [shift={(-1cm,1cm)},brown,scale=2,anchor=south east] (CSE)
        at (current page.south east) {\pgfornament[height=1cm,width=1cm,symmetry=c]{61}};
        \pgfornamentline[color=brown]{current page.north west}{current page.north east}{2}{87}
        \pgfornamentline{current page.south west}{current page.south east}{2}{87}
        \pgfornamentline{current page.north west}{current page.south west}{3}{87}
        \pgfornamentline{current page.north east}{current page.south east}{3}{87}
        \end{tikzpicture}%
\thispagestyle{empty}
\pagebreak
\begin{center}
  {\fontsize{30pt}{26pt}\selectfont
    \hypertarget{一元二次方程、多项式}{一元二次方程、多项式} \label{一元二次方程、多项式}
  }
\end{center}
\separator
\vspace{1pt}
\nopagecolor
\begin{questions}
    \question 已知方程 
    \[
    x^2 - m x - m + 3 = 0
    \]
    的两根满足下列条件,求 $m$ 的取值范围:
    \begin{parts}
    \part 一根大于 $1$,另一根小于 $1$  
    \begin{solution}
        设$f(x)=x^2 - m x - m + 3$,所求即 
        \[
        f(1)=1-m-m+3<0 \Rightarrow m>2
        \]
    \end{solution}
    \part 一根小于 $0$,另一根大于 $2$  
    \begin{solution}
        即
        \[
        \begin{cases}
        f(0)= -m+3<0 \\
        f(2)= 4-2m-m+3<0
        \end{cases} \Rightarrow m > 3
        \]
    \end{solution}
    \part 一根在 $0$ 与 $1$ 之间,另一根在 $1$ 与 $2$ 之间  
    \begin{solution}
        即 
        \[
        \begin{cases}
        f(0)= -m+3>0 \\
        f(1)= 1-m-m+3<0 \\
        f(2)= 4-2m-m+3>0
        \end{cases} \Rightarrow 2<m<\frac{7}{3}
        \]
    \end{solution}
    \part 两根都在 $-4$ 与 $0$ 之间  
    \begin{solution}
        即  
        \[
        \begin{cases}
        f(-4)=16+4m-m+3>0 \\
        f(0)=-m+3>0 \\[1mm]
        \displaystyle -4<\frac{m}{2}<0 \\[2mm]
        \displaystyle f\left(\frac{m}{2}\right)=\frac{m^2}{4}-\frac{m^2}{2}-m+3\le 0
        \end{cases} \Rightarrow -\frac{19}{3}<m<-6
        \]
    \end{solution}
    \part 两根都大于 $-5$  
    \begin{solution}
        即  
        \[
        \begin{cases}
        f(-5)=25+5m-m+3>0 \\[1mm]
        \displaystyle \frac{m}{2}>-5 \\[2mm]
        \displaystyle f\left(\frac{m}{2}\right)=\frac{m^2}{4}-\frac{m^2}{2}-m+3\le 0
        \end{cases} \Rightarrow -7<m\le -6 \quad \text{或} \quad m\ge 2
        \]
    \end{solution}
    \part 有且仅有一根在 $0$ 与 $2$ 之间  
    \begin{solution}
        即  
        \[
        \begin{cases}
            \displaystyle 0<\frac{m}{2}<2, \\[2mm]
            \displaystyle f\left(\frac{m}{2}\right)=\frac{m^2}{4}-\frac{m^2}{2}-m+3=0
        \end{cases}   
        \]
        或
        \[
        f(0) \cdot f(2)=(-m+3)(4-2m-m+3)<0
        \]
        解得  
        \[
        m=2 \quad \text{或} \quad \frac{7}{3}<m<3
        \]
    \end{solution}
    \end{parts}

    \question 已知$y = x^3 - x^2 + 3x - 4$与$y = a x^2 - x - 4$恰好相交于两点, 求 $a$ 的可能值。 
    \begin{solution}
        联立方程得 \[ 
            x[x^2-(a+1)x+4]=0
        \]
        已知两线交于恰好两点,其中一点为$(0,4)$,则意味$$x^2-(a+1)x+4=0$$有重根,其判别式为零:$$[-(a+1)]^2-4\cdot1\cdot4=0$$解得$a=3, -5$
    \end{solution}
    
    \question 已知由 $f(x)=-x^{2}+4x+1$ 和 $g(x)=-x+5$ 所围成的封闭区域 $A$,若作一直线 $L$ 垂直于 $x$ 轴,分别与封闭区域 $A$ 的边界交于 $P,Q$ 两点,求在封闭区域内的 $\overline{PQ}$ 长的最大值为多少？
    \begin{solution}
        令 $f(x) = g(x)$,则
        \[
         -x^2 + 4x + 1 = -x + 5 \Rightarrow (x-1)(x-4) = 0 \Rightarrow x = 1,\,4
        \]
        所以封闭区域的 $x$ 范围为 $[1,4]$,而 $\overline{PQ}$ 即
        \[
        f(x) - g(x) = (-x^2 + 4x + 1) - (-x + 5) = -x^2 + 5x - 4 =-\left(x - \frac{5}{2}\right)^2 + \frac{9}{4}
        \]
        在 $x = \dfrac{5}{2}$时有最大值$\dfrac{9}{4}$
    \end{solution}

    \question 已知$k$为有理数,且使得方程式 $kx^{2}+(k-1)x+(k+1)=0$ 只有整数解,求$k$的所有可能值。
    \begin{solution}
        情况一:$k=0,$则原方程式为
        \[
        -x+1=0 \Rightarrow x=1 \in \mathbb{Z}
        \]
        情况二:$k\ne 0$,原方程式变为
        \[
        k(x^2+x+1)=x-1 \Rightarrow k=\frac{x-1}{x^2+x+1} \tag{1}
        \]
        设 $\alpha, \beta$ 为 $kx^2+(k-1)x+(k+1)=0$ 的两根,则
        \[
        \alpha+\beta = \frac{1-k}{k} = \frac{1}{k}-1 
        \]
        将 (1) 代入上式,
        \[
        \alpha+\beta = \frac{x^2+x+1}{x-1} - 1 =x+1+\frac{3}{x-1}\in \mathbb{Z}
        \]
        解得
        \[
        x-1=\pm1,\pm3
        \Rightarrow
        \begin{cases}
        x=2,4 \Rightarrow k=\dfrac{1}{7} \\
        x=0,-2 \Rightarrow k=-1
        \end{cases}
        \]
        故$k$的所有可能值为$-1,0,\dfrac{1}{7}$
    \end{solution}

    \question 若抛物线 $y = mx^2 - 1$ 上必存在相异两点对称于直线 $x+y=0$,求 $m$ 的范围。
    \begin{solution}
        设 $A, B$ 在抛物线 $\Gamma:y = mx^2 - 1$ 上且对称于直线 $L_1:x+y=0$,则 $A, B$ 同在直线 $$L_2:y = x + k$$ 上,其中 $L_1 \perp L_2$,抛物线与 $L_2$ 联立得
        \[
        mx^2 - x - 1 - k = 0 \tag{1}
        \]
        设 (1) 的解为 $a, b$,则$A(a, a+k), B(b, b+k)$的中点 $C\Big(\frac{a+b}{2}, \frac{a+b}{2}+k\Big)$ 在 $L_1$ 上,得
        \[
        \frac{a+b}{2} + \frac{a+b}{2} + k = 0 \Rightarrow k = -(a+b)
        \]
        由韦达定理,
        \[
        k = -(a+b) = -\frac{1}{m}
        \]
        且两实根相异,则判别式大于 $0$:
        \[
        \Delta = (-1)^2 - 4m(-\frac{1}{m} + 1) > 0  \Rightarrow m > \frac{3}{4}
        \]
    \end{solution}

    \question 实系数二次多项方程 $f(x)=0$ 有一根为 $2$,且方程 $f(f(x))=0$ 恰只有一实根为 $5$,求 $f(0)$。
    \begin{solution}
        设
        \[
        f(x)=a(x-2)(x-b)
        \]
        则
        \begin{align*}
        g(x)=f(f(x))
        &=a\,[a(x-2)(x-b)-2]\,[a(x-2)(x-b)-b]\\
        &=a^3\left(x^2-(b+2)x+2b-\frac{2}{a}\right)\left(x^2-(b+2)x+2b-\frac{b}{a}\right)
        &\equiv a^3 f_1(x) f_2(x)
        \end{align*}
        若$f_1=(x-5)^2,\ f_2 \text{的判别式} < 0,$可得
        \[
        a=-\frac{2}{9},\ b=8
        \]
        若$f_2=(x-5)^2,\ f_1 \text{的判别式} < 0,$可得
        \[
        a=-\frac{8}{9},\ b=8
        \]
        但 $f_1$ 判别式 $>0$,故舍去;因此
        \[
        f(0)=2ab = 2 \cdot \left(-\frac{2}{9}\right) \cdot 8 = -\frac{32}{9}
        \]
    \end{solution}

    \question 求所有实数 $c$,使得 $f(x)=x^{2}+4x+c$ 满足$f(f(x))$ 恰好有三个相异实根。
    \begin{solution}
        若 $f$ 的根为重根,则 $f \circ f$ 至多只有两个相异根,因此 $f$ 必须有两个相异根。设 $r$ 为其中一根使得 $f(x)=r$ 有重根,解
        \[
        x^2 + 4x + c - r = (x+2)^2
        \]
        得$r=c-4$ 是 $f$ 的根,故
        \[
        (c-4)^2 + 4(c-4) + c = 0 \Rightarrow c=0 \text{ 或 } c=3.
        \]
        当 $c=0$,解$x^2+4x=0$及$x^2+4x=-4$得
        \[
        x=0,-4,-2,
        \]
        当 $c=3$,解$x^2+4x+3=-3$及$x^2+4x+3=-1$,发现其中
        \[
        x^2+4x+6
        \]
        判别式为负,故唯一满足条件的 $c$ 为 $0$。
    \end{solution}

    \question 求非零实数三元组 $(a,b,c)$,使得
    \[
    (x^2 + 2ax + b)^2 + 2a(x^2 + 2ax + b) - b = (x-c)^4
    \]
    是多项式恒等式。
    \begin{solution}
        设 
        \[
        P(x) = x^2 + 2ax + b,Q(x) = x^2 + 2ax - b,
        \] 
        且$Q(P(x))$ 的根满足 $P(x) = r_1$ 或 $P(x) = r_2$,其中 $r_1, r_2$ 是 $Q$ 的根。欲使$Q(P(x))$ 只有一个四重根,$Q$ 必须有重根,因此判别式为
        \[
        4a^2 + 4b = 0 \Rightarrow b = -a^2,
        \]
        其中根为 $-a$,此时$P(x) = -a$ 也必须有重根,因此 $P(x) + a = x^2 + 2ax + b + a = 0$ 的判别式为
        \[
        4a^2 - 4(a+b) = 0.
        \]
        代入 $b = -a^2$ 解得
        \[
        a = \frac{1}{2}, \quad b = -\frac{1}{4}.
        \]
        $c$ 是 $P(x)+a=0$ 的解,即 $c = -a = -\frac{1}{2}$,所以三元组为
        \[
        (a,b,c) = \left(\frac{1}{2}, -\frac{1}{4}, -\frac{1}{2}\right)
        \]
    \end{solution}

    \question 求所有实数 $k$,使方程
    \[
    4x^2 + 4(2-k)x - k^2 = 0
    \]
    有两个实数解 $x_1,x_2$ 满足 $|x_1| = 2 + |x_2|$。
    \begin{solution}
        由韦达定理,
        \[
        x_1+x_2 = k-2, \quad x_1x_2 = -\frac{k^2}{4} \le 0
        \]
        条件 $|x_1| = 2 + |x_2|$ 等价于
        \[
        4 = (|x_1|-|x_2|)^2 = x_1^2+x_2^2-2|x_1x_2| = (x_1+x_2)^2 - 2x_1x_2 + 2x_1x_2
        \]
        解得
        \[
        k=0 \text{ 或 } \;4
        \]
    \end{solution}

    \question 解方程
    \[ 
    (x+1)(x+2)(x+3)^2(x+4)(x+5)=360 
    \]
    \begin{solution}   
        排版后有
        \begin{align*}
        (x+1)(x+5)(x+2)(x+4)(x+3)(x+3)&=360 \\
        (x^2+6x+5)(x^2+6x+8)(x^2+6x+9)&=360
        \end{align*}
        设 $y=x^2+6x$,变为
        \begin{align*}
        (y+5)(y+8)(y+9)&=360\\
        y(y^2+22y+157)&=0
        \end{align*}
        其中$y^2+22y+157=0$无实数解,解 $x^2+6x=0$ 可得 
        \[
        x=0, -6
        \]
    \end{solution}

    \question 解方程
    \[
    9x^4 - 24x^3 - 2x^2 - 24x + 9 = 0 \quad x \in \mathbb{R}
    \]
    \begin{solution}
        此方程为倒数方程式,将原方程写成
        \begin{align*}
        9x^2 - 24x - 2 - \frac{24}{x} + \frac{9}{x^2} &= 0 \\
        9\left(x^2 + \frac{1}{x^2}\right) - 24\left(x + \frac{1}{x}\right) - 2 &= 0
        \end{align*}
        设 $y = x + \frac{1}{x}$,则
        \[
        x^2 + \frac{1}{x^2} = y^2 - 2
        \]
        代入方程得
        \[
        9(y^2-2) - 24y - 2 = 0 \Rightarrow y = \frac{10}{3},-\frac{2}{3}
        \]
        当$x + \dfrac{1}{x} = \dfrac{10}{3}$,解得
        \[
        x = \frac{1}{3}, \quad x = 3
        \]
        当$x + \dfrac{1}{x} = -\dfrac{2}{3}$,方程
        \begin{align*}
        3x^2 + 2x + 3 &= 0
        \end{align*}
        的判别式 $2^2-4\cdot3\cdot3 = -32 < 0$,故无实数解,因此原方程式的实根为
        \[
        x = \frac{1}{3}, \quad x = 3
        \]
    \end{solution}

    \question 解
        \[
        \frac{x^2+16x+54}{x^2+11x+35} = \frac{x^2+13x+35}{x^2+14x+54}
        \]
    \begin{solution}
        使得原式交叉相乘时能顺利平方差,写成
        \[
        \frac{x^2+15x+54 + x}{x^2+12x+35 - x} = \frac{x^2+12x+35 + x}{x^2+15x+54 - x}
        \]
        \[
        (x^2+15x+54)^2 - x^2 = (x^2+12x+35)^2 - x^2
        \]
        \[
        (x^2+15x+54)^2 - (x^2+12x+35)^2 = 0
        \]
        再平方差得
        \[
        (3x+19)(2x^2+27x+89) = 0
        \]
        经检验,原方程的解为 $x = -\dfrac{19}{3},\dfrac{-27\pm\sqrt{17}}{4}$
    \end{solution}

    \question 解
    \[
    (x+1)^5+(x+1)^4(x-1)+(x+1)^3(x-1)^2+(x+1)^2(x-1)^3+(x+1)(x-1)^4+(x-1)^5=0
    \]
    其中$x \in \mathbb{C}$。
    \begin{solution}
        利用恒等式
        \[
        \frac{a^6-b^6}{a-b}=a^5+a^4b+a^3b^2+a^2b^3+ab^4+b^5
        \]
        则原式可写为
        \[
        \frac{(x+1)^6-(x-1)^6}{(x+1)-(x-1)}=0 (x+1)^6-(x-1)^6=0
        \]
        展开并化简得
        \[
        x(3x^2+1)(x^2+3)=0
        \]
        解得
        \[
        x=0,\quad x=\pm\frac{\sqrt{3}}{3}i,\quad x=\pm\sqrt{3}i
        \]
    \end{solution}

    \question 已知 $\alpha,\beta$ 是方程
    \[
    x^2 + (m-2)x + 1 = 0
    \]
    的根,求
    \[
    (1 + m\alpha + \alpha^2)(1 + m\beta + \beta^2).
    \]
    \begin{solution}
        由于 $\alpha,\beta$ 是方程
        \[
        x^2 + mx - 2x + 1 = 0
        \]
        的根,故满足
        \[
        \alpha^2 + m\alpha + 1 = 2\alpha,\quad 
        \beta^2 + m\beta + 1 = 2\beta
        \]
        于是
        \[
        (1 + m\alpha + \alpha^2)(1 + m\beta + \beta^2)  = (2\alpha)(2\beta) = 4\alpha\beta
        \]
        由韦达定理,$\alpha\beta = 1$,因此
        \[
        (1 + m\alpha + \alpha^2)(1 + m\beta + \beta^2) = 4 \cdot 1 = 4
        \]
    \end{solution}

    \question 若三次多项式 $x^{3}+3x-2=0$ 的根为 $a,b,c$,求以 $(a-b)^{2},(b-c)^{2},(c-a)^{2}$ 为根且首项系数为 $1$ 的三次多项式。
    \begin{solution}
        由韦达定理,
        \[
        a+b+c=0,\quad ab+bc+ca=3,\quad abc=2.
        \]
        且有
        \[
        a^3+3a=b^3+3b=c^3+3c=2.
        \]
        因此
        \[
        a^3-b^3+3(a-b)=0\Rightarrow (a-b)\bigl(a^2+ab+b^2+3\bigr)=0\Rightarrow a^2+ab+b^2=-3
        \]
        于是
        \[
        (a-b)^2=-3-3ab=-3(ab+1)
        \]
        同理,
        \[
        (b-c)^2=-3(bc+1),\quad (c-a)^2=-3(ca+1)
        \]
        故
        \begin{align*}
        (a-b)^2+(b-c)^2+(c-a)^2 
        &= -3(3+ab+bc+ca) = -18
        \end{align*}
        \begin{align*}
        &(a-b)^2(b-c)^2+(b-c)^2(c-a)^2+(c-a)^2(a-b)^2 \\
        &= 9\big((ab+1)(bc+1)+(bc+1)(ca+1)+(ca+1)(ab+1)\big) \\
        &= 9\big(2(a+b+c)+2(ab+bc+ca)+3\big) = 81
        \end{align*}
        \begin{align*}
        (a-b)^2(b-c)^2(c-a)^2
        &= -27(ab+1)(bc+1)(ca+1) \\
        &= -27\big(abc(a+b+c)+(abc)^2+ab+bc+ca+1\big) = -216
        \end{align*}
        所求三次多项式为
        \[
        x^3 +18x^2 +81x +216
        \]
    \end{solution}

    \question 设 $a, b, c \in \mathbb{R}$, 证明方程式
    \[
    \frac{b+c}{x-a}+\frac{c+a}{x-b}+\frac{a+b}{x-c}=3
    \]
    的根都是实根。
    \begin{solution}
        有
        \begin{align*}
        \frac{b+c}{x-a} + \frac{c+a}{x-b} + \frac{a+b}{x-c} &= 3 \\
        1 - \frac{b+c}{x-a} + 1 - \frac{c+a}{x-b} + 1 - \frac{a+b}{x-c} &= 0 \\
        \frac{x-a-b-c}{x-a} + \frac{x-a-b-c}{x-b} + \frac{x-a-b-c}{x-c} &= 0 \\
        (x-a-b-c)\left( \frac{1}{x-a} + \frac{1}{x-b} + \frac{1}{x-c} \right) &= 0
        \end{align*}
        \[
        (x-a-b-c) \frac{ x^2 - (b+c)x + bc + x^2 - (a+c)x + ac + x^2 - (a+b)x + ab }{ (x-a)(x-b)(x-c) } = 0 
        \]
        \[
        (x-a-b-c)\left( 3x^2 - 2(a+b+c)x + ab + bc + ca \right) = 0
        \]
        故方程式
        \[
        \frac{b+c}{x-a} + \frac{c+a}{x-b} + \frac{a+b}{x-c} = 3
        \]
        有一实根 $x = a + b + c$,现探讨另两根,方程式
        \[
        3x^2 - 2(a+b+c)x + ab + bc + ca = 0.
        \]
        的判别式为
        \[
        \Delta = 4(a+b+c)^2 - 12(ab+bc+ca) = 4\left(a^2 + b^2 + c^2 - ab - bc - ca\right).
        \]
        由 AM-GM 不等式,
        \[
        ab + bc + ca \leq \frac{a^2 + b^2}{2} + \frac{b^2 + c^2}{2} + \frac{c^2 + a^2}{2} = a^2 + b^2 + c^2.
        \]
        故判别式为非负,于是方程式 $3x^2 - 2(a+b+c)x + ab + bc + ca = 0$的根都是实根,得证原方程式的根都是实根。
    \end{solution}

    \question 设 $f(x)=x^3+4x^2+8x+16$,计算
    \[
    f(x+7)-f(x+6)-f(x+5)+f(x+4)-f(x+3)+f(x+2)+f(x+1)-f(x).
    \]
    \begin{solution}
        设 $g(x)=f(x+1)-f(x)$,则 
        \[
        g(x)=3x^2+10x+12 
        \]
        同理设 $h(x)=g(x+2)-g(x)=12x+24$,原式可写成
        \[
        g(x+6)-g(x+4)-g(x+2)+g(x)=h(x+4)-h(x)=12(x+4)-12x=48
        \]
    \end{solution}

    \question 已知关于$x$的方程$x^3+ax^2+bx+c=0$的三个非零实数根成等比数列,求$a^3c-b^3$的值。
    \begin{solution}
        设方程$x^3+ax^2+bx+c=0$的三根为$\dfrac{\alpha}{r},\alpha,\alpha r,\alpha \ne 0 $,由韦达定理,
        \begin{align*}
            -a&=\frac{\alpha}{r}+\alpha+\alpha r=\alpha(\frac{1}{r}+1+r) \tag{1}\\
            b&=\frac{\alpha^2}{r}+\alpha^2+\alpha^2 r=\alpha^2(\frac{1}{r}+1+r) \tag{2}\\
            -c&=\alpha^3 \tag{3}
        \end{align*}
        由$(1),(2)$得$-\dfrac{b}{a}=\alpha$,代入(3)得
        \[
        -c=-\frac{b^3}{a^3} \Rightarrow a^3c-b^3=0
        \]
    \end{solution}

    \question 已知$\left(x + \dfrac{1}{x}\right)^2=3$,求
    \[
    x^{63} + x^{44} + x^{37} + x^{31} + x^{26} + x^9 + 6
    \]
    的值。
    \begin{solution}
        由已知得
        \[
        \left(x + \frac{1}{x}\right)^3 = x^3 + \frac{1}{x^3} + 3(1)(\sqrt{3}) = 3\sqrt{3}
        \]
        给出
        \[
        x^3 + \frac{1}{x^3} = 0
        \]
        同理可得,
        \[
        x^9 + \frac{1}{x^9} = 0, \quad x^{27} + \frac{1}{x^{27}} = 0
        \]
        故
        \begin{align*}
        &x^{63} + x^{44} + x^{37} + x^{31} + x^{26} + x^9 + 6 \\
        & =x^{36}\left(x^{27} + \frac{1}{x^{27}}\right) + x^{35} \left(x^9 + \frac{1}{x^9}\right) + x^{34} \left(x^3 + \frac{1}{x^3}\right) + 6 =6
        \end{align*}
    \end{solution}

    \question 已知实数$x$满足
    \[
    x^3 + \frac{1}{x^3} = 18
    \]
    求
    \[
    x^{11} + \frac{1}{x^{11}}.
    \]
    \begin{solution}
        设$a = x + \frac{1}{x}$,则
        \[
        a^3 - 3a = 18 \Rightarrow (a-3)(a^2 + 3a + 6) = 0
        \]
        所以
        \[
        a = x + \frac{1}{x} = 3 \in \mathbb{R}
        \]
        则
        \[
        x^2+\frac{1}{x^2} = 3^2-2=7
        \]
        同理得
        \[
        x^4+\frac{1}{x^4} = 47,\quad x^8+\frac{1}{x^8} = 2207
        \]
        于是
        \[
        \left(x^2+\frac{1}{x^2}\right)\left(x^3 + \frac{1}{x^3}\right)=x^5+\frac{1}{x^5} + x + \frac{1}{x}= 126 \Rightarrow x^5+\frac{1}{x^5}=123
        \]
        且
        \[
        \left(x^3+\frac{1}{x^3}\right)\left(x^8 + \frac{1}{x^8}\right)=x^{11}+\frac{1}{x^{11}} + x^5 + \frac{1}{x^5}= 39726
        \]
        故
        \[
        x^{11}+\frac{1}{x^{11}}=39726-123=39603
        \]
    \end{solution}

    \question 已知 $\alpha, \beta, \gamma$ 是方程
    \[
    x^3 - x - 1 = 0
    \]
    的三根, 计算
    \[
    \frac{1-\alpha}{1+\alpha} + \frac{1-\beta}{1+\beta} + \frac{1-\gamma}{1+\gamma}
    \]
    的值。 
    \begin{solution}
        由韦达定理,$\;\alpha+\beta+\gamma=0,\;\alpha\beta+\beta\gamma+\gamma\alpha=-1,\;\alpha\beta\gamma=1$,故
        \begin{align*}
        \frac{1-\alpha}{1+\alpha} + \frac{1-\alpha}{1+\beta} + \frac{1-\gamma}{1+\gamma} 
        &= \frac{2}{1+\alpha} + \frac{2}{1+\beta} + \frac{2}{1+\gamma}-3 \\
        &= 2\cdot\frac{(1+\alpha)(1+\beta)+(1+\beta)(1+\gamma)+(1+\gamma)(1+\alpha)}{(1+\alpha)(1+\beta)(1+\gamma)}-3 \\
        &= 2\cdot\frac{3+2(\alpha+\beta+\gamma)+\alpha\beta+\beta\gamma+\gamma\alpha}{1+\alpha+\beta+\gamma+\alpha\beta+\beta\gamma+\gamma\alpha+\alpha\beta\gamma}-3 \\
        &= 2\cdot\frac{3-1}{1-1+1}-3 \\
        &= 1
        \end{align*}
    \end{solution}
        
    \question 已知 $\alpha,\beta,\gamma,\delta$ 为 $x^{4}+2x^{3}+x^{2}+2x+1=0$ 的四根,试求 $$(\alpha^{2}+\alpha+1)(\beta^{2}+\beta+1)(\gamma^{2}+\gamma+1)(\delta^{2}+\delta+1)$$ 的值。
    \begin{solution}
        $\alpha$ 是方程
        $$x^{4}+2x^{3}+x^{2}+2x+1=0$$
        的根,则
        \[
        \alpha^4 + 2\alpha^3 + \alpha^2 + 2\alpha + 1 = 0
        \Rightarrow (\alpha^2 + \alpha + 1)^2 = 2\alpha^2
        \Rightarrow \alpha^2 + \alpha + 1 = \sqrt{2} \alpha
        \]
        同理可得
        \[
        \beta^2 + \beta + 1 = \sqrt{2} \beta,\quad
        \gamma^2 + \gamma + 1 = \sqrt{2} \gamma,\quad
        \delta^2 + \delta + 1 = \sqrt{2} \delta
        \]
        由韦达定理,
        \[
        (\alpha^2 + \alpha + 1)(\beta^2 + \beta + 1)(\gamma^2 + \gamma + 1)(\delta^2 + \delta + 1)
        = (\sqrt{2})^4 \cdot \alpha \beta \gamma \delta =4 \cdot 1 = 4
        \]
    \end{solution}

    \question 设方程 $x^{3}-4x+1=0$ 的三个相异复数根为 $a, b, c$,求 
    $$\frac{a+1}{(a-1)^{4}}+\frac{b+1}{(b-1)^{4}}+\frac{c+1}{(c-1)^{4}}$$ 的值。
    \begin{solution}
        发现
        \[\frac{a+1}{(a-1)^4} = \frac{1}{(a-1)^3} + \frac{2}{(a-1)^4}
        \],因此欲求以 $\displaystyle \frac{1}{a-1}, \frac{1}{b-1}, \frac{1}{c-1}$ 为三根的多项式,令 $x_1 = x - 1 $ 代入原方程式
        \[
        (x_1+1)^3 - 4(x_1+1) + 1 = 0 \Rightarrow x_1^3 + 3x_1^2 - x_1 - 2 = 0
        \]
        再令 $x_2 = \frac{1}{x_1}$ 代入可得
        \[
        \frac{1}{x_2^3} + \frac{3}{x_2^2} - \frac{1}{x_2} - 2 = 0 \Rightarrow x^3_2 + \frac12 x^2_2 - \frac32 x_2 - \frac12 = 0
        \]
        故
        \[
        \frac{a+1}{(a-1)^4} + \frac{b+1}{(b-1)^4} + \frac{c+1}{(c-1)^4}
        = (\alpha^3 + \beta^3 + \gamma^3) + 2(\alpha^4 + \beta^4 + \gamma^4) \tag{1}
        \]
        设 $\displaystyle f(x) = x^3 + \frac12 x^2 - \frac32 x - \frac12$, 则 $ f'(x) = 3x^2 + x - \dfrac32$,利用长除法计算 $\dfrac{f'(x)}{f(x)}$:
        \[
        \frac{f'(x)}{f(x)} = \frac{3}{x} - \frac12 \cdot \frac{1}{x^2} + \frac{13}{4} \cdot \frac{1}{x^3}
        - \frac{7}{8} \cdot \frac{1}{x^4} + \frac{81}{16} \cdot \frac{1}{x^5} + \cdots
        \]
        由系数比较得
        \[
        \alpha^3 + \beta^3 + \gamma^3 = -\frac78,\quad
        \alpha^4 + \beta^4 + \gamma^4 = \frac{81}{16}
        \]
        故所求为
        \[
        -\frac78 + 2 \cdot \frac{81}{16} = \frac{37}{4}
        \]
    \end{solution}

    \question 已知 $f(x) = x^4 + px^3 + qx^2 + rx + s$ 且 $f(1)=59,f(2)=118,f(3)=177$,求 $f(9)+f(-5)$ 。
    \begin{solution}
        考虑函数 $g(x) = 59x$,则
        \[
        f(1)=g(1)=59,\quad f(2)=g(2)=118,\quad f(3)=g(3)=177。
        \]
        定义 $h(x) = f(x)-g(x)$,则 $x=1,2,3$ 为 $h(x)$ 的根且 $h(x)$ 首项系数为 $1$,所以设
        \[
        h(x) = (x-1)(x-2)(x-3)(x-a)
        \]
        其中$a$是一实数,于是
        \begin{align*}
        f(9)+f(-5) &= h(9)+g(9) + h(-5)+g(-5) \\
        &= 336(9-a) + 9\cdot 59 + 336(5+a) - 5\cdot 59 \\
        &= 4940
        \end{align*}
    \end{solution}

    \question 已知 $x_1, x_2, x_3$ 是方程
    \[
    x^3 - 6x^2 + a x - a = 0
    \]
    的根,且满足
    \[
    (x_1 - 3)^3 + (x_2 - 3)^3 + (x_3 - 3)^3 = 0,
    \]
    求 $a$ 的值。
    \begin{solution}
        令 $r_i = x_i - 3$, 则 $r_1, r_2, r_3$ 是方程
        \[
        (x+3)^3 - 6(x+3)^2 + a(x+3) - a = 0
        \]
        即方程
        \[
        x^3 + 3x^2 + (a-9)x + (2a-27) = 0
        \]
        的根,由韦达定理,
        \[
        r_1 + r_2 + r_3 = -3, \quad r_1r_2 + r_2r_3 + r_3r_1 = a-9, \quad r_1r_2r_3 = 27-2a
        \]
        由立方和公式,
        \[
        r_1^3 + r_2^3 + r_3^3 = -3(9 - 3(a-9)) + 3(27 - 2a) = 0
        \]
        解得
        \[
        a = 9
        \]
    \end{solution}

    \question 已知 $x_{1}, x_{2}, x_{3}$ 为方程 $\sqrt{123}x^{3}-247x^{2}+2=0$ 的相异实根, 且 $x_{1}<x_{2}<x_{3}$, 求 $x_{2}(x_{3}^{2}-x_{1}^{2})$的值。
    \begin{solution}
        设 $\sqrt{123}=a$, 则方程 $\sqrt{123}x^{3}-247x^{2}+2=0$ 化为 $$ax^{3}-(2a^{2}+1)x^{2}+2=0\Rightarrow(ax-1)(x^{2}-2ax-2)=0$$
        所以方程的3个实数根为 $$\frac{1}{a},a+\sqrt{a^{2}+2},a-\sqrt{a^{2}+2}$$
        因为 $a=\sqrt{123}>1$,所以 $$a-\sqrt{a^{2}+2}<0<\frac{1}{a}<1<a+\sqrt{a^{2}+2}$$
        因此 $x_{1}=a-\sqrt{a^{2}+2},x_{2}=\dfrac{1}{a}, x_{3}=a+\sqrt{a^{2}+2}$,而
        $$x_{2}(x_{3}^{2}-x_{1}^{2})=\frac{1}{a}\cdot2\sqrt{a^{2}+2}\cdot 2a=20\sqrt5$$
    \end{solution}
    
    \question 已知正整数 $m,n$ 满足 $m^2 - n = 32$,且 $\sqrt[5]{m+\sqrt{n}} + \sqrt[5]{m-\sqrt{n}}$ 是方程 $$x^5 - 10x^3 + 20x - 40 = 0$$ 的一个实根,求 $m,n$ 的值。
    
    \begin{solution}
        令 $a = \sqrt[5]{m+\sqrt{n}},b = \sqrt[5]{m-\sqrt{n}}$,则 $x = a + b$ 是该方程的一个实根,且
        \[
        a^5 + b^5 = 2m, \quad a^5 b^5 = m^2 - n = 32
        \Rightarrow ab = 2
        \]
        考虑 $(a + b)^5$ 的展开式:
        \[
        x^5 = a^5 + b^5 + 5ab(a^3 + b^3) + 10a^2 b^2 (a + b)
        \]
        由$a^3 + b^3 = (a + b)\left[(a + b)^2 - 3ab\right] = x(x^2 - 6)$,得 
        \[
        x^5 = 2m + 10x(x^2 - 6) + 40x 
        \]
        整理得:
        \[
        x^5 - 10x^3 + 20x- 2m=0
        \]
        对比原方程式得
        \[
        m = 20 , n = 368
        \]
    \end{solution}

    \question 设方程式 $x^5 + x^4 - x^2 + 1 = 0$ 的五个根为 $\alpha_1, \alpha_2, \alpha_3, \alpha_4, \alpha_5$, 若 $P(x) = x^4 - 1$,求 $$P(\alpha_1) P(\alpha_2) P(\alpha_3)  P(\alpha_4) P(\alpha_5)$$ 的值。
    \begin{solution}
        设
        \[
        f(x) = x^5 + x^4 - x^2 + 1 = (x - \alpha_1)(x - \alpha_2)(x - \alpha_3)(x - \alpha_4)(x - \alpha_5) 
        \]
        令$x=1$得
        \[
        f(1) = (1 - \alpha_1)(1 - \alpha_2)(1 - \alpha_3)(1 - \alpha_4)(1 - \alpha_5) = 0
        \]
        由于
        \[P(x) = x^4 - 1 = (x^2 + 1)(x - 1)(x + 1) \]
        故
        \begin{align*}
        &P(\alpha_1) P(\alpha_2) P(\alpha_3)  P(\alpha_4) P(\alpha_5)\\
        &= -(1 - \alpha_1)(1 - \alpha_2)(1 - \alpha_3)(1 - \alpha_4)(1 - \alpha_5) \prod_{k=1}^5 (\alpha_{k}^2+1)(\alpha_{k}+1)\\ &= 0
        \end{align*}
    \end{solution}

    \question 已知 $p(x)$ 是一个整系数多项式,定义 $q(x)=\dfrac{p(x)}{x(1-x)}$,若对所有 $x\neq 0,1$ 都有
    \[
    q(x) = q\left(\frac{1}{1-x}\right),
    \]
    且 $p(2)=p(3)=5$,求 $p(4)$ 的值。
    \begin{solution}
        由 $q(x)=q\left(\dfrac{1}{1-x}\right)$ 可得
        \[
        \frac{p(x)}{x(1-x)} = \frac{p\left(\frac{1}{1-x}\right)}{\frac{1}{1-x}\left(1-\frac{1}{1-x}\right)} \Rightarrow p(x) = (1-x)^3 p\left(\frac{1}{1-x}\right)
        \]
        这说明 $\deg(p)\le 3$,设 $p(x)=ax^3+bx^2+cx+d$,则
        \[
        (1-x)^3 p\left(\frac{1}{1-x}\right)=d(1-x)^3+c(1-x)^2+b(1-x)+a
        \]
        比较$x^3,x^2$系数得
        \[
        a=d, \quad b=-c-3d
        \]
        且由 $p(2)=p(3)=5$ 得
        \[
        -3a - 2c = a - 6c =5 \Rightarrow a=c=-1
        \]
        于是
        \[
        p(x) = -x^3 +4x^2 - x -1 \Rightarrow p(4) = -5
        \]
    \end{solution}

    \question 已知多项式 $x^7 - 5$ 的七个相异根为 $r_1, \dots, r_7$,求
    \[
    \prod_{1 \le i < j \le 7} (r_i + r_j)^2
    \] 的值。
    \begin{solution}
        设所求乘积为 $P$,由韦达定理,考虑
        \begin{align*}
        2^7 \cdot 5 \cdot P 
        &= \prod_{i=1}^{7} 2 r_i \prod_{1 \le i < j \le 7} (r_i + r_j)^2\\
        &= \prod_{1 \le i \le j \le 7} (r_i + r_j) \prod_{1 \le i < j \le 7} (r_i + r_j) \prod_{1 \le j < i \le 7} (r_i + r_j)\\
        &= \prod_{i=1}^{7} \prod_{j=1}^{7} (r_i + r_j)
        \end{align*}
        且注意到
        \[
        \prod_{j=1}^7(x-r_j) = x^7 - 5 \Rightarrow \prod_{j=1}^7(x+r_j) = x^7 + 5,
        \]
        故
        \[
        \prod_{j=1}^{7} (r_i + r_j) = r_i^7 + 5 = 10
        \]
        因此
        \[
        2^7 \cdot 5 \cdot P = 10^7 \Rightarrow P = 5^6
        \]
    \end{solution}

    \question 设 $P(x)$ 为次数为 $10$ 的多项式,且满足 
    \[
    P(2^i)=i,0\le i\le 10,
    \]
    求 $P(x)$ 中 $x$ 项的系数。
    \begin{solution}
        令 $Q(x)=P(2x)-P(x)-1$,则 $Q(2^i)=0,0\le i\le9$,所以
        \[
        Q(x)=\alpha\prod_{k=0}^{9}(x-2^k)
        \]
        其中$\alpha$为一常数,又 $Q(0)=P(0)-P(0)-1=-1$,且
        \[
        \prod_{k=0}^{9}(0-2^k)=2^{45},
        \]
        因此 $\alpha=-\dfrac{1}{2^{45}}$,记 
        \[
        R(x)=\prod_{k=0}^{9}(x-2^k),
        \]
        其 $x$ 项系数为
        \[
        \text{coef}_{x}(R) = -\left(\prod_{k=0}^{9}2^k\right)\sum_{k=0}^{9}\frac{1}{2^k}
        = -2^{45}\left(2-\frac{1}{2^9}\right)
        \]
        因为$P(2x)-P(x)-1$ 的 $x$ 项系数恰好等于 $P(x)$ 的 $x$ 项系数,故
        \[
        \text{coef}_{x}(Q)=\alpha\cdot\text{coef}_{x}(R)
        = -\frac{1}{2^{45}}\cdot\left(-2^{45}\cdot\frac{1023}{512}\right)=\frac{1023}{512}
        \]
    \end{solution}

    \question 已知 $x_1, x_2, \ldots, x_{2015}$ 是方程
    \[
    x^{2015}+x^{2014}+\dots+x^2+x+1=0
    \]
    的根,求
    \[
    \frac{1}{1-x_1}+\frac{1}{1-x_2}+\dots+\frac{1}{1-x_{2015}}.
    \]
    \begin{solution}
        设 
        \[
        a_k = \frac{1}{1-x_k}, \quad k=1,2,\dots,2015.
        \]
        则
        \[
        x_k = \frac{a_k-1}{a_k}.
        \]
        代入原方程,得到
        \[
        \frac{(a_k-1)^{2015}}{a_k^{2015}} + \frac{(a_k-1)^{2014}}{a_k^{2014}} + \dots + \frac{a_k-1}{a_k} + 1 = 0.
        \]
        两边同时乘以 $a_k^{2015}$,得
        \[
        (a_k-1)^{2015} + a_k(a_k-1)^{2014} + \dots + a_k^{2014}(a_k-1) + a_k^{2015} = 0.
        \]
        因此
        \[
        a_1 + a_2 + \dots + a_{2015} = \frac{\comb{2015}{1}+\comb{2014}{1}+\comb{1}{1}}{2016} = \frac{2015}{2}
        \]
        而当$k=1,2,\dots,n,a_1 + a_2 + \dots + a_{2015}=\dfrac{n}{2}$
    \end{solution}
    \begin{solution}
        设
        \[
        f(x) = x^{2015}+x^{2014}+\dots+x+1,
        \]
        则
        \[
        f(x) = \prod_{k=1}^{2015}(x-x_k) \Rightarrow \frac{f'(x)}{f(x)} = \sum_{k=1}^{2015} \frac{1}{x-x_k}.
        \]
        取 $x=1$,得到
        \[
        \sum_{k=1}^{2015} \frac{1}{1-x_k} = \frac{f'(1)}{f(1)}= \frac{2015}{2}
        \]
        其中
        \[
        f(1) = 2016, \quad f'(1) = 2015+2014+\dots+1 = \frac{2015\cdot 2016}{2}
        \]
    \end{solution}

    \question 
    \begin{parts}
    \part 已知 $\alpha,\beta,\gamma$ 为三实数,设 $t=-(\alpha+\beta+\gamma),v=\alpha\beta+\beta\gamma+\gamma\alpha$,且满足
    \[
    \alpha\beta\gamma=-1, \quad t+v=-3,
    \]
    试证
    \[
    \alpha^{\frac{1}{3}}+\beta^{\frac{1}{3}}+\gamma^{\frac{1}{3}}=\sqrt[3]{(-t-6)+3\sqrt[3]{t^2+3t+9}}.
    \]
    \begin{solution}
        设
        \[
        f(x)=(x-\alpha)(x-\beta)(x-\gamma)=x^3+tx^2+vx+1,
        \]
        \[
        g(y)=(y-\alpha^{\frac{1}{3}})(y-\beta^{\frac{1}{3}})(y-\gamma^{\frac{1}{3}})
        =y^3+ay^2+by+1
        \]
        由$g(y)=0$得$y^3+1=-y(ay+b)$,于是
        \[
        (y^3+1)^3=-y^3\left(a^3y^3+b^3+3ab y(ay+b)\right)
        =-y^3\left(a^3y^3+b^3-3ab(y^3+1)\right)
        \]
        令 $x=y^3$,展开整理得
        \[
        x^3+(a^3-3ab+3)x^2+(b^3-3ab+3)x+1=0
        \]
        于是$a^3-3ab+3=t, b^3-3ab+3=v$,令 $z=ab$,则
        \begin{align*}
        z^3&=(t+3z-3)(v+3z-3) \\
        &=tv+3z(t+v)-3(t+v)+9z^2-18z+9 \\
        &=t(-3-t)+3z(-3)-3(-3)+9z^2-18z+9 \\
        &=9z^2-27z+(-t^2-3t+18)
        \end{align*}
        即
        \[
        (z-3)^3=-t^2-3t-9 \Rightarrow z=3-\sqrt[3]{t^2+3t+9}
        \]
        由韦达定理,
        \[
        \alpha^{\frac{1}{3}}+\beta^{\frac{1}{3}}+\gamma^{\frac{1}{3}}=-a=-\sqrt[3]{t+3z-3}=\sqrt[3]{(-t-6)+3\sqrt[3]{t^2+3t+9}}
        \]
        故证毕。
    \end{solution}
    \part 据此,试证
    \[
    \sqrt[3]{\cos\frac{2\pi}{9}}+\sqrt[3]{\cos\frac{4\pi}{9}}+\sqrt[3]{\cos\frac{8\pi}{9}}
    =\sqrt[3]{\frac{3}{2}(\sqrt[3]{9}-2)}
    \]
    \begin{solution}
        考虑
        \[
        \alpha=2\cos\frac{2\pi}{9},\quad
        \beta=2\cos\frac{4\pi}{9},\quad
        \gamma=2\cos\frac{8\pi}{9}.
        \]
        则由
        \begin{align*}
        \cos{\frac{2\pi}{9}} \cos{\frac{4\pi}{9}} \cos{\frac{8\pi}{9}} &= \frac{1}{2} \left(\cos{\frac{6\pi}{9}}+\cos{\frac{2\pi}{9}} \right)\cos{\frac{8\pi}{9}} \\
        &= \frac{1}{2} \left(-\frac{1}{2}+\cos{\frac{2\pi}{9}} \right)\cos{\frac{8\pi}{9}} \\
        &= \frac{1}{2} \left(-\frac{1}{2} \cos{\frac{8\pi}{9}}+ \cos{\frac{2\pi}{9}}\cos{\frac{8\pi}{9}} \right) \\
        &= \frac{1}{2} \left(-\frac{1}{2} \cos{\frac{8\pi}{9}}+ \frac{1}{2}\left(\cos{\frac{10\pi}{9}} +\cos{\frac{6\pi}{9}}\right) \right) \\
        &= \frac{1}{2} \left(\frac{1}{2} \cos{\frac{\pi}{9}}- \frac{1}{2}\cos{\frac{\pi}{9}} -\frac{1}{4} \right) = -\frac{1}{8} 
        \end{align*}
        满足已知$\alpha\beta\gamma=-1$,此时
        \begin{align*}
        t =-(\alpha+\beta+ \gamma) &=-2\left(\cos{\frac{2\pi}{9}}+ \cos{\frac{4\pi}{9}}+ \cos{\frac{8\pi}{9}} \right) \\
        &=-2\left(2\cos{\frac{3\pi}{9}}\cos{\frac{\pi}{9}}+ \cos{\frac{8\pi}{9}} \right) \\
        &=-2\left(\cos{\frac{\pi}{9}}+ \cos{\frac{8\pi}{9}} \right) \\
        &=-2\cdot 2\cos{\frac{\pi}{2}}\cos{\frac{7\pi}{18}} =0
        \end{align*}
        故
        \[
        \alpha^{\frac{1}{3}}+\beta^{\frac{1}{3}}+\gamma^{\frac{1}{3}}
        =\sqrt[3]{-6+3\sqrt[3]{9}}
        \]
        且
        \[
        \sqrt[3]{\cos\frac{2\pi}{9}}+\sqrt[3]{\cos\frac{4\pi}{9}}+\sqrt[3]{\cos\frac{8\pi}{9}}
        =\frac{1}{\sqrt[3]{2}}\bigl(\alpha^{\frac{1}{3}}+\beta^{\frac{1}{3}}+\gamma^{\frac{1}{3}}\bigr)
        =\sqrt[3]{\frac{3}{2}(\sqrt[3]{9}-2)}.
        \]
        证毕。
    \end{solution}
    \end{parts}

    \question 求 
    \[
    \frac{(1^{4}+18^{2})(11^{4}+18^{2})(23^{4}+18^{2})(35^{4}+18^{2})(47^{4}+18^{2})}{(5^{4}+18^{2})(7^{4}+18^{2})(17^{4}+18^{2})(29^{4}+18^{2})(41^{4}+18^{2})}
    \] 
    的值。
    \begin{solution}
        因式分解给出
        \[
        a^4+18^2=(a^2+18)^2-(6a)^2=(a^2+6a+18)(a^2-6a+18)=((a+3)^2+9)((a-3)^2+9)
        \] 
        故
        \begin{align*}
        &\frac{(1^{4}+18^{2})(11^{4}+18^{2})(23^{4}+18^{2})(35^{4}+18^{2})(47^{4}+18^{2})}{(5^{4}+18^{2})(7^{4}+18^{2})(17^{4}+18^{2})(29^{4}+18^{2})(41^{4}+18^{2})}\\
        &=\frac{(4^2+9)(2^2+9)(14^2+9)(8^2+9)(26^2+9)(20^2+9)(38^2+9)(32^2+9)(50^2+9)(44^2+9)}{(8^2+9)(2^2+9)(10^2+9)(4^2+9)(20^2+9)(14^2+9)(32^2+9)(26^2+9)(44^2+9)(38^2+9)} \\
        &=\frac{50^2+9}{10^2+9}=\frac{2509}{109}
        \end{align*}
    \end{solution}

    \question 已知三实数 $a,b,c$ 满足 $\sqrt{3}(a-b)+3(b-c)+(c-a)=0$,且 $b \neq c$,求 
    \[
    \frac{(a-b)(a-c)}{(b-c)^2}
    \]
    的值。
    \begin{solution}
        令$\alpha = a-b, \beta = b-c, \gamma = c-a$,则
        \[
        \alpha+\beta+\gamma= \sqrt{3}\alpha +3\beta +\gamma =0
        \]
        由此得
        \[
        \gamma = -(\alpha+\beta), \quad \alpha+\beta = \sqrt{3}\alpha + 3\beta \Rightarrow \beta = \frac{1-\sqrt{3}}{2}\alpha
        \]
        因此
        \[
        \frac{(a-b)(a-c)}{(b-c)^2} = \frac{-\alpha \gamma}{\beta^2} = \frac{\alpha(\alpha+\beta)}{\left(\frac{1-\sqrt{3}}{2}\right)^2 \alpha^2} = \frac{\frac{3-\sqrt{3}}{2}\alpha^2}{\frac{4-2\sqrt{3}}{4}\alpha^2} = \frac{3-\sqrt{3}}{2-\sqrt{3}} \cdot \frac{2+\sqrt{3}}{2+\sqrt{3}}= 3+\sqrt{3}
        \]
    \end{solution}
    \question 已知 $a+b+c=6$,且
    \[
    \frac{1}{a}+\frac{1}{b}+\frac{1}{c}=2,
    \]
    求
    \[
    \frac{b+c}{a}+\frac{c+a}{b}+\frac{a+b}{c}
    \]
    的值。
    \begin{solution}
        不难发现
        \begin{align*}
        &\frac{b+c}{a}+\frac{c+a}{b}+\frac{a+b}{c} \\
        &=\left(\frac{b+c}{a}+1\right)+\left(\frac{c+a}{b}+1\right)+\left(\frac{a+b}{c}+1\right)-3 \\
        &=\frac{a+b+c}{a}+\frac{a+b+c}{b}+\frac{a+b+c}{c}-3 \\
        &=6\cdot2-3=9
        \end{align*}
    \end{solution}

    \question 已知 $a,b,c$ 皆为实数,若
    \[
    \frac{a}{b+c}+\frac{b}{c+a}+\frac{c}{a+b}=1
    \]
    求
    \[
    \frac{a^{2}}{b+c}+\frac{b^{2}}{c+a}+\frac{c^{2}}{a+b}
    \]
    的值。
    \begin{solution}
        有
        \[
        \frac{a(a+b+c)}{b+c}+\frac{b(a+b+c)}{c+a} +\frac{c(a+b+c)}{a+b}=a+b+c
        \]
        于是
        \[
        \frac{a^2}{b+c}+ a+\frac{b^2}{c+a}+ b+\frac{c^2}{a+b}+c=a+b+c
        \]
        即
        \[
        \frac{a^2}{b+c}+ \frac{b^2}{c+a}+\frac{c^2}{a+b}=0
        \]
    \end{solution}

    \question 已知 $abc=1$,证明
    \begin{parts}
    \part 
    \[
    \frac{a}{ab+a+1} + \frac{b}{bc+b+1} + \frac{c}{ca+c+1}=1
    \]
    \begin{solution}
        由$abc=1$,
        \begin{align*}       
        &\frac{a}{ab+a+1} + \frac{b}{bc+b+1} + \frac{c}{ca+c+1}\\
        &= \frac{a}{ab+a+1} + \frac{ab}{abc+ab+a} + \frac{abc}{abca+abc+ab} \\
        &= \frac{a}{ab+a+1} + \frac{ab}{1+ab+a} + \frac{1}{a+1+ab} \\
        &= \frac{a+ab+1}{ab+a+1} \\
        &= 1
        \end{align*}
    \end{solution}
    \part 
    \[
    \frac{1}{ab+a+1} + \frac{1}{bc+b+1} + \frac{1}{ca+c+1}=1
    \]
    \begin{solution}
        由(a),
        \begin{align*}
        &\frac{1}{ab+a+1} + \frac{1}{bc+b+1} + \frac{1}{ac+c+1} \\
        &=\frac{c}{abc+ac+c} + \frac{a}{abc+ab+a} + \frac{b}{abc+bc+b} \\
        &=\frac{c}{1+ac+c} + \frac{a}{1+ab+a} + \frac{b}{1+bc+b} \\
        &=1
        \end{align*}
    \end{solution}
    \end{parts}

    \question 已知 $a, b, c$ 为非零相异实数,且 $a + b + c = 0$,求
    \[ 
    \left(\frac{b - c}{a} + \frac{c - a}{b} + \frac{a - b}{c}\right) \left(\frac{a}{b - c} + \frac{b}{c - a} + \frac{c}{a - b}\right) 
    \]
    的值。
    \begin{solution}
        首先有
        \begin{align*}
        \frac{b - c}{a} + \frac{c - a}{b} + \frac{a - b}{c} &= \frac{b^2 c - b c^2 + a c^2 - a^2 c + a^2 b - a b^2}{abc} \\
        &= -\frac{(a - b)(b - c)(c - a)}{abc}
        \end{align*}
        现令 $a' = b - c,b' = c - a,c' = a - b$,发现
        \[ 
        b' - c' = (c - a) - (a - b) = b + c - 2a = -3a
        \]
        由此得 $a = -\dfrac{b' - c'}{3}$,同理有
        \[ 
        b = -\frac{c' - a'}{3}, \quad c = -\frac{a' - b'}{3} \]
        于是
        \begin{align*}
        \frac{a}{b - c} + \frac{b}{c - a} + \frac{c}{a - b} &= \frac{a}{a'} + \frac{b}{b'} + \frac{c}{c'} \\
        &= -\frac{1}{3} \left(\frac{b' - c'}{a'} + \frac{c' - a'}{b'} + \frac{a' - b'}{c'}\right) \\
        &= -\frac{1}{3} \left[ -\frac{(a' - b')(b' - c')(c' - a')}{a' b' c'} \right] \\
        &= \frac{1}{3} \frac{(-3c)(-3a)(-3b)}{(b - c)(c - a)(a - b)} \\
        &= -9 \frac{abc}{(a - b)(b - c)(c - a)}
        \end{align*}
        故所求的值为
        \begin{align*}
        &\left(\frac{b - c}{a} + \frac{c - a}{b} + \frac{a - b}{c}\right) \left(\frac{a}{b - c} + \frac{b}{c - a} + \frac{c}{a - b}\right) \\
        &=\left[ -\frac{(a - b)(b - c)(c - a)}{abc} \right] \left[ -9 \frac{abc}{(a - b)(b - c)(c - a)} \right] \\
        &= 9
        \end{align*}
    \end{solution}

    \question 已知 $x^5=1$ 且 $x \neq 1$, 求
    \[
    \frac{x}{1+x^2} + \frac{x^2}{1+x^4} + \frac{x^3}{1+x} + \frac{x^4}{1+x^3} 
    \]
    的值。
    \begin{solution}
        由$x^5=1$,
        \begin{align*}
        \frac{x}{1+x^2} + \frac{x^2}{1+x^4} + \frac{x^3}{1+x} + \frac{x^4}{1+x^3} 
        &= \frac{x}{1+x^2} + \frac{x^3}{x+1} + \frac{x^3}{1+x} + \frac{x}{x^2+1} \\
        &= 2 \left(\frac{x}{1+x^2} + \frac{x^3}{x+1}\right) \\
        &= 2 \left(\frac{x+x^2+x^3+x^5}{1+x+x^2+x^3}\right) \\
        &= 2
        \end{align*}
    \end{solution}

    \question 设 $a,b,c,d \in \mathbb{R}$ 且 $abcd \neq 0$,又 $a+b+c+d=0$,求
    \[
    S=a\left(\frac{1}{b}+\frac{1}{c}+\frac{1}{d}\right)+b\left(\frac{1}{c}+\frac{1}{d}+\frac{1}{a}\right)+c\left(\frac{1}{d}+\frac{1}{a}+\frac{1}{b}\right)+d\left(\frac{1}{a}+\frac{1}{b}+\frac{1}{c}\right)。
    \]
    的值。
    \begin{solution}
        由$a+b+c+d=0$,得
        \begin{align*}
        S &= a\left(\frac{1}{b}+\frac{1}{c}+\frac{1}{d}\right) + b\left(\frac{1}{c}+\frac{1}{d}+\frac{1}{a}\right) + c\left(\frac{1}{d}+\frac{1}{a}+\frac{1}{b}\right) + d\left(\frac{1}{a}+\frac{1}{b}+\frac{1}{c}\right) \\
        &= -(b+c+d)\left(\frac{1}{b}+\frac{1}{c}+\frac{1}{d}\right) - (a+c+d)\left(\frac{1}{c}+\frac{1}{d}+\frac{1}{a}\right) \\
        &\quad - (a+b+d)\left(\frac{1}{d}+\frac{1}{a}+\frac{1}{b}\right) - (a+b+c)\left(\frac{1}{a}+\frac{1}{b}+\frac{1}{c}\right) \\
        &= -12-\frac{2(b+c+d)}{a}-\frac{2(a+c+d)}{b} -\frac{2(a+b+d)}{c} -\frac{2(a+b+c)}{d}\\
        &= -12-\frac{-2a}{a} -\frac{-2b}{b} -\frac{-2c}{c} -\frac{-2d}{d}\\
        &= -4
        \end{align*}
    \end{solution}

    \question 设 $a,b,c$ 满足 $a+b+c=a^{3}+b^{3}+c^{3}=0$, $n$ 为任意实数, 求 $a^{2n+1}+b^{2n+1}+c^{2n+1}$的值。
    \begin{solution}
        由恒等式
        \[
        a^3 + b^3 + c^3 - 3abc = (a + b + c)(a^2 + b^2 + c^2 - ab - bc - ca)
        \]
        将已知代入得
        \[
        0 - 3abc = 0 \Rightarrow abc = 0.
        \]
        若$c = 0$,由 $a + b + c = 0$ 可得 $b = -a$,则
        \[
        a^{2n+1} + b^{2n+1} + c^{2n+1} = a^{2n+1} + (-a)^{2n+1} + 0  = 0
        \]
        同理若$a = 0$或$b = 0,\ a^{2n+1} + b^{2n+1} + c^{2n+1}=0$
    \end{solution}

\question 已知 $x, y, p, q$ 为实数,且满足 $2x^{2}+3p^{2}=2y^{2}+3q^{2}=(xq-yp)^{2}=6$,求 $$(x^{2}+y^{2})(p^{2}+q^{2})$$ 的值。

\begin{solution}
\textcolor{red}{(待另三种解)}
\end{solution}

    \begin{solution}
        \textbf{解法一}:
        
        设
        \[
        x = \sqrt{3} \cos\alpha, \quad p = \sqrt{2} \sin\alpha, \quad y = \sqrt{3} \cos\beta, \quad q = \sqrt{2} \sin\beta
        \]
        则
        \[
        6=(xq - yp)^2 = (\sqrt{6}(\cos\alpha \sin\beta - \cos\beta \sin\alpha))^2= 6\sin^2(\beta - \alpha)
        \]
        于是$\sin^2(\beta - \alpha) = 1 \Rightarrow \beta - \alpha = \pm \frac{\pi}{2} + 2k\pi,$
        因此有
        \[
        \cos\beta = \mp \sin\alpha, \quad \sin\beta = \pm \cos\alpha
        \]
        故
        \[
        (x^2 + y^2)(p^2 + q^2) = (\underbrace{3\cos^2\alpha + 3\cos^2\beta}_{=3(\cos^2\alpha + \sin^2\alpha) = 3})(\underbrace{2\sin^2\alpha + 2\sin^2\beta}_{=2(\sin^2\alpha + \cos^2\alpha) = 2})
        = 6
        \]
    \end{solution}

    \question 解方程
    \[
    (x+1)^6-2(x-1)^6 = (x^2-1)^3
    \]
    \begin{solution}
        原方程两边除以 $(x-1)^6$,得
        \[
        \left(\frac{x+1}{x-1}\right)^6 - 2 = \left(\frac{x+1}{x-1}\right)^3
        \]
        设 $y = \left(\dfrac{x+1}{x-1}\right)^3$,得到二次方程:
        \[
        y^2 - y - 2 = 0 \Rightarrow y = 2 \quad \text{或} \quad y=-1
        \]
        解为
        \[
        \frac{x+1}{x-1} = 2^{\frac{1}{3}} \Rightarrow x = \frac{2^{\frac{1}{3}}+1}{2^{\frac{1}{3}}-1} = 3+2^{\frac{2}{3}}+2^{\frac{1}{3}}
        \]
        或
        \[
        \frac{x+1}{x-1} = -1 \Rightarrow x = 0
        \]
    \end{solution}

    \question 若$a,b$为正实数,解方程
    \[
    (a+b)(ax+b)(a-bx)=(a^2x-b^2)(a+bx), \quad x\in \mathbb{R}.
    \]
    \begin{solution}
        首先设
        \[
        f(x)=(a+b)(ax+b)(a-bx)-(a^2x-b^2)(a+bx)
        \]
        不难发现$f(1)=0$,故$(x-1)$是$f(x)$的因式。正所谓:解决问题仍需暴力,展开并整理得
        \begin{align*}
        f(x) &= a^3x + a^2bx^2 - a^2b - b^3x - (a^3x - a^2bx^2 + a^2b - ab^2x + a^2bx - ab^2x^2 + ab^2 - b^3x) \\
        &= (2a^2b + ab^2)x^2 + (ab^2 - a^2b)x - (a^2b + ab^2) \\
        &= ab[(2a+b)x^2 + (b-a)x - (a+2b)] \\
        &= ab(x-1)((2a+b)x+a+2b)
        \end{align*}
        故$f(x)=0$给出了
        \[
        x = 1 \quad \text{或} \quad x = -\frac{a+2b}{2a+b}
        \]
    \end{solution}
\end{questions}
\pagebreak

\begin{center}
  {\fontsize{30pt}{26pt}\selectfont
    \hypertarget{因式定理、余式定理}{因式定理、余式定理} \label{因式定理、余式定理}
  }
\end{center}
\separator
\vspace{1pt}

\begin{questions}
    \question 设 $f(x) = a x^3 + b x^2 - 18 x + 3, g(x) = a x^3 + 9 x^2 + b x - 9$。当$f(x)$ 和 $g(x)$ 除以 $2x - 1$ 时所得余数相同,且当 $f(x)$ 除以 $x - 2$ 时余数是 $-5$。
    \begin{parts}
    \part 求 $a,\;b$ 的值。  
    \begin{solution}
        令 $x=\tfrac{1}{2}$,有:
        \[
        f\left(\frac{1}{2}\right) = g\left(\frac{1}{2}\right) \Rightarrow a\cdot\frac{1}{8} + b\cdot\frac{1}{4} - 9 + 3
         = a\cdot\frac{1}{8} + \frac{9}{4} + \frac{b}{2} - 9
         \Rightarrow b = 3
        \]
        又
        \[
        f(2) = 8a + 4b - 36 + 3 = -5
        \Rightarrow a = 2
        \]
        $\therefore a=2,\;b=3$
    \end{solution}
        
    \part 若 $x+4$ 整除 $f(x) + k g(x)$,求 $k$。
    \begin{solution}
        有\[
        f(-4) + k g(-4) = 0 \Rightarrow -128a + 16b + 72 + 3 + k(-128a + 144 - 4b - 9)=0
        \]
        代入 $a=2,\;b=3$ 得:
        \[
            -133(1 + k) = 0 \Rightarrow k = -1
        \]
    \end{solution}
    \end{parts}
    
    \question 已知 $f(x)$ 为三次多项式,以 $x^2 + x + 2$ 除之得余式 $x + 3$,以 $x^2 + x - 2$ 除之得余式 $5x + 7$,求 $f(x)$。
    \begin{solution}
        据题意有$f(1)=5(1)+7=12,f(-2)=5(-2)+7=-3,$
        设$$f(x)=(x^2 + x + 2)(ax+b)+x+3$$
        令$x=1,$ $$f(1)=(1+1+2)(a+b)+1+3=12 \Rightarrow a+b=2$$
        令$x=-2,$ $$f(-2)=(4-2+2)(-2a+b)-2+3=-3 \Rightarrow -2a+b=-1$$
        解得 $a=1,\;b=1$,所以$$f(x)=(x^2 + x + 2)(x+1)+x+3=x^3 + 2x^2 + 4x + 5$$
    \end{solution}
    
    \question 若 $f(x)$ 以 $x^2 - 1$ 除余 $3x + 2$,$g(x)$ 以 $x^2 + 2x - 3$ 除余 $5x + 2$,求 $(x+3) f(x) + (5x^2 + 1) g(x)$ 除以 $x - 1$ 的余式。
    \begin{solution}
        即求$(1+3)f(1)+(5+1)g(1)=4f(1)+6g(1)$的值,由已知
        $$f(1)=3\cdot1+2=5,\;g(1)=5\cdot1+2=7$$故$(x+3) f(x) + (5x^2 + 1) g(x)$ 除以 $x - 1$ 的余式为$$4\cdot5+6\cdot7=62$$
    \end{solution}
    
    \question 以 $x^2 + 2x + 3$ 除 $f(x)$ 余 $x + 12$,以 $(x+1)^2$ 除 $f(x)$ 余 $5x + 4$,求以 $(x+1)(x^2 + 2x + 3)$ 除 $f(x)$ 的余式。
    \begin{solution}
        同上,设$$f(x)=(x+1)(x^2+2x+3)Q(x)+a(x^2+2x+3)+x+12$$
        据题意$f(-1)=-5+4=-1$,现令$x=-1$有$$f(-1)=a(1-2+3)-1+12=-1 \Rightarrow a=-6$$

        $\therefore$ 以 $(x+1)(x^2 + 2x + 3)$ 除 $f(x)$ 的余式为 $-6(x^2 + 2x + 3)+x+12=-6x^2 - 11x - 6$    
    \end{solution}
    
    \question $f(x)$ 以 $x(x - 1)$ 除之,余式为 $-x + 3$；以 $x(x+1)$ 除之,余式为 $-3x + 3$,则 $f(x)$ 除以 $x(x^2 - 1)$ 的余式为?
    \begin{solution}
        设$$f(x)=x(x^2-1)Q(x)+ax(x-1)-x+3$$
        据题意$f(-1)=3+3=6$,令$x=-1$有$$f(-1)=a(-1)(-2)+1+3=6 \Rightarrow a=1$$

        $\therefore$ $f(x)$ 除以 $x(x^2 - 1)$ 的余式为 $x(x-1)-x+3=x^2 - 2x + 3$  
    \end{solution}
    
    \question 设多项式 $f(x)$ 除以 $(x+1)^3$ 得余式 $2x^2 + 8$,除以 $(x - 2)^2$ 得余式 $15x + 40$,且 $\deg f(x) \geq 4$,则 $f(x)$ 除以 $(x+1)^3 (x-2)$ 的余式为?
    \begin{solution}
        设$$f(x)=(x+1)^3 (x-2)Q(x)+a(x+1)^3+2x^2 + 8$$
        据题意$f(2)=15\cdot2+40=70$,令$x=2$有$$f(2)=a(2+1)^3+2\cdot2^2 + 8 =70\Rightarrow a=2$$

        $\therefore$ $f(x)$ 除以 $(x+1)^3 (x-2)$ 的余式为 $2(x+1)^3+2x^2 + 8=2x^3 + 8x^2 + 6x + 10$  
    \end{solution}

    
    \question 设 $\deg f(x) \geq 3$, 且 $f(x)$ 以 $(x-1)^2$ 除之余 $3x + 2$,以 $(x+2)^2$ 除之余 $5x - 3$,求:
        
    \begin{parts}
    \part 以 $x-1$ 除之的余式。
    \begin{solution}
        即 $f(1) = 3 \cdot 1 + 2 = 5$
    \end{solution}
        
    \part 以 $(x-1)(x+2)$ 除之的余式。
    \begin{solution}
        设$$f(x)=(x-1)(x+2)Q(x)+ax+b$$
        令$x=1,$ $$f(1)=a+b=5$$
        令$x=-2,$ $$f(-2)=-2a+b=-13$$ 
        解得 $a = 6,\ b = -1$,故余式为$6x-1$
    \end{solution}
        
    \part 以 $(x-1)^2(x+2)$ 除之的余式。
    \begin{solution}
        设$$f(x)=(x-1)^2(x+2)Q(x)+a(x-1)^2+3x+2$$
        令$x=-2,$ $$f(-2)=9a+-4=-13$$ 
        解得 $a = -1 $,故余式为$-x^2 + 5x + 1$
    \end{solution}
    \end{parts}
    
    \question $f(x)$ 为一多项式,若 $(x+1) f(x)$ 除以 $x^2 + x + 1$ 的余式为 $5x + 3$,则 $f(x)$ 除以$x^2 + x + 1$的余式为何？
    \begin{solution}
        据题意有$$(x+1)f(x)=(x+1)(x^2 + x + 1)Q(x)+a(x^2+x+1)+5x+3$$
        且设$$f(x)=(x^2 + x + 1)Q(x)+ax+b$$
        两边乘$x+1$有$$(x+1)f(x)=(x+1)(x^2 + x + 1)Q(x)+(x+1)(ax+b)$$
        比较得 $$a(x^2+x+1)+5x+3=(x+1)(ax+b)$$
        即 $$ax^2+(a+b)x+b=ax^2+(a+5)x+a+3 $$
        得知 $a=2,\;b=5$,余式为 $2x+5$
    \end{solution}
    
    \question 以 $(x+1)^3$ 除 $f(x)$ 余式为 $x^2 - 2x + 3$,求以 $(x+1)^2$ 除 $f(x)$ 所得余式。
    \begin{solution}
        设\[
        f(x)=Q(x)(x+1)^{3}+x^{2}-2x+3
        \]
        有\[
        f(x)=Q(x)(x+1)^{3}+(x+1)^{2}+(-4x+2) =(x+1)^{2}[Q(x)(x+1)+1]+(-4x+2)  
        \]
        轻而易举得知余式为$-4x + 2$.
    \end{solution}

    \question 设 $f(x)$ 为实系数三次多项式,若 $f(x)$ 除以 $x-2$ 的余式为 $-5$,且 $(x+1)f(x)$ 除以 $x^3-3$ 的余式为 $3x-1$,求多项式 $f(x)$ 
    \begin{solution}
        设 
        \[
        f(x) = a x^3 + b x^2 + c x + d,
        \]
        则
        \[
        (x+1) f(x) = a x^4 + (a + b) x^3 + (b + c) x^2 + (c + d) x + d.
        \]
        由题意,
        \[
        f(2) = 8a + 4b + 2c + d = -5 \tag{1}
        \]
        且有
        \[
        (x+1) f(x) = (x^3 - 3)(a x + (a + b)) + 3x - 1,
        \]
        比较系数得:
        \[
        b + c = 0 \tag{2}
        \]
        \[
        3a + 3b + d = -1 \tag{3}
        \]
        \[
        3a + c + d = 3 \tag{4}
        \]
        由$(1)-(4)$联立解得
        \[
        a = -1, \quad b = -1, \quad c = 1, \quad d = 5
        \]
        因此
        \[
        f(x)=-x^3-x^2+x+5
        \]
    \end{solution}

    \question 已知一多项式 $f(x) = x^{2025}(x^{2} + ax + b)$,其中 $a, b$ 为实数,如果将 $f(x)$ 除以 $(x - 2)^2$得余式 $2^{2025}(x - 2)$,求 $a, b$。
    \begin{solution}
        设
        \[
        f(x) = x^{2025}(x^2 + ax + b) = (x - 2)^2 p(x) + 2^{2025}(x - 2)
        \]
        对$x$求导有
        \[
        f'(x) = 2025x^{2024}(x^2 + ax + b) + x^{2025}(2x + a)
        = 2(x - 2)p(x) + (x - 2)^2 p'(x) + 2^{2025}
        \]
        现令$x=2$,则
        \[
        \begin{cases}
         \ 2^{2025}(4 + 2a + b) = 0  \\
         \ 2025 \cdot 2^{2024}(4 + 2a + b) + 2^{2025}(4 + a) = 2^{2025}
        \end{cases}
        \]
        解得
        \[
         a = -3,b = 2
        \]
    \end{solution}
    
    \question 设 $f(x) = x^4 + 5x^3 + a x^2 + b x + c$ 可被 $(x-1)^3$ 整除,求 $a,b,c$ 的值。 
    \ifprintanswers
    \begin{table}[H]
        \centering
        \begin{tabular}{llllll}
        1 & 5  & $a$    & $b$    & \multicolumn{1}{l|}{$c$}  & 1 \\
        & 1   & 6    & $a+6$   & \multicolumn{1}{l|}{$a+b+6$} &   \\ \cline{1-5}
        1 & 6    & $a$+6    & \multicolumn{1}{l|}{$a+b+6$} & $a+b+c+6$       &   \\
        & 1    & 7    & \multicolumn{1}{l|}{$a+13$}  &  &   \\ \cline{1-4}
        1 & 7   & \multicolumn{1}{l|}{$a+13$} & $2a+b+19$  &   &  \\
          & 1  & \multicolumn{1}{l|}{8}  &  &  &   \\ \cline{1-3}
        1 & \multicolumn{1}{l|}{8} & $a+21$ & & &   \\
          & \multicolumn{1}{l|}{1} &  &  &   &   \\ \cline{1-2}
        1 & 9  &  &    &    &  
        \end{tabular}
    \end{table}
    \fi
    \begin{solution}
        连续综合除法得\[
        f(x)=(x-1)^4+9(x-1)^3+(a+21)(x-1)^2+(2a+b+19)(x-1)+(a+b+c+6)
        \]
        已知 $f(x)$ 可被 $(x-1)^3$ 整除,则
        \[
        \begin{cases}
            a+21=0\\2a+b+19=0\\a+b+c+6=0
        \end{cases} \Rightarrow (a,b,c)=(-21,23,-8)
        \]
    \end{solution}

    \question 已知 $x^2 - x + b$ 为多项式 $6x^4 - 7x^3 + a x^2 + 3x + 2$ 的因式,求 $a,b$ 的值。
    \begin{solution}
        设 $f(x) = 6x^4 - 7x^3 + ax^2 + 3x + 2= (x^2 - x + b)(6x^2 + px + q)$,展开右边得
        \[
        (x^2 - x + b)(6x^2 + px + q)
        = 6x^4 + (p - 6)x^3 + (q - p + 6b)x^2 + (pb - q)x + bq
        \]
        比较系数得
        \begin{align}
        p - 6 &= -7 \tag{1} \\
        q - p + 6b &= a \tag{2} \\
        pb - q &= 3 \tag{3} \\
        bq &= 2 \tag{4}
        \end{align}
        由(3)得$q = -b - 3$, 代入(4)得
        \[
        b(-b - 3) = 2 
        \Rightarrow (b + 1)(b + 2) = 0
        \Rightarrow b = -1 \text{ 或 } -2
        \]
        当 $b = -1$ 时,$\;q = -2,a = -7$;当 $b = -2$ 时,$\;q = -1,a = - 12$
        
        $\therefore (a,b) = (-7, -1) \text{ 或 } \;(-12, -2)$
    \end{solution}
    
    \question 若 $a,b,c$ 为相异实数,分别用 $x - a, x - b, x - c$ 除多项式 $p(x)$,所得的余数分别为 $a, b, c$,求以 $(x - a)(x - b)(x - c)$ 除 $p(x)$ 的余式。 
    \begin{solution}
        设
        \[
        p(x) = (x-a)(x-b)(x-c)q(x) + r(x-a)(x-b)+s(x-a)+a,
        \]
        令 $x = b$, 由 $a \ne b$,
        \[
        p(b) = s(b-a)+a = b  \Rightarrow s(b-a) = b-a \Rightarrow s=1
        \]
        令 $x = c$,由 $a \ne c,b \ne c,$
        \[
        p(c) = r(c-a)(c-b)+(c-a)+a = c  \Rightarrow r(c-a)(c-b)=0 \Rightarrow r=0
        \]
        因此以 $(x - a)(x - b)(x - c)$ 除 $p(x)$ 的余式为$x$
    \end{solution}
    
    \question 若 $g(x)$ 除以 $2x - 3$ 的余式为 $1$,且 $f(x) = g(x)(2x - 3) + 5$,求 $[f(x)]^2$ 除以 $(2x - 3)^2$ 的余式。
    \begin{solution}
        设\[
        g(x)=(2x-3)Q(x)+1
        \]
        由$f(x) = g(x)(2x - 3) + 5$,
        \begin{align*}
        [f(x)]^2 
        &=[g(x)]^2(2x-3)^2+10g(x)(2x-3)+25 \\
        &=[g(x)]^2(2x-3)^2+10[(2x-3)Q(x)+1](2x-3)+25\\
        &=[g(x)]^2(2x-3)^2+10Q(x)(2x-3)^2+10(2x-3)+25\\
        &=[g(x)]^2(2x-3)^2+10Q(x)(2x-3)^2+20x-5
        \end{align*}
        故$[f(x)]^2$ 除以 $(2x - 3)^2$ 的余式为$20x - 5$
    \end{solution}
    
    \question 若多项式 $p(x) = x^2 + b x + c$,其中 $b,c$ 为实数,且 $p(p(1)) = p(p(2)) = 0$,且 $p(1) \neq p(2)$,求 $b,c$ 的值。 
    \begin{solution}
        由 \(p(p(1)) =p(p(2)) = 0\),意味 \(p(1)\)和\(p(2) \)是方程 \(x^2 + bx + c = 0\) 的两根,且
        \[
        p(1)+p(2) = -b,\quad p(1)p(2) = c
        \]
        即
        \[
         (1 + b + c) + (4+2b+c)=-b \tag{1}
        \]
        \[
         (1 + b + c)(4+2b+c)=c \tag{2}
        \]
        由(1)得$2b+c = -\dfrac{5}{2} $,代入(2):
        \[
        (1 - \frac{5}{2} - b)(4  - \frac{5}{2}) = -\frac{5}{2} - 2b \Rightarrow  b = -\frac{1}{2},c=-\frac{3}{2}
        \]
    \end{solution}
    
    \question $a,b,c$ 为整数,且 $0 < a < b$,若 $x - c$ 是多项式 $x(x-a)( x - b) - 17$ 的因式,求 $(a,b,c)$。 
    \begin{solution}
        已知$f(x) = x(x - a)(x - b) - 17$且 \(x - c\) 是它的因式,于是
        \[
        f(c) = c(c - a)(c - b) - 17 = 0 \Rightarrow c(c - a)(c - b) = 17.
        \]
        由于 \(a, b, c\) 均为整数,且 \(0 < a < b\),我们只需枚举整数三元组使得左边乘积为 17。
    
        因为 17 是质数,其非零整数因式只有:
        \[
        (\pm1, \pm1, \pm17) \text{ 的排列组合。}
        \]
    
        枚举 \(c = 1\),
        \[
        1(1 - a)(1 - b) = 17 \Rightarrow (1 - a)(1 - b) = 17.
        \]
    
        由于 \(17 = 1 \times 17 = (-1)\times(-17)\),列出可能组合:
        \begin{itemize}
            \item \(1 - a = 1,\ 1 - b = 17 \Rightarrow a = 0,\ b = -16\),不满足 \(0 < a < b\)。
            \item \(1 - a = -1,\ 1 - b = -17 \Rightarrow a = 2,\ b = 18\),满足条件！
        \end{itemize}
    
        若 \(c = -1\):
        \[
        (-1)(-1 - a)(-1 - b) = 17 \Rightarrow -(a + 1)(b + 1) = 17
        \Rightarrow (a + 1)(b + 1) = -17
        \]
        没有满足 \(0 < a < b\) 的整数解,尝试 \(c = 17, -17\) 会得到更大的数,不满足 \(0 < a < b\)。
    
        因此,唯一符合题意的解为:
        \[
        (a, b, c) = (2, 18, 1)
        \]
    \end{solution}

    \question 设 $f(x)$ 与 $g(x)$ 均为实系数二次多项式且首项系数都是 $2$。已知 $(f(x))^2$ 除以 $g(x)$ 的余式为 $5x+3$,而 $(g(x))^2$ 除以 $f(x)$ 的余式为 $x+1$,求 $f(x),g(x)$。
    \begin{solution}
        设
        \[
        f(x) = 2x^2 + a x + b, \quad g(x) = f(x) - c x - d = 2x^2 + (a-c)x + b-d.
        \]
        则
        \[
        (f(x))^2 = (g(x) + c x + d)^2 = (g(x))^2 + 2(c x + d) g(x) + (c x + d)^2.
        \]
        由已知 $(f(x))^2$ 除以 $g(x)$ 的余式为 $5x+3$,
        \[
        (c x + d)^2 = c^2 x^2 + 2 c d x + d^2 = \frac{c^2}{2} g(x) + \Big(2 c d - \frac{c^2}{2}(a-c)\Big)x + \Big(d^2 - \frac{c^2}{2}(b-d)\Big).
        \]
        比较系数得
        \[
        \begin{cases}
        4 c d - a c^2 + c^3 = 10,\\
        2 d^2 - b c^2 + c^2 d = 6.
        \end{cases}
        \]
        同理,由$(g(x))^2$ 除以 $f(x)$ 的余式为 $x+1$得
        \[
        \begin{cases}
        4 c d - a c^2 = 2,\\
        2 d^2 - b c^2 = 2.
        \end{cases}
        \]
        解得
        \[
        c^3 = 8 \Rightarrow c = 2, \quad c^2 d = 4 \Rightarrow d = 1, \quad a = \frac{3}{2}, \quad b = 0.
        \]
        于是
        \[
        f(x) = 2 x^2 + \frac{3}{2}x, \quad g(x) = 2 x^2 - \frac{1}{2}x - 1.
        \]
    \end{solution}

    \question 求多项式 $(x + 1)^6$ 除以 $x^2 + 1$ 的余式。
    \begin{solution}
        有$x^2 \equiv -1 \pmod{x^2+1}$,则
        $$(x+1)^6 = (x^2 + 2x + 1)^3 \equiv (2x)^3
        = 8x \cdot x^2 \equiv 8x \cdot (-1)
        = -8x \pmod{x^2+1}$$
    \end{solution}  
    
    \question 设 $f(x) = x^4 + x^3 + x^2 + x + 1$,求 $f(x^5)$ 除以 $f(x)$ 所得余式。 
    \begin{solution}
        首先注意到
        \[
        f(x) = x^4 + x^3 + x^2 + x + 1 = \frac{x^5 - 1}{x - 1}, \quad x \ne 1,
        \]
        所以 \( f(x) \mid x^5 - 1 \),即
        \[
        x^5 \equiv 1 \pmod{f(x)}
        \]
        故
        \[
        f(x^5) = (x^5)^4 + (x^5)^3 + (x^5)^2 + x^5 + 1 \equiv 1 + 1 + 1 + 1 + 1 = 5 \pmod{f(x)}
        \]
    \end{solution}

    \question 求以 $x^2 + 2x + 3$ 除 $(x^2 + 3x + 4)^4$ 所得的余式。
    \begin{solution}
    有
        \[
        (x^2 + 3x + 4)^4 \equiv (x+1)^4 = (x^2+2x+1)^2 \equiv (-2)^2 = 4 \pmod{x^2 + 2x + 3}
        \]
    \end{solution}

    
    \question 设 $f(x) = x^{37} - 2 x^{26} + 4 x^7 - 3$,则
    \begin{parts}
    \part 求$f(x)$ 除以 $x^2 + 1$ 的余式。
    \begin{solution}
        有
        \begin{align*}
        f(x)&=x^{37}-2x^{26}+4x^{7}-3 \\
        &=(x^{2})^{18} \cdot x - 2(x^{2})^{13} + 4(x^{2})^{3} \cdot x - 3 \\
        &=(-1)^{18} \cdot x - 2(-1)^{13} + 4(-1)^{3} \cdot x - 3 \\
        &=x+2-4x-3 \\
        &=-3x-1 \pmod{x^2 + 1}
        \end{align*}
    \end{solution}
    \part 求$f(x)$ 除以 $x^2 + x + 1$ 的余式。
    \begin{solution}
        有 $x^{3}\equiv1 \pmod{x^{2}+x+1}$,于是
        \begin{align*}
        f(x)&=x^{37}-2x^{26}+4x^{7}-3 \\
        &=(x^{3})^{12} \cdot x - 2(x^{3})^{8} \cdot x^{2} + 4(x^{3})^{2} \cdot x - 3 \\
        &\equiv x-2x^{2}+4x-3 \\
        &=-2x^{2}+5x-3 \\
        &\equiv-2(-x-1)+5x-3 \\
        &=7x-1 \pmod{x^{2}+x+1}
        \end{align*}
    \end{solution}  
    \end{parts}

    \question 求以 $x^4 - x$ 除 $x^{87} - 2 x^{44} - x^3 + 3 x^2 + 1$ 所得的余式。
    \begin{solution}
        有$x^4\equiv x \Rightarrow x^3\equiv 1\pmod{x^4-x}$,故
        \[
        x^{87} - 2 x^{44} - x^3 + 3 x^2 + 1
        \equiv 1 - 2 x^2 - 1 + 3 x^2 + 1
        = x^2 + 1      
        \]
    \end{solution}
                
    \question 求$x^{100}$ 除以 $x^3 + 2 x^2 + 2 x + 1$ 的余式。
    \begin{solution}
        $x^{100}$ 除以 $x+1$ 的余式即$(-1)^{100}=1$,又\[
        x^3\equiv1\pmod{x^2+x+1}
        \]故$x^{100}$ 除以 $x^2+x+1$ 的余式为 $x^{100} \equiv x\pmod{x^2+x+1}$,现设$$x^{100}=(x+1)(x^2+x+1)Q(x)+a(x^2+x+1)+x$$
        令$x=-1$有$$1=a(1-1+1)-1\Rightarrow a=2$$

        $\therefore$ $x^{100}$ 除以 $x^3 + 2 x^2 + 2 x + 1$ 的余式为 $2x^2+3x+2$ 
    \end{solution}
    
    \question 求$x^{200}$ 除以 $(x - 1)^2$ 的余式。
    \begin{solution}
        发现$f(x)$ 除以 $(x-1)$ 的余式为 $f(1)=1^{200}=1$ ,将所求写成
        \begin{align*}
        x^{200}&=(x-1)^{2}Q(x)+a(x-1)+1 \\
        x^{200}-1&=(x-1)^{2}Q(x)+a(x-1) \\
        (x-1)(x^{199}+\dots+x+1)&=(x-1)^{2}Q(x)+a(x-1) \\
        x^{199}+\dots+x+1&=(x-1)Q(x)+a
        \end{align*}
        两边代 $x=1$ ,得到
        $$200=0+a$$
        故余式为 $200(x-1)+1=200x-199$
    \end{solution}

    \begin{solution}
        设 $t=x-1$ ,则题目变成 $x^{200}=(t+1)^{200}$ 除以 $t^{2}$ 的余式。由二项式定理,
        $$(t+1)^{200}=\comb{200}{200}t^{200}+\dots+\comb{200}{2}t^{2}+\comb{200}{1}t+\comb{200}{0}$$
            则显然, $(t+1)^{200}$ 除以 $t^{2}$ 的余式即为 $\comb{200}{1}t+\comb{200}{0}=200(x-1)+1=200x-199$
    \end{solution}
    
    \question 已知多项式 $f(x)=x^{130}-1,g(x)=x^{4}-x^{3}+2x^{2}-x+1$,求 $f(x)$ 除以 $g(x)$ 的余式。
    \begin{solution}
        设
        \[
        x^{130}-1= (x^2 + 1)(x^2 - x + 1)Q(x) + Ax^3 + Bx^2 + Cx + D.
        \]
        令 $x = i$,则
        \[
        -1 - 1 = Ai^3 + Bi^2 + Ci + D = (-B + D) + (-A + C)i,
        \]
        得
        \[
        \begin{cases}
        -B + D = -2, \\
        -A + C = 0
        \end{cases}
        \]
        令 $\omega = \dfrac{1 + \sqrt{3}i}{2}$,
        \[
        \omega^{130}-1=-\omega -1=-A+B(\omega-1)+C\omega+D =(B+C)\omega-A-B+D
        \]
        得到方程组
        \[
        \begin{cases}
        B+C=-1 \\
        -A-B+D=-1
        \end{cases}
        \]
        解得$A = -1, B = 0, C = -1, D = -2,$因此余式为$-x^{3} - x - 2$
    \end{solution}

    \question 设 $f(x)$ 是一个次数有限的多项式,且满足
    \[
    (x + 9) f(x + 1) = (x + 3) f(x + 3),
    \]
    已知 $f(0) = 1$,求 $f(1)$ 的值。
    \begin{solution}
        令 $x = -3$ 得 $f(-2) = 0$,所以 $f(x)$ 有因式 $(x+2)$,设 $f(x) = (x+2)g(x)$得
        \[
        (x+9)(x+3)g(x+1) = (x+3)(x+5)g(x+3) \Rightarrow (x+9)g(x+1) = (x+5)g(x+3)
        \]
        同理,令 $x = -5$ 得 $g(-4) = 0$,故 $g(x)$ 有因式 $(x+4)$,设 $g(x) = (x+4)h(x)$得
        \[
        (x+9)(x+5)h(x+1) = (x+5)(x+7)h(x+3) \Rightarrow (x+9)h(x+1) = (x+7)h(x+3)
        \]
        令 $x = -7$ 得 $h(-6) = 0$,故 $h(x)$ 有因式 $(x+6)$,设 $h(x) = (x+6)p(x)$得
        \[
        (x+9)(x+7)p(x+1) = (x+7)(x+9)p(x+3) \Rightarrow p(x+1) = p(x+3)
        \]
        即$p(x)$是周期为2的函数,因为 $p(x)$ 是次数有限的多项式,故 $p(x)=c,c\in\mathbb{R}$ 
        
        于是$f(x) = c(x+2)(x+4)(x+6)$,由 $f(0) = 1$ 得$c = \frac{1}{48}$,故
        \[
        f(1) = \frac{1}{48}(3)(5)(7) =\frac{35}{16}
        \]
    \end{solution}
\end{questions}
\pagebreak

\begin{center}
  {\fontsize{30pt}{26pt}\selectfont
    \hypertarget{根式、绝对值}{根式、绝对值} \label{根式、绝对值}
  }
\end{center}
\separator
\vspace{1pt}
\begin{questions}
    \question 求
    \[
    \frac{\sqrt{10+\sqrt{1}} + \sqrt{10+\sqrt{2}} + \cdots + \sqrt{10+\sqrt{99}}}{\sqrt{10-\sqrt{1}} + \sqrt{10-\sqrt{2}} + \cdots + \sqrt{10-\sqrt{99}}}
    \]
    之值。
    \begin{solution}
        由
        \[
        \left(\sqrt{10+\sqrt{a}} - \sqrt{10-\sqrt{a}}\right)^2 = 20 - 2\sqrt{100 - a}
        \]
        于是有
        \[
        \sqrt{10+\sqrt{a}} - \sqrt{10-\sqrt{a}} = \sqrt{2} \cdot \sqrt{10 - \sqrt{100 - a}}
        \]
        令
        \[
        L = \frac{\sqrt{10+\sqrt{1}} + \sqrt{10+\sqrt{2}} + \cdots + \sqrt{10+\sqrt{99}}}{\sqrt{10-\sqrt{1}} + \sqrt{10-\sqrt{2}} + \cdots + \sqrt{10-\sqrt{99}}}
        \]
        则
        \[
        L - 1 = \frac{
        (\sqrt{10+\sqrt{1}} - \sqrt{10-\sqrt{1}}) + \cdots + (\sqrt{10+\sqrt{99}} - \sqrt{10-\sqrt{99}})
        }{
        \sqrt{10-\sqrt{1}} + \cdots + \sqrt{10-\sqrt{99}}
        }
        \]
        \[
        = \frac{
        \sqrt{2} \left(\sqrt{10 - \sqrt{99}} + \sqrt{10 - \sqrt{98}} + \cdots + \sqrt{10 - \sqrt{1}}\right)
        }{
        \sqrt{10 - \sqrt{1}} + \sqrt{10 - \sqrt{2}} + \cdots + \sqrt{10 - \sqrt{99}}
        } = \sqrt{2}
        \]
        因此
        \[
        L = \sqrt{2} + 1
        \]
    \end{solution}
    
    \question 求根式方程
    \[
    \sqrt{6x-9} + \sqrt{2x-5} = x-1
    \]
    的实数根。
    \begin{solution}
        原方程式两边平方得
        \[
        6x-9+2x-5+2\sqrt{(6x-9)(2x-5)}=x^2-2x+1
        \]
        看来并不是明智之举。考虑
        \[
        (\sqrt{6x-9} + \sqrt{2x-5})(\sqrt{6x-9} - \sqrt{2x-5}) = (6x-9) - (2x-5) = 4x-4
        \]
        由于$x \neq 1$,化为
        \[
        \sqrt{6x-9} - \sqrt{2x-5} = 4
        \]
        消去$\sqrt{2x-5}$可得
        \[
        \sqrt{6x-9}=\frac{1}{2}(x-1+4)
        \]
        两边平方,
        \[
        4(6x-9)=x^2+6x+9 \Rightarrow (x-3)(x-15) = 0
        \]
        经检验得$x=3$是增根,且$x=15$是原方程式的解。
    \end{solution}

    \question 解方程
    \[
    4+\sqrt{x^2-6x+13} = x+\sqrt{2x-5}
    \]
    \begin{solution}
        由
        \[
        \sqrt{x^2-6x+13} = x-4 + \sqrt{2x-5}
        \]
        设 $x-4 = y$,化简得
        \[
        \sqrt{y^2 + 2y + 5} = y + \sqrt{2y+3}
        \]
        两边平方,
        \[
        y^2 + 2y + 5 = y^2 + 2y\sqrt{2y+3} + 2y + 3 \Rightarrow
        1 = y\sqrt{2y+3}
        \]
        再次平方得
        \[
        2y^3 + 3y^2 - 1 = 0 \Rightarrow (y+1)(2y^2+y-1) = 0
        \]
        解得
        \[
        y = -1, \quad y = \frac{1}{2} \Rightarrow x = 3, \quad x = \frac{9}{2}
        \]
        经检验$x=3$是增根,原方程式的解为
        \[
        x=\frac{9}{2}
        \]
    \end{solution}

    \question 求所有实数 $x$ 满足
    \[
    x = \sqrt{x - \frac{1}{x}} + \sqrt{1 - \frac{1}{x}}.
    \]
    \begin{solution}
        原方程式变为
        \[
        x - \sqrt{x - \frac{1}{x}} = \sqrt{1 - \frac{1}{x}}
        \]
        两边平方得
        \[
        (x^2-1)-2\sqrt{x(x^2-1)}+x=0 \implies (\sqrt{x^2-1}-\sqrt{x})^2=0
        \]
        解得
        \[
        x^2 - x - 1 = 0 \Rightarrow x = \frac{1+\sqrt{5}}{2} >0
        \]
    \end{solution}

    \question 解方程
    \[
    \sqrt{7x+7}+\sqrt{7x-6}+2\sqrt{49x^2+7x-42}=181-14x
    \]
    \begin{solution}
        由原方程
        \[
        \sqrt{7x+7}+\sqrt{7x-6}+2\sqrt{(7x+7)(7x-6)}=181-14x
        \]
        设$ya=\sqrt{7x+7}+\sqrt{7x-6}$,则
        \[
        y^2=14x+1+2\sqrt{(7x+7)(7x-6)}
        \]
        代回原方程得
        \[
        y^2-(14x+1)+y=181-14x \Rightarrow (y+14)(y-13)=0
        \]
        由于平方根之和为非负,舍去 $y=-14$,于是由
        \[
        \sqrt{7x+7}+\sqrt{7x-6}=13
        \]
        两边平方,
        \[
        14x+1+2\sqrt{(7x+7)(7x-6)}=169 \Rightarrow \sqrt{49x^2-6}=84-7x
        \]
        再次平方,
        \[
        49x^2-6=7056-1176x+49x^2 \Rightarrow x=6
        \]
        经检验,原方程式的解为
        \[
        x=6
        \]
    \end{solution}

    \question 解
    \[
    x - \sqrt{\frac{x}{2} + \frac{7}{8} - \sqrt{\frac{x}{8} + \frac{13}{64}}} = 179
    \]
    \begin{solution}
        将两边乘以4得
        \[
        4x - \sqrt{8x + 14 - 2\sqrt{8x + 13}} = 716
        \]
        观察到
        \[
        8x + 14 - 2\sqrt{8x + 13} = (8x + 13) - 2\sqrt{8x + 13} + 1 = (\sqrt{8x + 13} - 1)^2
        \]
        且由原方程知$x\ge 179$,故$\sqrt{8x + 13} - 1>0$,方程变为
        \[
        4x - (\sqrt{8x + 13} - 1) = 716 \Rightarrow 8x+13-2\sqrt{8x+13}-1443=0
        \]
        因式分解得
        \[
        (\sqrt{8x + 13} - 39)(\sqrt{8x + 13} + 37) = 0
        \]
        由于 $\sqrt{8x + 13} \ge 0$,所以 
        \[
        \sqrt{8x + 13} = 39 \Rightarrow x = 188.5
        \]
        经检验,原方程式的解为
        \[
        x=188.5
        \]
    \end{solution}

    \question 解
    \[
    \frac{x-7}{2+\sqrt{x-3}} + \frac{x-5}{1+\sqrt{x-4}} = \sqrt{10}
    \]
    \begin{solution}
        首先有
        \[
        -\frac{4-(x-3)}{2+\sqrt{x-3}} - \frac{1-(x-4)}{1+\sqrt{x-4}} = \sqrt{10}
        \]
        于是
        \[
        -(2-\sqrt{x-3}) - (1-\sqrt{x-4}) = \sqrt{10} \Rightarrow \sqrt{x-3}-\sqrt{x-4} = \sqrt{10}+3
        \]
        两边移项得到
        \[
        \sqrt{x-3} -\sqrt{10} = 3 + \sqrt{x-4}
        \]
        两边平方,
        \[
        x - 3 + 10 - 2\sqrt{10(x - 3)} = 9 + x - 4 - 6\sqrt{x - 4} \Rightarrow \sqrt{10(x - 3)} - 1 = 3\sqrt{x - 4}
        \]
        再次平方,
        \[
        10(x - 3) + 1 - 2\sqrt{10(x - 3)} = 9(x - 4) \Rightarrow x + 7 = 2\sqrt{10(x - 3)}
        \]
        再再次平方,
        \[
        x^2 + 14x + 49 = 40(x - 3) \Rightarrow (x - 13)^2 = 0
        \]
        经检验,原方程式的解为
        \[
        x=13
        \]
    \end{solution}

    \question 
    \begin{parts}
    \part 已知相异非零实数 $a,b,c,d$ 满足
    \[
    \frac{a}{b}=\frac{c}{d},
    \]
    证明
    \[
    \frac{a+b}{a-b}=\frac{c+d}{c-d}
    \]
    \begin{solution}
        由$\dfrac{a}{b}=\dfrac{c}{d}$,得
        \[
        \frac{a}{b}+1=\frac{c}{d}+1,\quad
        \frac{a}{b}-1=\frac{c}{d}-1
        \]
        即
        \[
        \frac{a+b}{b}=\frac{c+d}{d},\quad
        \frac{a-b}{b}=\frac{c-d}{d}
        \]
        两式相除,得
        \[
        \frac{a+b}{a-b}=\frac{c+d}{c-d}
        \]
        证毕。
    \end{solution}
    \part 利用 (a) 或其他方法,解方程
    \[
    \frac{\sqrt{x+1}+\sqrt{x-1}}{\sqrt{x+1}-\sqrt{x-1}}=\frac{4x-1}{2}
    \]
    \begin{solution}
        由(a),原方程化为
        \[
        \frac{(\sqrt{x+1}+\sqrt{x-1})+(\sqrt{x+1}-\sqrt{x-1})}
        {(\sqrt{x+1}+\sqrt{x-1})-(\sqrt{x+1}-\sqrt{x-1})}
        =\frac{(4x-1)+2}{(4x-1)-2}
        \]
        即
        \[
        \frac{2\sqrt{x+1}}{2\sqrt{x-1}}=\frac{4x+1}{4x-3}
        \]
        两边平方,
        \[
        \frac{x+1}{x-1}=\frac{(4x+1)^2}{(4x-3)^2}
        =\frac{16x^2+8x+1}{16x^2-24x+9}.
        \]
        再由(a)得
        \[
        \frac{(x+1)+(x-1)}{(x+1)-(x-1)}
        =
        \frac{(16x^2+8x+1)+(16x^2-24x+9)}
        {(16x^2+8x+1)-(16x^2-24x+9)}.
        \]
        化简得
        \[
        \frac{2x}{2}=\frac{32x^2-16x+10}{32x-8} \Rightarrow x=\frac{5}{4}
        \]
        经检验,原方程的解为
        \[
        x=\frac{5}{4}
        \]
    \end{solution}
    \end{parts}

    \question 解方程
    \[
    \sqrt[3]{x} + \sqrt[3]{2x-3} = \sqrt[3]{12(x-1)}.
    \]
    \begin{solution}
        两边立方得
        \[
        x+2x-3+3\sqrt[3]{x(2x-3)}(\sqrt[3]{x} + \sqrt[3]{2x-3})=12(x-1)
        \]
        由原方程$\sqrt[3]{x} + \sqrt[3]{2x-3} = \sqrt[3]{12(x-1)}$代入得
        \[
        \sqrt[3]{x(2x-3)12(x-1)} = 3(x-1)
        \]
        两边再立方,
        \[
        (x-1)(x-3)^2=0 \Rightarrow x=1,3
        \]
        经检验,原方程式的解为
        \[
        x = 1 \quad \text{或} \quad x = 3
        \]
    \end{solution}

    \question 证明
    \[
    \sqrt[3]{9+4\sqrt{5}}+\sqrt[3]{9-4\sqrt{5}}=3
    \]
    \begin{solution}
        设
        \[
        x=\sqrt[3]{9+4\sqrt{5}}+\sqrt[3]{9-4\sqrt{5}}
        \]
        完全立方展开得
        \[
        x^3=9+4\sqrt{5}+9-4\sqrt{5}+3\sqrt[3]{81-80}(x)
        \]
        即解得
        \[
        x^3-3x-18=0 \Rightarrow (x-3)(x^2 + 3x + 6) = 0 \Rightarrow x=3 \in \mathbb{R}
        \]
        因此得证
        \[
        \sqrt[3]{9+4\sqrt{5}}+\sqrt[3]{9-4\sqrt{5}}=3
        \]
    \end{solution}

    \question 解方程
    \[
    \sqrt{x+\sqrt{x}} - \sqrt{x-\sqrt{x}} = \frac{199}{100} \sqrt{\frac{x}{x+\sqrt{x}}}.
    \]
    \begin{solution}
        设 $a=x+\sqrt{x}, b=x-\sqrt{x}$,则 $a-b=\sqrt{a}-\sqrt{b}=2\sqrt{x}$,原方程化为
        \[
        \sqrt{a}+\sqrt{b} = \frac{200}{199}\sqrt{x+\sqrt{x}}, \quad
        \sqrt{a}-\sqrt{b} = \frac{199}{100}\frac{\sqrt{x}}{\sqrt{x+\sqrt{x}}}
        \]
        两式相加得
        \[
        2\sqrt{x+\sqrt{x}}= \frac{200}{199}\sqrt{x+\sqrt{x}}+\frac{199}{100}\frac{\sqrt{x}}{\sqrt{x+\sqrt{x}}}
        \]
        解得
        \[
        \sqrt{x}(19800\sqrt{x}-19801)=0 \Rightarrow \sqrt{x} = \frac{19801}{19800}>0 \Rightarrow x = \frac{19801^2}{19800^2}
        \]
    \end{solution}

    \question 设 $a,b$ 为正实数,且
    \[
    \frac{1}{a}+\frac{1}{b}=1,\quad 2022a^2=2023b^2,
    \]
    求 $\sqrt{2022a+2023b}$ 的值。
    \begin{solution}
        由 $2022a^2=2023b^2$ 得
        \[
        \frac{b}{a}=\sqrt{\frac{2022}{2023}}.
        \]
        记 $k=\sqrt{\dfrac{2022}{2023}}$,则 $b=ka$。由 $\dfrac{1}{a}+\dfrac{1}{b}=1$ 得
        \[
        \frac{1}{a}+\frac{1}{ka}=1 \Rightarrow a=1+\frac{1}{k}.
        \]
        于是
        \begin{align*}
        2022a+2023b
        &=2022a+2023ka\\
        &=a(2022+2023k)\\
        &=\left(1+\frac{1}{k}\right)(2022+2023k)\\
        &=2022+2023+2022\cdot\frac{1}{k}+2023k\\
        &=2022+2023+2022\sqrt{\frac{2023}{2022}}+2023\sqrt{\frac{2022}{2023}}\\
        &=2022+2023+\sqrt{2022\cdot2023}+\sqrt{2022\cdot2023}\\
        &=(\sqrt{2022}+\sqrt{2023})^2
        \end{align*}
        故
        \[
        \sqrt{2022a+2023b}=\sqrt{2022}+\sqrt{2023}.
        \]
    \end{solution}

    \question 若
    \[
    (x-\sqrt{x^2-2011})(y+\sqrt{y^2-2011})+2011=0,
    \]
    求$2x+y$的值。
    \begin{solution}
        原式两边乘于$x+\sqrt{x^2-2011}$可得
        \[
        2011(y+ \sqrt{y^2-2011})+2011(x+\sqrt{x^2-2011})=0
        \]
        \[
        2011(x+y+\sqrt{x^2-2011} +\sqrt{y^2-2011})=0
        \]
        即
        \[
        x+y= -(\sqrt{x^2-2011} +\sqrt{y^2-2011}) \tag{1}
        \]
        同理,原式两边乘于$y-\sqrt{y^2-2011}$可得 
        \[
        (x-\sqrt{x^2-2011})2011+ 2011(y-\sqrt{y^2-2011})=0 
        \]
        \[
        2011(x+y-(\sqrt{x^2-2011}+ \sqrt{y^2-2011})) =0 
        \]
        \[
        x+y=\sqrt{x^2-2011}+ \sqrt{y^2-2011} \tag{2}
        \]
        由(1)及(2)可得
        \[
        \sqrt{x^2-2011}+ \sqrt{y^2-2011}=0 \Rightarrow x=\pm \sqrt{2011},y=\mp \sqrt{2011}\quad (\,\because x+y=0)
        \]
        故\[
        2x+y=x+y+x= \pm \sqrt{2011}
        \]
    \end{solution}

    \question 求 
    \[
    \left|\sqrt{x + 1} - 2\right| + \left|\sqrt{x + 1} - 3\right| = 1
    \]
    的整数解集。
    \begin{solution}
        令 $t = \sqrt{x+1}$, 其中 $t \geq 0$, 则方程变为
        \[
        |t-2| + |t-3| = 1
        \]
        分三种情况讨论:

        情况1:当 $0 \leq t < 2$ 时,
        \[
        (2-t) + (3-t) = 1 \Rightarrow t = 2
        \]
        这与 $t < 2$ 矛盾, 无解。

        情况2:当 $2 \leq t \leq 3$ 时,
        \[
        (t-2) + (3-t) = 1 \Rightarrow 1 = 1
        \]
        恒成立, 因此 $t \in [2, 3]$。

        情况3:当 $t > 3$ 时,
        \[
        (t-2) + (t-3) = 1 \Rightarrow t = 3
        \]
        这与 $t > 3$ 矛盾, 无解。

        综合得$2 \leq \sqrt{x+1} \leq 3$,平方得
        \[
        4 \leq x+1 \leq 9 \Rightarrow 3 \leq x \leq 8
        \]
        因此整数解集为 $\{3, 4, 5, 6, 7, 8\}$。
    \end{solution}

    \question 解
    \[ 
    |x + 1|-|x|+3|x-1|-2|x-2|=x+2 
    \]
    \begin{solution}
        分情况讨论:

        情况1:当 $x < -1$ 时,
        \[
        -(x+1) + x - 3(x-1) + 2(x-2) = x+2 \Rightarrow x = -2
        \]
        得一解$x = -2$。

        情况2:当 $-1 \le x < 0$ 时,
        \[
        (x+1) + x - 3(x-1) + 2(x-2) = x+2 \Rightarrow x = x+2
        \]
        故无解。
        
        情况3:当 $0 \le x < 1$ 时,
        \[
        (x+1) - x - 3(x-1) + 2(x-2) = x+2 \Rightarrow x = -1 \]
        这与$0 \le x < 1$矛盾,故无解。

        情况4:当 $1 \le x < 2$ 时,
        \[
        (x+1) - x + 3(x-1) + 2(x-2) = x+2 \Rightarrow x = 2 \]
        这与$1 \le x < 2$矛盾,故无解。

        情况4:当 $x \ge 2$ 时,
        \[ 
        (x+1) - x + 3(x-1) - 2(x-2) = x+2 \Rightarrow x = x
        \]
        恒成立,因此$x \ge 2$都是解。

        结论:原方程的解为 $x = -2$ 以及所有满足 $x \ge 2$ 的实数。
    \end{solution}

\end{questions}
\pagebreak

\begin{center}
  {\fontsize{30pt}{26pt}\selectfont
    \hypertarget{指数与对数}{指数与对数} \label{指数与对数}
  }
\end{center}
\separator
\vspace{1pt}

\begin{questions}
    \question 设 \( a=2019^{2017}, b=2019^{2018}, c=2019^{2019} \),计算
    \[
    \frac{1}{2019^{a - a} + 2019^{a - b} + 2019^{a - c}} + \frac{1}{2019^{b - a} + 2019^{b - b} + 2019^{b - c}} + \frac{1}{2019^{c - a} + 2019^{c - b} + 2019^{c - c}} 
    \]
    \begin{solution}
        观察到原式为
        \[
        \frac{2019^{-a}}{2019^{- a} + 2019^{- b} + 2019^{- c}} + \frac{2019^{-b}}{2019^{- a} + 2019^{- b} + 2019^{- c}} + \frac{2019^{-c}}{2019^{- a} + 2019^{- b} + 2019^{- c}} 
        \]
        答案呼之欲出$\Rightarrow1$
    \end{solution}          

    \question 试证
    \[
    2^{\sqrt{\log_2 3}} = 3^{\sqrt{\log_3 2}}
    \]
    \begin{solution}
        注意到
        \[
        2^{\sqrt{\log_2 3}} 
        = 2^{\frac{\log_2 3}{\sqrt{\log_2 3}}} 
        = \left(2^{\log_2 3}\right)^{\sqrt{\log_3 2}} 
        = 3^{\sqrt{\log_3 2}}
        \]
        故左式等于右式,得证。
    \end{solution}
    \begin{solution}
        设$x=2^{\sqrt{\log_2 3}},y=3^{\sqrt{\log_3 2}}$,则\[
        \log x=\sqrt{\log_2 3}\cdot \log 2 = \sqrt{\frac{\log 3}{\log 2}}\cdot\log 2 = \sqrt{\log 2 \log 3}
        \]同理\[
        \log y=\sqrt{\log_3 2} \cdot\log 3 = \sqrt{\frac{\log 2}{\log 3}}\cdot\log 3 = \sqrt{\log 2 \log 3}
        \]故\[\log x = \log y \Rightarrow x=y \quad\text{（得证）}\]
    \end{solution}

    \question 解方程
    \[
    2^{x+2}5^{6-x} = 10^{x^{2}},
    \]
    \begin{solution}
        两边取$\log$得\[
        (x+2)\log 2+(6-x)\log5=x^2 \Rightarrow x^2-\left(\log \frac{2}{5}\right)x-2\log 250=0
        \]
        因式分解给出\[
        (x-2)(x+\log 250)=0 \Rightarrow x = 2, -\log 250
        \]
    \end{solution}

    \question 求解方程
    \[
    \frac{e^{2x}+16^x}{(4e)^x} = \frac{4+e}{2\sqrt{e}}.
    \]
    \begin{solution}
        原方程即
        \[
        \left(\frac{e}{4}\right)^x + \left(\frac{4}{e}\right)^x = \frac{\sqrt{e}}{2} + \frac{2}{\sqrt{e}}
        \]
        令
        \[
        y = \left(\frac{e}{4}\right)^x
        \]
        则由
        \[
        y^2 - \left(\left(\frac{e}{4}\right)^{\frac{1}{2}} + \left(\frac{e}{4}\right)^{-\frac{1}{2}}\right)y + 1 = 0
        \]
        因式分解得
        \[
        \left(y - \left(\frac{e}{4}\right)^{\frac{1}{2}}\right)\left(y - \left(\frac{e}{4}\right)^{-\frac{1}{2}}\right) = 0
        \]
        解得
        \[
        x = \frac{1}{2}, \quad x = -\frac{1}{2}
        \]
    \end{solution}

    \question 解 
    \[
    x^{x^{x}} = (x^{x})^{x}
    \]
    \begin{solution}
        原方程即
        \[
        x^{x^{x}} = x^{x^2}
        \]
        先考虑指数方程的特殊情况$x=-1,0,1$,可得解
        \[
        x=-1,1
        \]
        当$x \neq -1,0,1$时,由底数相等所以指数相等的性质得,
        \[
        x=2
        \]
        故原方程式的解为
        \[
        x=-1 \quad \text{或} \quad x=1 \quad \text{或} \quad x=2
        \]
    \end{solution}

    \question 已知 
    \[
    (x\sqrt{x}\sqrt[3]{x})^{x} = x^{x\sqrt{x}\sqrt[3]{x}},
    \]
    试求 $x^{5}$ 的值。
    \begin{solution}
        \begin{align*}
        (x\sqrt{x}\sqrt[3]{x})^{x} &= x^{x\sqrt{x}\sqrt[3]{x}} \\
        x^{(1+\frac{1}{2}+\frac{1}{3})x} &= x^{x^{1+\frac{1}{2}+\frac{1}{3}}}\\
        x^{(\frac{11}{6})x} &= x^{x^{\frac{11}{6}}}
        \end{align*}
        考虑特殊情况 $x=-1, 0, 1,$ 可得解 $x=1$,此时因底数相等所以指数相等,得
        \[
        \frac{11}{6}x = x^{\frac{11}{6}} \Rightarrow x^5 = \left(\frac{11}{6}\right)^{6} 
        \]
        故$x^5$的可能值为
        \[
        x^5=1 \quad \text{或} \quad x^5 = \left(\frac{11}{6}\right)^{6}
        \]
    \end{solution}

    \question 求满足方程 
    \[
    (x^{2}+2x)^{x^{2}-3x+2}=1
    \] 
    的实数解。
    \begin{solution}
        设 $A=x^{2}+2x$, $B=x^{2}-3x+2$。欲使 $A^{B}=1$,仅以下三种情况成立:
        \begin{itemize}
        \item $B=0$:
        \[
        x^{2}-3x+2=0 \;\;\Rightarrow\;\; (x-1)(x-2)=0\Rightarrow\;\;x=1,2
        \]
        \item $A=1$:
        \[
        x^{2}+2x=1 \;\;\Rightarrow\;\; x=-1\pm\sqrt{2}
        \]
        \item $A=-1$ 且 $B$ 为偶数: 
        \[
        x^{2}+2x=-1 \;\;\Rightarrow\;\; (x+1)^{2}=0 \;\;\Rightarrow\;\; x=-1
        \]  
        此时  
        \[
        B=(-1)^{2}-3(-1)+2=6
        \]  
        为偶数,成立。
        \end{itemize}
        综上,解为  
        \[
        x\in \{1,2,-1+\sqrt{2},-1-\sqrt{2},-1\}
        \]
    \end{solution}

    \question 计算
    \[
    5^{(\log 2)^3} \cdot 8^{(\log 5)^2} \cdot 5^{(\log 5)^3}
    \]
    \begin{solution}
        令$a=\log 2,b=\log 5$,则$a+b=1$,原式为
        \[
        L=5^{a^3}\cdot 8^{b^2} \cdot 5^{b^3}
        \]
        两边取对数,则
        \[
        \log L 
        = (a^3 + b^3)b + 3ab^2 
        = (a+b)(a^2-ab+b^2) b + 3ab^2 
        = b(a+b)^2 
        = \log 5
        \]
        因此
        \[
        L = 5
        \]
    \end{solution}

    \question 已知正实数 $x$ 满足
    \[
    \log_2 x\log_4 x\log_6 x = \log_2 x\log_4 x+\log_2 x\log_6 x+\log_4 x\log_6 x,
    \]
    求 $x$。
    \begin{solution}
        设$\log_2 x=a, \log_6 x=b$,则$\log_4 x=\dfrac{1}{2}a$,原式为
        \[
        \frac{1}{2}a^2 b = \frac{1}{2}a^2 + ab + \frac{1}{2}ab
        \Rightarrow ab = a + 3b
        \]
        将 $a=\dfrac{\log x}{\log 2},b=\dfrac{\log x}{\log 6}$ 代入得
        \[
        \frac{(\log x)^2}{(\log 2)(\log 6)}
        = \frac{\log x(\log 6+3\log 2)}{(\log 2)(\log 6)} \Rightarrow \log x(\log x-\log 48)=0
        \]
        所以
        \[
        x=1\quad\text{或}\quad x=48
        \]
    \end{solution}

    \question 求
    \[
    \sqrt{\log_3 \sqrt{6} + \sqrt{\log_3 2}} + \sqrt{\log_3 \sqrt{6} - \sqrt{\log_3 2}}
    \]
    之值。
    \begin{solution}
        令
        \[
        a = \log_3 \sqrt{6} = \frac{1}{2} (1 + \log_3 2), \quad 
        b = \sqrt{\log_3 2} < 1
        \]
        则有
        \[
        a = \frac{b^2 + 1}{2}
        \]
        因此
        \begin{align*}
        \sqrt{a+b} + \sqrt{a-b} &= \sqrt{\frac{b^2+1}{2} + b} + \sqrt{\frac{b^2+1}{2} - b} \\
        &= \sqrt{\frac{1}{2} (b+1)^2} + \sqrt{\frac{1}{2} (b-1)^2} \\
        &= \frac{1}{\sqrt{2}} (b+1) + \frac{1}{\sqrt{2}} (1-b) = \sqrt{2}
        \end{align*}
    \end{solution}

    \question 解不等式
    \[
    \sqrt{\log_2 x - 1} + \frac{1}{2} \log_{\frac{1}{2}} x^3 + 2 > 0
    \]
    \begin{solution}
        设 \( y = \log_2 x \),则不等式化为
        \[
        \sqrt{y - 1} - \frac{3}{2} y + 2 > 0, \quad \text{且} \;y \ge 1.
        \]
        设$z=\sqrt{y - 1}$,则$y=z^2+1$,方程变为
        \[
        z - \frac{3}{2} (z^2+1) + 2 > 0 \Rightarrow (3z+1)(z-1)<0 
        \]
        于是
        \[
        -\frac13<z<1 \Rightarrow 0\leq z^2<1 \Rightarrow1\leq y<2 
        \]
        再考虑$y$的定义域$y\geq1$,经检验有
        \[
        2\leq x <4
        \]
    \end{solution}
    
    \question 求不等式的解集
    \[
    \log (x - 40) + \log (60 - x) < 2
    \]
    \begin{solution}
        变为
        \[
            (x - 40)(60 - x) < 100
        \]
        注意到\((x - 40)(60 - x) = -(x - 50)^2 + 100\),所以不等式等价于
        \[
        -(x - 50)^2 + 100 < 100 \Rightarrow (x - 50)^2 > 0\Rightarrow x \ne 50
        \]
        且由定义域\[
        \begin{cases}
        x - 40 > 0 \\
        60 - x > 0
        \end{cases} \Rightarrow 40<x<60
        \]经检验得解集
        \[
        (40,\;50)\cup(50,\;60).
        \]
    \end{solution}

    \question 解不等式:
    \[
    \left(\log_{\frac{1}{3}} x - 1\right)\left(\log_{\frac{1}{4}} x + 2\right)\left(\log_{\frac{1}{5}} x - 3\right) > 0.
    \]
    \begin{solution}
    换底得
    \[
    \left(\frac{\log x}{-\log 3} - 1\right)\left(\frac{\log x}{-\log 4} + 2\right)\left(\frac{\log x}{-\log 5} - 3\right) > 0
    \]
    即
    \[
    (\log x + \log 3)(\log x - 2\log 4)(\log x + 3\log 5) < 0
    \]
    符号分析得
    \[
    -\log 3 < \log x < 2\log 4, \quad \log x < -3\log 5
    \]
    解集为
    \[
    \frac{1}{3} < x < 16 \quad \text{或} \quad 0 < x < \frac{1}{125}
    \]
    \end{solution}

    \question 求 
    \[
    |x - 1|^{\log_{2}(4 - x)} < |x - 1|^{\log_{2}(1 + x)}
    \]
    的解之范围。
    \begin{solution}
        由$\log_{2}(4 - x)$及$\log_{2}(1 + x)$的定义域,得
        \[  
        x < 4 \quad \text{且} \quad x > -1 
        \]
        即  $x \in (-1, 4)$,且$x \neq 0,1,2$,否则不等式不成立。分四种情况讨论:

        情况1:$x \in (-1, 0)$,则
        \begin{align*}
        (1-x)^{\log_2 (4-x)} < (1-x)^{\log_2 (1+x)} & \Rightarrow \log_2 (4-x) < \log_2 (1+x) \\
        & \Rightarrow 4-x < 1+x \\
        & \Rightarrow x > \frac{3}{2}
        \end{align*}
        此情况解集为 $(-1, 0) \cap \left(\dfrac{3}{2}, \infty\right) = \emptyset$,无解。

        情况2:$x \in (0, 1)$,则
        \begin{align*}
        (1-x)^{\log_2 (4-x)} < (1-x)^{\log_2 (1+x)} & \Rightarrow \log_2 (4-x) > \log_2 (1+x) \\
        & \Rightarrow 4-x > 1+x \\
        & \Rightarrow x < \frac{3}{2} 
        \end{align*}
        此情况解集为 $(0, 1) \cap \left(-\infty, \dfrac{3}{2}\right) = (0, 1)$。

        情况3:$x \in (1, 2)$,则
        \begin{align*}
        (x-1)^{\log_2 (4-x)} < (x-1)^{\log_2 (1+x)} & \Rightarrow \log_2 (4-x) > \log_2 (1+x) \\
        & \Rightarrow 4-x > 1+x \\
        & \Rightarrow x < \frac{3}{2}
        \end{align*}
        此情况解集为 $(1, 2) \cap \left(-\infty, \dfrac{3}{2}\right) = \left(1, \dfrac{3}{2}\right)$。

        情况4:$x \in (2, 4)$,则
        \begin{align*}
        (x-1)^{\log_2 (4-x)} < (x-1)^{\log_2 (1+x)} & \Rightarrow \log_2 (4-x) < \log_2 (1+x) \\
        & \Rightarrow 4-x < 1+x \\
        & \Rightarrow x > \frac{3}{2}
        \end{align*}
        此情况解集为 $(2, 4) \cap \left(\dfrac{3}{2},\infty \right) = (2, 4)$。

        综上所述,原方程式的解集为 $\left(0, \dfrac{3}{2}\right) \cup (2, 4) \setminus \{1\}$。
    \end{solution}

    \question 求所有实数 $x$,使得
    \[
    \sqrt{\log_2 x \cdot \log_2 (4x) + 1} + \sqrt{\log_2 x \cdot \log_2 \left(\frac{x}{64}\right) + 9} = 4
    \]
    \begin{solution}
        发现
        \begin{align*}
        f(x) 
        &= \sqrt{\log_2 x \cdot \log_2(4x) + 1} + \sqrt{\log_2 x \cdot \log_2\left(\frac{x}{64}\right) + 9} \\
        &= \sqrt{\log_2 x (2 + \log_2 x) + 1} + \sqrt{\log_2 x (\log_2 x - 6) + 9} \\
        &= \sqrt{(\log_2 x + 1)^2} + \sqrt{(\log_2 x - 3)^2} \\
        &= 
        \begin{cases}
        \begin{aligned}
        &-2\log_2 x+ 2&&, \log_2 x \le -1 \\[1mm]
        &4&&, -1 < \log_2 x \le 3 \\[1mm]
        &2\log_2 x- 2&&, \log_2 x > 3
        \end{aligned}
        \end{cases}
        \end{align*}
        欲使$f(x)=4$,当 $\log_2 x \le -1,\log_2x = -1$;
        当 $-1 < \log_2x \le 3,f(x) = 4$ 恒成立;
        当 $\log_2x > 3$,无解;因此解为
        \[
        -1 \le \log_2x \le 3 \Rightarrow \frac{1}{2} \le x \le 8
        \]
    \end{solution}

    \question 已知 
    \[
    \frac{\log a}{b - c} = \frac{\log b}{c - a} = \frac{\log c}{a - b},
    \]
    其中 $a \neq b \neq c$, 试求 $a^{a}b^{b}c^{c}$ 的值。
    \begin{solution}
        设
        \[
        \frac{\log a}{b-c} = \frac{\log b}{c-a} = \frac{\log c}{a-b} = k
        \]
        则
        \[
        \log a = k(b-c), \quad \log b = k(c-a), \quad \log c = k(a-b)
        \]
        因此
        \begin{align*}
        \log(a^a b^b c^c) &= a\log a + b\log b + c\log c \\
        &= ak(b-c) + bk(c-a) + ck(a-b) \\
        &= k[ab - ac + bc - ab + ca - cb] \\
        &= 0
        \end{align*}
        即 
        \[
        a^a b^b c^c = 1
        \]
    \end{solution}

    \question 设 $a,b$同号,且 $a^2 - 2ab - 9b^2 = 0$,求
    \[
    \log(a^2 + ab - 6b^2) - \log(a^2 + 4ab + 15b^2)
    \]
    的值。
    \begin{solution}
        由 $a^2 - 2ab - 9b^2 = 0 $解得
        $$a = (\sqrt{10} + 1)b$$
        故
        \[
        \begin{aligned}
        &\log(a^2 + ab - 6b^2) - \log(a^2 + 4ab + 15b^2) \\
        &= \log(3ab + 3b^2) - \log(6ab + 24b^2) \\
        &= \log\left(\frac{3b(a + b)}{6b(a + 4b)}\right)\\
        &= \log\left(\frac{1}{2} \cdot \frac{a + b}{a + 4b}\right) \\
        &= \log\left(\frac{1}{2} \cdot \frac{(\sqrt{10} + 2)b}{(\sqrt{10} + 5)b}\right) \\
        &= \log\left(\frac{1}{2} \cdot \frac{(\sqrt{10} + 2)(5 - \sqrt{10})}{(\sqrt{10} + 5)(5 - \sqrt{10})}\right) \\
        &= \log\left(\frac{1}{2} \cdot \frac{3\sqrt{10}}{15}\right)= \log\left(\frac{\sqrt{10}}{10}\right) = -\frac{1}{2}
        \end{aligned}
        \]
    \end{solution}

    \question 设正实数 \( x, y \)（\( x \neq 1, y \neq 1 \)）满足
    \[
    \log_2 x = \log_y 16, \quad xy = 64,
    \]
    求
    \[
    \left(\log_2 \frac{x}{y}\right)^2。
    \]
    \begin{solution}
        由 \(xy = 64\)得 \(x = \frac{64}{y}\),代入第一式:
        \[
        \log_2 \left( \frac{64}{y} \right) = \log_y 16
        \Rightarrow 6 - \log_2 y = \frac{4}{\log_2 y}
        \]
        设 \(a = \log_2 y\),则变为
        \[
        6 - a = \frac{4}{a} \Rightarrow a^2 - 6a + 4 = 0
        \]
        解得
        \[
        \log_2 y =a= 3 \pm \sqrt{5}
        \]
        所求为
        \[
        \left(\log_2 x - \log_2 y\right)^2 = (a - (6 - a))^2 = (2a - 6)^2 = (\pm 2\sqrt{5})^2 = 20
        \]
    \end{solution}
    
    \question 已知 $a^x = bc,b^y = ac,c^z = ab,$ 证明 $$xyz = x + y + z + 2$$

    \begin{solution}
        由 $a^x = bc$,两边取 $yz$ 次方
        \[
        (a^x)^{yz} = (bc)^{yz} \Rightarrow a^{xyz} = b^{yz} \cdot c^{yz}
        \]
        由 $b^y = ac$,得 $b^{yz} = (ac)^z = a^z c^z$ ;由 $c^z = ab$,得 $c^{yz} = (ab)^y = a^y b^y$;于是        
        \[
        a^{xyz} = (a^z c^z)(a^y b^y) = a^{z+y} b^y c^z = a^{z+y} \cdot ac \cdot ab = a^{z+y+2} bc
        \]
        又 $a^x = bc$,故
        \[
        a^{xyz} = a^{z+y+2} \cdot a^x = a^{x+y+z+2} \Rightarrow xyz = x + y + z + 2
        \]
    \end{solution}

    \question 设 \(a, b, c\) 均为异于 1 的正数,且满足 \(abc = 1\), 证明
    \[
    \log_a b + \log_a c + \log_b c + \log_b a + \log_c b + \log_c a = -3
    \]
    \begin{solution}
    有\begin{align*}
          &\log_a a+\log_a b + \log_a c + \log_b b+ \log_b c + \log_b a + \log_c c + \log_c b + \log_c a \\ &=\log_a (abc)+\log_b (abc)+\log_c (abc)\\&=0
      \end{align*}
      故原式得证。
    \end{solution}

    \question 已知
        \[
        \log_{10} \sin x + \log_{10} \cos x = -1,
        \]
        且整数 $n$ 满足
        \[
        \log_{10} (\sin x + \cos x) = \frac{1}{2} (\log_{10} n - 1),
        \]
        求 \( n\)。
    \begin{solution}
        由
        \[
        \log_{10} \sin x + \log_{10} \cos x = -1
        \Rightarrow \sin x \cos x = 10^{-1}= \frac{1}{10}
        \]
        又$\sin x + \cos x=\sqrt{1+2\cdot\dfrac{1}{10}}=\sqrt{\dfrac{6}{5}}$ ,则
        \[
        \log_{10}{n} = 2 \log_{10}{\sqrt{\frac{6}{5}}} + 1 = \log_{10}{\frac{6}{5}} + \log_{10}{10} = \log_{10}{12} \Rightarrow n=12
        \]
    \end{solution}

    \question 解方程式 
    \[
    (\log_{5}x)^{2} + \log_{5x}\left(\frac{5}{x}\right) = 1
    \]
    \begin{solution}
        设$\log_5 x = t$,则
        \[
        \log_{5x}\left(\frac{5}{x}\right) = \frac{\log_5\frac{5}{x}}{\log_5(5x)} = \frac{\log_5 5 - \log_5 x}{\log_5 5 + \log_5 x} = \frac{1-t}{1+t}
        \]
        原方程式变为
        \begin{align*}
        t^2 + \frac{1-t}{1+t} &= 1 \\
        t^2(1+t) + 1-t &= 1+t \\
        t^3 + t^2 - 2t &= 0 \\
        t(t+2)(t-1) &= 0
        \end{align*}
        解得 $t = 0, -2, 1$,即 
        \[
        x = 5^0 = 1, \quad \text{或} \quad x = 5^{-2} = \frac{1}{25}, \quad \text{或} \quad x = 5^1 = 5
        \]
        经检验得$x = 1, x = \dfrac{1}{25}, x=5$都是原方程的解。
    \end{solution}
    
    \question 解方程
    \[
    \log_{2x} \left( 48 \sqrt[3]{3} \right) = \log_{3x} \left( 162 \sqrt[3]{2} \right)
    \]
    \begin{solution}
        原式变为
        \[
        \frac{\log (48 \sqrt[3]{3})}{\log (2x)} = \frac{\log (162 \sqrt[3]{2})}{\log (3x)}.
        \]
        又因
        \[
        \log (48 \sqrt[3]{3}) = 4 \log 2 + \frac{4}{3} \log 3,\quad 
        \log (162 \sqrt[3]{2}) = \frac{4}{3} \log 2 + 4 \log 3.
        \]
        于是有
        \[
        \frac{4 \log 2 + \frac{4}{3} \log 3}{\log 2 +  \log x}
        =
        \frac{\frac{4}{3} \log 2 + 4 \log 3}{\log 3 +  \log x}.
        \]
        化简最终得到
        \[
        \log x = \frac{1}{2} \log 6 \Rightarrow x = \sqrt{6}.
        \]
    \end{solution}
    
    \question 求所有实数 $x$,使得
    \[
    \log_{5x+9}(x^2+6x+9) + \log_{x+3}(5x^2+24x+27) = 4.
    \]
    \begin{solution}
        原方程式化为
        \begin{align*}
        \frac{\log(x^2+6x+9)}{\log(5x+9)} + \frac{\log(5x^2+24x+27)}{\log(x+3)} &= 4 \\
        \frac{2\log(x+3)}{\log(5x+9)} + \frac{\log(5x+9)+\log(x+3)}{\log(x+3)} &= 4 \\
        2 \frac{\log(x+3)}{\log(5x+9)} + \frac{\log(5x+9)}{\log(x+3)} + 1 &= 4
        \end{align*}
        令 $t = \dfrac{\log(x+3)}{\log(5x+9)}$,解得
        \[
        2t + \frac{1}{t} = 3 \Rightarrow t = 1 \ \text{或}\ t = \frac{1}{2}
        \]
        当$t = 1$,解得
        \[
        \log(x+3) = \log(5x+9) \Rightarrow x = -\frac{3}{2}
        \]
        当$t = \frac{1}{2}$,则 $2\log(x+3) = \log(5x+9)$,即
        \[
        2\log(x+3) = \log(5x+9) \Rightarrow x = 0 \ \text{或}\ x = -1
        \]
        经检验得,原方程的实数解为
        \[
        x = 0, -1, -\frac{3}{2}
        \]
    \end{solution}

    \question 求解下列方程组:
    \[
    \begin{cases}
    (2x)^{\ln 2} = (3y)^{\ln 3} \\
    3^{\ln x} = 2^{\ln y}
    \end{cases}
    \]
    \begin{solution}
        两边取对数,
        \[
        (\ln 2)^2 + \ln 2 \cdot \ln x = (\ln 3)^2 + \ln 3 \cdot \ln y \tag{1}
        \]
        \[
        \ln x \cdot \ln 3 = \ln y \cdot \ln 2 \Rightarrow \ln x = \frac{\ln y \cdot \ln 2}{\ln 3}\tag{2} 
        \]
        将(2)代入(1),可解得
        \[
        x=\frac{1}{2},\quad y= \frac{1}{3}
        \]
    \end{solution}

    \question 求下列方程组的所有实数解:
    \[
    \begin{cases}
    x + \log_{10} x = y-1 \\
    y + \log_{10}(y-1) = z-1 \\
    z + \log_{10}(z-2) = x+2
    \end{cases}
    \]
    \begin{solution}
        将方程组改写为
        \begin{align*}
        x + \log_{10} x &= y-1 \\
        (y-1) + \log_{10}(y-1) &= z-2 \\
        (z-2) + \log_{10}(z-2) &= x
        \end{align*}
        令 $a = x,b = y-1,c = z-2$,方程组可改写为:
        \begin{align}
        a + \log_{10} a &= b \tag{1} \\
        b + \log_{10} b &= c \tag{2} \\
        c + \log_{10} c &= a \tag{3}
        \end{align}
        容易验证 $(a,b,c)=(1,1,1)$ 是一解,现考虑 $a>1$,则
        \[
        \log_{10} a > 0 \Rightarrow b = a+\log_{10}a > a > 1 \Rightarrow c = b+\log_{10}b > b > a > 1
        \]
        由 (3) 得 $a = c + \log_{10}c > c > b > a$,矛盾。而当 $0<a<1$,
        \[
        \log_{10} a < 0 \Rightarrow b = a+\log_{10}a < a < 1 \Rightarrow c = b+\log_{10}b < b < a < 1
        \]
        由 (3) 得 $a = c + \log_{10}c < c < b < a$,矛盾。
        因此$(a,b,c)=(1,1,1)$ 是唯一解,即 
        \[
        (x,y,z)=(1,2,3)
        \]
    \end{solution}

    \question 若正实数 $x,y,z$ 满足
    \[
    \begin{cases}
    \log_{2}x + \log_{4}y + \log_{4}z = 2 \\
    \log_{3}y + \log_{9}z + \log_{9}x = 2 \\
    \log_{4}z + \log_{16}x + \log_{16}y = 2
    \end{cases}
    \]
    求 $xyz$。
    \begin{solution}
        原方程组给出
        \[
        \begin{cases}
        \log_2 x + \log_2 \sqrt{y} + \log_2 \sqrt{z} = 2 \Rightarrow x\sqrt{yz} = 2^2 \\
        \log_3 y + \log_3 \sqrt{z} + \log_3 \sqrt{x} = 2 \Rightarrow y\sqrt{xz} = 3^2 \\
        \log_4 z + \log_4 \sqrt{x} + \log_4 \sqrt{y} = 2
        \Rightarrow z\sqrt{xy} = 4^2
        \end{cases}
        \]
        将三式相乘得
        \[
        (xyz)^2 = 2^6 \cdot 3^2 \Rightarrow xyz = 24
        \]
    \end{solution}

    \question 解联立方程组
    \[
    \begin{cases}
    \log x \log y - 3 \log 5 y - \log 8 x = -4 \\
    \log y \log z - 4 \log 5 y - \log 16 z = 4 \\
    \log z \log x - 4 \log 8 x - 3 \log 625 z = -18
    \end{cases}
    \]
    \begin{solution}
        三式化简得
        \begin{align*}
        \log x \log y - 3 \log  y - \log x &= -1 \\
        \log y \log z - 4 \log  y - \log  z &= 8 \\
        \log z \log x - 4 \log  x - 3 \log z &= -6
        \end{align*}
        因式分解:
        \begin{align}
        (\log x-3)( \log y -1 )&= 2 \\
        (\log y-1)( \log z -4 )&= 12 \\
        (\log z-4)( \log x -3 )&= 6 
        \end{align}
        三式相乘可得
        \begin{equation}
        (\log x-3)( \log y -1 )(\log z-4)= \pm\sqrt{2\cdot 6\cdot 12} =\pm12
        \end{equation}
        分别作$\dfrac{(4)}{(2)},\dfrac{(4)}{(3)},\dfrac{(4)}{(1)}$得解   $\log x-3= \pm1,\log y -1=\pm2,\log z-4=\pm6$,即\[
        (x,y,z) = (10^4, 10^3, 10^{10}) \quad \text{或} \quad (10^2, 10^{-1}, 10^{-2})
        \]
    \end{solution}
\end{questions}
\pagebreak

\begin{center}
  {\fontsize{30pt}{26pt}\selectfont
    \hypertarget{方程组}{方程组} \label{方程组}
  }
\end{center}
\separator
\vspace{1pt}
\begin{questions}
    \question 若 $x, y, z$ 都是正数且满足 
    \[
    \begin{cases}
    \displaystyle x+\frac{1}{y}=4,\\[8pt] 
    \displaystyle y+\frac{1}{z}=1,\\[8pt] 
    \displaystyle z+\frac{1}{x}=\frac{7}{3}, 
    \end{cases}
    \]
    求 $xyz$ 的值。
    \begin{solution} 
        \begin{align*} 
        \left(x+\frac{1}{y}\right)\left(y+\frac{1}{z}\right)\left(z+\frac{1}{x}\right)
        &= xyz + x+\frac{1}{y} + y+\frac{1}{z} + z+\frac{1}{x} + \frac{1}{xyz}\\
        &= xyz + 4 + 1 + \frac{7}{3} + \frac{1}{xyz} = \frac{28}{3}
        \end{align*}
        于是
        \[
        xyz + \frac{1}{xyz} = 2 
        \Rightarrow (xyz)^2 - 2(xyz) + 1 = 0 
        \Rightarrow xyz = 1
        \]
    \end{solution}

    \question 求满足方程组
    \[
    \begin{cases}
    x^2 + x^2y^2 + x^2y^4 = 525 \\
    x + xy + xy^2 = 35
    \end{cases}
    \]
    的所有实数序对$(x,y)$。
    \begin{solution}
        \begin{align}
        x^2 + x^2y^2 + x^2y^4 &= 525 \\
        x + xy + xy^2 &= 35
        \end{align}
        发现
        \[
        525=x^2 + x^2y^2 + x^2y^4=(x - xy + xy^2)(x +xy+xy^2)
        \]
        于是得到
        \[
        x - xy + xy^2  = 15 \tag{3}
        \]
        $(2)-(3)$得
        \[
        2xy = 20 \Rightarrow x = \frac{10}{y}
        \]
        代回(3)解得
        \[
        (x,y)=(5,2),\left(20,\frac{1}{2}\right)
        \]
    \end{solution}

    \question 若两正数 $a,b$ 满足
    \[
    \begin{cases}
    a\sqrt{a} + b\sqrt{b} = 50 \\
    a\sqrt{b} + b\sqrt{a} = 25
    \end{cases}
    \]
    求 $ab$ 之值。
    \begin{solution}
        \[
        a\sqrt{a} + b\sqrt{b} = 50 \tag{1}
        \]
        \[
        a\sqrt{b} + b\sqrt{a} = 25 \tag{2}
        \]
        由 $\dfrac{(1)}{(2)}$ 得
        \[
        \frac{(\sqrt{a})^3 + (\sqrt{b})^3}{\sqrt{ab}(\sqrt{a}+\sqrt{b})}
        = \frac{a - \sqrt{ab} + b}{\sqrt{ab}} = 2 \Rightarrow a + b = 3\sqrt{ab} \Rightarrow a^2 + b^2 = 7ab \tag{3}
        \]
        又由 $(1)\times(2)$得
        \[
        (a\sqrt{a} + b\sqrt{b})(a\sqrt{b} + b\sqrt{a})
        = \sqrt{ab}(a^2+b^2) + ab(a+b) = 1250 \tag{4}
        \]
        将 (3) 代入 (4),
        \[
        \sqrt{ab}(7ab) + ab(3\sqrt{ab}) =10(\sqrt{ab})^3= 1250 \Rightarrow ab = 25
        \]
    \end{solution}

    \question 求正实数对 $(x,y)$ 满足
    \[
    \begin{cases}
    x^{2}+x\sqrt[3]{xy^{2}}=208 \\
    y^{2}+y\sqrt[3]{yx^{2}}=1053
    \end{cases}
    \]
    \begin{solution}
        原方程即
        \begin{align}
        x^{\frac{4}{3}}(x^{\frac{2}{3}}+y^{\frac{2}{3}})&=208 \\
        y^{\frac{4}{3}}(x^{\frac{2}{3}}+y^{\frac{2}{3}})&=1053
        \end{align}
        两式相除得
        \[
        \left(\frac{y}{x}\right)^{\frac{4}{3}}=\left(\frac{3}{2}\right)^{4} \Rightarrow \frac{y}{x}=\left(\frac{3}{2}\right)^{3}=\frac{27}{8},
        \]
        将$y=\dfrac{27}{8}x$代入(1),解得
        \[
        x^2\left(1+\frac{9}{4}\right)=208 \Rightarrow (x,y)=(8,27).
        \]
    \end{solution}

    \question 设相异实数 $x,y$ 满足
    \[
    \begin{cases}
    x^2 + \sqrt{3} y = 4 \\
    y^2 + \sqrt{3} x = 4 
    \end{cases}
    \]
    求 $\dfrac{x}{y} + \dfrac{y}{x}$ 的值。
    \begin{solution}
        \[
        x^2 + \sqrt{3} y = 4 \tag{1}
        \]
        \[
        y^2 + \sqrt{3} x = 4 \tag{2}
        \]
        由于$x\neq y,(1)-(2)$得
        \[
        x^2 - y^2 = \sqrt{3} (x - y) \Rightarrow x + y = \sqrt{3}.
        \]
        且$(1)+(2)$得
        \[
        x^2 + y^2 + \sqrt{3} \cdot \sqrt{3} = 8 \Rightarrow x^2 + y^2 = (\sqrt3)^2-2xy=5 \Rightarrow xy = -1
        \]
        因此
        \[
        \frac{x}{y} + \frac{y}{x} = \frac{x^2 + y^2}{xy} = \frac{5}{-1} = -5.
        \]
    \end{solution}

    \question 已知实数 $x, y, z$ 满足
    \[
    \begin{cases}
    xy + 4z = 60 \\
    yz + 4x = 60 \\
    zx + 4y = 60 
    \end{cases}
    \]
    求$x$的可能值。
    \begin{solution}
        \[
        xy + 4z = 60 \tag{1}
        \]
        \[
        yz + 4x = 60 \tag{2}
        \]
        \[
        zx + 4y = 60 \tag{3}
        \]
        由 (1) -(2) 得
        \[
        y(x-z)+4(z-x)=0 \Rightarrow (x-z)(y-4)=0 \Rightarrow y=4 \; \text{或} \; x=z
        \]
        同理有
        \[
        z=4 \; \text{或} \; x=y
        \]
        及
        \[
        x=4 \; \text{或} \; y=z
        \]
        若$x=4$,由(1)及(2)得
        \[
        y+z=15,\quad yz=44 
        \]
        由此可知$y, z$ 是方程 
        \[
        t^2-15t+44=0
        \]
        的根,解得
        \[
        (x,y,z)=(4,4,11),(4,11,4)
        \]
        若$y=z$,由(2)-(1)得
        \[
        y^2-(x+4)y+4x=0 \Rightarrow (y-4)(y-x)=0 
        \]
        当$y=4$有
        \[
        (x,y,z)=(11,4,4)
        \]
        而当$y=x$有$x=y=z$,代入(1)得
        \[
        x^2+4x-60=0 \Rightarrow (x+10)(x-6)=0 
        \]
        给出
        \[
        (x,y,z)=(-10,-10,-10),(6,6,6)
        \]
        所以原方程组的解为
        \[
        (x,y,z)=(-10,-10,-10),(6,6,6),(11,4,4),(4,11,4),(4,4,11)
        \]
    \end{solution}

    \question 解方程组
    \[
    \begin{cases}
    x+y+z=6 \\
    x^{2}+y^{2}+z^{2}=18 \\
    \sqrt{x}+\sqrt{y}+\sqrt{z}=4
    \end{cases}
    \]
    \begin{solution}
        由
        \[
        (x+y+z)^{2} = x^{2}+y^{2}+z^{2}+2(xy+xz+yz)
        \]
        给出
        \[
        xy+xz+yz = \frac{1}{2}(6^2-18) = 9
        \]
        又由 
        \[
        (\sqrt{x}+\sqrt{y}+\sqrt{z})^{2} = x+y+z+2(\sqrt{xy}+\sqrt{xz}+\sqrt{yz})
        \]
        所以 
        \[
        4^{2} = 6+2(\sqrt{xy}+\sqrt{xz}+\sqrt{yz}) \Rightarrow \sqrt{xy}+\sqrt{xz}+\sqrt{yz} = 5
        \]
        对上式两边平方,有
        \[
        xy+xz+yz+2\sqrt{xyz}(\sqrt{x}+\sqrt{y}+\sqrt{z})=25
        \]
        即
        \[
        9+2\sqrt{xyz} \cdot 4=25 \Rightarrow xyz=4
        \]
        由韦达定理, $x, y, z$ 为三次方程
        \[
        t^3 - 6t^2 + 9t - 4 = 0
        \]
        的根,因式分解得
        \[
        (t-1)^2(t-4) = 0 \Rightarrow t=1, 4
        \]
        所以方程组的解为 
        \[
        (x,y,z) = (1, 1, 4), (1, 4, 1), (4, 1, 1)
        \]
    \end{solution}

    \question 若 $x,y,z$ 满足
    \[
    \begin{cases} 
    \dfrac{x}{3} + \dfrac{y}{3+\log 2} + \dfrac{z}{3+\log 5} = 1 \\[8pt]
    \dfrac{x}{7} + \dfrac{y}{7+\log 2} + \dfrac{z}{7+\log 5} = 1 \\[8pt]
    \dfrac{x}{11} + \dfrac{y}{11+\log 2} + \dfrac{z}{11+\log 5} = 1
    \end{cases}
    \]
    求 $x+y+z$ 之值。
    \begin{solution}
        不妨设
        \[
        \frac{a}{t}+\frac{b}{t+\log 2}+\frac{c}{t+\log 5}=1
        \]
        则
        \[
        f(t)=a(t+\log 2)(t+\log 5)+b\,t(t+\log 5)+c\,t(t+\log 2)-t(t+\log 2)(t+\log 5)=0
        \]
        由韦达定理,三根之和为
        \[
        a+b+c-1=3+7+11=21 \Rightarrow a+b+c=22
        \]
    \end{solution}

    \question 已知 $x, y \in \mathbb{R}$,求解方程组
    \[
    \begin{cases} 
    \sqrt{x}\left(1+\dfrac{1}{x+y}\right)=2 \\[8pt]
    \sqrt{y}\left(1-\dfrac{1}{x+y}\right)=\sqrt{2} 
    \end{cases}
    \]
    \begin{solution}
        \[
        1+\frac{1}{x+y} = \frac{2}{\sqrt{x}} \tag{1} 
        \]
        \[
        1-\frac{1}{x+y} = \frac{\sqrt{2}}{\sqrt{y}} \tag{2} 
        \]
        $(1)^2-(2)^2$:
        \[
        \left(1+\frac{1}{x+y}\right)^2 - \left(1-\frac{1}{x+y}\right)^2 = \frac{4}{x} - \frac{2}{y} \Rightarrow \frac{4}{x+y} = \frac{4y-2x}{xy}
        \]
        整理得
        \[
        x^2 + xy - 2y^2 = 0 \Rightarrow (x+2y)(x-y)=0
        \]
        由于 $x,y \ge 0$,取 $x=y$,代入(1)解得
        \[
        1 + \frac{1}{2x} = \frac{2}{\sqrt{x}} \Rightarrow 4x^2 - 12x +1 =0 \Rightarrow x = y = \frac{3+2\sqrt{2}}{2}
        \]
    \end{solution}

    \question 解方程组
    \[
    \begin{cases} 
    \sqrt{\dfrac{x+y}{x}}+\sqrt{\dfrac{x}{x+y}}=\dfrac{5}{2} \\[8pt]
    2x^{2}+y^{2}=176
    \end{cases}
    \]
    \begin{solution}
        \[
        \sqrt{\frac{x+y}{x}}+\sqrt{\frac{x}{x+y}}=\frac{5}{2} \tag{1}
        \]
        \[
        2x^{2}+y^{2}=176 \tag{2}
        \]
        设
        \[
        u = \sqrt{\frac{x+y}{x}}
        \]
        方程(1)变为
        \[
        u + \frac{1}{u} = \frac{5}{2} \Rightarrow (2u-1)(u-2)=0
        \]
        当$u=\frac{1}{2}$,
        \[
        \frac{x+y}{x} = \frac{1}{4} \Rightarrow y = -\frac{3}{4}x
        \]
        代入(2),解得
        \[
        x = \pm 16 \sqrt{\frac{11}{41}}, \quad y = -\frac{3}{4}x = \mp 12 \sqrt{\frac{11}{41}}
        \]
        而当$u=2$,
        \[
        \frac{x+y}{x} = 16 \Rightarrow y = 3x
        \]
        代入(2),解得
        \[
        x = \pm 4, \quad y = 3x = \pm 12
        \]
        故原方程组解为
        \[
        (x,y) = (4,12), (-4,-12), \left(16\sqrt{\frac{11}{41}}, -12\sqrt{\frac{11}{41}}\right), \left(-16\sqrt{\frac{11}{41}}, 12\sqrt{\frac{11}{41}}\right)
        \]
    \end{solution}

    \question 解方程组
    \[
    \begin{cases}
    (91-2x)^{3}=216xy^{2} \\
    (37-2y)^{3}=216x^{2}y
    \end{cases}
    \]
    \begin{solution}
        \[
        (91-2x)^{3}=216xy^{2} \tag{1}
        \]
        \[
        (37-2y)^{3}=216x^{2}y \tag{2}
        \]
        令
        \[
        x = u^3, \quad y = v^3.
        \]
        将(1),(2)两边开立方得
        \[
        91-2u^3 = 6 u v^2, \quad 37-2v^3 = 6 u^2 v.
        \]
        两式相加,
        \[
        2(u+v)^3=128 \Rightarrow u+v=4
        \]
        将 $v = 4-u$ 代入 (1),
        \[
        6 u (4-u)^2 + 2 u^3 = 91 \Rightarrow (2 u - 7) (4 u^2 - 10 u + 13)=0 \Rightarrow u=\frac{7}{2} \in \mathbb{R}
        \]
        因此
        \[
        v = 4 - u = \frac{1}{2}
        \]
        最终解为
        \[
        x = \frac{343}{8}, \quad y = \frac{1}{8}
        \]
    \end{solution}

    \question 解
    \[
    \begin{cases}
    3(a^2+b^2)^{\frac32}=125a\\[4pt]
    4(a^2+b^2)=-25b
    \end{cases}
    \]
    \begin{solution}
        \[
        3(a^2+b^2)^{\frac32}=125a \tag{1}
        \]
        \[
        4(a^2+b^2)=-25b \tag{2}
        \]
        由$\dfrac{(1)}{(2)}$得
        \[
        -\frac{20a}{3b}=(a^2+b^2)^{\frac12}
        \]
        平方得
        \[
        \frac{400a^2}{9b^2}=a^2+b^2 \Rightarrow a^2=\frac{9b^2}{400-9b^2}
        \]
        代入(2)得
        \[
        4\left(\frac{9b^2}{400-9b^2}+b^2\right)+25b=0 \Rightarrow b(b+4)(9b-100)=0 
        \]
        经检验得$b=\dfrac{100}{9}$是增根,原方程组的解为
        \[
        (a,b)=(0,0)\quad \text{或} \quad (3,-4)
        \]
    \end{solution}

    \question 解方程组
    \[
    \begin{cases}
    x^3 + 9x^2y = -28 \\
    y^3 + xy^2 = 1
    \end{cases}
    \]
    \begin{solution}
        \[
        x^3 + 9x^2y = -28 \tag{1}
        \]
        \[
        y^3 + xy^2 = 1 \tag{2}
        \]
        $(1)+27\cdot (2)$,得
        \[
        x^3 + 9x^2y + 27xy^2 + 27y^3 = -1 
        \]
        即
        \[
        (x + 3y)^3 = -1 \Rightarrow x + 3y = -1
        \]
        将 $x = -1 - 3y$ 代入(2),
        \[
        y^3 + (-1 - 3y)y^2 = 1 \Rightarrow (y+1)(2y^2 - y + 1)=0
        \]
        其中$2y^2 - y + 1$的判别式
        \[
        \Delta = (-1)^2 - 4\cdot 2 \cdot 1 = -7 < 0
        \]
        故$2y^2 - y + 1=0$无实数解,因此唯一实数解为 
        \[
        (x, y) = (2, -1)
        \]
    \end{solution}

    \question 解方程组
    \[
    \begin{cases}
    x^3 + 6xy^2 = 99 \\
    2y^3 + 3x^2y = 70 
    \end{cases}
    \]
    \begin{solution}
        \[
        x^3 + 6xy^2 = 99 \tag{1}
        \]
        \[
        2y^3 + 3x^2y = 70 \tag{2}
        \]
        观察到方程式齐次,不妨设 $y = m x$,其中 $m \neq 0$,代入方程得到
        \[
        x^3 (1 + 6 m^2) = 99, \quad x^3 (2 m^3 + 3 m) = 70.
        \]
        两式相除得
        \[
        \frac{1 + 6 m^2}{2 m^3 + 3 m} = \frac{99}{70} \Rightarrow (3m - 2)(66 m^2 - 96 m + 35)=0
        \]
        其中$66 m^2 - 96 m + 35$的判别式为
        \[
        \Delta = (-96)^2 - 4 \cdot 66 \cdot 35 = -24 < 0
        \]
        故$66 m^2 - 96 m + 35=0$无实根,因此唯一实根为
        \[
        m = \frac{2}{3} \Rightarrow  y = \frac{2}{3} x
        \]
        将之代入(1),
        \[
        x^3 + 6 x \left(\frac{2}{3} x\right)^2 = 99 \Rightarrow x = 3
        \]
        得到
        \[
        (x, y) = (3, 2)
        \]
    \end{solution}

    \question 求方程
    \[
    \begin{cases}
    36y^{2}(x+1)+36x^{2}(y+1)=7x^{2}y^{2} \\
    6x+6y+xy=0
    \end{cases}
    \]
    的实数解。
    \begin{solution}
        \[
        36y^{2}(x+1)+36x^{2}(y+1)=7x^{2}y^{2} \tag{1}
        \]
        \[
        6x+6y+xy=0
        \]
        由(1),
        \[
        \frac{1}{x} + \frac{1}{x^2} + \frac{1}{y} + \frac{1}{y^2} = \frac{7}{36} \tag{3}
        \]
        由(2),
        \[
        \frac{1}{x} + \frac{1}{y} = -\frac{1}{6} \tag{4}
        \]
        将(3)代入(4),得
        \[
        \frac{1}{x^2} + \left(-\frac{1}{x}-\frac{1}{6}\right)^2 = \frac{13}{36} \Rightarrow (x-3)(x+2)=0
        \]
        故原方程式的解为
        \[
        (x,y) = (3,6),\quad \left(-2,-\frac{3}{2}\right)
        \]
    \end{solution}

    \question 设 $a,b,c$ 为相异非零实数,且
    \[
    \frac{1+a^3}{a}=\frac{1+b^3}{b}=\frac{1+c^3}{c}.
    \]
    求 $a^3+b^3+c^3$ 的所有可能值。
    \begin{solution}
        设
        \[
        \frac{1+a^3}{a}=\frac{1+b^3}{b}=\frac{1+c^3}{c}=k
        \]
        则 $a,b,c$ 是方程
        \[
        x^3-kx+1=0
        \]
        的三根,故设
        \[
        x^3-kx+1=(x-a)(x-b)(x-c)=x^3-(a+b+c)x^2+(ab+ac+bc)x-abc
        \]
        比较系数得
        \[
        a+b+c=0,\quad abc=-1
        \]
        故
        \[
        a^3+b^3+c^3=(a+b+c)(a^2+b^2+c^2-ab-bc-ca)+3abc0+3(-1)=-3
        \]
    \end{solution}

    \question 若 $a,b,c \in \mathbb{R}^+$ 满足
    \[
    abc = 120, \quad a+b+c = a^2+b^2+c^2 = a^3+b^3+c^3 = \lambda,
    \]
    求 $\lambda \in \mathbb{Z}^+$。
    \begin{solution}
        已知
        \[
        a+b+c = a^2+b^2+c^2 = a^3+b^3+c^3 = \lambda.
        \]
        由$a^2+b^2+c^2 = (a+b+c)^2 - 2(ab+bc+ca)$,
        \[
        ab+bc+ca = \frac{\lambda^2 - \lambda}{2}.
        \]
        由$a^3+b^3+c^3 = (a+b+c)^3 - 3(a+b+c)(ab+bc+ca) + 3abc$,
        \[
        \lambda = \lambda^3 - 3 \lambda \left(\frac{\lambda^2 - \lambda}{2}\right) + 3 \cdot 120 
        \]
        解得
        \[
        (\lambda-10)(\lambda^2+7\lambda+72)=0 \Rightarrow \lambda = 10 \in \mathbb{Z}^+
        \]
    \end{solution}

    \question 已知
    \[
    \begin{cases} 
    \alpha+\beta+\gamma=6 \\ 
    \alpha^3+\beta^3+\gamma^3=87 \\ 
    (\alpha+1)(\beta+1)(\gamma+1)=33 
    \end{cases}
    \]
    求
    \[
    \frac{1}{\alpha}+\frac{1}{\beta}+\frac{1}{\gamma}.
    \]
    \begin{solution}
        由$\alpha^3+\beta^3+\gamma^3 - 3\alpha\beta\gamma = (\alpha+\beta+\gamma)\bigl((\alpha+\beta+\gamma)^2 - 3(\alpha\beta+\beta\gamma+\gamma\alpha)\bigr)$,则
        \[
        87 - 3\alpha\beta\gamma = 6\bigl(36 - 3(\alpha\beta+\beta\gamma+\gamma\alpha)\bigr) \Rightarrow 6(\alpha\beta+\beta\gamma+\gamma\alpha) - \alpha\beta\gamma = 43 \tag{1}
        \]
        又由$33=(\alpha+1)(\beta+1)(\gamma+1) = 1 + 6 + \alpha\beta+\beta\gamma+\gamma\alpha + \alpha\beta\gamma$,
        \[
        \alpha\beta+\beta\gamma+\gamma\alpha + \alpha\beta\gamma = 26 \tag{2}
        \]
        联立(1),(2)得,
        \[
        \alpha\beta+\beta\gamma+\gamma\alpha = \frac{69}{7}, \quad \alpha\beta\gamma = \frac{113}{7}
        \]
        因此
        \[
        \frac{1}{\alpha}+\frac{1}{\beta}+\frac{1}{\gamma} = \frac{\alpha\beta+\beta\gamma+\gamma\alpha}{\alpha\beta\gamma} = \frac{69}{113}.
        \]
    \end{solution}

    \question 已知 $x,y,z$ 为实数满足 
    \[
    \begin{cases}
    \displaystyle \frac{1}{x}+\frac{1}{y}+\frac{1}{z}=1 \\[6pt]
    x^2+y^2+z^2=\dfrac{3}{2} \\
    x^3+y^3+z^3=1
    \end{cases}
    \]
    且$x+y+z$ 为整数, 求 $x+y+z$ 的值。
    \begin{solution}
        \[
        \begin{cases}
        xy+yz+zx=xyz \\
        x^2+y^2+z^2=\dfrac{3}{2} \\
        x^3+y^3+z^3=1
        \end{cases}
        \]
        令 $a=xyz$, 则
        \[
        (x+y+z)^2 = \frac{3}{2} + 2a, \quad 1-3a = (x+y+z)\left(\frac{3}{2}-a\right) 
        \]
        联立得
        \[
        \left(\frac{1-3a}{\frac{3}{2}-a}\right)^2 = \frac{3}{2}+2a
        \]
        解得 
        \[
        (4a+1)(4a^2-28a+19)=0 \Rightarrow a=xyz=-\frac{1}{4}
        \]
        故
        \[
        x+y+z = \frac{1-3xyz}{\frac{3}{2}-xyz} = \frac{\frac{7}{4}}{\frac{7}{4}} = 1
        \]
    \end{solution}

    \question 已知存在实数 $a, b, x, y$ 满足以下方程:
    \[
    \begin{cases}
    a x^{2014} + b y^{2014} = 6 \\
    a x^{2015} + b y^{2015} = 7 \\
    a x^{2016} + b y^{2016} = 3 \\
    a x^{2017} + b y^{2017} = 50
    \end{cases}
    \]
    求 $a x^{2018} + b y^{2018}$ 的值。
    \begin{solution} 
        不妨设 $f(n)=ax^n+by^n$,于是
        \[
        f(2014)=6,f(2015)=7,f(2016)=3,f(2017)=50,
        \]
        发现
        \[
        (x+y)f(2015)=ax^{2016}+by^{2016}+ax^{2015}y + bxy^{2015} =f(2016)+xyf(2014)
        \]
        \[
        (x+y)f(2016)=ax^{2017}+by^{2017}+ax^{2016}y + bxy^{2016} =f(2017)+xyf(2015)
        \]
        解得
        \[
        x+y=-9,\;xy=-11
        \]
        故
        \[
        (x+y)f(2017)=f(2018)+xyf(2016) \Rightarrow a x^{2018} + b y^{2018}=f(2018)=-417
        \]
    \end{solution}

    \question 求联立方程
    \[
    \begin{cases}
    x\left(2x^2 + y - \dfrac{1}{2}\right) = 0 \\[3pt]
    y \left(x - y + \dfrac{5}{2}\right) = 0
    \end{cases}
    \]
    的所有实数解 $(x, y)$。
    \begin{solution}
        \begin{align}
        x\left(2x^2 + y - \dfrac{1}{2}\right) &= 0 \\[3pt]
        y\left(x - y - \dfrac{5}{2}\right) &= 0
        \end{align}
        \textbf{情况一}:若 $x = 0$,由(2)得
        \[
        y\left(0 - y - \dfrac{5}{2}\right) = 0
        \Rightarrow y = 0 \; \text{或} \; y = -\dfrac{5}{2}.
        \]
        解为$$(0, 0),\left(0, -\dfrac{5}{2}\right)$$
        \textbf{情况二}: 若 $y = 0$,由(1)得
        \[
        x\left(2x^2 + 0 - \dfrac{1}{2}\right) = 0
        \Rightarrow x = 0 \; \text{或} \; x = \pm \dfrac{1}{2}.
        \]
        其中 $x=0$ 已考虑,解为
        \[
        \left( \dfrac{1}{2}, 0 \right), \left( -\dfrac{1}{2}, 0 \right)
        \]
        \textbf{情况三}:若$x\neq0$且$y\neq0$,
        \[
        \begin{cases}
        2x^2 + y - \dfrac{1}{2} = 0 \\[3pt]
        x - y - \dfrac{5}{2} = 0
        \end{cases}
        \]
        两式相加得
        \[
        2x^2 + x - 3 = 0
        \Rightarrow x = \dfrac{3}{2} \; \text{或} \; x = -1
        \]
        解为
        \[
        \left(\dfrac{3}{2}, -1\right),\; \left(-1, -\dfrac{7}{2}\right)
        \]
        $\therefore\;$原方程组的所有解为
        \[
        (0, 0),\;
        \left(0, -\dfrac{5}{2}\right),\;
        \left(\dfrac{1}{2}, 0\right),\;
        \left(-\dfrac{1}{2}, 0\right),\;
        \left(\dfrac{3}{2}, -1\right),\;
        \left(-1, -\dfrac{7}{2}\right)
        \]
    \end{solution}

    \question 求联立方程
    \[
    \begin{cases}
    x^2 - y - 2z = 4 \\
    y^2 - 2z - 3x = -2 \\
    2z^2 - 3x - 5y = -22
    \end{cases}
    \]
    的所有实数解 $(x, y, z)$。 
    \begin{solution}
        三式相加得
        \[
        x^2 - 6x +y^2 - 6y+2z^2-4z = -20,
        \]
        经配方后变为
        \[
        (x-3)^2+(y-3)^2+2(z-1)^2=0
        \]
        故$x-3=0,y-3=0,z-1=0$,解为$(3, 3, 1)$
    \end{solution}
    
    \question 设实数 $x, y, z$ 满足以下方程
    \[
    \begin{cases}
    4x + 2 y z - 6 z + 9 x z^2 = 4 \\
    x y z = 1
    \end{cases}
    \]
    求 $x + y + z$ 的所有可能值。
    \begin{solution}
        眼光发现$xyz=1$提示了换元$x=\dfrac{a}{b},\;y=\dfrac{b}{c},\;z=\dfrac{c}{a}$,第一式变为
        \[
        4\left(\dfrac{a}{b}\right)+2\left(\dfrac{b}{c}\right)\left(\dfrac{c}{a}\right)-6\left(\dfrac{c}{a}\right)+9\left(\dfrac{a}{b}\right)\left(\dfrac{c}{a}\right)^2=4 
        \]
        化简得
        \[
        4a^2-4ab+2b^2−6bc+9c^2=0 \Rightarrow (2a-b)^2+(b+3c)^2=0
        \]
        于是$2a-b=0,b+3c=0$,即
        \[
        x=\frac{a}{b}=\frac{1}{2},\;y=\frac{b}{c}=-3,\;z=\frac{1}{xy}=-\frac{2}{3}
        \]
        解得
        \[
        (x,y,z)=\left(\frac{1}{2},-3,-\frac{2}{3}\right) \Rightarrow x+y+z = -\frac{13}{6}
        \]
    \end{solution}
    
    \question 求所有非零实数对 $(x, y)$, 使得满足
    \[
    \frac{x}{x^2 + y} + \frac{y}{x + y^2} = -1 \quad \text{且} \quad \frac{1}{x} + \frac{1}{y} = 1
    \]
    \begin{solution}
        令 $u = x + y,\; v = xy$,则第二式化为 $\dfrac{u}{v} = 1 \Rightarrow u = v$。
        
        由第一式得
        \[
        x(x + y^2) + y(x^2 + y) = -[(x^2 + y)(x + y^2)]
        \]
        展开并代换为 $u, v$ 得
        \[
        u^2 - 2v + uv = - (u^3 - 3uv + v + v^2)
        \Rightarrow u^3 + u^2 - 2v + v^2 - v = 0
        \]
        代入 $u = v$
        \[
        v^3 - v = 0 \Rightarrow v(v^2 - 1) = 0
        \]
        因 $xy \ne 0$,故 $v = \pm1$
        
        若 $v = 1$,则 $x + y = 1,\; xy = 1 \Rightarrow x^2 - x + 1 = 0$,无实根。
        
        若 $v = -1$,则 $x + y = -1,\; xy = -1 \Rightarrow x^2 + x -1 = 0$
        
        解得
        \[
        (x, y) = \left(\frac{-1 \pm \sqrt{5}}{2},\; \frac{-1 \mp \sqrt{5}}{2} \right)
        \]
    \end{solution}

    \question 解方程组
    \[
    \begin{cases}
    x+y+z = 9  \\
    x^2+y^2+z^2 = 41  \\
    x^2(y+z) + y^2(z+x) + z^2(x+y) = 180
    \end{cases}
    \]
    \begin{solution}
        \[
        x+y+z = 9 \tag{1}
        \]
        \[
        x^2+y^2+z^2 = 41 \tag{2}
        \]
        \[
        x^2(y+z) + y^2(z+x) + z^2(x+y) = 180 \tag{3}
        \]
        由(2),
        \[
        (x+y+z)^2 - 2(xy+xz+yz) = 41 \Rightarrow xy+xz+yz = 20
        \]
        由(3),
        \begin{align*}
        x^2(9-x) + y^2(9-y) + z^2(9-z) &= 180 \\
        9(x^2+y^2+z^2) - (x^3+y^3+z^3) &= 180 \\
        9(41) - [(x+y+z)(x^2+y^2+z^2 - xy - yz - zx) + 3xyz] &= 180 \\
        369 - [9(41-20) + 3xyz] &= 180 \\
        xyz &= 0
        \end{align*}
        由韦达定理,实数$x, y, z$ 为方程式
        \[
        t^3 - 9t^2 + 20t = 0 \Rightarrow t(t-5)(t-4) = 0.
        \]
        的三根,于是解为
        \[
        (0,5,4), (0,4,5), (4,5,0), (4,0,5), (5,0,4), (5,4,0)
        \]
    \end{solution}

    \question 已知实数 $x, y, z$ 满足方程组
    \[
    \begin{cases}
    x+y+z=-3 \\[1mm]
    \dfrac{1}{x}+\dfrac{1}{y}+\dfrac{1}{z}=-\dfrac{1}{3} \\[1mm]
    x^2(y+z)+y^2(z+x)+z^2(x+y)=-24
    \end{cases}
    \]
    求 $x^2+y^2+z^2$的值。
    \begin{solution}
        \[
        x+y+z=-3 \tag{1}
        \]
        \[
        \dfrac{1}{x}+\dfrac{1}{y}+\dfrac{1}{z}=-\dfrac{1}{3} \tag{2}
        \]
        \[
        x^2(y+z)+y^2(z+x)+z^2(x+y)=-24 \tag{3}
        \]
        设 $x^2+y^2+z^2 = a$,由 $(1)$ 可得
        \[
        (x+y+z)^2 = 9 = a + 2(xy+yz+zx) \Rightarrow xy+yz+zx = \frac{9-a}{2}
        \]
        由 $(2)$ 得
        \[
        \frac{1}{x}+\frac{1}{y}+\frac{1}{z} = \frac{xy+yz+zx}{xyz} = -\frac{1}{3} \Rightarrow xyz = \frac{3(a-9)}{2}
        \]
        由 $(3)$ 得
        \[
        \begin{aligned}
        &x^2(y+z) + y^2(z+x) + z^2(x+y) \\
        &= x^2(-3-x) + y^2(-3-y) + z^2(-3-z) \\ 
        &= -3(x^2+y^2+z^2) - (x^3+y^3+z^3) \\ 
        &= -3a - \Big[ (x+y+z)((x^2+y^2+z^2)-(xy+yz+zx)) + 3xyz \Big] \\
        &= -3a - \big[ (-3)\left(a - \frac{9-a}{2}\right) + 3 \cdot \frac{3(a-9)}{2} \big] \\
        &= -3a + \frac{-9a+27}{2} + \frac{9a-81}{2} \\
        &= -3a + 27=-24 \Rightarrow a = 17
        \end{aligned}
        \]
    \end{solution}

    \question 已知方程组
    \[
    \begin{cases}
    x+y+z=0 \\
    x^2+y^2+z^2=6 \\
    x^3+y^3+z^3=-3
    \end{cases}
    \]
    求 $x^4+y^4+z^4$ 之值。
    \begin{solution}
        \[
        x+y+z=0 \tag{1}
        \]
        \[
        x^2+y^2+z^2=6 \tag{2}
        \]
        \[
        x^3+y^3+z^3=-3 \tag{3}
        \]
        由(2),
        \[
        (x+y+z)^2 - 2(xy+yz+xz) = 6 \Rightarrow xy+yz+xz = -3
        \]
        由(3),
        \[
        x^3+y^3+z^3 = (x+y+z)(x^2+y^2+z^2 - xy - yz - xz) + 3xyz \Rightarrow xyz = -1
        \]
        因此 $x,y,z$ 为方程式
        \[
        t^3 - 3t + 1 = 0
        \]
        的三根,于是有
        \[
        x^3=3x-1,\quad y^3=3y-1,\quad z^3=3z-1
        \]
        故
        \[
        x^4+y^4+z^4 = 3(x^2+y^2+z^2) - (x+y+z) = 3\cdot 6 - 0 = 18
        \]
    \end{solution}

    \question 求序对 $(x,y,z)$ 满足
    \[
    \begin{cases}
    (x+y)^3 = z,\\
    (y+z)^3 = x,\\
    (z+x)^3 = y.
    \end{cases}
    \]
    \begin{solution}
        不失一般性,设 $x \ge y \ge z$,则有
        \[
        2x \ge x+y \ge x+z, \quad x+y \ge 2y\ge y+z, \quad x+z \ge y+z \ge 2z
        \]
        立方可得
        \[
        \begin{cases}
        8x^3\ge (x+y)^3 \ge (x+z)^3\\
        (x+y)^3 \ge 8y^3 \ge (y+z)^3\\
        (x+z)^3 \ge (y+z)^3 \ge 8z^3
        \end{cases}
        \Rightarrow
        \begin{cases}
        8x^3\ge z \ge y\\
        z \ge 8y^3 \ge x\\
        y \ge x \ge z^3
        \end{cases}
        \Rightarrow z \ge y \ge x
        \]
        因此$x=y=z$,解$(x+x)^3=x$得
        \[
        (x,y,z)=(0,0,0),\left(\frac{\sqrt2}{4},\frac{\sqrt2}{4},\frac{\sqrt2}{4}\right),\left(-\frac{\sqrt2}{4},-\frac{\sqrt2}{4},-\frac{\sqrt2}{4}\right)
        \]
    \end{solution}

    \question 求满足
    \[
    \begin{cases}
    (1+x)(1+x^2)(1+x^4) = 1+y^7,\\
    (1+y)(1+y^2)(1+y^4) = 1+x^7
    \end{cases}
    \]
    的实数序对 $(x,y)$ 个数。
    \begin{solution}
        情况一:$xy=0$,得解 $(0,0)$。  

        情况二:$xy<0$。不失一般性设 $x>0>y$,则 
        \[(1+x)(1+x^2)(1+x^4) > 1
        \]
        且 $1+y^7 < 1$,故无解。  

        情况三:$x,y>0, x \ne y$。不失一般性设 $x>y>0$,则
        \[
        (1+x)(1+x^2)(1+x^4) > 1+x^7 > 1+y^7
        \]
        无解。  

        情况四:$x,y<0, x \ne 0$。不失一般性设 $x<y<0$,则
        \[
        1-x^8=(1+y^7)(1-x)=1-x+y^7-xy^7 \tag{1}
        \]
        \[
        1-y^8=(1+x^7)(1-y)=1-y+x^7-x^7y \tag{2}
        \]
        $(2)-(1)$得,
        \[
        x^8-y^8=x-y+x^7-y^7-xy(x^6-y^6)
        \]
        由于$x<y<0$,则$x^8-y^8>0,x-y<0,x^7-y^7<0,-xy<0,x^6-y^6>0$,左式为正,右式为负,故无解。

        情况五:$x=y$,则
        \[
        1-x^8=1-x+y^7-xy^7=1-x+x^7-x^8 \Rightarrow x=-1,0,1
        \] 
        其中只有 $(0,0),(-1,-1)$ 成立。  

        综上,解为 $(0,0),(-1,-1)$,共有 $2$ 个实数序对。
    \end{solution}

    \question 若非零实数 $a,b,c$ 满足
    \[
    \begin{cases}
    a^2 + b^2 + c^2 = 1,\\[1mm]
    \displaystyle a\left(\frac{1}{b}+\frac{1}{c}\right)+b\left(\frac{1}{a}+\frac{1}{c}\right)+c\left(\frac{1}{a}+\frac{1}{b}\right) = -3,
    \end{cases}
    \]
    求 $a+b+c$ 可能的取值个数。
    \begin{solution}
        由第二个方程式,
        \[
        a^2(b+c) + b^2(a+c) + c^2(a+b) = -3abc \Rightarrow (a+b+c)(ab+bc+ca)=0.
        \]
        又$(a+b+c)^2=1-2(ab+bc+ca)$,所以
        \[
        \frac{1}{2}(a+b+c)((a+b+c)^2-1)=0 \Rightarrow a+b+c=-1,0,1
        \]
        因此 $a+b+c$ 共有 $3$ 个可能值。
    \end{solution}

    \question 求在区间 $[0,2]$ 内满足方程组
    \[
    \begin{cases}
    2x^2 - 4x + 7 = y,\\
    2y^2 - 4y + 7 = z,\\
    2z^2 - 4z + 7 = x
    \end{cases}
    \]
    的无序三元组 $(x,y,z)$个数。
    \begin{solution}
        令 $a=x-1, b=y-1, c=z-1$,则 $a,b,c\in[-1,1]$,方程组化为
        \[
        \begin{cases}
        2a^2-1=b,\\
        2b^2-1=c,\\
        2c^2-1=a.
        \end{cases}
        \]
        令 $a=\cos\theta, b=\cos 2\theta, c=\cos 4\theta,\theta\in[0,\pi]$,得到
        \[
        \cos\theta = \cos 8\theta \Rightarrow -2\sin\frac{9\theta}{2}\sin\frac{7\theta}{2} = 0.
        \]
        解得 
        \[
        \theta -2n\pi= 0, \frac{2\pi}{9}, \frac{4\pi}{9}, \frac{6\pi}{9}, \frac{8\pi}{9}, \frac{2\pi}{7}, \frac{4\pi}{7}, \frac{6\pi}{7},n \in \mathbb{Z}
        \] 
        在区间 $[-1,1]$ 内的无序三元组$(a,b,c)$为 
        \[
        (0,0,0), \left(\cos\frac{2\pi}{7},\cos\frac{4\pi}{7},\cos\frac{8\pi}{7}\right), \left(\cos\frac{2\pi}{9},\cos\frac{4\pi}{9},\cos\frac{8\pi}{9}\right), \left(\cos\frac{2\pi}{3},\cos\frac{2\pi}{3},\cos\frac{2\pi}{3}\right)
        \] 
        故原方程式无序三元组个数为$4$。
    \end{solution}

    \question 若 $x,y,z$ 为相异复数,满足
    \[
    \begin{cases}
    x+y+z=1,\\
    x^2+y=y^2+z=z^2+x,
    \end{cases}
    \]
    求 $(x-y)(y-z)(z-x)$ 的值。
    \begin{solution}
        由$x^2+y=y^2+z$,得
        \[
        x^2-y^2=(x+y)(x-y)=(1-z)(x-y)=z-y
        \]
        同理得
        \[
        (1-x)(y-z)=x-z,\quad (1-y)(z-x)=y-x
        \]
        由于$x \neq y \neq z$,三式相乘可得
        \[
        (1-x)(1-y)(1-z)=-1 \tag{1}
        \]
        且由$(1-z)(x-y)=z-y$展开得
        \[
        x-z=xz-yz=z(x-y)
        \]
        同理得
        \[
        y-x=x(y-z),\quad z-y=y(z-x)
        \]
        由于$x \neq y \neq z$,三式相乘可得
        \[
        xyz=-1
        \]
        现由$(1)$,得
        \[
        1-(x+y+z)+xy+yz+zx-xyz=-1 \Rightarrow xy+yz+zx=-2
        \]
        设 $k = x^2+y = y^2+z = z^2+x$,则
        \[
        3k=x^2+y^2+z^2+x+y+z=1-2(-2)+1=6 \Rightarrow x^2+y=k=2
        \]
        则
        \[
        x-y=x-(2-x^2)=x^2+x-2=(x-1)(x+2),
        \]
        同理
        \[
        y-z=(y-1)(y+2),\quad z-x=(z-1)(z+2)
        \]
        于是
        \begin{align*}
        &(x-y)(y-z)(z-x) \\
        &= (x-1)(y-1)(z-1)(x+2)(y+2)(z+2) \\
        &= [xyz-(xy+yz+zx)+x+y+z-1][xyz+2(xy+yz+zx)+4(x+y+z)+8]=7
        \end{align*}
    \end{solution}

    \question 已知 $xyz \neq 0,a, b, c$ 不全为零,且满足方程组
    \[
    \begin{cases}
    a = \dfrac{by}{z} + \dfrac{cz}{y} \\[2mm]
    b = \dfrac{cz}{x} + \dfrac{ax}{z} \\[2mm]
    c = \dfrac{ax}{y} + \dfrac{by}{x}
    \end{cases}
    \]
    \begin{parts}
    \part 证明 $a^3x^3 + b^3y^3 + c^3z^3 + abcxyz = 0$。
    \begin{solution}
        将方程组改写为
        \[
        \begin{cases}
        (-ax) \dfrac{1}{x} + (cz) \dfrac{1}{y} + (by) \dfrac{1}{z} = 0 \\[2mm]
        (cz) \dfrac{1}{x} + (-by) \dfrac{1}{y} + (ax) \dfrac{1}{z} = 0 \\[2mm]
        (by) \dfrac{1}{x} + (ax) \dfrac{1}{y} + (-cz) \dfrac{1}{z} = 0
        \end{cases}
        \]
        由与存在非零解 $\left(\dfrac{1}{x},\dfrac{1}{y},\dfrac{1}{z}\right)$,则
        \[
        \begin{vmatrix}
        -ax & cz & by \\
        cz & -by & ax \\
        by & ax & -cz
        \end{vmatrix} = 0
        \]
        展开可得
        \[
        a^3x^3 + b^3y^3 + c^3z^3 + abcxyz = 0
        \]
    \end{solution}
    \part 证明 $\displaystyle \frac{yz}{x^2} + \frac{zx}{y^2} + \frac{xy}{z^2} = -1$。
    \begin{solution}
        改写成
        \[
        \begin{cases}
        - a + \dfrac{y}{z} b + \dfrac{z}{y} c = 0 \\[2mm]
        \dfrac{z}{x} a - b + \dfrac{x}{z} c = 0 \\[2mm]
        \dfrac{x}{y} a + \dfrac{y}{x} b - c = 0
        \end{cases}
        \]
        同理,
        \[
        \begin{vmatrix}
        -1 & \dfrac{y}{z} & \dfrac{z}{y} \\[3mm]
        \dfrac{z}{x} & -1 & \dfrac{x}{z} \\[3mm]
        \dfrac{x}{y} & \dfrac{y}{x} & -1
        \end{vmatrix} = 0
        \]
        展开行列式可得
        \[
        \frac{yz}{x^2} + \frac{zx}{y^2} + \frac{xy}{z^2} = -1
        \]
    \end{solution}
    \part 证明 $a^3 + b^3 + c^3 - 5abc = 0$。
    \begin{solution}
        由原方程组,有
        \[
        a^3 = \left(\frac{by}{z}\right)^3 + \left(\frac{cz}{y}\right)^3 + 3abc, \:
        b^3 = \left(\frac{cz}{x}\right)^3 + \left(\frac{ax}{z}\right)^3 + 3abc, \:
        c^3 = \left(\frac{ax}{y}\right)^3 + \left(\frac{by}{x}\right)^3 + 3abc
        \]
        于是
        \[
        2(a^3+b^3+c^3) = (a^3x^3 + b^3y^3 + c^3z^3) \left(\frac{1}{x^3} + \frac{1}{y^3} + \frac{1}{z^3}\right) + 9abc
        \]
        由 $(a),(b)$ 得
        \[
        2(a^3+b^3+c^3) = -abcxyz \left(\frac{1}{x^3} + \frac{1}{y^3} + \frac{1}{z^3}\right) + 9abc = 10abc
        \]
        即得证
        \[
        a^3 + b^3 + c^3 - 5abc = 0
        \]
    \end{solution}
    \end{parts}

    \question 求实数解 $(a, b, c, d)$ 满足方程组
    \[
    \begin{cases}
    a + 4b + 8c + 4d = 53 \\
    3a^2 + 4b^2 + 12 c^2 + 2 d^2 = 159 \\
    9a^3 + 4b^3 + 18 c^3 + d^3 = 477
    \end{cases}
    \]
    \begin{solution}
        由柯西不等式, 
        \[
        53\cdot477=(a + 4b + 8c + 4d)(9a^3 + 4b^3 + 18 c^3 + d^3)\ge (3a^2 + 4b^2 + 12 c^2 + 2 d^2)^2=159^2
        \]
        此时等号成立,有
        \[
        \frac{a}{3a^2} = \frac{4b}{4b^2} = \frac{8c}{12c^2} = \frac{4d}{2d^2}.
        \]
        设
        \[
        \lambda=\frac{1}{3a} = \frac{1}{b} = \frac{2}{3c} = \frac{2}{d}.
        \]
        则
        \[
        a = \frac{1}{3\lambda}, \quad b = \frac{1}{\lambda}, \quad c = \frac{2}{3\lambda}, \quad d = \frac{2}{\lambda}
        \]
        代入第一式得
        \[
        \frac{1}{3\lambda} + \frac{4}{\lambda} + \frac{16}{3\lambda} + \frac{8}{\lambda} =  \frac{29}{3\lambda} = 53 \Rightarrow \lambda = \frac{29}{159}
        \]
        故实数解为
        \[
        \left( \frac{159}{87}, \frac{159}{29}, \frac{106}{87}, \frac{318}{29} \right)
        \]
    \end{solution}
    
    \question 设正实数 $x$, $y$, $z$ 满足
    \[
    \begin{cases}
        x=\sqrt{y^{2}-\dfrac{1}{49}}+\sqrt{z^{2}-\dfrac{1}{49}}\\[3mm]
        y=\sqrt{x^{2}-\dfrac{1}{64}}+\sqrt{z^{2}-\dfrac{1}{64}}\\[3mm]
        z=\sqrt{x^{2}-\dfrac{1}{81}}+\sqrt{y^{2}-\dfrac{1}{81}}
    \end{cases}
    \]
    求 $x+y+z$。
    \ifprintanswers
    \begin{figure}[H]
        \centering        
        \includegraphics[width=0.5\textwidth]{images/image97.jpg}
    \end{figure}
    \fi
    \begin{solution}
        考虑一边长为 $x,y,z$,对应高为 
        \[
        h_x=\frac{1}{7},\quad h_y=\frac{1}{8},\quad h_z=\frac{1}{9}
        \] 
        的$\triangle ABC$,则 $x,y,z$ 满足题意;$\triangle ABC$面积为 
        \[
        S = \frac{1}{2} x h_x = \frac{1}{2} y h_y = \frac{1}{2} z h_z \Rightarrow x:y:z = 7:8:9
        \] 
        设 $x=7k$, $y=8k$, $z=9k$,半周长 $s=\frac{1}{2}(x+y+z)=12k$,面积又为
        \[
        S = \sqrt{12k\cdot 5k \cdot 4k \cdot 3k} = 12\sqrt{5} k^2
        \]
        联立得
        \[S = \frac{1}{2} x h_x = \frac{k}{2} =12\sqrt{5} k^2 \Rightarrow k = \frac{1}{24\sqrt{5}}
        \]
        因此 
        \[
        x+y+z = 24 k = \frac{\sqrt{5}}{5}
        \]
    \end{solution}

    \question 已知正数 $x,y,z$ 满足
    \[
    \begin{cases}
    x^{2}+xy+y^{2}=2 \\
    y^{2}+yz+z^{2}=1 \\
    z^{2}+zx+x^{2}=3 
    \end{cases}
    \]
    求 $xy+yz+zx$。
    \ifprintanswers
    \begin{figure}[H]
        \centering
        \includegraphics[width=0.4\linewidth]{images/image149.png}
    \end{figure}
    \fi
    \begin{solution}
        方程组改写成
        \[
        \begin{cases}
        x^{2}-2xy\cos 120^\circ+y^{2}=(\sqrt{2})^2 \\
        y^{2}-2yz\cos 120^\circ+z^{2}=1^2 \\
        z^{2}-2zx\cos 120^\circ+x^{2}=(\sqrt{3})^2 
        \end{cases}
        \]
        考虑一边长为$AB=\sqrt{2},BC=1,CA=\sqrt{3}$的$\triangle ABC$,其中
        \[
        OA=x,OB=y,OC=z,\angle AOB = \angle BOC = \angle COA = 120^\circ 
        \]
        且满足
        \[
        AB^2+BC^2=CA^2
        \]
        故 $\triangle ABC$ 为直角三角形,面积为
        \[
        [ABC]=\frac12(xy+yz+zx)\sin120^\circ=\frac12\cdot 1\cdot\sqrt{2}
        \]
        解得
        \[
        xy+yz+zx=\frac{2\sqrt{6}}{3}
        \]
    \end{solution}

    \question 已知 \(x,y,z\) 为实数且满足
    \[
    \begin{cases}
    x^{2} + y^{2} = 18 \\
    y^{2} + \sqrt{3} y z + z^{2} = 13 \\
    x^{2} + x z + z^{2} = 19
    \end{cases}
    \]
    求 \(2 x y + y z + \sqrt{3} x z\)。
    \ifprintanswers
    \begin{figure}[H]
        \centering
        \includegraphics[width=0.5\linewidth]{images/image71.png}
    \end{figure}
    \fi
    \begin{solution}
        方程组改写成
        \[
        \begin{cases}
        x^{2} + y^{2} - 2 x y \cos 90^\circ = (3\sqrt{2})^{2} \\
        y^{2} + z^{2} - 2 y z \cos 150^\circ = (\sqrt{13})^{2} \\
        x^{2} + z^{2} - 2 x z \cos 120^\circ = (\sqrt{19})^{2}
        \end{cases}
        \]
        构造\(\triangle ABC\),满足边长:
        \[
        AB = \sqrt{18}, \quad BC = \sqrt{13}, \quad AC = \sqrt{19}, \quad OA = x, \quad OB = y, \quad OC = z
        \]
        及夹角
        \[
        \angle AOB = 90^\circ, \quad \angle BOC = 150^\circ, \quad \angle AOC = 120^\circ
        \]
        设半周长
        \[
        s = \frac{\sqrt{18} + \sqrt{13} + \sqrt{19}}{2},
        \]
        三角形面积为
        \[
        S = \sqrt{s (s - \sqrt{18})(s - \sqrt{13})(s - \sqrt{19})} = 3 \sqrt{\frac{11}{2}}.
        \]
        面积也等于
        \[
        \frac{1}{2} \left( x y \sin 90^\circ + y z \sin 150^\circ + z x \sin 120^\circ \right) = \frac{1}{2} \left( x y + \frac{1}{2} y z + \frac{\sqrt{3}}{2} x z \right).
        \]
        因此
        \[
        \frac{1}{2} \left( x y + \frac{1}{2} y z + \frac{\sqrt{3}}{2} x z \right) = 3 \sqrt{\frac{11}{2}} \implies 2 x y + y z + \sqrt{3} x z = 6 \sqrt{22}.
        \]
    \end{solution}

    \question 已知联立方程组
    \[
    \begin{cases}
    3x^2 + y^2 - 3xy = 3 + 2\sqrt{2} \\
    y^2 + z^2 - yz = 9 + 6\sqrt{2} \\
    z^2 + w^2 + \sqrt{3}zw = 3 + 2\sqrt{2} \\
    w^2 + 3x^2 + \sqrt{3}wx = 9 + 6\sqrt{2}
    \end{cases}, 
    \]
    求$\sqrt{3}xz + yw$之值。
    \ifprintanswers
    \begin{figure}[H]
        \centering
        \includegraphics[width=0.5\linewidth]{images/image115.jpg}
    \end{figure}
    \fi
    \begin{solution}
        设
        \[
        3x^2 + y^2 - 3xy = 3 + 2\sqrt{2} = (\sqrt{2}+1)^2 = a^2,
        \]
        \[
        y^2 + z^2 - yz = 9 + 6\sqrt{2} = (\sqrt{6}+\sqrt{3})^2 = b^2,
        \]
        \[
        z^2 + w^2 + \sqrt{3}zw = 3 + 2\sqrt{2} = a^2,
        \]
        \[
        w^2 + 3x^2 + \sqrt{3}wx = 9 + 6\sqrt{2} = b^2.
        \]
        考虑一圆内接四边形,边长分别为 $\sqrt{3}x$, $y$, $z$, $w$,对角线长为 $a,b$;由托勒密定理,
        \[
        \sqrt{3} xz + yw = (\sqrt{2}+1)(\sqrt{6}+\sqrt{3}) = 3\sqrt{3} + 2\sqrt{6}.
        \]
    \end{solution}

    \question 已知实数$x,y \in (0,1)$满足
    \[
    \begin{cases} 
    \dfrac{1 - \sqrt{1 - x^2}}{x} + \dfrac{2y}{1 + \sqrt{1 - y^2} + y} = 1, \\[2pt]
    25(1 - y^2) = 41 - 40\sqrt{1 - x^2}
    \end{cases}
    \]
    求$(x,y)$的所有解。
    \begin{solution}
        设$x =\sin \alpha,y=\sin \beta, 0 < \alpha,\beta < \dfrac{\pi}{2}$,则第一式变为
        \[
        \frac{1 - \cos \alpha}{\sin \alpha} + \frac{2 \sin \beta}{1 + \cos \beta + \sin \beta} = 1 
        \]
        继续化简得
        \[
        \tan \frac{\alpha}{2} 
        = 1 - \frac{2 \sin \beta}{1 + \cos \beta + \sin \beta}
        = \frac{1 + \cos\beta - \sin\beta}{1 + \cos\beta + \sin\beta}
        = \frac{1 - \tan \frac{\beta}{2}}{1 + \tan \frac{\beta}{2}}
        = \tan \left(\frac{\pi}{4} - \frac{\beta}{2}\right)
        \]
        故
        \[
        \frac{\alpha}{2}=\frac{\pi}{4} - \frac{\beta}{2} +n\pi, n \in \mathbb{Z} \Rightarrow \alpha + \beta = \frac{\pi}{2} \Rightarrow y=\sin\left(\frac{\pi}{2}-\alpha\right)=\cos \alpha
        \]
        于是第二方程可化为
        \[
        25(1 - y^2) = 41 - 40y
        \]
        解得
        \[
        x=\frac{3}{5},\quad y=\frac{4}{5}
        \]
    \end{solution}

    \question 已知实数 $x>0$ 且 $y, z$ 均为实数,求联立方程组
    \[
    \begin{cases}
    5\left(x + \dfrac{1}{x}\right) = 12\left(y + \dfrac{1}{y}\right) = 13\left(z + \dfrac{1}{z}\right), \\
    xy + yz + zx = 1
    \end{cases}
    \]
    的解 $(x,y,z)$。
    \begin{solution}
        由条件$x > 0, \:x,y,z \in \mathbb{R},$可设
        \[
        x = \tan A, \quad y = \tan B, \quad z = \tan C,
        \]
        其中$0 < A < \dfrac{\pi}{2}, \:-\dfrac{\pi}{2} < B,C < \dfrac{\pi}{2}$,于是
        \[
        5\left(x + \frac{1}{x}\right) = 12\left(y + \frac{1}{y}\right) = 13\left(z + \frac{1}{z}\right) \implies 5 \cdot \frac{2}{\sin 2A} = 12 \cdot \frac{2}{\sin 2B} = 13 \cdot \frac{2}{\sin 2C},
        \]
        且
        \[
        \tan A \tan B + \tan B \tan C + \tan C \tan A = 1.
        \]
        因此
        \[
        \frac{5}{\sin 2A} = \frac{12}{\sin 2B} = \frac{13}{\sin 2C}, \quad A + B + C = 90^\circ.
        \]
        由此得
        \[
        \tan 2A = \frac{5}{12}, \quad \tan 2B = \frac{12}{5}, \quad \tan 2C = \infty,
        \]
        故
        \[
        \tan A = \frac{1}{5}, \quad \tan B = \frac{2}{3}, \quad \tan C = 1.
        \]
        即
        \[
        (x,y,z) = \left(\frac{1}{5}, \frac{2}{3}, 1\right)
        \]
    \end{solution}

    \question 已知正实数 $x,y,z$ 满足
    \[
    \begin{cases}
    x+y+z=xyz\\[4pt]
    \dfrac{x^2}{16(1+x^2)}=\dfrac{y^2}{25(1+y^2)}=\dfrac{z^2}{36(1+z^2)}
    \end{cases}
    \]
    求
    \[
    \frac{x^2(1+x^2)^2}{z^2(1+z^2)^2}
    \]
    的值。
    \begin{solution}
        考虑换元$x=\tan A,y=\tan B,z=\tan C$,其中$A,B,C\in \left(0,\dfrac{\pi}{2}\right),$且注意到
        \[
        \tan A+\tan B+\tan C=\tan A\tan B\tan C \Longleftrightarrow A+B+C=\pi
        \]
        于是可以将 $A,B,C$ 视为某三角形的内角;由
        \[
        \frac{x^2}{16(1+x^2)}=\frac{y^2}{25(1+y^2)}=\frac{z^2}{36(1+z^2)}
        \]
        化简得
        \[
        \frac{4}{\sin A}=\frac{5}{\sin B}=\frac{6}{\sin C}
        \]
        由正弦定理,不妨设$\triangle ABC$边长为$a=4k,b=5k,c=6k,k\ne0$,则由余弦定理,
        \[
        \cos A=\frac{(5k)^2+(6k)^2-(4k)^2}{2\cdot5k\cdot6k}=\frac{3}{4},\quad \cos C=\frac{(4k)^2+(5k)^2-(6k)^2}{2\cdot4k\cdot5k}=\frac{1}{8}
        \]
        故\[
        \frac{x^2(1+x^2)^2}{z^2(1+z^2)^2}=\frac{\tan^2A\sec^4A}{\tan^2C\sec^4C}
        =\frac{(\frac{\sqrt{7}}{3})^2(\frac{4}{3})^4}{(\sqrt{63})^2\cdot 8^4}= \frac{1}{186624}
        \]
    \end{solution}

    \question 求所有 $(a, b, c) \in \mathbb{R}^3$ 满足方程组
    \[
    \begin{cases}
    a^2 = \dfrac{b^3 + 9\sqrt{3}}{3b} = \dfrac{c^3 + 16}{3c} \\[6pt]
    b^2 = \dfrac{a^3 - 10}{3a} = \dfrac{c^3 + 28}{3c} \\[6pt]
    c^2 = \dfrac{b^3 + 45\sqrt{3}}{3b} = \dfrac{a^3 - 88}{3a}
    \end{cases}
    \]
    \begin{solution}
        由$(1),(2)$
        \[
        3a^2b-b^3=9\sqrt{3}, \quad a^3-3ab^2=10
        \]
        此时展开式
        \[
        (a\pm b)^3=a^3 \pm 3a^2b + 3ab^2 \pm b^3
        \]
        毫无帮助,不妨考虑
        \[
        (a+bi)^3=a^3-3ab^2+(3a^2b-b^3)i=10+9\sqrt{3}i
        \]
        两边取模得
        \[
        (\sqrt{a^2+b^2})^3=\sqrt{10^2+(9\sqrt{3})^2} \Rightarrow a^2+b^2=7 \tag{3}
        \]
        同理可得
        \[
        b^2+c^2=19,\quad c^2+a^2=20 \tag{4}
        \]
        由$(3),(4)$解得$a^2=4,b^2=3,c^2=16$,经检验得原方程组的解为
        \[
        a=-2,\;b=\sqrt{3},\;c=-4
        \]
    \end{solution}
\end{questions}
\pagebreak
\begin{center}
  {\fontsize{30pt}{26pt}\selectfont
    \hypertarget{取整}{取整} \label{取整}
  }
\end{center}
\separator
\vspace{1pt}
\begin{questions}
    \question 求满足
    \[
    \floor{0.5+\floor{x}} = 20
    \]
    的所有$x$的取值范围。
    \begin{solution}
        由下取整定义,原方程即
        \[
        20 \le 0.5+\floor{x}< 21 \Rightarrow 19.5 \le \floor{x} < 20.5
        \]
        由于$\floor{x} \in \mathbb{Z}$,
        \[
        \floor{x}=20 \Rightarrow 20 \le x < 21
        \]
    \end{solution}

    \question 求满足
    \[
    \ceil{\ceil{x}-1.3} = 16
    \]
    的所有$x$。
    \begin{solution}
        由上取整定义,原方程即
        \[
        15 \le \ceil{x} -1.3 < 16 \Rightarrow 16.3 \le \ceil{x} < 17.3
        \]
        由于$\ceil{x} \in \mathbb{Z}$,
        \[
        \ceil{x}=17 \Rightarrow 16 < x \le 17
        \]
    \end{solution}

    \question 解方程 
    \[
    2x = \floor{x} + \frac{16}{5}
    \]
    \begin{solution}
        发现$x\neq \floor{x}$,由于 
        \[
        \floor{x} < x < \floor{x} + 1 \Rightarrow 2\floor{x} < 2x < 2\floor{x} + 2,
        \]
        观察右式的取值范围可知
        \[ 
        2\floor{x} < \floor{x} + \frac{16}{5} < 2\floor{x} + 2 \Rightarrow \frac{6}{5} < \floor{x} < \frac{16}{5} 
        \]
        因此 
        \[
        \floor{x} = 2, 3 \Rightarrow x = \frac{13}{5}, \frac{31}{10}
        \]
    \end{solution}

    \question 解方程
    \[
    2\lfloor x \rfloor + 3x = 4 - 5\{x\}
    \]
    \begin{solution}
        设 $x = n+r$, 其中 $n = \lfloor x \rfloor,r = \{x\}$,代入方程得
        \[
        2n + 3(n+r) = 4 - 5r \Rightarrow 5x=4-3r
        \]
        由 $0 \le r < 1$, 有
        \[
        1 < 4-3r \le 4 \Rightarrow \frac{1}{5} < x \le \frac{4}{5}
        \]
        由此可得$x=\{x\}$,即$\floor{x}=0$,此时
        \[
        2\cdot 0 +3x=4-5x \Rightarrow x=\frac{1}{2}
        \]
    \end{solution}

    \question 令 
    \[
    f(x) = \left\lfloor x + \frac{1}{2} \right\rfloor + \left\lfloor x - \frac{1}{2} \right\rfloor - \lfloor 2x \rfloor,
    \]
    证明只有当 $\dfrac{1}{2} \le \{x\} < 1$ 时,满足方程 $f(x) = 0$ 。
    \begin{solution}
        设 $x = n + \{x\}$ ,其中 $n$ 为整数,且$0 \le \{x\} < 1$,将$n$从取整符号中提出,
        \[ 
        f(x) = \left\lfloor \{x\} + \frac{1}{2} \right\rfloor + \left\lfloor \{x\} - \frac{1}{2} \right\rfloor - \lfloor 2\{x\} \rfloor 
        \]
        当 $0 \le \{x\} < \dfrac{1}{2}$ 时,
        \[ 
        f(x) = 0 + (-1) - 0 = -1 \neq 0
        \]
        而当 $\dfrac{1}{2} \le \{x\} < 1$ 时,
        \[
        f(x) = 1 + 0 - 1 = 0
        \]
        得证只有当 $\dfrac{1}{2} \le \{x\} < 1$ 时,$f(x) = 0$。
    \end{solution}

    \question 解 
    \[
    8\lfloor 3x \rfloor - 5\lfloor 2x \rfloor = 3
    \]
    \begin{solution}
        令 $x = n + r$ ,其中 $n \in \mathbb{Z}, 0 \le r < 1$,原方程变为 
        \[ 
        3n + 8\lfloor 3r \rfloor - 5\lfloor 2r \rfloor = 3 
        \]
        当 $0 \le r < \dfrac{1}{3}$ 时,$\lfloor 3r \rfloor = 0, \lfloor 2r \rfloor = 0$,有
        \[
        3n = 3 \Rightarrow n = 1
        \]
        当 $\dfrac{1}{3} \le r < \dfrac{1}{2}$ 时,$\lfloor 3r \rfloor = 1, \lfloor 2r \rfloor = 0$,有
        \[
        3n + 8 = 3 \Rightarrow n = -\frac{5}{3} \notin \mathbb{Z}
        \]
        当 $\dfrac{1}{2} \le r < \dfrac{2}{3}$ 时, $\lfloor 3r \rfloor = 1, \lfloor 2r \rfloor = 1$,有
        \[
        3n + 8 - 5 = 3 \Rightarrow n = 0
        \]
        当 $\dfrac{2}{3} \le r < 1$ 时,$\lfloor 3r \rfloor = 2, \lfloor 2r \rfloor = 1$,有
        \[
        3n + 16 - 5 = 3 \Rightarrow n = -\frac{8}{3} \notin \mathbb{Z}
        \]
        综上,解为
        \[
        n = 1, 0 \le r < \frac{1}{3} \Rightarrow 1 \le x < \frac{4}{3}
        \]
        或
        \[
        n = 0, \frac{1}{2} \le r < \frac{2}{3} \Rightarrow \frac{1}{2} \le x < \frac{2}{3}
        \]
    \end{solution}

    \question 已知
    \[
    \begin{cases}
    \lfloor b \rfloor + \lceil a \rceil + \{c\} = 16 \\
    \lfloor c \rfloor + \lceil b \rceil + \{a\} = 11.3 \\
    \lfloor a \rfloor + \lceil c \rceil + \{b\} = 9.7
    \end{cases}
    \]
    求$a+b+c$的值。
    \begin{solution}
        由题意可知,小数部分分别为
        \[
        \{a\}=0.3,\quad \{b\}=0.7,\quad \{c\}=0
        \]
        原方程组变为
        \[
        \begin{cases}
        \lfloor b \rfloor + \lfloor a \rfloor + 1 = 16 \\
        \lfloor c \rfloor + \lfloor b \rfloor + 1 = 11 \\
        \lfloor a \rfloor + \lfloor c \rfloor = 9
        \end{cases}
        \]
        由第二式与第三式可得
        \[
        \lfloor a \rfloor =7, \quad \lfloor b \rfloor = 9-1=8, \quad \lfloor c \rfloor = 2
        \]
        因此
        \[
        a=7.3,\quad b=8.7,\quad c=2 \Rightarrow a+b+c=18
        \]
    \end{solution}

    \question 已知 $\dfrac{1}{3-\sqrt{7}}$ 的整数部分为 $a$ ,小数部分为 $b$ ,求 $a^2 + (1+\sqrt{7})ab$ 。
    \begin{solution}
        有理化得
        \[ 
        \frac{1}{3-\sqrt{7}} = \frac{3+\sqrt{7}}{2} 
        \]
        由 
        \[
        2 < \sqrt{7} < 3 \Rightarrow \frac{5}{2} < \frac{3+\sqrt{7}}{2} < 3,
        \]
        得 
        \[
        a = 2 ,\quad b = \frac{3+\sqrt{7}}{2} - a = \frac{\sqrt{7}-1}{2}
        \]
        故
        \begin{align*}
        a^2 + (1+\sqrt{7})ab &= 4 + (1+\sqrt{7}) \cdot 2 \cdot \frac{\sqrt{7}-1}{2} \\
        &= 4 + (\sqrt{7}+1)(\sqrt{7}-1) \\
        &= 10
        \end{align*}
    \end{solution}

    \question 求大于 $(\sqrt{3}+\sqrt{2})^{6}$ 的最小整数。
    \begin{solution}
        令
        \[
        a=\sqrt{3}+\sqrt{2},\quad b=\sqrt{3}-\sqrt{2}
        \]
        则
        \[
        a^2 = 5+2\sqrt{6},\quad b^2 = 5-2\sqrt{6},\quad ab = 1
        \]
        所以
        \[
        a^6 + b^6 = (a^2 + b^2)(a^4 + b^4 - a^2b^2) = 10(10^2-3) = 970
        \]
        由于 $0 < b^2 < 1$,所以
        \[
        969 < a^6 < 970 \quad \Rightarrow \quad \ceil{a^6} = 970
        \]
    \end{solution}

    \question 求满足
    \[
    \lfloor \sqrt{x}\rfloor=\lfloor \sqrt{x+34}\rfloor
    \]
    的最小整数$x$。
    \begin{solution}
        设$y,y+1$为连续正整数,使得他们的平方根的差不小于34,则
        \[
        (y+1)^2-y^2>34 \Rightarrow y>16.5
        \]
        由于$y \in \mathbb{Z}$,最小的$x$为
        \[
        x=y^2=17^2=289
        \]
    \end{solution}

    \question 解
    \[
    \lceil x\rceil \lceil 2x\rceil = 15
    \]
    \begin{solution}
        设
        \[
        x=n-r, \quad n \in \mathbb{Z}, \quad 0<r\le 1
        \]
        当 $r < \dfrac{1}{2}$,原方程即
        \[
        n(2n) = 15 \Rightarrow n = \sqrt{\frac{15}{2}} \not \in \mathbb{Z}
        \]
        当  $r \ge \dfrac{1}{2}$,原方程即
        \[
        n(2n-1) = 15 \Rightarrow (n-3)(2n+5) = 0 \Rightarrow n = 3 \in \mathbb{Z}
        \]
        故$x$的取值范围为
        \[
        2 < x \le 2.5
        \]
    \end{solution}

    \question 求正实数$x$满足
    \[
    x^2+\{x\}^2=27
    \]
    \begin{solution}
        设
        \[
        x=n+r, \quad n \in \mathbb{Z}, \quad 0<r\le 1
        \]
        且注意到
        \[
        x^2 \le x^2+\{x\}^2 < x^2+1
        \]
        因此
        \[
        26<x^2\le27 \Rightarrow n=5
        \]
        代入原方程得
        \[
        (5+r)^2+r^2=27 \Rightarrow r=\frac{-5+\sqrt{29}}{2}
        \]
        于是
        \[
        x=5+r=\frac{5+\sqrt{29}}{2}
        \]
    \end{solution}

    \question 解
    \[
    x^2 - 6\lfloor x \rfloor + 5 = 0
    \]
    \begin{solution}
        原方程化为
        \[
        x^2 - 6\lfloor x \rfloor + 5 = 0 \Rightarrow x^2 - 6x + 5 = -6 \{x\}
        \]
        $-6 \{x\}$的值域为$(-6,0]$,且对于$x \in [1,5]$,有
        \[
        x^2 - 6x + 5 \in (-6,0]
        \]
        当$x \in [1,2)$,有$\floor{x}=1$,则
        \[
        x^2=6\cdot1-5 \Rightarrow x=1
        \]
        类似地,
        \begin{align*}
        \floor{x}&=2 \Rightarrow x^2 = 7 \Rightarrow x = \sqrt{7} \\
        \floor{x}&=3 \Rightarrow x^2 = 13 \Rightarrow x = \sqrt{13} \\
        \floor{x}&=4 \Rightarrow x^2 = 19 \Rightarrow x = \sqrt{19} \\
        \floor{x}&=5 \Rightarrow x^2 = 25 \Rightarrow x = 5 
        \end{align*}
        因此,方程的解为
        \[
        x = 1, \sqrt{7}, \sqrt{13}, \sqrt{19}, 5
        \]
    \end{solution}

    \question 解
    \[
    \lfloor x^2 \rfloor - \lfloor x \rfloor^2 = 1999
    \]
    \begin{solution}
        设 $x = n + r$, 其中 $n = \lfloor x \rfloor$, $r = \{x\}$,则
        \[
        \lfloor x^2 \rfloor - n^2 = \floor{(n+r)^2-n^2} =\lfloor 2 n r + r^2 \rfloor
        \]
        故
        \[
        \lfloor x^2 \rfloor - \lfloor x \rfloor^2 = 1999 \Rightarrow 2 n r + r^2 \ge 1999
        \]
        最小的 $n$ 滿足
        \[
        2n r \ge 1999 \implies n \ge \lceil 1999/2 \rceil = 1000.
        \]
        此时有
        \[
        r^2 + 2000 r - 1999 = 0 \Rightarrow r = -1000 + \sqrt{1001999}
        \]
        所以
        \[
        x = n + r = \sqrt{1001999}
        \]
    \end{solution}

    \question 解 
    \[
    \left \lfloor \frac{4x+3}{5} \right \rfloor = \frac{3(2x+1)}{4}
    \]    
    \begin{solution}
        令 
        \[
        m = \frac{3(2x+1)}{4} \Rightarrow x = \frac{4m-3}{6},
        \]
        代回原方程化简得到 
        \[ 
        \left \lfloor \frac{8m+3}{15} \right \rfloor= m 
        \]
        有 
        \[
        0 \le \frac{8m+3}{15} - m < 1  \Rightarrow -\frac{12}{7} < m \le \frac{3}{7} 
        \]
        由于$m \in \mathbb{Z}$于是
        \[ 
        m = 0, -1 \Rightarrow x = -\frac{1}{2}, -\frac{7}{6}
        \]
    \end{solution}

    \question 求满足以下条件的正整数 $x$ 的个数:
    \[
    \left\lfloor \frac{x}{99} \right\rfloor = \left\lfloor \frac{x}{101} \right\rfloor.
    \]
    \begin{solution}
        设
        \[
        \left\lfloor \frac{x}{99} \right\rfloor = \left\lfloor \frac{x}{101} \right\rfloor = m \in \mathbb{Z}
        \]
        由定义有
        \[
        m \le \frac{x}{99} < m+1 \Rightarrow 99m \le x < 99(m+1)
        \]
        \[
        m \le \frac{x}{101} < m+1 \Rightarrow 101m \le x < 101(m+1)
        \]
        因此
        \[
        101m \le x < 99(m+1)
        \]
        当 $m>49$ 时, $101m > 99(m+1)$, 无解。对于 $0 \le m \le 49$, 由于 $x=0$ 不是正整数,可行整数个数为
        \[
        \sum_{m=0}^{49} (99(m+1)-101m) =\sum_{m=0}^{49} (99-2m) -1 = 99 \cdot 50 - 2 \cdot \frac{49 \cdot 50}{2} -1 = 2499
        \]
    \end{solution}

    \question 求解 $x\in \mathbb{Z}$ 满足
    \[
    \left\lfloor \frac{x}{1!} \right\rfloor + \left\lfloor \frac{x}{2!} \right\rfloor + \left\lfloor \frac{x}{3!} \right\rfloor = 224
    \]
    \begin{solution}
        原方程给出
        \[
        \frac{x}{1} + \frac{x}{2} + \frac{x}{6} -3 < 224 < \frac{x}{1} + \frac{x}{2} + \frac{x}{6}
        \]
        化简得
        \[
        134.4 \le x < 136.2
        \]
        由于 $x \in \mathbb{Z}$,且 $x=136$ 不满足原方程,故原方程式的解只有
        \[
        x = 135
        \]
    \end{solution}

    \question 求满足
    \[
    \frac{\lfloor x \rfloor}{x} = \frac{9}{10}
    \]
    的最大实数 $x$。
    \begin{solution}
        由 $x = \lfloor x \rfloor + \{x\}$, 交叉相乘得
        \[
        10 \lfloor x \rfloor = 9 (\lfloor x \rfloor + \{x\}) \Rightarrow \lfloor x \rfloor = 9 \{x\}
        \]
        由于 $0 \le \{x\} < 1,\{x\}$的可能值为
        \[
        \frac{1}{9}, \frac{2}{9}, \dots, \frac{8}{9}, 
        \]
        取最大值 $\{x\} = \dfrac{8}{9}$, 于是
        \[
        x = \lfloor x \rfloor + \{x\} = \frac{80}{9}
        \]
    \end{solution}
    
    \question 18. 如 $x > 0$ 且满足 $\lfloor x \rfloor^2 = x(x - \lfloor x \rfloor)$ ,求 $x - \frac{1}{x}$ 。
    \begin{solution}
        设 $n = \lfloor x \rfloor$。由于 $x > 0$,故 $n \ge 0$。
        方程变为
        \[ 
        n^2 = x(x - n) \Rightarrow x^2 - nx - n^2 = 0
        \]
        由于 $x > 0$,取正根
        \[ 
        x = \frac{(1 + \sqrt{5})n}{2} 
        \]
        根据 $\lfloor x \rfloor = n$ 的定义,必须满足
        \[ 
        n \le \frac{(1 + \sqrt{5})n}{2} < n + 1 \tag{1}
        \]
        对于 $n=0$,(1) 成立但 $x=0$,与 $x>0$ 矛盾。对于 $n \ge 1$,(1)左侧 
        \[
        \frac{1+\sqrt{5}}{2} \approx 1.618 \ge 1
        \]
        恒成立,而右侧
        \[
        \frac{(1 + \sqrt{5})n}{2} < n + 1 \implies 0.618n < 1 \implies n < \frac{1 + \sqrt{5}}{2}
        \]
        故只能$n = 1$,此时
        \[
        x - \frac{1}{x} = \frac{1 + \sqrt{5}}{2} - \frac{2}{1 + \sqrt{5}} = \frac{1 + \sqrt{5}}{2} - \frac{\sqrt{5} - 1}{2} = 1
        \]
    \end{solution}

\question 证明方程 $\lfloor x \rfloor + \lfloor 2x \rfloor + \lfloor 4x \rfloor + \lfloor 8x \rfloor = 147$ 无实数解。
\begin{solution}
    设 $x = n + \alpha$,其中 $n = \lfloor x \rfloor$ 为整数部分,$\alpha \in [0, 1)$ 为小数部分。
    将 $\alpha$ 表示为二进制形式:
    \[ \alpha = (0.b_1 b_2 b_3 b_4 \dots)_2 = \frac{b_1}{2} + \frac{b_2}{4} + \frac{b_3}{8} + \frac{b_4}{16} + \dots \]
    其中 $b_i \in \{0, 1\}$。

    利用性质 $\lfloor m + y \rfloor = m + \lfloor y \rfloor$（$m$ 为整数）以及二进制移位性质:
    \begin{align*}
    \lfloor x \rfloor &= n \\
    \lfloor 2x \rfloor &= 2n + \lfloor 2\alpha \rfloor = 2n + b_1 \\
    \lfloor 4x \rfloor &= 4n + \lfloor 4\alpha \rfloor = 4n + 2b_1 + b_2 \\
    \lfloor 8x \rfloor &= 8n + \lfloor 8\alpha \rfloor = 8n + 4b_1 + 2b_2 + b_3
    \end{align*}
   
    将上述结果代入原方程左式（LHS）:
    \begin{align*}
    \text{LHS} &= (n + 2n + 4n + 8n) + (b_1 + 2b_1 + 4b_1) + (b_2 + 2b_2) + b_3 \\
    &= 15n + 7b_1 + 3b_2 + b_3
    \end{align*}

    已知方程为 $15n + 7b_1 + 3b_2 + b_3 = 147$。
    由于 $b_1, b_2, b_3 \in \{0, 1\}$,系数部分 $7b_1 + 3b_2 + b_3$ 的最大取值为:
    \[ 7(1) + 3(1) + 1(1) = 11 \]
    
    这意味着对于特定的 $n$,函数值的波动范围仅为 $[15n, 15n + 11]$。

    若 $n = 9$,最大值为 $15(9) + 11 = 135 + 11 = 146$。
    若 $n = 10$,最小值为 $15(10) + 0 = 150$。
    
    数值 147 恰好落在两个区间之间（$146 < 147 < 150$）,因此该方程在实数范围内无解。
    \textcolor{red}{(待验证)}
\end{solution}

    \question 解
    \[
    \lceil x \lfloor x \rfloor \rceil + \lfloor x \lceil x \rceil \rfloor = 111
    \]
    \begin{solution}
        令 $x = n + r$,其中 $n\in \mathbb{Z}, 0 \le r < 1$,则
        \[
        \lceil (n+r) \lfloor(n+r) \rfloor\rceil + \lfloor (n+r)\lceil(n+r)\rceil \rfloor = 111
        \]
        展开得
        \[
        2n^2 +n + \lceil nr \rceil + \lfloor (n+1)r \rfloor = 111
        \]
        解不等式组
        \[
        \begin{cases}
        2n^2+n+1 \le 111 \\
        2n^2+3n \ge 111
        \end{cases}
        \]
        得$n=7$,因此解
        \[
        \lceil 7r \rceil + \lfloor 8r \rfloor =6
        \]
        得
        \[
        \frac{3}{8} \le r \le \frac{3}{7} \Rightarrow \frac{59}{8} \le x \le \frac{52}{7}
        \]
    \end{solution}

    \question 
    \begin{parts}
    \part 求方程
    \[
    \left\lfloor \frac{n}{2} \right\rfloor + \left\lfloor \frac{n}{3} \right\rfloor + \left\lfloor \frac{n}{6} \right\rfloor = n
    \]
    的正整数解个数,其中$1 \le n \le 100$。
    \begin{solution}
        原方程化为
        \[
        \left\lfloor \frac{n}{2} \right\rfloor + \left\lfloor \frac{n}{3} \right\rfloor + \left\lfloor \frac{n}{6} \right\rfloor = \frac{n}{2} + \frac{n}{3} + \frac{n}{6}
        \]
        明示了
        \[
        \frac{n}{2},\frac{n}{3},\frac{n}{6}
        \]
        皆为整数, 于是$n$ 必须为 6 的倍数,因此正整数解共有 $\left\lfloor\dfrac{100}{6}\right\rfloor=16$ 个。
    \end{solution}
    \part 求方程
    \[
    \left\lfloor \frac{n}{2} \right\rfloor + \left\lfloor \frac{n}{3} \right\rfloor + \left\lfloor \frac{n}{6} \right\rfloor = n-1
    \]
    的正整数解个数,其中$1 \le n \le 100$。
    \begin{solution}
        设 
        \[
        f(n) = \left\lfloor \frac{n}{2} \right\rfloor + \left\lfloor \frac{n}{3} \right\rfloor + \left\lfloor \frac{n}{6} \right\rfloor, 
        \]
        则
        \[
        f(n+6) = f(n) + 6
        \]
        定义 $g(n) = f(n) - (n-1)$, 则
        \[
        g(n+6) = g(n).
        \]
        且有
        \[
        g(1) = g(2) = g(3) = g(4) = 0, \quad g(5) \neq 0, \quad g(6) \neq 0.
        \]
        因此正整数解共有 
        \[
        17+17+17+17=68
        \]
        个。
    \end{solution}
    \end{parts}

    \question 求数列
    \[
    \left\lfloor \frac{1^2}{2016} \right\rfloor, \left\lfloor \frac{2^2}{2016} \right\rfloor, \dots, \left\lfloor \frac{2016^2}{2016} \right\rfloor
    \]
    中相异整数的个数。
    \begin{solution}
        考虑连续两项的差
        \[
        \frac{n^2}{2016} - \frac{(n-1)^2}{2016} = \frac{2n-1}{2016} < 1 \implies n < 1007.5
        \]
        由于
        \[
        \frac{1007^2}{2016} \approx 503.0005 > 503
        \]  
        前 1007 项即整数 $0,1,2,\dots,503$,共有 $504$ 个不同整数。对于第 1008 项到 2016 项,
        每项至少比前一项大 1,因此产生 
        \[
        2016-1008+1 = 1009
        \]
        个不同整数,因此数列中相异整数的个数为
        \[
        504 + 1009 = 1513
        \]
    \end{solution}

    \question 求 
    \[
    \sum_{k=1}^{512} \lfloor \log_2 k \rfloor
    \]
    \begin{solution}
        一般地,对所有整数$n\ge 0,k \ge 1$,
        \[
        2^n \le k \le 2^{n+1} \lfloor \log_2 k \rfloor = n
        \]
        因此
        \begin{align*}
        \sum_{k=1}^{512} \lfloor \log_2 k \rfloor &= 0 + (4-2) \cdot 1 + (8-4) \cdot 2 + \dots + (512-256) \cdot 8 + 9 \\
        &= 2 \cdot 1 + 2^2 \cdot 2 + \dots + 2^8 \cdot 8 + 9 \\
        &= 3595
        \end{align*}
    \end{solution}

    \question 计算
    \[
    \sum_{r=1}^{34}\left\lfloor \frac{18r}{35}\right\rfloor
    \]
    \begin{solution}
        将整数部分与小数部分分开,
        \[
        \sum_{r=1}^{34}\left\lfloor \frac{18r}{35}\right\rfloor
        =\sum_{r=1}^{34}\frac{18r}{35}
        -\sum_{r=1}^{34}\left\{\frac{18r}{35}\right\}
        \]
        其中
        \[
        \sum_{r=1}^{34}\frac{18r}{35}
        =\frac{18}{35}\cdot\frac{34\cdot35}{2}
        =306
        \]
        注意到对任意正整数$1 \le r \le 34$,
        \[
        \left\{\frac{18r}{35}\right\}+\left\{\frac{18(35-r)}{35}\right\}=1
        \]
        当$r$从$1$到$34$时,可以配成$17$对,因此
        \[
        \sum_{r=1}^{34}\left\{\frac{18r}{35}\right\}=17
        \]
        于是
        \[
        \sum_{r=1}^{34}\left\lfloor \frac{18r}{35}\right\rfloor
        =289
        \]
    \end{solution}
    \begin{solution}
        对任意正整数$1 \le r \le 34$,有
        \[
        \frac{18r}{35}+\frac{18(35-r)}{35}=18
        \]
        因此
        \[
        \left\lfloor \frac{18r}{35}\right\rfloor
        +\left\lfloor \frac{18(35-r)}{35}\right\rfloor
        +\left\{\frac{18r}{35}\right\}
        +\left\{\frac{18(35-r)}{35}\right\}=18
        \]
        又因为
        \[
        \left\{\frac{18r}{35}\right\}
        +\left\{\frac{18(35-r)}{35}\right\}=1
        \]
        所以
        \[
        \left\lfloor \frac{18r}{35}\right\rfloor
        +\left\lfloor \frac{18(35-r)}{35}\right\rfloor
        =17
        \]
        因此
        \[
        \sum_{r=1}^{34}\left\lfloor \frac{18r}{35}\right\rfloor=17\cdot 17=289
        \]
    \end{solution}

    \question 设$r$为实数,且满足
    \[
    \left\lfloor r+\frac{19}{100}\right\rfloor
    +\left\lfloor r+\frac{20}{100}\right\rfloor
    +\cdots
    +\left\lfloor r+\frac{91}{100}\right\rfloor
    =546
    \]
    求$\lfloor100r\rfloor$的值。
    \begin{solution}
        从$\dfrac{19}{100}$到$\dfrac{91}{100}$,共有
        \[
        91-19+1=73
        \]
        项,这些取整值只可能为$7$或$8$。

        若全部等于$7$,则和为$73\cdot 7=521$;若全部等于$8$,则和为$73\cdot 8=581$,

        设前$k$项等于$7$,其余$79-k$项等于$8$,则
        \[
        7k+8(73-k)=546 \Rightarrow k=38
        \]
        这表示
        \[
        \left\lfloor r+\frac{56}{100}\right\rfloor=7,\quad
        \left\lfloor r+\frac{57}{100}\right\rfloor=8
        \]
        于是
        \[
        \floor{100r}=\floor{r}+\cdots + \left\lfloor r+\frac{56}{100}\right\rfloor+\left\lfloor r+\frac{57}{100}\right\rfloor+\cdots +\left\lfloor r+\frac{99}{100}\right\rfloor=57\cdot7+43\cdot8=743
        \]
    \end{solution}

    \question 证明 
    \[
    \lfloor \sqrt{1} \rfloor + \lfloor \sqrt{2} \rfloor + \dots + \lfloor \sqrt{n^2 - 1} \rfloor = \frac{n(n-1)(4n+1)}{6}
    \]
    \begin{solution}
        考虑满足 
        \[
        \lfloor \sqrt{k} \rfloor = m \tag{1}
        \]
        的正整数 $k$ 的个数。(1)等价于
        \[
        m \le \sqrt{k} < m+1 \Longleftrightarrow m^2 \le k < (m+1)^2 
        \]
        因此,对于每一个确定的 $m$ ,满足条件的 $k$ 有 
        \[
        (m+1)^2 - m^2 = 2m+1
        \]
        个,于是
        \begin{align*}
        \sum_{k=1}^{n^2-1} \lfloor \sqrt{k} \rfloor &= \sum_{m=1}^{n-1} m(2m+1) \\
        &= 2 \sum_{m=1}^{n-1} m^2 + \sum_{m=1}^{n-1} m \\
        &= 2 \cdot \frac{(n-1)n(2n-1)}{6} + \frac{(n-1)n}{2} \\
        &= \frac{n(n-1)(4n+1)}{6}
        \end{align*}
        得证。
    \end{solution}

    \question 计算
    \[
    \ceil{\sqrt{1}} + \ceil{\sqrt{2}} + \ceil{\sqrt{3}} + \dots + \ceil{\sqrt{2025}}
    \]
    \begin{solution}
        注意到$\sqrt{2025}=45$,当$(m-1)^2<k\le m^2$时有
        \[
        \lceil \sqrt{k}\rceil=m
        \]
        因此可按 $m=1,2,\dots,45$ 分组求和,
        \[
        \sum_{k=1}^{2025}\lceil \sqrt{k}\rceil
        =\sum_{m=1}^{45}\sum_{k=(m-1)^2+1}^{m^2} m
        \]
        对固定的 $m$, 求和的项数为$m^2-(m-1)^2=2m-1$,于是
        \[
        \sum_{k=1}^{2025}\lceil \sqrt{k}\rceil
        =\sum_{m=1}^{45} m(2m-1)
        =2\sum_{m=1}^{45}m^2-\sum_{m=1}^{45}m
        =2\cdot\frac{45\cdot46\cdot91}{6}-\frac{45\cdot46}{2}
        =61755
        \]
    \end{solution}

    \question 求
    \[
    \lceil 1\rceil+\lceil 1.7\rceil+\lceil 2.4\rceil+\lceil 3.1\rceil+\dots+\lceil 999.9\rceil
    \]
    \begin{solution}
        首先发现
        \[
        \lceil 1\rceil+\lceil 8\rceil+\dots+\lceil 995\rceil
        =\frac{143}{2}(1+995)
        \]
        \[
        \lceil 1.7\rceil+\lceil 8.7\rceil+\dots+\lceil 995.7\rceil
        =\frac{143}{2}(2+996)
        \]
        以此类推,原式即
        \begin{align*}
        &\lceil 1\rceil+\lceil 8\rceil+\dots+\lceil 995\rceil+\lceil 1.7\rceil+\lceil 8.7\rceil+\dots+\lceil 995.7\rceil \\
        &+\cdots + \lceil 6.6\rceil+\lceil 13.6\rceil+\dots+\lceil 993.6\rceil+\lceil 7.3\rceil+\lceil 14.3\rceil+\dots+\lceil 994.3\rceil \\
        &=\frac{143}{2}(1+995+2+996+3+997+4+998+4+998+5+999+6+1000+6+1000)\\
        &\quad+\frac{142}{2}(7+994+8+995)\\
        &=715285
        \end{align*}
    \end{solution}

    \question 求
    \[
    \sum_{n=0}^{1000} \left\lfloor \frac{2^n}{3} \right\rfloor
    \]
    \begin{solution}
        利用二项式展开,
        \begin{align*}
        a_n &= \left\lfloor \frac{(3-1)^n}{3} \right\rfloor \\
        &= \left\lfloor 3^{n-1} - n \cdot 3^{n-2} + n(n-1)3^{n-3} + \dots + (-1)^{n-1} + \frac{(-1)^n}{3} \right\rfloor \\
        &= 3^{n-1} - n \cdot 3^{n-2} + n(n-1)3^{n-3} + \dots + (-1)^{n-1} 
        \end{align*}
        且注意到 
        \[
        a_n 
        = \begin{cases}
        \dfrac{2^n-1}{3}, \quad n\text{为偶数} \\[8pt]
        \dfrac{2^n-2}{3}, \quad n\text{为奇数}
        \end{cases}
        = \frac{2^n}{3} - \frac{1}{2} + \frac{(-1)^n}{6}
        \]
        因此
        \begin{align*}
        \sum_{n=0}^{1000} \left\lfloor \frac{2^n}{3} \right\rfloor
        &= \sum_{n=0}^{1000} \left( \frac{2^n}{3} - \frac{1}{2} + \frac{(-1)^n}{6} \right) \\
        &= \frac{1}{3} \cdot \frac{2^{1001}-1}{2-1} - \frac{1001}{2} + \frac{1}{6} \\
        &= \frac{2^{1001}-1502}{3}
        \end{align*}
    \end{solution}

    \question 设 $m, n$ 是整数,$n \ge 1$ ,证明 
    \[
    \left\lfloor \dfrac{m+1}{n} \right\rfloor = 
    \begin{cases} 
    \left\lfloor \dfrac{m}{n} \right\rfloor &,n \nmid m+1 \\[8pt]
    \left\lfloor \dfrac{m}{n} \right\rfloor + 1 &,n \mid m+1
    \end{cases}
    \]
    \begin{solution}
        当 $n$ 不整除 $m+1$,有$m+1 = kn+r$ ,其中 $k \in \mathbb{Z},1 \le r \le n-1$,此时欲证 
        \[
        \left\lfloor \frac{m+1}{n} \right\rfloor = \left\lfloor \frac{m}{n} \right\rfloor \tag{1}
        \]
        观察左式为
        \[
        \left\lfloor \frac{m+1}{n} \right\rfloor = \left\lfloor \frac{kn+r}{n} \right\rfloor = \left\lfloor k + \frac{r}{n} \right\rfloor = k 
        \]
        右式为
        \[
        \left\lfloor \frac{m}{n} \right\rfloor = \left\lfloor \frac{kn+r-1}{n} \right\rfloor = \left\lfloor k + \frac{r-1}{n} \right\rfloor = k
        \]
        故(1)得证。当 $n$ 整除 $m+1$ 时,即 $m+1 = kn$,欲证
        \[
        \left\lfloor \frac{m+1}{n} \right\rfloor = \left\lfloor \frac{m}{n} \right\rfloor + 1 \tag{2}
        \]
        此时左式为
        \[
        \left\lfloor \frac{m+1}{n} \right\rfloor = \left\lfloor \frac{kn}{n} \right\rfloor = \left\lfloor k \right\rfloor = k
        \]
        右式为
        \[
        \left\lfloor \frac{m}{n} \right\rfloor = \left\lfloor \frac{kn-1}{n} \right\rfloor = \left\lfloor k - \frac{1}{n} \right\rfloor = k-1 
        \]
        故(2)得证。
    \end{solution}

    \question 证明对任意实数 $\alpha, \beta$ 有 
    \[
    \lfloor 2\alpha \rfloor + \lfloor 2\beta \rfloor \ge \lfloor \alpha \rfloor + \lfloor \beta \rfloor + \lfloor \alpha + \beta \rfloor
    \]
    \begin{solution}
        令 
        \[
        \alpha = n_1 + r_1,\beta = n_2 + r_2, \quad n_1, n_2 \in \mathbb{Z},\quad 0 \le r_1, r_2 < 1
        \]
        则原不等式等价于
        \[ 
        \lfloor 2r_1 \rfloor + \lfloor 2r_2 \rfloor \ge \lfloor r_1 + r_2 \rfloor \tag{1}
        \]
        若 $0 \le r_1 < \dfrac{1}{2}$ 且 $0 \le r_2 < \frac{1}{2}$ ,则左式 $= 0 + 0 = 0$ 。由于 $r_1 + r_2 < 1$ ,右式 $= 0$,(1)成立；

        若 $r_1, r_2$ 中有一个大于等于 $\frac{1}{2}$ ,不妨设 $r_1 \ge \dfrac{1}{2}$ ,则左式 $\ge 1 + 0 = 1$ 。而 $r_1 + r_2 < 2$ ,右式最大为 $1$ ,(1)成立；

        若 $r_1 \ge \dfrac{1}{2}$ 且 $r_2 \ge \frac{1}{2}$ ,则左式 $= 1 + 1 = 2$ 。由于 $r_1 + r_2 < 2$ ,右式最大为 $1$ ,(1)也成立。

        综上,原不等式成立。
    \end{solution}

    \question 证明对任意实数 $x, y$ 有 
    \[
    \lfloor x - y \rfloor \le \lfloor x \rfloor - \lfloor y \rfloor \le \lfloor x - y \rfloor + 1
    \]
    \begin{solution}
        由取整函数的性质可知 
        \[
        x - 1 < \lfloor x \rfloor \le x,\quad -(y - 1) > -\lfloor y \rfloor \ge -y
        \]
        由 $x = (x - y) + y$ 得 
        \[
        \lfloor x \rfloor = \lfloor (x - y) + y \rfloor \ge \lfloor x - y \rfloor + \lfloor y \rfloor
        \]
        移项得 
        \[
        \lfloor x \rfloor - \lfloor y \rfloor \ge \lfloor x - y \rfloor
        \]
        得证左边不等式,又因 $x - y = x + (-y)$ 得 
        \[
        \lfloor x - y \rfloor = \lfloor x + (-y) \rfloor \ge \lfloor x \rfloor + \lfloor -y \rfloor
        \]
        由于 
        \[
        \lfloor -y \rfloor = 
        \begin{cases}
        -\lfloor y \rfloor &,y \in \mathbb{Z} \\ 
        -\lfloor y \rfloor - 1 &,y \not\in \mathbb{Z}
        \end{cases}
        \]
        即 
        \[
        \lfloor -y \rfloor \ge -\lfloor y \rfloor - 1
        \]
        代入得 
        \[
        \lfloor x - y \rfloor \ge \lfloor x \rfloor - \lfloor y \rfloor - 1
        \]
        得证右边不等式,综上原不等式成立。
    \end{solution}

    \question 证明埃尔米特恒等式:对所有实数 $x$ 及所有正整数 $n$ ,以下恒等式成立:
    \[ 
    \sum_{k=0}^{n-1} \left\lfloor x + \frac{k}{n} \right\rfloor = \lfloor nx \rfloor 
    \]
    \begin{solution}
        将 $x$ 分解为整数部分和小数部分,即 
        \[
        x = \lfloor x \rfloor + \{x\}
        \]
        在 $\{1, \dots, n\}$ 中,恰好存在一个 $k'$ 使得:
        \[ 
        \lfloor x \rfloor = \left\lfloor x + \frac{k'-1}{n} \right\rfloor \le x < \left\lfloor x + \frac{k'}{n} \right\rfloor = \lfloor x \rfloor + 1 
        \]
        通过在不等式左右两边的取整符号内减去同一个整数 $\lfloor x \rfloor$ ,可以重写为:
        \[ 
        0 = \left\lfloor \{x\} + \frac{k'-1}{n} \right\rfloor \le \{x\} < \left\lfloor \{x\} + \frac{k'}{n} \right\rfloor = 1  
        \]
        因此有
        \[ 
        1 - \frac{k'}{n} \le \{x\} < 1 - \frac{k'-1}{n}  
        \]
        即求点
        \[ 
        n - k' \le n\{x\} < n - k' + 1 
        \]
        这意味着 $\lfloor n\{x\} \rfloor = n - k'$,于是
        \begin{align*}
        \sum_{k=0}^{n-1} \left\lfloor x + \frac{k}{n} \right\rfloor &= \sum_{k=0}^{k'-1} \lfloor x \rfloor + \sum_{k=k'}^{n-1} (\lfloor x \rfloor + 1) \\
        &= k' \lfloor x \rfloor + (n - k')(\lfloor x \rfloor + 1) \\
        &= n \lfloor x \rfloor + n - k' \\
        &= n \lfloor x \rfloor + \lfloor n\{x\} \rfloor \\
        &= \lfloor n \lfloor x \rfloor + n\{x\} \rfloor \\
        &= \lfloor nx \rfloor
        \end{align*}
        得证。
    \end{solution}
\end{questions}
\pagebreak

\begin{center}
  {\fontsize{30pt}{26pt}\selectfont
    \hypertarget{函数}{函数} \label{函数}
  }
\end{center}
\separator
\vspace{1pt}

\begin{questions}
    \question 求下列函数的值域:
    \begin{parts}
    \part $f(x) = x^2-2x+5$,其中$D_f= [-1,2]$
    \begin{solution}
        配方法得
        \[
        f(x) = x^2-2x+5 = (x-1)^2+4
        \]
        当$x=1,f_{\min}=4$;当$x=-1,f_{\max}=8$,故
        \[
        R_f = [4,8]
        \]
    \end{solution}
    \part $f(x) = x+\sqrt{x(2-x)}$
    \begin{solution}
        设$y=x+\sqrt{x(2-x)}$,则
        \[
        2x^2-2(y+1)x+y^2=0
        \]
        由于$x\in \mathbb{R}$,判别式为非负,
        \[
        4(y+1)^2-8y \ge 0 \Rightarrow 1-\sqrt{2} \le y \le 1+\sqrt{2}
        \]
        但$0 \le x \le 2, y=x+\sqrt{x(2-x)} \ge 0$,故$y_{\min}=0$。 

        而当$y=1+\sqrt{2},x_1=\dfrac{2+\sqrt{2}-2\sqrt[4]{2}}{\sqrt{2}} \in [0,2]$,故$y_{\max}=1+\sqrt{2}$。

        于是
        \[
        R_f = [0,1+\sqrt{2}]
        \]
    \end{solution}
    \part $f(x) = \dfrac{3x+4}{5x+6}$
    \begin{solution}
        设$y=\dfrac{3x+4}{5x+6}$,则$x=\dfrac{4-6y}{5y-3}$,故反函数
        \[
        f^{-1}(x)=\frac{4-6x}{5x-3}
        \]
        的定义域为
        \[
        D_{f^{-1}}=\left(-\infty,\frac{3}{5}\right)\cup\left(\frac{3}{5},\infty\right)
        \]
        即$R_f$。
    \end{solution}
    \part $f(x) = \dfrac{3x+2}{x+1}$
    \begin{solution}  
        发现
        \[
        f(x) = \dfrac{3x+2}{x+1} = 3 - \frac{1}{x+1} \neq 3
        \]
        故
        \[
        R_f=(-\infty,3) \cup (3,\infty)
        \]
    \end{solution}
    \part $f(x) = \dfrac{\cos x}{\sin x-3}$
    \begin{solution}
        设$y=\dfrac{\cos x}{\sin x-3}$,则
        \[
        3y=y \sin x -\cos x=\sqrt{y^2+1}\cos(x-\alpha)
        \Rightarrow \cos(x-\alpha)=\frac{3y}{\sqrt{y^2+1}} 
        \]
        由于$\cos(x-\alpha)\in [-1,1]$,解得
        \[
        -1 \le \frac{3y}{\sqrt{y^2+1}} \le 1 \Rightarrow -\frac{\sqrt{2}}{4} \le y \le \frac{\sqrt{2}}{4}
        \]
        即
        \[
        R_f = \left[-\frac{\sqrt{2}}{4},\frac{\sqrt{2}}{4}\right]
        \]
    \end{solution}
    \part $f(x) = 2^{x-5}+\log_3 \sqrt{x-1}$,其中$D_f= [2,10]$
    \begin{solution}
        由于$2^{x-5}$与$\log_3 \sqrt{x-1}$在$[2,10]$上皆为增函数,故$f(x)$也为增函数。

        当$x=2,f(x)=\dfrac{1}{8}$;当$x=10,f(x)=33$。故
        \[
        R_f=\left[\frac{1}{8},33\right]
        \]
    \end{solution}
    \part $f(x) = \sqrt{x+1}-\sqrt{x-1}$
    \begin{solution}
        发现
        \[
        f(x) = \frac{2}{\sqrt{x+1}+\sqrt{x-1}}
        \]
        由$x\ge 1$,得$\sqrt{x+1}+\sqrt{x-1} \ge 2$,于是
        \[
        0 < \frac{2}{\sqrt{x+1}+\sqrt{x-1}} \le \frac{2}{\sqrt{2}} = \sqrt{2}
        \]
        即
        \[
        R_f=(0,\sqrt{2}]
        \]
    \end{solution}
    \part $f(x) = x+2+\sqrt{1-(x+1)^2}$
    \begin{solution}
        由于$1-(x+1)^2 \ge 0$,有$(x+1)^2 \le 1$,不妨设 $x+1 = \cos \alpha,\alpha \in [0,\pi]$,则
        \[
        f(x) = \cos \alpha + 1 + \sqrt{1-\cos^2 \alpha} = \sin \alpha + \cos \alpha + 1 = \sqrt{2}\cos\left(\alpha-\frac{\pi}{4}\right)+1
        \]
        由$\alpha \in [0,\pi]$,可知
        \[
        -\frac{\sqrt{2}}{2} \le \cos\left(\alpha-\frac{\pi}{4}\right) \le 1 \Rightarrow 0 \le \sqrt{2}\cos\left(\alpha-\frac{\pi}{4}\right) +1\le 1+\sqrt{2}
        \]
        即
        \[
        R_f=[0,1+\sqrt{2}]
        \]
    \end{solution}
    \part $f(x) = \dfrac{x^3-x}{x^4+2x^2+1}$
    \begin{solution}
        首先有
        \[
        f(x)=\frac{1}{2}\cdot \frac{2x}{1+x^2} \cdot \frac{1-x^2}{1+x^2}
        \]
        令$x=\tan \beta$,则
        \[
        f(x)=\frac{1}{2}\sin 2\beta \cos 2\beta=\frac{1}{4}\sin 4\beta
        \]
        若$k \in \mathbb{Z}$,当$\beta=\dfrac{k\pi}{2}-\dfrac{\pi}{8},f_{\max}=\dfrac{1}{4}$;当$\beta=\dfrac{k\pi}{2}+\dfrac{\pi}{8},f_{\min}=-\dfrac{1}{4}$,且此时$\tan \beta$有意义,于是
        \[
        R_f=\left[-\frac{1}{4},\frac{1}{4}\right]
        \]
    \end{solution}
    \part $f(x) = x+4+\sqrt{5-x^2}$
    \begin{solution}
        由$5-x^2 \ge 0$得$|x| \le \sqrt{5}$,令$x=\sqrt{5}\cos\alpha, \alpha \in [0,\pi]$,则
        \[
        f(x)=\sqrt{5}\cos\alpha+4+\sqrt{5}\sin\alpha=4+\sqrt{10}\cos\left(\alpha-\frac{\pi}{4}\right)
        \]
        由于$\alpha \in [0,\pi],\alpha-\dfrac{\pi}{4}\in\left[-\dfrac{\pi}{4},\,\dfrac{3\pi}{4}\right]$,故
        \[
        f_{\max}=4+\sqrt{10}\cdot 1,\quad f_{\min}=4+\sqrt{10}\cdot \left(-\frac{\sqrt2}{2}\right)
        \]
        故
        \[
        R_f=[4-\sqrt{5},4+\sqrt{10}]
        \]
    \end{solution}
    \part $f(x)= \dfrac{\sqrt{x+2}}{x+3}$
    \begin{solution}
        设$t=\sqrt{x+2},t \ge 0$,则$x+3=t^2+1$。当$t>0$,
        \[
        f(x)=\frac{t}{t^2+1} = \frac{1}{t+\frac{1}{t}} \le \frac{1}{2}
        \]
        等号成立当且仅当$t=1$即$x=-1$,于是
        \[
        0 < f(x) \le \frac{1}{2}
        \]
        而当$t=0$时,$f(x)=0$,故
        \[
        R_f=\left[0,\frac{1}{2}\right]
        \]
    \end{solution}
    \part $f(x) = (\sin x + 1)(\cos x + 1),x \in \left[-\dfrac{\pi}{12},\dfrac{\pi}{2}\right]$
    \begin{solution}
        由
        \[
        f(x)=\sin x \cos x +\sin x +\cos x +1
        \]
        令$t=\sin x +\cos x$,则$\sin x \cos x=\dfrac{1}{2}(t^2-1)$,化为
        \[
        f(x)=\frac{1}{2}(t^2-1)+t+1=\frac{1}{2}(t+1)^2
        \]
        又由$t=\sin x +\cos x=\sqrt{2}\cos\left(x-\dfrac{\pi}{4}\right),x \in \left[-\dfrac{\pi}{12},\dfrac{\pi}{2}\right]$知
        \[
        \frac{\sqrt{2}}{2}\le t\le \sqrt{2}
        \]
        当$t=\sqrt{2},f_{\max}=\dfrac{3}{2}+\sqrt{2}$;当$t=\dfrac{\sqrt{2}}{2},f_{\min}=\dfrac{3}{4}+\dfrac{\sqrt{2}}{2}$,于是
        \[
        R_f=\left[\frac{3}{4}+\frac{\sqrt{2}}{2},\frac{3}{2}+\sqrt{2}\right]
        \]
    \end{solution}
    \part $f(x)=\left(\sin x + \dfrac{1}{\sin x}\right)^2+\left(\cos x + \dfrac{1}{\cos x}\right)^2$
    \begin{solution}
        展开得
        \[
        f(x)=\sin^2 x + \cos^2 x + \csc^2 x + \sec^2 x + 4 = 7 + \tan^2 x + \cot^2 x 
        \]
        由AM-GM不等式,
        \[
        f(x) \ge 7 + 2 \sqrt{1} = 9
        \]
        等号成立当且仅当$\tan x =\cot x$即$x=k\pi \pm \dfrac{\pi}{4}, k \in \mathbb{Z}$,故
        \[
        R_f=[9,\infty)
        \]
    \end{solution}
    \part $f(x)=2 \sin x \sin 2x$
    \begin{solution}
        由AM-GM不等式,
        \[
        [f(x)]^2=16\sin^4 x \cos^2 x = 8\sin ^2 x \sin ^2 x(2-2\sin ^2 x) \le 8\cdot \left(2\cdot \frac{1}{3}\right)^3=\frac{64}{27}
        \]
        等号成立当且仅当$\sin ^2 x =2-2\sin ^2 x$即$\sin ^2 x=\dfrac{2}{3}$,故
        \[
        R_f=\left[-\frac{8\sqrt{3}}{9},\frac{8\sqrt{3}}{9}\right]
        \]
    \end{solution}
    \part $f(x)=|x-2|+|x+8|$
    \begin{solution}
        在一维数轴上设$A=-8,B=2$,当动点$P$在线段$AB$上,
        \[
        f(x)=|x-2|+|x+8| = |AB| = 10
        \]
        当动点$P$在线段$AB$外,
        \[
        f(x)=|x-2|+|x+8| > |AB| = 10
        \]
        故
        \[
        R_f = [10,\infty)
        \]
    \end{solution}
    \part $f(x)=\sqrt{x^2-6x+13}+\sqrt{x^2+4x+5}$
    \begin{solution}
        写成
        \[
        f(x)=\sqrt{(x-3)^2+(0-2)^2}+\sqrt{(x+2)^2+(0+1)^2}
        \]
        即$x$轴上的动点$P(x,0)$到两定点$A(3,2),B(-2,-1)$的距离之和,而当$P$为线段与$x$轴的交点时,
        \[
        f_{\min}=|AB|=\sqrt{(3+2)^2+(2+1)^2}=\sqrt{43}
        \]
        于是
        \[
        R_f = [\sqrt{43},\infty)
        \]
    \end{solution}
    \part $f(x)=\sqrt{x^2-6x+13}-\sqrt{x^2+4x+5}$
    \begin{solution}
        写成
        \[
        f(x)=\sqrt{(x-3)^2+(0-2)^2}-\sqrt{(x+2)^2+(0-1)^2}
        \]
        即$AP$距离与$BP$距离之差,其中$A(3,2),B(-2,1),P(x,0)$。当$P$在$x$轴上且不是直线$AB$与$x$轴的交点时,即$P'$,则构成$\triangle ABP'$,有
        \[
        \bigl||AP'|-|BP'|\bigr|<|AB|=\sqrt{(3+2)^2+(2-1)^2}=\sqrt{26}
        \]
        即
        \[
        -\sqrt{26} < f(x) < \sqrt{26}
        \]
        而当$P$恰好在直线$AB$与$x$轴的交点时,有
        \[
        \bigl||AP'|-|BP'|\bigr|=|AB|=\sqrt{26}
        \]
        故
        \[
        R_f=(-\sqrt{26},\sqrt{26}]
        \]
    \end{solution}
    \part $f(x) = x^2+\sqrt{x^4-3x^2+2x+5}$
    \begin{solution}
        变形可得
        \[
        f(x)=x^2+\sqrt{(x^2-2)^2+(x+1)^2}
        \]
        发现$f(x)$表示抛物线$y=x^2$上动点$P(x,x^2)$到点$A(-1,2)$和$x$轴的距离之和。
        
        过$P$点作$PB \perp x$轴于$B$, 过$A$点作$AC \perp x$轴于$C$,
        $BC$交抛物线$y=x^2$于点$P_0$,
        故
        \[
        |PA|+|PB| \ge |P_0A|+|P_0C| = |AC| = 2
        \]
        所以
        \[
        R_f = [2,\infty)
        \]
    \end{solution}
    \end{parts}

    \question 设函数 $f:\mathbb{R}\setminus\{0,-b\}\to \mathbb{R}\setminus\{1\}$满足
    \[
    f(x) = \frac{x+a}{x+b}
    \]
    求出所有实数对 $(a,b)$ 使得
    \[
    f(f(x)) = -\frac{1}{x}
    \]
    \begin{solution}
        据题意,
        \[
        f(f(x)) = \frac{\frac{x+a}{x+b} + a}{\frac{x+a}{x+b} + b}
        = \frac{(1+a)x + a + ab}{(1+b)x + a + b^2} = -\frac{1}{x}
        \]
        整理得
        \[
        (1+a)x^2 + (a(1+b) + 1 + b)x + a + b^2 = 0.
        \]
        解得
        \[
        \begin{cases}
        1 + a = 0 \\
        a + b^2 = 0 
        \end{cases}
        \Rightarrow a = -1,b = \pm 1
        \]
        经检验,$(-1, -1)$ 不合题意,故解为$(-1, 1)$
    \end{solution}

    \question 已知函数 $f,g:\mathbb{R} \to \mathbb{R}$ 互为反函数,且函数
    \[
    F(x) = f(x+1)-2, \quad G(x) = g(2x+1),
    \]
    也互为反函数,若 $f(1)=4$,求 $f(100)$。
    \begin{solution}
        由于 $F,G$ 互为反函数,有
        \[
        G(F(x))=g(2F(x)+1)=g(2(f(x+1)-2)+1)=x
        \]
        又$f,g$ 互为反函数,则
        \[
        2f(x+1)-3=f(x) \Rightarrow f(x+1) = \frac{1}{2}f(x) + \frac{3}{2}
        \]
        由 $f(1)=4$,可依次得到
        \[
        f(2)=3+\frac{1}{2}, \quad f(3)=3+\frac{1}{4}, \dots
        \]
        归纳可得
        \[
        f(n) = 3 + \frac{1}{2^{\,n-1}}, \ n \in \mathbb{N}
        \]
        因此
        \[
        f(100) = 3 + \frac{1}{2^{99}}
        \]
    \end{solution}

    \question 已知函数 $g(x)=2x-4$, 其反函数为 $g^{-1}$,且函数 $f$ 对任意实数 $x$ 皆满足
    \[
    g(f(g^{-1}(x))) = 2x^2 + 16x + 26,
    \]
    求 $f(\pi)$ 的值。
    \begin{solution}
        已知 $g$ 可逆,令$x=g(y)$,则$y=g^{-1}(x)$,由题意得
        \[
        g(f(y)) = 2(g(y))^2 + 16g(y) + 26 
        = 2(2y-4)^2 + 16(2y-4) + 26 = 8y^2 - 6
        \]
        于是有
        \[
        f(y) = g^{-1}(g(f(y)))= g^{-1}(8y^2 - 6) = \frac{8y^2 -6 +4}{2} = 4y^2 - 1
        \]
        所以
        \[
        f(\pi) = 4\pi^2 - 1
        \]
    \end{solution}
\begin{solution}
    由$y=g(x)=2x-4$,得
    \[
    x=g^{-1}(y) = \frac{y+4}{2}
    \]
    由$g(f(g(x))) = 2x^2 + 16x + 26$,
    \[
    f(g(x)) = g^{-1}(2x^2 + 16x + 26) = x^2 + 8x + 15 = (x+4)^2 - 1
    \]
    且
    \[
    f(g(x)) = f\left(\textcolor{red}{\frac{x+4}{2}}\right) = (x+4)^2 - 1
    \]
    令 $\dfrac{x+4}{2} = \pi$, 则
    \[
    f(\pi) = 4\pi^2 - 1.
    \]
    \textcolor{red}{请问是哪一行有误?}
\end{solution}

    % \question 求 $k$ 的求值范围,使得函数 $f(x) = x^3 + 2x^2 + kx - 1$ 反函数存在。
    % \begin{solution}
    %     对$f(x)$ 求导得
    %     $$ f'(x) = 3x^2 + 4x + k = 3\left(x + \frac{2}{3}\right)^2 + k - \frac{4}{3} $$
    %     只有当$k> \dfrac{4}{3}$时, $f(x)$ 单调递增,故$f(x)$在$\mathbb{R}$上一对一,反函数$f^{-1}(x)$ 存在。
        
    %     现考虑$k=\dfrac{4}{3}$,因$f'(x)=0$有唯一解$x=-\dfrac23$,
    %     且$f$在$(-\infty,-\dfrac23)$及$(-\dfrac23,\infty)$上严格递增,故$x=-\dfrac23$是一驻点。由图像法可知,$f(x)$在$\mathbb{R}$上一对一,此时反函数$f^{-1}(x)$ 存在。\textcolor{red}{(待验证)}
    
    %     综上,当$k\ge \dfrac{4}{3}$时,反函数存在。
    % \end{solution}


    \question 已知$f(x) + 2 f(4-x) = x + 8,$求 $f(16)$。 
    \begin{solution}
        令$x=-12,16$,可得联立方程
        \[
        \begin{cases}
            f(-12)+2(16)=-4\\f(16)+2f(-12)=24
        \end{cases}
        \]
        解得 $$f(16)= -\dfrac{32}{3}$$
    \end{solution}
        
    \question 若 $f(x)$ 为实系数二次多项式,已知 $p,q,r$ 为三相异非零实数,使得 $f(p)=qr,f(q)=rp,f(r)=pq$,证明$f(p+q+r) = f(p)+f(q)+f(r)$。
    \begin{solution}
        取 $g(x) = x f(x) - pqr$,则 $g(p) = g(q) = g(r) = 0$,且有
        \[
        g(x) = x f(x) - pqr = a(x - p)(x - q)(x - r)
        \]
        于是
        \[
        f(x) = \frac{1}{x} \left( a x^3 - a(p + q + r)x^2 + a(pq + qr + rp)x - apqr + pqr \right)
        \]
        由于 $f(x)$ 为二次多项式,因此 $-apqr + pqr = 0 \Rightarrow a = 1$,得
        \[
        f(x) = x^2 - (p + q + r)x + pq + qr + rp
        \]
        故得证
        \[
        f(p + q + r) = pq + qr + rp = f(p)+f(q)+f(r)
        \]
    \end{solution}
    
    \question 设 $Q(x)$ 是一个 $2017$ 次多项式,且满足
    \[
    Q'(r)=\frac{2017!}{r} \quad ,r=1,2,3,\cdots, 2017
    \]
    另定义
    \[
    P(x)=xQ(x)-\int Q(x)\,dx.
    \]
    若多项式 $P(x)$ 的所有根之和为 $a$,求 $a \pmod{1000}$。
    \begin{solution}
        有
        \[
        P'(x)=xQ'(x) +Q(x)-Q(x) \Rightarrow P(x) = \int xQ'(x)\,dx
        \]
        考虑函数
        \[
        R(x)=xQ(x)-2017!
        \]
        则$R(r)=0,r=1,2,3,\cdots, 2017$,故可设
        \[
        R(x)=A(x-1)(x-2)(x-3)\cdots(x-2017)
        \]
        令$x=0$可得$A=1$,故
        \[
        Q'(x)=\frac{(x-1)(x-2)(x-3)\cdots(x-2017)+2017!}{x}
        \]
        且
        \begin{align*}
        P(x) = \int xQ'(x)\,dx 
        &= \int ((x-1)(x-2)(x-3)\cdots(x-2017)+2017!)\,dx \\
        &=\int \left(x^{2017}-\frac{2017\cdot2018}{2}x^{2016} +\cdots\right)\,dx \\
        &=\frac{1}{2018}x^{2018} - 1009 x^{2017}+\cdots
        \end{align*}
        由韦达定理,$\ P(x)$ 的所有根之和为 
        \[
        a=1009 \cdot 2018 \equiv 162 \pmod{1000}
        \]
    \end{solution}

    \question 定义函数 $f:\mathbb{R} \to \mathbb{R}$,使得对所有实数 $x$ 有
    \[
    f(f(x)) = x^2 - x + 1.
    \]
    求 $f(0)$。
    \begin{solution}
        设 $f(0) = b$,则
        \[
        f(b) = f(f(0)) = 0^2 - 0 + 1 = 1
        \]
        又有
        \[
        f(f(b)) = f(1) = b^2 - b + 1
        \]
        因为 $f(b) = 1$,所以
        \[
        f(f(b)) = f(1) = 1 = b^2 - b + 1 \implies b(b-1) = 0.
        \]
        若 $b=0$,则 $f(0)=0$,但 $f(f(0)) = f(0) = 0 \neq 1$,矛盾。故 $b = 1$,即
        \[
        f(0) = 1
        \]
    \end{solution}

    \question 已知 $f(x) = e(x) + o(x)$,其中 $e(x)$ 为偶函数,$o(x)$ 为奇函数,且
    \[
    e(x) + x^2 = o(x),
    \]
    求 $f(2)$。
    \begin{solution}
        由已知$e(x)=f(x)-o(x)$,换元得
        \[
        e(-x)=f(-x)-o(-x) \Rightarrow e(x)=f(-x)+o(x)
        \]
        故
        \[
        f(-x)=-x^2 \Rightarrow f(2)=4
        \]
    \end{solution}

    \question 定义 $f$ 在 $[0,1]$ 上满足
    \[
    2f\left(\frac{x}{3}\right) = f(x), \quad f(x) + f(1-x) = 1,
    \]
    求 $f\left(\dfrac{1}{3}\right)$。
    \begin{solution}
        由递推公式,则
        \[
        f\left(\frac{1}{13}\right) + f\left(\frac{12}{13}\right) = 1,
        \]
        且
        \[
        f\left(\frac{12}{13}\right) 
        = 2 f\left(\frac{4}{13}\right) 
        = 2 \left[1-f\left(\frac{9}{13}\right)\right]
        = 2 - 2 f\left(\frac{9}{13}\right)
        = 2 - 4 f\left(\frac{3}{13}\right)
        = 2 - 8 f\left(\frac{1}{13}\right)
        \]
        解得
        \[
        f\left(\frac{1}{13}\right) = \frac{1}{7}.
        \]
    \end{solution}

    \question 已知$f(x)$是定义在$\mathbb{R}$上的函数,且对任意实数$x$均有$2f(x)+f(x^2-1)=1$,试求$f(\sqrt{2})$的值。
    \begin{solution}
        分别令 $x=-1,0,1,\sqrt{2}$ 可得
        \begin{align*}
            2f(-1) + f(0) = 1 \tag{1}\\
            2f(0) + f(-1) = 1 \tag{2}\\
            2f(1) + f(0) = 1 \tag{3}\\
            2f(\sqrt{2}) + f(1) = 1 \tag{4}
        \end{align*}
        由 $(1),(2),(3)$ 解得
        \[
        f(-1) = f(0)= f(1)= \frac{1}{3}
        \]
        故由(4)得
        \[
        f(\sqrt{2}) = \frac{1}{3}
        \]
        \end{solution}

    \question 设函数 $f(x)=\cos x+\log_{2}x,$其中$x>0$ ,若正实数 $a$ 满足 $f(a)=f(2a)$ ,求$f(2a)-f(4a)$ 的值。
    \begin{solution}
        由条件得 $$\cos a+\log_{2}a=\cos 2a+\log_{2}2a=2\cos^{2}a-1+\log_{2}a+1$$
        所以
        $$\cos a(2\cos a-1) =0\Rightarrow (\cos a,\cos 2a) = (0,-1) \text{或}\; (\frac{1}{2},-\frac{1}{2})$$
        于是
        \begin{align*}
        f(2a) - f(4a) 
        &= \cos 2a + \log_2 2a - \cos 4a - \log_2 4a \\
        &= \cos 2a - 2 \cos^2 2a\\
        &= \begin{cases}
            -3, & \text{若 } \cos 2a = -1, \\
            -1, & \text{若 } \cos 2a = -\frac{1}{2}
        \end{cases}
        \end{align*}
    \end{solution}

    \question 已知函数 $f(x)$ 的定义域为 $\mathbb{R}$, 且 $f(x+2)-2$ 为奇函数, $f(2x+1)$ 为偶函数。 若 $f(1)=0$, 求 $f(1)+f(2)+\cdots+f(2023)$ 的值。
    \begin{solution}
        $f(x+2)-2$ 为奇函数, 则
        \[
        f(-x+2)-2+f(x+2)-2=0 \Rightarrow f(2-x)+f(2+x)=4
        \]
        于是 $f(2)=2$, $f(x)$ 的图象关于点 $(2,2)$ 对称,又由 $f(2x+1)$ 为偶函数, 
        $$f(-2x+1)=f(2x+1)\Rightarrow f(1-2x)=f(1+2x)$$ 
        即
        $$f(1-x)=f(1+x)$$
        $f(x)$ 图象关于直线 $x=1$ 对称,
        由上可得 
        $$f(x)=f(2-x)=4-f(2+x)=4-f(-x)=4-[4-f(x+4)]=f(x+4)$$
        所以 $f(x)$ 是周期为4的函数,且
        \[
        f(1)=0,f(2)=2,f(3)=4,f(4)=2
        \]
        故
        \begin{align*}   
        &f(1)+f(2)+\cdots+f(2023)\\
        &=505[f(1)+f(2)+f(3)+f(4)]+f(1)+f(2)+f(3)\\
        &=4046 
        \end{align*}
    \end{solution}

    \question 设 $f$ 为实值函数,且对任意实数 $x$ 满足
    \[
    f(10+x) = f(10-x), \quad f(20+x) = -f(20-x)
    \]
    \begin{parts}
    \part 证明 $f$ 是奇函数。
    \begin{solution}
        \[
        f(10+x) = f(10-x) \tag{1}
        \]
        \[
        f(20+x) = -f(20-x) \tag{2}
        \]
        由(1),令$x = y - 10$,
        \[
        f(y) = f(20 - y)
        \]
        由(2),令 $x = - y$,
        \[
        f(20 - y) = -f(20 + y)
        \]
        因此
        \[
        f(y) = -f(20 + y)
        \]
        由(1),令 $x = y + 10$,
        \[
        f(y + 20) = f(-y)
        \]
        结合上式可得
        \[
        f(y) = -f(20+y) = f(-y),
        \]
        所以得证 $f$ 是奇函数。
    \end{solution}
    \part 证明 $f$ 是周期函数。
    \begin{solution}
        由(2),令 $x = y - 20$ 在 (b) 中,则
        \[
        f(y) = -f(40 - y)
        \]
        由于 $f$ 是奇函数,有$-f(40-y) = f(y-40)$,所以
        \[
        f(y) = f(y-40),
        \]
        因此得证 $f$ 是周期函数。
    \end{solution}
    \end{parts}

    \question 函数 $f(x)$ 在 $\mathbb{R}$ 上定义,且满足
    \[
    f(-x)=f(x), f(0)=993,
    \]
    若另一函数 $g(x)$满足
    \[
    g(-x)=-g(x),g(x)=f(x-1)
    \]
    求 $f(1992)$的值。
    \begin{solution}
        由 $g(-x)=-g(x)$ 及 $g(x)=f(x-1)$ 得 
        \[
        f(-x-1) = -f(x-1)
        \]
        结合 $f(x)$ 为偶函数,有
        \[
        f(x+1) = f(-x-1) = -f(x-1)
        \]
        从而 $f(x+2) = -f(x)$,进而 $f(x+4) = f(x)$,发现$f(x)$是周期为4的函数,故
        \[
        f(1992) = f(4 \cdot 498) = f(0) = 993
        \]
    \end{solution}

    \question 若实数 \(x, y, z\) 满足
    \[
    \begin{cases}
    \displaystyle \frac{x}{1^2 + 4^2} + \frac{y}{1^2 + 5^2} + \frac{z}{1^2 + 6^2} = 1 \\[2mm]
    \displaystyle \frac{x}{2^2 + 4^2} + \frac{y}{2^2 + 5^2} + \frac{z}{2^2 + 6^2} = 1 \\[2mm]
    \displaystyle\frac{x}{3^2 + 4^2} + \frac{y}{3^2 + 5^2} + \frac{z}{3^2 + 6^2} = 1
    \end{cases}
    \]
    求 \(x + y + z\)。
    \begin{solution}
        设
        \[
        f(k) = \frac{x}{k + 4^2} + \frac{y}{k + 5^2} + \frac{z}{k + 6^2} - 1,
        \]
        据题意
        \[
        f(1^2) = f(1) = 0, \quad f(2^2) = f(4) = 0, \quad f(3^2) = f(9) = 0。
        \]
        因此 \(f(k)\) 有三个根\(k = 1, 4, 9\),所以
        \[
        (k - 1)(k - 4)(k - 9)
        \]
        与
        \[
        (k + 4^2)(k + 5^2)(k + 6^2) - x(k + 5^2)(k + 6^2) - y(k + 4^2)(k + 6^2) - z(k + 4^2)(k + 5^2)
        \]
        相等,比较 \(k^2\) 系数得
        \[
        - (1 + 4 + 9) = 4^2 + 5^2 + 6^2 - (x + y + z) \Rightarrow x + y + z = 91
        \]
    \end{solution}

    \question 已知实数 $\alpha,\beta$ 满足
    \[
    3^{\frac{\alpha}{2}} = \sqrt{3} - \alpha, \quad \log_3 \beta = 2\sqrt{3} - 2\beta
    \]
    求
    \[
    (\alpha + \beta)^2 - 2 (\sqrt{3})^\alpha - \log_3 \beta
    \]
    \begin{solution}
        将底改为 $\sqrt{3}$:
        \[
        \sqrt{3}^\alpha = \sqrt{3} - \alpha, \quad \log_{\sqrt{3}} \beta = \sqrt{3} - \beta
        \]
        设 $A(\alpha, \sqrt{3}-\alpha)$ 为两图形 $y = \sqrt{3}^x$ 与 $y = \sqrt{3} - x$ 的交点$,B(\beta, \sqrt{3}-\beta)$ 为两图形 $y = \log_{\sqrt{3}} x$ 与 $y = \sqrt{3}-x$ 的交点。  

        由于 $y = \sqrt{3}^x$ 与 $y = \log_{\sqrt{3}} x$ 关于 $y = x$ 对称,
        \[
        \alpha = \sqrt{3} - \beta \Rightarrow \alpha+\beta = \sqrt{3}
        \]
        因此
        \[
        (\alpha + \beta)^2 - 2 (\sqrt{3})^\alpha - \log_3 \beta = (\sqrt{3})^2 - 2(\sqrt{3}-\alpha) - 2(\sqrt{3}-\beta) = 3 - 2\sqrt{3}
        \]
    \end{solution}

    \question 设实数 $\alpha,\beta$ 满足
    \[
    \begin{cases}
    \alpha^3 - 6\alpha^2 + 13\alpha = 6 \\
    \beta^3 - 6\beta^2 + 13\beta = 14
    \end{cases}
    \]
    求 $\alpha + \beta$。

    \begin{solution}
        令 $$f(x) = x^3 - 6x^2 + 13x \Rightarrow f(\alpha) = 6,\ f(\beta) = 14$$

        又 $$f'(x) = 3x^2 - 12x + 13 \Rightarrow f''(x) = 6x - 12 = 0 \Rightarrow x = 2$$
        因此图形 $y = f(x)$ 的对称中心为 $(2, f(2) = 10)$,所以 $P(\alpha, 6), Q(\beta, 14)$的对称点为
        \[
        P' = (4 - \alpha, 14), Q' = (4 - \beta, 6)
        \]
        于是由$P = Q' ,\quad Q = P'$解得
        \[
        \alpha = 4 - \beta,\quad \beta = 4 - \alpha \Rightarrow \alpha + \beta = 4
        \]
    \end{solution}

    \question 解方程
    \[
    e^{x-1} \ln x + e^{y-1} \ln y = \ln(xy)
    \]
    \begin{solution}
        原式化为:
        \[
        (e^{x-1} - 1)\ln x + (e^{y-1} - 1)\ln y = 0
        \]
        考虑函数 \(f(t) = (e^{t-1} - 1)\ln t\),
        \begin{itemize}
        \item 当 \(t > 1\) 时,\(\;e^{t-1} - 1 > 0\),\(\;\ln t > 0\),故 \(f(t) > 0\)
        \item 当 \(0 < t < 1\) 时,\(\;e^{t-1} - 1 < 0\),\(\;\ln t < 0\),故 \(f(t) > 0\)
        \item 当 \(t = 1\) 时,\(\;e^{t-1} - 1 = 0\),\(\;\ln 1 = 0\),故 \(f(1) = 0\)
        \end{itemize}
        故 $$f(x) + f(y) = 0 \Rightarrow x = y = 1$$
    \end{solution}

    \question 已知$x,y \in \mathbb{R}$, 且满足
    \[
    \begin{cases}
    \;(x+1)^{3}+2023(x+1)=-2023 \\
    \;(y+1)^{3}+2023(y+1)=2023
    \end{cases}
    \]
    求$x+y$的值。
    \begin{solution}
    考虑函数 \( f(t) = t^3 + 2023t \),据题意有
        \[
        f(x+1) = -2023, \quad f(y+1) = 2023
        \]
        注意到 \( f(t) \) 是奇函数,
        \[
        f(x+1) = -f(y+1)= f(-(y+1))
        \]
        且导数为
        \[
        f'(t) = 3t^2 + 2023 > 0, \quad \forall\, t \in \mathbb{R}
        \]
        所以$f$在$\mathbb{R}$上严格递增,故$f$一对一,因此有
        \[
        x+1 = -(y+1) \Rightarrow x + y = -2
        \]
    \end{solution}

    \question 求定义于 $\mathbb{R} \setminus \{0\}$ 上的所有实值函数满足 
    \[
    \frac{1}{x} f(-x) + f\left( \frac{1}{x} \right) = x
    \]
    \begin{solution}
        令 $x = \dfrac{1}{y}$,
        \[
        y f\left( -\frac{1}{y} \right) + f(y) = \frac{1}{y} \tag{1}
        \]
        令 $x = -y$ 代入,得 
        \[
        -\frac{1}{y} f(y) + f\left( -\frac{1}{y} \right) = -y \tag{2}
        \]
        联立(1),(2),解得
        \[
        f(x) = \frac{x^3 + 1}{2x}, \quad x \neq 0
        \]
    \end{solution}

    \question 解函数方程 
    \[
    xf(y) + yf(x) = (x+y)f(x)f(y)
    \]
    \begin{solution}
        令 $x = y = 1$,得 $2f(1) = 2f^2(1)$,因此 $f(1) = 0$ 或 $1$。若 $f(1) = 0$,令 $y = 1$,得 
        \[
        xf(1) + 1 \cdot f(x) = (x+1)f(x)f(1) \Rightarrow f(x) = 0
        \]
        若 $f(1) = 1$,令 $y = 1$,得 
        \[
        x + f(x) = (x+1)f(x) \Rightarrow x[f(x) - 1] = 0
        \]
        故解为
        \[
        f(x) = 0 \quad \text{或} \quad f(x) = 
        \begin{cases} 
        1 & ,x \neq 0 \\ a & ,x = 0, a \in \mathbb{R}
        \end{cases}
        \]
    \end{solution}

    \question 已知函数$f(x)$满足 
    \[
    f(m+n) = f(m) + f(n) + mn, \quad \forall m,n \in \mathbb{N}
    \]
    且$f(1) = 1$,求 $f(x)$。
    \begin{solution}
        令 $m = 1$,则 
        \[
        f(n+1) = f(n) + f(1) + n = f(n) + (n+1)
        \]
        当 $n=1, \dots, n-1$ 时,由累加法,
        \[
        f(n) = f(1) + 2 + 3 + \dots + n = \frac{n(n+1)}{2}
        \]
        即 
        \[
        f(x) = \frac{x(x+1)}{2}, \quad x \in \mathbb{N}
        \]
    \end{solution}

    \question 已知函数 $f(x)$ 满足$f(0) = 1, f\left(\dfrac{\pi}{2}\right) = 2$,且对任意的 $x, y \in \mathbb{R}$,
    \[
    f(x+y) + f(x-y) = 2f(x)\cos y,
    \]
    求 $f(x)$。
    \begin{solution}
        令 $x = 0, y = t$,
        \[
        f(t) + f(-t) = 2f(0)\cos t = 2\cos t \tag{1}
        \]
        令 $x = \dfrac{\pi}{2} + t, y = \dfrac{\pi}{2}$,
        \[
        f(t+\pi) + f(t) = 2f\left(\frac{\pi}{2}\right)\cos\frac{\pi}{2} = 0 \Rightarrow f(t+\pi) = -f(t) \tag{2}
        \]
        再令 $x = \dfrac{\pi}{2}, y = \dfrac{\pi}{2} + t$,
        \[
        f(\pi+t) + f(-t) = 2f\left(\frac{\pi}{2}\right)\cos\left(\frac{\pi}{2}+t\right) = -4\sin t \tag{3}
        \]
        将$(2)$代入(3),
        \[
        -f(t) + f(-t) = -4\sin t \tag{4}
        \]
        (1) + (4),给出
        \[
        f(-t) = \cos t - 2\sin t
        \]
        于是
        \[
        f(x) = \cos(-x) - 2\sin(-x) = \cos x + 2\sin x
        \]
    \end{solution}

    \question 已知 $f(x)$ 为多项式,解 
    \[
    f(x+1) + f(x-1) = 2x^2 - 4x
    \]
    \begin{solution}
        若$f(x)$是一$n$次多项式,则不难看出$f(x+1) + f(x-1)$也为$n$次多项式,据此$f(x)$必为二次多项式,设 
        \[
        f(x) = ax^2 + bx + c
        \]
        代入原式,整理得
        \[
        2ax^2 + 2bx + 2a + 2c = 2x^2 - 4x
        \]
        比较系数,解得
        \[
        a = 1, b = -2, c = -1
        \]
        故 
        \[
        f(x) = x^2 - 2x - 1
        \]
    \end{solution}

    \question 设函数$f(x)$满足 $f(1)=1, f(2)=8$ 且 
    \[
    f(n) = \sqrt{f(n-1)f(n-2)}, \quad n \ge 3
    \]
    求 $f(n)$。
    \begin{solution}
        令 $g(n) = \log_2 f(n)$,对方程两边取以 2 为底的对数,则 
        \[
        g(n) = \frac{1}{2} [g(n-1) + g(n-2)],
        \]
        特征方程为 
        \[
        2r^2 - r - 1 = 0
        \]
        解得 $r_1 = 1, r_2 = -\dfrac{1}{2}$,故 
        \[
        g(n) = c_1 + c_2\left(-\frac{1}{2}\right)^{n-1}
        \]
        代入 $g(1) = \log_2 1 = 0,g(2) = \log_2 8=3$ 可解得 $c_1 = 2, c_2 = -2$,于是
        \[
        g(n) = 2 - 2\left(-\frac{1}{2}\right)^{n-1} = 2 + 4\left(-\frac{1}{2}\right)^n
        \]
        最后由 $f(n) = 2^{g(n)}$ 可得
        \[
        f(n) = 2^{2 + (-2)^{2-n}}
        \]
    \end{solution}

    \question 设 $ab \neq 0, a^2 \neq b^2$,求方程 
    \[
    af(x) + bf\left(\frac{1}{x}\right) = cx
    \]
    的解。
    \begin{solution}
        已知 
        \[
        af(x) + bf\left(\frac{1}{x}\right) = cx \tag{1}
        \]
        以 $\dfrac{1}{x}$ 代替 $x$ 得
        \[
        af\left(\frac{1}{x}\right) + bf(x) = \frac{c}{x} \tag{2}
        \]
        联立(1),(2)解得
        \[
        f(x) = \frac{c}{a^2 - b^2} \left(ax - \frac{b}{x}\right) ,\quad x \neq 0
        \]
        其中$ab \neq 0, a^2 \neq b^2$。
    \end{solution}

    \question 求函数 $f(x)$ 满足 
    \[
    af(x^n) + f(-x^n) = bx,
    \]
    其中 $a^2 \neq 1, n$ 是奇数。
    \begin{solution}
        已知 
        \[
        af(x^n) + f(-x^n) = bx \tag{1}
        \]
        以 $-x$ 代替 $x$,由于 $n$ 为奇数,得
        \[
        af(-x^n) + f(x^n) = -bx \tag{2}
        \]
        联立(1),(2)解得
        \[
        f(x^n) = \frac{bx(a+1)}{(a+1)(a-1)} = \frac{bx}{a-1}
        \]
        即 
        \[
        f(x) = \frac{bx^{\frac{1}{n}}}{a-1}
        \]
    \end{solution}

    \question 解函数方程 
    \[
    2f(x^2) + f\left(\frac{1}{x^2}\right) = x
    \]
    其中$x > 0$。
    \begin{solution}
        已知 
        \[
        2f(x^2) + f\left(\frac{1}{x^2}\right) = x \tag{1}
        \]
        以 $\frac{1}{x}$ 代替 $x$ 得
        \[
        2f\left(\frac{1}{x^2}\right) + f(x^2) = \frac{1}{x} \tag{2}
        \]
        联立(1),(2)得
        \[
        3f(x^2) = 2x - \frac{1}{x}
        \]
        令 $t = x^2 \Rightarrow x = \sqrt{t}$,则 
        \[
        f(t) = \frac{1}{3}\left(2\sqrt{t} - \frac{1}{\sqrt{t}}\right)
        \]
        即 
        \[
        f(x) = \frac{2x-1}{3\sqrt{x}}, \quad x > 0
        \]
    \end{solution}

    \question 已知 $f(1) = \dfrac{1}{5}$,且当 $n > 1$ 时,
    \[
    \frac{f(n-1)}{f(n)} = \frac{2nf(n-1)+1}{1-2f(n)},
    \]
    求 $f(n)$。
    \begin{solution}
        原方程式即 
        \[
        \frac{1}{f(n)} - \frac{1}{f(n-1)} = 2(n+1)
        \]
        利用累加法,
        \[
        \frac{1}{f(n)} = \frac{1}{f(1)} + \sum_{i=2}^{n} 2(i+1) = 5 + 2\left[\frac{(n+1)(n+2)}{2}-3\right]
        \]
        整理得 
        \[
        f(n) = \frac{1}{n^2 + 3n + 1},\quad n > 1
        \]
    \end{solution}

    \question 解函数方程 
    \[
    [f(x+y)]^2 = [f(x)]^2 + [f(y)]^2
    \]
    \begin{solution}
        令 $x = y = 0$,则有 
        \[
        [f(0)]^2 = 2[f(0)]^2 \Rightarrow f(0) = 0
        \]
        令 $y = -x$,得到 
        \[
        [f(x)]^2 + [f(-x)]^2 = [f(0)]^2 = 0
        \]
        由于平方数在$\mathbb{R}$是非负的,解为
        \[
        f(x) = 0,\quad  \forall x \in \mathbb{R}
        \]
    \end{solution}

    \question 设函数 $f:(0,\infty)\rightarrow\mathbb{R}$,满足
    \[
    f\left(1-\frac{1}{1+t}\right) + f\left(\frac{1+t}{t}\right)\log(1+t) = f\left(\frac{1+t}{t}\right)\log t + 2022,
    \]
    求 $f(1000)$。
    \begin{solution}
        令 
        \[x = \frac{t}{1+t} = 1 - \frac{1}{1+t},
        \]
        则原式变为
        \[
        f(x) + f\left(\frac{1}{x}\right) \log(1+t) = f\left(\frac{1}{x}\right) \log t + 2022
        \]
        即
        \[
        f(x) - f\left(\frac{1}{x}\right) \log x = 2022 \tag{1}
        \]
        再换元得
        \[
        f\left(\frac{1}{x}\right) - f(x)\log \frac{1}{x} = 2022 \Rightarrow f\left(\frac{1}{x}\right) + f(x) \log x = 2022
        \]
        代入 (1)得
        \[
        f(x) - 2022 \log x + f(x)(\log x)^2 = 2022 \Rightarrow f(x) = \frac{2022 (1 + \log x)}{1 + (\log x)^2}
        \]
        故
        \[
        f(1000) = \frac{2022 \cdot 4}{1 + 9} = 808.8
        \]
    \end{solution}

    \question 设函数 \( f:\mathbb{R}\setminus\{0,1\} \to \mathbb{R} \) 使得
        \[
        f(x) + f\left(\frac{1}{1-x}\right) = \frac{1}{x(1-x)}.
        \]
        求 $f(x)$。
    \begin{solution}
        代入 \( x \), \( \dfrac{1}{1-x} \), 及\( \dfrac{x-1}{x} \) 得
        \[
        f(x) + f\left(\frac{1}{1-x}\right) = \frac{1}{x(1-x)} \tag{1} 
        \]
        \[
        f\left(\frac{1}{1-x}\right) + f\left(\frac{x-1}{x}\right) = -\frac{(x-1)^2}{x}\tag{2}
        \]
        \[
        f\left(\frac{x-1}{x}\right) + f(x) = \frac{x^2}{x-1} \tag{3}
        \]
        $(1)+(3)-(2)$,得
        \begin{align*}
        2f(x) &= \frac{1}{x(1-x)}+\frac{x^2}{x-1}+\frac{(x-1)^2}{x}\\
        &=\frac{-1+x^3+(x-1)^3}{x(x-1)}\\
        &=\frac{(x^2+x+1)+(x^2-2x+1)}{x}\\
        &=\frac{2x^2-x+2}{x}
        \end{align*}
        故
        \[
        f(x)=x+\frac1x-\frac12
        \]
    \end{solution}

    \question 已知函数 $f(x)$ 定义在 $\mathbb{R}$ 上, 且满足 $f(1)=1$, 对所有实数 $x$ 有
    \[
    f(x+5) \geq x+5, \quad f(x+1) \leq f(x)+1.
    \]
    若定义函数 $g(x) = f(x)+1 - x$, 求 $g(2002)$。     
    %https://artofproblemsolving.com/community/c1076938h2031392p14321295
    \begin{solution}
        $\forall x \in \mathbb{R}$,由 $f(x+1) \leq f(x)+1$,
        \[
        g(x+1)=f(x+1)+1-(x+1)\leq f(x)+1-x= g(x) 
        \]
        即$g(x)$是递减函数,且由$f(x+5)\geq x+5$,
        \[
        g(x+5)=f(x+5)+1 - (x+5) \geq x+5+1 - (x+5)=1 
        \]
        因为 $g(1)=f(1)+1-1=1$,故$g(x)=1\;\forall x \in \mathbb{R}$,即$g(2002)=1$
    \end{solution}

    \question 设$f:R\setminus\{0\}\to \mathbb{R}$,使得对于所有非零实数$x, y$有
    \[
    xy(f(x) - f(y)) + 2x = xf(x) - yf(y) + 2y
    \]
    若$f'(1) = 2018$,求$f(2018)$。
    \begin{solution}
        令$y=1$,有
        \[
        xf(x) - xf(1) + 2x= xf(x) - f(1) + 2 \Rightarrow f(1)=2
        \] 
        令$y=-1$,有
        \[
        -x(f(x) - f(-1)) + 2x = xf(x) + f(-1) -2 \Rightarrow f(x) = \frac{f(-1) + 2}{2} + \frac{2 - f(-1)}{2x}
        \]
        求导得
        \[
        f'(x) = \frac{f(-1) - 2}{2x^2} 
        \]
        故由
        \[
        f'(1) = \frac{f(-1) - 2}{2} = 2018
        \]
        得$f(-1)=4038$,于是
        \[
        f(x)=2020-\frac{2018}{x}
        \]
        经检验,原方程式左式=右式$=2020x-2018y$,故
        \[
        f(2018)=2019
        \]
    \end{solution}
    \begin{solution}
        将原式写成
        \[
        xf(x) (y-1) + 2x = yf(y) (x-1) + 2y 
        \]
        \[
        xf(x)(y-1) + 2(x-1) = yf(y)(x-1) + 2(y-1)
        \]
        \[
        \frac{xf(x)}{x-1} + \frac{2}{y-1} = \frac{yf(y)}{y-1} + \frac{2}{x-1}
        \]
        \[
        \frac{xf(x)-2}{x-1} = \frac{yf(y)-2}{y-1} \quad \forall x, y \neq 0, 1
        \]
        设
        \[
        \frac{xf(x)-2}{x-1} = c \Rightarrow f(x) = \frac{c(x-1)+2}{x} = c + \frac{2-c}{x}
        \]  
        求导得
        \[
        f'(x) = \frac{c-2}{x^2} 
        \]
        由$f'(1) = 2018$得 $c=2020$,故
        \[
        f(x)=2020-\frac{2018}{x}
        \]
        经检验,原方程式左式=右式$=2020x-2018y$,故
        \[
        f(2018)=2019
        \]
    \end{solution}
    
    \question 设函数 $f(x), g(x)$ 在 $\mathbb{R}$ 上连续且满足
    \[
    \begin{cases}
    f(x+y) = f(x)g(y) + g(x)f(y) \\
    g(x+y) = f(x)f(y) + g(x)g(y)
    \end{cases}
    \]
    且 $g(x)$ 为偶函数,$\;g(x) > 0$ 恒成立,$\;g(0) \ne \dfrac{1}{2}$,已知 $f(2) + g(2) = \dfrac{\sqrt{5}}{15}$,求 $f(8) - g(8)$。
    \begin{solution}
        两式相加得
        \[
        f(x+y) + g(x+y) = (f(x) + g(x))(f(y) + g(y))
        \]
        设 $h(x) = f(x) + g(x)$,则 $h(x+y) = h(x)h(y)$,得 $h(x) = a^x$。
        两式相减得
        \[
        g(x+y) - f(x+y) = (g(x) - f(x))(g(y) - f(y))
        \]
        设 $k(x) = g(x) - f(x)$,同理得 $k(x) = b^x$,于是
        \[
        f(x) = \frac{h(x) - k(x)}{2} = \frac{a^x - b^x}{2},\quad
        g(x) = \frac{h(x) + k(x)}{2} = \frac{a^x + b^x}{2}
        \]
        由于 $g(x)$ 为偶函数,$\,g(x) = g(-x)$,即
        \[
        \frac{a^x + b^x}{2} = \frac{a^{-x} + b^{-x}}{2} \Rightarrow (a^x + b^x) a^x b^x = b^x + a^x \Rightarrow (ab)^x = 1 \Rightarrow ab = 1,\ b = \frac{1}{a}
        \]
        故
        \[
        f(x) = \frac{a^x - a^{-x}}{2},\quad
        g(x) = \frac{a^x + a^{-x}}{2}
        \]
        由 $f(2) + g(2) = \dfrac{\sqrt{5}}{15}$ 可知
        \[
        h(2) = a^2 = \frac{\sqrt{5}}{15}
        \Rightarrow a^8 = \left( \frac{\sqrt{5}}{15} \right)^4 = \frac{1}{2025}
        \]
        所以
        \[
        f(8) - g(8) = -k(8) = -b^8 = -\frac{1}{a^8} = -2025
        \]
    \end{solution}

    \question 设 $a,b>0$使得关于 $x$ 的方程 $\sqrt{|x|}+\sqrt{|x+a|}=b$ 恰有三个相异实数解 $x_1,x_2,x_3$, 且 $x_1<x_2<x_3=b$, 求 $a+b$ 的值。
    \begin{solution}
        令 $t=x+\frac{a}{2}$, 则关于 $t$ 的方程 $$\sqrt{\left|t-\frac{a}{2}\right|}+\sqrt{\left|t+\frac{a}{2}\right|}=b$$ 恰有三个不同的实数解 $$t_i=x_i+\frac{a}{2},i=1,2,3$$

        由于 $f(t)=\sqrt{|t-\frac{a}{2}|}+\sqrt{|t+\frac{a}{2}|}$ 为偶函数, 故方程 $f(t)=b$ 的三个实数解关于原点对称分布, 从而必有 $b=f(0)=\sqrt{2a}$
        
        现求方程 $f(t)=\sqrt{2a}$ 的实数解:
        \begin{itemize}
            \item 当 $|t|\le\frac{a}{2}$ 时, $f(t)=\sqrt{\frac{a}{2}-t}+\sqrt{\frac{a}{2}+t}=\sqrt{a+\sqrt{a^2-4t^2}}\le\sqrt{2a}$, 等号成立当且仅当 $t=0$
            \item 当 $t>\frac{a}{2}$ 时, $f(t)$ 单调递增, 且当 $t=\frac{5a}{8}$ 时 $f(t)=\sqrt{2a}$
            \item 当 $t<-\frac{a}{2}$ 时, $f(t)$ 单调递减, 且当 $t=-\frac{5a}{8}$ 时 $f(t)=\sqrt{2a}$
        \end{itemize}
        从而方程 $f(t)=\sqrt{2a}$ 恰有三个实数解 $$t_1=-\frac{5a}{8},\ t_2=0,\ t_3=\frac{5a}{8}$$
        由条件知 $b=x_3=t_3-\frac{a}{2}=\frac{a}{8}$, 结合 $b=\sqrt{2a}$ 得 $a=128$,于是 $a+b=\frac{9a}{8}=144$
    \end{solution}  
\end{questions}
\pagebreak

\begin{center}
  {\fontsize{30pt}{26pt}\selectfont
    \hypertarget{不等式}{不等式} \label{不等式}
  }
\end{center}
\separator
\vspace{1pt}

\begin{questions}
    \question 求满足不等式
    \[
    \frac{4x^{2}}{(1-\sqrt{1+2x})^{2}}<2x+9
    \]
    的实数 $x$ 的集合。
    \begin{solution}
        首先定义域要求 $1+2x \ge 0$,即 $x \ge -\dfrac{1}{2}$,且分母不为零,即 $1-\sqrt{1+2x} \ne 0$,得 $x \ne 0$,发现
        \[
        \frac{4x^{2}}{(1-\sqrt{1+2x})^{2}}=\left(\frac{1-(1+2x)}{1-\sqrt{1+2x}}\right)^2= (1+\sqrt{1+2x})^2
        \]
        由$(1+\sqrt{1+2x})^2 < 2x+9$解得
        \[
        \sqrt{1+2x} < \frac{7}{2} \Rightarrow x < \frac{45}{8}.
        \]
        故$x$ 的取值范围为
        \[
        \left[-\frac{1}{2},0\right) \cup \left(0,\frac{45}{8}\right).
        \]
    \end{solution}

    \question 已知
    \[
    \begin{cases}
    \displaystyle \sqrt{x(x-3)} + \frac{8}{y^2} + \frac{y^2}{\sqrt{x(x-3)}} = 6, \\
    x+y < 0
    \end{cases}
    \]
    求 $(x,y)$。
    \begin{solution}
        据题意,
        \[
        6=\sqrt{x(x-3)} + \frac{8}{y^2} + \frac{y^2}{\sqrt{x(x-3)}} \ge 2|y|+\frac{8}{y^2} \Rightarrow 3y^2 \ge y^2|y|+4 \tag{1}
        \]
        当$y>0$,化简得
        \[
        y^3-3y^2+4=(y+1)(y-2)^2\le0
        \]
        由于$(y-2)^2\ge0$且$y+1>0$,故
        \[
        y=2 \Rightarrow x=-1,4
        \]
        经检验得(-1,2),(4,2)均不合题意;同理当$y<0$可解得
        \[
        y=-2,\Rightarrow x=-1,4
        \]
        经检验得,原方程式只有唯一解$(-1,-2)$
    \end{solution}

    \question 解不等式
    \[
    |x|^3 - 2x^2 - 4|x| + 3 < 0,
    \]
    \begin{solution}
        令 \( t = |x| \),不等式化为
        \[
        t^3 - 2t^2 - 4t + 3 < 0
        \]
        求根得 \( t = 3 \) 和 \( t = \dfrac{\sqrt{5}-1}{2}  \)。
        于是
        \[
        \frac{\sqrt{5}-1}{2} < |x| < 3
        \]
        解得
        \[
        x \in \left( -3, -\frac{\sqrt{5}-1}{2} \right) \cup \left( \frac{\sqrt{5}-1}{2}, 3 \right)
        \]
    \end{solution}

    \question 假设 $a>1$ 且 $x,y>0$ 满足
    \[
    (\log_a x)^2 + (\log_a y)^2 - \log_a(x^2y^2) \le 2, \quad \log_a y \ge 1.
    \]
    求 $\log_a(x^2y)$ 的取值区间。
    \begin{solution}
        设 $u = \log_a x$, $v = \log_a y$。不等式化为
        \[
        u^2 + v^2 - 2(u+v) \le 2, \quad v \ge 1.
        \]
        这是以 $(1,1)$ 为圆心、半径为 $2$ 的半圆的上半部分。我们要求 $2u+v$ 的取值范围。

        最小值在左下角 $(-1,1)$,得 $2u+v = -1$。

        最大值出现在斜率为 $-2$ 的直线与半圆切点,即联立
        \[
        v-1 = \frac{1}{2}(u-1), \quad (u-1)^2 + (v-1)^2 = 4.
        \]
        解得
        \[
        (u,v) = \left(1+\frac{4}{\sqrt{5}},\,1+\frac{2}{\sqrt{5}}\right) \implies 2u+v = 3 + 2\sqrt{5}.
        \]
        因此区间为 $[-1, 3+2\sqrt{5}]$。
    \end{solution}

    \question 已知不等式 
    \[
    |x^2 + ax + 2| \ge |x + 1|
    \]
    对任意 $x \in \mathbb{R}$ 恒成立,求实数 $a$ 的取值范围。
    \begin{solution}
        将不等式两边平方得
        \[
        (x^2 + ax + 2)^2 \ge (x + 1)^2
        \]
        \[
        (x^2 + ax + 2)^2 - (x + 1)^2 \ge 0
        \]
        \[
        [(x^2 + ax + 2) - (x + 1)][(x^2 + ax + 2) + (x + 1)] \ge 0
        \]
        \[
        [x^2 + (a - 1)x + 1][x^2 + (a + 1)x + 3] \ge 0
        \]
        因为两个二次式领导系数为正,两个因式需对任意 $x$ 恒非负。
        
        对二次式 $x^2 + (a - 1)x + 1$,
        \[
        \Delta_1  = (a - 1)^2 - 4 \le 0
        \Rightarrow -1 \le a \le 3
        \]
        对二次式 $x^2 + (a + 1)x + 3$,
        \[
        \Delta_2  = (a + 1)^2 - 12 \le 0
        \Rightarrow -1 - 2\sqrt{3} \le a \le -1 + 2\sqrt{3}
        \]
        综上所述,
        \[
        a \in [-1, -1 + 2\sqrt{3}]
        \]
    \end{solution}
        
    \question 求 
    \[
    |2|x-2|-5| \le 3
    \]
    的正整数解。
    \begin{solution}
        情况一:$\ |2|x-2|-5| = 3$  
        \[
        \Rightarrow
        \begin{cases}
        2|x-2| - 5 = 3 \Rightarrow |x-2| = 4 \Rightarrow 
        \begin{cases}
        x-2 = 4 \Rightarrow x = 6 \\
        x-2 = -4 \Rightarrow x = -2 \notin \mathbb{N}
        \end{cases} \\
        2|x-2| - 5 = -3 \Rightarrow |x-2| = 1 \Rightarrow 
        \begin{cases}
        x = 3 \\
        x = 1
        \end{cases}
        \end{cases}
        \]
        情况二:$\ |2|x-2|-5| = 2$  
        \[
        \Rightarrow
        \begin{cases}
        2|x-2| - 5 = 2 \Rightarrow |x-2| = \frac{7}{2} \notin \mathbb{Z} \\
        2|x-2| - 5 = -2 \Rightarrow |x-2| = \frac{3}{2} \notin \mathbb{Z}
        \end{cases}
        \]
        情况三:$\ |2|x-2|-5| = 1$  
        \[
        \Rightarrow
        \begin{cases}
        2|x-2| - 5 = 1 \Rightarrow |x-2| = 3 \Rightarrow 
        \begin{cases}
        x = 5 \\
        x = -1 \notin \mathbb{N}
        \end{cases} \\
        2|x-2| - 5 = -1 \Rightarrow |x-2| = 2 \Rightarrow 
        \begin{cases}
        x = 4 \\
        x = 0 \notin \mathbb{N}
        \end{cases}
        \end{cases}
        \]
        情况四:$\ |2|x-2|-5| = 0$  
        \[
        \Rightarrow |x-2| = \frac{5}{2} \notin \mathbb{Z}
        \]
        综上,正整数解为
        \[
        x = 1, 3, 4, 5, 6
        \]
    \end{solution}

    \question 若实数 $a$ 可使得对任意实数 $x$,不等式 
    \[
    |4x-3a|+|5x-4a|\ge a^{2}
    \]
    恒成立,求 $a$ 的范围。
    \begin{solution}
        记
        \[
        f(x)=|4x-3a| +|5x-4a|= 4\left( \Big|x-\frac{3}{4}a\Big|+ \Big|x-\frac{4}{5}a\Big|\right)+\Big|x-\frac{4}{5}a\Big|
        \]
        由三角不等式得
        \[
        f(x) \ge 4\Big| \Big(x-\frac{3}{4}a\Big)- \Big(x-\frac{4}{5}a\Big)\Big|+\Big|x-\frac{4}{5}a\Big|
        = \frac{|a|}{5}+ \Big|x-\frac{4}{5}a\Big|
        \]
        取 $x=\dfrac{4}{5}a$,得 $f(x)$ 的最小值为 $\dfrac{|a|}{5}$,欲使不等式对任意 $x$ 恒成立需有
        \[
        \frac{|a|}{5} \ge |a|^2 \Rightarrow |a|(5|a|-1)\le 0  \Rightarrow \quad 0\le |a|\le \frac{1}{5}
        \]
    \end{solution}

    \question 已知实数 $a,b$满足 $a+b\ge 3$, 试证 
    \[
    |a-2b^2|+|b-2a^2| \ge 6
    \]
    \begin{solution}
        由三角不等式,
        \begin{align*}
        \left| a-2b^2\right|+ \left| b-2a^2\right| 
        &\ge \left| a-2b^2+b-2a^2\right| \\
        &=\left| 2(a^2+b^2)-(a+b)\right| \\
        &\ge 2(a^2+b^2)-(a+b) \\
        &=(a+b)^2+(a-b)^2-(a+b) \\
        &\ge 3^2-3=6
        \end{align*}
        故原不等式得证。    
    \end{solution}

    \question 若二次实系数多项式函数 $f(x)$ 满足
    \[
    \begin{cases}
    -1 \le f(1) \le 3 \\
    6 \le f(2) \le 10 \\
    2 \le f(4) \le 24
    \end{cases}
    \]
    求 $f(7)$ 的最大值。
    \begin{solution}
        设 $f(x)=ax^2+bx+c$,则
        \[
        f(1)=a+b+c,\quad f(2)=4a+2b+c,\quad f(4)=16a+4b+c
        \]
        解得
        \[
        a = \frac{1}{3}f(1) - \frac{1}{2}f(2) + \frac{1}{6}f(4),\;
        b = -2f(1) + \frac{5}{2}f(2) - \frac{1}{2}f(4),\;
        c = \frac{8}{3}f(1) - 2f(2) + \frac{1}{3}f(4)
        \]
        因此
        \[
        f(7)=49a+7b+c = 5f(1)-9f(2)+5f(4)
        \]
        取 $f(1)=3,f(2)=6,f(4)=24$,得
        \[
        f(7)_{\max} =5\cdot 3 - 9\cdot 6 +5\cdot 24 =81
        \]
    \end{solution}

    \question 若 $a<b$,已知对任意实数 $x$,均有 $ax^2+bx+c\ge 0$,且 $m<\dfrac{4a+3b+2c}{b-a}$ 恒成立,求实数 $m$ 的范围。
    \begin{solution}
        由 $ax^2+bx+c\ge0\ (\forall x)$,可知二次项系数满足
        \[
        a>0,\quad b^2-4ac\le0
        \]
        写成
        \[
        \frac{4a+3b+2c}{b-a}
        =\frac{4\cdot\frac{a}{b}+3+2\cdot\frac{c}{b}}{1-\frac{a}{b}}
        \]
        记 $t=\dfrac{a}{b}$,则 $0<t<1$。由 $b^2-4ac\le0$ 得 $\dfrac{c}{b}\ge \dfrac{b}{4a}=\dfrac{1}{4t}$,于是
        \[
        \frac{4a+3b+2c}{b-a}
        \ge \frac{4t+3+2\cdot\frac{1}{4t}}{1-t}
        = \frac{4t+3+\frac{1}{2t}}{1-t}
        \]
        记
        \[
        k=\frac{4t+3+\frac{1}{2t}}{1-t}
        \]
        其中$0<t<1$,整理得关于 $t$ 的二次方程
        \[
        (2k+8)t^2+(6-2k)t+1=0
        \]
        为使该方程在 $0<t<1$ 上有解,判别式必须非负,
        \[
        (6-2k)^2-4(2k+8)\ge0 \Rightarrow k^2-8k+1 \ge 0 
        \]
        解得
        \[
        k \ge 4+\sqrt{15}>0
        \]
        于是对任意符合条件的 $a,b,c$,都有
        \[
        \frac{4a+3b+2c}{b-a}\ge 4+\sqrt{15}
        \]
        由题意 $m$ 满足 $m<\dfrac{4a+3b+2c}{b-a}$ 恒成立,故 
        \[
        m<4+\sqrt{15}
        \]
    \end{solution}

    \question 已知 $a>0, b>0, a+b=1$,证明不等式
    \[
    \frac{1}{a} + \frac{1}{b} + \frac{1}{ab} \ge 8
    \]
    \begin{solution}
        由 AM-GM 不等式,
        \begin{align*}
        \frac{1}{a} + \frac{1}{b} + \frac{1}{ab} &= \frac{1}{a} + \frac{1}{b} + \frac{a+b}{ab} \\
        &= 2\left(\frac{1}{a} + \frac{1}{b}\right) \\
        &= 2 \cdot \frac{a+b}{ab} \\
        &= 4 + 2\left(\frac{a}{b} + \frac{b}{a}\right) \\
        &\ge 4 + 2\cdot 2 = 8
        \end{align*}
    \end{solution}

    \question 设 \(a,b\) 均为正数,求
    \[
    f(a,b)=ab (72 - 3a - 4b)
    \]
    的最大值。
    \begin{solution}
        由 AM-GM 不等式,
        \begin{align*}
        f(a,b) &= \frac{1}{12} \cdot 3a \cdot 4b \cdot (72 - 3a - 4b) \\
        &\leq \frac{1}{12}\left(\frac{3a+4b+72-3a-4b}{3}\right)^3 = 1152        
        \end{align*}
        当$a=b=\frac{1}{2}$时等号成立。
    \end{solution}

    \question 设 $a,b,c$ 为正实数且满足 $a+b^{2}+c^{3}=11$,求 $abc$ 的最大值。
    \begin{solution}
        由 AM-GM 不等式,
        \begin{align*}
        11&=\frac{a}{6}+\frac{a}{6}+\frac{a}{6}+\frac{a}{6}+\frac{a}{6}+\frac{a}{6}+\frac{b^{2}}{3}+\frac{b^{2}}{3}+\frac{b^{2}}{3}+\frac{c^{3}}{2}+\frac{c^{3}}{2}\\
        &=11\sqrt[11]{\left(\frac{a}{6}\right)^6\left(\frac{b^2}{3}\right)^3\left(\frac{c^3}{2}\right)^2} = 11\sqrt[11]{\frac{(abc)^6}{6^6\cdot 27\cdot 4}}
        \end{align*}
        即
        \[
        abc \le 6\sqrt{3}\,\sqrt[3]{2}
        \]
        当且仅当
        \[
        \frac{a}{6}=\frac{b^2}{3}=\frac{c^3}{2}=1,
        \]
        即 $a=6,\ b=\sqrt{3},\ c=\sqrt[3]{2}$ 时等号成立。
    \end{solution}

    \question 实数 $x, y, z$ 满足 $x^2 + y^2 + z^2 = 1$,求 $\sqrt{2}xy + yz$ 的最大值。
    \begin{solution}
        由 AM-GM 不等式,
        \[
        \frac{\sqrt{3}}{2} x^2 + \frac{1}{\sqrt{3}} y^2 \ge 2\sqrt{\frac{1}{2} x^2 y^2} = \sqrt{2} xy,
        \]
        \[
        \frac{1}{2\sqrt{3}} y^2 + \frac{\sqrt{3}}{2} z^2 \ge 2\sqrt{\frac{1}{4} y^2 z^2} = yz.
        \]
        将两式相加得
        \[
         \sqrt{2} xy + yz \le \frac{\sqrt{3}}{2}(x^2 + y^2 + z^2) = \frac{\sqrt{3}}{2}
        \]
    \end{solution}

    \question 设 $x,y,z$ 为不全为 $0$ 的实数,求分式
    \[
    \frac{xy+2yz}{x^{2}+y^{2}+z^{2}}
    \]
    的最大值。
    \begin{solution}
        由 AM-GM 不等式,
        \[
        x^2+\frac{1}{5}y^2 \ge 2sqrt{\frac{1}{5}x^2y^2}=\frac{2}{\sqrt{5}} xy
        \]
        同理,
        \[
        \frac{4}{5}y^2 + z^2 \ge 2\sqrt{\frac{4}{5}y^2 z^2} = \frac{4}{\sqrt{5}} yz 
        \]
        两式相加得
        \[
        x^2 + y^2 + z^2 \ge \frac{2}{\sqrt{5}} (xy + 2 yz)
        \quad \Rightarrow \quad \frac{xy+2yz}{x^2+y^2+z^2} \le \frac{\sqrt{5}}{2}.
        \]
    \end{solution}

    \question 求
    \[
    f(x)=\frac{x^4 + 2x^3 + 3x^2 + 2x + 10}{x^2 + x + 1}
    \]
    的最小值。
    \begin{solution}
        由 AM-GM 不等式,
        \begin{align*}
        f(x)&=x^2 + x + 1 + \frac{9}{x^2 + x + 1} \\
        &\geq 2 \sqrt{(x^2 + x + 1)\cdot  \frac{9}{x^2 + x + 1}}  =6
        \end{align*}
        当 $(x^2 + x + 1)^2=9$即$x=2$或$x=-1$时,$\ f(x)$ 取得最大值 $6$
    \end{solution}

    \question 求函数 $$y=\frac{4^{x}+1}{2^{x}+1}$$ 的最小值。
    \begin{solution}
        由 AM-GM 不等式,
        \begin{align*}
        y=\frac{4^{x}-1}{2^{x}+1}+\frac{2}{2^{x}+1}
        &=2^{x}+1+\frac{2}{2^{x}+1}-2\\
        &\ge2\sqrt{(2^{x}+1)\cdot\frac{2}{2^{x}+1}}-2=2\sqrt{2}-2
        \end{align*}

        当 $(2^x+1)^2=2$即$x=\log_{2}(\sqrt2-1)$时,$\ y$ 取得最大值 $2\sqrt{2}-2$
    \end{solution}
    
    \question 求函数 $$f(x)=\log 2 \log 5-\log 2x \log 5x$$ 的最大值。 
    \begin{solution}
        \begin{align*}
            f(x)&=\log 2\log 5-(\log 2+\log x)(\log 5+\log x)\\
            &=-(\log 2 + \log 5)\log x - \log^2 x \\
            &=-\log x - \log^2 x \\
            &=-\left(\log x + \frac{1}{2}\right)^2 + \frac{1}{4} \le \frac{1}{4}
        \end{align*}
        当 $\log x = -\dfrac{1}{2}$即 $x=\dfrac{\sqrt{10}}{10}$ 时,$\ f(x)$ 取得最大值 $\dfrac{1}{4}$.
    \end{solution}

    \question 求函数 $$f(x,y) = x^2 + 2xy + 2y^2 + 2x + 4y + 5$$ 的最小值。
    \begin{solution}
        配方求最小值:
        \begin{align*}
        f(x,y) &= x^2 + 2xy + 2y^2 + 2x + 4y + 5 \\
        &= x^2 + 2(y+1)x + (y+1)^2 + y^2 + 2y + 4 \\
        &= [x+(y+1)]^2 + (y+1)^2 + 3 \ge 3
        \end{align*}
        当 $x=0, y=-1$ 时$f(x,y)$取到最小值 $3$。
    \end{solution}
    
    \question \(x,y \in \mathbb{R}\),求
    \[
    x^4 + 4x^3 + 8x^2 + 4xy + 6x + 4y^2 + 10
    \]
    的最小值。 
    \begin{solution}
        先配$y$项,写成\[
        f(x,y)=(x+2y)^2 + x^4 + 4x^3 + 7x^2+ 6x + 10
        \]
        再用火眼金睛注意到\[
        x^4 + 4x^3 + 7x^2+ 6x + 2=(x + 1)^2 (x^2 + 2 x + 2)    \]故原式为\[
        (x+2y)^2+(x + 1)^2 (x^2 + 2 x + 2)+8 \geq 8
        \]
        当$x=-1,y=\dfrac12$时,$\ f(x,y)$ 取得最小值 $8$.
    \end{solution}

    \question 已知 $x,y>0$, 证明
    \[
    \sqrt{x^2+y^2} > \sqrt[3]{x^3+y^3}
    \]
    \begin{solution}
        发现$\sqrt{x^2+y^2} > \sqrt[3]{x^3+y^3}$等价于
        \[
        x^6+3x^4y^2+3x^2y^4+y^6 > x^6+2x^3y^3+y^6 \Leftrightarrow x^2y^2(3x^2-2xy+3y^2) > 0
        \]
        其中关于$x$的两次函数$3x^2-2xy+3y^2$判别式为
        \[
        \Delta = 4y^2-4(3)(3)(y^2) =-32y^2< 0
        \]
        故$3x^2-2xy+3y^2 > 0$,得证
        \[
        x^2y^2(3x^2-2xy+3y^2) > 0
        \]
    \end{solution}

    \question 已知实数 $a,b,c$ 满足 $c<b<a,a+b+c=1,a^2+b^2+c^2=1$,求证 
    \[
    1<a+b<\frac{4}{3}.
    \]
    \begin{solution}
        由 $a^2+b^2+c^2=1,a+b+c=1$ 可得
        \[
        ab = \frac{(a+b)^2-(a^2+b^2)}{2} 
        = \frac{(1-c)^2-(1-c^2)}{2} 
        = c^2-c.
        \]
        所以 $a+b=1-c$,且 $a,b$ 是方程
        \[
        x^2-(1-c)x+c^2-c=0
        \]
        的两个相异实根,其中判别式为
        \[
        \Delta=(1-c)^2-4(c^2-c) >0 \Rightarrow -\frac{1}{3}<c<1.
        \]
        又有
        \[
        (c-a)(c-b)=c^2-c(a+b)+ab
        = c^2-c(1-c)+c^2-c
        = 3c^2-2c.
        \]
        因为 $c<b<a$,所以 $c-a<0,\;c-b<0$,故 $(c-a)(c-b)>0$,即
        \[
        3c^2-2c>0 \;\;\Longrightarrow\;\; c<0 \;\; \text{或} \;\; c>\frac{2}{3}.
        \]
        综合 $-\frac{1}{3}<c<1$ 与上述条件,得到
        \[
        -\frac{1}{3}<c<0 \quad \text{或} \quad \frac{2}{3}<c<1.
        \]
        若 $\dfrac{2}{3}<c<1$,则
        \[
        0<a+b<\frac{1}{3}
        \]
        但由 $c<b<a$ 得 $a+b>c>\dfrac{2}{3}$,矛盾,因此只能有 $-\dfrac{1}{3}<c<0$,则
        \[
        1<a+b<\frac{4}{3}
        \]
    \end{solution}

    \question 已知$a,b,c$为正实数,求
    \[
    \left\lfloor \frac{a+b}{c} \right\rfloor + \left\lfloor \frac{b+c}{a} \right\rfloor + \left\lfloor \frac{c+a}{b} \right\rfloor
    \]
    的最小值。
    \begin{solution}
        由$\lfloor x\rfloor > x-1,\;x \notin \mathbb{Z}$及AM-GM不等式,
        \begin{align*}
        S=\left\lfloor \frac{a+b}{c} \right\rfloor + \left\lfloor \frac{b+c}{a} \right\rfloor + \left\lfloor \frac{c+a}{b} \right\rfloor &> \frac{a}{b} + \frac{b}{a} + \frac{b}{c} + \frac{c}{b} + \frac{c}{a} + \frac{a}{c}-3 \\
        &\ge 2+2+2-3 =3
        \end{align*}
        由于$S \in \mathbb{Z}$,最小值为$S_{\min}=4$,例如可在$a=6,b=8,c=9$时取得。
    \end{solution}

    \question 已知$x>0$,求
    \[
    x\lfloor x\rfloor + \left\lfloor\frac{1}{x}\right\rfloor + x + \frac{1}{x} + x\lceil x\rceil + \left\lceil\frac{1}{x}\right\rceil
    \]
    的最小值。
    \begin{solution}
        设
        \[
        f(x)=x\lfloor x\rfloor + \left\lfloor\frac{1}{x}\right\rfloor + x + \frac{1}{x} + x\lceil x\rceil + \left\lceil\frac{1}{x}\right\rceil
        \]
        观察到$f(1)=6$,若$x > 1$,即$0 < \dfrac{1}{x} \le 1$,
        \[
        x\lfloor x\rfloor + \left\lfloor\frac{1}{x}\right\rfloor \ge x >1,\quad
        x\lceil x\rceil + \left\lceil\frac{1}{x}\right\rceil \ge 2x+1 >3,\quad
        x + \frac{1}{x} > 2
        \]
        于是当$x > 1,f(x) > 6 =f(1)$。若$0 < x \le \dfrac{1}{2}$,即$2 \le \dfrac{1}{x}$
        \[
        x\lfloor x\rfloor + \left\lfloor\frac{1}{x}\right\rfloor \ge 2, \quad
        x\lceil x\rceil + \left\lceil\frac{1}{x}\right\rceil \ge x + 2 > 2, \quad
        x + \frac{1}{x} > 2  
        \]
        于是当$0 < x \le \dfrac{1}{2},f(x) > 6 =f(1)$。若$\dfrac{1}{2} < x < 1$,即$1 < \dfrac{1}{x} < 2$,
        \[
        f(x)=1+x+\frac{1}{x}+x+2=3+2x+\frac{1}{x} \ge 3+2\sqrt{2}
        \]
        等号成立当且仅当
        \[
        2x=\frac{1}{x}
        \]
        即$x=\dfrac{\sqrt{2}}{2}$。由于$3+2\sqrt{2}<6,f(x)$的最小值为$3+2\sqrt{2}$。
    \end{solution}

    \question 若实数 $x, y$ 满足 $4x^2 - 2xy + 2y^2 = 1$, 则 $3x^2 + xy + y^2$ 的最大值与最小值的和为
    \begin{solution}
        \textcolor{red}{(待解)}
    \end{solution}

    \question 求函数 
    \[
    y = \sqrt{x^2 - 3x + 2} + \sqrt{2 + 3x - x^2}
    \]
    的最大值和最小值。
    \begin{solution}
        考虑
        \[
        1 = \frac{\sqrt{x^2 - 3x + 2}}{y} + \frac{\sqrt{2 + 3x - x^2}}{y}
        \]
        令 
        \[
        \cos^2 \alpha = \frac{\sqrt{x^2 - 3x + 2}}{y}, \sin^2 \alpha = \frac{\sqrt{2 + 3x - x^2}}{y}
        \]
        则
        \[
        \cos^4 \alpha + \sin^4 \alpha = \frac{x^2 - 3x + 2}{y^2} + \frac{2 + 3x - x^2}{y^2}= \frac{4}{y^2}
        \]
        于是
        \[
        y^2 = \frac{4}{\cos^4 \alpha + \sin^4 \alpha} = \frac{8}{2 - \sin^2 2\alpha}
        \]
        由于 $0 \le \sin^2 2\alpha \le 1$,所以 
        \[
        4 \le y^2 \le 8 \Rightarrow 2 \le y \le 2\sqrt{2}
        \]
    \end{solution}

    \question 若 
    \[
    x^6 + \frac{y^3}{3} + \frac{z}{4} = 1,
    \]
    求 $x^{12}y^6z^3$ 的最大值,其中 $x, y, z \in \mathbb{R}^{+}$。
    \begin{solution}
        由AM-GM不等式,
        \begin{align*}
        1 &= x^6 + \frac{y^3}{3} + \frac{z}{4} = \frac{x^6}{2} + \frac{x^6}{2} + \frac{y^3}{6} + \frac{y^3}{6} + \frac{z}{12} + \frac{z}{12} + \frac{z}{12} \\
        &\ge 7 \sqrt[7]{\frac{x^{12}y^6z^3}{2^2 \cdot 6^2 \cdot 12^3}}
        \end{align*}
        因此
        \[
        x^{12}y^6z^3 \le \frac{12^5}{7^7},
        \]
        等号成立当且仅当
        \[
        \frac{x^6}{2} = \frac{y^3}{6} = \frac{z}{12} = \frac{1}{7}
        \]
        即
        \[
        x = \sqrt[6]{\frac{2}{7}}, \quad y = \sqrt[3]{\frac{6}{7}}, \quad z = \frac{12}{7}
        \]
    \end{solution}

    \question 已知正实数$x, y, z$,求函数
    \[
    f(x, y, z) = \frac{x}{y} + \sqrt{\frac{y}{z}} + \sqrt[3]{\frac{z}{x}},
    \]
    的最小值。
    \begin{solution}
        由AM-GM不等式,
        \begin{align*}
        f(x, y, z) &= \frac{x}{y} + \frac{1}{2}\sqrt{\frac{y}{z}} + \frac{1}{2}\sqrt{\frac{y}{z}} + \frac{1}{3}\sqrt[3]{\frac{z}{x}} + \frac{1}{3}\sqrt[3]{\frac{z}{x}} + \frac{1}{3}\sqrt[3]{\frac{z}{x}}\\
        &\ge 6 \cdot \sqrt[6]{\frac{x}{y} \cdot \frac{1}{2}\sqrt{\frac{y}{z}} \cdot \frac{1}{2}\sqrt{\frac{y}{z}} \cdot \frac{1}{3}\sqrt[3]{\frac{z}{x}} \cdot \frac{1}{3}\sqrt[3]{\frac{z}{x}} \cdot \frac{1}{3}\sqrt[3]{\frac{z}{x}}} \\ 
        &= 6 \cdot \sqrt[6]{\frac{1}{2^2 \cdot 3^3} \cdot \frac{xyz}{yzx}} = \sqrt[3]{4} \cdot \sqrt{3}
        \end{align*}
        等号成立当且仅当
        \[
        \frac{x}{y} = \frac{1}{2}\sqrt{\frac{y}{z}} = \frac{1}{3}\sqrt[3]{\frac{z}{x}} \Rightarrow x :y :z = 2 :2\sqrt[3]{2}\sqrt{3} :3\sqrt{3}
        \]
    \end{solution}

    \question 已知 $x, y, z$ 均为正实数,求 
    \[
    \frac{xy + yz}{x^2 + y^2 + z^2}
    \]
    的最大值。
    \begin{solution}
        由AM-GM不等式,
        \begin{align*}
        \frac{xy + yz}{x^2 + y^2 + z^2} 
        &= \frac{\frac{x}{y} + \frac{z}{y}}{(\frac{x}{y})^2 + 1 + (\frac{z}{y})^2} \\
        &= \frac{\frac{x}{y} + \frac{z}{y}}{[(\frac{x}{y})^2 + \frac{1}{2}] + [(\frac{z}{y})^2 + \frac{1}{2}]} \\
        &\le \frac{\frac{x}{y} + \frac{z}{y}}{2(\frac{x}{y})\sqrt{\frac{1}{2}} + 2(\frac{z}{y})\sqrt{\frac{1}{2}}} \\
        &= \frac{\sqrt{2}}{2}
        \end{align*}
        等号成立当且仅当
        \[
        \frac{1}{2} = (\frac{x}{y})^2 = (\frac{z}{y})^2 \Rightarrow x=z=\frac{\sqrt{2}}{2}y
        \]
    \end{solution}

    \question 证明 
    \[
    x + \frac{4}{x-1} \ge 2\sqrt{x+3},
    \]
    其中 $x > 1$。
    \begin{solution}
        利用比较法。令 $u = x-1 > 0$,原不等式即
        \[
        u + 1 + \frac{4}{u} \ge 2\sqrt{u+4},
        \]
        现在比较 $u + 1 + \frac{4}{u}$ 与 $2\sqrt{u+4}$ 的大小,
        \[
        \frac{u + 1 + \frac{4}{u}}{2\sqrt{u+4}} = \frac{u^2 + (u+4)}{2u\sqrt{u+4}} \ge \frac{2\sqrt{u^2(u+4)}}{2u\sqrt{u+4}} = 1
        \]
        等号在 $u^2 = u + 4$ 即$u=\dfrac{1+\sqrt{17}}{2}$时成立,故原不等式得证。
    \end{solution}

    \question 已知 $x_1 + x_2 = 2$,证 
    \[
    \frac{1}{1+x_1} + \frac{1}{1+x_2} \ge \frac{x_1}{1+x_1^2} + \frac{x_2}{1+x_2^2},
    \]
    其中 $x_1,x_2$ 都是正数。
    \begin{solution}
        由右式 
        \[
        \frac{x_1}{1+x_1^2} + \frac{x_2}{1+x_2^2} = \frac{1}{\frac{1}{x_1} + x_1} + \frac{1}{\frac{1}{x_2} + x_2} \le \frac{1}{2} + \frac{1}{2} = 1
        \]
        由左式 
        \[
        \frac{1}{1+x_1} + \frac{1}{1+x_2} \ge 2\sqrt{\frac{1}{(1+x_1)(1+x_2)}}
        \]
        考虑分母 
        \[
        (1+x_1)(1+x_2) = 1 + x_1 + x_2 + x_1x_2 = 1 + 2 + x_1(2-x_1) = -x_1^2 + 2x_1 + 3 = -(x_1-1)^2 + 4 \le 4
        \]
        于是
        \[
        \frac{1}{1+x_1} + \frac{1}{1+x_2} \ge 2\sqrt{\frac{1}{4}} = 1
        \]
        因此左式 $\ge 1 \ge$ 右式,原不等式成立,等号成立当且仅当$x_1=x_2=1$。
    \end{solution}

    \question 若 $a, b, c$ 是$\triangle ABC$的三边,证明 
    \[
    \frac{1}{4} < \frac{(a+b)(b+c)(c+a)}{(a+b+c)^3} \le \frac{8}{27}
    \]
    \begin{solution}
        由AM-GM不等式, 
        \[
        (a+b) + (b+c) + (c+a) \ge 3 \sqrt[3]{(a+b)(b+c)(c+a)}\]
        于是有
        \[
        \frac{(a+b)(b+c)(c+a)}{(a+b+c)^3} \le \frac{8}{27}
        \]
        等号成立当且仅当$a=b=c$,又 
        \[
        a+b > c \Rightarrow 2a+2b > a+b+c \Rightarrow \frac{a+b}{a+b+c} > \frac{1}{2},
        \]
        同理可得 
        \[
        \frac{b+c}{a+b+c} > \frac{1}{2},\quad \frac{c+a}{a+b+c} > \frac{1}{2}
        \]
        设 
        \[
        \frac{a+b}{a+b+c} = \frac{1+x}{2},\quad \frac{b+c}{a+b+c} = \frac{1+y}{2},\quad \frac{c+a}{a+b+c} = \frac{1+z}{2}
        \]
        由于 
        \[
        \frac{a+b}{a+b+c} + \frac{b+c}{a+b+c} + \frac{c+a}{a+b+c} = 2 \Rightarrow x+y+z = 1
        \]
        于是 
        \[
        \frac{(a+b)(b+c)(c+a)}{(a+b+c)^3} = \frac{(1+x)(1+y)(1+z)}{8} > \frac{1+x+y+z}{8} = \frac{1}{4}
        \]
        故得证。
    \end{solution}

    \question 若 $a > b > c$,求证
    \[
    frac{1}{a-b} + \frac{1}{b-c} \ge \frac{4}{a-c}
    \]
    \begin{solution}
        因为 $a-c > 0$,由柯西不等式,
        \[
        (a-c)\left(\frac{1}{a-b} + \frac{1}{b-c}\right) = [(a-b) + (b-c)]\left(\frac{1}{a-b} + \frac{1}{b-c}\right) \ge (1 + 1)^2 = 4
        \]
        故原不等式得证,当$a - b = b - c$,即$b = \dfrac{a + c}{2}$时等号成立。
    \end{solution}

    \question 证 
    \[
    (x+y+z)(\frac{1}{y+z} + \frac{1}{x+z} + \frac{1}{x+y}) \ge \frac{9}{2}
    \]
    \begin{solution}
        由柯西不等式,
        \begin{align*}
        &(x+y+z)(\frac{1}{y+z} + \frac{1}{x+z} + \frac{1}{x+y})\\
        &= \frac{1}{2}[(x+y) + (y+z) + (z+x)](\frac{1}{y+z} + \frac{1}{x+z} + \frac{1}{x+y}) \\
        &\ge \frac{1}{2}(1+1+1)^2 = \frac{9}{2}
        \end{align*}
        等号成立当且仅当$x = y = z$。
    \end{solution}

    \question 设 $a, b$ 为不相等的正数,试证 
    \[
    (a+b)(a^3+b^3) > (a^2+b^2)^2
    \]
    \begin{solution}
        由柯西不等式,左式改写为
        \[
        [(\sqrt{a})^2 + (\sqrt{b})^2][(\sqrt{a^3})^2 + (\sqrt{b^3})^2] \ge (\sqrt{a} \cdot \sqrt{a^3} + \sqrt{b} \cdot \sqrt{b^3})^2 = (a^2 + b^2)^2
        \]
        若等号成立,则必有一实数 $k$ 存在,使得 $\sqrt{a^3} = k\sqrt{a}, \sqrt{b^3} = k\sqrt{b}$ 于是 $a=b=k$,与 $a, b$ 不相等的假设矛盾,故得证
        \[
        (a+b)(a^3+b^3) > (a^2+b^2)^2
        \]
    \end{solution}

    \question 已知$a, b, c$为正实数,求证
    \[
    \frac{a}{b+c} + \frac{b}{c+a} + \frac{c}{a+b} \ge \frac{3}{2}
    \]
    \begin{solution}
        由柯西不等式,
        \begin{align*}
        &\frac{a}{b+c} + \frac{b}{c+a} + \frac{c}{a+b} + 3 \\
        &= \frac{a+b+c}{b+c} + \frac{b+c+a}{c+a} + \frac{c+a+b}{a+b} \\
        &= (a+b+c)(\frac{1}{b+c} + \frac{1}{c+a} + \frac{1}{a+b}) \\
        &= \frac{1}{2}[(b+c) + (c+a) + (a+b)](\frac{1}{b+c} + \frac{1}{c+a} + \frac{1}{a+b}) \\
        &\ge \frac{1}{2}(1+1+1)^2 = \frac{9}{2}
        \end{align*}
        故原不等式得证,等号成立当且仅当$a=b=c$。
    \end{solution}

    \question 已知$a, b, c$ 为 $\triangle ABC$ 的三边,求证
    \[
    \frac{a^2}{b+c-a} + \frac{b^2}{c+a-b} + \frac{c^2}{a+b-c} \ge a+b+c.
    \] 
    \begin{solution}
        由柯西不等式, 
        \begin{align*}
        &(a+b+c)\left(\frac{a^2}{b+c-a} + \frac{b^2}{c+a-b} + \frac{c^2}{a+b-c}\right) \\
        &= [(b+c-a) + (c+a-b) + (a+b-c)]\left(\frac{a^2}{b+c-a} + \frac{b^2}{c+a-b} + \frac{c^2}{a+b-c}\right) \\
        &\ge (a+b+c)^2
        \end{align*}
        整理即得原不等式,等号成立当且仅当$a=b=c$。
    \end{solution}

    \question 当 $\alpha, \beta \in \left[0, \dfrac{\pi}{4}\right]$,求 
    \[
    \sin(\alpha-\beta) + 2\sin(\alpha+\beta)
    \]
    的最大值。
    \begin{solution}
        由柯西不等式,
        \begin{align*}
        \sin(\alpha-\beta) + 2\sin(\alpha+\beta)
        &= 3\sin\alpha\cos\beta + \sin\beta\cos\alpha \\
        &\le \sqrt{(9\sin^2 \alpha + \cos^2 \alpha)(\cos^2 \beta + \sin^2 \beta)}\\
        &= \sqrt{8\sin^2 \alpha + 1} \\
        &\le \sqrt{8\sin^2 \frac{\pi}{4} + 1} = \sqrt{5}
        \end{align*}
        等号在 $\alpha = \dfrac{\pi}{4}, \beta = \tan^{-1} \dfrac{1}{3}$ 成立。
    \end{solution}

    \question 若$x, y, z \in \mathbb{R}^+$满足 
    \[
    \frac{1}{x^3} + \frac{3}{y^3} + \frac{12}{z^3} = 1,
    \]
    求 $xy^2z$ 的最小值。
    \begin{solution}
        由AM-GM不等式,
        \[
        1 = \frac{1}{x^3} + \frac{3}{2y^2} + \frac{3}{2y^2} + \frac{12}{z} \ge 4 \sqrt[4]{\frac{27}{(xy^2z)^3}}
        \]
        即
        \[
        xy^2z \ge 12 \sqrt[3]{4}
        \]
        等号成立当且仅当$x = \sqrt[3]{4}, y = \sqrt[3]{6}, z  = 2\sqrt[3]{6}$。
    \end{solution}

    \question 求多元函数 
    \[
    f(x, y, z) = \sqrt[4]{\frac{2x}{y}} + \sqrt{\frac{3y}{2z}} + \frac{z}{2x}
    \]
    的最小值,其中 $D_f = \mathbb{R}^+$。
    \begin{solution}
        由AM-GM不等式,
        \begin{align*}
        f(x, y, z) &= \frac{1}{4}\sqrt[4]{\frac{2x}{y}} + \frac{1}{4}\sqrt[4]{\frac{2x}{y}} + \frac{1}{4}\sqrt[4]{\frac{2x}{y}} + \frac{1}{4}\sqrt[4]{\frac{2x}{y}} + \frac{1}{2}\sqrt{\frac{3y}{2z}} + \frac{1}{2}\sqrt{\frac{3y}{2z}} + \frac{z}{2x} \\
        &= 7 \sqrt[7]{\frac{1}{4^4}\cdot\frac{1}{2^2}\cdot2\cdot\frac{3}{2}\cdot\frac{1}{2}} = \frac{7}{4}\sqrt[7]{24}
        \end{align*}
        等号成立当且仅当
        \[
        \frac{1}{4}\sqrt[4]{\frac{2x}{y}} = \frac{1}{2}\sqrt{\frac{3y}{2z}} = \frac{z}{2x} \Rightarrow x :y :z = 2 \cdot 24^{\frac{4}{7}} :4 :24^{\frac{5}{7}}
        \]
    \end{solution}

    \question 设 $x, y, z$ 为实数,若 
    \[
    \frac{(x-1)^2}{9} + \frac{(y+1)^2}{4} + \frac{(z-2)^2}{3} = 1,
    \]
    求 $x+y+z$ 之最大值及最小值。
    \begin{solution}
        由柯西不等式,
        \[
        1=\frac{(x-1)^2}{9} + \frac{(y+1)^2}{4} + \frac{(z-2)^2}{3} \ge \frac{(x-1+y+1+z-2)^2}{9+4+3}
        \]
        给出
        \[
        (x-1+y+1+z-2)^2 \le 16 \Rightarrow -2 \le x+y+z \le 6
        \]
        等号成立当且仅当比例相等,即
        \[
        \frac{x-1}{9} = \frac{y+1}{4} = \frac{z-2}{3} = k
        \]
        代入约束条件:
        \[
        9k^2 + 4k^2 + 3k^2 = 1 \Rightarrow k = \pm \frac{1}{4}
        \]
        当 $k = \dfrac{1}{4}$ 时,$x = \dfrac{13}{4}, y = 0, z = \dfrac{11}{4}$,取得最大值 6,当 $k = -\dfrac{1}{4}$ 时,$x = -\dfrac{5}{4}, y = -2, z = \dfrac{5}{4}$,此时取得最小值 -2。
    \end{solution}

\question 设 $0 < \theta < \dfrac{\pi}{2}$,试求 
\[
\frac{3}{\cos \theta} + \frac{2}{\sin \theta}
\]
的最小值。
\begin{solution}
    由赫尔德不等式,
    \[
    \left(\frac{3}{\cos \theta} + \frac{2}{\sin \theta}\right)^2(\cos^2 \theta+\sin^2 \theta) \ge (\sqrt[3]{9}+\sqrt[3]{4})^3
    \]
    故最小值为 $(\sqrt[3]{9}+\sqrt[3]{4})^\frac{3}{2}$,当$\tan \theta = \sqrt[3]{\dfrac{2}{3}}$ 时等号成立。
    \textcolor{red}{不安全,赫尔德不等式不给出最小值}
\end{solution}

    \question 设实数 $x, y$ 满足 
    \[
    x^2 + y^2 = 6x - 4y - 9,
    \]
    求 $2x - 3y$ 的最大值和最小值。
    \begin{solution}
        原方程配方得圆的方程:
        \[
        (x-3)^2 + (y+2)^2 = 4
        \]
        设 $x = 3 + 2\cos \alpha, y = -2 + 2\sin \alpha$,则
        \begin{align*}
        2x - 3y &= 2(3 + 2\cos \alpha) - 3(-2 + 2\sin \alpha) \\
        &= 12 + 4\cos \alpha - 6\sin \alpha \\
        &= 12 + 2\sqrt{13} \sin(\phi - \alpha)
        \end{align*}
        其中 $\phi$ 满足 
        \[
        \sin \phi = \frac{2}{\sqrt{13}}, \cos \phi = \frac{3}{\sqrt{13}},
        \]
        当 $\sin(\phi - \alpha) = 1$ 时,取得最大值$12 + 2\sqrt{13}$,此时
        \[
        x = 3 + \frac{4}{\sqrt{13}}, y = -2 - \frac{6}{\sqrt{13}}
        \]
        当 $\sin(\phi - \alpha) = -1$ 时,取得最小值$12 - 2\sqrt{13}$,此时
        \[
        x = 3 - \frac{4}{\sqrt{13}}, y = -2 + \frac{6}{\sqrt{13}}
        \]
    \end{solution}

\question 若 $a, b$ 均为实数且满足 
\[
a\sqrt{1-b^2} + b\sqrt{1-a^2} = 1,
\]
求 $a^2 + b^2$ 的值。
\begin{solution}
    由根号的定义可知 
    \[
    1-a^2 \ge 0,1-b^2 \ge 0 \Rightarrow a^2 \le 1, b^2 \le 1
    \]
    设 
    \[
    a = \sin \alpha,\quad b = \sin \beta, \quad \alpha, \beta \in \left[-\frac{\pi}{2}, \frac{\pi}{2}\right]
    \]
    原式变为
    \[
    \sin \alpha \cos \beta + \sin \beta \cos \alpha =\sin(\alpha + \beta) = 1
    \]
    由此得 
    \[
    \alpha + \beta = \frac{\pi}{2} \Rightarrow \beta = \frac{\pi}{2} - \alpha 
    \]
    则 
    \[
    b = \sin\left(\frac{\pi}{2} - \alpha\right) = \cos \alpha
    \]
    因此 
    \[
    a^2 + b^2 = \sin^2 \alpha + \cos^2 \alpha = 1
    \]
\textcolor{red}{(移走)}
\end{solution}
\begin{solution}
    由根号的定义可知 
    \[
    1-a^2 \ge 0,1-b^2 \ge 0 \Rightarrow a^2 \le 1, b^2 \le 1
    \]
    由柯西不等式,
    \[
    1=\left( a\sqrt{1-b^2} + b\sqrt{1-a^2} \right)^2 \le (a^2 + b^2)(1-b^2+1-a^2)
    \]
    设 $x = a^2 + b^2$,上式变为
    \[
    1 \le x(2 - x) \Rightarrow x^2 - 2x + 1 \le 0 \Rightarrow (x - 1)^2 \le 0
    \]
    由于实数的平方非负,因此上述不等式只有在等号成立时有解
    \[
    (x - 1)^2 = 0 \Rightarrow a^2 + b^2 = x = 1
    \]
\end{solution}

\question 当点 $P$ 沿着直线 $y = 2x + 6$ 移动,点 $Q$ 在椭圆 
\[
\dfrac{x^2}{6} + \dfrac{y^2}{4} = 1
\] 
运动时,求线段 $PQ$ 的长度最小值。
\begin{solution}
线段 $PQ$ 的最小值即为平行于已知直线的椭圆切线与该直线之间的距离。设与 $y = 2x + 6$ 平行的椭圆切线方程为 $y = 2x + k$,代入椭圆方程整理得,
\[
2x^2 + 3(4x^2 + 4kx + k^2) = 12 \Rightarrow 14x^2 + 12kx + (3k^2 - 12) = 0
\]
令判别式 $\Delta = (12k)^2 - 4(14)(3k^2 - 12) = 0$:
\[
\Delta = (12k)^2 - 4(14)(3k^2 - 12) = 144k^2 - 168k^2 + 672 = 0 \Rightarrow k = \pm 2\sqrt{7}
\]
取较近的切线 $y = 2x + 2\sqrt{7}$,故
\[
PQ_{min} = \frac{|6 - 2\sqrt{7}|}{\sqrt{2^2 + (-1)^2}} = \frac{6 - 2\sqrt{7}}{\sqrt{5}} = \frac{6\sqrt{5} - 2\sqrt{35}}{5}
\]
\textcolor{red}{(移走)}
\end{solution}

    \question 若$a, b$ 都是正数,证明 
    \[
    a + b \ge \sqrt[3]{a^2b} + \sqrt[3]{b^2a}
    \]
    \begin{solution}
        由因式分解得
        \begin{align*}
        &a + b - \left(\sqrt[3]{a^2b} + \sqrt[3]{b^2a}\right) \\
        &= (\sqrt[3]{a} + \sqrt[3]{b})(\sqrt[3]{a^2} - \sqrt[3]{ab} + \sqrt[3]{b^2}) - \sqrt[3]{ab}(\sqrt[3]{a} + \sqrt[3]{b}) \\
        &= (\sqrt[3]{a} + \sqrt[3]{b})(\sqrt[3]{a^2} - 2\sqrt[3]{ab} + \sqrt[3]{b^2}) \\
        &= (\sqrt[3]{a} + \sqrt[3]{b})(\sqrt[3]{a} - \sqrt[3]{b})^2 \ge 0
        \end{align*}
        故原不等式成立。
    \end{solution}
    \begin{solution}
        由赫尔德不等式,
        \[
        (a+b)(a+b)(b+a) \ge \left(\sqrt[3]{a^2b} + \sqrt[3]{b^2a}\right)^3
        \]
        故原不等式成立。
    \end{solution}

    \question 设 $a, b, c \in \mathbb{R}$,证明
    \[
    a^2 + ac + c^2 + 3b(a + b + c) \ge 0.
    \]
    \begin{solution}
        欲证
        \[
        3b^2 + 3(a+c)b + (a^2 + ac + c^2) \ge 0
        \]
        对于二次函数 $f(b) = 3b^2 + 3(a+c)b + (a^2 + ac + c^2)$,其判别式为
        \[
        \Delta = [3(a+c)]^2 - 4(3)(a^2 + ac + c^2) = -3(a-c)^2
        \]
        因为 $-3(a-c)^2 \le 0$ 且二次项系数 $3 > 0$,故该二次式恒大于或等于 $0$,故原不等式得证。
    \end{solution}
    \begin{solution}
        由配方法,
        \[
        a^2 + ac + c^2 + 3b(a + b + c)=3 \left( b + \frac{a+c}{2} \right)^2 + \frac{1}{4}(a-c)^2\ge 0
        \]
    \end{solution}

    \question 已知 $x, y, z$ 是正实数,求 
    \[
    \frac{x}{y} + \left(\frac{y}{z}\right)^2 + \left(\frac{z}{x}\right)^3
    \]
    的最小值。
    \begin{solution}
        由AM-GM不等式,原式即
        \[
        6 \cdot \frac{x}{6y} + 3 \cdot \frac{y^2}{3z^2} + 2 \cdot \frac{z^3}{2x^3}= \ge 11 \sqrt[11]{\frac{1}{6^6 \cdot 3^3 \cdot 2^2}} = \frac{11}{6} \sqrt[11]{72}
        \]
    \end{solution}

    \question 设 $a, b, c, d$ 为实数,满足 $a^2 + b^2 = c^2 + d^2$,求
    \[
    \frac{20bc + 14bd + 20ad - 14ac}{a^2 + b^2 + c^2 + d^2}
    \]
    的最大值。
    \begin{solution}
        由柯西不等式,
        \begin{align*}
        (a(20d - 14c) + b(20c + 14d))^2 
        &\le (a^2 + b^2)((20d - 14c)^2 + (20c + 14d)^2)\\
        &596(c^2 + d^2)
        \end{align*}
        由此可知
        \[
        |20bc + 14bd + 20ad - 14ac| \le \sqrt{a^2 + b^2} \cdot \sqrt{596(c^2 + d^2)}
        \]
        记 $a^2 + b^2 = c^2 + d^2 = S$,代入上式得
        \[
        \frac{20bc + 14bd + 20ad - 14ac}{2S} \le \frac{\sqrt{S} \cdot \sqrt{596S}}{2S} = \frac{\sqrt{596}S}{2S} = \frac{\sqrt{596}}{2}=\sqrt{149}
        \]
    \end{solution}

    \question 实数 $x,y,z,t$ 满足
    \[
    x^2+y^2=16,\quad z^2+t^2=25,\quad xt-yz=20,
    \]
    求 $xz$ 的最大值。
    \begin{solution}
        由柯西不等式,
        \[
        16\cdot 25=(x^2+y^2)(t^2+(-z)^2)\ge (xt-yz)^2=20^2
        \]
        此时等号成立,可解得
        \[
        \frac{x}{t}=\frac{y}{-z}=\frac{4}{5} \Rightarrow xz=-\frac{5}{4}xy
        \]
        由 $x^2+y^2=16$ 可得
        \[
        |xy|\le \frac{x^2+y^2}{2}=8
        \]
        所以
        \[
        -10= -\frac{5}{4}\cdot 8 \le xz \le -\frac{5}{4}\cdot(-8) =10
        \]
        因此 $x z$ 的最大值为$10$。
    \end{solution}

    \question 设 \(a_1,a_2,\dots,a_{18}\) 均为大于 1 的实数,求  
    \[
    \frac{\log_{a_1} 2023 + \log_{a_2} 2023 + \cdots + \log_{a_{18}} 2023}{\log_{a_1 a_2 \cdots a_{18}} 2023}
    \]
    的最小值。
    \begin{solution}
        利用换底公式,原式即
        \[
        \frac{\sum_{i=1}^{18} \frac{1}{\log_{2023} a_i}}{\frac{1}{\sum_{i=1}^{18} \log_{2023} a_i}} 
        = \left(\sum_{i=1}^{18} \frac{1}{\log_{2023} a_i}\right) \left(\sum_{i=1}^{18} \log_{2023} a_i\right).
        \]
        由柯西不等式,
        \[
        \left(\sum_{i=1}^{18} \log_{2023} a_i\right) \left(\sum_{i=1}^{18} \frac{1}{\log_{2023} a_i}\right) \ge \left(\sum_{i=1}^{18} 1\right)^2 = 18^2 = 324.
        \]
    \end{solution}

    \question 存在 2017 个正实数 \( x_1, x_2, \ldots, x_{2017} \)使得
    \[
    \sum_{i=1}^{2017} x_i = \sum_{i=1}^{2017} \frac{1}{x_i} = 2018,
    \]
    求
    \[
    x_1 + \frac{1}{x_1}
    \]
    的最大可能值。
    \begin{solution}
        由柯西不等式,
        \[
        (2018 - x_1)\left(2018 - \frac{1}{x_1} \right)=\left(\sum_{i=2}^{2017} x_i\right)\left(\sum_{i=2}^{2017} \frac{1}{x_i}\right)  \geq 2016^2.
        \]
        展开左边得
        \[
        2018^2 - 2018\left(x_1 + \frac{1}{x_1} \right) + 1 \geq 2016^2.
        \]
        于是
        \[
        x_1 + \frac{1}{x_1} \leq \frac{8069}{2018}.
        \]
        等号成立当且仅当 $x_i = \dfrac{1}{x_i}$,即$x_i=1,i=2,3,\dots,2017$
    \end{solution}

    \question 已知 $a, b, c$ 为正数,且 $a+b+c=1$。求 $$ \left(\frac{1}{a}-1\right)\left(\frac{1}{b}-1\right)\left(\frac{1}{c}-1\right)$$ 的最小值。
    \begin{solution}
    由AM-GM 不等式,
        \begin{align*}
        \left(\frac{1}{a}-1\right)\left(\frac{1}{b}-1\right)\left(\frac{1}{c}-1\right)
        &= \left(\frac{1-a}{a}\right)\left(\frac{1-b}{b}\right)\left(\frac{1-c}{c}\right)\\[1mm]
        &= \left(\frac{b+c}{a}\right)\left(\frac{a+c}{b}\right)\left(\frac{a+b}{c}\right)\\[1mm]
        &\ge \left(\frac{2\sqrt{bc}}{a}\right)\left(\frac{2\sqrt{ac}}{b}\right)\left(\frac{2\sqrt{ab}}{c}\right) =8 
        \end{align*}
        当且仅当 $a=b=c=\dfrac{1}{3}$ 时等号成立。
    \end{solution}

    \question 已知 $a,b,c$ 为三角形三边,证明不等式
    \[
    \frac{a}{b+c-a} + \frac{b}{c+a-b} + \frac{c}{a+b-c} \ge 3
    \]
    \begin{solution}
        发现
        \begin{align*}
        \frac{a}{b+c-a} + \frac{b}{c+a-b} + \frac{c}{a+b-c} &\ge 3 \\
        \iff \left(\frac{2a}{b+c-a}+1\right) + \left(\frac{2b}{c+a-b}+1\right) + \left(\frac{2c}{a+b-c}+1\right) &\ge 9 \\
        \iff (a+b+c) \left(\frac{1}{b+c-a} + \frac{1}{c+a-b} + \frac{1}{a+b-c}\right) &\ge 9
        \end{align*}
        由于 
        \[
        (b+c-a)+(c+a-b)+(a+b-c) = a+b+c,
        \]
        且
        \[
        b+c-a>0,\;c+a-b>0,\;a+b-c>0
        \]
        由AM-HM不等式或柯西不等式,原不等式得证。
    \end{solution}

    \question 设 \(x, y, z\) 为正实数且 \(xyz = 1\),证明:
    \[
    \frac{x}{y} + \frac{y}{z} + \frac{z}{x} \geq x + y + z.
    \]
    \begin{solution}
        由AM-GM 不等式,
        \[
        \frac{1}{3} \left( \frac{x}{y} + \frac{x}{y} + \frac{y}{z} \right) \geq \sqrt[3]{\frac{x^2}{yz}} = x,
        \]
        \[
        \frac{1}{3} \left( \frac{y}{z} + \frac{y}{z} + \frac{z}{x} \right) \geq \sqrt[3]{\frac{y^2}{xz}} = y, 
        \]
        \[
        \frac{1}{3} \left( \frac{z}{x} + \frac{z}{x} + \frac{x}{y} \right) \geq \sqrt[3]{\frac{z^2}{xy}} = z. 
        \]
        三式相加得
        \[
        \frac{x}{y} + \frac{y}{z} + \frac{z}{x} \geq x + y + z.
        \]
        故证毕。
    \end{solution}
    
    \question 若 $x^2+y^2+z^2=1$,求 $x+2y+3z$ 的最大值。
    \begin{solution}
        由柯西不等式,
        \[
        (x^2+y^2+z^2)(1^2+2^2+3^2) \ge (x+2y+3z)^2
        \Rightarrow x+2y+3z \le \sqrt{14}
        \]
        当且仅当 $(x,y,z)=\dfrac{1}{\sqrt{14}}(1,2,3)$ 时等号成立。
    \end{solution}

    \question 若 $x,y,z \in \mathbb{R}^+$ 且 $x+y+z=3$,求 $\dfrac{yz+4xz+9xy}{xyz}$ 的最小值。
    \begin{solution}
        由柯西不等式,
        \[
        \frac{x^2}{1^2} + \frac{y^2}{2^2} + \frac{z^2}{3^2} \ge \frac{(x+y+z)^2}{1+4+9} = \frac{9}{14}.
        \]
        当且仅当 $(x,y,z)=\left(\dfrac{1}{2},1,\dfrac{3}{2}\right)$ 时等号成立。
    \end{solution}

    \question 若 $a,b,c \in \mathbb{R}^+$ 且 $a+b+c=6$,求
    \[
    \left(a+\frac{1}{b}\right)^2 + \left(b+\frac{1}{c}\right)^2 + \left(c+\frac{1}{a}\right)^2
    \]
    的最小值。
    \begin{solution}
        由柯西不等式,
        \begin{align*}
        \left(a+\frac{1}{b}\right)^2 + \left(b+\frac{1}{c}\right)^2 + \left(c+\frac{1}{a}\right)^2 
        &\ge \frac{\left(a+b+c + \frac{1}{a}+\frac{1}{b}+\frac{1}{c}\right)^2}{1+1+1} \\
        &\ge \frac{\left(6 + \frac{9}{a+b+c}\right)^2}{1+1+1} = \frac{75}{4}
        \end{align*}
        等号成立当且仅当 $a=b=c=2$。
    \end{solution}

    \question 已知实数$a,b,c,d$满足 
    \[
    ab + bc + cd = 8,\quad b^{2} + c^{2} = 2,
    \]
    试求 $a^{2} + d^{2}$ 的最小值。
    \begin{solution}
        首先有
        \[
        b^2+c^2 \ge 2bc \Rightarrow bc \le 1
        \]
        由柯西不等式,
        \[ 
        (a^2+d^2)(b^2+c^2) \ge (ab+cd)^2 
        \]
        即
        \[ 
        (a^2+d^2) \ge \frac{(ab+cd)^2}{b^2+c^2} = \frac{(8-bc)^2}{2} \ge \frac{(8-1)^2}{2} = \frac{49}{2} 
        \]
        当 $bc=1$ 且 $\dfrac{a}{c} = \dfrac{d}{b}$, 即 $b=c=-1,a=d=-\dfrac{7}{2}$ 时等号成立,此时$a^2+d^2$ 取最小值 $\dfrac{49}{2}$。
    \end{solution}

    \question 若 $3a+4b=15$,求 $\sqrt{a^2+b^2}$ 的最小值。
    \begin{solution}
        由柯西不等式,
        \[
        5\sqrt{a^2+b^2} = \sqrt{3^2+4^2}\cdot\sqrt{a^2+b^2} \ge |3a+4b| = 15 \implies \sqrt{a^2+b^2} \ge 3
        \]
        当且仅当 $a=\dfrac{9}{5},b=\dfrac{12}{5}$ 时等号成立。
    \end{solution}

    \question 若 $abcd=1$,求 $(1+a)(1+b)(1+c)(1+d)$ 的最小值。
    \begin{solution}
        由 AM–GM不等式,
        \[
        (1+a)(1+b)(1+c)(1+d) \ge (2\sqrt{a})(2\sqrt{b})(2\sqrt{c})(2\sqrt{d}) =16
        \]
        等号成立当且仅当 $a=b=c=d=1$。
    \end{solution}

    \question 若 $x\in \mathbb{R}$,求 
    \[
    \frac{x^2+2-\sqrt{x^4+4}}{x}
    \] 
    的最大值。
    \begin{solution}
        设$x>0$,由AM–GM不等式,
        \begin{align*}
        \frac{x^2+2-\sqrt{x^4+4}}{x} 
        &= x+\frac{2}{x}-\sqrt{x^2+\frac{4}{x^2}}\\
        &= \sqrt{x^2+\frac{4}{x^2}+4}-\sqrt{x^2+\frac{4}{x^2}}\\
        &=\frac{4}{\sqrt{x^2+\frac{4}{x^2}+4}+\sqrt{x^2+\frac{4}{x^2}}}\\
        &\le \frac{4}{\sqrt{4+4}+\sqrt{4}} = 2\sqrt{2}-2
        \end{align*}
        当且仅当 $x=\sqrt{2}$ 时等号成立。
    \end{solution}

    \question 若 $a,b$ 为正实数且满足 $ab(a+b)=2000$,求 $\dfrac{1}{a} + \dfrac{1}{b} + \dfrac{1}{ab}$ 的最小值。
    \begin{solution}
        由 AM–GM不等式,
        \[
        2000=ab(a+b)\le \left(\frac{a+b}{2}\right)(a+b) \Rightarrow a+b \ge 20
        \]
        再由 AM–GM不等式,
        \begin{align*}
        \frac{1}{a} + \frac{1}{b} + \frac{1}{ab} 
        &= \frac{1}{2a}+\frac{1}{2a}+\frac{1}{2b}+\frac{1}{2b}+\frac{1}{a+b} \\
        &\ge 5\sqrt[5]{\frac{1}{16a^2b^2(a+b)}} \\
        &= 5 \sqrt[5]{\frac{a+b}{16a^2b^2(a+b)^2}} \\
        &\ge 5 \sqrt[5]{\frac{20}{16\cdot 2000^2}} = \frac{1}{4}
        \end{align*}
        等号成立当且仅当$a=b=10$。以下为错误示范:

        由 AM–GM不等式,
        \[
        \frac{1}{a} + \frac{1}{b} + \frac{1}{ab} \ge 3 \sqrt[3]{\frac{1}{ab(a+b)}}= 3 \sqrt[3]{\frac{1}{2000}}=\frac{3\sqrt[3]{4}}{20}
        \]
        但等号成立当且仅当$a=b=0$,不合题意。

        由 AM–GM不等式,
        \[
        \frac{1}{a} + \frac{1}{b} + \frac{1}{ab} = \frac{a+b}{ab} +\frac{1}{ab} = \frac{(a+b)^2}{2000}+\frac{1}{2(a+b)}+\frac{1}{2(a+b)}
        \ge 3\sqrt[3]{\frac{1}{8000}} =\frac{3}{20}
        \]
        但等号成立当且仅当$a+b=10$且$ab=200$,解得$a,b$为非实数,不合题意。
    \end{solution}

    \question 若 $a, b, c$ 皆为正实数满足 $a+b+c=100$,求 $\sqrt{a} + \sqrt{ab} + \sqrt{abc}$ 的最大值。
    \begin{solution}
        由AM–GM不等式,
        \[
        \sqrt{a} + \sqrt{ab} + \sqrt{abc} = a + 2\sqrt{\frac{a}{4}\cdot b} + 3\sqrt{\frac{a}{12}\cdot \frac{b}{3} \cdot \frac{4c}{3}} \le a+\frac{a}{4}+b+\frac{a}{12}+\frac{b}{3} + \frac{4c}{3} = \frac{4}{3}(a+b+c) =\frac{400}{3}
        \]  
        等号成立当且仅当 $a = \dfrac{160}{21}, b = \dfrac{40}{21}, c = \dfrac{10}{21}$。
    \end{solution}

    \question 若 $x_1, x_2, x_3, x_4, x_5$ 为实数,且
    \[
    x_1+x_2+x_3+x_4+x_5=8, \quad x_1^2+x_2^2+x_3^2+x_4^2+x_5^2=16
    \]  
    求 $x_5$ 的最大值。
    \begin{solution}
        由柯西不等式,
        \[
        (x_1^2+x_2^2+x_3^2+x_4^2)(1+1+1+1) \ge (x_1+x_2+x_3+x_4)^2 \Rightarrow 4(16-x_5^2) \ge (8-x_5)^2
        \]  
        解得 
        \[
        0 \le x_5 \le \frac{16}{5}
        \]  
        当 $x_1=x_2=x_3=x_4=\dfrac{6}{5},x_5 = \dfrac{16}{5}$时等号成立。
    \end{solution}

    \question 若 $a,b,c$ 为正实数且 $a+b+c \le 9$,求 $(a^2+b^2+c^2)(2ab+2bc+2ca+5)$ 的最大值。
    \begin{solution}
        由AM–GM不等式,
        \begin{align*}
        (a^2+b^2+c^2)(2ab+2bc+2ca+5) 
        &\le \left(\frac{a^2+b^2+c^2 + 2ab+2bc+2ca+5}{2}\right)^2 \\
        &= \left(\frac{(a+b+c)^2+5}{2}\right)^2 \\
        &\le \left(\frac{9^2+5}{2}\right)^2 = 1849
        \end{align*}
        当$a^2+b^2+c^2=43,ab+bc+ca=19$时等号成立。
    \end{solution}

    \question $x,y,z$ 为正实数且满足 $x+y+z=91$,求
    \[
    \frac{yz}{x} + \frac{xz}{y} + \frac{xy}{z}
    \]
    的最小值。
    \begin{solution}
        由排序不等式,
        \[
        \sum_{cyc} \frac{xy}{z} \ge \sum_{cyc} \frac{xy}{y} = x+y+z=91
        \]
    \end{solution}

    \question 已知正实数 $a,b$ 满足 $a + b = 1$,求 
    \[
    \left(a + \dfrac{1}{a}\right)^{2} + \left(b + \dfrac{1}{b}\right)^{2}
    \]
    的最小值。
    \begin{solution}
        由柯西不等式或QM-AM不等式, 
        \[
        \left(a+\frac{1}{a}\right)^2+\left(b+\frac{1}{b}\right)^2
        \ge
        \frac12\left(a+b+\frac{1}{a}+\frac{1}{b}\right)^2=\frac12\left(1+\frac{1}{ab}\right)^2
        \]
        又由AM-GM不等式,
        \[
        ab\le\left(\frac{a+b}{2}\right)^2=\frac14
        \]
        故
        \[
        \left(a+\frac{1}{a}\right)^2+\left(b+\frac{1}{b}\right)^2
        \ge
        \frac12\cdot (1+4)^2=\frac{25}{2}
        \]
        当且仅当 $a=b=\dfrac12$ 时等号成立,此时有最小值 $\dfrac{25}{2}$。
    \end{solution}

    \question 已知 $x,y,z$ 为正实数,求 $\dfrac{(x^4+1)(y^4+1)(z^4+1)}{xy^2z}$ 的最小值。
    \begin{solution}
        由AM–GM不等式,
        \begin{align*}
        \frac{(x^4y^4+1)(y^4z^4+1)(z^4x^4+1)}{xy^2z}
        &=\left(x^3+\frac{1}{3x}+\frac{1}{3x}+\frac{1}{3x}\right)\left(y^2+\frac{1}{y^2}\right)\left(z^3+\frac{1}{3z}+\frac{1}{3z}+\frac{1}{3z}\right)\\
        &\ge 4\sqrt[3]{\frac{1}{27}}\cdot 2 \cdot 4\sqrt[3]{\frac{1}{27}} = \frac{32\sqrt{3}}{9}
        \end{align*}
    \end{solution}

    \question 已知$x \in [0,2]$,求 $\sqrt{x}+4\sqrt{1-\dfrac{x}{2}}$ 的最大值。
    \begin{solution}
        发现
        \[
        \sqrt{x}+4\sqrt{1-\frac{x}{2}}=\sqrt{x}+\sqrt{2^3}\cdot \sqrt{2-x}
        \]
        由柯西不等式,
        \[
        (\sqrt{x}+\sqrt{2^3}\cdot \sqrt{2-x})^2\le(1+8)(x+2-x)=18 \Leftrightarrow \sqrt{x}+4\sqrt{1-\frac{x}{2}} \le 3\sqrt{2}
        \]
        等号成立当且仅当$x=\dfrac{2}{9}$
    \end{solution}

    \question 若 $x,y$ 为正数, 且 
    \[
    x^2+\frac{y^2}{45}=1, 
    \]
    试求 
    \[
    \frac{2}{1-x}+\frac{75}{10-y}
    \] 
    的最小值。
    \begin{solution}
        由AM-GM不等式,
        \[
        \frac{1-x+\frac{x}{2}+\frac{x}{2}}{3} \ge \sqrt[3]{\frac{x^2(1-x)}{4}}
        \]
        于是
        \[
        \frac{2}{1-x} = \frac{x^2}{2}\cdot \frac{4}{x^2(1-x)}\ge \frac{x^2}{2}\cdot\left( \frac{3}{1-x+\frac{x}{2}+\frac{x}{2}}\right)^3 = \frac{27x^2}{2} \tag{1}
        \]
        由AM-GM不等式,
        \[
        \frac{y+(10-y)}{2} \ge \sqrt{y(10-y)}
        \]
        故
        \[
        \frac{75}{10-y}
        = \frac{75}{10}\left( \frac{y}{10-y}+1\right)
        \ge \frac{15}{2}\cdot \frac{y^2}{\left( \frac{y+(10-y)}{2}\right)^2}+\frac{15}{2}
        = \frac{15}{2}\cdot \frac{y^2}{25}+ \frac{15}{2} \tag{2}
        \]
        由 $(1),(2)$,
        \[
        \frac{2}{1-x}+\frac{75}{10-y}
        \ge \frac{27}{2}\left( x^2+\frac{y^2}{45}\right)+\frac{15}{2}= 21
        \]
        等号成立当且仅当$x=\dfrac{2}{3},y=5$。
    \end{solution}

    \question 若正实数$a,b,c$满足 $a+b+c=21$,求 $a^2+3b^2+5c^2$ 的最小值。
    \begin{solution}
        由排序不等式,
        \begin{align*}
        a^2+3b^2+5c^2 
        &= a^2+b^2+c^2 +2b^2+2c^2+2c^2 \\
        &\ge a^2+b^2+c^2 + 2ab+2ac+2bc = (a+b+c)^2 = 441
        \end{align*}
        当$(a,b,c)=(21,0,0),\left(\dfrac{21}{2},\dfrac{21}{2},0\right),(7,7,7)$时等号成立。
    \end{solution}

    \question 已知 $x,y$ 为实数,求 $x^6+y^6-54xy$ 的最小值。
    \begin{solution}
        由AM–GM不等式,
        \[
        \frac{x^6+y^6+27+27+27+27}{6} \ge \sqrt[6]{x^6y^6 \cdot 3^{12}} \Rightarrow x^6y^6 - 54xy \ge -108
        \]  
        当$x=y=\pm\sqrt{3}$时等号成立。
    \end{solution}

\question 若 $x,y,z$ 为正实数,且满足 $x+y+z=1$,求 $x(x+y)^2(y+z)^3(z+x)^4$ 的最大值。
\begin{solution}
    由AM–GM不等式,
    \[
    \frac{5(x+y+z)}{10}=\frac{x + 2(x+y) + 3(y+z) + 4\left(\frac{z+x}{2}\right)}{10} \ge \sqrt[10]{\frac{x(x+y)^2(y+z)^3(z+x)^4}{2^4}} 
    \]  
    故
    \[
    x(x+y)^2(y+z)^3(z+x)^4 \le \frac{1}{64}
    \]
    \textcolor{red}{等号成立$y=0,x=z=\frac{1}{2}$？}
\end{solution}

    \question 设 $a,b,c,d$ 为正实数,且满足 $abcd=1$,求
    \[
    a^2+b^2+c^2+d^2+ab+ac+ad+bc+bd+cd
    \]
    的最小值。
    \begin{solution}
        由AM-GM不等式,有
        \[
        a^2+b^2+c^2+d^2+ab+ac+ad+bc+bd+cd \ge \sqrt[4]{(abcd)^2} + \sqrt[6]{(abcd)^3}= 4+6=10
        \]
        等号成立当且仅当$a=b=c=d=1$。
    \end{solution}

    \question 设 $a,b,c$ 为正实数,且满足 $a^2+b^2+c^2 = 1$,证明
    \[
    \frac{1}{a^2} + \frac{1}{b^2} + \frac{1}{c^2} \ge 3 + \frac{2(a^3+b^3+c^3)}{abc}.
    \]
    \begin{solution}
        考虑左式减右式,有
        \begin{align*}
        &\frac{1}{a^2} + \frac{1}{b^2} + \frac{1}{c^2} - 3 - \frac{2(a^3+b^3+c^3)}{abc} \\
        &= \frac{a^2+b^2+c^2}{a^2} + \frac{a^2+b^2+c^2}{b^2} + \frac{a^2+b^2+c^2}{c^2} - 3 - 2 \left(\frac{a^2}{bc} + \frac{b^2}{ca} + \frac{c^2}{ab}\right) \\
        &= \frac{b^2+c^2}{a^2} + \frac{a^2+c^2}{b^2} + \frac{a^2+b^2}{c^2} - 2 \left(\frac{a^2}{bc} + \frac{b^2}{ca} + \frac{c^2}{ab}\right) \\
        &= a^2\left(\frac{1}{b} - \frac{1}{c}\right)^2 + b^2\left(\frac{1}{c} - \frac{1}{a}\right)^2 + c^2\left(\frac{1}{a} - \frac{1}{b}\right)^2 \ge 0
        \end{align*}
        等号成立当且仅当 $a = b = c$。
    \end{solution}

    \question 若 $a,b,c$ 为实数且 $a+b+c=12$,且
    \[
    \frac{1}{a}+\frac{1}{b}+\frac{1}{c}+\frac{1}{abc}=1,
    \]  
    求 $abc-(a+2b-3c)$ 的最大值。
    \begin{solution}
        由
        \[
        \frac{1}{a}+\frac{1}{b}+\frac{1}{c}+\frac{1}{abc}=1,
        \] 
        则
        \[
        2abc=2ab+2bc+2ca+2=144-a^2-b^2-c^2+2
        \]
        所以
        \begin{align*}
        2abc - 2(a+2b-3c) 
        &= 146 - (a^2+2a)-(b^2+4b)-(c^2-6c) \\
        &= 160 - [(a+1)^2+(b+2)^2+(c-3)^2] \\
        &\le 160-3\left(\frac{a+1+b+2+c-3}{3}\right)^2=112        
        \end{align*}
        因此
        \[
        abc-(a+2b-3c) \le 56
        \]  
        等号成立当且仅当 $a=3, b=2, c=7$。
    \end{solution}

    \question 设 $a,b,c$ 为正实数,且 $a+b+c=1$,求
    \[
    \sqrt{a^2+b^2}+\sqrt{b^2+c^2}+\sqrt{c^2+a^2}
    \]
    的最小值。
    \begin{solution}
        由柯西不等式,
        \[
        (1^2+1^2)(a^2+b^2)\ge (a+b)^2 \Rightarrow \sqrt{a^2+b^2}\ge \frac{a+b}{\sqrt{2}},
        \]
        同理
        \[
        \sqrt{b^2+c^2}\ge \frac{b+c}{\sqrt{2}},\quad \sqrt{c^2+a^2}\ge \frac{c+a}{\sqrt{2}}.
        \]
        相加得
        \[
        \sqrt{a^2+b^2}+\sqrt{b^2+c^2}+\sqrt{c^2+a^2}
        \ge \frac{2(a+b+c)}{\sqrt{2}}=\sqrt{2}.
        \]
        当 $a=b=c=\dfrac13$ 时取等号,故最小值为 $\sqrt{2}$。
    \end{solution}

    \question 已知 $x>1,y>1,z>1$ 且 $\dfrac{1}{x}+\dfrac{1}{y}+\dfrac{1}{z}=2$, 试证
    \[
    \sqrt{x+y+z}\ge \sqrt{x-1}+\sqrt{y-1}+\sqrt{z-1}
    \]
    \begin{solution}
        首先有
        \[
        \frac{x-1}{x}+\frac{y-1}{y}+\frac{z-1}{z}
        =3-\left(\frac{1}{x}+\frac{1}{y}+\frac{1}{z}\right)=3-2=1
        \]
        由柯西不等式,
        \[
        \left(\frac{x-1}{x}+\frac{y-1}{y}+\frac{z-1}{z}\right)(x+y+z)
        \ge\left(\sqrt{x-1}+\sqrt{y-1}+\sqrt{z-1}\right)^2
        \]
        左式为 $1\cdot(x+y+z)=x+y+z$,因此
        \[
        \sqrt{x+y+z}\ge \sqrt{x-1}+\sqrt{y-1}+\sqrt{z-1}
        \]
        故证毕。
    \end{solution}

    \question 
    \begin{parts}
    \part 已知 $a,b,c>0$,证明
    \[
    \frac{a^{2}}{2a+b}+\frac{b^{2}}{2b+c}+\frac{c^{2}}{2c+a}\ge\frac{a+b+c}{3}
    \]
    \begin{solution}
        由柯西不等式,
        \[
        \left( \frac{a^2}{2a+b} + \frac{b^2}{2b+c} + \frac{c^2}{2c+a} \right) (2a+b+2b+c+2c+a) \ge (a+b+c)^2
        \]
        所以得证
        \[
        \frac{a^2}{2a+b} + \frac{b^2}{2b+c} + \frac{c^2}{2c+a} \ge \frac{a+b+c}{3}
        \]
    \end{solution}
    \part 已知 $a_{1},a_{2},b_{1},b_{2},c_{1},c_{2}>0$,证明
    \[
    (a_{1}^{3}+a_{2}^{3})(b_{1}^{3}+b_{2}^{3})(c_{1}^{3}+c_{2}^{3})\ge(a_{1}b_{1}c_{1}+a_{2}b_{2}c_{2})^{3}
    \]
    \begin{solution}
        令
        \[
        r_1 = \frac{a_2}{a_1}, \quad r_2 = \frac{b_2}{b_1}, \quad r_3 = \frac{c_2}{c_1}
        \]
        则
        \[
        (a_1^3+a_2^3)(b_1^3+b_2^3)(c_1^3+c_2^3) = a_1^3 b_1^3 c_1^3 (1+r_1^3)(1+r_2^3)(1+r_3^3)
        \]
        \[
        (a_1 b_1 c_1 + a_2 b_2 c_2)^3 = a_1^3 b_1^3 c_1^3 (1+r_1 r_2 r_3)^3
        \]
        由AM-GM不等式,
        \begin{align*}
        (1+r_1^3)(1+r_2^3)(1+r_3^3) 
        &= 1 + (r_1^3+r_2^3+r_3^3) + (r_1 r_2)^3+(r_2 r_3)^3+(r_3 r_1)^3 + r_1^3 r_2^3 r_3^3 \\
        &\ge 1 + 3 r_1 r_2 r_3 + 3 r_1^2 r_2^2 r_3^2 + r_1^3 r_2^3 r_3^3  \\
        &= (1+r_1 r_2 r_3)^3
        \end{align*}
        因此得证
        \[
        (a_1^3+a_2^3)(b_1^3+b_2^3)(c_1^3+c_2^3) \ge (a_1 b_1 c_1 + a_2 b_2 c_2)^3
        \]
    \end{solution}
    \end{parts}

    \question 设 \(a, b, c\) 均为正数,且 \(a+b+c=3\),证明
    \[
    \frac{a}{b^{2}+1}+\frac{b}{c^{2}+1}+\frac{c}{a^{2}+1} \ge \frac{3}{2}.
    \]
    \begin{solution}
        由不等式
        \[
        (b-1)^2 = b^2 + 1 - 2b \ge 0 \Rightarrow \frac{1}{b^2 + 1} \le \frac{1}{2b},
        \]
        因此
        \[
        \frac{a}{b^2 + 1} = a - \frac{ab^2}{b^2 + 1} \ge a - \frac{ab^2}{2b} = a - \frac{ab}{2}.
        \]
        同理得
        \[
        \frac{b}{c^2 + 1} \ge b - \frac{bc}{2}, \quad \frac{c}{a^2 + 1} \ge c - \frac{ac}{2}.
        \]
        又由柯西不等式,
        \[
        (a^2 + b^2 + c^2)(1^2 + 1^2 + 1^2) \ge (a + b + c)^2 \Rightarrow a^2 + b^2 + c^2 \ge 3.
        \]
        因为
        \[
        ab + bc + ca = \frac{(a + b + c)^2 - (a^2 + b^2 + c^2)}{2} \le \frac{9 - 3}{2} = 3
        \]
        故
        \[
        \frac{a}{b^2 + 1} + \frac{b}{c^2 + 1} + \frac{c}{a^2 + 1} \ge (a + b + c) - \frac{1}{2}(ab + bc + ca) \ge 3 - \frac{3}{2} = \frac{3}{2}.
        \]
        故证毕。
    \end{solution}

    \question 设 $a,b,c$ 为正实数,证明
    \[
    \sqrt{ab(a+b)} + \sqrt{bc(b+c)} + \sqrt{ca(c+a)} \le \frac{3}{2}\sqrt{(a+b)(b+c)(c+a)}.
    \]
    \begin{solution}
        由AM-GM不等式,
        \begin{align*}
        3&=\frac{b}{b+c} + \frac{c}{c+a} + \frac{a}{a+b} + \frac{c}{b+c} + \frac{b}{a+b} + \frac{a}{c+a} \\
        &\ge 2\left(\sqrt{\frac{ab}{(b+c)(c+a)}} + \sqrt{\frac{bc}{(a+b)(c+a)}} + \sqrt{\frac{ca}{(a+b)(b+c)}}\right)
        \end{align*}
        即得证
        \[
        \sqrt{ab(a+b)} + \sqrt{bc(b+c)} + \sqrt{ca(c+a)} \le \frac{3}{2}\sqrt{(a+b)(b+c)(c+a)}.
        \]
    \end{solution}

    \question 证明对任意实数$x, y, z$皆有 $$ \frac{y^{2}-x^{2}}{2x^{2}+1}+\frac{z^{2}-y^{2}}{2y^{2}+1}+\frac{x^{2}-z^{2}}{2z^{2}+1}\ge0 $$
    \begin{solution}
        由 AM-GM 不等式,
        \begin{align*}
            \frac{y^{2}-x^{2}}{2x^{2}+1}+\frac{z^{2}-y^{2}}{2y^{2}+1}+\frac{x^{2}-z^{2}}{2z^{2}+1}+\frac{3}{2}
            &=\frac{1}{2}\left(\frac{2y^{2}+1}{2x^{2}+1}+\frac{2z^{2}+1}{2y^{2}+1}+\frac{2x^{2}+1}{2z^{2}+1}\right) \\
            &\ge\frac{3}{2}\sqrt[3]{\frac{2y^{2}+1}{2x^{2}+1}\cdot \frac{2z^{2}+1}{2y^{2}+1}\cdot\frac{2x^{2}+1}{2z^{2}+1}}  =\frac{3}{2}
        \end{align*}
        故
        $$ \frac{y^{2}-x^{2}}{2x^{2}+1}+\frac{z^{2}-y^{2}}{2y^{2}+1}+\frac{x^{2}-z^{2}}{2z^{2}+1}\ge0 $$
    \end{solution}
    
    \question 设 $a>0, b>0, c>0$,求
    \[
    \frac{a+3c}{a+2b+c} + \frac{4b}{a+b+2c} - \frac{8c}{a+b+3c} + 17
    \]
    的最小值。
    \begin{solution}
        由于
        \[
        \begin{cases}
        a + 2b + c = x\\ a + b + 2c = y\\ a + b + 3c = z
        \end{cases}
        \Rightarrow
        \begin{cases}
        a = -x + 5y - 3z \\ b = x - 2y + z \\ c = -y + z
        \end{cases}
        \]
        故
        \begin{align*}
        &\frac{a+3c}{a+2b+c} + \frac{4b}{a+b+2c} - \frac{8c}{a+b+3c} + 17 \\[1mm]
        &= \frac{-x + 2y}{x} + \frac{4x - 8y + 4z}{y} - \frac{-8y + 8z}{z} + 17 \\
        &= \left(\frac{2y}{x} + \frac{4x}{y}\right) + \left(\frac{4z}{y} + \frac{8y}{z}\right) \\[1mm]
        &\geq 2\sqrt{\frac{2y}{x} \cdot \frac{4x}{y}} + 2\sqrt{\frac{4z}{y} \cdot \frac{8y}{z}} = 12\sqrt{2}
        \end{align*}
    \end{solution}

    \question 证明对所有自然数 $n$,有
    \[
    1 + \frac{2}{3n-2} \le \sqrt[n]{3} \le 1 + \frac{2}{n}.
    \]
    \begin{solution}
        由AM-GM不等式,
        \[
        \frac{n+2}{n} = \frac{3 + \overbrace{1 + \cdots + 1}^{n-1 \text{个1}}}{n} \ge \sqrt[n]{3\cdot 1^{n-1}} = \sqrt[n]{3} \Rightarrow \sqrt[n]{3} \le 1 + \frac{2}{n}
        \]
        且
        \[
        \frac{3n-2}{3n} = \frac{1/3 + \overbrace{1 + \cdots + 1}^{n-1 \text{个1}}}{n} \ge \sqrt[n]{(1/3)\cdot 1^{n-1}} = \frac{1}{\sqrt[n]{3}} \Rightarrow \sqrt[n]{3} \ge  1 + \frac{2}{3n-2}.
        \]
        综上得到
        \[
        1 + \frac{2}{3n-2} \le \sqrt[n]{3} \le 1 + \frac{2}{n}
        \]
    \end{solution}

    \question 已知 $0 \le x \le 1$, $n \in \mathbb{N}$,证明:
    \[
    0 \le x^n - x^{n+1} \le \frac{n^n}{(n+1)^{\,n+1}}.
    \]
    \begin{solution}
        由 $0 \le x \le 1$,有
        \[
        0 \le 1-x \quad \Rightarrow \quad 0 \le x^n(1-x) = x^n - x^{\,n+1} \tag{1}
        \]
        由AM-GM不等式,
        \[
        \frac{\overbrace{\frac{x}{n} + \cdots + \frac{x}{n}}^{n \text{ 个}} + (1-x)}{\,n+1} \ge \sqrt[n+1]{\left(\frac{x}{n}\right)^n (1-x)} 
        \]
        即
        \[
        \frac{n^n}{(n+1)^{\,n+1}} \ge x^n(1-x) = x^n - x^{\,n+1} \tag{2}
        \]
        由 (1) 与 (2), 得证
        \[
        0 \le x^n - x^{\,n+1} \le \frac{n^n}{(n+1)^{\,n+1}}.
        \]
    \end{solution}

    \question 已知 $\triangle ABC$ 的面积为 $27,R$ 为其外接圆半径,$r$ 为其内切圆半径。求 $Rr^{3}$ 的最大值。
    \begin{solution}
        设$\triangle ABC$边长为$a,b,c>0$,面积为
        \[
        [\triangle ABC]=\frac{1}{2}(a+b+c)r=\frac{abc}{4R}=27
        \]
        由AM-GM不等式,
        \[
        \frac{2\cdot27}{3r}=\frac{a+b+c}{3}\ge \sqrt[3]{abc} = \sqrt[3]{4R\cdot 27} \Rightarrow Rr^3 \le 54
        \]
        等号成立当且仅当$a=b=c$即$\triangle ABC$为等边三角形。
    \end{solution}

    \question 设实数 $x,y$ 满足 $\sin x + \cos y = 1$,求 $\cos x + \sin y$ 的最大值。
    \begin{solution}
        注意到
        \[
        (\sin x + \cos y)^2 + (\cos x + \sin y)^2 = 2 + 2\sin(x + y) \le 4
        \]
        由题设 $\sin x + \cos y = 1$,代入上式
        \[
        1^2 + (\cos x + \sin y)^2 \le 4 
        \Rightarrow \cos x + \sin y \le \sqrt{3}
        \]
        当 $x = \dfrac{\pi}{6},\; y = \dfrac{\pi}{3}$ 时,$\;\sin x + \cos y = \dfrac{1}{2} + \dfrac{1}{2} = 1,\cos x + \sin y = \dfrac{\sqrt{3}}{2} + \dfrac{\sqrt{3}}{2} = \sqrt{3}$ 取到最大值。
    \end{solution}

    \question 若 $x \in \mathbb{R} \setminus \left\{\dfrac{\pi}{2} + k\pi, k\pi \mid k \in \mathbb{Z} \right\}$,试求 
    \[
    f(x) = |\sin x + \cos x + \tan x + \cot x + \sec x + \csc x|
    \]
    的最小值。
    \begin{solution}
        设 $t = \cos x + \sin x $,有
        \[
        \sin x\cos x = \frac{t^2 - 1}{2}
        \]
        且
        \[
        \tan x + \cot x = \frac{\sin^2 x+\cos^2 x}{\sin x \cos x} = \frac{2}{t^2 - 1} 
        \]
        \[
        \sec x + \csc x = \frac{\sin x + \cos x}{\cos x \sin x} = \frac{2t}{t^2 - 1}
        \]
        故
        \[
        f(x) = g(t) = \left| t + \frac{2}{t^2 - 1} + \frac{2t}{t^2 - 1} \right| = \left| t - 1 + \frac{2}{t - 1} + 1 \right| \ge \left| -2\sqrt{2} + 1 \right| = 2\sqrt{2} - 1
        \]
    \end{solution}

    \question 已知 $x^2 + y^2 = 1$,求 $\dfrac{y-3}{x+2}$ 的最大值。
    \begin{solution}
        设 $k = \dfrac{y-3}{x+2}$,则直线方程为 
        \[
        kx - y + 2k + 3 = 0
        \]
        当直线与圆 $x^2 + y^2 = 1$ 相切时,$k$取得最大值或最小值,此时圆心 $(0,0)$ 到直线的距离等于半径:
        \[ 
        \frac{|2k+3|}{\sqrt{k^2+1}} = 1 \Rightarrow 3k^2 + 12k + 8 = 0 
        \]
        解得最大值为 
        \[
        k_{max} = \frac{-12 + \sqrt{144 - 96}}{6} =  \frac{-6 + 2\sqrt{3}}{3}
        \]
    \end{solution}

    \question 已知实数 $x, y$ 满足 $3x - 4y = 5$,求 $x^2 + y^2$ 的最小值。
    \begin{solution}
        $x^2 + y^2$ 的最小值即为原点 $(0,0)$ 到直线 $3x - 4y - 5 = 0$ 距离的平方:
        \[ 
        d = \frac{|-5|}{\sqrt{3^2 + 4^2}} = 1 
        \]
        所以最小值为$1^2=1$。
    \end{solution}

    \question 已知实数 $x, y$ 满足 $12x + 5y = 1$,求 $x^2 + y^2 + 4x - 6y + 20$ 的最小值。
    \begin{solution}
        原式可配方为 $(x+2)^2 + (y-3)^2 + 7$,其几何意义为点 $(-2,3)$ 到直线 $12x + 5y - 1 = 0$ 距离的平方加 $7$,该距离为
        \[ 
        d = \frac{|12(-2) + 5(3) - 1|}{\sqrt{12^2 + 5^2}} = \frac{10}{13} 
        \]
        最小值即 $d^2 + 7 = \frac{1283}{169}$。
    \end{solution}

    \question 已知实数 $x, y$ 满足 $x^2 + y^2 + 2x - 2y = 7$,求 $\sqrt{(x-5)^2 + (y-6)^2}$ 的最小值。
    \begin{solution}
        圆方程配方为 $(x+1)^2 + (y-1)^2 = 3^2$,圆心 $A(-1,1)$,半径 $r=3$。设点 $B(5,6)$,则 $\sqrt{(x-5)^2 + (y-6)^2}$ 的最小值即
        \[ 
        AB - r= \sqrt{(5+1)^2 + (6-1)^2} - 3 = \sqrt{61} - 3 
        \]
    \end{solution}

    \question 求 $\sqrt{x^2+6x+73} + \sqrt{x^2-4x+8}$ 的最小值。
    \begin{solution}
        原式可化为 
        \[
        \sqrt{(x+3)^2 + (0-8)^2} + \sqrt{(x-2)^2 + (0-2)^2},
        \]
        即为 $x$ 轴上的点 $(x,0)$ 到两点 $A(-3,8)$ 和 $B(2,2)$ 距离之和,找 $B$ 关于 $x$ 轴的对称点 $B'(2,-2)$,最小值为线段 $AB'$ 的长度
        \[ 
        AB' = \sqrt{(-3-2)^2 + (8+2)^2} = 5\sqrt{5} 
        \]
    \end{solution}

    \question 求$\dfrac{3-\sin x}{2+\cos x}$ 的最大值及最小值。
    \begin{solution}
        设 $y = \dfrac{3-\sin x}{2+\cos x}$,则 
        \[
        3 - 2y = y\cos x + \sin x = \sqrt{y^2+1} \cos(x-\alpha)
        \]
        其中$\alpha$为锐角,由 $|\cos(x-\alpha)| \le 1$得 
        \[
        \left| \frac{3-2y}{\sqrt{y^2+1}} \right| \le 1 \Rightarrow 3y^2 - 12y + 8 \le 0 
        \]
        解方程 $3y^2 - 12y + 8 = 0$ 得 
        \[
        y = 2 \pm \frac{2\sqrt{3}}{3}
        \]
        所以最大值为 $2 + \dfrac{2\sqrt{3}}{3}$,最小值为 $2 - \dfrac{2\sqrt{3}}{3}$。
    \end{solution}

    \question 求代数式 $\sqrt{x^2-12x+40} + \sqrt{x^2+y^2} + \sqrt{y^2-8y+20}$ 的最小值。
    \begin{solution}
        将各式进行配方,转化为两点间距离公式:
        \[
        \sqrt{(x-6)^2 + (0 + 2)^2} + \sqrt{(x-0)^2 + (0-y)^2} + \sqrt{(0 + 2)^2 + (y-4)^2}
        \]
        构造直角坐标系,考虑坐标内四点
        \[
        A = (6, -2),\quad B = (x, 0),\quad C = (0, y),\quad D = (-2, 4)
        \]
        则根式之和即$AB+BC+CD$。当 $A, B, C, D$ 四点共线时,长度之和取得最小值,最小距离即为点 $A = (6, -2)$ 与点 $D = (-2, 4)$ 之间的距离:
        \[
        \sqrt{(6 - (-2))^2 + (-2 - 4)^2} = 10
        \]
        因此,原式的最小值为 10。
    \end{solution}

    \question 已知 $a, b$ 为正实数,求 $\sqrt{49+a^2-7\sqrt{2}a} + \sqrt{a^2+b^2-\sqrt{2}ab} + \sqrt{50+b^2-10b}$ 的最小值。
    \begin{solution}
        由余弦定理,观察到
        \begin{align*}
        \sqrt{a^2 + 49 - 7\sqrt{2}a} &= \sqrt{a^2 + 7^2 - 2 \cdot a \cdot 7 \cos 45^\circ}\\
        \sqrt{a^2 + b^2 - \sqrt{2}ab} &= \sqrt{a^2 + b^2 - 2 \cdot a \cdot b \cos 45^\circ} \\
        \sqrt{b^2 + 50 - 10b} &= \sqrt{b^2 + (5\sqrt{2})^2 - 2 \cdot b \cdot 5\sqrt{2} \cos 45^\circ}
        \end{align*}
        在平面上构造三角形$\triangle OAB,\triangle OBC,\triangle OCD$,使得
        \[
        OA=7,\quad OB= a,\quad OC=b,\quad OD=5\sqrt{2},\quad \angle AOB = \angle BOC = \angle COD = 45^\circ
        \]
        故原式即$AB+BC+CD$的长度,欲使其为最小,调整正实数$a,b$使得$A,B,C,D$四点共线,在$\triangle OAD$中,由余弦定理,最小值为
        \[
        AD = \sqrt{7^2 + 50 - 2 \cdot 7 \cdot 5\sqrt{2} \cos 135^\circ} =13
        \]
    \end{solution}

    \question 试求 
    \[
    \sqrt{10-6\cos\theta}+\frac{1}{4}\sqrt{34-24\sqrt{2}\sin\theta}+\sqrt{19-2\sqrt{2}\cos\theta-8\sin\theta}
    \] 
    的最小值。
    \ifprintanswers
    \begin{figure}[H]
        \centering
        \includegraphics[width=0.5\linewidth]{images/image42.png}
    \end{figure}
    \fi
    \begin{solution}
        发现
        \[
        \sqrt{10 - 6\cos\theta} = \sqrt{(\cos\theta - 3)^2 + \sin^2\theta}
        \]
        \[
        \frac{1}{4}\sqrt{34 - 24\sqrt{2} \sin\theta} = \sqrt{\cos^2\theta - \left(\sin\theta - \frac{3}{4}\sqrt{2} \right)^2}
        \]
        \[
        \sqrt{19 - 2\sqrt{2} \cos\theta - 8\sin\theta} = \sqrt{(\cos\theta - \sqrt{2})^2 + (\sin\theta - 4)^2}
        \]
        构造坐标系,设 $P(\cos\theta, \sin\theta)$ 在单位圆上,$A(3,0),\ B(0,\frac{3\sqrt{2}}{4}),\ C(\sqrt{2},4)$,原式即为 
        \[
        PA + PB + PC
        \]
        直线 $AB$方程式为
        \[
        \sqrt{2}x + 4y = 3\sqrt{2},
        \]
        故圆心 $O$ 到 $AB$ 的距离为 $1$,且当 $P = \left(\dfrac{1}{3}, \dfrac{2\sqrt{2}}{3} \right),OPC$ 共线,于是$PA + PB + PC$最小值为
        \[
        AB + OC - 1
        = \sqrt{9 + \frac{18}{16}} + \sqrt{18} - 1 = \frac{21\sqrt{2}}{4}  - 1
        \]
    \end{solution}

    \question 若锐角 $A,B,C$ 满足 $\sin^2 A + \sin^2 B + \sin^2 C = 2$, 求 
    \[
    \frac{1}{\sin^2 A\cos^4 B} + \frac{1}{\sin^2 B \cos^4 C} + \frac{1}{\sin^2 C \cos^4 A}
    \]
    的最小值。
    \begin{solution}
        由柯西不等式,
        \[
        \begin{aligned}
        &\frac{1}{\sin^2 A \cos^4 B} + \frac{1}{\sin^2 B \cos^4 C} + \frac{1}{\sin^2 C \cos^4 A} \\
        &= \frac{1}{2} \left( \frac{1}{\sin^2 A \cos^4 B} + \frac{1}{\sin^2 B \cos^4 C} + \frac{1}{\sin^2 C \cos^4 A} \right) (\sin^2 A + \sin^2 B + \sin^2 C) \\
        &\ge \frac{1}{2} \left( \frac{1}{\cos^2 B} + \frac{1}{\cos^2 C} + \frac{1}{\cos^2 A} \right)^2 \\
        &= \frac{1}{2} \left[ \left( \frac{1}{\cos^2 A} + \frac{1}{\cos^2 B} + \frac{1}{\cos^2 C} \right) (\cos^2 A + \cos^2 B + \cos^2 C) \right]^2 \ge \frac{81}{2}
        \end{aligned}
        \]
        当 $\sin^2 A = \sin^2 B = \sin^2 C = \dfrac{2}{3},\cos^2 A = \cos^2 B = \cos^2 C = \dfrac{1}{3}$ 时等号成立,
        \[
        \frac{1}{\sin^2 A \cos^4 B} + \frac{1}{\sin^2 B \cos^4 C} + \frac{1}{\sin^2 C \cos^4 A} = \frac{27}{2} \cdot 3 = \frac{81}{2}.
        \]
    \end{solution}

    \question 设$\theta_1, \theta_2, \theta_3, \dots, \theta_{2025}$皆为锐角,且
    \[
    \sin^2 \theta_1 + \sin^2 \theta_2 + \cdots + \sin^2 \theta_{2025} = 1,
    \]
    求
    \[
    \frac{\sin \theta_1 + \sin \theta_2 + \cdots + \sin \theta_{2025}}{\cos \theta_1 + \cos \theta_2 + \cdots + \cos \theta_{2025}}
    \]
    的最大值。
    \begin{solution}
        由题意,
        \[
        \sin^2 \theta_2 + \cdots + \sin^2 \theta_{2025} = 1-\sin^2 \theta_1  = \cos^2 \theta_1
        \]
        由柯西不等式,
        \[
        \cos^2 \theta_1 = (\sin^2 \theta_2 + \cdots + \sin^2 \theta_{2025})(1^2 + 1^2 + \cdots + 1^2)
        \ge (\sin \theta_2 + \sin \theta_3 + \cdots + \sin \theta_{2025})^2
        \]
        即
        \[
        \cos \theta_1 \ge \frac{\sin \theta_2 + \cdots + \sin \theta_{2025}}{\sqrt{2024}}
        \]
        同理,
        \[
        \cos \theta_2 \ge \frac{\sin \theta_1 + \sin \theta_3 + \cdots + \sin \theta_{2025}}{\sqrt{2024}}, \ \cdots \ , \cos \theta_{2025} \ge \frac{\sin \theta_1 + \cdots + \sin \theta_{2024}}{\sqrt{2024}}
        \]
        将以上不等式相加得
        \[
        \cos \theta_1 + \cdots + \cos \theta_{2025} \ge \frac{2024}{\sqrt{2024}} (\sin \theta_1 + \cdots + \sin \theta_{2025})
        \]
        故
        \[
        \frac{\sin \theta_1 + \cdots + \sin \theta_{2025}}{\cos \theta_1 + \cdots + \cos \theta_{2025}} \le \frac{\sqrt{2024}}{2024}
        = \frac{\sqrt{506}}{1012}
        \]
    \end{solution}

    \question 设 $a, b, c$ 为三角形的三边长,$\Delta$ 为此三角形面积,试证:
    \[
    \Delta \le \frac{\sqrt{3}}{4} \left( \frac{a + b + c}{3} \right)^2
    \]
    且等号成立当且仅当 $a = b = c$。
    \begin{solution}
        由海伦公式,三角形面积为
        \[
        \Delta = \sqrt{s(s - a)(s - b)(s - c)}
        \]
        其中$s$为半周长,由AM-GM不等式,
        \[
        \frac{s}{3}=\frac{(s - a) + (s - b) + (s - c)}{3} \ge \sqrt[3]{(s - a)(s - b)(s - c)}
        \]
        故
        \[
        \Delta \le \sqrt{ s\cdot\left(  \frac{s}{3} \right)^3 } = \frac{s^2}{3\sqrt{3}} = \frac{(a + b + c)^2}{12\sqrt{3}} = \frac{\sqrt{3}}{4} \left( \frac{a + b + c}{3} \right)^2
        \]
        等号成立条件为 $s - a = s - b = s - c$,即 $a = b = c$ 时,该三角形为正三角形,面积为 $\dfrac{\sqrt{3}}{4} a^2$
    \end{solution}

    \question 设 $a, b, c$ 是一个三角形的三边,周长为 $2$。证明
    \[
    \frac{3}{2} < a^2+b^2+c^2 + 2abc < 2.
    \]
    \begin{solution}
        由于周长为 $2$,有$a,b,c <1$,则三角形面积
        \[
        A = \frac{1}{2} bc \sin \alpha < \frac{1}{2}
        \]
        由海伦公式,
        \[
        A^2 = (1-a)(1-b)(1-c),
        \]
        因此展开
        \[
        0 < (1-a)(1-b)(1-c) < \frac{1}{4}
        \]
        可得
        \[
        1 < (ab+ac+bc) - abc < \frac{5}{4}
        \]
        且由
        \[
        a^2+b^2+c^2 = (a+b+c)^2 - 2(ab+ac+bc),
        \]
        因此
        \[
        4 - \frac{5}{2} < a^2+b^2+c^2 + 2abc < 4 - 2 \Rightarrow \frac{3}{2} < a^2+b^2+c^2 + 2abc < 2
        \]
    \end{solution}
        
    \question 在 $\triangle ABC$ 中,角 $A, B, C$ 所对的边为 $a, b, c$,且
    \[
    a^2 + b^2 + 2c^2 = 8
    \]
    求 $\triangle ABC$ 面积的最大值。
    \begin{solution}
        由余弦定理
        \[
        \cos C = \frac{a^2 + b^2 - c^2}{2ab}= \frac{8 - 3c^2}{2ab}
        \]
        又
        \[
        S = \frac{1}{2}ab\sin C = \frac{1}{2}ab\sqrt{1 - \left(\frac{8-3c^2}{2ab}\right)^2} = \frac{1}{4} \sqrt{4a^2b^2 - (8-3c^2)^2}
        \]
        由不等式$a^2+b^2\ge2ab,$
        \[
        S \leq \frac{1}{4} \sqrt{(a^2+b^2)^2 - (8-3c^2)^2}= \frac{1}{4} \sqrt{(8-2c^2)^2 - (8-3c^2)^2} = \frac{1}{4} \sqrt{16c^2 - 5c^4}
        \]
        对于$c^2$的函数 $f(c^2) = 16c^2 - 5c^4$,令导数为零得 $c^2 = \dfrac{8}{5}$,代入得
        \[
        S_{\text{max}} = \frac{1}{4} \sqrt{\frac{64}{5}} = \frac{2\sqrt{5}}{5}
        \]
    \end{solution}
    \begin{solution}
        以 $AB$ 所在直线为 $x$ 轴,$AB$ 中垂线为 $y$ 轴,设 $A(-\dfrac{c}{2},0),B(\dfrac{c}{2},0),C(x,y)$,由条件
        \[
        a^2 + b^2 + 2c^2 = 8
        \]
        得
        \[
        (x - \frac{c}{2})^2 + y^2 + (x + \frac{c}{2})^2 + y^2 + 2c^2 = 8
        \]
        即
        \[
        x^2 + y^2 = 4 - \frac{5c^2}{4}
        \]
        $\triangle ABC$面积为
        \[
        S = \frac{1}{2}c|y|
        \]
        由 $y^2 \leq 4 - \dfrac{5c^2}{4}$,得
        \[
        S \leq \frac{1}{2}c\sqrt{4 - \frac{5c^2}{4}} = \frac{1}{4}\sqrt{16c^2 - 5c^4}
        \]
        解法同上。
    \end{solution}
    
    \question  在 $\triangle ABC$ 中,试求$$\frac{2\sin^2 A + \sin B \sin C}{\sin A \sin B \sin C}$$ 的最小值。
    \begin{solution}
        由余弦定理,
        \[
        \cos A = \dfrac{b^2 + c^2 - a^2}{2bc}\ge 1 - \dfrac{a^2}{2bc}\Rightarrow \frac{2a^2}{bc} \ge 4 - 4\cos A
        \]
        由正弦定理,
        \[
        \frac{2\sin^2 A + \sin B \sin C}{\sin A \sin B \sin C}
        = \frac{1}{\sin A} \left( \frac{2a^2}{bc} + 1 \right)
        \ge \frac{5 - 4\cos A}{\sin A}
        \]
        设 $\tan \dfrac{A}{2} = t > 0$,则
        \[
        \cos A = \frac{1 - t^2}{1 + t^2},\quad \sin A = \frac{2t}{1 + t^2}
        \]
        代入得
        \[
        \frac{5 - 4\cos A}{\sin A}
        = \frac{5 - 4\cdot \frac{1 - t^2}{1 + t^2}}{\frac{2t}{1 + t^2}}
        = \frac{9t^2 + 1}{2t}
        \]
        由AM-GM不等式:
        \[
        \frac{9t^2 + 1}{2t} \ge \frac{2\sqrt{9t^2 \cdot 1}}{2t} = 3
        \]
        当 $A=\arcsin\dfrac{3}{5}$, $B=C=\dfrac{\pi-A}{2}$ 时等号成立, 原式取最小值 $3$
    \end{solution}
    
    \question 在锐角三角形 $\triangle ABC$ 中,角 $A,B,C$ 所对的边为 $a,b,c$,且
        \[
        a^2 + 2ab\cos C = 3b^2
        \]
        求 $\tan A \tan B \tan C$ 的最小值。
    \begin{solution}
        由余弦定理,有
        \[
        2ab\cos C = a^2+b^2-c^2
        \]
        代入原式得
        \[
        a^2 + (a^2+b^2-c^2) = 3b^2 \Rightarrow c^2 = 2a^2 - 2b^2 \Rightarrow b^2 + c^2 - a^2 = \frac{c^2}{2} \Rightarrow \frac{b^2 + c^2 - a^2}{2bc} = \frac{c^2}{4bc}
        \]
        由正、余弦定理得
        \[
        \cos A = \frac{\sin C}{4\sin B} \Rightarrow 4\sin B \cos A  =\sin C = \sin(A+B) = \sin A \cos B + \cos A \sin B
        \]
        有 $3\sin B \cos A = \sin A \cos B$.
        
        $\because \triangle ABC$ 是锐角三角形, $\cos A\neq0,\cos B\neq0, \therefore \tan A=3\tan B $
        
        现设 $\tan B = t$, 又
        \[
        \tan C = -\tan(A+B) = -\frac{\tan A + \tan B}{1-\tan A\tan B} = \frac{4t}{3t^2-1}
        \]
        所以
        \[
        \tan A \tan B \tan C = 3t \cdot t \cdot \frac{4t}{3t^2-1} = \frac{12t^3}{3t^2-1}
        \]
        对函数
        \[
        f(t) = \frac{12t^3}{3t^2-1}, t>0
        \]
        求导
        \[
        f'(t) = \frac{36t^2(3t^2-1) - 12t^3 \cdot 6t}{(3t^2-1)^2} = \frac{36t^2(t^2-1)}{(3t^2-1)^2}
        \]
        令 $f'(t)=0$,解得 $t=1$; 当 $t\in (0,1),f'(t)<0,$当 $t\in (1,\infty),f'(t)>0,$

        $\therefore\; f(t)$在$t=1$有最小值 $f(1)=6$。
    \end{solution}
    \begin{solution}
        以 $AB$ 所在直线为 $x$ 轴,$\;AB$ 中垂线为 $y$ 轴, 设 $A(-\dfrac{c}{2},0),B(\dfrac{c}{2},0),C(x,y)$,       
        由距离公式
        \[
        a^2 = (x-\frac{c}{2})^2 + y^2, \quad b^2 = (x+\frac{c}{2})^2 + y^2
        \]
        代入已知条件
        \[
        2b^2 + c^2 - 2a^2 = 0 \implies 2\left[(x+\frac{c}{2})^2 + y^2\right] + c^2 - 2\left[(x-\frac{c}{2})^2 + y^2\right] = 0
        \]
        展开化简得
        \[
        4c x + c^2 = 0 \implies x = -\frac{c}{4}
        \]
        作 $CD \perp AB$ 于 $D$,则 $D$ 点坐标为 $(x,0) = \left(-\dfrac{c}{4},0\right)$。发现
        \[
        AD = \left|-\frac{c}{2} - \left(-\frac{c}{4}\right)\right| = \frac{c}{4} \quad \text{且} \quad BD = \left|\frac{c}{2} - \left(-\frac{c}{4}\right)\right| = \frac{3c}{4}
        \]
        所以 $BD = 3AD$,又
        \[
        \tan A = \frac{CD}{AD} = \frac{y}{\frac{c}{4}}
        \]
        \[
        \tan B = \frac{CD}{BD} = \frac{y}{\frac{3c}{4}}
        \]
        因此有
        \[
        \tan A = 3 \tan B
        \]
        解法同上。
    \end{solution}
    
    \question 已知非负实数 $a, b, c,d$ 满足 $a \le 1, a+b \le 5, a+b+c \le 14, a+b+c+d \le 30$,试证
    \[
    \sqrt{a}+\sqrt{b}+\sqrt{c}+\sqrt{d}\le10
    \]
    \begin{solution}
        由柯西不等式,
        \begin{align*}
        &(\sqrt{a}+\sqrt{b}+\sqrt{c}+\sqrt{d})^2\\
        &=\left(\sqrt{a}+\frac{\sqrt{b}}{2}+\frac{\sqrt{b}}{2}+\frac{\sqrt{c}}{3}+\frac{\sqrt{c}}{3}+\frac{\sqrt{c}}{3}+\frac{\sqrt{d}}{4}+\frac{\sqrt{d}}{4}+\frac{\sqrt{d}}{4}+\frac{\sqrt{d}}{4}\right)^{2} \\
        &\le\left(a+\frac{b}{2}+\frac{b}{2}+\frac{c}{9}+\frac{c}{9}+\frac{c}{9}+\frac{d}{16}+\frac{d}{16}+\frac{d}{16}+\frac{d}{16}\right)(1+...+1)=10(a+\frac{b}{2}+\frac{c}{3}+\frac{d}{4})
        \end{align*}
        而
        \begin{align*}
        a+\frac{b}{2}+\frac{c}{3}+\frac{d}{4}&=\frac{1}{4}(a+b+c+d)+\frac{1}{12}(a+b+c)+\frac{1}{6}(a+b)+\frac{1}{2}a\\
        &\le\frac{30}{4}+\frac{14}{12}+\frac{5}{6}+\frac{1}{2}=10
        \end{align*}
        故
        $$\sqrt{a}+\sqrt{b}+\sqrt{c}+\sqrt{d}\le10$$
    \end{solution}

    \question 已知 $a,b,c$ 为正实数,且满足 $abc=1$,试证明:
    \[
    \frac{1}{a^{3}(b+c)} + \frac{1}{b^{3}(a+c)} + \frac{1}{c^{3}(a+b)} \ge \frac{3}{2}
    \]
    \begin{solution}
        由柯西不等式,
        \[
        \left(\frac{1}{a^{3}(b+c)} + \frac{1}{b^{3}(a+c)} + \frac{1}{c^{3}(a+b)}\right) (a(b+c) + b(a+c) + c(a+b)) \ge \left(\frac{1}{a} + \frac{1}{b} + \frac{1}{c}\right)^2
        \]
        注意到
        \[
        a(b+c) + b(a+c) + c(a+b) = 2(ab + bc + ca),
        \]
        且
        \[
        \left(\frac{1}{a} + \frac{1}{b} + \frac{1}{c}\right)^2 = \left(\frac{ab + bc + ca}{abc}\right)^2 = (ab + bc + ca)^2
        \]
        故由AM-GM不等式,
        \[
        \frac{1}{a^{3}(b+c)} + \frac{1}{b^{3}(a+c)} + \frac{1}{c^{3}(a+b)} \ge \frac{1}{2}(ab + bc + ca)
        \ge \frac{3}{2}\sqrt[3]{(abc)^2} = \frac{3}{2}
        \]
    \end{solution}
    \begin{solution}
        由 $abc=1$,有
        \[
        \sum_{cyc} \frac{1}{a^3(b+c)} = \sum_{cyc} \frac{bc}{a^2(b+c)} = \sum_{cyc} \frac{(\frac{1}{a})^2}{\frac{1}{b}+\frac{1}{c}}
        \]
        由 AM–GM 不等式,可推出
        \[
        \frac{x^2}{y} \ge x - \frac{y}{4}, \quad x,y>0,
        \]
        据此得到
        \[
        \sum_{cyc} \frac{(\frac{1}{a})^2}{\frac{1}{b}+\frac{1}{c}} \ge \sum_{cyc} \left[ \frac{1}{a} - \frac{1}{4}\left( \frac{1}{b} + \frac{1}{c} \right) \right] = \frac{1}{2} \left( \frac{1}{a} + \frac{1}{b} + \frac{1}{c} \right)
        \]
        再由 AM–GM 不等式,
        \[
        \frac{1}{2} \left( \frac{1}{a} + \frac{1}{b} + \frac{1}{c} \right) \ge \frac{3}{2} \sqrt[3]{\frac{1}{abc}} = \frac{3}{2}
        \]
        因此原不等式成立。
    \end{solution}

    \question 已知 $0\le x,y,z\le1$ 且
    \[
    xyz=(1-x)(1-y)(1-z),
    \]
    证明
    \[
    x(1-z)+y(1-x)+z(1-y)\ge\frac{3}{4}.
    \]
    \begin{solution}
        由于 $0\le x\le1$,有
        \[
        (x-\frac12)^2 \ge 0 \Rightarrow 0\le x(1-x)\le\frac14,
        \]
        对 $y,z$ 同理成立。于是
        \[
        xyz(1-x)(1-y)(1-z)\le\frac{1}{64}
        \]
        由已知关系 $xyz=(1-x)(1-y)(1-z)$得
        \[
        (xyz)^2\le\frac{1}{64} \Rightarrow xyz\le\frac{1}{8}
        \]
        于是
        \[
        x(1-z)+y(1-x)+z(1-y)= x+y+z - (xy+yz+zx)= 1 - 2xyz \ge 1-2\cdot\frac{1}{8}=\frac{3}{4}
        \]
        当 $x=y=z=\dfrac12$ 时等号成立,
    \end{solution}

    \question 设正实数 $a,b,c,x,y,z$ 满足 $a+b+c = x+y+z$,证明
    \[
    \frac{2a^2}{a+x} + \frac{2b^2}{b+y} + \frac{2c^2}{c+z} \ge a+b+c
    \]
    \begin{solution}
    由柯西不等式,
        \[
        \left(\frac{2a^2}{a+x} + \frac{2b^2}{b+y} + \frac{2c^2}{c+z}\right)\left(2(a+x) +2(b+y) +2(c+z) \right) \ge (2a+2b+2c)^2 
        \]
        又 $a+b+c = x+y+z$,得
        \[
        \frac{2a^2}{a+x} + \frac{2b^2}{b+y} + \frac{2c^2}{c+z} \ge a+b+c.
        \]
    \end{solution}

    \question 已知 $a,b,c,d>0$,证明不等式
    \[
    \frac{a^{3}+b^{3}+c^{3}}{a^{2}+b^{2}+c^{2}}+\frac{b^{3}+c^{3}+d^{3}}{b^{2}+c^{2}+d^{2}}+\frac{c^{3}+d^{3}+a^{3}}{c^{2}+d^{2}+a^{2}}+\frac{d^{3}+a^{3}+b^{3}}{d^{2}+a^{2}+b^{2}}\ge a+b+c+d.
    \]
    \begin{solution}
        由 $a,b>0$,有
        \[
        0\le (a+b)(a-b)^2= (a^2-b^2)(a-b) = a^3 - a^2 b - b^2 a + b^3 
        \]
        可知
        \[
        a^3 + b^3 \ge a^2 b + a b^2, \quad b^3 + c^3 \ge b^2 c + b c^2, \quad c^3 + a^3 \ge c^2 a + c a^2.
        \]
        于是
        \begin{align*}
        3(a^3 + b^3 + c^3) &\ge a^3 + b^3 + c^3+ a^2 b + a b^2 + b^2 c + b c^2 + c^2 a + c a^2\\
        &=(a+b+c)(a^2 + b^2 + c^2)
        \end{align*}
        因此
        \[
        \frac{a^3 + b^3 + c^3}{a^2 + b^2 + c^2} \ge \frac{a+b+c}{3}
        \]
        故
        \begin{align*}
        &\frac{a^{3}+b^{3}+c^{3}}{a^{2}+b^{2}+c^{2}}+\frac{b^{3}+c^{3}+d^{3}}{b^{2}+c^{2}+d^{2}}+\frac{c^{3}+d^{3}+a^{3}}{c^{2}+d^{2}+a^{2}}+\frac{d^{3}+a^{3}+b^{3}}{d^{2}+a^{2}+b^{2}} \\[1mm]
        &\ge \frac{1}{3} \cdot 3(a+b+c+d) = a+b+c+d 
        \end{align*}
    \end{solution}

    \question 已知正实数 $a, b, c, d$,求
    \[
    \frac{ab + bc + cd}{a^2 + b^2 + c^2 + d^2}
    \]
    的最大值。
    \begin{solution}
        对任意正数 $t$,由 AM-GM 不等式,
        \[
        ab \le \frac{t}{2} a^2 + \frac{1}{2t} b^2, \quad 
        bc \le \frac{1}{2} b^2 + \frac{1}{2} c^2, \quad
        cd \le \frac{1}{2t} c^2 + \frac{t}{2} d^2.
        \]
        取 $t = \frac{1}{2}(1+\sqrt{5})$,使得 
        \[
        \frac{t}{2} = \frac{1}{2t} + \frac{1}{2}
        \]
        则
        \[
        ab + bc + cd \le \frac{t}{2} (a^2 + b^2 + c^2 + d^2).
        \]
        因此最大值为
        \[
        \frac{t}{2} = \frac{\sqrt{5}+1}{4}
        \]
        当 $ta = b = c = td$ 时等号成立。
    \end{solution}

    \question 已知正数 $a,b,c$ 满足 $abc > 1$,求 $\dfrac{abc(a + b + c + 8)}{abc - 1}$ 的最小值。
    \begin{solution}
        设 $\lambda = \dfrac{abc(a + b + c + 8)}{abc - 1} > 0$,整理得
        \[
        a + b + c + \dfrac{\lambda}{abc} = \lambda - 8
        \]
        由AM-GM不等式得
        \[
        a + b + c + \dfrac{\lambda}{abc} \ge 4\sqrt[4]{abc \cdot \dfrac{\lambda}{abc}} = 4\sqrt[4]{\lambda}
        \]
        于是
        \[
        \lambda - 8 \ge 4\sqrt[4]{\lambda}
        \]
        令 $k = \sqrt[4]{\lambda} > 0$,则不等式转化为
        \[
        k^4 - 4k - 8 \ge 0 \Rightarrow (k - 2)(k^3 + 2k^2 + 4k + 4) \ge 0
        \]
        由于 $k > 0,k^3 + 2k^2 + 4k + 4>0$ ,故$k\ge 2$,即  $$\lambda \ge 2^4 = 16$$
        等号成立当且仅当 $a = b = c = 2$
    \end{solution}

    \question 已知 $x^2+y^2+z^2 \leq 1$, 求 $x^2+2y-2z+3$ 的取值范围。
    \begin{solution}
        由柯西不等式, $$[y+(-z)]^2 \leq (1+1)[y^2+(-z)^2] \leq 2(x^2+y^2+z^2) \leq 2$$
        则 $y-z \geq -\sqrt{2}$, 且 $x^2 \geq 0$, 可得 $$x^2+2y-2z+3 \geq 0+2(-\sqrt{2})+3=3-2\sqrt{2}$$
        当且仅当 $x=0, y=-\dfrac{\sqrt{2}}{2}, z=\dfrac{\sqrt{2}}{2}$ 时,等号成立;
        
        又因为 $x^2+y^2+z^2 \leq 1$, 则 $x^2 \leq 1-y^2-z^2$,
        可得 $$x^2+2y-2z+3 \leq 4-(y^2+z^2-2y+2z)$$
        且 $$y^2 + z^2 - 2y + 2z = (y-1)^2 + (z+1)^2 - 2,$$
        设点 $A(-1,1)$ 和标准单位圆上一点 $P(y,z)$, 则 $$(y-1)^2 + (z+1)^2 - 2 = |PA|^2 - 2,$$
        又因为 $|PA|^2 \geq (|OA|-1)^2 = 3-2\sqrt{2}$, 可得 $(y-1)^2 + (z+1)^2 - 2 \geq 1-2\sqrt{2},$
        则 $$x^2+2y-2z+3 \leq 4-(y^2+z^2-2y+2z) \leq 3+2\sqrt{2},$$
        当且仅当 $x=0, y=\dfrac{\sqrt{2}}{2}, z=-\dfrac{\sqrt{2}}{2}$ 时,等号成立。
        
        $\therefore$取值范围是 $[3-2\sqrt{2}, 3+2\sqrt{2}].$
    \end{solution}  

    \question 已知 $x,y,z>0$ ,求 $$\frac{\sqrt{x^{2}+y^{2}}+\sqrt{y^{2}+4z^{2}}+\sqrt{z^{2}+16x^{2}}}{9x+3y+5z}$$ 的最小值。
    \begin{solution}
        \textcolor{red}{(待解)}
    \end{solution}

    \question 证明
    \[
    \sum_{n=1}^\infty \frac{1}{\sqrt{n}(n+1)} < 2.
    \]
    \begin{solution}
        首先证明对任意正整数 $n$,都有如下不等式成立
        \[
        \frac{1}{\sqrt{n}(n+1)} < \frac{2}{\sqrt{n}} - \frac{2}{\sqrt{n+1}} \tag{1}
        \]
        将不等式 (1) 两边同乘以 $\sqrt{n}(n+1)$, 可得其等价形式
        \[
        1 < 2(n+1) - 2\sqrt{n(n+1)} \Longleftrightarrow 2\sqrt{n(n+1)} < n + (n+1)
        \]
        由AM-GM不等式,显然上述等价形式成立,故       
        \[
        \sum_{n=1}^\infty \frac{1}{\sqrt{n}(n+1)} < \sum_{n=1}^\infty \left( \frac{2}{\sqrt{n}} - \frac{2}{\sqrt{n+1}} \right)
        \]
        右端为裂项级数,其前 $k$ 项部分和为
        \[
        \sum_{n=1}^k\left( \frac{2}{\sqrt{n}} - \frac{2}{\sqrt{n+1}} \right) = 2 \left( 1 - \frac{1}{\sqrt{k+1}} \right)
        \]
        当 $k \to \infty$ 时, $S_k \to 2$,因此
        \[
        \sum_{n=1}^\infty \frac{1}{\sqrt{n}(n+1)} < 2
        \]
    \end{solution}

    \question 设实数 $a,b,c \in [-1,1]$ 满足
    \[
    1 + 2abc \ge a^2 + b^2 + c^2.
    \]
    证明对所有正整数 $n$ 有
    \[
    1 + 2(abc)^n \ge a^{2n} + b^{2n} + c^{2n}.
    \]
    \begin{solution}
        条件可改写为
        \[
        (a - bc)^2 \le (1-b^2)(1-c^2). \tag{1}
        \]
        由柯西不等式,
        \begin{align*}
        (a^{n-1} + a^{n-2} bc + \cdots + b^{n-1} c^{n-1})^2 
        &\le (|a|^{n-1} + |a|^{n-2}|bc| + \cdots + |bc|^{n-1})^2 \\
        &\le (1 + |bc| + \cdots + |bc|^{n-1})^2 \\
        &\le (1 + |b|^2 + \cdots + |b|^{2(n-1)}) (1 + |c|^2 + \cdots + |c|^{2(n-1)})
        \end{align*}
        乘以不等式 (1),得到
        \begin{align*}
        (a - bc)^2 (a^{n-1} + a^{n-2} bc + \cdots + b^{n-1} c^{n-1})^2 
        &\le ((1-b^2)(1 + |b|^2 + \cdots + |b|^{2(n-1)})) \\
        &\quad \cdot ((1-c^2)(1 + |c|^2 + \cdots + |c|^{2(n-1)})) 
        \end{align*}
        即
        \[
        (a^n - b^n c^n)^2 \le (1 - b^{2n})(1 - c^{2n})
        \]
        从而得到
        \[
        1 + 2(abc)^n \ge a^{2n} + b^{2n} + c^{2n}
        \]
    \end{solution}

    \question 设 $a_1,\dots,a_n$ 为正实数,且设 $s=a_1+\cdots+a_n$,试证明
    \[
    (1+a_1)(1+a_2)\cdots(1+a_n) \le 1+\frac{s}{1!}+\frac{s^2}{2!}+\cdots+\frac{s^n}{n!}.
    \]
    \begin{solution}
        由AM-GM不等式,
        \[
        \prod_{j=1}^n(1+a_j)\le \left(\frac{1}{n}\sum_{j=1}^n(1+a_j)\right)^n=\left(1+\frac{s}{n}\right)^n
        \]
        二项展开得
        \[
        \Big(1+\frac{s}{n}\Big)^n=\sum_{j=0}^n\binom{n}{j}\frac{s^j}{n^j}
        =\sum_{j=0}^n\frac{n!}{(n-j)!\,j!}\,\frac{s^j}{n^j}.
        \]
        注意当 $0\le j\le n$ 时,
        \[
        \frac{n!}{(n-j)!\,n^j}=\prod_{k=0}^{j-1}\Big(1-\frac{k}{n}\Big)\le 1,
        \]
        因此每一项满足
        \[
        \frac{n!}{(n-j)!\,j!}\,\frac{s^j}{n^j}\le\frac{s^j}{j!}.
        \]
        代入二项和得到
        \[
        \prod_{j=1}^n(1+a_j)\le\sum_{j=0}^n\frac{s^j}{j!},
        \]
        即
        \[
        (1+a_1)\cdots(1+a_n)\le 1+\frac{s}{1!}+\frac{s^2}{2!}+\cdots+\frac{s^n}{n!},
        \]
        证毕。
    \end{solution}

    \question 设 $n$ 个正实数的最小值为 $m$, 最大值为 $M$, 记它们的算术平均为 $A$, 几何平均为 $G$。证明
    \[
    A - G \ge \frac{1}{n} \left( \sqrt{M} - \sqrt{m} \right)^2.
    \]
    \begin{solution}
        设 $n$ 个正实数为 $x_1, \dots, x_n$, 不妨令 $x_1 = m, x_n = M$。不等式等价于
        \begin{align*}
        & A - G \ge \frac{\left( \sqrt{M} - \sqrt{m} \right)^2}{n} \\
        \Longleftrightarrow & \sum_{j=1}^n x_j - n \left( \prod_{j=1}^n x_j \right)^{\frac{1}{n}} \ge M - 2\sqrt{Mm} + m \\
        \Longleftrightarrow & \sum_{j=2}^{n-1} x_j + 2\sqrt{Mm} \ge n \left( \prod_{j=1}^n x_j \right)^{\frac{1}{n}}
        \end{align*}
        对这 $n$ 个数
        \[
        \sqrt{Mm}, \; x_2, \; x_3, \; \dots, \; x_{n-1}, \; \sqrt{Mm}
        \]
        应用 AM-GM 不等式,有
        \[
        \frac{1}{n} \left( \sum_{j=2}^{n-1} x_j + 2\sqrt{Mm} \right) \ge \left( Mm \prod_{j=2}^{n-1} x_j \right)^{\frac{1}{n}} = \left( x_1 x_n \prod_{j=2}^{n-1} x_j \right)^{\frac{1}{n}} = \left( \prod_{j=1}^n x_j \right)^{\frac{1}{n}}
        \]
        即
        \[
        \sum_{j=2}^{n-1} x_j + 2\sqrt{Mm} \ge n \left( \prod_{j=1}^n x_j \right)^{\frac{1}{n}}
        \]
        故原不等式得证。
    \end{solution}

    \question 设 $x_1,\ldots,x_n\ge -1$,且
    \[
    \sum_{i=1}^n x_i^3=0,
    \]
    证明
    \[
    \sum_{i=1}^n x_i\le \frac{n}{3}.
    \]
    \begin{solution}
        当 $x\ge -1$ 时,有不等式
        \[
        0\le x^3-\frac34 x+\frac14=(x+1)\left(x-\frac12\right)^2
        \]
        将 $x_1,\ldots,x_n$ 代入上式并求和,得到
        \[
        0\le =\sum_{i=1}^n x_i^3-\frac34\sum_{i=1}^n x_i+\frac{n}{4}
        \]
        由已知条件 
        \[
        \sum_{i=1}^n x_i^3=0,
        \]
        可得
        \[
        0\le -\frac34\sum_{i=1}^n x_i+\frac{n}{4} \Rightarrow \sum_{i=1}^n x_i\le \frac{n}{3}
        \]
        当且仅当 $n=9k$,其中 $k$ 个 $x_i=-1$,其余 $8k$ 个 $x_i=\dfrac12$ 时取等号。
    \end{solution}

    \question 已知非负实数 $x_1, x_2, \ldots, x_5$ 满足 $x_1^2 + x_2^2 + \cdots + x_5^2 = 4$. 求下式的最大值:
    \[
    S = \sum_{i=1}^5 \frac{1}{x_i+1}\sum_{i=1}^5 x_i
    \]
    \begin{solution}
        当 $x_1 = x_2 = x_3 = x_4 = 1, x_5 = 0$ 时, $S = 12$,下面证明 $S \le 12$.
        注意到 
        \[
        \frac{x(x-1)^2}{x+1} \ge 0$$ 对 $x \ge 0$ 成立, 由此可得$$\sum_{i=1}^5 \frac{x_i(x_i-1)^2}{x_i+1} \ge 0
        \]
        展开可知
        \[
        \sum_{i=1}^5 \left(x_i^2 - 3x_i + 4 - \frac{4}{x_i+1}\right) \ge 0
        \]
        整理得
        \[
        \sum_{i=1}^5 \frac{1}{x_i+1} \le 6 - \frac{3}{4} \sum_{i=1}^5 x_i
        \]
        由上式及AM-GM不等式可知 
        \[
        S \le \frac{4}{3} \left(6 - \frac{3}{4}\sum_{i=1}^5 x_i\right) \left(\frac{3}{4}\sum_{i=1}^5 x_i\right) \le \frac{4}{3} \cdot 3^2 = 12
        \]
    \end{solution}
\end{questions}

\pagebreak

\begin{center}
  {\fontsize{30pt}{26pt}\selectfont
    \hypertarget{数列与级数}{数列与级数} \label{数列与级数}
  }
\end{center}
\separator
\vspace{1pt}

\begin{questions}
    \question 设 \(a, b, c\) 三数成等比数列,且满足 \(a + b + c = 9\) 及 \(a^2 + b^2 + c^2 = 189\),求 $b$。
    \begin{solution}
        关键:运用恒等式
        \[
        (a+b+c)^2=a^2+b^2+c^2+2(ab+bc+ac)
        \]
        又已知 $b^2=ac$,故
        \[
        81=189+2(b(9-b)+b^2) \Rightarrow b=6
        \]
    \end{solution}
    
    \question 已知 $a,b,c,d$ 成等差数列,且实数 $x,y,z,u$ 满足
    \[
    \begin{cases}
    a+b+c+d = 60 \\
    x+y+z+u = 12 \\
    az + bu + cx + dy = 168
    \end{cases}
    \]
    求 $ay + bx + cu + dz$。
    \begin{solution}
        已知$a,b,c,d$ 成等差数列,于是
        \begin{align*}
        ay + bx + cu + dz+168 
        &=ay + bx + cu + dz+(az + bu + cx + dy) \\
        &=(a+d)(y+z)+(b+c)(x+u)\\
        &=30(x+y+z+u)= 360
        \end{align*}
        故
        \[
        ay + bx + cu + dz=192
        \]
    \end{solution}
    
    \question 已知 $a_1=-1,a_{n+1}=2a_n-3$, 求 $a_{2017}$。
    \begin{solution}
        由递推式,
        \[
        a_{n+1}-3 = 2(a_n-3),
        \]
        所以 $a_1-3,a_2-3,\dots$ 成等比数列,公比为 $2$。
        \[
        a_{n+1}-3 = 2^n(a_1-3) = 2^n(-1-3) = -2^{n+2} \Rightarrow a_{2017} = -2^{2018}+3.
        \]
    \end{solution}
    \begin{solution}
        由
        \[
        a_{n+1}-a_n =2a_n-3-(2a_{n-1}-3)= 2(a_n-a_{n-1}),
        \]
        所以 $\{a_{n+1}-a_n\}$ 为等比数列,公比为 $2$,因此
        \begin{align*}
        a_n &= (a_2-a_1)+(a_3-a_2)+\cdots+(a_n-a_{n-1})+a_1 \\
        &= \frac{(a_2-a_1)(2^{n-1}-1)}{2-1}-1 \\ 
        &= (2a_1-3-a_1)(2^{n-1}-1) =-2^{n+1}+3
        \end{align*}
        所以
        \[
        a_{2017}=-2^{2018}+3
        \]
    \end{solution}

    \question 若整数 $m\ge 1$,函数 $f$ 满足
    \[
    f(m+1)=m(-1)^{m+1}-2f(m), \quad f(1)=f(2001),
    \]
    求 $f(1)+f(2)+\cdots+f(2000)$。
    \begin{solution}
        由递推式可得
        \[
        f(2)=1-2f(1),\quad f(3)=-2-2f(2),\quad f(4)=3-2f(3),\;\ldots,\quad f(2001)=2000-2f(2000),
        \]
        将 $f(2001)$ 替换为 $f(1)$ 并将所有式子相加,得
        \[
        \sum_{i=1}^{2000} f(i)=1-2+3-4+\cdots+1999-2000-2\sum_{i=1}^{2000} f(i)
        \]
        所以
        \[
        \sum_{i=1}^{2000} f(i)=\frac{1}{3}\left(\frac{2000}{2}-2000\right)=-\frac{1000}{3}
        \]
    \end{solution}

    \question 已知数列 $\{a_n\}$ 定义为
    \[
    a_{n+1} = a_n + \frac{2n}{n^4 + n^2 + 1}, \quad a_1 = 0, \quad n \in \mathbb{N},
    \]
    求 $a_{61}$ 的值。
    \begin{solution}
        由递推关系,对 $n = 1, 2, \dots, 60$ 累加得
        \[
        a_{61} = a_1 + \sum_{n=1}^{60} \frac{2n}{n^4 + n^2 + 1} = \sum_{n=1}^{60} \left(\frac{1}{n^2 - n + 1} - \frac{1}{n^2 + n + 1}\right)
        \]
        再由累加法,
        \[
        a_{61} = \frac{1}{1^2 - 1 + 1} - \frac{1}{60^2 + 60 + 1} = 1 - \frac{1}{3661} = \frac{3660}{3661}
        \]
    \end{solution}

    \question 数列 \(a_0, a_1, a_2, \ldots, a_n, \ldots\) 满足 \(a_0 = 3\),且
    \[
    (3 - a_{n-1})(6 + a_n) = 18,
    \]
    证明
    \[
    \sum_{k=1}^n \frac{1}{a_k} = \frac{1}{3}\left(2^{n+2}-n-4\right)
    \]
    \begin{solution}
        由递推关系
        \begin{equation}
            (3-a_{n+1})(a_n+6) = 18 \implies a_{n+1} = \frac{3a_n}{a_n+6}
        \end{equation}
        两边加3得到
        \begin{equation}
            a_{n+1}+3 = \frac{3a_n}{a_n+6}+3 = \frac{6(a_n+3)}{a_n+6}
        \end{equation}
        两式相比得
        \[
        \frac{a_{n+1}+3}{a_{n+1}} =  \frac{2(a_n+3)}{a_n}
        \]
        故数列 $\left\{\dfrac{a_n+3}{a_n}\right\}$ 是首项为2,公比为2的等比数列,解得
        \[
        \frac{a_n+3}{a_n} = 2^{n+1} \Rightarrow a_n = \frac{3}{2^{n+1}-1}
        \]
        因此
        \[
        \sum_{k=1}^n \frac{1}{a_k} = \sum_{k=1}^n \frac{2^{k+1}-1}{3} = \frac{1}{3}\left(2^{n+2}-n-4\right)
        \]
    \end{solution}
        
    \question 设一函数 $f$ 的定义域为所有正整数,如果 $f(1) = 101$,且对所有正整数 $n > 1$皆有  
    \[
    f(1) + f(2) + f(3) + \cdots + f(n) = n^2 f(n)
    \]  
    求 $f(100)$ 的值。
    \begin{solution}
        发现
        \[
        f(n) = \sum_{k=1}^n f(k)-\sum_{k=1}^{n-1} f(k)=n^2 f(n) - (n - 1)^2 f(n - 1)
        \]
        即
        \[
         f(n) = \frac{n - 1}{n + 1} f(n - 1)
        \]
        于是
        \[
        f(100) = \frac{99 \cdot 98 \cdots 1}{101 \cdot 100 \cdots 3} f(1) = \frac{2}{101 \cdot 100} \cdot 101 = \frac{1}{50}
        \]
    \end{solution}
    
    \question 已知数列$\{a_n\}$满足:$a_1 =1$, $a_{n+1} = a_n + \dfrac{1}{a_n}$, 则$\lfloor a_{2025}\rfloor=$
    \begin{solution}
        由题意可知$$a_{n+1}^2 = a_n^2 + \frac{1}{a_n^2} + 2$$所以有
        $$a_{2025}^2 = a_1^2 + 2 \times (2025-1) + \left(\frac{1}{a_1^2} + \frac{1}{a_2^2} + \dots + \frac{1}{a_{2024}^2}\right)$$又因为$a_1=1$, 且有
        $$\frac{1}{a_1^2} + \frac{1}{a_2^2} + \dots + \frac{1}{a_{2024}^2} < 47$$所以$$63=\sqrt{3969} < a_{2025} < \sqrt{4096}=64$$故$[a_{2025}]=63$.
    \end{solution}

    \question 设 $f:\mathbb{N} \to \mathbb{R}$满足 $f(1) = \frac{3}{2}$,且对所有$n \in \mathbb{N}$且满足$$\ f(n+1) = \left(1 + \frac{1}{n+1}\right) f(n) + \left(1 + \frac{n}{2}\right) f(1) + n^2 + 2n,$$ 求 $f(40)$。
    \begin{solution}
        将递推式改写为
        \[
        f(n+1) = \frac{n+2}{n+1} f(n) + g(n),
        \]
        其中 $\displaystyle g(n) = (n+2)\left(n + \frac{3}{4}\right)$,于是
        \begin{align*}
        f(40) &= \frac{41}{40} f(39) + g(39) \\
        &= \frac{41}{39} f(38) + \frac{41}{40} g(38) + g(39) \\
        &= \cdots \\
        &= \frac{41}{2} f(1) + \sum_{k=1}^{39} \frac{41}{k+2}\ g(k) \\
        &= \frac{41}{2} \cdot \frac{3}{2} + 41 \sum_{k=1}^{39} (k+\frac{3}{4}) \\
        &= \frac{123}{4} + \frac{39\cdot 40}{2}+39\cdot \frac{3}{4}\\
        &= 33210
        \end{align*}
    \end{solution}
    
    \question 已知数列 $\{a_{n}\}$ 满足$a_{1}=a_{2}=a_{3}=1$,令 $b_{n}=a_{n}+a_{n+1}+a_{n+2},\ n\in \mathbb{N}$,若 $\{b_{n}\}$ 是公比为3的等比数列,求 $a_{100}$ 的值。
    \begin{solution}
        由条件知 $$b_{n}=b_{1}\cdot3^{n-1}=3^{n}$$
        因此$$a_{n+3}-a_{n}=b_{n+1}-b_{n}=3^{n+1}-3^{n}=2\cdot3^{n}$$
        于是
        $$a_{100}=a_{1}+\sum_{k=1}^{33}(a_{3k+1}-a_{3k-2})=1+\sum_{k=1}^{33}2\cdot3^{3k-2}$$
        $$=1+6\cdot\frac{27^{33}-1}{27-1}=1+\frac{3}{13}(3^{99}-1)=\frac{3^{100}+10}{13}$$
    \end{solution}

    \question 已知递归数列满足
    \[
    a_1=\frac{1}{2},\quad a_{n+1}=\frac{a_n+3}{2a_n-4},
    \]
    求通项公式 $a_n$。
    \begin{solution}
        先求不动点,设 $x=\dfrac{x+3}{2x-4}$,得 $2x^2-5x-3=0$,所以不动点为 $x=-\dfrac12,\,3$,于是
        \[
        R_n=\frac{a_n+\tfrac12}{a_n-3}.
        \]
        且
        \[
        R_{n+1}=\frac{a_{n+1}+\tfrac12}{a_{n+1}-3}=-\frac{2}{5}R_n,
        \]
        即 $\{R_n\}$ 为公比$-\dfrac{2}{5}$的等比数列,由于$R_1=\frac{a_1+\frac12}{a_1-3}=-\frac{2}{5}$,故
        \[
        R_n=\left(-\frac{2}{5}\right)^n
        \]
        于是
        \[
        a_n=\frac{3R_n+\tfrac12}{R_n-1}=\frac{6\big(-\tfrac{2}{5}\big)^n+1}{2\Big(\big(-\tfrac{2}{5}\big)^n-1\Big)}
        =\frac{5^n+6(-2)^n}{2\big((-2)^n-5^n\big)}.
        \]
    \end{solution}

    \question 设实数数列 \(\{a_n\}\) 满足
    \[
    a_n a_{n+2} - a_{n+1}^2 - (n+1) a_n a_{n+1} = 0。
    \]
    已知 \(a_1 = 1, a_2 = 2018\),求
    \[
    \frac{a_{2018} \cdot a_{2016}}{a_{2017}^2}
    \]
    \begin{solution}
        由递推关系得
        \[
        \frac{a_{n+2}}{a_{n+1}} = \frac{a_{n+1}}{a_n} + (n+1)
        \]
        由$\dfrac{a_2}{a_1}=2018$,有
        \[
        \dfrac{a_{n+1}}{a_n} = 2018 + \sum_{k=2}^n k = 2017 + \frac{n(n+1)}{2}
        \]
        故
        \[
        \frac{a_{2018}a_{2016}}{a_{2017}^2} = 
        \frac{2017 \left(1 + \frac{2018}{2} \right)}{2017 \left(1 + \frac{2016}{2} \right)} = \dfrac{1010}{1009}
        \]
    \end{solution}

    \question 设数列 $a_1, a_2, \dots, a_n, \dots$ 满足 $a_1 = a_2 = 1, a_3 = 2$ ,且对任意自然数 $n$ ,都有 
    \[
    a_n  a_{n+1}  a_{n+2}  a_{n+3} = a_n + a_{n+1} + a_{n+2} + a_{n+3},
    \]
    且$a_n  a_{n+1}  a_{n+2} \neq 1$,求 $a_1 + a_2 + \dots + a_{100}$。
    \begin{solution}
        由 
        \[
        a_n a_{n+1} a_{n+2} a_{n+3} = a_n + a_{n+1} + a_{n+2} + a_{n+3}
        \]
        及 
        \[
        a_{n+1} a_{n+2} a_{n+3} a_{n+4} = a_{n+1} + a_{n+2} + a_{n+3} + a_{n+4}
        \]
        两式相减得 
        \[
        (a_{n+1} a_{n+2} a_{n+3} - 1)(a_{n+4} - a_n) = 0
        \]
        因为 $a_{n+1} a_{n+2} a_{n+3} \neq 1$ ,所以 $a_{n+4} = a_n$ ,数列$\{a_n\}$周期为 4,数列的前四项为 1, 1, 2, 4,于是
        \[ 
        a_1 + a_2 + \dots + a_{100} = 25 \cdot (1 + 1 + 2 + 4) = 200 
        \]
    \end{solution}

    \question 已知函数 
    \[
    f(x) = \frac{(x+1)^4 + (x-1)^4}{(x+1)^4 - (x-1)^4},\quad x \neq 0
    \]
    在数列 $\{a_n\}$ 中,$a_1 = 2,a_{n+1} = f(a_n)$ 。求数列 $\{a_n\}$ 的通项公式。
    \begin{solution}
        观察到
        \[ 
        \frac{a_{n+1}-1}{a_{n+1}+1} = \frac{f(a_n)-1}{f(a_n)+1} = \frac{(a_n-1)^4}{(a_n+1)^4} = \left( \frac{a_n-1}{a_n+1} \right)^4 
        \]
        令 $b_n = \dfrac{a_n - 1}{a_n + 1}$ ,则有 
        \[
        b_{n+1} = b_n^4
        \]
        已知 $b_1 = \dfrac{2-1}{2+1} = \dfrac{1}{3}$,
        由此可得 
        \[
        b_n = b_1^{4^{n-1}} = \left( \frac{1}{3} \right)^{4^{n-1}}
        \]
        由 
        \[
        b_n = \frac{a_n - 1}{a_n + 1} \Rightarrow a_n = \frac{1+b_n}{1-b_n},
        \]
        代入得
        \[ 
        a_n = \frac{3^{4^{n-1}} + 1}{3^{4^{n-1}} - 1} 
        \]
    \end{solution}

    \question 已知数列 $\{a_n\}$ 中,$a_1 = 2$,
    \[
    a_n = \frac{1+a_{n-1}}{1-a_{n-1}},
    \]
    求 $\{a_n\}$ 的通项公式。
    \begin{solution}
        通过计算数列的前几项寻找规律:
        \begin{align*}
        a_1 &= 2\\
        a_2 &= \frac{1+2}{1-2} = -3 \\
        a_3 &= \frac{1-3}{1-(-3)} = -\frac{1}{2}\\
        a_4 &= \frac{1-1/2}{1-(-1/2)} = \frac{1}{3}\\ 
        a_5 &= \frac{1+1/3}{1-1/3} = 2 = a_1 
        \end{align*}
        由此可知数列 $\{a_n\}$ 是以 4 为周期的周期数列,通项公式为
        \[ 
        a_n = 
        \begin{cases} 
        2 & n = 4k-3 \\ 
        -3 & n = 4k-2 \\[4pt]
        -\dfrac{1}{2} & n = 4k-1 \\[8pt]
        \dfrac{1}{3} & n = 4k 
        \end{cases} ,k \in \mathbb{N}
        \]
    \end{solution}

    \question 已知数列 $\{a_n\}$ 由递推关系定义如下:
    \[
    a_0 = 1, \quad a_{n+1} = \frac{a_n}{2(n+1)}, \quad n \ge 0,
    \]
    若函数 $f$ 定义为
    \[
    f(x) = \sum_{n=0}^{\infty} a_n x^n,
    \]
    求 $f(2)$ 的精确值。
    \begin{solution}
        根据递推关系,可以看出所有奇数项 $a_{2n+1} = 0$,偶数项为
        \[
        a_0 = 1, \quad a_2 = \frac{1}{2}, \quad a_4 = \frac{1}{8}, \quad a_6 = \frac{1}{48}, \dots
        \]
        满足
        \[
        a_{2n} = \frac{1}{2^n n!}
        \]
        因此函数 $f$ 可表示为
        \[
        f(x) = \sum_{n=0}^{\infty} a_n x^n = \sum_{n=0}^{\infty} a_{2n} x^{2n} = \sum_{n=0}^{\infty} \frac{1}{2^n n!} x^{2n} = \sum_{n=0}^{\infty} \frac{1}{n!} \left(\frac{x^2}{2}\right)^n = e^\frac{x^2}{2}
        \]
        所以
        \[
        f(2) = e^2
        \]
    \end{solution}

    \question 设 $\{b_n\}_{n=0}^\infty$ 为一列正实数,且 $b_0=1$,
    \[
    b_n=2+\sqrt{b_{n-1}}-2\sqrt{1+\sqrt{b_{n-1}}},
    \]
    计算
    \[
    \sum_{n=1}^\infty b_n 2^n.
    \]
    \begin{solution}
        令 $a_n=1+\sqrt{b_n}$,其中 $n\ge0$,则 $a_n>1,a_0=2$,并且
        \[
        a_n=1+\sqrt{1+a_{n-1}-2\sqrt{a_{n-1}}}=\sqrt{a_{n-1}}
        \]
        因此
        \[
        a_n=2^{2^{-n}}
        \]
        于是
        \begin{align*}
        \sum_{n=1}^N b_n 2^n
        &=\sum_{n=1}^N (a_n-1)^2 2^n \\
        &=\sum_{n=1}^N \bigl(a_n^2 2^n-a_n 2^{n+1}+2^n\bigr) \\
        &=\sum_{n=1}^N \bigl((a_{n-1}-1)^2 2^n-(a_n-1)^2 2^{n+1}\bigr) \\
        &=(a_0-1)^2 2-(a_N-1)^2 2^{N+1} \\
        &=2-2\frac{2^{2^{-N}}-1}{2^{-N}}
        \end{align*}
        令 $x=2^{-N}$,当 $N\to\infty$ 时有 $x\to0$,由洛必达法则,
        \[
        \sum_{n=1}^\infty b_n 2^n
        =\lim_{N\to\infty}\left(2-2\frac{2^{2^{-N}}-1}{2^{-N}}\right)
        =\lim_{x\to0}\left(2-2\frac{2^x-1}{x}\right)
        =2-2\ln2
        \]
    \end{solution}
    
    \question 一列 11 个正实数 $a_1, a_2, \dots, a_{11}$ 满足 $a_1 = 4$, $a_{11} = 1024$,并且对 $2 \le n \le 11$ 有
    \[
    a_n + a_{n-1} = \frac{5}{2} \sqrt{a_n a_{n-1}}.
    \]
    求满足条件的序列数量 $S$ 。
    \begin{solution}
        设 $a_n = x,a_{n-1} = y$,解
        \[
        x+y = \frac{5}{2}\sqrt{xy}
        \]
        得
        \[
        x = 4y \text{ 或 } \ x = \frac{y}{4}.
        \]
        考虑二叉树,每一步可以乘以 4 或除以 4。从 $a_1=4$ 到 $a_{11}=1024$ 共 10 步:
        \[
        4^m \cdot 4^{-(10-m)} = 4^4 \Rightarrow m =7
        \]
        即有 7 步乘以 4,3 步除以 4。可能序列数为选择 3 步除以 4 的位置数:
        \[
        S = \comb{10}{3} = 120
        \]
    \end{solution}

    \question 已知复数列 $\{z_{n}\}$ 满足
        $z_{1}=\dfrac{\sqrt{3}}{2},z_{n+1}=\overline{z_{n}}(1+z_{n}i)$,其中$n\in\mathbb{N}$,求 $z_{2021}$ 的值。
    \begin{solution}
        对 $n\in \mathbb{N}$ ,设 $z_{n}=a_{n}+b_{n}i,\ a_{n},b_{n}\in \mathbb{R}$ ,则
        \[
        a_{n+1}+b_{n+1}i=z_{n+1}=\overline{z_{n}}(1+z_{n}i)=\overline{z_{n}}+|z_{n}|^{2}i=a_{n}-b_{n}i+(a_{n}^{2}+b_{n}^{2})i
        \]
        因此 
        \[
        a_{n+1}=a_{n}=a_1=\frac{\sqrt{3}}{2}, b_{n+1}=a_{n}^{2}+b_{n}^{2}-b_{n}=b_{n}^{2}-b_{n}+\frac{3}{4}
        \]
        即 
        \[
        b_{n+1}-\frac{1}{2}=b_{n}^{2}-b_{n}+\frac{1}{4}=(b_{n}-\frac{1}{2})^{2}\]
        所以当 $n\ge2$,
        \[
        b_{n}=\frac{1}{2}+(b_{1}-\frac{1}{2})^{2^{n-1}}=\frac{1}{2}+(-\frac{1}{2})^{2^{n-1}}=\frac{1}{2}+\frac{1}{2^{2^{n-1}}}
        \]
        于是
        \[
        z_{2021}=a_{2021}+b_{2021}i=\frac{\sqrt{3}}{2}+(\frac{1}{2}+\frac{1}{2^{2^{2020}}})i
        \]
    \end{solution}

    \question 已知$a_0 = \dfrac{1}{2}$,
    \[
    a_n = \left(\frac{1 + a_{n-1}}{2}\right)^{\frac{1}{2}},\quad \forall n \in \mathbb{N},
    \]
    求 $\displaystyle \lim_{n \rightarrow \infty} 4^n \cdot (1 - a_n)$。
    \begin{solution}
        由
        \[
        \cos \theta = \left(\frac{1 + \cos 2\theta}{2}\right)^{\frac{1}{2}}
        \]
        可得 $a_0 = \cos \dfrac{\pi}{3} = \dfrac{1}{2}$,从而
        \[
        a_1 = \cos \frac{\pi}{3 \cdot 2} = \left(\frac{1 + \frac{1}{2}}{2}\right)^{\frac{1}{2}} = \frac{\sqrt{3}}{2},  a_2 = \cos \frac{\pi}{3 \cdot 2^2}, \cdots,  a_n = \cos \frac{\pi}{3 \cdot 2^n}
        \]
        由泰勒展开,
        \[
        a_n = \cos \frac{\pi}{3 \cdot 2^n} = 1 - \frac{\left(\frac{\pi}{3 \cdot 2^n}\right)^2}{2!} + \frac{\left(\frac{\pi}{3 \cdot 2^n}\right)^4}{4!} - \cdots
        \]
        因此
        \[
        \lim_{n \to \infty} 4^n \cdot (1 - a_n)
        = \lim_{n \to \infty} 4^n \cdot \left( \frac{\left(\frac{\pi}{3 \cdot 2^n}\right)^2}{2!} - \frac{\left(\frac{\pi}{3 \cdot 2^n}\right)^4}{4!} + \cdots \right)
        = \frac{\pi^2}{9} \cdot \frac{1}{2!}
        = \frac{\pi^2}{18}
        \]
    \end{solution}
    
    \question Let $a_1,a_2,\dots$ be a sequence of real numbers satisfying
$$\frac{a_{n+1}}{a_n}-\frac{a_{n+2}}{a_n}-\frac{a_{n+1}a_{n+2}}{a_n^2}=\frac{na_{n+2}a_{n+1}}{a_n}.$$Given that $a_1=-1$ and $a_2=-\tfrac{1}{2}$, find the value of $\tfrac{a_9}{a_{20}}$.
    \begin{solution}
        \textcolor{red}{(待解)}
    \end{solution}

    \question 数列 $1, 2, 2, 3, 3, 3, 4, 4, 4, 4, 5, \dots$ 的第 $100$ 项是?
    \begin{solution}
        一般地, 数字 $n$ 出现 $n$ 项,因此从数字 $1$ 到数字 $n$ 共出现
        \[
        1+2+\dots+n = \frac{n(n+1)}{2}
        \]
        项,由 $\dfrac{n(n+1)}{2} \leq 100$ 得 $n<14$,取 $n=13$ 时共有 $91$ 项,因此第 $100$ 项为 $14$。
    \end{solution}

    \question 等差数列 $1, 3, 5, 7, 9, 11,\dots$ 按如下方法分组
    \[
    (1),(3,5),(7,9,11),(13,15,17,19),\dots
    \]
    求第 $n$ 组中 $n$ 个数的和 $S_n$。
    \begin{solution}
        等差数列 $1,3,5,7,\dots$ 的通项公式为
        \[
        a_n = 2n-1.
        \]
        第 $n$ 组有 $n$ 个奇数, 因此第 $n$ 组的第 $n$ 个奇数是等差数列中的第
        \[
        1+2+\dots+n = \frac{n(n+1)}{2}
        \]
        个奇数,且第 $n$ 组的第$1$个奇数是等差数列中的第
        \[
        \frac{n(n+1)}{2}-(n-1)=\frac{n^2-n+2}{2}
        \]
        故第$n$组中第 $n$ 个奇数是
        \[
        2\cdot \frac{n(n+1)}{2}-1=n^2+n-1
        \]
        第$n$组中第 $1$ 个奇数是
        \[
        2\cdot \frac{n^2-n+2}{2}-1=n^2-n+1
        \]
        所以
        \[
        S_n=\frac{n[(n^2+n-1)-(n^2-n+1)]}{2}=n^3
        \]
    \end{solution}

    \question 已知数列 $\{a_n\}$, 且 $a_1=1$, 定义 $S_n=n^2 a_n$. 求 $a_n$ 和 $S_n$。
    \begin{solution}
        由
        \[
        a_n = S_n - S_{n-1} =n^2 a_n - (n-1)^2 a_{n-1} \Rightarrow (n^2-1)a_n = (n-1)^2 a_{n-1}
        \]
        于是
        \[
        \frac{a_n}{a_{n-1}} =\frac{n-1}{n+1}.
        \]
        因此
        \[
        a_n = \frac{a_n}{a_{n-1}} \cdot \frac{a_{n-1}}{a_{n-2}}\cdot\cdots\cdot \frac{2}{4}\cdot \frac{1}{3}\cdot 1 = \frac{2}{n(n+1)} 
        \]
        且
        \[
        S_n =n^2 a_n= \frac{2n}{n+1}.
        \]
    \end{solution}

    \question 已知等差数列 $\{a_n\}$ 公差 $d=2$, $a_1+a_2+\dots+a_{100}=100$, 求 $a_4+a_8+\dots+a_{100}$ 的值。
    \begin{solution}
        由
        \[
        a_1+a_2+\dots+a_{100}=100,
        \]
        解得
        \[
        100a_1 + \frac{100 \cdot 99}{2} \cdot 2 = 100\Rightarrow a_1=-98
        \]
        于是$a_4=-92,a_100=100$,
        \[
        a_4+a_8+\dots+a_{100} = \frac{25}{2}(-92+100)=100.
        \]
    \end{solution}

    \question 已知 $F_1 = 0, F_2 = 1$ ,求斐波那契数列 $F_n + F_{n+1} = F_{n+2}$ 的通项 $F_n$ 。
    \begin{solution}
        数列
        \[
        F_n + F_{n+1} = F_{n+2}
        \]
        的特征方程为 
        \[
        1 + \lambda = \lambda^2
        \]
        解得
        \[
        \lambda_1 = \frac{1+\sqrt{5}}{2}, \lambda_2 = \frac{1-\sqrt{5}}{2}
        \]
        于是 
        \[
        F_n = C_1 \left( \frac{1+\sqrt{5}}{2} \right)^{n-1} + C_2 \left( \frac{1-\sqrt{5}}{2} \right)^{n-1}
        \]
        代入初值 $F_1 = 0, F_2 = 1$ 解得
        \[
        C_1 = \frac{1}{\sqrt{5}},\quad C_2 = -\frac{1}{\sqrt{5}}
        \]
        因此 
        \[
        F_n = \frac{1}{\sqrt{5}} \left[ \left( \frac{1+\sqrt{5}}{2} \right)^{n-1} - \left( \frac{1-\sqrt{5}}{2} \right)^{n-1} \right]
        \]
    \end{solution}

    \question 已知数列 $a_1 = 1, a_2 = 5$ ,且 $a_{n+1} = 4a_n - 4a_{n-1}, n \ge 2$ ,求通项公式 $a_n$ 。
    \begin{solution}
        此数列对应特征方程为 
        \[
        \lambda^2 - 4\lambda + 4 = 0
        \]
        解得 $\lambda_1 = \lambda_2 = 2$,设此数列的通项公式为 
        \[
        a_n = (C_1 + nC_2) \cdot 2^n
        \]
        代入初值 $a_1 = 1, a_2 = 5$ 可知
        \[
        C_1 = -\frac{1}{4},\quad C_2 = \frac{3}{4}
        \]
        所以 
        \[
        a_n = (3n-1) \cdot 2^{n-2}
        \]
    \end{solution}

    \question 设数列 $\{a_n\}$ 满足 $a_0 = 1, a_1 = 4$ ,且 
    \[
    a_{n+2} - 14a_{n+1} + a_n + 6 = 0.
    \]
    求证$a_n$ 是完全平方数。
    \begin{solution}
        将递推式化为 
        \[
        \left(a_{n+2} - \frac{1}{2}\right) - 14\left(a_{n+1} - \frac{1}{2}\right) + \left(a_n - \frac{1}{2}\right) = 0
        \]
        令 $b_n = a_n - \dfrac{1}{2}$ ,且 $b_0 = \dfrac{1}{2}, b_1 = \dfrac{7}{2}$ ,则 
        \[
        b_{n+2} - 14b_{n+1} + b_n = 0
        \]
        由特征方程 $\lambda^2 - 14\lambda + 1 = 0$ 得两根 
        \[
        \lambda_1 = 7+4\sqrt{3}, \lambda_2 = 7-4\sqrt{3}
        \]
        故设
        \[
        b_n = C_1 \lambda_1^n + C_2 \lambda_2^n ,
        \]
        由 $n=0, 1$ 时条件解得 $C_1 = C_2 = \frac{1}{4}$,因此 
        \[
        b_n = \frac{1}{4}(7+4\sqrt{3})^n + \frac{1}{4}(7-4\sqrt{3})^n = \frac{1}{4}(2+\sqrt{3})^{2n} + \frac{1}{4}(2-\sqrt{3})^{2n}
        \]
        则 
        \[
        a_n = b_n + \frac{1}{2} = \frac{1}{4}(2+\sqrt{3})^{2n} + \frac{1}{4}(2-\sqrt{3})^{2n} + \frac{1}{2} = \left[ \frac{(2+\sqrt{3})^n + (2-\sqrt{3})^n}{2} \right]^2
        \]
        因为 $(2+\sqrt{3})^n + (2-\sqrt{3})^n$ 为偶数,所以 $a_n$ 是完全平方数,故得证。
    \end{solution}

    \question 设
    \[
    D_1 = 4,\quad D_2 = \begin{vmatrix} 4 & 2 \\ 2 & 4 \end{vmatrix},\quad D_3 = \begin{vmatrix} 4 & 2 & 0 \\ 2 & 4 & 2 \\ 0 & 2 & 4 \end{vmatrix}, \dots 
    \]
    求 $D_{10}$ 。
    \begin{solution}
        观察到
        \[
        D_n = 4D_{n-1} - 4D_{n-2}
        \]
        特征方程为
        \[
        \lambda^2 = 4\lambda - 4 
        \]
        解得$\lambda_1 = \lambda_2 = 2$,设 
        \[
        D_n = (C_1 + nC_2) \cdot 2^{n-1}
        \]
        由 $D_1 = 4, D_2 = 12$, 解得 $C_1 = 2, C_2 = 2$,因此 
        \[
        D_n = (2 + 2n) \cdot 2^{n-1} = (n+1) \cdot 2^n
        \]
        当 $n=10$ 时,
        \[
        D_{10} = 11 \cdot 1024 = 11264
        \]
    \end{solution}

    \question 已知对任意 $x$, 
    \[
    f(x+1) + f(x-1) = \sqrt{2} f(x),
    \]
    且 $f(1) = 2, f(2) = 1$,求 $f(2010)$ 。
    \begin{solution}
        将 $f(x+1) + f(x-1) = \sqrt{2} f(x)$ 化为 
        \[
        a_{n+2} + a_n = \sqrt{2} a_{n+1}
        \]
        且 $a_1 = 2, a_2 = 1$,特征方程为 
        \[
        \lambda^2 - \sqrt{2}\lambda + 1 = 0 
        \]
        解得
        \[
        \lambda = \frac{\sqrt{2}}{2} \pm i \frac{\sqrt{2}}{2} = \cos \frac{\pi}{4} \pm i \sin \frac{\pi}{4}
        \]
        则 
        \[
        a_n = C_1 \cos \frac{n\pi}{4} + C_2 \sin \frac{n\pi}{4}
        \]
        代入 $a_1=2, a_2=1$ 得 $C_1 = 2\sqrt{2}-1, C_2 = 1$,于是
        \[
        a_{2010} = (2\sqrt{2}-1) \cos \frac{2010\pi}{4} + \sin \frac{2010\pi}{4} = 1
        \]
    \end{solution}

    \question 一实数列 $\{u_1, u_2, \dots, u_n\}$ 定义为
    \[
    u_1 = \sqrt{2},\quad u_2 = \pi,\quad u_n = u_{n-1} - u_{n-2} \quad \forall n \ge 3
    \]
    求 $u_n, u_{2008}$ 及 $u_{2016}$。
    \begin{solution}
        由递推关系 $u_n - u_{n-1} + u_{n-2} = 0$,其特征方程为
        \[
        x^2 - x + 1 = 0
        \]
        求得特征根为
        \[
        x = \frac{1 \pm \sqrt{3}i}{2} = \cos \frac{\pi}{3} \pm i \sin \frac{\pi}{3}
        \]
        因此数列的通项公式形式为
        \[
        u_n = C_1 \cos \frac{n\pi}{3} + C_2 \sin \frac{n\pi}{3}
        \]
        代入初始条件 $u_1 = \sqrt{2}, u_2 = \pi$,解得
        \[
        C_1 = \sqrt{2} - \pi,\quad C_2 = \frac{\sqrt{2} + \pi}{\sqrt{3}}
        \]
        故通项公式为
        \[
        u_n = (\sqrt{2} - \pi) \cos \frac{n\pi}{3} + \frac{\sqrt{2} + \pi}{\sqrt{3}} \sin \frac{n\pi}{3}
        \]
        令 $n=2016$,可得
        \[
        u_{2016} = \sqrt{2} - \pi
        \]
    \end{solution}

    \question 已知数列 $\{a_n\}$ 中,$a_1 = 1$, 
    \[
    S_n = \frac{S_{n-1}}{2S_{n-1} + 1},
    \]
    求 $\{a_n\}$ 的通项公式。
    \begin{solution}
        有
        \[ 
        \frac{1}{S_n} = \frac{2S_{n-1} + 1}{S_{n-1}} = \frac{1}{S_{n-1}} + 2 
        \]
        说明数列 $\left\{\dfrac{1}{S_n}\right\}$ 是以 $\dfrac{1}{S_1} = \dfrac{1}{a_1} = 1$ 为首项, $2$ 为公差的等差数列,则 
        \[
        \frac{1}{S_n} = 1 + (n-1) \cdot 2 = 2n - 1 \Rightarrow S_n = \frac{1}{2n-1}
        \]
        当 $n \ge 2$ 时,
        \[ 
        a_n = S_n - S_{n-1} = \frac{1}{2n-1} - \frac{1}{2n-3} = -\frac{2}{(2n-1)(2n-3)} 
        \]
        故通项公式为 
        \[
        a_n = \begin{cases} 1 & ,n=1 \\[4pt] -\dfrac{2}{(2n-1)(2n-3)} & ,n \ge 2 \end{cases}
        \]
    \end{solution}

\question 已知正数数列 $\{a_n\}$ 的前 $n$ 项和 
\[
S_n = \frac{1}{2} \left(a_n + \frac{1}{a_n}\right) ,
\]
求 $\{a_n\}$ 的通项公式。
%曾龙文师-高二数理培训队笔记
\begin{solution}
当 $n=1$ 时,
\[
a_1 = S_1 = \frac{1}{2}\left(a_1 + \frac{1}{a_1}\right) \Rightarrow a_1^2 = 1 \Rightarrow a_1 = 1 >0
\]
当 $n \ge 2$ 时,有
\[ 
a_n = S_n - S_{n-1} = \frac{1}{2}\left(a_n + \frac{1}{a_n} - a_{n-1} - \frac{1}{a_{n-1}}\right) 
\]
\[ a_n + a_{n-1} = \frac{1}{a_n} - \frac{1}{a_{n-1}} = \frac{a_{n-1} - a_n}{a_n a_{n-1}} \]
由于 $a_n + a_{n-1} > 0$ ,等式两边同除以 $(a_n + a_{n-1})$ 得:
\[ 1 = \frac{a_{n-1} - a_n}{(a_n + a_{n-1}) a_n a_{n-1}} \]
经观察或递推可得 $S_n = \sqrt{n}$ 满足原式,则 $a_n = \sqrt{n} - \sqrt{n-1}$ 。
验证:当 $a_n = \sqrt{n} - \sqrt{n-1}$ 时,$\frac{1}{a_n} = \sqrt{n} + \sqrt{n-1}$ ,
则 $\frac{1}{2}(a_n + \frac{1}{a_n}) = \frac{1}{2}(\sqrt{n} - \sqrt{n-1} + \sqrt{n} + \sqrt{n-1}) = \sqrt{n} = S_n$ ,符合题意。
故 $a_n = \sqrt{n} - \sqrt{n-1}$ 。
\textcolor{red}{(待验证)}
\end{solution}

    \question 数列$\{a_n\}$中, $a_1=1$ 且 $a_n a_{n+1}=4^n$, 求前$n$项和$S_n$。
    \begin{solution}
        由题意,
        \[
        \frac{a_{n+2}}{a_n} = \frac{a_{n+2}}{a_{n+1}} \cdot \frac{a_{n+1}}{a_{n}} = 4,
        \]
        故 $\{a_1,a_3,a_5,\dots\}$ 与 $\{a_2,a_4,a_6,\dots\}$ 皆为公比为 $4$的等比数列,其中 
        \[
        a_1=1,a_2=4,
        \]
        当 $n$ 为奇数,
        \begin{align*}
        S_n 
        &= (a_1+a_3+\cdots+a_n) + (a_2+a_4+\cdots+a_{n-1}) \\
        &= \frac{1(1-4^{\frac{n+1}{2}})}{1-4} + \frac{4(1-4^{\frac{n-1}{2}})}{1-4} = \frac{2^{n+1}-1}{3} + \frac{2^{n+1}-4}{3} = \frac{4}{3}\cdot 2^n - \frac{5}{3}.
        \end{align*}
        当 $n$ 为偶数,
        \begin{align*}
        S_n 
        &= (a_1+a_3+\cdots+a_{n-1}) + (a_2+a_4+\cdots+a_n) \\
        &= \frac{1(1-4^{\frac{n}{2}})}{1-4} + \frac{4(1-4^{\frac{n}{2}})}{1-4} = \frac{5}{3}\cdot 2^n - \frac{5}{3}.
        \end{align*}
        综上,
        \[
        S_n = \frac{9+(-1)^n}{3}\cdot 2^{n-1} - \frac{5}{3}
        \]
    \end{solution}

    \question 已知数列 $\{x_{n}\}$ 的首项$x_{1}=3$,通项
    \[
    x_{n}=2^{n}p+nq, \quad n\in\mathbb{N},\quad p,q\in \mathbb{R},
    \]
    且$x_{1},x_{4},x_{5}$成等差数列,求:
    \begin{parts}
    \part $p,q$的值;
    \begin{solution}
        由$x_{1}=3$, 得
        \[
        2p+q=3 \tag{1}
        \]
        又$x_{1}+x_{5}=2x_{4}$,得
        \[
        3+2^{5}p+5q=2(2^{4}p+4q) \tag{2}
        \]
        解得
        \[
        p=1,\quad q=1
        \]
    \end{solution}
    \part 数列 $\{x_{n}\}$ 的前$n$项的和$S_{n}$的公式。
    \begin{solution}
        有
        \[
        S_{n}=(2^{1}+2^{2}+\cdots+2^{n})+(1+2+\cdots+n)= 2^{n+1}-2+\frac{n(n+1)}{2}
        \]
    \end{solution}
    \end{parts}

    \question 设数列 $\{a_n\}$ 的前 $n$ 项和为
    \[
    S_n=2a_n-2^n.
    \]
    \begin{parts}
    \part 求 $a_1,a_4$；
    \begin{solution}
        当 $n=1$ 时,解得
        \[
        S_1=a_1=2a_1-2 \Rightarrow a_1=2
        \]
        由$S_n=2a_n-2^n$,得
        \[
        2a_{n+1}=S_{n+1}+2^{n+1}=a_{n+1}+S_n+2^{n+1}
        \]
        从而
        \[
        a_{n+1}=S_n+2^{n+1}
        \]
        因此
        \[
        a_2=S_1+2^2=6,\ S_2=8,\
        a_3=S_2+2^3=16,\ S_3=24,\
        a_4=S_3+2^4=40.
        \]
    \end{solution}
    \part 证明$\{a_{n+1}-2a_n\}$ 是等比数列；
    \begin{solution}
        由上式可得
        \[
        a_{n+1}-2a_n=(S_n+2^{n+1})-(S_n+2^n)=2^n
        \]
        故数列 $\{a_{n+1}-2a_n\}$ 是首项为 $2$、公比为 $2$ 的等比数列。
    \end{solution}
    \part 求数列 $\{a_n\}$ 的通项公式。
    \begin{solution}
        由
        \[
        a_{n+1}-2a_n=2^n
        \]
        得
        \[
        a_n=(a_n-2a_{n-1})+2(a_{n-1}-2a_{n-2})+\cdots+2^{n-2}(a_2-2a_1)+2^{n-1}a_1.
        \]
        代入 $a_1=2,\ a_2=6$,并利用 $a_{k+1}-2a_k=2^k$,可得
        \[
        a_n=2^{n-1}(n+1)
        \]
    \end{solution}
    \end{parts}

    \question 已知数列 $\{a_n\}$ 的首项 $a_1=\dfrac{2}{3}$,且
    \[
    a_{n+1}=\frac{2a_n}{a_n+1}\quad (n=1,2,3,\dots).
    \]
    \begin{parts}
    \part 证明数列 $\left\{\dfrac{1}{a_n}-1\right\}$ 是等比数列；
    \begin{solution}
        由
        \[
        a_{n+1}=\frac{2a_n}{a_n+1}
        \]
        得
        \[
        \frac{1}{a_{n+1}}=\frac{a_n+1}{2a_n}
        =\frac{1}{2}+\frac{1}{2}\cdot\frac{1}{a_n}.
        \]
        于是
        \[
        \frac{1}{a_{n+1}}-1=\frac{1}{2}\left(\frac{1}{a_n}-1\right).
        \]
        又
        \[
        a_1=\frac{2}{3},\quad \frac{1}{a_1}-1=\frac{1}{2}.
        \]
        因此,数列 $\left\{\dfrac{1}{a_n}-1\right\}$ 是首项为 $\dfrac{1}{2}$、公比为 $\dfrac{1}{2}$ 的等比数列。
    \end{solution}
    \part 求数列 $\left\{\dfrac{n}{a_n}\right\}$ 的前 $n$ 项和 $S_n$。
    \begin{solution}
        由 (a) 可得
        \[
        \frac{1}{a_n}-1=\frac{1}{2}\ cdot \frac{1}{2^{n-1}}=\frac{1}{2^n} \Rightarrow \frac{n}{a_n} = \frac{n}{2^n} + n
        \]
        设 
        \[
        T_n = \frac{1}{2} + \frac{2}{2^2} + \frac{3}{2^3} + \dots + \frac{n}{2^n},
        \]
        则
        \[
        \frac{1}{2}T_n = \frac{1}{2^2} + \frac{2}{2^3} + \dots + \frac{n-1}{2^n} + \frac{n}{2^{n+1}}
        \]
        两式相减得
        \[
        \frac{1}{2}T_n = \frac{1}{2} + \frac{1}{2^2} + \dots + \frac{1}{2^n} - \frac{n}{2^{n+1}} = 1 - \frac{1}{2^n} - \frac{n}{2^{n+1}}  \Rightarrow T_n= 2-\frac{1}{2^{n-1}} - \frac{n}{2^{n}}
        \]
        又因为 
        \[
        1 + 2 + 3 + \dots + n = \frac{n(n+1)}{2},
        \]
        所以数列 $\left\{\dfrac{n}{a_n}\right\}$ 的前 $n$ 项和 $S_n$ 为
        \[
        S_n = 2 - \frac{n + 2}{2^n} + \frac{n(n+1)}{2} = \frac{n^2 + n + 4}{2} - \frac{n + 2}{2^n}
        \]
    \end{solution}
    \end{parts}

    \question 等差数列 $\{a_n\}$ 的各项均为正数, $a_1=3$, 前 $n$ 项和为 $S_n$, $\{b_n\}$ 为等比数列, $b_1=1$, 且
    \[
    b_2S_2=64,\quad b_3S_3=960.
    \]
    \begin{parts}
    \part 求 $a_n$ 与 $b_n$；
    \begin{solution}
        设等差数列 $\{a_n\}$ 的公差为 $d$, 等比数列 $\{b_n\}$ 的公比为 $q$。
        则
        \[
        a_n=3+(n-1)d,\quad b_n=q^{n-1}.
        \]
        由
        \[
        S_2=3+(3+d)=6+d,\quad S_3=3+(3+d)+(3+2d)=9+3d,
        \]
        代入已知条件得
        \[
        (6+d)q=64,\quad (9+3d)q^2=960
        \]
        解得
        \[
        d=2,\; q=8 \quad\text{或}\quad d=-\frac{6}{5},\; q=\frac{40}{3}
        \]
        因 $\{a_n\}$ 各项均为正数, 舍去第二组解,故
        \[
        a_n=3+2(n-1)=2n+1,\quad b_n=8^{n-1}.
        \]
    \end{solution}
    \part 求 
    \[
    \frac{1}{S_1}+\frac{1}{S_2}+\dots+\frac{1}{S_n}.
    \]
    \begin{solution}
        由 $a_n=2n+1$ 得
        \[
        S_n=3+5+\dots+(2n+1)=n(n+2).
        \]
        因此
        \begin{align*}
        \frac{1}{S_1}+\frac{1}{S_2}+\dots+\frac{1}{S_n}
        &=\frac{1}{1\cdot3}+\frac{1}{2\cdot4}+\dots+\frac{1}{n(n+2)} \\
        &=\frac{1}{2}\left(1-\frac{1}{3}+\frac{1}{2}-\frac{1}{4}+\dots+\frac{1}{n}-\frac{1}{n+2}\right)\\
        &=\frac{1}{2}\left(1+\frac{1}{2}-\frac{1}{n+1}-\frac{1}{n+2}\right)\\
        &=\frac{3}{4}-\frac{2n+3}{2(n+1)(n+2)}
        \end{align*}
    \end{solution}
    \end{parts}

    \question 设$S_n$是等差数列$\{a_n\}$前$n$项的和,已知$\dfrac{1}{3}S_3$与$\dfrac{1}{4}S_4$的等比中项为$\dfrac{1}{5}S_5$,等差中项为$1$,求等差数列的通项$a_n$。
    \begin{solution}
        由$(S_3)(S_4)=(S_5)^2$及$\dfrac{1}{3}S_3+\dfrac{1}{4}S_4=2$得
        \[
        \begin{cases} 
        \displaystyle \frac{1}{3}\left(3a+\frac{3\cdot 2}{2}d\right)\cdot\frac{1}{4}\left(4a+\frac{4\cdot3}{2}d\right)=\frac{1}{25}\left(5a+\frac{5\cdot4}{2}d\right)^2 \\[8pt] 
        \displaystyle \frac{1}{3}\left(3a+\frac{3\cdot2}{2}d\right)+\frac{1}{4}\left(4a+\frac{4\cdot3}{2}d\right)=2 
        \end{cases}
        \]
        解得
        \[
        d=0,a=1 \;\; \text{或} \;\; d=-\frac{12}{5},a=4
        \]
        经检验得通项解为$a_n=1$或
        \[
        a_n=4-\frac{12}{5}(n-1)=\frac{32}{5}-\frac{12}{5}n
        \]
    \end{solution}

    \question 已知数列 $\{a_n\}$ 中每一项均为正数,且数列前 $n$ 项之和为 $S_n$,若 
    \[
    \sum_{k=1}^{n}\frac{4S_k}{a_k+2}=S_n,
    \]
    试求 $a_n$ 及 $S_n$。
    \begin{solution}
        由递推关系得
        \[
        S_1 = a_1 = \frac{4a_1}{a_1+2} 
        \Rightarrow a_1^2 - 2a_1 = 0 
        \]
        因\ $a_1 \ne 0$,解得 $a_1 = 2$,于是
        \[
        a_n = S_n - S_{n-1} = \sum_{k=1}^n \frac{4S_k}{a_k+2} - \sum_{k=1}^{n-1} \frac{4S_k}{a_k+2} = \frac{4S_n}{a_n+2} 
        \Rightarrow S_n = \frac{1}{4}a_n(a_n+2)
        \]
        又
        \[
        S_{n-1} = S_n-a_n= \frac{1}{4}a_n^2 - \frac{1}{2}a_n = \frac{1}{4}a_{n-1}(a_{n-1}+2)
        \]
        因此
        \[
        a_n^2 - 2a_n = a_{n-1}^2 + 2a_{n-1}
        \Rightarrow (a_n - 1)^2 = (a_{n-1} + 1)^2
        \Rightarrow a_n - 1 = a_{n-1} + 1
        \]
        即
        \[
        a_n = a_{n-1} + 2 = \cdots = a_1 + 2(n-1) = 2n, S_n = n^2 + n
        \]
    \end{solution}

    \question 已知$n^2 (n \ge 4)$ 个正数排成$n$行、$n$列,如下:
    \[ 
    \begin{matrix}
    a_{11} & a_{12} & a_{13} & a_{14} & \dots & a_{1n} \\
    a_{21} & a_{22} & a_{23} & a_{24} & \dots & a_{2n} \\
    a_{31} & a_{32} & a_{33} & a_{34} & \dots & a_{3n} \\
    \dots & \dots & \dots & \dots & \dots & \dots \\
    a_{n1} & a_{n2} & a_{n3} & a_{n4} & \dots & a_{nn}
    \end{matrix} 
    \]
    且每一行的数成等差数列,每一列的数成等比数列,并且所有公比相等。已知 
    \[
    a_{24} = 1, a_{42} = \frac{1}{8}, a_{43} = \frac{3}{16},
    \]
    试求 
    \[
    a_{11} + a_{22} + \dots + a_{nn}
    \]
    的值。
    \begin{solution}
        设每一列的公比为 $q$ ,第 $i$ 行的等差数列公差为 $d_i$,
        由每一列成等比数列且公比相等可知
        \[
        a_{ij} = a_{1j} \cdot q^{i-1}
        \]
        对于第四行, $a_{42} = \dfrac{1}{8},a_{43} = \dfrac{3}{16}$ ,则公差为$d_4 = a_{43} - a_{42} = \dfrac{1}{16}$,据此有 
        \[
        a_{44} = a_{43} + d_4 = \dfrac{1}{4},
        \]
        又已知 $a_{24} = 1$ ,且第四列成等比数列,则
        \[ 
        a_{44} = a_{24} \cdot q^2 \Rightarrow q = \frac{1}{2}>0
        \]
        由 
        \[
        a_{1j} = a_{4j} \cdot 2^3 = 8 \cdot \frac{j}{16} = \frac{j}{2}
        \]
        因此 
        \[
        a_{ii} = a_{1i} \cdot q^{i-1} = \frac{i}{2} \cdot \left( \frac{1}{2} \right)^{i-1} = \frac{i}{2^i}
        \]
        所求和为
        \[ 
        S = \sum_{i=1}^n a_{ii} = \frac{1}{2^1} + \frac{2}{2^2} + \dots + \frac{n}{2^n} 
        \]
        利用错位相减法,
        \[ 
        \frac{1}{2}S = \frac{1}{2^2} + \frac{2}{2^3} + \dots + \frac{n}{2^{n+1}} 
        \]
        \[ 
        \frac{1}{2}S = \frac{1}{2^1} + \frac{1}{2^2} + \dots + \frac{1}{2^n} - \frac{n}{2^{n+1}} 
        \]
        \[ 
        \frac{1}{2}S = (1 - \frac{1}{2^n}) - \frac{n}{2^{n+1}} 
        \]
        得
        \[ 
        S = 2 - \frac{n+2}{2^n} 
        \]
    \end{solution}

    \question 已知数列$\{a_n\}$满足 $\begin{cases} a_0 = 1 \\ a_{2k+1} = a_k \\ a_{2k+2} = a_k + a_{k+1} \end{cases}, k \in N \cup \{0\}$, 求$\displaystyle \sum_{k=0}^{63} a_k$。
    \begin{solution}
        记
        \[
        S(n)=\sum_{k=0}^{2^n-1}a_k
        \]
        则
        \[
        S(n)=(a_0+a_2+a_4+\cdots+a_{2^n-2})+(a_1+a_3+\cdots+a_{2^n-1})
        \]
        其中偶数项为
        \[
        a_0+(a_0+a_1)+(a_1+a_2)+(a_2+a_3)+\cdots+(a_{2^{n-1}-2}+a_{2^{n-1}-1})
        =2(a_0+a_1+\cdots+a_{2^{n-1}-2})+a_{2^{n-1}-1}
        \]
        奇数项为
        \[
        a_0+a_1+a_2+\cdots+a_{2^{n-1}-1}
        \]
        所以
        \[
        S(n)=3(a_0+a_1+\cdots+a_{2^{n-1}-2})+2a_{2^{n-1}-1}
        =3(a_0+a_1+\cdots+a_{2^{n-1}-1})-a_{2^{n-1}-1}
        \]
        由 $a_{2^{n-1}-1}=1$,得
        \[
        S(n)=3S(n-1)-1
        \]
        于是
        \[
        S(8)=3S(7)-1=3^2S(6)-3-1=\cdots
        =3^5S(1)-(1+3+\cdots+3^4)
        \]
        又 $S(1)=a_0+a_1=2$,所以
        \[
        S(8)=243\cdot 2-121=365
        \]
    \end{solution}

    \question 已知各项皆为正整数的数列$\{ a_n \}$的前 $n$ 项和为$S_n$, 且对任意正整数$n$有
    \[
    \sqrt{S_n} = \lambda(a_n-1)+1,
    \]
    其中$\lambda$为正实数。若$2a_2=a_1+a_3$, 试求数列$\{a_n\}$的一般项。
    \begin{solution}
        由$2a_2=a_1+a_3$,设
        \[
        d=a_2-a_1=a_3-a_2
        \]
        则
        \[
        \sqrt{S_2}-\sqrt{S_1}=\lambda(a_2-a_1)=\lambda(a_3-a_2)=\sqrt{S_3}-\sqrt{S_2} \Rightarrow 2\sqrt{S_2}=\sqrt{S_1}+\sqrt{S_3}
        \]
        两边平方得
        \[
        4S_2=S_1+S_3+2\sqrt{S_1S_3} \Rightarrow 4(2a_1+d)=4a_1+3d+2\sqrt{a_1(3a_1+3d)}
        \]
        解得$d=2a_1$,则
        \[
        S_1=a_1,\; S_2=4a_1,\; S_3=9a_1
        \]
        解
        \[
        \lambda={\sqrt{S_2}-\frac{1}{ a_2-1}}
        ={\sqrt{S_3}-\frac{1}{a_3-1}}
        \]
        得$a_1=1,\lambda=\dfrac12$,于是
        \begin{align*}
        \sqrt{S_n}&=\frac12(a_n-1)+1 \\[1mm]
        a_n=S_n-S_{n-1}&=2\sqrt{S_n}-1 \\
        (\sqrt{S_n}-1)^2&=S_{n-1} \\
        \sqrt{S_n}&=\sqrt{S_{n-1}}+1 \\[1mm]
        \frac{a_n+1}{2}&=\frac{a_{n-1}+1}{2}+1
        \end{align*}
        故
        \[
        a_n=a_{n-1}+2=a_{n-2}+4=\cdots=a_1+2(n-1)=2n-1
        \]
    \end{solution}

    \question 已知数列 $\{a_n\}$ 满足 $a_1=1$, $a_{n+1}=\dfrac{a_n^2}{2a_n+1}$, 求通项公式并证明
    \[
    \sum_{k=1}^{n}\frac{a_k}{1+a_k}<\frac{7}{8}.
    \]
    \begin{solution}
        由递推式
        \[
        a_{n+1}=\frac{a_n^2}{2a_n+1},
        \]
        可得
        \[
        \frac{1}{a_{n+1}}=\frac{2a_n+1}{a_n^2}=\frac{2}{a_n}+\frac{1}{a_n^2},
        \]
        从而
        \[
        1+\frac{1}{a_{n+1}}=1+\frac{2}{a_n}+\frac{1}{a_n^2}=\left(1+\frac{1}{a_n}\right)^2.
        \]
        取对数,得
        \[
        \log\!\left(1+\frac{1}{a_{n+1}}\right)=2\log\!\left(1+\frac{1}{a_n}\right).
        \]
        因此数列 $\left\{\log\!\left(1+\tfrac{1}{a_n}\right)\right\}$ 是首项为 $\log 2$、公比 $2$ 的等比数列,于是
        \[
        \log\!\left(1+\frac{1}{a_n}\right)=2^{\,n-1}\log 2=\log\!\left(2^{2^{n-1}}\right),
        \]
        即
        \[
        1+\frac{1}{a_n}=2^{2^{\,n-1}}\Rightarrow a_n=\frac{1}{2^{2^{\,n-1}}-1}.
        \]
        且有
        \[
        \frac{a_n}{1+a_n}=\frac{1}{2^{2^{\,n-1}}}.
        \]
        所以
        \[
        \sum_{k=1}^{n}\frac{a_k}{1+a_k}=\frac{1}{2}+\frac{1}{2^2}+\frac{1}{2^4}+\cdots+\frac{1}{2^{2^{\,n-1}}}.
        \]
        注意当 $n\geq 4$ 时,有
        \[
        2^{n-1} > n+1 \Rightarrow \left(\frac{1}{2}\right)^{2^{n-1}} < \left(\frac{1}{2}\right)^{n+1}.
        \]
        因此
        \begin{align*}
        \sum_{k=1}^{n}\frac{a_k}{1+a_k}
        &< \frac{1}{2}+\frac{1}{4}+\frac{1}{16}+\sum_{k=5}^{n+1}\left(\frac{1}{2}\right)^k \\
        &= \frac{1}{2}+\frac{1}{4}+\frac{1}{16}+\frac{\left(\tfrac{1}{2}\right)^5\bigl(1-(\tfrac{1}{2})^{n-3}\bigr)}{1-\tfrac{1}{2}} \\
        &= \frac{7}{8}-\left(\frac{1}{2}\right)^{n+1} < \frac{7}{8}
        \end{align*}
        当 $n=1,2,3$ 时直接检验,亦满足不等式,故
        \[
        \sum_{k=1}^{n}\frac{a_k}{1+a_k}<\frac{7}{8},\quad \forall n\in\mathbb{N}.
        \]
    \end{solution}

    \question 设各项均为正数的数列 $\{a_{n}\}$ 满足 
    \[
    a_{1}=2, a_{n} = a_{n+1}^{\frac{3}{2}}a_{n+2}
    \]
    其中$n\ge2$。
    \begin{parts}
    \part 若 $a_{2}=\dfrac{1}{4}$, 求$a_{3}, a_{4}$, 并猜想$a_{2008}$的值(不需证明);
    \begin{solution}
        因
        \[
        a_{1}=2\left(-2\right)^{0}, a_{2}=2\left(-2\right)^{\frac{1}{4}}, a_{3}=2\left(-2\right)^{\frac{1}{2}}, a_{4}=2\left(-2\right)^{\frac{3}{4}}
        \]
        从而猜想$a_{n}$的通项为
        \[
        a_{n} = 2\left(-2\right)^{\frac{n-1}{4}}
        \]
        据此有
        \[
        a_{2008}=2\left(-2\right)^{\frac{2007}{4}}
        \]
    \end{solution}
    \part 若 
    \[
    2\sqrt{2}\leq a_{1}a_{2}\cdots a_{n} < 4
    \]
    对$n\geq 2$恒成立, 求$a_{2}$的值.
    \begin{solution}
        令 $x_n = \log_2 a_n$,则 $a_2 = 2^{x_2}$,求 $x_2$ 即可。
        设 $S_n$ 表示 $x_n$ 的前 $n$ 项和,则 
        \[
        a_1 a_2 \cdots a_n = 2^{S_n}
        \]
        由 $2\sqrt{2} \leq a_1 a_2 \cdots a_n < 4$ 得
        \[
        \frac{3}{2} \leq S_n = x_1 + x_2 + \cdots + x_n < 2, \quad n \geq 2
        \]
        由于上式对 $n=2$ 成立,可得 
        \[
        \frac{3}{2} \leq x_1 + x_2 < 2
        \]
        由 $a_1 = 2$ 得 $x_1 = 1$,故 $x_2 \geq \dfrac{1}{2}$,又由$a_1 = 2,a_{1}=2, a_{n} = a_{n+1}^{\frac{3}{2}}a_{n+2}$得
        \[
        x_n=\frac{3}{2}x_{n+1}+x_{n+2}, \quad n \geq 2
        \]
        即
        \[
        x_{n+2} + 2x_{n+1} = x_{n+2} + \frac{3}{2}x_{n+1} + \frac{1}{2}x_{n+1} = \frac{1}{2}(x_{n+1} + 2x_n)
        \]
        由此可知数列 $\{x_{n+1} + 2x_n\}$ 是首项为 $ x_2 + 2$,公比为 $\dfrac{1}{2}$ 的等比数列,故
        \[
        x_{n+1} + 2x_n = (x_2 + 2) \left(\frac{1}{2}\right)^{n-1} 
        \]
        对 $n$ 求和得
        \[
        S_{n+1} - x_1 + 2S_n = (x_2 + 2) \left( 1 + \frac{1}{2} + \cdots + \frac{1}{2^{n-1}} \right) = (x_2 + 2)\left( 2 - \frac{1}{2^{n-1}} \right) 
        \]
        因 $S_n < 2,S_{n+1} < 2$,且 $x_1 = 1$,故
        \[
        (x_2 + 2)\left( 2 - \frac{1}{2^{n-1}} \right) = S_{n+1} + 2S_n - 1 < 2 + 4 - 1 = 5 
        \]
        由此整理得:
        \[
        2x_2 - 1 < \frac{x_2 + 2}{2^{n-1}},\quad n\ge 2
        \]
        下证 $x_2 \leq \dfrac{1}{2}$。假设 $x_2 > \dfrac{1}{2}$,则由上式可知
        \[
        2^{n-1} < \frac{x_2 + 2}{2x_2 - 1}
        \]
        对 $n \geq 2$ 恒成立,但这对于趋于无穷的 $n$ 是不可能的,因此必有 
        \[
        x_2 \leq \frac{1}{2}
        \]
        又已知 $x_2 \geq \dfrac{1}{2}$,故 $x_2 = \dfrac{1}{2}$,所以 
        \[
        a_2 = 2^{x_2} = \sqrt{2}
        \]
    \end{solution}
    \end{parts}

    \question 在数列 $\{a_n\}$,$\{b_n\}$ 是各项均为正数的等比数列,设
    \[
    c_n=\frac{b_n}{a_n},\quad n\in \mathbb{N}
    \]
    \begin{parts}
    \part 数列 $\{c_n\}$ 是否为等比数列？证明你的结论。
    \begin{solution}
        设 $\{a_n\}$ 的公比为 $q_1>0,\{b_n\}$ 的公比为 $q_2>0$,则
        \[
        \frac{c_{n+1}}{c_n}
        =\frac{b_{n+1}}{a_{n+1}}\cdot\frac{a_n}{b_n}
        =\frac{b_{n+1}}{b_n}\cdot\frac{a_n}{a_{n+1}}
        =\frac{q_2}{q_1}\ne 0,
        \]
        故 $\{c_n\}$ 是公比为$\dfrac{q_2}{q_1}$的等比数列。
    \end{solution}
    \part 设数列 $\{\ln a_n\},\{\ln b_n\}$ 的前 $n$ 项和分别为 $S_n,T_n$。若 $a_1=2$,
    \[
    \frac{S_n}{T_n}=\frac{n}{2n+1},
    \]
    求数列 $\{c_n\}$ 的前 $n$ 项和。
    \begin{solution}
    因 $\{a_n\}$,$\{b_n\}$ 为等比数列,
    \[
    S_n=n\ln a_1+\frac{n(n-1)}{2}\ln q_1,\quad
    T_n=n\ln b_1+\frac{n(n-1)}{2}\ln q_2
    \]
    由题意
    \[
    \frac{n\ln a_1+\frac{n(n-1)}{2}\ln q_1}
    {n\ln b_1+\frac{n(n-1)}{2}\ln q_2}
    =\frac{n}{2n+1}
    \]
    化简得对一切 $n$ 成立的多项式恒等式:
    \[
    (2\ln q_1-\ln q_2)n^2
    +(4\ln a_1-\ln q_1-2\ln b_1+\ln q_2)n
    +(2\ln a_1-\ln q_1)=0
    \]
    于是由
    \[
    \begin{cases}
    2\ln q_1-\ln q_2=0\\
    4\ln a_1-\ln q_1-2\ln b_1+\ln q_2=0\\
    2\ln a_1-\ln q_1=0
    \end{cases}
    \]
    及 $a_1=2$ 解得
    \[
    q_1=4,\quad q_2=16,\quad b_1=8
    \]
    故
    \[
    c_n=\frac{b_n}{a_n}
    =\frac{8\cdot16^{\,n-1}}{2\cdot4^{\,n-1}}
    =4^n
    \]
    所以数列 $\{c_n\}$ 的前 $n$ 项和为
    \[
    c_1+c_2+\cdots+c_n
    =4+4^2+\cdots+4^n
    =\frac{4}{3}\left(4^n-1\right)
    \]
    \end{solution}
    \end{parts}

    \question 将数列 $\{a_n\}$ 中的所有项按每一行比上一行多一项的规则排成如下
    \[
    \begin{array}{l}
    a_1\\
    a_2,\ a_3\\
    a_4,\ a_5,\ a_6\\
    a_7,\ a_8,\ a_9,\ a_{10}\\
    \end{array}
    \]
    记表中的第一列数 $a_1,a_2,a_4,a_7,\cdots$ 构成的数列为 $\{b_n\}$,且$b_1=a_1=1$。$S_n$ 为数列 $\{b_n\}$ 的前 $n$ 项和,且满足
    \[
    \frac{2b_n}{b_nS_n-S_n^2}=1,\quad n\ge2
    \]
    \begin{parts}
    \part 证明数列 $\left\{\dfrac{1}{S_n}\right\}$ 成等差数列,并求数列 $\{b_n\}$ 的通项公式；
    \begin{solution}
        由$b_n=S_n-S_{n-1}$,当$n\ge2$,
        \[
        \frac{2(S_n-S_{n-1})}{(S_n-S_{n-1})S_n-S_n^2}=1
        \]
        化简得
        \[
        \frac{1}{S_n}-\frac{1}{S_{n-1}}=\frac{1}{2}.
        \]
        又 $S_1=b_1=1$,故数列 $\left\{\dfrac{1}{S_n}\right\}$ 是首项为 $1$,公差为 $\dfrac12$ 的等差数列,
        \[
        \frac{1}{S_n}=1+\frac12(n-1)=\frac{n+1}{2} \Rightarrow S_n=\frac{2}{n+1}
        \]
        则
        \[
        b_n=S_n-S_{n-1}=\frac{2}{n+1}-\frac{2}{n}=-\frac{2}{n(n+1)}
        \]
        因此
        \[
        b_n=
        \begin{cases}
        1&, n=1,\\[4pt]
        -\dfrac{2}{n(n+1)}&, n\ge2
        \end{cases}
        \]
    \end{solution}
    \part 若从第三行起,每一行中的数按从左到右的顺序均构成等比数列,且公比为同一个正数。当 $a_{81}=-\dfrac{4}{91}$ 时,求上表中第 $k\ge3$ 行所有项的和。
    \begin{solution}
        设从第三行起,每一行的公比为 $q>0$。因为
        \[
        1+2+\cdots+12=\frac{12 \cdot 13}{2}=78,
        \]
        所以表中第1行至第12行共含有数列 $\{a_n\}$ 的前 $78$ 项,故 $a_{81}$ 在表中第 $13$ 行第 $3$ 列,于是
        \[
        a_{81}=b_{13}q^2=-\frac{4}{91}.
        \]
        又
        \[
        b_{13}=-\frac{2}{13\cdot 14},
        \]
        解得 $q=2$,记表中第 $k\ge3$ 行所有项的和为 $S$,则
        \[
        S=-\frac{2}{k(k+1)}\cdot\frac{1-2^{k}}{1-2}
        =\frac{2}{k(k+1)}(1-2^{k}),\quad k\ge3
        \]
    \end{solution}
    \end{parts}

    \question 数列 $\{a_n\}$ 满足
    \[
    a_1=0,\ a_2=2,\ a_{n+2}=\left(1+\cos^2\frac{n\pi}{2}\right)a_n+4\sin^2\frac{n\pi}{2},\quad n=1,2,3,\dots
    \]
    \begin{parts}
    \part 求 $a_3,a_4$, 并求数列 $\{a_n\}$ 的通项公式；
    \begin{solution}
        已知 $a_1=0,a_2=2$,则
        \[
        a_3=\left(1+\cos^2\frac{\pi}{2}\right)a_1+4\sin^2\frac{\pi}{2}=4,\quad        a_4=\left(1+\cos^2\pi\right)a_2+4\sin^2\pi=4
        \]
        一般地, 当 $n=2k-1,\ k\in\mathbb{N}$ 时,
        \[
        a_{2k+1}=\left(1+\cos^2\frac{(2k-1)\pi}{2}\right)a_{2k-1}
        +4\sin^2\frac{(2k-1)\pi}{2}
        =a_{2k-1}+4
        \]
        即
        \[
        a_{2k+1}-a_{2k-1}=4
        \]
        因此 $\{a_{2k-1}\}$ 为首项为 $0$, 公差为 $4$ 的等差数列,
        \[
        a_{2k-1}=4(k-1)
        \]
        当 $n=2k,\ k\in\mathbb{N}$ 时,
        \[
        a_{2k+2}=\left(1+\cos^2\frac{2k\pi}{2}\right)a_{2k}+4\sin^2\frac{2k\pi}{2}
        =2a_{2k}
        \]
        因此 $\{a_{2k}\}$ 为首项为 $2$, 公比为 $2$ 的等比数列,
        \[
        a_{2k}=2^k
        \]
        综上,
        \[
        a_n=
        \begin{cases}
        4(k-1), & n=2k-1,\ k\in\mathbb{N}\\
        2^k, & n=2k,\ k\in\mathbb{N}
        \end{cases}
        \]
    \end{solution}
    \part 设
    \[
    S_k=a_1+a_3+\dots+a_{2k-1},\quad
    T_k=a_2+a_4+\dots+a_{2k},\quad
    W_k=\frac{2S_k}{2+T_k},\quad k\in\mathbb{N},
    \]
    求使 $W_k>1$ 的所有 $k$ 的值。
    \begin{solution}
        由(1)得
        \[
        S_k=0+4+\dots+4(k-1)=2k(k-1),\quad T_k=2+2^2+\dots+2^k=2^{k+1}-2
        \]
        于是
        \[
        W_k=\frac{2S_k}{2+T_k}
        =\frac{4k(k-1)}{2^{k+1}}
        =\frac{k(k-1)}{2^{k-1}}
        \]
        计算得
        \[
        W_1=0,\quad W_2=1,\quad W_3=\frac{3}{2},\quad
        W_4=\frac{3}{2},\quad W_5=\frac{5}{4},\quad W_6=\frac{15}{16}
        \]
        当 $k\ge6$ 时,
        \[
        W_{k+1}-W_k
        =\frac{(k+1)k}{2^k}-\frac{k(k-1)}{2^{k-1}}=\frac{k(3-k)}{2^k}<0
        \]
        故 $\{W_k\}$ 单调递减,又 $W_6<1$, 因此当 $k\ge6$ 时 $W_k<1$。综上, 满足 $W_k>1$ 的所有 $k$ 为
        \[
        k=3,4,5
        \]
    \end{solution}
    \end{parts}

    \question 证明由正实数组成的数列\(a_0, a_1, a_2, \ldots, a_n\) 是等比数列当且仅当
    \[
    \left(a_0 a_1 + a_1 a_2 + a_2 a_3 + \cdots + a_{n-1} a_n\right)^2 = \left(a_0^2 + a_1^2 + \cdots + a_{n-1}^2\right)\left(a_1^2 + a_2^2 + \cdots + a_n^2\right).
    \]
    \begin{solution}
        充分条件:已知 $a_0,\, a_1,\, \ldots,\, a_n$ 是等比数列,设$a_k = a_0 r^k,\; 1 \leq k \leq n$, 于是
        \[
        a_0 a_1 + a_1 a_2 + \ldots + a_{n-1} a_n 
        = a_0^2 (r + r^3 + \ldots + r^{2n-1}) = a_0^2 r \left(1 + r^2 + \ldots + r^{2n-2}\right)
        \]
        且
        \begin{align*}
        (a_0^2 + a_1^2 + \ldots + a_{n-1}^2)(a_1^2 + a_2^2 + \ldots + a_n^2)
        &= a_0^4 (1 + r^2 + \ldots + r^{2n-2})(r^2 + r^4 + \ldots + r^{2n}) \\
        &= a_0^4 r^2 \left(1 + r^2 + \ldots + r^{2n-2}\right)^2
        \end{align*}
        故得证
        \[
        (a_0 a_1 + a_1 a_2 + \ldots + a_{n-1} a_n)^2 = (a_0^2 + a_1^2 + \ldots + a_{n-1}^2)(a_1^2 + a_2^2 + \ldots + a_n^2)
        \]
        必要条件:已知
        \[
        (a_0 a_1 + a_1 a_2 + \ldots + a_{n-1} a_n)^2 = (a_0^2 + a_1^2 + \ldots + a_{n-1}^2)(a_1^2 + a_2^2 + \ldots + a_n^2)
        \]
        考虑函数
        \[
        f(r) = \sum_{k=1}^n (a_k - r a_{k-1})^2 = r^2 \sum_{k=1}^n a_{k-1}^2 - 2r \sum_{k=1}^n a_k a_{k-1} + \sum_{k=1}^n a_k^2
        \]
        观察得 $f(r) \geq 0 \quad \forall r$,又关于$r$方程式的判别式为
        \[
        4 \left[ (a_0 a_1 + a_1 a_2 + \ldots + a_{n-1} a_n)^2 - (a_0^2 + a_1^2 + \ldots + a_{n-1}^2)(a_1^2 + a_2^2 + \ldots + a_n^2) \right] = 0,
        \]
        所以得知$f(r) = 0$ 有重根 $r = r_0$,且只有当  
        \[
        a_k = r_0 a_{k-1},\quad 1 \leq k \leq n
        \]
        才成立,意味 $a_0,\, a_1,\, \ldots,\, a_n$ 是一个公比为$r_0$ 的等比数列。
    \end{solution}

    \question 求和
    \[
    \sum_{n=2}^{20} \frac{n^3-n^2+1}{n^2-n}
    \]
    \begin{solution}
        注意到
        \[
        \frac{n^3-n^2+1}{n^2-n} = n + \frac{1}{n(n-1)} = n + \frac{1}{n-1} - \frac{1}{n}
        \]
        于是由累加法,
        \begin{align*}
        \sum_{n=2}^{20} \frac{n^3-n^2+1}{n^2-n}
        &= \sum_{n=2}^{20} n + \sum_{n=2}^{20} \left(\frac{1}{n-1} - \frac{1}{n} \right) \\
        &= \frac{19}{2}(2+20) + 1 - \frac{1}{20} \\
        &= \frac{4199}{20}
        \end{align*}
    \end{solution}

    \question 试求级数
    \[
    \frac{1\cdot 3\cdot 5}{3^6-64} + \frac{3\cdot 5\cdot 7}{5^6-64} + \frac{5\cdot 7\cdot 9}{7^6-64} + \cdots + \frac{19\cdot 21\cdot 23}{21^6-64}
    \] 
    之和。
    \begin{solution}
        注意到
        \begin{align*}
        &\frac{(2k-1)(2k+1)(2k+3)}{(2k+1)^6 - 64}\\
        &= \frac{(2k-1)(2k+1)(2k+3)}{((2k+1)^3-8)((2k+1)^3+8)} \\
        &= \frac{(2k+1)}{((2k+1)^2 + 2(2k+1) + 4)((2k+1)^2 - 2(2k+1) + 4)}\\
        &= \frac{2k+1}{(4k^2 + 3)(4k^2 + 8k + 7)} \\
        &= \frac{2k+1}{(4k^2 + 3)(4(k+1)^2 + 3)}\\
        &= \frac{1}{4} \left( \frac{1}{4k^2 + 3} - \frac{1}{4(k+1)^2 + 3} \right)
        \end{align*}
        故
        \[
        \sum_{k=1}^{10} \frac{(2k-1)(2k+1)(2k+3)}{(2k+1)^6 - 64}= \frac{1}{4} \left( \frac{1}{7} - \frac{1}{487} \right)
        = \frac{120}{3409}
        \]
    \end{solution}

    \question 求和
    \[
    \sum_{n=1}^{24} \frac{1}{\sqrt{n + \sqrt{n^2 - 1}}} 
    \]
    \begin{solution}
        不难发现
        \[
        \frac{1}{\sqrt{n + \sqrt{n^2 - 1}}} = \sqrt{n - \sqrt{n^2 - 1}} = \frac{1}{\sqrt{2}}(\sqrt{n+1} - \sqrt{n-1}),
        \]
        由累加法,
        \begin{align*}
        \sum_{n=1}^{24} \frac{1}{\sqrt{n + \sqrt{n^2 - 1}}} &= \frac{1}{\sqrt{2}} \sum_{n=1}^{24} (\sqrt{n+1} - \sqrt{n-1})\\
        &= \frac{1}{\sqrt{2}} (\sqrt{25} + \sqrt{24} - \sqrt{1} - \sqrt{0})\\
        &=2\sqrt{2} + 2\sqrt{3}
        \end{align*}
    \end{solution}

    \question 若 $f(n)=(n^{2}-2n+1)^{\frac{1}{3}}+(n^{2}-1)^{\frac{1}{3}}+(n^{2}+2n+1)^{\frac{1}{3}},$ 求 
    \[
    \sum_{k=1}^{500}\frac{1}{f(2k-1)}
    \]
    \begin{solution}
        首先有
        \begin{align*}
        \frac{1}{f(n)} &= \frac{1}{\sqrt[3]{(n-1)^2} +\sqrt[3]{(n+1)(n-1)} +\sqrt[3]{(n+1)^2}} \\
        &= \frac{\sqrt[3]{(n+1)} - \sqrt[3]{(n-1)}}{(n+1)-(n-1)} \\
        &= \frac{1}{2}(\sqrt[3]{(n+1)} - \sqrt[3]{(n-1)})
        \end{align*}
        故
        \[
        \sum_{k=1}^{500} \frac{1}{f(2k-1)} = \frac{1}{2} \sum_{k=1}^{500} \left( \sqrt[3]{2k} - \sqrt[3]{2k-2} \right) = \frac{1}{2} \cdot \sqrt[3]{1000} = 5
        \]
    \end{solution}

    \question 设数列 $\{a_n\}$ 的通项公式为
    \[
    a_n = \frac{1}{(\sqrt{n}+\sqrt{n+1})(\sqrt{n+1}+\sqrt{n+2})(\sqrt{n}+\sqrt{n+2})},
    \]
    求 
    \[
    \sum_{k=1}^{\infty} a_k
    \]
    的值。
    \begin{solution}
        发现
        \begin{align*}
        a_n &= \frac{1}{\sqrt{n+2}-\sqrt{n}} \left( \frac{1}{(\sqrt{n}+\sqrt{n+1})(\sqrt{n}+\sqrt{n+2})} - \frac{1}{(\sqrt{n}+\sqrt{n+2})(\sqrt{n+1}+\sqrt{n+2})} \right) \\
        &= \frac{1}{(\sqrt{n+2}-\sqrt{n})(\sqrt{n}+\sqrt{n+2})} \left( \frac{1}{\sqrt{n}+\sqrt{n+1}} - \frac{1}{\sqrt{n+1}+\sqrt{n+2}} \right) \\
        &= \frac{1}{2} \left( \frac{1}{\sqrt{n}+\sqrt{n+1}} - \frac{1}{\sqrt{n+1}+\sqrt{n+2}} \right)
        \end{align*}
        由$\lim_{n \to \infty} \dfrac{1}{\sqrt{n}+\sqrt{n+1}}=0,$故
        \[
        \sum_{k=1}^{\infty} a_k = \frac{1}{2} \left( \frac{1}{\sqrt{1}+\sqrt{2}} - \frac{1}{\sqrt{2}+\sqrt{3}} + \frac{1}{\sqrt{2}+\sqrt{3}} - \frac{1}{\sqrt{3}+\sqrt{4}} + \cdots \right)
        = \frac{1}{2} \cdot \frac{1}{1+\sqrt{2}} = \frac{\sqrt{2}-1}{2}
        \]
    \end{solution}
    
    \question 求 
    \[
    \sqrt{1+\frac{1}{1^{2}}+\frac{1}{2^{2}}}+\sqrt{1+\frac{1}{2^{2}}+\frac{1}{3^{2}}}+\dots+\sqrt{1+\frac{1}{2025^{2}}+\frac{1}{2026^{2}}}
    \]
    的值。
    \begin{solution}
        注意到
        \begin{align*}
        &\sqrt{1+\frac{1}{n^2} +\frac{1}{(n+1)^2}} 
        = \sqrt{1+\frac{(n+1)^2+n^2}{n^2(n+1)^2} } \\
        &= \sqrt{1+\frac{2n(n+2) +1}{n^2(n+1)^2} } 
        = \sqrt{1+\frac{2}{n(n+1)} +\frac{1}{n^2(n+1)^2} } \\
        &= \sqrt{\left(1+\frac{1}{n(n+1)}\right)^2} 
        = 1+\frac{1}{n(n+1)} 
        = 1+\frac{1}{n}-\frac{1}{n+1}
        \end{align*}
        因此原式为
        \[
        \sum_{n=1}^{2025} \left(1 + \frac{1}{n} - \frac{1}{n+1}\right)= 2025 + \left(1 - \frac{1}{2026} \right)=2016-\frac{1}{2026}
        \]
    \end{solution}

    \question 求
    \[
    \sum_{r=2}^{\infty}\left[\frac{4r-1}{r(r-1)}\left(-\frac{1}{3}\right)^r\right]
    \]
    \begin{solution}
        作部分分式分解,有
        \[
        \frac{4r-1}{r(r-1)}=\frac{1}{r}+\frac{3}{r-1}
        \]
        因此
        \begin{align*}
        \frac{4r-1}{r(r-1)}\left(-\frac{1}{3}\right)^r
        &=\frac{1}{r}\left(-\frac{1}{3}\right)^r+\frac{3}{r-1}\left(-\frac{1}{3}\right)^r\\
        &=\frac{1}{r}\left(-\frac{1}{3}\right)^r
        -\frac{1}{r-1}\left(-\frac{1}{3}\right)^{r-1}
        \end{align*}
        故由累加法,
        \[
        \sum_{r=2}^{n}\left[\frac{4r-1}{r(r-1)}\left(-\frac{1}{3}\right)^r\right]
        =\frac{1}{n}\left(-\frac{1}{3}\right)^n+\frac{1}{3}
        \]
        令$n\to\infty$,由于$\dfrac{1}{n}\left(-\dfrac{1}{3}\right)^n\to 0$,从而
        \[
        \sum_{r=2}^{\infty}\left[\frac{4r-1}{r(r-1)}\left(-\frac{1}{3}\right)^r\right]
        =\frac{1}{3}
        \]
    \end{solution}

    \question 求以下级数的和
    \[
    \frac{1}{4\cdot 2!} + \frac{1}{5\cdot 3!} + \frac{1}{6\cdot 4!} + \frac{1}{7\cdot 5!} + \frac{1}{8\cdot 6!} + \cdots
    \]
    \begin{solution}
        级数即
        \[
        \sum_{r=1}^{\infty} \frac{1}{(r+3)(r+1)!}
        \]
        发现
        \[
        \frac{1}{(r+3)(r+1)!} = \frac{1}{(r+2)!} - \frac{1}{(r+3)!}
        \]
        因此由累加法,
        \[
        \sum_{r=1}^{N} \frac{1}{(r+3)(r+1)!}= \frac{1}{3!} - \frac{1}{(N+3)!}
        \]
        令 $N \to \infty$,由于$\dfrac{1}{(N+3)!} \to 0$,则无穷级数为
        \[
        \sum_{r=1}^{\infty} \frac{1}{(r+3)(r+1)!} = \frac{1}{3!} = \frac{1}{6}
        \]
    \end{solution}

    \question 证明
    \[
    (2\cdot 1!)+(5\cdot 2!)+(10\cdot 3!)+(17\cdot 4!)+\cdots+(n^2+1)n!=n(n+1)!
    \]
    \begin{solution}
        先将数列写成求和形式
        \[
        S_n=\sum_{r=1}^{n}(r^2+1)r!
        \]
        尝试将 $(r^2+1)r!$ 表示成阶乘的差,注意到
        \[
        (r+2)!=(r+2)(r+1)r!=(r^2+3r+2)r!,
        \]
        因此
        \[
        (r+2)!-r!=(r^2+3r+1)r!
        \]
        又因为$(r+1)!-r!=r\cdot r!$,于是
        \[
        (r+2)!-r!=(r^2+1)r!+3r\cdot r!=(r^2+1)r!+3[(r+1)!-r!]
        \]
        即
        \[
        (r^2+1)r!=(r+2)!-3(r+1)!+2r!
        \]
        故由累加法,
        \begin{align*}
        \sum_{r=1}^{n}(r^2+1)r!
        &=\sum_{r=1}^{n}\big[(r+2)!-3(r+1)!+2r!\big]\\
        &=(n+2)!-2(n+1)! \\
        &=n(n+1)!
        \end{align*}
    \end{solution}

    \question 求
    \[
    \sum_{r=1}^{n}\left[\frac{r^2+9r+19}{(r+5)!}\right]
    \]
    的化简表达式,并据此求
    \[
    \sum_{r=1}^{\infty}\left[\frac{r^2+7r+11}{(r+4)!}\right]
    \]
    \begin{solution}
        注意到
        \[
        \frac{r^2+9r+19}{(r+5)!}
        =\frac{1}{(r+3)!}-\frac{1}{(r+5)!}
        \]
        由累加法,
        \begin{align*}
        \sum_{r=1}^{n}\left[\frac{r^2+9r+19}{(r+5)!}\right]
        &=\left(\frac{1}{4!}+\frac{1}{5!}\right)
        -\left(\frac{1}{(n+4)!}+\frac{1}{(n+5)!}\right)\\
        &=\frac{1}{20}-\left(\frac{1}{(n+4)!}+\frac{1}{(n+5)!}\right)
        \end{align*}
        改写变量,令 $k=r-1$,则
        \begin{align*}
        \sum_{r=1}^{n}\left[\frac{r^2+7r+11}{(r+4)!}\right]
        &=\sum_{k=0}^{n-1}\left[\frac{(k+1)^2+7(k+1)+11}{(k+5)!}\right]\\
        &=\sum_{k=0}^{n-1}\left[\frac{k^2+9k+19}{(k+5)!}\right]\\
        &= \frac{19}{5!} + \frac{1}{20}-\left(\frac{1}{(n+4)!}+\frac{1}{(n+5)!}\right)\\
        &= \frac{5}{24}-\frac{n+6}{(n+5)!}
        \end{align*}
        令 $n\to\infty$,因$\dfrac{n+6}{(n+5)!}\to 0$,故
        \[
        \sum_{r=1}^{\infty}\left[\frac{r^2+7r+11}{(r+4)!}\right]
        =\frac{5}{24}
        \]
    \end{solution}

    \question 求极限
    \[
    \lim_{n \to \infty} \sum_{k=1}^{n} \left( \sin  \frac{\pi}{2k} - \cos  \frac{\pi}{2k} - \sin  \frac{\pi}{2(k+2)} + \cos \frac{\pi}{2(k+2)} \right).
    \]
    \begin{solution}
        令 
        \[
        f(k) = \sin \left( \frac{\pi}{2k} \right) - \cos \left( \frac{\pi}{2k} \right)
        \]
        由累加法,部分和为 
        \begin{align*}
        S_n &= \sum_{k=1}^{n} [f(k) - f(k+2)] \\
        &= f(1) + f(2) - f(n+1) - f(n+2) \\
        &= 1-0 - \left( \sin \frac{\pi}{2(n+1)} - \cos \frac{\pi}{2(n+1)} \right) - \left( \sin \frac{\pi}{2(n+2)} - \cos \frac{\pi}{2(n+2)} \right)
        \end{align*}
        当 $n \to \infty$ 时,$\dfrac{\pi}{2(n+1)} \to 0,\dfrac{\pi}{2(n+2)} \to 0$:
        \begin{align*}
        \lim_{n \to \infty} S_n &= 1 - (\sin 0 - \cos 0) - (\sin 0 - \cos 0) \\
        &= 1 - (0 - 1) - (0 - 1) \\
        &= 3.
        \end{align*}
        故级数的极限值为 3。
    \end{solution}

    \question 求级数
    \[
    \sum_{n=0}^{\infty} \frac{1}{2011^{2^{n}} - 2011^{-2^{n}}}
    = \frac{1}{2011^{1} - 2011^{-1}} + \frac{1}{2011^{2} - 2011^{-2}} + \frac{1}{2011^{4} - 2011^{-4}} + \cdots
    \]
    \begin{solution}
        更一般地,设
        \[
        S(x) = \sum_{n=0}^{\infty} \frac{x^{2^{n}}}{1 - x^{2^{n+1}}}, \quad
        S_N(x) = \sum_{n=0}^{N} \frac{x^{2^{n}}}{1 - x^{2^{n+1}}}
        \]
        所求即 $S\left(\dfrac{1}{2011}\right)$。下面证明当 $0<x<1$ 时,
        \[
        S(x) = \frac{x}{1-x}
        \]
        注意到
        \[
        \frac{x^{2^{n}}}{1 - x^{2^{n+1}}}
        = \frac{1}{1 - x^{2^{n}}} - \frac{1}{1 - x^{2^{n+1}}},
        \]
        代入部分和 $S_N(x)$,由累加法
        \[
        S_N(x) = \frac{1}{1-x} - \frac{1}{1-x^{2^{N+1}}}
        \]
        当 $N \to \infty$ 时,由于 $0<x<1$,有 $x^{2^{N+1}} \to 0$,从而
        \[
        \lim_{N\to\infty} \frac{1}{1-x^{2^{N+1}}} = 1
        \]
        于是
        \[
        S(x) = \lim_{N\to\infty} S_N(x) = \frac{1}{1-x} - 1 = \frac{x}{1-x}
        \]
        因此
        \[
        \sum_{n=0}^{\infty} \frac{1}{2011^{2^{n}} - 2011^{-2^{n}}}
        = S\!\left(\frac{1}{2011}\right)
        = \frac{\frac{1}{2011}}{1-\frac{1}{2011}}
        = \frac{1}{2010}
        \]
    \end{solution}

    \question 求下列无穷乘积的值
    \[
    \left(\frac{7}{9}\right) \cdot \left(\frac{26}{28}\right) \cdot \left(\frac{63}{65}\right) \cdot \dots = \lim_{n \to \infty} \,\prod_{k=2}^{n} \frac{k^3 - 1}{k^3 + 1}
    \]
    \begin{solution}
        首先注意到部分和为
        \[
        \prod_{k=2}^{n} \frac{k^3 - 1}{k^3 + 1} = \prod_{k=2}^{n} \frac{k-1}{k+1} \;\prod_{k=2}^{n} \frac{k^2+k+1}{k^2-k+1}
        \]
        其中
        \[
        \prod_{k=2}^{n} \frac{k^2+k+1}{k^2-k+1} = \prod_{k=2}^{n} \frac{k^2+k+1}{(k-1)^2 + (k-1) + 1}
        \]
        同时
        \[
        \prod_{k=2}^{n} \frac{k-1}{k+1} = \frac{1\cdot 2}{3\cdot 4} \cdot \frac{2\cdot 3}{4\cdot 5} \cdots \frac{(n-1)n}{(n+1)n} = \frac{2}{n(n+1)}
        \]
        由累乘法,
        \[
        \prod_{k=2}^{n} \frac{k^3 - 1}{k^3 + 1} = \frac{2}{n(n+1)} \cdot \frac{n^2+n+1}{3} = \frac{2}{3} \left(1 + \frac{1}{n(n+1)}\right)
        \]
        因此
        \[
        \prod_{k=2}^{\infty} \frac{k^3 - 1}{k^3 + 1} = \lim_{n \to \infty} \frac{2}{3} \left(1 + \frac{1}{n(n+1)}\right) = \frac{2}{3}
        \]
    \end{solution}

    \question 计算无穷乘积
    \[
    \prod_{n=3}^\infty \frac{(n^3+3n)^2}{n^6-64}
    \]
    \begin{solution}
        设
        \[
        a_n=\frac{(n^3+3n)^2}{n^6-64}
        \]
        注意到
        \begin{align*}
        a_n
        &=\frac{(n^3+3n)^2}{(n^3-8)(n^3+8)}
        =\frac{n^2(n^2+3)^2}{(n-2)(n^2+2n+4)(n+2)(n^2-2n+4)} \\
        &=\frac{n}{n-2}\cdot\frac{n}{n+2}\cdot
        \frac{n^2+3}{(n-1)^2+3}\cdot
        \frac{n^2+3}{(n+1)^2+3}
        \end{align*}
        因此当 $N\ge3$ 时,
        \begin{align*}
        \prod_{n=3}^N a_n
        &=\left(\prod_{n=3}^N\frac{n}{n-2}\right)
        \left(\prod_{n=3}^N\frac{n}{n+2}\right)
        \left(\prod_{n=3}^N\frac{n^2+3}{(n-1)^2+3}\right)
        \left(\prod_{n=3}^N\frac{n^2+3}{(n+1)^2+3}\right) \\
        &=\frac{N(N-1)}{1\cdot2}\cdot
        \frac{3\cdot4}{(N+1)(N+2)}\cdot
        \frac{N^2+3}{2^2+3}\cdot
        \frac{3^2+3}{(N+1)^2+3} \\
        &=\frac{72}{7}\cdot
        \frac{N(N-1)(N^2+3)}{(N+1)(N+2)((N+1)^2+3)} \\
        &=\frac{72}{7}\cdot
        \frac{(1-\frac1N)(1+\frac{3}{N^2})}
        {(1+\frac1N)(1+\frac{2}{N})((1+\frac1N)^2+\frac{3}{N^2})}
        \end{align*}
        令 $N\to\infty$,得到
        \[
        \prod_{n=3}^\infty a_n
        =\lim_{N\to\infty}\prod_{n=3}^N a_n
        =\frac{72}{7}
        \]
    \end{solution}

    \question 设
    \[
    f(r)=\sum_{j=2}^{2013}\frac{1}{j^r}
    =\frac{1}{2^r}+\frac{1}{3^r}+\dots+\frac{1}{2013^r},
    \]
    求
    \[
    \sum_{k=2}^{\infty} f(k).
    \]
    \begin{solution}
        交换求和次序,并注意到内层和是收敛的等比级数,有
        \begin{align*}
        \sum_{k=2}^{\infty}\sum_{j=2}^{2013}\frac{1}{j^k}
        &=\sum_{j=2}^{2013}\sum_{k=2}^{\infty}\frac{1}{j^k}\\
        &=\sum_{j=2}^{2013}\frac{1/j^2}{1-1/j}\\
        &=\sum_{j=2}^{2013}\frac{1}{j(j-1)}\\
        &=\sum_{j=2}^{2013}\left(\frac{1}{j-1}-\frac{1}{j}\right)\\
        &=\frac{1}{1}-\frac{1}{2}+\frac{1}{2}-\frac{1}{3}+\dots+\frac{1}{2012}-\frac{1}{2013}\\
        &=\frac{2012}{2013}
        \end{align*}
    \end{solution}

    \question 试证 
    \[
    8 + 88 + 888 + \ldots + \underbrace{888\ldots88}_{\text{88's}} = \frac{8}{81}(10^{89} - 802)
    \] 
    \begin{solution}
        考虑到
        \[
        \underbrace{99\ldots99}_{\text{n's}}= 10^n - 1
        \]
        于是
        \begin{align*}
        8 + 88 + 888 + \cdots + \underbrace{888\ldots88}_{\text{88's}} 
        &= \frac{8}{9}(9 + 99 + 999 + \cdots + \underbrace{99\ldots99}_{\text{88's}}) \\
        &= \frac{8}{9}(10 + 10^2 + 10^3 + \cdots + 10^{88} - 88) \\
        &= \frac{8}{9}\left[\frac{10(10^{88}-1)}{10-1}-88\right] \\
        &= \frac{8}{81}(10^{89} - 802)
        \end{align*}
        故得证。
    \end{solution}

    \question 求无穷级数
    \[
    \frac{1^{2}}{11} + \frac{5^{2}}{11^{2}} + \frac{9^{2}}{11^{3}} + \frac{13^{2}}{11^{4}} + \cdots + \frac{(4n - 3)^{2}}{11^{n}} + \cdots
    \]
    \begin{solution}
        设原式为 
        \[
        S = \sum_{n=1}^{\infty}\frac{(4n-3)^2}{11^n}
        \]
        由 $(4n-3)^2 = 16n^2 - 24n + 9$,有
        \[
        S = 16\sum_{n=1}^{\infty}\frac{n^2}{11^n} - 24\sum_{n=1}^{\infty}\frac{n}{11^n} + 9\sum_{n=1}^{\infty}\frac{1}{11^n}
        \]
        当 $|x| < 1$ 时,有
        \[
        \sum_{n=1}^{\infty}x^n = \frac{x}{1-x}, \quad
        \sum_{n=1}^{\infty}nx^n = \frac{x}{(1-x)^2}, \quad
        \sum_{n=1}^{\infty}n^2x^n = \frac{x(1+x)}{(1-x)^3}
        \]
        令 $x = \dfrac{1}{11}$, 则
        \[
        \sum_{n=1}^{\infty}\frac{1}{11^n} = \frac{1}{10}, \quad \sum_{n=1}^{\infty}\frac{n}{11^n} = \frac{11}{100}, \quad \sum_{n=1}^{\infty}\frac{n^2}{11^n} = \frac{33}{250}
        \]
        于是
        \[
        S = 16 \cdot \frac{33}{250} - 24 \cdot \frac{11}{100} + 9 \cdot \frac{1}{10} = \frac{93}{250}
        \]
    \end{solution}
    \begin{solution}
        设
        \[
        S=\frac{1^{2}}{11} + \frac{5^{2}}{11^{2}} + \frac{9^{2}}{11^{3}} + \frac{13^{2}}{11^{4}} + \frac{17^{2}}{11^{5}} + \cdots \tag{1}
        \]
        则
        \[
        \frac{1}{11}S=\frac{1^{2}}{11^2} + \frac{5^{2}}{11^{3}} + \frac{9^{2}}{11^{4}} + \frac{13^{2}}{11^{5}} + \frac{17^{2}}{11^{6}} + \cdots \tag{2}
        \]
        (1) - (2) 得
        \[
        \frac{10}{11} S = \frac{1}{11} + \frac{5^2-1^2}{11^2} + \frac{9^2-5^2}{11^3} + \frac{13^2-9^2}{11^4} + \frac{17^2-13^2}{11^5} + \cdots \tag{3}
        \]
        \[
        \frac{10}{121} S = \frac{1}{11^2} + \frac{5^2-1^2}{11^3} + \frac{9^2-5^2}{11^4} + \frac{13^2-9^2}{11^5} + \frac{17^2-13^2}{11^6} + \cdots \tag{4}
        \]
        (3) - (4) 得
        \[
        \frac{100}{121} S = \frac{1}{11} + \frac{23}{11^2} + \frac{32}{11^3} + \frac{32}{11^4} + \frac{32}{11^5} + \cdots = \frac{1}{11} + \frac{23}{11^2} + \frac{32}{11^3} \cdot \frac{1}{1 - \frac{1}{11}}
        \]
        故可得
        \[
        S = \frac{93}{250}
        \]
    \end{solution}

    \question 已知
    \[
    1 + \frac{3}{x} + \frac{5}{x^2} + \frac{7}{x^3} + \frac{9}{x^4} + \cdots = 91,
    \]
    求
    \[
    S=1 + \frac{4}{x} + \frac{9}{x^2} + \frac{16}{x^3} + \frac{25}{x^4} + \cdots
    \]
    \begin{solution}
        发现
        \[
        91 = 1 + \frac{3}{x} + \frac{5}{x^2} + \frac{7}{x^3} + \cdots = \left(1 + \frac{1}{x} + \frac{1}{x^2} + \cdots \right) + 2\left( \frac{1}{x} + \frac{2}{x^2} + \frac{3}{x^3} + \cdots \right)
        \]
        考虑
        \[
        \sum_{n=1}^{\infty} n y^{n-1} =\left(\sum_{n=1}^{\infty} y^n\right)'=\left(\frac{y}{1-y}\right)'= \frac{1}{(1 - y)^2}
        \Rightarrow \sum_{n=1}^{\infty} n y^n = \frac{y}{(1 - y)^2}\]
        记 \(y = \dfrac{1}{x}\),则
        \[
        \sum_{n=1}^{\infty} \frac{n}{x^n} = \frac{1}{x} \cdot \frac{1}{(1 - \frac{1}{x})^2} = \frac{x}{(x - 1)^2}
        \]
        所以
        \[
        91= \frac{x}{x - 1} + \frac{2x}{(x - 1)^2} = \frac{x^2 + x}{(x - 1)^2} \Rightarrow 90x^2 - 183x + 91 = 0 \Rightarrow x = \frac{7}{6} \ \text{或} \ \frac{13}{15}
        \]
        级数收敛需满足 \(|x| > 1\),故取 \(x = \dfrac{7}{6}\),现考虑
        \[
        S=\sum_{n=1}^\infty \frac{n^2}{x^{n-1}}=\sum_{n=1}^{\infty} n^2 y^{n-1}
        \]
        由上,
        \[
        \sum_{n=1}^{\infty} n y^n = \frac{y}{(1 - y)^2}
        \Rightarrow \sum_{n=1}^{\infty} n^2 y^{n-1} = \left( \frac{y}{(1 - y)^2} \right)' = \frac{1 + y}{(1 - y)^3}
        \]
        取$y=\dfrac67$即得$S=637$;又解:发现$S=91+\dfrac{6}{7}S \Rightarrow S=637$
    \end{solution}

    \question 求
    \[
    \sum_{k=0}^{\infty} \frac{(-1)^k}{2k+1}
    \]
    \begin{solution}
        发现
        \begin{align*}
        1-\frac{1}{3}+\frac{1}{5}-\frac{1}{7}+\frac{1}{9}- \cdots
        &= \left[x-\frac{1}{3}x^3 +\frac{1}{5}x^5 -\frac{1}{7}x^7+\frac{1}{9}x^9-\cdots\right] _0^1 \\
        &=\int_0^1 (1-x^2+x^4-x^6+x^8+\cdots)\,dx \\
        &=\int_0^1 \frac{1}{1+x^2}\,dx \\
        &= \left[ \tan^{-1} x \right]_0^1 = \frac{\pi}{4}
        \end{align*}
    \end{solution}

    \question 求下列级数的值或证明其发散:
    \[
    \sum_{k=1}^{\infty} \arctan\left(-\frac{2}{k^{2}}\right)
    \]
    \begin{solution}
        由
        \[
        \arctan\left(-\frac{2}{k^{2}}\right)
        =\arctan\left(\frac{(k-1)-(k+1)}{1+(k-1)(k+1)}\right)
        =\arctan(k-1)-\arctan(k+1),
        \]
        知原级数是一列项求和,于是
        \begin{align*}
        &\sum_{k=1}^{\infty}\left[\arctan(k-1)-\arctan(k+1)\right] \\
        &=\lim_{N\to\infty}(\arctan 0+\arctan 1-\arctan N-\arctan(N+1)) \\
        &=0+\frac{\pi}{4}-\frac{\pi}{2}-\frac{\pi}{2} =-\frac{3\pi}{4}
        \end{align*}
        级数收敛,其值为$-\dfrac{3\pi}{4}$。
    \end{solution}

    \question 求
    \[
    \sum_{n=1}^{\infty} \frac{1}{2^n} \cdot \frac{1}{1 + 2^{2^{-n}}}
    \]
    \begin{solution}
        注意到:
        \[
        \frac{1}{1 - x^2} = \frac{1}{(1 - x)(1 + x)} = \frac{1}{2} \left( \frac{1}{1 - x} + \frac{1}{1 + x} \right)
        \]
        令 $x = 2^{2^{-n}}$,则
        \[
        \frac{1}{1 + 2^{2^{-n}}} = \frac{2}{1 - 2^{2^{1-n}}} - \frac{1}{1 - 2^{2^{-n}}}
        \]
        因此原式化为
        \begin{align*}
        &\sum_{n=1}^{\infty} \frac{1}{2^n} \cdot \left( \frac{2}{1 - 2^{2^{1-n}}} - \frac{1}{1 - 2^{2^{-n}}} \right)\\
        &= \sum_{n=1}^{\infty} \left( \frac{1}{2^{n-1}} \cdot \frac{1}{1 - 2^{2^{-(n-1)}}} - \frac{1}{2^n} \cdot \frac{1}{1 - 2^{2^{-n}}} \right)\\
        &= -1 - \lim_{N \to \infty} \left(\frac{1}{2^N} \cdot \frac{1}{1 - 2^{2^{-N}}} \right)
        \end{align*}
        现设 $x = 2^{-N}$,则当$N \to \infty, x \to 0^+$,且
        \[
        \frac{1}{2^N} \cdot \frac{1}{1 - 2^{2^{-N}}} = \frac{x}{1 - 2^x} = -\frac{x}{2^x - 1}
        \]        
        由洛必达法则,
        \[
        \lim_{x \to 0^+} \frac{x}{2^x - 1} = \lim_{x \to 0^+} \frac{1}{\ln 2 \cdot 2^x} = \frac{1}{\ln 2}
        \]
        故
        \[
        \sum_{n=1}^{\infty} \frac{1}{2^n} \cdot \frac{1}{1 + 2^{2^{-n}}} = -1 + \frac{1}{\ln 2}
        \]
    \end{solution}

    \question 设
    \[
    f(x) = \frac{2016^x}{2016^x + \sqrt{2016}},
    \]
    求
    \[
    \sum_{k=0}^{2016} f\left(\frac{k}{2016}\right) 
    \]
    \begin{solution}
        记$a=2016$,观察到
        \[
        f(x)+f(1-x)=\frac{a^{x}}{a^{x}+a^\frac12}+\frac{a^{1-x}}{a^{1-x}+a^\frac12}
        =\frac{a^{x}(a^{1-x}+a^{\frac12})+a^{1-x}(a^{x}+a^\frac12)}{(a^{x}+a^\frac12)(a^{1-x}+a^\frac12)}=1
        \]
        于是
        \[
        \sum_{k=0}^{2016} f\!\left(\frac{k}{2016}\right)
        =\underbrace{\sum_{k=0}^{1007}\left[f\!\left(\frac{k}{2016}\right)+f\!\left(1-\frac{k}{2016}\right)\right]}_{1008\text{ 个“1”}}+f\!\left(\frac{1008}{2016}\right)
        =1008+f\!\left(\frac12\right)
        \]
        而$f\!\left(\dfrac12\right)=\dfrac{a^\frac12}{a^\frac12+a^\frac12}=\dfrac12,$ 故
        \[
        \sum_{k=0}^{2016} f\!\left(\frac{k}{2016}\right)=1008+\frac12=\frac{2017}{2}
        \]
    \end{solution}

    \question 求级数
    \[
    \sum_{r=0}^{\infty}\frac{\sin^4(\pi 2^{r-2})}{4^r}
    \]
    \begin{solution}
        利用恒等式
        \[
        \sin^4\theta=\sin^2\theta-\frac{1}{4}\sin^2 2\theta
        \]
        考虑前 $n$ 项和:
        \[
        \sum_{r=0}^{n}\frac{\sin^4(\pi  2^{r-2})}{4^r}
        =\sum_{r=0}^{n}\left[
        \frac{1}{4^r}\sin^2(\pi 2^{r-2})
        -\frac{1}{4^{r+1}}\sin^2(\pi 2^{r-1})
        \right]
        \]
        由累加法,
        \[
        \sum_{r=0}^{n}\frac{\sin^4(\pi 2^{r-2})}{4^r}
        = \sin^2\frac{\pi }{4}
        -\frac{1}{4^{n+1}}\sin^2(\pi 2^{n-1})
        \]
        由于
        \[
        0\le \frac{1}{4^{n+1}}\sin^2(\pi 2^{n-1}) \le \frac{1}{4^{n+1}}
        \]
        且$\lim_{n \to \infty} \frac{1}{4^{n+1}}=0$,由夹挤定理,
        \[
        \lim_{n \to \infty} \frac{1}{4^{n+1}}\sin^2(\pi 2^{n-1}) = 0
        \]
        于是
        \[
        \sum_{r=0}^{\infty}\frac{\sin^4(\pi 2^{r-2})}{4^r}= \sin^2\frac{\pi}{4}= \frac{1}{2}
        \]
    \end{solution}

    \question 求下列级数的和
    \[
    \sum_{n=0}^{\infty} \frac{\cos^2\frac{n\pi}{6}}{2^n}
    \]
    \begin{solution}
        先由三角恒等式降幂,再由欧拉公式,
        \begin{align*}
        S=\sum_{n=0}^{\infty} \frac{\cos^2\frac{n\pi}{6}}{2^n} 
        &= \frac{1}{2} \sum_{n=0}^{\infty} \frac{1+\cos\frac{n\pi}{3}}{2^n}\\
        &=\frac{1}{2} \left(\sum_{n=0}^{\infty} \frac{1}{2^n} + \sum_{n=0}^{\infty} \Re \left[\frac{e^{in\pi/3}}{2^n} \right]\right) \\
        &=\frac{1}{2} \left(\frac{1}{1-\frac{1}{2}}+\Re \left[\frac{1}{1 - \frac{e^{i\pi/3}}{2}}\right]\right)
        \end{align*}
        其中
        \begin{align*}
        \frac{1}{1 - \frac{e^{i\pi/3}}{2}} 
        &= \frac{2 - e^{-i\pi/3}}{4 - 2(e^{i\pi/3} + e^{-i\pi/3}) + 1} \\
        &= \frac{2 - e^{-i\pi/3}}{5 - 4\cos\frac{\pi}{3}} \\
        &= \frac{2 - (\frac{1}{2} - i\frac{\sqrt{3}}{2})}{3} \\
        &= \frac{1}{2}+\frac{\sqrt{3}}{6}i
        \end{align*}
        故
        \[
        S = \frac{1}{2}\left( 2 + \frac{1}{2}\right) = \frac{3}{2}
        \]
    \end{solution}

    \question 求 
    \[
    \sum_{n=1}^{\infty}\frac{n}{2^{n}}\cos\frac{n\pi}{3}
    \]
    \begin{solution}
        设
        \[
        z=\frac{1}{2}e^{i\pi/3}=\frac{1}{4}(1+\sqrt 3\, i)
        \]
        且
        \[
        f(z) =\sum_{n=1}^\infty z^n = \frac{z}{1-z}
        \]
        考虑
        \[
        \sum_{n=1}^\infty \frac{n}{2^n} e^{n\pi i/3} = \sum_{n=1}^\infty n z^{n} = \frac{z}{(1-z)^2} = zf'(z)
        \]
        由欧拉公式,
        \[
        \sum_{n=1}^\infty \frac{n}{2^n}\cos\frac{n\pi}{3}=\Re(zf'(z)) 
        \]
        其中
        \[
        zf'(z)
        = \frac{\frac{1}{4}(1+\sqrt{3}\,i)}{\left(\frac{3}{4}-\frac{\sqrt 3}{4} i\right)^2}
        = \frac{2}{3}\cdot \frac{1+\sqrt 3\, i}{1-\sqrt 3\, i}
        = -\frac{1}{3}+\frac{\sqrt 3}{3} i
        \]
        故得
        \[
        \sum_{n=1}^\infty \frac{n}{2^n}\cos\frac{n\pi}{3} = -\frac{1}{3}
        \]
    \end{solution}

    \question 证明不等式
    \[ 
    \frac{1}{2} < \frac{1}{n+1} + \frac{1}{n+2} + \cdots + \frac{1}{2n} < \frac{3}{4}
    \]
    \begin{solution}
        由放缩法,
        \[ 
        \frac{1}{n+1} + \frac{1}{n+2} + \cdots + \frac{1}{2n} > \underbrace{\frac{1}{2n} + \frac{1}{2n} + \cdots + \frac{1}{2n}}_{n\ \text{个}} = \frac{1}{2} 
        \]
        故左侧不等式得证,同由放缩法,且注意到
        \begin{align*}
        \frac{1}{n} + \frac{1}{n+1} + \cdots + \frac{1}{2n} 
        &= \frac{1}{2} \left[ \left(\frac{1}{n} + \frac{1}{2n}\right) + \left(\frac{1}{n+1} + \frac{1}{2n-1}\right) + \cdots + \left(\frac{1}{2n} + \frac{1}{n}\right) \right] \\
        &= \frac{1}{2}  \left[\frac{3n}{2n^2} + \frac{3n}{2n^2 + (n-1)} + \frac{3n}{2n^2 + 2(n-2)} + \cdots + \frac{3n}{2n^2}\right] \\
        &< \frac{1}{2} \underbrace{\left[ \frac{3n}{2n^2} + \frac{3n}{2n^2} + \cdots + \frac{3n}{2n^2} \right]}_{(n+1)\ \text{个}} \\
        &= \frac{1}{2} (n+1) \frac{3}{2n} \\
        &= \frac{3}{4} + \frac{1}{4n} \\
        &< \frac{3}{4} + \frac{1}{n}
        \end{align*}
        于是右侧不等式得证。
    \end{solution}

    \question 计算 
    \[
    1 + \frac{1}{\sqrt{2}} + \frac{1}{\sqrt{3}} + \frac{1}{\sqrt{4}} + \cdots + \frac{1}{\sqrt{1,000,000}}
    \]
    的整数部分。
    \begin{solution}
        由放缩法,注意到
        \[ 
        2\sqrt{n+1} - 2\sqrt{n} = \frac{2(n+1-1)}{\sqrt{n+1} + \sqrt{n}} < \frac{2}{\sqrt{n} + \sqrt{n}} = \frac{1}{\sqrt{n}} \tag{1}
        \]
        同理,
        \[
        \frac{1}{\sqrt{n}} = \frac{2}{\sqrt{n} + \sqrt{n}} < \frac{2(n-(n-1))}{\sqrt{n} + \sqrt{n-1}} = 2\sqrt{n} - 2\sqrt{n-1} \tag{2}
        \]
        对于下界,由(1)得
        \begin{align*}
        &1 + \frac{1}{\sqrt{2}} + \cdots + \frac{1}{\sqrt{1,000,000}} \\
        &> 1 + 2 \left[(\sqrt{3} - \sqrt{2}) + (\sqrt{4} - \sqrt{3}) + \cdots + (\sqrt{1,000,001} - \sqrt{1,000,000})\right] \\
        &= 1 + 2(\sqrt{1,000,001} - \sqrt{2}) \\
        &> 2 \cdot 1000 - \sqrt{8} + 1 \\
        &> 2000 - 3 + 1 = 1998
        \end{align*}
        对于上界,由(2)得
        \begin{align*}
        &1 + \frac{1}{\sqrt{2}} + \cdots + \frac{1}{\sqrt{1,000,000}} \\
        &< 1 + 2 \left[(\sqrt{2} - 1) + (\sqrt{3} - \sqrt{2}) + \cdots + (\sqrt{1,000,000} - \sqrt{999,999})\right] \\
        &= 1 + 2(\sqrt{1,000,000} - 1) \\
        &= 1 + 2 \cdot 999 = 1999
        \end{align*}
        综上,该数的整数部分为 1998。
    \end{solution}

    \question 求
    \[
    \frac{1}{\sqrt[3]{4}} + \frac{1}{\sqrt[3]{5}} + \frac{1}{\sqrt[3]{6}} + \cdots + \frac{1}{\sqrt[3]{1,000,000}}
    \]
    的整数部分。
    \begin{solution}
        需要前置作业,观察到二项展开为
        \[
        \left(1 + \frac{1}{n}\right)^2 = 1 + \frac{2}{n} + \frac{1}{n^2},\quad \left(1 + \frac{2}{3} \cdot \frac{1}{n}\right)^3 = 1 + \frac{2}{n} + \frac{4}{3n^2} + \frac{8}{27n^3}
        \]
        由上可知
        \[ 
        \left(1 + \frac{2}{3} \cdot \frac{1}{n}\right)^3 > \left(1 + \frac{1}{n}\right)^2 
        \]
        两边开立方根并乘以$n^{\frac{2}{3}}$,得
        \[
        n^{\frac{2}{3}} + \frac{2}{3} n^{-\frac{1}{3}} > (n+1)^{\frac{2}{3}}
        \]
        即
        \[ 
        \frac{1}{\sqrt[3]{n^2}} > \frac{3}{2} \left[ \sqrt[3]{(n+1)^2} - \sqrt[3]{n^2} \right] 
        \]
        同理,
        \[
        \left( 1 - \frac{2}{3n} \right)^3 
        = 1 - \frac{2}{n} + \frac{4}{3n^2} - \frac{8}{27n^3} > 1 - 2\frac{1}{n} + \frac{1}{n^2} 
        = \left( 1 - \frac{1}{n} \right)^2
        \]
        其中 
        \[
        \frac{1}{3n^2} - \frac{8}{27n^3} = \frac{1}{3n^2} - \frac{1}{3n^3} \ge 0,
        \]
        于是
        \[
        1 - \frac{2}{3n} > \left( 1 - \frac{1}{n} \right)^{\frac{2}{3}}
        \]
        \[
        n^{\frac{2}{3}} - \frac{2}{3}n^{-1/3} > (n - 1)^{\frac{2}{3}}
        \]
        \[ 
        \frac{1}{\sqrt[3]{n^2}} < \frac
        {3}{2} \left[ \sqrt[3]{n^2} - \sqrt[3]{(n-1)^2} \right] 
        \]
        由此对原和式进行估计。对于下界:
        \begin{align*}
        \frac{1}{\sqrt[3]{4}} + \cdots + \frac{1}{\sqrt[3]{1,000,000}} 
        &> \frac{3}{2} \left[(\sqrt[3]{5^2} - \sqrt[3]{4^2}) + \cdots + (\sqrt[3]{1,000,001^2} - \sqrt[3]{1,000,000^2})\right] \\
        &= \frac{3}{2} (\sqrt[3]{1,000,002,000,001} - \sqrt[3]{16}) \\
        &> \frac{3}{2} \cdot 10,000 - \sqrt[3]{54} \\
        &> 15,000 - 4 = 14,996
        \end{align*}   
        对于上界:
        \begin{align*}
        \frac{1}{\sqrt[3]{4}} + \cdots + \frac{1}{\sqrt[3]{1,000,000}} 
        &< \frac{3}{2} \left[(\sqrt[3]{4^2} - \sqrt[3]{3^2}) + \cdots + (\sqrt[3]{1,000,000^2} - \sqrt[3]{999,999^2})\right] \\
        &= \frac{3}{2} (\sqrt[3]{1,000,000,000,000} - \sqrt[3]{9}) \\
        &< \frac{3}{2} (10,000 - 2) = 14,997
        \end{align*}
        综上,该数的整数部分等于 14,996。
    \end{solution}

    \question 求
    \[
    \sum_{m=1}^{19} \sum_{n=m}^{19} (2m+n).
    \]
    \begin{solution}
        将求和顺序交换,得
        \begin{align*}
        S &=\sum_{m=1}^{19} \sum_{n=m}^{19} (2m+n) \\
        &= \sum_{n=1}^{19} \sum_{m=1}^{n} (2m+n)\\
        &= \sum_{n=1}^{19} \left[ 2 \sum_{m=1}^{n} m + \sum_{m=1}^{n} n \right] \\
        &= \sum_{n=1}^{19} [2 \cdot \frac{n(n+1)}{2} + n^2] \\ 
        &= \sum_{n=1}^{19} (2n^2+n) \\
        &= 2 \cdot \frac{19 \cdot 20 \cdot 39}{6} + \frac{19 \cdot 20}{2} \\
        &= 5130
        \end{align*}
    \end{solution}

    \question 令
    \[
    S(x) = \sum_{k=1}^\infty \frac{(-1)^k x^{k+1}}{3^k(k+1)},
    \]
    \begin{parts}
    \part 试证
    \[
    S'(x) = \frac{3}{3 + x} - 1.
    \]
    \begin{solution}
        对各项微分,有
        \[
        S'(x) = \sum_{k=1}^\infty \frac{(-1)^k x^k}{3^k} = \sum_{k=1}^\infty \left(-\frac{x}{3}\right)^k= \frac{-\frac{x}{3}}{1+\frac{x}{3}}= \frac{-x}{3+x}= \frac{3}{3+x} - 1
        \]
    \end{solution}
    \part 据此,证
    \[
    S(x) = 3 \ln(3 + x) - 3 \ln 3 - x.
    \]
    \begin{solution}
        \[
        S(x) = \int \left( \frac{3}{3+x} - 1 \right) dx = 3\ln(3+x) - x + C
        \]
        由 \( S(0)=0 \)可得 \( C = -3\ln 3 \),故
        \[
        S(x) = 3\ln(3+x) - 3\ln 3 - x
        \]
    \end{solution}
    \part 求
    \[
    \sum_{k=1}^\infty \frac{(-1)^k}{k \cdot 3^k}
    \]
    \begin{solution}改写求和,
        \begin{align*}
        \sum_{k=1}^\infty \frac{(-1)^k}{k \cdot 3^k}
        &= \sum_{k=0}^\infty \frac{(-1)^{k+1}}{(k+1)3^{k+1}}\\
        &= -\frac{1}{3}\sum_{k=0}^\infty \frac{(-1)^{k}}{(k+1)3^{k}}\\
        &= -\frac{1}{3} (1+S(1)) \\
        &= -\frac{1}{3}(1+3 \ln4 - 3 \ln 3 - 1)\\
        &= \ln\frac{3}{4}
        \end{align*}
    \end{solution}
    \end{parts}

    \question 
    \begin{parts}
    \part 定义
    \[
    f(n)=\ln\left(\frac{n+1}{n}\right),\quad g(n)=\sum_{k=n}^{2n-1} [f(k)]^3, \quad n\ge1,
    \]
    证明对于任意正整数 $n$,都有
    \[
    g(n)-g(n+1) = 3f(n)f(2n)f(2n+1)
    \]
    \begin{solution}
        已知
        \[
        f(n)=\ln\left(\frac{n+1}{n}\right),\quad n\ge1
        \]
        注意到
        \[
        f(2n)+f(2n+1)= \ln\left(\frac{2n+1}{2n}\cdot \frac{2n+2}{2n+1}\right) = \ln\left(\frac{n+1}{n}\right) = f(n)
        \]
        因此
        \begin{align*}
        g(n)-g(n+1)
        &=[f(n)]^3-[f(2n)]^3-[f(2n+1)]^3 \\
        &=[f(2n)+f(2n+1)]^3-[f(2n)]^3-[f(2n+1)]^3 \\
        &=3(f(2n)+f(2n+1))f(2n)f(2n+1) \\
        &=3f(n)f(2n)f(2n+1)
        \end{align*}
    \end{solution}
    \part 据此,求
    \[
    \sum_{n=1}^\infty 
    \ln\left(1+\frac{1}{n}\right)
    \ln\left(1+\frac{1}{2n}\right)
    \ln\left(1+\frac{1}{2n+1}\right)
    \]
    的值。
    \begin{solution}
        由累加法,部分和即
        \[
        \sum_{n=1}^N f(n)f(2n)f(2n+1)
        =\frac{1}{3}\sum_{n=1}^N\bigl(g(n)-g(n+1)\bigr)
        =\frac{1}{3}\bigl(g(1)-g(N+1)\bigr).
        \]
        由不等式 $\ln(1+x)\le x$ 可得
        \[
        f(n)= \ln\left(1+\frac{1}{n}\right)\le\frac{1}{n}
        \]
        由放缩法,知
        \[
        g(n)<n [f(n)]^3\le\frac{1}{n^2}
        \]
        因此当 $N\to\infty$ 时,$g(N+1)\to0$,于是
        \[
        \sum_{n=1}^\infty f(n)f(2n)f(2n+1)
        =\frac{1}{3}g(1)
        =\frac{1}{3}(\ln 2)^3
        \]
        这正是所求级数的值。
    \end{solution}
    \end{parts}

    \question 考虑正整数 $m, n$,且 $m \ge 2$。$(m,n)$-锯齿数列是从 1 开始的连续整数序列,有 $n$ 个“齿”,每个齿从 2 上升到 $m$ 再下降到 1,如$(3,4)$-锯齿数列为
    \[
    \begin{array}{ccccccccccccccccccccc}
    &&3&&&&3&&&&3&&&&3&& \\
    &2&&2&&2&&2&&2&&2&&2&&2& \\
    1&&&&1&&&&1&&&&1&&&&1 \\
    \end{array}
    \]
    该序列共有 17 项,平均数为 $\dfrac{33}{17}$。
    \begin{parts}
    \part 求 $(4,2)$-锯齿数列的项和。
    \begin{solution}
        $(4,2)$-锯齿数列的项为
        \[
        1,2,3,4,3,2,1,2,3,4,3,2,1
        \]
        其和为 $31$。
    \end{solution}
    \part 对任意正整数 $m \ge 2$,求 $(m,3)$-锯齿数列中所有数字之和的通项。
    \begin{solution}
        $(m,3)$-锯齿数列由初始 1 和 3 个齿组成,每个齿为 $2,3,\dots,m-1,m,m-1,\dots,2,1$,项和为
        \[
        2 + 3 + \dots + m + (m-1) + \dots + 1 = 2(1+2+\dots+m)-1-m = m^2-1.
        \]
        因此序列和为 $1 + 3(m^2-1) = 3m^2 - 2$。
    \end{solution}
    \part 求所有使 $(m,n)$-锯齿数列的项和为 145 的 $(m,n)$ 序对。
    \begin{solution}
        每个齿的和为 $m^2-1$,整个序列和为
        \[
        1 + n(m^2-1) = 145 \Rightarrow n(m^2-1) = 144.
        \]
        考虑$144$的因数分解,当$m \ge 2$ 时,$m^2-1$可能值为
        \[
        3,8,15,24,35,48,63,80,99,120,143,
        \]
        其中 $3,8,24,48$ 能整除 144,因此符合条件的序对为
        \[
        (m,n) = (2,48), (3,18), (5,6), (7,3).
        \]
    \end{solution}
    \part 证明对于所有正整数对 $(m,n)$ 且 $m \ge 2,(m,n)$-锯齿数列的平均数不是整数。
    \begin{solution}
        平均数为
        \[
        \frac{1 + n(m^2-1)}{1 + n(2m-2)}
        \]
        假设平均数为整数 $k$,可得一关于$m$的二次方程式
        \[
        \frac{1 + n(m^2-1)}{1 + n(2m-2)} = k \Rightarrow m^2 n - 2mn k + (2nk - n - k + 1) = 0
        \]
        由于$m$为整数,故判别式
        \[
        \Delta = (-2nk)^2 - 4n(2nk - n - k +1) = (2n(k-1)+1)^2 - 1.
        \]
        必须为完全平方数,且两完全平方数差为 1 的只有当 0 和 1,因此 
        \[
        2n(k-1)+1 = 1 \Rightarrow k=1
        \]
        但若平均数为 1,则
        \[
        n(m^2-1)+1 = n(2m-2)+1 \Rightarrow n(m-1)^2 = 0,
        \]
        得 $n>0$ 且 $m \ge 2$,矛盾,因此$(m,n)$-锯齿数列的平均数不可能为整数。
    \end{solution}
    \end{parts}

    \question 已知正整数 $n$,考虑一个边长为 $n$ 、顶点朝上的等边三角形,将其划分为单位等边三角形。对于每个 $n$,记 $f(n)$ 为所有大小的顶点朝下等边三角形的总数,例如$f(3)=3,f(4)=7$。
    %2008 euclid P10
    \begin{figure}[H]
        \centering        
        \includegraphics[width=0.6\textwidth]{images/image197.png}
    \end{figure}
    \begin{figure}[H]
        \centering        
        \includegraphics[width=0.6\textwidth]{images/image198.png}
    \end{figure}
    \begin{parts}
    \part 求 $f(5)$ 和 $f(6)$。
    \begin{solution}
        当 $n=5$ 时,有$1+2+3+4=10$个边长为1、$1+2=3$个边长为 2 的三角形,所以
        \[
        f(5) = 10 + 3 = 13
        \]
        当 $n=6$ 时,有$1+2+3+4+5=15$个边长为 1、$1+2+3=6$个边长为2、$1$个边长为 3 的三角形 ,所以
        \[
        f(6) = 15 + 6 + 1 = 22
        \]
    \end{solution}
    \part 证明对每个正整数 $k\ge 1$,有 $f(2k) = f(2k-1) + k^2$。
    \ifprintanswers
    \begin{figure}[H]
        \centering        
        \includegraphics[width=0.4\textwidth]{images/image199.png}
    \end{figure}
    \fi
    \begin{solution}
        设第 $i$ 行有 $i+1$ 个单位点,从左至右记为$0,1,\cdots,i$,考虑边长为 $m$ 、顶点朝下的等边三角形,且其底顶点位于行 $i$ 时,注意到每个这样的三角形都被其底顶点确定,故可能的底点数为
        \[
        (i-m)-m+1 = i+1-2m,
        \]
        要求 $2m \le i \le n$ 才能形成三角形,因此边长为 $m$ 、顶点朝下的三角形总数为
        \[
        \sum_{i=2m}^{n} (i+1-2m) = \frac{(n+1-2m)(n+2-2m)}{2}
        \]
        若 $n=2k$,则 $m=1,2,\dots,k$,所以
        \begin{align*}
        f(2k) &= \sum_{m=1}^{k} \frac{(2k+1-2m)(2k+2-2m)}{2} \\
        &= \sum_{l=1}^{k} (2l-1)l \quad (\text{令 } l=k+1-m) \\
        &= 2\sum_{l=1}^{k} l^2 - \sum_{l=1}^{k} l \\
        &= 2\cdot \frac{k(k+1)(2k+1)}{6} - \frac{k(k+1)}{2} \\
        &= \frac{k(k+1)(4k-1)}{6}
        \end{align*}
        若 $n=2k-1$,则 $m=1,2,\dots,k-1$,所以
        \begin{align*}
        f(2k-1) &= \sum_{m=1}^{k-1} \frac{(2k-2m)(2k+1-2m)}{2} \\
        &= \sum_{l=1}^{k-1} l(2l+1) \quad (\text{令 } l=k-m) \\
        &= 2\sum_{l=1}^{k-1} l^2 + \sum_{l=1}^{k-1} l \\
        &= 2\cdot \frac{(k-1)k(2k-1)}{6} + \frac{(k-1)k}{2} \\
        &= \frac{k(k-1)(4k+1)}{6}
        \end{align*}
        因此,
        \[
        f(2k) - f(2k-1) = \frac{k(k+1)(4k-1)}{6} - \frac{k(k-1)(4k+1)}{6} = k^2
        \]
        证毕。
    \end{solution}
    \end{parts}
    
    \question 如下图,$O(0,0), A_1(8,0),A_1A_2$ 与 $x$ 轴正向夹 $45^\circ$ 角,又 
    \[
    A_1A_2 \parallels A_3A_4 \parallels A_5A_6 \parallels \dots,\quad OA_1 \parallels A_2A_3 \parallels A_4A_5 \parallels \dots.
    \]
    已知 $A_1A_2 = 8$,且对所有 $k \in \mathbb{N}$,有 
    \[
    A_kA_{k+1} = 2A_{k+1}A_{k+2}.
    \]
    若点 $A_n$ 的坐标为 $(x_n, y_n)$,求 $\displaystyle \lim_{n \to \infty} (x_n + y_n)$。
    \begin{figure}[H]
        \centering        
        \includegraphics[width=0.4\textwidth]{images/image90.png}
    \end{figure}
    \begin{solution}
        令 $a_n = A_nA_{n+1} = 2^{4-n}, n \in \mathbb{N},$有
        \begin{align*}
        x_\infty &= 8 + \frac{a_1}{\sqrt 2} - a_2 - \frac{a_3}{\sqrt 2} + a_4 + \frac{a_5}{\sqrt 2} - a_6 - \frac{a_7}{\sqrt 2} + a_8 + \cdots \\
        &= 8 + (a_4 + a_8 + a_{12} + \cdots) + \frac{1}{\sqrt 2}(a_1 + a_5 + a_9 + \cdots) \\
        &\quad - \frac{1}{\sqrt 2}(a_3 + a_7 + a_{11} + \cdots) - (a_2 + a_6 + a_{10} + \cdots) \\
        &= 8 + \frac{1}{1 - \frac{1}{16}} + \frac{1}{\sqrt 2} \cdot \frac{8}{1 - \frac{1}{16}} - \frac{1}{\sqrt 2} \cdot \frac{2}{1 - \frac{1}{16}} - \frac{4}{1 - \frac{1}{16}} = 8 + \frac{16}{5} (\sqrt 2 - 1) \\
        y_\infty &= \frac{a_1}{\sqrt 2} - \frac{a_3}{\sqrt 2} + \frac{a_5}{\sqrt 2} - \frac{a_7}{\sqrt 2} + \cdots \\
        &= \frac{1}{\sqrt 2} (a_1 + a_5 + a_9 + \cdots) - \frac{1}{\sqrt 2} (a_3 + a_7 + a_{11} + \cdots) \\
        &= \frac{1}{\sqrt 2} \cdot \frac{8}{1 - \frac{1}{16}} - \frac{1}{\sqrt 2} \cdot \frac{2}{1 - \frac{1}{16}} = \frac{16 \sqrt 2}{5} 
        \end{align*}
        所以
        \[
        x_\infty + y_\infty = \frac{24 + 32 \sqrt 2}{5}
        \]
    \end{solution}
\end{questions}
\pagebreak

\begin{center}
  {\fontsize{30pt}{26pt}\selectfont
    \hypertarget{二项展开式}{二项展开式} \label{二项展开式}
  }
\end{center}
\separator
\vspace{1pt}

\begin{questions}
    \question 求\((x^2 + x + y)^5\) 展开式中\( x^5 y^2 \) 的系数。  
    \begin{solution}
        系数为
        \[
        \comb{5}{2}\comb{3}{1}\comb{2}{2}=30
        \]
    \end{solution}

    \question 求\(\left(x \sqrt{y} - y \sqrt{x}\right)^4\) 的展开式中\( xy \) 的系数。
    \begin{solution}
        观察发现
        \[
        (x\sqrt{y}-y\sqrt{x})^4 = x^2y^2(\sqrt{x}-\sqrt{y})^4,
        \]
        即$(\sqrt{x}-\sqrt{y})^4$展开式中含 $xy$ 项的系数 $\comb{4}{2}=6$。
    \end{solution}

    \question 求在 $\left(x+\dfrac{4}{x}+4\right)^6$ 的展开式中$x^4$ 的系数。
    \begin{solution}
        原式可化为
        \[
        \left(x + \dfrac{4}{x} + 4\right)^6 = \left(\dfrac{x^2 + 4x + 4}{x}\right)^6 = \dfrac{(x+2)^{12}}{x^6}
        \]
        即求 $(x+2)^{12}$ 展开式中$x^{10}$ 的系数
        \[
        \comb{12}{2} \cdot 2^2 = 264
        \]
    \end{solution}

    \question 求 \( (1 + x^2) + (1 + x^2)^2 + \cdots + (1 + x^2)^{10} \) 展开式中 \( x^6 \) 的系数。  
    \begin{solution}
        原式为
        \[
        \frac{(1+x^2) \left[ (1+x^2)^{10}-1 \right]}{(1+x^2)-1} = \frac{(1+x^2)^{11}-(1+x^2)}{x^2}
        \]
        而展开式$(1+x^2)^{11}$中 $x^8$ 系数为 
        \[
        \comb{10}{4} = \dfrac{11 \cdot 10 \cdot 9 \cdot 8}{4 \cdot 3 \cdot 2} = 330
        \]
        故原式中 $x^6$ 系数为 $330$。
    \end{solution}

    \question 求$(1+x)^{1000} + x(1+x)^{999} + x^2(1+x)^{998} + \dots + x^{1000}$ 展开式中 $x^{50}$ 的系数。
    \begin{solution}
        由等比数列求和公式,原式可化为
        \begin{align*}
        &(1+x)^{1000} + x(1+x)^{999} + \dots + x^{1000} \\
        &= \frac{(1+x)^{1001} - x^{1001}}{(1+x) - x} \\
        &= (1+x)^{1001} - x^{1001} \\
        &= 1 + \comb{1001}{1}x + \comb{1001}{2}x^2 + \dots + \comb{1001}{1001}x^{1001} - x^{1001}
        \end{align*}
        所求的系数为
        \[ 
        \comb{1001}{50} = \frac{1001!}{50!\,951!} 
        \]
    \end{solution}

    \question 求$(1+x) + 2(1+x)^2 + 3(1+x)^3 + \dots + 1000(1+x)^{1000}$ 展开式中 $x^{50}$ 的系数。
    \begin{solution}
        设已知多项式为 $P(x)$,可以写出
        \begin{align*}
        x P(x) &= (1+x)P(x) - P(x) \\
        &= \left[(1+x)^2 + 2(1+x)^3 + \dots + 999(1+x)^{1000} + 1000(1+x)^{1001}\right] \\
        &\quad - \left[(1+x) + 2(1+x)^2 + 3(1+x)^3 + \dots + 1000(1+x)^{1000}\right] \\
        &= 1000(1+x)^{1001} - \left[(1+x) + (1+x)^2 + (1+x)^3 + \dots + (1+x)^{1000}\right] \\
        &= 1000(1+x)^{1001} - \frac{(1+x)[(1+x)^{1000} - 1]}{(1+x) - 1} \\
        &= 1000(1+x)^{1001} - \frac{(1+x)^{1001} - (1+x)}{x}
        \end{align*}
        从而得出
        \[ 
        P(x) = \frac{1000(1+x)^{1001}}{x} - \frac{(1+x)^{1001} - (1+x)}{x^2} 
        \]
        $P(x)$ 中 $x^{50}$ 的系数即为 $\dfrac{1000(1+x)^{1001}}{x}$ 中 $x^{50}$ 的系数减去 $\dfrac{(1+x)^{1001}}{x^2}$ 中 $x^{50}$ 的系数,即 $(1+x)^{1001}$ 中 $x^{51}$ 的系数减去$(1+x)^{1001}$ 中 $x^{52}$ 的系数:
        \[
        1000 \comb{1001}{51} - \comb{1001}{52} = \frac{1000 \cdot 1001!}{51!\,950!} - \frac{1001!}{52!\,950!} = \frac{51050 \cdot 1001!}{52!\,950!}
        \]
    \end{solution}

    \question 求\(\left[a + (b + c)^2\right]^8\) 展开式中 \( a^5 b^4 c^2 \) 的系数。
    \begin{solution}
        令$(b+c)^2=x,$展开式为
        $$(a+x)^8 = \sum_{k=0}^{8} \frac{8!}{k!(8-k)!} a^k x^{8-k} $$
        欲求$a^5b^4c^2$ 系数,于是$k=5$,有
        $$\frac{8!}{5!3!} a^5 [(b+c)^2]^3 = 56a^5 (b+c)^6 $$
        观察
        $$(b+c)^6 = \sum_{t=0}^{6} \frac{6!}{t!(6-t)!} b^t c^{6-t}$$
        欲求$b^4c^2$ 系数,于是 $t=4$,有
        $$\frac{6!}{4!2!} b^4 c^2 = 15b^4 c^2 $$
        $\therefore a^5b^4c^2$ 系数为 $56 \cdot 15 = 840$
    \end{solution}

    \question 求 $(-xy+2x+3y-6)^{6}$的展开式中$x^3y^3$的系数。
    \begin{solution}
        考虑
        \[
        \frac{6!}{n_1!n_2!n_3!n_4!} (-xy)^{n_1}(2x)^{n_2}(3y)^{n_3}(-6)^{n_4}=Ax^3y^3,
        \]
        其中$n_1+n_2+n_3+n_4=6$
        \begin{itemize}
        \item 若 \(n_1=0\),则 \(n_2=3, n_3=3, n_4=0\),
        \[
        \frac{6!}{0!3!3!0!}(-xy)^0(2x)^3(3y)^3(-6)^0 = 4320x^3y^3
        \]
        \item 若 \(n_1=1\),则 \(n_2=2, n_3=2, n_4=1\),
        \[
        \frac{6!}{1!2!2!1!}(-xy)^1(2x)^2(3y)^2(-6)^1 = 38880x^3y^3
        \]
        \item 若 \(n_1=2\),则 \(n_2=1, n_3=1, n_4=2\),
        \[
        \frac{6!}{2!1!1!2!}(-xy)^2(2x)^1(3y)^1(-6)^2 = 38880x^3y^3
        \]
        \item 若 \(n_1=3\),则 \(n_2=0, n_3=0, n_4=3\),
        \[
        \frac{6!}{3!0!0!3!}(-xy)^3(2x)^0(3y)^0(-6)^3  = 4320x^3y^3
        \]
        \end{itemize}
        因此$x^3y^3$的系数为$4320 + 38880 + 38880 + 4320 = 86400$
    \end{solution}

    \question \(\left(x + \dfrac{a}{x}\right) \left(2x - \dfrac{1}{x}\right)^5\) 的展开式中各项系数的和为 2,则该展开式中的常数项为  
    \begin{solution}
        \textbf{解法一}
        
        令$x=1$ 得 $a=1$.故原式为 $$(x+\frac{1}{x})(2x-\frac{1}{x})^5$$ 
        通项
        $$T_{r+1}=\comb{5}{r}(2x)^{5-r}(-x^{-1})^r=\comb{5}{r}(-1)^r 2^{5-r}x^{5-2r}$$
        由$5-2r=1$得$r=2$,对应的常数项$=80$,由$5-2r=-1$得
        $r=3$,对应的常数项$=-40$,故所求的常数项为$40$

        \textbf{解法二}
        
        用组合提取法,把原式看做$6$个因式相乘,若第$1$个括号提出$x$,从余下的$5$个括号中选$2$个提出$x$,选$3$个提出$\dfrac{1}{x}$,
        
        若第$1$个括号提出$\dfrac{1}{x}$,从余下的$5$个括号中选$2$个提出$\dfrac{1}{x}$,选$3$个提出$x$.
        
        故常数项$$x \cdot \comb{5}{2}(2x)^2 \cdot \comb{3}{3}(-\frac{1}{x})^3 + \frac{1}{x} \cdot \comb{5}{2}(-\frac{1}{x})^2 \cdot \comb{3}{3}(2x)^3 =-40+80=40$$
    \end{solution}

    \question 试求 
    \[
    (\comb{2}{2}+\comb{3}{2}x+\comb{4}{2}x^2+\comb{5}{2}x^3+\comb{6}{2}x^4+\comb{7}{2}x^5)^4
    \] 
    展开式中 \(x^5\) 的系数。
    \begin{solution}
        考虑 
        \[
        f(x) = \left(\sum_{n=0}^\infty \comb{n+2}{2} x^n \right)^4
        \]
        设 
        \[
        g(x)=\frac{x^2}{1-x} = \sum_{n=0}^\infty x^{n+2}
        \]
        则
        \[
        g'(x) = \frac{x(2-x)}{(x-1)^2}= \sum_{n=0}^\infty (n+2)x^{n+1}
        \]
        \[
        g''(x) = -\frac{2}{(x-1)^3}=\sum_{n=0}^\infty (n+2)(n+1)x^n = 2 \sum_{n=0}^\infty \comb{n+2}{2} x^n
        \]
        于是
        \[
        f(x) = \left( \sum_{n=0}^\infty \comb{n+2}{2} x^n \right)^4 = \left(- \frac{1}{(1 - x)^3} \right)^4 = \frac{1}{(1 - x)^{12}}
        \]
        故\((\comb{2}{2}+\comb{3}{2}x+\comb{4}{2}x^2+\comb{5}{2}x^3+\comb{6}{2}x^4+\comb{7}{2}x^5)^4\)  展开式中 \(x^5\) 的系数为
        \[
        \comb{12 + 5 - 1}{5} = 4368
        \]
    \end{solution}

    \question 若 
    \[
    (1 + x + x^2)^6 = a_0 + a_1 x + a_2 x^2 + \cdots + a_{12} x^{12},
    \]
    求\(a_2 + a_4 + \cdots + a_{12}\)的值。
    \begin{solution}
        设$f(x)=(1 + x + x^2)^6 = a_0 + a_1 x + a_2 x^2 + \cdots + a_{12} x^{12}$,则
        \[
        a_2 + a_4 + \cdots + a_{12}=\frac{f(1)+f(-1)}{2}-a_0=\frac{3^6+1}{2}-1=364
        \]
    \end{solution}

    \question 已知非零数列$\{a_n\}$满足$a_{1}=1,a_{2}=3,a_{n}(a_{n+2}-1)=a_{n+1}(a_{n+1}-1),n\in N,$求 
    \[
    \comb{2023}{0}a_{1}+\comb{2023}{1}a_{2}+\comb{2023}{2}a_{3}+\cdot\cdot\cdot+\comb{2023}{2023}a_{2024}
    \]
    的值。
    \begin{solution}
        由递推关系得
        \[
        \frac{a_{n+2}-1}{a_{n+1}} = \frac{a_{n+1}-1}{a_n}= \frac{a_2-1}{a_1} = \frac{3-1}{1} = 2
        \]
        即
        \[
         a_{n+1}-1 = 2a_n \Rightarrow a_{n+1}+1 = 2(a_n+1) 
        \]
        数列 $\{a_n+1\}$ 是首项为$2$,公比为$2$的等比数列, 故$a_n = 2^n - 1$,则
        \begin{align*}
        &\comb{2023}{0}a_{1}+\comb{2023}{1}a_{2}+\comb{2023}{2}a_{3}+\comb{2023}{3}a_{4}+\cdots+\comb{2023}{2023}a_{2024} \\
        &=\ \comb{2023}{0}(2^{1}-1)+\comb{2023}{1}(2^{2}-1)+\comb{2023}{2}(2^{3}-1)+\cdots+\comb{2023}{2023}(2^{2024}-1) \\
        &=\ 2\left(\comb{2023}{0}2^{0}+\comb{2023}{1}2^{1}+\comb{2023}{2}2^{2}+\cdots+\comb{2023}{2023}2^{2023}\right) \\
        &\qquad- \left(\comb{2023}{0}+\comb{2023}{1}+\comb{2023}{2}+\cdots+\comb{2023}{2023}\right) \\
        &=\ 2(1+2)^{2023} - 2^{2023} \\
        &=\ 2 \cdot 3^{2023} - 2^{2023}
        \end{align*}
    \end{solution}

    \question 求级数 
    \[
    1 + \frac{2}{4} + \frac{2 \cdot 3}{4 \cdot 8} + \frac{2 \cdot 3 \cdot 4}{4 \cdot 8 \cdot 12} + \frac{2 \cdot 3 \cdot 4 \cdot 5}{4 \cdot 8 \cdot 12 \cdot 16} + \dots
    \]
    的和。
    \begin{solution}
        通过提取公因数,整理得
        \begin{align*}
        S &= 1 + \frac{2}{4} + \frac{2 \cdot 3}{4 \cdot 8} + \frac{2 \cdot 3 \cdot 4}{4 \cdot 8 \cdot 12} + \frac{2 \cdot 3 \cdot 4 \cdot 5}{4 \cdot 8 \cdot 12 \cdot 16} + \dots \\
        &= 1 + \frac{2}{1!}\left(\frac{1}{4}\right) + \frac{2 \cdot 3}{2!}\left(\frac{1}{4}\right)^2 + \frac{2 \cdot 3 \cdot 4}{3!}\left(\frac{1}{4}\right)^3 + \frac{2 \cdot 3 \cdot 4 \cdot 5}{4!}\left(\frac{1}{4}\right)^4 + \dots \\
        &= 1 + \frac{-2}{1!}\left(-\frac{1}{4}\right)^1 + \frac{(-2)(-3)}{2!}\left(-\frac{1}{4}\right)^2 + \frac{(-2)(-3)(-4)}{3!}\left(-\frac{1}{4}\right)^3 \\
        &\quad + \frac{(-2)(-3)(-4)(-5)}{4!}\left(-\frac{1}{4}\right)^4 + \dots
        \end{align*}
        发现这是广义二项式 $(1 - x)^{-2}$ 在 $x=\dfrac{1}{4}$ 时的展开式,因此
        \[
        S = \left(1 - \frac{1}{4}\right)^{-2} = \frac{16}{9}
        \]
    \end{solution}

    \question 求 
    \[
    1 + \frac{1}{24} + \frac{1 \cdot 4}{24 \cdot 48} + \frac{1 \cdot 4 \cdot 7}{24 \cdot 48 \cdot 72} + \dots
    \]
    \begin{solution}
        整理得
        \begin{align*}
        S &= 1 + \frac{1}{24} + \frac{1 \cdot 4}{24 \cdot 48} + \frac{1 \cdot 4 \cdot 7}{24 \cdot 48 \cdot 72} + \dots \\
        &= 1 + \frac{3\cdot\frac{1}{3}}{24(1)} + \frac{3^2\cdot\frac{1}{3} \cdot \frac{4}{3}}{24^2 \cdot 1 \cdot 2} + \frac{3^3\cdot\frac{1}{3} \cdot \frac{4}{3} \cdot \frac{7}{3}}{24^3 \cdot 1 \cdot 2 \cdot 3} + \dots \\
        &= 1 + \frac{\left(-\frac{1}{3}\right)}{1}\left(-\frac{1}{8}\right) + \frac{\left(-\frac{1}{3}\right)\left(-\frac{4}{3}\right)}{1 \cdot 2}\left(-\frac{1}{8}\right)^2 + \frac{\left(-\frac{1}{3}\right)\left(-\frac{4}{3}\right)\left(-\frac{7}{3}\right)}{1 \cdot 2 \cdot 3}\left(-\frac{1}{8}\right)^3 + \dots
        \end{align*}
        发现这是广义二项式 $(1 - x)^{-\frac{1}{3}}$ 在 $x=\dfrac{1}{8}$ 时的展开式,因此
        \[
        S = \left(1 - \frac{1}{8}\right)^{-\frac{1}{3}} = \frac{2}{\sqrt[3]{7}}
        \]
    \end{solution}

    \question 求
    \[
    S = 1 - \frac{1}{4} + \frac{1 \cdot 3}{4 \cdot 8} - \frac{1 \cdot 3 \cdot 5}{4 \cdot 8 \cdot 12} + \dots
    \]
    \begin{solution}
        首先有
        \begin{align*}
        S &= 1 - \frac{1}{4} + \frac{1 \cdot 3}{4 \cdot 8} - \frac{1 \cdot 3 \cdot 5}{4 \cdot 8 \cdot 12} + \dots \\
        &= 1 - \frac{2^1\left(\frac{1}{2}\right)}{4^1 \cdot 1} + \frac{2^2\left(\frac{1}{2}\right)\left(\frac{3}{2}\right)}{4^2(1 \cdot 2)} - \frac{2^3\left(\frac{1}{2}\right)\left(\frac{3}{2}\right)\left(\frac{5}{2}\right)}{4^3(1 \cdot 2 \cdot 3)} + \dots \\
        &= 1 + \frac{-\frac{1}{2}}{1}\left(\frac{1}{2}\right)^1 + \frac{-\frac{1}{2}\left(-\frac{3}{2}\right)}{1 \cdot 2}\left(\frac{1}{2}\right)^2 + \frac{-\frac{1}{2}\left(-\frac{3}{2}\right)\left(-\frac{5}{2}\right)}{1 \cdot 2 \cdot 3}\left(\frac{1}{2}\right)^3 + \dots 
        \end{align*}
        这是广义二项式 $(1+x)^{-\frac{1}{2}}$ 在 $x = \dfrac{1}{2}$ 时的展开式,因此
        \[
        S = \left(1 + \frac{1}{2}\right)^{-\frac{1}{2}} = \sqrt{\frac{2}{3}}
        \]
    \end{solution}

    \question 求下列无穷级数的和
    \[
    \frac{3}{8}+\frac{3\cdot9}{8\cdot16}+\frac{3\cdot9\cdot15}{8\cdot16\cdot24}
    +\frac{3\cdot9\cdot15\cdot21}{8\cdot16\cdot24\cdot32}+\cdots
    \]
    \begin{solution}
        将一般项改写为
        \begin{align*}
        S&=\frac{3}{8}+\frac{3\cdot9}{8\cdot16}+\frac{3\cdot9\cdot15}{8\cdot16\cdot24}
        +\frac{3\cdot9\cdot15\cdot21}{8\cdot16\cdot24\cdot32}+\cdots \\
        &=\sum_{k=1}^{\infty}
        \frac{3^k(1\cdot3\cdot5\cdots(2k-1))}{8^{k+1}(1\cdot2\cdot3\cdots k)} \\
        &= \sum_{k=1}^{\infty}
        \frac{\left(\frac12\right)\left(\frac32\right)\cdots\left(\frac{2k-1}{2}\right)}{k!}\left(\frac34\right)^k \\
        &=\sum_{k=1}^{\infty}
        \frac{\left(-\frac12\right)\left(-\frac32\right)\cdots\left(-\frac{2k-1}{2}\right)}{k!}\left(-\frac34\right)^k
        \end{align*}
        考虑广义二项式展开,
        \[
        1+S=\sum_{k=0}^{\infty}
        \frac{\left(-\frac12\right)\left(-\frac32\right)\cdots\left(-\frac{2k-1}{2}\right)}{k!}\left(-\frac34\right)^k=\left(1-\frac34\right)^{-\frac12}
        \]
        给出
        \[
        S=1
        \]
    \end{solution}

    \question 求下列无穷和
    \[
    \sum_{k=1}^{\infty}\left[\prod_{r=1}^{k}\left(\frac{8r-7}{40r}\right)\right].
    \]
    \begin{solution}
        观察到
        \begin{align*}
        S&=\sum_{k=1}^{\infty}\left[\prod_{r=1}^{k}\left(\frac{8r-7}{40r}\right)\right] \\
        &= \frac{1}{40}
        + \frac{1}{40}\cdot\frac{9}{80}
        + \frac{1}{40}\cdot\frac{9}{80}\cdot\frac{17}{120}
        + \frac{1}{40}\cdot\frac{9}{80}\cdot\frac{17}{120}\cdot\frac{25}{160}
        + \cdots \\
        &= \frac{1}{(-8)(-5)\,1!}
        + \frac{1\cdot 9}{(-8)^2(-5)^2\,2!}
        + \frac{1\cdot 9\cdot 17}{(-8)^3(-5)^3\,3!}
        + \frac{1\cdot 9\cdot 17\cdot 25}{(-8)^4(-5)^4\,4!}
        + \cdots \\
        &= \frac{-\frac{1}{8}}{1!}\left(-\frac{1}{5}\right)
        + \frac{\left(-\frac{1}{8}\right)\left(-\frac{9}{8}\right)}{2!}\left(-\frac{1}{5}\right)^2 + \frac{\left(-\frac{1}{8}\right)\left(-\frac{9}{8}\right)\left(-\frac{17}{8}\right)}{3!}\left(-\frac{1}{5}\right)^3
        + \cdots 
        \end{align*}
        这是二项式级数展开式
        \[
        (1+x)^{\alpha}
        = \sum_{k=0}^{\infty}\frac{\alpha(\alpha-1)\cdots(\alpha-k+1)}{k!}x^k
        \]
        在 \(\alpha=-\dfrac{1}{8},x=-\dfrac{1}{5}\) 下去掉首项后的结果,因此
        \[
        S= \left(1-\frac{1}{5}\right)^{-\frac{1}{8}} - 1 = \sqrt[8]{\frac{5}{4}} - 1
        \]
    \end{solution}

    \question 已知级数
    \[
    S=\frac{5}{18}-\frac{5\cdot8}{18\cdot24}
    +\frac{5\cdot8\cdot11}{18\cdot24\cdot30}
    -\frac{5\cdot8\cdot11\cdot14}{18\cdot24\cdot30\cdot36}
    +\cdots
    \]
    收敛,求其和。
    \begin{solution}
        将原式写成
        \begin{align*}
        S &=\frac{5}{18}-\frac{5\cdot8}{18\cdot24}
        +\frac{5\cdot8\cdot11}{18\cdot24\cdot30}
        -\frac{5\cdot8\cdot11\cdot14}{18\cdot24\cdot30\cdot36}
        +\cdots \\
        &= \frac{5}{6 \cdot 3} - \frac{5 \cdot 8}{6^2(3 \cdot 4)} + \frac{5 \cdot 8 \cdot 11}{6^3(3 \cdot 4 \cdot 5)} - \frac{5 \cdot 8 \cdot 11 \cdot 14}{6^4(3 \cdot 4 \cdot 5 \cdot 6)} + \dots \\
        &= \frac{\frac{5}{3}}{3}\left(\frac{1}{2}\right) - \frac{\frac{5}{3}\left(\frac{8}{3}\right)}{4 \cdot 3}\left(\frac{1}{2}\right)^2 + \frac{\frac{5}{3}\left(\frac{8}{3}\right)\left(\frac{11}{3}\right)}{5 \cdot 4 \cdot 3}\left(\frac{1}{2}\right)^3 - \frac{\frac{5}{3}\left(\frac{8}{3}\right)\left(\frac{11}{3}\right)\left(\frac{14}{3}\right)}{6 \cdot 5 \cdot 4 \cdot 3}\left(\frac{1}{2}\right)^4 + \dots
        \end{align*}
        为了构建二项式展开的形式,写成
        \begin{align*}
        &S \cdot \frac{1}{2} \cdot \left(\frac{1}{2}\right)^2 \cdot \frac{2}{3} \cdot \frac{1}{3} \\
        &= \frac{\frac{1}{3}\left(\frac{2}{3}\right)\left(\frac{5}{3}\right)}{3!}\left(\frac{1}{2}\right)^3 - \frac{\frac{1}{3}\left(\frac{2}{3}\right)\left(\frac{5}{3}\right)\left(\frac{8}{3}\right)}{4!}\left(\frac{1}{2}\right)^4 + \dots \\
        &= \frac{\frac{1}{3}\left(-\frac{2}{3}\right)\left(-\frac{5}{3}\right)}{3!}\left(\frac{1}{2}\right)^3 + \frac{\frac{1}{3}\left(-\frac{2}{3}\right)\left(-\frac{5}{3}\right)\left(-\frac{8}{3}\right)}{4!}\left(\frac{1}{2}\right)^4 + \dots \\
        &= \left[ 1 + \frac{\frac{1}{3}}{1!}\left(\frac{1}{2}\right)^1 + \frac{\frac{1}{3}\left(-\frac{2}{3}\right)}{2!}\left(\frac{1}{2}\right)^2 + \dots \right] - \left[ 1 + \frac{\frac{1}{3}}{1!}\left(\frac{1}{2}\right)^1 + \frac{\frac{1}{3}\left(-\frac{2}{3}\right)}{2!}\left(\frac{1}{2}\right)^2 \right]
        \end{align*}
        第一部分即为 $\left(1 + \dfrac{1}{2}\right)^{\frac{1}{3}}$ 的二项式展开:
        \[
        \frac{1}{36}S = \left(1 + \frac{1}{2}\right)^{\frac{1}{3}} - \left(1 + \frac{1}{6} - \frac{1}{36}\right) 
        \]
        故
        \[
        S = 9\sqrt[3]{96} - 41
        \]
    \end{solution}

    \question 根据组合数的定义,证明下列恒等式:
    \begin{parts}
    \part
    \[
    \comb{n}{r} = \comb{n}{n-r}
    \]
    \begin{solution}
        有
        \[
        \comb{n}{n-r} = \frac{n!}{(n-r)!(n-(n-r))!} = \frac{n!}{(n-r)!r!} = \comb{n}{r}
        \]
    \end{solution}
    \part 
    \[
    k \cdot \comb{n}{k} = n \cdot \comb{n-1}{k-1}
    \]
    \begin{solution}
        有
        \begin{align*}
        k \cdot \comb{n}{k} 
        &= k \cdot \frac{n!}{k!(n-k)!} \\
        &= \frac{n!}{(k-1)!(n-k)!} \\
        &= n \cdot \frac{(n-1)!}{(k-1)!\,[ (n-1)-(k-1) ]!} \\
        &= n \cdot \comb{n-1}{k-1}
        \end{align*}
    \end{solution}
    \part 帕斯卡公式:
    \[
    \comb{n}{r-1} + \comb{n}{r} = \comb{n+1}{r}
    \]
    \begin{solution}
        有
        \begin{align*}
        \comb{n}{r-1} + \comb{n}{r} &= \frac{n!}{(r-1)!(n-(r-1))!} + \frac{n!}{r!(n-r)!} \\
        &= \frac{n!}{(r-1)!(n-r+1)!} + \frac{n!}{r!(n-r)!} \\
        &= \frac{n! \cdot r}{r!(n-r+1)!} + \frac{n! \cdot (n-r+1)}{r!(n-r+1)!} \\
        &= \frac{n! [r + (n-r+1)]}{r!(n-r+1)!} \\
        &= \frac{n! (n+1)}{r!(n-r+1)!} \\
        &= \frac{(n+1)!}{r!(n+1-r)!} \\
        &= \comb{n+1}{r}
        \end{align*}
    \end{solution}
    \part 
    \[
    \comb{n}{m}\cdot\comb{m}{k} = \comb{n}{k}\cdot\comb{n-k}{m-k}
    \]
    \begin{solution}
        有
        \begin{align*}
        \comb{n}{m}\cdot\comb{m}{k} &= \frac{n!}{m!(n-m)!} \cdot \frac{m!}{k!(m-k)!} \\
        &= \frac{n!}{k!(m-k)!(n-m)!} \\
        &= \frac{n!}{k!(n-k)!} \cdot \frac{(n-k)!}{(m-k)!(n-m)!} \\
        &= \comb{n}{k}\cdot\comb{n-k}{m-k}
        \end{align*}
    \end{solution}
    \end{parts}

    \question 设 $n,k \in \mathbb{N}$, 根据帕斯卡公式,试证明
    \begin{parts}
    \part
    \[
    \comb{n}{0} + \comb{n+1}{1} + \dots + \comb{n+k}{k} = \comb{n+k+1}{k}
    \]
    \begin{solution}
        反复对最后一个二项式系数运用帕斯卡公式,得
        \begin{align*}
        \comb{n}{k} &= \comb{n-1}{k} + \comb{n-1}{k-1} \\
        &= \comb{n-1}{k} + \comb{n-2}{k-1} + \comb{n-2}{k-2} \\
        &\;\vdots \\
        &= \comb{n-1}{k} + \comb{n-2}{k-1} + \dots +  \comb{n-k-2}{0} \\
        &= \comb{n-1}{k} + \comb{n-2}{k-1} + \dots +  \comb{n-k-1}{0} 
        \end{align*}
        以 $n+k+1$ 取代 $n$,即得证。
    \end{solution}
    \part 
    \[
    \comb{k}{k} + \comb{k+1}{k} + \dots + \comb{n}{k} = \comb{n+1}{k+1}
    \]
    \begin{solution}
        反复对第一个二项式系数运用帕斯卡公式,得
        \begin{align*}
        \comb{n}{k} &= \comb{n-1}{k} + \comb{n-1}{k-1} \\
        &= \comb{n-2}{k} + \comb{n-2}{k-1} + \comb{n-1}{k-1} \\
        &\;\vdots \\
        &= \comb{0}{k} + \comb{0}{k-1} + \comb{1}{k-1} + \dots + \comb{n-2}{k-1} + \comb{n-1}{k-1}
        \end{align*}
        而 $\comb{0}{k} = 0$, 以 $n+1$ 取代 $n$,且以 $k+1$ 取代 $k$,即得证。
    \end{solution}
    \end{parts}

    \question 证明恒等式
    \[
    \sum_{r=0}^{n} (\comb{n}{r})^2 = \comb{2n}{n}, \quad n \ge 0
    \]
    \begin{solution}
        设 $S$ 是一个有 $2n$ 个元素的集合,$\comb{2n}{n}$ 为计数 $S$ 的 $n$ 子集的数目。又把 $S$ 划分成 $A$ 和 $B$ 两个子集,每一个子集都有 $n$ 个元素。利用 $S$ 的这个划分去划分 $S$ 的 $n$ 子集。$S$ 的每一个 $n$ 子集包含 $k$ 个 $A$ 的元素和 $n-k$ 个 $B$ 的元素,$k$ 是介于 $0$ 和 $n$ 之间的任意整数。我们把 $S$ 的 $n$ 子集划分成 $n+1$ 个部分:
        \[
        C_0, C_1, \dots, C_n
        \]
        其中 $C_k$ 是由包含 $k$ 个 $A$ 的元素和 $n-k$ 个 $B$ 的元素组成的 $n$ 子集。根据加法原理,有
        \[
        \comb{2n}{n} = |C_0| + |C_1| + \dots + |C_n|
        \]
        $C_k$ 中的 $n$ 子集可以通过从 $A$ 中选择 $k$ 个元素,再从 $B$ 中选择 $n-k$ 个元素而得到。根据乘法原理,
        \[
        |C_k| = \comb{n}{k} \comb{n}{n-k} = \comb{n}{k}^2, \quad k \ge 0.
        \]
        把上式代入得
        \[
        \sum_{k=0}^{n} (\comb{n}{k})^2 = (\comb{n}{0})^2 + (\comb{n}{1})^2 + \dots + (\comb{n}{n})^2 = \comb{2n}{n}
        \]
    \end{solution}

    \question 根据二项式定理,证明下列恒等式:
    \begin{parts}
    \part 
    \[
    \comb{n}{0} + \comb{n}{2} + \dots = \comb{n}{1} + \comb{n}{3} + \dots = 2^{\,n-1}
    \]
    \begin{solution}
        由二项式定理,
        \[
        (1-1)^n=1^n+\comb{n}{1}1^{n-1}(-1)+\comb{n}{2}1^{n-2}(-1)^2+\dots+(-1)^n
        \]
        即
        \[
        \comb{n}{0} - \comb{n}{1} + \comb{n}{2} - \dots + (-1)^n \cdot \comb{n}{n} = 0
        \]
        由此可得
        \[
        \comb{n}{0} + \comb{n}{2} + \dots = \comb{n}{1} + \comb{n}{3} + \dots
        \]
        由二项式定理,
        \[
        (1+1)^n=1^n+\comb{n}{1}1^{n-1}1+\comb{n}{2}1^{n-2}1^2+\dots+1^n
        \]
        给出
        \[
        \comb{n}{0} + \comb{n}{1} + \comb{n}{2} + \dots + \comb{n}{n} = 2^n
        \]
        因此
        \[
        \comb{n}{0} + \comb{n}{2} + \dots = \comb{n}{1} + \comb{n}{3} + \dots = \frac{1}{2}(2^n) = 2^{\,n-1}
        \]
    \end{solution}
    \part 
    \[
    1 \cdot \comb{n}{1} + 2 \cdot \comb{n}{2} + \dots + n \cdot \comb{n}{n}=n 2^{n-1}
    \]
    \begin{solution}
        考虑二项式展开,
        \[
        \sum_{k=0}^n \comb{n}{k} x^k = (1+x)^n
        \]
        两边对 $x$ 求导,
        \[
        \sum_{k=1}^n k \cdot \comb{n}{k} x^{k-1} = n(1+x)^{n-1}
        \]
        令 $x=1$,
        \[
        \sum_{k=1}^n k \cdot \comb{n}{k} = n 2^{n-1}
        \]
    \end{solution}
    \begin{solution}
        由恒等式
        \[
        k \cdot \comb{n}{k} = n \cdot \comb{n-1}{k-1},\quad \sum_{k=0}^{n} \comb{n}{k} = 2^n
        \]
        有
        \begin{align*}
        \sum_{k=1}^{n} k \cdot \comb{n}{k} &= 1\comb{n}{1} + 2\comb{n}{2} + \dots + n\comb{n}{n} \\
        &= n \left[ \comb{n-1}{0} + \comb{n-1}{1} + \dots + \comb{n-1}{n-1} \right] \\
        &= n 2^{n-1}
        \end{align*}
    \end{solution}
    \part 
    \[
    \sum_{k=1}^n k^2 \cdot \comb{n}{k} = n(n+1) 2^{n-2}
    \]
    \begin{solution}
        考虑二项式展开,
        \[
        \sum_{k=0}^n \comb{n}{k} x^k = (1+x)^n
        \]
        两边求导,并乘以 $x$ 得
        \[
        \sum_{k=1}^{n} k \comb{n}{k} x^k = x \frac{d}{dx}(1+x)^n = x n (1+x)^{n-1}
        \]
        再次对两边求导,
        \[
        \sum_{k=1}^{n} k^2 \comb{n}{k} x^{k-1} = n (1+x)^{n-1} + x n (n-1) (1+x)^{n-2}
        \]
        代入 $x=1$,化简得
        \[
        \sum_{k=1}^{n} k^2 \comb{n}{k} = n \cdot 2^{n-1} + n(n-1) \cdot 2^{n-2} = n(n+1) 2^{n-2}
        \]
    \end{solution}
    \part 
    \[
    \comb{n}{0} + \frac{1}{2} \comb{n}{1} + \frac{1}{3} \comb{n}{2} + \cdots + \frac{1}{n+1} \comb{n}{n} = \frac{1}{n+1} \left(2^{n+1} - 1\right)
    \]
    \begin{solution}
        考虑二项式展开,
        \[
        (1+x)^n = \sum_{k=0}^{n} \comb{n}{k} x^k = \comb{n}{0} + \comb{n}{1}x + \comb{n}{2}x^2 + \dots + \comb{n}{n}x^n
        \]
        对等式两边在区间 $[0, 1]$ 上关于 $x$ 进行定积分:
        \[
        \int_{0}^{1} (1+x)^n \, dx = \int_{0}^{1} \left( \comb{n}{0} + \comb{n}{1}x + \comb{n}{2}x^2 + \dots + \comb{n}{n}x^n \right) \, dx
        \]
        左式为
        \[
        \int_{0}^{1} (1+x)^n \, dx = \left[ \frac{(1+x)^{n+1}}{n+1} \right]_{0}^{1} = \frac{2^{n+1} - 1}{n+1}
        \]
        右式为
        \begin{align*}
        \int_{0}^{1} \sum_{k=0}^{n} \comb{n}{k} x^k \, dx &= \sum_{k=0}^{n} \comb{n}{k} \int_{0}^{1} x^k \, dx \\
        &= \sum_{k=0}^{n} \comb{n}{k} \left[ \frac{x^{k+1}}{k+1} \right]_{0}^{1} \\
        &= \sum_{k=0}^{n} \frac{1}{k+1} \comb{n}{k} 
        \end{align*}
        于是得证
        \[
        \comb{n}{0} + \frac{1}{2} \comb{n}{1} + \frac{1}{3} \comb{n}{2} + \dots + \frac{1}{n+1} \comb{n}{n} = \frac{1}{n+1} (2^{n+1} - 1)
        \]
    \end{solution}
    \begin{solution}
        \begin{align*}
        \sum_{k=0}^{n} \frac{\comb{n}{k}}{k+1} 
        &= \sum_{k=0}^{n} \frac{1}{k+1} \cdot \frac{n!}{k!(n-k)!} \\
        &= \sum_{k=0}^{n} \frac{n!}{(k+1)!(n-k)!} \\
        &= \frac{1}{n+1} \sum_{k=0}^{n} \frac{(n+1)!}{(k+1)!(n-k)!} \\
        &= \frac{1}{n+1} \sum_{k=0}^{n} \comb{n+1}{k+1} \\
        &= \frac{1}{n+1} \left( \comb{n+1}{1} + \comb{n+1}{2} + \dots + \comb{n+1}{n+1} \right) \\
        &= \frac{1}{n+1} (2^{n+1}-1)
        \end{align*}
        于是得证。
    \end{solution}
    \part 
    \[
    \sum_{k \text{ even}} k \comb{n}{k} = n 2^{n-2}
    \]
    \begin{solution}
        由二项式定理,
        \[
        (1+x)^n = \comb{n}{0} + \comb{n}{1}x + \comb{n}{2}x^2 + \dots + \comb{n}{n}x^n
        \]
        两边求导得
        \[
        n(1+x)^{n-1} = \comb{n}{1} + 2\comb{n}{2}x + 3\comb{n}{3}x^2 + \dots + n\comb{n}{n}x^{n-1}
        \]
        将 $x$ 替换为 $-x$,得
        \[
        n(1-x)^{n-1} = \comb{n}{1} - 2\comb{n}{2}x + 3\comb{n}{3}x^2 - \dots + (-1)^{n-1} n \comb{n}{n}x^{n-1}
        \]
        两式相加
        \[
        n(1+x)^{n-1} + n(1-x)^{n-1} = 2 \sum_{k \text{ even}} k \comb{n}{k} x^{k-1}
        \]
        代入 $x=1$,得到
        \[
        \sum_{k \text{ even}} k \comb{n}{k} = n 2^{n-2}
        \]
    \end{solution}
    \end{parts}

    \question 证明范德蒙恒等式
    \[
    \comb{n+m}{k} = \sum_{i=0}^{k} \comb{n}{i}\comb{m}{k-i}
    \]
    \begin{solution}
        注意到
        \[
        (1+x)^n (1+x)^m = (1+x)^{n+m}
        \]
        左式根据二项展开得
        \[
        (1+x)^n (1+x)^m = \left(\sum_{i=0}^{n} \comb{n}{i}x^i\right)\left(\sum_{j=0}^{m} \comb{m}{j}x^j\right) = \sum_{k=0}^{m+n} \left(\sum_{i=0}^{k} \comb{n}{i}\comb{m}{k-i}\right)x^k
        \]
        而右式为
        \[
        (1+x)^{n+m} = \sum_{k=0}^{n+m} \comb{n+m}{k} x^k
        \]
        比较 $x^k$ 系数, 得到
        \[
        \comb{n+m}{k} = \sum_{i=0}^{k} \comb{n}{i}\comb{m}{k-i}
        \]
    \end{solution}

    \question 证明  
    \[
    \left(\comb{n}{0}-\comb{n}{2}+\comb{n}{4}-\cdots\right)^{2}+\left(\comb{n}{1}-\comb{n}{3}+\comb{n}{5}-\cdots\right)^{2}=2^{n}
    \]
    \begin{solution}
        设 \(f(x)=(x+1)^n = \comb{n}{0} + \comb{n}{1} x + \comb{n}{2} x^2 + \cdots + \comb{n}{n} x^n\),取 \(x=i\) 得 
        \[
        f(i) = (i+1)^n = (\sqrt{2}\, e^{\frac{\pi i}{4}})^n = 2^{\frac{n}{2}} e^{\frac{\pi i}{4}}
        \]
        且有
        \[
        f(i) = (\comb{n}{0} - \comb{n}{2} + \comb{n}{4} - \cdots) + i (\comb{n}{1} - \comb{n}{3} + \comb{n}{5} - \cdots)
        \]
        比较实部与虚部得 
        \[
        \comb{n}{0} - \comb{n}{2} + \comb{n}{4} - \cdots = 2^{\frac{n}{2}} \cos\frac{\pi i}{4},\quad \comb{n}{1} - \comb{n}{3} + \comb{n}{5} - \cdots = 2^{\frac{n}{2}} \sin\frac{\pi i}{4}
        \]
        因此得证
        \[
        (\comb{n}{0} - \comb{n}{2} + \cdots)^2 + (\comb{n}{1} - \comb{n}{3} + \cdots)^2 
        = 2^n\cos^2\left(\frac{\pi i}{4}+\sin^2\frac{\pi i}{4}\right) = 2^n
        \]
    \end{solution}

    \question 若将 $(x+1)^{2000}$ 展开,问有多少个系数是奇数？
    \begin{solution}
        一般地,若 $(x+1)^n$ 展开式中奇数系数的个数为 $2^{a}$,其中 $a$ 是 $n$ 的二进制表示中 1 的个数,将 $2000$ 转换为二进制:
        \[
        2000 = 11111010000_2
        \]
        其中有 6 个 1,因此 $(x+1)^{2000}$ 中奇数系数的个数为
        \[
        2^6 = 64.
        \]
        又解:对 2 取模,有
        \begin{align*}
        (x+1)^{2000} 
        &= (x+1)^{1024}(x+1)^{512}(x+1)^{256}(x+1)^{128}(x+1)^{64}(x+1)^{16} \\
        &\equiv (x^{1024}+1)(x^{512}+1)\cdots(x^{16}+1) \pmod{2},
        \end{align*}
        展开后共有 $2^6=64$ 项系数为奇数。
    \end{solution}

    \question 设 \( n \) 为大于或等于 1 的正整数,考虑数列 \( a_n = \left(1 + \dfrac{1}{n} \right)^n\),试用二项式定理证明对所有 \( n \) 皆满足  
    \[
    a_n < e
    \]
    \begin{solution}
        设 \( a_n = \left(1 + \dfrac{1}{n} \right)^n \),考虑对 \( a_n \) 进行二项式展开
        \[
        \left(1 + \frac{1}{n} \right)^n = \sum_{k=0}^{n} \binom{n}{k} \cdot \left(\frac{1}{n}\right)^k
        = \sum_{k=0}^{n} \frac{n(n-1)\cdots(n-k+1)}{k! \cdot n^k}
        \]
        注意到
        \[
        \binom{n}{k} \cdot \left(\frac{1}{n}\right)^k \le \frac{1}{k!}
        \Rightarrow a_n=\left(1 + \frac{1}{n} \right)^n < \sum_{k=0}^{n} \frac{1}{k!} < \sum_{k=0}^{\infty} \frac{1}{k!} = e 
        \]
    \end{solution}

    \question 定义 \( a_n \) 为 \( (3 - \sqrt{x})^n \) 展开式中 \( x \) 项的系数,其中 \( n \) 是正整数。求极限  
    \[
    \lim_{n \to \infty} \left( \frac{3^2}{a_2} + \frac{3^3}{a_3} + \cdots + \frac{3^n}{a_n} \right)
    \]
    \begin{solution}
        观察到 \( (3 - \sqrt{x})^n \) 的二项展开式为:
        \[
        (3 - \sqrt{x})^n = \sum_{k=0}^n \binom{n}{k} 3^{n - k} (-\sqrt{x})^k
        \]
        \( x \) 项的系数对应 \( k = 2 \):
        \[
        a_n = \binom{n}{2} 3^{n - 2}
        \]
        所以
        \[
        \frac{3^n}{a_n} = \frac{3^n}{\binom{n}{2} 3^{n - 2}} = \frac{3^2}{\binom{n}{2}} = \frac{9}{\binom{n}{2}}
        \]
        原式变为
        \[
        9 \sum_{n=2}^{\infty} \frac{1}{\binom{n}{2}}=
        18\sum_{n=2}^{\infty} \frac{2}{n(n-1)} = 18\sum_{n=2}^{\infty} \left( \frac{1}{n-1} - \frac{1}{n} \right)
        \]
        由累加法,
        \[
        \sum_{n=2}^{\infty} \left( \frac{1}{n-1} - \frac{1}{n} \right) =  1 - \frac{1}{2} + \frac{1}{2} - \frac{1}{3} + \cdots = 1
        \]
        因此原极限为 $18$。
    \end{solution}

    \question 用 $|S|$ 表示集合 $S$ 中元素的个数。已知集合
    \[
    S = \left\{ \frac{1}{2}, \frac{1}{3}, \frac{1}{4}, \ldots, \frac{1}{113} \right\},\quad
    T = \left\{ A \subseteq S \,\middle|\, |A| = 2n,\, n \in \mathbb{N} \right\},
    \]
    试回答下列问题:
    \begin{parts}
    \part 求$|T|$。
    \begin{solution}
        集合 $S$ 中的元素个数为 $|S| = 112$,所以:
        \[
        |T| = \comb{112}{2} + \comb{112}{4} + \cdots + \comb{112}{112}
        \]
        由二项式定理,
        \begin{align*}
        (1 + 1)^{112} &= \comb{112}{0} + \comb{112}{1} + \cdots + \comb{112}{112} = 2^{112} \\
        (1 - 1)^{112} &= \comb{112}{0} - \comb{112}{1} + \comb{112}{2} - \cdots + (-1)^{112} \cdot \comb{112}{112} = 0
        \end{align*}
        两式相加得
        \[
        2^{112} = 2 \left( \comb{112}{0} + \comb{112}{2} + \comb{112}{4} + \cdots + \comb{112}{112} \right)
        \Rightarrow \sum_{k=0}^{56} \comb{112}{2k} = 2^{111}
        \]
        于是
        \[
        |T| = 2^{111} - 1
        \]
    \end{solution}
    \part 对任意 $A_i \in T$,将 $A_i$ 中所有元素相乘的乘积记为 $m_i$,再将所有的 $m_i$ 相加,其和为 $M$,求 $M$ 的值。
    \end{parts}
    \begin{solution}
        设
        \[
        f(x) = \left(x + \frac{1}{2}\right)\left(x + \frac{1}{3}\right)\cdots\left(x + \frac{1}{113}\right)
        \]
        展开 $f(x)$ 得
        \begin{itemize}
            \item $x^{110}$ 的系数为所有 $|A_i| = 2$ 的子集乘积之和；
            \item $x^{108}$ 的系数为所有 $|A_i| = 4$ 的子集乘积之和；
            \item $\cdots$；
            \item 常数项为所有 $|A_i| = 112$ 的子集乘积之和。
        \end{itemize}
        故所有偶数个元素子集的乘积之和为
        \[
        M = \frac{1}{2} \left( f(1) + f(-1) \right) - 1
        \]
        而
        \[
        f(1) = \left(1 + \frac{1}{2}\right)\left(1 + \frac{1}{3}\right)\cdots\left(1 + \frac{1}{113}\right)
        = \frac{3}{2} \cdot \frac{4}{3} \cdot \frac{5}{4} \cdots \frac{114}{113} = 57
        \]
        \[
        f(-1) = \left(-1 + \frac{1}{2}\right)\left(-1 + \frac{1}{3}\right)\cdots\left(-1 + \frac{1}{113}\right)
        = \left(-\frac{1}{2}\right)\left(-\frac{2}{3}\right)\cdots\left(-\frac{112}{113}\right) = \frac{1}{113}
        \]
        因此
        \[
        M = \frac{1}{2} \left( 57 + \frac{1}{113} \right) - 1 = \frac{3108}{113}
        \]
    \end{solution}

    \question 函数 $f$ 由实数 $a,b,c$ 定义为
    \[
    f(x)=(a+bx+cx^2)(1-x)^{-3},\quad x\in\mathbb{R},\ |x|<1.
    \]
    \begin{parts}
    \part 证明
    \[
    f(x)=a+(3a+b)x+\frac{1}{2}\sum_{n=2}^{\infty}\left[a(n+1)(n+2)+bn(n+1)+cn(n-1)\right]x^n.
    \]
    \begin{solution}
        考虑二项展开式
        \begin{align*}
        (1-x)^{-3}
        &= 1+\frac{-3}{1!}(-x)+\frac{-3(-3-1)}{2!}(-x)^2+\frac{-3(-3-1)(-3-2)}{3!}(-x)^3+\cdots \\
        &=\sum_{n=0}^{\infty}\frac{1}{2}\cdot \frac{(n+2)!}{n!}x^n \\
        &=\sum_{n=0}^{\infty}\frac{1}{2}(n+1)(n+2)x^n
        \end{align*}
        于是
        \begin{align*}
        f(x)
        &=(a+bx+cx^2)(1-x)^{-3} \\
        &=(a+bx+cx^2)\sum_{n=0}^{\infty}\frac{1}{2}(n+1)(n+2)x^n \\
        &=\frac{1}{2}a\sum_{n=0}^{\infty}(n+1)(n+2)x^n
        +\frac{1}{2}b\sum_{n=0}^{\infty}(n+1)(n+2)x^{n+1} \\
        &\quad +\frac{1}{2}c\sum_{n=0}^{\infty}(n+1)(n+2)x^{n+2}.
        \end{align*}
        将求和指标统一从 $n=2$ 开始,可得
        \[
        f(x)= a+(3a+b)x+\frac{1}{2}\sum_{n=2}^{\infty}\left[a(n+1)(n+2)+bn(n+1)+cn(n-1)\right]x^n
        \]
    \end{solution}
    \part 据此,求
    \[
    \sum_{n=1}^{\infty}\frac{n^2}{2^n}.
    \]
    \begin{solution}
        强迫展开式中 $x^n$ 的系数
        \[
        \frac{1}{2}\left[(a+b+c)n^2+(3a+b-c)n+2a\right]
        \]
        等于 $n^2$,比较系数得
        \[
        \frac{1}{2}(a+b+c)=1,\quad
        \frac{1}{2}(3a+b-c)=0,\quad
        a=0.
        \]
        解得
        \[
        a=0,\quad b=1,\quad c=1
        \]
        于是
        \[
        f(x)=(x+x^2)(1-x)^{-3}= x+\sum_{n=2}^{\infty}n^2x^n
        \]
        令 $x=\dfrac{1}{2}$,有
        \[
        f\left(\frac{1}{2}\right)
        =\left(\frac{1}{2}+\frac{1}{4}\right)\left(1-\frac{1}{2}\right)^{-3}=\frac{1}{2}+\sum_{n=2}^{\infty}\frac{n^2}{2^n}
        \]
        因此
        \[
        \sum_{n=2}^{\infty}\frac{n^2}{2^n}
        =6-\frac{1}{2}=\frac{11}{2}.
        \]
        再加上 $n=1$ 的项 $\dfrac{1}{2}$,得到
        \[
        \sum_{n=1}^{\infty}\frac{n^2}{2^n}=6
        \]
    \end{solution}
    \end{parts}

    \question 设函数
    \[
    f(x)=\frac{1}{\sqrt{1-x}},\quad -1<x<1
    \]
    \begin{parts}
    \part 证明
    \[
    f(x)=\sum_{r=0}^{\infty}\comb{2r}{r}\left(\frac{1}{4}x\right)^r
    \]
    \begin{solution}
        由广义二项式定理,
        \begin{align*}
        (1-x)^{-\frac12}
        &=1+\frac{-\frac12}{1!}(-x)+\frac{-\frac12\left(-\frac32\right)}{2!}(-x)^2+\cdots \\
        &=1+\sum_{r=1}^{\infty} \frac{\left(-\frac12\right)\left(-\frac32\right)\cdots\left(-\frac{2r-1}{2}\right)}{r!}(-x)^r \\
        &=1+\sum_{r=1}^{\infty}\frac{1\cdot3\cdot5\cdots(2r-1)}{r!2^r}x^r
        \end{align*}
        利用
        \[
        1\cdot3\cdot5\cdots(2r-1)=\frac{(2r)!}{2^r r!}
        \]
        可得
        \[
        (1-x)^{-\frac12}
        =\sum_{r=0}^{\infty}\frac{(2r)!}{r!r!}\left(\frac{x}{4}\right)^r
        =\sum_{r=0}^{\infty}\comb{2r}{r}\left(\frac{x}{4}\right)^r
        \]
    \end{solution}
    \part 求 
    \[
    \frac{1}{\sqrt{16-x^2}}
    \]
    的展开式,并进一步证明
    \[
    \frac{x}{(16-x^2)^{\frac{3}{2}}}
    =\sum_{r=1}^{\infty}\comb{2r}{r}\left(\frac{1}{16}\right)^r\left(\frac{1}{8}x\right)^{2r-1}
    \]
    \begin{solution}
        由(a),二项展开得,
        \begin{align*}
        (16-x^2)^{-\frac12}
        &=16^{-\frac12}\left(1-\frac{x^2}{16}\right)^{-\frac12} \\
        &=\frac14\sum_{r=0}^{\infty}\comb{2r}{r}\left(\frac{x^2}{64}\right)^r \\
        &=\frac14\sum_{r=0}^{\infty}\comb{2r}{r}\left(\frac{x}{8}\right)^{2r}
        \end{align*}
        对两边求导,化简得
        \[
        \frac{x}{(16-x^2)^{\frac32}}
        =\frac14 \sum_{r=1}^{\infty}\comb{2r}{r}(2r)\left(\frac{x}{8}\right)^{2r-1}\cdot \frac18
        \]
    \end{solution}
    \part 求下列级数的值
    \[
    \sum_{r=1}^{\infty}\comb{2r}{r}\left(\frac{5}{32}\right)^r\left(\frac{4}{25}\right)^r
    \]
    \begin{solution}
        原级数化简为
        \[
        \sum_{r=1}^{\infty}\comb{2r}{r}\left(\frac{1}{16}\right)^r\left(\frac{2}{5}\right)^{2r-1}
        \]
        由(b),令
        \[
        \frac{x}{8}=\frac{2}{5}\Rightarrow x=\frac{16}{5}
        \]
        则
        \[
        \sum_{r=1}^{\infty}\comb{2r}{r}\left(\frac{1}{16}\right)^r\left(\frac{2}{5}\right)^{2r-1}
        =\frac{\frac{16}{5}}{\bigl(16-\left(\frac{16}{5}\right)^2\bigr)^{\frac32}}=\frac{25}{108}
        \]
    \end{solution}
    \end{parts}

\end{questions}
\pagebreak

\begin{center}
  {\fontsize{30pt}{26pt}\selectfont
    \hypertarget{泰勒展开式}{泰勒展开式} \label{泰勒展开式}
  }
\end{center}
\separator
\vspace{1pt}
\begin{questions}
    \question 设 
    \[
    y = e^{\tan x},\quad x \in \mathbb{R}
    \]
    \begin{parts}
    \part 证明
    \[
    \frac{d^2y}{dx^2} = (1+\tan^2 x + 2 \tan x) \frac{dy}{dx}.
    \]
    \begin{solution}
        首先对 $y$ 求导
        \[
        \frac{dy}{dx} = e^{\tan x} \sec^2 x = y \sec^2 x.
        \]
        再次求导得
        \begin{align*}
        \frac{d^2y}{dx^2} 
        &= \frac{dy}{dx} \sec^2 x + y \cdot 2 \sec^2 x \tan x \\
        &= \frac{dy}{dx} \sec^2 x + 2 \frac{dy}{dx} \tan x \\
        &= \frac{dy}{dx} ( \sec^2 x + 2 \tan x ) \\
        &= \frac{dy}{dx} (1 + \tan^2 x + 2 \tan x ) 
        \end{align*}
    \end{solution}
    \part 求 $e^{\tan x}$ 展开至 $x^3$项的麦克劳林展开式。
    \begin{solution}
        在 $x=0$处的泰勒展开式为
        \[
        y = y_0 + y'_0 x + \frac{y''_0}{2!} x^2 + \frac{y'''_0}{3!} x^3 + O(x^4),
        \]
        其中$y_0 = 1,y'_0 = 1,y''_0 = 1$,计算三阶导数:
        \[
        \frac{d^3y}{dx^3} = 2(1+\tan x) \sec^2 x \frac{dy}{dx} + (1+\tan x)^2 \frac{d^2y}{dx^2},
        \]
        代入 $x=0$ 得 $y'''_0 = 3$,于是
        \[
        e^{\tan x} = 1 + x + \frac{1}{2}x^2 + \frac{1}{2}x^3 + O(x^4)
        \]
    \end{solution}
    \end{parts}

    \question 设 
    \[
    f(x) = \sin[\ln(1+x)],\quad x \in \mathbb{R},\quad x > -1
    \]
    \begin{parts}
    \part 求关于 $f(x)$ 的微分方程。
    \begin{solution}
        首先求导得
        \[
        f'(x) = \cos[\ln(1+x)] \frac{1}{1+x} \Rightarrow 
        (1+x) f'(x) = \cos[\ln(1+x)]
        \]
        再次求导
        \[
        (1+x) f''(x) + f'(x) = - \sin[\ln(1+x)] \frac{1}{1+x}
        \]
        移项得微分方程
        \[
        (1+x)^2 f''(x) + (1+x) f'(x) + f(x) = 0 \tag{1}
        \]
    \end{solution}
    \part 求$f(x)$的麦克劳林展开前 3 个非零项。
    \begin{solution}
        利用微分方程可求各阶导数在 $x=0$ 的值,首先有
        \[
        f(0) = \sin(\ln 1) = 0, \quad f'(0) = \cos(\ln 1) = 1.
        \]
        由(1)得,
        \[
        f''(0) + f'(0) + f(0) = 0 \implies f''(0) = -1.
        \]
        对(1)再求导得
        \[
        (1+x)^2 f''' + 3(1+x) f'' + 2 f' = 0
        \]
        于是
        \[
        f'''(0) + 3 f''(0) + 2 f'(0) = 0 \Rightarrow f'''(0) = 1.
        \]
        因此麦克劳林展开为
        \[
        f(x) = f(0) + f'(0) x + \frac{f''(0)}{2!} x^2 + \frac{f'''(0)}{3!} x^3 + \dots
        = x - \frac{1}{2} x^2 + \frac{1}{6} x^3 + \dots
        \]
    \end{solution}
    \part 据此,求 $\cos[\ln(1+x)]$ 的麦克劳林展开前 2 个非零项。
    \begin{solution}
    由(b),
        \[ 
        \sin[\ln(1+x)] = x - \frac{1}{2}x^2 + \frac{1}{6}x^3 + \dots 
        \]
        两边关于 $x$ 求导,
        \[ 
        \cos[\ln(1+x)] \cdot \frac{1}{1+x} = 1 - x + \frac{1}{2}x^2 + \dots 
        \]
        于是
        \begin{align*}
        \cos[\ln(1+x)] 
        &= (1+x)(1 - x + \frac{1}{2}x^2 + \dots) \\
        &= 1 - x + \frac{1}{2}x^2 + x - x^2 + \dots \\
        &= 1 - \frac{1}{2}x^2 + \dots
        \end{align*}
    \end{solution}
    \end{parts}

    \question 设 
    \[
    y = \tan x,\quad 0 \le x < \frac{\pi}{2}
    \]
    \begin{parts}
    \part 证明下列公式:
    \begin{subparts}
    \subpart 
    \[
    \frac{d^2y}{dx^2} = 2y \frac{dy}{dx}
    \]
    \begin{solution}
        对$x$求导得,
        \[
        \frac{dy}{dx} = \sec^2 x = 1 + \tan^2 x = 1+y^2
        \]
        再求导得
        \[
        y'' = 2y y'
        \]
    \end{solution}
    \subpart 
    \[
    \frac{d^5y}{dx^5} = 6 \left(\frac{d^2y}{dx^2}\right)^2 + 8 \frac{dy}{dx} \frac{d^3y}{dx^3} + 2y \frac{d^4y}{dx^4}
    \]
    \begin{solution}
        再次求导得,
        \[
        y''' = 2(y')^2 + 2yy''
        \]
        \[
        y^{(4)} = 4y'y''+2y'y'' + 2yy''' =6y'y''+ 2yy'''
        \]
        \[
        y^{(5)} = 6(y'')^2 + 6y'y''' + 2y'y''' + 2yy^{(4)} = 6(y'')^2 + 8y'y''' + 2y y^{(4)}
        \]
    \end{solution}
    \end{subparts}
    \part 求 $y = \tan x$ 的麦克劳林展开前 3 个非零项。
    \begin{solution}
        在 $x=0$ 处求各阶导数为
        \[
        y_0 = \tan(0) = 0, \quad y'_0 =1, \quad y''_0 = 0, \quad y'''_0 = 2, \quad y^{(4)}_0 = 0, \quad y^{(5)}_0 = 16.
        \]
        故麦克劳林展开式为
        \[
        \tan x = x + \frac{1}{3}x^3 + \frac{2}{15}x^5 + O(x^6)
        \]
    \end{solution}
    \end{parts}

    \question 设 
    \[
    y = \ln(2-e^x),\quad x<\ln 2
    \]
    \begin{parts}
    \part 证明下列公式:
    \[
    e^y \left[ \frac{d^3 y}{dx^3} + 3 \frac{dy}{dx} \frac{d^2 y}{dx^2} + \left(\frac{dy}{dx}\right)^3 \right] + e^x = 0
    \]
    \begin{solution}
        由 $y = \ln(2-e^x)$ 得
        \[
        e^y = 2-e^x
        \]
        两边对 $x$ 求导,
        \[
        e^y y' = -e^x \Rightarrow e^y y' + e^x = 0
        \]
        再次求导,
        \[
        e^y[(y')^2 + y''] + e^x = 0
        \]
        再求一次导数,
        \[
        e^y[y''' + 3y'y'' + (y')^3] + e^x = 0
        \]
    \end{solution}
    \part 求麦克劳林展开式的前三个非零项。
    \begin{solution}
        在 $x=0$ 处求各阶导数:
        \[
        y_0 = \ln(2-e^0)= 0
        \]
        \[
        y'_0 + 1 = 0 \Rightarrow y'_0 = -1
        \]
        \[
        (-1)^2 + y''_0 + 1 = 0 \Rightarrow y''_0 = -2
        \]
        \[
        (-1)^3 + 3(-1)(-2) + y'''_0 + 1 = 0 \Rightarrow y'''_0 = -6
        \]
        故麦克劳林展开式为
        \[
        \ln(2-e^x) = -x - x^2 - x^3 + O(x^4)
        \]
    \end{solution}
    \end{parts}

    \question 函数 $f,g$ 定义如下
    \[
    f(x) = \arctan\left(\frac{2}{3}x\right), \quad x\in \mathbb{R},
    \]
    \[
    g(y) = \frac{1}{1+y}, \quad y\in \mathbb{R}, \quad -1<y<1.
    \]
    \begin{parts}
    \part 将 $g(y)$ 展开为二项级数,并保留 $y^3$ 项。
    \begin{solution}
        二项展开得,
        \begin{align*}
        g(y) = (1+y)^{-1} 
        &= 1 + (-1)y + \frac{(-1)(-2)}{2!}y^2 + \frac{(-1)(-2)(-3)}{3!}y^3 + O(y^4) \\
        &= 1 - y + y^2 - y^3 + O(y^4)
        \end{align*}
    \end{solution}
    \part 利用 (a) 的结果,证明
    \[
    \arctan\left(\frac{2}{3}x\right) \approx \frac{2}{3}x - \frac{8}{81}x^3 + \frac{32}{1215}x^5 - \frac{128}{15309}x^7
    \]
    \begin{solution}
        对 $x$ 求导,
        \[
        f'(x) = \frac{\frac{2}{3}}{1+(\frac{2}{3}x)^2} = \frac{2}{3}\left(1+\frac{4}{9}x^2\right)^{-1}
        \]
        令$y = \dfrac{4}{9}x^2$,
        \begin{align*}
        f'(x) 
        &= \frac{2}{3}\left[1 - \frac{4}{9}x^2 + \left(\frac{4}{9}x^2\right)^2 - \left(\frac{4}{9}x^2\right)^3 + O(x^8)\right]\\
        &= \frac{2}{3} - \frac{8}{27}x^2 + \frac{32}{243}x^4 - \frac{128}{2187}x^6 + O(x^8)
        \end{align*}
        两边积分得
        \begin{align*}
        \arctan\left(\frac{2}{3}x\right) 
        &=\int \left[\frac{2}{3} - \frac{8}{27}x^2 + \frac{32}{243}x^4 - \frac{128}{2187}x^6 + O(x^8) \right] dx \\
        &= \frac{2}{3}x - \frac{8}{81}x^3 + \frac{32}{1215}x^5 - \frac{128}{15309}x^7 + O(x^9)+ C
        \end{align*}
        由 $x=0$ 可得 $C=0$,因此
        \[
        \arctan\left(\frac{2}{3}x\right) \approx \frac{2}{3}x - \frac{8}{81}x^3 + \frac{32}{1215}x^5 - \frac{128}{15309}x^7
        \]
    \end{solution}
    \end{parts}

    \question 已知 
    \[
    y = \tan x
    \]
    \begin{parts}
    \part 证明
    \[
    \frac{d^3 y}{dx^3} = 2y \frac{d^2 y}{dx^2} + 2 \left(\frac{dy}{dx}\right)^2
    \]
    \begin{solution}
        求导,
        \[
        \frac{dy}{dx} = \sec^2 x = 1 + \tan^2 x = 1 + y^2
        \]
        再求两次导数得
        \[
        \frac{d^2 y}{dx^2} = 2y \frac{dy}{dx},\quad \frac{d^3 y}{dx^3} = 2y \frac{d^2 y}{dx^2} + 2\left(\frac{dy}{dx}\right)^2
        \]
    \end{solution}
    \part 求 $\tan x$ 关于 $\left(x - \dfrac{\pi}{4}\right)$ 的泰勒级数展开式的前四项。
    \begin{solution}
        计算 $f(x)$ 及其各阶导数在 $x = \dfrac{\pi}{4}$ 的值:
        \[ 
        f\left(\frac{\pi}{4}\right) = \tan\frac{\pi}{4} = 1 
        \]
        \[ 
        f'(x) = \sec^2 x , \quad f'\left(\frac{\pi}{4}\right) = \sec^2\frac{\pi}{4} = (\sqrt{2})^2 = 2 
        \]
        \[ 
        f''(x) = 2\sec x \cdot (\sec x \tan x) = 2\sec^2 x \tan x, \quad f''\left(\frac{\pi}{4}\right) = 2(2)(1) = 4 
        \]
        \[ 
        f'''(x) = 4\sec x(\sec x \tan x)\tan x + 2\sec^2 x(\sec^2 x) = 4\sec^2 x \tan^2 x + 2\sec^4 x 
        \]
        \[ 
        f'''\left(\frac{\pi}{4}\right) = 4(2)(1)^2 + 2(2)^2 = 16 
        \]
        将上述数值代入泰勒公式:
        \begin{align*}
        \tan x &= 1 + 2\left(x - \frac{\pi}{4}\right) + \frac{4}{2!}\left(x - \frac{\pi}{4}\right)^2 + \frac{16}{3!}\left(x - \frac{\pi}{4}\right)^3 + \dots \\
        &= 1 + 2\left(x - \frac{\pi}{4}\right) + 2\left(x - \frac{\pi}{4}\right)^2 + \frac{8}{3}\left(x - \frac{\pi}{4}\right)^3 + \dots 
        \end{align*}
    \end{solution}
    \part 据此,求证
    \[
    \tan \frac{5\pi}{18} \approx 1+\frac{\pi}{18}+\frac{\pi^2}{648}+\frac{\pi^3}{17496}
    \]
    \begin{solution}
        令$x=\dfrac{5\pi}{18}$,则
        \[
        x-\frac{\pi}{4} = \frac{\pi}{36}
        \]
        代入泰勒展开得
        \[
        \tan \frac{5\pi}{18} \approx 1 + 2\cdot \frac{\pi}{36} + 2\cdot \left(\frac{\pi}{36}\right)^2 + \frac{8}{3} \cdot \left(\frac{\pi}{36}\right)^3
        \]
        即
        \[
        \tan \frac{5\pi}{18} \approx 1 + \frac{\pi}{18} + \frac{\pi^2}{648} + \frac{\pi^3}{17496}
        \]
    \end{solution}
    \end{parts}

    \question 求 $\arctan x$ 的麦克劳林展开式,并利用该展开证明
    \[
    \pi=\sum_{n=0}^{\infty} f(n)
    \]
    其中 $f(n)$ 为某函数。
    \begin{solution}
        考虑
        \[
        \frac{d}{dx}(\arctan x)=\frac{1}{1+x^{2}}
        \]
        当 $|x|<1$ 时,有几何级数展开
        \[
        \frac{1}{1+x^{2}}=1-x^{2}+x^{4}-x^{6}+x^{8}-\cdots=\sum_{n=0}^{\infty}(-1)^{n}x^{2n}
        \]
        对上式逐项积分,有
        \[
        \arctan x=\sum_{n=0}^{\infty}\int(-1)^{n}x^{2n}\,dx =\sum_{n=0}^{\infty}\frac{(-1)^{n}x^{2n+1}}{2n+1}+C
        \]
        令 $x=0$,则 $C=0$,于是
        \[
        \arctan x=\sum_{n=0}^{\infty}\frac{(-1)^{n}x^{2n+1}}{2n+1}
        \]
        代入 $x=1$ 得
        \[
        \arctan 1=\frac{\pi}{4}=\sum_{n=0}^{\infty}\frac{(-1)^{n}}{2n+1}
        \]
        即
        \[
        \pi=\sum_{n=0}^{\infty}\frac{4(-1)^{n}}{2n+1}
        \]
        其中
        \[
        f(n)=\frac{4(-1)^{n}}{2n+1}
        \]
    \end{solution}

    \question 选择适当的二项展开式,证明
    \[
    \arcsin x = \sum_{r=0}^{\infty} \left[ \binom{2r}{r}\frac{2}{2r+1}\left(\frac{x}{2}\right)^{2r+1} \right]
    \]
    \begin{solution}
        从二项式展开 $(1-x^2)^{-\frac{1}{2}}$ 出发,
        \[
        \frac{1}{\sqrt{1-x^2}} = (1-x^2)^{-\frac{1}{2}} = 1 + \frac{-\frac{1}{2}}{1}(-x^2) + \frac{-\frac{1}{2}(-\frac{3}{2})}{1\cdot 2}(-x^2)^2 + \frac{-\frac{1}{2}(-\frac{3}{2})(-\frac{5}{2})}{1\cdot 2\cdot 3}(-x^2)^3 + O(x^8)
        \]
        整理得
        \begin{align*}
        \frac{1}{\sqrt{1-x^2}} 
        &= 1 + \frac{2!}{(1!)^2} \frac{x^2}{2^2} + \frac{4!}{(2!)^2} \frac{x^4}{2^4} + \frac{6!}{(3!)^2} \frac{x^6}{2^6} + O(x^8)  \\
        &= \sum_{r=0}^{\infty} \left[ \frac{(2r)!}{(r!)^2} \left( \frac{x}{2} \right)^{2r} \right] \\
        &= \sum_{r=0}^{\infty} \left[ \binom{2r}{r} \left( \frac{x}{2} \right)^{2r} \right] 
        \end{align*}
        在收敛半径内,对等式两边积分:
        \[ 
        \int \frac{1}{\sqrt{1-x^2}} dx = \int \sum_{r=0}^{\infty} \left[ \binom{2r}{r} \frac{x^{2r}}{2^{2r}} \right] dx 
        \]
        得到
        \[ 
        \arcsin x = \sum_{r=0}^{\infty} \left[ \binom{2r}{r} \frac{x^{2r+1}}{2r+1} \cdot \frac{1}{2^{2r}} \right] + C =\sum_{r=0}^{\infty} \left[ \binom{2r}{r} \frac{2}{2r+1} \left( \frac{x}{2} \right)^{2r+1} \right] +C
        \]
        当 $x=0$ 时,故 $C=0$,因此
        \[ 
        \arcsin x = \sum_{r=0}^{\infty} \left[ \binom{2r}{r} \frac{2}{2r+1} \left( \frac{x}{2} \right)^{2r+1} \right] 
        \]
    \end{solution}

    \question 求下列无穷级数的和
    \[
    1+\frac{1}{3\cdot4}+\frac{1}{5\cdot4^2}+\frac{1}{7\cdot4^3}+\frac{1}{9\cdot4^4}+\cdots
    \]
    \begin{solution}
        考虑$\ln(1+x)$及$\ln(1-x)$的幂展开:
        \[
        \ln(1+x) = x - \frac{1}{2}x^2 + \frac{1}{3}x^3 - \frac{1}{4}x^4 + \dots
        \]
        \[
        \ln(1-x) = -x - \frac{1}{2}x^2 - \frac{1}{3}x^3 - \frac{1}{4}x^4 + \dots
        \]
        两式相减得
        \[
        \ln(1+x) - \ln(1-x) = 2\left( x + \frac{x^3}{3} + \frac{x^5}{5} + \frac{x^7}{7} + \dots \right)
        \]
        即
        \[
        \ln\left(\frac{1+x}{1-x}\right) = 2 \sum_{k=0}^{\infty} \frac{x^{2k+1}}{2k+1}
        \]
        在收敛半径内令 $x = \frac{1}{2}$,则
        \[
        \ln 3 = 2 \sum_{k=0}^{\infty} \frac{1}{(2k+1)2^{2k+1}} = \sum_{k=0}^{\infty} \frac{1}{(2k+1)4^k}
        \]
        即所求级数的和为$\ln 3$。
    \end{solution}

    \question 求下列交错级数的和
    \[
    \frac{1}{1} - \frac{1}{1+4} + \frac{1}{1+4+9} - \frac{1}{1+4+9+16} + \frac{1}{1+4+9+16+25} - \cdots
    \]
    \begin{solution}
        写成
        \[
        \sum_{n=1}^{\infty} \frac{(-1)^{n+1}}{1^2 + 2^2 + \cdots + n^2} = \sum_{n=1}^{\infty} \frac{(-1)^{n+1}}{\frac{1}{6}n(n+1)(2n+1)} = \sum_{n=1}^{\infty} \frac{6(-1)^{n+1}}{n(n+1)(2n+1)}
        \]
        部分分式分解得
        \[
        \frac{1}{n(n+1)(2n+1)} = \frac{1}{n} + \frac{1}{n+1} - \frac{4}{2n+1}
        \]
        因此级数可写为
        \[
        6\sum_{n=1}^{\infty} \frac{(-1)^{n+1}}{n} + 6\sum_{n=1}^{\infty} \frac{(-1)^{n+1}}{n+1} - 24\sum_{n=1}^{\infty} \frac{(-1)^{n+1}}{2n+1}
        \]
        分别处理各个级数
        \[
        \sum_{n=1}^{\infty} \frac{(-1)^{n+1}}{n} = \ln 2, \quad 6\sum_{n=1}^{\infty} \frac{(-1)^{n+1}}{n+1} = 6 - 6 \ln 2
        \]
        而由 $\arctan x$ 的展开,
        \[
        \arctan x = \sum_{n=0}^{\infty} \frac{(-1)^n x^{2n+1}}{2n+1}  
        \]
        令$x=1$,易知
        \[
        \sum_{n=0}^{\infty} \frac{(-1)^n}{2n+1} = \frac{\pi}{4}
        \]
        因此原级数的和为
        \[
        \frac{1}{1} - \frac{1}{1+4} + \frac{1}{1+4+9} - \cdots = 6\ln 2 + (6 - 6\ln 2) - 24\left(\frac{\pi}{4} - 1\right) = 6(\pi - 3)
        \]
    \end{solution}

    \question 证明
    \[
    \sum_{r=1}^{\infty}\left[\frac{2r+3}{(r+1)3^r}\right]=3\ln\frac{3}{2}
    \]
    \begin{solution}
        首先注意到
        \[
        \frac{2r+3}{(r+1)3^r}
        =2\left(\frac{1}{3}\right)^r+\frac{1}{r+1}\left(\frac{1}{3}\right)^r
        \]
        因此
        \[
        \sum_{r=1}^{\infty}\left[\frac{2r+3}{(r+1)3^r}\right]
        =\sum_{r=1}^{\infty}2\left(\frac{1}{3}\right)^r
        +\sum_{r=1}^{\infty}\frac{1}{r+1}\left(\frac{1}{3}\right)^r
        \]
        第一项是首项为 $\dfrac{2}{3}$,公比为 $\dfrac{1}{3}$的等比级数,故
        \[
        \sum_{r=1}^{\infty}2\left(\frac{1}{3}\right)^r
        =\frac{\frac{2}{3}}{1-\frac{1}{3}}=1
        \]
        对第二项,考虑
        \[
        \ln(1-x)=-x-\frac{1}{2}x^2-\frac{1}{3}x^3-\frac{1}{4}x^4-\cdots
        \]
        两边除以 $-x$,得
        \[
        -\frac{1}{x}\ln(1-x)
        =1+\frac{1}{2}x+\frac{1}{3}x^2+\frac{1}{4}x^3+\cdots
        =1+\sum_{r=1}^{\infty}\frac{1}{r+1}x^r
        \]
        取 $x=\dfrac{1}{3}$,由于 $|x|<1$,级数收敛,于是
        \[
        \sum_{r=1}^{\infty}\frac{1}{r+1}\left(\frac{1}{3}\right)^r
        =-3\ln\frac{2}{3}-1
        \]
        故
        \[
        \sum_{r=1}^{\infty}\left[\frac{2r+3}{(r+1)3^r}\right]
        =1-3\ln\frac{2}{3}-1
        =3\ln\frac{3}{2}
        \]
    \end{solution}

    \question 求下列级数的值:
    \[
    \sum_{r=1}^{\infty} \left[\left(\frac{1}{2r}+\frac{1}{2r+1}\right)\left(\frac{1}{4}\right)^r\right]
    \]
    \begin{solution}
        由
        \[
        \ln(1+x)=x-\frac{1}{2}x^2+\frac{1}{3}x^3-\cdots,\quad |x|<1
        \]
        \[
        \ln(1-x)=-x-\frac{1}{2}x^2-\frac{1}{3}x^3-\cdots,\quad |x|<1
        \]
        两式相加得
        \[
        \ln(1-x^2)= -x^2-\frac{1}{2}x^4-\frac{1}{3}x^6-\cdots
        \]
        即
        \[
        -\frac{1}{2}\ln(1-x^2)=\frac{1}{2}x^2+\frac{1}{4}x^4+\frac{1}{6}x^6+\cdots
        \]
        令 $x=\dfrac{1}{2}$,得
        \[
        \sum_{r=1}^{\infty} \frac{1}{2r} \cdot \left(\frac{1}{4}\right)^r
        =-\frac{1}{2}\ln\frac{3}{4}
        \]
        同理,两式相减得
        \[
        \ln\left(\frac{1+x}{1-x}\right)
        =2\left(x+\frac{x^3}{3}+\frac{x^5}{5}+\cdots\right)
        \]
        从而
        \[
        \frac{1}{2x}\ln\left(\frac{1+x}{1-x}\right)
        =1+\frac{x^2}{3}+\frac{x^4}{5}+\cdots
        \]
        令 $x=\dfrac{1}{2}$,得
        \[
        \sum_{r=1}^{\infty} \frac{1}{2r+1} \cdot \left(\frac{1}{4}\right)^r
        =\ln 3-1
        \]
        将两部分相加,
        \[
        \sum_{r=1}^{\infty} \left[\left(\frac{1}{2r}+\frac{1}{2r+1}\right)\left(\frac{1}{4}\right)^r\right]
        =-\frac{1}{2}\ln\frac{3}{4}+\ln 3-1=\frac{1}{2}\ln 12-1
        \]
    \end{solution}

    \question
    \begin{parts} 
    \part 以适当的积分方法计算
    \[
    \int_{0}^{1} x^{3}\arctan x\,dx
    \]
    \begin{solution}
        先用分部积分法,
        \[
        \int_{0}^{1} x^{3}\arctan x\,dx
        =\left[\frac{x^{4}}{4}\arctan x\right]_{0}^{1}
        -\frac{1}{4}\int_{0}^{1}\frac{x^{4}}{1+x^{2}}\,dx
        \]
        其中
        \begin{align*}
        \int_{0}^{1}\frac{x^{4}}{1+x^{2}}\,dx
        &=\int_{0}^{1}\left(x^{2}-1+\frac{1}{1+x^{2}}\right)\,dx \\
        &= \left[\frac{x^{3}}{3}-x+\arctan x\right]_{0}^{1} \\
        &=\frac{1}{3}-1+\frac{\pi}{4} 
        \end{align*}
        代回得
        \[
        \int_{0}^{1} x^{3}\arctan x\,dx
        =\frac{\pi}{16}-\frac{1}{4}\left(\frac{1}{3}-1+\frac{\pi}{4}\right)=\frac{1}{6}
        \]
    \end{solution}
    \part 求 $\arctan x$ 的无穷级数展开,并利用该展开验证 (a) 的结果
    \begin{solution}
        由
        \[
        \frac{1}{1+x^{2}}=1-x^{2}+x^{4}-x^{6}+\cdots
        \]
        两边积分得
        \[
        \arctan x
        =x-\frac{x^{3}}{3}+\frac{x^{5}}{5}-\frac{x^{7}}{7}+\cdots
        =\sum_{n=0}^{\infty}\frac{(-1)^{n}x^{2n+1}}{2n+1}
        \]
        根据富比尼定理（在满足一致收敛或绝对收敛的条件下）,可以将积分号与求和号交换。代入积分式得
        \begin{align*}
        \int_{0}^{1} x^{3}\arctan x\,dx
        &=\int_{0}^{1}x^{3}\sum_{n=0}^{\infty}\frac{(-1)^{n}x^{2n+1}}{2n+1}\,dx\\
        &=\sum_{n=0}^{\infty}\frac{(-1)^{n}}{2n+1}\int_{0}^{1}x^{2n+4}\,dx\\
        &=\sum_{n=0}^{\infty}\frac{(-1)^{n}}{(2n+1)(2n+5)}
        \end{align*}
        部分分式分解,
        \[
        \frac{1}{(2n+1)(2n+5)}
        =\frac{1}{4}\left(\frac{1}{2n+1}-\frac{1}{2n+5}\right)
        \]
        于是由累加法,
        \[
        \int_{0}^{1} x^{3}\arctan x\,dx
        =\frac{1}{4}\sum_{n=0}^{\infty}(-1)^{n}
        \left(\frac{1}{2n+1}-\frac{1}{2n+5}\right)=\frac{1}{4}\left(1-\frac{1}{3}\right)
        =\frac{1}{6}
        \]
        与 (a) 的结果一致。
    \end{solution}
    \end{parts}

    \question
    已知
    \[
    \sum_{n=0}^{\infty}\frac{(-1)^n}{2n+1}=\frac{\pi}{4},
    \quad
    \sum_{n=0}^{\infty}\frac{(-1)^n}{n+1}=\ln 2,
    \quad
    \sum_{n=0}^{\infty}\frac{(-1)^n}{(n+1)^2}=\frac{\pi^2}{12}.
    \]
    求
    \[
    \int_0^1 (\ln x)(\arctan x)\,dx
    \]
    的值。
    \begin{solution}
        考虑$\arctan x$的幂级数展开,由
        \[
        \frac{d}{dx}(\arctan x)=\frac{1}{1+x^2}=1-x^2+x^4-x^6+\cdots
        \]
        逐项积分得
        \[
        \arctan x=x-\frac{1}{3}x^3+\frac{1}{5}x^5-\frac{1}{7}x^7+\cdots
        =\sum_{n=0}^{\infty}\frac{(-1)^n x^{2n+1}}{2n+1},
        \quad 0\le x\le1.
        \]
        将其代入积分,并交换求和与积分次序:
        \[
        \int_0^1 (\arctan x)(\ln x)\,dx
        =\sum_{n=0}^{\infty}\frac{(-1)^n}{2n+1}
        \int_0^1 x^{2n+1}\ln x\,dx.
        \]
        作分部积分,取$u=\ln x,dv=x^{2n+1}dx$,则
        \[
        \int_0^1 x^{2n+1}\ln x\,dx
        =\left[\frac{x^{2n+2}\ln x}{2n+2}\right]_0^1
        -\frac{1}{2n+2}\int_0^1 x^{2n+1}\,dx.
        \]
        化简得
        \[
        \int_0^1 x^{2n+1}\ln x\,dx=-\frac{1}{(2n+2)^2}.
        \]
        于是
        \[
        \int_0^1 (\arctan x)(\ln x)\,dx
        =-\sum_{n=0}^{\infty}\frac{(-1)^n}{(2n+1)(2n+2)^2}.
        \]
        作部分分式分解
        \[
        \frac{1}{(2n+1)(2n+2)^2}
        =\frac{1}{2n+1}-\frac{1}{2n+2}-\frac{1}{(2n+2)^2}.
        \]
        代回得
        \[
        \int_0^1 (\arctan x)(\ln x)\,dx=-\sum_{n=0}^{\infty}(-1)^n\left[\frac{1}{2n+1}-\frac{1}{2n+2}-\frac{1}{(2n+2)^2}\right]
        \]
        将其拆分为三个已知级数:
        \[
        \sum_{n=0}^{\infty}\frac{(-1)^n}{2n+1}=\frac{\pi}{4},
        \quad
        \sum_{n=0}^{\infty}\frac{(-1)^n}{n+1}=\ln 2,
        \quad
        \sum_{n=0}^{\infty}\frac{(-1)^n}{(n+1)^2}=\frac{\pi^2}{12}.
        \]
        因此
        \[
        \int_0^1 (\arctan x)(\ln x)\,dx
        =-\left[\frac{\pi}{4}-\frac{1}{2}\ln 2-\frac{1}{4}\cdot\frac{\pi^2}{12}\right]
        =\frac{\pi^2-12\pi+24\ln 2}{48}
        \]
    \end{solution}
\begin{solution}
    使用莱布尼兹积分法则。设
    \[
    I(a)=\int_0^1 (\ln x)\arctan(ax)\,dx,
    \quad a\ge0.
    \]
    所求积分为 $I(1)$,对参数 $a$ 求导,有
    \[
    I'(a)=\int_0^1 (\ln x)\frac{x}{1+a^2x^2}\,dx.
    \]
    交换积分顺序,令 $u=ax$,得
    \[
    I'(a)=\int_0^1 x\ln x
    \int_0^{a}\frac{du}{1+u^2x^2}\,dx.
    \]
    改为先对 $x$ 积分。注意到
    \[
    \int_0^1 \frac{x\ln x}{1+u^2x^2}\,dx
    =-\frac{1}{2u^2}\ln(1+u^2),
    \]
    因此
    \[
    I'(a)
    =-\frac12\int_0^{a}\frac{\ln(1+u^2)}{u^2}\,du.
    \]
    对该积分作分部积分,取
    $U=\ln(1+u^2)$,$dV=\frac{du}{u^2}$,
    则
    \[
    I'(a)
    =-\frac12
    \left[
    -\frac{\ln(1+u^2)}{u}
    +\int_0^{a}\frac{2u}{1+u^2}\cdot\frac{1}{u}\,du
    \right].
    \]
    化简得
    \[
    I'(a)
    =\frac12\frac{\ln(1+a^2)}{a}
    -\int_0^{a}\frac{du}{1+u^2}.
    \]
    于是
    \[
    I'(a)=\frac12\frac{\ln(1+a^2)}{a}-\arctan a.
    \]
    再对 $a$ 从 $0$ 积分到 $1$,注意 $I(0)=0$,得到
    \[
    I(1)=\int_0^1
    \left(
    \frac12\frac{\ln(1+a^2)}{a}-\arctan a
    \right)da.
    \]
    分项计算。首先
    \[
    \int_0^1 \arctan a\,da
    =\frac{\pi}{4}-\frac12\ln2.
    \]
    其次
    \[
    \int_0^1 \frac{\ln(1+a^2)}{a}\,da
    =\int_0^1\int_0^{a^2}\frac{1}{a(1+t)}\,dt\,da
    =\int_0^1\frac{1-a^2}{a(1+a^2)}\,da
    =\frac{\pi^2}{12}
    \]
    综上
    \[
    I(1)
    =\frac12\cdot\frac{\pi^2}{12}
    -\left(\frac{\pi}{4}-\frac12\ln2\right)
    \]
    因此
    \[
    \int_0^1 (\ln x)(\arctan x)\,dx
    =\frac{\pi^2}{48}+\frac12\ln2-\frac{\pi}{4}
    =\frac{\pi^2-12\pi+24\ln2}{48}
    \]
    \textcolor{red}{(待验证)}
\end{solution}

\end{questions}
\pagebreak

\begin{center}
  {\fontsize{30pt}{26pt}\selectfont
    \hypertarget{矩阵}{矩阵} \label{矩阵}
  }
\end{center}
\separator
\vspace{1pt}
\begin{questions}
    \question 已知
    \[
    A=\begin{bmatrix}1&0&0\\ -2&3&0\\ 0&-4&5\end{bmatrix},\quad
    I=\begin{bmatrix}1&0&0\\ 0&1&0\\ 0&0&1\end{bmatrix},\quad
    B=(I+A)^{-1}(I-A),
    \]
    求矩阵 $(I+B)^{-1}$。
    \begin{solution}
        由 $B=(I+A)^{-1}(I-A)$ 可得
        \[
        (I+A)B=I-A
        \]
        两边加上 $I+A$ 得
        \[
        I+A+B+AB =(I+A)(I+B)=2I
        \]
        因此
        \[
        (I+B)^{-1}=\frac{1}{2}(I+A)=\frac{1}{2}\begin{bmatrix}2&0&0\\ -2&4&0\\ 0&-4&6\end{bmatrix}
        =\begin{bmatrix}1&0&0\\ -1&2&0\\ 0&-2&3\end{bmatrix}
        \]
    \end{solution}

    \question 已知
    \[
    \begin{bmatrix} a & b \\ c & d \end{bmatrix} 
    \begin{bmatrix} p & q \\ r & s \end{bmatrix} 
    = \begin{bmatrix} 5 & 3 \\ 3 & 2 \end{bmatrix}
    \]
    且联立方程
    \[
    \begin{cases} ax+by=5\\ cx+dy=-3 \end{cases}
    \]
    恰有一组解 $(x,y)=(1,2)$,求联立方程
    \[
    \begin{cases} pu+qv=1\\ ru+sv=2 \end{cases}
    \]
    的解 $(u,v)$。
    \begin{solution}
        已知
        \[
        \begin{bmatrix} a & b \\ c & d \end{bmatrix} \begin{bmatrix} 1 \\ 2 \end{bmatrix} = \begin{bmatrix} 5 \\ -3 \end{bmatrix} 
        \]
        且
        \[
        \begin{bmatrix} a & b \\ c & d \end{bmatrix} 
        \begin{bmatrix} p & q \\ r & s \end{bmatrix} 
        = \begin{bmatrix} 5 & 3 \\ 3 & 2 \end{bmatrix} 
        \Rightarrow 
        \begin{bmatrix} p & q \\ r & s \end{bmatrix}^{-1} = 
        \begin{bmatrix} 2 & -3 \\ -3 & 5 \end{bmatrix} 
        \begin{bmatrix} a & b \\ c & d \end{bmatrix}
        \]
        因此
        \[
        \begin{bmatrix} p & q \\ r & s \end{bmatrix}^{-1} \begin{bmatrix} 1 \\ 2 \end{bmatrix} 
        = \begin{bmatrix} 2 & -3 \\ -3 & 5 \end{bmatrix} \begin{bmatrix} 5 \\ -3 \end{bmatrix} 
        = \begin{bmatrix} 19 \\ -30 \end{bmatrix}
        \]
        即
        \[
        \begin{bmatrix} p & q \\ r & s \end{bmatrix} \begin{bmatrix} 19 \\ -30 \end{bmatrix} = \begin{bmatrix} 1 \\ 2 \end{bmatrix}
        \]
        故联立方程的解为
        \[
        (u,v) = (19,-30)
        \]
    \end{solution}

    \question 设三阶方阵 
    \[
    A = \begin{bmatrix}1 & a & b \\ 0 & 1 & a \\ 0 & 0 & 1\end{bmatrix},\quad 
    A^{10} = \begin{bmatrix}1 & ka & pa^2 + qb \\ 0 & 1 & ka \\ 0 & 0 & 1\end{bmatrix},
    \]
    其中$k,p,q$为常数,试求 $k+p+q$。
    \begin{solution}
        将 $A$ 分解为 $A = I + B$,其中
        \[
        I = \begin{bmatrix}1 & 0 & 0 \\ 0 & 1 & 0 \\ 0 & 0 & 1\end{bmatrix},\quad
        B = \begin{bmatrix}0 & a & b \\ 0 & 0 & a \\ 0 & 0 & 0\end{bmatrix}.
        \]
        注意到 $B^3 = 0$,且
        \[
        B^2 = \begin{bmatrix}0 & 0 & a^2 \\ 0 & 0 & 0 \\ 0 & 0 & 0\end{bmatrix}.
        \]
        由二项式定理,
        \[
        A^{10} = (I + B)^{10} = \sum_{k=0}^{10} \comb{10}{k} B^k I^{10-k} = I + 10B + 45 B^2=\begin{bmatrix}1 & 10a & 45a^2+10b \\ 0 & 1 & 10a \\ 0 & 0 & 1\end{bmatrix}
        \]
        所以
        \[
        k = 10,\quad p = 45,\quad q = 10 \quad \Rightarrow \quad k+p+q = 65.
        \]
    \end{solution}

    \question 设实数 $a>b$,且有二阶方阵 $X,Y$ 满足
    \[
    X + Y = I, \quad XY = O,
    \]
    其中
    \[
    I = \begin{bmatrix}1 & 0 \\ 0 & 1 \end{bmatrix}, \quad O = \begin{bmatrix}0 & 0 \\ 0 & 0 \end{bmatrix},
    \]
    且
    \[
    A = \begin{bmatrix} 2 & 4 \\ 1 & -1 \end{bmatrix} = aX + bY,
    \]
    求 $a,b$ 的值。据此,求$X^{2025}-Y^{2025}$。
    \begin{solution}
        首先有
        \[
        A = aX + bY = aX + b(I - X) = bI + (a - b)X,
        \]
        故
        \[
        X = \frac{A - bI}{a - b}, \quad Y = \frac{aI - A}{a - b}.
        \]
        所以
        \[
        XY = \frac{1}{(a - b)^2} (A - bI)(aI - A) = O,
        \]
        即
        \[
        \begin{bmatrix} 2 - b & 4 \\ 1 & -1 - b \end{bmatrix}
        \begin{bmatrix} a - 2 & -4 \\ -1 & a + 1 \end{bmatrix}
        = \begin{bmatrix} 0 & 0 \\ 0 & 0 \end{bmatrix}.
        \]
        得到方程组
        \[
        \begin{cases}
        (2 - b)(a - 2) + 4(-1) = 0, \\
        (2 - b)(-4) + 4(a + 1) = 0, \\
        1 \cdot (a - 2) + (-1 - b)(-1) = 0, \\
        1 \cdot (-4) + (-1 - b)(a + 1) = 0.
        \end{cases}
        \]
        由$a>b$,解得
        \[
        a = 3 ,b = -2 
        \]
        因此
        \[
        X=\frac{1}{5}\begin{bmatrix}4&4\\1&1\end{bmatrix},\quad
        Y=I-X=\frac{1}{5}\begin{bmatrix}1&-4\\-1&4\end{bmatrix}
        \]
        且发现 $X^2=X,\ Y^2=Y$,即$X,Y$是幂等矩阵,故对任意正整数 $n$,
        \[
        X^n=X,\quad Y^n=Y
        \]
        因此
        \[
        X^{2025}-Y^{2025}=X-Y=\frac{1}{5}\begin{bmatrix}3&8\\2&-3\end{bmatrix}
        \]
    \end{solution}

    \question 设 $$A=\begin{bmatrix} 4 & 4 \\ -1 & -1 \end{bmatrix}$$满足 $$A+A^2+\cdots+A^n = \begin{bmatrix} 2(3^n-1) & a \\ b & c \end{bmatrix},$$求 $b+c$。
    \begin{solution}
        发现
        \[
        A^2=\begin{bmatrix}4&4\\-1&-1\end{bmatrix}^2
        =\begin{bmatrix}12&12\\-3&-3\end{bmatrix}=3A,
        \]
        于是对任意 $k\ge1$ 有 $A^k=3^{\,k-1}A$,因此
        \[
        A+\cdots+A^n
        =A\sum_{k=0}^{n-1}3^k
        =\frac{3^n-1}{2}\,A
        =\begin{bmatrix} 2(3^n-1) & 2(3^n-1) \\ \dfrac12(1-3^n) & \dfrac12(1-3^n) \end{bmatrix}
        \]
        故
        \[
        b+c = \frac12(1-3^n) + \frac12(1-3^n) = 1 - 3^n
        \]
    \end{solution}
    \begin{solution}
        又解:对角化$A$可得
        \[
        A=\begin{bmatrix}4 & 4 \\ -1 & -1 \end{bmatrix} 
        =\begin{bmatrix}-1 & -4 \\ 1 & 1 \end{bmatrix} 
        \begin{bmatrix}0 & 0 \\ 0 & 3 \end{bmatrix} 
        \begin{bmatrix} \frac13 & \frac43 \\ -\frac13 & -\frac13 \end{bmatrix} \equiv P D P^{-1}
        \]
        由$A^n = P D^n P^{-1}$,
        \[
        A + A^2 + \cdots + A^n = P(D + D^2 + \cdots + D^n)P^{-1}
        \]
        \[
        = P \begin{bmatrix}0 & 0 \\ 0 & 3+3^2+\cdots+3^n \end{bmatrix} P^{-1} 
        = P \begin{bmatrix}0 & 0 \\ 0 & \dfrac{3^{n+1}-3}{2} \end{bmatrix} P^{-1}
        = \begin{bmatrix} 2(3^n-1) & 2(3^n-1) \\ \dfrac12(1-3^n) & \dfrac12(1-3^n) \end{bmatrix}
        \]
        同上。
    \end{solution}

    \question 设矩阵 $$A=\begin{pmatrix} 7 & -8 \\ -7 & 8 \end{pmatrix}$$ 满足 $$(I+A)^n = I + a_n A,$$其中 $n$ 为自然数,$$I=\begin{pmatrix} 1 & 0 \\ 0 & 1 \end{pmatrix},$$且 $\{a_n\}$ 为一个数列,求 $a_n$ 的通项公式。
    \begin{solution}
        求 $A$ 的特征值:
        \[
        \det(A - \lambda I) = \begin{vmatrix} 7 - \lambda & -8 \\ -7 & 8 - \lambda \end{vmatrix} = \lambda^2 - 15\lambda = 0 \Rightarrow \lambda = 0,\,15
        \]
        设 
        \[
        f(\lambda) = (1+\lambda)^n = (\lambda^2 - 15\lambda)p(\lambda) + a\lambda + b,
        \]令 $\lambda=0,15$,解得
        \[
        a = \frac{16^n - 1}{15},\quad b = 1
        \]
        由凯莱-哈密顿定理,
        \[
        (I + A)^n = aA + bI = \frac{16^n - 1}{15}A + I
        \Rightarrow a_n = \frac{1}{15}(16^n - 1)
        \]
    \end{solution}

    \question 设矩阵 
    \[
    A = \begin{bmatrix} -5 & -4 \\ 9 & 7 \end{bmatrix}, 
    \] 
    求 
    \[
    A^{51} - A^{50} + A^3 - 3A^2 - 2A + 4I_2,
    \] 
    其中 \(I_2 = \begin{bmatrix} 1 & 0 \\ 0 & 1 \end{bmatrix}\)。
    \begin{solution}
        发现 
        \[
        (A-I)^2 = {\begin{bmatrix}-6 & -4 \\ 9 & 6\end{bmatrix}}^2=\begin{bmatrix}0 & 0 \\ 0 & 0\end{bmatrix} \Rightarrow A^2=2A-I
        \] 
        且考虑方程
        \[
        \lambda^{50} = q(\lambda)(\lambda-1)^2 + c_1 \lambda + c_0 \tag{1}
        \]
        对 \(\lambda\) 求导,可得
        \[
        50\lambda^{49}=q'(\lambda)(\lambda-1)^2+2q(\lambda)(\lambda-1)+c_1 \tag{2}
        \]
        由$(1),(2),$代入 \(\lambda=1\) 得 
        \[
        c_0=-49, \ c_1=50
        \]
        由凯莱-哈密顿定理,
        \[
        A^{50} = 50A - 49I
        \] 
        因此 
        \[
        A^{51} - A^{50} = (50A^2-49A) - (50A - 49I) = 100A-50I -99A - 49I=A-I
        \]
        又 
        \[
        A^3 - 3A^2 - 2A + 4I = (A-I)^3-5(A-I)=-5(A-I)
        \] 
        故
        \[
        A^{51} - A^{50} + A^3 - 3A^2 - 2A + 4I = -4(A-I) = \begin{bmatrix} 24 & 16 \\ -36 & -24 \end{bmatrix}
        \]
    \end{solution}

    \question 已知
    \[
    \begin{bmatrix}
    a & b & c \\
    d & e & f \\
    g & h & i
    \end{bmatrix}^{101}
    =
    \begin{bmatrix}
    -1 & -2 & -2 \\
    1 & 2 & 1 \\
    -1 & -1 & 0
    \end{bmatrix}\\[2mm]
    \]
    求 $a + b + c + d + e + f + g + h + i$。
    \begin{solution}
        设
        \[
        A = \begin{bmatrix}
        -1 & -2 & -2 \\
        1 & 2 & 1 \\
        -1 & -1 & 0
        \end{bmatrix}
        \]
        解 $\det(A - \lambda I) = 0$,即
        \[
        \begin{vmatrix}
        -1-\lambda & -2 & -2 \\
        1 & 2-\lambda & 1 \\
        -1 & -1 & -\lambda
        \end{vmatrix}
        = -(\lambda - 1)^2(\lambda + 1) =0 \Rightarrow \lambda = 1 \ (m=2), \ \lambda = -1
        \]
        对于 $\lambda = 1$,解 $(A - I)\textbf{x} = \textbf{0}$:
        \[
        A - I =
        \begin{bmatrix}
        -2 & -2 & -2 \\
        1 & 1 & 1 \\
        -1 & -1 & -1
        \end{bmatrix}
        \Rightarrow
        \begin{bmatrix}
        1 & 1 & 1 \\
        0 & 0 & 0 \\
        0 & 0 & 0
        \end{bmatrix}
        \Rightarrow x_1 = -x_2 - x_3,\ x_2,\ x_3 \in \mathbb{R}
        \]
        取两个线性无关解:
        \[
        \vec{v}_1 =
        \begin{bmatrix}
        -1 \\ 1 \\ 0
        \end{bmatrix},\quad
        \vec{v}_2 =
        \begin{bmatrix}
        -1 \\ 0 \\ 1
        \end{bmatrix}
        \]
        对于 $\lambda = -1$,解 $(A + I)\textbf{x} = \textbf{0}$:
        \[
        A + I =
        \begin{bmatrix}
        0 & -2 & -2 \\
        1 & 3 & 1 \\
        -1 & -1 & 1
        \end{bmatrix}
        \Rightarrow
        \begin{bmatrix}
        1 & 3 & 1 \\
        0 & 1 & 1 \\
        0 & 0 & 0
        \end{bmatrix} \Rightarrow x_1 = x_3,\ x_2 = -x_3
        \]
        取
        \[
        \vec{v}_3 =
        \begin{bmatrix}
        2 \\ -1 \\ 1
        \end{bmatrix}
        \]
        于是
        \[
        P =
        \begin{bmatrix}
        -1 & -1 & 2 \\
        1 & 0 & -1 \\
        0 & 1 & 1
        \end{bmatrix}, \quad
        D =
        \begin{bmatrix}
        1 & 0 & 0 \\
        0 & 1 & 0 \\
        0 & 0 & -1
        \end{bmatrix}
        \]
        现计算逆矩阵 \(P^{-1}\),由$|P|=2$得
        \[
        P^{-1} 
        =\frac{1}{2}
        \begin{bmatrix}
        \begin{vmatrix} 0 & -1 \\ 1 & 1 \end{vmatrix} & -\begin{vmatrix} 1 & -1 \\ 0 & 1 \end{vmatrix} & \begin{vmatrix} 1 & 0 \\ 0 & 1 \end{vmatrix} \\
        -\begin{vmatrix} -1 & 2 \\ 1 & 1 \end{vmatrix} & \begin{vmatrix} -1 & 2 \\ 0 & 1 \end{vmatrix} & -\begin{vmatrix} -1 & -1 \\ 0 & 1 \end{vmatrix} \\
        \begin{vmatrix} -1 & 2 \\ 0 & -1 \end{vmatrix} & -\begin{vmatrix} -1 & 2 \\ 1 & -1 \end{vmatrix} & \begin{vmatrix} -1 & -1 \\ 1 & 0 \end{vmatrix}
        \end{bmatrix}^T
        =\begin{bmatrix}
        \frac{1}{2} & \frac{3}{2} & \frac{1}{2} \\
        -\frac{1}{2} & -\frac{1}{2} & \frac{1}{2} \\
        \frac{1}{2} & \frac{1}{2} & \frac{1}{2}
        \end{bmatrix}
        \]
        于是
        \[
        A = P D P^{-1} =
        \begin{bmatrix}
        -1 & -1 & 1 \\
        1 & 0 & -1 \\
        0 & 1 & 1
        \end{bmatrix}
        \begin{bmatrix}
        1 & 0 & 0 \\
        0 & 1 & 0 \\
        0 & 0 & -1
        \end{bmatrix}
        \begin{bmatrix}
        \frac{1}{2} & \frac{3}{2} & \frac{1}{2} \\
        -\frac{1}{2} & -\frac{1}{2} & \frac{1}{2} \\
        \frac{1}{2} & \frac{1}{2} & \frac{1}{2}
        \end{bmatrix}
        \]
        所以
        \[
        \begin{bmatrix}
        a & b & c \\
        d & e & f \\
        g & h & i
        \end{bmatrix}
        =
        \begin{bmatrix}
        -1 & -1 & 2 \\
        1 & 0 & -1 \\
        0 & 1 & 1
        \end{bmatrix}
        \begin{bmatrix}
        1^{\frac{1}{101}} & 0 & 0 \\
        0 & 1^{\frac{1}{101}} & 0 \\
        0 & 0 & (-1)^{\frac{1}{101}}
        \end{bmatrix}
        \begin{bmatrix}
        \frac{1}{2} & \frac{3}{2} & \frac{1}{2} \\
        -\frac{1}{2} & -\frac{1}{2} & \frac{1}{2} \\
        \frac{1}{2} & \frac{1}{2} & \frac{1}{2}
        \end{bmatrix}
        =
        \begin{bmatrix}
        -1 & -2 & -2 \\
        1 & 2 & 1 \\
        -1 & -1 & 0
        \end{bmatrix}
        \]
        恰好也是原矩阵,因此
        \[
        a + b + c + d + e + f + g + h + i = -3
        \]
    \end{solution}

    \question 设 $A,B$ 为 $n\times n$ 实矩阵,且满足
    \[
    AB + A + B = 0,
    \]
    证明 $AB = BA$。
    \begin{solution}
        注意到
        \[
        (A+I)(B+I) = AB + A + B + I = I
        \]
        其中 $I$ 为单位矩阵,因此 $A+I$ 及 $B+I$ 互为逆矩阵,据此有
        \[
        (A+I)(B+I) = (B+I)(A+I) = I
        \]
        展开等式得到
        \[
        AB + A + B + I = BA + B + A + I \Rightarrow AB = BA
        \]
    \end{solution}

    \question 若$n$阶方阵 $A,B,A+B$ 皆可逆,证明 $A^{-1}+B^{-1}$ 也是可逆矩阵。
    \begin{solution}
        需证明存在一个矩阵 $C$,使得
        \[
        (A^{-1}+B^{-1})C=I
        \]
        设
        \[
        C=B(A+B)^{-1}A
        \]
        下面验证该选择可行,观察
        \begin{align*}
        (A^{-1}+B^{-1})C
        &=(A^{-1}+B^{-1})B(A+B)^{-1}A \\
        &=A^{-1}B(A+B)^{-1}A+(A+B)^{-1}A \\
        &=A^{-1}[(A+B)-A](A+B)^{-1}A+(A+B)^{-1}A \\
        &=[A^{-1}(A+B)-I](A+B)^{-1}A+(A+B)^{-1}A \\
        &=A^{-1}A-(A+B)^{-1}A+(A+B)^{-1}A \\
        &=I
        \end{align*}
        因此存在矩阵 $C=B(A+B)^{-1}A$ 使得 $(A^{-1}+B^{-1})C=I$,从而 $A^{-1}+B^{-1}$ 可逆,其逆矩阵为
        \[
        (A^{-1}+B^{-1})^{-1}=B(A+B)^{-1}A.
        \]
    \end{solution}

    \question 设 $A,B$ 均为 $n$ 阶方阵,且 $I+AB$ 可逆,证明矩阵$I+BA$也可逆。
    \begin{solution}
        首先发现到
        \[
        (I+BA)B=B + BAB=B(I + AB)
        \]
        由此观察
        \begin{align*}
        &(I+BA)\bigl[I - B(I+AB)^{-1}A\bigr] \\
        &= I + BA - (I+BA)B(I+AB)^{-1}A \\
        &= I + BA - B(I + AB)(I+AB)^{-1}A \\
        &= I + BA - BA \\
        &= I
        \end{align*}
        故得证$I+BA$也可逆。
    \end{solution}

    \question 设 $A,B,C$ 为同阶实方阵,且 $A$ 可逆,证明若
    \[
    (A-B)C = BA^{-1},
    \]
    则有
    \[
    C(A-B) = A^{-1}B.
    \]
    \begin{solution}
        由假设
        \[
        (A-B)C = BA^{-1}
        \]
        在两边加上 $(A-B)A^{-1}$,得到
        \[
        (A-B)C + (A-B)A^{-1} = BA^{-1} + (A-B)A^{-1} = I
        \]
        于是由
        \[
        (A-B)(C + A^{-1}) = I 
        \]
        知$A+B$可逆,有
        \[
        (C + A^{-1})(A-B) = I
        \]
        展开即得到
        \[
        C(A-B) = A^{-1}B
        \]
    \end{solution}

    \question 对任意整数 $n\ge 2$,设 $A,B$ 为两个 $n\times n$ 实矩阵,且满足
    \[
    A^{-1}+B^{-1}=(A+B)^{-1}.
    \]
    证明 
    \[
    \det A=\det B.
    \]
    若 $A,B$ 为复矩阵,结论是否仍成立？
    \begin{solution}
        由条件
        \[
        A^{-1}+B^{-1}=(A+B)^{-1}
        \]
        有
        \[
        I=(A+B)(A+B)^{-1} =(A+B)(A^{-1}+B^{-1}) =I+AB^{-1}+BA^{-1}+I
        \]
        于是
        \[
        AB^{-1}+BA^{-1}+I=0
        \]
        令 $X=AB^{-1}$,则 $A=XB$,且$BA^{-1}=X^{-1}$,因此
        \[
        X+X^{-1}+I=0 
        \]
        于是
        \[
        X^3-I=(X-I)X(X+X^{-1}+I)=0 \Rightarrow X^3=I
        \]
        取行列式可得
        \[
        (\det X)^3=\det(X^3)=\det I=1
        \]
        由于 $X$ 为实矩阵,故$\det X=1$,于是
        \[
        \det A=\det(XB)=\det X\det B=\det B,
        \]
        故得证。下面说明在复矩阵情形下结论不成立。设
        \[
        \omega=-\frac{1}{2}+\frac{\sqrt{3}}{2}i \notin\mathbb{R}
        \]
        其中 $\omega^3=1$,且有
        \[
        1+\omega+\omega^2=0.
        \]
        取 $A=I$,且$B$ 为对角矩阵,其对角元均取为 $\omega$ 或 $\omega^2$使得 $\det B\neq1$(若 $n$ 不是 $3$ 的倍数可直接取 $B=\omega I$),此时
        \[
        A^{-1}=I,\quad B^{-1}=\overline{B}
        \]
        且
        \[
        I+B+\overline{B}=0
        \]
        因此
        \[
        (A+B)^{-1}=(-\overline{B})^{-1}=-B=I+\overline{B}
        =A^{-1}+B^{-1}
        \]
        但由构造可知
        \[
        \det A=1\ne\det B
        \]
    \end{solution}

    \question 设 $A,B$ 为 $n\times n$ 实矩阵,满足
    \[
    A^2+B^2=AB.
    \]
    证明若 $BA-AB$ 为可逆矩阵,则 $n$ 能被 $3$ 整除。
    \begin{solution}
        设$S=A+\omega B$,其中$\omega=-\dfrac12+\dfrac{\sqrt3}{2}i$,则
        \[
        S\bar S=(A+\omega B)(A+\bar\omega B) =A^2+\omega BA+\bar\omega AB+B^2=AB+\omega BA+\bar\omega AB
        \]
        注意到 $\bar\omega+1=-\omega$,于是
        \[
        S\bar S=\omega(BA-AB)
        \]
        两边取行列式,有
        \[
        \det S\cdot\det\bar S=\det(\omega(BA-AB))=\omega^n\det(BA-AB)。
        \]
        由于 $\det S\cdot\det\bar S \in \mathbb{R}$,且$BA-AB$ 可逆,故 $\det(BA-AB)\ne0$,从而 $\omega^n$ 必须为实数,因此 $n$ 必须是 $3$ 的倍数。
    \end{solution}

    \question 求所有复数 $\lambda$,使得存在正整数 $n$ 及实 $n \times n$ 矩阵 $A$,满足 
    \[
    A^2 = A^T
    \]
    且 $\lambda$ 是 $A$ 的一个特征值。
    \begin{solution}
        有
        \[ 
        A^4 = (A^2)^2 = (A^T)^2 = (A^2)^T = (A^T)^T = A
        \]
        所以
        \[ 
        A^4 - A = 0
        \]
        由此可知$A$ 的所有特征值都是多项式 $X^4 - X$ 的根,而$X^4 - X = X(X^3 - 1)$ 的根是 
        \[
        0, 1,\frac{-1 \pm \sqrt{3}i}{2}
        \]
        考虑以下矩阵
        \[ 
        A_0 = (0) , \quad A_1 = (1) , \quad A_2 = \begin{pmatrix} -\frac{1}{2} & \frac{\sqrt{3}}{2} \\ -\frac{\sqrt{3}}{2} & -\frac{1}{2} \end{pmatrix} , \quad A_4 = \begin{pmatrix} 0 & 0 & 0 & 0 \\ 0 & 1 & 0 & 0 \\ 0 & 0 & -\frac{1}{2} & \frac{\sqrt{3}}{2} \\ 0 & 0 & -\frac{\sqrt{3}}{2} & -\frac{1}{2} \end{pmatrix} 
        \]
        $0$ 和 $1$ 分别是 $1 \times 1$ 矩阵 $A_0$ 和 $A_1$ 的特征值, $\frac{-1 \pm \sqrt{3}i}{2}$ 是 $A_2$ 的特征值,不难看出
        \[ 
        A_2^2 = \begin{pmatrix} -\frac{1}{2} & -\frac{\sqrt{3}}{2} \\ \frac{\sqrt{3}}{2} & -\frac{1}{2} \end{pmatrix} = A_2^T
        \]
        故矩阵 $A_4$ 包含了所有四个可能特征值。
    \end{solution}

    \question 已知 $A$ 是 $4 \times 2$ 实矩阵,$\ B$ 是 $2 \times 4$ 实矩阵,且
    \[
    AB = \begin{pmatrix}
    1 & 0 & -1 & 0 \\
    0 & 1 & 0 & -1 \\
    -1 & 0 & 1 & 0 \\
    0 & -1 & 0 & 1
    \end{pmatrix}.
    \]
    求 $BA$。
    \begin{solution}
        将矩阵分块成
        \[
        A = \begin{pmatrix} A_1 \\ A_2 \end{pmatrix}, \quad B = \begin{pmatrix} B_1 & B_2 \end{pmatrix},
        \]
        其中 $A_1, A_2, B_1, B_2$ 为 $2 \times 2$ 矩阵,于是
        \[
        AB = \begin{pmatrix} A_1 \\ A_2 \end{pmatrix} \begin{pmatrix} B_1 & B_2 \end{pmatrix} = 
        \begin{pmatrix} A_1B_1 & A_1B_2 \\ A_2B_1 & A_2B_2 \end{pmatrix}
        \]
        比较得
        \[
        A_1B_1 = A_2B_2 = I_2, \quad A_1B_2 = A_2B_1 = -I_2
        \]
        即
        \[
        B_1 = A_1^{-1}, \quad B_2 = -A_1^{-1}, \quad A_2 = -A_1
        \]
        于是
        \[
        BA = \begin{pmatrix} B_1 & B_2 \end{pmatrix} \begin{pmatrix} A_1 \\ A_2 \end{pmatrix} 
        = B_1A_1 + B_2A_2 = A_1^{-1}A_1 + (-A_1^{-1})(-A_1) = 2 I_2
        \]
    \end{solution}

    \question 已知 
    \[
    A=\begin{bmatrix} 
    2 & 4 & 0 & 0 \\ 
    1 & 2 & 0 & 0 \\ 
    0 & 0 & 2 & 4 \\ 
    0 & 0 & 2 & 2 
    \end{bmatrix}, 
    \] 
    求 $A^n$。
    \begin{solution}
        将 $A$ 分块为 
        \[
        A=\begin{bmatrix} B & 0 \\ 0 & C \end{bmatrix}, \quad 
        B=\begin{bmatrix} 2 & 4 \\ 1 & 2 \end{bmatrix}, \quad 
        C=\begin{bmatrix} 2 & 4 \\ 0 & 2 \end{bmatrix}
        \]
        注意到 $B = \begin{bmatrix} 2 \\ 1 \end{bmatrix} \begin{bmatrix} 1 & 2 \end{bmatrix}$,所以可递推得
        \[
        B^n = 4^{\,n-1} B
        \]
        且由于 $\begin{bmatrix} 0 & 4 \\ 0 & 0 \end{bmatrix}^2 = 0$,则
        \[
        C^n = \bigl(2I + \begin{bmatrix} 0 & 4 \\ 0 & 0 \end{bmatrix}\bigr)^n = 2^n I + n2^{\,n-1} \begin{bmatrix} 0 & 4 \\ 0 & 0 \end{bmatrix} = \begin{bmatrix} 2^n & 4n2^{\,n-1} \\ 0 & 2^n \end{bmatrix}
        \]
        因此
        \[
        A^n = \begin{bmatrix} B^n & 0 \\ 0 & C^n \end{bmatrix} 
        = \begin{bmatrix}
        2 \cdot 4^{\,n-1} & 4 \cdot 4^{\,n-1} & 0 & 0 \\
        1 \cdot 4^{\,n-1} & 2 \cdot 4^{\,n-1} & 0 & 0 \\
        0 & 0 & 2^n & n 2^{\,n+1} \\
        0 & 0 & 0 & 2^n
        \end{bmatrix} 
        = \begin{bmatrix}
        2^{\,2n-1} & 2^{\,2n+1} & 0 & 0 \\
        2^{\,2n-2} & 2^{\,2n} & 0 & 0 \\
        0 & 0 & 2^n & n 2^{\,n+1} \\
        0 & 0 & 0 & 2^n
        \end{bmatrix}
        \]
    \end{solution}

    \question 设 $A$ 和 $B$ 是 $n \times n$ 实矩阵,且满足
    \[
    \mathrm{rank}(AB - BA + I) = 1,
    \]
    其中 $I$ 是 $n\times n$ 单位矩阵。证明
    \[
    \mathrm{tr}(ABAB) - \mathrm{tr}(A^2B^2) = \frac{1}{2}n(n-1).
    \]
    \begin{solution}
        令 $X = AB - BA$,注意到
        \begin{align*}
        \mathrm{tr}(X^2) &= \mathrm{tr}(ABAB - ABBA - BAAB + BABA) \\
        &= 2\mathrm{tr}(ABAB) - 2\mathrm{tr}(A^2B^2),
        \end{align*}
        因为迹具有循环性。因此只需证明 $\mathrm{tr}(X^2) = n(n-1)$。由假设,$X+I$ 的秩为 $1$,所以
        \[
        X+I = v^t w
        \]
        其中$v,w$为向量,于是
        \[
        X^2 = (v^t w - I)^2 = I - 2 v^t w + v^t w v^t w = I + (w v^t - 2)v^t w.
        \]
        且有 $\mathrm{tr}(X) = 0$,因此 
        \[
        w v^t=\mathrm{tr}(w v^t) = \mathrm{tr}(v^t w) = n
        \]
        于是
        \[
        \mathrm{tr}(X^2) = n + (n-2)n = n(n-1)
        \]
        即得证。
    \end{solution}
    \begin{solution}
        由于 $X+I$ 的秩为 $1$,它有特征值$0$,重数为 $n-1$。因此 $X$ 有特征值 $-1$,重数也为 $n-1$。由于 $\mathrm{tr}(X) = 0$,剩余的特征值只能为为 $n-1$,于是
        \[
        \mathrm{tr}(X^2) = (n-1)^2 + (n-1)\cdot 1^2 = n(n-1)
        \]
        同样得证。
    \end{solution}

    \question 确定是否存在奇正整数 $n$ 和具有整数元素的 $n \times n$ 矩阵 $A,B$,满足以下条件:
    \begin{itemize}
        \item $\det(B) = 1$;
        \item $AB = BA$;
        \item $A^4 + 4A^2B^2 + 16B^4 = 2019I$
    \end{itemize}
    \begin{solution}
        我们证明不存在这样的矩阵。注意到 $A^4 + 4A^2B^2 + 16B^4$ 可以因式分解为
        \[
        A^4 + 4A^2B^2 + 16B^4 = (A^2 + 2AB + 4B^2)(A^2 - 2AB + 4B^2)
        \]
        令 $C = A^2 + 2AB + 4B^2,D = A^2 - 2AB + 4B^2$ ,则
        \[
        \det C \cdot \det D = \det(CD) = \det(2019I) = 2019^n
        \]
        矩阵 $C, D$ 的元素皆为整数,因此它们的行列式也为整数,此外由 $C \equiv D \pmod 4$ 可知
        \[
        \det C \equiv \det D \pmod 4
        \]
        这意味 $\det C \cdot \det D \equiv (\det C)^2 \pmod 4$,但 $2019^n \equiv 3$ 是模 $4$ 的二次非剩余,矛盾。
    \end{solution}
    \begin{solution}
        注意
        \[ 
        A^4 \equiv A^4 + 4A^2B^2 + 16B^4 = 2019I \pmod{4} 
        \]
        所以
        \[ 
        (\det A)^4 = \det A^4 \equiv \det(2019I) = 2019^n \pmod{4}. 
        \]
        但是 $2019^n \equiv 3$ 是模 $4$ 的二次非剩余,矛盾。
    \end{solution}

    \question 设 $A$ 是 $n \times n$ 实矩阵,且 
    \[
    A^3 = 0
    \]
    证明存在唯一的 $n \times n$ 实矩阵 $X$ 满足
    \[
    X + AX + XA^2 = A,
    \]
    且以 $A$ 表示 $X$。
    \begin{solution}
        首先假设某个矩阵 $X$ 满足该方程。如果我们在给定方程的左侧乘以 $A$ 的某个幂,右侧乘以 $A$ 的另一个幂,就可以得到新的方程。例如,
        \[ A^2(X + AX + XA^2)A^2 = A^2XA^2 + A^3 \cdot XA^2 + A^2XA \cdot A^3 = A^2XA^2. \]
        所以
        \begin{align*}
        A^2XA^2 &= A^2(X + AX + XA^2)A^2 = A^5 = 0 \\
        A^2X &= A^2(X + AX + XA^2) = A^3 = 0 \\
        AXA &= A(X + AX + XA^2)A = A^3 = 0 \\
        XA^2 &= (X + AX + XA^2)A^2 = A^3 = 0 \\
        AX &= A(X + AX + XA^2)A = A^2 \\
        X &= A - AX - XA^2 = A - A^2
        \end{align*}
        因此,除 $A - A^2$ 以外,没有其他矩阵能满足该方程。注意上述论证并未证明矩阵 $X = A - A^2$ 满足方程,因为这些步骤不能逆推。事实上,
        \[ 
        X + AX + XA^2 = (A - A^2) + A(A - A^2) + (A - A^2)A^2 = A - A^4 = A
        \]
        因此,$X = A - A^2$ 是方程的唯一解。
    \end{solution}
    \begin{solution}
        我们使用另一种不同的方法来用 $A$ 表示 $X$。如果存在某个矩阵 $X$ 满足方程,那么
        \[ 
        X = A - AX - XA^2. 
        \]
        让我们在右侧反复代入此恒等式,直到 $X$ 在各处都被消去。注意到根据条件 $A^3 = 0$,我们有 $A^3 = A^4 = A^5 = A^3X = XA^4 = AXA^4 = A^3XA^2 = 0$,所以
        \begin{align*}
        X &= A - AX - XA^2 \\
        &= A - A(A - AX - XA^2) - (A - AX - XA^2)A^2 \\
        &= A - (A^2 - A^2X - AXA^2) - (A^3 - AXA^2 - XA^4) \\
        &= A - A^2 + A^2X + 2AXA^2 \\
        &= A - A^2 + A^2(A - AX - XA^2) + 2A(A - AX - XA^2)A^2 \\
        &= A - A^2 + (A^3 - A^3X - A^2XA^2) + 2(A^4 - A^2XA^2 - AXA^4) \\
        &= A - A^2 - 3A^2XA^2 \\
        &= A - A^2 - 3A^2(A - AX - XA^2)A^2 \\
        &= A - A^2 - 3(A^5 - A^3XA^2 - A^2XA^4) \\
        &= A - A^2.
        \end{align*}
        同上,仍需验证 $X = A - A^2$ 确实是一个解,这一步与解 1 相同。
    \end{solution}
    \begin{solution}
        设 $B = I - A + A^2$,使得 $B$ 是 $I + A$ 的逆矩阵。在左侧乘以 $B$,方程等价于
        \[ 
        X + BXA^2 = BA. \tag{1} 
        \]
        现在假设 $X$ 满足方程。在右侧乘以 $A^2$ 并利用 $A^3 = 0$,我们得到 $XA^2 = 0$。因此方程简化为 
        \[
        X = BA = A - A^2
        \]
        另一方面,$X = BA$ 显然满足 (1)。
    \end{solution}

    \question 设 $A,B,C$ 为 $n \times n$ 复矩阵,满足
    \[
    A^2 = B^2 = C^2, \quad B^3 = ABC + 2I.
    \]
    证明 $A^6 = I$。
    \begin{solution}
        发现到$B^3=A^2B$可得
        \[
        B^3 - ABC = A^2B - ABC = 2I \Rightarrow A(AB - BC) = 2I
        \]
        同理由$B^3=BC^2$可得
        \[
        BC^2 - ABC = 2I \implies (BC - AB)C = 2I
        \]
        这说明 $A$ 是 $\dfrac{AB - BC}{2}$ 的左逆矩阵,而 $-C$ 是右逆矩阵,所以 $A = -C$。因此
        \[
        ABA = 2I- B^3
        \]
        于是 $ABA$ 与 $B$ 交换,由此得 
        \[
        (AB)^2 = (BA)^2
        \]
        而
        \begin{align*}
        (AB - BA)(AB + BA) &= (AB)^2 + AB^2A - BA^2B - (BA)^2 \\
        &= (AB)^2 + A^4 - B^4 - (AB)^2 = 0
        \end{align*}
        又因为 $AB - BC = AB + BA$ 可逆,故 $AB = BA$,于是 $ABA = A^2B = B^3$,由此得
        \[
        B^3 = I \Rightarrow A^6 = B^6 = (B^3)^2 = I^2 = I
        \]
    \end{solution}

    \question 设$a=\begin{pmatrix} 1 \\ 2 \\ 1 \end{pmatrix},b=\begin{pmatrix} 1 \\ \frac12 \\ 0 \end{pmatrix}, c=\begin{pmatrix} 0 \\ 0 \\ 8 \end{pmatrix}$, 且$A=ab^T$, $B=ba^T$, 解矩阵方程$$2B^2 A^2 X = A^4 X + B^4 X + c$$\textcolor{red}{(题目有误,不符合矩阵乘法定义)}

\end{questions}
\pagebreak

\begin{center}
  {\fontsize{30pt}{26pt}\selectfont
    \hypertarget{行列式}{行列式} \label{行列式}
  }
\end{center}
\separator
\vspace{1pt}

\begin{questions}
    \question 若 $a+b+c=x+y+z=0$,求行列式 
    \[
    \begin{vmatrix}
    xa & yb & zc \\
    yc & za & xb \\
    zb & xc & ya
    \end{vmatrix}
    \]
    \begin{solution}
        展开行列式得
        \[
        \begin{vmatrix}
        xa & yb & zc \\
        yc & za & xb \\
        zb & xc & ya
        \end{vmatrix}=xyz(a^3+b^3+c^3)+abc(x^3+y^3+z^3)
        \]由于
        \[
        a^3+b^3+c^3-3abc=(a+b+c)(a^2+b^2+c^2-ab-bc-ca)=0
        \] 
        \[
        x^3+y^3+z^3-3xyz=(x+y+z)(x^2+y^2+z^2-xy-yz-zx)=0
        \]
        故
        \[
        \begin{vmatrix}
        xa & yb & zc \\
        yc & za & xb \\
        zb & xc & ya
        \end{vmatrix}=xyz(3abc)-abc(3xyz)=0\]
    \end{solution}

    \question 计算行列式
    \[
    \begin{vmatrix} 
    \log a & \log b & \log \dfrac{1}{ab} \\[8pt]
    \log b & \log c & \log \dfrac{1}{bc} \\[8pt]
    \log c & \log a & \log \dfrac{1}{ca} 
    \end{vmatrix}
    \]
    \begin{solution}
        令\[
        D=\begin{vmatrix} 
        \log a & \log b & \log \dfrac{1}{ab} \\[8pt]
        \log b & \log c & \log \dfrac{1}{bc} \\[8pt]
        \log c & \log a & \log \dfrac{1}{ca} 
        \end{vmatrix}
        \]
        化为
        \[
        D=\begin{vmatrix} 
        \log a & \log b & -\log a - \log b \\ 
        \log b & \log c & -\log b - \log c \\ 
        \log c & \log a & -\log c - \log a 
        \end{vmatrix}
        \]
        进行列变换 $C_3 \leftarrow C_3 - C_1 - C_2$,得到
        \[
        D=\begin{vmatrix} 
        \log a & \log b & 0 \\ 
        \log b & \log c & 0 \\ 
        \log c & \log a & 0 
        \end{vmatrix} = 0
        \]
    \end{solution}

    \question  若 $m$为正整数,求行列式 \[
        \begin{vmatrix}
        1 & 1 & 1 \\
        \comb{m}{1} & \comb{m+1}{1} & \comb{m+2}{1} \\
        \comb{m}{2} & \comb{m+1}{2} & \comb{m+2}{2}
        \end{vmatrix}
        \]
    \begin{solution}
        令\[
        D=\begin{vmatrix}
        1 & 1 & 1 \\
        \comb{m}{1} & \comb{m+1}{1} & \comb{m+2}{1} \\
        \comb{m}{2} & \comb{m+1}{2} & \comb{m+2}{2}
        \end{vmatrix}
        \]
        由性质 $\comb{n-1}{k-1}+\comb{n-1}{k}=\comb{n}{k}$,有\[
        D=\begin{vmatrix}
        1 & 1 & 1 \\
        \comb{m}{1} & \comb{m+1}{1} & \comb{m+1}{0}+\comb{m+1}{1} \\
        \comb{m}{2} & \comb{m+1}{2} & \comb{m+1}{1}+\comb{m+1}{2}
        \end{vmatrix} 
        \]
        进行列变换$C_3\rightarrow C_3-C_2,$\[
        D=\begin{vmatrix}
        1 & 1 & 0 \\
        \comb{m}{1} & \comb{m+1}{1} & \comb{m+1}{0} \\
        \comb{m}{2} & \comb{m+1}{2} & \comb{m+1}{1}
        \end{vmatrix} 
        \]
        再由性质 $\comb{n-1}{k-1}+\comb{n-1}{k}=\comb{n}{k}$,\[
        D=\begin{vmatrix}
        1 & 1 & 0 \\
        \comb{m}{1} & \comb{m}{0}+\comb{m}{1} & \comb{m+1}{0} \\
        \comb{m}{2} & \comb{m}{1}+\comb{m}{2} & \comb{m+1}{1}
        \end{vmatrix} 
        \]
        进行列变换$C_2\rightarrow C_2-C_1,$\[
        D=\begin{vmatrix}
        1 & 0 & 0 \\
        \comb{m}{1} & \comb{m}{0} & \comb{m+1}{0} \\
        \comb{m}{2} & \comb{m}{1} & \comb{m+1}{1}
        \end{vmatrix} 
        \]
        按第一行展开得\[
        D= \comb{m}{0}\comb{m+1}{1}-\comb{m}{1}\comb{m+1}{0}=1\cdot(m+1)-m\cdot1=1
        \]
    \end{solution}

    \question 已知$A,B,C$为 $\triangle ABC$的内角,求  
    \[
    \begin{vmatrix} 
    \tan A & 1 & 1 \\ 
    1 & \tan B & 1 \\ 
    1 & 1 & \tan C 
    \end{vmatrix}
    \]
    \begin{solution}
        令\[
        D=\begin{vmatrix} 
        \tan A & 1 & 1 \\ 
        1 & \tan B & 1 \\ 
        1 & 1 & \tan C 
        \end{vmatrix}
        \]
        展开行列式:
        \[
        D = \tan A \tan B \tan C - (\tan A + \tan B + \tan C) + 2
        \]
        由于 $A + B + C = \pi$,有
        \[
        \tan(A + B) = \tan(\pi - C) = -\tan C
        \]
        又
        \[
        \tan(A + B) = \frac{\tan A + \tan B}{1 - \tan A \tan B}
        \]
        因此可得
        \[
        \frac{\tan A + \tan B}{1 - \tan A \tan B} \Rightarrow \tan A + \tan B + \tan C = \tan A \tan B \tan C
        \]
        故得到
        \[
        D = 2
        \]
    \end{solution}

    \question 计算行列式
    \[
    \begin{vmatrix}
    \tan 40^{\circ} & \tan 10^{\circ} & \tan 50^{\circ} \\
    \tan 20^{\circ} & \tan 50^{\circ} & \tan 70^{\circ} \\
    \tan 10^{\circ} & \tan 70^{\circ} & \tan 80^{\circ}
    \end{vmatrix}
    \]
    \begin{solution}
        令
        \[
        D=
        \begin{vmatrix}
        \tan 40^\circ & \tan 10^\circ & \tan 50^\circ \\
        \tan 20^\circ & \tan 50^\circ & \tan 70^\circ \\
        \tan 10^\circ & \tan 70^\circ & \tan 80^\circ
        \end{vmatrix}
        \]
        进行列变换 $C_3 \to -C_1 - C_2 + C_3 $,
        \[
        D=
        \begin{vmatrix}
        \tan 40^\circ & \tan 10^\circ & -\tan 10^\circ - \tan 40^\circ + \tan 50^\circ \\
        \tan 20^\circ & \tan 50^\circ & -\tan 20^\circ - \tan 50^\circ + \tan 70^\circ \\
        \tan 10^\circ & \tan 70^\circ & -\tan 10^\circ - \tan 70^\circ + \tan 80^\circ
        \end{vmatrix}
        \]
        由
        \[
        \tan(A+B) = \frac{\tan A + \tan B}{1 - \tan A \tan B} 
        \]
        可得
        \[
        \tan A \tan B \tan (A+B) = \tan(A+B) - \tan A - \tan B
        \]
        且若 $A + B = 90^\circ$,则 
        \[
        \tan A \tan B = 1
        \]
        因此
        \[
        D=\begin{vmatrix}
        \tan 40^\circ & \tan 10^\circ & \tan 10^\circ \cdot \tan 40^\circ \cdot \tan 50^\circ \\
        \tan 20^\circ & \tan 50^\circ & \tan 20^\circ \cdot \tan 50^\circ \cdot \tan 70^\circ \\
        \tan 10^\circ & \tan 70^\circ & \tan 10^\circ \cdot \tan 70^\circ \cdot \tan 80^\circ
        \end{vmatrix}
        =
        \begin{vmatrix}
        \tan 40^\circ & \tan 10^\circ & \tan 10^\circ \\
        \tan 20^\circ & \tan 50^\circ & \tan 50^\circ \\
        \tan 10^\circ & \tan 70^\circ & \tan 70^\circ
        \end{vmatrix}
        =0
        \]
    \end{solution}

    \question 已知$A,B,C$为 $\triangle ABC$的内角,计算行列式
    \[
    \begin{vmatrix} 
    a & a^2 & \cos A \\ 
    b & b^2 & \cos B \\ 
    c & c^2 & \cos C 
    \end{vmatrix}
    \]
    \begin{solution}
        令
        \[
        D=\begin{vmatrix} 
        a & a^2 & \cos A \\ 
        b & b^2 & \cos B \\ 
        c & c^2 & \cos C 
        \end{vmatrix}
        \]
        由余弦定理,
        \[
        D= \frac{1}{2abc} 
        \begin{vmatrix} 
        a & a^2 & a(b^2+c^2-a^2) \\ 
        b & b^2 & b(a^2+c^2-b^2) \\ 
        c & c^2 & c(a^2+b^2-c^2) 
        \end{vmatrix}
        = \frac{1}{2} 
        \begin{vmatrix} 
        1 & a & b^2+c^2-a^2 \\ 
        1 & b & a^2+c^2-b^2 \\ 
        1 & c & a^2+b^2-c^2 
        \end{vmatrix}
        \]
        进行行变换$R_2 \to R_2-R_1, R_3 \to R_3-R_1$
        \[
        D=\frac{1}{2} 
        \begin{vmatrix} 
        1 & a & b^2+c^2-a^2 \\ 
        0 & b-a & 2a^2-2b^2 \\ 
        0 & c-a & 2a^2-2c^2 
        \end{vmatrix}
        \]
        按$C_1$展开,
        \[
        D= \frac{1}{2} 
        \begin{vmatrix} 
        b-a & 2a^2-2b^2 \\ 
        c-a & 2a^2-2c^2 
        \end{vmatrix}
        = -(b-a)(c-a) 
        \begin{vmatrix} 
        1 & a+b \\ 
        1 & a+c 
        \end{vmatrix}
        = (a-b)(a-c)(b-c)
        \]
    \end{solution}

    \question 在坐标平面上,设 $\triangle ABC$ 经二阶方阵 $M = \begin{bmatrix} a & b \\ c & d \end{bmatrix}$ 作线性变换后成 $\triangle A'B'C'$。若 $\triangle ABC$ 的面积为 $\triangle,\triangle A'B'C'$ 的面积为 $\triangle'$,试证$$\triangle' =  \begin{vmatrix} a & b \\ c & d \end{vmatrix} \cdot \triangle$$
    \begin{solution}
        设$A(x_a, y_a),B(x_b, y_b),C(x_c, y_c)$
        经变换后为
        \[
        A'(a x_a + b y_a,\; c x_a + d y_a),B'(a x_b + b y_b,\; c x_b + d y_b),C'(a x_c + b y_c,\; c x_c + d y_c)
        \]
        则
        \[
        \Delta' = \frac{1}{2}
        \begin{vmatrix}
        a x_a + b y_a & c x_a + d y_a & 1 \\
        a x_b + b y_b & c x_b + d y_b & 1 \\
        a x_c + b y_c & c x_c + d y_c & 1
        \end{vmatrix}
        \]
        变为
        \[
        \Delta' = \frac{1}{2} \begin{vmatrix}
        a x_a & d y_a & 1 \\
        a x_b & d y_b & 1 \\
        a x_c & d y_c & 1
        \end{vmatrix}
        +
        \frac{1}{2} \begin{vmatrix}
        b y_a & c x_a & 1 \\
        b y_b & c x_b & 1 \\
        b y_c & c x_c & 1
        \end{vmatrix} +0+0
        = \frac{1}{2}(ad-bc)
        \begin{vmatrix}
        x_a & y_a & 1 \\
        x_b & y_b & 1 \\
        x_c & y_c & 1
        \end{vmatrix}
        = \begin{vmatrix} a & b \\ c & d \end{vmatrix} \cdot \triangle
        \]
    \end{solution}

    \question 求下列行列式的立方根:
    \[ 
    \begin{vmatrix}
    -\dfrac{2019^2}{\sqrt[5]{80}} & \dfrac{2020^2}{\sqrt[5]{80^2}} & -\dfrac{2021^2}{\sqrt[5]{80^3}} \\[8pt]
    \dfrac{2022^2}{\sqrt[5]{80^4}} & -\dfrac{2023^2}{\sqrt[5]{80^5}} & \dfrac{2024^2}{\sqrt[5]{80^6}} \\[8pt]
    -\dfrac{2025^2}{\sqrt[5]{80^7}} & \dfrac{2026^2}{\sqrt[5]{80^8}} & -\dfrac{2027^2}{\sqrt[5]{80^9}}
    \end{vmatrix}
    \]
    \begin{solution}
        设$a=2019,r=-\dfrac{1}{\sqrt[5]{80}}$,原行列式为
        \[
        D=
        \begin{vmatrix}
        a^2 r & (a+1)^2 r^2 & (a+2)^2 r^3 \\
        (a+3)^2 r^4 & (a+4)^2 r^5 & (a+5)^2 r^6 \\
        (a+6)^2 r^7 & (a+7)^2 r^8 & (a+8)^2 r^9
        \end{vmatrix}
        \]
        进行列变换$C_2 \rightarrow C_2-C_1,C_3 \rightarrow C_3-C_2$,
        \[
        D=\begin{vmatrix}
        a^2 r & (2a+1) r^2 & (2a+3) r^3 \\
        (a+3)^2 r^4 & (2a+7) r^5 & (2a+9) r^6 \\
        (a+6)^2 r^7 & (2a+13) r^8 & (2a+15) r^9
        \end{vmatrix}
        \]
        进行列变换$C_3 \rightarrow C_3-C_2$,
        \[
        D=\begin{vmatrix}
        a^2 r & (2a+1) r^2 & 2r^3 \\
        (a+3)^2 r^4 & (2a+7) r^5 & 2r^6 \\
        (a+6)^2 r^7 & (2a+13) r^8 & 2r^9
        \end{vmatrix}
        \]
        进行行变换$R_3 \rightarrow R_3-r^3R_2,R_2 \rightarrow R_2-r^3R_1$
        \[
        D=\begin{vmatrix}
        a^2 r & (2a+1) r^2 & 2r^3 \\
        (6a+9) r^4 & 6 r^5 & 0 \\
        (6a+27) r^7 & 6 r^8 & 0
        \end{vmatrix}
        \]
        进行行变换$R_3 \rightarrow R_3-r^3R_2$,
        \[
        D=\begin{vmatrix}
        a^2 r & (2a+1) r^2 & 2r^3 \\
        (6a+9) r^4 & 6 r^5 & 0 \\
        18 r^7 & 0 & 0
        \end{vmatrix}
        \]
        进行列变换$C_1 \leftrightarrow C_3$,
        \[
        D=-\begin{vmatrix}
        2r^3 & (2a+1) r^2 & a^2 r \\
        0 & 6 r^5 & (6a+9) r^4 \\
        0 & 0 & 18 r^7
        \end{vmatrix}
        =-(2r^3)(6r^5)(18r^7)=-216r^{15}=\frac{216}{80^3}
        \]
        故原行列式的立方根为$\dfrac{3}{40}$
    \end{solution}

    \question 若 $a^2+b^2+c^2=47$, 求
    \[
    \begin{vmatrix}
    a^2-1 & ab & ca \\
    ab & b^2-1 & bc \\
    ca & bc & c^2-1
    \end{vmatrix}
    \]
    \begin{solution}
        令
        \[
        D=\begin{vmatrix}
        a^2-1 & ab & ca \\
        ab & b^2-1 & bc \\
        ca & bc & c^2-1
        \end{vmatrix}
        \]
        进行行变换$R_1\to aR_1, R_2\to bR_2, R_3\to cR_3$,
        \[
        D=\frac{1}{abc}\begin{vmatrix}
        a(a^2-1) & a^2b & ca^2 \\
        ab^2 & b(b^2-1) & b^2c \\
        c^2a & bc^2 & c(c^2-1)
        \end{vmatrix}
        \]
        进行列变换$C_1\to \dfrac{C_1}{a}, C_2\to \dfrac{C_2}{b}, C_3\to \dfrac{C_3}{c}$,
        \[
        D=\begin{vmatrix}
        a^2-1 & a^2 & a^2 \\
        b^2 & b^2-1 & b^2 \\
        c^2 & c^2 & c^2-1
        \end{vmatrix}
        \]
        进行行变换 $R_1\to R_1+R_2+R_3$,
        \begin{align*}
            D&=\begin{vmatrix}
            a^2+b^2+c^2-1 & a^2+b^2+c^2-1 & a^2+b^2+c^2-1 \\
            b^2 & b^2-1 & b^2 \\
            c^2 & c^2 & c^2-1
            \end{vmatrix}\\
            &=(a^2+b^2+c^2-1)
            \begin{vmatrix}
            1 & 1 & 1 \\
            b^2 & b^2-1 & b^2 \\
            c^2 & c^2 & c^2-1
            \end{vmatrix}
        \end{align*}
        进行列变换 $C_2\to -C_1+C_2,C_3\to -C_1+C_3$,
        \[
        D=(a^2+b^2+c^2-1) 
        \begin{vmatrix}
        1 & 0 & 0 \\
        b^2 & -1 & 0 \\
        c^2 & 0 & -1
        \end{vmatrix}
        =(47-1)\cdot1=46
        \]
    \end{solution}

    \question 证明
    \[
    \begin{vmatrix} 
    b+c & c+a & a+b \\ 
    q+r & r+p & p+q \\ 
    y+z & z+x & x+y 
    \end{vmatrix}
    = 2\begin{vmatrix} 
    a & b & c \\ 
    p & q & r \\ 
    x & y & z 
    \end{vmatrix}
    \]
    \begin{solution}
        令\[
        D=\begin{vmatrix} 
        b+c & c+a & a+b \\ 
        q+r & r+p & p+q \\ 
        y+z & z+x & x+y 
        \end{vmatrix}
        \]
        进行列变换$C_3\to C_3-C_2$,
        \[
        D = \begin{vmatrix} 
        b+c & c+a & b-c \\ 
        q+r & r+p & q-r \\ 
        y+z & z+x & y-z 
        \end{vmatrix}
        \]
        进行列变换$C_1\to C_1-C_3$,
        \[
        D = 2 \begin{vmatrix} 
        b & c+a & b-c \\ 
        q & r+p & q-r \\ 
        y & z+x & y-z 
        \end{vmatrix}
        \]
        进行列变换$C_3\to C_3-C_1$,
        \[
        D = 2 \begin{vmatrix} 
        b & c+a & -c \\ 
        q & r+p & -r \\ 
        y & z+x & -z 
        \end{vmatrix}
        \]
        进行列变换$C_2\to C_2+C_3$,
        \[
        D = 2 \begin{vmatrix} 
        b & a & -c \\ 
        q & p & -r \\ 
        y & x & -z 
        \end{vmatrix}
        =2 \begin{vmatrix} 
        a & b & c \\ 
        p & q & r \\ 
        x & y & z 
        \end{vmatrix}
        \]
        故得证。
    \end{solution}

    \question 证明
    \[
    \begin{vmatrix} 
    b+c & a & a \\ 
    b & c+a & b \\ 
    c & c & a+b 
    \end{vmatrix} = 4abc
    \]
    \begin{solution}
        令
        \[
        D=\begin{vmatrix} 
        b+c & a & a \\ 
        b & c+a & b \\ 
        c & c & a+b 
        \end{vmatrix}
        \]
        进行行变换$R_1 \to R_1 - R_2 - R_3$,
        \[
        D=\begin{vmatrix} 
        0 & -2c & -2b \\ 
        b & c+a & b \\ 
        c & c & a+b 
        \end{vmatrix}
        =-2 \begin{vmatrix} 
        0 & c & b \\ 
        b & c+a & b \\ 
        c & c & a+b 
        \end{vmatrix}
        \]
        进行行变换$R_2 \to R_2 - R_1,R_3 \to R_3 - R_1$,
        \[
        D = -2 \begin{vmatrix} 
        0 & c & b \\ 
        b & a & 0 \\ 
        c & 0 & a 
        \end{vmatrix} = -2 (-2abc) = 4abc
        \]
    \end{solution}

    \question 证明$$ \begin{vmatrix}
        a^2 & b^2 + c^2 & bc \\
        b^2 & c^2 + a^2 & ca \\
        c^2 & a^2 + b^2 & ab
        \end{vmatrix}= -(a - b)(b - c)(c - a)(a + b + c)(a^2 + b^2 + c^2)$$
    \begin{solution}
        令
        \[
        D = \begin{vmatrix}
        a^2 & b^2 + c^2 & bc \\
        b^2 & c^2 + a^2 & ca \\
        c^2 & a^2 + b^2 & ab
        \end{vmatrix}
        \]
        进行行变换 $R_2\rightarrow R_2-R_1,\;R_3\rightarrow R_3-R_1,$
        \[D=
        \begin{vmatrix}
        a^2 & b^2 + c^2 & bc \\
        b^2 - a^2 & a^2 - b^2 & ca - bc \\
        c^2 - a^2 & a^2 - c^2 & ab - bc
        \end{vmatrix}
        \]
        化简得
        \begin{align*}
        D&=\begin{vmatrix}
        a^2 & b^2 + c^2 & bc \\
        (b+a)(b-a)& (a+b)(a-b) & c(a - b) \\
        (c+a)(c-a) & (a+c)(a-c) & b(a - c)
        \end{vmatrix} \\
        &= (a-b)(a-c)\begin{vmatrix}
        a^2 & b^2 + c^2 & bc \\
        -a-b& a+b & c \\
        -a-c & a+c & b
        \end{vmatrix} 
        \end{align*}
        进行列变换 $C_1 \rightarrow C_1+C_2,$
        \[
        D=(a-b)(a-c)\begin{vmatrix}
        a^2+b^2+c^2 & b^2 + c^2 & bc \\
        0& a+b & c \\
        0 & a+c & b
        \end{vmatrix}         
        \]
        按第一列展开得
        \begin{align*}
        D&=(a-b)(a-c)(a^2+b^2+c^2)(ab+b^2-ac-c^2)\\
        &=(a-b)(a-c)(a^2+b^2+c^2)(a(b-c)+(b+c)(b-c))\\
        &=-(a - b)(b - c)(c - a)(a + b + c)(a^2 + b^2 + c^2)
        \end{align*}
    \end{solution}
        
    \question 证明
    \[
    \begin{vmatrix}
    (b+c)^2 & a^2 & a^2 \\
    b^2 & (c+a)^2 & b^2 \\
    c^2 & c^2 & (a+b)^2 \\
    \end{vmatrix}
    = 2abc(a+b+c)^3
    \]
    \begin{solution}
        令
        \[
        D = \begin{vmatrix}
        (b+c)^2 & a^2 & a^2 \\
        b^2 & (c+a)^2 & b^2 \\
        c^2 & c^2 & (a+b)^2 \\
        \end{vmatrix}
        \]
        进行列变换 \( C_2 \rightarrow C_2 - C_1,\; C_3 \rightarrow C_3 - C_1 \),得
        \[
        D = \begin{vmatrix}
        (b+c)^2 & a^2 - (b+c)^2 & a^2 - (b+c)^2 \\
        b^2 & (c+a)^2 - b^2 & 0 \\
        c^2 & 0 & (a+b)^2 - c^2 \\
        \end{vmatrix}
        \]
        化简得
        \begin{align*}
        D &= \begin{vmatrix}
        (b+c)^2 & (a+b+c)(a-b-c) & (a+b+c)(a-b-c) \\
        b^2 & (a+b+c)(a-b+c) & 0 \\
        c^2 & 0 & (a+b+c)(a+b-c) \\
        \end{vmatrix} \\
        &= (a+b+c)^2 \begin{vmatrix}
        (b+c)^2 & a-b-c & a-b-c \\
        b^2 & a-b+c & 0 \\
        c^2 & 0 & a+b-c \\
        \end{vmatrix}
        \end{align*}
        进行行变换 \( R_1 \rightarrow R_1 - R_2 - R_3 \),得
        \[
        D = (a+b+c)^2 \begin{vmatrix}
        2bc & -2c & -2b \\
        b^2 & a-b+c & 0 \\
        c^2 & 0 & a+b-c \\
        \end{vmatrix}
        \]
        再进行列变换 \( C_2 \rightarrow C_2 + \dfrac{1}{b}C_1,\; C_3 \rightarrow C_3 + \dfrac{1}{c}C_1 \),得
        \[
        D = (a+b+c)^2 \begin{vmatrix}
        2bc & 0 & 0 \\
        b^2 & a+c & \dfrac{b^2}{c} \\
        c^2 & \dfrac{c^2}{b} & a+b \\
        \end{vmatrix}
        \]
        按第一行展开行列式得
        \begin{align*}
        D &= (a+b+c)^2 \cdot 2bc \left[(a+c)(a+b) - \frac{b^2}{c} \cdot \frac{c^2}{b} \right] \\
        &= 2bc(a+b+c)^2 \left(a^2 + ab + ac + bc - bc \right) \\
        &= 2abc(a+b+c)^3
        \end{align*}
    \end{solution}

    \question 证明
    \[
    \left|
    \begin{matrix}
    ax + by & ay + bz & az + bx \\
    ay + bz & az + bx & ax + by \\
    az + bx & ax + by & ay + bz \\
    \end{matrix}
    \right|
    =
    \left(a^3 + b^3\right)
    \left|
    \begin{matrix}
    x & y & z \\
    y & z & x \\
    z & x & y \\
    \end{matrix}
    \right|
    \]
    \begin{solution}
        令\[
        D=\begin{vmatrix}
        ax + by & ay + bz & az + bx \\
        ay + bz & az + bx & ax + by \\
        az + bx & ax + by & ay + bz \\
        \end{vmatrix}
        \]
        按$C_1$拆分
        \begin{align*}
            D&=\begin{vmatrix}
            ax  & ay + bz & az + bx \\
            ay  & az + bx & ax + by \\
            az  & ax + by & ay + bz \\
            \end{vmatrix} \;+\;       
            \begin{vmatrix}
            by & ay + bz & az + bx \\
            bz & az + bx & ax + by \\
            bx & ax + by & ay + bz \\
            \end{vmatrix}\\
            &= a\begin{vmatrix}
            x  & ay + bz & az + bx \\
            y  & az + bx & ax + by \\
            z  & ax + by & ay + bz \\
            \end{vmatrix} \;+\;b       
            \begin{vmatrix}
            y & ay + bz & az + bx \\
            z & az + bx & ax + by \\
            x & ax + by & ay + bz \\
            \end{vmatrix}
        \end{align*}
        第一个行列式进行$C_3 \rightarrow C_3-bC_1$,第二个行列式进行$C_2 \rightarrow C_2-aC_1$得
        \begin{align*}
            D&=a\begin{vmatrix}
            x & ay + bz & az \\
            y & az + bx & ax \\
            z & ax + by & ay \\
            \end{vmatrix}
            +b\begin{vmatrix}
            y & bz & az + bx \\
            z & bx & ax + by \\
            x & by & ay + bz \\
            \end{vmatrix}\\
            &=a^2\begin{vmatrix}
            x & ay + bz & z \\
            y & az + bx & x \\
            z & ax + by & y \\
            \end{vmatrix}
            +b^2\begin{vmatrix}
            y & z & az + bx \\
            z & x & ax + by \\
            x & y & ay + bz \\
            \end{vmatrix}
        \end{align*}第一个行列式进行$C_2 \rightarrow C_2-bC_3$,第二个行列式进行$C_3 \rightarrow C_3-aC_2$得
        \begin{align*}
            D&=a^2\begin{vmatrix}
            x & ay & z \\
            y & az & x \\
            z & ax & y \\
            \end{vmatrix}
            +b^2\begin{vmatrix}
            y & z & bx \\
            z & x & by \\
            x & y & bz \\
            \end{vmatrix}\\
            &=(a^3+b^3)\begin{vmatrix}
                x & y & z \\
                y & z & x \\
                z & x & y \\
                \end{vmatrix}
        \end{align*}
    \end{solution}
    
    \question 证明
    \[
    \begin{vmatrix} 
    a+bx & c+dx & p+qx \\ 
    ax+b & cx+d & px+q \\ 
    u & v & w 
    \end{vmatrix} 
    = (1-x^2) 
    \begin{vmatrix} 
    a & c & p \\ 
    b & d & q \\ 
    u & v & w 
    \end{vmatrix}
    \]
    \begin{solution}
        令
        \[
        D=
        \begin{vmatrix} 
        a+bx & c+dx & p+qx \\ 
        ax+b & cx+d & px+q \\ 
        u & v & w 
        \end{vmatrix} 
        \]
        进行列行变换 \( R_1 \rightarrow R_1 - R_2 \),得
        \begin{align*}
        D &=
        \begin{vmatrix} 
        a+bx-x(ax+b) & c+dx-x(cx+d) & p+qx-x(px+q) \\ 
        ax+b & cx+d & px+q \\ 
        u & v & w 
        \end{vmatrix} \\
        &= 
        \begin{vmatrix} 
        a-ax^2 & c-cx^2 & p-px^2 \\ 
        ax+b & cx+d & px+q \\ 
        u & v & w 
        \end{vmatrix} \\
        &= (1-x^2) 
        \begin{vmatrix} 
        a & c & p \\ 
        ax+b & cx+d & px+q \\
        u & v & w 
        \end{vmatrix}
        \end{align*}
        按$R_2$拆分,得
        \[
        D= 0+(1-x^2) 
        \begin{vmatrix} 
        a & c & p \\ 
        b & d & q \\
        u & v & w 
        \end{vmatrix}
        = (1-x^2)
        \begin{vmatrix} 
        a & c & p \\ 
        b & d & q \\
        u & v & w 
        \end{vmatrix}
        \]
    \end{solution}

    \question 证明
    \[
    \begin{vmatrix} 
    a & a^2 & bc \\ 
    b & b^2 & ca \\ 
    c & c^2 & ab 
    \end{vmatrix} 
    = 
    \begin{vmatrix} 
    1 & a^2 & a^3 \\ 
    1 & b^2 & b^3 \\ 
    1 & c^2 & c^3 
    \end{vmatrix}
    \]
    \begin{solution}
        令
        \[
        D=
        \begin{vmatrix} 
        a & a^2 & bc \\ 
        b & b^2 & ca \\ 
        c & c^2 & ab 
        \end{vmatrix} 
        \]
        将$a,b,c$分别乘入$R_1,R_2,R_3$
        \[
        D= \frac{1}{abc} 
        \begin{vmatrix} 
        a^2 & a^3 & abc \\ 
        b^2 & b^3 & abc \\ 
        c^2 & c^3 & abc 
        \end{vmatrix} 
        = \frac{abc}{abc} 
        \begin{vmatrix} 
        a^2 & a^3 & 1 \\ 
        b^2 & b^3 & 1 \\ 
        c^2 & c^3 & 1 
        \end{vmatrix} 
        \]
        进行$C_2 \leftrightarrow C_3$,再$C_1 \leftrightarrow C_2$得
        \[
        D= - 
        \begin{vmatrix} 
        a^2 & 1 & a^3 \\ 
        b^2 & 1 & b^3 \\ 
        c^2 & 1 & c^3 
        \end{vmatrix} 
        = 
        \begin{vmatrix} 
        1 & a^2 & a^3 \\ 
        1 & b^2 & b^3 \\ 
        1 & c^2 & c^3 
        \end{vmatrix}
        \]
    \end{solution}

    \question 证明
    \[
    \begin{vmatrix} 
    \sin\alpha & \cos\alpha & \cos(\alpha+\delta) \\ 
    \sin\beta & \cos\beta & \cos(\beta+\delta) \\ 
    \sin\gamma & \cos\gamma & \cos(\gamma+\delta) 
    \end{vmatrix} = 0
    \]
    \begin{solution}
        令\[
        D=    
        \begin{vmatrix} 
        \sin\alpha & \cos\alpha & \cos(\alpha+\delta) \\ 
        \sin\beta & \cos\beta & \cos(\beta+\delta) \\ 
        \sin\gamma & \cos\gamma & \cos(\gamma+\delta) 
        \end{vmatrix}
        \]
        由三角公式 $\cos(\theta+\delta) = \cos\theta\cos\delta - \sin\theta\sin\delta$,将$C_3$展开得
        \[
        D=\begin{vmatrix} \sin\alpha & \cos\alpha & \cos(\alpha+\delta) \\ \sin\beta & \cos\beta & \cos(\beta+\delta) \\ \sin\gamma & \cos\gamma & \cos(\gamma+\delta) \end{vmatrix} 
        = \cos\delta \begin{vmatrix} \sin\alpha & \cos\alpha & \cos\alpha \\ \sin\beta & \cos\beta & \cos\beta \\ \sin\gamma & \cos\gamma & \cos\gamma \end{vmatrix} 
        - \sin\delta \begin{vmatrix} \sin\alpha & \cos\alpha & \sin\alpha \\ \sin\beta & \cos\beta & \sin\beta \\ \sin\gamma & \cos\gamma & \sin\gamma \end{vmatrix} 
        \]
        其中两个行列式都有两列成比例,因此
        \[
        D=0
        \]
    \end{solution}
    
    \question 证明
    \[
    \begin{vmatrix}
    \cos^2 a & \sin a & \sin^2 a \\
    -\cos 2a & \cos a - \sin a & \cos 2a \\
    -\sin 2a - \cos 2a & 2 \cos a & 1 + \sin 2a + \cos 2a
    \end{vmatrix} = \sin a \cos a (\cos a - \sin a)
    \]
    \begin{solution}
        令
        \[
        D=\begin{vmatrix}
        \cos^2 a & \sin a & \sin^2 a \\
        -\cos 2a & \cos a - \sin a & \cos 2a \\
        -\sin 2a - \cos 2a & 2 \cos a & 1 + \sin 2a + \cos 2a
        \end{vmatrix}
        \]
        进行列变换 $C_1 \rightarrow C_1 - C_3$,
        \[
        D=\begin{vmatrix}
        1 & \sin a & \sin^2 a \\
        0 & \cos a - \sin a & \cos 2a \\
        1 & 2 \cos a & 1 + \sin 2a + \cos 2a
        \end{vmatrix}
        \]
        进行行变换 $R_3 \rightarrow R_3 - R_2$,
        \[
        D= \begin{vmatrix}
        1 & \sin a & \sin^2 a \\
        0 & \cos a - \sin a & \cos 2a \\
        1 & \cos a + \sin a & 1 + \sin 2a
        \end{vmatrix}
        \]
        进行行变换 $R_2 \rightarrow R_2 + R_1$,化为一范德蒙行列式,
        \[
        D= \begin{vmatrix}
        1 & \sin a & \sin^2 a \\
        1 & \cos a & \cos^2 a \\
        1 & \cos a + \sin a & (\cos a + \sin a)^2
        \end{vmatrix} 
        = \sin a \cos a (\cos a - \sin a)
        \]  
    \end{solution}

    \question 证明
        \[
        \begin{vmatrix}
        x-1 & x-2 & x-3 & x-4 \\
        x^2 - 1 & x^2 - 2^2 & x^2 - 3^2 & x^2 - 4^2 \\
        x^3 - 1 & x^3 - 2^3 & x^3 - 3^3 & x^3 - 4^3 \\
        x^4 - 1 & x^4 - 2^4 & x^4 - 3^4 & x^4 - 4^4 \\
        \end{vmatrix}
        =12(x - 1)(x - 2)(x - 3)(x - 4)
        \]
    \begin{solution}
        令
        \[
        D=\begin{vmatrix}
        x-1 & x-2 & x-3 & x-4 \\
        x^2 - 1 & x^2 - 2^2 & x^2 - 3^2 & x^2 - 4^2 \\
        x^3 - 1 & x^3 - 2^3 & x^3 - 3^3 & x^3 - 4^3 \\
        x^4 - 1 & x^4 - 2^4 & x^4 - 3^4 & x^4 - 4^4 \\
        \end{vmatrix}        
        \]
        提取公因式得$D=(x - 1)(x - 2)(x - 3)(x - 4)M$,其中
        \[
        M=
        \begin{vmatrix}
        1 & 2 & 3 & 4 \\
        x+1 & x+2 & x+3 & x+4 & \\
        x^2+x+1 & x^2+2x+2^2 & x^2+3x+3^2 & x^2+4x+4^2 & \\
        x^3+x^2+x+1 & x^3+2x^2+2^2x+2^3 & x^3+3x^2+3^2x+3^3 & x^3+4x^2+4^2x+4^3 & \\
        \end{vmatrix}        
        \]
        依次进行行变换$R_4 \rightarrow R_4-xR_3,R_3 \rightarrow R_3-xR_2,R_2 \rightarrow R_2-xR_1,$
        \[
        M=
        \begin{vmatrix}
        1 & 2 & 3 & 4 \\
        1^2 & 2^2 & 3^2 & 4^2 \\
        1^3 & 2^3 & 3^3 & 4^3 \\
        1^4 & 2^4 & 3^4 & 4^4 \\
        \end{vmatrix}
        \]
        是一范德蒙行列式,
        \[
        M= \prod_{1\leq i<j\leq4}(j−i)=(2−1)(3-1)(3−2)(4−1)(4−2)(4−3)=1\cdot2\cdot1\cdot3\cdot2\cdot1= 12
        \]
        于是得证
        \[
        D=12(x - 1)(x - 2)(x - 3)(x - 4)
        \]
    \end{solution}

    \question 设 $V_n$ 为 $n$ 阶范德蒙行列式,定义如下:
    \[
    V_n = \begin{vmatrix}
    1 & x_1 & x_1^2 & \dots & x_1^{n-1} \\
    1 & x_2 & x_2^2 & \dots & x_2^{n-1} \\
    \vdots & \vdots & \vdots & \ddots & \vdots \\
    1 & x_n & x_n^2 & \dots & x_n^{n-1}
    \end{vmatrix}
    \]
    证明
    \[
    V_n = \prod_{1 \le i < j \le n} (x_j - x_i).
    \]
    \begin{solution}
        设
        \[
        V_n = \begin{vmatrix}
        1 & x_1 & x_1^2 & \dots & x_1^{n-1} \\
        1 & x_2 & x_2^2 & \dots & x_2^{n-1} \\
        1 & x_3 & x_3^2 & \dots & x_3^{n-1} \\
        \vdots & \vdots & \vdots & \ddots & \vdots \\
        1 & x_{n-1} & x_{n-1}^2 & \dots & x_{n-1}^{n-1} \\
        1 & x_n & x_n^2 & \dots & x_n^{n-1}
        \end{vmatrix}
        \]
        进行行变换$R_2\rightarrow R_2-R_1,R_3\rightarrow R_3-R_1,\cdots R_n\rightarrow R_n-R_1$,得
        \[
        V_n = \begin{vmatrix}
        1 & x_1 & x_1^2 & \dots & x_1^{n-2} & x_1^{n-1} \\
        0 & x_2-x_1 & x_2^2-x_1^2 & \dots & x_2^{n-2}-x_1^{n-2} & x_2^{n-1}-x_1^{n-1} \\
        0 & x_3-x_1 & x_3^2-x_1^2 & \dots & x_3^{n-2}-x_1^{n-2} & x_3^{n-1}-x_1^{n-1} \\
        \vdots & \vdots & \vdots & \ddots & \vdots & \vdots \\
        0 & x_{n-1}-x_1 & x_{n-1}^2-x_1^2 & \dots & x_{n-1}^{n-2}-x_1^{n-2} & x_{n-1}^{n-1}-x_1^{n-1} \\
        0 & x_n-x_1 & x_n^2-x_1^2 & \dots & x_n^{n-2}-x_1^{n-2} & x_n^{n-1}-x_1^{n-1}
        \end{vmatrix}
        \]
        进行列变换$C_2\rightarrow C_2-C_1,C_3\rightarrow C_3-C_2,\cdots C_n\rightarrow C_n-C_{n-1}$,得
        \[
        V_n =
        \begin{vmatrix}
        1 & 0 & 0 & \cdots & 0 & 0 \\
        0 & x_2-x_1 & (x_2-x_1)x_2 & \cdots & (x_2-x_1)x_2^{n-3} & (x_2-x_1)x_2^{n-2} \\
        0 & x_3-x_1 & (x_3-x_1)x_3 & \cdots & (x_3-x_1)x_3^{n-3} & (x_3-x_1)x_3^{n-2} \\
        \vdots & \vdots & \vdots & \ddots & \vdots & \vdots \\
        0 & x_{n-1}-x_1 & (x_{n-1}-x_1)x_{n-1} & \cdots & (x_{n-1}-x_1)x_{n-1}^{n-3} & (x_{n-1}-x_1)x_{n-1}^{n-2} \\
        0 & x_n-x_1 & (x_n-x_1)x_n & \cdots & (x_n-x_1)x_n^{n-3} & (x_n-x_1)x_n^{n-2}
        \end{vmatrix} 
        \]
        提取公因式,
        \[
        V_n= \prod_{k=2}^{n} (x_k-x_1)
        \begin{vmatrix}
        1 & 0 & 0 & \cdots & 0 & 0 \\
        0 & 1 & x_2 & \cdots & x_2^{n-3} & x_2^{n-2} \\
        0 & 1 & x_3 & \cdots & x_3^{n-3} & x_3^{n-2} \\
        \vdots & \vdots & \vdots & \ddots & \vdots & \vdots \\
        0 & 1 & x_{n-1} & \cdots & x_{n-1}^{n-3} & x_{n-1}^{n-2} \\
        0 & 1 & x_n & \cdots & x_n^{n-3} & x_n^{n-2}
        \end{vmatrix}
        \]
        按$R_1$展开得
        \[
        V_n = \prod_{k=2}^{n} (x_k-x_1)
        \begin{vmatrix}
        1 & x_2 & \cdots & x_2^{n-3} & x_2^{n-2} \\
        1 & x_3 & \cdots & x_3^{n-3} & x_3^{n-2} \\
        \vdots & \vdots & \ddots & \vdots & \vdots \\
        1 & x_{n-1} & \cdots & x_{n-1}^{n-3} & x_{n-1}^{n-2} \\
        1 & x_n & \cdots & x_n^{n-3} & x_n^{n-2}
        \end{vmatrix}
        = \prod_{k=2}^{n}(x_k-x_1)\,V_{n-1}
        \]
        依次递推得证
        \[
        V_n = \prod_{1 \le i < j \le n} (x_j - x_i)
        \]
        其中
        \[
        V_2=\begin{vmatrix}1 & x_{n-1}\\1 & x_n\end{vmatrix}=x_n-x_{n-1}
        \]
    \end{solution}

    \question 设
    \[
    \begin{cases}
    k x_1 + x_2 + x_3 = 5, \\
    3 x_1 + 2 x_2 + k x_3 = 18 - 5k, \\
    x_2 + 2 x_3 = 2,
    \end{cases}
    \]
    问 $k$ 取何值时,方程组无解、有唯一解、有无穷解？在有无穷解时,求全部解。
    \begin{solution}
        原方程组为$ \mathbf{A} \mathbf{x}=\mathbf{b}$,其中
        \[
        \mathbf{A} = \begin{bmatrix}
        k & 1 & 1 \\
        3 & 2 & k \\
        0 & 1 & 2
        \end{bmatrix},
        \quad
        \mathbf{b} = \begin{bmatrix} 5 \\ 18 - 5k \\ 2 \end{bmatrix}.
        \]
        解
        \[
        \det(\mathbf{A}) = 4k-3-k^2-6= -k^2 + 4k - 3 = 0
        \]
        得$k=1,3$, 因此:
        \begin{itemize}
        \item 当 $k \ne 1, 3$ 时,$\det(\mathbf{A}) \ne 0$,方程组有唯一解；
        \item 当 $k = 1$ 时,代入方程组得:
        \[
        \begin{cases}
        x_1 + x_2 + x_3 = 5, \\
        3x_1 + 2x_2 + x_3 = 13, \\
        x_2 + 2x_3 = 2.
        \end{cases}
        \]
        用第三式解出 $x_2 = 2 - 2x_3$,代入前两式得
        \[
        \begin{aligned}
        x_1 + (2 - 2x_3) + x_3 &= 5 \Rightarrow x_1 = 3 + x_3, \\
        3x_1 + 2x_2 + x_3 &= 13 \Rightarrow 3(3 + x_3) + 2(2 - 2x_3) + x_3 = 13.
        \end{aligned}
        \]
        检验等式成立:
        \[
        9 + 3x_3 + 4 - 4x_3 + x_3 = 13 \Rightarrow 13 = 13.
        \]
        成立,说明方程组有无穷多解,设 $x_3 = \lambda$,则
        \[
        x_1 = 3 + \lambda, \quad
        x_2 = 2 - 2\lambda, \quad
        x_3 = \lambda \quad (\lambda \in \mathbb{R}).
        \]
        \item 当 $k = 3$ 时,代入原方程组得
        \[
        \begin{cases}
        3x_1 + x_2 + x_3 = 5, \\
        3x_1 + 2x_2 + 3x_3 = 3, \\
        x_2 + 2x_3 = 2.
        \end{cases}
        \]
        由前两式相减得:
        \[
        (3x_1 + 2x_2 + 3x_3) - (3x_1 + x_2 + x_3) = -2 \Rightarrow x_2 + 2x_3 = -2,
        \]
        与第三式 $x_2 + 2x_3 = 2$ 矛盾,因此方程组无解。
        \end{itemize}
    \end{solution}

\question 将帕斯卡三角形（Pascal's Triangle）写成如下无穷矩阵的形式:
\[
\begin{pmatrix}
1 & 1 & 1 & 1 & 1 & \dots \\
1 & 2 & 3 & 4 & 5 & \dots \\
1 & 3 & 6 & 10 & 15 & \dots \\
1 & 4 & 10 & 20 & 35 & \dots \\
1 & 5 & 15 & 35 & 70 & \dots \\
\vdots & \vdots & \vdots & \vdots & \vdots & \ddots
\end{pmatrix}
\]
其中第一行和第一列全部由 1 组成,其他每个元素由其左侧元素与上方元素之和构成。对于每个正整数 $n$,令 $D_n$ 表示由该阵列的前 $n$ 行和前 $n$ 列组成的 $n \times n$ 矩阵。求 $D_n$ 的行列式并证明你的结论。

\begin{solution}
对较小 $n$ 值的直接计算表明,对于所有 $n$,恒有 $\det(D_n) = 1$。

\textbf{证明大纲:}

对 $D_n$ 按顺序执行以下初等行变换,依次处理第 $n$ 行至第 2 行:
\begin{enumerate}
    \item 从第 $n$ 行减去第 $n-1$ 行；
    \item 从第 $n-1$ 行减去第 $n-2$ 行；
    \item[\dots]
    \item[$(n-1).$] 从第 2 行减去第 1 行。
\end{enumerate}

由于该阵列的构造满足 $a_{i,j} = a_{i-1,j} + a_{i,j-1}$,上述操作后,$D_n$ 的第 $i$ 行第 $j$ 列元素变为:
\[
a_{i,j}' = a_{i,j} - a_{i-1,j} = a_{i,j-1}
\]
这会将每一列向右“平移”一个位置,并使第一列变为 $[1, 0, 0, \dots, 0]^T$。

对剩余的子矩阵重复此过程（即对第 $n$ 行至第 3 行执行减法,以此类推）,最终会将 $D_n$ 转化为一个对角线元素均为 1 的上三角矩阵。由于上三角矩阵的行列式等于主对角线元素的乘积,因此 $\det(D_n) = 1^n = 1$。

\textbf{另一种方法（递推法）:}
在进行第一组行变换后,按第一列进行余子式展开,可以直接得到递推关系:
\[
\det(D_n) = 1 \cdot \det(D_{n-1})
\]
由于 $D_1 = [1]$ 且 $\det(D_1) = 1$,通过数学归纳法可知 $\det(D_n) = 1$。
\end{solution}

\question 证明
    \[
    \left|
    \begin{matrix}
    a & b & c & d \\
    -b & a & -d & c \\
    -c & d & a & -b \\
    -d & -c & b & a \\
    \end{matrix}
    \right|
    =
    \left(a^2 + b^2 + c^2 + d^2\right)^2
    \]  
\begin{solution}
    \textcolor{red}{(待验证)}
\end{solution}  

\question 求行列式$D_n = 
    \begin{vmatrix}
    1 & 1 & 1 & \cdots & 1 & 1 & -n \\
    1 & 1 & 1 & \cdots & 1 & -n & 1 \\
    1 & 1 & 1 & \cdots & -n & 1 & 1 \\
    \vdots & \vdots & \vdots & \ddots & \vdots & \vdots & \vdots \\
    -n & 1 & 1 & \cdots & 1 & 1 & 1
    \end{vmatrix}$
    \begin{solution}
    这是一个 $n\times n$ 的行列式,每一行有 $n{-}1$ 个 $1$ 和一个 $-n$,且 $-n$ 分别出现在每行的不同列。
    
    设该行列式为 $D_n$,我们将其简化。
    
    \textbf{第一步:行变换}。
    
    令第 $i$ 行减去第 $n$ 行,$i = 1,2,\ldots,n-1$,则得到一个新行列式 $D_n'$,变换后前 $n{-}1$ 行只有两个非零元素（$1{-}1=0$,$-n{-}1=-(n{+}1)$,$1{-}(-n)=n{+}1$ 等等）,便于处理。
    
    经此变换,前 $n{-}1$ 行在对角线上为 $n{+}1$,其余为 $0$。
    
    \textbf{第二步:观察结构}。
    
    经过上述行变换,$D_n'$ 为一个上三角矩阵,其前 $n{-}1$ 个对角元为 $n+1$,最后一行为原来的第 $n$ 行,未变动。
    
    于是可得:
    \[
    \begin{aligned}
    D_n &= \begin{vmatrix}
    n+1 & 0   & 0   & \cdots & 0   & a_1 \\
    0   & n+1 & 0   & \cdots & 0   & a_2 \\
    0   & 0   & n+1 & \cdots & 0   & a_3 \\
    \vdots & \vdots & \vdots & \ddots & \vdots & \vdots \\
    0   & 0   & 0   & \cdots & n+1 & a_{n-1} \\
    - n & 1   & 1   & \cdots & 1   & 1
    \end{vmatrix}
    \end{aligned}
    \]
    
    将该行列式按最后一列展开,只需考虑对角线乘积（因为其余部分为 0）,于是有:
    \[
    D_n = (n+1)^{n-1} \cdot \text{余子式项}.
    \]
    
    由于我们做了 $(n-1)$ 次行变换,每次减去第 $n$ 行,所以符号为 $(-1)^{\frac{n(n+1)}{2}}$（由置换中逆序数判断）。
    
    \textbf{最终结果为}:
    \[
    \boxed{D_n = (-1)^{\frac{n(n+1)}{2}} (n+1)^{n-1}}.
    \]
    \textcolor{red}{(待验证)}
\end{solution}

\question 设 $M$ 为一个任意 $3 \times 3$ 矩阵,其每个项 $m_{ij}$ 均选自集合 $\{-1, 1\}$,记 $\det M$ 为其行列式。证明 $|\det M| \leq 4$。
%missouri 2011 II p3
\begin{solution}
一个项为 $\pm 1$ 的 $2 \times 2$ 矩阵,其行列式必然为 0 或 $\pm 2$。对 $3 \times 3$ 的 $\det M$ 按第一行进行余子式展开:
\[
\det M = m_{11} \begin{vmatrix} m_{22} & m_{23} \\ m_{32} & m_{33} \end{vmatrix} - m_{12} \begin{vmatrix} m_{21} & m_{23} \\ m_{31} & m_{33} \end{vmatrix} + m_{13} \begin{vmatrix} m_{21} & m_{22} \\ m_{31} & m_{32} \end{vmatrix}
\]
由于每个余子式都是偶数,左侧的结果必然也是偶数。同时由上式可知 $|\det M|$ 不可能超过 6。只有当所有余子式都不为 0 时才可能达到 6。但在由 1 和 -1 组成的 $2 \times 3$ 矩阵中,总能找到一个 $2 \times 2$ 的子矩阵其行列式为 0。因此,上述展开式中至少有一项为 0,从而得出 $|\det M| \leq 4$。
\end{solution}

    \question 计算 $n \times n$ 矩阵 $A = [a_{ij}]$ 的行列式,其中
    \[
    a_{ij} = 
    \begin{cases} 
    (-1)^{|i-j|}, & i \neq j, \\
    2, &  i = j.
    \end{cases}
    \]
    \begin{solution}
        进行行变换$R_1 \to R_1 + R_2,R_2 \to R_2 + R_3,\cdots R_{n-1} \to R_{n-1} + R_n$,有
        \[
        \det(A)=
        \begin{vmatrix}
        2 & -1 & +1 & \cdots & \pm1 & \mp1 \\
        -1 & 2 & -1 & \cdots & \pm1 & \mp1 \\
        +1 & -1 & 2 & \cdots & \pm1 & \mp1 \\
        \vdots & \vdots & \vdots & \ddots & \vdots & \vdots \\
        \mp1 & \pm1 & \mp1 & \cdots & 2 & -1 \\
        \pm1 & \mp1 & \pm1 & \cdots & -1 & 2
        \end{vmatrix}
        =
        \begin{vmatrix}
        1 & 1 & 0 &  \cdots & 0 & 0 \\
        0 & 1 & 1 &  \cdots & 0 & 0 \\
        0 & 0 & 1 &  \cdots & 0 & 0 \\
        \vdots & \vdots & \vdots & \ddots & \vdots & \vdots \\
        0 & 0 & 0 &  \cdots & 1 & 1 \\
        \pm1 & \mp1 & \pm1 & \cdots & -1 & 2
        \end{vmatrix}
        \]
        进行列变换$C_2 \to C_2 - C_1,C_3 \to C_3 - C_2, \cdots C_n \to C_n - C_{n-1}$,得到
        \[
        \det(A)=
        \begin{vmatrix}
        1 & 0 & 0 & \cdots & 0 & 0 \\
        0 & 1 & 0 & \cdots & 0 & 0 \\
        0 & 0 & 1 & \cdots & 0 & 0 \\
        \vdots & \vdots & \vdots & \ddots & \vdots & \vdots \\
        0 & 0 & 0 & \cdots & 1 & 0 \\
        0 & 0 & 0 & \cdots & 0 & n+1
        \end{vmatrix}
        = n+1
        \]
    \end{solution}

    \question 设对所有 $j=0,\dots,n$皆有 
    \[
    a_j = a_0 + jd,
    \]
    其中 $a_0, d$ 为实数,令
    \[
    A = \begin{pmatrix}
    a_0 & a_1 & a_2 & \cdots & a_n \\
    a_1 & a_0 & a_1 & \cdots & a_{n-1} \\
    a_2 & a_1 & a_0 & \cdots & a_{n-2} \\
    \vdots & \vdots & \vdots & \ddots & \vdots \\
    a_n & a_{n-1} & a_{n-2} & \cdots & a_0
    \end{pmatrix}
    \]
    计算 $\det(A)$, 其中 $\det(A)$ 表示 $A$ 的行列式。
    \begin{solution}
        进行列变换$C_{n+1} \to C_{n+1} + C_1$,得到
        \[
        \det(A) = (a_0 + a_n) \det \begin{pmatrix}
        a_0 & a_1 & a_2 & \cdots & 1 \\
        a_1 & a_0 & a_1 & \cdots & 1 \\
        a_2 & a_1 & a_0 & \cdots & 1 \\
        \vdots & \vdots & \vdots & \ddots & \vdots \\
        a_n & a_{n-1} & a_{n-2} & \cdots & 1
        \end{pmatrix}
        \]
        进行行变换$R_{n+1} \to R_{n+1} -R_n,R_n \to R_n -R_{n-1}, \cdots R_2 \to R_2-R_1$,得到
        \[
        \det(A) = (a_0 + a_n) \det \begin{pmatrix}
        a_0 & a_1 & a_2 & \cdots & 1 \\
        d & -d & -d & \cdots & 0 \\
        d & d & -d & \cdots & 0 \\
        \vdots & \vdots & \vdots & \ddots & \vdots \\
        d & d & d & \cdots & 0
        \end{pmatrix}
        \]
        按$C_{n+1}$展开得
        \[
        \det(A) = (-1)^n (a_0 + a_n) \det \begin{pmatrix}
        d & -d & -d & \cdots & -d \\
        d & d & -d & \cdots & -d \\
        d & d & d & \cdots & -d \\
        \vdots & \vdots & \vdots & \ddots & \vdots \\
        d & d & d & \cdots & d
        \end{pmatrix}.
        \]
        将上述矩阵的最后一行加至其他各行,有
        \[
        \det(A) = (-1)^n (a_0 + a_n) \det \begin{pmatrix}
        2d & 0 & 0 & \cdots & 0 \\
        2d & 2d & 0 & \cdots & 0 \\
        2d & 2d & 2d & \cdots & 0 \\
        \vdots & \vdots & \vdots & \ddots & \vdots \\
        d & d & d & \cdots & d
        \end{pmatrix} = (-1)^n (a_0 + a_n) 2^{n-1} d^n
        \]
    \end{solution}
    
    \question 已知$n > 1$ 是一个奇正整数,设 $A = (a_{ij})_{i,j=1\dots n}$ 是一个 $n \times n$ 矩阵,其元素定义为:
    \[
    a_{ij} = 
    \begin{cases} 
    2 & ,\;i = j \\ 
    1 & ,\;i - j \equiv \pm 2 \pmod n \\ 
    0 & ,\;\text{反之} 
    \end{cases}
    \]
    求 $\det A$。
    \begin{solution}
        注意到 $A = B^2$,其中矩阵 $B$ 的元素定义为:
        \[
        b_{ij} = 
        \begin{cases} 
        1 & ,\;i - j \equiv \pm 1 \pmod n \\ 
        0 & ,\;\text{反之} 
        \end{cases}
        \]
        计算 $\det B$,先按$R_1$展开行列式,然后对得到的两项分别按$C_1$展开。
        \begin{align*}
        \det B &= 
        \begin{vmatrix} 
        0 & 1 & & & & & 1 \\ 
        1 & 0 & 1 & & & \\ 
        & 1 & 0 & 1 & & \\ 
        & & 1 & \ddots & \ddots & \\ 
        & & & \ddots & 0 & 1 \\ 
        & & & & 1 & 0 & 1 \\ 
        1 & & & & & 1 & 0 
        \end{vmatrix} 
        = - \begin{vmatrix} 
        1 & 1 & & & & \\ 
        & 0 & 1 & & &  \\ 
        & 1 & \ddots & \ddots & & \\ 
        & & \ddots & 0 & 1 & \\ 
        & & & 1 & 0 & 1 \\ 
        1 & & & & 1 & 0 
        \end{vmatrix} 
        + \begin{vmatrix} 
        1 & 0 & 1 & & & \\ 
        & 1 & 0 & 1 & & \\ 
        & & 1 & \ddots & \ddots & \\ 
        & & & \ddots & 0 & 1 \\ 
        & & & & 1 & 0 \\ 
        1 & & & & & 1 
        \end{vmatrix} \\
        &= - \left( \begin{vmatrix} 
        0 & 1 & & & \\ 
        1 & \ddots & \ddots & & \\ 
        & \ddots & 0 & 1 & \\ 
        & & 1 & 0 & 1 \\ 
        & & & 1 & 0 
        \end{vmatrix} - 
        \begin{vmatrix} 
        1 & & & & \\ 
        0 & 1 & & & \\ 
        1 & \ddots & \ddots & & \\ 
        & \ddots & 0 & 1 \\ 
        & & 1 & 0 & 1 
        \end{vmatrix} \right) + 
        \left( \begin{vmatrix} 
        1 & 0 & 1 & & \\ 
        & 1 & \ddots & \ddots & \\ 
        & & \ddots & 0 & 1 \\ 
        & & & 1 & 0 \\ 
        & & & & 1 
        \end{vmatrix} - \begin{vmatrix} 
        0 & 1 & & & \\ 
        1 & 0 & 1 & & \\ 
        & 1 & \ddots & \ddots & \\ 
        & & \ddots & 0 & 1 \\ 
        & & & 1 & 0 \end{vmatrix} \right) \\
        &= -(0 - 1) + (1 - 0) = 2
        \end{align*}
        其中第二和第三个矩阵是下/上三角矩阵,而第一个和第四个矩阵满足 
        \[
        R_1 - R_3 + R_5 - \dots \pm R_{n-2} = (0)_{1\times(n-2)},
        \]
        故 $\det B = 2$,从而 
        \[
        \det A = (\det B)^2 = 4
        \]
    \end{solution}
\end{questions}
\pagebreak

\begin{center}
  {\fontsize{30pt}{26pt}\selectfont
    \hypertarget{复数}{复数} \label{复数}
  }
\end{center}
\separator
\vspace{1pt}

\begin{questions}
    \question 已知 $i = \sqrt{-1}$, 若 
    \[
    \frac{1}{i^{2025}} - \frac{2}{i^{2024}} + \frac{3}{i^{2023}} - \frac{4}{i^{2022}} + \dots - \frac{2024}{i^2} + \frac{2025}{i} = a+bi
    \]
    其中 $a,b$ 为实数, 求 $a-b$。
    \begin{solution}
        \begin{align*}
        &\frac{1}{i^{2025}} - \frac{2}{i^{2024}} + \frac{3}{i^{2023}} - \frac{4}{i^{2022}} + \cdots - \frac{2024}{i^2} + \frac{2025}{i} \\
        &= \frac{1}{i} - 2 - \frac{3}{i} + 4 + \frac{5}{i} - 6 - \frac{7}{i} + 8 + \cdots + 2024 + \frac{2025}{i} \\
        &= \frac{1}{i}\,(1 + 5 + \cdots + 2025) - (2 + 6 + \cdots + 2022) - \frac{1}{i}\,(3 + 7 + \cdots + 2023) + (4 + 8 + \cdots + 2024) \\
        &= \frac{1}{i} \cdot \frac{2026 \cdot 507}{2} - \frac{2024 \cdot 506}{2} - \frac{1}{i} \cdot \frac{2026 \cdot 506}{2} + \frac{2028 \cdot 506}{2} \\
        &= \frac{1013}{i} + 1012 \\
        &= 1012 - 1013 i \quad \Rightarrow \quad a - b = 2025
        \end{align*}
    \end{solution}

    \question 设 $z$ 是 $1$ 的七次方根,且 $z \neq 1$,试求 $z + z^2 + z^4$的值。
    \begin{solution}
        发现
        \[
        z^7 = 1 \Rightarrow 1 + z + \cdots + z^6 = 0
        \Rightarrow z + z^2 + \cdots + z^6 = -1
        \]
        设$\alpha = z + z^2 + z^4 ,\beta = z^3 + z^5 + z^6$,则
        \[
        \alpha + \beta = z + z^2 + \cdots + z^6 = -1
        \]
        又
        \[
        \alpha\beta = (z + z^2 + z^4)(z^3 + z^5 + z^6) = z^4(1 + z + z^2 + \cdots + z^6 + 2z^3)= z^4(0 + 2z^3) = 2z^7 = 2
        \]
        因此$\alpha, \beta$是方程$x^2 + x + 2 = 0$ 的两根,解得
        \[
        \alpha = z + z^2 + z^4 = \frac{-1 \pm \sqrt{7}i}{2}
        \]
    \end{solution}

    \question 若$z \in \mathbb{C}$满足$z+\dfrac1z=\sqrt 3$,求$z^{2025}+\dfrac{1}{z^{2025}}$的值。
    \begin{solution}
        先求 $z$,
        \[
        z + \frac{1}{z} = \sqrt{3} \quad\Rightarrow\quad z^2 - \sqrt{3} z + 1 = 0 \quad\Rightarrow\quad z = \frac{\sqrt{3}+i}{2} = e^{\frac{\pi i}{6}}
        \]
        于是
        \[
        z^{2025} = e^{\frac{2025 \pi i}{6}} = e^{\frac{3 \pi i}{2}} = -i \quad\Rightarrow\quad z^{2025} + \frac{1}{z^{2025}} = -i + \frac{1}{-i} = 0
        \]
    \end{solution}

    \question 求
    \[
    \left(\frac{\sqrt{3}+1}{2\sqrt{2}} + \frac{\sqrt{3}-1}{2\sqrt{2}} i \right)^6
    \]
    \begin{solution}
        观察
        \[
        \left(\frac{\sqrt{3}+1}{2\sqrt{2}} + \frac{\sqrt{3}-1}{2\sqrt{2}} i \right)^2 
        = \frac{\sqrt{3}}{2} + \frac{1}{2} i 
        = \cos 30^\circ + i \sin 30^\circ.
        \]
        由棣莫弗定理,原式为
        \[
        (\cos 30^\circ + i \sin 30^\circ)^3 = \cos 90^\circ + i \sin 90^\circ = i.
        \]
    \end{solution}

    \question 已知 $a,b,c \in \mathbb{C}$, 且 $a+b+c=a^2+b^2+c^2=3, a^3+b^3+c^3=6$, 试求 
    \[
    (a-1)^{2023}+(b-1)^{2023}+(c-1)^{2023}
    \]
    的值。
    \begin{solution}
        由
        \[
        3=3^{2}-2(ab+bc+ca), \quad 6-3abc=3(3-3)
        \]
        解得
        \[
        ab+bc+ca=3, abc=2
        \]
        因此 $a,b,c$ 是方程
        \[
        x^{3}-3x^{2}+3x-2=0
        \]
        即
        \[
        (x-1)^{3}=1
        \]
        的根,因此
        \[
        a-1,b-1,c-1 \in \{1,\ \omega,\ \omega^{2}\}, \omega=e^{i\frac{2\pi}{3}}
        \]
        不失一般性,
        \[
        (a-1)^{2023}+(b-1)^{2023}+(c-1)^{2023}
        =1^{2023}+\omega^{2023}+(\omega^{2})^{2023}=1+\omega+\omega^{2}=0
        \]
    \end{solution}

    \question 证明
    \[
    1 + 2i + 3i^2 + 4i^3 + \dots + (4n+1)i^{4n}=(2n+1) - 2n i, \quad n \in \mathbb{N}.
    \]
    \begin{solution}
        设
        \[
        S = 1 + 2i + 3i^2 + 4i^3 + \dots + (4n+1)i^{4n}
        \]
        则
        \[
        iS = i + 2i^2 + 3i^3 + 4i^4 + \dots + (4n)i^{4n} + (4n+1)i^{4n+1}
        \]
        两式相减得
        \begin{align*}
        (1-i)S &= 1 + i + i^2 + i^3 + \dots + i^{4n} - (4n+1)i^{4n+1} \\
        &= \frac{i^{4n+1}-1}{i-1} - (4n+1)i = 1 - (4n+1)i
        \end{align*}
        故
        \[
        S = \frac{(1 - (4n+1)i)(1+i)}{(1-i)(1+i)} = (2n+1) - 2n i
        \]
    \end{solution}

    \question 考虑复平面作图 \(z=\sqrt{2}(1+i),\,w=\sqrt{3}-i\),证
    \[
    \tan\frac{\pi}{24}=\sqrt{6}-\sqrt{3}+\sqrt{2}-2
    \]
    \ifprintanswers
    \begin{figure}[H]
    \centering
    \includegraphics[width=0.4\textwidth]{images/image14.png}
    \end{figure}
    \fi
    \begin{solution} 
        如图作$z=\sqrt{2}(1+i),\,w=\sqrt{3}-i,$则$$\arg(z)=\arctan \frac{\sqrt2}{\sqrt2}=\dfrac{\pi}{4},\arg(w)=\arctan \frac{\sqrt3}{-1}=-\dfrac{\pi}{6}$$
        于是 \begin{align*}
        \arg(z+w)&=\angle QOR-|\arg w|\\
        &=\frac{1}{2}\angle POR-|\arg w|\\
        &=\frac{1}{2}(\angle POQ+|\arg w|)-|\arg w|\\
        &=\frac{1}{2}(\frac{\pi}{4}+\frac{\pi}{6})-\frac{\pi}{6}=\frac{\pi}{24}
        \end{align*}
        又$z+w=\sqrt2+\sqrt3+(\sqrt2-1)i$,所以有
        \[
        \tan\frac{\pi}{24}=\frac{\sqrt2-1}{\sqrt2+\sqrt3}=\frac{\sqrt2-1}{\sqrt2+\sqrt3}\cdot\frac{\sqrt3-\sqrt2}{\sqrt3-\sqrt2}=\sqrt{6}-\sqrt{3}+\sqrt{2}-2
        \]
    \end{solution}

    \question 复平面上有三点 $P, Q, R$ 对应三个复数 $z_1, z_2, z_3$,且 $|z_1| = \sqrt{2},|z_2| = \sqrt{5},|z_3| = 3$。若原点 $O$ 为 $\triangle PQR$ 的重心,求 $\Re(\overline{z_1} z_2)$ 。
    \begin{solution}
        设
        \[
        z_1 = a_1 + b_1i, \quad z_2 = a_2 + b_2i, \quad z_3 = a_3 + b_3i
        \]
        由于点 $P, Q, R$ 的重心在原点,故有
        \[
        a_1 + a_2 + a_3 = 0, \quad b_1 + b_2 + b_3 = 0
        \]
        且$a_1^2 + b_1^2 = 2,a_2^2 + b_2^2 = 5,a_3^2 + b_3^2 = 9$,又因为$a_3 = -a_1 - a_2, \ b_3 = -b_1 - b_2,$
        \begin{align*}
        a_3^2 + b_3^2 &= (a_1 + a_2)^2 + (b_1 + b_2)^2\\
        &= a_1^2 + a_2^2 + 2a_1a_2 + b_1^2 + b_2^2 + 2b_1b_2\\
        &= 7 + 2(a_1a_2 + b_1b_2)=9
        \end{align*}
        故
        \[
        \Re(\overline{z_1}z_2)= a_1a_2 + b_1b_2=1
        \]
    \end{solution}

    \question 设 $\omega$ 为复数,且 $|\omega| = 5$。存在一正实数 $\lambda > 1$,使得 $\omega, \omega^2, \lambda\omega$ 这三个复数在复平面上构成一个正三角形,试求 $\lambda$ 的值。
    \ifprintanswers
        \begin{figure}[H]
        \centering
        \includegraphics[width=0.4\textwidth]{images/image19.png}
        \end{figure}
    \fi
    \begin{solution}
        设$O(0, 0),A(\omega),B(\omega^2),C(\lambda \omega)$,则
        \[
        \begin{aligned}
        OA &= |\omega| = 5 \\
        OB &= |\omega^2| = |\omega|^2 = 25 \\
        AB &= AC = |\lambda\omega - \omega| = 5(\lambda - 1) \\
        \end{aligned}
        \]
        又由于 $\triangle ABC$ 是正三角形,因此 
        \[
        \angle OAB = 180^\circ - 60^\circ = 120^\circ
        \]        
        在 $\triangle OAB$中,由余弦定理,
        \[
        \cos \angle OAB = \frac{OA^2 + AB^2 - OB^2}{2 \cdot OA \cdot AB}
        = \frac{5^2 + (5(\lambda - 1))^2 - 25^2}{2 \cdot 5 \cdot 5(\lambda - 1)}=-\frac{1}{2}
        \]
        解得
        \[
        \lambda = \frac{1 + \sqrt{97}}{2}>1
        \]
    \end{solution}

    \question 如图,在坐标平面上有一个半径为 2 的圆,其圆心 $O$ 为原点,且正七边形 $ABCDEFG$ 内接于此圆。若 $A(2,0),P(1,1)$,求 $\overline{PA}\cdot \overline{PB}\cdot\overline{PC}\cdot\overline{PD}\cdot\overline{PE}\cdot\overline{PF}\cdot\overline{PG}$。
    \begin{figure}[H]
        \centering
        \includegraphics[width=0.4\textwidth]{images/image30.png}
    \end{figure}
    \begin{solution}
        设$x^7=2^7$ 的七根分别为 $2,\omega,\omega^2,\dots,\omega^6,$其中 $\omega =2e^{\frac{2\pi i}{7}}$,则
        \[ 
        f(x)=x^7-2^7=(x-2)(x-\omega)(x-\omega^2)\cdots(x-\omega^6)
        \]
        令$x=1+i$,可得
        \[
        \overline{PA}\cdot \overline{PB}\cdot\overline{PC}\cdot\overline{PD}\cdot\overline{PE}\cdot\overline{PF}\cdot\overline{PG}=|f(1+i)|=|(1+i)^7-2^7|=|2i(1+i)-128| =8\sqrt{226}
        \]
    \end{solution}

    \question 已知 $\omega=\cos\dfrac{2\pi}{9}+i\sin\dfrac{2\pi}{9}$,求 $|2-\omega|^2+|2-\omega^2|^2+\cdots+|2-\omega^8|^2$。
    \ifprintanswers
    \begin{figure}[H]
        \centering
        \includegraphics[width=0.5\linewidth]{images/image77.png}
    \end{figure}
    \fi
    \begin{solution}
        即求点 $P(2,0)$ 至单位圆上正九边形各顶点距离的平方和:
        \[
        |2-\omega|^2+|2-\omega^2|^2+\cdots+|2-\omega^8|^2
        =\sum_{n=1}^8 A_nP^2
        \]
        在$\triangle OA_nP$中,由余弦定理,
        \[
        A_nP^2=1^2+2^2-2\cdot 2\cdot 1\cdot \cos\angle A_nOP=5-4\cos(40^\circ\cdot n)
        \]
        于是
        \begin{align*}
        \sum_{n=1}^8 A_nP^2
        &=\sum_{n=1}^8\left(5-4\cos(40^\circ\cdot n)\right)\\
        &=40-4(\cos40^\circ+\cos80^\circ+\cdots+\cos320^\circ)\\
        &=40-8(\cos40^\circ+\cos80^\circ+\cos120^\circ+\cos160^\circ)\\
        &=40-8\left(2\cos60^\circ\cos20^\circ-\frac12+\cos160^\circ\right)\\
        &=40-8\left(\cos20^\circ+\cos160^\circ-\frac12\right)\\
        &=40-8\left(2\cos90^\circ\cos70^\circ-\frac12\right)\\
        &=40-8\left(-\frac12\right)=44
        \end{align*}
    \end{solution}

    \question $\omega^{503}=1,\ \omega \neq 1$,求
        $$
        \frac{\omega^2}{\omega-1}+\frac{\omega^4}{\omega^2-1}+\frac{\omega^6}{\omega^3-1}+\cdots+\frac{\omega^{1004}}{\omega^{502}-1}
        $$

        \begin{solution}
        由 $\omega^{503}=1$ 得
        \[
        (1-\omega)(1+\omega+\omega^2+\cdots+\omega^{502})=0
        \]
        所以
        \[
        \sum_{k=0}^{502}\omega^k=0
        \]
        又
        \[
        \frac{1}{\omega^k-1}+\frac{1}{\omega^{503-k}-1}
        =\frac{1}{\omega^k-1}+\frac{1}{\omega^{-k}-1}
        =\frac{1}{\omega^k-1}+\frac{\omega^k}{1-\omega^k}=-1
        \]
        原式
        \[
        =\sum_{k=1}^{502}\frac{\omega^{2k}}{\omega^k-1}
        =\sum_{k=1}^{502}\left( \frac{\omega^{2k}-1}{\omega^k-1}+\frac{1}{\omega^k-1} \right)
        =\sum_{k=1}^{502} \left( \omega^k+1+\frac{1}{\omega^k-1} \right)
        \]
        \[
        =-1+502+\sum_{k=1}^{502}\frac{1}{\omega^k-1}
        \]
        \[
        =501+\left(\frac{1}{\omega-1}+\frac{1}{\omega^{502}-1}\right)
        +\left(\frac{1}{\omega^2-1}+\frac{1}{\omega^{501}-1}\right)
        +\cdots
        +\left(\frac{1}{\omega^{251}-1}+\frac{1}{\omega^{252}-1}\right)
        \]
        \[
        =501+(-1)\times 251=250
        \]
    \end{solution}

    \question 证明
        \[
        |z+w|^2-|z+\overline{w}|^2=4\Re(z)\Re(w)
        \]
    \begin{solution}
        设$z = x + iy,\;w = u + iv,$其中$x,y,u,i\in \mathbb{R}$,则
        \begin{align*}
        |z+w|^2 - |z-\overline{w}|^2 
        &= |(x+iy) + (u+iv)|^2 - |(x+iy) - (u-iv)|^2 \\
        &= |(x+u) + i(y+v)|^2 - |(x-u) + i(y+v)|^2 \\
        &= (x+u)^2 + (y+v)^2 - (x-u)^2 - (y+v)^2 \\
        &= 4xu \\
        &= 4 \Re(z) \Re(w)
        \end{align*}
    \end{solution}
    \begin{solution}
        又解:
        \begin{align*} 
        |z + w|^2 - |z - w|^2 
        &= (z + w)(\overline{z+w}) - (z - \overline{w})(\overline{z - \overline{w}}) \\
        &= (z + w)(\overline{z}+\overline{w}) - (z - \overline{w})(\overline{z} - w)\\
        &= z\overline{z} + z\overline{w} + w\overline{z} + w\overline{w} - \left( z\overline{z} - zw - \overline{w}\:\overline{z} + w\overline{w} \right) \\
        &= zw+z\overline{w}+w\overline{z}+\overline{w}\:\overline{z}\\
        &=(z+\overline{z})(w+\overline{w})\\
        &=2\Re(z)\cdot 2\Re(w)\\
        &=4 \Re(z) \Re(w)
        \end{align*}
    \end{solution}    
    
    \question 已知 \(z=\cos\theta+i\sin\theta,\ 0\leq\theta<2\pi\),证明:
    \[
    \frac{2}{1+z}=1-i\tan\frac{\theta}{2}
    \]
    \begin{solution}
        运用各位三角恒等式化简
        \begin{align*}
        \frac{2}{1+z} 
        &= \frac{2}{1+\cos\theta+i\sin\theta} \\
        &= \frac{2(1+\cos\theta - i\sin\theta)}{(1+\cos\theta)^2 + \sin^2\theta} \\
        &= \frac{2(1+\cos\theta - i\sin\theta)}{1 + 2\cos\theta + \cos^2\theta + \sin^2\theta} \\
        &= \frac{2(1+\cos\theta - i\sin\theta)}{2 + 2\cos\theta} \\
        &= \frac{1+\cos\theta}{1+\cos\theta} - i\cdot\frac{\sin\theta}{1+\cos\theta} \\
        &= 1 - i \cdot \frac{2\sin\frac{\theta}{2}\cos\frac{\theta}{2}}{2\cos^2\frac{\theta}{2}} \\
        &= 1 - i \tan\frac{\theta}{2}
        \end{align*}
    \end{solution}  
    
    % \question 若 \(m,n\in\mathbb{N}\),证明方程 \(z^m=2-i\) 与 \(z^n=1+i\) 无公共复根。  
    % \begin{solution}
    %     设$m,n\in\mathbb{N}$,假设存在复数 \(z\)使得
    %     \[
    %     z^m = 2 - i,\quad z^n = 1 + i
    %     \]
    %     取模得到
    %     \[
    %     |z^m| = |2 - i| = \sqrt{2^2 + 1^2} = \sqrt{5},\quad |z^n| = |1 + i| = \sqrt{1^2 + 1^2} = \sqrt{2}
    %     \]
    %     因此,
    %     \[
    %     |z|^m = \sqrt{5} \Rightarrow |z| = 5^{\frac{1}{2m}},\quad |z|^n = \sqrt{2} \Rightarrow |z| = 2^{\frac{1}{2n}}
    %     \]
    %     于是,
    %     \[
    %     5^{\frac{1}{2m}} = 2^{\frac{1}{2n}}
    %     \Rightarrow \left(5^{\frac{1}{2m}}\right)^{2mn} = \left(2^{\frac{1}{2n}}\right)^{2mn}
    %     \Rightarrow 5^n = 2^m
    %     \]
    %     左边是奇数,右边是偶数,矛盾!
    % \end{solution}

    \question 设复数 $z$ 满足 $|z|=1$ 且 $z = e^{i\theta}$,求证
    \[
    \Re \left[ \frac{z(1-\bar{z})}{\bar{z}(1+z)} \right]=-2 \sin^2 \frac{\theta}{2}
    \]
    \begin{solution}
        先将表达式化简,
        \begin{align*}
        \frac{z(1-\bar{z})}{\bar{z}(1+z)} 
        &= \frac{z - |z|^2}{\bar{z} + |z|^2}  \\
        &= \frac{e^{i\theta} - 1}{e^{-i\theta} + 1} \\
        &=  \frac{(e^{i\theta}-1)(e^{i\theta}+1)}{(e^{-i\theta}+1)(e^{i\theta}+1)}  \\
        &=  \frac{e^{i2\theta}-1}{2 + e^{i\theta}+e^{-i\theta}} \\
        &=  \frac{e^{i\theta} (e^{i\theta}-e^{-i\theta})}{2 + 2\cos\theta}  \\
        &=  \frac{e^{i\theta} \cdot 2i\sin\theta}{2(1+\cos\theta)} \\
        &=  \frac{e^{i\theta} \cdot i\sin\theta}{1+\cos\theta} 
        \end{align*}
        故
        \begin{align*}
        \Re \left[ \frac{z(1-\bar{z})}{\bar{z}(1+z)} \right]=\Re \left[ \frac{e^{i\theta} \cdot i\sin\theta}{1+\cos\theta} \right] &= \frac{1}{1+\cos\theta} \Re \left[ i \sin\theta (\cos\theta + i\sin\theta) \right] \\
        &= \frac{1}{1+\cos\theta} \Re \left[ i \sin\theta \cos\theta - \sin^2\theta \right] \\
        &= -\frac{\sin^2\theta}{1+\cos\theta} \\
        &= -\frac{(2\sin\frac{\theta}{2}\cos\frac{\theta}{2})^2}{2\cos^2\frac{\theta}{2}} \\
        &= -2 \sin^2 \frac{\theta}{2}
        \end{align*}
    \end{solution}
    
    \question 已知两个相异复数 \(z_1,z_2\),且 \(|z_1|=|z_2|\ne 0\),证明
    \[
    \frac{z_1+z_2}{z_1-z_2}
    \]
    是纯虚数。  
    \begin{solution}
        设$w = \dfrac{z_1 + z_2}{z_1 - z_2}$ ,欲证 \(\overline{w} = -w\),因为 \(|z_1| = |z_2| = r\),且 \(z_1, z_2 \ne 0\),根据性质
        \[
        \overline{z} = \frac{r^2}{z}
        \Rightarrow \overline{z_1} = \frac{r^2}{z_1},\quad \overline{z_2} = \frac{r^2}{z_2}
        \]
        因此
        \begin{align*}
        \overline{w} 
        &= \overline{\left(\frac{z_1 + z_2}{z_1 - z_2}\right)} 
        = \frac{\overline{z_1} + \overline{z_2}}{\overline{z_1} - \overline{z_2}} 
        = \frac{\dfrac{r^2}{z_1} + \dfrac{r^2}{z_2}}{\dfrac{r^2}{z_1} - \dfrac{r^2}{z_2}} 
        = \frac{\dfrac{1}{z_1} + \dfrac{1}{z_2}}{\dfrac{1}{z_1} - \dfrac{1}{z_2}} 
        = \frac{z_2 + z_1}{z_2 - z_1} = -\frac{z_1 + z_2}{z_1 - z_2} = -w
        \end{align*}
    \end{solution}
    
    % \question 复数 \(z\) 满足关系式:
    %     \[
    %     5(z+i)^n = (4+3i)(1+iz)^n, \quad n \in \mathbb{R}
    %     \]
    %     试证 \(z \in \mathbb{R}\)。
        
    % \begin{solution}
    %     原式写成
    %     \[
    %     5(z+i)^n = i^n (4+3i)(z - i)^n
    %     \]
    %     两边取模有
    %     \begin{align*}
    %     5 |z+i|^n &= |i^n| \cdot |4+3i| \cdot |z-i|^n \\
    %     5 |z+i|^n &= 1 \cdot \sqrt{4^2 + 3^2} \cdot |z - i|^n = 5 |z - i|^n\\
    %     |z+i|^n &= |z-i|^n
    %     \end{align*}
    %     设 \(z = x + iy\),代入得
    %     \begin{align*}
    %     |x + i(y+1)|^n &= |x + i(y-1)|^n \\
    %     \left(\sqrt{x^2 + (y+1)^2}\right)^n &= \left(\sqrt{x^2 + (y-1)^2}\right)^n
    %     \end{align*}
    %     故
    %     \[
    %     x^2 + (y+1)^2 = x^2 + (y-1)^2 \Rightarrow y = 0
    %     \]
    %     $\therefore \;z \in \mathbb{R}$
    % \end{solution}
    
    % \question 已知复数 \(z, w\) 满足 \(|z| = |w| = 1\),证明:
    %     \[
    %     \frac{z + w}{1 + zw} \in \mathbb{R}
    %     \]
        
    % \begin{solution}
    %     设
    %     \[
    %     v = \frac{z + w}{1 + zw}
    %     \]
    %     即证$v = \overline{v}$:
    %     \begin{align*}
    %     \overline{v} &= \overline{\left(\frac{z + w}{1 + zw}\right)} = \frac{\overline{z} + \overline{w}}{1 + \overline{z} \cdot \overline{w}}
    %     \end{align*}
        
    %     由于 \[|z| = |w| = 1 \Rightarrow \overline{z} = \frac{1}{z},\; \overline{w} = \frac{1}{w}
    %     \]
    %     则
    %     \begin{align*}
    %     \overline{v} &= \frac{\frac{1}{z} + \frac{1}{w}}{1 + \frac{1}{zw}} 
    %     = \frac{w + z}{zw + 1} = v
    %     \end{align*}
        
    %     \(\therefore v \in \mathbb{R}\)。
        
    % \end{solution}

    % \question 若 \(z,w\in\mathbb{C},\,|w|>1,\,z+w=|z|w\),试以 \(|w|\) 表示 \(|z|\)。  

    % \begin{solution}
    %     由题意
    %     \[
    %     z + w = |z|w \Rightarrow z = (|z| - 1)w
    %     \]
    %     设 \(w = a + bi,\)
    %     \[
    %     z = (|z| - 1)(a + bi) = a(|z| - 1) + b(|z| - 1)i
    %     \]
    %     取模得
    %     \[
    %     |z| = \left|a(|z| - 1) + b(|z| - 1)i\right| = \sqrt{a^2(|z| - 1)^2 + b^2(|z| - 1)^2}
    %     \]
    %     故
    %     \[
    %     |z| = \pm |w|(|z| - 1)
    %     \Rightarrow |z| = \pm |w||z| \mp |w|
    %     \]
    %     即
    %     \[
    %     |z| = \frac{\mp |w|}{1 \mp |w|}
    %     \]
    % \end{solution}

    \question 设复数 $z_1, z_2$ 满足 $|z_1|=2,|z_2|=3,|z_1+z_2|=4$, 求 $\dfrac{z_1}{z_2}$。 
    \begin{solution}
        设 $w = \dfrac{z_1}{z_2}$,则 $z_1 = w z_2$。由
        \[
        |z_1 + z_2|^2 = (z_1 + z_2)(\overline{z_1} + \overline{z_2}) = |z_1|^2 + |z_2|^2 + z_1\overline{z_2} + z_2\overline{z_1}
        \]
        代入已知得
        \[
        z_1\overline{z_2} + z_2\overline{z_1} = 3
        \]
        将 $z_1 = w z_2$ 代入得
        \[
        wz_2 \cdot \overline{z_2} + z_2 \cdot \overline{w z_2} = w |z_2|^2 + \overline{w} |z_2|^2 = (w + \overline{w}) \cdot 9=3 \Rightarrow w + \overline{w} = \frac{1}{3}
        \]
        设 $w = x + iy,x,y\in \mathbb{R}$,则 $w + \overline{w} = 2x = \dfrac{1}{3} \Rightarrow x = \dfrac{1}{6}$;又因 $|w| = \dfrac{|z_1|}{|z_2|} = \dfrac{2}{3}$,所以
        \[
        |w|^2 = \frac{1}{36} + y^2 = \frac{4}{9}
        \Rightarrow y = \pm \frac{\sqrt{15}}{6} \Rightarrow   \frac{z_1}{z_2} = \frac{1}{6} \pm i \frac{\sqrt{15}}{6}
        \]
    \end{solution}

    \question 若复数 $z$ 使得$\dfrac{z-3i}{z+i}$ 为负实数, $\dfrac{z-3}{z+1}$ 为纯虚数,求 $z$。
    \begin{solution}
        设 $z=a+bi,a,b\in \mathbb{R}$,由于 
        \[
        \frac{z-3i}{z+i}=\frac{a+(b-3)i}{a+(b+1)i}
        \]
        为负实数,则 
        \[
        (a+(b-3)i)(a-(b+1)i)<0
        \]
        所以 $a^{2}+(b-3)(b+1)<0$ 且 $-4a=0$,即 $a=0$ 且 $b-3,b+1$ 异号; 又此时
        \[
        \frac{z-3}{z+1}=\frac{-3+bi}{1+bi}
        \]
        为纯虚数,故 
        \[
        \Re((-3+bi)(1-bi))=b^{2}-3=0
        \]
        又 $b-3,b+1$异号知 $b=\sqrt{3}$,所以 $z=\sqrt{3}i$.
    \end{solution}
    
    \question  已知虚数 $z$ 使得 $z_1 = \dfrac{z}{1+z^2}$ 和 $z_2 = \dfrac{z^2}{1+z}$ 都是实数,求 $z$。
    \begin{solution}
        由已知$(z^2 + 1)z_1 = z, (1 + z)z_2 = z^2$,联立两式
        \[
        z = (z^2 + 1)z_1 = ((1+z)z_2 + 1)z_1
        \]
        整理得
        \[
        z_1 + z_1 z_2 = z(1 - z_1 z_2)
        \]
        由于 $z$ 为虚数, $z_1 + z_1 z_2$ 与 $1 - z_1 z_2$ 为实数,故
        \[
        1 - z_1 z_2 = 0 \Rightarrow z_1 z_2 = 1
        \]
        代入 $z_1 = \dfrac{z}{1 + z^2},z_2 = \dfrac{z^2}{1 + z}$ 得
        \[
        z_1 z_2 = \frac{z}{1+z^2} \cdot \frac{z^2}{1+z} = \frac{z^3}{(1+z)(1+z^2)} = 1
        \]
        给出
        \[
        z^3 = (1+z)(1+z^2)
        \Rightarrow z^2 + z + 1 = 0
        \]
        即
        \[
        z = -\frac{1}{2} \pm \frac{\sqrt{3}}{2}i
        \]
    \end{solution}
    
    \question 已知复数 \(z\ne 0\),满足方程 \(z^2=z+i|z|\),求 \(|z|\) 的值。
    \begin{solution}
        设 $z = x + iy,\;x,y\in \mathbb{R},$则
        \[
        z^2 = z + i|z| \Rightarrow x^2 - y^2 + 2ixy  = x + i\left(y + \sqrt{x^2 + y^2}\right)
        \]
        比较实部与虚部得
        \[
        x^2 - y^2 = x \tag{1}
        \]
        \[
        2xy = y + \sqrt{x^2 + y^2} \tag{2}
        \]
        将(2)移项并平方
        \[
        y^2(2x - 1)^2 = x^2 + y^2 
        \]
        由(1)得
        \[
        (x^2 - x)(2x - 1)^2 = 2x^2 - x \Rightarrow x^2(2x-3)(2x-1)=0
        \]
        解得 $x=0,\dfrac{1}{2},\dfrac{3}{2}$,其中$x=0$不合理$(z\neq 0)$,且$x=\dfrac{1}{2}$不合理$(y^2=-\dfrac{1}{4}<0)$,故
        \[
        x = \dfrac{3}{2},\; y^2 = \dfrac{3}{4}\Rightarrow|z| = \sqrt{\frac{9}{4} + \frac{3}{4}} = \sqrt3
        \]
    \end{solution}

    \question 确定所有满足 $|a|=|b|=1$ 且 $a + b + a\bar{b} \in \mathbb{R}$ 的复数对 $(a,b)$。
    \begin{solution}
        设 $a = e^{ix},b = e^{iy}$,其中 $x, y \in [0, 2\pi)$,由欧拉公式,
        \begin{align*}
        \Im(a+b+a\bar{b}) &= (\sin x + \sin y) + \sin(x-y) \\
        &= 2 \sin \frac{x+y}{2} \cos \frac{x-y}{2} + 2 \sin \frac{x-y}{2} \cos \frac{x-y}{2} \\
        &= 2 \left(\sin \frac{x+y}{2} + \sin \frac{x-y}{2}\right) \cos \frac{x-y}{2} \\
        &= 4 \sin \frac{x}{2} \cos \frac{y}{2} \cos \frac{x-y}{2}.
        \end{align*}
        因此$a+b+a\bar{b}$ 为实数当且仅当
        \[
        \sin \frac{x}{2} = 0, \quad \cos \frac{y}{2} = 0, \quad \text{或} \quad \cos \frac{x-y}{2} = 0,
        \]
        分别对应
        \[
        x = 2k\pi, \quad y = (2k+1)\pi, \quad x = y + (2k+1)\pi, \quad k \in \mathbb{Z}.
        \]
        所以解集为
        \[
        (a,b) = (1, b), \quad (a,-1), \quad (a,-a),
        \]
        其中$|a|=|b|=1$。
    \end{solution}
    \begin{solution}
        注意到
        \[
        a+b+a\bar{b}\in\mathbb{R} \iff 1+a+b+a\bar{b} \in \mathbb{R}.
        \]
        设 $c \in \mathbb{C}$ 且 $a=c^2$,则
        \[
        \overline{c}(1+a+b+a\bar{b}) = \overline{c} + \overline{c}c^2 + \overline{c}b + \overline{c}c^2\overline{b} = \overline{c}+c+\overline{c}b+c\overline{b} \in \mathbb{R},
        \]
        因此 $c \in \mathbb{R}$或 $1+a+b+a\overline{b}=0$。第一种情况 $c=\pm 1$,于是 $a=1$。第二种情况将等式因式分解为
        \[
        (a+b)(1+\overline{b}) = 1+a+b+a\overline{b}=0,
        \]
        得到 $a=-b$ 或 $b=-1$,所以得到解为
        \[
        (a,b) = (1,b), \quad (a,-1), \quad (a,-a), \quad |a|=|b|=1.
        \]
    \end{solution}

    \question 已知两复数 $z_1, z_2$ 满足 $|z_1 - (3 + 3i)| = 2,|iz_2 - 1| = 1$,求 $|z_1 - z_2|$ 的最小值。
    \begin{solution}
        由$|z_1 - (3 + 3i)| = 2,P(z_1)$在圆
        \[
        C_1 :(x - 3)^2 + (y - 3)^2 = 2^2 
        \]
        上,设$z_2 = a + bi$,则$iz_2 - 1 = (-1 - b) + ai$,故$ |iz_2 - 1| = 1 $意味$Q(z_2)$在圆 
        \[
        C_2 :x^2 + (y + 1)^2 = 1 
        \]
        上,故$|z_1 - z_2|=PQ$ 的最小值为
        \[
        \text{两圆心距离} - \text{两圆半径和} = \sqrt{3^2 + 4^2} - (2 + 1) = 2
        \]
    \end{solution}

    \question 设 $\alpha$ 为 $\left|(z-1)\left(z+\dfrac{1}{2}\right)\right|$ 在圆盘上 $\{z \in \mathbb{C}:|z| \leq 1\}$ 的最大值,则 $\alpha^2=$ ?
    \ifprintanswers
    \begin{figure}[H]
        \centering
        \includegraphics[width=0.5\linewidth]{images/image75.png}
    \end{figure}
    \fi
    \begin{solution}
        如图,$P$ 在圆周上,$A(1,0),B\left(-\dfrac12,0\right)$,则
        \[
        PA \cdot PB = \left|(z-1)\left(z+\frac{1}{2}\right)\right|
        \]
        设$\theta=\angle POA$,在$\triangle POA$及$\triangle POB$中,由余弦定理,
        \[
        PA^2 = 1^2+1^2-2\cdot 1 \cdot 1 \cdot \cos \theta= 2-2\cos\theta
        \]
        \[
        PB^2 = 1^2+ \left(\frac{1}{2}\right)^2 -2\cdot 1 \cdot \left(\frac{1}{2}\right) \cdot \cos(\pi-\theta)=\cos\theta+\frac54
        \]
        因此
        \[
        f(\theta) = (2-2\cos\theta)\left(\cos\theta+\frac54\right) = -2\cos^2\theta - \frac12\cos\theta + \frac52  
        \]
        且
        \[
        \alpha^2 = f\left(-\frac18\right) = \frac{81}{32}
        \]
    \end{solution}

    \question 设有一虚部不为零的复数 $z$, 其长度为 $2$, 且在复数平面上与 $-2$ 及 $z^2$ 刚好在同一直线上, 与 $1$ 及 $z^3$ 也同在另一直线上, 试求以 $z,z^2,z^3$ 所围成的三角形面积。
    \begin{solution}
        由 $1,z,z^3$ 共线可得
        \[
        \frac{z^3-z}{z-1}=k \in \mathbb R \Rightarrow z^2+z-k=0 \Rightarrow z=\frac{-1\pm \sqrt{\,1+4k\,}}{2}
        \]
        其中$z$实部为$-\dfrac{1}{2}$,由 $-2,z,z^2$ 共线可得
        \[
        \frac{z^2-z}{z+2}=t \in \mathbb R
        \Rightarrow z^2-(t+1)z-2t=0
        \Rightarrow z=\frac{t+1\pm \sqrt{(t+1)^2+8t}}{2}
        \]
        因此
        \[
        \frac{t+1}{2}=-\frac{1}{2}\Rightarrow t=-2
        \]
        代入得
        \[
        z=\frac{-1\pm \sqrt{15}\, i}{2}.
        \]
        取$z=\dfrac{-1+\sqrt{15}\, i}{2}$,则
        \[
        z^2=\frac{-7-\sqrt{15}\, i}{2},\quad
        z^3=\frac{11-3\sqrt{15}\, i}{2}.
        \]
        $z,z^2,z^3$在复平面上的坐标为
        \[
        A(z)=\left(-\frac12,\frac{\sqrt{15}}{2}\right),\quad
        B(z^2)=\left(-\frac72,-\frac{\sqrt{15}}{2}\right),\quad
        C(z^3)=\left(\frac{11}{2},-\frac{3\sqrt{15}}{2}\right).
        \]
        故所求面积为
        \[
        \frac{1}{2}
        \begin{vmatrix}
        -\dfrac12 & \dfrac{\sqrt{15}}{2} & 1\\[6pt]
        -\dfrac72 & -\dfrac{\sqrt{15}}{2} & 1\\[6pt]
        \dfrac{11}{2} & -\dfrac{3\sqrt{15}}{2} & 1
        \end{vmatrix}
        =6\sqrt{15}
        \]
    \end{solution}

    \question 若实数 $m,n$ 使得关于 $x$ 的方程 $x^3 + mx + n = 0$ 有模为 $3$ 的虚根,求 $m+n$ 的取值范围。
    \begin{solution}
        由于系数皆为实数,设三根为 $a + bi,\ a - bi,\ -2a$,其中 $a,b \in \mathbb{R},\,a^2 + b^2 = 9,\ b \ne 0$。
        
        由韦达定理,
        \[
        m = (a + bi)(a - bi) + (-2a)(a + bi + a - bi)
        = 9 - 4a^2
        \]
        \[
        n=(a + bi)(a - bi)(-2a)= -18a 
        \]
        所以
        \[
        m + n = 9 - 4a^2 + 18a = -4\left(a - \dfrac{9}{4} \right)^2 + \dfrac{117}{4}
        \]
        又$a^2<9,-3<a<3$,则\[
        m+n \in \left(-81,\ \dfrac{117}{4}\right]
        \]
    \end{solution}

    \question 关于 $z$ 的方程 $z^{n+1} - \sqrt{3}z^n - 1 = 0$ 存在一个模为 $1$ 的虚根,求正整数 $n$ 的最小值。
    \begin{solution}
        设该虚根为 $z_0= e^{i\theta} = \cos\theta + i\sin\theta$,有        
        \[
        z_0^n(z_0 - \sqrt{3}) = 1
        \]
        两边取模得
        \[
         |z_0^n||z_0 - \sqrt{3}| = 1
        \]
        由于 $|z_0^n| = 1$,从而
        \[
        |z_0 - \sqrt{3}| = 1
        \]
        将$z_0= e^{i\theta} = \cos\theta + i\sin\theta$代入得
        \[
        |(\cos\theta - \sqrt{3}) + i\sin\theta| = 1
        \]
        于是
        \[
        (\cos\theta - \sqrt{3})^2 + \sin^2\theta = 1 \Rightarrow \cos\theta = \frac{\sqrt{3}}{2}
        \]
        得 $\theta = \dfrac{\pi}{6}$,即 $z_0 = e^{i\frac{\pi}{6}}$,代入 $z_0^n(z_0 - \sqrt{3}) = 1$ 得:
        \[
        \left(e^{i\frac{\pi}{6}}\right)^n \left(e^{i\frac{\pi}{6}} - \sqrt{3}\right) =
        e^{i\frac{n\pi}{6}} \left( \frac{\sqrt{3}}{2} + \frac{1}{2}i - \sqrt{3} \right) = e^{i\frac{n\pi}{6}} \left( -\frac{\sqrt{3}}{2} + \frac{1}{2}i \right) = 1
        \]
        注意到$-\frac{\sqrt{3}}{2} + \frac{1}{2}i = e^{i\frac{5\pi}{6}},$变为
        \[
        e^{i\frac{n\pi}{6}} \cdot e^{i\frac{5\pi}{6}}= e^{i\frac{(n+5)\pi}{6}}= 1 \Rightarrow \frac{(n+5)\pi}{6} = 2k\pi \Rightarrow n = 12k - 5
        \]
        令 $k=1$ 得最小正整数解为 $n = 7$.
    \end{solution}

    \question 已知复数 \(z_1, z_2, z_3\) 满足
    \[
    \begin{cases}
    |z_1| = |z_2| = |z_3| = 1 \\[1mm]
    \displaystyle \frac{z_1}{z_2} + \frac{z_2}{z_3} + \frac{z_3}{z_1} = 1
    \end{cases}
    \]
    求 \(|z_1 + 2 z_2 + 3 z_3|\) 最大可能值。
    \begin{solution}
        设存在一三次多项式,根为 \(\dfrac{z_1}{z_2}, \dfrac{z_2}{z_3}, \dfrac{z_3}{z_1}\),则
        \[
        \frac{z_1}{z_2} + \frac{z_2}{z_3} + \frac{z_3}{z_1} = 1 
        \]
        \[
        \frac{z_1}{z_2} \cdot \frac{z_2}{z_3} \cdot \frac{z_3}{z_1} = 1
        \]
        \begin{align*}
        &\frac{z_1}{z_2} \cdot \frac{z_2}{z_3} + \frac{z_2}{z_3} \cdot \frac{z_3}{z_1} + \frac{z_3}{z_1} \cdot \frac{z_1}{z_2} = \frac{z_1}{z_3} + \frac{z_2}{z_1} + \frac{z_3}{z_2} \\
        &= \overline{\left(\frac{z_3}{z_1}\right)} + \overline{\left(\frac{z_1}{z_2}\right)} + \overline{\left(\frac{z_2}{z_3}\right)} = \overline{\left(\frac{z_1}{z_2} + \frac{z_2}{z_3} + \frac{z_3}{z_1}\right)} = 1
        \end{align*}
        因此该多项式为
        \[
        f(x) = x^3 - x^2 + x - 1 = (x^2 + 1)(x - 1),
        \]
        根为 \(x = 1, \pm i\);为使 \(|z_1 + 2 z_2 + 3 z_3|\) 最大,取
        \[
        \frac{z_2}{z_3} = 1, \quad \frac{z_3}{z_1} = i, \quad \frac{z_1}{z_2} = -i,
        \]
        则
        \[
        z_2 = z_3 = 1, \quad z_1 = -i,
        \]
        所以
        \[
        |z_1 + 2 z_2 + 3 z_3| = |-i + 5| = \sqrt{26}
        \]
    \end{solution}

    \question 设 $z$ 为复数,且满足 $|z+1|>2$。证明
    \[
    |z^3+1|>1.
    \]
    \begin{solution}
        注意到
        \[
        z^3+1=(z+1)(z^2-z+1),
        \]
        因此只需证明
        \[
        |z^2-z+1|> \frac{1}{2}.
        \]
        设$z+1=re^{i\varphi}$,其中 $r=|z+1|>2$,$\varphi$ 为某个实数。于是
        \[
        z^2-z+1=(re^{i\varphi}-1)^2-(re^{i\varphi}-1)+1 =r^2e^{2i\varphi}-3re^{i\varphi}+3
        \]
        于是
        \begin{align*}
        |z^2-z+1|^2
        &=(r^2e^{2i\varphi}-3re^{i\varphi}+3)
        (r^2e^{-2i\varphi}-3re^{-i\varphi}+3) \\
        &=r^4+9r^2+9-(6r^3+18r)\cos\varphi+6r^2\cos 2\varphi \\
        &=r^4+9r^2+9-(6r^3+18r)\cos\varphi+6r^2(2\cos^2\varphi-1) \\
        &=12\left(r\cos\varphi-\frac{r^2+3}{4}\right)^2
        +\frac{1}{4}(r^2-3)^2.
        \end{align*}
        由于 $r>2$,可得
        \[
        |z^2-z+1|^2>\frac{1}{4}\Rightarrow |z^2-z+1|>\frac{1}{2}
        \]
        因此
        \[
        |z^3+1|=|z+1|\,|z^2-z+1|>2\cdot\frac{1}{2}=1
        \]
        证毕。
    \end{solution}

    \question 复数 $z_1,z_2,z_3,z_4,z_5$ 满足
    \begin{equation*}
    \begin{cases}
    \;|z_1| \leq 1 \\
    \;|z_2| \leq 1 \\
    \;|2z_3 - (z_1 + z_2)| \leq |z_1 - z_2| \\
    \;|2z_4 - (z_1 + z_2)| \leq |z_1 - z_2| \\
    \;|2z_5 - (z_3 + z_4)| \leq |z_3 - z_4|
    \end{cases}
    \end{equation*}
    求 $|z_5|$ 的最大值。
    \begin{solution}
        由\[
        |z_1 - z_2| \geq |2z_3 - (z_1 + z_2)| \geq |2|z_3| - |z_1 + z_2||
        \] 
        我们有
        \[
        |z_1 + z_2| - |z_1 - z_2| \leq 2|z_3| \leq |z_1 + z_2| + |z_1 - z_2|
        \] 
        由AM-QM不等式,
        \[
        |z_3| \leq \frac{|z_1 + z_2| + |z_1 - z_2|}{2} \leq \sqrt{\frac{|z_1 + z_2|^2 + |z_1 - z_2|^2}{2}} = \sqrt{\frac{2|z_1|^2 + 2|z_2|^2}{2}} \leq \sqrt{1^2 + 1^2} = \sqrt{2}
        \] 
        同理可得 $|z_4| \leq \sqrt{2}$,故
        $$ |z_5| \leq \frac{|z_3 + z_4|}{2} \leq \frac{|z_3| + |z_4|}{2} \leq \frac{\sqrt{2} + \sqrt{2}}{2} = \sqrt{2}$$ 
        考虑等号何时成立:当$z_1=1, z_2=i$时, \(z_{3}=1+i\);当$z_1=-1,z_2=i$时, \(z_{4}=-1+i\); 当$z_3=1+i, z_4=-1+i$时, \(z_{5}=2i\),这时 \(|z_{5}|=2\). 故$|z_5|_{\max}=2$.
    \end{solution}
    
    \question 设复数 $z$ 满足 $|z|=1$,则 $|z^7 + \bar{z}^5 - 3z^3 - 3\bar{z}|$ 的最大值为
    \begin{solution}
        试 $z = \frac{1+i}{\sqrt{2}}$,算得值为:
        \[
        |z^7 + \bar{z}^5 - 3z^3 - 3\bar{z}| = |2i + 2 - 6i - 6| = |-4i - 4| = \sqrt{16 + 16} = \boxed{4\sqrt{2}}
        \]
        \textcolor{red}{(待解)(设$z=\cos \theta + i \sin \theta,...$感觉是尝试因式分解再AM-GM)}

    \end{solution}

    \question 
    \begin{parts}
    \part 已知 $z = \cos \theta + i \sin \theta$,证明
    \[
    z^n + z^{-n} = 2\cos n\theta, \quad z^n - z^{-n} = 2i \sin n\theta.
    \]
    \begin{solution}
        由棣莫弗定理,
        \[
        z^n = \cos n\theta + i \sin n\theta, \quad z^{-n} = \cos n\theta - i \sin n\theta
        \]
        相加得到
        \[
        z^n + z^{-n} = (\cos n\theta + i \sin n\theta) + (\cos n\theta - i \sin n\theta) = 2\cos n\theta
        \]
        而相减得到
        \[
        z^n - z^{-n} = (\cos n\theta + i \sin n\theta) - (\cos n\theta - i \sin n\theta) = 2i \sin n\theta
        \]
    \end{solution}
    \part 证明
    \[
    16 \sin^5\theta = \sin 5\theta - 5 \sin 3\theta + 10 \sin \theta,
    \]
    且
    \[
    32 \cos^6\theta = \cos 6\theta + 6 \cos 4\theta + 15 \cos 2\theta + 10.
    \]
    \begin{solution}
        令$n=1$,则 
        \[
        2i \sin \theta = z - z^{-1}, \quad 2 \cos \theta = z + z^{-1}
        \]
        由第一式,二项式展开得
        \begin{align*}
        (2i \sin \theta)^5 &= (z - z^{-1})^5 \\
        32 i^5 \sin^5 \theta &= z^5 - 5 z^3 z^{-2} + 10 z z^{-4} - 10 z^2 z^{-3} + 5 z z^{-4} - z^{-5} \\
        32 i \sin^5 \theta &= (z^5 - z^{-5}) - 5(z^3 - z^{-3}) + 10 (z - z^{-1}) \\
        32 i \sin^5 \theta &= 2i \sin 5\theta - 10 i \sin 3\theta + 20 i \sin \theta \\
        16 \sin^5 \theta &= \sin 5\theta - 5 \sin 3\theta + 10 \sin \theta
        \end{align*}
        由第二式,
        \begin{align*}
        (2 \cos \theta)^6 &= (z + z^{-1})^6 \\
        64 \cos^6 \theta &= (z^6 + z^{-6}) + 6(z^4 + z^{-4}) + 15 (z^2 + z^{-2}) + 20 \\
        64 \cos^6 \theta &= 2 \cos 6\theta + 12 \cos 4\theta + 30 \cos 2\theta + 20 \\
        32 \cos^6 \theta &= \cos 6\theta + 6 \cos 4\theta + 15 \cos 2\theta + 10
        \end{align*}
        故得证。
    \end{solution}
    \end{parts}

    \question 证明
    \[
    \cos^5\theta \sin^3\theta = \frac{1}{128}\Bigl(6\sin 2\theta + 2\sin 4\theta - 2\sin 6\theta - \sin 8\theta\Bigr)
    \]
    \begin{solution}
        设 $z = \cos\theta + i \sin\theta$,由棣莫弗定理,
        \[
        z^n = \cos n\theta + i\sin n\theta, \quad z^{-n} = \cos n\theta - i\sin n\theta,
        \]
        且
        \[
        z^n + z^{-n} = 2\cos n\theta, \quad z^n - z^{-n} = 2i \sin n\theta.
        \]
        由 $2\cos\theta = z + z^{-1}$ 得
        \begin{align*}
        (2\cos\theta)^5 &= (z + z^{-1})^5 \\
        32 \cos^5\theta &= z^5 + 5 z^4 z^{-1} + 10 z^3 z^{-2} + 10 z^2 z^{-3} + 5 z z^{-4} + z^{-5} \\
        32 \cos^5\theta &= (z^5 + z^{-5}) + 5(z^3 + z^{-3}) + 10(z + z^{-1}) \\
        16 \cos^5\theta &= \cos 5\theta + 5\cos 3\theta + 10\cos\theta
        \end{align*}
        由 $4\sin^3\theta = 3\sin\theta - \sin 3\theta$,得到
        \[
        16\cos^5\theta \cdot 4 \sin^3\theta = (\cos 5\theta + 5\cos 3\theta + 10\cos\theta)(3\sin\theta - \sin 3\theta).
        \]
        展开并和差化积,
        \begin{align*}
        64 \cos^5\theta \sin^3\theta &= 3\cos 5\theta \sin\theta - \cos 5\theta \sin 3\theta + 15 \cos 3\theta \sin\theta - 5 \cos 3\theta \sin 3\theta \\
        &\quad + 30 \cos\theta \sin\theta - 10 \cos\theta \sin 3\theta \\
        &= \frac{3}{2}(\sin 6\theta - \sin 4\theta) - \frac{1}{2}(\sin 8\theta - \sin 2\theta) + \frac{15}{2}(\sin 4\theta - \sin 2\theta) \\
        &\quad - \frac{5}{2} \sin 6\theta + \frac{30}{2} \sin 2\theta - \frac{10}{2} \sin 4\theta \\
        128 \cos^5\theta \sin^3\theta &= 6 \sin 2\theta + 2 \sin 4\theta - 2 \sin 6\theta - \sin 8\theta
        \end{align*}
        因此得证
        \[
        \cos^5\theta \sin^3\theta = \frac{1}{128}\left(6\sin 2\theta + 2\sin 4\theta - 2\sin 6\theta - \sin 8\theta\right)
        \]
    \end{solution}

    \question 定义无穷级数
        \[
        C=\cos\theta+\frac{1}{2}\cos5\theta+\frac{1}{4}\cos9\theta+\cdots
        \]
        \[
        S=\sin\theta+\frac{1}{2}\sin5\theta+\frac{1}{4}\sin9\theta+\cdots
        \]
        证明
        \[
        C+iS = \frac{2e^{i\theta}}{2 - e^{4i\theta}}, \quad 
        S = \frac{4\sin\theta + 2\sin3\theta}{5 - 4\cos4\theta}
        \]
    \begin{solution}
        记
        \begin{align*}
        C+iS &= \cos\theta + i\sin\theta + \frac{1}{2}(\cos5\theta + i\sin5\theta) + \frac{1}{4}(\cos9\theta + i\sin9\theta) + \cdots\\
        &= e^{i\theta} + \frac{1}{2}e^{i5\theta} + \frac{1}{4}e^{i9\theta} + \cdots
        = \sum_{k=0}^{\infty} \frac{1}{2^k} e^{i(4k+1)\theta}
        = e^{i\theta} \sum_{k=0}^{\infty} \left(\frac{e^{4i\theta}}{2}\right)^k        \end{align*}
        这是一个公比为 \(\dfrac{e^{4i\theta}}{2}\)的等比级数,其模为 \(\dfrac{1}{2}<1\),因此级数收敛:
        \[
        C+iS = \frac{e^{i\theta}}{1 - \frac{1}{2}e^{4i\theta}} = \frac{2e^{i\theta}}{2 - e^{4i\theta}}
        \]
        又
        \[
        C+iS =\frac{2e^{i\theta}}{2 - e^{4i\theta}} \cdot \frac{2 - e^{-4i\theta}}{2 - e^{-4i\theta}} = \frac{4e^{i\theta} - 2e^{-3i\theta}}{|2 - e^{4i\theta}|^2}=\frac{(4\cos\theta - 2\cos3\theta) + i(4\sin\theta + 2\sin3\theta)}{(2 - \cos4\theta)^2 + \sin^2 4\theta}
        \]
        因此
        \[
        S = \Im(C + iS) = \frac{4\sin\theta + 2\sin3\theta}{5 - 4\cos4\theta}
        \]
    \end{solution}

        \question 证明
    \begin{align*}
    \cos \alpha + \comb{n}{1} \cos 2\alpha + \comb{n}{2} \cos 3\alpha + \dots + \cos (n + 1)\alpha &= 2^n \cos^n \frac{\alpha}{2} \cos \frac{n+2}{2} \alpha \\
    \sin \alpha + \comb{n}{1} \sin 2\alpha + \comb{n}{2} \sin 3\alpha + \dots + \sin (n + 1)\alpha &= 2^n \cos^n \frac{\alpha}{2} \sin \frac{n+2}{2} \alpha
    \end{align*}
    \begin{solution}
        考虑下列复数和的实部与虚部系数
        \[ 
        (\cos \alpha + i \sin \alpha) + \comb{n}{1}(\cos 2\alpha + i \sin 2\alpha) + \dots + [\cos (n + 1)\alpha + i \sin (n + 1)\alpha] 
        \]
        令 
        \[
        x = \cos \alpha + i \sin \alpha
        \]
        由棣莫弗定理及二项式定理,
        \begin{align*}
        x + \comb{n}{1} x^2 + \comb{n}{2} x^3 + \dots + x^{n+1} &= x(1 + x)^n \\
        &= (\cos \alpha + i \sin \alpha)(1 + \cos \alpha + i \sin \alpha)^n \\
        &= (\cos \alpha + i \sin \alpha)\left(2\cos^2 \frac{\alpha}{2} + 2i \sin \frac{\alpha}{2} \cos \frac{\alpha}{2}\right)^n \\
        &= (\cos \alpha + i \sin \alpha) \left[ 2\cos \frac{\alpha}{2} \left(\cos \frac{\alpha}{2} + i \sin \frac{\alpha}{2}\right) \right]^n \\
        &= 2^n \cos^n \frac{\alpha}{2} (\cos \alpha + i \sin \alpha)\left(\cos \frac{n\alpha}{2} + i \sin \frac{n\alpha}{2}\right) \\
        &= 2^n \cos^n \frac{\alpha}{2} \left(\cos \frac{n+2}{2} \alpha + i \sin \frac{n+2}{2} \alpha\right)
        \end{align*}
        比较实部与虚部,可得
        \begin{align*}
        \cos \alpha + \comb{n}{1} \cos 2\alpha + \comb{n}{2} \cos 3\alpha + \dots + \cos (n + 1)\alpha &= 2^n \cos^n \frac{\alpha}{2} \cos \frac{n+2}{2} \alpha \\
        \sin \alpha + \comb{n}{1} \sin 2\alpha + \comb{n}{2} \sin 3\alpha + \dots + \sin (n + 1)\alpha &= 2^n \cos^n \frac{\alpha}{2} \sin \frac{n+2}{2} \alpha
        \end{align*}
    \end{solution}

    \question
    \begin{parts}
    \part 证明对于任意正整数 $n$ 有
    \begin{align*}
    \sin n\theta &= \binom{n}{1}\sin\theta \cos^{n-1}\theta - \binom{n}{3}\sin^{3}\theta \cos^{n-3}\theta + \binom{n}{5}\sin^{5}\theta \cos^{n-5}\theta - \dots, \\
    \cos n\theta &= \cos^{n}\theta - \binom{n}{2}\sin^{2}\theta \cos^{n-2}\theta + \binom{n}{4}\sin^{4}\theta \cos^{n-4}\theta - \dots
    \end{align*}
    \begin{solution}
        根据棣莫弗定理,我们有
        \[
        \cos n\theta + i \sin n\theta = (\cos\theta + i \sin\theta)^n \tag{1}
        \]
        利用二项式定理展开等式右边
        \[
        (\cos\theta + i \sin\theta)^n = \sum_{k=0}^{n} \binom{n}{k} (i \sin\theta)^k (\cos\theta)^{n-k}
        \]
        展开各项并提取 $i$ 的幂次:
        \begin{align*}
        (\cos\theta + i \sin\theta)^n = &\binom{n}{0}\cos^{n}\theta + i\binom{n}{1}\cos^{n-1}\theta \sin\theta + i^2\binom{n}{2}\cos^{n-2}\theta \sin^2\theta \\
        &+ i^3\binom{n}{3}\cos^{n-3}\theta \sin^3\theta + i^4\binom{n}{4}\cos^{n-4}\theta \sin^4\theta + \dots
        \end{align*}
        由于 $i^2 = -1, i^3 = -i, i^4 = 1, i^5 = i,\dots$,可将展开式整理为
        \begin{align*}
        &(\cos\theta + i \sin\theta)^n \\
        &= \left( \cos^{n}\theta - \binom{n}{2}\cos^{n-2}\theta \sin^{2}\theta + \binom{n}{4}\cos^{n-4}\theta \sin^{4}\theta - \dots \right) \\
        &+ i \left( \binom{n}{1}\cos^{n-1}\theta \sin\theta - \binom{n}{3}\cos^{n-3}\theta \sin^{3}\theta + \binom{n}{5}\cos^{n-5}\theta \sin^{5}\theta - \dots \right)
        \end{align*}
        由(1),比较实部和虚部即得
        \begin{align*}
        \cos n\theta &= \cos^{n}\theta - \binom{n}{2}\sin^{2}\theta \cos^{n-2}\theta + \binom{n}{4}\sin^{4}\theta \cos^{n-4}\theta - \dots \\
        \sin n\theta &= \binom{n}{1}\sin\theta \cos^{n-1}\theta - \binom{n}{3}\sin^{3}\theta \cos^{n-3}\theta + \binom{n}{5}\sin^{5}\theta \cos^{n-5}\theta - \dots
        \end{align*}
        于是得证。
    \end{solution}
    \part 证明对于区间 $[-1,1]$ 上的所有 $x$ 和任意正整数 $n$,函数
    \[
    T_n(x) = \cos(n \cos^{-1} x)
    \]
    是 $x$ 的次数为 $n$ 的多项式,且首项系数为 $2^{\,n-1}$。
    \begin{solution}
        取 $\theta = \cos^{-1} x$,则 $\cos\theta = x$,$\sin\theta = \sqrt{1-x^2}$,应用(a)可得
        \[
        \cos(n \cos^{-1} x) = x^n - \binom{n}{2}(1-x^2)x^{\,n-2} + \binom{n}{4}(1-x^2)^2 x^{\,n-4} - \dots
        \]
        显然 $T_n(x)$ 是 $x$ 的次数为 $n$ 的多项式。其首项系数为
        \[
        1 + \binom{n}{2} + \binom{n}{4} + \dots + \binom{n}{2k}
        \]
        由于
        \[
        2^n = (1+1)^n = 1 + \binom{n}{1} + \binom{n}{2} + \dots + \binom{n}{n},
        \]
        \[
        0 = (1-1)^n = 1 - \binom{n}{1} + \binom{n}{2} - \dots + (-1)^n \binom{n}{n},
        \]
        两式相加,得到
        \[
        1 + \binom{n}{2} + \binom{n}{4} + \dots + \binom{n}{2k} = 2^{\,n-1}.
        \]
        因此首项系数为 $2^{\,n-1}$。
    \end{solution}
    \end{parts}

    \question 证明下列无穷级数的和:
    \[
    \sum_{n=1}^{\infty} (-1)^{n+1} \frac{\sin(n\theta)}{n!} = e^{-\cos\theta}\sin(\sin\theta)
    \]
    \begin{solution}
        定义
        \[
        C = \sum_{n=1}^{\infty} (-1)^{n+1} \frac{\cos(n\theta)}{n!},\quad 
        S = \sum_{n=1}^{\infty} (-1)^{n+1} \frac{\sin(n\theta)}{n!}
        \]
        则
        \[
        C+iS = \sum_{n=1}^{\infty} (-1)^{n+1} \frac{\cos(n\theta)+i\sin(n\theta)}{n!} 
        = \sum_{n=1}^{\infty} (-1)^{n+1} \frac{(e^{i\theta})^n}{n!}
        \]
        观察到指数函数的展开式为
        \[
        e^{-z} = 1 - z + \frac{z^2}{2!} - \frac{z^3}{3!} + \cdots
        \]
        于是
        \[
        \sum_{n=1}^{\infty} \frac{(-1)^{n+1} z^n}{n!} = 1 - e^{-z}
        \]
        取 $z = e^{i\theta}$,得
        \begin{align*}
        C+iS &= 1 - e^{-e^{i\theta}} = 1 - e^{-(\cos\theta + i\sin\theta)} = 1 - e^{-\cos\theta} e^{-i\sin\theta}\\
        &= 1 - e^{-\cos\theta} [\cos(\sin\theta) - i \sin(\sin\theta)] \\
        &= [1 - e^{-\cos\theta}\cos(\sin\theta)] + i [e^{-\cos\theta}\sin(\sin\theta)]
        \end{align*}
        取虚部得到
        \[
        \sum_{n=1}^{\infty} (-1)^{n+1} \frac{\sin(n\theta)}{n!} = \Im(C+iS) = e^{-\cos\theta}\sin(\sin\theta)
        \]
    \end{solution}

    \question 求下列多项式方程的所有实数解:
    \[
    x^7-7x^6-21x^5+35x^4+35x^3-21x^2-7x+1=0.
    \]
    \begin{solution}
        设$\cos\theta=c, \sin\theta=s$,由棣莫弗定理,
        \[
        (c+is)^7=(\cos\theta+i\sin\theta)^7=\cos7\theta+i\sin7\theta
        \]
        且由二项展开得
        \[
        (c+is)^7=c^7+7ic^6s-21c^5s^2-35ic^4s^3+35c^3s^4+21ic^2s^5-7cs^6-is^7
        \]
        比较实部与虚部,得到
        \begin{align*}
        \cos7\theta&=c^7-21c^5s^2+35c^3s^4-7cs^6,\\
        \sin7\theta&=7c^6s-35c^4s^3+21c^2s^5-s^7.
        \end{align*}
        设$t=\tan\theta=\dfrac{s}{c}$,则
        \begin{align*}
        \tan7\theta=\frac{\sin7\theta}{\cos7\theta}=\frac{7t-35t^3+21t^5-t^7}{1-21t^2+35t^4-7t^6}.
        \end{align*}
        令 $\tan7\theta=1$,得到
        \[
        \frac{7t-35t^3+21t^5-t^7}{1-21t^2+35t^4-7t^6}=1,
        \]
        整理得
        \[
        t^7-7t^6-21t^5+35t^4+35t^3-21t^2-7t+1=0.
        \]
        这正是题目所给的多项式方程,其中 $t=x$。由 $\tan7\theta=1$,
        \[
        7\theta=\frac{\pi}{4}+n\pi,
        \]
        从而
        \[
        \theta=\frac{(4n+1)\pi}{28},\quad n=0,1,2,3,4,5,6.
        \]
        因此方程的所有实数解为
        \[
        \tan\frac{\pi}{28},\tan\frac{5\pi}{28},\tan\frac{9\pi}{28},\tan\frac{13\pi}{28},\tan\frac{17\pi}{28},\tan\frac{21\pi}{28},\tan\frac{25\pi}{28}.
        \]
    \end{solution}

    \question 
    \begin{parts}
    \part 证明:
    \[
    \sin7\theta = 7\sin\theta - 56\sin^3\theta + 112\sin^5\theta - 64\sin^7\theta
    \]
    \begin{solution}
        从棣莫弗定理出发,
        \[
        (\cos\theta + i\sin\theta)^7 = \cos 7\theta + i\sin 7\theta
        \]
        由二项式定理展开左式,
        \[
        (\cos\theta + i\sin\theta)^7 = \sum_{k=0}^{7} \comb{7}{k} \cos^{7-k} \theta (i\sin\theta)^k
        \]
        比较虚部项系数得
        \[
        \sin 7\theta =  \comb{7}{1} \cos^6\theta \sin\theta - \comb{7}{3} \cos^4\theta \sin^3\theta + \comb{7}{5} \cos^2\theta \sin^5\theta - \comb{7}{7} \sin^7\theta
        \]
        由恒等式 \(\cos^2\theta = 1 - \sin^2\theta\),令 \(s = \sin\theta\),则有
        \[
        \sin 7\theta = 7(1 - s^2)^3 s - 35(1 - s^2)^2 s^3 + 21(1 - s^2) s^5 - s^7=7s - 56s^3 + 112s^5 - 64s^7
        \]
        故得证。
    \end{solution}
    \part 据此,解方程
    \[
    1 + 7x - 56x^3 + 112x^5 - 64x^7 = 0
    \]
    \begin{solution}
        原方程
        \[
        1 + 7x - 56x^3 + 112x^5 - 64x^7 = 0
        \]
        可写为
        \[
        1 + \sin 7\theta = 0 \Rightarrow \sin 7\theta = -1
        \]
        解得
        \[
        7\theta = \dots,-\frac{5\pi}{2},-\frac{\pi}{2},\frac{3\pi}{2},\frac{7\pi}{2},\dots
        \Rightarrow  \theta = \dots,-\frac{5\pi}{14},-\frac{\pi}{14},\frac{3\pi}{14},\frac{7\pi}{14},\dots
        \]
        所以
        \[
        x=\sin \theta = \sin\left(-\frac{5\pi}{14}\right),\sin\left(-\frac{\pi}{14}\right),\sin \frac{3\pi}{14},\sin \frac{7\pi}{14}=-0.901, -0.223, 0.623, 1
        \]
    \end{solution}
    \part 通过构造合适的多项式方程,证明
    \[
    \csc^2\frac{\pi}{7}+\csc^2\frac{2\pi}{7}+\csc^2\frac{3\pi}{7}=8
    \]
    \begin{solution}
        考虑解为
        \[
        \theta=\cdots,0, \frac{\pi}{7}, \frac{2\pi}{7}, \frac{3\pi}{7},\cdots2
        \]
        的方程
        \[
        \sin 7\theta=0,
        \]
        由三角恒等式
        \[
        \sin7\theta
        =7\sin\theta-56\sin^3\theta+112\sin^5\theta-64\sin^7\theta,
        \]
        方程变为
        \[
        -\sin\theta\bigl(64\sin^6\theta-112\sin^4\theta+56\sin^2\theta-7\bigr)=0.
        \]
        注意到
        \[
        \sin\frac{\pi}{7}=\sin\frac{6\pi}{7},\quad
        \sin\frac{2\pi}{7}=\sin\frac{5\pi}{7},\quad
        \sin\frac{3\pi}{7}=\sin\frac{4\pi}{7},
        \]
        则对非零解 $\sin\theta\neq0$,令$z=\sin^2\theta$,则方程
        \[
        64z^3-112z^2+56z-7=0
        \]
        的三根为
        \[
        \alpha=\sin^2\frac{\pi}{7},\quad
        \beta=\sin^2\frac{2\pi}{7},\quad
        \gamma=\sin^2\frac{3\pi}{7}.
        \]
        由韦达定理,
        \[
        \alpha+\beta+\gamma = \frac{112}{64}=\frac{7}{4},\quad
        \alpha\beta+\beta\gamma+\gamma\alpha = \frac{56}{64}=\frac{7}{8},\quad
        \alpha\beta\gamma = \frac{7}{64}
        \]
        于是
        \[
        \csc^2\frac{\pi}{7}+\csc^2\frac{2\pi}{7}+\csc^2\frac{3\pi}{7}
        =\frac{1}{\alpha}+\frac{1}{\beta}+\frac{1}{\gamma}
        =\frac{\alpha\beta+\beta\gamma+\gamma\alpha}{\alpha\beta\gamma}
        =8
        \]
        证毕。
    \end{solution}
    \end{parts}

    \question
    \begin{parts}
    \part 证明
    \[
    (1+i\tan\theta)^4+(1-i\tan\theta)^4=\frac{2\cos4\theta}{\cos^4\theta}.
    \]
    \begin{solution}
        由
        \begin{align*}
        (1+i\tan\theta)^4 + (1-i\tan\theta)^4 &= \left(1 + \frac{i\sin\theta}{\cos\theta}\right)^4 + \left(1 - \frac{i\sin\theta}{\cos\theta}\right)^4 \\
        &= \left(\frac{\cos\theta + i\sin\theta}{\cos\theta}\right)^4 + \left(\frac{\cos\theta - i\sin\theta}{\cos\theta}\right)^4 \\
        &= \frac{\cos 4\theta + i\sin 4\theta}{\cos^4\theta} + \frac{\cos 4\theta - i\sin 4\theta}{\cos^4\theta} \\
        &= \frac{2\cos 4\theta}{\cos^4\theta}
        \end{align*}
        故左式=右式,等式得证。
    \end{solution}
    \part 通过构造合适的多项式方程,证明
    \[
    \tan^2\frac{\pi}{8}\tan^2\frac{3\pi}{8}=1, \quad
    \tan^2\frac{\pi}{8}+\tan^2\frac{3\pi}{8}=6
    \]
    \begin{solution}
        考虑 $\cos 4\theta = 0$的解
        \[ 
        4\theta = \frac{\pi}{2}, \frac{3\pi}{2}, \frac{5\pi}{2}, \frac{7\pi}{2}, \dots \Rightarrow \theta = \frac{\pi}{8}, \frac{3\pi}{8}, \frac{5\pi}{8}, \frac{7\pi}{8}, \dots 
        \]
        则 
        \[
        z=i\tan\frac{\pi}{8},\; i\tan\frac{3\pi}{8},\; i\tan\frac{5\pi}{8},\; i\tan\frac{7\pi}{8}
        \]
        为方程
        \[
        (1+z)^4 + (1-z)^4=0 \Rightarrow z^4 + 6z^2 + 1 = 0
        \]
        的解,由韦达定理,
        \[ 
        i\tan\frac{\pi}{8} \cdot i\tan\frac{3\pi}{8} \cdot i\tan\frac{5\pi}{8} \cdot i\tan\frac{7\pi}{8} = 1 
        \]
        又因为 
        \[
        \tan\frac{7\pi}{8} = -\tan\frac{\pi}{8},\quad \tan\frac{5\pi}{8} = -\tan\frac{3\pi}{8}
        \]
        故得证
        \[ 
        \tan^2\frac{\pi}{8} \tan^2\frac{3\pi}{8} = 1 
        \]
        又由韦达定理,
        \[
        i\tan\frac{\pi}{8} + i\tan\frac{3\pi}{8} + i\tan\frac{5\pi}{8} + i\tan\frac{7\pi}{8} = 0
        \]
        两边平方得
        \[
        \left(i\tan\frac{\pi}{8}\right)^2 + \left(i\tan\frac{3\pi}{8}\right)^2 + \left(i\tan\frac{5\pi}{8}\right)^2 + \left(i\tan\frac{7\pi}{8}\right)^2 + 2\cdot 6=0
        \]
        给出
        \[ 
        \tan^2\frac{\pi}{8} + \tan^2\frac{3\pi}{8} = 6
        \]
    \end{solution}
    \end{parts}

    \question 设 
    \[
    w = \cos\frac{2\pi}{5} + i\sin\frac{2\pi}{5},
    \]
    \begin{parts}
    \part 证明 
    \[
    1+w+w^2+w^3+w^4=0.
    \] 
    \begin{solution}
        由棣莫弗定理, 
        \[
        w^5 = (\cos\frac{2\pi}{5} + i\sin\frac{2\pi}{5})^5 = \cos 2\pi + i \sin 2\pi = 1,
        \]
        由于$w \neq 1$,所以由$w^5 - 1 = 0$得
        \[
       (w-1)(1+w+w^2+w^3+w^4) = 0 \Rightarrow 1+w+w^2+w^3+w^4 = 0.
        \]
    \end{solution}
    \part 求 
    \[
    (1-w)(1-w^2)(1-w^3)(1-w^4)
    \]
    的值。
    \begin{solution}
        有
        \begin{align*}
        &(1-w)(1-w^2)(1-w^3)(1-w^4) \\
        &= (1-w)(1-w^4)(1-w^2)(1-w^3) \\
        &= [2-(w+w^4)][2-(w^2+w^3)] \\
        &= 4 - 2(w+w^4+w^2+w^3) + (w+w^4)(w^2+w^3) \\
        &= 4 - 2(-1) + w^3+w^4+w^6+w^7 \\
        &= 6 + w^3+w^4+w+w^2 \\
        &= 6 - 1 \\
        &= 5
        \end{align*}
    \end{solution}
    \part 证明 
    \[
    (1-w)(1-w^4) = 4\sin^2\frac{\pi}{5}.
    \] 
    \begin{solution}
        由于$w$与$w^4$共轭,有
        \begin{align*}
        (1-w)(1-w^4) &= 1 - w - w^4 + w^5 \\
        &= 2 - (w + w^4) \\
        &= 2 - 2\cos\frac{2\pi}{5} \\
        &= 4\sin^2\frac{\pi}{5}
        \end{align*}
    \end{solution}
    \part 证明 
    \[
    (1-w^2)(1-w^3) = 4\sin^2\frac{2\pi}{5}.
    \] 
    \begin{solution}
        同理由于$w^2$与$w^3$共轭,有
        \begin{align*}
        (1-w^2)(1-w^3) &= 1 - w^2 - w^3 + w^5 \\ 
        &= 2 - (w^2 + w^3) \\
        &= 2 - 2 \cos\frac{4\pi}{5} \\
        &= 4\sin^2\frac{2\pi}{5}
        \end{align*}
    \end{solution}
    \part 证明 
    \[
    \sin\frac{\pi}{5}\sin\frac{2\pi}{5} = \frac{\sqrt{5}}{4}.
    \]
    \begin{solution}
        由(b),(c),(d)得
        \[
        16 \sin^2\frac{\pi}{5} \sin^2\frac{2\pi}{5} = (1-w)(1-w^2)(1-w^3)(1-w^4) = 5, 
        \]
        即       
        \[
        \sin^2\frac{\pi}{5} \sin^2\frac{2\pi}{5} = \frac{5}{16},
        \]
        由于 $\frac{\pi}{5},\frac{2\pi}{5}$为锐角,所以$\sin\dfrac{\pi}{5},\sin\dfrac{2\pi}{5}>0$,故取正值
        \[
        \sin\frac{\pi}{5} \sin\frac{2\pi}{5} = \frac{\sqrt{5}}{4}
        \]
    \end{solution}
    \end{parts}

    \question 设 \(S_n=a^n+b^n\),其中 \(a,b\) 是方程 \(x^2+x+1=0\) 的根,求
    \[
    \sum_{n=0}^{1729}(-1)^n S_n
    \]
    \begin{solution}
        方程根为 \(a = \omega,\ b = \omega^2\),其中 \(\omega = e^{2\pi i/3}\),满足 \(\omega^3 = 1\)。设 
        \[
        S_n = \omega^n + \omega^{2n},
        \]
        发现 \(S_n\) 是周期为 $3$ 的数列:
        \[
        S_0 = 2,\quad S_1 = -1,\quad S_2 = -1
        \]
        将原式写成
        \[
        \sum_{k=0}^{576} \left[(-1)^{3k} S_{3k} + (-1)^{3k+1} S_{3k+1} + (-1)^{3k+2} S_{3k+2}\right] + (-1)^{1728} S_{1728} + (-1)^{1729} S_{1729}
        \]
        前面 $576$ 组和为 $0,$剩余
        \[
        (-1)^{1728} S_{1728} + (-1)^{1729} S_{1729} = 1 \cdot 2 + (-1) \cdot (-1) = 3
        \]
    \end{solution}

    \question 设 \(\omega \in \mathbb{C},\omega \ne 1\)、且 \(\omega^7 = 1\),计算
    \[
    \prod_{k=0}^{6} (\omega^{2k} + 2\omega^k + 4)
    \]
    \begin{solution}
        设方程 \(x^7 - 1 = 0\) 的根为 \(1,\omega,\omega^2,\dots,\omega^6\),则
        \[
        \prod_{k=0}^{6} \left((\omega^k)^2 + 2\omega^k + 4\right)
        =\prod_{k=0}^{6} \frac{(\omega^k)^3 - 8}{\omega^k - 2}
        \]
        由于 $(\omega^0,\omega^3,\omega^6,\omega^9,\omega^{12},\omega^{15},\omega^{18})=(\omega^0,\omega^3,\omega^6,\omega^2,\omega^{5},\omega^{1},\omega^{4})$,因此:
        \[
        \prod_{k=0}^{6} ((\omega^k)^3 - 8) = \prod_{j=0}^{6} (\omega^j - 8)
        \]
        又因为 \(x^7 - 1 = \prod_{j=0}^{6} (x - \omega^j)\),令 \(x = 8\),得
        \[
        \prod_{j=0}^{6} (\omega^j - 8) = -(8^7 - 1) = 1 - 8^7
        \]
        同理,令 \(x = 2\)得
        \[
        \prod_{j=0}^{6} (\omega^j - 2) = -(2^7 - 1) = 1 - 2^7
        \]
        因此原式为:
        \[
        \frac{1 - 8^7}{1 - 2^7}=16513
        \]
    \end{solution}

    \question 令 $\omega = \cos\dfrac{2\pi}{111} + i \sin\dfrac{2\pi}{111}$, 其中 $i = \sqrt{-1}$, 求 \[ 
    \sum_{k=1}^{110} \frac{\omega^{2k}}{\omega^k - 1}
    \]
    \begin{solution}
        由$\omega = \cos\dfrac{2\pi}{111} + i \sin\dfrac{2\pi}{111}$可知
        \[
        \omega^{111}-1 = 0 \Rightarrow \sum_{k=0}^{110} \omega^k = 0
        \]
        因此
        \[
        \sum_{k=1}^{110} \frac{\omega^{2k}}{\omega^k - 1} = \sum_{k=1}^{110} \left(\omega^k + 1 + \frac{1}{\omega^k - 1}\right) = -1 + 110 + \sum_{k=1}^{110} \frac{1}{\omega^k - 1} \tag{1}
        \]
        设 
        \[f(x) = \sum_{k=0}^{110} x^k = \prod_{k=1}^{110} (x - \omega^k)
        \]
        则
        \[ f'(x) = \sum_{k=1}^{110} k x^{k-1} = \sum_{m=1}^{110} \prod_{k=1, k\ne m}^{110} (x - \omega^k)
        \]
        于是
        \[
        g(x) = \frac{f'(x)}{f(x)} = \sum_{k=1}^{110} \frac{1}{x - \omega^k}
        \]
        令$x=1$,
        \[
        g(1) = \frac{111 \cdot \frac{110}{2}}{111} = \sum_{k=1}^{110} \frac{1}{1 - \omega^k}\Rightarrow \sum_{k=1}^{110} \frac{1}{\omega^k - 1} = -55
        \]
        代入 (1) 得
        \[
        \sum_{k=1}^{110} \frac{\omega^{2k}}{\omega^k - 1} = -1 + 110 - 55 = 54
        \]
    \end{solution}

    \question 设 \(z=\cos\dfrac{\pi}{6}+i\sin\dfrac{\pi}{6}\),求
    \[
    \left|1+2z+3z^2+\cdots+2016z^{2015}\right|
    \]
    \begin{solution}
        记
        \[
        S = 1 + 2z + 3z^2 + \cdots + 2016z^{2015} 
        \]
        考虑函数
        \[
        f(x) = x + x^2 + \cdots + x^{2016} = \frac{x(x^{2016} - 1)}{x - 1},\quad 
        \]
        则有$$S = f'(z)$$
        求导:
        \[
        f'(x) = \frac{[(x^{2016} - 1) + 2016x^{2016}](x - 1) - x(x^{2016} - 1)}{(x - 1)^2}
        \]
        设 \(z = \cos\dfrac{\pi}{6} + i\sin\dfrac{\pi}{6} = e^{\frac{i\pi}{6}}\),由棣莫弗定理得
        \[
        z^{2016} = \left(e^{i\pi/6}\right)^{2016} = e^{i\cdot336\pi} = 1
        \]
        于是
        \[
        f'(z) = \frac{2016(z - 1)}{(z - 1)^2} = \frac{2016}{z - 1}
        \Rightarrow |S| = \frac{2016}{|z - 1|}
        \]
        而
        \[
        z - 1 = \left(\frac{\sqrt{3}}{2} - 1\right) + \frac{1}{2}i
        \]
        \[
        |z - 1| = \sqrt{\left(\frac{\sqrt{3}}{2} - 1\right)^2 + \left(\frac{1}{2}\right)^2} = \sqrt{2 - \sqrt{3}} = \frac{\sqrt6+\sqrt2}{2}
        \]
        所以
        \[
        |S| = \frac{2016\cdot2}{\sqrt6+\sqrt2}\cdot \frac{\sqrt6-\sqrt2}{\sqrt6-\sqrt2} = 1008\sqrt{2} + 1008\sqrt{6}
        \]
    \end{solution}

    \question 设 $\omega = e^{\frac{2\pi i}{n}}$,求
    \[ 
    -\sum_{k=1}^{n-1} \prod_{\substack{j=1 \\ j \neq k}}^{n-1} (\omega^k - \omega^j). 
    \]
    \begin{solution}
        设
        \[
        P(x)=\prod_{k=1}^{n-1}(x-\omega^k)=\frac{x^n-1}{x-1}=x^{n-1}+x^{n-2}+\cdots +x+1
        \]
        则
        \[
        P'(x)=(n-1)x^{n-2}+(n-2)x^{n-3}+\cdots + 1
        \]
        于是
        \[
        \sum_{k=1}^{n-1} \prod_{\substack{j=1 \\ j \neq k}}^{n-1} (\omega^k - \omega^j)=-\sum_{k=1}^{n-1} P'(\omega^k)
        \]
        观察到由于$\omega^n=1$,
        \[
        \sum_{k=0}^{n-1} \omega^{(n-1)k}
        =\begin{cases}
        n, \quad k=0 \\[2pt]
        \dfrac{1-\omega^{nk}}{1-\omega^k}=0,\quad k \neq 0 \\
        \end{cases}
        \]
        因此
        \[
        \sum_{k=0}^{n-1} P'(\omega^k)=0+\cdots+0+n=n
        \]
        故
        \[
        \sum_{k=1}^{n-1} \prod_{\substack{j=1 \\ j \neq k}}^{n-1} (\omega^k - \omega^j)=-\sum_{k=1}^{n-1} P'(\omega^k)=-(n-P'(1))=\frac{n(n-1)}{2}-n=\frac{n(n-3)}{2}
        \]
    \end{solution}

    \question 求解方程 $\sin x = \pi$。
\begin{solution}
利用恒等式 $\sin x = \frac{e^{ix} - e^{-ix}}{2i}$：
\[ \frac{e^{ix} - e^{-ix}}{2i} = \pi \implies e^{ix} - e^{-ix} = 2i\pi \]
设 $u = e^{ix}$，构造二次方程 $u^2 - 2i\pi u - 1 = 0$。
解得：
\[ u = i\pi \pm \sqrt{-4\pi^2 + 4} / 2 = i(\pi \pm \sqrt{\pi^2 - 1}) \]
对 $u$ 取复对数 $\ln u = \ln |u| + i\text{arg}(u)$。
由于 $u$ 位于虚轴正半轴，其幅角为 $\frac{\pi}{2} + 2k\pi$。
最终得到解为：
\[ x = \frac{4n+1}{2}\pi - i \ln(\pi \pm \sqrt{\pi^2 - 1}), \quad n \in \mathbb{Z} \]
\end{solution}

    \question 设复数平面上三点 $A(\alpha)$, $B(\beta)$, $C(\gamma)$ 可连成正三角形 $ABC$,  
已知 $\alpha$, $\beta$, $\gamma$ 满足
\[
\alpha^4 - 2\alpha^3\beta + (\beta^2 - 4)\alpha^2 + 8\alpha\gamma - 4\gamma^2 = 0
\]
且 $\alpha$ 的实部和虚部均为正数,  
当 $\triangle ABC$ 的重心 $G$ 为 $\frac{\alpha}{2^{110}}$ 时,求 $\beta$ 及 $\gamma$ 各为何？
\begin{solution}
求解 $\alpha$ 的值
由原式因式分解得:
\[ (\alpha^2 - \alpha\beta)^2 - 4(\alpha - \gamma)^2 = 0 \implies \alpha(\alpha - \beta) = \pm 2(\alpha - \gamma) \]
由于 $\triangle ABC$ 为正三角形,满足 $\alpha - \beta = (\alpha - \gamma)e^{\pm i\frac{\pi}{3}}$,代入上式:
\[ \alpha (\alpha - \gamma) e^{\pm i\frac{\pi}{3}} = \pm 2(\alpha - \gamma) \implies \alpha = \pm 2 e^{\mp i\frac{\pi}{3}} \]
已知 $\alpha$ 的实部与虚部均为正,故:
\[ \alpha = 2e^{i\frac{\pi}{3}} = 1 + \sqrt{3}i \]

利用重心与旋转性质求解 $\beta, \gamma$

已知重心 $G = \frac{\alpha}{2^{1110}}$。根据正三角形的几何性质,顶点 $\beta$ 和 $\gamma$ 是由顶点 $\alpha$ 绕重心 $G$ 分别旋转 $120^\circ$ 和 $-120^\circ$ 得到的:
\begin{align*}
\beta &= G + (\alpha - G)e^{i\frac{2\pi}{3}} = \frac{\alpha}{2^{1110}} + \alpha\left(1 - \frac{1}{2^{1110}}\right)e^{i\frac{2\pi}{3}} \\
\gamma &= G + (\alpha - G)e^{-i\frac{2\pi}{3}} = \frac{\alpha}{2^{1110}} + \alpha\left(1 - \frac{1}{2^{1110}}\right)e^{-i\frac{2\pi}{3}}
\end{align*}

具体数值计算
将 $\alpha = 2e^{i\frac{\pi}{3}}$ 代入 $\beta$:
\begin{align*}
\beta &= \frac{2e^{i\frac{\pi}{3}}}{2^{1110}} + 2e^{i\frac{\pi}{3}}\left(1 - \frac{1}{2^{1110}}\right)e^{i\frac{2\pi}{3}} \\
&= \frac{1 + \sqrt{3}i}{2^{1110}} + 2\left(1 - \frac{1}{2^{1110}}\right)e^{i\pi} \\
&= \frac{1 + \sqrt{3}i}{2^{1110}} - 2 + \frac{2}{2^{1110}} \\
&= \left( \frac{3}{2^{1110}} - 2 \right) + \frac{\sqrt{3}}{2^{1110}}i
\end{align*}

将 $\alpha = 2e^{i\frac{\pi}{3}}$ 代入 $\gamma$:
\begin{align*}
\gamma &= \frac{2e^{i\frac{\pi}{3}}}{2^{1110}} + 2e^{i\frac{\pi}{3}}\left(1 - \frac{1}{2^{1110}}\right)e^{-i\frac{2\pi}{3}} \\
&= \frac{1 + \sqrt{3}i}{2^{1110}} + 2\left(1 - \frac{1}{2^{1110}}\right)e^{-i\frac{\pi}{3}} \\
&= \frac{1 + \sqrt{3}i}{2^{1110}} + 2\left(1 - \frac{1}{2^{1110}}\right)\left(\frac{1}{2} - \frac{\sqrt{3}}{2}i\right) \\
&= \frac{1 + \sqrt{3}i}{2^{1110}} + \left(1 - \frac{1}{2^{1110}}\right) - \sqrt{3}\left(1 - \frac{1}{2^{1110}}\right)i \\
&= 1 + \sqrt{3}\left(\frac{2}{2^{1110}} - 1\right)i
\end{align*}

综上所述:
\[ \beta = \left( \frac{3}{2^{1110}} - 2 \right) + \frac{\sqrt{3}}{2^{1110}}i \]
\[ \gamma = 1 + \left( \frac{\sqrt{3}}{2^{1109}} - \sqrt{3} \right)i \]
\textcolor{red}{(待验证)}
\end{solution}

    \question 设 $P(z)$ 是一个 $n$ 次复系数多项式,其所有零点都位于复平面的单位圆上。证明多项式
    \[
    \tilde{P}(z)=2zP'(z)-nP(z)
    \]
    的所有零点也都位于同一单位圆上。
    \begin{solution}
        不失一般性,考虑首项系数为 $1$ 的多项式。设
        \[
        P(z)=(z-\alpha_1)(z-\alpha_2)\cdots(z-\alpha_n),
        \]
        其中 $|\alpha_j|=1,j=1,2,\ldots,n$,并允许这些复数 $\alpha_1,\alpha_2,\ldots,\alpha_n$ 相同,则
        \begin{align*}
        \tilde{P}(z)&=2zP'(z)-nP(z)\\
        &=(z+\alpha_1)(z-\alpha_2)\cdots(z-\alpha_n)+\cdots+(z-\alpha_1)(z-\alpha_2)\cdots(z+\alpha_n)
        \end{align*}
        因此
        \[
        \frac{\tilde{P}(z)}{P(z)}=\sum_{k=1}^n \frac{z+\alpha_k}{z-\alpha_k}。
        \]
        注意到,对任意复数 $z,\alpha$ 且 $z\neq\alpha$,皆有
        \[
        \Re\left[\frac{z+\alpha}{z-\alpha}\right]
        =\frac{|z|^2-|\alpha|^2}{|z-\alpha|^2}
        \]
        于是
        \[
        \Re\left[\frac{\tilde{P}(z)}{P(z)}\right]
        =\sum_{k=1}^n \frac{|z|^2-|\alpha_k|^2}{|z-\alpha_k|^2}
        \]
        由于 $|\alpha_k|=1$,可知当 $|z|\neq1$ 时,上式不为零,因此若 $\tilde{P}(z)=0$,则必有
        \[
        \Re\left[\frac{\tilde{P}(z)}{P(z)}\right]=0,
        \]
        从而推出 $|z|=1$,这说明 $\tilde{P}(z)$ 的所有零点也都位于单位圆上。
    \end{solution}

    \question 复数平面上以原点为中心的单位圆中,有一内接四边形,其顶点为 $z_1, z_2, z_3, z_4$,设
    \[
    S_n = z_1^n + z_2^n + z_3^n + z_4^n
    \]
    且 $S_1 = 0$ 且 $S_2 = 1$,
    \begin{parts}
    \part 证明:若
    \[
    |z_1|=|z_2|=|z_3|=|z_4|=1,\quad z_1+z_2+z_3+z_4=0
    \]
    则 $z_1,z_2,z_3,z_4$ 为一个矩形的顶点。
    \begin{solution}
        设
        \[
        f(z) = (z-z_1)(z-z_2)(z-z_3)(z-z_4) = z^4 + a z^3 + b z^2 + c z + d
        \]
        由 $S_1 = z_1+z_2+z_3+z_4 = 0$ 得 $a=0$,又由于 $|z_k|=1$,有 $\bar z_k = \dfrac{1}{z_k}$,于是
        \[
        \overline{z_1} + \overline{z_2} + \overline{z_3} + \overline{z_4} = 0
        \quad\Rightarrow\quad
        \frac1{z_1}+\frac1{z_2}+\frac1{z_3}+\frac1{z_4} = 0
        \]
        即
        \[
        z_2 z_3 z_4 + z_1 z_3 z_4 + z_1 z_2 z_4 + z_1 z_2 z_3 = 0
        \]
        所以 $c=0$,因此
        \[
        f(z) = z^4 + b z^2 + d
        \]
        为偶函数,四根关于原点成对对称,于是可设
        \[
        z_1+z_3=0,z_2+z_4=0
        \]
        此时两对顶点互为相反数,对角线过原点,且 $|z_1|=|z_2|=1$,故为内接矩形。
    \end{solution}
    \part 计算该矩形的面积。
    \begin{solution}
        设
        \[
        z_1 = a+bi,\quad z_2 = -a+bi,\quad z_3 = -a-bi,\quad z_4 = a-bi
        \]
        由 $S_2 = z_1^2+z_2^2+z_3^2+z_4^2 = 4(a^2-b^2) = 1$ 得
        \[
        a^2 - b^2 = \frac14
        \]
        又因 $|z_1|=1$,有 $a^2+b^2=1$,解得
        \[
        a = \frac{\sqrt5}{2\sqrt2},\quad b = \frac{\sqrt3}{2\sqrt2}
        \]
        故矩形面积为
        \[
        2a \cdot 2b = \sqrt{\frac52} \cdot \sqrt{\frac32} = \frac{\sqrt{15}}{2}
        \]
    \end{solution}
    \end{parts}

\end{questions}
\pagebreak
\begin{center}
  {\fontsize{30pt}{26pt}\selectfont
    \hypertarget{数学归纳法}{数学归纳法} \label{数学归纳法}
  }
\end{center}
\separator
\vspace{1pt}

\begin{questions}
    \question 证明对任意正整数 $n$, 
    \[
    (n + 1)(n + 2)(n + 3) \cdots (2n - 1)(2n) = 2^n \cdot 1 \cdot 3 \cdot 5 \cdots (2n - 1)
    \]
    \begin{solution}
        定义命题
        \[
        P_n:\prod_{i=1}^{n} (n + i) = 2^n \cdot \prod_{i=1}^{n} (2i - 1), \quad n \in \mathbb{N}
        \]
        基础步骤:观察当 $n=1$ 时, 左式为
        \[
        (1 + 1) = 2
        \]
        而右式为
        \[
        2^1 \cdot 1 = 2
        \]
        此时左式等于右式, 即 $P_1$ 成立。

        归纳步骤:设 $k \in \mathbb{Z}^{+}$, 假设归纳假设 $P_k$ 成立, 有
        \[
        (k + 1)(k + 2) \cdots (2k) = 2^k \cdot 1 \cdot 3 \cdots (2k - 1)
        \]
        欲证 $P_{k+1}$ 成立, 即证
        \[
        (k + 2)(k + 3) \cdots (2k)(2k + 1)(2k + 2) = 2^{k+1} \cdot 1 \cdot 3 \cdots (2k - 1)(2k + 1)
        \]
        由归纳假设, 观察 $P_{k+1}$ 的左式:
        \begin{align*}
        (k + 2)(k + 3) \cdots (2k)(2k + 1)(2k + 2) 
        &= \frac{(k + 1)(k + 2) \cdots (2k) \cdot (2k + 1) \cdot (2k + 2)}{k + 1} \\
        &= \frac{ \left[ 2^k \cdot 1 \cdot 3 \cdots (2k - 1) \right] \cdot (2k + 1) \cdot 2(k + 1)}{k + 1} \\
        &= \frac{2^k \cdot 1 \cdot 3 \cdots (2k - 1) \cdot (2k + 1) \cdot 2 \cdot (k + 1)}{k + 1} \\
        &= 2^{k} \cdot 2 \cdot 1 \cdot 3 \cdots (2k - 1)(2k + 1) \\
        &= 2^{k+1} \cdot 1 \cdot 3 \cdots (2k + 1)
        \end{align*}
        此时左式等于 $P_{k+1}$ 的右式, 即命题 $P_{k+1}$ 成立, 由此蕴含 $P_k \implies P_{k+1}$ 成立。

        由数学归纳法, 命题 $P_n$ 对所有正整数 $n$ 皆成立。
    \end{solution}

    \question 证明对任意正整数 $n \ge p$,
    \[
    \sum_{i=1}^{n-p+1} \frac{1}{i(i+1)\cdots(i+p-1)} = \frac{1}{p-1} \left[ \frac{1}{(p-1)!} - \frac{1}{(n-p+2)(n-p+3)\cdots n} \right]
    \]
    \begin{solution}
        定义命题
        \[
        P_n:\sum_{i=1}^{n-p+1} \frac{1}{i(i+1)\cdots(i+p-1)} = \frac{1}{p-1} \left[ \frac{1}{(p-1)!} - \frac{1}{(n-p+2)(n-p+3)\cdots n} \right],
        \]
        其中$n \ge p >1 ,n,p \in \mathbb{N}$。

        基础步骤:当 $n = p$ 时, 左式为
        \[
        \sum_{i=1}^{1} \frac{1}{i(i+1)\cdots(i+p-1)} = \frac{1}{1 \cdot 2 \cdots p} = \frac{1}{p!}
        \]
        右式为
        \[
        \frac{1}{p-1} \left[ \frac{1}{(p-1)!} - \frac{1}{2 \cdot 3 \cdots p} \right] = \frac{1}{p-1} \left[ \frac{p}{p!} - \frac{1}{p!} \right] = \frac{1}{p-1} \cdot \frac{p-1}{p!} = \frac{1}{p!}
        \]
        此时左式等于右式, 即 $P_p$ 成立。

        归纳步骤:设 $k \ge p$, 假设归纳假设 $P_k$ 成立, 即
        \[
        \sum_{i=1}^{k-p+1} \frac{1}{i(i+1)\cdots(i+p-1)} = \frac{1}{p-1} \left[ \frac{1}{(p-1)!} - \frac{1}{(k-p+2)\cdots k} \right]
        \]
        观察当 $n = k + 1$ 时, $P_{k+1}$ 的左式为
        \begin{align*}
        \sum_{i=1}^{k-p+2} \frac{1}{i \cdots (i+p-1)} 
        &= \sum_{i=1}^{k-p+1} \frac{1}{i \cdots (i+p-1)} + \frac{1}{(k-p+2)(k-p+3) \cdots (k+1)} \\
        &= \frac{1}{p-1} \left[ \frac{1}{(p-1)!} - \frac{1}{(k-p+2) \cdots k} \right] + \frac{1}{(k-p+2) \cdots (k+1)} \\
        &= \frac{1}{p-1} \left[ \frac{1}{(p-1)!} - \left( \frac{1}{(k-p+2) \cdots k} - \frac{p-1}{(k-p+2) \cdots (k+1)} \right) \right]
        \end{align*}
        对括号内各项进行通分, 分母取 $(k-p+2)(k-p+3) \cdots (k+1)$:
        \begin{align*}
        \frac{1}{(k-p+2) \cdots k} - \frac{p-1}{(k-p+2) \cdots (k+1)} 
        &= \frac{(k+1) - (p-1)}{(k-p+2)(k-p+3) \cdots (k+1)} \\
        &= \frac{k-p+2}{(k-p+2)(k-p+3) \cdots (k+1)} \\
        &= \frac{1}{(k-p+3) \cdots (k+1)}
        \end{align*}
        代入原式得$P_{k+1}$左式为
        \[
        \frac{1}{p-1} \left[ \frac{1}{(p-1)!} - \frac{1}{(k-p+3) \cdots (k+1)} \right]
        \]
        此时左式等于 $P_{k+1}$ 的右式, 即命题 $P_{k+1}$ 成立, 此时蕴含 $P_k \implies P_{k+1}$。

        由数学归纳法, 命题 $P_n$ 对所有整数 $n \ge p$ 皆成立。
    \end{solution}

    \question 证明:对于任意整数 $n > 1$, 有
    \[ 
    \frac{1}{2} \cdot \frac{3}{4} \cdot \frac{5}{6} \cdots \frac{2n - 1}{2n} < \frac{1}{\sqrt{3n + 1}} 
    \]
    据此, 试证
    \[
    \frac{1}{2} \cdot \frac{3}{4} \cdot \frac{5}{6} \cdots \frac{99}{100} < \frac{1}{12}
    \]
    \begin{solution}
        定义命题 
        \[ 
        P_n: \frac{1}{2} \cdot \frac{3}{4} \cdot \frac{5}{6} \cdots \frac{2n - 1}{2n} < \frac{1}{\sqrt{3n + 1}},\quad n > 1, n \in \mathbb{N}
        \]
        基础步骤:当 $n = 2$ 时, 左式为 $\frac{1}{2} \cdot \frac{3}{4} = 0.375$, 右式为 $\frac{1}{\sqrt{7}} \approx 0.37796$。
        显然 $0.375 < 0.378$, 故 $P_2$ 成立。

        归纳步骤:假设当 $n = k$时, 命题 $P_k$ 成立, 即:
        \[ 
        \frac{1}{2} \cdot \frac{3}{4} \cdots \frac{2k - 1}{2k} \le \frac{1}{\sqrt{3k + 1}} 
        \]
        欲证 $P_{k+1}$ 成立, 观察其左式:
        \[ 
        \left( \frac{1}{2} \cdots \frac{2k - 1}{2k} \right) \cdot \frac{2k + 1}{2k + 2} \le \frac{1}{\sqrt{3k + 1}} \cdot \frac{2k + 1}{2k + 2} 
        \]
        只需证
        \[ 
        \frac{2k + 1}{(2k + 2)\sqrt{3k + 1}} < \frac{1}{\sqrt{3k + 4}} 
        \]
        将上式两边平方, 即证
        \[ 
        \frac{(2k + 1)^2}{(2k + 2)^2(3k + 1)} < \frac{1}{3k + 4} \iff (2k + 1)^2 (3k + 4) < (2k + 2)^2 (3k + 1)
        \]
        展开左式得$12k^3 + 28k^2 + 19k + 4$,展开右式得$12k^3 + 28k^2 + 20k + 4$。
        由于左式>右式对 $k > 1$ 显然成立, 故 $P_{k+1}$ 成立, 由此蕴含 $P_k \implies P_{k+1}$成立。
        
        由数学归纳法可知, 原不等式对所有 $n > 1$ 均成立。

        现令 $n = 50$, 代入上述不等式得
        \[
        \frac{1}{2} \cdot \frac{3}{4} \cdot \frac{5}{6} \cdots \frac{99}{100} < \frac{1}{\sqrt{3(50) + 1}} = \frac{1}{\sqrt{151}}
        \]
        由于 $12^2 = 144 < 151$, 因此 $\sqrt{151} > 12$,由此可得:
        \[
        \frac{1}{\sqrt{151}} < \frac{1}{12}
        \]
        故原不等式 
        \[
        \frac{1}{2} \cdot \frac{3}{4} \cdots \frac{99}{100} < \frac{1}{12}
        \]
        成立。
    \end{solution}

    \question 已知
    \[
    f(n)=3^{2n+4}-2^{2n},\quad n\in\mathbb{N},
    \]
    用数学归纳法证明 $f(n)$ 能被 $5$ 整除。
    \begin{solution}
        定义命题 $P_n$ 为 $5 \mid f(n)$,即 $f(n)$ 能被 $5$ 整除。

        基础步骤:当 $n=1$ 时,
        \begin{align*}
        f(1) &= 3^{2(1)+4}-2^{2(1)} \\
        &= 3^6-2^2 \\
        &= 729-4 = 725
        \end{align*}
        由于 $725 = 5 \times 145$, 显然 $P_1$ 成立。

        归纳步骤:假设当 $n=k$ ($k \in \mathbb{N}$) 时结论成立, 即存在 $m \in \mathbb{Z}$, 使得
        \[
        f(k) = 3^{2k+4}-2^{2k} = 5m
        \]
        欲证 $P_{k+1}$ 成立, 观察 $f(k+1)$ 的表达式:
        \begin{align*}
        f(k+1) &= 3^{2(k+1)+4} - 2^{2(k+1)} \\
        &= 3^{2k+6} - 2^{2k+2} \\
        &= 3^2 \cdot 3^{2k+4} - 2^2 \cdot 2^{2k} \\
        &= 9 \cdot 3^{2k+4} - 4 \cdot 2^{2k}
        \end{align*}
        利用归纳假设 $2^{2k} = 3^{2k+4} - 5m$, 代入上式得:
        \begin{align*}
        f(k+1) &= 9 \cdot 3^{2k+4} - 4(3^{2k+4} - 5m) \\
        &= 9 \cdot 3^{2k+4} - 4 \cdot 3^{2k+4} + 20m \\
        &= 5 \cdot 3^{2k+4} + 20m \\
        &= 5(3^{2k+4} + 4m)
        \end{align*}
        由于 $3^{2k+4} + 4m$ 为整数, 故 $5 \mid f(k+1)$, 即命题 $P_{k+1}$ 成立。
        由此蕴含 $P_k \implies P_{k+1}$ 成立。

        由数学归纳法可知, 对所有 $n \in \mathbb{N}$, $f(n) = 3^{2n+4}-2^{2n}$ 均能被 $5$ 整除。
    \end{solution}

    \question 证明对所有正整数 $n$, 都有
    \[
    \frac{1}{\sqrt{1}} + \frac{1}{\sqrt{2}} + \dots + \frac{1}{\sqrt{n}} \ge \sqrt{n}.
    \]
    \begin{solution}
        定义命题 $P_n:\sum_{i=1}^{n} \frac{1}{\sqrt{i}} \ge \sqrt{n}, \quad n \in \mathbb{Z}^{+}$。

        基础步骤:当 $n=1$ 时, 
        左式为 $\frac{1}{\sqrt{1}} = 1$, 右式为 $\sqrt{1} = 1$。
        此时左式等于右式, 即 $P_1$ 成立。

        归纳步骤:设 $k \in \mathbb{Z}^{+}$, 假设归纳假设 $P_k$ 成立, 有
        \[
        \frac{1}{\sqrt{1}} + \frac{1}{\sqrt{2}} + \dots + \frac{1}{\sqrt{k}} \ge \sqrt{k}
        \]
        欲证 $P_{k+1}$ 成立, 即
        \[
        \frac{1}{\sqrt{1}} + \frac{1}{\sqrt{2}} + \dots + \frac{1}{\sqrt{k}} + \frac{1}{\sqrt{k+1}} \ge \sqrt{k+1}
        \]
        由归纳假设, $P_{k+1}$ 的左式满足
        \[
        \text{左式} \ge \sqrt{k} + \frac{1}{\sqrt{k+1}}
        \]
        只需证
        \begin{align*}
        &\iff \sqrt{k} + \frac{1}{\sqrt{k+1}} \ge \sqrt{k+1} \\
        &\iff \sqrt{k(k+1)} + 1 \ge k+1 \\
        &\iff \sqrt{k^2+k} + 1 \ge k+1 \\
        &\iff \sqrt{k^2+k} \ge k \\
        &\iff k^2+k \ge k^2
        \end{align*}
        由于 $k^2+k \ge k^2$ 对所有 $k \in \mathbb{Z}^{+}$ 显然成立, 故得 $P_{k+1}$ 成立, 此时蕴含 $P_k \implies P_{k+1}$。

        由数学归纳法, 命题 $P_n$ 对所有正整数 $n$ 均成立。
    \end{solution}

    \question 用数学归纳法证明:若 $n \in \mathbb{N}, n \geq 3$, 则
    \[
    n^{n+1} > (n+1)^n,
    \]
    并由此推导出若 $n \in \mathbb{N}, n \geq 3$, 则
    \[
    \sqrt[n]{n} > \sqrt[n+1]{n+1}.
    \]
    \begin{solution}
        定义命题 
        \[
        P_n:n^{n+1} > (n+1)^n, \quad n \ge 3,n \in \mathbb{N}
        \]
        基础步骤:当 $n=3$ 时, 
        左式为 $3^{3+1} = 3^4 = 81$, 右式为 $(3+1)^3 = 4^3 = 64$。
        由于 $81 > 64$, 故 $P_3$ 成立。

        归纳步骤:设 $k \ge 3$, 假设归纳假设 $P_k$ 成立, 即 $k^{k+1} > (k+1)^k$。欲证 $P_{k+1}$ 成立, 即证:
        \[
        (k+1)^{k+2} > (k+2)^{k+1} \iff \frac{(k+1)^{k+2}}{(k+2)^{k+1}} > 1
        \]
        由归纳假设可知 $k+1 > \frac{(k+1)^{k+1}}{k^{k+1}}$, 考虑目标不等式的左端:
        \begin{align*}
        \frac{(k+1)^{k+2}}{(k+2)^{k+1}} &= \frac{(k+1)^{k+1} \cdot (k+1)}{(k+2)^{k+1}} \\
        &> \frac{(k+1)^{k+1} \cdot \frac{(k+1)^{k+1}}{k^{k+1}}}{(k+2)^{k+1}} \quad (\text{利用 } P_k) \\
        &= \frac{(k+1)^{2k+2}}{k^{k+1}(k+2)^{k+1}} \\
        &= \left[ \frac{(k+1)^2}{k(k+2)} \right]^{k+1} \\
        &= \left[ \frac{k^2+2k+1}{k^2+2k} \right]^{k+1}
        \end{align*}
        由于 
        \[
        \frac{k^2+2k+1}{k^2+2k} = 1 + \frac{1}{k^2+2k} > 1, 
        \]
        则其任意正整数次幂均大于 $1$。因此 $(k+1)^{k+2} > (k+2)^{k+1}$ 成立, 即 $P_{k+1}$ 成立, 蕴含 $P_k \implies P_{k+1}$。

        由数学归纳法, 原不等式对所有 $n \ge 3$ 均成立。
        
        由上述已证不等式可知, 当 $n \ge 3$ 时:
        \[
        n^{n+1} > (n+1)^n
        \]
        对不等式两边同时开 $n(n+1)>0$ 次方, 则
        \[
        (n^{n+1})^{\frac{1}{n(n+1)}} > ((n+1)^n)^{\frac{1}{n(n+1)}} \\
        \]
        即
        \[
        \sqrt[n]{n} > \sqrt[n+1]{n+1}
        \]
        得证。
    \end{solution}

    \question 设 $x_0=5$, 且 $x_{n+1}=x_n+\frac{1}{x_n}$。证明对一切自然数 $n \ge 1$ 都有
    \[
    2n < x_n^2 - 25 < \frac{47n}{23}.
    \]
    \begin{solution}
        定义命题 
        \[
        P_n:2n < x_n^2 - 25 < \frac{47n}{23}, \quad n \ge 1,n \in \mathbb{N}
        \]
        基础步骤:当 $n=1$ 时,$x_1 = 5.2, x_1^2 - 25 = 2.04$。观察不等式:
        \[
        2(1) = 2 < 2.04 < \frac{47(1)}{23} \approx 2.04348
        \]
        此时左式与右式均成立, 即 $P_1$ 成立。

        归纳步骤:设 $k \ge 1$, 假设归纳假设 $P_k$ 成立, 即 
        \[
        2k < x_k^2 - 25 < \frac{47k}{23}
        \]
        递推关系平方可得
        \[
        x_{k+1}^2 = \left( x_k + \frac{1}{x_k} \right)^2 = x_k^2 + 2 + \frac{1}{x_k^2}
        \]
        即
        \[
        x_{k+1}^2 - 25 = (x_k^2 - 25) + 2 + \frac{1}{x_k^2}
        \]
        证明下界:由归纳假设 
        \[
        x_k^2 - 25 > 2k,
        \]
        且由于 $x_k^2 > 0$, 必有 $\frac{1}{x_k^2} > 0$,则
        \[
        x_{k+1}^2 - 25 > 2k + 2 + 0 = 2(k+1)
        \]
        证明上界:由归纳假设 
        \[
        x_k^2 - 25 < \frac{47k}{23}
        \]
        注意到数列 $\{x_n\}$ 单调递增, 故对于 $k \ge 1$, 有 $x_k \ge x_1 = 5.2$。
        由此可得 $x_k^2 \ge 27.04 > 23$, 故 $\frac{1}{x_k^2} < \frac{1}{23}$。代入表达式:
        \[
        x_{k+1}^2 - 25 < \frac{47k}{23} + 2 + \frac{1}{23} = \frac{47(k+1)}{23}
        \]
        由此可见 $P_{k+1}$ 成立, 即蕴含 $P_k \implies P_{k+1}$成立。

        由数学归纳法, 命题 $P_n$ 对所有正整数 $n$ 均成立。
    \end{solution}

    \question 设数列 $\{a_n\}$ 定义为 $a_1 = 1,2 = 1$, 且 
    \[
    a_{m+1} = a_m + a_{m-1} \quad \forall m \ge 2, m \in \mathbb{N}
    \]
    证明
    \[
    a_n = \frac{1}{\sqrt{5}}\left[\left(\frac{1+\sqrt{5}}{2}\right)^n - \left(\frac{1-\sqrt{5}}{2}\right)^n\right]
    \]
    对所有正整数 $n$ 皆成立。
    \begin{solution}
        定义命题 
        \[
        P_n:a_n = \frac{1}{\sqrt{5}}\left(\phi^n - \psi^n\right), 
        \]
        其中 
        \[
        \phi = \frac{1+\sqrt{5}}{2},\quad \psi = \frac{1-\sqrt{5}}{2}
        \]
        注意到 $\phi$ 和 $\psi$ 是特征方程 $x^2 - x - 1 = 0$ 的两个根, 满足 
        \[
        \phi^2 = \phi + 1, \quad\psi^2 = \psi + 1
        \]
        基础步骤:当 $n=1$ 时, 右式为 
        \[
        \frac{1}{\sqrt{5}}\left(\frac{1+\sqrt{5}}{2} - \frac{1-\sqrt{5}}{2}\right) = \frac{\sqrt{5}}{\sqrt{5}} = 1 = a_1, 
        \]
        当 $n=2$ 时, 右式为 
        \[
        \frac{1}{\sqrt{5}}\left[\left(\frac{1+\sqrt{5}}{2}\right)^2 - \left(\frac{1-\sqrt{5}}{2}\right)^2\right] = \frac{1}{\sqrt{5}}\left(\frac{6+2\sqrt{5}}{4} - \frac{6-2\sqrt{5}}{4}\right) = \frac{\sqrt{5}}{\sqrt{5}} = 1 = a_2,
        \]
        故 $P_1,P_2$ 成立。

        归纳步骤:假设 $P_{k-1}$ 与 $P_k$ 成立,即
        \[
        a_{k-1} = \frac{1}{\sqrt{5}}(\phi^{k-1} - \psi^{k-1}), \quad a_k = \frac{1}{\sqrt{5}}(\phi^k - \psi^k)
        \]
        欲证 $P_{k+1}$ 成立, 根据递推关系 $a_{k+1} = a_k + a_{k-1}$:
        \begin{align*}
        a_{k+1} &= \frac{1}{\sqrt{5}}(\phi^k - \psi^k) + \frac{1}{\sqrt{5}}(\phi^{k-1} - \psi^{k-1}) \\
        &= \frac{1}{\sqrt{5}} \left[ (\phi^k + \phi^{k-1}) - (\psi^k + \psi^{k-1}) \right] \\
        &= \frac{1}{\sqrt{5}} \left[ \phi^{k-1}(\phi + 1) - \psi^{k-1}(\psi + 1) \right]
        \end{align*}
        利用特征根的性质 $\phi + 1 = \phi^2$ 及 $\psi + 1 = \psi^2$, 代入上式得:
        \[
        a_{k+1} = \frac{1}{\sqrt{5}} \left[ \phi^{k-1}(\phi^2) - \psi^{k-1}(\psi^2) \right] 
        = \frac{1}{\sqrt{5}} (\phi^{k+1} - \psi^{k+1})
        \]
        此时命题 $P_{k+1}$ 成立, 由此蕴含 $(P_{k-1} \land P_k) \implies P_{k+1}$。

        由数学归纳法, 命题 $P_n$ 对所有正整数 $n$ 均成立。
    \end{solution}

    \question 已知数列由递推关系
    \[
    a_{n+1}=\frac{a_n-5}{3a_n-7},\quad a_1=-1,\quad n\in\mathbb{N}
    \]
    定义, 用数学归纳法证明
    \[
    a_n=\frac{2^{\,n+1}-5}{2^{\,n+1}-3}.
    \]
    \begin{solution}
        定义命题 
        \[
        P_n:a_n = \frac{2^{n+1}-5}{2^{n+1}-3} \quad, n \in \mathbb{N}
        \]
        基础步骤:当 $n=1$ 时,右式为 
        \[
        \frac{2^{1+1}-5}{2^{1+1}-3} = \frac{4-5}{4-3} = -1=a_1
        \]
        此时左式等于右式, 即 $P_1$ 成立。

        归纳步骤:设 $k \in \mathbb{N}$, 假设 $P_k$ 成立, 有
        \[
        a_k = \frac{2^{k+1}-5}{2^{k+1}-3}
        \]
        欲证 $P_{k+1}$ 成立, 即证 $a_{k+1} = \dfrac{2^{k+2}-5}{2^{k+2}-3}$,由递推关系及归纳假设:
        \begin{align*}
        a_{k+1} &= \frac{a_k - 5}{3a_k - 7} \\
        &= \frac{\frac{2^{k+1}-5}{2^{k+1}-3} - 5}{3\left(\frac{2^{k+1}-5}{2^{k+1}-3}\right) - 7} \\
        &= \frac{(2^{k+1}-5) - 5(2^{k+1}-3)}{3(2^{k+1}-5) - 7(2^{k+1}-3)} \\
        &= \frac{2^{k+1} - 5 - 5 \cdot 2^{k+1} + 15}{3 \cdot 2^{k+1} - 15 - 7 \cdot 2^{k+1} + 21} \\
        &= \frac{-4 \cdot 2^{k+1} + 10}{-4 \cdot 2^{k+1} + 6}\\
        &= \frac{2^{k+2} - 5}{2^{k+2} - 3}
        \end{align*}
        此时命题 $P_{k+1}$ 成立, 蕴含 $P_k \implies P_{k+1}$成立。

        由数学归纳法, 命题 $P_n$ 对所有正整数 $n$ 均成立。
    \end{solution}

    \question 证明对所有整数 $n \ge 1$, 有
    \[
    (a+b)^n = \comb{n}{0} a^n + \comb{n}{1} a^{n-1}b + \comb{n}{2} a^{n-2}b^2 + \dots + \comb{n}{n} b^n.
    \]
    \begin{solution}
        定义命题 
        \[
        P_n:(a+b)^n = \sum_{r=0}^{n} \comb{n}{r} a^{n-r}b^r \quad, n \in \mathbb{N}
        \]
        基础步骤:当 $n=1$ 时, 
        左式为 $(a+b)^1 = a+b$。
        右式为 $\comb{1}{0} a^1 + \comb{1}{1} b^1 = a+b$。
        此时左式等于右式, 即 $P_1$ 成立。

        归纳步骤:设 $k \in \mathbb{N}$, 假设归纳假设 $P_k$ 成立, 即
        \[
        (a+b)^k = \sum_{r=0}^{k} \comb{k}{r} a^{k-r}b^r
        \]
        欲证 $P_{k+1}$ 成立, 考虑其左式:
        \begin{align*}
        (a+b)^{k+1} &= (a+b)(a+b)^k \\
        &= (a+b) \sum_{r=0}^{k} \comb{k}{r} a^{k-r}b^r \\
        &= a \sum_{r=0}^{k} \comb{k}{r} a^{k-r}b^r + b \sum_{r=0}^{k} \comb{k}{r} a^{k-r}b^r \\
        &= \sum_{r=0}^{k} \comb{k}{r} a^{k-r+1}b^r + \sum_{r=0}^{k} \comb{k}{r} a^{k-r}b^{r+1}
        \end{align*}
        在第一个求和中提取 $r=0$ 的项, 在第二个求和中提取 $r=k$ 的项, 并对剩余部分合并:
        \begin{align*}
        (a+b)^{k+1} &= \comb{k}{0} a^{k+1} + \sum_{r=1}^{k} \comb{k}{r} a^{k-r+1}b^r + \sum_{r=0}^{k-1} \comb{k}{r} a^{k-r}b^{r+1} + \comb{k}{k} b^{k+1} \\
        &= \comb{k}{0} a^{k+1} + \sum_{r=1}^{k} \comb{k}{r} a^{k-r+1}b^r + \sum_{r=1}^{k} \comb{k}{r-1} a^{k-r+1}b^r + \comb{k}{k} b^{k+1} \\
        &= \comb{k}{0} a^{k+1} + \sum_{r=1}^{k} \left( \comb{k}{r} + \comb{k}{r-1} \right) a^{k-r+1}b^r + \comb{k}{k} b^{k+1}
        \end{align*}
        利用帕斯卡恒等式 
        \[
        \comb{k}{r} + \comb{k}{r-1} = \comb{k+1}{r}, 
        \]
        且注意到 $\comb{k}{0} = \comb{k+1}{0} = 1$, $\comb{k}{k} = \comb{k+1}{k+1} = 1$, 得:
        \[
        (a+b)^{k+1} = \comb{k+1}{0} a^{k+1} + \sum_{r=1}^{k} \comb{k+1}{r} a^{k+1-r}b^r + \comb{k+1}{k+1} b^{k+1} = \sum_{r=0}^{k+1} \comb{k+1}{r} a^{k+1-r}b^r
        \]
        此时命题 $P_{k+1}$ 成立, 蕴含 $P_k \implies P_{k+1}$成立。

        由数学归纳法, 命题 $P_n$ 对所有正整数 $n$ 均成立。
    \end{solution}

    \question 设数列 $(a_n)_{n=0}^\infty$ 满足对所有 $n \ge 0$ 都有 $\dfrac{1}{2} < a_n < 1$。定义数列 $(x_n)$ 为
    \[
    x_0=a_0,\quad x_{n+1}=\frac{a_{n+1}+x_n}{1+a_{n+1}x_n} \quad (n \ge 0).
    \]
    求 $\displaystyle \lim_{n \to \infty} x_n$ 的所有可能取值, 并判断该数列是否可能发散。
    \begin{solution}
        定义命题 
        \[
        P_n: 0 < 1 - x_n < \frac{1}{2^{n+1}} \quad, n \in \mathbb{N} \cup \{0\}
        \]
        基础步骤: 当 $n=0$ 时, 
        由题设 $\frac{1}{2} < a_0 < 1$ 且 $x_0 = a_0$, 有
        \[
        1 - x_0 = 1 - a_0 < 1 - \frac{1}{2} = \frac{1}{2}.
        \]
        又因 $a_0 < 1$, 故 $1 - x_0 > 0$。
        此时 $0 < 1 - x_0 < \frac{1}{2^{0+1}}$ 成立, 故 $P_0$ 成立。

        归纳步骤: 设 $k \in \mathbb{N} \cup \{0\}$, 假设归纳假设 $P_k$ 成立, 即
        \[
        0 < 1 - x_k < \frac{1}{2^{k+1}}
        \]
        欲证 $P_{k+1}$ 成立, 考虑其表达式:
        \begin{align*}
        1 - x_{k+1} &= 1 - \frac{a_{k+1} + x_k}{1 + a_{k+1}x_k} \\
        &= \frac{1 + a_{k+1}x_k - a_{k+1} - x_k}{1 + a_{k+1}x_k} \\
        &= \frac{(1 - a_{k+1}) - x_k(1 - a_{k+1})}{1 + a_{k+1}x_k} \\
        &= \frac{1 - a_{k+1}}{1 + a_{k+1}x_k} (1 - x_k)
        \end{align*}
        由于 $\frac{1}{2} < a_{k+1} < 1$ 且由 $P_k$ 知 $0 < x_k < 1$, 观察系数项:
        \[
        0 < \frac{1 - a_{k+1}}{1 + a_{k+1}x_k} < \frac{1 - \frac{1}{2}}{1 + \frac{1}{2} \cdot 0} = \frac{1}{2}
        \]
        由归纳假设,得
        \[
        0 < 1 - x_{k+1} < \frac{1}{2} \cdot (1 - x_k) < \frac{1}{2} \cdot \frac{1}{2^{k+1}} = \frac{1}{2^{k+2}}
        \]
        此时命题 $P_{k+1}$ 成立, 蕴含 $P_k \implies P_{k+1}$ 成立。由数学归纳法, 命题 $P_n$ 对所有非负整数 $n$ 均成立。

        由于对所有 $n,0 < 1 - x_n < \dfrac{1}{2^{n+1}}$, 且 
        \[
        \lim_{n \to \infty} \frac{1}{2^{n+1}} = 0, 
        \]
        由夹挤定理,
        \[
        \lim_{n \to \infty} (1 - x_n) = 0 \implies \lim_{n \to \infty} x_n = 1.
        \]
        因此, 极限的所有可能取值仅为 $1$, 且该数列在给定条件下必然收敛, 不可能发散。
    \end{solution}

    \question 对于实数 $a>0$, 定义数列 $\{x_n\}$ 为
    \[
    x_{n+1}=a(x_n^2+4),\quad x_0=0.
    \]
    求 $\displaystyle \lim_{n\to\infty}x_n$ 存在且有限的充要条件。
    \begin{solution}
        首先分析极限存在的必要条件。若 $\displaystyle \lim_{n\to\infty}x_n=x$ 存在, 对递推式两边取极限得
        \[
        x=a(x^2+4) \implies ax^2 - x + 4a = 0
        \]
        根据判别式 
        \[
        \Delta = (-1)^2 - 4(a)(4a) = 1 - 16a^2 \ge 0, 
        \]
        且 $a>0$, 得到 $0 < a \le \dfrac{1}{4}$。在此条件下, 方程的根为 
        \[
        x = \frac{1 \pm \sqrt{1-16a^2}}{2a}
        \]

        下面证明当 $0 < a \le \dfrac{1}{4}$ 时, 数列收敛。令 
        \[
        L = \frac{1 - \sqrt{1-16a^2}}{2a}
        \]
        为较小的根, 注意到 $L = a(L^2+4)$ 且 $L \le 2$ (当 $a=\dfrac{1}{4}$ 时 $L=2$)。

        定义命题 
        \[
        P_n: x_n \le x_{n+1} < L \quad, n \in \mathbb{N} \cup \{0\}
        \]
        基础步骤: 当 $n=0$ 时,$x_0 = 0,x_1 = a(0^2+4) = 4a$。由于 $0 < a \le \frac{1}{4}$, 则 
        \[
        x_0 = 0 \le x_1 = 4a \le 1 < L 
        \]
        故 $P_0$ 成立。

        归纳步骤: 设 $k \in \mathbb{N} \cup \{0\}$, 假设归纳假设 $P_k$ 成立, 即
        \[
        x_k \le x_{k+1} < L
        \]
        欲证 $P_{k+1}$ 成立, 考虑
        \begin{align*}
        x_{k+2} - x_{k+1} &= a(x_{k+1}^2 + 4) - a(x_k^2 + 4) \\
        &= a(x_{k+1}^2 - x_k^2) \\
        &= a(x_{k+1} + x_k)(x_{k+1} - x_k)
        \end{align*}
        由 $a>0$ 及 $P_k$ 知 $x_{k+1} - x_k \ge 0$ 且 $x_k, x_{k+1} \ge 0$, 故 $x_{k+2} - x_{k+1} \ge 0$。再考虑上界:由于 $f(t) = a(t^2+4)$ 在 $t \ge 0$ 时单调递增, 且由 $P_k$ 知 $x_{k+1} < L$, 则:
        \[
        x_{k+2} = a(x_{k+1}^2 + 4) < a(L^2 + 4) = L
        \]
        此时命题 $P_{k+1}$ 成立, 蕴含 $P_k \implies P_{k+1}$ 成立。

        由数学归纳法, 命题 $P_n$ 对所有非负整数 $n$ 均成立。

        因此, 数列 $\{x_n\}$ 单调递增且有上界 $L$, 根据单调有界原理, 极限 $\displaystyle \lim_{n\to\infty}x_n$ 必然存在且有限。

        综上所述, 所求的充要条件为 
        \[
        0 < a \le \frac{1}{4}
        \]
    \end{solution}

    \question 设 $n$ 为正整数, 且 $a_k, b_k$ 满足
    \[
    a_{k}=\frac{1}{\binom{n}{k}}, \quad b_{k}=2^{k-n}, \quad k=1,2,\dots,n.
    \]
    试证:
    \[
    \frac{a_1-b_1}{1}+\frac{a_2-b_2}{2}+\cdots+\frac{a_n-b_n}{n}=0 
    \]
    \begin{solution}
        欲证
        \[
        \sum_{k=1}^{n} \frac{1}{k \binom{n}{k}} = \sum_{k=1}^{n} \frac{2^{k-n}}{k}
        \]
        由于对所有 $k \ge 1$ 均有恒等式 
        \[
        k \binom{n}{k} = n \binom{n-1}{k-1},
        \]
        得等价命题:
        \[
        P_n:\frac{2^n}{n} \sum_{k=0}^{n-1} \frac{1}{\binom{n-1}{k}} = \sum_{k=1}^{n} \frac{2^k}{k}
        \]
        基础步骤:当 $n=1$ 时,左式为 $\frac{2^1}{1} \cdot \frac{1}{\binom{0}{0}} = 2$,右式为 $\frac{2^1}{1} = 2$。左式等于右式,$P_1$ 成立。

        归纳步骤:假设归纳假设 $P_n$ 成立, 则对于 $n+1$:
        \begin{align*}
        \frac{2^{n+1}}{n+1} \sum_{k=0}^{n} \frac{1}{\binom{n}{k}} 
        &= \frac{2^n}{n+1} \left( 1 + \sum_{k=0}^{n-1} \left( \frac{1}{\binom{n}{k}} + \frac{1}{\binom{n}{k+1}} \right) + 1 \right) \\
        &= \frac{2^n}{n+1} \sum_{k=0}^{n-1} \frac{\frac{n-k}{n}+\frac{k+1}{n}}{\binom{n-1}{k}} + \frac{2^{n+1}}{n+1} \\
        &= \frac{2^n}{n} \sum_{k=0}^{n-1} \frac{1}{ \binom{n-1}{k}} + \frac{2^{n+1}}{n+1} \\
        &= \frac{2^n}{n} \sum_{k=0}^{n-1} \frac{1}{\binom{n-1}{k}} + \frac{2^{n+1}}{n+1} \\
        &= \sum_{k=1}^{n+1} \frac{2^k}{k}
        \end{align*}
        即命题 $P_{n+1}$ 成立, 蕴含 $P_n \implies P_{n+1}$ 成立。

        由数学归纳法, 命题 $P_n$ 对所有正整数 $n$ 均成立。
    \end{solution}

    \question 定义数列 $\{x_n\}$ 递推关系为 $x_1 = \sqrt{5}$ 且 $x_{n+1} = x_n^2 - 2$ 对自然数 $n \ge 1$。求
    \[
    \lim_{n \to \infty} \frac{x_1 \cdot x_2 \cdot \cdots \cdot x_n}{x_{n+1}}.
    \]
    \begin{solution}
        令 $y_n = x_n^2$。根据递推关系 $x_{n+1} = x_n^2 - 2$,平方得 $y_{n+1} = (y_n - 2)^2$。
        由此可得关键恒等式:
        \[
        y_{n+1} - 4 = (y_n - 2)^2 - 4 = y_n^2 - 4y_n = y_n(y_n - 4).
        \]
        即有 
        \[
        \frac{y_n}{y_{n+1}-4} = \frac{1}{y_n-4}
        \]

        首先,我们用归纳法证明对所有 $n \ge 1$, 有 $y_n \ge 5$ 且 $y_n \to \infty$。
        定义命题 $P_n: y_n \ge 5$。
        基础步骤:当 $n=1$ 时,$y_1 = (\sqrt{5})^2 = 5 \ge 5$,故 $P_1$ 成立。
        归纳步骤:假设 $y_k \ge 5$ 成立,则 $y_{k+1} = (y_k - 2)^2 \ge (5-2)^2 = 9 > 5$。
        故由数学归纳法知 $y_n \ge 5$ 对所有 $n \ge 1$ 成立。
        进一步地,
        \[
        y_{n+1} - y_n = y_n^2 - 5y_n + 4 = (y_n - 1)(y_n - 4)
        \]
        由于 $y_n \ge 5$,则 $y_{n+1} - y_n \ge (5-1)(5-4) = 4$,故数列 $\{y_n\}$ 趋于无穷大。

        接下来,我们用归纳法证明命题 
        \[
        Q_n: \prod_{i=1}^n y_i = \frac{y_{n+1}-4}{y_1-4} \cdot y_1, \quad n \in \mathbb{N}\tag{1}
        \]
        基础步骤:当 $n=1$ 时,左式 $= y_1$；右式 $= \frac{y_2-4}{y_1-4} \cdot y_1 = \frac{y_1(y_1-4)}{y_1-4} \cdot y_1 = y_1$。$Q_1$ 成立。
        归纳步骤:假设 $Q_k$ 成立,则对于 $n=k+1$:
        \begin{align*}
        \prod_{i=1}^{k+1} y_i &= \left( \prod_{i=1}^{k} y_i \right) \cdot y_{k+1} \\
        &= \frac{y_{k+1}-4}{y_1-4} \cdot y_1 \cdot y_{k+1} \\
        &= \frac{y_1}{y_1-4} \cdot [y_{k+1}(y_{k+1}-4)] \\
        &= \frac{y_1}{y_1-4} (y_{k+2}-4)
        \end{align*}
        此时 $Q_{k+1}$ 成立。由数学归纳法知(1)对所有 $n \in \mathbb{N}$ 成立。

        现在计算目标极限的平方
        \[
        \left( \frac{x_1 x_2 \cdots x_n}{x_{n+1}} \right)^2 = \frac{\prod_{i=1}^n y_i}{y_{n+1}} 
        = \frac{y_1}{y_1-4} \cdot \frac{y_{n+1}-4}{y_{n+1}} 
        = 5 \left( 1 - \frac{4}{y_{n+1}} \right)
        \]
        当 $n \to \infty$ 时,$y_{n+1} \to \infty$,故
        \[
        \lim_{n \to \infty} \left( \frac{x_1 x_2 \cdots x_n}{x_{n+1}} \right)^2 = 5(1 - 0) = 5
        \]
        由于 $x_n$ 显然为正,故
        \[
        \lim_{n \to \infty} \frac{x_1 \cdot x_2 \cdot \cdots \cdot x_n}{x_{n+1}} = \sqrt{5}
        \]
    \end{solution}

    \question 设数列 $\{a_n\}_{n=1}^{\infty}$ 定义为 $a_1=1, a_2=1, a_3=2$,并且当 $n\ge3$ 时,满足递推式:
    \[
    a_{n+1}=\frac{n a_n a_{n-2}}{a_{n-1}}.
    \]
    \begin{parts}
    \part 证明对所有 $n\ge1$, $a_n$ 都是正整数。
    \begin{solution}
        首先推导数列的一个更简便的递推关系。
        由题设, 当 $n \ge 3$ 时, 有 $a_{n+1} a_{n-1} = n a_n a_{n-2}$。
        则对于 $n+1$ 项, 有:
        \[
        a_{n+2} a_{n} = (n+1) a_{n+1} a_{n-1} = (n+1) \cdot n a_n a_{n-2} \implies a_{n+2} = n(n+1) a_{n-2}.
        \]
        将下标平移, 得到对所有 $n \ge 1$:
        \[
        a_{n+4} = (n+3)(n+2) a_n. \tag{1}
        \]

        下面用数学归纳法证明 $a_n$ 均为正整数。
        定义命题 $P_n: a_{4n-3}, a_{4n-2}, a_{4n-1}, a_{4n}$ 均为正整数。

        基础步骤: 当 $n=1$ 时, 
        $a_1 = 1, a_2 = 1, a_3 = 2$ 为已知正整数。
        由递推式 $a_4 = \frac{3 a_3 a_1}{a_2} = \frac{3 \cdot 2 \cdot 1}{1} = 6$, 亦为正整数。
        故 $P_1$ 成立。

        归纳步骤: 设 $k \in \mathbb{N}$, 假设归纳假设 $P_k$ 成立, 即 $a_{4k-3}, a_{4k-2}, a_{4k-1}, a_{4k}$ 均为正整数。
        考虑 $n=k+1$ 时的项, 根据公式 (1):
        \begin{align*}
        a_{4(k+1)-3} &= a_{4k+1} = (4k)(4k-1) a_{4k-3} \\
        a_{4(k+1)-2} &= a_{4k+2} = (4k+1)(4k) a_{4k-2} \\
        a_{4(k+1)-1} &= a_{4k+3} = (4k+2)(4k+1) a_{4k-1} \\
        a_{4(k+1)}   &= a_{4k+4} = (4k+3)(4k+2) a_{4k}
        \end{align*}
        由于 $a_{4k-3}, \dots, a_{4k}$ 为正整数, 且系数 $(4k+i)$ 均为正整数, 故它们的积也必为正整数。
        即命题 $P_{k+1}$ 成立, 蕴含 $P_k \implies P_{k+1}$ 成立。

        由数学归纳法, 命题 $P_n$ 对所有 $n \in \mathbb{N}$ 成立, 故对所有 $n \ge 1$, $a_n$ 均为正整数。
    \end{solution}
    \part 定义数列 $\{b_n\}_{n=1}^{\infty}$ 为 
    \[
    b_n=\frac{a_n}{\sqrt{(n+1)!}}.
    \]
    证明 $\{b_n\}$ 有界。
    \begin{solution}
        由 (3) 式,有
        \begin{align*}
        b_{n+4}
        &=\frac{a_{n+4}}{\sqrt{(n+5)!}} \\
        &=\frac{(n+3)(n+2)a_n}{\sqrt{(n+1)!(n+2)(n+3)(n+4)(n+5)}} \\
        &=\frac{a_n}{\sqrt{(n+1)!}}
        \sqrt{\frac{(n+3)(n+2)}{(n+4)(n+5)}} \\
        &\le \frac{a_n}{\sqrt{(n+1)!}} \\
        &=b_n
        \end{align*}
        即归纳步骤成立。注意到基础步骤也成立 $b_1,b_2,b_3,b_4\le1$,由数学归纳法,可得对所有 $n\ge1$,
        \[
        b_n\le1
        \]
        因此数列 $\{b_n\}$ 有界。
    \end{solution}
    \end{parts}

    \question 设 $x_1=1$, 且对 $n \ge 1$,
    \[
    x_{n+1}=\frac{1}{x_n}\left(\sqrt{1+x_n^2}-1\right).
    \]
    证明数列 $\{2^n x_n\}$ 收敛, 并求其极限。
    \begin{solution}
        定义命题 
        \[
        P_n: x_n = \tan\left(\frac{\pi}{2^{n+1}}\right) \quad, n \in \mathbb{N}
        \]
        基础步骤: 当 $n=1$ 时, 
        左式为 $x_1 = 1$。
        右式为 $\tan\left(\frac{\pi}{2^{1+1}}\right) = \tan\left(\frac{\pi}{4}\right) = 1$。
        此时左式等于右式, 即 $P_1$ 成立。

        归纳步骤: 设 $k \in \mathbb{N}$, 假设归纳假设 $P_k$ 成立, 即
        \[
        x_k = \tan\left(\frac{\pi}{2^{k+1}}\right)
        \]
        欲证 $P_{k+1}$ 成立, 考虑其递推表达式:
        \begin{align*}
        x_{k+1} &= \frac{1}{x_k}\left(\sqrt{1+x_k^2}-1\right) \\
        &= \frac{1}{\tan\left(\frac{\pi}{2^{k+1}}\right)}\left(\sqrt{1+\tan^2\left(\frac{\pi}{2^{k+1}}\right)}-1\right) \\
        &= \frac{\sec\left(\frac{\pi}{2^{k+1}}\right)-1}{\tan\left(\frac{\pi}{2^{k+1}}\right)} \\
        &= \frac{1-\cos\left(\frac{\pi}{2^{k+1}}\right)}{\sin\left(\frac{\pi}{2^{k+1}}\right)}
        \end{align*}
        由半角公式 
        \[
        \tan\left(\frac{\theta}{2}\right) = \frac{1-\cos\theta}{\sin\theta},
        \]
        可得
        \[
        x_{k+1} = \tan\left(\frac{1}{2} \cdot \frac{\pi}{2^{k+1}}\right) = \tan\left(\frac{\pi}{2^{k+2}}\right)
        \]
        此时命题 $P_{k+1}$ 成立, 蕴含 $P_k \implies P_{k+1}$ 成立。

        由数学归纳法, 命题 $P_n$ 对所有正整数 $n$ 均成立。

        接下来求数列 $\{2^n x_n\}$ 的极限:
        \[
        \lim_{n\to\infty} 2^n x_n = \lim_{n\to\infty} 2^n \tan\left(\frac{\pi}{2^{n+1}}\right)
        \]
        利用等价无穷小 $\tan \theta \sim \theta$ (当 $\theta \to 0$ 时):
        由于当 $n \to \infty$ 时, $\frac{\pi}{2^{n+1}} \to 0$, 则
        \begin{align*}
        \lim_{n\to\infty} 2^n \tan\left(\frac{\pi}{2^{n+1}}\right) &= \lim_{n\to\infty} 2^n \cdot \frac{\pi}{2^{n+1}} \\
        &= \lim_{n\to\infty} \frac{\pi}{2} \\
        &= \frac{\pi}{2}
        \end{align*}
        因此数列 $\{2^n x_n\}$ 收敛, 其极限为 $\dfrac{\pi}{2}$。
    \end{solution}

    \question 设无穷数列$\{a_n\}$符合 $a_0=0$ 且当 $n \ge 1$ 时,
    \[
    a_n - a_{n-1} =
    \begin{cases}
    \left(\dfrac{1}{5}\right)^n, & n \:\text{为偶数} \\[2mm]
    \left(\dfrac{1}{5}\right)^n - \left(\dfrac{1}{3}\right)^n, & n \:\text{为奇数}
    \end{cases}
    \]
    \begin{parts}
    \part 求 $\lim_{n \to \infty} a_{2n}$ 的值。
    \begin{solution}
        由于
        \begin{align*}
        a_{2n} - a_0 
        &= a_{2n} - a_{2n-1}+a_{2n-1} - a_{2n-2}+ \cdots +a_2 - a_1+a_1 - a_0 \\
        &= \left(\frac15 + \frac1{5^2} + \cdots + \frac1{5^{2n}}\right)
        - \left(\frac13 + \frac1{3^3} + \cdots + \frac1{3^{2n-1}}\right)
        \end{align*}
        所以
        \[
        \lim_{n\to\infty} a_{2n} = \frac{\frac15}{1 - \frac15} - \frac{\frac13}{1 - \frac19}
        = -\frac18
        \]
    \end{solution}
    \part 证明当 $n \ge 0,a_{2n+2} - a_{2n} < 0$,并依此证明对于所有正整数 $n$,不等式 $$-\frac{1}{8} \le a_{2n} < 0$$ 恒成立。
    \begin{solution}
        首先有
        \begin{align*} 
        a_{2n+2} - a_{2n} &=a_{2n+2} - a_{2n+1}+a_{2n+1} - a_{2n}\\
        &=\left(\frac15\right)^{2n+2}+\left(\frac15\right)^{2n+1} - \left(\frac13\right)^{2n+1}\\
        &= \frac{\frac65 \cdot 3^{2n+1} - 5^{2n+1}}{15^{2n+1}}
        \end{align*}
        当 $n=0$ 时,$\:a_2 - a_0 = \frac{18/5 - 5}{15} < 0$,成立;假设 $n=k$ 时成立,即
        \[
        \frac65 \cdot 3^{2k+1} < 5^{2k+1}
        \]
        观察当 $n=k+1$ 时,
        \[
        a_{2k+4} - a_{2k+2} = \frac{\frac65 \cdot 3^{2k+3} - 5^{2k+3}}{15^{2k+3}}
        \]
        由归纳假设,
        \[
        \frac65 \cdot 3^{2k+3} < 3^2 \cdot 5^{2k+1} < 5^2 \cdot 5^{2k+1} = 5^{2k+3}
        \]
        因此 $a_{2k+4} - a_{2k+2} < 0$,即当 $n=k+1$ 时也成立;由数学归纳法,$\:a_{2n+2} - a_{2n} < 0$;又 $$a_{2n} < a_{2n-2} < \cdots < a_2 < a_0 = 0,$$故 $a_{2n} < 0$,结合 $\displaystyle \lim_{n\to\infty} a_{2n} = -\frac18$得
        \[
        -\frac18 \le a_{2n} < 0,\quad n \in \mathbb{N}
        \]
    \end{solution}
    \end{parts}

\question 设正实数 $a_0, a_1, \dots, a_n$ 满足 $a_{k+1} - a_k \ge 1$ 对所有 $k = 0,1,\dots,n-1$ 成立。证明
\[
1 + \frac{1}{a_0}\prod_{k=1}^{n} \left(1+\frac{1}{a_k-a_0}\right) \le \prod_{k=0}^{n} \left(1+\frac{1}{a_k}\right).
\]
%IMC 2010 D2P2
\begin{solution}
    定义命题 
\[
P_n: 1 + \frac{1}{a_0} \prod_{i=1}^n \left(1 + \frac{1}{a_i - a_0}\right) \le \prod_{i=0}^n \left(1 + \frac{1}{a_i}\right), \quad n \in \mathbb{N} \cup \{0\}
\]
其中 $a_k$ 是正实数且满足 $a_{k+1} - a_k \ge 1$。

基础步骤: 当 $n=0$ 时, 
左式为 $1 + \frac{1}{a_0} \cdot (1) = 1 + \frac{1}{a_0}$（此处空乘积定义为 1）。
右式为 $1 + \frac{1}{a_0}$。
此时左式等于右式, 即 $P_0$ 成立。

归纳步骤: 设 $n \in \mathbb{N} \cup \{0\}$, 假设归纳假设 $P_n$ 成立。
欲证 $P_{n+1}$ 成立, 即证明：
\[
1 + \frac{1}{a_0} \prod_{i=1}^{n+1} \left(1 + \frac{1}{a_i - a_0}\right) \le \prod_{i=0}^{n+1} \left(1 + \frac{1}{a_i}\right)
\]
注意到上式左边可以拆分为 $P_n$ 的左项与增量项。根据 $P_n$ 已成立，只需证明增量部分的差异满足：
\[
\frac{1}{a_0} \prod_{i=1}^{n} \left(1 + \frac{1}{a_i - a_0}\right) \cdot \frac{1}{a_{n+1} - a_0} \le \prod_{i=0}^{n} \left(1 + \frac{1}{a_i}\right) \cdot \frac{1}{a_{n+1}} \quad \cdots (1)
\]
针对不等式 (1)，定义子命题 $Q_n$ 并利用第二数学归纳法证明：

基础步骤: 当 $n=0$ 时, (1) 式变为:
\begin{align*}
\frac{1}{a_0} \cdot \frac{1}{a_1 - a_0} &\le \left(1 + \frac{1}{a_0}\right) \frac{1}{a_1} \\
a_1 &\le (a_0 + 1)(a_1 - a_0) \\
a_0 &\le a_0 a_1 - a_0^2 \\
1 &\le a_1 - a_0
\end{align*}
由已知条件 $a_1 - a_0 \ge 1$ 知其成立。

归纳步骤: 假设 $Q_n$ 成立, 欲证 $Q_{n+1}$ 成立, 需证明：
\[
\left(1 + \frac{1}{a_{n+1} - a_0}\right) \frac{a_{n+1} - a_0}{a_{n+2} - a_0} \le \left(1 + \frac{1}{a_{n+1}}\right) \frac{a_{n+1}}{a_{n+2}}
\]
展开整理得：
\begin{align*}
(a_{n+1} - a_0 + 1) a_{n+2} &\le (a_{n+1} + 1)(a_{n+2} - a_0) \\
a_0 &\le a_0 a_{n+2} - a_0 a_{n+1} \\
1 &\le a_{n+2} - a_{n+1}
\end{align*}
由于已知 $a_{n+2} - a_{n+1} \ge 1$, 故 $Q_{n+1}$ 成立。

由此可知 (1) 式恒成立，进而蕴含 $P_n \implies P_{n+1}$ 成立。

由数学归纳法, 命题 $P_n$ 对所有非负整数 $n$ 均成立。
\textcolor{red}{(待验证)}
\end{solution}

    \question 用数学归纳法证明:对所有正整数 $n$, 有
    \[
    \sum_{r=1}^{n} r(r+1)\left(\frac{1}{2}\right)^{r-1} = 16 - \left(\frac{1}{2}\right)^{n-1}(n^2+5n+8).
    \]
    \begin{solution}
        定义命题
        \[
        P_n:\sum_{r=1}^{n} r(r+1)\left(\frac{1}{2}\right)^{r-1} = 16 - \left(\frac{1}{2}\right)^{n-1}(n^2+5n+8), \quad n \in \mathbb{N}.
        \]
        基础步骤:当 $n=1$ 时, 左式为
        \[
        \sum_{r=1}^{1} r(r+1)\left(\frac{1}{2}\right)^{r-1} = 1 \cdot 2 \cdot \left(\frac{1}{2}\right)^0 = 2.
        \]
        而右式为
        \[
        16 - \left(\frac{1}{2}\right)^{1-1}(1^2 + 5 \cdot 1 + 8) = 16 - 1 \cdot 14 = 2.
        \]
        此时左式等于右式, 即命题 $P_1$ 成立。

        归纳步骤:设 $k \in \mathbb{N}$, 假设归纳假设 $P_k$ 成立, 有
        \[
        \sum_{r=1}^{k} r(r+1)\left(\frac{1}{2}\right)^{r-1} = 16 - \left(\frac{1}{2}\right)^{k-1}(k^2+5k+8).
        \]
        观察当 $n=k+1$ 时,
        \begin{align*}
        \sum_{r=1}^{k+1} r(r+1)\left(\frac{1}{2}\right)^{r-1} &= \sum_{r=1}^{k} r(r+1)\left(\frac{1}{2}\right)^{r-1} + (k+1)(k+2)\left(\frac{1}{2}\right)^k \\
        &= 16 - \left(\frac{1}{2}\right)^{k-1}(k^2+5k+8) + (k+1)(k+2)\left(\frac{1}{2}\right)^k \\
        &= 16 + \left(\frac{1}{2}\right)^k \left[ (k+1)(k+2) - 2(k^2+5k+8) \right] \\
        &= 16 - \left(\frac{1}{2}\right)^k (k^2+7k+14) \\
        &= 16 - \left(\frac{1}{2}\right)^{(k+1)-1} \left[ (k+1)^2 + 5(k+1) + 8 \right]
        \end{align*}
        发现命题 $P_{k+1}$ 成立, 此时蕴含 $P_k \implies P_{k+1}$ 成立。

        由数学归纳法, 命题 $P_n$ 对所有正整数 $n$ 皆成立。
    \end{solution}

    \question 对每个 $n \ge 1$, 令
    \[
    a_n = \sum_{k=0}^\infty \frac{k^n}{k!}, \quad b_n = \sum_{k=0}^\infty \frac{(-1)^k k^n}{k!}.
    \]
    证明 $a_n \cdot b_n$ 是整数。
    \begin{solution}
        定义命题 $P_n$ 为:对给定的 $n \ge 0$, 存在整数 $A_n, B_n$ 使得 $a_n = A_n e$ 且 $b_n = B_n e^{-1}$ (约定 $0^0=1$)。

        基础步骤:当 $n=0$ 时, 由指数函数的幂级数展开 $e^x = \sum_{k=0}^\infty \frac{x^k}{k!}$ 可知
        \[
        a_0 = \sum_{k=0}^\infty \frac{1}{k!} = e, \quad b_0 = \sum_{k=0}^\infty \frac{(-1)^k}{k!} = e^{-1}.
        \]
        此时 $A_0 = 1, B_0 = 1$ 均为整数, 即命题 $P_0$ 成立。

        归纳步骤:假设归纳假设对所有 $m \le n$ 均成立, 即 $a_m = A_m e, b_m = B_m e^{-1}$ 且 $A_m, B_m \in \mathbb{Z}$。
        对于 $a_{n+1}$, 利用指标变换 $j = k-1$:
        \begin{align*}
        a_{n+1} &= \sum_{k=1}^\infty \frac{k^{n+1}}{k!} = \sum_{k=1}^\infty \frac{k^n}{(k-1)!} = \sum_{j=0}^\infty \frac{(j+1)^n}{j!} \\
        &= \sum_{j=0}^\infty \frac{1}{j!} \sum_{m=0}^n \binom{n}{m} j^m = \sum_{m=0}^n \binom{n}{m} \left( \sum_{j=0}^\infty \frac{j^m}{j!} \right) \\
        &= \sum_{m=0}^n \binom{n}{m} a_m.
        \end{align*}
        根据归纳假设 $a_m = A_m e$, 则 $a_{n+1} = \left( \sum_{m=0}^n \binom{n}{m} A_m \right) e$。由于 $A_m$ 与组合数均为整数, 故其和 $A_{n+1}$ 亦为整数。
        同理, 对于 $b_{n+1}$:
        \begin{align*}
        b_{n+1} &= \sum_{k=1}^\infty \frac{(-1)^k k^{n+1}}{k!} = \sum_{k=1}^\infty \frac{(-1)^k k^n}{(k-1)!} = \sum_{j=0}^\infty \frac{(-1)^{j+1} (j+1)^n}{j!} \\
        &= -\sum_{j=0}^\infty \frac{(-1)^j}{j!} \sum_{m=0}^n \binom{n}{m} j^m = -\sum_{m=0}^n \binom{n}{m} b_m.
        \end{align*}
        根据归纳假设 $b_m = B_m e^{-1}$, 则 $b_{n+1} = \left( -\sum_{m=0}^n \binom{n}{m} B_m \right) e^{-1}$。显然 $B_{n+1}$ 亦为整数。
        由此命题 $P_{n+1}$ 成立, 此时蕴含 $P_0, \dots, P_n \implies P_{n+1}$ 成立。

        由数学归纳法, 命题 $P_n$ 对所有自然数 $n$ 皆成立。此时有
        \[
        a_n b_n = (A_n e)(B_n e^{-1}) = A_n B_n.
        \]
        由于 $A_n, B_n$ 均为整数, 故其积 $a_n b_n$ 是整数。
    \end{solution}

    \question 试证对任意正整数 $n$, 
    \[
    \lfloor\sqrt{n}\rfloor + \sum_{k=1}^{n} \left\lfloor\frac{n}{k}\right\rfloor
    \]
    必为偶数, 其中 $\lfloor x \rfloor$ 为向下取整函数。
    \begin{solution}
        定义命题 
        \[
        P_n:a_n = \lfloor\sqrt{n}\rfloor + \sum_{k=1}^n \left\lfloor\frac{n}{k}\right\rfloor \text{ 为偶数}, \quad n \in \mathbb{Z}^+.
        \]
        基础步骤:观察当 $n=1$ 时, 
        \[
        a_1 = \lfloor\sqrt{1}\rfloor + \left\lfloor\frac{1}{1}\right\rfloor = 1 + 1 = 2
        \]
        为偶数, 即 $P_1$ 成立。

        归纳步骤:假设 $P_{\ell}$ 成立, 即当 $n=\ell$ 时, 
        \[
        a_{\ell} = \lfloor\sqrt{\ell}\rfloor + \sum_{k=1}^{\ell} \lfloor\frac{\ell}{k}\rfloor
        \]
        为偶数。观察当 $n=\ell+1$ 时, $\sum_{k=1}^{n} \lfloor\frac{n}{k}\rfloor$ 在数值上等于从 $1$ 到 $n$ 所有正整数的正因数个数之和。因此:
        \[
        \sum_{k=1}^{\ell+1} \left\lfloor\frac{\ell+1}{k}\right\rfloor = \sum_{k=1}^{\ell} \left\lfloor\frac{\ell}{k}\right\rfloor + d(\ell+1)
        \]
        其中 $d(\ell+1)$ 表示 $\ell+1$ 的正因数个数。分以下两种情况讨论:
        \begin{itemize}
            \item 若 $\ell+1$ 为完全平方数, 则存在正整数 $m$ 使得 $\ell+1 = m^2$。此时有
            \[
            \lfloor\sqrt{\ell+1}\rfloor = m = \lfloor\sqrt{\ell}\rfloor + 1.
            \]
            由于完全平方数的正因数个数必为奇数, 即 $d(\ell+1)$ 为奇数。于是
            \[
            a_{\ell+1} = (\lfloor\sqrt{\ell}\rfloor + 1) + \left( \sum_{k=1}^{\ell} \left\lfloor\frac{\ell}{k}\right\rfloor + d(\ell+1) \right) = a_\ell + 1 + d(\ell+1).
            \]
            因为 $a_\ell$ 为偶数且 $1 + d(\ell+1)$ 为偶数, 故 $a_{\ell+1}$ 为偶数。

            \item 若 $\ell+1$ 不是完全平方数, 则有
            \[
            \lfloor\sqrt{\ell+1}\rfloor = \lfloor\sqrt{\ell}\rfloor.
            \]
            且非完全平方数的正因数个数必为偶数, 即 $d(\ell+1)$ 为偶数。于是
            \[
            a_{\ell+1} = \lfloor\sqrt{\ell}\rfloor + \left( \sum_{k=1}^{\ell} \left\lfloor\frac{\ell}{k}\right\rfloor + d(\ell+1) \right) = a_\ell + d(\ell+1).
            \]
            因为 $a_\ell$ 为偶数且 $d(\ell+1)$ 为偶数, 故 $a_{\ell+1}$ 为偶数。
        \end{itemize}
        综上所述, 无论何种情况均有 $a_{\ell+1}$ 为偶数, 即命题 $P_{\ell+1}$ 成立, 此时蕴含 $P_\ell \implies P_{\ell+1}$ 成立。

        由数学归纳法, 命题 $P_n$ 对任意正整数 $n$ 都成立。
    \end{solution}

    \question 票箱中有甲、乙两人的选票分别为 $m$ 张和 $n$ 张且 $m>n$。令 $P_{m,n}$ 表示开票的过程中甲的选票会一路领先乙的选票的概率,
    \begin{parts}
    \part 计算 $P_{m,1}$ 和 $P_{m,2}$。
    \begin{solution}
        视甲得$1$票为向上走一步,乙得$1$票为向右走一步,
        则 $P_{m,n}$ 表示从原点 $O(0,0)$ 走格点至 $P(n,m)$ 但不经过 $(a,a)$ 的概率。

        $m$ 个上、1 个右的排列数为 $\dfrac{(m+1)!}{m!}=m+1$,其中
        \[
        \begin{cases}
        \text{开头为右的数列有 1 个}\\
        \text{开头为上右的数列也只有 1 个}
        \end{cases}
        \]
        因此符合条件的数列有 $m+1-2=m-1$ 个,故$$P_{m,1}=\frac{m-1}{m+1}$$

        $m$ 个上、2 个右的排列数为 $\dfrac{(m+2)!}{m!2!}=\dfrac{(m+2)(m+1)}{2}$,其中
        \[
        \begin{cases}
        \text{开头为右的数列有 $m+1$ 个}\\
        \text{开头为上右的数列有 $m$ 个}\\
        \text{开头为上上右右的数列有 1 个}
        \end{cases}
        \]
        因此符合条件的数列有
        \[
        \frac{(m+2)(m+1)}{2}-(m+1)-m-1=\frac{(m+1)(m-2)}{2}
        \]
        所以 $$P_{m,2}=\frac{m-2}{m+2}$$
    \end{solution}
    \part 证明 $P_{m,n} = \dfrac{m}{m+n} P_{m-1,n} + \dfrac{n}{m+n} P_{m,n-1}$。
    \begin{solution}
        从 $(0,0)$ 至 $(n,m)$ 的方法数等于从 $(0,0)$ 至 $(n-1,m)$ 的方法数加上从 $(0,0)$ 至 $(n,m-1)$ 的方法数,因此
        \[
        P_{m,n}=\frac{\frac{(m+n-1)!}{(m-1)!n!}P_{m-1,n}+\frac{(m+n-1)!}{m!(n-1)!}P_{m,n-1}}{\frac{(m+n)!}{m!n!}}=\frac{m}{m+n}P_{m-1,n}+\frac{n}{m+n}P_{m,n-1}
        \]
    \end{solution}
    \part 先猜测 $P_{m,n}$ 的答案,再利用(b)用归纳法证明你的猜测。
    \begin{solution}
        由(a)可猜 $$P_{m,n}=\frac{m-n}{m+n},\ m\ge n$$
        令 $k=m+n$,观察当 $k=2$ 时,
        \[
        \begin{cases}
        m=2,n=0 \Rightarrow P_{2,0}=\frac{2-0}{2+0}=1\\
        m=n=1 \Rightarrow P_{1,1}=\frac{1-1}{1+1}=0
        \end{cases}
        \]
        显然猜测成立。现假设 $k=N$ 时猜测成立,观察当 $k=N+1$ 时,
        \begin{align*} 
        P_{m,n}&=\frac{m}{m+n}P_{m-1,n}+\frac{n}{m+n}P_{m,n-1}\\
        &=\frac{m}{m+n}\cdot\frac{m-n-1}{m+n-1}+\frac{n}{m+n}\cdot\frac{m-n+1}{m+n-1}\\
        &=\frac{(m-n)(m+n-1)}{(m+n)(m+n-1)}=\frac{m-n}{m+n}
        \end{align*}
        猜测亦成立,故猜测对所有$m,n\in \mathbb{N}\cup\{0\}$且$m\ge n$皆成立。
    \end{solution}
    \end{parts}

    \question 用数学归纳法证明:对所有正整数 $n$, 有
    \[
    \cos x + \cos 3x + \cos 5x + \cdots + \cos[(2n-1)x] = \frac{\sin(2nx)}{2\sin x}.
    \]
    \begin{solution}
        定义命题 $P_n$ 为等式
        \[
        P_n:\sum_{r=1}^{n} \cos[(2r-1)x] = \frac{\sin(2nx)}{2\sin x}, \quad n \in \mathbb{N}.
        \]
        基础步骤:当 $n=1$ 时, 左式为
        \[
        \sum_{r=1}^{1} \cos[(2r-1)x] = \cos(2 \cdot 1 - 1)x = \cos x.
        \]
        而右式为
        \[
        \frac{\sin(2 \cdot 1 \cdot x)}{2\sin x} = \frac{\sin 2x}{2\sin x} = \frac{2\sin x \cos x}{2\sin x} = \cos x.
        \]
        此时左式等于右式, 即 $P_1$ 成立。

        归纳步骤:设 $k \in \mathbb{N}$, 假设归纳假设 $P_k$ 成立, 即有
        \[
        \sum_{r=1}^{k} \cos[(2r-1)x] = \frac{\sin(2kx)}{2\sin x}.
        \]
        由归纳假设, 观察当 $n=k+1$ 时,
        \begin{align*}
        \sum_{r=1}^{k+1} \cos[(2r-1)x] &=  \sum_{r=1}^{k} \cos[(2r-1)x] + \cos[(2k+1)x] \\
        &= \frac{\sin(2kx)}{2\sin x} + \cos[(2k+1)x] \\
        &= \frac{\sin(2kx) + 2\sin x \cos[(2k+1)x]}{2\sin x} \\
        &= \frac{\sin(2kx) + \sin(2kx+2x) - \sin(2kx)}{2\sin x} \\
        &= \frac{\sin[2(k+1)x]}{2\sin x}.
        \end{align*}
        即命题 $P_{k+1}$ 成立, 此时蕴含 $P_k \implies P_{k+1}$ 成立。

        由数学归纳法, 命题 $P_n$ 对所有正整数 $n$ 皆成立。
    \end{solution}

    \question 用数学归纳法证明:对所有正整数 $n$, 有
    \[
    \frac{d^n}{dx^n} [e^x \cos x] = 2^{\frac{n}{2}} e^x \cos\left(x + \frac{n\pi}{4}\right).
    \]
    \begin{solution}
        定义命题 $P_n$ 为等式
        \[
        P_n:\frac{d^n}{dx^n} [e^x \cos x] = 2^{\frac{n}{2}} e^x \cos\left(x + \frac{n\pi}{4}\right), \quad n \in \mathbb{N}.
        \]

        基础步骤:当 $n=1$ 时, 左式为一阶导数, 由乘积法则得
        \[
        \frac{d}{dx} [e^x \cos x] = e^x \cos x + e^x (-\sin x) = e^x (\cos x - \sin x).
        \]
        而右式为
        \[
        2^{\frac{1}{2}} e^x \cos\left(x + \frac{\pi}{4}\right) = \sqrt{2} e^x \left( \frac{\sqrt{2}}{2}\cos x - \frac{\sqrt{2}}{2}\sin x \right) = e^x (\cos x - \sin x)
        \]
        此时左式等于右式, 即 $P_1$ 成立。

        归纳步骤:设 $k \in \mathbb{N}$, 假设归纳假设 $P_k$ 成立, 即有
        \[
        \frac{d^k}{dx^k} [e^x \cos x] = 2^{\frac{k}{2}} e^x \cos\left(x + \frac{k\pi}{4}\right).
        \]
        欲证命题 $P_{k+1}$ 成立, 对 $P_k$ 的等式两边关于 $x$ 求导:
        \begin{align*}
        \frac{d^{k+1}}{dx^{k+1}} [e^x \cos x] &= \frac{d}{dx} \left[ 2^{\frac{k}{2}} e^x \cos\left(x + \frac{k\pi}{4}\right) \right] \\
        &= 2^{\frac{k}{2}} \left[ e^x \cos\left(x + \frac{k\pi}{4}\right) + e^x \left( -\sin\left(x + \frac{k\pi}{4}\right) \right) \right] \\
        &= 2^{\frac{k}{2}} e^x \left[ \cos\left(x + \frac{k\pi}{4}\right) - \sin\left(x + \frac{k\pi}{4}\right) \right].
        \end{align*}
        利用辅助角公式 
        \[
        \cos A - \sin A = \sqrt{2} \cos(A + \frac{\pi}{4}),
        \]
        有
        \begin{align*}
        \frac{d^{k+1}}{dx^{k+1}} [e^x \cos x] &= 2^{\frac{k}{2}} e^x \cdot \sqrt{2} \cos\left(x + \frac{k\pi}{4} + \frac{\pi}{4}\right) \\
        &= 2^{\frac{k+1}{2}} e^x \cos\left(x + \frac{(k+1)\pi}{4}\right).
        \end{align*}
        即命题 $P_{k+1}$ 成立, 此时蕴含 $P_k \implies P_{k+1}$ 成立。

        由数学归纳法, 命题 $P_n$ 对所有正整数 $n$ 皆成立。
    \end{solution}

    \question 设 $n$ 为正整数, 计算定积分:
    \[
    \int_{0}^{\frac{\pi}{2}} \frac{\sin [(2n + 1) x]}{\sin x} \,dx.
    \]
    \begin{solution}
        记 
        \[
        I_n = \int_{0}^{\frac{\pi}{2}} \frac{\sin [(2n+1)x]}{\sin x} \,dx
        \]
        下面通过数学归纳法证明 $I_n = \dfrac{\pi}{2}$。

        基础步骤:当 $n=1$ 时, 利用三倍角公式 $\sin 3x = 3\sin x - 4\sin^3 x$, 可得:
        \[
        I_1 = \int_0^{\frac{\pi}{2}} \frac{3\sin x - 4\sin^3 x}{\sin x} \,dx = \int_0^{\frac{\pi}{2}} (3 - 4\sin^2 x) \,dx.
        \]
        利用降幂公式 $4\sin^2 x = 2(1 - \cos 2x)$, 积分化为
        \[
        I_1 = \int_0^{\frac{\pi}{2}} (1 + 2\cos 2x) \,dx = \left[ x + \sin 2x \right]_0^{\frac{\pi}{2}} = \frac{\pi}{2}.
        \]
        故 $n=1$ 时结论成立。

        归纳步骤:假设对于 $n=k-1$ 时结论成立, 即 $I_{k-1} = \frac{\pi}{2}$。考虑 $I_k$ 与 $I_{k-1}$ 的差值:
        \[
        I_k - I_{k-1} = \int_0^{\frac{\pi}{2}} \frac{\sin(2k+1)x - \sin(2k-1)x}{\sin x} \,dx.
        \]
        利用和差化积:
        \[
        \sin(2k+1)x - \sin(2k-1)x = 2\cos(2kx)\sin x.
        \]
        于是
        \[
        I_k - I_{k-1} = \int_0^{\frac{\pi}{2}} \frac{2\cos(2kx)\sin x}{\sin x} \,dx = 2 \int_0^{\frac{\pi}{2}} \cos(2kx) \,dx = \frac{\sin(k\pi)}{k}.
        \]
        由于 $k$ 是正整数, $\sin(k\pi) = 0$, 因此 $I_k - I_{k-1} = 0$, 即 $I_k = I_{k-1}$。
        根据归纳假设 $I_{k-1} = \frac{\pi}{2}$, 故 $I_k = \frac{\pi}{2}$ 也成立。

        由数学归纳法可知, 对所有正整数 $n$, 该积分的值恒为 $\frac{\pi}{2}$。
    \end{solution}

    \question 已知矩阵
    \[
    \mathbf{A} = \begin{pmatrix} 1 & 0 \\ 2 & 1 \end{pmatrix},
    \]
    用数学归纳法证明:对所有正整数 $n$, 有
    \[
    \mathbf{A}^n = n\mathbf{A} - (n-1)\mathbf{I}, \quad n \in \mathbb{N}.
    \]
    \begin{solution}
        定义命题
        \[
        P_n:\mathbf{A}^n = n\mathbf{A} - (n-1)\mathbf{I}.
        \]
        基础步骤:当 $n=1$ 时, 左式为 $\mathbf{A}^1 = \mathbf{A}$, 右式为 $1\cdot\mathbf{A} - (1-1)\mathbf{I} = \mathbf{A}$。
        此时左式等于右式, 即 $P_1$ 成立。

        归纳步骤:设 $k \in \mathbb{N}$, 假设归纳假设 $P_k$ 成立, 即
        \[
        \mathbf{A}^k = k\mathbf{A} - (k-1)\mathbf{I}.
        \]
        欲证命题 $P_{k+1}$ 成立, 考虑 $\mathbf{A}^{k+1}$:
        \[
        \mathbf{A}^{k+1} = \mathbf{A}^k \cdot \mathbf{A} = [k\mathbf{A} - (k-1)\mathbf{I}] \mathbf{A} = k\mathbf{A}^2 - (k-1)\mathbf{A}.
        \]
        观察到
        \[
        \mathbf{A}^2 = \begin{pmatrix} 1 & 0 \\ 2 & 1 \end{pmatrix} \begin{pmatrix} 1 & 0 \\ 2 & 1 \end{pmatrix} = \begin{pmatrix} 1 & 0 \\ 4 & 1 \end{pmatrix}=2\mathbf{A} - \mathbf{I}
        \]
        因此
        \begin{align*}
        \mathbf{A}^{k+1} &= k(2\mathbf{A} - \mathbf{I}) - (k-1)\mathbf{A} \\
        &= 2k\mathbf{A} - k\mathbf{I} - k\mathbf{A} + \mathbf{A} \\
        &= (k+1)\mathbf{A} - k\mathbf{I} \\
        &= (k+1)\mathbf{A} - [(k+1)-1]\mathbf{I}.
        \end{align*}
        即命题 $P_{k+1}$ 成立, 此时蕴含 $P_k \implies P_{k+1}$ 成立。

        由数学归纳法, 命题 $P_n$ 对所有正整数 $n$ 皆成立。
    \end{solution}

    \question 设 $V_n$ 为 $n$ 阶范德蒙行列式,定义如下:
    \[
    V_n = \begin{vmatrix}
    1 & x_1 & x_1^2 & \dots & x_1^{n-1} \\
    1 & x_2 & x_2^2 & \dots & x_2^{n-1} \\
    \vdots & \vdots & \vdots & \ddots & \vdots \\
    1 & x_n & x_n^2 & \dots & x_n^{n-1}
    \end{vmatrix}
    \]
    试以数学归纳法证明
    \[
    V_n = \prod_{1 \le i < j \le n} (x_j - x_i).
    \]
    \begin{solution}
        定义命题
        \[
        P_n:V_n = 
        \begin{vmatrix}
        1 & x_1 & x_1^2 & \dots & x_1^{n-1} \\
        1 & x_2 & x_2^2 & \dots & x_2^{n-1} \\
        \vdots & \vdots & \vdots & \ddots & \vdots \\
        1 & x_n & x_n^2 & \dots & x_n^{n-1}
        \end{vmatrix}= \prod_{1 \le i < j \le n} (x_j - x_i)\quad,n \in \mathbb{N}
        \]
        基础步骤:显然$P_1$成立,因$V=1=|1|=1$,且显然$P_2$亦成立,因为
        \[
        V_2=
        \begin{vmatrix} 
        1 & x_1 \\
        1 & x_2 
        \end{vmatrix} 
        =x_2-x_1
        \]
        归纳步骤:假设命题$P_k$成立,即
        \[
        V_k = 
        \begin{vmatrix}
        1 & x_1 & x_1^2 & \dots & x_1^{k-1} \\
        1 & x_2 & x_2^2 & \dots & x_2^{k-1} \\
        \vdots & \vdots & \vdots & \ddots & \vdots \\
        1 & x_k & x_k^2 & \dots & x_k^{k-1}
        \end{vmatrix}= \prod_{1 \le i < j \le k} (x_j - x_i)
        \]
        观察$V_{k+1}$,将 $x_1$写成变量$x$,即
        \[
        V_{k+1} = \begin{vmatrix}
        1 & x & x^2 & \dots & x^{k} \\
        1 & x_2 & x_2^2 & \dots & x_2^{k} \\
        \vdots & \vdots & \vdots & \ddots & \vdots \\
        1 & x_k & x_k^2 & \dots & x_k^{k} \\
        1 & x_{k+1} & x_{k+1}^2 & \dots & x_{k+1}^{k}
        \end{vmatrix}
        \]
        $V_{k+1}$按$R_1$展开后是次数不大于$k$的多项式,设$f(x)=V_{k+1}$,发现当$x=x_2,x_3,\cdots,x_{k+1}$,行列式$V_{k+1}$有相同的两行,故
        \[
        f(x_2)=f(x_3)=\;\dots\;=f(x_{k+1})=0 
        \]
        故 $f(x)$ 是一个次数为 $k$且以 $x_2,x_3,\cdot,x_{k+1}$ 为根的多项式,由余氏定理,
        \[
        f(x) = C(x-x_2)(x-x_3)\cdots(x-x_{k+1}) \tag{1}
        \]
        且$C$是一常数。又$V_{k+1}$按$R_1$展开式中$x^k$的系数为
        \[
        \begin{vmatrix}
        1 & x_1 & x_1^2 & \dots & x_1^{k-1} \\
        1 & x_2 & x_2^2 & \dots & x_2^{k-1} \\
        \vdots & \vdots & \vdots & \ddots & \vdots \\
        1 & x_k & x_k^2 & \dots & x_k^{k-1}
        \end{vmatrix}
        \]
        由归纳假设知此系数为
        \[
        \prod_{2 \le i < j \le k+1} (x_j - x_i)
        \]
        故由$(1)$知
        \[
        C=\prod_{2 \le i < j \le k+1} (x_j - x_i)
        \]
        即
        \[
        f(x)=(x-x_2)(x-x_3)\cdots(x-x_{k+1})\prod_{2 \le i < j \le k+1} (x_j - x_i)
        \]
        将$x$置换回$x_1$,即得命题$P_{k+1}$成立,此时蕴含$P_k \implies P_{k+1}$成立。
        
        由数学归纳法,命题$P_n$对所有正整数$n$皆成立。
        \textcolor{red}{验证是$x_j - x_i$还是$x_i - x_j$}
    \end{solution}

    \question 试证莱布尼茨公式:设函数 $f,g$ 定义在开区间 $I$ 上,$n$ 为正整数,$x \in I$ 为$I$内一点使得 $f,g$ 均可导 $n$ 次,则有
    \[
    (f(x)g(x))^{(n)} = \sum_{k=0}^{n} \binom{n}{k} f^{(k)}(x) g^{(n-k)}(x),
    \]
    其中 $(n)$ 表示导数的阶数。
    \begin{solution}
        定义命题
        \[
        P_n:(f(x)g(x))^{(n)} = \sum_{k=0}^{n} \binom{n}{k} f^{(k)}(x) g^{(n-k)}(x) \quad,n \in \mathbb{N}
        \]
        基础步骤:观察当$n=1$时,由乘积求导法则,左式为
        \[
        (f(x)g(x))' = f(x)g'(x) + f'(x)g(x)
        \]
        而右式为
        \[
        \sum_{k=0}^{1} \binom{1}{k} f^{(k)}(x) g^{(1-k)}(x)
        = \binom{1}{0} f^{(0)}(x) g^{(1)}(x) + \binom{1}{1} f^{(1)}(x) g^{(0)}(x)
        = f(x)g'(x) + f'(x)g(x),
        \]
        此时左式等于右式,即$P_1$成立。

        归纳步骤:设 $n \in \mathbb{N}$,假设归纳假设$P_n$成立,有
        \[
        (f(x)g(x))^{(n)} = \sum_{k=0}^{n} \binom{n}{k} f^{(k)}(x) g^{(n-k)}(x)
        \]
        欲证
        \[
        (f(x)g(x))^{(n+1)} = \sum_{k=0}^{n+1} \binom{n+1}{k} f^{(k)}(x) g^{(n+1-k)}(x)
        \]
        由归纳假设,
        \begin{align*}
        (f(x)g(x))^{(n+1)}
        &= \left( \sum_{k=0}^{n} \binom{n}{k} f^{(k)}(x) g^{(n-k)}(x) \right)' \\
        &= \sum_{k=0}^{n} \binom{n}{k} \left(f^{(k)}(x) g^{(n-k)}(x) \right)' \\
        &= \sum_{k=0}^{n} \binom{n}{k} \left( f^{(k+1)}(x) g^{(n-k)}(x) + f^{(k)}(x) g^{(n-k+1)}(x) \right) \\
        &= \sum_{k=0}^{n} \binom{n}{k} f^{(k+1)}(x) g^{(n-k)}(x)
        + \sum_{k=0}^{n} \binom{n}{k} f^{(k)}(x) g^{(n-k+1)}(x)
        \end{align*}
        将第一个求和中 $k=n$ 的项分离,第二个求和中 $k=0$ 的项分离,得
        \begin{align*}
        (f(x)g(x))^{(n+1)} 
        &= \sum_{k=0}^{n-1} \binom{n}{k} f^{(k+1)}(x) g^{(n-k)}(x)
        + \sum_{k=1}^{n} \binom{n}{k} f^{(k)}(x) g^{(n-k+1)}(x) \\
        &\quad + \binom{n}{n} f^{(n+1)}(x) g^{(0)}(x)
        + \binom{n}{0} f^{(0)}(x) g^{(n+1)}(x) \\
        &= \sum_{k=1}^{n-1} \binom{n}{k-1} f^{(k)}(x) g^{(n+1-k)}(x)
        + \sum_{k=1}^{n} \binom{n}{k} f^{(k)}(x) g^{(n-k+1)}(x) \\
        &\quad + \binom{n}{n} f^{(n+1)}(x) g^{(0)}(x)
        + \binom{n}{0} f^{(0)}(x) g^{(n+1)}(x)
        \end{align*}
        由帕斯卡恒等式,
        \begin{align*}
        &(f(x) g(x))^{(n+1)} \\
        &= \sum_{k=1}^{n} \binom{n+1}{k} f^{(k)}(x) g^{(n+1-k)}(x) + \binom{n}{0} f^{(0)}(x) g^{(n+1)}(x) + \binom{n}{n} f^{(n+1)}(x) g^{(0)}(x) \\
        &= \sum_{k=1}^{n} \binom{n+1}{k} f^{(k)}(x) g^{(n+1-k)}(x) + \binom{n+1}{0} f(x) g^{(n+1)}(x) + \binom{n+1}{n+1} f^{(n+1)}(x) g(x) \\
        &= \sum_{k=0}^{n+1} \binom{n+1}{k} f^{(k)}(x) g^{(n+1-k)}(x)
        \end{align*}
        即命题$P_{n+1}$成立,此时蕴含$P_n \implies P_{n+1}$成立。
        
        由数学归纳法,命题$P_n$对所有正整数$n$皆成立。
    \end{solution}

    \question 设 $a_0, a_1, a_2, \dots$ 为一无限实数数列, 且满足 $$\frac{a_{n-1}+a_{n+1}}{2} \ge a_n, \forall \,n\in \mathbb{N}$$ 证明 $$\frac{a_0+a_{n+1}}{2} \ge \frac{a_1+a_2+\dots+a_n}{n},\forall \,n\in \mathbb{N}$$
    \begin{solution}
        先证明:
        \[
        a_{0}+na_{n+1}\ge(n+1)a_{n}\tag{*}
        \]
        当$n=1$时, $a_{0}+a_{2}\ge2a_{1}$,显然(*)成立;假设$n=k$时(*)成立, 即
        \[
        a_{0}+ka_{k+1}\ge(k+1)a_{k}\tag{1}
        \]
        又
        \[
        a_{k+2}+a_{k}\ge2a_{k+1} \tag{2}
        \]
        由$(1),(2)$得$$a_{0}+(k+1)a_{k+2}\ge (k+1)(a_{k+2}+a_k)-ka_{k+1}\ge2(k+1)a_{k+1}-ka_{k+1}=(k+2)a_{k+1}$$即$n=k+1$时(*)也成立,由数学归纳法知(*)对所有$n\in \mathbb{N}$均成立。
        
        现证:
        \[
        n(a_{n}+a_{n+1})\ge2\sum_{i=1}^{n}a_{i}\tag{**}
        \]
        当$n=1$时,$\ a_{0}+a_{2}\ge2a_{1}$,(**)成立;假设$n=k$时(**)成立,即 
        \[
        k(a_{0}+a_{k+1})\ge2\sum_{i=1}^{k}a_{i}\tag{3}
        \]
        由$(*)$知
        \[
        a_{0}+(k+1)a_{k+2}\ge(k+2)a_{k+1}\tag{4}
        \]
        由$(3),(4)$得$$(k+1)(a_{0}+a_{k+2})\ge(k+2)a_{k+1}+2\sum_{i=1}^{k}a_{i}-ka_{k+1}=2\sum_{i=1}^{k+1}a_{i}$$ 即$n=k+1$时(**)也成立,由数学归纳法知(**)对所有$n\in \mathbb{N}$均成立
        即$$\frac{a_{0}+a_{n+1}}{2}\ge\frac{a_{1}+a_{2}+\cdot\cdot\cdot+a_{n}}{n},\forall \,n\in \mathbb{N}$$
        证毕。
    \end{solution}
    \begin{solution}
        设\(a_{n+1}-a_{n}=d_{n}\),由题意可知
        \[
        a_{n+1}-a_{n}\ge a_{n}-a_{n-1}\Rightarrow d_{n}\ge d_{n-1}
        \]
        且 $a_{n}=a_{0}+d_{1}+d_{2}+\cdot\cdot\cdot+d_{n-1}$,欲证等价于
        \begin{align*}
        &\frac{a_{0}+a_{n+1}}{2}\ge\frac{a_{1}+a_{2}+\cdot\cdot\cdot+a_{n}}{n}\\
        \Longleftrightarrow \ &n(a_{0}+a_{n+1})\ge2(a_{1}+a_{2}+\cdot\cdot\cdot+a_{n})\\
        \Longleftrightarrow \ &na_{0}+n(a_{0}+d_{0}+d_{1}+\cdot\cdot\cdot+d_{n-1})\\
        &\ge 2[(a_{0}+d_{0})+(a_{0}+d_{0}+d_{1})+\cdot\cdot\cdot+(a_{0}+d_{0}+d_{1}+\cdot\cdot\cdot+d_{n-1})]\\
        \Longleftrightarrow \ &n d_{n}\ge n d_{0}+(n-2)d_{1}+(n-4)d_{2}+(n-6)d_{3}+\cdot\cdot\cdot+(4-n)d_{n-2}+(2-n)d_{n-1}
        \end{align*} 
        当$n$为奇数,即证 \[n(d_{0}-d_{n})+(n-2)(d_{1}-d_{n-1})+(n-4)(d_{2}-d_{n-2})+\cdot\cdot\cdot+(d_{\frac{n-1}{2}}-d_{\frac{n+1}{2}})\le0 \tag{1}\]
        而 $d_{n}\ge d_{n-1}$, 则$$d_{0}-d_{n}\le0, d_{1}-d_{n-1}\le0, \cdot\cdot\cdot, d_{\frac{n-1}{2}}-d_{\frac{n+1}{2}}\le0$$ 此时(1)显然成立。
        
        当$n$为偶数,即证 \[n(d_{0}-d_{n})+(n-2)(d_{1}-d_{n-1})+(n-4)(d_{2}-d_{n-2})+\cdot\cdot\cdot+2(d_{\frac{n-2}{2}}-d_{\frac{n+2}{2}})\le0\tag{2}\]
        上式中$d_{\frac{n}{2}}$这一项没有,而$d_{n}\ge d_{n-1}$
        则$$d_{0}-d_{n}\le0, d_{1}-d_{n-1}\le0,\cdot\cdot\cdot,d_{\frac{n-2}{2}}-d_{\frac{n+2}{2}}\le0$$ 此时(2)也成立。
        
        故得证 $$\frac{a_0+a_{n+1}}{2} \ge \frac{a_1+a_2+\dots+a_n}{n},\forall \,n\in \mathbb{N}$$
    \end{solution}

    \question 已知多项式 $P(x) = x^2 + 2019x + 1$,试证对任意正整数 $n$,函数 $P^{(n)}(x) = 0$ 至少有一个实数根,其中 $P^{(n)}(x)$ 表示 $P$ 自身复合 $n$ 次。
    \begin{solution}
        定义命题:对任意正整数 $n,P^{(n)}(x)$ 有一个负实根 $-c_n$,且 $0 < c_n < 2$。
        
        基础步骤:当 $n=1$ 时,
        \[
        P^{(1)}(x) = P(x) = x^2 + 2019x + 1
        \]
        其判别式为 $\Delta_1 = 2019^2 - 4 > 0$,所以 $P^{(1)}(x)$ 有两个实数根,设为 $-b_1, -c_1$,其中 $c_1 \le b_1$,由韦达定理,
        \[
        b_1 + c_1 = 2019,\quad b_1 c_1 = 1
        \]
        因此 $b_1, c_1$ 都为正数,且因为乘积是 $1$,可得 $c_1 \le 1 < 2$,所以 $-c_1$ 是一个负实根,满足 $0 < c_1 < 2$,故命题在$n=1$ 时成立。
        
        归纳步骤:假设命题在$n=k$ 时成立,即 $P^{(k)}(x)$ 有一负实根 $-c_k$,其中 $0 < c_k < 2$,即
        \[
        P^{(k)}(x) = (x + c_k) Q_k(x)
        \]
        其中$Q_k(x)$为某多项式 ,则有
        \[
        P^{(k+1)}(x) = P^{(k)}(P(x)) = (P(x) + c_k) Q_k(P(x))
        \]
        其中
        \[
        P(x) + c_k = x^2 + 2019x + 1 + c_k
        \]
        的判别式为
        \[
        \Delta_{k+1} = 2019^2 - 4(1 + c_k) > 0
        \]
        即$P(x) + c_k$有两个实根,设为 $-b_{k+1}, -c_{k+1}$,其中 $c_{k+1} \le b_{k+1}$,由韦达定理,
        \[
        b_{k+1} + c_{k+1} = 2019,\quad b_{k+1} c_{k+1} = 1 + c_k > 0
        \]
        可知 $b_{k+1}, c_{k+1}$ 都为正数,且
        \[
        c_{k+1} \le \sqrt{1 + c_k} < \sqrt{3} < 2
        \]
        所以 $-c_{k+1}$ 是 $P^{(k+1)}(x)$ 的负实根,满足 $0 < c_{k+1} < 2$,故命题在$n=k+1$ 时也成立。
        
        由数学归纳法,对任意正整数 $n,P^{(n)}(x)$ 至少有一个实数根。
    \end{solution}
\end{questions}


\pagecolor{PageColor}
\
\vfil
\hfil  {\fontsize{50pt}{36pt}\selectfont{组合数学}} \hfil
\vfil
\begin{tikzpicture}[remember picture,overlay,every node/.style={inner sep=0pt}]
        \node [shift={(1cm,-1cm)},brown,scale=2,anchor=north west] (CNW)
        at (current page.north west) {\pgfornament[height=1cm,width=1cm]{61}};
        \node [shift={(-1cm,-1cm)},brown,scale=2,anchor=north east] (CNE)
        at (current page.north east) {\pgfornament[height=1cm,width=1cm,symmetry=v]{61}};
        \node [shift={(1cm,1cm)},brown,scale=2,anchor=south west] (CSW)
        at (current page.south west) {\pgfornament[height=1cm,width=1cm,symmetry=h]{61}};
        \node [shift={(-1cm,1cm)},brown,scale=2,anchor=south east] (CSE)
        at (current page.south east) {\pgfornament[height=1cm,width=1cm,symmetry=c]{61}};
        \pgfornamentline[color=brown]{current page.north west}{current page.north east}{2}{87}
        \pgfornamentline{current page.south west}{current page.south east}{2}{87}
        \pgfornamentline{current page.north west}{current page.south west}{3}{87}
        \pgfornamentline{current page.north east}{current page.south east}{3}{87}
        \end{tikzpicture}%
\thispagestyle{empty}
\pagebreak

\begin{center}
  {\fontsize{30pt}{26pt}\selectfont
    \hypertarget{排列与组合}{排列与组合} \label{排列与组合}
  }
\end{center}
\separator
\vspace{1pt}
\nopagecolor
\begin{questions}
    \question 求满足以下条件的三位数的个数:  
    \begin{itemize}
        \item 三个数字各不相同
        \item 数字按递减顺序排列
        \item 其中一个数字是5
    \end{itemize}  
    \begin{solution}
        即先选数字$5$,再从$9$个数字中任选$2$个,所以有 
        \[
        \comb{9}{2}=36
        \]
        个这样的三位数。注意到每组都能唯一地按递减顺序排成一个三位数,而且首位不会是 $0$。
    \end{solution}

    \question 已知 $a,b,c$ 为相异正整数且满足 $abc=2310$,求所有可能相异集合 $\{a,b,c\}$的个数。
    \begin{solution}
        发现$2310 = 2\times 3\times 5\times 7\times 11$,从 $5$ 个因数中,
        \begin{itemize}
            \item 挑 3 个连乘积作为 $a$,剩下 2 个作为 $b,c$,有
            $\comb{5}{3} = 10$种可能。
            \item 挑 2 个连乘积作为 $a$,剩下 3 个再挑 2 个乘积作为 $b$,有 $\comb{5}{2} \cdot \comb{3}{2} = 30$
            种可能。
        \end{itemize}
        故共有$10 + 30 = 40$种可能。
    \end{solution}

    \question $A$ 有一本共有 $2017$ 页的书，页码从 $1$ 到 $2017$。问有多少个页码同时包含至少一个数字 $1$ 和至少一个数字 $9$？例如 $91,1921,191$ 都符合条件。
    \begin{solution}
        情况一:页码为二位数。只有$19,91$满足要求,共2个页码符合要求。

        情况二:页码为三位数。页码中必有数字$1,9$,排列数为$3\cdot 2\cdot 10$,但此时多算了页码$119,191,199,911,919,991,019,091$,故符合要求的页码共有$3\cdot 2\cdot 10-6=52$。

        情况二:页码为三位数。千位必为数字1,而数字9可位于百位、十位或各位,排列数为$3\cdot 10\cdot 10$,但此时多考虑了含两个数字9的页码，共$3\cdot 10$。减去其之后,发现又少算了页码1999。故符合要求的页码共有$3\cdot 10\cdot 10-3\cdot 10+1=271$。

        故共有
        \[
        2+52+271=325
        \]
        个页码符合条件。
    \end{solution}

    \question 设 $N$ 为满足 $x < y < z$ 且
    \[
    xyz = 2^2 \cdot 3^2 \cdot 5^2 \cdot 7^2 \cdot 11^2 \cdot 13^2 \cdot 17^2 \cdot 19^2
    \]
    的正整数三元组 $(x, y, z)$ 的个数。求 $N$。
    \begin{solution}
        忽略$x < y < z$,先计算无序三元组 $(a,b,c)$的个数,使得
        \[
        abc = 2^2 \cdot 3^2 \cdot 5^2 \cdot 7^2 \cdot 11^2 \cdot 13^2 \cdot 17^2 \cdot 19^2
        \]
        每个平方的质因数可以分配给 $a,b,c$ 的方式有 6 种:两个都给 $a$,两个都给 $b$,两个都给 $c$,或者每个数各取一个。共有 8 个平方质因数,因此共有 $6^8$ 种分配方式。

        现考虑$a,b,c$中两个数相等的情况。观察到积$abc$不是完全立方数,所以$a=b=c$不成立。计算恰有一对相等的三元组:以 $a=b$ 为例,每个平方质因数分配给 $c$,亦或分配给 $a$ 和 $b$ 各一个,共 $2^8$ 种三元组需要排除,同理$a=c,b=c$ 也各有 $2^8$ 种,所以相异的三元组数为:
        \[
        6^8 - 3\cdot 2^8
        \]
        将无序三元组 $(a,b,c)$ 转换为有序三元组 $(x,y,z)$,其中$x<y<z$,每个三元组对应 6 个无序三元组,因此
        \[
        N = \frac{1}{6}(6^8 - 3\cdot 2^8) = 279808
        \]
    \end{solution}

    \question 已知 $a_0, a_1, a_2, a_3, a_4$ 的值正好都是 $-1, 0, 1$ 中的数, 则 $a_0 + 3a_1 + 3^2a_2 + 3^3a_3 + 3^4a_4$ 的值是正整数共有多少个？
    \begin{solution}
        形如
        \[
        1\triangle \triangle \triangle \triangle,01\triangle \triangle \triangle,
        001\triangle \triangle,
        0001\triangle,
        00001
        \]
        的数共有
        \[
        3^4,3^3,3^2,3^1,1
        \]
        个,总计
        \[
        81 + 27 + 9 + 3 + 1 = 121
        \]
    \end{solution}

    \question 求正整数有序三元组 $(x,y,z)$ 的个数,使得
    \[
    xyz = 4000.
    \]
    \begin{solution}
        质因数分解得$4000 = 2^5 \cdot 5^3$,设
        \[
        x = 2^a 5^d, \quad y = 2^b 5^e, \quad z = 2^c 5^f,
        \]
        其中 $a,b,c,d,e,f$ 为非负整数,则需满足
        \[
        a+b+c = 5, \quad d+e+f = 3
        \]
        非负整数解的个数为
        \[
        \#(a,b,c) = \comb{5+3-1}{3-1} = 21, \quad \#(d,e,f) = \comb{3+3-1}{3-1} = 10
        \]
        因此有序三元组 $(x,y,z)$ 的个数为
        \[
        21 \cdot 10 = 210
        \]
    \end{solution}

    \question 将 $1,2,3,4,5,6,7$ 排成一列,若规定排列后不得出现 $12,23,34,45,56,67$(如：1273546 不合题意,7362154 符合题意),则有多少种排法？
    \begin{solution}
    7 个相异数字任意排列共有 $7!$ 种。

    从$12,23,34,45,56,67$选$12$后有6个板块,出现$12$的排列数为$\comb{6}{1}6!$。

    依此类推,由容斥原理,总排列数为
    \[
    7!-\comb{6}{1}6!+\comb{6}{2}5!-\comb{6}{3}4!+\comb{6}{4}3!-\comb{6}{5}2!+1=2119
    \]
    \end{solution}

    \question 如果一个数的每一位都大于前一位,则称它为上升数。例如 457 是上升数,但 447 不是。问 400 到 5000 之间共有多少上升数？
    \begin{solution}
        情况一:上升数为三位数。三位数必须在 400 到 999 之间,百位可以是 4,5,6,7,8,9。4 到 9 的 6 个数字中任选 3 个,按升序排列即可得到上升数。排列数为$\comb{6}{3}$。

        情况二:上升数为四位数。四位数小于 5000,千位必须为 1,2,3,4。1 到 9 的数字中任选 4 个,再扣除首位 $\ge 5$ 的数。排列数为$\comb{9}{4} - \comb{5}{4}$。

        因此共有 
        \[
        \comb{6}{3} + (\comb{9}{4} - \comb{5}{4}) = 141
        \]
        个这样的上升数。
    \end{solution}

    \question Ricardo 想要将三个 1、三个 2、两个 3 和一个 4 排成一个九位正整数,且满足以下条件:
    \begin{itemize}
        \item 从左到右,至少有一个 1 在第一个 2 之前,至少有一个 2 在第一个 3 之前,至少有一个 3 在 4 之前;
        \item 任意两个数字 2 不能相邻。
    \end{itemize}
    求总共有多少种符合条件的九位数。
    \begin{solution}
        设 $N$ 为满足条件的整数。$N$ 的首位必须是 1,因此 $N$ 可以以 1、11 或 111 开头。又因为第一个非 1 的数字必须是 2,所以 $N$ 只能以 12、112 或 1112 开头。
        
        情况 1:$N$ 以 12 开头。剩余两个 2 不能相邻,可放在下列位置组合(从左数第几位):
        \[
        (4,6),(4,7),(4,8),(4,9),(5,7),(5,8),(5,9),(6,8),(6,9),(7,9)
        \]
        共有 10 种可能。剩余两个 1 可以放在剩下的 5 个空位中的任意两位,有 $\comb{5}{2}=10$ 种方法。剩余的两个 3 和一个 4 需放在最后 3 个位置中。为保证至少有一个 3 在 4 之前,第一个空位必须放 3,剩下两个位置放另一个 3 和一个 4,有 2 种排列。因此本情况共有:
        \[
        10 \cdot 10 \cdot 2 = 200.
        \]
        情况 2:$N$ 以 112 开头。剩余两个 2 不能相邻,可放在:
        \[
        (5,7),(5,8),(5,9),(6,8),(6,9),(7,9)
        \]
        共有 6 种可能。剩余一个 1 可放在 4 个空位中的任意一位,有 4 种方法。剩余两个 3 和一个 4 需放在最后 3 个位置中。第一个空位放 3,剩下两个位置放另一个 3 和一个 4,有 2 种排列。本情况共有:
        \[
        6 \cdot 4 \cdot 2 = 48.
        \]
        情况 3:$N$ 以 1112 开头。剩余两个 2 不能相邻,可放在
        \[
        (6,8),(6,9),(7,9)
        \]
        共有 3 种可能。剩余的两个 3 和一个 4 需放在最后 3 个位置中。第一个空位放 3,剩下两个位置放另一个 3 和一个 4,有 2 种排列。本情况共有:
        \[
        3 \cdot 2 = 6.
        \]
        综上,共有 $200 + 48 + 6 = 254$ 种符合条件的九位数。
    \end{solution}

    \question
3) 从数字1至9中选出7位数，要求每个数字不重复，并且5和6不能连续出现，求共有多少种排列方式。

\begin{solution}
考虑不同情况：

(a) 5和6都出现：

- 情况一：第1位是5，第2位不是6：
\[
5 \times \perm{7}{5} = 12600
\]

- 情况二：第7位是5，第6位不是6：
\[
5 \times \perm{7}{5} = 12600
\]

- 情况三：5出现在首尾之外的其他位置：
\[
5 \times 4 \times \perm{7}{5} = 50400
\]

因此，5和6都出现的排列数：
\[
12600 + 12600 + 50400 = 75600
\]

(b) 5和6都不出现：
\[
\perm{7}{7} = 7! = 5040
\]

(c) 5出现但6不出现：
\[
7 \times \perm{7}{6} = 7 \times 7! = 35280
\]

(d) 5不出现但6可以出现：
\[
7 \times \perm{7}{6} = 7 \times 7! = 35280
\]

根据加法原理，总排列数为：
\[
75600 + 5040 + 2 \times 35280 = 151200
\]
\end{solution}


    \question 令 $a_n$ 为第 $n$ 小的各位数字之和为 $3$ 的正整数。例如 $a_1=3, a_2=12, a_3=21, a_4=30$。求 $a_{2012}$ 有多少位数。  
    \begin{solution}
        若一个数最多有 $d$ 位,那么其各位和为 $3$ 的这样的数的个数为
        \[
        \comb{d+2}{3}
        \]
        这是因为我们可以用隔板法来理解:将 $3$ 个相同的球(代表数字和为 3)分配到 $d$ 个位置(代表 $d$ 位数),等价于在 $d+2$ 个位置中选择 $3$ 个位置放置隔板。具体地,我们在一排 $d+2$ 个位置中放置恰好 $3$ 个 O(其余为 X),则第 $i$ 位的数字等于第 $i-1$ 个 X 与第 $i$ 个 X 之间的 O 的个数。因此需找最小的 $d$ 使得
        \[
        \comb{d+2}{3}\ge 2012.
        \]
        计算得
        \[
        \comb{23}{3}=1771,\quad \comb{24}{3}=2024.
        \]
        因此当 $d=22$ 时,$\comb{24}{3}=2024\ge2012$,且 $d=21$ 时不足。故 $a_{2012}$ 恰有 $22$ 位。
    \end{solution}

    \question 将$A,B,C,D,E,F,G,H$八个字母排成一列,使得$B$在$A$之右方,$E$在$C$与$D$之间,且$F,G$不相邻,试问符合条件的排法有多少种？
    \begin{solution}
        在所有的排列中,$B$ 在 $A$ 的右方与 $A$ 在 $B$ 的右方各占一半,因此满足"$B$ 在 $A$ 的右方"的排列占全部排列的 $\dfrac{1}{2}$。
        
        同理,在 $C,D,E$ 三个字母的相对位置中,$E$ 在 $C$ 与 $D$ 之间、$C$ 在 $E$ 与 $D$ 之间、$D$ 在 $E$ 与 $C$ 之间各占三分之一,因此满足"$E$ 在 $C$ 与 $D$ 之间"的排列占全部排列的 $\dfrac{1}{3}$。
        
        现计算满足"$F,G$ 不相邻"的排列数:
        
        八个字母的全排列有 $8!$ 种。若 $F,G$ 相邻,可将它们视为一个整体,有 $7!$ 种排法,而 $F,G$ 内部有 $2$ 种排列,因此 $F,G$ 相邻的排法有 $7! \cdot 2$ 种。
        
        所以 $F,G$ 不相邻的排法有 $8! - 7! \cdot 2$ 种。
        
        由于这三个条件相互独立,符合所有条件的排法数为:
        \[
        \frac{8! - 7! \cdot 2}{2 \cdot 3} = 5040
        \]
    \end{solution}

    \question 有 6 个人有网络账号，已知每个人都恰好与自己以外的 2 个人互为好友，则共有种不同的组成方法?
    \begin{solution}
        用六个顶点表示六个人，两顶点有连线表示两人互为好友。

        情况一:一个六边形。共有$\dfrac{5!}{2}$种。

        情况二:两个三角形。共有$\dfrac{\comb{6}{3}}{2}$种。

        因此共有
        \[
        \frac{5!}{2}+\frac{\comb{6}{3}}{2}=70
        \]
        种组合方式。
    \end{solution}

\question 求不大于 $2018$ 的正整数中,二进制中 1 出现比 0 多的个数。
\begin{solution}
    考虑不大于 $2047=2^{11}-1$且满足条件的正整数个数,再减去 2019 至 2047 的 29 个数,因为这29 个数的二进制中至少有6个1。

    对于偶数位数（去掉开头的 1 后有偶数位）,一半的数满足 1 的个数多于 0。具体公式为:
    \[
    \frac{1}{2}\left(2^{2k} + \comb{2k}{k}\right)
    \]
    因为恰好 $k$ 个 1 的情况算一半,其余一半满足条件,计算:
    \[
    \frac{1}{2}\Big(2047 + \comb{0}{0} + \comb{2}{1} + \comb{4}{2} + \comb{6}{3} + \comb{8}{4} + \comb{10}{5}\Big) = 1199
    \]
    故答案为$1199-29=1170$。
    \textcolor{red}{Half of the numbers with an even number of bits have this property, since they have an odd number of bits after their initial 1. Of the $2^{2k}$ numbers with $2k$ bits following the initial 1, the number with this property is $\frac{1}{2}(2^{2k} + \binom{2k}{k})$ since those with $k$ 1's will be included, and half of the others will.}
\end{solution}

    \question 小王有八个编号为 1 到 8 的盒子和八个编号为 1 到 8 的球。问他把球放入盒子,使每个盒子恰好有一个球,且球 1 不在盒子 1,球 2 不在盒子 2,球 3 不在盒子 3,有多少种方法？
    \begin{solution}
        设
        \[
        A_1 = \{\text{球 1 在盒子 1}\}, \quad
        A_2 = \{\text{球 2 在盒子 2}\}, \quad
        A_3 = \{\text{球 3 在盒子 3}\}.
        \]
        根据容斥原理,不允许的排列数为
        \begin{align*}
        &|A_1 \cup A_2 \cup A_3| \\
        &= |A_1| + |A_2| + |A_3| - |A_1 \cap A_2| - |A_1 \cap A_3| - |A_2 \cap A_3| + |A_1 \cap A_2 \cap A_3| \\
        &=3\cdot 7! - 3\cdot 6! + 5! = 13080           
        \end{align*}
        因此满足条件的排列数为
        \[
        S = 8! - |A_1 \cup A_2 \cup A_3| = 27240
        \]
    \end{solution}

    \question 在一排有 20 张椅子的座位区中, 要安排甲、乙、丙、丁、戊 5 人入坐, 一人坐一张椅子, 要求第一张与最后一张椅子不能安排人入坐, 且每相邻的 5 张椅子至少要有一人入坐, 任两人不能坐在相邻的椅子上。试问 5 人入坐的方法有多少种可能?
    \begin{solution}
    20 张椅子 5 人先入坐, 有 $5!$ 排法, 剩下 15 张椅子。依规定头尾各放一张, 剩下 13 张; 两人中间各摆一张, 剩下 9 张。

    9 张椅子可以在 5 人的六个间隔中摆放, 有 $\comb{9+6-1}{6-1}$ 摆法。

    为符合「相邻 5 张椅子至少要有一人坐」, 需扣除某个间隔塞入四张或以上的情形,因此实际摆法有
    \[
    5!(\comb{9+6-1}{6-1} - \comb{6}{1} \cdot \comb{6+5-1}{5-1} + \comb{6}{2} \cdot \comb{6+1-1}{1-1})= 69600
    \]
    \end{solution}

    \question 由字母 $A$ 和 $B$ 组成的 9 个字母的字母串中,有多少个不包含连续字母串 $ABBA$？
    \begin{solution}
        总字母串数为 $2^9 = 512$,我们先计算至少包含一个 $ABBA$ 的字符串数量,再用总数减去它。

        $ABBA$ 可以出现在前 6 个位置中的任意一个位置,剩余 5 个位置可以随意填充 $A$ 或 $B$,因此每个起始位置有 $2^5 = 32$ 种方法,共 $6 \cdot 32 = 192$ 种字母串。  

        这 192 种字母串中有些被重复计算,因为 $ABBA$ 可以重叠如下:
        \[
        xyABBABBA, \quad xABBABBAy, \quad ABBABBAxy
        \]
        其中 $x, y$ 可以是 $A$ 或 $B$,共有 $12$ 种。还有 6 种形为
        \[
        xABBAABBA, \quad ABBAxABBA, \quad ABBAABBAx
        \]
        这些被重复计算了 2 次。  

        根据容斥原理,因此至少包含一个 $ABBA$ 的字母串数为 $192-12-6 = 174$。故不包含 $ABBA$ 的字母串数为:
        \[
        512 - 174 = 338
        \]
    \end{solution}

    \question 设 Bauman 字串满足以下条件:
    \begin{itemize}
        \item 每个字母只能是 $A, B, C, D, E$ 之一；
        \item 相邻两个字母不能相同。
    \end{itemize}
    例如 $AECD, BDCEC$ 是 Bauman 字串,而 $ABBC, DAEEE$ 不是。
    \begin{parts}
    \part 长度为 5 的 Bauman 字串中,首尾字母都是 $A$ 的有多少个？
    \begin{solution}
        若首尾均为 $A$,则第二位和第四位不能是 $A$。按第三位分类讨论:

        情形 1:第三位是 $A$。字串形如 $A\_A\_A$,第二位有 4 种选择($B,C,D,E$),第四位有 4 种选择,共 $4 \cdot 4 = 16$ 种。

        情形 2:第三位不是 $A$。第三位有 4 种选择($B,C,D,E$),第二位需与第三位不同且不能为 $A$,共 3 种选择;第四位同理有 3 种选择,共 $4 \cdot 3 \cdot 3 = 36$ 种。

        总数为 $16 + 36 = 52$。
    \end{solution}
    \begin{solution}
        按第二位和第四位是否相同分类:

        情形 1:第二位和第四位相同。字串形如 $Ax\_xA$(其中 $x$ 表示第二位和第四位的字母),第二位有 4 种选择,第三位需与第二位不同,有 4 种选择,共 $4 \cdot 4 = 16$ 种。

        情形 2:第二位和第四位不同。字串形如 $Axy\_A$(其中 $x$, $y$ 分别表示第二位和第四位的字母),第二位有 4 种选择,第四位需与第二位不同,有 3 种选择,第三位需与第二位和第四位都不同,有 3 种选择,共 $4 \cdot 3 \cdot 3 = 36$ 种。

        总数为 $16 + 36 = 52$。
    \end{solution}
    \part 长度为 6 的 Bauman 字串中,包含超过一个 B 的有多少个？
    \begin{solution}
        总字串数为 $5 \cdot 4^5 = 5120$(首位 5 种选择,其余每位 4 种选择)。

        无 $B$ 的字串:首位 4 种选择($A,C,D,E$),其余每位 3 种选择(需与前一位不同且不能是 $B$),共 $4 \cdot 3^5 = 972$ 种。

        恰有 1 个 $B$ 的字串:
        \begin{itemize}
            \item 若 $B$ 在首位或末位:
            \begin{itemize}
                \item $B$ 在首位:第二位 4 种,其余每位 3 种,共 $1 \cdot 4 \cdot 3^4 = 324$ 种
                \item $B$ 在末位:首位 4 种,第二到第五位中第二位 3 种,其余每位 3 种,共 $4 \cdot 3^4 = 324$ 种
            \end{itemize}
            \item 若 $B$ 在第 2 到第 5 位(共 4 个位置):每个位置,$B$ 前一位 4 种,$B$ 后一位 4 种,其余 3 个位置各 3 种,每个位置贡献 $4 \cdot 1 \cdot 4 \cdot 3^3 = 432$ 种,共 $4 \cdot 432 = 1728$ 种
        \end{itemize}
        恰有 1 个 $B$ 的总数为 $324+ 324 + 1728 = 2376$。

        超过 1 个 $B$ 的数目为:
        \[
        5120 - 972 - 2376 = 1772.
        \]
    \end{solution}
    \begin{solution}
       长度为 6 的 Bauman 字串中,由于相邻字母不能相同,$B$ 最多出现 3 次。故只需计算恰有 2 个 $B$ 或恰有 3 个 $B$ 的字串数。

        情形 1:恰有 3 个 $B$。三个 $B$ 必须间隔放置。可能的位置模式有:
        \begin{itemize}
            \item $B\_B\_B\_$(第 1, 3, 5 位):第 2, 4, 6 位各 4 种选择,共 $4^3 = 64$ 种
            \item $\_B\_B\_B$(第 2, 4, 6 位):第 1, 3, 5 位各 4 种选择,共 $4^3 = 64$ 种
            \item $B\_B\_B$(第 1, 3, 5 位):第 2 位 4 种,第 4 位 4 种,第 6 位 3 种,共 $4 \cdot 4 \cdot 3 = 48$ 种
            \item $\_B\_B\_$(第 2, 4, 6 位):第 1 位 4 种,第 3 位 4 种,第 5 位 3 种,共 $4 \cdot 4 \cdot 3 = 48$ 种
        \end{itemize}
        合计 $64 + 64 + 48 + 48 = 224$ 种。

        情形 2:恰有 2 个 $B$。从 6 个位置中选 2 个放 $B$,有 $\comb{6}{2} = 15$ 种选择,但需排除相邻的 5 种,剩余 10 种位置模式。

        逐一计数(用下划线表示非 $B$ 位置):
        \begin{itemize}
            \item $B\_\_\_B\_$:第 2 位 4 种,第 4 位 4 种,第 5, 6 位各 3 种,共 $4 \cdot 4 \cdot 3 \cdot 3 = 144$ 种
            \item $\_\_B\_\_B$:对称,同样 $144$ 种
            \item $B\_\_B\_\_$:第 2, 4 位各 4 种,第 5 位需与第 4 位不同有 3 种,第 6 位 3 种,共 $4 \cdot 4 \cdot 3 \cdot 3 = 144$ 种
            \item $\_\_B\_B\_$:对称,同样 $144$ 种
            \item $B\_B\_\_\_$:第 2 位 4 种,第 4, 5, 6 位各 3 种,共 $4 \cdot 3 \cdot 3 \cdot 3 = 108$ 种
            \item $\_\_\_B\_B$:对称,第 1 位 4 种,第 2, 3 位各 3 种,第 5 位 3 种,共 $4 \cdot 3 \cdot 3 \cdot 3 = 108$ 种
            \item $B\_B\_\_\_$:第 2 位 4 种,第 4 位需与第 3 位不同有 4 种,第 5, 6 位各 3 种,共 $4 \cdot 4 \cdot 3 \cdot 3 = 144$ 种
            \item $\_\_B\_B$:对称,同样 $144$ 种
            \item $B\_\_\_\_B$:第 2 位 4 种,第 3, 4, 5 位各 3 种,共 $4 \cdot 3^3 = 108$ 种
            \item $\_B\_B\_\_$:第 1, 3, 5 位各 4 种,第 6 位 3 种,共 $4^3 \cdot 3 = 192$ 种
        \end{itemize}
        合计 $4 \cdot 144 + 2 \cdot 108 + 2 \cdot 144 + 108 + 192 = 1548$ 种。

        恰 2 个 $B$ 的有 1548 种,恰 3 个 $B$ 的有 224 种,总和为 $1548 + 224 = 1772$。
    \end{solution}
    \end{parts}


    \question 将字母串 $AAAABBBCCC$ 排成一列,且相同字母不相邻的排法有多少种？
    \begin{solution}
        按 $B,C$ 的排列分类,再插入 $A$。先排列 3 个 $B$ 和 3 个 $C$ 的相对次序,记为「$BC$ 串」。共有
        \[
        \comb{6}{3}=\frac{6!}{3!3!}=20
        \]
        种 $BC$ 串。按 $BC$ 串中出现相邻同字母的"块数"分类,列举如下:
        \begin{center}
        \begin{tabular}{c|l}
        编号 & 字串 \\ \hline 
        1 & $BCBCBC$ \\ 
        2 & $CBCBCB$ \\ \hline 
        3 & $BCCBCB$ \\ 
        4 & $CBCBBC$ \\ 
        5 & $BCBCCB$ \\ 
        6 & $CBBCBC$ \\ \hline 
        7 & $BCCBBC$ \\ 
        8 & $CCBBCB$ \\ 
        9 & $BCBBCC$ \\ 
        10 & $CCBCBB$ \\ 
        11 & $BBCCBC$ \\ 
        12 & $CBBCCB$ \\ 
        13 & $BBCBCC$ \\ 
        14 & $CBCCBB$ \\ \hline 
        15 & $CBBBCC$ \\ 
        16 & $CCBBBC$ \\ 
        17 & $BBCCCB$ \\ 
        18 & $BCCCBB$ \\ \hline 
        19 & $CCCBBB$ \\ 
        20 & $BBBCCC$ \\ \hline
        \end{tabular}
        \end{center}
        现在将 4 个 $A$ 插入 $BC$ 串的空位中,使得任何相同字母都不相邻:
        \begin{enumerate}
            \item $BC$ 串中没有相邻同字母（编号 1--2）:有 7 个空位插入 4 个 A,方式数 $\comb{7}{4}=35$,对应 2 种 $BC$ 串,共 $2\cdot 35=70$。
            \item 有 1 处相邻同字母（编号 3--6）:此时剩 6 个空位插入 4 个 A（等价插入 3 个 A 以避免产生相邻相同字母）,方式数 $\comb{6}{3}=20$,对应 4 种 $BC$ 串,共 $4\cdot20=80$。
            \item 有 2 处相邻同字母（编号 7--14）:剩 5 个空位插入 2 个 A,方式数 $\comb{5}{2}=10$,对应 8 种 $BC$ 串,共 $8\cdot10=80$。
            \item 有 3 处相邻同字母（编号 15--18）:剩 4 个空位插入 1 个 A,方式数 $\comb{4}{1}=4$,对应 4 种 $BC$ 串,共 $4\cdot4=16$。
            \item 有 4 处相邻同字母（编号 19--20）:A 恰好插满一组,计 $2$ 种。
        \end{enumerate}
        合计:
        \[
        70+80+80+16+2=248
        \]
    \end{solution}
    \begin{solution}
        根据容斥原理,相同字母不相邻的排法数
        \[
        = \text{4个A不相邻} - (\text{至少2个B相邻} \cup \text{至少2个C相邻})
        \]
        4个A不相邻:先排 3 个 $B$ 和 3 个 $C$,有 $\comb{6}{3}=20$ 种方式,产生 7 个空位,在其中插入 4 个 $A$,有 $\comb{7}{4}=35$ 种方式,共 $20 \cdot 35=700$ 种。

        至少2个B相邻:将 2 个 $B$ 看作一个整体 $BB$,与剩余 1 个 $B$ 和 3 个 $C$ 排列,有 $\frac{5!}{3!}=20$ 种方式,产生 6 个空位插入 4 个 $A$,有 $\comb{6}{4}=15$ 种方式,共 $20 \cdot 15=300$ 种。

        但这样会重复计算 3 个 $B$ 都相邻的情况。当 3 个 $B$ 都相邻时,将它们看作 $BBB$,与 3 个 $C$ 排列有 4 种方式,产生 5 个空位插入 4 个 $A$,有 $\comb{5}{4}=5$ 种方式,共 $4 \cdot 5=20$ 种。

        因此至少 2 个 $B$ 相邻为 $300-20=280$ 种。

        由对称性,至少 2 个 $C$ 相邻也是 $280$ 种。

        至少2个B相邻且至少2个C相邻:将 2 个 $B$ 看作 $BB$,2 个 $C$ 看作 $CC$,与剩余 1 个 $B$ 和 1 个 $C$ 排列,有 $4!=24$ 种方式,产生 5 个空位插入 4 个 $A$,有 $\comb{5}{4}=5$ 种方式,共 $24 \cdot 5=120$ 种。

        需扣除 3 个 $B$ 都相邻的情况:$BBB,CC,C$ 排列有 $3!=6$ 种,产生 4 个空位插入 4 个 $A$,有 $\comb{4}{4}=1$ 种,共 $6 \cdot 1=6$ 种。

        同样扣除 3 个 $C$ 都相邻的情况:$6$ 种。

        因此至少 2 个 $B$ 相邻且至少 2 个 $C$ 相邻为 $120-6-6=108$ 种。

        最终答案为 
        \[
        \comb{6}{3}\comb{7}{4} - \left[2 \times \left(\frac{5!}{3!}\comb{6}{4}-4\comb{5}{4}\right) - \left(4!\comb{5}{4}-2 \times 3!\comb{4}{4}\right)\right]=248
        \]
    \end{solution}

    \question 要爬 12 级台阶,但每步只能上 1 级或 2 级。第 8 级有一条蛇,所以不能踩。问共有多少种爬法？
    \begin{solution}
        由于第 8 级不能踩,必须从第 7 级迈 2 级直接到第 9 级。于是问题可以拆成两段:
        \begin{itemize}
            \item 爬前 7 级的方式数
            \item 从第 9 级到第 12 级的方式数
        \end{itemize}
        总数为
        \[
        (\comb{7}{0}+\comb{6}{1}+\comb{5}{2}+\comb{4}{3})\cdot 3=63
        \]
    \end{solution}

    \question 在集合 $\{1,2,3,4,5,6,7\}$ 的所有非空子集 S 中,有多少个子集不包含数 $|S|$($|S|$ 表示子集元素个数)？例如 $\{3,4\}$ 是这样的子集,因为它不包含数字 2。
    \begin{solution}
        非空子集总数为 $2^7-1=127$。若子集大小为 $k$ 且包含数字 $k$,则剩下的 $k-1$ 个元素必须从另外的 6 个数中选取,共有
        \[
        \comb{6}{k-1}
        \]
        种。对 $k=1,2,\dots,7$ 求和得包含其大小的子集总数为
        \[
        \sum_{k=1}^7 \comb{6}{k-1} = \sum_{j=0}^6 \comb{6}{j} = 2^6 = 64.
        \]
        因此不包含其大小的子集个数为
        \[
        127-64=63.
        \]
    \end{solution}

    \question 有多少个非负整数有序四元组 $(a,b,c,d)$ 满足 $a+b+c+d\le15$?
    \begin{solution}
        将不等式转化为等式
        \[
        a+b+c+d+e=15,
        \]
        其中 $e\ge0$,则原题等价于求此方程的非负整数解个数。根据隔板法,有
        \[
        \comb{15+5-1}{5-1}=\comb{19}{4}=3876
        \]
        个非负整数有序四元组。
    \end{solution}

    \question 以一个正方体的顶点为顶点的四面体共有多少个？
    \begin{solution}
        从正方体的 8 个顶点中任取 4 点:
        \[
        \comb{8}{4} = 70
        \]
        需扣除 4 点共平面的情形,包括正方体的六个面,共 6 种;斜平面,共 6 种。所以四面体个数为
        \[
        70 - 12 = 58
        \]
    \end{solution}

    \question 在 5 只蜗牛进行的比赛中,最多只会出现一次打平,但可以接受任意数量的蜗牛打平。例如,比赛结果可能是 Dazzler 获得第一名;Abby、Cyrus 和 Elroy 并列第二名,而 Bruna 获得第五名。问这种比赛共有多少种不同的结果? 
    \begin{solution}
        考虑没有、两只、三只、四只、五只蜗牛打平,总共有
        \[
        5! + \comb{5}{2} \cdot 4! + \comb{5}{3} \cdot 3! + \comb{5}{4} \cdot 2! + 1 = 431
        \]
        个不同的结果。
    \end{solution}

    \question 在一个 $3 \times 3$ 的网格中,放置 $3$ 个红色棋子、$3$ 个蓝色棋子和 $3$ 个绿色棋子,所有棋子同色不可相邻（横或竖方向）,且棋子同色不可区分。请问满足条件的放置方法有多少种？
        \ifprintanswers
        \begin{figure}[H]
        \centering
        \includegraphics[width=0.4\textwidth]{images/image13.png}
        \end{figure}
        \fi
    \begin{solution}
        观察发现,只存在两种基本的放置方法（如上图所示）,其他放置方法只能通过这两种结构旋转或颜色变化得到,所以总合法方案数为
        \[
        (4 + 2) \cdot 3! = 36
        \]
    \end{solution}

    \question 小明在注册账号时可以使用字符$1,2,3,a,b,c,A,B,C$来组成五位密码,但要求必须包含数字、小写字母和大写字母,且不可以出现两个相同的字符相邻,例如密码可以设置为$123aA$或$laA12$,但不能设置为$123ab$或$112aA$,试求可以设置不同的密码的个数。
    \begin{solution}
        先考虑相邻字符不同的密码,共有$9 \cdot 8^4 = 36864$种,这里面不满足密码要求的有两类:
        \begin{itemize}
            \item 仅包含单一字符类型(如全数字),这类共有$3 \cdot (3 \cdot 2^4) = 144$种
            \item 仅包含两种字符类型(如数字和小写字母),只满足相邻不同的密码有$6 \cdot 5^4 = 3750$种,但此时我们多算了$2 \cdot (3 \cdot 2^4) = 96$种单一字符类型, 故第二类共有$3 \cdot (3750 - 96) = 10962$种
        \end{itemize}
        不同密码的总个数为$36864 - 144 - 10962 = 25758$.
    \end{solution}

    \question 阿绿想在她的表演服装上缝 6 颗相同的红色钮扣、3 颗相同的绿色钮扣和 3 颗相同的黄色钮扣。若所有钮扣需竖直地排成一直线,且相邻钮扣不同色,则阿绿有多少种排列方法?
    \begin{solution}
        假设红、绿、黄钮扣为$R,G,Y$,先将 $6R$ 排成一列,共 $7$ 个间隔。由于同色不相邻,必须将中间 $5$ 个间隔放入 $3G$ 与 $3Y$ 中的五个。

        情况一: 5 个间隔放 $3G2Y$ 且剩下 $1Y$ 放头或尾端。排列数有$\comb{5}{3}\cdot 2$。

        情况二: 5 个间隔放 $2G3Y$ 且剩下 $1G$ 放头或尾端。排列数有$\comb{5}{3}\cdot 2$。

        情况三: 5 个间隔放 $3G2Y$ 且剩下 $1Y$ 不放头尾。可能情况有$G\cdot G\cdot Y\cdot Y\cdot \boxed{GY},G\cdot G\cdot Y\cdot Y\cdot \boxed{YG}$,排列数各有 $\dfrac{5!}{2!2!}$。

        故排列方法共有
        \[
        \comb{5}{3}\cdot 2+\comb{5}{3}\cdot 2+\frac{5!}{2!2!}+\frac{5!}{2!2!}=100
        \]
    \end{solution}

    
    \question 在一个 $3\times3$ 的方格中,每个小方格被涂成红、白、蓝、绿四种颜色中的一种,要求任何 $2\times2$ 的小方格块都恰好包含四种不同的颜色。共有多少种不同的涂色方案？
    \begin{solution}
        先固定中心格:中心格有 $4$ 种颜色可选。再考虑中心行的左右两格:
        \begin{itemize}
        \item 两格同色。共有 $3$ 种颜色可供选择,此时上下行个别有$2$种填色方法,共 $2^2=4$ 种合法填法。  
        \item 两格异色。左格 $3$ 种选,右格再 $2$ 种选,共 $3\cdot2=6$ 种。在此情况下,其余 4 个角格的颜色被唯一地确定。
        \end{itemize}
        故方法数共有 
        \[
        4 \cdot (3 \cdot 4 + 6) =72
        \]
    \end{solution}

    \question Rita 正在给一朵花涂色。她已经涂好了花心和花茎。接下来,她将用红色、橙色、黄色和蓝色给六片花瓣涂色,每片花瓣只用一种颜色。规则如下:
    \begin{itemize}
        \item 相邻的花瓣不能涂相同的颜色；
        \item 不一定要用到所有四种颜色。
    \end{itemize}
    求共有多少种不同的涂法数?
    \begin{figure}[H]
    \centering
    \includegraphics[width=0.3\textwidth]{images/image170.png}
    \end{figure}
    \begin{solution}
        不失一般性,假设顶端花瓣为红色。

        情况 1:红色出现三次。由于相邻花瓣不能同色,红色只能出现在特定的三片花瓣上,其余三片花瓣可以任意用其他3种颜色涂色,因此这种情况下共有
        \[
        3^3 = 27 \text{ 种涂法。}
        \]
        情况 2:红色出现两次。分情况讨论,
        \begin{itemize}
            \item 两红色花瓣之间相隔两个花瓣。与顶端红色花瓣相邻的两片花瓣可以各自用3种颜色涂色,剩余两片花瓣可以各自用2种颜色涂色,总共有
            \[
            3 \cdot 3 \cdot 2 \cdot 2 = 36 \text{ 种方式。}
            \]
            \item 两红色花瓣之间相隔一个花瓣。两红色花瓣可能在顶端及左侧或右侧,红色花瓣之间的花瓣可以用3种颜色,其他三片花瓣根据两种或三种颜色组合有 $6+6=12$ 种方式,因此该配置共有
            \[
            2 \cdot 3 \cdot 12 = 36 \text{ 种方式。}
            \]
        \end{itemize}
        情况 3:红色出现一次。按顺时针给花瓣涂色（从红色顶端开始）:第一片花瓣有 3 种选择,接下来的四片花瓣每片有 2 种选择（不能和相邻花瓣相同,也不能用红色）,因此共有
        \[
        3 \cdot 2 \cdot 2 \cdot 2 \cdot 2 = 48 \text{ 种方式。}
        \]
        顶端花瓣可以选择 4 种颜色,故
        \[
        N = 4 \cdot (27 + (36+36) + 48) = 732
        \]
    \end{solution}

    \question 一位农夫拥有一个矩形农田,它被划分成 $2\times2$ 的四块矩形区域。在每块区域中,农夫将种植一种作物:玉米、小麦、大豆或土豆。农夫不希望玉米和小麦出现在相邻的两块地中,也不希望大豆和土豆出现在相邻的两块地中。在这些限制下,农夫有多少种不同的方法为这四块地选择作物？
    \begin{solution}
        注意到:对于每种作物,有恰好一种作物是不能与其相邻的。不失一般性,设左上角种的是小麦。

        情况一:左上角右边和下边两个相邻地块种的是相同的作物。这两块地不能种玉米（因为玉米不能与小麦相邻）,因此它们可以是小麦、大豆或土豆中的一种,共有 $3$ 种选择,对于右下角地块,只要与周围的两块地都不冲突即可,有 $3$ 种选择。
        
        情况二:右边和下边两块地种的是不同的作物。可从 $3$ 种合法作物中选择两个不同的,并安排在右和下两个位置,共有 $3 \cdot 2 = 6$ 种方式,此时右下角地块的选择只有 $2$ 种。
        
        综合两种情况相加并乘以左上角作物的 $4$ 种选择,总方法数为
        \[
        4 \cdot (3 \cdot 3 + 6 \cdot 2) = 84 
        \]
    \end{solution}
            
\question 三双不同的鞋子被排成一列,要求不能有一只左脚鞋与另一双的右脚鞋相邻。问共有多少种排法？
\begin{solution}
    设鞋子编号如下:右脚鞋:$1,\ 2,\ 3$,左脚鞋:$4,\ 5,\ 6$,其中 $n$ 与 $n+3$ 配对（例如 $1$ 和 $4$ 是一对）。

    不妨设右脚鞋 $1$ 出现在排列中的第二个位置。此时我们枚举前三个位置可能的情况:
    \begin{itemize}
        \item 剩下的三个鞋子只能是
            \[
            654 \quad \text{或} \quad 645
            \]
            因为 $5$ 或 $4$ 不可以和 $3$（另一双的右脚鞋）相邻。
        \item 剩下的三个鞋子只能是
            \[
            563
            \]
            （不能选 $365$,因为 $6$ 和 $5$ 是不同双鞋的左、右脚,会相邻）
        \item 剩下的三个鞋子可以是
            \[
            256 \quad \text{或} \quad 652
            \]
    \end{itemize}
    所以总方法数为
    \[
    (2 + 1 + 2) \times 12 = \boxed{60}
    \]
    其中乘上 $12$ 是因为我们一开始假设了鞋子 $1$ 出现在第二位,实际有 $2$ 种选哪个右脚鞋先出现,$6$ 种选配对顺序,所以总共乘 $12$。
    \textcolor{red}{实际有 $2$ 种选哪个右脚鞋先出现?}
\end{solution}

    \question 竹东高中的多元选修课程共开设了六门选修课:$A,B,C $为第一类选修课,$D,E,F$ 为第二类选修课,要求每名同学须从中选修三门课,第一类选修课至少要选两门。现有甲、乙、丙三位同学选课,则任意一位同学与其他两位同学均至少有两门相同选修课的选法共有几种？
    \begin{solution}
        每个同学可能的选法有$$ABC, ABD, ABE, ABF, ACD, ACE, ACF, BCD, BCE, BCF,$$共 10 种；甲、乙、丙三人共有 $10^3=1000$ 种选法。

        不符规定的选法分析如下:
        \begin{itemize}
            \item 第一人选了 $ABC$:第二人在第一类有 $\comb{3}{2}=3$ 种选法,第三人在第一类只能与第二人重复一门课,因此有 2 种选法；第二人在第二类的课程有 3 种选择,第三人有 2 种选择,共有 $3\cdot2 \cdot 3\cdot2 =108$ 种。
            \item 第一人、第二人在第一类完全相同:第三人在第一类只有 2 种选择,在第二类课程中三人任选再扣除三人完全相同,即 $3^3-3$；共有 $\comb{3}{2} \cdot 3 \cdot 2 \cdot (3^3-3) =432$ 种。
            \item 三人在第一类彼此最多重复一门课:第一类的选择有 $3!=6$,第二类与情况(2)相同,即 $3^3-3$；共有 $6 \cdot (3^3-3) =144$ 种。
        \end{itemize}
        因此符合要求的选法共有
        \[
        1000-108-432-144 = 316
        \]
    \end{solution}

    \question 下图显示了一个宽为3格、高为3格的点阵,共包含9个小正方形。Carl 在这些正方形的边上放置长为 \(\dfrac{1}{2}\) 英寸的牙签,要求形成一个不相交的闭合环。图中某些格子内标注了一个数字,表示该格子被牙签覆盖的边数。若无数字,则该格子可以被任意数量的牙签覆盖。求 Carl 放置牙签的可能方法数。  
    \begin{figure}[H]
        \centering
        \includegraphics[width=0.4\textwidth]{images/image7.png}
    \end{figure}
    \begin{solution}
        由于牙签回路不能穿过中间的五列,所以整个回路只有两种大致布局:

        情况一:不穿过第二行。只剩下 $2$ 种回路（对称）。
        
        情况一:必须穿过第二行（左、右各穿一次）。则
        \begin{itemize}
          \item 中间四个标 $1$ 的小方格,各有两种放牙签方式,共 $2^4$;
          \item 回路在左右两侧穿过第二行,各有 $3$ 种走法,共 $3^2$。
        \end{itemize}
        因此所有合法回路总数为
        \[
        2^4 \cdot 3^2 + 2 = 146
        \]
    \end{solution}

    \question 有一张由 $5\times6$ 个正方形组成的格线纸,如右图。小强想沿着实线以向左、向右及向下的方向将格线纸剪成两张面积相等的纸张,并且先由 $A$ 点向下剪一格,最后从 $B,C,D,E,F$ 中某一点剪断纸张。问有多少种不同的剪法?
    \begin{figure}[H]
        \centering
        \includegraphics[width=0.3\textwidth]{images/image95.png}
    \end{figure}
    \ifprintanswers
    \begin{figure}[H]
        \centering
        \includegraphics[width=0.4\textwidth]{images/image96.jpg}
    \end{figure}
    \fi
    \begin{solution}
        格线纸共有 $5 \times 6 = 30$ 个方格。设剪线左侧从上往下第 $1,2,3,4,5$ 行的方格数分别为 $2,x,y,z,w$,需满足
        \[
        2 + x + y + z + w = 15,\quad 1 \le x,y,z,w \le 5, \quad x,y,z,w \in \mathbb{N}
        \]
        令 $x' = x - 1, y' = y - 1, z' = z - 1, w' = w - 1$,则
        \[
        x' + y' + z' + w' = 9, \quad 0 \le x',y',z',w' \le 4
        \]
        总方法数即不考虑$0 \le x',y',z',w' \le 4$的情况扣除 $x',y',z',w'$ 中至少有一个大于 $4$ 的情况:
        \[
        \comb{9+4-1}{4-1} - 4 \cdot \comb{4+4-1}{4-1} = 80
        \]
    \end{solution}

    \question 考虑集合 $\{1,2,3,\dots,2024\}$ 的所有恰有 1000 个元素的子集。对于每个这样的子集 $S$,记 $m(S)$ 为 $S$ 中的最小元素。求所有 $m(S)$ 的算术平均数。
    \begin{solution}
        设最小元素为 $a$,则有选择余下 $999$ 个元素的方法数为
        \[
        \comb{2024 - a}{999}, \quad a=1,2,\dots,1025 \quad (1025=2024-999).
        \]
        算术平均数为
        \[
        \frac{\displaystyle \sum_{a=1}^{1025} a \cdot \comb{2024 - a}{999}}{\displaystyle\sum_{a=1}^{1025} \comb{2024 - a}{999}}=\frac{\comb{2025}{1001}}{\comb{2024}{1000}} = \frac{2025}{1001}
        \]
        利用组合恒等式 $\sum_{k=r}^{n} \comb{k}{r} = \comb{n+1}{r+1}$,令 $k = 2024-a$,当 $a$ 从 1 到 1025 时,$k$ 从 2023 到 999:
        \[
        \sum_{a=1}^{1025} \comb{2024 - a}{999} = \sum_{k=999}^{2023} \comb{k}{999} = \comb{2024}{1000}
        \]
        利用恒等式:
\[
\sum_{i=r}^{n} \comb{i}{r} = \comb{n+1}{r+1}
\]

以及权重求和恒等式:
\[
\sum_{k=0}^{n} (k+1) \cdot \comb{n-k}{r} = \comb{n+1}{r+2}
\]

令 $k = a-1$,则 $a = k+1$,当 $a$ 从 1 到 1025 时,$k$ 从 0 到 1024:
\[
\sum_{a=1}^{1025} a \cdot \comb{2024-a}{999} = \sum_{k=0}^{1024} (k+1) \cdot \comb{2024-(k+1)}{999} = \sum_{k=0}^{1024} (k+1) \cdot \comb{2023-k}{999}
\]

利用恒等式 $\sum_{k=0}^{n-r} (k+1) \cdot \comb{n-k}{r} = \comb{n+2}{r+2}$:

令 $n = 2023$, $r = 999$:
\[
\sum_{k=0}^{1024} (k+1) \cdot \comb{2023-k}{999} = \comb{2025}{1001}
\]
        \textcolor{red}{(待解)}
    \end{solution}

    \question 用数字 1、2、3 组成 10 位数,要求其中数字 1 出现次数为偶数。求这样的 10 位数的个数。
    \begin{solution}
        设 \(a_n\) 表示长度为 \(n\) 的、由 1、2、3 组成且数字 1 出现偶数次的序列个数。
        \begin{itemize}
            \item 若首位为 2 或 3（共 2 种选择）,则余下 \(n-1\) 位仍需 1 出现偶数次,方案数为 \(2a_{n-1}\)。
            \item 若首位为 1（1 种选择）,则已有一个 1,余下 \(n-1\) 位中 1 需出现奇数次。余下 \(n-1\) 位的总序列数为 \(3^{n-1}\),其中 1 出现偶数次的为 \(a_{n-1}\),故出现奇数次的为 \(3^{n-1} - a_{n-1}\)。
        \end{itemize}
        因此,当 \(n > 1\) 时,有递推关系：
        \[
        a_n = 2a_{n-1} + (3^{n-1} - a_{n-1}) = a_{n-1} + 3^{n-1}, \quad a_1 = 2
        \]
        求通项
        \[
        a_n = a_1 + \sum_{k=1}^{n-1} 3^k =2 + \frac{3(3^{n-1} - 1)}{2}=\frac{3^n + 1}{2}
        \]
        代入 \(n=10\)得
        \[
        a_{10} = \frac{3^{10} + 1}{2} = 29525
        \]
    \end{solution}

    \question 某语言只使用字母 $A, B, C, D, E$,其中 $A$ 和 $E$ 为元音,$B, C, D$ 为辅音。一个字母序列称为单词,当且仅当它不含有相邻两个相同字母,且不含有相邻两个元音。问该语言中长度为 10 且以元音开头的单词有多少个？
    \begin{solution}
        设 $v_n$ 为长度为 $n$ 且以元音开头的单词数,$c_n$ 为长度为 $n$ 且以辅音开头的单词数。据题意,
        \[
        v_1=2,\ c_1=3
        \]
        若 $n\ge2$,则  
        \[
        v_n=2c_{n-1},
        \]
        因为在辅音开头的 $(n-1)$ 字母单词前可加 $A$ 或 $E$。且有  
        \[
        c_n=3v_{n-1}+2c_{n-1},
        \]
        其中 $3v_{n-1}$ 意味可在元音开头的 $(n-1)$ 字母单词前加 $B, C, D,2c_{n-1}$ 意味可在辅音开头的 $(n-1)$ 字母单词前加不同于首字母的两个辅音。递推如下表:
        \begin{center}
        \begin{tabular}{|c|c|c|}
        \hline
        $n$ & $v_n$ & $c_n$ \\
        \hline
        1 & 2 & 3 \\
        2 & 6 & 12 \\
        3 & 24 & 42 \\
        4 & 84 & 156 \\
        5 & 312 & 564 \\
        6 & 1128 & 2064 \\
        7 & 4128 & 7512 \\
        8 & 15024 & 27408 \\
        9 & 54816 & 99888 \\
        10 & 199776 & 364224 \\
        \hline
        \end{tabular}
        \end{center}
        因此,长度为 10 且以元音开头的单词共有 $199776$ 个。
    \end{solution}

    \question 有六个格子,初始时三个硬币从左到右依次占据前三个格子（记为 $X,Y,Z$）。一次移动是将某一枚硬币向右移动一格,前提是目标格子为空,且硬币之间不能相互跳跃（因此三枚硬币的相对顺序始终保持 $X$ 在 $Y$ 左边,$Y$ 在 $Z$ 左边）。问：有多少种不同的移动序列,能使三枚硬币最终全部移动到最右边的三个格子？
    \begin{figure}[H]
    \centering
    \includegraphics[width=0.3\textwidth]{images/image162.png}
    \end{figure}
    \begin{solution}
        我们将每个允许的移动序列看作由 $X,Y,Z$ 组成的字符串。例如,字符串 $ZZYXZ$ 表示先移动 $Z$ 一格,再移动 $Z,Y,X,Z$。  

        对于每个整数三元组 $(x,y,z)$,其中 $0 \le x,y,z \le 3$,定义 $S(x,y,z)$ 为使 X 移动 $x$ 格,Y 移动 $y$ 格,Z 移动 $z$ 格的移动序列数。欲求 $S(3,3,3)$。  

        由于硬币不能交叉且只能向右移动,当 $x > y$ 或 $y > z$ 或 $x > z$,有$S(x, y, z) = 0$。因此仅需考虑 $0 \le x \le y \le z \le 3$ 的情况。此时有
        \[
        S(x,y,z) = S(x-1,y,z) + S(x,y-1,z) + S(x,y,z-1)
        \]
        其中当 $x=0$ 时 $S(x-1,y,z)=0,y=0$ 时 $S(x,y-1,z)=0,z=0$ 时 $S(x,y,z-1)=0$,且  
        \[
        S(0,0,0) = 1, \quad S(1,0,0) = 0, \quad S(0,1,0) = 0, \quad S(0,0,1) = 1
        \]
        接下来通过表格逐步填入 $S(x, y, z)$,其中 $z = 1,2,3$,纵轴为 $y$,横轴为 $x$。
        \[  
        \begin{tabular}{c c c}
            \begin{tabular}{|c|cccc|}
            \multicolumn{1}{c}{} & \multicolumn{4}{c}{$x$} \\
            \cline{2-5}
            \multicolumn{1}{c|}{$y$} & 0 & 1 & 2 & 3 \\
            \cline{1-5}
            0 & 1 & 0 & 0 & 0 \\
            1 & 1 & 1 & 0 & 0 \\
            2 & 0 & 0 & 0 & 0 \\
            3 & 0 & 0 & 0 & 0 \\
            \cline{1-5}
            \end{tabular} &
            \begin{tabular}{|c|cccc|}
            \multicolumn{1}{c}{} & \multicolumn{4}{c}{$x$} \\
            \cline{2-5}
            \multicolumn{1}{c|}{$y$} & 0 & 1 & 2 & 3 \\
            \cline{1-5}
            0 & 1 & 0 & 0 & 0 \\
            1 & 2 & 3 & 0 & 0 \\
            2 & 2 & 5 & 5 & 0 \\
            3 & 0 & 0 & 0 & 0 \\
            \cline{1-5}
            \end{tabular} &
            \begin{tabular}{|c|cccc|}
            \multicolumn{1}{c}{} & \multicolumn{4}{c}{$x$} \\
            \cline{2-5}
            \multicolumn{1}{c|}{$y$} & 0 & 1 & 2 & 3 \\
            \cline{1-5}
            0 & 1 & 0 & 0 & 0 \\
            1 & 3 & 6 & 0 & 0 \\
            2 & 5 & 16 & 21 & 0 \\
            3 & 5 & 21 & 42 & 42 \\
            \cline{1-5}
            \end{tabular} \\
            $z=1$ & $z=2$ & $z=3$
        \end{tabular}
        \]
        由表得 
        \[
        S(0,0,1)=1(Z),\quad S(0,1,1)=1(ZY),\quad S(1,1,1)=1
        (ZYX)
        \]
        因此序列的数量为$S(3,3,3)=42$。
    \end{solution}

    \question 给定数字集合 $\{1, 2, 4, 8, 16, 32, 64, 128\}$,每个数字最多可使用 $3$ 次,问：用这些数字之和恰好等于 $530$ 的方案数有多少种？
    \begin{solution}
        每个数字 $2^k$ 可以出现 $0,1,2,3$ 次,对应多项式
        \[
        1 + x^{2^k} + x^{2\cdot 2^k} + x^{3\cdot 2^k}.
        \]
        考虑生成函数
        \[
        (1+x+x^2+x^3)(1+x^2+x^4+x^6)(1+x^4+x^8+x^{12}) \cdots (1+x^{128}+x^{256}+x^{384}).
        \]
        因式分解得
        \[
        \frac{x^4-1}{x-1}\cdot\frac{x^8-1}{x^2-1}\cdots \frac{x^{512}-1}{x^{128}-1}
        = (1+x+x^2+\cdots+x^{255})(1+x^2+x^4+\cdots+x^{510}).
        \]
        要得到 $x^{530}$,记作 $x^{2i}\cdot x^{530-2i}$,其中 $10 \le i \le 127$,共有 118 个 $i$ 的取值,因此方法数为 $118$。
    \end{solution}

    \question 求正整数有序三元组 $(a_1, a_2, a_3)$ 的个数,使得
    \[
    a_1 + a_2 + a_3 = 2020, \quad a_1 \not\equiv 0 \pmod 2, \quad a_2 \not\equiv 0 \pmod 3, \quad a_3 \not\equiv 0 \pmod 4。
    \]
    \begin{solution}
        $a_1$ 为奇数,生成函数为
        \[
        x + x^3 + x^5 + \cdots = \frac{1}{1-x} - \frac{1}{1-x^2} = \frac{x}{1-x^2}
        \]
        $a_2$ 不可被 3 整除,生成函数为
        \[
        x + x^2 + x^4 + x^5 + \cdots = \frac{1}{1-x} - \frac{1}{1-x^3} = \frac{x(1-x^2)}{(1-x)(1-x^3)}
        \]
        $a_3$ 不可被 4 整除,生成函数为
        \[
        x + x^2 + x^3 + x^5 + x^6 + x^7 + \cdots = \frac{1}{1-x} - \frac{1}{1-x^4} = \frac{x(1-x^3)}{(1-x)(1-x^4)}
        \]
        三式相乘得总生成函数
        \[
        \frac{x}{1-x^2} \cdot \frac{x(1-x^2)}{(1-x)(1-x^3)} \cdot \frac{x(1-x^3)}{(1-x)(1-x^4)} = \frac{x^3}{(1-x)^2(1-x^4)}
        \]
        因此所求为 $x^{2020}$ 的系数,即 $x^{2017}$ 在
        \[
        (1 + 2x + 3x^2 + \cdots)(1 + x^4 + x^8 + \cdots)
        \]
        中的系数,即
        \[
        2 + 6 + 10 + \cdots + 2018 = 505 \cdot 1010 = 510050
        \]
    \end{solution}

    \question 给定正立方体 $ABCD-EFGH$,其中 $ABCD$ 为底面,$EFGH$ 为顶面,$A$ 与 $E$ 对应,$B$ 与 $F$ 对应,$C$ 与 $G$ 对应,$D$ 与 $H$ 对应。一只小虫从顶点 $A$ 出发,每分钟沿一条棱移动到相邻顶点（每次移动一条棱）。问：经过恰好 $9$ 分钟后,小虫到达顶点 $G$ 的不同路径数为多少？
    \begin{solution}
        转换矩阵为
        \[
        A=
        \begin{bmatrix}
        0 & 1 & 0 & 1 & 1 & 0 & 0 & 0 \\
        1 & 0 & 1 & 0 & 0 & 1 & 0 & 0 \\
        0 & 1 & 0 & 1 & 0 & 0 & 1 & 0 \\
        1 & 0 & 1 & 0 & 0 & 0 & 0 & 1 \\
        1 & 0 & 0 & 0 & 0 & 1 & 0 & 1 \\
        0 & 1 & 0 & 0 & 1 & 0 & 1 & 0 \\
        0 & 0 & 1 & 0 & 0 & 1 & 0 & 1 \\
        0 & 0 & 0 & 1 & 1 & 0 & 1 & 0
        \end{bmatrix}
        \]
        发现
        \[
        A^9=
        \begin{bmatrix}
        0 & 4921 & 0 & 4921 & 4921 & 0 & 4920 & 0 \\
        4921 & 0 & 4921 & 0 & 0 & 4921 & 0 & 4920 \\
        0 & 4921 & 0 & 4921 & 4920 & 0 & 4921 & 0 \\
        4921 & 0 & 4921 & 0 & 0 & 4920 & 0 & 4921 \\
        4921 & 0 & 4920 & 0 & 0 & 4921 & 0 & 4921 \\
        0 & 4921 & 0 & 4920 & 4921 & 0 & 4921 & 0 \\
        4920 & 0 & 4921 & 0 & 0 & 4921 & 0 & 4921 \\
        0 & 4920 & 0 & 4921 & 4921 & 0 & 4921 & 0
        \end{bmatrix}
        \]
        故由
        \[
        A^9
        \begin{bmatrix}
        1 \\ 0 \\ 0 \\ 0 \\ 0 \\ 0 \\ 0 \\ 0
        \end{bmatrix}
        =
        \begin{bmatrix}
        0 \\ 4921 \\ 0 \\ 4921 \\ 4921 \\ 0 \\ 4920 \\ 0
        \end{bmatrix}
        \]
        可知$A \to G$ 有 $4920$ 种路线。
        \textcolor{red}{(待递推解)}
    \end{solution}

\end{questions}

\pagebreak

\begin{center}
  {\fontsize{30pt}{26pt}\selectfont
    \hypertarget{概率、期望值}{概率、期望值} \label{概率、期望值}
  }
\end{center}
\separator
\vspace{1pt}

\begin{questions}
    \question 正八面体有 $8$ 个全等的正三角形面。任取 $3$ 个面,其面心组成的三角形是正三角形的概率是多少？
    \begin{solution}
        正八面体的对偶多面体是正方体,其面心对应正方体的 $8$ 个顶点。  
        等边三角形对应正方体中形如正四面体侧面的三角形,共 $8$ 个,因此概率为
        \[
        \frac{8}{\comb83} = \frac{1}{7}
        \]
    \end{solution}

    \question 给定正八边形 $PQRSTUVW$。从中随机选择 $4$ 条边并将其延长为直线,问：这 $4$ 条直线相交所形成的四边形恰好包含原八边形的概率是多少？
    \begin{solution}
        选择的四条边相邻,或未选的三条相邻边及一条不相邻的边,皆不能形成包含原八边形的四边形,这样的组合数有
        \[
        8 + 8 \cdot 3 = 32
        \]
        所求概率为
        \[
        \frac{\comb{8}{4}-32}{\comb{8}{4}}=\frac{19}{35}
        \]
    \end{solution}
    
    \question $X$是从$\{1,2,\dots,9\}$中随机抽取$3$个不同的数排列出的最大的三位数,
        $Y$是从$\{1,2,\dots,8\}$中随机抽取$3$个不同的数排列出的最大的三位数,求$X>Y$的概率。
    \begin{solution}
        若 $X$ 的三个数字中包含 $9$,则由其组成的最大三位数一定大于由 $1$ 到 $8$ 中任意三个数字组成的最大三位数,此时 $X > Y$ 恒成立,这样的$X$排列数有$\comb82$,因此
        \[
        P(X > Y) = \frac{\comb82 \cdot \comb83}{\comb93 \cdot \comb83} = \frac{1}{3}
        \]
    \end{solution}
    
    \question 在数字 $1,2,3,4,5,6,7$ 的所有排列中,求至少有两个相邻数字不互质的概率。
    \begin{solution}
        先对奇数 $1,3,5,7$ 排列,有 $\perm44$ 种排法。
        
        数字 $6$ 不能与 $3$ 相邻,$4$ 个奇数间有 $5$ 个空位,去除 $3$ 的左右两个空位,剩 $3$ 个空位可放 $6$,有 $3$ 种放法。
        
        数字 $2$ 和 $4$ 放在剩余 $4$ 个空位,共有 $\perm42 $ 种排法。
        
        故至少有两个相邻数字不互质的概率为
        \[
        1 - \frac{\perm{4}{4} \cdot 3 \cdot \perm{4}{2}}{7!} = \frac{29}{35}
        \]
    \end{solution}

    \question 设 $A$ 和 $B$ 是从集合 $\{1,2,3,4,5,6\}$ 中随机选出的子集,允许 $A=B$。求事件“$A$ 包含于 $B$ 或 $B$ 包含于 $A$”的概率。
    \begin{solution}
        $A \subseteq B$ 当且仅当不存在元素 $x$ 满足 $x\in A$ 且 $x\notin B$,即概率为
        \[
        P(A\subseteq B) = \left(1-\left(\frac{1}{2}\right)^2\right)^6 = \left(\dfrac{3}{4}\right)^6
        \]
        同理 $P(B\subseteq A) = \left(\dfrac{3}{4}\right)^6$,且
        \[
        P(A=B) = \left(\frac{1}{2}\right)^6
        \]
        根据容斥原理,所求概率为
        \[
        2\left(\frac{3}{4}\right)^6 - \left(\frac{1}{2}\right)^6 = \frac{697}{2048}
        \]
    \end{solution}

\question Alice 和 Bob 各自有一副相同的牌,包含 3 张红牌、3 张白牌和 3 张蓝牌。他们轮流从自己的牌中随机抽取一张牌,且不放回。Alice 先抽。求 Alice 在 Bob 抽到任何红牌之前抽完她的所有红牌的概率。

\begin{solution}
每个人都将 9 张牌随机排列,考虑每个人的 3 张红牌的位置。共有
\[
\comb{9}{3}^2
\]
种可能。

考虑一条长度为 10 的序列,由 Alice 抽牌直到她抽到第 3 张红牌为止,再加上 Bob 在 Alice 抽到第 3 张红牌后接下来抽的牌。为了满足条件,这条序列中必须有 6 张红牌,我们在枚举这些情况。

因此概率为
\[
\frac{\comb{10}{6}}{\comb{9}{3}^2} = \frac{5}{168}.
\]
        \textcolor{red}{(待验证, lehigh 2025 q15)}
\end{solution}

    \question 掷一个公平的六面骰子连续十次。求出现恰好三个 6 连在一起的概率。（四个或更多连续的 6 不计。）
    \begin{solution}
        情况一:三个 6 在序列的一端,且相邻数字不是 6:有 2 种方式,对应概率为
        \[
        2 \cdot \left(\frac{1}{6}\right)^3 \cdot \frac{5}{6}
        \]
        情况二:三个 6 在序列中间,两边相邻数字都不是 6:有 6 种方式,对应概率为
        \[
        6 \cdot \left(\frac{1}{6}\right)^3 \cdot \left(\frac{5}{6}\right)^2
        \]
        总概率为
        \[
        \left(\frac{1}{6}\right)^3 \left(2 \cdot \frac{5}{6} + 6 \cdot \frac{25}{36}\right) = \frac{35}{1296}
        \]
    \end{solution}

    \question 投掷一枚均匀硬币 8 次。已知在前 3 次投掷中至少出现一次正面,求 8 次投掷中恰好出现 4 次正面的概率。
    \begin{solution} 
        前 3 次投掷中至少出现一次正面且8次投掷中恰好有4次正面的情形,即前3次及后5次有1正3正、2正2正、3正1正,方法数为
        \[
        \comb{3}{1}\comb{5}{3},\comb{3}{2}\comb{5}{2},\comb{3}{3}\comb{5}{1}
        \]
        8 次投掷中前3次都是反面的方法数为$2^5$,前3次至少1次正面的方法数为$2^8 - 2^5$,因此所求概率为
        \[
        \frac{\comb{3}{1}\comb{5}{3} + \comb{3}{2}\comb{5}{2} + \comb{3}{3}\comb{5}{1}}{2^8 - 2^5}
        = \frac{65}{224}
        \]
    \end{solution}

    \question 有三个碗,每个碗里各有 6 个球。现随机选择一个碗,再选择一个不同的碗,把第一个碗中的一个球移动到第二个碗。经过 5 次这样的移动后,三个碗再次各有 6 个球的概率是多少？
    \begin{solution}
        设碗标为 $A, B, C$。为了最终三个碗各有 6 个球,必须有两个碗各被拿走 2 个球,第三个碗只被拿走 1 个球。选择只拿走 1 个球的碗有 3 种方法,假设为碗 $A$。

        这个移球的回合有 5 种选择,并且球要移到的碗有 2 种选择。

        若从 $A$ 移到 $B$,则必须从 $C$ 移回 $A$,否则无法使 $B$ 和 $C$ 最终数量相等。这个移动的回合有 4 种选择。

        剩下 3 次移动必须是 2 次从 $B$ 到 $C$,1 次从 $C$ 到 $B$,这三次的顺序有 3 种。

        综上,有 $3 \cdot 5 \cdot 2 \cdot 4 \cdot 3$ 种序列使三碗最终各有 6 个球,每次移动有 $3 \cdot 2 = 6$ 种可能,因此总共有 $6^5$ 种可能的移动序列,因此所求概率为
        \[
        \frac{3 \cdot 5 \cdot 2 \cdot 4 \cdot 3}{6^5} = \frac{5}{108}
        \]
    \end{solution}

    \question 一个正立方体骰子六个面上分别有 $2,3,4,5,6,7$ 个点。随机移除一个点(每个点被移除的概率相等)后投掷这个骰子,求朝上的面的点数为奇数的概率。
    \begin{solution}
        移除一个点后,如果从偶数点的面移除一个点,该面变为奇数点。如果从奇数点的面移除一个点,该面变为偶数点。

        情况 1:从偶数点的面移除。则有 4 个奇数点面和 2 个偶数点面,掷出奇数点的概率为 $\dfrac{4}{6}$。

        情况 2:从奇数点的面移除。则有 2 个奇数点面和 4 个偶数点面,掷出奇数点的概率为 $\dfrac{2}{6}$。

        移除某个面的点的概率为该面点数除以 $2+3+4+5+6+7=27$,所求概率为
        \[
        \frac{2}{27}\cdot \frac{4}{6} + \frac{3}{27}\cdot \frac{2}{6} + \frac{4}{27}\cdot \frac{4}{6} + \frac{5}{27}\cdot \frac{2}{6} + \frac{6}{27}\cdot \frac{4}{6} + \frac{7}{27}\cdot \frac{2}{6} = \frac{13}{27}
        \]
    \end{solution}

    \question $A$ 和 $B$ 玩一场卡牌游戏。$A$ 有 6 张牌：2 张红色、2 张黄色、2 张绿色。$B$ 有 4 张牌：2 张紫色、2 张白色。两人轮流出牌,$A$ 先出。每回合,玩家随机选择一张自己手中的牌放到桌上。若某玩家在桌上放出两张同色牌,则该玩家获胜。求 $A$ 获胜的概率。
    \begin{solution}
        游戏一定在 $B$ 的第三回合之前结束,故$A$ 只能在自己的第二回合或第三回合获胜,且不能在第一回合直接获胜。

        情况 1:$A$ 第二回合获胜。$A$ 第二张牌必须和第一张同色。此时 $A$ 还剩 5 张牌,其中 1 张与第一张同色,所以概率为$\dfrac{1}{5}$。$A$ 和 $B$ 的第一张牌没有限制。

        情况 2:$A$ 第三回合获胜。需满足$A$ 第二张牌颜色与第一张不同,$B$ 第二张牌颜色与她的第一张不同且$A$ 第三张牌与他前两张牌之一同色,概率为
        \[
        \frac{4}{5} \cdot \frac{2}{3} \cdot \frac{1}{2} = \frac{4}{15}
        \]
        综上,$A$ 获胜的概率为
        \[
        \frac{1}{5} + \frac{4}{15} = \frac{7}{15}
        \]
    \end{solution}

    \question $C$参加一个比赛,比赛中没有平局,她会一直比赛直到输掉 2 场比赛为止,此时她将被淘汰,不再继续比赛。已知:
    \begin{itemize}
        \item 第一场比赛获胜的概率为 $\dfrac{1}{2}$
        \item 若上一场获胜,则下一场获胜的概率为 $\dfrac{3}{4}$
        \item 若上一场失败,则下一场获胜的概率为 $\dfrac{1}{3}$
    \end{itemize}
    求 $C$ 在被淘汰之前(即输掉第 2 场之前)恰好赢得 3 场比赛的概率。
    \begin{solution}
        $C$ 赢 3 场比赛且输少于 2 场比赛的概率,可能的情况有赢 3场,输 0场或赢 3场,输 1场。以 $W$ 表示胜利,$L$ 表示失败,比赛结果可能的序列为
        \[
        WWW, \quad LWWW, \quad WLWW, \quad WWLW.
        \]
        故所求概率为
        \[
        \frac{1}{2}\cdot\frac{3}{4}\cdot\frac{3}{4} + \frac{1}{2}\cdot\frac{1}{3}\cdot\frac{3}{4}\cdot\frac{3}{4} + \frac{1}{2}\cdot\frac{1}{4}\cdot\frac{1}{3}\cdot\frac{3}{4} + \frac{1}{2}\cdot\frac{3}{4}\cdot\frac{1}{4}\cdot\frac{1}{3} = \frac{7}{16}
        \]
    \end{solution}

    \question $A$ 有 12 个桶:3 个绿色、3 个红色、3 个蓝色和 3 个黄色。他将球随机放入桶中:
    \begin{itemize}
        \item 将 4 个球随机放入 3 个绿色桶(每个球独立地等概率选择一个绿色桶)
        \item 将 3 个球随机放入 3 个红色桶
        \item 将 2 个球随机放入 3 个蓝色桶
        \item 将 1 个球随机放入 3 个黄色桶
    \end{itemize}
    求存在一个绿色桶,其中的球数皆多于其他 11 个桶中任何一个桶的球数的概率。
    \begin{solution}
        以$(a,b,c)$ 表示 3 个桶中球的无序分布。黄色桶分布只能是 $(1,0,0)$。

        蓝色桶:将 2 个球放入 3 个桶,总共有 $3^2=9$ 种方法。
        \begin{itemize}
            \item $(2,0,0)$: 2 个球在同一个桶,有 3 种方式,
            \item $(1,1,0)$: 2 个球分在不同桶,有 $9-3=6$ 种方式。
        \end{itemize}
        红色桶:将 3 个球放入 3 个桶,总共有 $3^3=27$ 种方法。
        \begin{itemize}
            \item $(3,0,0)$: 3 个球在同一个桶,有 3 种方式,
            \item $(1,1,1)$: 每个桶 1 个球,有 $\comb{3}{2}=6$ 种方式,
            \item $(2,1,0)$: 其余 18 种方式。
        \end{itemize}
        绿色桶:将 4 个球放入 3 个桶,总共有 $3^4=81$ 种方法。
        \begin{itemize}
            \item $(4,0,0)$: 3 种方式,
            \item $(3,1,0)$: $4\cdot 3!=24$ 种方式,
            \item $(2,1,1)$: $4\cdot 3 \cdot \dfrac{3!}{2!}=36$ 种方式,
            \item $(2,2,0)$: 其余 18 种方式。
        \end{itemize}
        要使绿色桶中的球比其他 11 个桶都多,可能的情况及概率如下:
        \begin{center}
        \begin{tabular}{|c|c|c|c|c|}
        \hline
        \text{绿色} & \text{红色} & \text{蓝色} & \text{黄色} & \text{概率} \\
        \hline
        $(4,0,0)$ & 任意 & 任意 & 任意 & $\frac{3}{81}$ \\
        \hline
        $(3,1,0)$ & 不是 $(3,0,0)$ & 任意 & 任意 & $\frac{24}{81} \cdot \frac{24}{27}$ \\
        \hline
        $(2,1,1)$ & $(1,1,1)$ & $(1,1,0)$ & 任意 & $\frac{36}{81} \cdot \frac{6}{27} \cdot \frac{6}{9}$ \\
        \hline
        \end{tabular}
        \end{center}
        其中分布 $2/2/0$ 不符合条件,因为没有单个绿色桶的球数比其他桶多,故所求概率为
        \[
        \frac{3}{81} + \frac{24}{81} \cdot \frac{24}{27} + \frac{36}{81} \cdot \frac{6}{27} \cdot \frac{6}{9} = \frac{89}{243}
        \]
    \end{solution}

    \question 有三个盒子:
    \begin{itemize}
        \item 盒子 1 里有 1 个金球和 1 个黑球；
        \item 盒子 2 里有 1 个金球和 2 个黑球；
        \item 盒子 3 里有 1 个金球和 3 个黑球。
    \end{itemize}
    每次从一个盒子里随机抽取一个球,盒子里每个球被选中的概率相等。过程如下:
    \begin{enumerate}
        \item 从盒子 1 中随机抽一个球,放入盒子 2；
        \item 然后从盒子 2 中随机抽一个球,放入盒子 3；
        \item 最后从盒子 3 中随机抽一个球。
    \end{enumerate}
    求从盒子 3 中抽到金球的概率。
        \ifprintanswers
        \begin{figure}[H]
            \centering        
            \includegraphics[width=0.5\textwidth]{images/image171.png}
        \end{figure}
        \fi
    \begin{solution}
        从盒子 1 抽到金球及黑球的概率皆为$\dfrac{1}{2}$。
        
        若从盒子 1 抽到金球,盒子 2 将有 2 个金球和 2 个黑球；若从盒子 1 抽到黑球,盒子 2 将有 1 个金球和 3 个黑球。

        若盒子 2 中有 2 金 2 黑,抽到金球的概率为 $\dfrac{1}{2}$。若盒子 2 中有 1 金 3 黑,抽到金球的概率为 $\dfrac{1}{4}$。
        
        放入盒子 3后,盒子 3 有 2 金 3 黑的概率为
        \[
        \frac{1}{2}\cdot \frac{1}{2} + \frac{1}{2}\cdot \frac{1}{4} = \frac{3}{8}, 
        \]
        此时从盒子 3 抽到金球的概率为
        \[
        \frac{3}{8} \cdot \frac{2}{5} + \frac{5}{8} \cdot \frac{1}{5} =\frac{11}{40}
        \]
    \end{solution}

    \question 八个人分属三个家庭,围坐在圆桌旁。其中两个家庭各有 3 人,另一个家庭有 2 人。求每个人都有至少一个来自不同家庭的邻座的概率。
    \begin{solution}
        对座位编号,并将同一家庭的成员视为不可区分。8 个人的座位排列总数为:
        \[
        \frac{8!}{3! 3!  2!} = 560
        \]
        现探讨至少有一人使得其邻座都来自自己家庭,即3 人家庭坐在一起的排列数。设家庭 $A$ 和 $B$ 各有 3 人,家庭 $C$ 有 2 人。要求家庭 $A$ 的 3 人坐在一起,视他们为一个整体,有 8 种方式选择这个整体的起始位置,剩余 5 个人的排列数为 $\dfrac{5!}{3! 2!}$。因此家庭 $A$ 坐在一起的排列数为
        \[
        8 \cdot \frac{5!}{3! 2!} = 80
        \]
        同理,家庭 $B$ 坐在一起的排列数也是 80。若家庭 $A$ 和 $B$ 都坐在一起,则将两个家庭各看作一个整体,排列数有
        \[
        8 \cdot \frac{3!}{2!} = 24
        \]
        由容斥原理,至少有一人使得其邻座都来自自己家庭的排列数为
        \[
        80 + 80 - 24 = 136
        \]
        所求概率为
        \[
        \frac{560 - 136}{560} = \frac{53}{70}
        \]
    \end{solution}

    \question 八支队伍参加一项比赛,采用单循环赛制(即任意两支队伍之间恰好进行一场比赛)。每场比赛没有平局,两支队伍各有 $\dfrac{1}{2}$ 的获胜概率。求每支队伍都至少输一场且至少赢一场的概率。
    \begin{solution}
        共有 $\comb{8}{2}= 28$场比赛,因此总共有$2^{28}$种可能的比赛结果。

        当一支队伍全胜,同时另一支队伍全败时,选择全胜队有 8 种,全败队有 7 种,它们之间的比赛及与其他队伍的比赛结果已经确定,剩下的 $28-7-6=15$ 场比赛未确定,因此至少有一支队伍全胜或全败的结果数为
        \[
        8 \cdot 2^{21} + 8 \cdot 2^{21} - 8 \cdot 7 \cdot 2^{15}=2^{15}\cdot 968
        \]
        则每支队伍都至少输一场且至少赢一场的概率为
        \[
        1 - \frac{2^{15}\cdot 968}{2^{28}}  = \frac{903}{1024}
        \]
    \end{solution}

    \question 盒子中有大小,形状完全相同的3个红球,3个白球,现抛掷一枚质地均匀的骰子, 掷出几点就从盒子中取出几个球,求取出的球中红球个数大于白球个数的概率。
    \begin{solution}
        记从盒子中取出 $i$ 个球,取出的球中红球个数大于白球个数的概率为 $P_i$ ($i=1, 2, ..., 6$),则 
        \[
        \begin{aligned}
        P_1 &= \frac{\comb{3}{1}\comb{3}{0}}{\comb{6}{1}} = \frac{1}{2}, \quad
        P_2 = \frac{\comb{3}{2} \comb{3}{0}}{\comb{6}{2}} = \frac{1}{5}, \quad
        P_3 = \frac{\comb{3}{2}\comb{3}{1} + \comb{3}{3}\comb{3}{0}}{\comb{6}{3}}  = \frac{1}{2}, \\
        P_4 &= \frac{\comb{3}{3} \comb{3}{1}}{\comb{6}{4}} = \frac{1}{5}, \quad
        P_5 = \frac{\comb{3}{3} \comb{3}{2}}{\comb{6}{5}} = \frac{1}{2}, \quad
        P_6 = 0.
        \end{aligned}
        \]
        所以取出的球中红球个数大于白球个数的概率为 $$\frac{1}{6}\left(\frac{1}{2}+\frac{1}{5}+\frac{1}{2}+\frac{1}{5}+\frac{1}{2}+0\right)=\frac{19}{60}$$
    \end{solution}

    \question $A,B,C$ 独立参加考试。已知:
    \begin{itemize}
        \item $A$ 及格且 $B$ 不及格的概率为 $\dfrac{3}{20}$；
        \item $B$ 及格且 $B$ 不及格的概率为 $\dfrac{1}{4}$；
        \item $A$ 和 $C$ 都及格的概率为 $\dfrac{2}{5}$。
    \end{itemize}
    求至少有一人不及格的概率。
    \begin{solution}
        设$A,B,C$及格的概率为$a,b,c$,由于及格与否为独立事件,据题意有
        \[
        a(1-b) = \frac{3}{20}, \quad b(1-c) = \frac{1}{4}, \quad ac = \frac{2}{5}
        \]
        解得
        \[
        (8b+1)(4b-3) = 0 \Rightarrow b = \frac{3}{4} >0
        \]
        因此至少有一人不及格的概率为
        \[
        1 - abc  = 1 - \frac{3}{5} \cdot \frac{3}{4} \cdot \frac{2}{3}= \frac{7}{10}
        \]
    \end{solution}

    \question 将一块正方形厚纸板划分成 $5 \times 5 = 25$ 个面积相等的小正方格。现将三枚不同颜色的棋子随机放置在小正方格的中心,每个小正方格最多放一枚棋子。求这三枚棋子的位置不共线(即构成三角形)的概率。
    \begin{solution}
        三点共线的方法数有
        \[
        16 \cdot \comb{3}{3} + 4 \cdot \comb{4}{3} + 12 \cdot \comb{5}{3} =152
        \]
        因此三点不共线（即能构成三角形）的概率为
        \[
        \frac{\comb{25}{3} - 152}{\comb{25}{3}} = \frac{537}{575}
        \]
    \end{solution}

\question 从圆内随机选择四个不同的点,这四个点确定四个三角形(包括可能出现的退化三角形)。随机选择其中两个三角形,求圆心位于所选两个三角形的并集内的概率。
\begin{solution}
首先计算圆心位于四个三角形并集内的概率。圆心不在四个三角形的并集内,当且仅当四点位于某条直径的同一侧,即它们在圆上的径向投影也在直径同侧。此时恰好存在一个点,其余三个点位于该点顺时针方向 180 度的弧内。对于给定点,其余三个点落在该弧内的概率为
\[
\left(\frac{1}{2}\right)^3 = \frac{1}{8},
\]
因此四个点中任一点作为该点的概率为 $4 \cdot \frac{1}{8} = \frac{1}{2}$。因此圆心不在凸包内的概率为 $\frac{1}{2}$,所以圆心在凸包内的概率也为 $\frac{1}{2}$。

接下来考虑随机选择两三角形。四点确定四个三角形,它们的交集互不重叠。共有 $\comb{4}{2} = 6$ 对三角形。圆心在随机选择的一对三角形交集内的概率为 $\frac{1}{6} = \frac{2}{12}$。因此圆心在所选两三角形并集内的概率为
\[
\frac{1}{2} - \frac{1}{12} = \frac{5}{12}.
\]
We first calculate the probability that the center is in the union of the four triangles. The center does not lie in the union of the four triangles iff the four points lie on the same side of some diameter of the circle iff this is true of their radial projections onto the circle iff for one of these, the other three are within 180 degrees clockwise of the point, and such a point is essentially unique. Given a point, the probability that the other three lie in this arc is $\left(\frac{1}{2}\right)^3=\frac{1}{8}$, so the probability that the center does not lie in the convex hull is $4 \cdot \frac{1}{8}=\frac{1}{2}$.Ignoring possibilities with probability 0, the center lies in the union of the four triangles iff it lies in the intersection of two of the triangles, and these possibilities are disjoint. Since there are six pairs, the probability of the center lying in a randomly selected intersection is $\frac{1}{6} = \frac{2}{12}$. The center is in the union of the two selected triangles iff it is in the union of the four but not in the intersection of the two which we did not select. This probability is $\frac{1}{2} - \frac{1}{12}$.
\end{solution}

    \question $A$ 和 $B$ 轮流进行游戏,$A$ 先开始。每轮中:
    \begin{itemize}
        \item $A$ 投掷一枚公平硬币,若出现正面则 $A$ 获胜,游戏结束
        \item 若 $A$ 未获胜,则 $B$ 投掷一枚公平骰子,若出现 3 则 $B$ 获胜,游戏结束
        \item 若 $B$ 也未获胜,则进入下一轮,$A$ 再次投掷
    \end{itemize}
    游戏持续进行直到某人获胜。求 $B$ 获胜的概率。
    \begin{solution}
        设事件 $N$ 表示 $B$ 掷出的数不是 3,则$B$ 获胜的可能顺序为
        \[
        T3, \quad TNT3, \quad TNTNT3, \;\ldots
        \]
        每种顺序的概率分别为:
        \[
        \frac{1}{2}\cdot \frac{1}{6} = \frac{1}{12},\quad
        \frac{1}{2}\cdot\frac{5}{6}\cdot\frac{1}{2}\cdot \frac{1}{6} = \frac{5}{144}, \;\ldots
        \]
        这是一个等比数列,公比为 $\displaystyle \frac{1}{2}\cdot \frac{5}{6}=\frac{5}{12}$,因此 $B$ 获胜的概率为
        \[
        \frac{\frac{1}{12}}{1 - \frac{5}{12}} = \frac{1}{7}
        \]
    \end{solution}

    \question 圆圈中有 25 人,随机选出 3 人,求所选三人互不相邻的概率。
    \begin{solution}
        先固定被选中的一人 $A$,其余两人必须从不与 $A$ 相邻的 $22$ 人中选出,且这 $22$ 人在圆上形成一条线段,因此其中相邻的成对共有 $21$ 对,于是可行对数为
        \[
        \comb{22}{2}-21
        \]
        由于三人中谁充当 $A$ 都可以,上式被重复计数 $3$ 次,故满足条件的三人组数为
        \[
        \frac{25}{3}(\comb{22}{2}-21)
        \]
        所选三人互不相邻的概率为
        \[
        \frac{\frac{25}{3}(\comb{22}{2}-21)}{\comb{25}{3}}=\frac{35}{46}
        \]
    \end{solution}

    \question 三对已婚夫妇随机围坐在一张圆桌旁,求没有夫妻相邻的概率。
    \begin{solution}
        先固定 $A$ 先生的座位,其他五人的排列共有 $5!$ 种。

        情况 1:$A$ 女士坐在 $A$ 先生对面。此时$B$ 先生 有 4 种选择,$B$ 女士不坐他旁边有 2 种选择,剩下两人顺序 2 种选择,共 $4\cdot2\cdot2$ 种。

        情况 2:$A$ 女士与 $A$ 先生之间间隔一人。这个人有 4 种选择,其配偶必须坐在其他三个座位的中间位置,$C$夫妇有 2 种选择,且$A$ 女士与 $A$ 先生可互换位置,共 $4\cdot2\cdot2$ 种。

        因此概率为
        \[
        \frac{4\cdot2\cdot2+4\cdot2\cdot2}{5!} = \frac{4}{15}
        \]
    \end{solution}

    \question 伯克利市每天的天气要么是雨天,要么是雾天,其中天气与前一天相同的概率为 \(\dfrac{3}{4}\)。如果我们只知道今天是雨天,求 7 天后仍是雨天的概率。
    \begin{solution}
        7 天后仍是雨天的概率为
        \begin{align*}
        &P(\text{天气完全没变}) + P(\text{天气变恰好2次})+ P(\text{天气变恰好4次})+ P(\text{天气变恰好6次}) \\
        &= \comb 70 \left(\frac14\right)^0\left(\frac34\right)^7
        + \comb 72 \left(\frac14\right)^2\left(\frac34\right)^5
        + \comb 74 \left(\frac14\right)^4\left(\frac34\right)^3 
        + \comb 76 \left(\frac14\right)^6\left(\frac34\right)^1 \\
        &= \frac{129}{256}
        \end{align*}
    \end{solution}

\question 袋中有红球 5 个、白球 3 个、黑球 4 个,若每球被选取的机会均等,每次由袋中取一球,取后不放回,取完为止,则黑球最先取完的概率为？

\begin{solution}
\[
{5\over 5+4}+{3\over 3+4}-{5+3\over 5+3+4}= \frac{20}{63}
\]
%https://math.ntnu.edu.tw/~horng/letter/hpm17010.pdf
\end{solution}

    \question 一对实数 $(a,b)$ 在闭单位圆盘内均匀随机选取,即满足
    \[
    a^{2} + b^{2} \leq \frac{1}{4}.
    \]
    求二次函数
    \[
    y = ax^{2} + 2bx - a
    \]
    与曲线
    \[
    y = x^{2}
    \]
    相交的概率。
    \ifprintanswers
    \begin{figure}[H]
        \centering        
        \includegraphics[width=0.6\textwidth]{images/image174.png}
    \end{figure}
    \fi
    \begin{solution}
        曲线相交当且仅当方程
        \[
        (a-1)x^2 + 2bx - a = 0
        \]
        有实根,此时判别式为非负:
        \[
        (2b)^2 - 4(a-1)(-a) \ge 0 \Rightarrow \left(a - \frac{1}{2}\right)^2 + b^2 \ge \left(\frac{1}{2}\right)^2
        \]
        实数 $(a,b)$满足在圆 $a^2+b^2 = \left(\dfrac{1}{2}\right)^2$ 内但在圆$\left(a - \dfrac{1}{2}\right)^2 + b^2 = \left(\dfrac{1}{2}\right)^2$外,以阴影部分表示,面积为
        \[
        \pi\left(\frac{1}{2}\right)^2-2\left(2\cdot \frac{1}{2}\left(\frac{1}{2}\right)^2\left(\frac{\pi}{3}\right)-\frac{\sqrt{3}}{4}\left(\frac{1}{2}\right)^2\right) = \frac{\pi}{12}+\frac{\sqrt{3}}{8}
        \]
        所求概率为
        \[
        \frac{\frac{\pi}{12}+\frac{\sqrt{3}}{8}}{\frac{\pi}{4}}
        = \frac{1}{3} + \frac{\sqrt{3}}{2\pi} \approx 0.609
        \]
    \end{solution}

    \question 一只蚂蚁从立方体的一个顶点出发,每次可以走到与当前位置相邻的三个顶点中的任意一个。求四步后,蚂蚁回到出发点的概率。
        \begin{solution}
        设蚂蚁起点为 $(0,0,0)$,三步后,蚂蚁抵达顶点 $(1,1,1)$的概率为
        \[
        \frac{3!}{3^3} = \frac{2}{9},
        \]
        否则,蚂蚁会停在 $(1,0,0),(0,1,0)$ 或 $(0,0,1)$。

        若三步后位于 $(1,1,1)$,蚂蚁下一步无法回到 $(0,0,0)$；但如果三步后在 $(1,0,0),(0,1,0)$ 或 $(0,0,1)$,则第四步回到 $(0,0,0)$ 的概率均为 $\dfrac{1}{3}$。

        故四步后回到出发点的概率为
        \[
        \frac{7}{9} \cdot \frac{1}{3} = \frac{7}{27}
        \]
    \end{solution}

    \question 一只蚂蚁在正方体 $ABCD-EFGH$ 的顶点 $A$ 处, 每次等概率地爬行到相邻三个顶点中的一个, 那么六次爬行之后回到顶点 $A$ 处的概率为
\begin{solution}
将立方体顶点坐标设为 $(0,0,0)$ 或 $1$,蚂蚁从 $(0,0,0)$ 出发。每步蚂蚁改变一个坐标的值（0 $\leftrightarrow$ 1）。要在六步后回到原点,每个坐标必须被改变偶数次。设三个坐标变化次数分别为 $x,y,z\ge 0$,且 $x+y+z=6$,且 $x,y,z$ 都是偶数。可行的组合为
\[
(0,0,6),\ (0,6,0),\ (6,0,0),\ (0,2,4),\ (0,4,2),\ (2,0,4),\ (2,4,0),\ (4,0,2),\ (4,2,0),\ (2,2,2)。
\]

每个组合对应的步序列数由多项式系数计算:
\[
(2,2,2) \rightarrow \frac{6!}{2!2!2!}=90,\quad
(0,2,4)\ \text{及类似组合} \rightarrow \frac{6!}{0!2!4!}=15。
\]

前六种 $(0,2,4)$ 类型组合共有 6 种,贡献步序列总数 $6\cdot 15=90$。加上 $(2,2,2)$ 的 90,总共有 $180$ 种有效序列。总步数序列数为 $3^6=729$。因此概率为
\[
\frac{180}{729}=\frac{20}{81}.
\]
\textcolor{red}{(chatgpt,待解)}
\end{solution}

    \question 一只蚂蚁从正八面体的一个顶点出发,每一步随机走到相邻顶点。求 10 步后它回到出发点的概率。  
    \begin{solution}
        设 $a_n,b_n,c_n$ 分别为走 $n$ 步后蚂蚁在 (a) 出发点、(b) 与出发点相邻的四个顶点、(c) 对面顶点 的概率。有递推关系
        \[
        a_n=\frac{1}{4}b_{n-1}, \quad c_n=\frac{1}{4}b_{n-1}, \quad b_n=a_{n-1}+\frac{1}{2}b_{n-1}+c_{n-1}
        \]
        消去 $a_n,c_n$ 得
        \[
        b_n=\frac{1}{2}b_{n-1}+\frac{1}{2}b_{n-2}
        \]
        特征方程为 $x^2=\frac{1}{2}x+\frac{1}{2}$,解得 $x=1,-\dfrac{1}{2}$,故
        \[
        b_n=\alpha\cdot 1^n+\beta\left(-\frac{1}{2}\right)^n
        \]
        初值$a_0=1,b_0=c_0=0$,所以 $b_1=1$,可得 $\alpha=\dfrac{2}{3}, \beta=-\dfrac{2}{3}$,即
        \[
        b_n=\frac{2}{3}\left(1-\left(-\frac{1}{2}\right)^n\right)
        \]
        于是
        \[
        a_{10}=\frac{1}{4}b_{9}=\frac{1}{4}\cdot \frac{2}{3}\left(1-\left(-\frac{1}{2}\right)^9\right)
        =\frac{513}{3072}
        \]
    \end{solution}

    \question 一只蜜蜂在三维空间中移动。掷一个公平的六面骰子,骰子的六个面分别标记为 \( \pm x, \pm y, \pm z \)。假设蜜蜂当前所在点为 \( (x,y,z) \)。若骰子显示 \( +x \),则蜜蜂移动到点 \( (x+1, y, z) \)；若骰子显示 \( -x \),则蜜蜂移动到点 \( (x-1, y, z) \)。其他四个面以此类推。已知蜜蜂从点 \( (0,0,0) \) 出发,掷骰子四次。求蜜蜂经过的四条边恰好是某单位立方体上的四条不同边的概率。 
    \begin{solution}
        由乘法原理及想象力,概率为
        \[
        1\cdot \dfrac{4}{6}\cdot \left(\dfrac{1}{6}\cdot \dfrac{3}{6}+\dfrac{2}{6}\cdot \dfrac{2}{6}\right)=\frac{7}{54}
        \]
    \end{solution}

    \question 两支队伍正在进行“三局两胜”的系列赛:最多进行 3 场比赛,率先赢下 2 场的队伍获胜。第一场比赛在 A 队的主场进行,后两场在 B 队的主场进行。A 队在主场赢的概率是 $\dfrac{2}{3}$,在客场赢的概率是 $p$。各场比赛的结果是独立的。若 A 队赢得整个系列赛的概率是 $\dfrac{1}{2}$,求 $p$ 的值。
    \begin{solution}
        已知 A 队主场胜率为 $\dfrac23$,客场胜率为 $p$。三局两胜赛制下,A 队赢下系列赛的情形只有  
        \[
        \text{AA},\; \text{ABA},\; \text{BAA}
        \]
        于是有
        \[
        \frac23 p+\frac23(1-p)p+\frac13p^{2}=\frac12
        \]
        解得  
        \[
        p=\dfrac12\bigl(4-\sqrt{10}\bigr)\le1
        \]
    \end{solution}
    
    \question  一个不均匀的骰子,掷出 $1,2,3,4,5,6$ 点的概率依次成等差数列. 独立地先后掷该骰子两次,所得的点数分别记为 $a,b$. 若事件 “$a+b=7$” 发生的概率为 $\dfrac{1}{7}$,求事件 “$a=b$” 发生的概率。 
    \begin{solution}     
        设掷出$1,2,\dots,6$点的概率分别为$p_1,p_2,\dots,p_6$。由于$p_1,p_2,\dots,p_6$成等差数列,且$p_1+p_2+\dots+p_6=1$,故
        \[
        p_1+p_6=p_2+p_5=p_3+p_4=\frac{1}{3}
        \]
        事件“$a+b=7$”发生的概率为
        \[
        P_1=p_1p_6+p_2p_5+\dots+p_6p_1
        \]
        事件“$a=b$”发生的概率为
        \[
        P_2=p_1^2+p_2^2+\dots+p_6^2
        \]
        于是
        \[
        P_1+P_2=(p_1+p_6)^2+(p_2+p_5)^2+(p_3+p_4)^2=3\cdot\left(\frac{1}{3}\right)^2=\frac{1}{3}
        \]
        由于$P_1=\dfrac{1}{7}$,所以
        \[
        P_2=\frac{1}{3}-\frac{1}{7}=\frac{4}{21}
        \]
    \end{solution}
    \begin{solution}
        设掷出点数$1,2,3,4,5,6$的概率依次为$p-5d,p-3d,p-d,p+d,p+3d,p+5d$,据题意有
        \begin{align}
        (p-5d)+(p-3d)+(p-d)+(p+d)+(p+3d)+(p+5d)&=1 \tag{1}\\
        2[(p-5d)(p+5d)+(p-3d)(p+3d)+(p-d)(p+d)]&=\dfrac{1}{7} \tag{2} 
        \end{align} 
        由$(1)$得$p=\dfrac{1}{6}$,由$(2)$化简得
        \[
        6p^2 - 70 d^2 = \frac{1}{7} \Rightarrow 70d^2=\frac{1}{42}
        \] 
        于是
        \begin{align*}
            P(a=b)&=(p-5d)^2+(p-3d)^2+(p-d)^2+(p+d)^2+(p+3d)^2+(p+5d)^2 \\
            &=6p^2+70d^2 =\frac{1}{6}+\frac{1}{42}=\frac{4}{21}
        \end{align*}
    \end{solution}

    \question 设甲、乙、丙三位射手之命中率依次为 $p,q,r$，其中 $p\ge q\ge r$。今三人同打一靶且互不影响，各发一弹时，此靶不中弹之概率为 $\dfrac{1}{4}$，恰中一弹之概率为 $\dfrac{11}{24}$，恰中二弹之概率为 $\dfrac{1}{4}$，求序组 $(p,q,r)$。
    \begin{solution} 
        据题意,
        \[
        \begin{cases}
        pqr=1-\left(\dfrac{1}{4}+\dfrac{11}{24}+\dfrac{1}{4}\right)=\dfrac{1}{24}, \\[2pt]
        p(1-q)(1-r)+q(1-r)(1-p)+r(1-p)(1-q)=\dfrac{11}{24}, \\[2pt]
        pq(1-r)+qr(1-p)+rp(1-q)=\dfrac{1}{4},
        \end{cases}
        \]
        其中 $p\ge q\ge r$,解得
        \[
        pqr=\frac{1}{24},\quad pq+qr+rp=\frac{3}{8},\quad p+q+r=\frac{13}{12},
        \]
        因此 $p,q,r$ 为三次方程
        \[
        t^3-\frac{13}{12}t^2+\frac{3}{8}t-\frac{1}{24}=0 \Rightarrow (2t-1)(3t-1)(4t-1)=0
        \]
        的三个根,解得
        \[
        (p,q,r)=\left(\frac{1}{2},\frac{1}{3},\frac{1}{4}\right)
        \]
    \end{solution}

    \question 在区间 $[0,10]$ 上等概率地随机选取一个数 $y$。以点 $(0,y)$ 为斜边一端点,斜边另一端点位于非负 $x$ 轴上,并与原点 $(0,0)$ 共同构成一个斜边长度为 10 的直角三角形。求该三角形面积大于 15 的概率。
    \begin{solution}
        由斜边长度 10 可得面积条件为
        \[
        \frac{1}{2} xy = \frac{1}{2} y\sqrt{100-y^2} > 15 
        \]
        解得
        \[
        y < \sqrt{10} \quad \text{或} \quad y> \sqrt{90}
        \]
        所求概率为
        \[
        \frac{\sqrt{90}-\sqrt{10}}{10} = \frac{\sqrt{10}}{5}
        \]
    \end{solution}

    \question 设甲袋中有 5 颗白球、2 颗黑球,乙袋中有 3 颗白球。先自甲袋中任取 4 颗球放入乙袋,再从乙袋中任取 5 颗球放入甲袋,完成一次操作称为一局。已知每颗球被抽到的机会均等,若第一局结束时甲袋中有黑球,求过程中过甲袋取得 2 颗白球、2 颗黑球放入乙袋的概率。
    \begin{solution}
        情况一:甲取 4 白球入乙,再从乙任取 5 球入甲。  
        \[
        P_1 = \frac{\comb{5}{4}}{\comb{7}{4}} = \frac{1}{7}
        \]
        情况二:甲取 3 白 1 黑入乙,再从乙任取 5 球入甲。  
        \[
        P_2 = \frac{\comb{5}{3}\comb{2}{1}}{\comb{7}{4}} = \frac{4}{7}
        \]
        情况三:甲取 2 白 2 黑入乙,再从乙任取 4 白 1 黑球入甲。  
        \[
        P_3 = \frac{\comb{5}{2}\comb{2}{2}}{\comb{7}{4}} \cdot \frac{\comb{5}{4}\comb{2}{1}}{\comb{7}{5}} = \frac{20}{147}
        \]
        情况四:甲取 2 白 2 黑入乙,再从乙任取 3 白 2 黑球入甲。  
        \[
        P_4 = \frac{\comb{5}{2}\comb{2}{2}}{\comb{7}{4}} \cdot \frac{\comb{5}{3}\comb{2}{2}}{\comb{7}{5}} = \frac{20}{147}
        \]
        已知第一局结束时甲袋中有黑球,因此条件概率为
        \[
        \frac{P_3+P_4}{P_1+P_2+P_3+P_4} = \frac{8}{29}
        \]
    \end{solution}

    \question 三个人玩一个游戏,轮流掷一枚公平的六面骰子。若掷出 5 或 6 则该玩家立即获胜,游戏结束；若掷出 1 或 2 则该玩家被淘汰。游戏继续进行,直到有人获胜,或者只剩下一人,则该人获胜。问第一个掷骰子的人获胜的概率是多少？
    \begin{solution}
        设第一个掷骰子的人为 $A$,任意一次掷出 1,2,5,6 的概率是 $\dfrac{2}{3}$,故$A$ 第一个掷出 1,2,5,6 的概率为
        \[
        \frac{2}{3}\left(1+\frac{1}{3}+\frac{1}{3^2}+\cdots\right)=\frac{\frac{2}{3}}{1-\frac{1}{9}}=\frac{9}{13}
        \]
        而此时$A$有$\dfrac{1}{2}$的概率掷出 5 或 6 从而立即获胜。

        若 $A$ 不是第一个掷出 1,2,5,6 的人(概率为 $\dfrac{4}{13}$),那么第一个掷出 1,2,5,6 的人有$\dfrac{1}{2}$的概率掷出 1 或 2 被淘汰。此时只剩下两人,$A$ 与另一人各有 $\dfrac{1}{2}$ 的获胜机会。

        综上,第一个掷骰子的人获胜的概率为
        \[
        \frac{9}{13}\cdot \frac{1}{2}+\frac{4}{13}\cdot \frac{1}{2}\cdot \frac{1}{2}=\frac{11}{26}
        \]
    \end{solution}

    \question 反复抛掷一枚公平硬币,直到第一次出现序列 $HTH$ 为止。求在此过程中,序列 $THTH$ 从未出现过的概率。(例如,事件包含 $HHTH$,但不包含 $TTHTH$。)
    \begin{solution}
        设事件 $E$ 表示在序列 $THTH$ 出现之前先出现 $HTH$。若当前已出现的前缀为 $s$,则记 $P(E\mid s)$ 为在以 $s$ 开始的条件下事件 $E$ 发生的概率。定义
        \[
        x = P(E\mid H), \quad y = P(E\mid T)
        \]
        根据一步分析法,有
        \[
        x = \frac{1}{2} P(E\mid HH) + \frac{1}{4} P(E\mid HTT) + \frac{1}{4} P(E\mid HTH) = \frac{1}{2}x + \frac{1}{4}y + \frac{1}{4},
        \]
        \[
        y = \frac{1}{2} P(E\mid TT) + \frac{1}{4} P(E\mid THH) + \frac{1}{8} P(E\mid THTT) + \frac{1}{8} P(E\mid THTH) = \frac{1}{2}y + \frac{1}{4}x + \frac{1}{8}y
        \]
        解得
        \[
        x = \frac{3}{4}, \quad y = \frac{1}{2}
        \]
        因此所求概率为
        \[
        \frac{1}{2}x + \frac{1}{2}y = \frac{5}{8}
        \]
    \end{solution}

    \question 反复抛掷一枚公平硬币,直到出现连续四个正面 $HHHH$ 或连续六个反面 $TTTTTT$ 为止。求事件“$HHHH$ 先于 $TTTTTT$ 出现”的概率。
    \begin{solution}
        设事件 $E$ 表示在序列 $TTTTTT$ 出现之前先出现 $HHHH$。对任意当前末尾已出现的序列 $s$,记 $P(E\mid s)$ 为在条件 $s$ 下事件 $E$ 发生的概率。定义
        \[
        x = P(E\mid H),\quad y = P(E\mid T).
        \]
        根据后续可能的状态转移,可得到方程组（按下一次或若干次投掷展开）：
        \[
        x = \frac{7}{8}y + \frac{1}{8}, \quad
        y = \frac{31}{32}x
        \]
        因此
        \[
        x = \frac{32}{39}, \quad y = \frac{31}{39}
        \]
        由于初始时无既定前缀（在第一次投掷前）,所求概率为 
        \[
        \frac12(x+y)= \frac{21}{26}
        \]
    \end{solution}

    \question 四位玩家各自掷一个标准六面骰子,点数最大者获胜。如果出现平手,则平手者继续掷骰,直到一人胜出。$H$ 是其中之一。已知他最终获胜,求他第一轮掷出点数为 $5$ 的条件概率。
    \begin{solution}
        设事件 $A$ 为“$H$ 首掷为 $5$”,事件 $B$ 为“$H$ 获胜”,根据贝叶斯公式,
        \[
        P(A \mid B) = \frac{P(B \mid A) \cdot P(A)}{P(B)}.
        \]
        已知 $P(A) = \dfrac{1}{6}, P(B) = \dfrac{1}{4}$, 而$P(B \mid A)$ 如下分情况:
        \begin{itemize}
            \item 其他人都掷出 $\le 4$,概率为 $\left( \dfrac{2}{3} \right)^3 = \dfrac{8}{27}$
            \item 恰一人也掷 $5$,概率 $3 \cdot \dfrac{1}{6} \cdot \left( \dfrac{2}{3} \right)^2 \cdot \dfrac{1}{2} = \dfrac{1}{9}$
            \item 恰两人也掷 $5$,概率 $3 \cdot \left( \dfrac{1}{6} \right)^2 \cdot \dfrac{2}{3} \cdot \dfrac{1}{3} = \dfrac{1}{54}$
            \item 三人都掷 $5$,概率 $\left( \dfrac{1}{6} \right)^3 \cdot \dfrac{1}{4} = \dfrac{1}{864}$
        \end{itemize}
        故
        \[
        P(B \mid A) = \frac{8}{27} + \frac{1}{9} + \frac{1}{54} + \frac{1}{864} = \frac{41}{96} ,\quad
        P(A \mid B) = \frac{\frac{41}{96} \cdot \frac{1}{6}}{\frac{1}{4}} = \frac{41}{144}
        \]
    \end{solution}

    \question 某种生物不会与其他个体产生互动,但能够自行繁殖。若单独放置一小时,它会以相等的概率变成 $0,1,2,3$ 个个体,分别对应死亡或不同的繁殖结果。其后新产生的个体在每个后续小时中也以同样的方式独立行动。若初始时只有一个个体,设 $p$ 为这一族群能够持续存在（即永不灭绝）的概率,求$p$。
    \begin{solution}
        设 $q$ 为族群最终灭绝的概率。若当前有 $n$ 个个体,则最终灭绝的概率为 $q^n$。题意给出
        \[
        q = \frac14 \left(1 + q + q^2 + q^3\right)
        \]
        化简得
        \[
        (q-1)(q^2 + 2q - 1) = 0
        \]
        在 $[0,1]$ 区间内唯一满足要求的解为
        \[
        q = \sqrt{2} - 1
        \]
        因此
        \[
        p = 1 - q = 2 - \sqrt{2} \approx 0.586
        \]
    \end{solution}

    \question 设 $P(n)$ 表示在 $n$ 次抛掷一枚公平硬币时,出现至少连续三个正面的概率。求满足 $P(n)\ge \dfrac12$ 的最小整数 $n$。
    \begin{solution}
        设事件 $E_i$ 表示首次出现连续三个正面恰好发生在第 $i$ 次抛掷时,则
        \[
        P(n) = \sum_{i=1}^{n-2} \Pr(E_i),
        \]
        其中
        \[
        \Pr(E_i) = 
        \begin{cases}
        \dfrac{1}{8}, & i=1,\\[2mm]
        \dfrac{1}{16}(1-P(i-2)), & i\ge 2.
        \end{cases}
        \]
        当 $i\ge 2$ 时,$E_i$ 发生的条件是第 $i$ 次抛出 H 前必须是 T,并且之前没有出现过三个连续 H。注意 $P(i-2)=0$ 当 $i<5$,逐步计算:
        \[
        P(3)=\frac{1}{8},\quad P(4)=\frac{3}{16},\quad P(5)=\frac{4}{16},\quad P(6)=\frac{5}{16},
        \]
        \[
        P(7)=\frac{6}{16}-\frac{1}{16}\cdot\frac{1}{8},\quad P(8)=\frac{7}{16}-\frac{1}{16}\left(\frac{1}{8}+\frac{3}{16}\right),
        \]
        \[
        P(9)=\frac{8}{16}-\frac{1}{16}\left(\frac{1}{8}+\frac{3}{16}+\frac{4}{16}\right),\quad 
        P(10)=\frac{9}{16}-\frac{1}{16}\left(\frac{1}{8}+\frac{3}{16}+\frac{4}{16}+\frac{5}{16}\right)=\frac{65}{128} > \frac{1}{2}
        \]
        因此最小的 $n$ 满足 $P(n)\ge \dfrac{1}{2}$ 为 10。
    \end{solution}

    \question 设甲、乙两箱中,甲箱内有 1 白球 1 红球,乙箱内有 1 白球 2 红球。现在每次先从甲箱中随机取一球,放入乙箱内,再从乙箱随机取一球放入甲箱,这样视为一局。试求  
    \begin{parts}
    \part 在第二局结束后,有 2 红球在甲箱内的概率。
    \begin{solution}
        设三状态如下表所示：
        \[
        \begin{array}{c|c|c}
        \hline
        \text{状态} & \text{甲箱} & \text{乙箱} \\ \hline
        S_1 & 1\text{白},\ 1\text{红} & 1\text{白},\ 2\text{红} \\ \hline
        S_2 & 2\text{红} & 2\text{白},\ 1\text{红} \\ \hline
        S_3 & 2\text{白} & 3\text{红} \\ \hline
        \end{array}
        \]
        因此转移矩阵为
        \begin{align*}
        A&=\begin{bmatrix}
        P(S_1\to S_1) & P(S_2\to S_1) & P(S_3\to S_1) \\
        P(S_1\to S_2) & P(S_2\to S_2) & P(S_3\to S_2) \\
        P(S_1\to S_3) & P(S_2\to S_3) & P(S_3\to S_3)
        \end{bmatrix} \\
        &=\begin{bmatrix}
        \frac{1}{2}\cdot \frac{1}{2}+ \frac{1}{2}\cdot \frac{3}{4} & 1\cdot\frac{1}{2} & 1\cdot\frac{3}{4} \\
        \frac{1}{2}\cdot\frac{1}{2} & 1\cdot\frac{1}{2} & 0 \\
        \frac{1}{2}\cdot\frac{1}{4} & 0 & 1\cdot\frac{1}{4}
        \end{bmatrix}
        =\begin{bmatrix}
        \frac{5}{8} & \frac{1}{2} & \frac{3}{4} \\
        \frac{1}{4} & \frac{1}{2} & 0 \\
        \frac{1}{8} & 0 & \frac{1}{4}
        \end{bmatrix}
        \end{align*}
        于是
        \[
        A^2
        \begin{bmatrix} 1\\ 0 \\0 \end{bmatrix}
        =
        \left[\begin{matrix}
        \frac{39}{64} & \frac{9}{16} & \frac{21}{32} \\
        \frac{9}{32} & \frac{3}{8} & \frac{3}{16} \\
        \frac{7}{64} & \frac{1}{16} & \frac{5}{32}
        \end{matrix}\right]
        \begin{bmatrix} 1\\ 0 \\0 \end{bmatrix}
        =
        \begin{bmatrix}
        \frac{39}{64} \\
        \frac{9}{32} \\
        \frac{7}{64}
        \end{bmatrix}
        \]
        即第二局结束后,有 2 红球在甲箱内的概率为$\dfrac{9}{32}$。
    \end{solution}
    \part 经过长时间交换后,有 2 红球在甲箱内的概率。
    \begin{solution}
        对角化转移矩阵得
        \[
        A=
        \left[\begin{matrix}
        -2 & 1 & 6 \\
        1 & -2 & 3 \\
        1 & 1 & 1
        \end{matrix}\right]
        \left[\begin{matrix}
        0 & 0 & 0 \\
        0 & \frac{3}{8} & 0 \\
        0 & 0 & 1
        \end{matrix}\right]
        \left[\begin{matrix}
        -\frac{1}{6} & \frac{1}{6} & \frac{1}{2} \\
        \frac{1}{15} & -\frac{4}{15} & \frac{2}{5} \\
        \frac{1}{10} & \frac{1}{10} & \frac{1}{10}
        \end{matrix}\right]
        \]
        于是
        \[
        A^\infty
        =
        \left[\begin{matrix}
        -2 & 1 & 6 \\
        1 & -2 & 3 \\
        1 & 1 & 1
        \end{matrix}\right]
        \left[\begin{matrix}
        0 & 0 & 0 \\
        0 & 0 & 0 \\
        0 & 0 & 1
        \end{matrix}\right]
        \left[\begin{matrix}
        -\frac{1}{6} & \frac{1}{6} & \frac{1}{2} \\
        \frac{1}{15} & -\frac{4}{15} & \frac{2}{5} \\
        \frac{1}{10} & \frac{1}{10} & \frac{1}{10}
        \end{matrix}\right]
        =
        \left[\begin{matrix}
        \frac{3}{5} & \frac{3}{5} & \frac{3}{5} \\
        \frac{3}{10} & \frac{3}{10} & \frac{3}{10} \\
        \frac{1}{10} & \frac{1}{10} & \frac{1}{10}
        \end{matrix}\right]
        \]
        且
        \[
        A^\infty
        \begin{bmatrix} 1\\ 0\\ 0\end{bmatrix}
        =\begin{bmatrix}
        \frac{3}{5} \\
        \frac{3}{10} \\
        \frac{1}{10}
        \end{bmatrix}
        \]
        经过长时间交换后,有 2 红球在甲箱内的概率为
        \[
        P(S_2)=\frac{3}{10}
        \]
    \end{solution}
    \end{parts}

    \question 小明给阿花 7 颗巧克力当礼物,然后跟阿花说:「你每天至少要吃一颗巧克力,表示你有想念我,当然你要一次吃好多颗也可以喔。」

    设随机变量 $X$ 代表阿花吃完巧克力的天数,求期望值 $E(X)$。
    \begin{solution}
        考虑将 7 颗巧克力在 $X$ 天内吃完,且每天至少吃一颗。这等价于将 7 个相同的球放入 $X$ 个不同的盒子,每个盒子至少一个球。由间隔法,
        \[
        P(X=k)=\frac{\comb{6}{k-1}}{2^6}
        \]
        故
        \begin{align*}
        E(X)&=\sum_{k=1}^{7}k\cdot P(X=k) \\
        &=\frac{1}{64}\left(1\cdot \comb{6}{0}+2\cdot \comb{6}{1}+3\cdot \comb{6}{2}+4\cdot \comb{6}{3}+5\cdot \comb{6}{4}+6\cdot \comb{6}{5}+7\cdot \comb{6}{6}\right)
        =4
        \end{align*}
    \end{solution}

    \question 在下图所示的大正方形中随机选择一个点。图中嵌套的小正方形无限向内延续。求该随机点所处正方形数量的期望值。例如,图中所示的样例点位于其中的 2 个正方形内。
    \begin{figure}[H]
        \centering
        \includegraphics[width=0.5\textwidth]{images/image154.png}
    \end{figure} 
    \begin{solution}
        由于四个三角形面积和皆占正方形的一半,所求为
        \[
        1 \cdot \frac{1}{2} + 2 \cdot \frac{1}{2^2} + 3 \cdot \frac{1}{2^3} + \cdots
        \]
        即
        \[
        \left(\frac{1}{2} + \frac{1}{4} + \cdots \right)\left(1 + \frac{1}{2} + \frac{1}{4} + \cdots \right)
        = 1 \cdot 2 = 2
        \]
    \end{solution}

    \question 一个人出生在一周七天中的任意一天且概率相同。一大群人依次报出自己的出生星期,直到七个星期几都至少出现一次为止。试求第一次凑齐全部七个星期几时,对应的那个人的期望序号。
    \begin{solution}
        这是经典的“优惠券收集（Coupon Collector）”问题。设当前已经出现了 $i-1$ 个不同的星期,则下一位出现一个新星期的概率为
        \[
        p=\frac{7-(i-1)}{7},
        \]
        因此获得一个新星期所需的期望人数为
        \[
        \frac{1}{p}=\frac{7}{7-(i-1)}.
        \]
        于是凑齐全部七天所需的期望总人数为
        \[
        1+\frac{7}{6}+\frac{7}{5}+\frac{7}{4}+\frac{7}{3}+\frac{7}{2}+\frac{7}{1}=\frac{363}{20}
        \]
        其中,第 $i$ 项表示在已有 $i-1$ 种不同星期之后,等到第 $i$ 个新星期出现的期望等待人数。
    \end{solution}

    \question 假设投掷某金币一枚出现正面与反面的概率皆为 $\dfrac{1}{2}$,试问平均要连续掷几次金币,才会出现连续三次正面？
    \begin{solution}
        设 $E(n)$ 为出现连续 $n$ 次正面的掷币次数。由一步分析法,
        \[
        E(n) = p(E(n-1) + 1) + (1-p)(E(n-1) + 1 + E(n))
        \]
        代入 $p = \dfrac{1}{2}$,
        \[
        E(1) = 1 + \frac{1}{2} E(1), \quad E(2) = E(1) + 1 + \frac{1}{2} E(2), \quad E(3) = E(2) + 1 + \frac{1}{2} E(3)
        \]
        得
        \[
        E(1) = 2, \quad E(2) = 6, \quad E(3) = 14
        \]
    \end{solution}

    \question 已知一个非公正硬币掷出正面机率为 $\dfrac{1}{3}$,反面机率为 $\dfrac{2}{3}$,今连续掷此硬币,记录每次掷出的结果,每次结果互不影响,令随机变量 $X$表示第一次看到正面、反面、正面依序出现所需的投掷次数,求 $X$的期望值。
    \begin{solution}
        假设期望值为 $E(X)$。若
        \begin{itemize}
        \item 出现反面,期望值变为$E(X)+1$
        \item 出现正面、正面,期望值变为$E(X)+1$
        \item 出现正面、反面、反面,期望值变为$E(X)+3$
        \end{itemize}
        由一步分析法,
        \[
        E(X)=\frac{2}{3}(E(X)+1)+\left(\frac{1}{3}\right)^2(E(X)+1)+\left(\frac{1}{3}\right)\left(\frac{2}{3}\right)^2(E(X)+3)
        \]
        解得
        \[
        E(X)=\frac{33}{2}
        \]
    \end{solution}

    \question 袋子里有 $n+3$ 颗球，其中红球有 3 颗，黑球有 $n$ 颗。每次取一球，取后不放回，直到取到两颗红球为止。设随机变量 $X$ 表示取到两颗红球停止时的次数，求：
    \begin{parts}
    \part 概率 $P(X=k)$，其中 $k=2,3,4,\dots,n+2$。
    \begin{solution}
        三颗红球、$n$ 颗黑球任意排列，共有
        \[
        \comb{n+3}{3}=\frac{(n+3)(n+2)(n+1)}{6}
        \]
        种排列方法。第 $k$ 次取到第 2 颗红球：第 $k$ 次为红球，前 $k-1$ 次中有 1 红球和 $k-2$ 黑球，共 $k-1$ 种排列；后 $n+3-k$ 颗球中有 1 红球和 $n+2-k$ 黑球，共 $n+3-k$ 种排列。因此
        \[
        P(X=k)=\frac{(k-1)(n+3-k)}{(n+3)(n+2)(n+1)/6}
        = \frac{6(k-1)(n-k+3)}{(n+3)(n+2)(n+1)}, \; k=2,3,4,\dots,n+2
        \]
    \end{solution}
    \part 随机变量 $X$ 的期望值 $E(X)$。
    \begin{solution}
        期望值为
        \begin{align*}
        E(X)
        &=\sum_{k=2}^{n+2} k P(X=k) \\
        &= \frac{6}{(n+3)(n+2)(n+1)} \sum_{k=2}^{n+2} k(k-1)(n-k+3)\\
        &= \frac{6}{(n+3)(n+2)(n+1)} \sum_{k=2}^{n+2} \left[-k^3 + (n+4)k^2 - (n+3)k \right] \\
        &= \frac{6}{(n+3)(n+2) (n+1)}\left[(n+3)(n+4) \left(\frac{1}{6}(n+4)(2n+7)-\frac{1}{4}(n+3)(n+6) \right) \right] \\
        &=\frac{n+4}{2}
        \end{align*}
    \end{solution}
    \end{parts}

    \question 在一个由 0 和 1 组成的序列中,"段"指的是由连续相同数字组成的一串,包括长度为 1 的段。例如序列 00100011 中共有四段。对于一个包含 15 个 0 和 9 个 1 的随机排列,求其段数的期望值。
    \begin{solution}
        令随机变量 $X_i$ 表示第 $i$ 个位置是否为某一段的起点。若是,则 $X_i=1$;否则 $X_i=0$。

        第一个位置一定是一段的起点,因此 $E(X_1)=1$。当 $i>1$ 时,第 $i$ 个位置成为段起点当且仅当前一位与当前位不同。其概率为
        \[
        \frac{15}{24}\cdot \frac{9}{23} + \frac{9}{24}\cdot \frac{15}{23}= \frac{45}{92}
        \]
        因此段数的期望值为
        \[
        E(X_1)+E(X_2)+\cdots+E(X_{24})= 1 + 23 \cdot \frac{45}{92}= \frac{49}{4}
        \]
    \end{solution}

    \question $A$ 从帕斯卡三角形的顶端开始。每次移动,她都会向下一层走,并且以相等概率选择向左或向右。完成 6 次移动后,$A$ 经过的数字（包括起点和终点）的期望和是多少？
    
    例如,在下图所示的路径中,$\;A$ 访问的数值之和为 $1 + 1 + 1 + 3 + 6 + 10 + 20 = 42$。 
    \begin{figure}[H]
        \centering
        \includegraphics[width=0.5\textwidth]{images/image8.png}
    \end{figure} 
    \begin{solution}
        设走到第 $n$ 层向右走了 $k$ 步,则 $k\sim\mathrm{Binom}\left(n,\dfrac12\right)$,该层数值为 $\dbinom{n}{k}$,则第 $n$ 层的期望值
        \[
        \mathbb E_n
        =\sum_{k=0}^{n}\binom{n}{k}^2\left(\frac12\right)^n
        \]
        由性质
        \[
        \sum_{k=0}^{n}\binom{n}{k}^2=\binom{2n}{n}
        \]
        总期望值为
        \[
        \sum_{n=0}^{6}\left(\frac12\right)^n\binom{2n}{n}
        =1 + 1 + \frac{3}{2} + \frac{5}{2} + \frac{35}{8} + \frac{63}{8} + \frac{231}{16}
        =\frac{523}{16}
        \]
    \end{solution}

\end{questions}
\pagebreak

\begin{center}
  {\fontsize{30pt}{26pt}\selectfont
    \hypertarget{统计}{统计} \label{统计}
  }
\end{center}
\separator
\vspace{1pt}

\begin{questions}
    \question 连续投掷一枚公平硬币,直到首次出现连续两个反面为止。设随机变量 $X$ 为所需的投掷次数,求 $X$ 的期望与方差。
    \begin{solution}
        令 $p(n)$ 表示投掷 $n$ 次铜板才出现连续两次反面的概率。若第一次出现正面,后 $n-1$ 次才出现连续两次反面。若第一次出现反面,第二次出现正面,后 $n-2$ 次才出现连续两次反面。有递推式:
        \[
        p(n) = \frac{1}{2} p(n-1) + \frac{1}{4} p(n-2),\quad n\ge 3,\quad p(1)=0,\quad p(2)=\frac{1}{4}
        \]
        于是
        \begin{align*}
        E(X) &= \sum_{k=1}^\infty k p(k) = 1\cdot 0 + 2\cdot \frac{1}{4} + \sum_{k=3}^\infty k\, p(k) \\
        &= \frac{1}{2} + \sum_{k=3}^\infty \left(\frac{1}{2} k p(k-1) + \frac{1}{4} k p(k-2)\right) \\
        &= \frac{1}{2} + \frac{1}{2} (1+E(X)) + \frac{1}{4} (2+E(X)) \Rightarrow E(X) = 6
        \end{align*}
        且
        \begin{align*}
        E(X^2) &= \sum_{k=1}^\infty k^2 p(k) 
        = 1^2 \cdot 0 + 2^2 \cdot \frac{1}{4} + \sum_{k=3}^\infty k^2 p(k) \\
        &= 1 + \sum_{k=3}^\infty \left(\frac{1}{2} k^2 p(k-1) + \frac{1}{4} k^2 p(k-2)\right) \\
        &= 1 + \frac{1}{2} (E(X^2) + 2 E(X) + 1) + \frac{1}{4} (E(X^2) + 4 E(X) + 4) \Rightarrow E(X^2) = 58
        \end{align*}
        故方差为
        \[
        Var(X) = E(X^2) - (E(X))^2 = 22
        \]
    \end{solution}

    \question 一有 $n$ 项的等差数列 $\{a_n\}$,公差为 $d=\dfrac{\sqrt{13}}{2}$,此数列的方差为 $260$,求 $n$。
    \begin{solution}
        设该等差数列为 $a,a+d,\dots,a+(n-1)d$。平均数为
        \[
        \overline{X}=a+\frac{n-1}{2}d
        \]
        且
        \[
        E(X^2)=\frac{1}{n}\sum_{k=0}^{n-1}(a+kd)^2
        =a^2+ad(n-1)+d^2\frac{(n-1)(2n-1)}{6}
        \]
        因此方差为
        \[
        \operatorname{Var}(X)=E(X^2)-\overline{X}^2
        =d^2\left(\frac{(n-1)(2n-1)}{6}-\frac{(n-1)^2}{4}\right)
        =\frac{d^2(n^2-1)}{12}
        \]
        代入 $d^2=\dfrac{13}{4}$ 且 $\operatorname{Var}(X)=260$得
        \[
        \frac{13}{4}\cdot\frac{n^2-1}{12}=260 \Rightarrow n=31
        \]
    \end{solution}

    \question 重复操作一个成功概率为 $p$ 的伯努利试验,且 $$S_n = \sum_{k=1}^{n} k^2 (1-p)^k,$$若 $\displaystyle\lim_{n \to \infty} S_n = 180$,试求第三次才出现第一次成功的概率。
    \begin{solution}
        令 $X \sim \text{Geometric}(p)$表示第一次成功出现所需的试验次数,则有
        \[
        E(X^2) = \text{Var}(X) + (E(X))^2 = \dfrac{1-p}{p^2} + \left(\dfrac{1}{p}\right)^2 = \dfrac{2 - p}{p^2}
        \]
        且据题意有
        \[
        E(X^2) = \dfrac{p}{1 - p} \cdot S_\infty
        \Rightarrow \dfrac{2 - p}{p^2} = \dfrac{p}{1 - p} \cdot 180
        \]
        解得
        \[
        180p^3 - p^2 + 3p - 2 =(5p - 1)(36p^2 + 7p + 2) = 0
        \Rightarrow p = \dfrac{1}{5} \quad (\text{舍去虚根})
        \]
        因此
        \[
        P(X = 3) = (1 - p)^2p  = \dfrac{16}{125}
        \]
    \end{solution}

    \question 某实验测得20组样本点 $(x_1,y_1), (x_2,y_2), \dots, (x_{20},y_{20})$,已知
    \[
    \sum_{i=1}^{20} x_i = 400, \quad \sum_{i=1}^{20} y_i = 900.
    \]
    利用最小平方法求得 $y$ 对 $x$ 的回归直线方程为 $y = a x + b$。若
    \[
    \sum_{i=1}^{20} (y_i - a x_i - b)^2 = 0,
    \]
    且 $(x_1,y_1) = (30,40)$,设
    \[
    x' = 2x - 4, \quad y' = -3y + 5, \quad i=1,2,\dots,20,
    \]
    求数据 $(x', y')$ 的回归直线方程。
    \begin{solution}
        有
        \[
        E(X) = \frac{400}{20} = 20, \quad E(Y) = \frac{900}{20} = 45,
        \]
        \[
        E(X') = E(2X - 4) = 2 \cdot 20 - 4 = 36, \quad E(Y') = E(-3Y + 5) = -3 \cdot 45 + 5 = -130.
        \]
        原回归直线经过 $(30,40)$ 与 $(20,45)$,斜率与截距为
        \[
        y = -\frac{1}{2} x + 55
        \]
        故
        \[
        m = \frac{\mathrm{Cov}(2X - 4, -3Y + 5)}{\mathrm{Var}(2X - 4)} = \frac{-6\,\mathrm{Cov}(X,Y)}{4\,\mathrm{Var}(X)} = -\frac{6}{4} \cdot \left(-\frac{1}{2}\right) = \frac{3}{4}.
        \]
        回归直线过点 $(E(X'), E(Y')) = (36, -130)$,故方程为
        \[
        y' - (-130) = \frac{3}{4}(x' - 36) \Rightarrow y' = \frac{3}{4} x' - 157.
        \]
    \end{solution}
    
    \question 已知5组数据如下,
    \[
    \begin{tabular}{c|ccccc}
    \hline
    $X$ & 1 & 2 & 5 & 3 & 4 \\
    \hline
    $Y$ & 3 & $a$ & 5 & $b$ & 6 \\
    \hline
    \end{tabular}
    \]
    且 $Y$ 对 $X$ 的回归直线为 
    \[
    Y = \frac{3}{5}X + \frac{11}{5},
    \]
    求数对 $(a,b)$。
    \begin{solution}
        有
        \[
        \bar x = \frac{1+2+5+3+4}{5} = 3.
        \]
        由于$\bar x,\bar y$在回归直线上,
        \[
        \bar y = \frac{3}{5} \bar x + \frac{11}{5} =\frac{3 + a + 5 + b + 6}{5} \Rightarrow a + b = 6 \tag{1}
        \]
        由回归直线斜率公式,
        \[
        \frac{3}{5} = \frac{\sum x_i y_i - n \bar x \bar y}{\sum x_i^2 - n \bar x ^2} = \frac{3+2a+25+3b+24-5 \cdot 3 \cdot 4}{55-5 \cdot 3^2} \tag{2}
        \]
        由$(1),(2)$解得
        \[
        (a,b)=(4,2)
        \]
    \end{solution}

    \question 有一笔资料包含 $n$ 个数 $x_1, x_2, x_3, \dots, x_n$,其平均数为 $\overline{x}$,试证
    \[
    \sqrt{\frac{1}{n}\sum_{i=1}^{n}x_{i}^{2}-\overline{x}^{2}}\ge\frac{1}{n}\sum_{i=1}^{n}|x_{i}-\overline{x}|
    \]
    \begin{solution}
        由柯西不等式,
        \[
        \left(\sum_{i=1}^n(x_i-\bar x)^2 \right) \left(\sum_{i=1}^n 1^2\right) \ge \left(\sum_{i=1}^n |x_i-\bar x| \right)^2
        \]
        得
        \[
        \frac{1}{n} \sum_{i=1}^n (x_i - \bar x)^2 \ge \left( \frac{1}{n} \sum_{i=1}^n |x_i - \bar x| \right)^2
        \]
        又
        \begin{align*}
        \frac{1}{n} \sum_{i=1}^n (x_i - \bar{x})^2 &= \frac{1}{n} \left( \sum_{i=1}^n x_i^2 - 2\bar{x} \sum_{i=1}^n x_i + \sum_{i=1}^n \bar{x}^2 \right)\\
        &= \frac{1}{n} \left( \sum_{i=1}^n x_i^2 - 2n\bar{x}^2 + n\bar{x}^2 \right)\\
        &= \frac{1}{n} \sum_{i=1}^n x_i^2 - \bar{x}^2
        \end{align*}
        即得证
        \[
        \sqrt{\frac{1}{n}\sum_{i=1}^{n}x_{i}^{2}-\overline{x}^{2}}\ge\frac{1}{n}\sum_{i=1}^{n}|x_{i}-\overline{x}|
        \]
    \end{solution}
\end{questions}
\pagecolor{PageColor}
\
\vfil
\hfil  {\fontsize{50pt}{36pt}\selectfont{几何}} \hfil
\vfil
\begin{tikzpicture}[remember picture,overlay,every node/.style={inner sep=0pt}]
        \node [shift={(1cm,-1cm)},brown,scale=2,anchor=north west] (CNW)
        at (current page.north west) {\pgfornament[height=1cm,width=1cm]{61}};
        \node [shift={(-1cm,-1cm)},brown,scale=2,anchor=north east] (CNE)
        at (current page.north east) {\pgfornament[height=1cm,width=1cm,symmetry=v]{61}};
        \node [shift={(1cm,1cm)},brown,scale=2,anchor=south west] (CSW)
        at (current page.south west) {\pgfornament[height=1cm,width=1cm,symmetry=h]{61}};
        \node [shift={(-1cm,1cm)},brown,scale=2,anchor=south east] (CSE)
        at (current page.south east) {\pgfornament[height=1cm,width=1cm,symmetry=c]{61}};
        \pgfornamentline[color=brown]{current page.north west}{current page.north east}{2}{87}
        \pgfornamentline{current page.south west}{current page.south east}{2}{87}
        \pgfornamentline{current page.north west}{current page.south west}{3}{87}
        \pgfornamentline{current page.north east}{current page.south east}{3}{87}
        \end{tikzpicture}%
\thispagestyle{empty}
\pagebreak

\begin{center}
  {\fontsize{30pt}{26pt}\selectfont
    \hypertarget{解三角形}{解三角形} \label{解三角形}
  }
\end{center}
\separator
\vspace{1pt}
\nopagecolor
\begin{questions}
    \question 已知钝角三角形的边长为 \(10,\,17,\,x\),求 \(x\) 的所有正整数解之和。  
    \begin{solution}
        钝角三角形需满足某边的平方大于其余两边平方和:
        
        \textbf{情况一:} 若 $x < 17$,最大边是 $17$,成立条件为:
        \[
        \begin{cases}
        17^2 > 10^2 + x^2 \Rightarrow x <\sqrt{189}\approx 13.75\\
        x + 10 > 17 \Rightarrow x > 7
        \end{cases}
        \]
        得 $x \in \{8,9,10,11,12,13\}$
        
        \textbf{情况二:} 若 $x > 17$,最大边是 $x$,成立条件为:
        \[
        \begin{cases}
        x^2 > 10^2 + 17^2 = 389 \Rightarrow x > \sqrt{389} \approx 19.7\\
        10 + 17 > x \Rightarrow x < 27
        \end{cases}
        \]
        得 $x \in \{20,21,22,23,24,25,26\}$
        
        故 $$\sum x=224$$
    \end{solution}
    
    \question 若$\triangle ABC$ 的三个高分别为 $\dfrac{1}{2},\ \dfrac{1}{3},\ \dfrac{1}{4}$,求$\triangle ABC$ 的周长。
    \begin{solution}
        设三角形三边长为 $a,\ b,\ c$,对应的高分别为 $\dfrac{1}{2},\ \dfrac{1}{3},\ \dfrac{1}{4}$,则$\triangle ABC$
        面积为\[
        \frac{1}{2}ah_a = \frac{1}{2}bh_b = \frac{1}{2}ch_c 
        \Rightarrow \frac{a}{2} = \frac{b}{3} = \frac{c}{4}
        \]
        现设$a = 2t, b = 3t, c = 4t$,则半周长
        \[
        s = \frac{a + b + c}{2} = \frac{9t}{2}
        \]
        由海伦公式,
        \[
        [\triangle ABC] 
        = \sqrt{\frac{9t}{2} \cdot \frac{5t}{2} \cdot \frac{3t}{2} \cdot \frac{t}{2}}
        =\frac{1}{2} \cdot 2t \cdot \frac{1}{2}
        \]
        解得$t=\dfrac{2\sqrt{15}}{45}$,因此周长为
        \[
        2s = 9t = \frac{2\sqrt{15}}{5}
        \]
    \end{solution}

    \question 设 \(\triangle ABC\) 的内角 \(A,B,C\) 对应边长为 \(a,b,c\)。若 \(b+c=2a,\; 3\sin A = 5\sin B\),求 \( C\)。  
    \begin{solution}
        由正弦定理,
        \[
        3\sin A = 5\sin B \Rightarrow 3a = 5b \Rightarrow b=\frac{3}{5}a
        \]
        又\[
        b+c=2a \Rightarrow c=\frac{8}{5}a
        \]
        由余弦定理,
        \[
        \left(\frac{8}{5}a\right)^2=a^2+\left(\frac{3}{5}a\right)^2-2a\left(\frac{3}{5}a\right)\cos C
        \]
        给出 
        \[
        C=120^\circ
        \]
    \end{solution}
    
    \question 若 \(\triangle ABC\) 的面积为 \(\dfrac{3\sqrt{15}}{4}\),且 \(\cos C = -\dfrac14\),且 \(\sin^2A + \sin^2B = \dfrac{13}{16}\sin^2C\),求边长 \(AB\)。  
    \begin{solution}
        已知面积
        \[
        \frac{1}{2} ab \sin C = \frac{3\sqrt{15}}{4}
        \]
        又由 \[
        \cos C = -\frac{1}{4} \Rightarrow \sin C = \frac{\sqrt{15}}{4}
        \]
        代入上式得 $ab = 6$, 且由正弦定理
        \[
        \sin^2 A + \sin^2 B = \frac{13}{16} \sin^2 C  \Rightarrow a^2 + b^2 = \frac{13}{16} c^2 
        \]
        又由余弦定理:
        \[
        c^2=a^2 + b^2-2ab\cos C =\frac{13}{16} c^2 -2 \cdot 6 \cdot \left(-\dfrac14\right) 
        \]
        得 $c=4$
    \end{solution}
    
    \question 已知 \(\triangle ABC\) 的周长为 \(14\,\text{cm}\),面积为 \(3\sqrt7\,\text{cm}^2\),且 \(\cos B = -\dfrac18\)。求 \(a,b,c\) 的边长。  
    \begin{solution}
        解法同上,已知周长\(a + b + c = 14\),面积\[
         \frac{1}{2} ac \sin B = 3\sqrt7
        \] 
        由\[
        \cos B =  -\dfrac18 \Rightarrow \sin B = \frac{3\sqrt{7}}{8}
        \]
        有 $ac=16$, 且由余弦定理, \[
        b^2=a^2+c^2-2ac \cos B =(a+c)^2-2ac-2ac \cos B =(14-b)^2-2\cdot 16-2\cdot 16 \cdot  \left(-\dfrac18\right)
        \]
        解得 $b=6$,即有$a+c=8$,所以$a,c$是关于$x$的方程式
        \[
        x^2-8x+16=0
        \]
        的两根,于是解得$a=4,\;c=4$,故解为$ a=4,\; b=6,\; c=4$
    \end{solution}
    
    \question 在锐角 \(\triangle ABC\) 中,已知 \(b=5,\; \sin A = \dfrac{\sqrt7}{4}\),面积为 \(\dfrac{15\sqrt7}{4}\),求 \(c\) 及 \(\sin C\)。  
    \begin{solution}
        由面积公式:
        \[
        S = \frac{1}{2}bc\sin A = \frac{1}{2} \cdot 5 \cdot c \cdot \frac{\sqrt{7}}{4} = \frac{15\sqrt{7}}{4}
        \]
        得 $c = 6$,由余弦定理,\[
            a^2=b^2+c^2-2bc \cos A
        \]
        其中$\sin A = \dfrac{\sqrt7}{4} \Rightarrow \cos A = \dfrac{3}{4}$(注意是锐角),上式变成\[
        a^2=5^2+6^2-2\cdot 5 \cdot 6 \cdot \dfrac{3}{4} \Rightarrow a=4
        \]再由正弦定理,\[
        \frac{a}{\sin A}= \frac{c}{\sin C} \Rightarrow \frac{4}{\frac{\sqrt7}{4}}= \frac{6}{\sin C} \Rightarrow \sin C=\dfrac{3\sqrt7}{8}
        \]
    \end{solution}

    \question 已知 $\triangle ABC$ 中,$\;C = 60^\circ$,求
    \[
    \frac{a}{b+c} + \frac{b}{a+c}
    \]
    的值。
    \begin{solution}
        由余弦定理,
        \[
        a^2 + b^2 = c^2+ 2ab \cos C = c^2+2ab \cdot \frac{1}{2} = c^2+ ab
        \]
        于是
        \[
        \frac{a}{b+c} + \frac{b}{a+c}
        = \frac{a^2 + ac + b^2 + bc}{ab+bc+ac+c^2} 
        = \frac{a^2 + b^2 + ac + bc}{bc+ac+a^2+b^2} 
        = 1
        \]
    \end{solution}

    \question 已知 $\triangle ABC$ 中,
    \[
    2a \cos B = c, \quad \sin A \sin B (2-\cos C) = \sin^2 \frac{C}{2} + \frac{1}{2},
    \]
    证明$\triangle ABC$ 为等腰直角三角形。
    \begin{solution}
        据题意,
        \[
        2a \cos B = c \Rightarrow 2 \sin A \cos B = \sin C = \sin(A+B) = \sin A \cos B + \cos A \sin B
        \]
        得
        \[
        \sin(A-B)= 0 \Rightarrow A=B \in (0,\pi)
        \]
        于是
        \[
        \sin A \sin B (2 - \cos C) = \sin^2 \frac{C}{2} + \frac{1}{2}
        \]
        化为
        \begin{align*}
            \sin^2 A (2 + \cos 2A) &= \cos^2 A + \frac{1}{2}, \\
            \sin^2 A (3 - 2 \sin^2 A) &= 1 - \sin^2 A + \frac{1}{2}, \\
            (2 \sin^2 A - 1)(2 \sin^2 A - 3) &= 0,
        \end{align*}
        舍去$\sin^2 A=\dfrac{3}{2}$,解得
        \[
        \sin A = \frac{\sqrt{2}}{2} >0\Rightarrow 
        A = B = \frac{\pi}{4}, C = \frac{\pi}{2}
        \]
        因此,$\triangle ABC$ 为等腰直角三角形。
    \end{solution}

    \question 单位圆上有三点 $A,B,C$,且
    \[
    AB^2+BC^2+CA^2=8
    \]
    证明:$\triangle ABC$ 为直角三角形。
    \begin{solution}
        三点$A,B,C$在单位圆上,外接圆半径为 $1$,由正弦定理,
        \[
        AB=2\sin C,\quad BC=2\sin A,\quad CA=2\sin B
        \]
        代入已知等式得
        \begin{align*}
        \sin^2 A + \sin^2 B + \sin^2 C = 2 &\Rightarrow \frac{1-\cos 2A}{2} + \frac{1-\cos 2B}{2} + \frac{1-\cos 2C}{2} = 2 \\ 
        &\Rightarrow \cos 2A + \cos 2B + \cos 2C = -1
        \end{align*}
        和差化积得
        \[
        -2\cos C\cos(A-B)+2\cos^2C-1=-1 \Rightarrow 2\cos C(\cos C-\cos(A-B))=0
        \]
        若 $\cos C=0$,则 $C=\dfrac{\pi}{2}$。若 $\cos C=\cos(A-B)$,则 $C=A-B$ 或 $C=B-A$,从而 $A=\dfrac{\pi}{2}$ 或 $B=\dfrac{\pi}{2}$。于是得证$\triangle ABC$ 为直角三角形。
    \end{solution}

    \question 设某三角形三边长成等差数列,公差为 $d$,若 $r,R$ 分别表示此三角形的内切圆半径与外接圆半径,试证公差
    \[
    d = \sqrt{2Rr - 4r^2}
    \]
    \begin{solution}
        设三角形三边为 $a<b<c$,三边长成等差,故设
        \[
        a = b-d, c = b+d \quad \Rightarrow s = \frac{a+b+c}{2}= \frac{3b}{2},
        \]
        其中$d>0$,则三角形面积为
        \[
        \triangle = r s = \sqrt{s(s-a)(s-b)(s-c)} = \frac{abc}{4R}
        \]
        即
        \[
        \frac{3br}{2} = \sqrt{\frac{3b}{2} \cdot \frac{b}{2}  \cdot \left(\frac{b^2}{4} - d^2\right)}= \frac{b(b^2 - d^2)}{4R} \Rightarrow 
        \begin{cases}
            b^2 = 4(d^2 + 3 r^2)\\ b^2 - d^2 = 6 R r
        \end{cases}
        \]
        因此,
        \[
        d^2 = 2Rr - 4 r^2 \Rightarrow d = \sqrt{2Rr - 4 r^2}
        \]
    \end{solution}

    \question 已知三角形 $ABC$ 及其 $BC$ 边上的任意一点 $P$ (其中 $P$ 异于 $B,C$ 两点),试证斯德瓦特定理: 
    \[
    (AP^2 + BP\cdot CP)BC= AB^2\cdot CP + AC^2\cdot BP.
    \]
    \begin{solution}
        在$\triangle APB$及$\triangle APC$中,由余弦定理,
        \[
        \begin{cases}
        \cos \angle APB = \dfrac{AP^2 + BP^2 - AB^2}{2AP\cdot BP}, \\[2mm]
        \cos \angle APC = \dfrac{AP^2 + CP^2 - AC^2}{2AP\cdot CP}.
        \end{cases}
        \]
        由于 $\angle APB + \angle APC = \pi \Rightarrow \cos \angle APB = -\cos \angle APC$,于是  
        \[
        CP(AP^2 + BP^2 - AB^2) = -BP(AP^2 + CP^2 - AC^2)
        \]
        整理得  
        \[
        AP^2(CP+BP) + CP\cdot BP(CP+BP) = AB^2\cdot CP + AC^2\cdot BP
        \]
        又 $CP+BP=BC$,所以  
        \[
        (AP^2 + BP\cdot CP)BC = AB^2\cdot CP + AC^2\cdot BP
        \]
    \end{solution}

    \question 如图,\;\(\triangle ABC\) 中 \(AB=10,\,AC=14,\,B=\dfrac{\pi}{3}\)。\(D\) 在 \(BC\) 上且 \(DC=6\)。  
    \begin{figure}[H]
        \centering
        \includegraphics[width=0.3\textwidth]{images/image1.png}
    \end{figure}
    \begin{parts}
        \part 求 \(\angle ADB\)。 
        \begin{solution}
            $\triangle ABC$中,由余弦定理,\[
            14^2=10^2+BC^2-2\cdot10\cdot BC \cdot \cos \frac{\pi}{3} \Rightarrow BC=16
            \]又$BD=16-6 = 10$,则$\triangle ABD$ 为等边三角形,故\(\angle ADB=60^\circ\)
        \end{solution}
        \part 求 \(\sin\angle DAC\)。
        \begin{solution}
            $\triangle ADC$中,由余弦定理,\[
            6^2=10^2+14^2-2\cdot10\cdot14\cdot \cos \angle DAC \Rightarrow \cos \angle DAC = \frac{13}{14}
            \]给出\[
                \sin \angle DAC = \dfrac{3\sqrt3}{14} 
            \]
        \end{solution}
    \end{parts}
    
    \question 在 \(\triangle ABC\) 中,点 \(D\) 在 \(BC\) 上且 \(AD\) 为 \(\angle BAC\) 的角平分线。  
        已知 \(AB=12,\;AD=9,\;AC=15\),令$\theta=\angle BAD=\angle DAC$, 求 $\cos \theta$。 
    \begin{figure}[H]
        \centering
        \includegraphics[width=0.3\textwidth]{images/image2.png}
    \end{figure}
    \begin{solution}
        考虑面积法,由正弦定理, 
        \[
        \frac{1}{2}\cdot12\cdot9\cdot \sin \theta + \frac{1}{2}\cdot9\cdot15\cdot \sin \theta=\frac{1}{2}\cdot12\cdot15\cdot \sin 2\theta
        \]
        化简得
        \[
        243\sin\theta = 360\sin\theta\cos\theta \Rightarrow \cos\theta = \dfrac{27}{40}
        \]
    \end{solution}

    \question 在$\triangle ABC$ 中,过 $A$ 与过 $B$ 的中线互相垂直。已知 $AC=6,BC=7$,求 $AB$。
    \ifprintanswers
    \begin{figure}[H]
        \centering        
        \includegraphics[width=0.4\textwidth]{images/image146.png}
    \end{figure}
    \fi
    \begin{solution}
        设$D,E$在$BC,AC$上使得中线$AD,BE$交于重心$G$,在$\triangle BDG$ 及$\triangle AEG$中,由毕氏定理,
        \[
        x^{2}+4y^{2}=\frac{49}{4},
        \quad
        4x^{2}+y^{2}=9.
        \]
        因此
        \[
        AB^{2}=4x^{2}+4y^{2}=\frac{4}{5}\Bigl(\frac{49}{4}+9\Bigr)=17 \Rightarrow AB=\sqrt{17}
        \]
    \end{solution}

    \question $\triangle ABC$ 中,$AB = AC = 3$, $AD = 1$, $\angle BCD = 60^\circ$, 求 $\cos A$。  
    \ifprintanswers
    \begin{figure}[H]
        \centering        
        \includegraphics[width=0.3\textwidth]{images/image35.png}
    \end{figure}
    \fi
    \begin{solution}
        设 $\angle ACD = \theta$, 则
        \[
        \angle ADC = 120^\circ + \theta, \angle B = 60^\circ + \theta
        \]
        在 $\triangle ADC$ 中,由正弦定理,
        \[
        \frac{1}{\sin \theta} = \frac{3}{\sin(120^\circ + \theta)} \Rightarrow \frac{\sqrt3}{2}\cos \theta - \frac{1}{2}\sin\theta=3\sin \theta \Rightarrow \tan \theta = \frac{\sqrt{3}}{7}
        \]
        由此可得
        \[
        \sin \theta = \frac{\sqrt{3}}{2\sqrt{13}},\quad \cos \theta = \frac{7}{2\sqrt{13}}
        \]
        在 $\triangle BCD$ 中,
        \[
        \frac{BC}{\sin(60^\circ - \theta)} = \frac{2}{\sin 60^\circ} \Rightarrow BC = \frac{4}{\sqrt{3}} \sin(60^\circ - \theta) = \frac{6}{\sqrt{13}}
        \]
        在 $\triangle ABC$ 中,由余弦定理,
        \[
        \cos A = \frac{9 + 9 - \frac{36}{13}}{2 \cdot 9} = \frac{11}{13}
        \]
    \end{solution}

    \question 在$\triangle ABC$中,在 $BC$ 边上取一点 $D$ 使得 $BD = AC$。若 $\angle DAC=60^\circ,\angle ACD=40^\circ$,求 $\angle BAD$。
    \ifprintanswers
    \begin{figure}[H]
        \centering
        \includegraphics[width=0.5\linewidth]{images/image44.png}
    \end{figure}
    \fi
    \begin{solution}
        设$BD = AC = a,AD = b, \angle BAD = \theta$,则
        \[
        \angle ADC = 80^\circ \Rightarrow \angle B = 80^\circ - \theta
        \]
        在$\triangle ADC$及$\triangle ABD$中,由正弦定理,
        \[
        \frac{a}{\sin 80^\circ} = \frac{b}{\sin 40^\circ}, \quad \frac{a}{\sin \theta} = \frac{b}{\sin (80^\circ - \theta)}
        \]
        消去 \( a \)得
        \[
        \frac{\sin \theta}{\sin(80^\circ - \theta)} = \frac{\sin 80^\circ}{\sin 40^\circ} = 2\cos 40^\circ
        \]
        整理得
        \[
        \sin \theta = 2 \cos 40^\circ \sin(80^\circ - \theta) 
        = \sin(120^\circ - \theta) + \sin(40^\circ - \theta)
        \]
        \[
        \sin \theta + \sin(\theta - 120^\circ) =2 \sin(\theta - 60^\circ) \cos 60^\circ = \sin(40^\circ - \theta)
        \]
        \[
        \sin(\theta - 60^\circ) = \sin(40^\circ - \theta)
        \Rightarrow \angle BAD=\theta = 50^\circ 
        \]
    \end{solution}

    \question 已知 $\triangle ABC$ 中, $AB=AC$, 且 $\angle B$ 的角平分线交 $AC$ 于 $D$, 且 $BC=BD+AD$, 求$\angle A$ 。
    \ifprintanswers
    \begin{figure}[H]
        \centering
        \includegraphics[width=0.35\linewidth]{images/image226.png}
    \end{figure}
    \fi
    \begin{solution}
        设 $\angle B=\angle C=2\theta$,则  
        \[
        \angle ABD=\angle DBC=\theta, \quad
        \angle A=\pi-4\theta,\quad
        \angle BDC=\pi-3\theta
        \]
        在$\triangle BDC$及$\triangle ABD$中,由正弦定理,
        \[
        \frac{BC}{BD}
        =\frac{\sin(\pi-3\theta)}{\sin 2\theta}
        =\frac{\sin 3\theta}{\sin 2\theta}, \quad
        \frac{AD}{BD}
        =\frac{\sin\theta}{\sin(\pi-4\theta)}
        =\frac{\sin\theta}{\sin 4\theta}
        \]
        由 $BC=BD+AD$,有
        \begin{align*}
        \frac{BC}{BD} &=1+\frac{AD}{BD} \\
        \frac{\sin 3\theta}{\sin 2\theta}&=1+\frac{\sin\theta}{\sin 4\theta}=1+\frac{\sin\theta}{2\sin 2\theta \cos 2\theta}\\
        2\cos 2\theta \sin 3\theta&=2\sin 2\theta \cos 2\theta+\sin\theta\\
        \sin 5\theta+\sin\theta&=\sin 4\theta+\sin\theta \\
        \sin 5\theta&=\sin 4\theta
        \end{align*}
        故
        \[
        \theta=\frac{\pi}{9} \Rightarrow \angle A=\frac{5\pi}{9}
        \]
    \end{solution}

    \question 设 $D$ 为 $\triangle ABC$ 的 $BC$ 上之一点,且 $BD=AC=1$,若 $\angle BAD=30^\circ$、$\angle CAD=90^\circ$,求 $CD$ 之长。 
    \ifprintanswers
    \begin{figure}[H]
        \centering        
        \includegraphics[width=0.4\textwidth]{images/image121.jpg}
    \end{figure}
    \fi
    \begin{solution}
        设 $CD=a$,由等高性质,
        \[
        \frac{[ABD]}{[ACD]}=\frac{BD}{CD}=\frac{1}{a} \tag{1}
        \]
        故
        \[
        \frac{[ABD]}{[ACD]}=\frac{\frac{1}{2}\,AD\cdot AB\sin30^\circ}{\frac{1}{2}\,AD\cdot AC}=\frac{AB}{2} \tag{2}
        \]
        由 $(1)=(2)$ 得 $AB=\dfrac{2}{a}$,由余弦定理,
        \[
        (a+1)^2=1+\left(\dfrac{2}{a}\right)^2-2\cdot \dfrac{2}{a}\cdot 1 \cdot \cos120^\circ
        \]
        化简得
        \[
        (a^3-2)(a+2)=0
        \]
        取正根 $a^3=2$,故 $a=\sqrt[3]{2}$。
    \end{solution}

    \question 已知一个半径为 1 的圆中,有两条相等的弦互相垂直,且互相分割成 $1:4$的线段。求弦长。
    \begin{solution}
        构造坐标系,使得圆心在原点上,两弦为 $x=-a$及 $y=-a$,且 $a>0$。由题意,
        \[
        4(\sqrt{1-a^{2}}-a)=\sqrt{1-a^{2}}+a
        \]
        解得
        \[
        a^{2}=\frac{9}{34}
        \]
        因此弦长为
        \[
        2\sqrt{1-a^{2}}=\frac{5\sqrt{34}}{17}
        \]
    \end{solution}

    \question 圆内接四边形 \(ABCD\) 中,\(\;AB=BC=3,\;CD=5,\;DA=8\)。求对角线 \(BD\)。  
    \begin{solution}
        圆内接四边形对角互补,有
        \[
        BD^2=3^2+8^2-2\cdot3\cdot8\cdot \cos A = 3^2+5^2-2\cdot3\cdot5\cdot \cos (\pi- A)
        \]
        解得 
        \[
        \cos A = \frac{1}{2} \Rightarrow BD=7
        \]
    \end{solution}
    
    \question 已知四边形 $ABCD$ 内接于圆,且 $BA=AD=1,\cos\angle BAD=-\dfrac{1}{3}$。求证 $BC$为外接圆的直径
    \ifprintanswers
    \begin{figure}[H]
        \centering        
        \includegraphics[width=0.4\textwidth]{images/image196.png}
    \end{figure}
    \fi
\begin{solution}
在$\triangle BAD$中,由余弦定理,
\[
BD^2 = 1^2 + 1^2 - 2\cdot 1 \cdot 1 \cdot \left(-\frac{1}{3}\right) \Rightarrow BD= \sqrt{\frac{8}{3}}
\]
设 \textcolor{red}{$x=\cos\angle ABC$,则 $DC=x$}。由于 $ABCD$ 为圆内接四边形,$\angle ADC=180^\circ-\angle ABC$,故
\[
\cos\angle ADC=-\cos\angle ABC=-x.
\]
同理,$\cos\angle BCD=-\cos\angle BAD=\dfrac{1}{3}$。

在 $\triangle ADC$ 与 $\triangle ABC$ 中应用余弦定理,因 $AC$ 公共,得:
\begin{align*}
1^2+x^2-2(1)(x)\cos\angle ADC &= 1^2+BC^2-2(1)(BC)\cos\angle ABC,\\
1+x^2-2x(-x) &= 1+BC^2-2BCx,\\
0 &= BC^2-2BCx-3x^2,\\
0 &= (BC-3x)(BC+x).
\end{align*}

因 $x>0$,故 $BC=3x$。

又因 $\cos\angle BCD=\tfrac{1}{3}$,且 $DC:BC=1:3$,由$\angle BDC$ 为直角。于是 $BC$ 为圆的直径。
\end{solution}

    \question 已知$\triangle ABC$ 中,$\angle B=90^\circ,\angle A=30^\circ$。点 $P, Q, R$ 分别在 $AB, BC, CA$ 上使得$\triangle PQR$ 为正三角形。若$BC=4$,且 $Q$ 是 $BC$ 的中点,求 $PR$。
    \begin{solution}
        已知$BQ=QC=2$,设 $PB=x,PQ=y,\angle PQB=\angle CRQ=\theta$,在$\triangle PBQ$及$\triangle QRC$中, 由正弦定理,
        \[
        \frac{x}{\sin\theta}=y, \quad \frac{2}{\sin\theta}=\frac{y}{\frac{\sqrt{3}}{2}}
        \]
        解得 $x=\sqrt{3}$,在 $\triangle PBQ$ 中,由毕氏定理
        \[
        y^2 = 3 + 2^2 =7 \Rightarrow PR=y=\sqrt{7}
        \]
    \end{solution}

    \question 已知在$\triangle ABC$中,三边 $AB,BC,CA$ 长度比为 $4:6:5$,三角形内一点 $P$ 满足 $PA=8,PB=4,PC=13$,试求 $\cos\angle APB$ 的值。
    \ifprintanswers
    \begin{figure}[H]
        \centering        
        \includegraphics[width=0.5\textwidth]{images/image25.png}
    \end{figure}
    \fi
    \begin{solution}
        在 ${BC}$ 外取一点 $Q$使得 ${BQ}=6$ 且 $\angle ABC=\angle PBQ$,由相似三角形:
        \[
        \frac{{AB}}{{BC}} = \frac{{BP}}{{BQ}} = \frac{4}{6} \Rightarrow \triangle ABC \backsim \triangle PBQ \  (SAS)\Rightarrow {PQ}=5
        \]
        且有
        \[
        \angle ABP=\angle CBQ \Rightarrow \triangle PBA \backsim \triangle QBC \ (SAS) \Rightarrow {QC}=12
        \]
        由余弦定理,
        \[
        \cos \angle BQP = \frac{5^2 + 6^2 - 4^2}{2 \cdot 5 \cdot 6} = \frac{3}{4} \Rightarrow \sin \angle BQP = \frac{\sqrt{7}}{4}
        \]
        又 $\triangle PQC$ 三边长为 $5,12,13$,故为直角三角形,$\angle PQC=90^\circ$,于是
        \[
        \cos \angle APB = \cos \angle BQC = \cos(\angle BQP + 90^\circ) = -\sin \angle BQP = -\frac{\sqrt{7}}{4}
        \]
    \end{solution}

    \question 在 $\triangle ABC$ 中,$AB=4,AC=6,\cos(B-C)=\dfrac{2}{3}$,求 $BC$。
    \ifprintanswers
    \begin{figure}[H]
        \centering        
        \includegraphics[width=0.35\textwidth]{images/image86.jpg}
    \end{figure}
    \fi
    \begin{solution}
        $AC$ 上找一点 $P$使得 $\angle PBC = \angle C = \theta$,令 $PB = PC = a$,则 $AP = 6 - a$, 由余弦定理,
        \[
        (6-a)^2=4^2 + a^2 - 2 \cdot 4 \cdot a \cdot \frac{2}{3} \Rightarrow a = 3
        \]
        设 $BC = b$,则
        \[
        \cos \angle PBC = \cos \theta = \frac{a^2 + b^2 - a^2}{2ab} ,\quad
        \cos \angle ACB = \cos \theta = \frac{b^2 + 6^2 - 4^2}{12b}
        \]
        联立得
        \[
        \frac{b}{6} = \frac{b^2 + 20}{12b} \Rightarrow b = 2\sqrt{5}
        \]
    \end{solution}

    \question 在 $\triangle ABC$ 中, $AB=1$, $AC=2$, $B-C=\dfrac{2\pi}{3}$ ,求 $\triangle ABC$ 的面积。
    \begin{solution}
        由正弦定理知 
        \[
        \frac{\sin B}{\sin C}=\frac{AC}{AB}=2
        \]
        又 $B-C=\dfrac{2\pi}{3}$,故
        \[
        2\sin C=\sin B=\sin(C+\frac{2}{3}\pi)=-\frac{1}{2}\sin C+\frac{\sqrt{3}}{2}\cos C
        \]
        即 
        \[
        \frac{5}{2}\sin C=\frac{\sqrt{3}}{2}\cos C\Rightarrow\tan C=\frac{\sqrt{3}}{5}
        \]
        注意到 $A=\pi-B-C=\dfrac{\pi}{3}-2C$ ,故$\triangle ABC$ 的面积为
        \[
        [\triangle ABC]=\frac{1}{2}AB\cdot AC\cdot \sin A=\sin \left(\frac{\pi}{3}-2C\right)=\frac{\sqrt{3}}{2}\cos 2C-\frac{1}{2}\sin 2C
        \]
        由 $\tan C=\dfrac{\sqrt{3}}{5}$ 知 
        \[
        \cos 2C=\frac{1-\tan^2 C}{1+\tan^2 C}=\frac{11}{14},\ \sin 2C=\frac{2\tan C}{1+\tan^2 C}=\frac{5\sqrt{3}}{14}
        \]
        从而
        \[
        [\triangle ABC]=\frac{\sqrt{3}}{2}\cdot\frac{11}{14}-\frac{1}{2}\cdot\frac{5\sqrt{3}}{14}=\frac{3\sqrt{3}}{14}\]
    \end{solution}

    \question 设 $\triangle ABC$ 的三边长为 $AB=4, BC=5, CA=6,D,E,F$分别在边$BC,CA,AB$上使得三高为 $AD, BE, CF$,试求三面积比 $[\triangle AEF] : [\triangle BDF] : [\triangle CDE]$。
    \begin{solution}
        由
        \[
        \frac{[\triangle AEF]}{[\triangle ABC]} 
        =\frac{\frac12 AE \cdot AF \sin  A}{\frac12 AB \cdot AC \sin  A}
        =\frac{AB \cos  A \cdot AC \cos  A}{AB \cdot AC} 
        =\cos^2  A
        \]
        同理可得
        \[
        \frac{[\triangle BDF]}{[\triangle ABC]} = \cos^2  B,\quad \frac{[\triangle CDE]}{[\triangle ABC]} = \cos^2 C
        \]
        因此
        \[
        [\triangle AEF] : [\triangle BDF] : [\triangle CDE] = \cos^2  A : \cos^2  B : \cos^2  C
        \]
        由余弦定理,
        \[
        \cos  A = \frac{4^2+6^2-5^2}{2\cdot4\cdot6} = \frac{9}{16},\;
        \cos  B = \frac{4^2+5^2-6^2}{2\cdot4\cdot5} = \frac{1}{8},\;
        \cos  C = \frac{5^2+6^2-4^2}{2\cdot5\cdot6} = \frac{3}{4}
        \]
        所以
        \[
        [\triangle AEF] : [\triangle BDF] : [\triangle CDE] = \frac{81}{256} : \frac{1}{64} : \frac{9}{16} = 81 : 4 : 144
        \]
    \end{solution}

    \question $\triangle ABC$ 中,在 $BC$ 边上取 $D,E$,使得  
    \[
    \angle BAD = \angle DAE = \angle EAC.
    \] 
    已知 $BD=3,DE=4,EC=8$,求 $\triangle ABC$ 的面积。
    \begin{solution}
        设 $\angle BAD = \angle DAE = \angle DAC = \theta, AB=b, AC=c, AD=d, AE=e$,由等高性质,
        \[
        [\triangle ABD] : [\triangle ADE] : [\triangle AEC] = BD : DE : EC = 3:4:8
        \]
        又有
        \[
        [\triangle ABD] = \frac12 bd \sin \theta, \quad 
        [\triangle ADE] = \frac12 de \sin \theta, \quad
        [\triangle AEC] = \frac12 ec \sin \theta.
        \]
        故$b:e = 3:4, d:c = 1:2$,令 $b = 3q,\; c = 2p, \; d = p, \; e = 4q$,由余弦定理,
        \[
        \cos \theta = \frac{b^2+d^2-9}{2bd} = \frac{d^2+e^2-16}{2de} = \frac{e^2+c^2-64}{2ec}.
        \]
        解得
        \[
        \frac{9q^2+p^2-9}{6pq} = \frac{p^2+16q^2-16}{8pq} = \frac{16q^2+4p^2-64}{16pq} \Rightarrow p=6\sqrt 2, \; q=\sqrt 7
        \]
        于是
        \[
        \cos \theta = \frac{\sqrt{14}}{4} \Rightarrow \sin \theta = \frac{\sqrt 2}{4}.
        \]
        $\triangle ABC$ 面积为
        \[
        [\triangle ABC] = \frac12 bc \sin 3\theta = \frac12 (3\sqrt 7)(12\sqrt 2) \bigl(3\sin \theta - 4\sin^3\theta \bigr) 
        = \frac{45\sqrt 7}{2} 
        \]
    \end{solution}

    \question 设 \(a,b,c\) 分别为 \(\triangle ABC\) 的 \(A, B, C\) 对边长,且 \(\angle B = 60^\circ\),证明
    \[
    (a+b+c)\left(\frac{1}{a+b}+\frac{1}{b+c}\right)=3
    \]
    \begin{solution}
        展开左式得
        \[
        (a+b+c)\left(\frac{1}{a+b} + \frac{1}{b+c}\right) = 1 + \frac{c}{a+b} + 1 + \frac{a}{b+c} = 2 + \frac{a^2 + c^2 + ab + bc}{(a+b)(b+c)}
        \]
        由余弦定理,
        \[
        b^2=a^2+c^2-2ac\left(\frac{1}{2}\right) \Rightarrow a^2 + c^2 = b^2 + ac 
        \]
        因此左式为 
        \[
        2 + \frac{b^2 + ac + ab + bc}{(a+b)(b+c)} = 2 + \frac{(a+b)(b+c)}{(a+b)(b+c)} = 3
        \]
    \end{solution}

    \question 在$\triangle ABC$中,其内角$A,B,C$所对应的边分别为$a,b,c$,且
    \[
    \dfrac{b}{a}+\dfrac{a}{b}=4\cos C
    \]
    求$\tan C(\cot A+\cot B)$的值。
    \begin{solution}
        由已知$$\frac{b}{a}+\frac{a}{b}=4\cos C$$及余弦定理,
        \[
        \cos C= \frac{a^2+b^2}{4ab} = \frac{a^2+b^2-c^2}{2ab} \Rightarrow a^2+b^2 = 2c^2
        \]
        于是由余弦、正弦定理,
        \begin{align*}
        &\tan C(\cot A+\cot B) \\
        &= \frac{\sin C}{\cos C}\left( \frac{\cos A}{\sin A} + \frac{\cos B}{\sin B} \right)\\
        &= \frac{\cos A}{\cos C} \cdot \frac{\sin C}{\sin A} + \frac{\cos B}{\cos C} \cdot \frac{\sin C}{\sin B}\\
        &= \left( \frac{b^2+c^2-a^2}{2bc} \cdot \frac{2ab}{a^2+b^2-c^2} \right) \cdot \frac{c}{a} + \left( \frac{a^2+c^2-b^2}{2ac} \cdot \frac{2ab}{a^2+b^2-c^2} \right) \cdot \frac{c}{b} \\
        &= \frac{b^2+c^2-a^2}{a^2+b^2-c^2} + \frac{a^2+c^2-b^2}{a^2+b^2-c^2}
        = \frac{2c^2}{2c^2-c^2} = 2
        \end{align*}
    \end{solution}

    \question 设 $\triangle ABC$ 中 $\angle A,\angle B,\angle C$ 的对边分别为 $a,b,c$,若 $a^2+b^2=8c^2$,求
    \[
    \frac{\tan C}{\tan A} + \frac{\tan C}{\tan B}
    \]
    的值。
    \begin{solution}
        设 $\triangle ABC$ 面积为 $S$,则
        \[
        S = \frac{1}{2}bc \sin A \Rightarrow \sin A = \frac{2S}{bc}, \quad \cos A = \frac{b^2+c^2-a^2}{2bc}
        \]
        因此
        \[
        \tan A = \frac{4S}{b^2+c^2-a^2}, \quad \tan B = \frac{4S}{c^2+a^2-b^2}, \quad \tan C = \frac{4S}{a^2+b^2-c^2}
        \]
        所以
        \[
        \frac{\tan C}{\tan A} + \frac{\tan C}{\tan B} = \frac{b^2+c^2-a^2}{a^2+b^2-c^2} + \frac{c^2+a^2-b^2}{a^2+b^2-c^2} = \frac{2c^2}{a^2+b^2-c^2}
        \]
        代入 $a^2+b^2 = 8c^2$得
        \[
        \frac{2c^2}{8c^2 - c^2} = \frac{2}{7}
        \]
    \end{solution}

    \question 设 $\triangle ABC$ 的三边长 $AB=c,\,BC=a,\,CA=b$, 且 
    \[
    |b-c|\cos\frac{A}{2}=5,\,(b+c)\sin\frac{A}{2}=10,
    \]求 $a$ 之值。 
    \begin{solution}
        两式平方得
        \[
        (b-c)^2 \cos^2 \frac{A}{2} + (b+c)^2 \sin^2 \frac{A}{2} = 125
        \]
        整理得
        \[
        (b^2+c^2)\cos^2 \frac{A}{2} + (b^2+c^2)\sin^2 \frac{A}{2} - 2bc(\cos^2 \frac{A}{2} - \sin^2 \frac{A}{2}) = 125
        \]
        故
        \[
        b^2+c^2-2bc\cos A = a^2 = 125 \Rightarrow a = 5\sqrt{5}
        \]
    \end{solution}

    \question 设 $\triangle ABC$ 中,$A,B,C$ 所对应的边长分别为 $a,b,c$,已知 \[\frac{b}{a}=\frac{|b^2+c^2-a^2|}{bc},\ \frac{c}{b}=\frac{|c^2+a^2-b^2|}{ca},\ \frac{a}{c}=\frac{|a^2+b^2-c^2|}{ab}\] 且 $a<b<c$,求 $A$。
    \begin{solution}
        首先由$a<b<c \Rightarrow b^2+c^2>a^2$,则
        \[
        \frac{b}{a} = \frac{|b^2+c^2-a^2|}{bc} = \frac{b^2+c^2-a^2}{bc} 
        \]
        于是由余弦、正弦定理,
        \[
        \frac{b}{2a} = \frac{b^2+c^2-a^2}{2bc} \Rightarrow \frac{\sin B}{2\sin A} = \cos A \Rightarrow \sin B = \sin2A \Rightarrow B = 2 A
        \]
        同理,
        \[
        \frac{c}{b} = \frac{|c^2+a^2-b^2|}{ca} \Rightarrow C = 2 B
        \]
        因此
        \[
        A +  B +  C =  7 A = \pi \Rightarrow A = \frac{\pi}{7}
        \]
        \textcolor{red}{(待验证$c^2+a^2-b^2>0$)}
    \end{solution}

    \question 在 $\triangle ABC$ 中,角 $A,B,C$ 所对的边分别为 $a,b,c$,其中 $b \ge a$。若 
    \[
    2a\cos(B+C)+c\cos B+b\cos C=0,
    \] 
    且 $\triangle ABC$ 的外接圆半径为 $2$,求 $2b - c$ 的取值范围。
    \begin{solution}
        由条件得
        \[
        2a\cos(B+C)+c\cos B+b\cos C = -2a\cos A + a = 0 \Rightarrow A = 60^\circ
        \]
        由正弦定理,
        \[
        \frac{a}{\sin A} = \frac{a}{\frac{\sqrt{3}}{2}} = 2 \cdot 2 \Rightarrow a = 2\sqrt{3}
        \]
        又因 $b \ge a$,所以
        \[
        4\sin B \ge 2\sqrt{3} \Rightarrow \sin B \ge \frac{\sqrt{3}}{2}
        \]
        且有
        \[
        2b - c = 2 \cdot 4\sin B - 4\sin C = 8\sin B - 4\sin(120^\circ - B)
        \]
        \[
        = 8\sin B - 4\left(\frac{\sqrt{3}}{2}\cos B + \frac{1}{2}\sin B\right)
        = 6\sin B - 2\sqrt{3}\cos B
        = 4\sqrt{3}\sin(B - 30^\circ)
        \]
        因为$60^\circ \le B < 120^\circ$故
        \[2b - c \in {[2\sqrt{3},\; 4\sqrt{3})}
        \]
    \end{solution}

    \question 在$\triangle ABC$中, 若$a^{2}+b^{2}+c^{2}=2\sqrt{3}ab\sin C$, 求 $\cos\dfrac{A}{2}+\cos\dfrac{B}{2}+\cos\dfrac{C}{2}$。
    \begin{solution}
        由余弦定理$\ c^2 = a^2 + b^2 - 2ab\cos C$,代入原式得:
        \[
        a^2 + b^2 + (a^2 + b^2 - 2ab\cos C) =2a^2 + 2b^2 - 2ab\cos C= 2\sqrt{3}ab\sin C
        \]
        整理得
        \[
        \frac{a}{b} + \frac{b}{a} = \cos C + \sqrt{3}\sin C =2\sin\left(C + \frac{\pi}{6}\right)
        \]
        由 AM-GM 不等式可知
        \[
        \frac{a}{b} + \frac{b}{a} \ge 2
        \]
        当且仅当 $a = b$ 时等号成立,此时
        \[
        \sin\left(C + \frac{\pi}{6}\right) = 1 \Rightarrow C = \frac{\pi}{3}
        \]
        因为 $a = b$ 且 $C = \dfrac{\pi}{3}$,三角形 $\triangle ABC$ 为等边三角形,故 $A = B = C = \dfrac{\pi}{3}$,于是
        \[
        \cos\frac{A}{2} + \cos\frac{B}{2} + \cos\frac{C}{2} = 3\cos\frac{\pi}{6} = \frac{3\sqrt{3}}{2}
        \]
    \end{solution}

    \question 设$\triangle ABC$的三个内角$A,B,C$所对应的边分别为$a,b,c$.若$a^4+b^4+c^4=2c^2(a^2+b^2)$且$A=63^\circ$,求$B$。
    \begin{solution}
        由展开式
        $$(a^2+b^2-c^2)^2 = a^4+b^4+c^4+2a^2b^2-2a^2c^2-2b^2c^2 = 2a^2b^2$$
        得$$a^2+b^2-c^2 = \pm\sqrt{2}ab$$
        由余弦定理得$$ \cos C = \pm\frac{\sqrt{2}}{2}$$
        当$\cos C=\frac{\sqrt{2}}{2}$,即$C=45^\circ$,故$B=72^\circ$
        
        当$\cos C=-\frac{\sqrt{2}}{2}$,即$C=135^\circ,A+C=198^\circ>180^\circ$,不合题意.
        
        $\therefore \;B=72^\circ$
    \end{solution}

    \question 设$a,b,c$为一直角三角形的三条边,$\;\theta$是该三角形最小的角。若$\dfrac{1}{a},\;\dfrac{1}{b},\;\dfrac{1}{c}$也构成另一直角三角形,证明$\sin\theta=\dfrac{\sqrt{5}-1}{2}$。
    \begin{solution}
        不妨设$a<b<c$,则$\dfrac{1}{c}<\dfrac{1}{b}<\dfrac{1}{a}$。因此有
        \[
        a^2+b^2=c^2 \quad \text{以及} \quad \frac{1}{a^2}=\frac{1}{b^2}+\frac{1}{c^2}
        \]
        于是有
        \[
        \frac{1}{a^2}-\frac{1}{c^2}=\frac{1}{c^2-a^2}
        \]
        进一步整理可得
        \[
        c^{4}+a^{4}-3c^2a^2=0
        \]
        求解得到
        \[
        a^2=\frac{3c^2-\sqrt{5}c^2}{2}
        \]
        由于$\theta$是最小角,$\ \sin\theta=\dfrac{a}{c}$,故
        \[
        \frac{a}{c}=\sqrt{\frac{6-2\sqrt{5}}{4}}=\frac{\sqrt{5}-1}{2}
        \]
        证毕。
    \end{solution}
    
    \question 已知三角形 $\triangle ABC$ 满足 $\sin A = \cos B = \dfrac{16}{21}\tan C$,且 ${AC} = 1$,求三角形面积。
    \begin{solution}
        因为 $\cos B = \sin\left(\dfrac{\pi}{2} \pm B\right)$,故
        \[
        \sin A = \sin\left(\frac{\pi}{2} \pm B\right)\Rightarrow A=\frac{\pi}{2} \pm B
        \]
        
        若 $ A = \dfrac{\pi}{2} - B$,则 $ C = \dfrac{\pi}{2}$,使 $\tan C$ 无定义,不合题意;于是
        \[
         A = \frac{\pi}{2} + B ,C = \pi -  A -  B = \frac{\pi}{2} - 2B
        \]
        所以
        \[
        \cos B = \frac{16}{21} \tan\left(\frac{\pi}{2} - 2B\right) = \frac{16}{21} \cot 2B = \frac{16}{21} \frac{\cos 2B}{\sin 2B}=\frac{16}{21}\cdot\frac{1-2\sin^{2}B}{2\sin B\cos B}
        \]
        解得
        $$(3\sin B-1)(7\sin^{2}B-3\sin B-8)=0,\sin B=\frac{1}{3}\ \text{或}\ \frac{3\pm\sqrt{233}}{14}$$ 
        注意到 $\sin B=\frac{3+\sqrt{233}}{14}>1$(不合题意) 及 $\sin B=\frac{3-\sqrt{233}}{14}<0\Rightarrow B>\frac{\pi}{2}$(不合题意), 故
        \[
        \sin B=\frac{1}{3}\Rightarrow \cos B = \sin A = \frac{2\sqrt{2}}{3}, \ \tan C = \frac{7\sqrt{2}}{8}
        \]
        由正弦定理
        \[
        \frac{{AC}}{\sin B} = \frac{{BC}}{\sin A} \Rightarrow {BC} = 2\sqrt{2}
        \]
        三角形面积为
        \[
        \frac{1}{2} \cdot {AC} \cdot {BC} \cdot \sin C = \frac{1}{2} \cdot 1 \cdot 2\sqrt{2} \cdot \frac{7}{9} = \frac{7\sqrt{2}}{9}
        \]
    \end{solution}
    
    \question 已知$\triangle ABC$内角$A,B,C$的对边分别为$a,\;b,\;c$,满足$b^2=ac$,点$D$在边$AC$上使得$AD = 2DC$,且$BD \sin\angle ABC = a \sin C$,求$\cos\angle ABC$。
    \begin{solution}
        \textbf{解法一}
        
        由题意知$\;BD = b,AD = \dfrac{2b}{3},DC = \dfrac{b}{3}$,且$\angle ADB,\angle CDB$互补,由余弦定理,
        \[
        \cos\angle ADB = -\cos\angle CDB\Rightarrow \frac{b^2 + \frac{4b^2}{9} - c^2}{2b \cdot \frac{2b}{3}}=-\frac{b^2 + \frac{b^2}{9} - a^2}{2b \cdot \frac{b}{3}}\Rightarrow2a^2 + c^2 = \dfrac{11b^2}{3}
        \]
        又$b^2 = ac$,所以$$6a^2 - 11ac + 3c^2 = 0\Rightarrow c = 3a\ \text{或}\ c = \dfrac{2}{3}a$$
        当$c = 3a$时,$\;b^2 = 3a^2$,由余弦定理得$\cos\angle ABC = \dfrac{7}{6}>1$,不合题意。
        
        当$c = \dfrac{2}{3}a$时,$\;b^2 = \dfrac{2}{3}a^2$,由余弦定理,
        \[
        \cos\angle ABC = \frac{a^2 + c^2 - b^2}{2ac} = \frac{a^2 + \frac{4}{9}a^2 - \frac{2}{3}a^2}{2a \cdot \frac{2}{3}a} = \frac{7}{12}
        \]
    \end{solution}
    \begin{solution}
        \textbf{解法二}
        
        已知 $AD = 2DC$,则 $[\triangle ABD] = \dfrac{2}{3} [\triangle ABC]$,即
        \[
        \frac{1}{2} \cdot \frac{2}{3} b^2 \sin\angle ADB = \frac{2}{3} \cdot \frac{1}{2} ac \sin\angle ABC,
        \]
        而 $b^2 = ac$,即 $\sin\angle ADB = \sin\angle ABC$,故有 $\angle ADB = \angle ABC$,从而 $\angle ABD = \angle C$。
        
        由 $b^2 = ac$,即 $\dfrac{b}{a} = \dfrac{c}{b}$,即 $\dfrac{CA}{CB} = \dfrac{BA}{BD}$,即 $\triangle ACB \backsim \triangle ABD \ (SAS)$,故
        \[
        \frac{AD}{AB} = \frac{AB}{AC} \implies \frac{2b}{c} = \frac{c}{b},
        \]
        又 $b^2 = ac$,所以 $c = \dfrac{2}{3}a$,则
        \[
        \cos\angle ABC = \frac{c^2 + a^2 - b^2}{2ac} = \frac{7}{12}.
        \]
    \end{solution}

    \question 已知三角形 $ABC$ 内切圆切$BC,CA,AB$于 $D,E,F$,若内切圆半径为 $r$,外接圆半径为 $R$,试证 
    \[
    \frac{[\triangle DEF]}{[\triangle ABC]} = \frac{r}{2R}.\]
    \begin{figure}[H]
        \centering
        \includegraphics[width=0.5\linewidth]{images/image43.png}
    \end{figure}
    \begin{solution}
        设三角形 $ABC$ 的边长为 $a, b, c$,内心 $I$,有 $ID = IE = IF = r,\triangle DEF$ 面积为
        \begin{align*}
        [\triangle DEF] 
        &= [\triangle IDE]+[\triangle IEF]+ [\triangle IFD] \\
        &= \frac{1}{2} r^2 (\sin (\pi-A) + \sin (\pi-B) + \sin (\pi-C)) \\
        &= \frac{1}{2} r^2 (\sin A + \sin B + \sin C)
        \end{align*}
        由正弦定理,
        \[
        \sin A = \frac{a}{2R},\quad \sin B = \frac{b}{2R},\quad \sin C = \frac{c}{2R}
        \]
        代入得
        \[
        [\triangle DEF] = \frac{1}{2} r^2 \left( \frac{a+b+c}{2R} \right) = \frac{r^2 s}{2R}
        \]
        因此:
        \[
        \frac{[\triangle DEF]}{[\triangle ABC]} = \frac{\frac{r^2 s}{2R}}{rs} = \frac{r}{2R}
        \]
    \end{solution}

    \question 设 $\triangle ABC$ 的三边长为 $a,b,c$,面积为 $S$,试证明
    \[
    \frac{a^{2} + b^{2} + c^{2}}{S} \ge 4\sqrt{3}.
    \]
    \begin{solution}
        已知三角形的三边为 $a,b,c$,其对角分别为 $A,B,C$。利用余弦定理与面积公式可得
        \[
        c^{2}=a^{2}+b^{2}-2ab\cos C, \quad S=\frac{1}{2}ab\sin C
        \]
        因此
        \begin{align*}
        a^{2}+b^{2}+c^{2}-4S\sqrt{3}&=2[a^{2}+b^{2}-ab(\cos C+\sqrt{3}\sin C)] \\
        &=2\left[a^{2}+b^{2}-2ab\sin\left(\frac{\pi}{6}+C\right)\right] \\
        &\ge 2(a^{2}+b^{2}-2ab) \\
        &=2(a-b)^{2} \\
        &\ge 0
        \end{align*}
        等号成立当且仅当 $a=b$ 且 $\sin\left(\dfrac{\pi}{6}+C\right)=1$即 $a=b,C=\dfrac{\pi}{3}$,亦即 $a=b=c$。
    \end{solution}

    \question 一个梯形的两底比为 $2:1$,两腰比为 $2:1$,两对角线比为 $2:1$。求短底与短腰的比。
    \ifprintanswers
    \begin{figure}[H]
        \centering
        \includegraphics[width=0.4\textwidth]{images/image147.png}
    \end{figure}
    \fi
    \begin{solution}
        设短腰为 $1$,短底为 $x$,高度为 $h$,与短腰相连的底边段长为 $y$,则
        \[
        1-y^{2}=h^{2}=4-(x-y)^{2}
        \]
        且对角线比给出
        \[
        4\bigl(h^{2}+(x+y)^{2}\bigr)=h^{2}+(2x-y)^{2}.
        \]
        解得
        \[
        x=\frac{\sqrt{10}}{2}.
        \]
        因此短底与短腰的比为
        \[
        \sqrt{10}:2
        \]
    \end{solution}

    \question 梯形 \(ABCD\) 中,\(AD\parallels BC,\;BC=3,\;CD=4\),且  
    \(\cos\angle ADC=\dfrac13,\;\angle ABC = 2\angle ADC\)。求 \(AC\)。  
    \begin{solution}
        \textcolor{red}{(待解)}
    \end{solution}
    
    \question 一匀速转动的摩天轮直径 $18$米,底端距地 $1$米。  
        康康在 \(P\) 处距地 $16$米,用 $4$ 秒到达顶端 \(T\),再用 $8$ 秒到达 \(Q\)。 求他抵达 \(Q\) 时离地面的距离。 
    \begin{figure}[H]
        \centering
        \includegraphics[width=0.3\textwidth]{images/image3.png}
    \end{figure}
    \begin{solution}
        设摩天轮圆心为$O$,\;$\angle POT=\theta$,则$\angle TOQ=2\theta$,发现到\[
        \cos \theta=\frac{6}{9}=\frac{2}{3}
        \]与此同时,作$QR$垂直于过圆心$O$的水平线于点$R$,有\[
            \sin (2\theta - 90^\circ)=\frac{QR}{9}
        \]又\[
            \sin (2\theta - 90^\circ)=-\cos2\theta =1-2\cos^2 \theta 
        \]则$QR=1$,故Q距地$(9-1)+1=9$米。
    \end{solution}
    
    \question 设圆心 \(O\)、半径 \(r\) 的圆的切线 \(AB\),点 \(C\) 在线 \(BD\) 上,且  
        \(AB=p,\;BC=CD=DO=q\)。 证明  
        \[
          p^{2}=q^{2}+r^{2}.
        \]
    \begin{figure}[H]
        \centering
        \includegraphics[width=0.25\textwidth]{images/image5.png}
    \end{figure}
    \begin{solution}
        在$\triangle ODC$中,由余弦定理
        \[
        r^2=q^2+q^2-2q^2\cos \angle ODC \Rightarrow \cos \angle ODC=\dfrac{2q^2-r^2}{2q^2}
        \]
        在$\triangle ODB$中,由余弦定理
        \[
        OB^2=q^2+(2q)^2-2 \cdot q \cdot 2q \cdot \dfrac{2q^2-r^2}{2q^2}=q^2+2r^2
        \]
        来到尾声,在$\triangle OAB$中,由毕氏定理
        \[
        OA^2+AB^2=OB^2 \Rightarrow r^2+p^2=q^2+2r^2
        \]
        于是得证 $p^{2}=q^{2}+r^{2}$。
    \end{solution}

    \question 一平行四边形$ABCD$面积为 $36$,对角线$AC,BD$长度为 $10,12$。求$AD$的长度。
    \begin{solution}
        设$M$在$BD$上使得$AM \perp BD$,于是
        \[
        \frac{1}{2}\cdot AM \cdot BD = 36 \Rightarrow AM = 3
        \]
        设对角线$AC,BD$交于点$O$,则 $AO = \dfrac{BD}{2} = 6$,在$\triangle AOM$ 中,由毕氏定理,
        \[
        OM = \sqrt{6^2 - 3^2} = 3\sqrt{3}
        \]
        在$\triangle AOM$ 中,由毕氏定理,
        \[
        AD = \sqrt{3^2 + (3\sqrt{3}+5)^2} = \sqrt{61 + 30\sqrt{3}}
        \]
    \end{solution}

    \question 在平行四边形 $ABCD$ 中,将边 $AD$ 延长至点 $F$,线段 $FB$ 与边 $CD$ 相交于 $G$,与对角线 $AC$ 相交于 $E$,且满足 $FG=4,GE=1$。求 $EB$ 的长度。
    \begin{figure}[H]
        \centering
        \includegraphics[width=0.4\textwidth]{images/image151.png}
    \end{figure}
    \begin{solution}
        设 $EB = x,DG = y,DC = b$,由$\triangle EBA \backsim \triangle EGC \; \text{(AAA)}$,
        \[
        \frac{x}{1} = \frac{b}{b-y}
        \]
        又由 $\triangle GFD \backsim \triangle GBC\; \text{(AAA)}$ 得
        \[
        \frac{4}{y} = \frac{1+x}{b-y} 
        \]
        解得
        \[
        EB =x= \sqrt{5}
        \]
    \end{solution}

    \question 平面上 $AC=AD,\angle ABC=90^\circ, \angle CAD=\alpha,\angle CBD=\beta, \angle CAB=\gamma$, 若 $\cos\alpha=\dfrac{4}{5}, \cos\beta=\dfrac{8}{17}$, 求 $\tan\gamma$。
    \ifprintanswers
    \begin{figure}[H]
        \centering
        \includegraphics[width=0.4\textwidth]{images/image91.jpg}
    \end{figure}
    \fi
    \begin{solution}
        在 $\triangle ABD$ 中, $\angle DAB = \gamma - \alpha, \angle ABD = 90^\circ + \beta,$则
        \[
        \angle ADB = 90^\circ + \alpha - \beta - \gamma
        \]
        由正弦定理,
        \[
        \frac{AB}{AD} = \frac{\sin(90^\circ+\alpha-\beta-\gamma)}{\sin(90^\circ+\beta)} = \frac{\cos(\alpha-\beta-\gamma)}{\cos\beta} = \frac{\cos(\alpha-\beta-\gamma)}{\frac{8}{17}}
        \]
        由于 $AD=AC=AB \cos \gamma$, 所以
        \begin{align*}
        \frac{8}{17}\cos \gamma 
        &= \cos(\alpha-\beta-\gamma)= \cos(\alpha-\beta)\cos\gamma + \sin(\alpha-\beta)\sin\gamma\\
        &= (\cos\alpha \cos\beta + \sin\alpha \sin\beta)\cos\gamma + (\sin\alpha \cos\beta - \sin\beta \cos\alpha)\sin\gamma\\
        &= \left(\frac{4}{5}\cdot \frac{8}{17} + \frac{3}{5}\cdot \frac{15}{17}\right)\cos\gamma + \left(\frac{3}{5}\cdot \frac{8}{17} - \frac{15}{17}\cdot \frac{4}{5}\right)\sin\gamma\\
        &= \frac{77}{85}\cos\gamma - \frac{36}{85}\sin\gamma
        \end{align*}
        得到
        \[
        \tan \gamma = \frac{37}{36}
        \]
    \end{solution}

    \question 四边形 $ABCD$ 的对角线 $AC,BD$ 交于 $E$ 点,且 $$BE=DE=\dfrac{1}{2}AD=3,\angle ADB=90^\circ,\angle ACD=45^\circ,$$求 $\triangle BCD$ 的面积。
    \ifprintanswers
    \begin{figure}[H]
        \centering        
        \includegraphics[width=0.5\textwidth]{images/image123.jpg}
    \end{figure}
    \fi
    \begin{solution}
        在直角 $\triangle ADE$ 中,
        \[
        AE^2=6^2+3^2=45 \Rightarrow AE=3\sqrt{5}
        \]
        在直角$\triangle ADB$ 中,
        \[
        AB^2=6^2+6^2=72 \Rightarrow AB=6\sqrt{2}
        \]
        在 $\triangle AEB$ 中,由余弦定理,
        \[
        \cos \angle AEB = \frac{45+9-72}{18\sqrt{5}} = -\frac{1}{\sqrt{5}} \Rightarrow \sin \angle AEB = \frac{2}{\sqrt{5}}
        \]
        在 $\triangle CDE$ 中,由正弦定理,
        \[
        \frac{3}{\frac{1}{\sqrt{2}}} = \frac{CD}{\frac{2}{\sqrt{5}}} \Rightarrow CD = 6\sqrt{\frac{2}{5}}
        \]
        又
        \[
        \sin \angle BDC = \sin(45^\circ+\angle DEC) = \frac{1}{\sqrt{2}}\cdot \left(-\frac{1}{\sqrt{5}}\right) + \frac{2}{\sqrt{5}}\cdot \frac{1}{\sqrt{2}} = \frac{1}{\sqrt{10}}
        \]
        因此$\triangle BCD$ 面积为
        \[
        \frac{1}{2}\cdot 6\sqrt{\frac{2}{5}}\cdot 6\cdot \frac{1}{\sqrt{10}} = \frac{18}{5}
        \]
    \end{solution}

    \question 有一正三角形 $\triangle ABC$ 的艺术品,边长为 6 公尺,以 $AB$ 为边斜靠在墙上,墙角为 $O$ 点,形成直角 $\triangle OAB,\angle AOB=90^\circ,A$ 点在地面上,$B$ 点在墙上。过 $B$ 点作与地面平行的直线交 $AC$ 于点 $D$,已知 $CD=2$ 公尺,求此艺术品的最高点离地面的高度。
    \ifprintanswers
    \begin{figure}[H]
        \centering        
        \includegraphics[width=0.5\textwidth]{images/image85.jpg}
    \end{figure}
    \fi
    \begin{solution}
        由余弦定理,
        \[
        BD^2=6^2+2^2-2 \cdot 2 \cdot 6 \cdot \frac{1}{2} \Rightarrow BD = 2\sqrt{7}
        \]
        在 $\triangle CBD$ 中,
        \[
        \frac{2}{\sin \angle CBD} = \frac{2\sqrt{7}}{\frac{\sqrt{3}}{2}} \Rightarrow \sin \angle CBD = \frac{\sqrt{3}}{2\sqrt{7}}
        \]
        于是
        \[
        CP = BC \sin \angle CBD = 6 \cdot \frac{\sqrt{3}}{2\sqrt{7}} = 3\sqrt{\frac{3}{7}}.
        \]
        又 $DP \parallels AQ$,由相似三角形,
        \[
        \frac{2}{6} = \frac{3\sqrt{\frac{3}{7}}}{CQ} \Rightarrow CQ = \frac{9}{7}\sqrt{21}\text{ 公尺}
        \]
    \end{solution}

    \question 某人在地面 $A$ 点,测得山峰的仰角为 $\theta$。此人向山脚前进 $110$ 公尺到达 $B$ 点,测得山峰仰角为 $2\theta$,再向山脚前进 $30$ 公尺到达 $C$ 点,又测得山峰仰角为 $90^\circ - \theta$。求山高 $h$。
    \ifprintanswers
    \begin{figure}[H]
        \centering        
        \includegraphics[width=0.6\textwidth]{images/image127.jpg}
    \end{figure}
    \fi
    \begin{solution}
        设从 $C$ 点到山脚的水平距离为 $a$ 公尺,则:
        \[
        h = a \tan (90^\circ - \theta) = \frac{a}{\tan \theta} \tag{1}
        \]
        \[
        h = (30+a) \tan 2\theta \tag{2}
        \]
        \[ 
        h = (140+a) \tan \theta \tag{3}
        \]
        由 $(1)=(3)$ 和 $(1)=(2)$,
        \[
        \frac{a}{140+a} = \tan^2 \theta \tag{4}
        \]
        \[
        \frac{a}{30+a} = \frac{2 \tan^2 \theta}{1 - \tan^2 \theta} = -2 + \frac{2}{1 - \tan^2 \theta} \tag{5}
        \]
        将 (4) 代入 (5),
        \[
        \frac{a}{30+a} = -2 + \frac{2}{1 - \frac{a}{140+a}} \Rightarrow a = 40
        \]
        于是
        \[
        \tan^2 \theta = \frac{40}{180} \Rightarrow \tan \theta = \frac{\sqrt{2}}{3}
        \]
        则山高为
        \[
        h = (140 + 40) \tan \theta = 180 \cdot \frac{\sqrt{2}}{3} = 60 \sqrt{2}
        \]
    \end{solution}

    \question 已知 \(\triangle ABC\) 满足 $2\angle BAC=3\angle ABC$,且 \(K\) 为 \(BC\) 上一点使得 \(\angle KAC = 2\,\angle KAB\),证明
    \begin{figure}[H]
        \centering
        \includegraphics[width=0.4\textwidth]{images/image4.png}
    \end{figure}
    \begin{parts}
    \part $AK=\dfrac{bc}{a},\;BK=\dfrac{a^{2}-b^{2}}{a}$
    \begin{solution}
        记$\angle KAB=\theta$,则$\angle KAC=2\theta,\angle ABC=2\theta,\angle CKA=\theta+2\theta=3\theta$,
        所以
        \[
        \triangle CAK \backsim \triangle CBA \ (\text{AAA}) \Rightarrow AK=\dfrac{bc}{a}
        \]
        由余弦定理,
        \[
        CK^2 =b^2+d^2-2bd \cos 2\theta 
        = b^2+\left(\dfrac{bc}{a}\right)^2-2b\left(\dfrac{bc}{a}\right) \left(\dfrac{a^2+c^2-b^2}{2ac}\right)
        = \dfrac{b^4}{a^2}
        \]
        于是 $CK= \dfrac{b^2}{a}$,且 
        \[
        BK=a-CK=\dfrac{a^{2}-b^{2}}{a}
        \]
    \end{solution}
    \part $(a^{2}-b^{2})(a^{2}-b^{2}+ac)=b^{2}c^{2}$
    \begin{solution}
        记$AK=x,BK=y$,由正弦定理,\[
        \frac{x}{\sin 2\theta}=\frac{y}{\sin \theta} \Rightarrow x=2y \cos \theta
        \]
        又由余弦定理,
        \[
        \cos \theta=\dfrac{c^2+x^2-y^2}{2cx}
        \]
        于是
        \[
        cx^2=y(x^2+c^2-y^2) \Rightarrow x^2=y(y+c)
        \]
        代入 $x=\dfrac{bc}{a}, y=\dfrac{a^{2}-b^{2}}{a}$可得\[
        (a^{2}-b^{2})(a^{2}-b^{2}+ac)=b^{2}c^{2}
        \]
    \end{solution}
    \end{parts}

    \question 已知正方形 $DEFG$ 内接于 $\triangle ABC,EF$ 在 $BC$ 上且 $D,G$ 分别在 $AB,AC$ 上,已知 $\triangle ADG,\triangle BDE,\triangle CGF$ 的面积分别为 $1,3,1$,求 $\triangle ABC$ 的面积。 
    \ifprintanswers
    \begin{figure}[H]
        \centering
        \includegraphics[width=0.5\textwidth]{images/image93.jpg}
    \end{figure}
    \fi
    \begin{solution}
        作 $AQ \perp BC$交 $DG$ 于 $P$,因此$\triangle APG \backsim \triangle GFC \ (\text{AAA}),\triangle APD \backsim \triangle DEB \ (\text{AAA})$,且有
        \[
        \frac{[\triangle APG]}{[\triangle APD]} = \frac{[\triangle GFC]}{[\triangle DEB]} = \frac{1}{3}
        \Rightarrow [\triangle APD]=\frac{3}{4},\quad [\triangle APG]=\frac{1}{4}
        \]
        又
        \[
        \frac{[\triangle BDE]}{[\triangle APD]}=\frac{3}{\frac{3}{4}} = \frac{DE^2}{AP^2}\Rightarrow DE = 2AP
        \]
        因此
        \[
        [\triangle ADG] = 1 = \frac12 \cdot DG \cdot AP = \frac12 \cdot 2(AP)^2 \Rightarrow AP = 1 , DG = 2 
        \]
        $\triangle ABC$面积为
        \[
        1+3+1+2^2 = 9
        \]
    \end{solution}

    \question 设 $P$ 为正 $\triangle ABC$ 内部一点,若 $P$ 点依序到三边 $BC,AC,AB$ 之距离比为 $1:3:2$,求 $PA^2 : PB^2 : PC^2$。
    \ifprintanswers
    \begin{figure}[H]
        \centering
        \includegraphics[width=0.4\textwidth]{images/image79.jpg}
    \end{figure}
    \fi
    \begin{solution}
        作
        \[
        MN \parallels AB, \quad ST \parallels BC, \quad UV \parallels AC,
        \]
        则 $\triangle PUM, \triangle PSV, \triangle PTN$ 均为正三角形,且 $1, 2, 3$ 分别为正三角形的高,因此
        \[
        PU = PM = UM = \frac{2}{\sqrt{3}}, \quad
        PV = PS = VS = \frac{4}{\sqrt{3}}, \quad
        PN = PT = NT = \frac{6}{\sqrt{3}}.
        \]
        由毕氏定理,
        \[
        PA^2 = 2^2 + \left(\frac{6}{\sqrt{3}} + \frac{2}{\sqrt{3}}\right)^2 = \frac{76}{3}, \quad
        PB^2 = 1^2 + \left(\frac{4}{\sqrt{3}} + \frac{1}{\sqrt{3}}\right)^2 = \frac{28}{3},
        \]
        \[
        PC^2 = 3^2 + \left(\frac{2}{\sqrt{3}} + \frac{3}{\sqrt{3}}\right)^2 = \frac{52}{3}.
        \]
        故
        \[
        PA^2 : PB^2 : PC^2 = 19 : 7 : 13
        \]
    \end{solution}

\question 圆内接四边形 $ABCD$ 满足 $AB=4,BC=1,CD=2,DA=3$。设 $P,Q$ 分别为 $BC$ 与 $DA$ 的中点,求 $PQ^2$。
\ifprintanswers
\begin{figure}[H]
    \centering
    \includegraphics[width=0.4\textwidth]{images/image157.png}
\end{figure}
\fi
\begin{solution}
    设对角线交点为 $R$,令 $CR=x$。由$\triangle BCR \backsim \triangle ADR \; \text{(AAA)}$ 及 $\triangle CDR \backsim \triangle BAR \; \text{(AAA)}$可得
    \[
    BR=2x,\quad DR=3x,\quad AR=6x,
    \]
    因此
    \[
    AC=7x,\quad BD=5x
    \]
    在四边形 $ABCD$中,由托勒密定理,
    \[
    AC\cdot BD=AB\cdot CD+BC\cdot DA 
    \]
    得
    \[
    x^2=\frac{11}{35}
    \]
    在$\triangle ABD$及$\triangle ACD$中,由中线定理,
    \[
    AB^2+BD^2=\frac{1}{2}DA^2+2BQ^2, \quad AC^2+CD^2=\frac{1}{2}DA^2 + 2CQ^2
    \]
    得
    \[
    \textcolor{red}{BQ^2=\frac{271}{28},\;CQ^2=\frac{571}{20}}
    \]
    故在$\triangle BCQ$中,由中线定理,
    \[
    BQ^2+CQ^2 = \frac{1}{2}BC^2+2PQ^2 \Rightarrow PQ^2=\frac{291}{35}
    \]
    \textcolor{red}{(待验证)}
\end{solution}

    \question 已知三角形 $ABC$ 和点 $X_1$ 在 $AC$ 上,满足 $AB=AC=91,BC=70,X_1C=65$。设 $X_0 = B$,点 $X_2, X_4, X_6, \ldots$ 在 $AB$ 上,$X_3, X_5, X_7, \ldots$ 在 $AC$ 上,且 $X_{n+1}X_n \perp X_nX_{n-1}$ 对 $n \ge 1$ 成立。求 $\displaystyle \sum_{n=1}^{\infty} X_{n-1}X_n$。
    \ifprintanswers
    \begin{figure}[H]
        \centering
        \includegraphics[width=0.4\textwidth]{images/image153.png}
    \end{figure}
    \fi
    \begin{solution}
        设 $M,N$ 分别为 $A,X_1$ 在 $BC$ 上的垂足,由$\triangle AMC \backsim \triangle X_1NC \; \text{(AAA)}$
        \[
        \frac{NC}{MC} = \frac{X_1C}{AC} \Rightarrow NC = \frac{65}{91} \cdot 35 = 25.
        \]
        由毕氏定理,
        \[
        X_1N = \sqrt{65^2 - 25^2} = 60, \quad X_0X_1 = \sqrt{60^2 + 45^2} = 75.
        \]
        则
        \[
        \tan \angle X_1BA = \tan(\angle ABC - \angle X_1BC) = \frac{\frac{12}{5} - \frac{4}{3}}{1 + \frac{12}{5} \cdot \frac{4}{3}} = \frac{16}{63}
        \]
        因此
        \[
        X_1X_2 = X_0X_1 \cdot \frac{16}{63} = \frac{400}{21}, \quad X_0X_2 = \sqrt{75^2 + \left(\frac{400}{21}\right)^2} = 25 \cdot \frac{65}{21}
        \]
        且$AX_n,AX_{n+2},\dots$成等比数列,公比为
        \[
        r = \frac{AX_2}{AX_0} = \frac{91 - \frac{25 \cdot 65}{21}}{91} = 1 - \frac{125}{147}
        \]
        因此无穷和为
        \[
        \sum_{n=1}^{\infty} X_{n-1}X_n = \frac{X_0X_1 + X_1X_2}{1-r} = \left(75 + \frac{400}{21}\right) \cdot \frac{147}{125} = \frac{553}{5}
        \]
    \end{solution}

    \question 一个等边三角形中有 $n$ 排全等的小圆（如下图 $n=4$ 所示）。求当 $n\to\infty$ 时,圆的面积与三角形面积的比值极限。
    \ifprintanswers
    \begin{figure}[H]
        \centering
        \includegraphics[width=0.4\textwidth]{images/image143.png}
    \end{figure}
    \fi
    \begin{solution}
        设三角形边长为 $1$,小圆半径为 $r$,由 $30$-$60$-$90$ 直角三角形可得
        \[
        2(n-1)r + 2 r\sqrt{3} = 1 \Rightarrow r = \frac{1}{2(n+\sqrt{3}-1)}.
        \]
        共有 $\frac{n(n+1)}{2}$ 个圆,因此圆的总面积为
        \[
        \frac{n(n+1)}{2} \cdot \frac{\pi}{4(n+\sqrt{3}-1)^2}= \frac{\pi}{8} \frac{n(n+1)}{(n+\sqrt{3}-1)^2}.
        \]
        当 $n\to\infty$ 时,$\dfrac{n(n+1)}{(n+\sqrt{3}-1)^2}\to 1$,因此面积比极限为
        \[
        \frac{\dfrac{\pi}{8}}{\dfrac{\sqrt{3}}{4}} = \frac{\pi\sqrt{3}}{6}
        \]
    \end{solution}

    \question 已知 $C_1$ 是半径为 1 的圆,$P_1$ 是 $C_1$ 的外接正方形。对于 $n \ge 2$,令 $C_n$ 为 $P_{n-1}$ 的外接圆,$P_n$ 为 $C_n$ 的外接正 $2^{n+1}$ 边形。设 $r_n$ 为 $C_n$ 的半径。求 $\lim_{n \to \infty} r_n$。
    \ifprintanswers
    \begin{figure}[H]
        \centering
        \includegraphics[width=0.4\textwidth]{images/image155.png}
    \end{figure}
    \fi
    \begin{solution}
        由图可知,
        \[
        r_2 = \frac{r_1}{\cos\frac{\pi}{4}}, \quad r_n = \frac{r_{n-1}}{\cos\frac{\pi}{2^n}}, \quad n\ge 2
        \]
        于是
        \[
        r_n = \frac{1}{\displaystyle \prod_{k=2}^{n} \cos \frac{\pi}{2^k}}
        \]
        且由
        \[
        \sin x = 2 \sin \frac{x}{2} \cos \frac{x}{2} = 2^2 \sin \frac{x}{4} \cos \frac{x}{4} \cos \frac{x}{2} = \cdots = 2^n \sin \frac{x}{2^n} \prod_{k=1}^{n} \cos \frac{x}{2^k}
        \]
        两边除以 $x$ 并取极限 $n\to\infty$ 得
        \[
        \frac{\sin x}{x} = \prod_{k=1}^{\infty} \cos \frac{x}{2^k}
        \]
        令 $x = \dfrac{\pi}{2}$,得到
        \[
        \prod_{k=1}^{\infty} \cos \frac{\pi}{2^{k+1}} = \frac{2}{\pi}
        \]
        因此
        \[
        \lim_{n\to\infty} r_n = \frac{\pi}{2}
        \]
    \end{solution}

    \question 圆内有两条平行弦,其长分别为 12,16,且它们之间的距离为 7。另一条弦位于两条弦中间且与两线等距,求该弦的长度。
    \begin{solution}
        设圆半径为 $r$,两条弦到圆心的距离分别为 $d_1,d_2$,由毕氏定理,
        \[
        \left(\frac{12}{2}\right)^2 + d_1^2 = r^2, \quad \left(\frac{16}{2}\right)^2 + d_2^2 = r^2 \implies 36 + d_1^2 = 64 + d_2^2.
        \]
        假设两弦在圆心的两侧,则
        \[
        d_1 + d_2 = 7, \quad d_1 - d_2 = 4 \implies d_1 = \frac{11}{2}>0, \quad d_2 = \frac{3}{2}>0
        \]
        于是
        \[
        r^2 = 36 + d_1^2 = \frac{265}{4}
        \]
        中间弦到圆心的距离为
        \[
        d = \frac{d_1 - d_2}{2} = 2
        \]
        设中间弦长度为 $x$,则由毕氏定理,
        \[
        \left(\frac{x}{2}\right)^2 + 2^2 = r^2 \Rightarrow x = \sqrt{249}
        \]
    \end{solution}

    \question 在图中,两圆在点 $P$ 相切。$QP$ 和 $SU$ 是大圆的垂直直径,交于 $O$。点 $V$ 在 $QP$ 上使得$VP$ 是小圆的直径。小圆与 $SU$ 相交于 $T$。已知 $QV = 9,ST = 5$,求两圆直径之和。
    \begin{figure}[H]
        \centering
        \includegraphics[width=0.4\textwidth]{images/image173.png}
    \end{figure}
    \begin{solution}
        设大圆直径为 $D$,小圆直径为 $d$,小圆圆心为 $C$,在$\triangle TOC$中,由毕氏定理,
        \[
        TO^2 + OC^2 = CT^2 \Rightarrow \left(\frac{D}{2} - 5\right)^2 + \left(\frac{D}{2} - \frac{d}{2}\right)^2 = \left(\frac{d}{2}\right)^2
        \]
        由 $QV=D-d = 9$,得
        \[
        (D-10)^2 + 9^2 = d^2 
        \]
        于是
        \[
        81 = (d-(D-10))(d+(D-10)) = 1\cdot(d+D-10)
        \]
        所以
        \[
        d + D = 91
        \]
    \end{solution}

    \question 如图所示,三个圆的圆心分别为 $X,Y,Z$,每个圆都与另外两个圆相切。圆心为 $X$ 的圆与矩形 $PQRS$ 相切于三条边,圆心为 $Z$ 的圆与矩形 $PQRS$ 相切于两条边。如果 $XY=30,YZ=20,XZ=40$,求矩形 $PQRS$ 的面积。  
    \begin{figure}[H]
        \centering
        \includegraphics[width=0.4\textwidth]{images/image165.png}
    \end{figure}
    \ifprintanswers
    \begin{figure}[H]
        \centering
        \includegraphics[width=0.4\textwidth]{images/image166.png}
    \end{figure}
    \fi
    \begin{solution}
        圆心之间的距离等于两个相切圆半径之和。设圆心 $X,Y,Z$ 的半径分别为 $x,y,z$,则有:
        \[
        XY = x + y = 30, \quad XZ = x + z = 40, \quad YZ = y + z = 20.
        \]
        解得
        \[
        x=25,\;y=5,\;z=15
        \]
        矩形高度为$2x = 50$,宽度为
        \[
        SR = SU + UV + VR = 25 + \sqrt{40^2 - 10^2} + 15 = 40 + 10\sqrt{15}
        \]
        故矩形面积为
        \[
        [PQRS]=50(40 + 10\sqrt{15}) = 2000 + 500\sqrt{15}
        \]
    \end{solution}

    \question 两个圆 $C_1$ 和 $C_2$ 外切,直线 $l$ 为它们的公切线。直线 $m$ 平行于 $l$ 并与 $C_1$ 和 $C_3$ 相切。三个圆互相切。已知 $C_2$ 的半径为 $9$,$C_3$ 的半径为 $4$,求 $C_1$ 的半径。
    \begin{figure}[H]
        \centering
        \includegraphics[width=0.4\textwidth]{images/image201.png}
    \end{figure}
    \ifprintanswers
    \begin{figure}[H]
        \centering
        \includegraphics[width=0.4\textwidth]{images/image202.png}
    \end{figure}
    \fi
    \begin{solution}
        设 $C_1$ 的半径为 $r$,有
        \[
        C_1C_2 = r+9, \quad C_1C_3 = r+4, \quad C_2C_3 = 13.
        \]
        如图所示,考虑矩形$ABC_2D$,由毕氏定理,在 $\triangle AC_3C_1$ 中,
        \[
        C_3A^2 = (r+4)^2-(r-4)^2 = 16r \Rightarrow C_3A = 4\sqrt{r}.
        \]
        在 $\triangle DC_2C_1$ 中,
        \[
        DC_2^2 = (r+9)^2-(r-9)^2 = 36r \Rightarrow DC_2 = 6\sqrt{r}
        \]
        在 $\triangle BC_3C_2$ 中,
        \[
        CB^2 = 13^2-(2r-13)^2 = -4r^2 + 52r \Rightarrow C_3B = \sqrt{-4r^2 + 52r}
        \]
        由$DC_2 = C_3A + C_3B$得
        \[
        6\sqrt{r} = 4\sqrt{r} + \sqrt{-4r^2+52r} 
        \]
        解得
        \[
        r=12>0
        \]
    \end{solution}

    \question 已知$\triangle ABC$中,$\;AB=13,BC=14,AC=15$,设 $D,E$ 分别为从 $A,B$ 作的高的垂足,求 $\triangle CDE$ 的外接圆直径。
    \begin{solution}
        在$\triangle ABC$中,由余弦定理,
        \[
        13^2 = 14^2 + 15^2 - 2\cdot 14 \cdot 15 \cos C \Rightarrow \cos C = \frac{3}{5}
        \]
        设 $AD$ 与 $BE$ 交于 $X$,由于 $\angle CDX = \angle CEX = 90^\circ$,有 
        \[
        \triangle CDA \backsim \triangle CEB \backsim \triangle XDB \ \text{(AAA)},
        \]
        得到
        \[
        CD = 9, \quad BD = 5, \quad DX = \frac{15}{4}
        \]
        在$\triangle CDX$中,由毕氏定理,
        \[
        CX = \sqrt{CD^2+DX^2}=\frac{39}{4}
        \]
        发现 $CDXE$ 是一圆内接四边形,因此 $\triangle CDE$ 的外接圆直径即为
        \[
        CX=\frac{39}{4}
        \]
    \end{solution}

    \question 已知点 $A,B,C,D$共圆, $AC$ 是直径,且 $\angle CBD=\angle DBA$。若 $BC=2, AB=4$,求 $BD$ 的长度。
    \begin{solution}
        $AC$ 是直径,所以 $\angle ABC=90^\circ$,因此 
        \[
        \angle CBD=\angle DBA=45^\circ
        \]
        又 $\angle CAD$ 与 $\angle CBD$ 为同弧所对的角,所以 
        \[
        \angle CAD=\angle CBD=45^\circ
        \]
        在$\triangle ABC$中,由毕氏定理,
        \[
        AC=\sqrt{AB^2+BC^2}=2\sqrt{5}
        \]
        故
        \[
        \sin \angle BAC=\frac{BC}{AC}=\frac{1}{\sqrt{5}}
        \]
        在$\triangle BAD$中,由正弦定理,
        \begin{align*}
        BD=AC \sin \angle BAD &= AC \sin \left(\angle BAC+\frac{\pi}{4}\right) \\
        &= 2\sqrt{5}\cdot \frac{\sqrt{2}}{2}(\sin\angle BAC+\cos\angle BAC)=3\sqrt{2}            
        \end{align*}
    \end{solution}

    \question 在 $\triangle ABC$ 中,已知 $AB=9,BC=10,CA=11$,且内切圆切 $BC,AB,AC$ 于 $D,E,F$,求 $AD$。
    \ifprintanswers
    \begin{figure}[H]
        \centering
        \includegraphics[width=0.4\textwidth]{images/image124.jpg}
    \end{figure}
    \fi
    \begin{solution}
        设$AE=AF=a, BD=BE=b, CF=CD=c,$则
        \[
        c = s - AB = \frac{9+10+11}{2} - 9 = 6
        \]
        在 $\triangle CAD$ 与 $\triangle CAB$中,由余弦定理,
        \[
        \cos C 
        = \frac{11^2 + 6^2 - AD^2}{2\cdot 11 \cdot 6}
        = \frac{11^2 + 10^2 - 9^2}{2\cdot 11 \cdot 10}
        \]
        解得
        \[
        AD = \sqrt{73}
        \]
    \end{solution}

    \question 在 $\triangle ABC$ 中,$AB=10$, $M$ 为 $AB$ 中点,$\;\triangle ABC$ 内切圆恰将线段 $CM$ 三等分,求 $\triangle ABC$ 面积。
    \ifprintanswers
    \begin{figure}[H]
        \centering
        \includegraphics[width=0.6\textwidth]{images/image110.jpg}
    \end{figure}
    \fi
    \begin{solution}
        设三切点为 $D,E,F$,中线 $CM$ 与内切圆交于 $P,Q$ 两点,据题意有
        \[
        CP = PQ = QA = a.
        \]
        由$\triangle MQD \backsim \triangle MDP \ (\text{AAA})$,
        \[
        \frac{MQ}{MD} = \frac{MD}{MP} \Rightarrow  MD^2 = MP \cdot MQ = 2a^2,
        \]
        同理由$\triangle CPF \backsim \triangle CFQ \ (\text{AAA})$得$CF^2 = 2a^2$,则
        \[
        CE = CF = MD = b  \Rightarrow b = \sqrt{2}a.
        \]
        由中线定理,
        \[
        CA^2 + CB^2 = 2(CM^2 + AM^2) \Rightarrow  (5+2b)^2 + 5^2 = 2(9a^2 + 5^2)
        \]
        代入 $b = \sqrt{2}a$ 解得
        \[
        a = 2\sqrt{2}, \quad b = 4.
        \]
        因此三边长为$AB = 10, BC = 5,  AC = 5 + 2b = 13$,半周长 $s = \dfrac{10+5+13}{2} = 14$,面积为
        \[
        \sqrt{14 \cdot 4 \cdot 9 \cdot 1} = 6\sqrt{14}
        \]
    \end{solution}

    \question 已知$PQRS$ 是边长为 $4$ 的正方形,点 $T$ 在 $QR$ 上,点 $U$ 在 $RS$ 上,且 $\angle UPT = 45^\circ$,求 $\triangle RUT$ 的最大周长。
    \ifprintanswers
    \begin{figure}[H]
        \centering
        \includegraphics[width=0.3\textwidth]{images/image200.png}
    \end{figure}
    \fi
    \begin{solution}
        将 $\triangle PSU$ 绕点 $P$ 逆时针旋转 $90^\circ$,得到全等的 $\triangle PQV$,此时 $V$ 位于 $RQ$ 的延长线上,且有
        \[
        \triangle PTU \congr \triangle PTV \ \text{(SAS)}
        \]
        于是$\triangle RUT$ 的周长
        \[
        UR + RT + UT = UR + RT + TV = UR + RQ + SU = SR + RQ = 8
        \]
        故最大周长为 $8$。
    \end{solution}
    \begin{solution}
        设 $\angle SPU = \theta,0^\circ < \theta < 45^\circ$,则
        \[
        SU = 4 \tan \theta, \quad UR = SR - SU = 4 - 4\tan \theta.
        \]
        由 $\angle UPT = 45^\circ$ 可得 $\angle QPT = 45^\circ - \theta$,于是
        \[
        QT = PQ \tan(45^\circ - \theta) = 4 \cdot \frac{1 - \tan \theta}{1 + \tan \theta}, \quad RT = QR - QT = \frac{8 \tan \theta}{1 + \tan \theta}.
        \]
        由毕氏定理,
        \[
        UT = \sqrt{UR^2 + RT^2} = 4 \cdot \frac{1 + \tan^2 \theta}{1 + \tan \theta}
        \]
        因此,$\triangle RUT$ 的周长为
        \[
        UR + RT + UT = 4 - 4 \tan \theta + \frac{8 \tan \theta}{1 + \tan \theta} + 4 \cdot \frac{1 + \tan^2 \theta}{1 + \tan \theta} 
        = 4 \frac{2 + 2 \tan \theta}{1 + \tan \theta} = 8
        \]
        无论 $\theta$ 取何值周长恒为 $8$,故最大可能为 $8$。
    \end{solution}

    \question 如图,$PRTY$ 和 $WRSU$ 是正方形,点 $Q,X$ 在 $PR,TY$ 上使得 $PQXY$ 是矩形,点 $T,W$ 在 $SU,QX$ 上,$UW$ 与 $TY$ 交于$V$。若矩形 $PQXY$ 面积为 $30$,求 $ST$ 的长度。
    \begin{figure}[H]
        \centering        
        \includegraphics[width=0.35\textwidth]{images/image175.png}
    \end{figure}
    \begin{solution}
        设 $ST = a,\angle STR = \theta$。则
        \[
        \angle TRS = 90^\circ - \theta ,\quad \angle QWR = \theta
        \]
        由 $\triangle RST$,
        \[
        PY = PR = RT = \frac{a}{\cos \theta}, \quad RW = RS = a \tan \theta
        \]
        且由 $\triangle QRW$,
        \[
        QR = RW \sin \theta = a \tan \theta \sin \theta, \quad
        PQ = PR - QR = \frac{a}{\cos \theta} - a \tan \theta \sin \theta = a \cos \theta
        \]
        矩形面积为 $PQ \cdot PY = 30$,即
        \[
        a \cos \theta \cdot \frac{a}{\cos \theta} = a^2 = 30 \Rightarrow a = \sqrt{30} 
        \]
    \end{solution}

    \question 如图,在直角 $\triangle ABC$ 中,$\angle ACB = 90^\circ, AC=6$,以 $AB$ 为一边向三角形外作正方形 $ABEF$。已知正方形的中心为 $O$,且 $OC = 8\sqrt{2}$,求 $BC$ 的长度。
    \begin{figure}[H]
        \centering        
        \includegraphics[width=0.35\textwidth]{images/image134.png}
    \end{figure}
    \ifprintanswers
    \begin{figure}[H]
        \centering        
        \includegraphics[width=0.5\textwidth]{images/image135.jpg}
    \end{figure}
    \fi
    \begin{solution}
        $\triangle OAB$ 为等腰直角三角形,设$G$ 为 $AB$ 中点,则 $OG$ 为中垂线,设$$GA = GB = OG = a,$$同弧所对的圆周角相等,有$\angle ACO = \angle ABO = 45^\circ$,由余弦定理,
        \[
        (\sqrt{2} a)^2=6^2 + (8\sqrt{2})^2-2 \cdot 6 \cdot 8\sqrt{2}\cdot \cos \angle 45^\circ
        \Rightarrow a^2 = 34
        \]
        在直角 $\triangle ABC$ 中,由毕氏定理,
        \[
        (2a)^2 = 6^2 + CB^2 \Rightarrow CB = 10
        \]
    \end{solution}

    \question 在直角$\triangle ABC$中,$D,E,F$分别为三边的中点,$G$为重心,且
    \[
    GD^2+GE^2+GF^2=\frac{50}{3}
    \]
    求斜边$AC$的长度。
    \begin{solution}
        由于$G$为重心,
        \[
        \frac{50}{3}=GD^2+GE^2+GF^2
        =\left(\frac13 BD\right)^2+\left(\frac13 CE\right)^2+\left(\frac13 AF\right)^2
        \]
        从而
        \[
        BD^2+CE^2+AF^2=150
        \]
        在直角$\triangle ABC$中,由毕氏定理,
        \[
        BD=\frac12 AC, \quad CE^2=BC^2+\left(\frac12 AB\right)^2, \quad AF^2=AB^2+\left(\frac12 BC\right)^2
        \]
        于是
        \[
        BD^2+CE^2+AF^2
        =\frac14 AC^2+\frac14 AB^2+BC^2+AB^2+\frac14 BC^2
        \]
        又由毕氏定理$AB^2+BC^2=AC^2$,
        \[
        150=BD^2+CE^2+AF^2=\frac14 AC^2+\frac14 AC^2+AC^2=\frac32 AC^2 \Rightarrow AC=10
        \]
    \end{solution}
    \begin{solution}
        由重心的性质可得
        \[
        GA^2+GB^2+GC^2=\frac13(AB^2+BC^2+AC^2)
        \]
        又$D,E,F$分别为三边的中点,$G$为重心,有
        \[
        GD^2+GE^2+GF^2=\frac14(GA^2+GB^2+GC^2)=\frac1{12}(AB^2+BC^2+AC^2)
        \]
        且由毕氏定理$AB^2+BC^2=AC^2$,因此
        \[
        GD^2+GE^2+GF^2=\frac16 AC^2=\frac{50}{3}
        \]
        解得
        \[
        AC=10
        \]
    \end{solution}

    \question 已知 $\triangle ABC$ 满足 $AB=9,BC=10,CA=17,B'$ 为 $B$ 关于 $AC$ 的对称点,$G,G'$为 $\triangle ABC,\triangle AB'C$的重心,求 $GG'$ 的长度。
    \ifprintanswers
    \begin{figure}[H]
        \centering        
        \includegraphics[width=0.5\textwidth]{images/image158.png}
    \end{figure}
    \fi
    \begin{solution}
        设 $M$ 为 $AC$ 的中点,由$GMG' \backsim BMB' \ \text{(SAS)}$,可得
        \[
        GG' = \frac{1}{3}BB'
        \]
        三角形面积为
        \[
        [\triangle ABC] = \frac{1}{4} BB' \cdot AC =\sqrt{18(18-9)(18-10)(18-17)} = 36 \Rightarrow BB' = \frac{144}{17}
        \]
        于是
        \[
        GG' = \frac{48}{17}
        \]
    \end{solution}

    \question 已知 $\triangle ABC$ 满足 $AB=4,BC=5, CA=6,\triangle ABC$ 的外接圆同时是 $\triangle A'B'C'$ 的内切圆,且 $A,B,C$ 分别在 $B'C',A'C',A'B'$ 上,求 $B'C'$ 的长度。
    \begin{figure}[H]
    \centering        
    \includegraphics[width=0.5\textwidth]{images/image159.png}
    \end{figure}
    \begin{solution}
        由弦切角定理,
        \[
        \angle C'AB = \angle C, \quad \angle B'AC = \angle B
        \]
        由直角三角形关系得
        \[
        C'A = \frac{2}{\cos\angle C'AB} = \frac{2}{\cos C}, \quad B'A = \frac{3}{\cos\angle B'AC} = \frac{3}{\cos B}
        \]
        由余弦定理,得
        \[
        \cos B = \frac{4^2+5^2-6^2}{2\cdot 4\cdot 5}=\frac{1}{8}, \quad \cos C = \frac{5^2+6+2-4^2}{2 \cdot 5 \cdot 6}=\frac{3}{4}
        \] 
        因此
        \[
        B'C' = C'A + B'A = 2\cdot \frac{4}{3} + 3\cdot 8 = \frac{80}{3}
        \]
    \end{solution}

    \question 在 $\triangle ABC$ 中,点 $D, E$ 分别在 $AB, AC$ 上,且 $AD=3, DB=1, AE=2, EC=4$。$BE$ 和 $CD$ 相交于 $P$ 点,若 $AP \perp BC$,求 $\cos \angle BAC$。
    \ifprintanswers
    \begin{figure}[H]
        \centering        
        \includegraphics[width=0.4\textwidth]{images/image128.jpg}
    \end{figure}
    \fi
    \begin{solution}
        设 $Q$ 为 $BC$ 上的垂足,由梅涅劳斯定理,
        \[
        \frac{CE}{EA} \cdot \frac{AD}{DB} \cdot \frac{BQ}{QC}= \frac{4}{2} \cdot \frac{3}{1} \cdot \frac{BQ}{QC} = 1 \Rightarrow \frac{BQ}{QC} = \frac{1}{6}
        \]
        设 $BQ = a,QC = 6a$,在直角$\triangle AQB$ 及 $\triangle AQC$中,由毕氏定理,
        \[
        QA^2 = 16 - a^2 = 36 - (6a)^2 \Rightarrow a^2 = \frac{4}{7}
        \]
        在$\triangle ABC$ 中,由余弦定理,
        \[
        \cos \angle BAC = \frac{4^2 + 6^2 - (7a)^2}{2 \cdot 4 \cdot 6} = \frac{1}{2}
        \]
    \end{solution}

    \question 如下图,三角形 $ABC$ 中,三线段 $AD,BE,CF$ 交于一点 $O$,若 $OD=OE=OF=4$ 且 $OA+OB+OC=37$,求 $OA \cdot OB \cdot OC$ 的值。 
    \begin{figure}[H]
        \centering        
        \includegraphics[width=0.4\textwidth]{images/image82.png}
    \end{figure}
    \begin{solution}
        由等高性质,
        \[
        \frac{S_{\triangle OBC}}{[\triangle ABC]}+ \frac{S_{\triangle OAC}}{[\triangle ABC]}+ \frac{S_{\triangle OAB}}{[\triangle ABC]}= \frac{OD}{AD} +\frac{OE}{BE} +\frac{OF}{CF} =1
        \]
        即
        \[
        \frac{4}{a}+\frac{4}{b}+ \frac{4}{c}=1
        \]
        其中 $a=AD=OA+4$, $b=BE=OB+4$, $c=CF=OC+4$,令
        \[
        ab+bc+ca=t,abc=4(ab+bc+ca)=4t,a+b+c=OA+OB+OC+12=49
        \]
        则设三次多项式 $f(x)$三根为 $a,b,c$,则
        \[
        f(x)=(x-a)(x-b)(x-c)=x^3-49x^2+tx-4t
        \]
        令 $x=4$ 得
        \[
        -(OA \cdot OB \cdot OC)=(4-a)(4-b)(4-c)=64-784+4t-4t=-720
        \]
        即
        \[
        OA \cdot OB \cdot OC=720
        \]
    \end{solution}

    \question 已知 $\triangle ABC$ 为等边三角形,$D,E,F$在 $BC,CA,AB$上使得$AF=BD=CE=\dfrac{1}{3}AB$。若$AD,BE,CF$分别交$CF,AD,BE$于$G,H,I$,求 
    \[
    \frac{[\triangle GHI]}{[\triangle ABC]}
    \]
    \ifprintanswers
    \begin{figure}[H]
        \centering        
        \includegraphics[width=0.4\textwidth]{images/image142.png}
    \end{figure}
    \fi
    \begin{solution}
        设
        \[
        \alpha = [\triangle AFG]=[\triangle BDH]=[\triangle CEI], \beta = [AEIG]=[BHGF]=[CIHD], \gamma = [\triangle GHI]
        \]
        由等高性质,得
        \[
        2(2\alpha+\beta) = 2\beta + \alpha + \gamma \Rightarrow \gamma = 3\alpha
        \]
        在$\triangle ACF$中,由余弦定理,
        \[
        CF^2 = AF^2 + (3AF)^2 - 2 \cdot AF \cdot (3AF)\cos 60^\circ = 7 AF^2.
        \]
        由于 $\triangle ACF \backsim \triangle ICE \ \text{(AAA)}$,
        \[
        \frac{[\triangle ACF]}{[\triangle ICE]}=\frac{2\alpha + \beta}{\alpha} = \frac{CF^2}{AF^2} = 7 \Rightarrow \beta = 5\alpha
        \]
        因此
        \[
        \frac{[\triangle GHI]}{[\triangle ABC]} = \frac{\gamma}{\gamma + 3\alpha + 3\beta} = \frac{1}{7}
        \]
    \end{solution}

    \question $\triangle ABC$ 中, $AB=6, BC=8, CA=12, D, E$ 在 $CA$ 上, 且 $AD = DE = EC,$ 求 $BD,BE$的长。
    \ifprintanswers
    \begin{figure}[H]
        \centering        
        \includegraphics[width=0.55\textwidth]{images/image89.jpg}
    \end{figure}
    \fi
    \begin{solution}
        令 $BD = a, BE = b,$由余弦定理,
        \[
        \cos \angle BDA = \frac{a^2 + 16 - 36}{8a}, \quad \cos \angle BDE = \frac{a^2 + 16 - b^2}{8a}
        \]
        由 $\cos \angle BDA = -\cos \angle BDE$得
        \[
        2a^2 - b^2 = 4 \tag{1}
        \]
        同理,由余弦定理,
        \[
        \cos \angle BED = \frac{b^2 + 16 - a^2}{8b}, \quad \cos \angle BEC = \frac{b^2 + 16 - 64}{8b}
        \]
        由 $\cos \angle BED = -\cos \angle BEC$可得
        \[
        a^2 - 2b^2 = -32 \tag{2}
        \]
        联立解得 
        \[
        a = \frac{2\sqrt{30}}{3},b=\frac{2\sqrt{17}}{3}
        \]
    \end{solution}

    \question 在 $\triangle ABC$ 中,$D,E$ 在 $AB$ 上使得 $AD:DE:EB=1:2:1,M$在$BC$上使得 $AM$为中线,$CD,CE$交$AM$于$G,H$,求 $AG:GH:HM$。
    \ifprintanswers
    \begin{figure}[H]
        \centering        
        \includegraphics[width=0.4\textwidth]{images/image150.png}
    \end{figure}
    \fi
    \begin{solution}
        设 $AG:GH:HM = x:y:z$,在$\triangle ABM$中,以 $CE$ 为横截线,由梅涅劳斯定理,
        \[
        \frac{AH}{HM} \cdot \frac{MC}{CB} \cdot \frac{BE}{EA} =
        \frac{x+y}{z} \cdot \frac{1}{2} \cdot \frac{1}{3} = 1 \Rightarrow x+y = 6z
        \]
        同理,以 $CD$ 为横截线,由梅涅劳斯定理,
        \[
        \frac{AG}{GM} \cdot \frac{MC}{CB} \cdot \frac{BD}{DA} =\frac{x}{y+z} \cdot \frac{1}{2} \cdot \frac{3}{1} = 1 \Rightarrow x = \frac{2}{3}(y+z)
        \]
        设 $z=1$,解得 $x=\dfrac{14}{5},\ y=\dfrac{16}{5}$,于是
        \[
        AG:GH:HM = 14:16:5
        \]
    \end{solution}

    \question $\triangle ABC$ 满足 $AB=20,AC=12$, 若在 $BC$ 边上取两点 $P$ 和 $Q$, 使得 $AQ$ 为 $\angle BAC$ 的角平分线且 $BP=QC$, 求 $\sqrt{AP^2-AQ^2}$ 的值。
    \ifprintanswers
    \begin{figure}[H]
        \centering        
        \includegraphics[width=0.5\textwidth]{images/image88.jpg}
    \end{figure}
    \fi
    \begin{solution}
        令 $BP = CQ = a, PQ = b$,在 $AB$ 上取点 $R$使得 $AR = AC = 12$, 则 
        \[
        \triangle AQR \congr \triangle AQC \; (SSS) \Rightarrow QR = QC = a
        \]
        又 $AQ$ 为 $\angle A$ 的角平分线, 因此 
        \[
        \frac{20}{12} = \frac{a+b}{a} \Rightarrow a+b = \frac{5}{3}a
        \]
        在$\triangle BQR$及$\triangle BAP$中,由余弦定理,
        \[
        \cos \angle B = \frac{(\frac{5a}{3})^2 + 8^2 - a^2}{2\cdot8 \cdot \frac{5a}{3}} = \frac{20^2 + a^2 - AP^2}{2\cdot 20 \cdot a}
        \]
        且在$\triangle BAQ$中,由余弦定理,
        \[
        \cos \angle B = \frac{20^2 + (\frac{5a}{3})^2 - AQ^2}{2\cdot 20 \cdot \frac{5a}{3}} 
        \]
        解得
        \[
        h^2 = 240 - \frac{5}{3}a^2, w^2 = 304 - \frac{5}{3}a^2 \Rightarrow \sqrt{w^2 - h^2} = 8
        \]
    \end{solution}

    \question 如图,$\triangle ABC$ 中,$\angle C=90^\circ$ 且 $AD = DE = EB$。已知 $\angle ACD = \alpha,\angle DCE = \beta,\angle ECB = \gamma$,求 
    \[
    \frac{\sin \alpha \cdot \sin \gamma}{\sin \beta}.
    \]
    \begin{figure}[H]
        \centering
        \includegraphics[width=0.4\linewidth]{images/image137.png}
    \end{figure}
    \ifprintanswers
    \begin{figure}[H]
        \centering
        \includegraphics[width=0.6\linewidth]{images/image138.jpg}
    \end{figure}
    \fi
    \begin{solution}
        过 $A$ 作水平线且过 $B$ 作垂直线,两线交于 $F$,则 $ACBF$ 为矩形。延长 $CE$ 交 $BF$ 于 $H$,延长 $CD$ 交 $AF$ 于 $G$,则有
        \[
        BH = HF,\quad AG = GF.
        \]
        设$BH = HF = a, \quad AG = GF = b$,
        在 $\triangle CAG$及 $\triangle CBH$中,由毕氏定理,
        \[
        \sin \alpha = \frac{b}{\sqrt{4a^2+b^2}}, \quad \sin \gamma = \frac{a}{\sqrt{a^2+4b^2}}.
        \]
        在 $\triangle CGH$ 中,由余弦定理,
        \[
        \cos \beta 
        = \frac{(4a^2+b^2)+(a^2+4b^2)-(a^2+b^2)}{2\sqrt{(4a^2+b^2)(a^2+4b^2)}}
        = \frac{2(a^2+b^2)}{\sqrt{(4a^2+b^2)(a^2+4b^2)}},
        \]
        于是
        \[
        \sin \beta = \sqrt{1-\cos^2 \beta} = \frac{3ab}{\sqrt{(4a^2+b^2)(a^2+4b^2)}}
        \]
        因此
        \[
        \frac{\sin \alpha \cdot \sin \gamma}{\sin \beta}= \frac{1}{3}
        \]
    \end{solution}

    \question 如下图,直角三角形 $ABC$ 中,$\angle BAC=90^\circ$,点 $D, E, F$ 分别为 $BC, CA, AB$ 之中点,$BE$ 与 $CF$ 交于点 $G,BC=18,\angle BGC=150^\circ$,求 $\triangle GBC$ 的面积。
    \begin{figure}[H]
        \centering
        \includegraphics[width=0.5\linewidth]{images/image74.png}
    \end{figure}
    \begin{solution}
        设$AF = FB = a,AE = EC = b,$则
        \[
        a^2 + b^2 = \frac{1}{4}\cdot 18^2=81
        \]
        设直线$BE,CF$斜率为
        \[
        m_1 = -\frac{b}{2a},\quad
        m_2 = -\frac{2b}{a}
        \]
        由
        \[
        \tan 150^\circ = \frac{m_2 - m_1}{1 + m_1 m_2}\Rightarrow
        -\frac{1}{\sqrt{3}} = \frac{-\frac{2b}{a} + \frac{b}{2a}}{1 + \left(-\frac{b}{2a}\right)\left(-\frac{2b}{a}\right)}
        = -\frac{3ab}{2(a^2 + b^2)}
        \]
        代入 $a^2 + b^2 = 81$, 有
        \[
        ab = 18 \sqrt{3}
        \]
        故$\triangle ABC$面积为
        \[
        \frac{1}{2}\cdot 2a \cdot 2b = 36 \sqrt{3}
        \]
        发现点$G$是重心,所以
        \[
        [\triangle GBC] = \frac{1}{3} [\triangle ABC]= 12 \sqrt{3}
        \]
    \end{solution}

    \question 在 $\triangle ABC$中,$\angle C=90^{\circ},G$为$\triangle ABC$的重心,且$G$到$BC,CA$的距离和为$6$。若 $AB=15$,试求 $\triangle ABC$ 的内切圆面积。
    \ifprintanswers
    \begin{figure}[H]
        \centering
        \includegraphics[width=0.5\linewidth]{images/image122.jpg}
    \end{figure}
    \fi
    \begin{solution}
        设$G$在$AC,BC$上的垂足分别为$H,I$,且$GI=a,GH=6-a$,由毕氏定理,
        \[
        GC^2 = a^2+(6-a)^2
        \]
        又$G$为重心,则 
        \[
        GC=\frac{2}{3}CD=\frac{2}{3}\cdot\frac{1}{2}AB=5
        \]
        解得
        \[
        a= \frac{6+\sqrt{14}}{2} \Rightarrow 6-a= \frac{6- \sqrt{14}}{2}
        \]
        又$\triangle AGH \backsim \triangle AFC \ \text{(AAA)}$,
        \[ 
        \frac{2}{3}=\frac{6-a}{FC}\Rightarrow FC= \frac{3}{2}(6-a) \Rightarrow BC =3(6-a)=\frac{18- 3\sqrt{14}}{2}
        \]
        同理,
        \[
        AC= 3a = \frac{18+3\sqrt{14}}{2}
        \]  
        所以 $\triangle ABC$面积为
        \[
        [\triangle ABC]=\frac{1}{2}\cdot \frac{18- 3\sqrt{14}}{2} \cdot \frac{18+ 3\sqrt{14}}{2}
        =\frac{1}{2}r\left(15+\frac{18- 3\sqrt{14}}{2} + \frac{18+ 3\sqrt{14}}{2}\right)  
        \]
        得$r=\dfrac{3}{2}$,故内切圆面积为$\dfrac{9\pi}{4}$。
    \end{solution}

    \question 已知$\triangle ABC$ 满足 $AB=AC=3,BC=1$,点 $D$ 在 $AB$ 上使得 $\triangle ACD$ 与 $\triangle BCD$ 的内切圆半径相等,求该内切圆半径。
    \begin{solution}
        设 $BD = x,CD = y$,三角形面积为内切圆半径乘半周长,于是
        \[
        \frac{[\triangle ACD]}{[\triangle BCD]} = \frac{s_{ACD}}{s_{BCD}} = \frac{6 - x + y}{y + 1 + x}.
        \]
        由等高性质,
        \[
        \frac{[\triangle ACD]}{[\triangle BCD]} = \frac{AD}{BD} = \frac{3-x}{x}.
        \]
        联立两式可得
        \[
        y = \frac{4x - 3}{3 - 2x} \tag{1}
        \]
        在$\triangle BCD$ 中,由余弦定理,
        \[
        y^2 = x^2 + 1^2 - 2 \cdot x \cdot 1 \cdot \frac{1}{6} = x^2 + 1 - \frac{1}{3}x \tag{2}
        \]
        联立$(1),(2)$得
        \[
        x(x-3)(12x^2 - 4x - 9) = 0 \Rightarrow x = \frac{1 + 2\sqrt{7}}{6}, y = \frac{\sqrt{7}}{2}.
        \]
        故 $\triangle ABC$ 面积为
        \[
        [\triangle ABC] = \frac{1}{2} r \cdot \left( \frac{6 - x + y}{2} + \frac{y + 1 + x}{2} \right) = \frac{1}{2} \sqrt{3^2 - \left(\frac{1}{2}\right)^2} = \frac{\sqrt{35}}{4}
        \]
        解得
        \[
        r = \frac{\sqrt{35} - \sqrt{5}}{12}
        \]
    \end{solution}

    \question $\triangle ABC$ 中, $AB=\sqrt{3},\ AC=3\sqrt{3},\ BC=\sqrt{21}$, 由各边分别向外作正三角形。已知此三个向外所作正三角形的重心会形成另一个新的正三角形, 求此新的三角形的面积。 
    \ifprintanswers
    \begin{figure}[H]
        \centering
        \includegraphics[width=0.5\linewidth]{images/image98.jpg}
    \end{figure}
    \fi
    \begin{solution}
        在$\triangle ABC$ 中,由余弦定理,
        \[
        (\sqrt{21})^2=(\sqrt 3)^2 + (3\sqrt 3)^2-2 \cdot \sqrt 3 \cdot 3\sqrt 3 \cdot \cos \angle CAB
        \Rightarrow \angle CAB = 60^\circ
        \]
        由正三角形$\triangle ACD$及$\triangle ABE$,
        \[
        AG_1=1,AG_3=3,\ \angle G_1AB=\angle G_3AC=30^\circ,
        \]
        在$\triangle AG_1G_3$ 中,由余弦定理,
        \[
        (G_1G_3)^2 = 1^2 + 3^2 - 2 \cdot 1 \cdot 3 \cdot \cos \angle G_1AG_3 \Rightarrow G_1G_3 = \sqrt{13}
        \]
        故
        \[
        [\triangle G_1G_2G_3]=\frac{13}{4}\sqrt 3
        \]
    \end{solution}

    \question 已知等腰三角形 $ABC$ 的腰长为 $AB=AC=13$,底边为 $BC=10$。设三角形的内心为 $I$,外心为 $O$,求 $OI$ 的长度。
    \ifprintanswers
    \begin{figure}[H]
        \centering
        \includegraphics[width=0.4\linewidth]{images/image133.jpg}
    \end{figure}
    \fi
    \begin{solution}
        设 $D$ 为$BC$ 中点,则 $AD \perp BC$,且 $BD = DC = 5$,在直角 $\triangle ADC$ 中,由毕氏定理,
        \[
        AD^2 + 5^2 = 13^2 \Rightarrow AD = 12.
        \]
        设 $G$ 为$AC$上的垂足,则
        \[
        \frac12 \cdot 12 \cdot 10 = \frac12 \cdot 13 \cdot DG \Rightarrow DG = \frac{60}{13}
        \]
        在直角 $\triangle ADG$ 中,由毕氏定理,
        \[
        12^2 = AG^2 + \left(\frac{60}{13}\right)^2 \Rightarrow AG = \frac{144}{13}
        \]
        由$\triangle AEO \backsim \triangle AGD \ (AAA)$,
        \[
        \frac{AO}{12} = \frac{\frac{13}{2}}{\frac{144}{13}} \Rightarrow AO = \frac{169}{24}
        \]
        且由$\triangle AFI \backsim \triangle AGD \ \text{(AAA)}$,
        \[
        \frac{ID}{\frac{60}{13}} = \frac{12-ID}{12} \Rightarrow ID = \frac{10}{3}
        \]
        所以
        \[
        OI = AD - AO - ID = 12 - \frac{169}{24} - \frac{10}{3} = \frac{13}{8}
        \]
    \end{solution}

%\question 已知$\triangle ABC$中, $\angle BAC=90^{\circ}$ , $AB=1$ , $AC=\sqrt{3}$ ,若$P$点在$\triangle ABC$内部且满足 $\frac{\vec{PA}}{|\vec{PA}|}+\frac{\vec{PB}}{|\vec{PB}|}+\frac{\vec{PC}}{|\vec{PC}|} = \vec{0}$ , 求序对$(PA,PB,PC)=$ \underline{\hspace{2cm}} . 
%\begin{solution}
% 因为 ${\vec u\over |\vec u|}$ 为单位向量,所以
% \[
% {\overrightarrow{PA} \over |\overrightarrow{PA}|} +{\overrightarrow{PB} \over |\overrightarrow{PB}|} +{\overrightarrow{PC} \over |\overrightarrow{PC}|} =\vec 0
% \]
% 可得
% \[
% \angle APB=\angle BPC=\angle APC = \tfrac{360^\circ}{3} =120^\circ .
% \]

% 令
% \[
% PA=a,\quad PB=b,\quad PC=c ,
% \]
% 则
% \[
% \cos \angle APB= \frac{a^2+b^2-1}{2ab},\quad 
% \cos \angle BPC= \frac{b^2+c^2-4}{2bc},\quad 
% \cos \angle APC= \frac{a^2+c^2-3}{2ac}.
% \]
% 由 $\cos 120^\circ=-\tfrac12$ 得
% \[
% \begin{cases}
% a^2+b^2+ab=1 \quad (1)\\
% b^2+c^2+bc=4 \quad (2)\\
% c^2+a^2+ca=3 \quad (3)
% \end{cases}
% \]

% 将(1)+(2)+(3)得
% \[
% 2(a^2+b^2+c^2)+ab+bc+ca=8 \quad (4)
% \]

% 另一方面,
% \[
% \triangle ABC 面积 = \triangle PAB+\triangle PBC+\triangle PCA
% \]
% \[
% = \tfrac12 ab\sin120^\circ + \tfrac12 bc\sin120^\circ + \tfrac12 ac\sin120^\circ
% = \tfrac{\sqrt{3}}{4}(ab+bc+ca).
% \]
% 而
% \[
% \triangle ABC 面积 = \tfrac12 AB\times AC = \tfrac{\sqrt{3}}{2}.
% \]
% 所以
% \[
% ab+bc+ca=2.
% \]

% 代入(4)得
% \[
% a^2+b^2+c^2=3.
% \]
% 因此
% \[
% (a+b+c)^2 = a^2+b^2+c^2+2(ab+bc+ca) = 3+4=7 \Rightarrow a+b+c=\sqrt{7}.
% \]

% 又由
% \[
% (1)+(2)\Rightarrow a^2+2b^2+c^2+ab+bc=5,
% \]
% \[
% (2)+(3)\Rightarrow a^2+b^2+2c^2+bc+ca=7,
% \]
% \[
% (1)+(3)\Rightarrow 2a^2+b^2+c^2+ab+ac=4,
% \]
% 可得
% \[
% \begin{cases}
% b^2+ab+bc=2 \\
% c^2+bc+ca=4 \\
% a^2+ab+ac=1
% \end{cases}
% \]

% 所以
% \[
% \begin{cases}
% b(a+b+c)=2 \\
% c(a+b+c)=4 \\
% a(a+b+c)=1
% \end{cases}
% \Rightarrow
% \begin{cases}
% a= \tfrac{1}{\sqrt{7}} \\
% b= \tfrac{2}{\sqrt{7}} \\
% c= \tfrac{4}{\sqrt{7}}
% \end{cases}
% \]

% 故答案为
% \[
% (PA,PB,PC)=\left( \tfrac{1}{\sqrt{7}}, \tfrac{2}{\sqrt{7}}, \tfrac{4}{\sqrt{7}} \right).
% \]
%\end{solution}

\question 在 $\triangle ABC$ 中,$AB=6,BC=4,CA=5$,圆 $O_1, O_2, O_3$ 为 $\triangle ABC$ 的三旁切圆,$O_1$ 与 $BC$ 相切于 $D$,$O_2$ 与 $CA$ 相切于 $E$,$O_3$ 与 $AB$ 相切于 $F$。求面积比 
\[
\frac{[\triangle DEF]}{[\triangle ABC]}.
\]
\begin{figure}[H]
    \centering
    \includegraphics[width=0.35\linewidth]{images/image108.png}
\end{figure}
\ifprintanswers
\begin{figure}[H]
    \centering
    \includegraphics[width=0.5\linewidth]{images/image109.jpg}
\end{figure}
\fi
\begin{solution}
    设$a = BC = 4, b = CA = 5, c = AB = 6, s = \dfrac{a+b+c}{2} = \dfrac{15}{2}$,由旁切圆性质,
    \[
    AF = CD = s - b = \frac{5}{2}, \quad
    BF = CE = s - a = \frac{7}{2}, \quad
    AE = BD = s - c = \frac{3}{2}
    \]
    由$\triangle ABC \backsim \triangle AEF \backsim \triangle BDF \backsim \triangle CDE$,
    \begin{align*}
    \frac{S_{\triangle DEF}}{[\triangle ABC]}
    &= 1- \frac{S_{\triangle AEF}}{[\triangle ABC]} - \frac{S_{\triangle BDF}}{[\triangle ABC]} - \frac{S_{\triangle CDE}}{[\triangle ABC]}\\[2mm]
    &= 1 - \left(
    \frac{AE \cdot AF}{AC \cdot AB} +
    \frac{BD \cdot BF}{BC \cdot AB} +
    \frac{CD \cdot CE}{AC \cdot BC}
    \right)\\[2mm]
    &= 1 - \left(
    \frac{\frac{3}{2} \cdot \frac{5}{2}}{5 \cdot 6} +
    \frac{\frac{3}{2} \cdot \frac{7}{2}}{4 \cdot 6} +
    \frac{\frac{5}{2} \cdot \frac{7}{2}}{5 \cdot 4}
    \right)= \frac{7}{32}
    \end{align*}
    \textcolor{red}{(待验证相似三角形,顺序,旁切圆性质)}
\end{solution}

    \question 设钝角三角形 $ABC$ 的内切圆半径为 $\sqrt{3}$,已知内切圆圆心 $O$ 到顶点 $A$ 的距离为 $2$,而 $O$ 到顶点 $B$ 的距离为 $2\sqrt{7}$,求 $\triangle ABC$ 的面积。
    \ifprintanswers
    \begin{figure}[H]
        \centering
        \includegraphics[width=0.3\linewidth]{images/image129.jpg}
    \end{figure}
    \fi
    \begin{solution}
        设 $D, E, F$ 分别为在$AB,BC,AC$上的内切圆切点,在$\triangle ADO,\triangle OBD$中,由毕氏定理,
        \[
        2^2 = (\sqrt{3})^2 + AD^2, (2\sqrt{7})^2 = (\sqrt{3})^2 + BD^2
        \]
        则
        \[
        AD = AF = 1, \quad BD = BE = 5, \quad CE = CF = a
        \]
        $\triangle ABC$半周长为$s=a+6$,则面积为
        \[
        [\triangle ABC]= \sqrt{3} (a + 6) =\sqrt{a+6\cdot a\cdot1\cdot5} \Rightarrow a = 9
        \]
        故
        \[
        [\triangle ABC]= 15 \sqrt{3}
        \]
    \end{solution}

    \question 在 $\triangle ABC$ 中,已知 $\cos\angle BAC = \dfrac{4\sqrt{3}}{7},\angle B = 30^\circ,O$ 为 $\triangle ABC$ 的内心,直线 $AO$ 与 $BC$ 相交于 $D$ 点,求 $\dfrac{BD}{CD}$。
    \ifprintanswers
    \begin{figure}[H]
        \centering
        \includegraphics[width=0.6\linewidth]{images/image116.jpg}
    \end{figure}
    \fi
    \begin{solution}
        设$\angle CAD = \angle DAB = \theta, CD = 1, BD = k$,  由 $\cos \angle A = \dfrac{4\sqrt{3}}{7}$,得 
        \[
        \sin \angle A = \dfrac{1}{7}
        \]
        在 $\triangle ABD$ 及$\triangle ACD$中,由正弦定理,
        \[
        \frac{k}{\sin \theta} = \frac{AD}{\sin 30^\circ}, \quad \frac{1}{\sin \theta} = \frac{AD}{\sin \angle C} 
        \]
        于是
        \[
        AD=\frac{k}{2\sin \theta} = \frac{\sin(150^\circ - \angle A)}{\sin \theta} \Rightarrow k = 2\sin(150^\circ - \angle A)
        \]
        即
        \[
        \frac{BD}{CD}=k = 2\sin(150^\circ - \angle A) = 2\left( \frac{1}{2} \cdot \frac{4\sqrt{3}}{7} - \left(-\frac{\sqrt{3}}{2}\right) \cdot \frac{1}{7} \right) = \frac{5\sqrt{3}}{7}
        \]
    \end{solution}

    \question 若等腰三角形$\triangle ABC$的外接圆半径 $R$ 与内切圆半径 $r$ 之比
    \[
    \frac{R}{r} = 1+\sqrt{2},
    \]
    且顶角$A$ 不为直角,求 $\sin\dfrac{A}{2}$。
    \ifprintanswers
    \begin{figure}[H]
        \centering
        \includegraphics[width=0.5\linewidth]{images/image126.jpg}
        \end{figure}
    \fi
    \begin{solution}
        设 $O$ 为外接圆圆心, $P$ 为内切圆圆心,$D,E$为$P$在$AC,AB$上的垂足,$G$为$AB$中点,则
        \[
        AB = 2R \cos \frac{A}{2}, \quad
        AF = 2R \cos^2 \frac{A}{2}, \quad
        OP = 2R \cos^2 \frac{A}{2} - R - r
        \]
        又由$\triangle AGO\backsim \triangle AEP \ (AAA)$,
        \[
        \frac{GO}{EP} = \frac{AO}{AP} \Rightarrow \frac{R \sin \frac{A}{2}}{r} = \frac{R}{2R\cos^2 \frac{A}{2} - r}
        \]
        令$x = \sin \dfrac{A}{2}$,将$\dfrac{R}{r} = 1+\sqrt{2}$代入得
        \[
        (x+1)\big((2+2\sqrt{2}) x^2 - (2+2\sqrt{2}) x + 1 \big) = 0
        \]
        舍去 $x=-1$,得
        \[
        \sin \frac{A}{2} = \frac{\sqrt{2}}{2} \quad \text{或} \quad \sin \frac{A}{2} = \frac{2-\sqrt{2}}{2}
        \]
    \end{solution}

    \question 已知四边形 $ABCD$ 中,$\; AB=BC=CD,\angle A=54^\circ,\angle B=108^\circ,\angle C$ 为钝角,求 $\angle C$。
    \ifprintanswers
    \begin{figure}[H]
        \centering
        \includegraphics[width=0.5\linewidth]{images/image92.jpg}
    \end{figure}
    \fi
    \begin{solution}
        作 $BD$,令 $\angle CBD=\angle CDB=\theta$, $AB=BC=CD=a$, 则
        \[
        \angle C=180^\circ-2\theta, \quad
        \angle ABD=108^\circ-\theta, \quad
        \angle ADB=18^\circ+\theta,  
        \] 
        在$\triangle ABD$及$\triangle BCD$中,由正弦定理,
        \[
        \frac{a}{\sin(18^\circ+\theta)}=\frac{BD}{\sin54^\circ},\quad
        \frac{a}{\sin\theta}=\frac{BD}{\sin(180^\circ-2\theta)}.
        \]
        故
        \[
        BD=\frac{\sin54^\circ}{\sin(18^\circ+\theta)}=\frac{\sin(180^\circ-2\theta)}{\sin\theta}= \frac{\sin2\theta}{\sin\theta}=2\cos\theta.
        \]
        于是
        \[
        \sin54^\circ=2\cos\theta\,\sin(18^\circ+\theta)=\sin18^\circ\cos2\theta+\cos18^\circ\sin2\theta+\sin18^\circ,
        \]
        即
        \[
        \sin(2\theta+18^\circ)=\sin54^\circ-\sin18^\circ=\frac{\sqrt5+1}{4}-\frac{\sqrt5-1}{4}=\frac{1}{2}
        \]
        故 
        \[
        2\theta+18^\circ=30^\circ \Rightarrow \angle C=168^\circ>90^\circ
        \]
    \end{solution}

    \question 已知边长为 $2$ 的正方形 $ABCD,E,F$ 分别为 $AD,CD$ 的中点,$AF$ 与 $EB,BD$ 相交于 $G,H$,求四边形 $DEGH$ 的面积。
    \ifprintanswers
    \begin{figure}[H]
        \centering
        \includegraphics[width=0.5\linewidth]{images/image145.png}
    \end{figure}
    \fi
    \begin{solution}
        首先有
        \[
        1=[\triangle ADF]=[\triangle ADH]+[\triangle DFH] =\frac{1}{2} \cdot DH \cdot 1 \cdot \sin 45^\circ +  \frac{1}{2} DH \cdot \frac{1}{2} \cdot \sin 45^\circ
        \]
        解得
        \[
        DH = \frac{2\sqrt{2}}{3}, \quad [\triangle DFH] = \frac{1}{3}
        \]
        设 $\alpha = \angle DAF$,由$\triangle ADF$得
        \[
        \sin \alpha = \frac{1}{\sqrt{5}}, \quad \sin(90^\circ-\alpha) = \frac{2}{\sqrt{5}}.
        \]
        同理有
        \[
        1=[\triangle ABE]=[\triangle ABG]+[\triangle AEG] 
        = \frac{1}{2} \cdot AG \cdot 1 \cdot \sin (90^\circ-\alpha) + \frac{1}{2} AG \cdot \frac{1}{2} \cdot \alpha
        \]
        解得
        \[
        AG = \frac{2\sqrt{5}}{5}, \quad [\triangle AEG] = \frac{1}{5}
        \]
        故四边形 $DEGH$ 的面积为
        \[
        [DEGH] = [\triangle ADF] - [\triangle DFH] - [\triangle AEG] = 1 - \frac{1}{3} - \frac{1}{5} = \frac{7}{15}
        \]
    \end{solution}

    \question 长方形 $ABCD$ 边长为 $2,3$,以 $A$ 为旋转中心逆时针旋转 $30^\circ$。求重叠部分面积。
    \begin{figure}[H]
        \centering
        \includegraphics[width=0.5\linewidth]{images/image67.png}
    \end{figure}
    \ifprintanswers
    \begin{figure}[H]
        \centering
        \includegraphics[width=0.5\linewidth]{images/image69.png}
    \end{figure}
    \fi
    \begin{solution}
        如图, $AB' = 3, \angle B'AB = 30^\circ$,则
        \[
        B'R = \frac{3}{2}, \quad AR = \frac{3\sqrt{3}}{2}
        \]
        又$B'Q = 2 - B'R = \dfrac{1}{2},\angle QB'C' = 30^\circ$,则
        \[ 
        QS = \frac{\sqrt{3}}{6} 
        \]
        重叠面积为
        \[
        [ABCD]- [B'PCS] - [ABPB'] =6 - \left(\frac{3}{2} - \frac{17\sqrt{3}}{24} \right) - \left(\frac{9}{2} - \frac{9\sqrt{3}}{8} \right) = \frac{11\sqrt{3}}{6} 
        \]
    \end{solution}

    \question 一边长为 $1$ 的正方形 $\Gamma_1$ 绕一顶点逆时针旋转 $\theta$ 得正方形 $\Gamma_2$,若两正方形重叠部分的面积为 $\dfrac{2}{3}$,求 $\tan\theta$。
    \ifprintanswers
    \begin{figure}[H]
        \centering
        \includegraphics[width=0.4\linewidth]{images/image72.png}
    \end{figure}
    \fi
    \begin{solution}
        令 $\angle AOP = \theta, \angle AOC = 90^\circ - \theta$,又 $\triangle OAB \congr \triangle OCB \ \text{(RHS)}$,  
        \[
        AB = OA \tan \angle AOB 
        = \tan\left( 45^\circ - \frac{\theta}{2} \right) 
        = \frac{1 - \tan\frac{\theta}{2}}{1 + \tan\frac{\theta}{2}}
        \]
        重叠面积为
        \[ 
        2[\triangle OAB] = OA \cdot AB = AB = \frac{1 - \tan\frac{\theta}{2}}{1 + \tan\frac{\theta}{2}} = \frac{2}{3}
        \]
        解得
        \[
        \tan\frac{\theta}{2} = \frac{1}{5} \Rightarrow \tan\theta = \frac{2\tan\frac{\theta}{2}}{1 - \tan^2\frac{\theta}{2}} = \frac{5}{12}
        \]
    \end{solution}

    \question $ABCD$为平行四边形,点$E,F$分别在边$AB,BC$上。已知
    \[
    [\triangle AED] = 7, \quad [\triangle EBF] = 3, \quad [\triangle CDF] = 6,
    \]
    求\(\triangle DEF\) 的面积。
    \ifprintanswers
    \begin{figure}[H]
        \centering
        \includegraphics[width=0.4\linewidth]{images/image66.png}
    \end{figure}
    \fi
    \begin{solution}
        不妨设$ABCD$为矩形,且设$AD = a, BE = b$,则
        \[
        AE = \frac{14}{a}, \quad BF = \frac{6}{b}, \quad CF = a - \frac{6}{b}, \quad CD = \frac{12}{a - \frac{6}{b}} = \frac{12b}{ab - 6}
        \]
        解 $AB = CD$得,
        \[
        b + \frac{14}{a} = \frac{12b}{ab - 6} \Rightarrow ab = 2 + 2\sqrt{22}
        \]
        \(\triangle DEF\) 的面积为
        \[
        [ABCD] - [\triangle AED] - [\triangle EBF] - [\triangle CDF] = (ab + 14) - 16 = 2\sqrt{22}
        \]
    \end{solution}
    \ifprintanswers
    \begin{figure}[H]
        \centering
        \includegraphics[width=0.5\linewidth]{images/image68.png}
    \end{figure}
    \fi
    \begin{solution}
        如果不假设为矩形,如上图,步骤相同,只是算出来的是 \(ab \sin \theta = 2 + 2\sqrt{22}\),结果是一样的。
    \end{solution}

    \question 如图,已知 $DE=10, EF=8, BF=6$,求正方形 $ABCD$ 面积。
    \begin{figure}[H]
        \centering
        \includegraphics[width=0.4\linewidth]{images/image60.png}
    \end{figure}
    \begin{solution}
        在$\triangle BFE$中,由毕氏定理,
        \[
        BE=\sqrt{6^2+8^2}=10
        \]
        设 $\angle BEF=\theta,\sin\theta=\dfrac{3}{5}$,则
        \[
        \cos \angle BED=\cos(90^\circ+\theta)=-\sin \theta=-\frac{3}{5}
        \]
        在$\triangle BED$中,由余弦定理,
        \[
        BD^2 = 10^2+10^2 - 2\cdot 10\cdot 10 \cdot \left(-\frac{3}{5}\right) \Rightarrow BD=\sqrt{320}
        \]
        正方形面积为
        \[
        AD^2=160
        \]
    \end{solution}

    \question 已知$P$为$\triangle ABC$内一点,满足$\angle PBA=80^{\circ},\angle PBC=20^{\circ},\angle PCB=10^{\circ}$,且$\angle PCA=30^{\circ}$,求$\angle PAC$。
    \ifprintanswers
    \begin{figure}[H]
        \centering
        \includegraphics[width=0.5\linewidth]{images/image76.png}
    \end{figure}
    \fi
    \begin{solution}
        延长$BP$交$AC$于$D$,设$\angle PAD=\theta$, $PD=a$,则  
        \[
        \angle PDA=60^\circ, \angle DPC=30^\circ \Rightarrow DC=DP=a
        \]  
        在$\triangle PDC$中,由余弦定理,
        \[
        CP^2=a^2+a^2-2a^2\cos120^\circ
        \Rightarrow CP=\sqrt3\,a
        \]
        在$\triangle ADP$中,由正弦定理,
        \[
        \frac{a}{\sin\theta}=\frac{AP}{\sin60^\circ} \Rightarrow AP=\frac{\sqrt3\,a}{2\sin\theta} \tag{1}
        \]
        在$\triangle ABP$中,由正弦定理,
        \[
        \frac{AP}{\sin80^\circ}=\frac{BP}{\sin(40^\circ-\theta)} \Rightarrow AP=\frac{\sin80^\circ}{\sin(40^\circ-\theta)}BP \tag{2}
        \]
        在$\triangle BCP$中,由正弦定理,
        \[
        \frac{\sqrt3\,a}{\sin20^\circ}=\frac{BP}{\sin10^\circ} \Rightarrow BP=\frac{\sin10^\circ}{\sin20^\circ}\sqrt3\,a \tag{3}
        \]
        由式$(1),(2),(3)$,
        \[
        AP=\frac{\sin80^\circ}{\sin(40^\circ-\theta)}\cdot\frac{\sin10^\circ}{\sin20^\circ}\sqrt3\,a=\frac{\sqrt3\,a}{2\sin\theta}
        \] 
        整理得
        \[
        \sin\theta=\frac{\sin20^\circ}{2\sin10^\circ\sin80^\circ}\sin(40^\circ-\theta)
        =\frac{\sin20^\circ}{\cos70^\circ-\cos90^\circ}\sin(40^\circ-\theta)
        \]
        故
        \[
        \theta=20^\circ
        \]
    \end{solution}

    \question 已知 $\triangle ABC$ 中, $\angle ACB=90^{\circ}$, $AB=2AC$, 又此三角形内部有一点 $P$,满足 $PA=\sqrt{2}, PB=\sqrt{10},PC=1$,求 $\triangle ABC$ 的面积。
    \ifprintanswers
    \begin{figure}[H]
        \centering
        \includegraphics[width=0.5\linewidth]{images/image56.png}
    \end{figure}
    \fi
    \begin{solution}
        设$\angle PCB = \theta, A(0,a), B(\sqrt{3}a, 0), C(0,0), P(\cos\theta, \sin\theta)$,则
        \[
        PA^2 = \cos^2\theta + (\sin\theta - a)^2 = 2,       PB^2 = (\cos\theta - \sqrt{3}a)^2 + \sin^2\theta = 10,
        \]
        展开得
        \[
        \sin\theta = \frac{a^2 - 1}{2a}, \quad \cos\theta = \frac{3a^2 - 9}{2\sqrt{3}a}
        \]
        解得
        \[
        \frac{(a^2 - 1)^2}{4a^2} + \frac{(3a^2 - 9)^2}{12 a^2} = 1 \Rightarrow a^2 = 3 + \sqrt{2}
        \]
        故$\triangle ABC$面积为
        \[
        \frac{1}{2} \cdot a \cdot \sqrt{3}a = \frac{3 \sqrt{3} + \sqrt{6}}{2}
        \]
    \end{solution}

    \question 设 $P$ 为正 $\triangle ABC$ 内一点,满足 $AP=2, BP=3, CP=\sqrt{7}$,求正 $\triangle ABC$ 的边长。
    \ifprintanswers
    \begin{figure}[H]
        \centering
        \includegraphics[width=0.5\linewidth]{images/image94.jpg}
    \end{figure}
    \fi
    \begin{solution}
        将 $\triangle CPB$ 绕顶点 $B$ 逆时针旋转 $60^\circ$,此时点 $C$ 与点 $A$ 重合,且有
        \[
        P'B = PB =3,\quad P'A = PC = \sqrt{7},\quad \angle PBP' =60^\circ
        \]
        所以 $\triangle PP'B$ 是正三角形,得 $PP'=3$,在$\triangle APP'$中,由余弦定理,
        \[
        7=2^2+3^2-2\cdot 2 \cdot 3 \cdot \cos \angle APP' \Rightarrow \angle APP' =60^\circ
        \]
        因此 $\angle APB = 60^\circ + 60^\circ =120^\circ$,在$\triangle APB$中,由余弦定理,
        \[
        AB^2 = 2^2 + 3^2 - 2\cdot 2\cdot 3 \cdot \cos \angle APB \Rightarrow AB = \sqrt{19}
        \]
    \end{solution}

    \question 如图,一点$P$在 $\triangle ABC$ 内使得$\angle PAB=\angle PBC=\angle PCA=\alpha$且 $AB=7,BC=8,CA=9$。求 $\tan\alpha$。
    \begin{figure}[H]
        \centering
        \includegraphics[width=0.4\linewidth]{images/image148.png}
    \end{figure}
    \begin{solution}
        设 $AP=x,BP=y,CP=z$,在 $\triangle ABP,\triangle BCP,\triangle CAP$中,由余弦定理,
        \[
        y^{2}=7^{2}+x^{2}-14x\cos\alpha,
        \quad
        z^{2}=8^{2}+y^{2}-16y\cos\alpha,
        \quad
        x^{2}=9^{2}+z^{2}-18z\cos\alpha.
        \]
        三式相加得
        \[
        (14x+16y+18z)\cos\alpha=194 \tag{1}
        \]
        又
        \[
        [\triangle ABC]=[\triangle ABP]+[\triangle BCP]+[\triangle CAP]=\frac12(7x+8y+9z)
        \]
        由海伦公式,
        \[
        [\triangle ABC]=\sqrt{12\cdot3\cdot4\cdot5}=12\sqrt{5}
        \]
        因此
        \[
        (7x+8y+9z)\sin\alpha=24\sqrt{5}\tag{2}
        \]
        $\dfrac{(2)}{(1)}$得
        \[
        \tan\alpha=\frac{24\sqrt{5}}{97}
        \]
    \end{solution}

    \question 在正方形 $ABCD$ 中,点 $E,F$ 在对角线 $AC$ 上,且 $\angle EDF=45^\circ$。若 $AE=x,EF=y,FC=z$,证明 $y^2=x^2+z^2$。
    \ifprintanswers
    \begin{figure}[H]
        \centering
        \includegraphics[width=0.4\linewidth]{images/image178.png}
    \end{figure}
    \fi
    \begin{solution}
        将 $\triangle DFC$ 绕点 $D$ 逆时针旋转 $90^\circ$得 $\triangle DF'A$,使得 $DF=DF'$。由$\angle F'AD=\angle FCD=45^\circ$,
        \[
        \angle F'AE = \angle EAD + \angle F'AD = \angle EAD +\angle FCD=90^\circ
        \]
        在 $\triangle F'AE$ 中,由毕氏定理,
        \[
        F'E^2 = F'A^2 + AE^2 =  x^2+z^2
        \]
        且由于 $\triangle F'DE \congr \triangle FDE \ \text{(SAS)},$有$F'E=FE=y$,故得证
        \[
        y^2 = x^2+z^2
        \]
    \end{solution}
    \ifprintanswers
    \begin{figure}[H]
        \centering
        \includegraphics[width=0.4\linewidth]{images/image179.png}
    \end{figure}
    \fi
    \begin{solution}
        设 $\angle ADE=\theta$,则
        \[
        \angle AED = 135^\circ-\theta,\quad \angle DEF=45^\circ+\theta,\quad \angle EFD=90^\circ-\theta
        \]
        在 $\triangle AED,\triangle DEF,\triangle DFC$ 中,由正弦定理,
        \[
        \frac{x}{\sin\theta}=\frac{ED}{\sin 45^\circ}, \quad
        \frac{y}{\sin 45^\circ} = \frac{FD}{\sin(45^\circ+\theta)},\quad
        \frac{z}{\sin(45^\circ-\theta)} = \frac{FD}{\sin 45^\circ}
        \]
        化简得
        \[
        \frac{x}{y}=\sin 2\theta,\qquad \frac{z}{y}=\cos 2\theta
        \]
        于是
        \[
        \frac{x^2}{y^2}+\frac{z^2}{y^2}=\sin^2 2\theta+\cos^2 2\theta=1 \Rightarrow x^2+z^2=y^2
        \]
        证毕。
        \textcolor{red}{怎样化简得？}
    \end{solution}

    \question 设 $P$ 是正方形 $ABCD$ 内部一点,且 $P$ 到 $A,B,C$ 三顶点的距离分别为 $1,2,3$,求此正方形的面积。
    \ifprintanswers
    \begin{figure}[H]
        \centering
        \includegraphics[width=0.4\linewidth]{images/image51.png}
    \end{figure}
    \fi
    \begin{solution}
        设正方形$ABCD$边长为 $a$,在$\triangle APB$及$\triangle BPC$中,由余弦定理,
        \[
        \cos \angle PBA = \frac{a^{2} + 2^2-1^2}{2\cdot 2 \cdot a}, \quad \cos \angle PBC = \frac{a^{2} +2^2 - 3^2}{2\cdot 2 \cdot a}
        \]
        由$\cos \angle PBC= \cos (90^\circ - \angle PBA)=\sin \angle PBA$,有
        \[
        \cos^{2} \angle PBA + \sin^2 \angle PBA = \left(\frac{a^{2} + 3}{4a}\right)^{2} + \left(\frac{a^{2} - 5}{4a}\right)^{2} = 1
        \]
        解得
        \[
        [ABCD]=a^{2} = 5 + 2 \sqrt{2}.
        \]
    \end{solution}

    \question 平面上,$P$为$AB$上一点满足$AP=5,BP=3,Q$为平面上一点满足$AQ=7,BQ=3$。若以$AB$为直径作一圆$C$,自$P$向$Q$作射线$PQ$交圆$C$于点$R$,试求$PR$的长度。
    \ifprintanswers
    \begin{figure}[H]
        \centering
        \includegraphics[width=0.5\linewidth]{images/image61.png}
    \end{figure}
    \fi
    \begin{solution}
        在$\triangle ABQ$中,由余弦定理,
        \[
        7^2 = 8^2+3^2 -2\cdot 3 \cdot 8 \cdot \cos \angle B
        \Rightarrow \angle B=60^\circ
        \]
        又$BP=BQ=3$,故$\triangle PQB$为等边三角形,于是$PQ=3$;在$\triangle OPR$中,由余弦定理,
        \[
        4^2=1^2 + PR^2-2 \cdot 1 \cdot PR \cdot \cos 60^\circ 
        \]
        解得
        \[
        PR = \frac{-1 + \sqrt{61}}{2}
        \]
    \end{solution}

    \question 两圆 $O_1,O_2$ 都经过 $\triangle ABC$ 的顶点 $A$,且分别与 $BC$ 边相切于点 $B,C$ 。已知圆 $O_1,O_2$ 的面积分别为 $m,n$,试以 $m,n$ 表示 $\triangle ABC$ 外接圆的面积。
    \ifprintanswers
    \begin{figure}[H]
        \centering
        \includegraphics[width=0.5\linewidth]{images/image224.png}
    \end{figure}
    \fi
    \begin{solution}
        设$O_1,O_2,\triangle ABC$外接圆的半径为$R_1,R_2,R_3;\angle ABC = \beta, \angle ACB = \gamma$。在 $\wideparen{AB},\wideparen{AC}$ 分别取点 $P,Q$,由弦切角定理,
        \[
        \angle APB = \angle ABC = \beta, \quad \angle AQC = \angle ACB = \gamma.
        \]
        由正弦定理,
        \[
        4 R_3^2 = \frac{\overline{AC}}{\sin \beta} \cdot \frac{\overline{AB}}{\sin \gamma} = \frac{2 R_2 \sin \gamma}{\sin \beta} \cdot \frac{2 R_1 \sin \beta}{\sin \gamma} = 4 R_1 R_2 \Rightarrow R_3^2 = R_1 R_2
        \]
        所以 $\triangle ABC$ 外接圆的面积为
        \[
        \pi R_3^2 = \pi R_1 R_2 = \sqrt{\frac{m}{\pi}} \cdot \sqrt{\frac{n}{\pi}} \cdot \pi = \sqrt{mn}
        \]
    \end{solution}

    \question 一半径为 4 的大圆里面有三个两两外切,且皆与大圆内切的小圆,其中两小圆半径分别为 1 和 3,且它们的圆心在大圆的一条直径上,求第三个小圆的半径。
    \ifprintanswers
    \begin{figure}[H]
        \centering
        \includegraphics[width=0.5\linewidth]{images/image140.png}
    \end{figure}
    \fi
    \begin{solution}
        设$B=(0,0)$为大圆圆心 ,$A=(-1,0)$为半径为 3 的小圆圆心,$C=(3,0)$为半径为 1 的小圆圆心,$D$为所求半径为 $r$ 的小圆圆心,由海伦公式及等高性质,
        \[
        3=\frac{[\triangle BDC]}{[\triangle ABD]}=
        \frac{\sqrt{4(3-r)\cdot r \cdot1}}{\sqrt{4(1-r)\cdot r\cdot3}}
        \]
        其中$\triangle ABD,\triangle BDC$半周长为
        \[
        [\triangle ABD]=\frac{1}{2}(3+r+(4-r)+1)=4, \quad s_{\triangle BDC}=\frac{1}{2}(1+r+(4-r)+3)=4
        \]
        解得第三个小圆的半径为
        \[
        r=\frac{12}{13}
        \]
    \end{solution}

    \question 在边长为 4 的正三角形 $ABC$ 内接圆中,点 $P$ 在弧 $AC$ 上满足 $AP \cdot PC = 5$。求 $BP$ 的长度。
    \begin{figure}[H]
        \centering
        \includegraphics[width=0.4\linewidth]{images/image152.png}
    \end{figure}
    \begin{solution}
        在四边形 $ABPC$中,由托勒密定理,
        \[
        AB \cdot PC + AP \cdot BC = AC \cdot BP 
        \]
        由$AB=AC=BC$,得
        \[
        BP = AP + PC
        \]
        易知$\angle APC = 120^\circ$,在$\triangle APC$中,由余弦定理,
        \[
        4^2 = AP^2 + PC^2 - 2 \cdot AP \cdot PC \cos 120^\circ \Rightarrow AP^2 + PC^2 = 11
        \]
        于是
        \[
        BP^2 = (AP+PC)^2 = AP^2 + PC^2 + 2 \cdot AP \cdot PC = 21 \Rightarrow BP = \sqrt{21}
        \]
    \end{solution}

    \question 已知三角形 $ABC$ 的外心为 $O$,且 $\angle B$ 为钝角。延长$BC$与外接圆在点 $A$处的切线交于 $D$使得$AO = AD = 2$。$F$ 在 $BC$ 上满足 $AF \perp OD$且交于点 $E$。过点 $A,D,O$ 的圆与 $BC$ 的交于点$H$。已知 $FH = \dfrac{\sqrt{3}}{3}$,求 $\triangle OEF$ 的面积。
    % 23. [6 pts] [7] Let $ABC$ be a triangle with circumcenter $O$ and $\angle B$ obtuse. Suppose the tangent to its circumcircle at $A$ intersects ray $CB$ at $D$, and that $AO=AD=2$. Suppose that $F$ lies on $BC$ with $AF \perp OD$, and let $E$ be the point where $AF$ meets $OD$. Let $H$ be the point where the circle passing through $ADO$ meets $BC$. What is the area of triangle $OEF$ if $FH=\frac{\sqrt{3}}{3}$?23. $\frac{1}{2}\sqrt{2}$. Let $J$ be the point where the extension of $AF$ meets circle $ADO$. Let $x = DF$ and $y = EF$. Power of a point using chords $DH$ and $AJ$ yields $x\frac{\sqrt{3}}{3}=(\sqrt{2}+y)(\sqrt{2}-y)$ while Pythagoras on triangle $DEF$ yields $x^{2}=y^{2}+2$. Combining, we get $x\frac{\sqrt{3}}{3}=2-(x^{2}-2)$ so $x^{2}+x\frac{\sqrt{3}}{3}-4=0$. The solution is $x=\sqrt{3}$ and $y=1$, so the area of the right triangle $OEF$ is $\frac{1}{2}\sqrt{2}\cdot1$.
    \ifprintanswers
    \begin{figure}[H]
        \centering
        \includegraphics[width=0.6\linewidth]{images/image160.png}
    \end{figure}
    \fi
    \begin{solution}
延长 $AF$ 与过点 $A,D,O$ 的圆交于点$J$,设 $x = DF,y = EF$。由幂点定理（$DH$ 与 $AJ$）有
\[
x \cdot \frac{\sqrt{3}}{3} = (\sqrt{2}+y)(\sqrt{2}-y) = 2 - y^2.
\]
由直角三角形 $DEF$ 得
\[
x^2 = y^2 + 2.
\]
联立得
\[
x \cdot \frac{\sqrt{3}}{3} = 2 - (x^2 - 2) \implies x^2 + x\frac{\sqrt{3}}{3} - 4 = 0.
\]
解得 $x = \sqrt{3}$,$y = 1$。

因此直角三角形 $OEF$ 的面积为
\[
\frac{1}{2} \cdot OE \cdot EF = \frac{1}{2} \cdot \sqrt{2} \cdot 1 = \frac{1}{2}\sqrt{2}.
\]
\textcolor{red}{$AE=\sqrt{2}$?怎样证明$\angle DAE = 45^\circ$?}
\end{solution}

    \question 已知$\triangle ABC$中,$AB=5,BC=7,AC=6$。点$D$在$BC$上移动,从$D$向$AB,AC$作垂线,垂足分别为$E,F$,求$EF$的最小值。
    \ifprintanswers
    \begin{figure}[H]
        \centering
        \includegraphics[width=0.45\linewidth]{images/image54.png}
    \end{figure}
    \fi
    \begin{solution}
        由$\angle DFA = \angle DEA = 90^\circ$,故点$A, E, D, F$共圆,设该圆半径为$R$,在$\triangle AEF$中,由正弦定理,
        \[
        EF = 2R \sin \angle A = AD \sin \angle A
        \]
        当$AD \perp BC$时,$AD$取得最小值;在$\triangle ABC$中,由余弦定理,
        \[
        \cos \angle A = \frac{5^2 + 6^2 - 7^2}{2 \cdot 5 \cdot 6} = \frac{1}{5} \Rightarrow \sin \angle A = \frac{2\sqrt{6}}{5}
        \]
        $\triangle ABC$面积为
        \[
        [\triangle ABC] = \frac{1}{2} \cdot 5 \cdot 6 \cdot \sin \angle A = \frac{1}{2} AD_{\min} \cdot BC \Rightarrow AD_{\min}=\frac{12 \sqrt{6}}{7}
        \]
        此时
        \[
        EF_{\min} = \frac{12 \sqrt{6}}{7} \cdot \frac{2 \sqrt{6}}{5} = \frac{144}{35}.
        \]
    \end{solution}

    \question 已知三角形 $ABC$ 的重心为 $G$,且 $GB=7,GC=3,G$ 到直线 $BC$ 的距离为 $2$, 求 $GA$ 的长度。
    \ifprintanswers
    \begin{figure}[H]
        \centering
        \includegraphics[width=0.5\linewidth]{images/image57.png}
    \end{figure}
    \fi
    \begin{solution}
        设$GH \perp BC, BC = a,$则半周长为$5+\dfrac{a}{2}$,则
        \[
        [\triangle GBC]=\sqrt{\left(5+\frac{a}{2}\right)\left(\frac{a}{2}-2\right)\left(\frac{a}{2}+2\right)\left(5-\frac{a}{2}\right)} = \frac{1}{2} \cdot a \cdot 2
        \]
        解得
        \[
        a = 4 \sqrt{5} \quad \text{或} \quad a = 2 \sqrt{5}
        \]
        设$P$为$BC$中点,由中线定理,
        \[
        7^2+3^2=2\left(GP^2+\left(\frac{a}{2}\right)^2\right) \Rightarrow GP = 3\ \text{或} \ 2 \sqrt{6}
        \]
        故
        \[
        GA=2GP=6 \ \text{或} \ 4 \sqrt{6}
        \]
    \end{solution}

    \question 已知 $G$ 为 $\triangle ABC$ 的重心, $BC=10,AG=4,\angle BGC=\dfrac{3\pi}{4}$,求 $\triangle ABC$ 的面积。
    \ifprintanswers
    \begin{figure}[H]
        \centering        
        \includegraphics[width=0.5\textwidth]{images/image225.png}
    \end{figure}
    \fi
    \begin{solution}
        设直线 $AG$ 与 $BC$ 相交于 $D$。由于 $G$ 为重心,
        \[
        GD=\frac{AG}{2}=2,\quad BD=DC=\frac{10}{2}=5.
        \]
        在 $\triangle GBC$ 中,由中线定理,
        \[
        GB^2+GC^2=2(2^2+5^2)=58
        \]
        又由余弦定理,
        \[
        10^2=GB^2+GC^2-2\,GB\cdot GC \cos\frac{3\pi}{4} \Rightarrow GB\cdot GC=21\sqrt{2}
        \]
        于是$\triangle GBC$ 面积为
        \[
        [\triangle GBC]
        =\frac12\cdot 21\sqrt{2}\cdot\sin\frac{3\pi}{4}
        =\frac{21}{2}
        \]
        由等高性质,$\ \triangle ABC$ 面积为
        \[
        [\triangle ABC]=3[\triangle GBC]
        =\frac{63}{2}
        \]
    \end{solution}

    \question $I$ 为 $\triangle ABC$ 内心,$D$ 在 $AB$ 上,$E$ 在 $AC$ 上,且 $D,I,E$ 三点共线。已知 $AD=DE=5,AE=6$, 求 $AI$。
    \ifprintanswers
    \begin{figure}[H]
        \centering
        \includegraphics[width=0.4\linewidth]{images/image47.png}
    \end{figure}
    \fi
    \begin{solution}
        设 $\triangle ABC$ 内切圆半径为 $h$,由 $AD=DE$ 可知 
        \[
        \angle AED=\angle DAE=2\angle IAE=2\theta
        \]
        $\triangle ADE$面积为
        \[
        [\triangle ADE]=\sqrt{8\cdot3\cdot3\cdot2}=12=\frac{1}{2}\cdot 5h+\frac{1}{2}\cdot 6h \Rightarrow h=\frac{24}{11}
        \]
        在$\triangle ADE$ 中,由余弦定理,
        \[
        \cos 2\theta=\frac{6^2+5^2-5^2}{2\cdot 6\cdot 5}=\frac{3}{5} \Rightarrow \sin\theta=\sqrt{\frac{1-\cos 2\theta}{2}}=\sqrt{\frac{1}{5}}
        \]
        故
        \[
        AI=\frac{h}{\sin\theta}=\frac{24\sqrt{5}}{11}
        \]
    \end{solution}

    \question $\triangle ABC$满足$AB:AC:BC=3:5:7$。令$I$为$\triangle ABC$的内心, 直线$AI$交$\triangle ABC$外接圆于另一点$D$, 试求$\dfrac{ID}{AD}$的值。
    \ifprintanswers
    \begin{figure}[H]
        \centering
        \includegraphics[width=0.4\textwidth]{images/image33.png}
    \end{figure}
    \fi
    \begin{solution}
        设$E$为$AD$与$BC$的交点,因为$AD$为$\angle A$角平分线,由角平分线定理,
        \[
        \frac{BE}{EC} = \frac{AB}{AC} = \frac{3}{5} \Rightarrow 
        BE = \frac{21}{8}, CE= \frac{35}{8}
        \]
        在$\triangle ABC$ 中,由余弦定理,
        \[
        (7k)^2=(3k)^2 + (5k)^2 - 2 \cdot 3k \cdot 5k \cdot 
        \cos \angle A, k>0 \Rightarrow \angle A = 120^\circ, \angle BAE = 60^\circ
        \]
        在$\triangle ABE$ 中,由余弦定理,
        \[
        \left(\frac{21}{8}\right)^2=3^2 + AE^2-2 \cdot 3 \cdot AE \cdot \cos \angle BAE \Rightarrow AE = \frac{15}{8}
        \]
        又$\angle ABE = \angle ADC$,于是有$\triangle ABE \backsim \triangle CDE \ (AAA)$,则
        \[
        \frac{\frac{15}{8}}{\frac{21}{8}} = \frac{\frac{35}{8}}{DE} \Rightarrow DE = \frac{49}{8}, AD = \frac{15}{8} + \frac{49}{8} = 8
        \]
        又由角平分线定理,
        \[
        \frac{\frac{21}{8}}{3} = \frac{EI}{AI} \Rightarrow AI = \frac{8}{15} \cdot AE = 1
        \]
        故
        \[
        \frac{ID}{AD} = \frac{8 - 1}{8} = \frac{7}{8}
        \]
    \end{solution}

    \question 在 $\triangle ABC$ 的边 $AB$ 与 $AC$ 的外侧分别作正三角形 $\triangle ABE$ 及 $\triangle ACF$。已知 $AC = 1$ 且 $EF = 2$,求 $\triangle ABC$ 面积的最大可能值。
    \begin{solution}
        设 $\angle BAC = \theta,AB = c$,则 $\angle EAF = 240^\circ - \theta$,在$\triangle AEF$中,由余弦定理,
        \[
        2^2 = c^2 + 1^2 - 2 c \cdot 1 \cdot \cos(240^\circ - \theta)
        \Rightarrow c^2 + c \cos \theta + \sqrt{3} c \sin \theta =3
        \]
        设 $x = c \cos \theta$, $y = c \sin \theta$,上式变为
        \[
        x^2 + y^2 + x + \sqrt{3} y = 3 \Rightarrow \left( x + \frac12 \right)^2 + \left( y + \frac{\sqrt{3}}{2} \right)^2 = 4
        \]
        $\triangle ABC$的面积为
        \[
        \frac12 \cdot 1 \cdot c \sin \theta = \frac{y}{2} \le \frac12 \left( 2 - \frac{\sqrt{3}}{2} \right) = 1 - \frac{\sqrt{3}}{4}
        \]
        等号成立当且仅当$x=-\dfrac{1}{2},y=2-\dfrac{\sqrt{3}}{2}$,即
        \[
        c=\sqrt{x^2+y^2}=\sqrt{\frac{17}{4}-2\sqrt{3}},
        \]
        \[
        \theta=\arctan\frac{y}{x}=\pi-\arctan(4-\sqrt{3})
        \]
    \end{solution}

    \question 如图,点 $M$ 是长方形 $ABCD$ 的边 $AD$ 的中点,$\triangle BME$ 是正三角形,$N$ 为线段 $EC$ 的中点。已知 $\angle AMB = 48^\circ$,${BM} = 1$,求 ${DN}$ 的长。 
        \begin{figure}[H]
            \centering
            \includegraphics[width=0.4\textwidth]{images/image26.png}
        \end{figure}
    \begin{solution}
        设$AB = a = \sin 48^\circ,MA = MB = b = \sin 42^\circ=\cos 48^\circ,$取
        \[A(0,a), B(0,0), C(2b,0),
        D(2b,a), E(\cos 12^\circ, -\sin 12^\circ),
        \]
        则$CE$中点 $N$ 为
        \[
        N\left( b + \frac{1}{2}\cos 12^\circ,\ -\frac{1}{2}\sin 12^\circ \right)
        \]
        故
        \[
        \begin{aligned}
        {DN}^2 &= \left(b - \frac{1}{2} \cos 12^\circ\right)^2 + \left(a + \frac{1}{2} \sin 12^\circ\right)^2 \\
        &= a \sin 12^\circ - b \cos 12^\circ + \frac{1}{4} (\cos^2 12^\circ + \sin^2 12^\circ) + a^2 + b^2 \\
        &= -\cos 60^\circ + \frac{1}{4} + 1  = \frac{3}{4} \Rightarrow DN = \frac{\sqrt{3}}{2}
        \end{aligned}
        \]
    \end{solution}

    \question 已知 $AC$ 为半圆的直径,若将弧 $AB$ 沿弦 $AB$ 向下折使其与 $AC$ 相交于 $D$ 点,且 $AD=9,DC=7$,则 $AB$ 的长度为多少？  
    \ifprintanswers
    \begin{figure}[H]
        \centering        
        \includegraphics[width=0.5\textwidth]{images/image80.jpg}
    \end{figure}
    \fi
    \begin{solution}
        作 $D'D'' \parallels BC$ 交圆于 $D',D''$ 及直线 $AB$ 于 $P$ 点,
        $AB$ 为 $DD'$ 的中垂线,
        \[
        AD' = AD = 9,
        \]
        在$\triangle AOD'$ 中,由余弦定理,
        \[
        \cos \angle AOD' = \frac{8^2 + 8^2 - 9^2}{2 \cdot 8 \cdot 8} = \frac{47}{128}
        \] 
        在$\triangle DOD'$ 中,由余弦定理,
        \[ 
        \cos \angle D'OD = -\frac{47}{128} = \frac{8^2 + 1^2 - DD'^2}{2 \cdot 8 \cdot 1} \Rightarrow DD' = \frac{9\sqrt{14}}{4}  ,PD = PD' = \frac{9\sqrt{14}}{8} 
        \]
        由圆幂定理,
        \[
        9 \cdot 7 = \frac{9\sqrt{14}}{4}  \cdot DD'' \Rightarrow DD'' = 2\sqrt{14}
        \]
        又 $\triangle APD \backsim \triangle ABC \ (AAA)$,
        \[
        \frac{AP}{PB} = \frac{AD}{DC} = \frac{9}{7}
        \]
        令 $AP = 9k,PB = 7k,k > 0$,由圆幂定理,
        \[
        \frac{9\sqrt{14}}{8}  \cdot \left( \frac{9\sqrt{14}}{8}  + 2\sqrt{14} \right) = 9k \cdot 7k \Rightarrow k = \frac{5\sqrt{2}}{8}
        \] 
        故
        \[
        AB = 16k = 10\sqrt{2}
        \]
    \end{solution}

    \question 如图,在正$\triangle ABC$中,将顶点$A$折至$D$,使得$D$在$BC$上,若$BD=1,CD=2$,求$EF$之长。
    \begin{figure}[H]
        \centering
        \includegraphics[width=0.4\linewidth]{images/image104.png}
    \end{figure}
    \ifprintanswers
    \begin{figure}[H]
        \centering
        \includegraphics[width=0.4\linewidth]{images/image105.jpg}
    \end{figure}
    \fi
    \begin{solution}
        设
        \[
        EA = ED = a, FA = FD = b, \angle EDA = \angle A = 60^\circ,
        \]
        在$\triangle BDE$ 及$\triangle CDF$中,由余弦定理,
        \[
        a^2=(3-a)^2 + 1-2(3-a)\cos 60^\circ ,\quad
        b^2=(3-b)^2 + 4-4(3-b)\cos 60^\circ
        \]
        解得
        \[
        a = \frac75,\quad b = \frac74
        \]
        同理,在$\triangle EDF$中,由余弦定理,
        \[
        EF^2 = a^2 + b^2 -2ab \cos 60^\circ \Rightarrow EF = \frac{7\sqrt{21}}{20}
        \]
    \end{solution}

    \question 如下图,将长 $AB=240$、宽 $BC=288$ 的长方形纸张对折,让顶点 $C$ 刚好落在 $AB$ 的中点 $M$ 上。若 $EF$ 是折线,求折线 $EF$ 的长度。
    \begin{figure}[H]
        \centering
        \includegraphics[width=0.35\linewidth]{images/image48.png}
    \end{figure}
    \ifprintanswers
    \begin{figure}[H]
        \centering
        \includegraphics[width=0.45\linewidth]{images/image139.png}
    \end{figure}
    \fi
    \begin{solution}
        如上图,设 $\angle EFC=\angle EFM=\theta$,则 
        \[
        \angle MFB=180^{\circ}-2\theta,\angle FMB=2\theta-90^{\circ}
        \]
        过 $E$ 点作水平线交 $BC$ 于 $G$ 点,由直角三角形$EFG$得 $EF = \dfrac{240}{\sin\theta}$,又  
        \[
        288 = BC = BF + FC = BF + FM = \frac{120}{\cos(2\theta-90^{\circ})} + 120\tan(2\theta-90^{\circ})
        \]
        化简得
        \[
        288 = \frac{120}{\sin 2\theta} - 120\cot 2\theta
        = 120\cdot \frac{1-\cos 2\theta}{\sin 2\theta}
        = 120\cdot \frac{2\sin^{2}\theta}{2\sin\theta\cos\theta}
        = 120\tan\theta \Rightarrow \tan\theta = \frac{12}{5}
        \]
        故 
        \[
        \sin\theta = \frac{12}{13} \Rightarrow EF = 240 \cdot \frac{13}{12} = 260
        \]
    \end{solution}
    \begin{solution}
        设 $MF = FC = a$,在直角 $\triangle BMF$ 中,由毕氏定理,
        \[
        a^2 = 120^2 + (288 - a)^2 \Rightarrow a = 169
        \]
        设 $DE = D'E = b$,在直角 $\triangle AEM$ 及$\triangle CDE$中,由毕氏定理,
        \[
        ME^2 = 120^2 + (288 - b)^2, \quad EC^2 = b^2 + 240^2
        \]
        由 $EM = EC$ 解得
        \[
        b = 69
        \]
        在直角 $\triangle EFG$ 中,由毕氏定理,
        \[
        EF^2 = 240^2 + (a - b)^2 = 260^2,
        \]
        故折线长度为 $EF = 260$。
    \end{solution}

    \question 如图,一个平行四边形 $PQRS$ 沿着线段 $MN$ 将平行四边形对折,恰好让 $P$ 点与 $R$ 点重合,且让 $Q$ 点坐落于 $Q'$ 的位置,而 $M,N$ 分别在 $PS,QR$ 上。若 $PQ=6,PM = 4\sqrt{3},\cos Q = -\dfrac{\sqrt{11}}{4}$,求折痕长 $MN$。
    \begin{figure}[H]
        \centering
        \includegraphics[width=0.6\textwidth]{images/image29.png}
    \end{figure}
    \begin{solution}
        设$\angle SMR = \angle MRQ = \alpha, SR = PQ = 6, \angle PMN = \angle MNR = \theta,$

        又 $MN$ 为折线,所以有 $MR = MP = 4\sqrt{3},\angle RMN = \angle PMN = \theta,$

        在 $\triangle MSR$ 中,由正弦定理,
        \[
        \frac{6}{\sin \alpha} = \frac{4\sqrt{3}}{1-\cos^2 Q} \Rightarrow \sin \alpha = \frac{\sqrt{15}}{8}
        \]
        于是
        \[
        \cos \alpha = \frac{7}{8}, \cos \frac{\alpha}{2} = \frac{\sqrt{15}}{4}
        \]
        又 $2\theta + \alpha = 180^\circ$,所以
        \[
        \sin \theta = \sin\left(90^\circ - \frac{\alpha}{2}\right) = \cos \frac{\alpha}{2} = \frac{\sqrt{15}}{4}
        \]
        在 $\triangle RMN$ 中,由正弦定理,
        \[
        \frac{MN}{\frac{\sqrt{15}}{8}} = \frac{4\sqrt{3}}{\frac{\sqrt{15}}{4}} \Rightarrow 
        MN = 2\sqrt{3}
        \]
    \end{solution}
    
    \question 以两种方式将矩形 $PQRS$放置于矩形 $ABCD$内:方式一为 $Q,R$ 分别在 $B,C$上,方式二为 $P,Q,R,S$ 分别在 $AB,BC,CD,DA$ 上。若 $AB=718,PQ=250$,求 $BC$ 的长度。
    \begin{figure}[H]
        \centering        
        \includegraphics[width=0.6\textwidth]{images/image186.png}
    \end{figure}
    \begin{solution}
        已知$AB=718,PQ=250$,在图二中,设 $BC=x,BQ=a,PB=b$,则
        \[
        AD=PS=QR=x,\quad QC=x-a,\quad AP=718-b
        \]
        由于$\triangle RDS \congr \triangle PBQ \ \text{(ASA)},\triangle SAP \congr \triangle QCR \ \text{(ASA)}$,因此
        \[
        \dfrac{SA}{PB}=\dfrac{AP}{BQ}=\dfrac{SP}{PQ} \Rightarrow
        \dfrac{x-a}{b}=\dfrac{718-b}{a}=\dfrac{x}{250} \tag{1}
        \]
        在 $\triangle PBQ$ 及 $\triangle SAP$ 中,由毕氏定理,
        \[
        a^2+b^2=250^2, \quad (x-a)^2+(718-b)^2=x^2 \tag{2,3}
        \]
        由$(1),(2),(3)$解得
        \[
        a=88,b=234,BC=x=1375
        \]
    \end{solution}

    \question 正方形 $PQRS$ 的边长为 $4$。点 $U$ 在对角线 $PR$ 上使得$PR=4UR$。以 $U$ 为圆心作圆,使得圆与正方形的两条边相切。$PW$ 是圆的切线,且 $W$ 在 $QR$ 上。求 $PW$ 的长度。
    \ifprintanswers
    \begin{figure}[H]
        \centering        
        \includegraphics[width=0.5\textwidth]{images/image214.png}
    \end{figure}
    \fi
    \begin{solution}
        由 $PR=4UR=4\sqrt{2}$,得
        \[
        PU=\frac{3}{4}PR=3\sqrt{2}, \quad UR=\sqrt{2}.
        \]
        设圆分别与 $PW,WR, RS$ 相切于点 $X,Y,Z$,由$\triangle UZR$知圆的半径为 $1$,在$\triangle PXU$中,由毕氏定理,
        \[
        PX=\sqrt{(3\sqrt{2})^2-1^2}=\sqrt{17}
        \]
        设$x=PW$,在$\triangle PQW$中,由毕氏定理,
        \[
        4^2+(3-x)^2=(\sqrt{17}+x)^2 \Rightarrow x=\frac{4}{\sqrt{17}+1}
        \]
        故
        \[
        PW=\sqrt{17}+\frac{4}{\sqrt{17}+1}=\frac{5\sqrt{17}-1}{4}
        \]
    \end{solution}

    \question 如下图所示, $AB=8$, 以 $AB$ 为直径的半圆上有 $C, D$ 两点, 且 $AC=2,BD=7$, 求 $CD$ 的长度。 
    \begin{figure}[H]
        \centering        
        \includegraphics[width=0.5\textwidth]{images/image99.png}
    \end{figure}
    \ifprintanswers
    \begin{figure}[H]
        \centering        
        \includegraphics[width=0.5\textwidth]{images/image100.jpg}
    \end{figure}
    \fi
    \begin{solution}
        $AB$ 为直径,在 $\triangle ABC$ 及 $\triangle ABD$中,由毕氏定理,
        \[
        BC=\sqrt{8^2-2^2}=2\sqrt{15}, AD=\sqrt{8^2-7^2}=\sqrt{15}
        \]
        设$\theta_1=\angle ABC,\theta_2=\angle ABD,$则
        \[
        \tan \theta_1 =\frac{1}{\sqrt{15}},\quad
        \tan (\theta_1+\theta_2) = \frac{\frac{1}{\sqrt{15}} + \tan \theta_2}{1 - \frac{1}{\sqrt{15}} \tan \theta_2} = \frac{\sqrt{15}}{7} \Rightarrow \tan \theta_2 = \frac{1}{\sqrt{15}}
        \]
        在 $\triangle BCD$中,由余弦定理,
        \[
        \cos \theta_2 = {\sqrt{15}\over 4} 
        = {(2\sqrt{15})^2 + 7^2 - CD^2 \over 2\cdot 7 \cdot 2\sqrt{15}} \Rightarrow CD = 2
        \]
    \end{solution}

    \question 三角形 $ABC$ 内接于一直径为 $5\sqrt{13}$ 的圆, $D$ 点在 $AC$ 上, 且 $\angle ABD = 90^\circ$. 已知 $BC = 10,AD = 13$, 求 $CD$。
    \ifprintanswers
    \begin{figure}[H]
        \centering        
        \includegraphics[width=0.5\textwidth]{images/image87.jpg}
    \end{figure}
    \fi
    \begin{solution}
        由于 $\angle ABD$ 为直角, 延长 $BD$ 交圆周于 $B'$, 则 $AB'$ 为直径,由$\triangle DBC \backsim \triangle DAB' \ (\text{AAA})$,
        \[
        \frac{DB}{DA} = \frac{BC}{AB'} \Rightarrow \frac{DB}{13} = \frac{10}{5\sqrt{13}} 
        \Rightarrow DB = 2\sqrt{13}
        \]
        在 $\triangle ABD$ 中,由毕氏定理,
        \[
        AB^2 + (2\sqrt{13})^2 = 13^2 \Rightarrow AB = 3\sqrt{13}
        \]
        在 $\triangle ABB'$ 中,由毕氏定理,
        \[
        BB' = \sqrt{(5\sqrt{13})^2 - (3\sqrt{13})^2} = 4\sqrt{13} \Rightarrow DB' = BB' - BD = 2\sqrt{13}
        \]
        再由 $\triangle DBC \backsim \triangle DAB'$,
        \[
        \frac{BC}{AB'} = \frac{CD}{DB'} \Rightarrow \frac{10}{5\sqrt{13}} = \frac{CD}{2\sqrt{13}} \Rightarrow CD = 4
        \]
    \end{solution}

    \question 已知圆内接六边形 $ABCDEF$ 的边长依次为 $5,5,7,7,5,7$,求该六边形的面积。
    \ifprintanswers
    \begin{figure}[H]
        \centering
        \includegraphics[width=0.5\linewidth]{images/image78.png}
    \end{figure}
    \fi
    \begin{solution}
        设圆心为 $O$,则
        \[
        \angle AOB : \angle BOC : \angle COD : \angle DOE : \angle EOF : \angle FOA = 5:5:7:7:5:7
        \]
        即
        \[
        \angle AOB = \angle BOC = \angle EOF = 50^\circ,\quad \angle COD = \angle DOE = \angle FOA = 70^\circ \Rightarrow \angle BAF = 120^\circ
        \]
        在$\triangle BAF$ 中,由余弦定理,
        \[
        BF^2 =5^2 + 7^2 - 2\cdot 5\cdot 7 \cdot \cos 120^\circ= 109
        \]
        六边形面积为
        \[
        [ABCDEF]= [\triangle BDF] + 3 [\triangle ABF]
        = \frac{\sqrt{3}}{4} \cdot 109 + \frac{3}{2} \cdot 5 \cdot 7 \cdot \sin 120^\circ
        = \frac{107\sqrt{3}}{2}
        \]
    \end{solution}

    \question 已知$\triangle ABC$中,$AB=25,AC=30,\angle B=2\angle C$, 且$\angle B$的内角平分线交$\triangle ABC$的外接圆于$D$点,且$B,D$为相异点。试求$\triangle BCD$的面积。 
    \ifprintanswers
    \begin{figure}[H]
        \centering
        \includegraphics[width=0.5\linewidth]{images/image70.png}
    \end{figure}
    \fi
    \begin{solution}
        设$\angle C = \theta$,则$\angle B = 2\theta$,$\angle A = \pi - 3\theta$, 在$\triangle ABC$ 中,由正弦定理,
        \[
        \frac{25}{\sin \theta} = \frac{30}{\sin 2\theta} \implies \cos \theta = \frac{3}{5}.
        \]
        因同弧所对圆周角相等,$\angle D = \angle A,\angle B = \angle DCB,$故$\triangle ABC \cong \triangle DCB$\ (ASA),因此
        \[
        [\triangle BCD] = [\triangle ABC] = \frac{1}{2} \cdot 25 \cdot 30 \cdot \sin(\pi - 3\theta) = 375 \sin 3\theta.
        \]
        由$\sin \theta = \dfrac{4}{5}$,
        \[
        [\triangle BCD] = 375 \left(3 \cdot \frac{4}{5} - 4 \cdot \left(\frac{4}{5}\right)^3 \right) = 132
        \]
    \end{solution}

    \question 半圆 $\Gamma$ 的直径 $AB$ 之长为 $14$, 圆 $\Omega$ 切 $AB$ 于点 $P$ 且交 $\Gamma$ 于点 $Q$ 与点 $R$。若 $QR=3\sqrt{3}$ 且 $\angle QPR=60^{\circ}$, 求 $\triangle PQR$ 的面积。
    \ifprintanswers
    \begin{figure}[H]
        \centering
        \includegraphics[width=0.5\linewidth]{images/image59.png}
    \end{figure}
    \fi
    \begin{solution}
        设半圆心为 $C$, 小圆心为 $O$, 小圆半径为 $r$。由$\angle QPR=60^\circ$ 得 $\angle QOR=120^\circ$,在$\triangle ABC$ 中,由正弦定理,
        \[
        \frac{QR}{\sin\angle QOR}=\frac{3\sqrt3}{\sin 120^\circ}=2r \Rightarrow OQ=OR=r=3
        \]
        且发现$\triangle COR \congr \triangle COQ$ \ (SSS),故 $\angle COQ=\angle COR=120^\circ$。在$\triangle COQ$中,由余弦定理,
        \[
        7^2 = CO^2+3^2-2\cdot CO\cdot 3 \cdot \cos 120^\circ \Rightarrow CO=5
        \]
        在$\triangle COP$中,由毕氏定理,$CP=\sqrt{5^2-3^2}=4$,设 $\angle COP=\theta$, 有
        \[
        \sin\theta=\frac{4}{5},\quad \cos\theta=\frac{3}{5}
        \]
        因此
        \begin{align*}  
        [\triangle PQR]
        &=[\triangle POQ]+[\triangle POR]+[\triangle OQR]\\
        &=\frac12\cdot 3^2\bigl(\sin(120^\circ-\theta)+\sin(120^\circ+\theta)\bigr)+\frac12\cdot 3^2\sin 120^\circ
        =\frac{99\sqrt3}{20}
        \end{align*}
    \end{solution}

    \question 已知圆内接四边形 $ABCD$ 中可以作一个内切圆,且 $E,F,G,H$ 分别在$AB,BC,CD,DA$上,且为此四边形 $ABCD$ 与其内切圆相切的四个切点,若 $FG=6,EF=7,EH=8$,求 $GH$ 的长度。
    \ifprintanswers
    \begin{figure}[H]
        \centering
        \includegraphics[width=0.4\textwidth]{images/image23.png}
    \end{figure}
    \fi
    \begin{solution}
        由圆内接四边形性质有
        \[
        \angle A + \angle C =\angle GOF + \angle C = 180^\circ 
        \Rightarrow \angle A = \angle GOF
        \]
        所以
        \[
        \triangle AHE \backsim \triangle OGF \ \text{(AAA)} \Rightarrow {AH} = {AE} = \dfrac{4}{3}r
        \]
        在$\triangle AHE$及$\triangle OHE$中,由余弦定理,
        \[
        \cos \angle A = -\cos \angle HOE \Rightarrow \frac{2r^2 - 8^2}{2r^2}= -\frac{\left(\frac{4r}{3}\right)^2 + \left(\frac{4r}{3}\right)^2 - 8^2}{2 \cdot \left(\frac{4r}{3}\right)^2}\Rightarrow r = 5
        \]
        同理可得$\angle D = \angle EOF$,在$\triangle OEF$中,由余弦定理,
        \[
        \cos \angle D = \cos \angle EOF = \frac{2r^2 - 7^2}{2r^2}  = \frac{1}{50}
        \]
        又 $\cos \angle GOH = -\cos \angle D = -\dfrac{1}{50}$
        ,在$\triangle OGH$中,由余弦定理,
        \[
        \cos \angle GOH = \frac{2r^2 - {GH}^2}{2r^2}
        \Rightarrow {GH} =\sqrt{51}
        \]
    \end{solution}

    \question 若圆 $O_1$ 与圆 $O_2$ 的半径比为 $2:3$,且圆 $O_1$ 与圆 $O_2$ 交于 $A,B$ 两点。过 $B$ 点作一直线分别交圆 $O_1$ 及圆 $O_2$ 于 $C,D$ 两点,且 $\angle CAD=\dfrac{2\pi}{3}$,求 $\tan \angle ACD$。 
    \ifprintanswers
    \begin{figure}[H]
    \centering        
    \includegraphics[width=0.7\textwidth]{images/image83.png}
    \end{figure}
    \fi
    \begin{solution}
        调整两圆心距使得  $C$ 在 $O_1O_2$ 上,令 $\angle ACD=\theta$,则
        \[
        \angle AO_2O_1 = \frac{1}{2}\angle AO_2B = \angle ADC = 60^\circ-\theta
        \]
        又 $\angle ACB=180^\circ-\theta$,所以
        \[
        \angle AO_1B = 360^\circ-2(180^\circ-\theta)=2\theta \Rightarrow \angle AO_1O_2=\theta
        \]
        在$\triangle APO_1$及$\triangle APO_2$中,设$r_1=2k,r_2=3k,$则
        \[
        \sin \angle AO_1O_2 = \sin \theta = \frac{AP}{2k}, \sin \angle AO_2O_1 = \sin (60^\circ-\theta) = \frac{AP}{3k}
        \]
        解得
        \[
        2\sin \theta = 3\left(\frac{\sqrt{3}}{2}\cos\theta - \frac{1}{2}\sin\theta\right)
        \Rightarrow \tan \theta = \frac{3\sqrt{3}}{7}
        \]
    \end{solution}
    
    \question 已知 $\triangle ABC$ 内接于半径为 $1$ 的圆。对于$\triangle ABC$的每一条边,作一圆使得与该边切于其中点,且该圆与大圆相切。已知其中两个小圆的半径分别为 $\dfrac{2}{3},\dfrac{2}{11}$,求第三个小圆的半径。
    \ifprintanswers
    \begin{figure}[H]
        \centering
        \includegraphics[width=0.5\textwidth]{images/image156.png}
    \end{figure}
    \fi
    \begin{solution}
        设大圆圆心为 $O,BC,AC,AB$中点为 $D,E,F$,对应小圆半径为 \[
        r_1=\frac{2}{3},r_2=\frac{2}{11},r_3
        \]  
        在$\triangle BDO$中,由毕氏定理,
        \[
        BC^2=4(OB^2-OD^2)
        =4\Bigl(1-(1-r_1)^2\Bigr)=\frac{32}{9}
        \]
        在$\triangle CEO$中,由毕氏定理,
        \[
        AC^2=4(OC^2-OE^2)
        =4\Bigl(1-(1-r_2)^2\Bigr)=\frac{288}{121}
        \]
        设 $\angle BOC=\beta,\ \angle COA=\gamma$,在$\triangle BOC$及$\triangle AOC$中,由余弦定理,
        \[
        \cos\beta=\frac{OB^2+OC^2-BC^2}{2OB\cdot OC}=-\frac{7}{9},\quad
        \cos\gamma=\frac{OC^2+OA^2-AC^2}{2OC\cdot OA}=-\frac{23}{121}.
        \]
        由此得
        \[
        \sin\beta=\frac{4\sqrt2}{9},\quad \sin\gamma=\frac{84\sqrt2}{121}.
        \]
        于是
        \[
        \cos\angle AOB=\cos(\beta-\gamma)=\cos\beta\cos\gamma+\sin\beta\sin\gamma
        =\frac{833}{1089}.
        \]
        在$\triangle AOB$中,由余弦定理,
        \[
        AB^2=OA^2+OB^2-2OA\cdot OB\cos\angle AOB=\frac{512}{1089}
        \]
        因此
        \[
        OF=\sqrt{OA^2-\frac14 AB^2}=\frac{31}{33} \Rightarrow r_3=\frac12(1-OF)=\frac{1}{33}
        \]
    \end{solution}

    \question $P,Q$ 为正方形 $ABCD$ 内两点使得 $DP \parallels BQ$ 且 $DP = PQ = BQ$,求 $\angle ADP$ 的最小可能值。
    \begin{figure}[H]
        \centering
        \includegraphics[width=0.3\textwidth]{images/image39.png}
    \end{figure}
    \ifprintanswers
    \begin{figure}[H]
        \centering
        \includegraphics[width=0.35\textwidth]{images/image40.png}
    \end{figure}
    \fi
    \begin{solution}
        设$DP = PQ = QB = a,\angle DPR = \angle QBC = \theta, \angle RPQ = \alpha$,则:
        \[
        CD = DR + MQ + QN = 2a \sin \theta + a \sin \alpha
        \]
        \[
        BC = BN + BN - PM = 2a \cos \theta - a \cos \alpha
        \]
        由$CD = BC$得
        \[
        2(\sin \theta - \cos \theta) = -(\sin \alpha + \cos \alpha)
        \]
        两边平方得
        \[
        4 - 4 \sin 2\theta = 1 + \sin 2\alpha
        \Rightarrow \sin 2\theta = \frac{3 - \sin 2\alpha}{4}
        \]
        由 $\sin 2\alpha \in [-1, 1]$,
        \[
        \frac{1}{2} \le \sin 2\theta \le 1
        \Rightarrow 15^\circ \le \theta \le 45^\circ
        \]
        因此$\theta$ 的最小值为 $15^\circ$。
    \end{solution}
    \ifprintanswers
    \begin{figure}[H]
        \centering
        \includegraphics[width=0.35\textwidth]{images/image190.png}
    \end{figure}
    \begin{figure}[H]
        \centering
        \includegraphics[width=0.35\textwidth]{images/image191.png}
    \end{figure}
    \fi
    \begin{solution}
        以$D$为原点构造坐标系,设 $P(a,b),Q(2-a,2-b)$,由 $PD=PQ$ 得
        \[
        a^2+b^2=(2-2a)^2+(2-2b)^2 \Rightarrow \left(a-\frac{4}{3}\right)^2+\left(b-\frac{4}{3}\right)^2=\frac{8}{9}
        \]
        因此 $P$ 在圆心 $\left(\dfrac{4}{3},\dfrac{4}{3}\right)$,半径 $\dfrac{2\sqrt{2}}{3}$ 的圆上。当 $DP$ 与此圆相切时,$\theta=\angle ADP$ 取最小值, 由$\triangle DOP$,
        \[
        \sin(45^\circ-\theta)=\frac{OP}{OD}=\frac{1}{2}
        \]
        此时 $45^\circ-\theta=30^\circ$,即$\theta=15^\circ$。
    \end{solution}
    \ifprintanswers
    \begin{figure}[H]
        \centering
        \includegraphics[width=0.35\textwidth]{images/image192.png}
    \end{figure}
    \fi
    \begin{solution}
        设正方形$ABCD$边长为 $1$,令 $\angle ADP=\theta,PD=a,PB=b$。在 $\triangle PDB$ 中,由余弦定理,
        \[
        \cos(45^\circ-\theta)=\frac{a^2+2-b^2}{2\sqrt{2}a} \tag{1}
        \]
        在 $\triangle PDB$ 中,由余弦定理,
        \[
        \cos(45^\circ-\theta)=\frac{a^2+\left(\frac{\sqrt2}{2}\right)^2-\left(\frac{a}{2}\right)^2}{2a\left(\frac{\sqrt2}{2}\right)}=\frac{\frac{3}{4}a^2+\frac{1}{2}}{\sqrt{2}a} \tag{2}
        \]
        由$(1),(2)$得
        \[
        b^2=\frac{1}{2}(2-a^2)
        \]
        代入$(1)$得
        \[
        \cos(45^\circ-\theta)=\frac{a^2+2-\frac{1}{2}(2-a^2)}{2\sqrt{2}a} = \frac{1}{4\sqrt{2}}\left(3a+\frac{2}{a}\right) \ge \frac{1}{4\sqrt{2}} \cdot 2\sqrt{6}=\frac{\sqrt{3}}{2},
        \]
        等号成立当且仅当$a^2=\frac{2}{3}$,当$\cos(45^\circ-\theta)$为最小时,$\theta=15^\circ$为最小。
    \end{solution}
    \ifprintanswers
    \begin{figure}[H]
        \centering
        \includegraphics[width=0.35\textwidth]{images/image193.png}
    \end{figure}
    \fi
    \begin{solution}
        连接 $BD$交 $PQ$ 于 $M$,由于$\triangle PDM \congr \triangle QBM \ \text{(ASA)}$,有$DM=BM$,即$M$为$BD$中点。不失一般性,设 $PM=1,PD=2$,记 $\angle ADP=\theta,\angle PDM=\alpha$,则由于 $\theta+ \alpha=45^\circ$,当$\alpha$取最大时$\theta$取最小。
        在$\triangle PMD$中,由正弦定理,
        \[
        \frac{\sin \alpha}{1}=\frac{\sin \angle PMD}{2}
        \]
        欲使$\sin \alpha$最大,取$\sin \angle PMD=1$,此时$\sin \alpha = \dfrac{1}{2}$,即 $\alpha = 30^\circ$,所以 $\theta$ 的最小值为 $15^\circ$。  
    \end{solution}
    \ifprintanswers
    \begin{figure}[H]
        \centering
        \includegraphics[width=0.35\textwidth]{images/image194.png}
    \end{figure}
    \fi
    \begin{solution}
        以$D$为原点构造坐标系,设 $PD=PQ=QB=a,\angle ADP=\theta$,则$P(a\sin\theta,a\cos\theta),Q(1-a\sin\theta,1-a\cos\theta)$,且有
        \[
        PQ^2=(1-2a\sin\theta)^2+(1-2a\cos\theta)^2=a^2,
        \]
        化简得
        \[
        \frac{2+3a^2}{4a}=\sin\theta+\cos\theta=\sqrt{2}\sin(\theta+45^\circ)
        \]
        从而
        \[
        \sin(\theta+45^\circ)=\frac{2+3a^2}{4\sqrt{2}a}=\frac{1}{2\sqrt{2}a}+\frac{3a}{4\sqrt{2}}\ge \frac{\sqrt{3}}{2},
        \]
        由于$0^\circ \le \theta \le 90^\circ,\sin(\theta+45^\circ)$在$\theta=15^\circ$有最小,故$\theta=15^\circ$为最小。
    \end{solution}

    \question 右图为月偏食的示意图,满月被地球的影锥遮蔽一部分。假设满月和地球影锥截面都是正圆,在图中标记圆弧 $\wideparen{AP} = \wideparen{PB}$ 均为满月的圆周一部分,圆弧 $\wideparen{AQ} = \wideparen{QB}$ 均为地球影锥截圆的圆周一部分。若 $AB$ 与 $PQ$ 分别是满月圆半径的 $\sqrt{3}$ 与 $2-\sqrt{3}$ 倍,求 
    \[
    \frac{\text{月偏食亮面面积}}{\text{满月圆面积}}.
    \]
    \begin{figure}[H]
        \centering
        \includegraphics[width=0.35\textwidth]{images/image36.png}
    \end{figure}
    \ifprintanswers 
    \begin{figure}[H]
        \centering
        \includegraphics[width=0.5\textwidth]{images/image37.png}
    \end{figure}
    \fi
    \begin{solution}
        设小圆(月亮)圆心为 $O_1$,半径$r = 1$ ,大圆(太阳)圆心为$O_2$,半径$R$,$C$为$O_1Q$ 与$AB$的交点,则
        \[
        BC = \frac{1}{2} AB = \frac{1}{2} \sqrt{3} \Rightarrow \angle CO_1B = 60^\circ,\angle O_1BC = 30^\circ
        \]
        又 $CQ = O_1P - PQ - O_1C = \sqrt{3} - \frac{3}{2}$,
        \[
        \tan \angle QBC = \frac{CQ}{BC} = 2 - \sqrt{3} \Rightarrow \angle QBC = 15^\circ,\angle CQB = 75^\circ
        \]
        由于 $O_2B = O_2Q = R$,因此
        \[
        \angle BO_2C = 180^\circ - 75^\circ \cdot 2 = 30^\circ \Rightarrow R = 2 BC = \sqrt{3}
        \]
        故弓形$AQB$面积为
        \[
        [\text{弓形}AQB] 
        = [\text{扇形}AQBO_2]-[\triangle ABO_2] 
        = \frac{1}{6} (\sqrt{3})^2 \pi - \frac{1}{2} (\sqrt{3})^2\sin 60^\circ
        =\frac{\pi}{2} - \frac{3\sqrt{3}}{4}
        \]
        同理,
        \[
        [\text{弓形}APB]
        = [\text{扇形}APBO_1]-[\triangle ABO_1] 
        = \frac{1}{3} (1)^2 \pi - \frac{1}{2} (1)^2 \sin 120^\circ
        =\frac{\pi}{3} - \frac{\sqrt{3}}{4}
        \]
        故
        \[
        \frac{\text{月偏食亮面面积}}{\text{满月圆面积}} = \frac{\left( \frac{\pi}{3} - \frac{\sqrt{3}}{4} \right) - \left( \frac{\pi}{2} - \frac{3\sqrt{3}}{4} \right)}{\pi(1)^2}=\frac{\sqrt{3}}{2\pi} - \frac{1}{6}
        \]
    \end{solution}

    \question 下图中,已知 $\triangle ABC$ 为正三角形,$DEFGHIJK$ 为正八边形,且 $E$ 为 $BC$ 上一点,$CE=2,A,C,D$ 三点共线,$A,B,F,G$ 四点共线,求 $AF$。
    \begin{figure}[H]
        \centering
        \includegraphics[width=0.5\textwidth]{images/image125.png}
    \end{figure}
    \begin{solution}
        正八边形内角为$(8-2)\times 180^\circ \div 8=135^\circ$,则
        \[
        \angle EFB = 45^\circ, \angle BEF = 15^\circ, \angle CED = 30^\circ, \angle CDE = 30^\circ \Rightarrow CD = CE = 2
        \]
        在 $\triangle CDE$ 中,由余弦定理,
        \[
        DE^2=2^2 + 2^2 - 2\cdot2\cdot2\cdot\cos 120^\circ \Rightarrow DE = 2\sqrt{3}
        \]
        在 $\triangle BEF$ 中,由正弦定理,
        \[
        \frac{2\sqrt{3}}{\sin 120^\circ} = \frac{EB}{\sin 45^\circ} = \frac{BF}{\sin 15^\circ} \Rightarrow EB = 2\sqrt{2}, BF = \sqrt{6} - \sqrt{2},
        \]
        其中$\sin 15^\circ=\dfrac{\sqrt{6}-\sqrt{2}}{4}$,因此
        \[
        AF = AB + BF = CE + EB + BF = 2 + \sqrt{6} + \sqrt{2}
        \]
    \end{solution}

    \question 某圆内接六边形 $ABCDEF$,其中 $AB=BC=CD=1,DE=EF=FA=2$,求此六边形的面积。
    \ifprintanswers
    \begin{figure}[H]
        \centering
        \includegraphics[width=0.5\textwidth]{images/image136.jpg}
    \end{figure}
    \fi
    \begin{solution}
        设圆心为 $O$,半径为$r$,且
        \[
        \angle AOB = \angle BOC = \angle COD = b, \quad \angle DOE = \angle EOF = \angle FOA = a, 
        \]
        由圆周角和
        \[
        3a + 3b = 360^\circ \Rightarrow \angle EDC = a+b=120^\circ 
        \]
        在 $\triangle EDC$ 中,由余弦定理,
        \[
        EC^2 = 1^2 + 2^2 - 2\cdot 1 \cdot 2 \cdot \cos 120^\circ \Rightarrow EC = \sqrt{7}
        \]
        在 $\triangle EOC$ 中,由正弦定理,
        \[
        \frac{\sqrt{7}}{\sin 120^\circ} = 2r \Rightarrow r = \sqrt{\frac{7}{3}}
        \]
        在 $\triangle EOD$ 及$\triangle DOC$中,由余弦定理,
        \[
        \cos a = \frac{r^2 + r^2 - 1^2}{2 r^2} = \frac{11}{14}, \cos b = \frac{r^2 + r^2 - 2^2}{2 r^2} = \frac{1}{7} \Rightarrow \sin a = \frac{5\sqrt{3}}{14}, \ \sin b = \frac{4\sqrt{3}}{7}.
        \]
        故六边形$ABCDEF$面积为
        \[
        [ABCDEF] = \frac{3}{2} r^2 \sin a + \frac{3}{2} r^2 \sin b = \frac{7}{2} \left( \frac{5\sqrt{3}}{14} + \frac{4\sqrt{3}}{7} \right) = \frac{13\sqrt{3}}{4}
        \]
    \end{solution}

    \question 边长为 $1$ 的正五边形内部,去掉同时与五个顶点距离都小于 $1$ 的点后,求剩余部分的面积。
    \ifprintanswers
    \begin{figure}[H]
        \centering
        \includegraphics[width=0.5\textwidth]{images/image132.jpg}
    \end{figure}
    \fi
    \begin{solution}
        如图,此题相当于求正五边形$ABCDE$扣除青色部分的面积,该部分由一个小正五边形和五个弓形组成,首先求小正五边形边长 $PQ$,易知
        \[
        \angle PAQ = \angle QAE + \angle PAB - \angle A = 60^\circ + 60^\circ - 108^\circ = 12^\circ, \angle QAB = 60^\circ - 12^\circ = 48^\circ.
        \]
        在 $\triangle APQ$ 中,由余弦定理,
        \[
        PQ^2 = 1^2+1^2-2 \cdot 1 \cdot 1 \cdot \cos 12^\circ = 2 - 2\cos 12^\circ
        \]
        弓形面积为
        \[
        [\text{弓形}]=[\text{扇形}APQ]-[\triangle APQ] = \frac{12}{360} \cdot 1^2 \cdot \pi \cdot-\frac{1}{2} \cdot 1^2 \cdot\sin 12^\circ=\frac{\pi}{30} - \frac{1}{2}\sin 12^\circ.
        \]
        小正五边形与五个弓形组成区域 $R$ 面积
        \[
        R = \frac{5}{4} (2 - 2\cos 12^\circ) \tan 54^\circ + 5\left(\frac{\pi}{30} - \frac{1}{2}\sin 12^\circ \right) 
        \]
        所以剩余部分面积为
        \begin{align*}
        [\text{剩余部分}] &= [ABCDE]-R \\
        &= \frac{5}{4}\tan 54^\circ - \frac{5}{4} (2 - 2\cos 12^\circ) \tan 54^\circ - \frac{\pi}{6} + \frac{5}{2}\sin 12^\circ \\
        &= \frac{5}{4} \tan 54^\circ (2\cos 12^\circ - 1) + \frac{5}{2}\sin 12^\circ - \frac{\pi}{6} \\
        &={5\over 4}{\cos 36^\circ \over \sin 36^\circ}(2\cos 12^\circ-1)+{5\over 2}\sin 12^\circ -{\pi \over 6} \\
        &={5\over 4} \times {2\cos 24^\circ-\cos 36^\circ \over \sin 36^\circ}-{\pi \over 6} \\
        &= \frac{5\sqrt{3}}{4} - \frac{\pi}{6}
        \end{align*}
        其中
        \[
        \cos24^\circ=\cos60^\circ\cos36^\circ+\sin60^\circ\sin36^\circ
        =\frac12\cos36^\circ+\frac{\sqrt3}{2}\sin36^\circ
        \]
    \end{solution}
\end{questions}

\pagebreak

\begin{center}
  {\fontsize{30pt}{26pt}\selectfont
    \hypertarget{三角函数}{三角函数} \label{三角函数}
  }
\end{center}
\separator
\vspace{1pt}

\begin{questions}
    \question 已知$\cos x+\cos^2 x+\cos^3x=1$,求$\cos^3 x+\sec x$的值。
    \begin{solution}
    反复操弄得 
    \begin{align*}
        \cos^3x+\dfrac{1}{\cos x}&=\dfrac{\cos^4x+1}{\cos x}\\
        &=\dfrac{\cos^4x+\cos x+\cos^2x+\cos^3x}{\cos x}\\
        &=\cos^3x+\cos^2x+\cos x+1\\
        &=1+1\\
        &=2
    \end{align*}
    \end{solution}
    \question 已知 \(\cos\theta + 8\sin\theta = 4\)。求 $\dfrac{1+\sin\theta}{\cos\theta}$ 的值。  
    \begin{solution}
        两边平方$\cos\theta =4- 8\sin\theta$有
        \[
        1-\sin^2 \theta=16-64\sin\theta+64\sin^2\theta \Rightarrow \sin \theta= \frac{3}{5}\; \text{或} \;\frac{5}{13}
        \]
        当$\sin \theta= \dfrac{3}{5},\cos \theta= \dfrac{3}{5},\dfrac{1+\sin\theta}{\cos\theta}=-2$;
        当$\sin \theta= \dfrac{5}{13},\cos \theta= \dfrac{12}{13},\dfrac{1+\sin\theta}{\cos\theta}=\dfrac{3}{2}$,故
        \[
        \dfrac{1+\sin\theta}{\cos\theta}=-2\;\text{或}\;\dfrac{3}{2}
        \]
    \end{solution}
    
    \question 若 \(\sec x + \tan x = 2018\),求 \(\csc x + \cot x\)。  
    \begin{solution}
        已知 \(\sec x + \tan x = 2018\),则 \(\sec x - \tan x =\dfrac{\sec^2 x - \tan^2 x}{\sec x + \tan x}= \dfrac{1}{2018}\),解得
        \[
        \sec x = \frac{2018^2+1}{2\cdot 2018},\; \tan x = \frac{2018^2-1}{2\cdot 2018}
        \]
        于是
        \[
        \csc x + \cot x=\frac{\sec x}{\tan x}+\frac{1}{\tan x}=\frac{2018^2+1+2\cdot2018}{2018^2-1}=\frac{2019^2}{2018^2-1}
        \]
    \end{solution}

    \question 已知 $\sin (1+\cos^2 x+\sin^4 x)=\dfrac{13}{14}$,求 $\cos (1+\sin^2 x+\cos^4 x)$的值。
    \begin{solution}
        记$\alpha=1+\cos^2 x+\sin^4 x, \beta=1+\sin^2 x+\cos^4 x$,发现
        \begin{align*}
        \alpha-\beta&=(1+\cos^2 x+\sin^4 x)-(1+\sin^2 x+\cos^4 x)\\
        &=\cos^2 x(1-\cos^2 x)+\sin^2 x(\sin^2 x-1)\\
        &=\cos^2 x\sin^2 x-\sin^2 x\cos^2 x\\
        &=0
        \end{align*}
        于是
        \[
        \cos^2 \beta = 1-\sin^2 \beta=1-\sin^2 \alpha= 1- \left(\frac{13}{14}\right)^2=\frac{27}{196} \Rightarrow \cos \beta=\pm \frac{3\sqrt3}{14}
        \]
    \end{solution}  

    \question 设 $a,b$ 为实数,且满足 
    \[
    \sin a+\sin b=\frac{\sqrt{2}}{2},\cos a+\cos b=\frac{\sqrt{6}}{2},
    \]
    求 $\sin(a+b)$ 的值。
    \begin{solution}
        由
        \[
        \frac{\sin a+\sin b}{\cos a+\cos b}
        =\frac{2\sin\frac{a+b}{2}\cos\frac{a-b}{2}}{2\cos\frac{a+b}{2}\cos\frac{a-b}{2}}
        =\tan\frac{a+b}{2}=\frac{1}{\sqrt{3}}
        \]
        因此
        \[
        \sin(a+b)=\frac{2\tan\frac{a+b}{2}}{1+\tan^2\frac{a+b}{2}}=\frac{\sqrt{3}}{2}
        \]
    \end{solution}
    \begin{solution}
        设复数$z_1=\cos a+i\sin a,\quad z_2=\cos b+i\sin b$,则
        \[
        z=z_1+z_2
        =\frac{\sqrt{6}}{2}+i\frac{\sqrt{2}}{2}
        =\sqrt{2}\left(\cos 30^\circ+i\sin 30^\circ\right)
        \Rightarrow \frac{a+b}{2}=30^\circ \Rightarrow a+b=60^\circ
        \]
        所以
        \[
        \sin(a+b)=\frac{\sqrt{3}}{2}
        \]
    \end{solution}

    \question 已知 $x\neq y,x\neq 0,y\neq 0$, 且$\tan(x+y)=2\tan(x-y)$,求
    \[
    \frac{\sin 2x}{\sin 2y}
    \]
    的值。
    \begin{solution}
        展开$\tan(x+y)=2\tan(x-y)$得
        \[
        \frac{\tan x+\tan y}{1-\tan x\tan y}
        =
        2\left(\frac{\tan x-\tan y}{1+\tan x\tan y}\right)\Rightarrow \tan x(1+\tan^2 y)=3\tan y(1+\tan^2 x)
        \]
        利用恒等式 $1+\tan^2\theta=\sec^2\theta$,得
        \[
        \frac{\tan x\sec^2 y}{\tan y\sec^2 x}=\frac{\sin x\cos x}{\sin y\cos y}=\frac{\sin 2x}{\sin 2y}=3
        \]
    \end{solution}

    \question
    已知三角方程
    \[
    \frac{\sin(x-\alpha)}{\cos(x-\alpha)-2\tan\alpha\sin(x-\alpha)}=\tan\alpha
    \]
    证明
    \[
    \tan x=2\tan\alpha
    \]
    \begin{solution}
        据题意有
        \[
        \frac{\sin(x-\alpha)}{\cos(x-\alpha)-2\tan\alpha\sin(x-\alpha)}=\tan\alpha
        \]
        \[
        \sin(x-\alpha)+2\tan^2\alpha\sin(x-\alpha)=\tan\alpha\cos(x-\alpha)
        \]
        \[
        \sin(x-\alpha)\left(1+2\tan^2\alpha\right)=\tan\alpha\cos(x-\alpha)
        \]
        \[
        \tan(x-\alpha)\left(1+2\tan^2\alpha\right)=\tan\alpha
        \]
        \[
        \frac{\tan x-\tan\alpha}{1+\tan x\tan\alpha}\left(1+2\tan^2\alpha\right)=\tan\alpha
        \]
        \[
        (\tan x-\tan\alpha)(1+2\tan^2\alpha)=\tan\alpha(1+\tan x\tan\alpha)
        \]
        \[
        \tan x+2\tan^2\alpha\tan x-\tan\alpha-2\tan^3\alpha
        =\tan\alpha+\tan^2\alpha\tan x
        \]
        \[
        (\tan x-2\tan\alpha)(\tan^2\alpha+1)=0
        \]
        由于 $\tan^2\alpha+1\neq0$,故
        \[
        \tan x=2\tan\alpha
        \]
        证毕。
    \end{solution}

    \question 已知 $0<x<\dfrac{\pi}{2},0<y<\dfrac{\pi}{2}$,且
    \[
    \sin(x+y)\sin(x-y)=\frac{5}{36}, \quad \cos x+\cos y=\frac{5}{6}.
    \]
    求 $\cos(x-y)$的值。
    \begin{solution}
        由
        \begin{align*}
        \frac{5}{36}=\sin(x+y)\sin(x-y) &= (\sin x\cos y + \cos x \sin y)(\sin x\cos y - \cos x \sin y) \\
        &= \sin^2 x \cos^2 y - \cos^2 x \sin^2 y \\
        &= (1-\cos^2 x)\cos^2 y - \cos^2 x (1-\cos^2 y) \\
        &= \cos^2 y - \cos^2 x 
        \end{align*}
        再由平方差公式,得
        \[
        \cos y - \cos x = \frac{1}{6}.
        \]
        联立可得
        \[
        \cos x = \frac{1}{3}, \quad \cos y = \frac{1}{2}
        \]
        由于 $x, y$ 皆为锐角,
        \[
        \sin x = \frac{2\sqrt{2}}{3}, \quad \sin y = \frac{\sqrt{3}}{2}
        \]
        使用差角公式:
        \[
        \cos(x-y) = \cos x \cos y + \sin x \sin y = \frac{1}{3}\cdot \frac{1}{2} + \frac{2\sqrt{2}}{3} \cdot \frac{\sqrt{3}}{2} = \frac{1+\sqrt{6}}{6}
        \]
    \end{solution}

    \question 已知方程
    \[
    \frac{\sin x}{\sin y} = \frac{1}{2}, \quad 
    \frac{\cos x}{\cos y} = \frac{3}{2},
    \]
    求 $\tan^2(x+y)$。
    \begin{solution}
        由
        \[
        \frac{\sin x}{\sin y} = \frac{1}{2}, \quad 
        \frac{\cos x}{\cos y} = \frac{3}{2},
        \]
        可得
        \[
        \frac{\sin x + \sin y}{\sin x - \sin y }=\frac{2\sin \frac{x+y}{2}\cos \frac{x-y}{2}}{2\sin \frac{x-y}{2}\cos \frac{x+y}{2}}=-3,\quad \frac{\cos x + \cos y}{\cos x - \cos y }=\frac{2\cos \frac{x+y}{2}\cos \frac{x-y}{2}}{-2\sin \frac{x-y}{2}\sin \frac{x+y}{2}}=5 ,
        \]
        故
        \[
        \tan^2 \frac{x+y}{2}=\frac{3}{5}
        \]
        又
        \[
        \tan (x+y)=\frac{2\tan \frac{x+y}{2}}{1-\frac{3}{5}}=5\tan\frac{x+y}{2}
        \]
        所以
        \[
        \tan^2 (x+y)=25\cdot \frac{3}{5}=15
        \]
    \end{solution}

    \question 设 $\beta \in \left[\dfrac{3\pi}{4}, \pi\right]$ 且满足
    \[
    \cos\alpha + \sin(\alpha+\beta) + \cos(\alpha+\beta) = \sqrt{3}
    \]
    求 $\beta$ 的值。
    \begin{solution}
        \begin{align*}
        \sqrt{3} &= \cos\alpha + \sin(\alpha+\beta) + \cos(\alpha+\beta)\\
        &= \cos\alpha + \sin\alpha \cos\beta + \cos\alpha \sin\beta + \cos\alpha \cos\beta - \sin\alpha \sin\beta \\
        &= \cos\alpha(1 + \sin\beta + \cos\beta) + \sin\alpha(\cos\beta - \sin\beta) \\
        &= \sqrt{(1 + \sin\beta + \cos\beta)^2 + (\cos\beta - \sin\beta)^2} \sin(\alpha + \gamma)\\
        &= \sqrt{3 + 2\sin\beta + 2\cos\beta}\sin(\alpha + \gamma)
        \end{align*}
        由于 $\sin(\alpha + \gamma) \le 1$,因此
        \[
        \sqrt{3 + 2\sin\beta + 2\cos\beta} \ge \sqrt3 \Rightarrow \sin\beta + \cos\beta \ge 0 \Rightarrow \sin\left(\beta + \frac{\pi}{4}\right) \ge 0
        \]
        而 $\beta + \frac{\pi}{4} \in [\pi, \frac{5\pi}{4}]\Rightarrow \sin\left(\beta + \frac{\pi}{4}\right) \le 0$,故\[\sin\left(\beta + \frac{\pi}{4}\right)=0 \Rightarrow \beta=\dfrac{3\pi}{4}\]
    \end{solution}

    \question 求解三角方程
    \[
    \sin x\sin 2x+\sin 2x\sin 3x+\sin 3x\sin 4x=0
    \]
    \begin{solution}
        原方程化为
        \[
        \sin 2x(\sin x+\sin 3x)+\sin 3x\sin 4x=0
        \]
        \[
        \sin 2x(2\sin 2x\cos x)+\sin 3x(2\sin 2x\cos 2x)=0
        \]
        \[
        2\sin 2x(\sin 2x\cos x+\sin 3x\cos 2x)=0
        \]
        \[
        2\sin 2x \sin x \left[2\cos^2 x + (3-4\sin^2 x)(1-2\sin^2 x)\right]=0
        \]
        \[
        \sin 2x \sin x \left[1+(3-4\sin^2 x)^2\right]=0
        \]
        此时$1+(3-4\sin^2 x)^2>0$,解$\sin 2x=0$即可,通解为
        \[
        x=\frac{n\pi}{2},\quad n\in\mathbb{Z}
        \]
    \end{solution}

    \question 解三角方程
    \[
    2\cos x\sin^2 x-2\cos^2 x\sin x+\cos^2 x-4\sin^2 x+3\cos x\sin x+2\sin x-2\cos x=0
    \]
    \begin{solution}
        原方程化为
        \[
        2\cos x\sin x(\sin x -\cos x)+(\cos x -\sin x)(\cos x +4\sin x)+2(\sin x-\cos x)=0
        \]
        \[
        (\sin x-\cos x)\left[2\cos x\sin x-(\cos x +4\sin x)+2\right]=0
        \]
        \[
        (\sin x-\cos x)(2\sin x-1)(\cos x-2)=0
        \]
        注意到方程$\cos x=2$无实数解,于是
        \[
        \tan x = 1\quad \text{或} \quad \sin x=\frac{1}{2}
        \]
        通解为
        \[
        x=n\pi+\frac{\pi}{4} \quad \text{或} \quad x=n\pi +(-1)^n \left(\frac{\pi}{6}\right),\quad n\in\mathbb{Z}
        \]
    \end{solution}

    \question 解方程 
    \[
    \sin^2 x + \sin^2 2x = \sin^2 3x.
    \]
    \begin{solution}
        首先有
        \begin{align*}
        \sin^2 3x - \sin^2 x 
        &= (\sin 3x + \sin x)(\sin 3x - \sin x)\\
        &=(2 \sin 2x \cos x)(2 \cos 2x \sin x)\\
        &=2 \sin^2 2x \cos 2x
        \end{align*}
        故原方程式变为
        \[
        \sin^2 2x (2 \cos 2x - 1) = 0
        \]
        解得
        \[
        x = \frac{n\pi}{2} \quad \text{或} \quad x = n\pi \pm \frac{\pi}{6},\quad n \in \mathbb{Z}
        \]
    \end{solution}

    \question 解方程 
    \[
    \cos^2 x + 3 \cos^2 2x = \cos^2 3x.
    \]
    \begin{solution}
        同上题,有
        \begin{align*}
        \cos^2 3x - \cos^2 x 
        &= (\cos 3x + \cos x)(\cos 3x - \cos x) \\
        &= (2 \cos 2x \cos x)(-2 \sin 2x \sin x) \\
        &= -2 \cos 2x \sin^2 2x \\
        &= -2 \cos 2x (1-\cos^2 2x)
        \end{align*}
        故原方程式变为
        \[
        \cos 2x (\cos 2x-2)(2\cos 2x+1) =0
        \]
        其中$\cos 2x=2$无实数解,故通解为
        \[
        x = \frac{n\pi}{2} + \frac{\pi}{4} \quad \text{或} \quad x = n\pi \pm \frac{\pi}{3},\quad n \in \mathbb{Z}
        \]
    \end{solution}
    
    \question 求解
    \[
    \sin\left(3x-\frac{\pi}{3}\right)+\sin\left(x-\frac{\pi}{3}\right)=\sin x
    \]
    \begin{solution}
        和差化积得
        \[
        2\sin\left(2x-\frac{\pi}{3}\right)\cos x=\sin x 
        \]
        \[
        \cos x (\sin2x-\sqrt3\cos2x)=\sin x
        \]
        \[ 
        \sin x(2\cos^2 x-1)=\sqrt{3}\cos 2x\cos x 
        \]
        \[
        \cos2x(\tan x-\sqrt3)=0
        \]
        故通解为
        \[
        x=n\pi \pm \frac{\pi}{4} \quad \text{或} \quad x=n\pi+\frac{\pi}{3}, \quad n\in\mathbb{Z}
        \]
    \end{solution}

    \question 解方程
    \[
    (\cos 4x + \cos x)^2 + (\sin 4x + \sin x)^2 = 2\sqrt{3}\sin 3x, \quad 0 \le x < \pi
    \]
    \begin{solution}
        首先展开左式,
        \[
        (\cos 4x+\cos x)^2+(\sin 4x+\sin x)^2 = 1 + 1 + 2\cos(4x-x) = 2 + 2\cos 3x
        \]
        原方程即
        \[
        2 + 2\cos 3x = 2\sqrt{3}\sin 3x \Rightarrow \sqrt{3}\sin 3x - \cos 3x = 1 \Rightarrow \sin\left(3x-\frac{\pi}{6}\right) = \frac{1}{2}
        \]
        故通解为
        \[
        3x - \frac{\pi}{6} = n\pi + (-1)^n\left(\frac{\pi}{6}\right) \Rightarrow x= \frac{n\pi}{3} + \frac{1+(-1)^n}{18}\pi, \quad n \in \mathbb{Z}
        \]
    \end{solution}

    \question 解三角方程式 
    \[
    16\sin^{2}\theta + \tan^{2}\theta + 9(\csc^{2}\theta + \cot^{2}\theta) = 30, 
    \]
    其中 $0^\circ \leq \theta \leq 360^\circ$。
    \begin{solution}
        原方程式即
        \[
        (4 \sin\theta - 3 \csc\theta)^2 + (\tan\theta - 3 \cot\theta)^2 = 0
        \]
        于是
        \[
        4 \sin\theta - 3 \csc\theta = 0 \Rightarrow \sin\theta = \pm\frac{\sqrt{3}}{2} \Rightarrow \theta =60^\circ, 120^\circ, 240^\circ, 300^\circ
        \]
        且
        \[
        \tan\theta - 3 \cot\theta = 0 \Rightarrow \tan\theta = \pm\sqrt{3} \Rightarrow \theta = 60^\circ, 120^\circ, 240^\circ, 300^\circ
        \]
        故原方程式的解为
        \[
        \theta = 60^\circ, 120^\circ, 240^\circ, 300^\circ
        \]
    \end{solution}

    \question 求解方程
    \[
    (x+y)^2 + 4(x+y)\cos(x-y) + 4 = 0
    \]
    \begin{solution}
        注意到
        \begin{align*}
        (x+y)^2 + 4(x+y)\cos(x-y) + 4 &= \left[(x+y) + 2\cos(x-y)\right]^2 - (2\cos(x-y))^2 + 4 \\
        &= \left[(x+y) + 2\cos(x-y)\right]^2 + \left[2\sin(x-y)\right]^2
        \end{align*}
        由于两个平方项均为实数,原方程等价于
        \[
        (x+y) + 2\cos(x-y) = 0 \tag{1}
        \]
        且
        \[
        \sin(x-y) = 0 \Rightarrow x-y = n\pi, \quad n \in \mathbb{Z}
        \]
        若 $n$ 为偶数,则$x-y = 2k\pi \Rightarrow \cos(x-y) = 1$,代入 (1) 得
        \[
        x+y = -2
        \]
        若 $n$ 为奇数,则$x-y = (2k+1)\pi \Rightarrow \cos(x-y) = -1$,代入 (1) 得
        \[
        x+y = 2
        \]
        故可解得
        \[
        x = -1+k\pi,  y = -1-k\pi, \quad \text{或} \quad
        x = 1+\frac{(2k+1)\pi}{2}, \, y = 1-\frac{(2k+1)\pi}{2}, \quad k \in \mathbb{Z}
        \]
    \end{solution}

    \question 解方程
    \[
    4\cos x\cos 2x\cos 5x+1=0.
    \]
    \begin{solution}
        原方程等价于
        \[
        4\cos x\cos 2x\cos 5x=-1
        \]
        注意到 $x=n\pi,n\in\mathbb{Z}$ 不是原方程的解,两边乘 $\sin x$得
        \[
        4\sin x\cos x\cos 2x\cos 5x=-\sin x
        \]
        即
        \[
        2\sin 2x\cos 2x\cos 5x=-\sin x \Rightarrow \sin 4x\cos 5x=-\sin x
        \]
        积化和差得
        \[
        \frac12\sin 9x-\frac12\sin x=-\sin x \Rightarrow
        \sin 9x=-\sin x
        \]
        即
        \[
        \sin 9x=\sin(-x)
        \]
        通解为
        \[
        9x=-x+2n\pi \Rightarrow x=\frac{n\pi}{5}, \quad n\in\mathbb{Z}, \quad n \not\equiv 0 \pmod 5
        \]
        或
        \[
        9x=\pi-(-x)+2n\pi \Rightarrow x=\frac{n\pi}{4}+\frac{\pi}{8}, \quad n\in\mathbb{Z}
        \]
    \end{solution}

    \question 解 
    \[
    \sin^{8}\theta + \cos^{8}\theta = \frac{17}{32}.
    \] 
    \begin{solution}
        化简左式得
        \begin{align*}
        \sin^8 \theta + \cos^8 \theta
        &=(\sin^4 \theta + \cos^4 \theta)^2 - 2 \sin^4 \theta \cos^4 \theta \\
        &=\left((\sin^2 \theta + \cos^2 \theta)^2 - 2 \sin^2 \theta \cos^2 \theta\right)^2 - 2 \sin^4 \theta \cos^4 \theta \\
        &=\left(1 - \frac{1}{2} \sin^2 2\theta\right)^2 - \frac{1}{8} \sin^4 2\theta 
        \end{align*}
        故原方程即
        \[
        4 \sin^2 2\theta - 32 \sin^2 2\theta + 15 = 0 \Rightarrow (2 \sin^2 2\theta - 1)(2 \sin^2 2\theta - 15) = 0
        \]
        于是由$\sin^2 2\theta = \dfrac{1}{2} \Rightarrow \sin 2\theta = \pm \dfrac{\sqrt{2}}{2}$,解为
        \[
        \theta = \frac{n\pi}{2} + \frac{\pi}{8}, \quad n \in \mathbb{Z}
        \]
        且由于$0 \leq \sin^2 \alpha \leq 1, \alpha \in \mathbb{R}$,方程式$\sin^2 2\theta = \dfrac{15}{2}$无解。
    \end{solution}

    \question 解方程  
    \[
    \cos\frac{4x}{3}=\cos^{2}x.
    \]
    \begin{solution}
        令 \(x = \dfrac{3t}{2}\),则 \(\dfrac{4x}{3} = 2t\),原式变为
        \[
        \cos2t = \cos^2\frac{3t}{2}
        \]
        无中生有,
        \[
        2\cos2t - 1 = 2\cos^2\frac{3t}{2} - 1 = \cos3t
        \]
        于是
        \[
        2(2\cos^2t-1) - 3 = 4\cos^3 t - 3\cos t
        \Rightarrow (4\cos^2 t - 3)(\cos t - 1) = 0
        \]
        若 \(\cos t = 1\),则 \(t = 2k\pi \Rightarrow x  = 3k\pi\);若 \(\cos^2 t = \dfrac{3}{4}\),则$t = \pm \dfrac{\pi}{6} + 2k\pi \Rightarrow x = \pm\dfrac{\pi}{4} + 3k\pi $,通解为
        \[
        x = 3k\pi \quad \text{或} \quad x = \pm\dfrac{\pi}{4} + 3k\pi,\quad k \in \mathbb{Z}
        \]
    \end{solution}

    \question 设 $x$ 满足 $0 < x < \dfrac{\pi}{2}$ 且 
    \[
    \cos \left(\frac{3}{2} \cos x\right) = \sin \left(\frac{3}{2} \sin x\right),
    \]
    求 $\sin 2x$ 的所有可能值。
    \begin{solution}
        由于 $x \in \left(0,\dfrac{\pi}{2}\right)$,所以 
        \[
        \frac{3}{2}\cos x,\frac{3}{2}\sin x \in \left(0,\frac{3}{2}\right) \subset \left(0,\frac{\pi}{2}\right)
        \]
        若 $Y,Z \in \left(0,\dfrac{\pi}{2}\right)$,则 $\cos Y = \sin Z$ 当且仅当 $Y+Z=\dfrac{\pi}{2}$,因此
        \[
        \cos\left(\frac{3}{2}\cos x\right) = \sin\left(\frac{3}{2}\sin x\right)
        \quad \Longleftrightarrow \quad
        \frac{3}{2}\cos x + \frac{3}{2}\sin x = \frac{\pi}{2}
        \]
        即
        \[
        \cos x + \sin x = \frac{\pi}{3}
        \]
        两边平方得
        \[
        \sin 2x = \frac{\pi^2}{9} - 1
        \]
        为唯一可能值。
    \end{solution}
    
    \question 设 $0^\circ \le \alpha \le 45^\circ$,若已知 $\cot 2\alpha - \sqrt{3} = \sec \alpha$,求 $\alpha$。
    \begin{solution}
        化简得
        \[
        \cot(2\alpha) - \sec \alpha 
        = \frac{\cos 2\alpha}{\sin 2\alpha} - \frac{1}{\cos \alpha} 
        = \frac{1-2\sin^2\alpha}{2\sin \alpha\cos \alpha} - \frac{1}{\cos \alpha}
        \]
        原方程即
        \[
        1 - 2\sin^2\alpha - 2\sin\alpha = 2\sqrt{3} \sin \alpha \cos \alpha 
        \]
        两边平方得
        \[
        (1 - 2\sin^2\alpha - 2\sin\alpha)^2 = 12\sin^2\alpha(1-\sin^2\alpha)
        \]
        \[
        16\sin^4\alpha + 8\sin^3\alpha - 12\sin^2\alpha - 4\sin\alpha + 1 = 0
        \]
        \[
        -4\sin\alpha(3\sin\alpha - 4\sin^3\alpha) - (-8\sin^3\alpha + 6\sin\alpha) + 2\sin\alpha + 1 = 0
        \]
        \[
        -4\sin\alpha \sin 3\alpha - 2\sin 3\alpha + 2\sin\alpha + 1 = 0
        \]
        \[
        -2\sin 3\alpha (2\sin\alpha + 1) + 2\sin\alpha + 1 = 0
        \]
        \[
        (1 - 2\sin 3\alpha)(2\sin\alpha + 1) = 0
        \]
        故
        \[
        \sin 3\alpha = \frac{1}{2} \Rightarrow 3\alpha = 30^\circ \Rightarrow \alpha = 10^\circ
        \]
    \end{solution}

    \question 
    \begin{parts}
    \part 已知$\theta \neq 2n\pi-\dfrac{\pi}{2},\ n\in \mathbb{Z}$,证明
    \[
    \tan^2\left(\frac{\pi}{4}-\frac{\theta}{2}\right) \equiv \frac{1-\sin\theta}{1+\sin\theta}
    \]
    \begin{solution}
        有
        \begin{align*}
        \tan^2\left(\frac{\pi}{4}-\frac{\theta}{2}\right) 
        &= \left(\frac{1-\tan\frac{\theta}{2}}{1+\tan\frac{\theta}{2}}\right)^2 \\
        &= \left(\frac{\cos\frac{\theta}{2}-\sin\frac{\theta}{2}}{\cos\frac{\theta}{2}+\sin\frac{\theta}{2}}\right)^2 \\
        &= \frac{1-2\sin\frac{\theta}{2}\cos\frac{\theta}{2}}{1+2\sin\frac{\theta}{2}\cos\frac{\theta}{2}} \\
        &= \frac{1-\sin\theta}{1+\sin\theta}.
        \end{align*}
        故左式=右式,恒等式得证。
    \end{solution}
    \part 解方程
    \[
    \tan\left(\frac{\pi}{4}-2x\right)=\sqrt{7+4\sqrt{3}},\quad 0\le x<\pi.
    \]
    \begin{solution}
        由(a),取 $\theta = 4x$,则
        \[
        \tan^2\left(\frac{\pi}{4}-2x\right) = \frac{1-\sin 4x}{1+\sin 4x}
        \]
        原方程两边平方得
        \[
        \frac{1-\sin 4x}{1+\sin 4x} = 7+4\sqrt{3} \Rightarrow \sin 4x = -\frac{\sqrt{3}}{2}
        \]
        所以通解为
        \[
        4x = n\pi+(-1)^n\left(-\frac{\pi}{3}\right) \Rightarrow x = \frac{n\pi}{4}+(-1)^n\left(-\frac{\pi}{12}\right), \quad n\in\mathbb{Z}
        \]
        且慢,注意到
        \[
        x=\frac{n\pi}{2}+\frac{\pi}{3}, \quad n\in\mathbb{Z}
        \]
        为增根,舍去之,于是只取
        \[
        x=\frac{n\pi}{2}-\frac{\pi}{12}, \quad n\in\mathbb{Z}
        \]
        在区间 $0\le x<\pi$ 内的解有
        \[
        x = \frac{5\pi}{12},\ \frac{11\pi}{12}
        \]
    \end{solution}
    \end{parts}

    \question 
    \begin{parts}
    \part 已知 $\theta \neq 2n\pi+\dfrac{\pi}{2},k\in\mathbb{Z}$,证明三角恒等式
    \[
    \tan\left(\frac{\theta}{2}+\frac{\pi}{4}\right)\equiv\tan\theta+\sec\theta
    \]
    \begin{solution}
        从左式开始,有
        \[
        \tan\left(\frac{\theta}{2}+\frac{\pi}{4}\right)
        =\frac{\tan\frac{\theta}{2}+1}{1-\tan\frac{\theta}{2}}=\frac{\sin\frac{\theta}{2}+\cos\frac{\theta}{2}}{\cos\frac{\theta}{2}-\sin\frac{\theta}{2}}
        \]
        将分子分母同乘 $\cos\dfrac{\theta}{2}+\sin\dfrac{\theta}{2}$,得
        \[
        =\frac{\sin^2\frac{\theta}{2}+2\sin\frac{\theta}{2}\cos\frac{\theta}{2}+\cos^2\frac{\theta}{2}}
        {\cos^2\frac{\theta}{2}-\sin^2\frac{\theta}{2}}
        =\frac{1+\sin\theta}{\cos\theta}=\sec\theta+\tan\theta
        \]
        于是左式等于右式,恒等式得证。
    \end{solution}
    \part 已知
    \[
    \tan x-\tan(x-\alpha)=2\tan\alpha
    \]
    其中 $\alpha$ 为常数。以 $\alpha$ 的三角函数表示 $\tan x$。
    \begin{solution}
        可得
        \[
        \tan x-\frac{\tan x-\tan\alpha}{1+\tan x\tan\alpha}=2\tan\alpha \Rightarrow \tan^2x-2\tan x\tan\alpha-1=0
        \]
        配方得
        \[
        (\tan x-\tan\alpha)^2=1+\tan^2\alpha=\sec^2\alpha
        \]
        因此
        \[
        \tan x=\tan\alpha\pm\sec\alpha
        \]
    \end{solution}
    \part 解三角方程
    \[
    \tan x-\tan\left(x-\frac{3\pi}{5}\right)=2\tan\frac{3\pi}{5},\quad 0\le x<2\pi
    \]
    \begin{solution}
        取 $\alpha=\dfrac{3\pi}{5}$,由 (b) 得
        \[
        \tan x=\tan\frac{3\pi}{5}\pm\sec\frac{3\pi}{5}
        \]
        且由(a),以$\theta-\pi$代替$\theta$有恒等式
        \[
        \tan\left(\frac{\theta}{2}-\frac{\pi}{4}\right)\equiv\tan\theta-\sec\theta
        \]
        因此
        \[
        \tan\frac{3\pi}{5}+\sec\frac{3\pi}{5}
        =\tan\left(\frac{3\pi}{10}+\frac{\pi}{4}\right)
        =\tan\frac{11\pi}{20}
        \]
        \[
        \tan\frac{3\pi}{5}-\sec\frac{3\pi}{5}
        =\tan\left(\frac{3\pi}{10}-\frac{\pi}{4}\right)
        =\tan\frac{\pi}{20}
        \]
        于是
        \[
        \tan x=\tan\frac{11\pi}{20}\quad \text{或}\quad \tan x=\tan\frac{\pi}{20}
        \]
        在 $0\le x<2\pi$ 内解得
        \[
        x=\frac{\pi}{20},\ \frac{11\pi}{20},\ \frac{21\pi}{20},\ \frac{31\pi}{20}
        \]
    \end{solution}
    \end{parts}

    \question 设 $0 \le x \le \pi$,求 
    \[
    1+\sqrt{\sin x} - \sqrt{x} = \cos 2x + 2x^2
    \] 
    的实根个数。
    \begin{solution}
        有
        \begin{align*}
        1+\sqrt{\sin x}- \sqrt x = 1-2\sin^2 x+2x^2& \\
        \sqrt{\sin x}-\sqrt x +2(\sin^2 x-x^2)&=0 \\
        \sqrt{\sin x}-\sqrt x +2(\sin x+x)(\sin x-x) &=0 \\
        \sqrt{\sin x}-\sqrt x +2(\sin x+x)(\sqrt{\sin x}-\sqrt x)(\sqrt{\sin x}+\sqrt x)&=0 \\
        (\sqrt{\sin x}-\sqrt x)\left[2(\sin x+x) (\sqrt{\sin x}+\sqrt x)+1\right]&=0
        \end{align*}
        由 $\sqrt{\sin x}=\sqrt x$ 得 $x=0$;由 $0\le x\le \pi \Rightarrow \sin x \ge 0$,可知 
        \[
        2(\sin x+x) (\sqrt{\sin x}+\sqrt x)+1>0
        \] 恒成立,因此原方程式只有 $1$ 个实根。
    \end{solution}

    \question 已知 $\theta,\alpha,\beta$ 为相异实数,且满足
    \[
    \tan(\theta-\alpha)+\tan(\theta-\beta)=x, \quad 
    \cot(\theta-\alpha)+\cot(\theta-\beta)=y.
    \]
    试以$x,y$表示 $\tan(\alpha-\beta)$。
    \begin{solution}
        不妨设$A = \tan(\theta-\alpha), B = \tan(\theta-\beta)$,由已知得
        \[
        A+B = x, \quad \frac{1}{A}+\frac{1}{B} = y \Rightarrow AB = \frac{x}{y}
        \]
        故
        \begin{align*}
        \tan(\alpha-\beta) 
        &= \tan[(\theta-\beta) - (\theta-\alpha)] \\
        &= \frac{\tan(\theta-\beta)-\tan(\theta-\alpha)}{1+\tan(\theta-\beta)\tan(\theta-\alpha)} \\
        &= \frac{B-A}{1+AB} \\
        &= \pm \frac{\sqrt{(B-A)^2}}{1+AB} \\
        &= \pm \frac{\sqrt{(A+B)^2 - 4AB}}{1+AB} \\
        &= \pm \frac{\sqrt{x^2 - \frac{4x}{y}}}{1 + \frac{x}{y}} \\
        &= \pm \frac{\sqrt{xy(xy - 4)}}{x+y}
        \end{align*}
    \end{solution}

    \question 已知正数 $x, y$ 及角 $\theta \neq \dfrac{n\pi}{2},n \in \mathbb{Z}$满足
    \[
    \begin{cases}
    \displaystyle \frac{\sin \theta}{x} = \frac{\cos \theta}{y} \\[6pt]
    \displaystyle \frac{\cos^4 \theta}{x^4} + \frac{\sin^4 \theta}{y^4} = \frac{7 \sin 2\theta}{x^3y + xy^3}
    \end{cases}
    \]
    求 $\dfrac{x}{y} + \dfrac{y}{x}$。
    \begin{solution}
        令 
        \[
        \frac{\sin \theta}{x} = \frac{\cos \theta}{y}=\frac{1}{k}
        \]
        即$x = k \sin \theta, y = k \cos \theta$,则
        \[
        \frac{\cos^4 \theta}{\sin^4 \theta} + \frac{\sin^4 \theta}{\cos^4 \theta} = \frac{14 \sin \theta \cos \theta}{\sin \theta \cos \theta (\sin^2 \theta + \cos^2 \theta)} = 14
        \]
        设 
        \[
        S = \frac{x}{y} + \frac{y}{x} = \tan \theta + \cot \theta
        \]
        注意到
        \[
        (S^2-2)^2-2=\tan^4 \theta + \cot^4 \theta=14
        \]
        故
        \[
        S = \sqrt{6} >0
        \]
    \end{solution}

    \question 定义 $f(x) = \sin^6 x + \cos^6 x + k (\sin^4 x + \cos^4 x)$,其中 $k \in \mathbb{R}$。
    \begin{parts}
    \part 求所有$k$使得 $f(x)$ 对任意 $x$ 恒为常数。
    \begin{solution}
        利用 $\sin^2x+\cos^2x=1$,设 $u=\sin^2x$,则
        \begin{align*}
        f(x)&=\sin^6x+(1-u)^3+k(\sin^4x+(1-u)^2)\\
        &=(1+k)-(3+2k)u+(3+2k)u^2。
        \end{align*}
        仅当 $3+2k=0$,即 $k=-\dfrac{3}{2}$时,$f(x)=1+k=-\dfrac{1}{2}$ 为常数。
        \end{solution}
    \begin{solution}
        由分解
        \begin{align*}
        f(x)&=(\sin^2x+\cos^2x)((\sin^2x+\cos^2x)^2-3\sin^2x\cos^2x) \\
        &+k((\sin^2x+\cos^2x)^2-2\sin^2x\cos^2x)
        \end{align*}
        整理得
        \[
        f(x)=(1+k)-(3+2k)\sin^2x\cos^2x
        \]
        仅当 $3+2k=0$,即 $k=-\dfrac{3}{2}$,有 $f(x)=-\dfrac{1}{2}, \forall x \in \mathbb{R}$
    \end{solution}
    \begin{solution}
        直接对 $f(x)$ 求导:
        \[
        f'(x)=2\sin x\cos x(\sin^2x-\cos^2x)(3+2k)。
        \]
        欲使 $f'(x)\equiv 0$,必须有 $3+2k=0$,即 $k=-\dfrac{3}{2}$。  
    \end{solution}
    \part 若 $k=-0.7$,求方程 $f(x)=0$ 的所有解。
    \begin{solution}
       当 $k=-0.7$ 时,由(a)解法二,
        \[
        f(x)=(1+k)-(3+2k)\sin^2x\cos^2x=0.3-1.6\sin^2x\cos^2x
        \]
        令 $u=\sin^2x$,则 $\cos^2x=1-u$,方程$f(x)=0$即
        \[
        1.6u(1-u)=0.3,
        \]
        解得
        \[
        u=\sin^2x=\frac{1}{4},\ \frac{3}{4} \Rightarrow \sin x=\pm \frac{1}{2},\ \pm \frac{\sqrt{3}}{2}
        \]
        解集为
        \[
        x=\frac{\pi}{6}+n\pi,\ \frac{5\pi}{6}+n\pi,\ \frac{\pi}{3}+n\pi,\ \frac{2\pi}{3}+n\pi,\quad n\in\mathbb{Z}。
        \]
    \end{solution}
    \part 求所有 $k$使得$f(x)=0$有实数解。
    \begin{solution}
        由 (a) 解法一,
        \[
        f(x)=(3+2k)u^2-(3+2k)u+(1+k),\quad u=\sin^2x \in [0,1]。
        \]
        若 $k=-\tfrac{3}{2}$,则$f(x)\equiv -\tfrac{1}{2}$无解,则由方程
        \[
        u^2-u+\frac{1+k}{3+2k}=0
        \]
        解得
        \[
        u=\frac{1}{2}\pm \frac{1}{2}\sqrt{-\frac{1+2k}{3+2k}}
        \]
        且判别式满足
        \[
        \Delta=(3+2k)^2-4(3+2k)(1+k)=(3+2k)(-1-2k) \ge 0
        \]
        即 
        \[
        -\frac{3}{2}<k\leq -\frac{1}{2}
        \]
        由于 $u\in[0,1]$,经检验得当 
        \[
        -1\leq k\leq -\frac{1}{2}
        \] 
        时方程$f(x)=0$有实数解。 
        \textcolor{red}{感觉有更简洁的解法}
    \end{solution}
    \end{parts}

    \question 证明恒等式
    \[
    \frac{\sin 2\theta}{1+\cos 2\theta} \cdot \frac{\cos \theta}{1+\cos \theta} \equiv \tan \frac{\theta}{2}
    \]
        \begin{solution}
        有
        \begin{align*}
        \frac{\sin 2\theta}{1+\cos 2\theta} \cdot \frac{\cos \theta}{1+\cos \theta} 
        &=\frac{2 \sin \theta \cos \theta}{2\cos^2 \theta} \cdot \frac{\cos \theta}{1+\cos \theta} \\
        &= \frac{\sin \theta}{1+\cos \theta} \\
        &= \frac{2 \sin \frac{\theta}{2} \cos \frac{\theta}{2}}{2\cos^2 \frac{\theta}{2}} \\
        &= \tan \frac{\theta}{2}
        \end{align*}
        故左式=右式,恒等式得证。
    \end{solution}

    \question 证明三角恒等式
    \[
    \frac{1 + \tan \theta \tan 3\theta}{1 + \tan 2\theta \tan 3\theta} \equiv \frac{\cos^2 2\theta}{\cos^2 \theta}.
    \]
    \begin{solution}
        由
        \begin{align*}
        \text{左式} &= \frac{1 + \tan \theta \tan 3\theta}{1 + \tan 2\theta \tan 3\theta} \\
        &= \frac{\cos 2\theta \cos 3\theta}{\cos \theta \cos 3\theta} \cdot \frac{\cos \theta \cos 3\theta+\sin \theta \sin 3\theta}{\cos 2\theta \cos 3\theta+ \sin 2\theta \sin 3\theta} \\
        &= \frac{\cos 2\theta}{\cos \theta} \cdot \frac{\cos(3\theta - \theta)}{\cos(3\theta - 2\theta)}  \\
        &= \frac{\cos^2 2\theta}{\cos^2 \theta} \\
        &= \text{右式}
        \end{align*}
        因此恒等式成立。
    \end{solution}
    
    \question 证明恒等式
    \[
    \frac{1}{\tan 3\theta - \tan \theta} - \frac{1}{\cot 3\theta - \cot \theta} \equiv \cot 2\theta
    \]
    \begin{solution}
        化简得
        \begin{align*}
        &\frac{1}{\tan 3\theta - \tan \theta} - \frac{1}{\cot 3\theta - \cot \theta} \\
        &= \frac{\cos 3\theta \cos \theta}{\sin 3\theta \cos \theta - \sin \theta \cos 3\theta} - \frac{\sin 3\theta \sin \theta}{\sin 3\theta \cos \theta - \cos 3\theta \sin \theta} \\
        &= \frac{\cos 3\theta \cos \theta - \sin 3\theta \sin \theta}{\sin 3\theta \cos \theta - \sin \theta \cos 3\theta} \\
        &= \frac{\cos(3\theta - \theta)}{\sin(3\theta - \theta)} \\
        &= \cot 2\theta
        \end{align*}
        故左式=右式,恒等式得证。
    \end{solution}

    \question 证明恒等式
    \[
    \frac{\sec 16\theta - 1}{\sec 8\theta - 1} \equiv \frac{\tan 16\theta}{\tan 4\theta}.
    \]
    \begin{solution}
        \begin{align*}
        \frac{\sec 16\theta - 1}{\sec 8\theta - 1}
        &= \frac{(1-\cos 16\theta)\cos 8\theta}{(1-\cos 8\theta)\cos 16\theta} \\
        &= \frac{1-\cos^2 16\theta}{1+\cos 16\theta} \cdot \frac{\cos 8\theta}{\cos 16\theta} \cdot \frac{1+\cos 8\theta}{1-\cos^2 8\theta} \\
        &= \frac{\sin^2 16\theta}{2\cos^2 8\theta} \cdot \frac{\cos 8\theta}{\cos 16\theta} \cdot \frac{2\cos^2 4\theta}{\sin^2 8\theta}\\
        &= \frac{\sin^2 16\theta}{\sin 16\theta} \cdot \frac{1}{\sin 8\theta} \cdot \frac{2\cos^2 4\theta}{\cos 16\theta}\\
        &= \tan 16 \theta \cdot \frac{2 \cos^2 4\theta}{2\cos 4 \theta \sin 4 \theta}\\
        &= \frac{\tan16\theta}{\tan 4\theta}
        \end{align*}
        故左式=右式,恒等式得证。
    \end{solution}

    \question 证明恒等式
    \[
    \cot^2 \theta \left( \frac{\sec \theta - 1}{\sin \theta - 1} \right) + \sec^2 \theta \left( \frac{\sin \theta - 1}{\sec \theta - 1} \right) \equiv 0
    \]
    \begin{solution}
        由于
        \begin{align*}
        &\cot^2 \theta \left( \frac{\sec \theta - 1}{\sin \theta - 1} \right) + \sec^2 \theta \left( \frac{\sin \theta - 1}{\sec \theta - 1} \right) \\
        &= -\cot^2 \theta \left( \frac{(\sec \theta - 1)(1+\sin \theta)}{1 - \sin^2 \theta} \right) + \sec^2 \theta \left( \frac{(\sin \theta - 1)(1-\sec \theta)}{\sec^2 \theta-1} \right) \\
        &= -\cot^2 \theta \left( \frac{(1 - \sec \theta)(\sin \theta - 1)}{\cos^2 \theta} \right) + \sec^2 \theta \left( \frac{(1 - \sec \theta)(\sin \theta - 1)}{\tan^2 \theta} \right) \\
        &= \left( -\frac{1}{\sin^2 \theta} + \frac{1}{\sin^2 \theta} \right) (1 - \sec \theta)(\sin \theta - 1) \\
        &= 0
        \end{align*}
        即左式=右式,故恒等式得证。
    \end{solution}

    \question 证明 
    \[
    \sin 5x \equiv 16\sin^5 x - 20 \sin^3 x + 5 \sin x
    \]
    \begin{solution}
        展开得
        \begin{align*}
        \sin 5x
        &= \sin(2x+3x) \\
        &= \sin 2x \cos 3x + \cos 2x \sin 3x \\
        &= 2 \sin x \cos x(4\cos^3 x - 3\cos x) + (1-2\sin^2 x)(3\sin x - 4\sin^3 x) \\
        &= 2 \sin x (1-\sin^2 x)\left(4(1-\sin^2 x)-3\right) + (1-2\sin^2 x)(3\sin x - 4\sin^3 x) \\
        &= 16\sin^5 x - 20\sin^3 x + 5\sin x
        \end{align*}
        于是左式=右式,故得证。
    \end{solution}

    \question 证明
    \[
    4 (\cos 3\theta \sin^3 \theta + \sin 3\theta \cos^3 \theta) \equiv 3 \sin 4\theta
    \]
    \begin{solution}
        \begin{align*}
        4 (\cos 3\theta \sin^3 \theta + \sin 3\theta \cos^3 \theta)
        &= 4 \bigl((4\cos^3\theta - 3\cos\theta)\sin^3\theta + (3\sin\theta - 4\sin^3\theta)\cos^3\theta\bigr) \\
        &= 4 \bigl(-3\cos\theta\sin^3\theta + 3\sin\theta\cos^3\theta\bigr) \\
        &= 12 \sin\theta\cos\theta(\cos^2\theta - \sin^2\theta) \\
        &= 6 \sin 2\theta \cos 2\theta \\
        &= 3 \sin 4\theta
        \end{align*}
        于是左式=右式,恒等式得证。
    \end{solution}

    \question 证明
    \[
    \sin^4\left(\frac{\pi}{4}-\theta\right)+\sin^4\left(\frac{\pi}{4}+\theta\right)
    \equiv \frac{3}{4}-\frac{1}{4}\cos 4\theta
    \]
    \begin{solution}
        有
        \begin{align*}
        \sin^4\left(\frac{\pi}{4}-\theta\right)+\sin^4\left(\frac{\pi}{4}+\theta\right)
        &=\frac{1}{4}\left[(1-\cos\left(\frac{\pi}{2}-2\theta\right))^2+(1+\cos\left(\frac{\pi}{2}+2\theta\right))^2\right] \\
        &=\frac{1}{4}\left[2+2\sin^2 2\theta\right] \\
        &=\frac{1}{2}\left(1+\frac{1-\cos 4\theta}{2}\right) \\
        &=\frac{3}{4}-\frac{1}{4}\cos 4\theta
        \end{align*}
        即左式=右式,恒等式得证。
    \end{solution}

    \question 已知
    \[
    8 \sin^3 x \sin 2x = a \cos 5x + b \cos 3x + c \cos x,
    \]
    求实数 $a, b, c$ 的值。
    \begin{solution}
    \noindent
        对左式化简,
        \begin{align*} 
        8 \sin^3 x \sin 2x 
        &= 4 \sin 2x \sin x \cdot 2 \sin^2 x \\ 
        &= 4 \sin 2x \sin x (1 - \cos 2x) \\ 
        &= 4 \sin 2x \sin x - 4 \sin 2x \cos 2x \sin x \\ 
        &= -2 (-2 \sin 2x \sin x) - 2 \sin 4x \sin x \\ 
        &= -2 (\cos 3x - \cos x) + \cos 5x - \cos 3x \\ 
        &= \cos 5x - 3 \cos 3x + 2 \cos x
        \end{align*}
        于是得到
        \[
        a = 1, \quad b = -3, \quad c = 2
        \]
    \end{solution}

    \question 证明恒等式
    \[
    32 \sin^4 x \cos^2 x \equiv 2 - \cos 2x - 2 \cos 4x + \cos 6x,
    \]
    \begin{solution}
        化简右边得
        \begin{align*}
        2 - \cos 2x - 2 \cos 4x + \cos 6x 
        &= 2(1 - \cos 4x) - 2 \sin 4x \sin 2x \\
        &= 2(2 \sin^2 2x) - 2 (2 \sin 2x \cos 2x) \sin 2x \\
        &= 4 \sin^2 2x (1 - \cos 2x) \\
        &= 4 (2 \sin x \cos x)^2 \cdot 2 \sin^2 x \\
        &= 32 \sin^4 x \cos^2 x
        \end{align*}
        所以恒等式得证。
    \end{solution}

    \question 将$16 \sin^5 \theta $表示成$a \sin \theta + b \sin 3\theta + c \sin 5\theta$的形式,其中 $a, b, c$ 为实数。
    \begin{solution}
        有
        \begin{align*}
        16 \sin^5 \theta 
        &= 16 \sin^3 \theta (1 - \cos^2 \theta) \\
        &= 16 \sin^3 \theta - 16 \sin^3 \theta \cos^2 \theta \\
        &= 16 \sin \theta (1 - \cos^2 \theta) - 4 \sin \theta \sin^2 2\theta \\
        &= 16 \sin \theta - 16 \sin \theta \cos^2 \theta - 4 \sin \theta \sin^2 2\theta \\
        &= 14 \sin \theta - 8 \sin 2\theta \cos \theta + 2 \sin \theta (1 - 2 \sin^2 2\theta) \\
        &= 14 \sin \theta - 4 (\sin 3\theta + \sin \theta) + 2 \sin \theta \cos 2\theta \\
        &= 14 \sin \theta - 4 \sin 3\theta - 4 \sin \theta + \sin 3\theta - \sin \theta \\
        &= 10 \sin \theta - 5 \sin 3\theta + \sin 5\theta
        \end{align*}
    \end{solution}

    \question 证明恒等式
    \[
    \tan \theta + \tan(\theta + 120^\circ) + \tan(\theta + 240^\circ) = 3 \tan 3\theta.
    \]
    \begin{solution}
        由左式化简得
        \begin{align*}
        &\tan \theta + \tan(\theta + 120^\circ) + \tan(\theta + 240^\circ) \\
        &= \tan \theta + \frac{\tan \theta - \sqrt{3}}{1 + \sqrt{3} \tan \theta} + \frac{\tan \theta + \sqrt{3}}{1 - \sqrt{3} \tan \theta} \\
        &= \tan \theta + \frac{(\tan \theta - \sqrt{3})(1 - \sqrt{3} \tan \theta) + (\tan \theta + \sqrt{3})(1 + \sqrt{3} \tan \theta)}{1 - 3 \tan^2 \theta} \\
        &= \tan \theta + \frac{8 \tan \theta}{1 - 3 \tan^2 \theta} \\
        &= \frac{9 \tan \theta - 3 \tan^3 \theta}{1 - 3 \tan^2 \theta} \\
        &= 3 \left( \frac{3 \tan \theta - \tan^3 \theta}{1 - 3 \tan^2 \theta} \right) \\
        &= 3 \tan 3\theta
        \end{align*}
        故恒等式得证。
    \end{solution}

    \question 已知
    \[
    x=\csc\theta-\sin\theta,\quad y=\sec\theta-\cos\theta,\quad 0<\theta<\frac{\pi}{2}
    \]
    利用三角恒等式证明
    \[
    y^2x^2\left(x^{\frac{2}{3}}+y^{\frac{2}{3}}\right)^3=1
    \]
    \begin{solution}
        发现
        \[
        x=\csc\theta-\sin\theta=\frac{1}{\sin\theta}-\sin\theta=
        \frac{1-\sin^2\theta}{\sin\theta}
        =\frac{\cos^2\theta}{\sin\theta}
        \]
        同理有
        \[
        y=\sec\theta-\cos\theta=\frac{1}{\cos\theta}-\cos\theta
        =\frac{1-\cos^2\theta}{\cos\theta}
        =\frac{\sin^2\theta}{\cos\theta}
        \]
        于是
        \[
        \frac{y}{x}
        =\frac{\sin^3\theta}{\cos^3\theta}
        =\tan^3\theta \Rightarrow \tan\theta=\left(\frac{y}{x}\right)^{\frac{1}{3}}
        \]
        再由$y=\tan\theta\sin\theta$两边平方得
        \[
        y^2=\tan^2\theta\sin^2\theta
        =\tan^2\theta\left(1-\frac{1}{\sec^2\theta}\right)
        =\tan^2\theta\left(1-\frac{1}{1+\tan^2\theta}\right)
        =\frac{\tan^4\theta}{1+\tan^2\theta}
        \]
        将 $\tan\theta=\left(\dfrac{y}{x}\right)^{\frac{1}{3}}$代入得,
        \[
        y^2=\frac{\left(\frac{y}{x}\right)^{\frac{4}{3}}}{1+\left(\frac{y}{x}\right)^{\frac{2}{3}}}
        \]
        整理得
        \[
        x^{\frac{4}{3}}y^{\frac{2}{3}}(x^{\frac{2}{3}}+y^{\frac{2}{3}})=1
        \]
        两边立方得
        \[
        y^2x^2\left(x^{\frac{2}{3}}+y^{\frac{2}{3}}\right)^3=1
        \]
        证毕。
    \end{solution}

    \question 证明以下恒等式:
    \begin{parts}
    \part 
    \[
    \sin x + \sin\left(x + \frac{2\pi}{3}\right) + \sin\left(x + \frac{4\pi}{3}\right) \equiv 0
    \]
    \begin{solution}
        使用和角公式,
        \begin{align*}
        \sin x + \sin\left(x + \frac{2\pi}{3}\right) + \sin\left(x + \frac{4\pi}{3}\right) 
        &= \sin x + 2 \sin(x + \pi) \cos\frac{\pi}{3} \\
        &= \sin x - \sin x  \\
        &= 0
        \end{align*}
    \end{solution}
    \part 
    \[
    \sin^3 x + \sin^3\left(x + \frac{2\pi}{3}\right) + \sin^3\left(x + \frac{4\pi}{3}\right) \equiv -\frac{3}{4} \sin 3x
    \]
    \end{parts}
    \begin{solution}
        使用立方正弦公式 
        \[
        \sin^3 \theta = \frac{1}{4}\left(3 \sin \theta - \sin 3\theta\right),
        \]
        \begin{align*}
        & \sin^3 x + \sin^3\left(x + \frac{2\pi}{3}\right) + \sin^3\left(x + \frac{4\pi}{3}\right) \\
        &= \frac{1}{4} \left[3 \left[\sin x + \sin\left(x + \frac{2\pi}{3}\right) + \sin\left(x + \frac{4\pi}{3}\right)\right] - (\sin 3x + \sin (3x+2\pi) + \sin (3x+4\pi)) \right] \\
        &= \frac{1}{4} \left[3 \cdot 0 - 3 \sin 3x \right] \\
        &= -\frac{3}{4} \sin 3x
        \end{align*}
    \end{solution}

    \question 证明下列恒等式:
    \begin{parts}
    \part 
    \[
    \cos \theta + \cos\left(\theta + \frac{2\pi}{3}\right) + \cos\left(\theta + \frac{4\pi}{3}\right) \equiv 0
    \]
    \begin{solution}
        和差化积得
        \begin{align*}
        \cos \theta + \cos\left(\theta + \frac{2\pi}{3}\right) + \cos\left(\theta + \frac{4\pi}{3}\right) 
        &= \cos \theta + 2 \cos\left(\theta + \pi\right)\cos\frac{\pi}{3} \\
        &= \cos \theta - \cos \theta \\
        &= 0
        \end{align*}
        故恒等式得证。
    \end{solution}
    \part 
    \[
    \cos^2 \theta + \cos^2\left(\theta + \frac{2\pi}{3}\right) + \cos^2\left(\theta + \frac{4\pi}{3}\right) \equiv \frac{3}{2}
    \]
    \begin{solution}
        由平方公式 
        \[
        \cos^2 x = \frac{1}{2}(1+\cos 2x),
        \]
        则
        \begin{align*}
        &\cos^2 \theta + \cos^2\left(\theta + \frac{2\pi}{3}\right) + \cos^2\left(\theta + \frac{4\pi}{3}\right)\\
        &= \frac{1}{2} \left[3 + \cos 2\theta + \cos\left(2\theta + \frac{4\pi}{3}\right) + \cos\left(2\theta + \frac{8\pi}{3}\right) \right] \\
        &= \frac{1}{2} \left[3 + \cos 2\theta + 2\cos (2\theta+2\pi)\cos \frac{2\pi}{3} \right] \\
        &= \frac{1}{2} \left[3 + \cos 2\theta - \cos 2\theta \right] \\ 
        &= \frac{3}{2}
        \end{align*}
        因此恒等式得证。
    \end{solution}
    \part 
    \[
    \cos^3 \theta + \cos^3\left(\theta + \frac{2\pi}{3}\right) + \cos^3\left(\theta + \frac{4\pi}{3}\right) \equiv \frac{3}{4} \cos 3\theta
    \]
    \begin{solution}
        由三次余弦公式,
        \[
        \cos^3 x = \frac{1}{4} \left(3\cos x + \cos 3x\right).
        \]
        因此:
        \begin{align*}
        &\cos^3 \theta + \cos^3\left(\theta + \frac{2\pi}{3}\right) + \cos^3\left(\theta + \frac{4\pi}{3}\right) \\ 
        &= \frac{1}{4} \left[3\left[\cos \theta + \cos\left(\theta + \frac{2\pi}{3}\right) + \cos\left(\theta + \frac{4\pi}{3}\right)\right]+\cos 3\theta + \cos (3\theta+2\pi) + \cos (3\theta+4\pi) \right] \\
        &= \frac{1}{4} \left[3 \cdot 0 + 3 \cos 3\theta \right] \\
        &= \frac{3}{4} \cos 3\theta
        \end{align*}
        故恒等式得证。
    \end{solution}
    \end{parts}

    \question 已知$A,B,C$为$\triangle ABC$的内角,证明恒等式
    \begin{parts}
    \part 
    \[
    \sin 2A - \sin 2B + \sin 2C = 4 \cos A \sin B \cos C
    \]
    \begin{solution}
        已知 $A + B + C = 180^\circ$,所以 
        \[
        \sin(A+B)=\sin C, \quad \cos(A + B) = -\cos C
        \]
        于是
        \begin{align*}
        \sin 2A - \sin 2B + \sin 2C
        &= 2\cos(A+B)\sin(A-B) + 2\sin C\cos C \\
        &= -2\cos C\sin(A-B) + 2\sin C\cos C \\
        &= 2\cos C[\sin C - \sin(A-B)] \\
        &= 2\cos C[\sin(A+B) - \sin(A-B)] \\
        &= 2\cos C \cdot 2\cos A\sin B \\
        &= 4\cos A\sin B\cos C
        \end{align*}
        故得证。
    \end{solution}
    \part 
    \[
    \sin A + \sin B + \sin C = 4 \cos \frac{A}{2} \cos \frac{B}{2} \cos \frac{C}{2}
    \]
    \begin{solution}
        已知 $A + B + C = 180^\circ$,所以 
        \[
        \sin\frac{A+B}{2}=\cos \frac{C}{2}, \quad \cos\frac{A+B}{2}=\sin \frac{C}{2}
        \]
        故
        \begin{align*}
        \sin A + \sin B + \sin C &= 2 \sin \frac{A+B}{2} \cos \frac{A-B}{2} + \sin C \\
        &= 2 \cos \frac{C}{2} \cos \frac{A-B}{2} + 2 \sin \frac{C}{2} \cos \frac{C}{2} \\
        &= 2 \cos \frac{C}{2} \left(\cos \frac{A-B}{2} + \sin \frac{C}{2}\right) \\
        &= 2 \cos \frac{C}{2} \left(\cos \frac{A-B}{2} + \cos \frac{A+B}{2}\right) \\
        &= 2 \cos \frac{C}{2} \cdot 2 \cos \frac{A}{2} \cos \frac{B}{2} \\
        &= 4 \cos \frac{A}{2} \cos \frac{B}{2} \cos \frac{C}{2}
        \end{align*}
        于是恒等式得证。
    \end{solution}
    \part 
    \[
    \cos A - \cos B + \cos C = 4 \cos \frac{A}{2} \sin \frac{B}{2} \cos \frac{C}{2} - 1
    \]
    \begin{solution}
        已知 $A + B + C = 180^\circ$,所以 
        \[
        \cos\frac{A+B}{2}=\sin \frac{C}{2}, \quad \sin\frac{A+B}{2}=\cos \frac{C}{2}
        \]
        故
        \begin{align*}
        \cos A - \cos B + \cos C 
        &= -2 \sin \frac{A+B}{2} \sin \frac{A-B}{2} + \cos C \\
        &= -2 \cos \frac{C}{2} \sin \frac{A-B}{2} + 2 \cos^2 \frac{C}{2} - 1 \\
        &= 2 \cos \frac{C}{2} \left(-\sin \frac{A-B}{2} + \cos \frac{C}{2}\right) - 1 \\
        &= 2 \cos \frac{C}{2} \left(-\sin \frac{A-B}{2} + \sin \frac{A+B}{2}\right) - 1 \\
        &= 2 \cos \frac{C}{2} \cdot 2 \cos \frac{A}{2} \sin \frac{B}{2} - 1 \\
        &= 4 \cos \frac{A}{2} \sin \frac{B}{2} \cos \frac{C}{2} - 1
        \end{align*}
        得证。
    \end{solution}
    \part 
    \[
    (\cot B + \cot C)(\cot C + \cot A)(\cot A + \cot B) = \csc A \csc B \csc C
    \]
    \begin{solution}
        由$B+C = 180^\circ - A$,所以 $\sin(B+C) = \sin A$,于是
        \[
        \cot B + \cot C = \frac{\cos B \sin C + \sin B \cos C}{\sin B \sin C} = \frac{\sin(B+C)}{\sin B \sin C}=\frac{\sin A}{\sin B \sin C} 
        \]
        同理可得
        \[
        \cot C + \cot A = \frac{\sin B}{\sin C \sin A},\quad
        \cot A + \cot B = \frac{\sin C}{\sin A \sin B}
        \]
        三项相乘得
        \[
        (\cot B + \cot C)(\cot C + \cot A)(\cot A + \cot B) 
        = \frac{1}{\sin A \sin B \sin C}
        = \csc A \csc B \csc C
        \]
        故得证。
    \end{solution}
    \part 
    \[
    \sin^2 \frac{A}{2} + \sin^2 \frac{B}{2} + \sin^2 \frac{C}{2} = 1 - 2 \sin \frac{A}{2} \sin \frac{B}{2} \sin \frac{C}{2}
    \]
    \begin{solution}
        由半角公式,
        \[
        \sin^2 \frac{A}{2} + \sin^2 \frac{B}{2} + \sin^2 \frac{C}{2} = \frac{1}{2}\left[3 - (\cos A + \cos B + \cos C)\right]
        \]
        而
        \begin{align*}
        \cos A + \cos B + \cos C 
        &= 2 \cos \frac{A+B}{2} \cos \frac{A-B}{2} + \cos C \\
        &= 2 \sin \frac{C}{2} \cos \frac{A-B}{2} + 1 - 2 \sin^2 \frac{C}{2} \\
        &= 1 + 2 \sin \frac{C}{2} (\cos \frac{A-B}{2} - \sin \frac{C}{2}) \\
        &= 1 + 2 \sin \frac{C}{2} (\cos \frac{A-B}{2} - \cos \frac{A+B}{2}) \\
        &= 1 + 2 \sin \frac{C}{2} \cdot 2\sin \frac{A}{2} \sin \frac{B}{2} \\
        &= 1 + 4 \sin \frac{A}{2} \sin \frac{B}{2} \sin \frac{C}{2}
        \end{align*}
        于是得证
        \[
        \sin^2 \frac{A}{2} + \sin^2 \frac{B}{2} + \sin^2 \frac{C}{2} = \frac{1}{2}\left(3 -1 - 4 \sin \frac{A}{2} \sin \frac{B}{2} \sin \frac{C}{2}\right) = 1 - 2 \sin \frac{A}{2} \sin \frac{B}{2} \sin \frac{C}{2}
        \]
    \end{solution}
    \part 
    \[
    \cot A + \cot B + \cot C = \csc A \csc B \csc C + \cot A \cot B \cot C.
    \]
    \begin{solution}
        首先有
        \begin{align*}
        \cot A + \cot B + \cot C 
        &= \frac{\cos A}{\sin A} + \frac{\cos B}{\sin B} + \frac{\cos C}{\sin C} \\
        &= \frac{\cos A \sin B \sin C + \sin A \cos B \sin C + \sin A \sin B \cos C}{\sin A \sin B \sin C}
        \end{align*}
        其中分子为
        \begin{align*}
        & \cos A \sin B \sin C + \sin A \cos B \sin C + \sin A \sin B \cos C \\
        &= \sin C (\cos A \sin B + \sin A \cos B) - \frac{1}{2} (\cos (A+B) - \cos (A-B)) \cos C \\
        &= \sin C \sin (A+B) - \frac{1}{2} (-\cos C - \cos (A-B)) \cos C \\
        &= \sin^2 C + \frac{1}{2} \cos^2 C + \frac{1}{2} \cos (A-B) \cos C \\
        &= \sin^2 C + \cos^2 C - \frac{1}{2} \cos^2 C + \frac{1}{2} \cos (A-B) \cos C \\
        &= 1 + \frac{1}{2} \cos C (\cos (A-B) - \cos C) \\
        &= 1 + \frac{1}{2} \cos C \left(-2 \sin \frac{A-B+C}{2} \sin \frac{A-B-C}{2}\right) \\
        &= 1 - \cos C \sin \frac{A-B+C}{2} \sin \frac{A-B+C}{2} \\
        &= 1 - \cos C \sin \frac{180^\circ-2B}{2} \sin \frac{-180^\circ+2A}{2} \\
        &= 1 + \cos C \sin (90^\circ-B) \sin (90^\circ-A) \\
        &= 1 + \cos A \cos B \cos C
        \end{align*}
        于是
        \[
        \cot A + \cot B + \cot C = \frac{1 + \cos A \cos B \cos C}{\sin A \sin B \sin C} = \csc A \csc B \csc C + \cot A \cot B \cot C.
        \]
        即左式=右式,恒等式得证。\textcolor{red}{有更简洁的证明?}
    \end{solution}
    \end{parts}

    \question 若$A,B,C$为任意角,证明
    \[
    \sin A + \sin B + \sin C - \sin(A+B+C) = 4 \sin \frac{A+B}{2} \sin \frac{B+C}{2} \sin \frac{C+A}{2} 
    \]
    \begin{solution}
        \begin{align*}
        &\sin A + \sin B + \sin C - \sin(A+B+C) \\
        &= 2 \sin \frac{A+B}{2} \cos \frac{A-B}{2} + \sin C - \sin(A+B)\cos C - \cos(A+B)\sin C \\
        &= 2 \sin \frac{A+B}{2} \cos \frac{A-B}{2} - 2 \sin \frac{A+B}{2} \cos \frac{A+B}{2} \cos C + \sin C (1-\cos(A+B)) \\
        &= 2 \sin \frac{A+B}{2} \cos \frac{A-B}{2} - 2 \sin \frac{A+B}{2} \cos \frac{A+B}{2} \cos C + \sin C \left(2\sin^2 \frac{A+B}{2}\right) \\
        &= 2 \sin \frac{A+B}{2} \left( \cos \frac{A-B}{2} - \cos \frac{A+B}{2} \cos C + \sin \frac{A+B}{2} \sin C \right) \\
        &= 2 \sin \frac{A+B}{2} \left( \cos \frac{A-B}{2} - \cos \left(\frac{A+B}{2}+C\right) \right) \\
        &= 2 \sin \frac{A+B}{2} \left( 2 \sin \frac{A+C}{2} \sin \frac{B+C}{2} \right) \\
        &= 4 \sin \frac{A+B}{2} \sin \frac{B+C}{2} \sin \frac{C+A}{2}
        \end{align*}
        故得证。
    \end{solution}

    \question 若$A,B,C$为任意角,证明
    \begin{align*}
    &\cos^2 A + \cos^2 B + \cos^2 C - 2 \cos A \cos B \cos C - 1 \\
    &= -4 \sin \frac{A+B+C}{2} \sin \frac{A+B-C}{2} \sin \frac{B+C-A}{2} \sin \frac{C+A-B}{2}.
    \end{align*}
    \begin{solution}
        \begin{align*}
        &\cos^2 A + \cos^2 B + \cos^2 C - 2 \cos A \cos B \cos C - 1 \\
        &= \frac{1}{2}(2+\cos 2A+\cos 2B) + \cos^2 C - (\cos(A+B) + \cos(A-B)) \cos C - 1 \\
        &= \frac{1}{2}(\cos 2A+\cos 2B) + \cos^2 C - (\cos(A+B) + \cos(A-B)) \cos C \\
        &= \cos(A+B)\cos(A-B) + \cos^2 C - \cos C \cos(A+B) - \cos C \cos(A-B) \\
        &= (\cos(A+B) - \cos C)(\cos(A-B) - \cos C) \\
        &= \left(-2 \sin \frac{A+B+C}{2} \sin \frac{A+B-C}{2}\right) \left(2 \sin \frac{A-B+C}{2} \sin \frac{A-B-C}{2}\right) \\
        &= -4 \sin \frac{A+B+C}{2} \sin \frac{A+B-C}{2} \sin \frac{B+C-A}{2} \sin \frac{C+A-B}{2}
        \end{align*}
        故得证。
    \end{solution}

    \question 证明
    \[
    \sin 50^\circ (1 + \sqrt{3} \tan 10^\circ)=1.
    \]
    \begin{solution}
        有
        \begin{align*}
        \sin 50^\circ (1 + \sqrt{3} \tan 10^\circ) &= \sin 50^\circ \left(\frac{\cos 10^\circ+\sqrt{3} \sin 10^\circ}{\cos 10^\circ}\right) \\
        &= 2 \sin 50^\circ \left(\frac{\cos 60^\circ \cos 10^\circ + \sin 60^\circ \sin 10^\circ}{\cos 10^\circ}\right) \\
        &= 2 \sin 50^\circ \frac{\cos 50^\circ}{\cos 10^\circ} \\
        &= \frac{\sin 100^\circ}{\cos 10^\circ} \\
        &= 1
        \end{align*}
    \end{solution}

    \question 证明
    \[
    4 \sin 40^{\circ} - \tan 40^{\circ} = \sqrt{3}.
    \]
    \begin{solution}
        \begin{align*}
        \tan 40^{\circ} + \sqrt{3} 
        & = \frac{\sin 40^{\circ}}{\cos 40^{\circ}} + \frac{\sin 60^{\circ}}{\cos 60^{\circ}} \\
        & = \frac{\sin 40^{\circ} \cos 60^{\circ} + \cos 40^{\circ} \sin 60^{\circ}}{\cos 40^{\circ} \cos 60^{\circ}} \\
        & = \frac{\sin 100^{\circ}}{\cos 40^{\circ} \cdot \frac{1}{2}} \\
        & = \frac{2 \sin 100^{\circ}}{\cos 40^{\circ}} \\
        & = \frac{2 \sin 80^{\circ}}{\cos 40^{\circ}} \\
        & = 4 \sin 40^{\circ} 
        \end{align*}
        于是得证
        \[
        4 \sin 40^{\circ} - \tan 40^{\circ} = \sqrt{3}
        \]
    \end{solution}
    
    \question 证明
    \[
    \tan 24^\circ (\sqrt{3} \sec 36^\circ + \tan 6^\circ)=1.
    \]
    \begin{solution}
        先证
        \[
        \cos 36^\circ - \cos 72^\circ = \frac{1}{2}
        \]
        有
        \begin{align*}
        2\cos^2 72^\circ - 1 &= \cos 144^\circ = -\cos 36^\circ \tag{1} \\
        2\cos^2 36^\circ - 1 &= \cos 72^\circ \tag{2}
        \end{align*}
        $(2)-(1)$得
        \[
        2(\cos 36^\circ + \cos 72^\circ)(\cos 36^\circ - \cos 72^\circ) = \cos 72^\circ + \cos 36^\circ
        \]
        即
        \begin{align*}
        \cos 36^\circ - \cos 72^\circ &= \frac{1}{2} = \cos 60^\circ 
        \end{align*}
        于是
        \[
        \cos 36^\circ = \cos 72^\circ + \cos 60^\circ = 2\cos 66^\circ \cos 6^\circ
        \]
        故
        \begin{align*}
        \tan 24^\circ (\sqrt{3} \sec 36^\circ + \tan 6^\circ) 
        &= \tan 24^\circ \left( \frac{2\sin 60^\circ}{2\cos 66^\circ \cos 6^\circ} + \tan 6^\circ \right) \\
        &= \tan 24^\circ \left( \frac{\sin 66^\circ \cos 6^\circ - \cos 66^\circ \sin 6^\circ}{\cos 66^\circ \cos 6^\circ} + \tan 6^\circ \right) \\
        &= \tan 24^\circ (\tan 66^\circ - \tan 6^\circ + \tan 6^\circ) \\
        &= 1
        \end{align*}
    \end{solution}
    
    \question 证明 
    \[
    \cos 20^\circ \cos 40^\circ \cos 80^\circ = \frac{1}{8}.
    \]
    \begin{solution}
        积化和差、再和差化积,
        \begin{align*}
        \cos 20^\circ \cos 40^\circ \cos 80^\circ
        &= \frac{1}{2}\cos 20^\circ (\cos 120^\circ + \cos 40^\circ) \\
        &= \frac{1}{2}\left(-\frac{1}{2}\right)\cos 20^\circ + \frac{1}{2}\cos 40^\circ \cos 20^\circ \\
        &= -\frac{1}{4}\cos 20^\circ + \frac{1}{4}\left(\frac{1}{2}+\cos 20^\circ\right) \\
        &= \frac{1}{8}
        \end{align*}
    \end{solution}
    \begin{solution}
        由恒等式$\sin2\theta=2\sin\theta\cos\theta$,
        \begin{align*}
        \cos 20^\circ \cos 40^\circ \cos 80^\circ
        &= \frac{1}{2\sin 20^\circ}\sin 40^\circ \cos 40^\circ \cos 80^\circ \\
        &= \frac{1}{4\sin 20^\circ}\sin 80^\circ \cos 80^\circ \\
        &= \frac{1}{8\sin 20^\circ}\sin 160^\circ \\
        &= \frac{1}{8}
        \end{align*}
    \end{solution}

    \question 证明 
    \[
    \sin 20^\circ \sin 40^\circ \sin 80^\circ = \frac{\sqrt{3}}{8}.
    \]
    \begin{solution}
        积化和差、再和差化积,
        \begin{align*}
        \sin 20^\circ \sin 40^\circ \sin 80^\circ
        &= \frac{1}{2} \sin 20^\circ (\cos 40^\circ+\cos 60^\circ) \\
        &= \frac{1}{4}(\sin 60^\circ-\sin 20^\circ)+\frac{1}{4}\sin 20^\circ \\
        &= \frac{1}{4}\cdot\frac{\sqrt{3}}{2} \\
        &= \frac{\sqrt{3}}{8}
        \end{align*}
    \end{solution}
    \begin{solution}
        由恒等式 
        \[
        \sin x \sin(60^\circ-x)\sin(60^\circ+x)=\frac{1}{4}\sin 3x
        \]
        取 $x=20^\circ$,则
        \[
        \sin 20^\circ \sin 40^\circ \sin 80^\circ
        = \frac{1}{4}\sin 60^\circ
        = \frac{\sqrt{3}}{8}
        \]
    \end{solution}

    \question 计算
    \[
    \sin 10^\circ \sin 30^\circ \sin 50^\circ \sin 70^\circ
    \]
    的值。
    \begin{solution}
        由恒等式$\sin2\theta=2\sin\theta\cos\theta$,
        \begin{align*}
        \sin 10^\circ \sin 30^\circ \sin 50^\circ \sin 70^\circ
        &= \sin 10^\circ \cdot \frac{1}{2} \cdot \cos 40^\circ \cos 20^\circ \\
        &= \frac{1}{4\cos 10^\circ} \sin 20^\circ \cos 20^\circ \cos 40^\circ \\
        &= \frac{1}{8\cos 10^\circ} \sin 40^\circ \cos 40^\circ \\
        &= \frac{1}{16\cos 10^\circ} \sin 80^\circ \\
        &= \frac{1}{16}
        \end{align*}
    \end{solution}

    \question 计算
    \[
    \cos 10^\circ \cos 30^\circ \cos 50^\circ \cos 70^\circ
    \]
    的值。
    \begin{solution}
        积化和差、再和差化积,
        \begin{align*}
        \cos 10^\circ \cos 30^\circ \cos 50^\circ \cos 70^\circ
        &= \frac{1}{2}(\cos60^\circ + \cos40^\circ) \cdot \frac{\sqrt{3}}{2} \cos 70^\circ  \\
        &= \frac{\sqrt{3}}{4}\left(\frac{1}{2}\cos 70^\circ+\frac{1}{2}(\cos110^\circ + \cos30^\circ)\right) \\
        &= \frac{\sqrt{3}}{8}(\cos 70^\circ-\cos70^\circ + \frac{\sqrt{3}}{2}) \\
        &= \frac{3}{16}
        \end{align*}
    \end{solution}

    \question 证明 
    \[
    \tan 50^\circ + \tan 60^\circ + \tan 70^\circ = \tan 80^\circ.
    \]
    \begin{solution}
        若 $A+B+C=180^\circ$,则有
        \[
        \tan A \tan B \tan C = \tan A + \tan B + \tan C。
        \]
        于是
        \begin{align*}
        \tan 50^\circ + \tan 60^\circ + \tan 70^\circ
        &= \tan 50^\circ \tan 60^\circ \tan 70^\circ \\[1pt]
        &= \frac{\sin 50^\circ \sin 60^\circ \sin 70^\circ}{\cos 50^\circ \cos 60^\circ \cos 70^\circ} \\[2pt]
        &= \frac{\sin 60^\circ \left(\cos 20^\circ + \frac{1}{2}\right)}{\cos 60^\circ \left(\cos 20^\circ - \frac{1}{2}\right)} \\[2pt]
        &= \frac{\frac{1}{2}(\sin 80^\circ + \sin 40^\circ) + \frac{1}{2}\sin 60^\circ}{\frac{1}{2}(\cos 80^\circ + \cos 40^\circ) + \frac{1}{2}\cos 60^\circ} \\[2pt]
        &= \frac{\sin 40^\circ + \sin 60^\circ + \sin 80^\circ}{\cos 40^\circ - \cos 60^\circ + \cos 80^\circ} \\[2pt]
        &= \frac{2 \sin 50^\circ \cos 10^\circ + \sin 80^\circ}{2 \sin 50^\circ \sin 10^\circ + \cos 80^\circ} \\[2pt]
        &= \frac{\sin 80^\circ (2 \sin 50^\circ + 1)}{\cos 80^\circ (2 \sin 50^\circ + 1)} \\[2pt]
        &= \tan 80^\circ
        \end{align*}
    \end{solution}

    \question 计算 
    \[
    \tan 9^\circ - \tan 27^\circ - \tan 63^\circ + \tan 81^\circ
    \]
    的值。
    \begin{solution}
        首先有
        \begin{align*}
        S &=\tan 9^\circ - \tan 27^\circ - \tan 63^\circ + \tan 81^\circ \\
        &= \tan 9^\circ + \cot 9^\circ - (\tan 27^\circ + \cot 27^\circ) \\[1pt]
        &= \frac{\sin^2 9^\circ + \cos^2 9^\circ}{\sin 9^\circ \cos 9^\circ}
        - \frac{\sin^2 27^\circ + \cos^2 27^\circ}{\sin 27^\circ \cos 27^\circ} \\[2pt]
        &= \frac{1}{\sin 9^\circ \cos 9^\circ} - \frac{1}{\sin 27^\circ \cos 27^\circ} \\[2pt]
        &= \frac{2}{\sin 18^\circ} - \frac{2}{\sin 54^\circ}
        \end{align*}
        现求$\sin 18^\circ$及$\sin 54^\circ$,设$\theta=18^\circ$,解
        \[
        \sin 2\theta = \sin (90^\circ-3\theta)=\cos 3\theta
        \]
        可得
        \[
        \sin 18^\circ=\frac{\sqrt{5}-1}{4}>0,\quad \sin 54^\circ = 3\sin 18^\circ-4\sin^3 18^\circ = \frac{\sqrt{5}+1}{4}
        \]
        故
        \[
        S = \frac{2}{\dfrac{\sqrt{5}-1}{4}} - \frac{2}{\dfrac{\sqrt{5}+1}{4}}= 4
        \]
    \end{solution}

    \question 试证 
    \[\dfrac{3}{\sin^2 40^\circ} - \dfrac{1}{\cos^2 40^\circ}=32 \sin 10^\circ.
    \]
    \begin{solution}
        有
        \begin{align*}
        \frac{3}{\sin^2 40^\circ} - \frac{1}{\cos^2 40^\circ} 
        &= \frac{3 \cos^2 40^\circ - \sin^2 40^\circ}{\sin^2 40^\circ \cos^2 40^\circ} \\[6pt]
        &= 16 \cdot \frac{\left(\displaystyle\frac{\sqrt{3}}{2} \cos 40^\circ + \frac{1}{2} \sin 40^\circ\right)
        \left(\displaystyle\frac{\sqrt{3}}{2} \cos 40^\circ - \frac{1}{2} \sin 40^\circ\right)}{\sin^2 80^\circ} \\[6pt]
        &= 16 \cdot \frac{\sin(60^\circ + 40^\circ)\,\sin(60^\circ - 40^\circ)}{\sin 80^\circ \cos 10^\circ} \\[6pt]
        &= 32 \sin 10^\circ
        \end{align*}
    \end{solution}

    \question 计算 $$\sin^2 37^\circ + \sin^2 8^\circ + \sqrt{2}\sin 37^\circ \sin 8^\circ$$ 之值。
    \begin{solution}
        由
        \[
        \sin 37^\circ = \sin(45^\circ - 8^\circ) = \frac{\sqrt{2}}{2} \cos 8^\circ - \frac{\sqrt{2}}{2} \sin 8^\circ
        \]
        则
        \[
        \begin{aligned}
        \sin^2 37^\circ &= \frac{1}{2} \cos^2 8^\circ - \sin 8^\circ \cos 8^\circ + \frac{1}{2} \sin^2 8^\circ \\
        \sqrt{2}\sin 37^\circ \sin 8^\circ &= \sin 8^\circ \cos 8^\circ - \sin^2 8^\circ
        \end{aligned}
        \]
        因此
        \[
        \sin^2 37^\circ + \sin^2 8^\circ + \sqrt{2}\sin 37^\circ \sin 8^\circ = \frac{1}{2} \cos^2 8^\circ + \frac{1}{2} \sin^2 8^\circ = \frac{3}{4}
        \]
    \end{solution}
    \begin{solution}
        考虑 $\triangle ABC$ ,其中外接圆半径 $R = \dfrac{1}{2}$,且
        \[
        \angle A = 135^\circ,\quad \angle B = 37^\circ,\quad \angle C = 8^\circ
        \]
        由正弦定理,
        \[
        a = \sin 135^\circ = \frac{\sqrt{2}}{2},\quad b = \sin 37^\circ, \quad c = \sin 8^\circ
        \]
        由余弦定理,
        \[
        \cos \angle A = -\frac{\sqrt{2}}{2} = \frac{\sin^2 37^\circ + \sin^2 8^\circ - \sin^2 135^\circ}{2 \cdot \sin 37^\circ \cdot \sin 8^\circ}
        \]
        即
        \[
        \sin^2 37^\circ + \sin^2 8^\circ + \sqrt{2} \sin 37^\circ \sin 8^\circ = \frac{1}{2}
        \]
    \end{solution}
    
    \question 在 $\triangle ABC$ 中,设 ${BC}=a,\ {AC}=b,\ {AB}=c$;若 $a,b,c$ 成等差数列,试求 $\tan\dfrac{A}{2}\tan\dfrac{C}{2}$ 的值。
    \begin{solution}
        据题意有$a + c = 2b$,由正弦定理,
        \[
        \sin A + \sin C = 2\sin B 
        \]
        又因 $A + B + C = 180^\circ,$
        \[
        2\sin B = 2\sin(A + C) = \sin A + \sin C
        \]
        于是
        \[
        4\sin\frac{A + C}{2}\cos\frac{A + C}{2} = 2\sin\frac{A + C}{2}\cos\frac{A - C}{2}
        \Rightarrow 2\cos\frac{A + C}{2} = \cos\frac{A - C}{2}
        \]
        将余弦展开得
        \[
        2\left( \cos\frac{A}{2}\cos\frac{C}{2} - \sin\frac{A}{2}\sin\frac{C}{2} \right)
        =
        \cos\frac{A}{2}\cos\frac{C}{2} + \sin\frac{A}{2}\sin\frac{C}{2}
        \]
        即
        \[
        3\sin\frac{A}{2}\sin\frac{C}{2} = \cos\frac{A}{2}\cos\frac{C}{2}
        \Rightarrow
        \tan\frac{A}{2} \tan\frac{C}{2} = \frac{1}{3}
        \]
    \end{solution}

    \question 已知$\triangle ABC$中,$AB=5,BC=9$,且$\cot A, \cot B, \cot C$成等差数列,试求$\tan^{2}\dfrac{B}{2}$。
    \begin{solution}
        据题意,有
        \[
        2 \cot B = \cot A + \cot C \Rightarrow \frac{\cos A}{\sin A} + \frac{\cos C}{\sin C} = \frac{\sin (A+C)}{\sin A \sin C} = \frac{\sin B}{\sin A \sin C} = \frac{2 \cos B}{\sin B}
        \]
        由余弦定理及正弦定理,
        \[
        \cos B = \frac{a^2 + c^2 - b^2}{2 a c}, \quad \frac{\sin^2 B}{2 \sin A \sin C} = \frac{b^2}{2 a c}
        \]
        因此
        \[
        2 b^2 = a^2 + c^2 = 9^2+5^2 \Rightarrow b^2 = 53 
        \]
        故
        \[
        \cos B = \frac{53}{90} \Rightarrow \tan^2 \frac{B}{2} = \frac{1 - \cos B}{1 + \cos B}= \frac{37}{143}
        \]
    \end{solution}

    \question 若 $\triangle ABC$ 的内角 $A,B,C$ 满足 $\sin A=\cos B=\tan C$ ,求 $\cos^{3}A+\cos^{2}A-\cos A$ 的值.
    \begin{solution}
        由 $\sin A=\cos B$ 知 $$A=\dfrac{\pi}{2}\pm B$$但 $\tan C$ 有意义,故 $C$ 不为直角,从而只
        能是 $$A=\frac{\pi}{2}+B$$
        进而有 $C=\pi-A-B=\dfrac{3\pi}{2}-2A$,所以 $$\sin A=\tan C=\tan\left(\frac{3\pi}{2}-2A\right)=\cot 2A$$
        于是
        \[1 = \sin A \cdot \tan 2A = \sin A \cdot \frac{\sin 2A}{\cos 2A} = \sin A \cdot \frac{2 \sin A \cos A}{2 \cos^2 A - 1}
        = \frac{2(1 - \cos^2 A) \cdot \cos A}{2 \cos^2 A - 1}\]
        即 
        \[
        2 \cos^2 A - 1 = 2(1 - \cos^2 A) \cdot \cos A \Rightarrow \cos^3 A + \cos^2 A - \cos A = \frac{1}{2}
        \]
    \end{solution}
    
    \question 已知 \(\triangle ABC\) 满足  
        \[
          \cos C = \frac{\sin A + \cos A}{2} = \frac{\sin B + \cos B}{2},
        \]
        求 \(\cos C\)。  
    \begin{solution}
        \textbf{解法一}
        
        由条件知 $$\cos C=\frac{\sqrt{2}}{2}\sin(A+\frac{\pi}{4})=\frac{\sqrt{2}}{2}\sin(B+\frac{\pi}{4})$$
        假如 $$(A+\frac{\pi}{4})+(B+\frac{\pi}{4})=\pi$$
        则 $C=\dfrac{\pi}{2},\cos C=0$, 但 $\sin(A+\dfrac{\pi}{4})>0$, 矛盾.
        所以只可能 $$A+\frac{\pi}{4}=B+\frac{\pi}{4}$$
        此时 $A=B\in(0,\dfrac{\pi}{2}),C=\pi-2A$,且注意到 $$\cos C=\frac{\sqrt{2}}{2}\sin(A+\frac{\pi}{4})>0$$故 $C<\dfrac{\pi}{2}$. 所以 $A=B\in(\dfrac{\pi}{4},\dfrac{\pi}{2})$, 结合条件得
        \begin{align*} 
        \cos C &= -\cos 2A = -\sin\left(2A+\frac{\pi}{2}\right) = -2\sin\left(A+\frac{\pi}{4}\right)\cos\left(A+\frac{\pi}{4}\right) \\ &= -2\sqrt{2}\cos C\cdot\left(-\sqrt{1-(\sqrt{2}\cos C)^{2}}\right)
        \end{align*}
        又$\cos C>0$, 化简得 $$8(1-2\cos^{2}C)=1\Rightarrow\cos C=\dfrac{\sqrt{7}}{4}$$
    \end{solution}  
    \begin{solution}
        \textbf{解法二}
        
        将第二个等式重写为
        \[
        \sin(A+45^{\circ}) = \sin(B+45^{\circ})
        \]
        假设 $A \neq B$,则由于 $A+45^{\circ}$ 和 $B+45^{\circ}$ 均在区间 $(45^{\circ}, 225^{\circ})$ 内, 
        \[
        (A+45^{\circ}) + (B+45^{\circ}) = 180^{\circ} \implies A+B = 90^{\circ}
        \]
        此时 $\cos(C) = 0$,故必须满足 $A=B=135^{\circ}$,矛盾;因此$A=B$,故
        \[
        \cos(C) = -\cos(A+B) = \sin^2 A - \cos^2 A = \frac{\sin A + \cos A}{2}
        \iff 
        \]
        \[
        \sin A - \cos A = \frac{1}{2} \iff (\cos A)^2 + \left(\cos A + \frac{1}{2}\right)^2 = 1
        \iff
        \]
        \[
        \cos A = \frac{-1 + \sqrt{7}}{4} \iff \cos C = \frac{\sqrt{7}}{4}
        \]
    \end{solution}
    \question 已知 $\triangle ABC$ 中, $AB=3$, 且外接圆半径为 $6$, 求
        \[
        \begin{vmatrix} 
        1 & \cos A & \cos B \\ 
        \cos A & -1 & \cos C \\ 
        \cos B & \cos C & -1 
        \end{vmatrix}
        \]
    \begin{solution}
        由正弦定理
        \[
        \frac{3}{\sin C} = 2 \cdot 6 \Rightarrow \sin C = \frac{1}{4}, \cos C = \frac{\sqrt{15}}{4}
        \]
        令
        \[
        \triangle = 
        \begin{vmatrix} 
        1 & \cos A & \cos B\\ 
        \cos A & -1 & \cos C\\ 
        \cos B & \cos C & -1
        \end{vmatrix} 
        = 1 + 2\cos A \cos B \cos C + \cos^2 A + \cos^2 B - \cos^2 C
        \]
        由 $ \cos C = -\cos(A+B)$,
        \begin{align*}
        \cos^2 C &= (\cos A \cos B - \sin A \sin B)^2 \\
        &= \cos^2 A \cos^2 B - 2 \cos A \cos B \sin A \sin B + (1-\cos^2 A)(1-\cos^2 B) \\
        &= 1 + 2 \cos^2 A \cos^2 B - 2 \cos A \cos B \sin A \sin B - \cos^2 A - \cos^2 B \\
        &= 1 + 2 \cos A \cos B (\cos A \cos B - \sin A \sin B) - \cos^2 A - \cos^2 B \\
        &= 1 + 2 \cos A \cos B \cos(A+B) - \cos^2 A - \cos^2 B\\
        \end{align*}
        故
        \[
        \triangle = 2 - 2 \cos^2 C= \frac{1}{8}
        \]
    \end{solution}

   \question 令 $A, B, C$ 为任意三角形 $\triangle ABC$ 的内角,满足
    \[
    \cos^{2}A + \cos^{2}B + 2 \sin A \sin B \cos C = \frac{15}{8},
    \]
    \[
    \cos^{2}B + \cos^{2}C + 2 \sin B \sin C \cos A = \frac{14}{9}.
    \]
    求
    \[
    \cos^{2}C + \cos^{2}A + 2 \sin C \sin A \cos B.
    \]
    \begin{solution}
        由$\sin (A + C)=\sin B$,两式相加得
        \[
        \cos^{2}A + 2 \cos^{2}B + \cos^{2}C + 2 \sin B \sin (A + C) = \cos^{2}A + 2 \cos^{2}B + \cos^{2}C + 2 \sin^{2} B = \frac{247}{72}
        \]
        因此
        \[
        \cos^{2}A + \cos^{2}C = \frac{247}{72} - 2 = \frac{103}{72} \tag{1}
        \]
        同理,两式相减得
        \[
        \cos^{2}A - \cos^{2}C + 2 \sin (A+C) \sin (A - C) = \cos^{2}A - \cos^{2}C + (\cos 2C - \cos 2A) =  \frac{23}{72}
        \]
        因此
        \[
        \cos^{2}A - \cos^{2}C + 2 \cos^{2}C - 1 - (2 \cos^{2}A - 1) = -\cos^{2}A + \cos^{2}C = \frac{23}{72} \tag{2}
        \]
        由 (1),(2) 解得
        \[
        \cos^{2}A = \frac{5}{9}, \quad \cos^{2}C = \frac{7}{8}
        \]
        因此
        \begin{align*}
        &\cos^{2}C + \cos^{2}A + 2 \sin C \sin A \cos B \\
        &= \frac{103}{72} + 2 \sin C \sin A (-\cos (A + C))\\
        &= \frac{103}{72} + 2 \sin C \sin A (\sin A \sin C - \cos A \cos C) \\
        &= \frac{103}{72} + 2 \sin^{2} A \sin^{2} C - 2 \sin C \cos C \sin A \cos A\\
        &= \frac{103}{72} + 2 \cdot \frac{4}{9} \cdot \frac{1}{8} \pm 2 \cdot \frac{1}{\sqrt{8}} \cdot \sqrt{\frac{7}{8}} \cdot \frac{2}{3} \cdot \frac{\sqrt{5}}{3} = \frac{111\pm4 \sqrt{35}}{72}
        \end{align*}
    \end{solution}

    \question 已知方程 $x \in \left(0, \dfrac{\pi}{2}\right)$
    \[
    \left(\tan^2 \frac{x}{2}\right) (\cot^4 x + 1)(\csc^2 x + \tan^2 x) = 1
    \] 
    有唯一实数解。求正整数$k$满足 
    \[
    \cos^{2019} x = \sin^k x
    \]
    \begin{solution}
        由
        \[
        \left(\tan^2 \frac{x}{2}\right) (\cot^4 x + 1)(\csc^2 x + \tan^2 x) = 1
        \]
        化为
        \[
        \frac{\sin^2 x}{(1+\cos x)^2}\left(\frac{\cos^4 x}{\sin^4 x}+1\right)\left(\frac{1}{\sin^2 x}+\frac{\sin^2 x}{\cos^2 x}\right)=1 
        \]
        \[
        (\cos^4 x +\sin^4 x)(\cos^2 x +\sin^4 x)=\sin^4 x \cos^2 x (1+\cos x)^2 
        \]
        展开化简得
        \[
        (\cos^3 x -\sin^4 x)^2=0 \Rightarrow \cos^3 x = \sin^4 x
        \]
        于是
        \[
        k=4 \cdot 673=2692
        \]
    \end{solution}
    \begin{solution}
        由柯西不等式,
        \[
        \cot^2 \frac{x}{2}=(\cot^4 x + 1)( \tan^2 x+\csc^2 x) \ge (\cot x +\csc x)^2 = \left(\frac{\cos x +1}{\sin x}\right)^2=\cot^2 \frac{x}{2}
        \]
        此时等号成立,于是有
        \[
        \frac{\cot^2 x}{\tan x}=\frac{1}{\csc x} \Rightarrow \cos^3 x = \sin^4 x
        \]
        同上,
        \[
        k=4 \cdot 673=2692
        \]
    \end{solution}

    \question 设 $\theta$ 为一锐角,满足
    \[
    \frac{16}{\sin^6 \theta} + \frac{1}{\cos^6 \theta} = 81,
    \]
    求 $\tan \theta$。
    \begin{solution}
        由柯西不等式,
        \[
        \left(\frac{16}{\sin^6 \theta} + \frac{1}{\cos^6 \theta}\right)(\sin^2 \theta + \cos^2 \theta) \ge \left(\frac{4}{\sin^2 \theta} + \frac{1}{\cos^2 \theta}\right)^2 \tag{1}
        \]
        再由柯西不等式,
        \[
        \left(\frac{4}{\sin^2 \theta} + \frac{1}{\cos^2 \theta}\right)(\sin^2 \theta + \cos^2 \theta) \ge (2+1)^2 \tag{2}
        \]
        由(1)与(2)得
        \[
        \frac{16}{\sin^6 \theta} + \frac{1}{\cos^6 \theta} \ge 9^2 = 81.
        \]
        等号成立当且仅当
        \[
        \frac{4}{\sin^2 \theta} = \frac{1}{\cos^2 \theta} \Rightarrow \tan \theta = \sqrt{2}>0 
        \]
    \end{solution}

    \question 设 $\theta,\phi$ 为锐角,且
    \[
    \frac{\sin^{2024}\theta}{\cos^{2022}\phi}+\frac{\cos^{2024}\theta}{\sin^{2022}\phi}=1,
    \]
    则 $\sin^{2023}\theta-\cos^{2023}\phi=?$

    \begin{solution}
        由柯西不等式,
\[
\left(\frac{\sin^{2024}\theta}{\cos^{2022}\phi}+\frac{\cos^{2024}\theta}{\sin^{2022}\phi}\right)\left(\frac{\cos^{2022} \phi}{\sin^{2020} \theta} + \frac{\sin^{2022} \phi}{\cos^{2020} \theta}\right) \ge (\sin^2 \theta + \cos^2 \theta)^2 =1
\]
\[
\Rightarrow 1\cdot \left( \left(\frac{\cos^{1011} \phi}{\sin^{1010} \theta} \right)^2+ \left(\frac{\sin^{1011} \phi}{\cos^{1010} \theta} \right)^2 \right) \ge 1
\]

\[
\Rightarrow \left(\frac{\cos^{1011} \phi}{\sin^{1010} \theta} \right)^2+ \left(\frac{\sin^{1011} \phi}{\cos^{1010} \theta} \right)^2 = 1
\]

此时
\[
\frac{\sin ^{2022}\theta}{\cos^{2022}\phi} = \frac{\cos^{2022} \theta}{\sin^{2022} \phi} \Rightarrow \sin\theta \sin \phi = \cos \theta \cos \phi \Rightarrow \cos(\theta+\phi)=0
\]

\[
\Rightarrow \theta+\phi =\frac{\pi}{2} \Rightarrow \sin \theta= \sin\left(\frac{\pi}{2}-\phi\right)= \cos \phi
\]

\[
\Rightarrow \sin^{2023} \theta-\cos^{2023} \phi = 0
\]
\textcolor{red}{为什么第三行=1?}
\end{solution}

    \question 证明对任意\(\triangle ABC\)均有  
        \[
        a\cos A + b\cos B + c\cos C = 2a\sin B\sin C\]
    \begin{solution}
        由正弦定理,设$a=k\sin A, b=k\sin B, c=k \sin C,$其中$k\in \mathbb{R}$,左式变为\begin{align*}
            a\cos A + b\cos B + c\cos C
            &=a\cos A+ k(\sin B \cos B+\sin C \cos C)\\
            &=a\cos A+ \frac{1}{2}k(\sin 2B+\sin 2C)\\
            &=a\cos A+ k\sin(B+C)\cos(B-C)\\
            &=a\cos A+ k\sin A\cos(B-C)\\
            &=a\cos A+ a\cos(B-C)\\
            &=a(\cos(B-C)-\cos (B+C))\\
            &=2a\sin B\sin C \text{(得证)}
        \end{align*}
    \end{solution}

    \question 已知一锐角三角形 $\triangle ABC$ 的边长分别为 $a,b,c$。设 $r$ 为 $\triangle ABC$ 内切圆的半径,$R$ 为 $\triangle ABC$ 外接圆的半径,试证
    \begin{parts}
    \part $$r = 4R \sin \frac{A}{2} \sin \frac{B}{2} \sin \frac{C}{2}$$
    \begin{solution}
        设 $S$ 为 $\triangle ABC$ 面积,则  
        \[
        S = \dfrac{1}{2}(a + b + c)r = \dfrac{abc}{4R}\Rightarrow r = \dfrac{abc}{2R(a + b + c)}
        \]
        又由于 $a = 2R\sin A,b = 2R\sin B,c = 2R\sin C$,代入得
        \[
        \begin{aligned}
        r &= \dfrac{2R\sin A \cdot 2R\sin B \cdot 2R\sin C}{2R(2R\sin A + 2R\sin B + 2R\sin C)} \\
        &= 2R \cdot \dfrac{\sin A \sin B \sin C}{\sin A + \sin B + \sin C} \\[1mm]
        &= 2R \cdot \dfrac{2\sin\frac{A}{2}\cos\frac{A}{2} \cdot 2\sin\frac{B}{2}\cos\frac{B}{2} \cdot 2\sin\frac{C}{2}\cos\frac{C}{2}}{2\sin\frac{A+B}{2}\cos\frac{A-B}{2} + 2\sin\frac{C}{2}\cos\frac{C}{2}} \\[1mm]
        &= 8R \cdot \dfrac{\sin\frac{A}{2}\sin\frac{B}{2}\sin\frac{C}{2}\cos\frac{A}{2}\cos\frac{B}{2}\cos\frac{C}{2}}{\cos\frac{A-B}{2}\cos\frac{C}{2} + \cos\frac{A+B}{2}\cos\frac{C}{2}} \\[1mm]
        &= 8R \cdot \dfrac{\sin\frac{A}{2}\sin\frac{B}{2}\sin\frac{C}{2}\cos\frac{A}{2}\cos\frac{B}{2}}{\cos\frac{A-B}{2} + \cos\frac{A+B}{2}} \\[1mm]
        &= 8R \cdot \sin\frac{A}{2}\sin\frac{B}{2}\sin\frac{C}{2} \cdot \dfrac{\frac{1}{2}(\cos\frac{A-B}{2} + \cos\frac{A+B}{2})}{\cos\frac{A-B}{2} + \cos\frac{A+B}{2}} \\[1mm]
        &= 4R \sin\frac{A}{2}\sin\frac{B}{2}\sin\frac{C}{2}
        \end{aligned}
        \]
    \end{solution}
    \part $$\frac{abc}{\sqrt{2(a^2 + b^2)(b^2 + c^2)(c^2 + a^2)}} \geq \frac{r}{2R}$$
    \begin{solution}
        由余弦定理,  
        \[
        \cos A = \dfrac{b^2 + c^2 - a^2}{2bc} \ge \dfrac{b^2 + c^2 - a^2}{b^2 + c^2} = 1 - \dfrac{a^2}{b^2 + c^2}
        \]
        而 $\cos A = 1 - 2\sin^2\frac{A}{2}$,所以
        \[
        1 - 2\sin^2\frac{A}{2} \ge 1 - \dfrac{a^2}{b^2 + c^2} \Rightarrow \dfrac{a^2}{b^2 + c^2} \ge 2\sin^2\frac{A}{2}
        \]
        同理可得
        \[
        \dfrac{b^2}{a^2 + c^2} \ge 2\sin^2\frac{B}{2},\quad
        \dfrac{c^2}{a^2 + b^2} \ge 2\sin^2\frac{C}{2}
        \]
        于是
        \[
        \dfrac{abc}{\sqrt{(a^2 + b^2)(b^2 + c^2)(c^2 + a^2)}} \ge \sqrt{2\sin^2\frac{A}{2} \cdot 2\sin^2\frac{B}{2}\cdot 2\sin^2\frac{C}{2}}= 2\sqrt{2} \sin\frac{A}{2} \sin\frac{B}{2} \sin\frac{C}{2}
        \]
        由 $(a)$ 得$$\sin\frac{A}{2} \sin\frac{B}{2} \sin\frac{C}{2} = \dfrac{r}{4R},$$
        故
        \[
        \dfrac{abc}{\sqrt{2(a^2 + b^2)(b^2 + c^2)(c^2 + a^2)}} \ge \frac{1}{\sqrt2}\cdot 2\sqrt{2}\cdot \dfrac{r}{4R}=\dfrac{r}{2R} 
        \]
    \end{solution}
    \end{parts}
    
    \question 已知 $A, B, C$ 是一锐角三角形的内角,证明
    \[
    (\tan A + \tan B + \tan C)^2 \ge (\sec A + 1)^2 + (\sec B + 1)^2 + (\sec C + 1)^2
    \]
    \begin{solution}
        由于 $A, B, C$ 都是锐角,故这些角的三角函数均为正。注意到
        \begin{align*}
        &(\cos A - \cos B)^2 \\
        &= \cos^2 A + \cos^2 B - 2 \cos A \cos B \\
        &= \cos(B+C)\cos A - \cos(A+C)\cos B - 2 \cos A \cos B \\
        &= \cos B \cos C \cos A + \sin B \sin C \cos A - \cos A \cos C \cos B + \sin A \sin C \cos B - 2 \cos A \cos B \\
        &= \cos A \cos B \cos C (\tan B \tan C + \tan A \tan C - 2 \sec C - 2)
        \end{align*}
        因为左边非负,有
        \[
        \tan B \tan C + \tan A \tan C \ge 2 \sec C + 2
        \]
        同理,
        \[
        \tan A \tan B + \tan A \tan C \ge 2 \sec A + 2, \quad
        \tan A \tan B + \tan B \tan C \ge 2 \sec B + 2
        \]
        将三式相加得
        \[
        \tan A \tan B + \tan A \tan C + \tan B \tan C \ge \sec A + \sec B + \sec C + 3
        \]
        因此,
        \begin{align*}
        &(\tan A + \tan B + \tan C)^2 \\
        &= \tan^2 A + \tan^2 B + \tan^2 C + 2(\tan A \tan B + \tan A \tan C + \tan B \tan C) \\
        &\ge \tan^2 A + \tan^2 B + \tan^2 C + 2(\sec A + \sec B + \sec C) + 6 \\
        &= \sec^2 A + \sec^2 B + \sec^2 C + 2(\sec A + \sec B + \sec C) + 3 \\
        &= (\sec A + 1)^2 + (\sec B + 1)^2 + (\sec C + 1)^2
        \end{align*}
    \end{solution}

    \question 已知 $x$ 为一实数, $0 < x < \pi$, 证明: 对于所有的自然数 $n$,
    \[
    \sin x + \frac{\sin 3x}{3} + \frac{\sin 5x}{5} + \dots + \frac{\sin(2n-1)x}{2n-1}
    \]
    的值为正数。
    \begin{solution}
        令 $$f(x) = \sin x + \frac{\sin 3x}{3} + \frac{\sin 5x}{5} + \dots + \frac{\sin(2n-1)x}{2n-1}$$
        由
        $$2\sin x \sin(2k-1)x = \cos(2k-2)x - \cos 2kx$$
        可得
        \begin{align*}
            &2f(x)\sin x\\
            &=1-\cos2x+\frac{\cos2x-\cos4x}{3}+\frac{\cos4x-\cos6x}{5}+\cdots+\frac{\cos(2n-2)x-\cos2nx}{2n-1} \\ 
            &=1-\left(1-\frac{1}{3}\right)\cos2x-\left(\frac{1}{3}-\frac{1}{5}\right)\cos4x-\cdots-\left(\frac{1}{2n-3}-\frac{1}{2n-1}\right)\cos(2n-2)x-\frac{\cos2nx}{2n-1} \\ 
            &\ge 1-\left[\left(1-\frac{1}{3}\right)+\left(\frac{1}{3}-\frac{1}{5}\right)+\cdots+\left(\frac{1}{2n-3}-\frac{1}{2n-1}\right)+\frac{1}{2n-1}\right]=0
        \end{align*}
        若等号成立,则有$\cos2kx=1(k=1,2,\dots,n)$,但因$0<x<\pi$,故$\cos2x\ne1$.于是得 $f(x)\sin x>0$.又因$\sin x>0$,所以$f(x)>0.$
    \end{solution}

    \question 证明  
    \[
    \cot\theta-\cot2\theta=\csc2\theta
    \]
    据此,若已知  
    \[
    \frac1{\sin8^\circ}+ \frac1{\sin16^\circ}+\cdots+\frac1{\sin4096^\circ}+ \frac1{\sin8192^\circ}
    =\frac1{\sin\alpha},
    \]
    其中$\alpha\in (0,90^\circ)$, 求$\alpha$。
    \begin{solution}
        有
        \[
        \cot \theta - \cot 2\theta 
        = \frac{\sin 2\theta \cos \theta - \sin \theta \cos 2\theta}{\sin \theta \sin 2\theta}
        = \frac{\sin(2\theta - \theta)}{\sin \theta \sin 2\theta}
        = \csc 2\theta \quad(\text{得证})
        \]
        故原式是一裂项和
        \[
        \sum_{k=0}^{10}\frac1{\sin\left(2^k8^\circ\right)}=\sum_{k=0}^{10}\left(\cot (2^{k-1}8^\circ) - \cot (2^k8^\circ)\right)=\cot 4^\circ - \cot 8192^\circ
        \]
        而 \(8192 \bmod 90 = 2\),故
        \begin{align*}
        \sum_{k=0}^{10}\frac1{\sin\left(2^k8^\circ\right)}
        &=\frac1{\tan4^\circ}-\frac1{\tan8192^\circ}\\
        &=\frac1{\tan4^\circ}+\tan2^\circ\\
        &=\frac{\cos4^\circ}{\sin4^\circ}+\frac{1-\cos4^\circ}{\sin4^\circ}\\
        &=\frac1{\sin4^\circ} \Rightarrow \alpha=4^\circ
        \end{align*}
        其中利用了恒等式$\tan \theta=\dfrac{1-\cos 2\theta}{\sin 2\theta}$。
    \end{solution}  

    \question 证明 
    \[
    \sum_{k=1}^{90} 2k\sin (2k)^\circ=90\cot 1^\circ
    \]
    \begin{solution}
        由 $k\sin (2k)^\circ+(90-k)\sin (180-2k)^\circ=90\sin (2k)^\circ,k=1,2,\cdots,90$,
        \begin{align*}
        \sum_{k=1}^{90} 2k\sin 2k^\circ
        &= 2(90\sin 2^\circ +90\sin 4^\circ +\cdots +90\sin 88^\circ)+ 2\cdot 45\sin 90^\circ \\
        &= \frac{90}{\sin 1^\circ}(2\sin 2^\circ \sin 1^\circ +2\sin 4^\circ \sin 1^\circ +\cdots +2\sin 88^\circ \sin 1^\circ)+90 \\
        &= \frac{90}{\sin 1^\circ}( \cos 1^\circ-\cos 3^\circ +\cos 3^\circ-\cos 5^\circ +\cdots +\cos 87^\circ-\cos 89^\circ)+90 \\
        &= \frac{90}{\sin 1^\circ} (\cos 1^\circ-\cos 89^\circ)+ 90 \\
        &= \frac{90}{\sin 1^\circ} (\cos 1^\circ-\sin 1^\circ)+ 90 \\
        &= 90\cot 1^\circ
        \end{align*}
    \end{solution}

    \question 计算  
    \[
    \sin^{6}1^\circ+\sin^{6}2^\circ+\cdots+\sin^{6}89^\circ
    \]
    \begin{solution}
        由$\sin x=\cos \left(\dfrac{\pi}{2}-x\right)$,设 
        \[
        S=\sin^{6}1^\circ+\sin^{6}2^\circ+\cdots+\sin^{6}89^\circ=\cos^{6}1^\circ+\cos^{6}2^\circ+\cdots+\cos^{6}89^\circ,
        \]
        则 
        \[
        2S= \sum_{k=1}^{44} (\sin^6 k^\circ+\cos^6 k^\circ) + 2\cdot\left(\frac{\sqrt2}{2}\right)^6
        \]
        其中
        \begin{align*}
        \sin^6 x + \cos^6 x &= (\sin^2 x + \cos^2 x)(\sin^4 x + \cos^4 x-\sin x \cos x) \\
        &=1 \cdot ((\sin^2 x + \cos^2 x)^2-3\sin^2 x \cos^2 x) \\
        &=1-\frac{3}{4}\sin^2 2x 
        \end{align*}且
        \[
        \sum_{k=1}^{44} \sin^2 (2k)^\circ=\sum_{k=1}^{22} (\sin^2 (2k)^\circ+\cos^2 (2k)^\circ)=22
        \]
        故
        \[
        2S=44-\frac{3}{4}\cdot 22+\frac{1}{4} \Rightarrow S=\frac{111}{8}
        \]
    \end{solution}
        
    \question 求 
    \[
    \tan^2 \frac {\pi}{16} \cdot \tan^2 \frac {3\pi}{16} 
    +\tan^2 \frac {\pi}{16} \cdot \tan^2 \frac {5\pi}{16}
    +\tan^2 \frac {3\pi}{16} \cdot \tan^2 \frac {7\pi}{16}
    +\tan^2 \frac {5\pi}{16} \cdot \tan^2 \frac {7\pi}{16}\] 
    的值。
    \begin{solution}
        将原式写成
        \[
        \left(\tan^2 \frac{\pi}{16}+\tan^2\frac{7\pi}{16}\right)\left(\tan^2 \frac{3\pi}{16}+\tan^2 \frac{5\pi}{16}\right) = \left(\tan^2 \frac{\pi}{16}+\cot^2\frac{\pi}{16}\right)\left(\tan^2 \frac{3\pi}{16}+\cot^2 \frac{3\pi}{16}\right).
        \]
        注意到        
        \[
        \tan^2\theta+\cot^2\theta=\frac{\sin^4\theta+\cos^4\theta}{\sin^2\theta \cos^2\theta}
        =\frac{1}{\sin^2\theta \cos^2 \theta}-2
        = \frac{4}{\sin^2 2\theta} -2
        = \frac{8}{1-\cos 4\theta}-2
        \]
        故 
        \[
        \tan^2\dfrac{\pi}{16}+\cot^2\dfrac{\pi}{16} = -2+\dfrac{8}{1-\frac{\sqrt 2}{2}} = 14+8\sqrt 2
        \]
        同理得 
        \[
        \tan^2\dfrac{3\pi}{16}+\cot^2\dfrac{3\pi}{16}=14-8\sqrt 2
        \]
        故原式 $=196-2\cdot 64=68$。
    \end{solution}    

    \question 求表达式 
    \[
    \log_{8}[(\tan 1^\circ+\sqrt{3})(\tan 2^\circ+\sqrt{3})(\tan 3^\circ+\sqrt{3})\cdots(\tan 29^\circ+\sqrt{3})]
    \]
    的值。
    \begin{solution}
        若 $A + B = 30^\circ$,则
        \[
        \tan 30^\circ = \frac{1}{\sqrt{3}} = \frac{\tan A + \tan B}{1 - \tan A \tan B} \Rightarrow \tan A + \tan B + \sqrt{3}(\tan A + \tan B) = 1
        \]
        于是
        \[
        (\tan A + \sqrt{3})(\tan B + \sqrt{3}) = \tan A \tan B + \sqrt{3}(\tan A + \tan B) + 3 = 4
        \]
        因此
        \[
        P=(\tan 1^\circ+\sqrt{3})(\tan 2^\circ+\sqrt{3})(\tan 3^\circ+\sqrt{3})\cdots(\tan 29^\circ+\sqrt{3})=4^{14} \cdot 2 =2^{29}
        \]
        故
        \[
        \log_{8}P= \frac{29}{3}
        \]
    \end{solution}

    \question 已知 $\sin\theta=\dfrac{3}{5}$,求
    \[
    \left( \frac{\sqrt{2}\cos\theta}{\sqrt{3}}+ \frac{\cos 2\theta}{\sqrt{6}}+ \frac{\cos 3\theta}{\sqrt{24}}+\cdots \right)
    \left( \frac{\sqrt{2}\sin\theta}{\sqrt{3}}+ \frac{\sin 2\theta}{\sqrt{6}}+ \frac{\sin 3\theta}{\sqrt{24}}+\cdots \right)
    \]
    \begin{solution}
        发现
        \[
        \frac{\sqrt{2}}{\sqrt{3}} \left( \cos\theta + \frac{1}{2}\cos 2\theta + \frac{1}{4}\cos 3\theta + \cdots \right)
        = \frac{1}{3} \sqrt{6} \cdot \Re \left( \frac{e^{i\theta}}{1 - \frac{1}{2} e^{i\theta}} \right)
        = \frac{2}{3} \sqrt{6} \cdot \Re \left( \frac{e^{i\theta}}{2 - e^{i\theta}} \right)        
        \]
        已知 \(\sin\theta = \dfrac{3}{5}\),则 \(\cos\theta = \dfrac{4}{5},\)代入 \(e^{i\theta} = \cos\theta + i\sin\theta = \dfrac{4}{5} + \dfrac{3}{5}i\) 得
        \[
        \frac{2\sqrt{6}}{3} \cdot \Re \left( \frac{1}{3} + \frac{2}{3}i \right)
        = \frac{2\sqrt{6}}{3} \cdot \frac{1}{3}
        = \frac{2\sqrt{6}}{9} 
        \]
        同理
        \[
        \frac{\sqrt{2}}{\sqrt{3}} \left( \sin\theta + \frac{1}{2}\sin 2\theta + \frac{1}{4}\sin 3\theta + \cdots \right) 
        = \frac{\sqrt{6}}{3}  \cdot \Im \left( \frac{e^{i\theta}}{1 - \frac{1}{2} e^{i\theta}} \right)
        = \frac{2\sqrt{6}}{3}  \cdot \Im \left( \frac{1}{3} + \frac{2}{3}i \right)
        = \frac{4\sqrt{6}}{9} 
        \]
        两式相乘得
        \[
        \left( \frac{\sqrt{2}\cos\theta}{\sqrt{3}}+ \frac{\cos 2\theta}{\sqrt{6}}+ \frac{\cos 3\theta}{\sqrt{24}}+\cdots \right)
        \left( \frac{\sqrt{2}\sin\theta}{\sqrt{3}}+ \frac{\sin 2\theta}{\sqrt{6}}+ \frac{\sin 3\theta}{\sqrt{24}}+\cdots \right)=\frac{2\sqrt{6}}{9}  \cdot \frac{4\sqrt{6}}{9} = \frac{16}{27}
        \]  
    \end{solution}

    \question 证明
    \[
    \cos \frac{\pi}{11} + \cos \frac{3\pi}{11} + \cos \frac{5\pi}{11} + \cos \frac{7\pi}{11} + \cos \frac{9\pi}{11} = \frac{1}{2}.
    \]
    \begin{solution}
        设 $z = \cos \theta + i \sin \theta$,其中 $\theta = \dfrac{\pi}{11}$,则由棣莫弗定理,
        \[
        z^k = \cos k\theta + i \sin k\theta.
        \]
        考虑和
        \[
        S = z + z^3 + z^5 + z^7 + z^9=z \frac{(z^2)^5 - 1}{z^2 - 1}=\frac{z^{11} - z}{z^2 - 1}
        \]
        由于 $z^{11} = \cos 11\theta + i \sin 11\theta = \cos \pi + i \sin \pi = -1$,所以
        \[
        S = \frac{-1 - z}{z^2 - 1} = \frac{1}{1-z}
        \]
        代入 $z = \cos \theta + i \sin \theta$:
        \begin{align*}
        S &= \frac{1}{1-\cos \theta - i \sin \theta} \\
        &= \frac{(1-\cos \theta) + i \sin \theta}{(1-\cos \theta)^2 + \sin^2 \theta} \\
        &= \frac{1}{2} + \frac{i \sin \theta}{2(1-\cos \theta)} \\
        &= \frac{1}{2} + \frac{i \cdot 2 \sin \frac{\theta}{2} \cos \frac{\theta}{2}}{2 \cdot 2 \sin^2 \frac{\theta}{2}} \\
        &= \frac{1}{2} + \frac{i}{2} \cot \frac{\theta}{2}
        \end{align*}
        故所求即
        \[
        \cos \frac{\pi}{11} + \cos \frac{3\pi}{11} + \cos \frac{5\pi}{11} + \cos \frac{7\pi}{11} + \cos \frac{9\pi}{11} = \Re(S)= \frac{1}{2}
        \]
    \end{solution}

    \question 求值 
    \[
    \frac{1}{\tan^2{\dfrac{\pi}{18}}}+\frac{1}{\tan^2{\dfrac{5\pi}{18}}}+\frac{1}{\tan^2{\dfrac{7\pi}{18}}}
    \]
    \begin{solution}
        设 \(\theta = \dfrac{\pi}{18} \Rightarrow 3\theta = \dfrac{\pi}{6}\),两边取$\tan,$
        \[
        \tan(3\theta) = \tan\frac{\pi}{6} \Rightarrow \frac{3\tan\theta - \tan^3\theta}{1 - 3\tan^2\theta} = \frac{1}{\sqrt{3}}
        \]
        整理得
        \[
        \sqrt{3}\tan^3\theta - 3\tan^2\theta - 3\sqrt{3}\tan\theta + 1 = 0
        \]
        发现该三次方程的三个根分别是
        \[
        x_1 = \tan\frac{\pi}{18},\; x_2 = -\tan\frac{5\pi}{18},\; x_3 = \tan\frac{7\pi}{18}
        \]
        由韦达定理
        \[
        x_1 + x_2 + x_3 = \sqrt 3,\;
        x_1x_2 + x_2x_3 + x_3x_1 = -3,\;
        x_1x_2x_3 = -\frac{1}{\sqrt3}
        \]
        则所求为
        \begin{align*}
        \frac{1}{x_1^2} + \frac{1}{x_2^2} + \frac{1}{x_3^2} &= \left( \frac{1}{x_1} + \frac{1}{x_2} + \frac{1}{x_3} \right)^2 - 2\left( \frac{1}{x_1x_2} + \frac{1}{x_2x_3} + \frac{1}{x_3x_1} \right)  \\
        &=\left(\frac{x_1x_2 + x_2x_3 + x_3x_1}{x_1x_2x_3}\right)^2-2\left(\frac{x_1 + x_2 + x_3}{x_1x_2x_3}\right)\\
        &= \left( \frac{-3}{-\frac{1}{\sqrt{3}}} \right)^2 - 2 \cdot \frac{\sqrt{3}}{-\frac{1}{\sqrt{3}}} \\
        &= 33
        \end{align*}
    \end{solution}

    \question 计算
    \[
    \cos^7 \frac{\pi}{9} + \cos^7 \frac{5\pi}{9} + \cos^7 \frac{7\pi}{9}
    \]
    \begin{solution}
        因 \( \theta = \dfrac{\pi}{9},\dfrac{5\pi}{9},\dfrac{7\pi}{9} \)皆满足
        \[
        \cos 3\theta = 4\cos^3 \theta - 3\cos \theta
        \]
        故 \(  \cos \dfrac{\pi}{9},\cos \dfrac{5\pi}{9} , \cos \dfrac{7\pi}{9} \) 是方程 \( 8x^3 - 6x - 1 = 0 \) 的三根。
        
        现设$x=\cos \dfrac{\pi}{9},y=\cos \dfrac{5\pi}{9} , z=\cos \dfrac{7\pi}{9}, S_n = x^n + y^n + z^n$,
        我们将用牛顿恒等式递推计算 \( S_7 \)。根据韦达定理可得
        \[
        \begin{aligned}
        S_1 &= x + y + z = 0 \\
        S_2 &= (x + y + z)^2 - 2(xy + yz + zx) = 0^2 - 2\cdot\left(-\frac{3}{4}\right) = \frac{3}{2} \\
        S_3 &= x^3 + y^3 + z^3 = 3xyz = \frac{3}{8}
        \end{aligned}
        \]
        递推公式为
        \begin{align*}
        S_n &= (x+y+z)S_{n-1} - (xy+yz+zx)S_{n-2} + xyz \cdot S_{n-3} \\
        &=0 \cdot S_{n-1} - (-\dfrac{3}{4}) S_{n-2} + \dfrac{1}{8} \cdot S_{n-3}
        \end{align*}
        于是
        \[
        \begin{aligned}
        S_4 &= 0 \cdot S_3 - (-\dfrac{3}{4}) S_2 + \dfrac{1}{8} \cdot S_1 = \dfrac{3}{4} \cdot \dfrac{3}{2} = \dfrac{9}{8} \\
        S_5 &= 0 \cdot S_4 - (-\dfrac{3}{4}) S_3 + \dfrac{1}{8} \cdot S_2 = \dfrac{3}{4} \cdot \dfrac{3}{8} + \dfrac{1}{8} \cdot \dfrac{3}{2} = \dfrac{9}{32} + \dfrac{3}{16} = \dfrac{15}{32} \\
        S_6 &= 0 \cdot S_5 - (-\dfrac{3}{4}) S_4 + \dfrac{1}{8} \cdot S_3 = \dfrac{3}{4} \cdot \dfrac{9}{8} + \dfrac{1}{8} \cdot \dfrac{3}{8} = \dfrac{27}{32} + \dfrac{3}{64} = \dfrac{57}{64} \\
        S_7 &= 0 \cdot S_6 - (-\dfrac{3}{4}) S_5 + \dfrac{1}{8} \cdot S_4 = \dfrac{3}{4} \cdot \dfrac{15}{32} + \dfrac{1}{8} \cdot \dfrac{9}{8} = \dfrac{45}{128} + \dfrac{9}{64} = \dfrac{63}{128}
        \end{aligned}
        \]
    \end{solution}

    \question 证明
    \[
    \tan \dfrac{2\pi}{15} \tan \dfrac{4\pi}{15} \tan \dfrac{8\pi}{15} \tan \dfrac{16\pi}{15} = -1
    \]
    \begin{solution}
        原式写成
        \[
        V=\frac{\sin 24^\circ \sin 48^\circ \sin 96^\circ \sin 192^\circ}{\cos 24^\circ \cos 48^\circ \cos 96^\circ \cos 192^\circ}
        \]
        记 \(T = \sin 24^\circ \sin 48^\circ \sin 96^\circ \sin 192^\circ,N = \cos 24^\circ \cos 48^\circ \cos 96^\circ \cos 192^\circ\)
        
        利用三角恒等式 \(\sin(60^\circ - \alpha)\sin \alpha \sin(60^\circ + \alpha) = \dfrac{1}{4}\sin 3\alpha\),整理 \(T\) 为
        \begin{align*}
        T &= -\sin 24^\circ \sin 48^\circ \sin 96^\circ \sin 12^\circ\\
        &= -\frac{\sin 12^{\circ}\sin 48^{\circ}\sin 108^{\circ}}{\sin 108^{\circ}} \cdot \sin 24^{\circ}\sin 96^{\circ}\\
        &=-\frac{1}{4}\frac{\sin 144^{\circ}}{\sin 108^{\circ}} \cdot \sin 24^{\circ}\sin 96^{\circ} \\
        &=-\frac{1}{4} \cdot \frac{\sin 24^\circ \sin 36^\circ \sin 96^\circ}{\sin 108^\circ} \\
        &= -\frac{1}{4} \cdot \frac{\frac{1}{4} \sin 108^\circ}{\sin 108^\circ} \\
        &= -\frac{1}{16}
        \end{align*}
        连续使用恒等式 $\sin2\alpha=2\sin \alpha\cos \alpha$有
        \[
        16\sin 24^\circ \cdot N = \sin 384^\circ = \sin 24^\circ \Rightarrow N = \frac{1}{16}
        \]
        故
        \[
        V = \frac{T}{N} = \frac{-\frac{1}{16}}{\frac{1}{16}} = \boxed{-1}
        \]
    \end{solution}
    \begin{solution}
        设 \(a = \dfrac{\pi}{15}\),有
        \begin{align*}
        \tan 2a \tan 4a \tan 8a \tan 16a &= \frac{2\sin a \cos a}{\cos 2a} \cdot \frac{2\sin 2a \cos 2a}{\cos 4a} \cdot \frac{2\sin 4a \cos 4a}{\cos 8a} \cdot \frac{2\sin 8a \cos 8a}{\cos 16a} \\&=
        16 \cdot \frac{\cos a \sin a \sin 2a \sin 4a \sin 8a}{\cos 16a}
        \end{align*}
        因为$\cos 16a=-\cos a,$于是
        \[
        \tan 2a \tan 4a \tan 8a \tan 16a = -16 \sin a \sin 2a \sin 4a \sin 8a
        \]
        我们来计算这个乘积。设 \(x = \cos 3a\),因为
        \[
        4\sin a \sin 4a = 2\cos3a - 2\cos5a=2\cos 3a-1
        \]
        \[
        4\sin 2a \sin 8a = 2\cos6a - 2\cos10a=2\cos 6a+1
        \]
        它们的乘积
        \[
        (2\cos 3a-1)(2\cos 6a+1)=(2x-1)(2(2x^2-1)+1)=2x(4x^2-2x-1)+1
        \]
        其中$x = \cos 3a$,又因为 \(\cos 9a = -\cos 6a\),可以推出
        \[
        4x^3-3x=1-2x^2 \Rightarrow (x + 1)(4x^2 - 2x - 1) = 0 \Rightarrow 2x\cdot 0+1=1
        \]
        故
        \[
        16 \sin a \sin 2a \sin 4a \sin 8a = -1
        \Rightarrow \tan \dfrac{2\pi}{15} \tan \dfrac{4\pi}{15} \tan \dfrac{8\pi}{15} \tan \dfrac{16\pi}{15} = -1
        \]
    \end{solution}

    \question 求  
    \[
    S = \sum_{n=1}^7 \tan^2\left(\frac{n\pi}{16}\right)
    \]
    \begin{solution}
        注意到 $$\tan\left(-\frac{7\pi}{16}\right), \tan\left(-\frac{6\pi}{16}\right), \dots, \tan\left(\frac{7\pi}{16}\right)$$ 是方程 $\tan16x = 0$ 的解,由恒等式
        \[
        \tan(n\theta) = \frac{\comb{n}{1} \tan\theta - \comb{n}{3} \tan^3\theta + \cdots}{1 - \comb{n}{2} \tan^2\theta + \cdots}
        \]也即是
        \[
        \comb{16}{1}x - \comb{16}{3}x^3 + \cdots + \comb{16}{13}x^{13} - \comb{16}{15}x^{15} = 0
        \]
        的根,去掉零根 $\tan0=0$ 后,
        $$\tan\left(\pm \frac{7\pi}{16}\right), \tan\left(\pm \frac{6\pi}{16}\right), \dots, \tan\left(\pm \frac{\pi}{16}\right)$$ 是
        \[
        \comb{16}{1} - \comb{16}{3}x^2 + \comb{16}{5}x^4 - \cdots + \comb{16}{13}x^{12} - \comb{16}{15}x^{14} = 0
        \]
        的根,因此$$\tan^2\left(\frac{7\pi}{16}\right), \tan^2\left(\frac{6\pi}{16}\right), \dots, \tan^2\left(\frac{\pi}{16}\right)$$ 
        是
        \[
        \comb{16}{1} - \comb{16}{3}y + \comb{16}{5}y^2 - \cdots + \comb{16}{13}y^6 - \comb{16}{15}y^7 = 0
        \]
        的根;由韦达定理,根之和为
        \[
        \frac{\comb{16}{13}}{\comb{16}{15}} = 35
        \]
    \end{solution}

    \question 若 $\alpha = \dfrac{2\pi}{1999}$, 试求 
    \[
    \cos\alpha\cos 2\alpha\cos 3\alpha\ldots\cos 998\alpha\cos 999\alpha
    \]
    的值。
    \begin{solution}
        记$S = \cos \alpha \cos 2\alpha \cdots \cos 999\alpha, T = \sin \alpha \sin 2\alpha \cdots \sin 999\alpha$,则
        \[
        ST = \sin \alpha \cos \alpha \sin 2\alpha \cos 2\alpha \cdots \sin 999\alpha \cos 999\alpha 
        \]
        即
        \[
        2^{999} ST = \sin 2\alpha \sin 4\alpha \cdots \sin 1998\alpha
        \]
        且注意到
        \[
        \sin 1998\alpha = - \sin(2\pi - 1998\alpha) = - \sin\frac{2\pi}{1999} = - \sin \alpha
        \]
        同理,
        \[
        \sin 1996\alpha = - \sin 3\alpha, \quad \sin 1994\alpha = - \sin 5\alpha, \cdots
        \]
        因此
        \[
        2^{999} ST = \sin \alpha \sin 4\alpha \cdots (-\sin 3\alpha)(-\sin \alpha) = \sin \alpha \sin 4\alpha \cdots \sin 3\alpha \sin \alpha = T 
        \]
        由于 $T \neq 0$,故
        \[
        S = \cos \alpha \cos 2\alpha \cos 3\alpha \cdots \cos 998\alpha \cos 999\alpha = \frac{1}{2^{999}}
        \]
    \end{solution}

    \question 证明:对于所有大于1的自然数 $n$,有
    \[
    \sin \frac{\pi}{n} \cdot \sin \frac{2\pi}{n} \cdots \sin \frac{(n-1)\pi}{n} = \frac{n}{2^{n-1}}.
    \]
    \begin{solution}
        设
        \[
        \omega = \cos \frac{2\pi}{n} + i \sin \frac{2\pi}{n},
        \]
        则多项式
        \[
        x^{n-1} + x^{n-2} + \cdots + 1 = 0
        \]的 $n-1$ 个根为$\omega^k, k=1,2,\dots,n-1,$且$\omega^n = 1,$令
        \[
        f(x) = x^{n-1} + x^{n-2} + \cdots + 1 = (x-\omega^1)(x-\omega^2)\cdots (x-\omega^{n-1}),
        \]
        则
        \[
        f(1) = n = (1-\omega^1)(1-\omega^2)\cdots (1-\omega^{n-1}).
        \]
        且发现
        \[
        1 - \omega^k = 2 \sin^2 \frac{k\pi}{n} - 2 i \sin \frac{k\pi}{n} \cos \frac{k\pi}{n} = 2 \sin \frac{k\pi}{n} \left(\sin \frac{k\pi}{n} - i \cos \frac{k\pi}{n}\right).
        \]
        故
        \[
        |1 - \omega^k| = 2 \sin \frac{k\pi}{n}.
        \]
        取模可得
        \[
        n = |1 - \omega^1||1 - \omega^2| \cdots |1 - \omega^{n-1}| = 2^{n-1} \sin \frac{\pi}{n} \cdot \sin \frac{2\pi}{n} \cdots \sin \frac{(n-1)\pi}{n}.
        \]
        即得证
        \[
        \sin \frac{\pi}{n} \cdot \sin \frac{2\pi}{n} \cdots \sin \frac{(n-1)\pi}{n} = \frac{n}{2^{n-1}}.
        \]
    \end{solution}
\question
计算下列乘积:

(a) $\sin \frac{\pi}{2n + 1} \sin \frac{2\pi}{2n + 1} \sin \frac{3\pi}{2n + 1} \dots \sin \frac{n\pi}{2n + 1}$
以及
$\sin \frac{\pi}{n} \sin \frac{2\pi}{n} \sin \frac{3\pi}{n} \dots \sin \frac{(n - 1)\pi}{n}$

(b) $\cos \frac{\pi}{2n + 1} \cos \frac{2\pi}{2n + 1} \cos \frac{3\pi}{2n + 1} \dots \cos \frac{n\pi}{2n + 1}$
以及
$\cos \frac{\pi}{2n} \cos \frac{2\pi}{2n} \cos \frac{3\pi}{2n} \dots \cos \frac{(n - 1)\pi}{2n}$

\begin{solution}
    (a) 设 $\omega = \cos \frac{2\pi}{n} + i \sin \frac{2\pi}{n}$ 为单位的 $n$ 次复根。
    考虑方程 $x^{n-1} + x^{n-2} + \cdots + 1 = 0$,其根为 $\omega^k$ ($k=1,2,\dots,n-1$)。
    因此有恒等式:
    \[ f(x) = x^{n-1} + x^{n-2} + \cdots + 1 = (x-\omega^1)(x-\omega^2)\cdots (x-\omega^{n-1}) \]
    令 $x=1$,得:
    \[ f(1) = n = (1-\omega^1)(1-\omega^2)\cdots (1-\omega^{n-1}) \]
    由于 $|1 - \omega^k| = |1 - (\cos \frac{2k\pi}{n} + i \sin \frac{2k\pi}{n})| = |2 \sin^2 \frac{k\pi}{n} - 2i \sin \frac{k\pi}{n} \cos \frac{k\pi}{n}| = 2 \sin \frac{k\pi}{n}$。
    取模可得:
    \[ n = 2^{n-1} \sin \frac{\pi}{n} \sin \frac{2\pi}{n} \dots \sin \frac{(n-1)\pi}{n} \]
    由此得出第一个重要结果:
    \[ \sin \frac{\pi}{n} \sin \frac{2\pi}{n} \dots \sin \frac{(n-1)\pi}{n} = \frac{n}{2^{n-1}} \]
    同理可证
    （通过替换 $n$ 为 $2n+1$ 并利用对称性）:
    \[ \sin \frac{\pi}{2n + 1} \sin \frac{2\pi}{2n + 1} \dots \sin \frac{n\pi}{2n + 1} = \frac{\sqrt{2n + 1}}{2^n} \]

    (b) 对于余弦乘积,我们利用倍角公式 $\sin 2\theta = 2 \sin \theta \cos \theta$。
    对于 $2n+1$ 的情形:
    \[ \sin \frac{2k\pi}{2n + 1} = 2 \sin \frac{k\pi}{2n + 1} \cos \frac{k\pi}{2n + 1} \]
    将 $k=1, 2, \dots, n$ 的等式全部相乘:
    \[ \prod_{k=1}^{n} \sin \frac{2k\pi}{2n+1} = 2^n \left( \prod_{k=1}^{n} \sin \frac{k\pi}{2n+1} \right) \left( \prod_{k=1}^{n} \cos \frac{k\pi}{2n+1} \right) \]
    由于 $\sin \frac{2k\pi}{2n+1}$ 的集合（当 $k$ 超过 $n/2$ 时利用 $\sin(\pi-\theta)=\sin\theta$）与 $\sin \frac{k\pi}{2n+1}$ 的集合完全相同,两边的正弦乘积项约去,得:
    \[ \cos \frac{\pi}{2n + 1} \cos \frac{2\pi}{2n + 1} \dots \cos \frac{n\pi}{2n + 1} = \frac{1}{2^n} \]

    对于 $\cos \frac{k\pi}{2n}$ 的乘积:
    利用 $\sin \theta \cos \theta = \frac{1}{2} \sin 2\theta$:
    \[ \left[ \prod_{k=1}^{n-1} \cos \frac{k\pi}{2n} \right] \left[ \prod_{k=1}^{n-1} \sin \frac{k\pi}{2n} \right] = \frac{1}{2^{n-1}} \prod_{k=1}^{n-1} \sin \frac{k\pi}{n} \]
    已知右侧乘积为 $\frac{1}{2^{n-1}} \cdot \frac{n}{2^{n-1}} = \frac{n}{2^{2n-2}}$。
    又因为 $\sin \frac{k\pi}{2n} = \cos \frac{(n-k)\pi}{2n}$,所以左侧两个括号内的数值相等。
    取平方根（项均为正）得:
    \[ \cos \frac{\pi}{2n} \cos \frac{2\pi}{2n} \dots \cos \frac{(n - 1)\pi}{2n} = \frac{\sqrt{n}}{2^{n-1}} \]

    注:若将 (a) 的结果除以 (b) 的结果,可得正切乘积:
    \[ \tan \frac{\pi}{2n + 1} \tan \frac{2\pi}{2n + 1} \dots \tan \frac{n\pi}{2n + 1} = \sqrt{2n + 1} \]
    \[ \tan \frac{\pi}{2n} \tan \frac{2\pi}{2n} \dots \tan \frac{(n - 1)\pi}{2n} = 1 \]
\end{solution}
    \question 证明:对于任意正整数 $n$,有
        \[
        \tan \frac{\pi}{2n+1} \tan \frac{2\pi}{2n+1} \cdots \tan \frac{n\pi}{2n+1} = \sqrt{2n+1}.
        \]
    \begin{solution}
        注意到
        \[
        \left( \cos \frac{k\pi}{2n+1} + i \sin \frac{k\pi}{2n+1} \right)^{2n+1} = \left( e^{i \frac{k\pi}{2n+1}} \right)^{2n+1} = e^{ik\pi} = (-1)^k.
        \]
        由二项式定理,比较虚部可得
        \[
        \sum_{j=0}^n (-1)^j \binom{2n+1}{2j+1} \left( \cos \frac{k\pi}{2n+1} \right)^{2n-2j} \left( \sin \frac{k\pi}{2n+1} \right)^{2j+1} = 0.
        \]
        将两边除以$\left( \cos \dfrac{k\pi}{2n+1} \right)^{2n+1}$得到
        \[
        \sum_{j=0}^n (-1)^j \binom{2n+1}{2j+1} \left( \tan \frac{k\pi}{2n+1} \right)^{2j+1} = 0.
        \]
        即\(\tan \dfrac{k\pi}{2n+1},1 \leq k \leq 2n+1\) 是次数为 \(2n+1\)的多项式
        \[
        \sum_{j=0}^n (-1)^j \binom{2n+1}{2j+1} z^{2j+1} = 0,
        \]
        的根;不考虑根 \(z = 0\) 可知, $\tan\left(\dfrac{k\pi}{2n+1}\right),\; 1 \leq k \leq 2n$是多项式
        \[
        \sum_{j=0}^n (-1)^j \binom{2n+1}{2j+1} z^{2j} = (-1)^n z^{2n} + \cdots + (2n+1) = 0.
        \]
        的根,因此由韦达定理,
        \[
        \prod_{k=1}^{2n} \tan\left(\frac{k\pi}{2n+1}\right) = (-1)^n \cdot (2n+1).
        \]
        且观察到
        \[
        \tan\left(\frac{(2n+1-k)\pi}{2n+1}\right) = -\tan\left(\frac{k\pi}{2n+1}\right),\;1 \leq k \leq n
        \]
        所以
        \[
        \left( \prod_{k=1}^n \tan\left(\frac{k\pi}{2n+1}\right) \right)^2 = 2n+1.
        \]
        又因\(\tan\left(\dfrac{k\pi}{2n+1}\right) > 0,\;1 \leq k \leq n\),故
        \[
        \tan\left(\frac{\pi}{2n+1}\right) \tan\left(\frac{2\pi}{2n+1}\right) \cdots \tan\left(\frac{n\pi}{2n+1}\right) = \sqrt{2n+1}.
        \]
    \end{solution}

    \question 设
    \[
    f_n(\theta)=\sin\theta\sin2\theta\sin4\theta\cdots\sin(2^{\,n-1}\theta), \quad n \in \mathbb{N}
    \]
    证明对任意实数 $\theta$ 及正整数 $n$皆有
    \[
    |f_n(\theta)|\le \frac{2}{\sqrt{3}}\,\left|f_n\left(\frac{\pi}{3}\right)\right|
    \]
    \begin{solution}
        由AM-GM不等式,
        \begin{align*}
        |g(\theta)| &= |\sin \theta| |\sin(2\theta)|^{1/2} = \frac{\sqrt{2}}{\sqrt[4]{3}} \left( \sqrt[4]{|\sin \theta| \cdot |\sin \theta| \cdot |\sin \theta| \cdot |\sqrt{3}\cos \theta|} \right)^2 \\
        &\le \frac{\sqrt{2}}{\sqrt[4]{3}} \cdot \frac{3 \sin^2 \theta + 3 \cos^2 \theta}{4} = \left(\frac{\sqrt{3}}{2}\right)^{3/2}
        \end{align*}
        故函数$g(\theta)$的最大值为 $\left(\dfrac{\sqrt{3}}{2}\right)^{3/2}$,并且在 $\theta=\dfrac{2^k\pi}{3},k \in \mathbb{N}$处取得。注意到
        \[
        f_n(\theta)
        = g(\theta)\cdot g(2\theta)^{\frac{1}{2}}\cdot g(4\theta)^{\frac{3}{4}}\cdots
        g(2^{n-1}\theta)^E,
        \]
        其中
        \[
        E=\frac{2}{3}\left(1-\left(-\frac{1}{2}\right)^n\right)
        \]
        于是
        \begin{align*}
        \left|\frac{f_n(\theta)}{f_n(\frac{\pi}{3})}\right|
        &=\left|\frac{g(\theta)\cdot g(2\theta)^{\frac{1}{2}}\cdots g(2^{n-1}\theta)^E}{g(\frac{\pi}{3})\cdot g(\frac{2\pi}{3})^{\frac{1}{2}}\cdots g(2^{n-1}\frac{\pi}{3})^E}\right|
        \cdot \left|\frac{\sin(2^n\theta)}{\sin(2^n\frac{\pi}{3})}\right|^{1-\frac{E}{2}}\\
        &\le \left|\frac{\sin(2^n\theta)}{\sin(2^n\frac{\pi}{3})}\right|^{1-\frac{E}{2}}
        \le \left(\frac{1}{\sqrt{3}/2}\right)^{1-\frac{E}{2}}
        \le \frac{2}{\sqrt{3}}
        \end{align*}
        因此得证对所有实数 $\theta$ 及任意正整数$n$ 皆有
        \[
        |f_n(\theta)|\le \frac{2}{\sqrt{3}}\,\left|f_n\left(\frac{\pi}{3}\right)\right|
        \]
    \end{solution}
\end{questions}
\pagebreak

\begin{center}
  {\fontsize{30pt}{26pt}\selectfont
    \hypertarget{反三角函数}{反三角函数} \label{反三角函数}
  }
\end{center}
\separator
\vspace{1pt}

\begin{questions}
    \question 若 $|x|\le 1$,证明
    \[
    \tan^{-1}\sqrt{\frac{1-x}{1+x}}=\frac{1}{2}\cos^{-1}x
    \]
    \begin{solution}
        设$\theta=\tan^{-1}\sqrt{\dfrac{1-x}{1+x}}$,于是
        \[
        1+\tan^2\theta=1+\frac{1-x}{1+x}=\frac{2}{1+x}
        \]
        从而
        \[
        \cos^2\theta=\frac{1}{\sec^2\theta}=\frac{1+x}{2} \Rightarrow \cos2\theta=2\cdot\frac{1+x}{2}-1=x.
        \]
        由于 $|x|\le 1$,且 $\theta \in \left[0,\dfrac{\pi}{2}\right]$,可知 $2\theta\in[0,\pi]$,因此
        \[
        2\theta=\cos^{-1}x \Rightarrow \theta=\frac{1}{2}\cos^{-1}x
        \]
        于是得证
        \[
        \tan^{-1}\sqrt{\frac{1-x}{1+x}}=\frac{1}{2}\cos^{-1}x
        \]
    \end{solution}

    \question 已知锐角 $\theta,\phi$ 满足
    \[
    2\cos\theta = \cos\phi, \quad 2\sin\theta = 3\sin\phi,
    \]
    证明
    \[
    \theta + \phi = \pi - \arctan \sqrt{15}.
    \]
    \begin{solution}
        将已知两式平方并相加,可得
        \[
        4(\cos^2\theta + \sin^2\theta) = \cos^2\phi + 9\sin^2\phi \Rightarrow \sin^2\phi = \frac{3}{8} \Rightarrow \sin\phi = \sqrt{\frac{3}{8}}>0
        \]
        同时可得
        \[
        4\sin^2\theta = 9\sin^2\phi = \frac{27}{8} \Rightarrow \sin\theta = \sqrt{\frac{27}{32}}
        \]
        于是$\tan\theta =\sqrt{\dfrac{27}{5}},\tan\phi=\sqrt{\dfrac{3}{5}}$,且
        \[
        \tan(\theta+\phi) = \frac{\tan\theta + \tan\phi}{1 - \tan\theta\tan\phi} = -\sqrt{15}
        \]
        由于 $0<\theta+\phi<\pi$,取
        \[
        \theta+\phi = \pi - \arctan\sqrt{15}
        \]
    \end{solution}

    \question 已知
    \[
    \sin^{-1}(\sin\alpha+\sin\beta)+\sin^{-1}(\sin\alpha-\sin\beta)=\frac{\pi}{2},
    \]
    求 $\sin^2\alpha+\sin^2\beta$ 的值。
    \begin{solution}
        由
        \[
        \sin^{-1}(\sin\alpha+\sin\beta)=\frac{\pi}{2}-\sin^{-1}(\sin\alpha-\sin\beta)=\cos^{-1}(\sin\alpha-\sin\beta)
        \]
        得到
        \[
        \sin\alpha+\sin\beta=\sin (\cos^{-1}(\sin\alpha-\sin\beta))=\sqrt{1-(\sin\alpha-\sin\beta)^2}
        \]
        因此
        \[
        (\sin\alpha+\sin\beta)^2+(\sin\alpha-\sin\beta)^2=1 \Rightarrow \sin^2\alpha+\sin^2\beta=\frac{1}{2}
        \]
    \end{solution}

    \question 已知
    \[
    \sin^{-1}x+\sin^{-1}y=\frac{2\pi}{3},\quad
    \cos^{-1}x-\cos^{-1}y=\frac{\pi}{3},
    \]
    求 $x,y$ 的值。
    \begin{solution}
        由反三角恒等式 
        \[
        \sin^{-1}x+\cos^{-1}x=\frac{\pi}{2}
        \]
        故
        \[
        \cos^{-1}x-\cos^{-1}y
        =\left(\frac{\pi}{2}-\sin^{-1}x\right)-\left(\frac{\pi}{2}-\sin^{-1}y\right)
        =\sin^{-1}y-\sin^{-1}x=\frac{\pi}{3}
        \]
        于是可解得
        \[
        \sin^{-1}x=\frac{\pi}{6},\;\sin^{-1}y=\frac{\pi}{2}\Rightarrow x=\frac{1}{2},\;y=1
        \]
    \end{solution}

    \question 解方程
    \[
    \tan^{-1}\!\left[x\cos\!\left(2\sin^{-1}\frac{1}{x}\right)\right]=\frac{\pi}{4}.
    \]
    \begin{solution}
        两边取正切,得
        \[
        x\cos\left(2\sin^{-1}\frac{1}{x}\right)=1 \Rightarrow
        \cos\left(2\sin^{-1}\frac{1}{x}\right)=\frac{1}{x}
        \]
        于是有
        \[
        2\sin^{-1}\frac{1}{x}=\pm\cos^{-1}\frac{1}{x}\pm 2n\pi,\quad n=0,1,2,3,\dots
        \]
        利用恒等式$\sin^{-1}t+\cos^{-1}t=\dfrac{\pi}{2}$,可得
        \[
        2\sin^{-1}\frac{1}{x}
        =\pm\left(\frac{\pi}{2}-\sin^{-1}\frac{1}{x}\right)\pm 2n\pi.
        \]
        给出
        \[
        \sin^{-1}\frac{1}{x}=\frac{\pi}{6}\pm\frac{2n\pi}{3}
        \quad \text{或} \quad
        \sin^{-1}\frac{1}{x}=-\frac{\pi}{2}\pm 2n\pi
        \]
        由于$-\dfrac{\pi}{2}\leq \sin^{-1}t\leq \dfrac{\pi}{2}$,故取
        \[
        \sin^{-1}\frac{1}{x}=\frac{\pi}{6}\quad\text{或}\quad
        \sin^{-1}\frac{1}{x}=-\frac{\pi}{2}.
        \]
        即
        \[
        x=2 \quad \text{或} \quad x=-1
        \]
        经检验,可知都是原方程式的解。
    \end{solution}

    \question 解方程
    \[
    \tan^{-1}\left(\sqrt{2}(x+1)\right)-\tan^{-1}\left(\frac{x-1}{\sqrt{2}}\right)=\tan^{-1}\!\frac{\sqrt{2}}{3}.
    \]
    \begin{solution}
        原方程式两边取正切,
        \[
        \frac{\sqrt{2}(x+1)-\dfrac{x-1}{\sqrt{2}}}{1+\sqrt{2}(x+1)\cdot\dfrac{x-1}{\sqrt{2}}}=\frac{\sqrt{2}}{3}
        \]
        解得
        \[
        x=3\quad\text{或}\quad x=-\frac{3}{2}
        \]
        经检验都是方程的实数解。
    \end{solution}

    \question 解反三角方程
    \[
    \sin^{-1}\frac{2x}{1+x^2}=\tan^{-1}\frac{3}{x+7}
    \]
    \begin{solution}
        两边取正弦得
        \[
        \frac{2x}{1+x^2}=\sin\left(\tan^{-1}\frac{3}{x+7}\right)=\pm \frac{3}{\sqrt{3^2+(x+7)^2}}
        \]
        解得
        \[
        x=\frac{1}{5},\quad x=7+2\sqrt{13}
        \]
        经检验皆为方程的实数解。
    \end{solution}

    \question 解方程
    \[
    3\sin^{-1}2x+3\sin^{-1}x = 2\pi
    \]
    \begin{solution}
        原方程式即
        \[
        \sin^{-1}2x+\sin^{-1}x = \frac{2\pi}{3}
        \]
        对两边取正弦得
        \[
        2x\cos(\sin^{-1}x)+x\cos(\sin^{-1}2x) = \frac{\sqrt{3}}{2}
        \]
        即
        \[
        2x\sqrt{1-x^2}+x\sqrt{1-(2x)^2}=\frac{\sqrt{3}}{2}
        \]
        解得$x=\pm \dfrac{1}{2}$,经检验,原方程式的解为
        \[
        x=\frac{1}{2}
        \]
    \end{solution}

    \question 解方程
    \[
    \sin(\cot^{-1}(x+1)) = \cos(\tan^{-1} x)
    \]
    \begin{solution}
        由 $\cos A = \sin\left(\frac{\pi}{2} - A\right)$,
        \[ 
        \sin[\cot^{-1}(x+1)] = \sin\left[\frac{\pi}{2} - \tan^{-1} x\right] 
        \]
        于是解为
        \[
        \cot^{-1}(x+1) = \frac{\pi}{2} - \tan^{-1} x \quad \text{或} \quad \cot^{-1}(x+1) = \pi - \left(\frac{\pi}{2} - \tan^{-1} x\right) 
        \]
        又由 $\cot^{-1} A = \tan^{-1}\dfrac{1}{A}$,
        \[ 
        \tan^{-1}\left(\frac{1}{x+1}\right) + \tan^{-1} x = \frac{\pi}{2} \quad \text{或} \quad \tan^{-1}\left(\frac{1}{x+1}\right) - \tan^{-1} x = \frac{\pi}{2} 
        \]
        两边取正切,对于左侧方程
        \begin{align*}
        \tan\left[\tan^{-1}\left(\frac{1}{x+1}\right) + \tan^{-1} x\right] &= \tan \frac{\pi}{2} \\
        \frac{\frac{1}{x+1} + x}{1 - \frac{1}{x+1} \cdot x} &= \infty \\
        x^2 + x + 1 &= \infty
        \end{align*}
        得 $x = \pm \infty$,而对于右侧方程,
        \begin{align*}
        \tan\left[\tan^{-1}\left(\frac{1}{x+1}\right) - \tan^{-1} x\right] &= \tan \frac{\pi}{2} \\
        \frac{\frac{1}{x+1} - x}{1 + \frac{1}{x+1} \cdot x} &= \infty \\
        \frac{1 - x^2 - x}{2x + 1} &= \infty \\
        2x + 1 &= 0 \\
        x &= -\frac{1}{2}
        \end{align*}
        所以解为 $x = -\dfrac{1}{2}$,经检验都是方程的解。
    \end{solution}

    \question 解方程
    \[
    (\tan^{-1} x)^2+(\cot^{-1} x)^2=\frac{5\pi^2}{8}.
    \]
    \begin{solution}
        设$\theta=\tan^{-1}x$,原方程式变为
        \[
        \theta^2+\left(\frac{\pi}{2}-\theta\right)^2=\frac{5\pi^2}{8}.
        \]
        整理得
        \[
        (4\theta-3\pi)(4\theta+\pi)=0 \Rightarrow \theta=\frac{3\pi}{4}\quad \text{或}\quad \theta=-\frac{\pi}{4}.
        \]
        注意到$-\dfrac{\pi}{2}\le \tan^{-1}x\le \dfrac{\pi}{2}$,因此舍去 $\theta=\dfrac{3\pi}{4}$,于是
        \[
        x=\tan\left(-\frac{\pi}{4}\right)=-1
        \]
        经检验是原方程的解。
    \end{solution}

    \question 若 
    \[
    \tan^{-1} x + \tan^{-1} y + \tan^{-1} z = \frac{\pi}{2},
    \] 
    证明 
    \[
    xy + yz + zx = 1
    \]
    \begin{solution}
        设$\alpha = \tan^{-1} x, \beta = \tan^{-1} y, \gamma = \tan^{-1} z$,则
        \[
        \alpha + \beta + \gamma = \frac{\pi}{2}.
        \]
        利用三角恒等式
        \[
        \tan(\alpha+\beta+\gamma)
        =\frac{\tan\alpha + \tan\beta + \tan\gamma - \tan\alpha \tan\beta \tan\gamma}
        {1 - \tan\alpha \tan\beta - \tan\beta \tan\gamma - \tan\gamma \tan\alpha}.
        \]
        代入 $\tan\alpha=x,\ \tan\beta=y,\ \tan\gamma=z$,得
        \[
        \tan(\alpha+\beta+\gamma)
        =\frac{x + y + z - xyz}{1 - xy - yz - zx}.
        \]
        由于$\alpha+\beta+\gamma = \dfrac{\pi}{2},\tan(\alpha+\beta+\gamma)$不存在,则
        \[
        1 - xy - yz - zx = 0 \Rightarrow xy + yz + zx = 1
        \]
        注:其逆命题不成立。
    \end{solution}

    \question 若 $x>1$, 求证
    \[
    \tan^{-1} x + \tan^{-1} \frac{1-x}{1+x} = \frac{\pi}{4}.
    \]
    \begin{solution}
        设 $\alpha = \tan^{-1} x, \beta = \tan^{-1} \dfrac{1-x}{1+x}$,于是
        \[
        \tan(\alpha + \beta) = \frac{\tan \alpha + \tan \beta}{1 - \tan \alpha \tan \beta} = \frac{x + \frac{1-x}{1+x}}{1 - x \cdot \frac{1-x}{1+x}}= \frac{x^2+1}{1+x^2} = 1 
        \]
        由于 $x>1$, 
        \[
        \frac{\pi}{4} < \alpha < \frac{\pi}{2}, \quad -1 < \frac{1-x}{1+x} < 0 \Rightarrow -\frac{\pi}{4} < \beta < 0
        \]
        因此 
        \[
        0 < \alpha + \beta < \frac{\pi}{2}
        \]
        这表示
        \[
        \tan^{-1} x + \tan^{-1} \frac{1-x}{1+x} =\alpha + \beta= \frac{\pi}{4}
        \]
    \end{solution}

    \question 证明
    \[
    \tan^{-1}\frac{1}{2}+\tan^{-1}\frac{1}{5}+\tan^{-1}\frac{1}{8}=\frac{\pi}{4}.
    \]
    \begin{solution}
        设
        \[
        A=\tan^{-1}\frac{1}{2},\quad B=\tan^{-1}\frac{1}{5},\quad C=\tan^{-1}\frac{1}{8}
        \]
        利用正切的三角和公式,
        \[
        \tan(A+B+C)
        =\frac{\tan A+\tan B+\tan C-\tan A\cdot\tan B\cdot\tan C}
        {1-(\tan A\tan B+\tan B\tan C+\tan C\tan A)}
        \]
        代入 $\tan A=\dfrac{1}{2},\ \tan B=\dfrac{1}{5},\ \tan C=\dfrac{1}{8}$,得
        \[
        \tan(A+B+C)=1=\tan\frac{\pi}{4}
        \]
        现证明 $\tan^{-1} t<t$:设 $f(x)=\tan x - x$,则 $f'(x)=\sec^2 x - 1 = \tan^2 x > 0$,且 $f(0)=0$,故对所有$x>0$有$f(x)>0$,即 $\tan x > x$对$x=\tan^{-1}t \in \left(0,\dfrac{\pi}{2}\right)$恒成立,从而 $\tan^{-1} t < t$。因此
        \[
        A+B+C < \frac{1}{2}+\frac{1}{5}+\frac{1}{8}=\frac{33}{40}<\frac{\pi}{2}
        \]
        又 $A,B,C>0$,所以 $A+B+C\in\left(0,\dfrac{\pi}{2}\right)$。在此区间上 $\tan\theta$ 严格递增,且仅在 $\theta=\frac{\pi}{4}$ 时 $\tan\theta=1$,故
        \[
        A+B+C=\frac{\pi}{4}
        \]
        证毕。
    \end{solution}
    \begin{solution}
        由恒等式
        \[
        \tan^{-1} x + \tan^{-1} y = \tan^{-1} \frac{x+y}{1-xy}, \quad x>0,\;y>0,\;xy<1
        \]
        有
        \[
        \tan^{-1}\frac{1}{2}+\tan^{-1}\frac{1}{5} = \tan^{-1}\frac{\frac{1}{2}+\frac{1}{5}}{1-\frac{1}{2}\cdot\frac{1}{5}} = \tan^{-1}\frac{7}{9}
        \]
        故原式为
        \[
        \tan^{-1}\frac{7}{9} + \tan^{-1}\frac{1}{8} =\tan^{-1}\frac{\frac{7}{9}+\frac{1}{8}}{1-\frac{7}{9}\cdot\frac{1}{8}} = \tan^{-1}1 = \frac{\pi}{4}
        \]
        即得证。
    \end{solution}

    \question 化简
    \[
    \tan^{-1} \left( \frac{\sqrt{1+\sin x} - \sqrt{1-\sin x}}{\sqrt{1+\sin x} + \sqrt{1-\sin x}} \right), 
    \]
    其中$0 < x < \dfrac{\pi}{4}$。
    \begin{solution}
        有理化分母,可得
        \begin{align*}
        \tan^{-1} \left( \frac{\sqrt{1+\sin x} - \sqrt{1-\sin x}}{\sqrt{1+\sin x} + \sqrt{1-\sin x}} \right) 
        &= \tan^{-1} \left( \frac{(\sqrt{1+\sin x} - \sqrt{1-\sin x})^2}{(1+\sin x) - (1-\sin x)} \right) \\
        &= \tan^{-1} \left( \frac{2 - 2\sqrt{1-\sin^2 x}}{2\sin x} \right) \\
        &= \tan^{-1} \left( \frac{1 - \cos x}{\sin x} \right) \\
        &= \tan^{-1} \left( \frac{2\sin^2 \frac{x}{2}}{2\sin \frac{x}{2} \cos \frac{x}{2}} \right) \\
        &= \tan^{-1} \left(\tan \frac{x}{2}\right) \\
        &= \frac{x}{2}
        \end{align*}
    \end{solution}

    \question 证明
    \[ 
    f(x) = \tan^{-1} (3x) + \sin^{-1} \left[(9x^2+1)^{-\frac{1}{2}}\right] 
    \]
    为一常数,并求此常数。
    \begin{solution}
        对 $x$ 求导,
        \begin{align*}
        f'(x)
        &=\frac{3}{1+(3x)^2}
        +\frac{1}{\sqrt{1-(9x^2+1)^{-1}}}
        \cdot\left(-\frac{1}{2}\right)(18x)(9x^2+1)^{-\frac{3}{2}} \\
        &=\frac{3}{1+9x^2}
        -\frac{9x}{\sqrt{\frac{9x^2}{9x^2+1}}}\cdot\frac{1}{(9x^2+1)^{\frac{3}{2}}} \\
        &=\frac{3}{1+9x^2}
        -3\sqrt{9x^2+1}\cdot\frac{1}{(9x^2+1)^{\frac{3}{2}}} \\
        &=\frac{3}{1+9x^2}-\frac{3}{9x^2+1} \\
        &=0.
        \end{align*}
        因此 $f(x)$ 为常数,取 $x=0$得
        \[
        f(0)=\tan^{-1}0+\sin^{-1}(1)=\frac{\pi}{2}.
        \]
        于是
        \[
        \tan^{-1}(3x)+\sin^{-1}\!\left(\frac{1}{\sqrt{9x^2+1}}\right)=\frac{\pi}{2}
        \]
    \end{solution}

    \question 若 $0 \le x \le 1$, 求证
    \[
    \cos(\sin^{-1} x) < \sin^{-1}(\cos x)
    \]
    \begin{solution}
        设 $\alpha = \sin^{-1} x$,则 $\sin\alpha = x$. 又 $x \in [0, 1],\alpha \in \left[0, \dfrac{\pi}{2}\right]$,从而有 
        \[
        \sqrt{1-x^2} = \cos\alpha
        \]
        故
        \[ x+\sqrt{1-x^2} = \sin\alpha + \cos\alpha = \sqrt{2} \sin\left(\alpha + \frac{\pi}{4}\right) \le \sqrt{2} < \frac{\pi}{2}
        \]
        即 
        \[
        \sqrt{1-x^2} < \frac{\pi}{2} - x
        \]
        由于$x \in [0, 1],x=\cos^{-1}(\cos x) =x$,故得证 
        \[
        \cos(\sin^{-1} x) < \sin^{-1}(\cos x),\quad 0 \le x \le 1
        \]
    \end{solution}

    \question 若 $\alpha, \beta \in (0, 2\pi), \ \alpha \neq \beta$, 且 $\alpha, \beta$ 满足方程
    \[
    \sin x + \sqrt{3} \cos x + a = 0,
    \]
    求实数 $a$ 的取值范围,以及 $\alpha + \beta$ 的值。
    \begin{solution}
        注意到原方程为
        \[
        \sin\left(x + \frac{\pi}{3}\right) = -\frac{a}{2}
        \]
        所以 
        \[
        -\frac{a}{2} \in [-1,1] \Rightarrow a \in [-2,2]
        \]
        设 $\alpha,\beta$ 为原方程的解,有
        \[
        \begin{cases}
        \sin \alpha + \sqrt{3} \cos \alpha + a = 0, \\
        \sin \beta + \sqrt{3} \cos \beta + a = 0.
        \end{cases}
        \]
        两式相减得
        \[
        (\sin\alpha - \sin\beta) + \sqrt{3} (\cos\alpha - \cos\beta) = 0
        \]
        和差化积得
        \[
        2 \cos \frac{\alpha+\beta}{2} \sin \frac{\alpha-\beta}{2} - 2\sqrt{3} \sin \frac{\alpha+\beta}{2} \sin \frac{\alpha-\beta}{2} = 0
        \]
        由于 $\alpha \ne \beta$,有 $\sin \dfrac{\alpha-\beta}{2} \ne 0$,于是
        \[
        \cos \frac{\alpha+\beta}{2} - \sqrt{3} \sin \frac{\alpha+\beta}{2} = 0 \Rightarrow \tan \frac{\alpha+\beta}{2}  = \frac{\sqrt{3}}{3}
        \]
        又 $\alpha+\beta \in (0,4\pi)$,因此
        \[
        \alpha + \beta = \frac{\pi}{3} \quad \text{或} \quad  \frac{7\pi}{3}
        \]
    \end{solution}

    \question 已知方程
    \[
    (\sin^{-1} x)^3 + (\cos^{-1} x)^3 = k\pi^3, |x| \le 1,
    \]
    其中 $k$ 为常数。
    \begin{parts}
    \part 证明该方程有解的一个必要非充分条件是 $k \ge \frac{1}{32}$。
    \begin{solution}
        由恒等式 $\sin^{-1} x + \cos^{-1} x = \dfrac{\pi}{2}$,
        \[
        (\sin^{-1} x)^3 + \left(\frac{\pi}{2} - \sin^{-1} x\right)^3 = k\pi^3 
        \]
        整理得
        \[
        (\sin^{-1} x)^2 - \frac{\pi}{2}(\sin^{-1} x) + \frac{\pi^2}{12}(1-8k) = 0 \tag{1}
        \]
        欲使(1)有实数解,判别式为非负:
        \begin{align*}
        \left(-\frac{\pi}{2}\right)^2 - 4 \cdot 1 \cdot \frac{\pi^2}{12}(1-8k) &\ge 0 \\
        \frac{\pi^2}{4} - \frac{\pi^2}{3}(1-8k) &\ge 0 \\
        \frac{1}{4} - \frac{1}{3} + \frac{8}{3}k &\ge 0 \\
        32k &\ge 1 \\
        k &\ge \frac{1}{32}
        \end{align*}
        由于条件$k \ge \dfrac{1}{32}$可能会生成 $|\sin^{-1} x| > 1$的解即非实数解,故此条件为必要非充分条件。
    \end{solution}
    \part 若方程只有一个解,求该解。
    \begin{solution}
        若只有一个解,则 $k = \dfrac{1}{32}$,方程变为
        \[
        (\sin^{-1} x)^2 - \frac{\pi}{2}(\sin^{-1} x) + \frac{\pi^2}{12}\left(1 - 8 \cdot \frac{1}{32}\right) = 0 
        \]
        即
        \[
        \left(\sin^{-1} x - \frac{\pi}{4}\right)^2 = 0 
        \]
        解得
        \[
        \sin^{-1} x = \frac{\pi}{4} \Rightarrow x = \frac{\sqrt{2}}{2}
        \]
    \end{solution}
    \part 若 $k = \dfrac{7}{96}$,求方程的两个解。
    \begin{solution}
        若 $k = \dfrac{7}{96}$,将方程配方可得
        \[
        \left(\sin^{-1} x - \frac{\pi}{4}\right)^2 = \frac{\pi^2}{36} 
        \]
        于是解为
        \[
        \sin^{-1} x - \frac{\pi}{4} = \pm \frac{\pi}{6} \Rightarrow
        \sin^{-1} x = \frac{5\pi}{12} \quad \text{或} \quad \sin^{-1} x = \frac{\pi}{12} 
        \]
        即
        \[
        x = \sin\frac{5\pi}{12} \quad \text{或} \quad  x = \sin\frac{\pi}{12}
        \]
    \end{solution}
    \end{parts}
\end{questions}

\pagebreak

\begin{center}
  {\fontsize{30pt}{26pt}\selectfont
    \hypertarget{平面向量}{平面向量} \label{平面向量}
  }
\end{center}
\separator
\vspace{1pt}

\begin{questions}
    \question 在坐标平面上的平行四边形 \(ABCD\) 中,若 \(\overrightarrow{AB} = (4,8)\),\(\;\overrightarrow{AD} = (1,4)\),求 $|\overrightarrow{AC}| + |\overrightarrow{BD}|$。
    \begin{solution}
        发现
        \[
        \overrightarrow{AC} = \overrightarrow{AB} + \overrightarrow{AD} = (4 + 1, 8 + 4) = (5,12)  
        \]
        且
        \[
        \overrightarrow{BD} = \overrightarrow{AD} - \overrightarrow{AB} = (1 - 4, 4 - 8) = (-3,-4)  
        \]
        因此
        \[
        |\overrightarrow{AC}| + |\overrightarrow{BD}| =\sqrt{5^2 + 12^2}+\sqrt{(-3)^2 + (-4)^2} = 18
        \]
    \end{solution}

    \question 在 \(\triangle ABC\) 中,已知 \(AB=4,BC=5,AC=6\),求 $\overrightarrow{AB} \cdot \overrightarrow{BC}$ 。
    \begin{solution}
    由余弦定理,
    \[
    6^2 = 4^2 + 5^2 - 2 \cdot 4 \cdot 5 \cdot \cos{\angle ABC}
    \Rightarrow \cos{\angle ABC} = \frac{1}{8}
    \]
    所以
    \begin{align*}
        \overrightarrow{AB} \cdot \overrightarrow{BC} &= |\overrightarrow{AB}| \cdot |\overrightarrow{BC}| \cdot \cos{\theta} \\
        &= 4 \cdot 5 \cdot \cos{(\pi - \angle ABC)} \\
        &= 4 \cdot 5 \cdot (-\cos{\angle ABC}) \\
        &= -\frac{5}{2}
    \end{align*}
    \end{solution}
    
    \question 已知 \(\mathbf{a} + \mathbf{b} + \mathbf{c} = \mathbf{0}\),且 \(|\mathbf{a}|=3\), \(|\mathbf{b}|=5\), \(|\mathbf{c}|=7\),试求 $|2\mathbf{a} + \mathbf{b}| $。
    \begin{solution}
        由已知
        \[
        \mathbf{a} + \mathbf{b} + \mathbf{c} = \mathbf{0}
        \Rightarrow \mathbf{a} + \mathbf{b} = -\mathbf{c} \Rightarrow |\mathbf{a} + \mathbf{b}| = |-\mathbf{c}| = |\mathbf{c}| \Rightarrow |\mathbf{a} + \mathbf{b}|^2 = |\mathbf{c}|^2
        \]
        展开得
        \[ |\mathbf{a}|^2 + 2\mathbf{a} \cdot \mathbf{b} + |\mathbf{b}|^2 = |\mathbf{c}|^2
        \Rightarrow 3^2 + 2\mathbf{a} \cdot \mathbf{b} + 5^2 = 7^2
        \Rightarrow \mathbf{a} \cdot \mathbf{b} = \frac{15}{2}
        \]
        于是
        \[
         |2\mathbf{a} + \mathbf{b}|^2 = 4|\mathbf{a}|^2 + 4\mathbf{a} \cdot \mathbf{b} + |\mathbf{b}|^2
        = 4 \cdot 3^2 + 4 \cdot \frac{15}{2} + 5^2 = 91
        \Rightarrow |2\mathbf{a} + \mathbf{b}| = \sqrt{91}
        \]
    \end{solution}

    \question 已知 $|\mathbf{a}|=1$, 且 $\mathbf{b}\cdot(\mathbf{b}-\mathbf{a})=0$, 求 $|\mathbf{b}-\mathbf{a}||\mathbf{b}+\mathbf{a}|$ 的最大值。
    \begin{solution}
        由 $\mathbf{b}\cdot(\mathbf{b}-\mathbf{a})=0$ 得
        \[
        |\mathbf{b}|^2 - \mathbf{a}\cdot \mathbf{b} = 0 \implies \mathbf{a}\cdot \mathbf{b} = |\mathbf{b}|^2
        \]
        于是
        \begin{align*}
        |\mathbf{b}-\mathbf{a}|^2|\mathbf{b}+\mathbf{a}|^2 
        &= [(\mathbf{b}-\mathbf{a})\cdot(\mathbf{b}-\mathbf{a})][(\mathbf{b}+\mathbf{a})\cdot(\mathbf{b}+\mathbf{a})] \\
        &= (|\mathbf{b}|^2 - 2 \mathbf{a}\cdot \mathbf{b} + |\mathbf{a}|^2)(|\mathbf{b}|^2 + 2 \mathbf{a}\cdot \mathbf{b} + |\mathbf{a}|^2) \\
        &= (|\mathbf{b}|^2 - 2|\mathbf{b}|^2 + 1)(|\mathbf{b}|^2 + 2|\mathbf{b}|^2 + 1) \\
        &= (1 - |\mathbf{b}|^2)(3|\mathbf{b}|^2 + 1)
        \end{align*}
        设 $|\mathbf{b}|^2 = m$, 配方得
        \[
        |\mathbf{b}-\mathbf{a}|^2|\mathbf{b}+\mathbf{a}|^2 = (1 - m)(3m + 1) = -3m^2 + 2m + 1 = -3\left(m - \frac{1}{3}\right)^2 + \frac{4}{3}
        \]
        当 $m = \dfrac{1}{3}$ 时取得最大值$\dfrac{4}{3}$,此时
        \[
        |\mathbf{b}-\mathbf{a}||\mathbf{b}+\mathbf{a}| = \frac{2\sqrt{3}}{3}
        \]
    \end{solution}

    \question 平面上非零向量 $\mathbf{a},\mathbf{b},\mathbf{c}$ 满足 $\mathbf{a}\perp\mathbf{b}$, 且 $\mathbf{b}\cdot\mathbf{c}=20|\mathbf{a}|\cdot|\mathbf{b}|,\mathbf{c}\cdot\mathbf{a}=25|\mathbf{b}|\cdot|\mathbf{c}|$, 求 $|\mathbf{c}|$ 的最小值。
    \begin{solution}
        设向量 $\mathbf{b}$ 与 $\mathbf{c}$ 的夹角为 $\theta$, 又因 $\mathbf{a}\perp\mathbf{b}$, 所以 $\mathbf{a}$ 与 $\mathbf{c}$ 的夹角为 $90^\circ-\theta$。由内积定义,
        \[
        \cos\theta=\frac{\mathbf{b}\cdot\mathbf{c}}{|\mathbf{b}||\mathbf{c}|},\quad
        \cos(90^\circ-\theta)=\frac{\mathbf{c}\cdot\mathbf{a}}{|\mathbf{c}||\mathbf{a}|}.
        \]
        据题意有,
        \[
        |\mathbf{c}|\cos\theta=\frac{20|\mathbf{a}|}{|\mathbf{b}|},\quad
        |\mathbf{c}|\sin\theta=\frac{25|\mathbf{b}|}{|\mathbf{a}|}.
        \]
        由AM-GM不等式,
        \[
        |\mathbf{c}|^2
        =\frac{400|\mathbf{a}|^2}{|\mathbf{b}|^2}
        +\frac{625|\mathbf{b}|^2}{|\mathbf{a}|^2}
        \ge 2\sqrt{400\cdot 625}=1000
        \]
        所以
        \[
        |\mathbf{c}|\ge 10\sqrt{10}
        \]
    \end{solution}
    
    \question 设 $\vec{a},\vec{b}$ 为两个垂直的平面向量, 且 $|\vec{a}|=2|\vec{b}|=10$. 当 $0 \le t \le 1$ 时, 记向量 $t\vec{a}+(1-t)\vec{b}$ 与向量 $(t-\frac{1}{5})\vec{a}+(1-t)\vec{b}$ 最大夹角为 $\theta$, 则 $\cos\theta=$
    \begin{solution}
        \textcolor{red}{(待解)}
    \end{solution}

    \question 已知 $\triangle ABC$ 中,$\ AB=3,AC=6,BC=7,O$ 为外心,$I$ 为内心,求 $\overrightarrow{OI} \cdot \overrightarrow{BC}$。
    \ifprintanswers
    \begin{figure}[H]
        \centering
        \includegraphics[width=0.4\textwidth]{images/image32.png}
    \end{figure}
    \fi
    \begin{solution}
        内心$I$的向量表示为
        \[
        \overrightarrow{OI} 
        = \frac{7\overrightarrow{OA}+6\overrightarrow{OB}+3\overrightarrow{OC}}{7+6+3}
        \]
        于是
        \begin{align*}
        \overrightarrow{OI} \cdot \overrightarrow{BC} 
        &= \frac{7}{16} \overrightarrow{OA} \cdot \left( \overrightarrow{BA} + \overrightarrow{AC} \right)
        + \frac{6}{16} \overrightarrow{OB} \cdot \overrightarrow{BC}
        + \frac{3}{16} \overrightarrow{OC} \cdot \overrightarrow{BC}\\
        &= \frac{7}{16} \left( \frac{1}{2} \cdot 3^2 - \frac{1}{2} \cdot 6^2 \right)
        + \frac{6}{16} \left( -\frac{1}{2} \cdot 7^2 \right)
        + \frac{3}{16} \left( \frac{1}{2} \cdot 7^2 \right)
        = -\frac{21}{2}
        \end{align*}
        其中$O$为外心时,有
        \begin{align*}  
        \overrightarrow{OX} \cdot \overrightarrow{YX}
        &=\overrightarrow{OX} \cdot (\overrightarrow{OX}-\overrightarrow{OY})\\
        &=|\overrightarrow{OX}|^2-\overrightarrow{OX} \cdot \overrightarrow{OY}\\
        &=|\overrightarrow{OX}|^2-\frac{1}{2}(|\overrightarrow{OX}|^2+|\overrightarrow{OY}|^2-|\overrightarrow{XY}|^2)\\
        &=\frac{1}{2}|\overrightarrow{XY}|^2
        \end{align*}
    \end{solution}

\question 若 $H$ 为 $\triangle ABC$ 的垂心,且
\[
\overrightarrow{AH} = \frac{1}{3}\overrightarrow{AB} + \frac{1}{4}\overrightarrow{AC},
\]
求 $\cos \angle BAC$。
\ifprintanswers
\begin{figure}[H]
    \centering
    \includegraphics[width=0.5\linewidth]{images/image53.png}
\end{figure}
\fi
\begin{solution}
因 $A,D,C$ 共线,且
\[
\overrightarrow{AH} = \frac{1}{3}\overrightarrow{AB} + \frac{1}{4}\overrightarrow{AC} = \frac{1}{3}\overrightarrow{AB} + \frac{2}{3}\overrightarrow{AD} \implies \overline{AD} : \overline{AC} = 3 : 8,????
\]
设
\[
A(0,0),\quad D(3k,0),\quad C(8k,0),\quad H(3k,t),\quad B(3k,s),
\]
则
\[
(3k,t) = \frac{1}{3}(3k,s) + \frac{1}{4}(8k,0) \implies s=3t,
\]
所以
\[
B(3k,3t).
\]

由垂心定义,有
\[
\overrightarrow{AH} \perp \overrightarrow{BC} \implies (3k,t) \cdot (5k, -3t) = 15k^2 - 3t^2 = 0 \implies t = \sqrt{5}k.
\]

计算
\[
\cos \angle BAC = \frac{\overrightarrow{AB} \cdot \overrightarrow{AC}}{|\overrightarrow{AB}||\overrightarrow{AC}|} = \frac{24k^2}{\sqrt{9k^2 + 9t^2} \cdot 8k} = \frac{24k^2}{\sqrt{54k^2} \cdot 8k} = \frac{1}{\sqrt{6}} = \frac{\sqrt{6}}{6}.
\]
\textcolor{red}{(待验证)}
\end{solution}

    \question 已知三角形 \(ABC\) 的外心 \(O\) 满足
    \[
    2\overrightarrow{OA}+3\overrightarrow{OB}+4\overrightarrow{OC}=0.
    \]
    求\(\sin A\).
    \begin{solution}
        由题意
        \[        
        3\overrightarrow{OB}+4\overrightarrow{OC}=-2\overrightarrow{OA}.
        \]
        两边平方得
        \[        
        9|\overrightarrow{OB}|^2+24\overrightarrow{OB}\cdot\overrightarrow{OC}+16|\overrightarrow{OC}|^2=4|\overrightarrow{OA}|^2
        \]
        注意到外心到各顶点距离相等,设 \(R=|\overrightarrow{OA}|=|\overrightarrow{OB}|=|\overrightarrow{OC}|\),代入得
        \[        
        9R^2+24\overrightarrow{OB}\cdot\overrightarrow{OC}+16R^2=4R^2.
        \]
        整理得
        \[
        \overrightarrow{OB}\cdot\overrightarrow{OC}=-\frac{7}{8}R^2.
        \]
        $\text{又因圆心角}=2\times\text{圆周角}$,
        \[        
        \overrightarrow{OB}\cdot\overrightarrow{OC}=|\overrightarrow{OB}||\overrightarrow{OC}|\cos\angle BOC=R^2\cos 2A,
        \]
        有
        \[
        \cos 2A=1-2\sin^2A=-\frac{7}{8}
        \Rightarrow \sin A=\frac{\sqrt{15}}{4}\quad(\sin A>0).
        \]
    \end{solution}
    
    \question 设 $\triangle ABC$ 的外心为$O$,重心为$H$,若 $|\overrightarrow{OA}+\overrightarrow{OB}+\overrightarrow{OC}|=\sqrt{3}$,求 $|\overrightarrow{HA}+\overrightarrow{HB}+\overrightarrow{HC}|$。
    \begin{solution}
        设 $O$ 为外心,$G$ 为重心,$H$ 为垂心,由于 $G$ 是重心,有 $\overrightarrow{GA} + \overrightarrow{GB} + \overrightarrow{GC} = \overrightarrow{0}$,于是
        \[
        |\overrightarrow{OA} + \overrightarrow{OB} + \overrightarrow{OC}| 
        = |3\overrightarrow{OG} + (\overrightarrow{GA} + \overrightarrow{GB} + \overrightarrow{GC})|=3|\overrightarrow{OG}| = \sqrt{3} \Rightarrow |\overrightarrow{OG}| = \frac{\sqrt{3}}{3}
        \]
        且
        \[
        |\overrightarrow{HA} + \overrightarrow{HB} + \overrightarrow{HC}|
        = |3\overrightarrow{HG} + (\overrightarrow{GA} + \overrightarrow{GB} + \overrightarrow{GC})|
        = 3|\overrightarrow{HG}|
        \]
        由欧拉线定理,
        \[
        |\overrightarrow{HG}| = 2|\overrightarrow{OG}| = \frac{2\sqrt{3}}{3}
        \]
        所以
        \[
        |\overrightarrow{HA} + \overrightarrow{HB} + \overrightarrow{HC}| = 2\sqrt{3}
        \]
    \end{solution}

    \question 
    \begin{parts}
    \part 在 $\triangle ABC$ 中,有 $\overrightarrow{AB} \cdot \overrightarrow{BC} = \overrightarrow{BC} \cdot \overrightarrow{CA} = \overrightarrow{CA} \cdot \overrightarrow{AB}$, 求$\cos A$。 
    \begin{solution}
        设 $a=|\overrightarrow{BC}|,b=|\overrightarrow{AC}|,c=|\overrightarrow{AB}|$,由余弦定理,
        \[
        a^2=b^2+c^2-2bc\cos A=b^2+c^2-2|\overrightarrow{AC}||\overrightarrow{AB}|\cos A
        \]
        故
        \[
        |\overrightarrow{AC}||\overrightarrow{AB}|\cos A
        =-\overrightarrow{CA}\cdot\overrightarrow{AB}
        =\frac{b^2+c^2-a^2}{2}.
        \]
        同理,
        \[
        |\overrightarrow{AB}||\overrightarrow{BC}|\cos B
        =-\overrightarrow{AB}\cdot\overrightarrow{BC}
        =\frac{a^2+c^2-b^2}{2},
        \]
        \[
        |\overrightarrow{AC}||\overrightarrow{BC}|\cos C
        =-\overrightarrow{BC}\cdot\overrightarrow{CA}
        =\frac{a^2+b^2-c^2}{2}.
        \]
        因此
        \[
        \frac{a^2+b^2-c^2}{2}=\frac{b^2+c^2-a^2}{2}=\frac{a^2+c^2-b^2}{2}
        \Rightarrow a=b=c
        \]
        即$\triangle ABC$ 为等边三角形,有
        \[
        \cos A=\frac{1}{2}
        \]
    \end{solution} 
    \part 在$\triangle ABC$ 中, $2\,\overrightarrow{AB} \cdot \overrightarrow{BC} = \overrightarrow{BC} \cdot \overrightarrow{CA} = 3\,\overrightarrow{CA} \cdot \overrightarrow{AB}$, 求 $\cos A$。
    \begin{solution}
        由上,
        \[
        3\,\overrightarrow{CA}\cdot\overrightarrow{AB}=-\frac{3}{2}(b^2+c^2-a^2)
        \]
        \[
        \overrightarrow{BC}\cdot\overrightarrow{CA}=-\frac12(a^2+b^2-c^2)
        \]
        \[
        2\,\overrightarrow{AB}\cdot\overrightarrow{BC}=-(a^2+c^2-b^2)
        \]
        设
        \[
        \frac{3}{2}(b^2+c^2-a^2)
        =\frac12(a^2+b^2-c^2)
        =a^2+c^2-b^2=3t
        \]
        可解得
        \[
        a^2=\frac{9t}{2},b^2=4t,c^2=\frac{5t}{2}
        \]
        于是
        \[
        \cos A=\frac{b^2+c^2-a^2}{2bc}
        =\frac{2t}{2\sqrt{10}\,t}
        =\frac{\sqrt{10}}{10}
        \]
    \end{solution}
    \end{parts}

    \question 已知 $G$ 为 $\triangle ABC$ 的重心,且 
    \[
    6\overrightarrow{GA}\cdot\overrightarrow{AB} = 3\overrightarrow{GB}\cdot\overrightarrow{BC} = 2\overrightarrow{GC}\cdot\overrightarrow{CA},
    \] 
    求 $\sin A$。
    \begin{solution}
        设 $AB=c,BC=a,CA=b$,由余弦定理
        \[
        \overrightarrow{GA}\cdot \overrightarrow{AB}=-\overrightarrow{AG}\cdot \overrightarrow{AB}
        =-\frac{1}{3}(\overrightarrow{AB}+\overrightarrow{AC})\cdot \overrightarrow{AB}
        =-\frac{1}{3} \left(c^2+\frac{b^2+c^2-a^2}{2}\right)
        =-\frac{b^2+3c^2-a^2}{6},
        \]
        同理
        \[
        \overrightarrow{GB}\cdot\overrightarrow{BC}
        =-\frac{c^2+3a^2-b^2}{6},\quad
        \overrightarrow{GC}\cdot\overrightarrow{CA}
        =-\frac{a^2+3b^2-c^2}{6}.
        \]
        由题意,可解得
        \[
        b^2+3c^2-a^2=\frac{1}{2}\left(c^2+3a^2-b^2\right)
        =\frac{1}{3}\left(a^2+3b^2-c^2\right) \Rightarrow a^2=b^2=\frac{5}{2}c^2
        \]
        由余弦定理,
        \[
        \cos A=\frac{b^2+c^2-a^2}{2bc}=\frac{1}{2}\cdot\frac{c}{b}
        =\frac{1}{2}\cdot\sqrt{\frac{2}{5}}=\frac{\sqrt{10}}{10}
        \]
        所以
        \[
        \sin A=\sqrt{1-\cos^2 A}=\frac{3\sqrt{10}}{10}.
        \]
    \end{solution}

    \question 圆之内接四边形 $ABCD$满足 
    \[
    \overrightarrow{AC} = \frac{3}{2}\overrightarrow{AB} + \frac{5}{2}\overrightarrow{AD}
    \]
    求 $\dfrac{\sin \angle DAB}{\sin \angle ABC}$。
    \ifprintanswers
    \begin{figure}[H]
        \centering
        \includegraphics[width=0.4\linewidth]{images/image45.png}
    \end{figure}
    \fi
    \begin{solution}
        设 $AC$ 与 $BD$ 交于 $P$,并令
        \[
        \overrightarrow{AP} = t\overrightarrow{AC} = \frac{3t}{2}\overrightarrow{AB} + \frac{5t}{2}\overrightarrow{AD} 
        \]
        由于$P$在$BD$上,解得
        \[
        \frac{3t}{2} + \frac{5t}{2} = 1 \Rightarrow t = \frac{1}{4}
        \]
        于是
        \[
        \overrightarrow{AP} = \frac{3}{8}\overrightarrow{AB} + \frac{5}{8}\overrightarrow{AD} \Rightarrow
        DP = 3k,PB = 5k,k>0
        \]
        同理,
        \[
        \overrightarrow{BD} = \overrightarrow{BA} + \overrightarrow{AD} = \overrightarrow{BA} + \frac{2}{5}(\overrightarrow{AC} - \frac{3}{2}\overrightarrow{AB}) = \frac{8}{5}\overrightarrow{BA} + \frac{2}{5}\overrightarrow{AC} = \frac{6}{5}\overrightarrow{BA} + \frac{2}{5}\overrightarrow{BC}
        \]
        解得
        \[
        \overrightarrow{BP} = t\overrightarrow{BD} = \frac{6t}{5}\overrightarrow{BA} + \frac{2t}{5}\overrightarrow{BC} \Rightarrow t = \frac{5}{8}
        \Rightarrow \overrightarrow{BP} = \frac{3}{4}\overrightarrow{BA} + \frac{1}{4}\overrightarrow{BC}
        \Rightarrow AP = m, PC = 3m
        \]
        由圆幂定理,
        \[
        AP \cdot PC = BP \cdot PD \Rightarrow 3m^2 = 15k^2 \Rightarrow \frac{k}{m} = \frac{1}{\sqrt{5}}
        \]
        由正弦定理,
        \[
        \frac{\sin \angle DAB}{\sin \angle ABC} = \frac{AC}{BD} = \frac{8k}{4m} = \frac{2\sqrt{5}}{5}
        \]
    \end{solution}

    \question 设 $\triangle ABC$ 的三边 $BC,CA,AB$ 的长度分别为 $a,b,c$, 在边 $BC,CA,AB$ 上分别取点 $L,M,N$使得 $BL:LC=c:b,CM:MA=a:c,AN:NB=b:a$。若 
    \[
    b\overrightarrow{BM}+c\overrightarrow{CN}+a\overrightarrow{AL}=\vec{0},
    \]
    试证$\triangle ABC$ 为正三角形。
    \begin{solution}
        据题意有
        \begin{align*}
        \overrightarrow{AL}&=\frac{b}{b+c}\overrightarrow{AB}+\frac{c}{b+c}\overrightarrow{AC} \\[1mm]
        \overrightarrow{BM}&=-\overrightarrow{AB}+\overrightarrow{AM}
        =-\overrightarrow{AB}+\frac{c}{a+c}\overrightarrow{AC} \\
        \overrightarrow{CN}
        &=\overrightarrow{CB}+\overrightarrow{BN}
        =\overrightarrow{AB}-\overrightarrow{AC}-\frac{a}{a+b}\overrightarrow{AB}
        =\frac{b}{a+b}\overrightarrow{AB}-\overrightarrow{AC}
        \end{align*}
        故$b\overrightarrow{BM}+c\overrightarrow{CN}+a\overrightarrow{AL}$ 为
        \[
        b\left(-\overrightarrow{AB}+\frac{c}{a+c}\overrightarrow{AC}\right)
        +c\left(\frac{b}{a+b}\overrightarrow{AB}-\overrightarrow{AC}\right)
        +a\left(\frac{b}{b+c}\overrightarrow{AB}+\frac{c}{b+c}\overrightarrow{AC}\right)
        \]
        \[
        =\left(-b+\frac{bc}{a+b}+\frac{ab}{b+c}\right)\overrightarrow{AB}
        +\left(\frac{bc}{a+c}-c+\frac{ac}{b+c}\right)\overrightarrow{AC}=\overrightarrow{0}
        \]
        于是由余弦定理,
        \[
        -b+\frac{bc}{a+b}+\frac{ab}{b+c}=0
        \Rightarrow \frac{c}{a+b}+\frac{a}{b+c}=1
        \Rightarrow b^2=a^2+c^2-ac
        \Rightarrow \angle B=60^\circ
        \]
        同理
        \[
        \frac{bc}{a+c}-c+\frac{ac}{b+c}=0
        \Rightarrow \angle C=60^\circ
        \]
        因此得证 $\triangle ABC$ 为正三角形。
    \end{solution}

    \question 点 \(O\) 在三角形 \(\triangle ABC\) 内,且满足 $\overrightarrow{OA} + 2\overrightarrow{OB} + 3\overrightarrow{OC} = \overrightarrow{0}$,求三角形面积的比值 $\frac{[\triangle ABC]}{[\triangle AOC]} =$ \ans{$3$}
    \begin{solution}
        \textcolor{red}{(待解)}
    \end{solution}

    \question 已知两向量 $|\textbf{a}| = |\textbf{b}| = 1$, 且 $\textbf{a}\cdot\textbf{b}=\dfrac12$, 若向量 $\textbf{c}$ 满足 $\textbf{a}-\textbf{c}$ 与 $\textbf{b}-\textbf{c}$ 夹角为 $120^\circ$, 求 $|\textbf{c}|$ 的最小值。
    \ifprintanswers
    \begin{figure}[H]
        \centering
        \includegraphics[width=0.4\linewidth]{images/image228.png}
    \end{figure}
    \fi
    \begin{solution}
        设$\textbf{a}=\overrightarrow{OA}, \textbf{b}=\overrightarrow{OB}, \textbf{c}=\overrightarrow{OC}$,其中 $O$ 为原点。由
        \[
        |\textbf{a}|=|\textbf{b}|=1,\quad \textbf{a}\cdot\textbf{b}=\frac12
        \]
        可知 $A,B$ 在单位圆上,且 $\angle AOB=60^\circ$。当 $|\textbf{c}|$ 最小时,有$|\textbf{a}-\textbf{c}|=|\textbf{b}-\textbf{c}|$,此时 $C$ 在 $AB$ 的中垂线上,且靠近原点 $O$。设$A(1,0), B\left(\dfrac12,\dfrac{\sqrt3}{2}\right)$, 且$P$ 为 $AB$ 中点,$OP$ 为等边 $\triangle OAB$ 的高,有
        \[
        OP=\frac{\sqrt3}{2}.
        \]
        $\triangle ACP$ 为 $30^\circ-60^\circ-90^\circ$ 三角形,故
        \[
        CP=\frac{1}{2\sqrt3}
        \]
        因此
        \[
        |\textbf{c}|=OC=OP-CP=\frac{\sqrt3}{3}
        \]
    \end{solution}

    \question 已知 $\textbf{a}$, $\textbf{b}$, $\textbf{c}$ 为三个非零向量,其中 $|\textbf{a}|=4$, $|\textbf{b}|=6$, $\textbf{a}$ 在 $\textbf{b}$ 上的正射影长为 1,若 $(\textbf{c}-\textbf{a})\cdot(\textbf{c}-\textbf{b})=0$,试求 $|\textbf{c}|$ 的最大可能值。
    \begin{solution}
        设
        \[
        \begin{cases}
        B(6,0) \\
        A(1,\sqrt{15}) \\
        C(a,b)
        \end{cases}
        \Rightarrow
        \begin{cases}
        \textbf a = \overrightarrow{OA} \\
        \textbf b = \overrightarrow{OB} \\
        \textbf c = \overrightarrow{OC}
        \end{cases}
        \Rightarrow
        \begin{cases}
        \textbf c - \textbf a = (a-1, b - \sqrt{15}) \\
        \textbf c - \textbf b = (a - 6, b)
        \end{cases}
        \]
        则
        \[
        (\textbf c - \textbf a) \cdot (\textbf c - \textbf b) = (a - 1)(a - 6) + b(b - \sqrt{15}) = 0
        \]
        由拉格朗日乘子法,令
        \[
        \begin{cases}
        f(a,b) = a^2 + b^2 \\
        g(a,b) = (a - 1)(a - 6) + b(b - \sqrt{15})
        \end{cases}
        \Rightarrow
        \begin{cases}
        f_a = \lambda g_a \\
        f_b = \lambda g_b
        \end{cases}
        \Rightarrow
        \begin{cases}
        2a = (2a - 7)\lambda \\
        2b = (2b - \sqrt{15})\lambda
        \end{cases}
        \]
        则
        \[
        \frac{a}{b} = \frac{2a - 7}{2b - \sqrt{15}} \Rightarrow b = \frac{\sqrt{15}}{7} a \quad 
        \]
        代入 $g(a,b) = 0$得
        \[
        32a^2 - 224a + 147 = 0 \Rightarrow a = \frac{28 + 7\sqrt{10}}{8}
        \]
        故
        \[
        |\textbf c| = \sqrt{f(a,b)} = \sqrt{a^2 + \frac{15}{49}a^2} = \frac{8}{7} \cdot \frac{28 + 7\sqrt{10}}{8} = 4 + \sqrt{10}
        \]
        \textcolor{red}{(非通解)}
    \end{solution}

    \question 已知 $\triangle ABC$ 为边长为 1 的正三角形,设 $BC$ 边上有 $n-1$ 个等分点,由 $B$ 点到 $C$ 点的顺序为$P_1, P_2, \dots, P_{n-1}$,且令 $B=P_0, C=P_n$。若
    \[
    S_n = \overrightarrow{AB} \cdot \overrightarrow{AP_0} + \overrightarrow{AP_1} \cdot \overrightarrow{AP_2} + \overrightarrow{AP_3} \cdot \overrightarrow{AP_4} + \dots + \overrightarrow{AP_{n-1}} \cdot \overrightarrow{AC},
    \]
    求 $\displaystyle \lim_{n \to \infty} \frac{S_n}{n}$。
    \begin{solution}
        由 $B=P_0, C=P_n$ 可得
        \[
        |\overrightarrow{AP_0}| = |\overrightarrow{AB}| =1, \quad 
        \overrightarrow{P_0P_k} = \frac{k}{n} \overrightarrow{P_0P_n} = \frac{k}{n}\overrightarrow{BC}, \quad k=0,1,\dots,n
        \]
        则
        \[
        S_n = \sum_{k=0}^{n-1} \overrightarrow{AP_k} \cdot \overrightarrow{AP_{k+1}}
        = \sum_{k=0}^{n-1} (\overrightarrow{AP_0} + \overrightarrow{P_0P_k}) \cdot (\overrightarrow{AP_0} + \overrightarrow{P_0P_{k+1}})
        \]
        展开得
        \[
        S_n = \sum_{k=0}^{n-1} \Big( |\overrightarrow{AP_0}|^2 + \overrightarrow{AP_0} \cdot \overrightarrow{P_0P_{k+1}} + \overrightarrow{P_0P_k} \cdot \overrightarrow{AP_0} + \overrightarrow{P_0P_k} \cdot \overrightarrow{P_0P_{k+1}} \Big)
        \]
        代入 $\overrightarrow{P_0P_k} = \dfrac{k}{n} \overrightarrow{BC}$ 及 $\overrightarrow{AP_0} = \overrightarrow{AB}$得
        \[
        S_n = \sum_{k=0}^{n-1} \Big( 1 + \overrightarrow{AB} \cdot \frac{2k+1}{n} \overrightarrow{BC} + \frac{k(k+1)}{n^2} \Big)
        = \sum_{k=0}^{n-1} \Big( 1 + \frac{2k+1}{n}\left(-\frac{1}{2}\right) + \frac{k(k+1)}{n^2} \Big)
        \]
        化简得
        \[
        S_n = \sum_{k=0}^{n-1} \left(1 - \frac{2k+1}{2n} + \frac{k(k+1)}{n^2}\right)
        = n - \frac{n^2 - n +1}{2n} + \frac{n^2 -1}{3n}
        \]
        所以
        \[
        \lim_{n \to \infty} \frac{S_n}{n} = \lim_{n \to \infty} \left( 1 - \frac{n^2 - n +1}{2n^2} + \frac{n^2 -1}{3n^2} \right)
        = 1 - \frac{1}{2} + \frac{1}{3} = \frac{5}{6}
        \]
    \end{solution}

\end{questions}

\pagebreak

\begin{center}
  {\fontsize{30pt}{26pt}\selectfont
    \hypertarget{直角坐标}{直角坐标} \label{直角坐标}
  }
\end{center}
\separator
\vspace{1pt}

\begin{questions}
    \question $\triangle ABC$ 的三边 $BC, AC, AB$, 其中点坐标分别为 $(1,5),(-2,1),(4,3)$,试求 $A, B, C$ 的坐标。
    \begin{solution}
        设$A(a_1,a_2),\; B(b_1,b_2),\; C(c_1,c_2),$据题意有
        \[
        \begin{cases}
        \dfrac{b_1+c_1}{2}=1 \\[6pt] \dfrac{a_1+c_1}{2}=-2 \\[6pt] \dfrac{a_1+b_1}{2}=4
        \end{cases}
        \quad
        \text{及}
        \quad
        \begin{cases}
        \dfrac{b_2+c_2}{2}=5 \\[6pt] \dfrac{a_2+c_2}{2}=1 \\[6pt] \dfrac{a_2+b_2}{2}=3
        \end{cases}
        \]
        解得 $A(1,-1),\; B(7,7),\; C(-5,3)$
    \end{solution}

    \question 已知 $A(1,2),B(3,3)$, 线段 $AB$ 绕某点 $P$ 旋转后变成了线段 $A'B'$, 其中 $A'(3,1),B'(4,3)$。求 $P$ 的坐标。
    \begin{solution}
        发现旋转中心 $P$ 即线段 $AA'$ 和 $BB'$ 垂直平分线的交点。$A(1,2)$,$A'(3,1)$中点为$\left(2,\dfrac{3}{2}\right)$,于是垂直平分线为 
        \[
        y - \frac{3}{2} = 2(x - 2) \Rightarrow y = 2x - \frac{1}{2}
        \] 
        同理可得 $B(3,3),B'(4,3)$垂直平分线为
        \[
        x = \frac{7}{2}
        \] 
        联立方程解得 $P\left( \dfrac{7}{2}, \dfrac{9}{2} \right)$
    \end{solution}
    
    \question 直线 $y = mx + 2$ 与 $|x| + |y| = 1$ 相交, 求 $m$ 的范围。
    \begin{solution}
        即求方程组 
        \[
        \begin{cases}
        y = mx + 2 \\ |x| + |y| = 1
        \end{cases}
        \]
        有实数解时$m$的范围。发现,由作图知即求\[
        \begin{cases}
        y = mx + 2 \\ x + |y| = 1
        \end{cases}
        \]有实数解时$m$的范围。分情况再以判别式的性质可求得 $m \ge 2 , m \le -2$
    \end{solution}
          
    \question 设 $P_1(1,1),\; P_2(-2,-1)$, 直线 $x + y + 1 = 0$ 与 $P_1P_2$ 相交于点 $P$, 求 $\dfrac{P_1P}{PP_2}$。
    \begin{solution}
        是否在考虑作$P_1P_2$的直线方程式$\rightarrow$与直线 $x + y + 1 = 0$联立解得点$P\rightarrow$用距离公式？其实只需这样:由相似三角形,
        \[
        \begin{aligned}
        \dfrac{P_1P}{PP_2} 
        &= \frac{\left|\dfrac{1+1+1}{\sqrt{1^2+1^2}}\right|}{\left|\dfrac{-2-1+1}{\sqrt{1^2+1^2}}\right|}
        &= \dfrac{3}{2}
        \end{aligned}
        \]轻轻松松。
    \end{solution}
    
    \question 平面上有两定点 $A(1,4)$ 与 $B(5,2)$,若点 $P$ 在直线 $L: x+y-6=0$ 上移动,则 $|PA-PB|$ 的最大值是多少？
    \begin{solution}
        设 $P(t, 6 - t)$,则
        \begin{align*}
        f(t) = |PA - PB| 
        &= \left| \sqrt{(t - 1)^2 + (2 - t)^2} - \sqrt{(t - 5)^2 + (4 - t)^2} \right|\\
        &= \left| \frac{((t - 1)^2 + (2 - t)^2) -((t - 5)^2 + (4 - t)^2)}{
        \sqrt{(t - 1)^2 + (2 - t)^2} + \sqrt{(t - 5)^2 + (4 - t)^2}} \right|\\
        &= \left| \frac{12t - 36}{\sqrt{2t^2 - 6t + 5} + \sqrt{2t^2 - 18t + 41}} \right| \\
        &= \left| \frac{12 - \frac{36}{t}}{\sqrt{2 - \frac6t + \frac{5}{t^2}} + \sqrt{2 - \frac{18}{t} + \frac{41}{t^2}}} \right|
        \end{align*}
        故
        \[
        \lim_{t \to \pm\infty} f(t) = \frac{12}{2\sqrt{2}} = 3\sqrt{2}
        \]
    \end{solution}

    \question 已知坐标平面上三条直线 $L, L_1, L_2$, 若 $L$ 为水平线, $L_1, L_2$ 的斜率分别为 $\dfrac{2}{3}$ 与 $-\dfrac{3}{2}$, 且 $L$ 被 $L_1, L_2$ 截出的线段长度为 $26$,求三线所围成的三角形面积。
    \begin{solution}
        观察到$ m_{L_1}\cdot m_{L_2}=\dfrac{2}{3} \cdot \left(-\dfrac{3}{2}\right)=-1$,即$L_1 \perp L_2$,故欲求三角形的高。设$L_1, L_2$分别交$L$于点$(p,q),(p+26,q)$,解
        \[
        \begin{cases}
        y-q=\dfrac{2}{3}(x-p)\\ y-q=-\dfrac{3}{2}(x-p-26)
        \end{cases}
        \]
        可得$L_1, L_2$的交点为$(p+18,q+12)$。故三角形面积为$\dfrac{1}{2} \cdot 26 \cdot 12=156$。
    \end{solution}    
        
    \question 已知直线
    \[
    L_1: 3x + 3y = 2,\qquad
    L_2: 2x - 3y = 3,\qquad
    L_3: x - ay = -2
    \]
    将平面分成六个区域, 求$a$的所有可能值。
    \begin{solution}
        情况一:$L_1 \parallels L_3$,有 $a=-1$。

        情况二:$L_2 \parallels L_3$,有 $a=\dfrac{3}{2}$。

        情况三:$L_1,L_2$与$L_3$有唯一公共点,解$L_1,L_2$得 $\left(1,-\dfrac{1}{3}\right)$,故$a=-9$。

        经检验,$a$的所有可能值为 
        \[
        a = -1, \frac{3}{2}, -9
        \]
    \end{solution} 

    \question 若三直线
    \[
    L_1: 2x-y-3=0, \quad L_2: 2x+ay+3=0, \quad L_3: ax-y-6=0,
    \]
    不能围成一三角形, 求实数$a$之值。
    \begin{solution}
        情况一:$L_1 \parallels L_2$,有 $a=-1$。

        情况二:$L_1 \parallels L_3$,有 $a=2$。

        情况三:$L_2 \parallels L_3$,有 $a^2=-2$,此时无实数解。

        情况四:$L_1,L_2$与$L_3$有唯一公共点,联立$L_1,L_3$得
        \[
        x=\frac{3}{a-2},\quad y=\frac{12-3a}{a-2}, \quad a \neq 2
        \]
        代入$L_2$解得
        \[
        -3a(a-5)=0 \Rightarrow a=0 \ \text{或} \ a=5
        \]
        经检验,实数$a$之可能值为
        \[
        a=-1,0,2,5
        \]
    \end{solution}

    \question 三直线
    \[
    L_1: x - y + 2 = 0, \quad L_2: 2x + 3y + 9 = 0, \quad L_3: 8x + 3y - 27 = 0,
    \]
    围成三角形 \(\triangle ABC\),若点 \(P = (3,k)\) 在三角形内部,求 \(k\) 的范围。
    \begin{solution}
        欲使\(P = (3,k)\) 在三角形内,需满足\[
        \begin{cases}
            3-k+2>0 \\ 2(3)+3k+9>0\\8x+3y-27<0
        \end{cases}\quad \Rightarrow \quad 
        \begin{cases}
            k<5 \\ k>-5\\k<1 
        \end{cases}\quad \Rightarrow \quad 
        -5<k<1\]
    \end{solution}

    \question 已知点 $P(2,1)$，$Q(5,5)$，点 $R(x, 3x)$。若 $\triangle PQR$ 分别以 $P$ 或 $Q$ 为直角顶点，求点 $R$ 的坐标。
\begin{solution}
情况一：若 $\angle RPQ = 90^\circ$，则 $k_{PR} \cdot k_{PQ} = -1$。
已知 $k_{PQ} = \frac{5-1}{5-2} = \frac{4}{3}$，则 $k_{PR} = -\frac{3}{4}$：
\[ \frac{3x - 1}{x - 2} = -\frac{3}{4} \implies 4(3x - 1) = -3(x - 2) \]
\[ 12x - 4 = -3x + 6 \implies 15x = 10 \implies x = \frac{2}{3} \]
此时 $y = 3(\frac{2}{3}) = 2$，点 $R$ 坐标为 $(\frac{2}{3}, 2)$。

情况二：若 $\angle PQR = 90^\circ$，则 $k_{QR} \cdot k_{PQ} = -1$。
则 $k_{QR} = -\frac{3}{4}$：
\[ \frac{3x - 5}{x - 5} = -\frac{3}{4} \implies 4(3x - 5) = -3(x - 5) \]
\[ 12x - 20 = -3x + 15 \implies 15x = 35 \implies x = \frac{7}{3} \]
此时 $y = 3(\frac{7}{3}) = 7$，点 $R$ 坐标为 $(\frac{7}{3}, 7)$。
综上所述，所有可能的点 $R$ 坐标为 $(\frac{2}{3}, 2)$ 或 $(\frac{7}{3}, 7)$。
\end{solution}
    
    \question 设 \(x,y\) 为二元一次联立不等式图形上的任一点,已知
    \[
    \begin{cases}
    x + y \geq -1 \\
    x - y \geq -3 \\
    4x - y \leq 6
    \end{cases}
    \]
    函数 \(P = x + k y\) 在点 \((1,-2)\) 取得唯一最小值,求实数 \(k\) 的范围。
    \begin{solution}
        由作图得知求出可行解之顶点分别为 $(1,-2),(3,6),(-2,1)$, 据题意有
        \[
        \begin{cases}
            P(1,-2)<P(-2,1)\\P(1,-2)<P(3,6)
        \end{cases}\quad \Rightarrow \quad
        \begin{cases}
            1-2k<-2+k\\1-2k<3+6k
        \end{cases}\quad \Rightarrow \quad
            k>1
        \]
    \end{solution}
    
    \question 在坐标平面上,求二元一次联立不等式
    \[
    |x - 2y| \leq 2, \quad |x + 2y| \leq 2,
    \]
    的解所成区域面积。
    \begin{solution}
        由作图得知该区域为一长为$4$,宽为$2$的菱形,故面积为$8$。
    \end{solution}

    \question 已知 $k>0$,直线 $y = 3kx + 4k^2$ 与抛物线 $y = x^2$ 相交于两点 $P$ 和 $Q$。若 $O$ 是原点,且 $\triangle OPQ$ 的面积为 $80$,求该直线的斜率。
    \ifprintanswers
    \begin{figure}[H]
        \centering
        \includegraphics[width=0.5\linewidth]{images/image172.png}
    \end{figure}
    \fi
    \begin{solution}
        联立直线 $y = 3kx + 4k^2$ 与抛物线 $y = x^2$,解得
        \[
        P=(-k,k^2),Q=(4k,16k^2)
        \]
        设 $S,T$ 为 $P,Q$ 在$x$ 轴的垂足,则 
        \[
        SP=k^2,TQ=16k^2,ST=4k-(-k)=5k
        \]
        且由
        \begin{align*}
        80=[\triangle OPQ]
        &=[\text{梯形} PSTQ]-[\triangle PSO]-[\triangle QTO] \\
        &=\frac{1}{2}(k^2+16k^2)\cdot 5k-\frac{1}{2} \cdot k^2 \cdot k-\frac{1}{2} \cdot 16k^2 \cdot 4k=10k^3
        \end{align*}
        可得 $k = 2$,因此直线斜率为 $3k = 6$。
    \end{solution}

    \question 设 $P_1, P_2, P_3$ 为抛物线 $y=x^2$ 上的三点,$l_1, l_2, l_3$ 分别为抛物线在这些点处的切线。切线两两相交,设 $Q_{12}=l_1 \cap l_2, Q_{13}=l_1 \cap l_3, Q_{23}=l_2 \cap l_3$。求 $\triangle P_1P_2P_3$ 与 $\triangle Q_{12}Q_{13}Q_{23}$ 面积之比。
    \begin{solution}
        设 $P_1=(a,a^2), P_2=(b,b^2), P_3=(c,c^2)$ 且 $a \neq b \neq c$,切线方程为
        \[
        l_1: y = 2ax - a^2, \quad l_2: y = 2bx - b^2, \quad l_3: y = 2cx - c^2。
        \]
        可解得
        \[
        Q_{12} = \left(\frac{a+b}{2}, ab\right), \quad Q_{13} = \left(\frac{a+c}{2}, ac\right), \quad Q_{23} = \left(\frac{b+c}{2}, bc\right)。
        \]
        于是
        \[
        [\triangle P_1P_2P_3] = \frac{1}{2}
        \left|\begin{vmatrix}
        1 & 1 & 1 \\
        a & b & c \\
        a^2 & b^2 & c^2
        \end{vmatrix}\right|
        =\frac{1}{2}|(a-b)(b-c)(c-a)|
        \]
        \[
        [\triangle Q_{12}Q_{13}Q_{23}] = \frac{1}{2} \cdot \frac{1}{2}
        \left|\begin{vmatrix}
        1 & 1 & 1 \\
        a+b & a+c & b+c \\
        ab & ac & bc
        \end{vmatrix}\right|
        =\frac{1}{4}|(a-b)(b-c)(c-a)|
        \]
        故
        \[
        \frac{[\triangle P_1P_2P_3]}{[\triangle Q_{12}Q_{13}Q_{23}]} = 2
        \]
    \end{solution}

    \question 已知抛物线 $P_1: f(x) = x^2 + bx + c$ 的顶点为 $P$,抛物线 $P_2: g(x) = -x^2 + dx + e$ 的顶点为 $Q$。$P,Q$ 不重合,且都在$P_1,P_2$上。
    \begin{parts}
    \part 证明 $bd=2(e - c)$。
    \begin{solution}
        抛物线$P_1$顶点 $P$ 在 $x = -\dfrac{1}{2}b$上。由于 $P$ 在两条抛物线上,
        \[
        \frac{1}{4}b^2 + b\Bigl(-\frac{1}{2}b\Bigr) + c = -\frac{1}{4}b^2 + d\Bigl(-\frac{1}{2}b\Bigr) + e 
        \]
        化简得
        \[
        bd = 2(e - c)
        \]
        因此得证。
    \end{solution}
    \part 证明经过 $P$ 与 $Q$ 的直线斜率为 $\dfrac{1}{2}(b+d),y$ 截距为 $\dfrac{1}{2}(c+e)$。
    \begin{solution}
        顶点 $P,Q$ 坐标为 
        \[
        P\left(-\frac{1}{2}b, -\frac{1}{4}b^2 + c\right),Q\left(-\frac{1}{2}d, -\frac{1}{4}d^2 + e\right)
        \]  
        直线 $PQ$ 的斜率为
        \[
        \frac{(-\frac{1}{4}b^2 + c) - (-\frac{1}{4}d^2 + e)}{-\frac{1}{2}b - (-\frac{1}{2}d)} 
        = \frac{1}{2}(b + d)
        \]
        故直线方程为
        \[
        y = \frac{1}{2}(b+d)x + \frac{1}{2}(c+e)
        \]
        因此斜率为 $\dfrac{1}{2}(b+d),y$ 截距为 $\dfrac{1}{2}(c+e)$。
    \end{solution}
    \begin{solution}
        两抛物线方程相加得到 $2y = (b+d)x + (c+e)$,即
        \[
        y = \frac{1}{2}(b+d)x + \frac{1}{2}(c+e)。
        \]
        因此 $P$ 与 $Q$ 在同一条直线上,直线斜率为 $\dfrac{1}{2}(b+d),y$ 截距为 $\dfrac{1}{2}(c+e)$。
    \end{solution}
    \end{parts}

    \question 已知 $a > \dfrac{1}{2}$,抛物线方程为 $y = ax^2 + 2$,顶点为 $V$。抛物线与直线 $y = -x + 4a$ 相交于 $B$ 和 $C$ 两点,若 $\triangle VBC$ 的面积为 $\dfrac{72}{5}$,求 $a$。
    \ifprintanswers
    \begin{figure}[H]
        \centering
        \includegraphics[width=0.5\linewidth]{images/image187.png}
    \end{figure}
    \fi
    \begin{solution}
        抛物线 $y = ax^2 + 2$ 的顶点在 $x=0$上,因此 $V=(0,2)$。联立抛物线与直线解得
        \[
        x = -2 \text{ 或 } \ x = \frac{2a-1}{a}.
        \]
        因此
        \[
        B(-2,4a+2),\quad C\left(2-\frac{1}{a},4a-2 + \frac{1}{a}\right)
        \]
        设 $P,Q$ 在通过 $V$ 的水平线上,使得 $BP,CQ$ 垂直于 $PQ$,则
        \[
        P(-2,2),Q\left(2-\dfrac{1}{a},2\right),
        \]
        且
        \[
        BP = 4a, \quad CQ = 4a -4 + \frac{1}{a}, \quad PQ = 4 - \frac{1}{a}
        \]
        因此 $\triangle VBC$ 面积为
        \begin{align*}
        \frac{72}{5}=[\triangle VBC]
        &=[\text{梯形} PBCQ]-[\triangle BPV]-[\triangle CQV] \\
        &=\frac{1}{2}\left(4a+4a -4 + \frac{1}{a}\right)\left(4 - \frac{1}{a}\right)-
        \frac{1}{2}\cdot 4a \cdot 2 - \frac{1}{2}\left(4a -4 + \frac{1}{a}\right)\left(2-\frac{1}{a}\right)
        \end{align*}
        解方程可得
        \[
        a = \frac{5}{2} > \frac{1}{2}
        \]
    \end{solution}

    \question 一个等边三角形位于第一象限,顶点为 $(0,0),(x_1,4),(x_2,11)$,求序对 $(x_1,x_2)$。
    \begin{solution}
        构造复平面,复数$x_2+11i$逆时针旋转 $60^\circ$后变为复数$x_1+4i$,即
        \[
        x_2+11i=\left(\frac12+\frac{\sqrt3}{2}i\right)(x_1+4i)
        \]
        比较实部与虚部得
        \[
        x_2=\frac12 x_1-2\sqrt3,\quad 11=\frac{\sqrt3}{2}x_1+2
        \]
        解得
        \[
        (x_1,x_2)=(6\sqrt3,\sqrt3)
        \]
    \end{solution}

    \question 已知正方形 $WXYZ$ 的对角线 $WY$ 的斜率为 $2$,求直线 $WX$ 及 $XY$ 的斜率之和。
    \ifprintanswers
    \begin{figure}[H]
        \centering
        \includegraphics[width=0.35\linewidth]{images/image184.png}
    \end{figure}
    \fi
    \begin{solution}
        设 $WY$ 与水平线成夹角 $\theta$, 有$\tan\theta = 2$,因此
        \[
        m_{WX}+m_{XY}
        =m_{WX}+m_{WZ}
        =\tan(\theta+45^\circ)+\tan(\theta-45^\circ)
        =\frac{2+1}{1-2\cdot1}+\frac{2-1}{1+2\cdot1}
        =-\frac{8}{3}
        \]
    \end{solution}
    \ifprintanswers
    \begin{figure}[H]
        \centering
        \includegraphics[width=0.35\linewidth]{images/image185.png}
    \end{figure}
    \fi
    \begin{solution}
        将正方形平移使得 $W$为原点。设 $Y(2a,2b)$,有
        \[
        m_{WY}=\frac{2b-0}{2a-0}=2 \Rightarrow b=2a \Rightarrow Y(2a,4a)
        \]
        对角线 $WY$ 的中点为 $C(a,2a)$,由于$WC \perp CX$,
        \[
        m_{WC}=2 \Rightarrow m_{XC}=-\frac{1}{2}
        \]
        于是由图可得 $X(-a,3a)$,因此
        \[
        m_{WX}=\frac{3a-0}{-a-0}=-3.
        \]
        又因 $XY\perp WX$,所以 $m_{XY}=\dfrac{1}{3}$,斜率之和为
        \[
        -3+\frac{1}{3}=-\frac{8}{3}
        \]
    \end{solution}

    \question 矩形 $ABCD$ 的顶点为 $A(0,0),B(0,12),C(6,12),D(6,0)$。一直线经过点 $U(0,u),V(2,4),W(6,w)$,且将矩形 $ABCD$ 分成两个梯形,使得该两个梯形面积之比为 $5:3$,求 $U,W$。
    \ifprintanswers
    \begin{figure}[H]
        \centering
        \includegraphics[width=0.7\linewidth]{images/image161.png}
    \end{figure}
    \fi
    \begin{solution}
        矩形 $ABCD$ 面积为 $6 \cdot 12 = 72$,因此两个梯形面积分别为
        \[
        \frac{5}{8}\cdot 72 = 45,\quad \frac{3}{8}\cdot 72 = 27.
        \]
        \textbf{情况 1:直线交于 $AB$ 与 $CD$。}  
        直线 $UV$ 与 $VW$ 共线,有
        \[
        \frac{4-u}{2-0} = \frac{w-4}{6-2} \Rightarrow w = 12 - 2u.
        \]
        其中 $0<u,w<12$,梯形 $ADWU$ 的面积为
        \[
        \frac{1}{2}\cdot 6 \cdot (w+u) = 3(w+u)=3(12-u)
        \]
        若 $3(12-u)=27$,解得 $U(0,3), W(6,6)$; 
        若 $3(12-u)=45$,解得 $u=-3$,不合题意。

        \textbf{情况 2:直线交于 $AD$ 与 $BC$}  
        设交点为 $E(e,0), F(f,12)$,此时
        \[
        \frac{4}{2-e} = \frac{8}{f-2} \Rightarrow f = 6 - 2e.
        \]
        梯形 $BFEA$ 的面积为
        \[
        \frac{1}{2}\cdot 12 \cdot (f+e) = 6(f+e) = 6(6-e)
        \]
        若 $6(6-e)=27$,得 $e=\dfrac{3}{2}, f=3$,此时直线斜率为 $8$,由 $V(2,4)$ 得 $U(0,-12), W(6,36)$;  
        若 $6(6-e)=45$,得 $e=-\dfrac{3}{2}$ 不合题意。

        故解为
        \[
        U(0,3), W(6,6) \quad \text{或} \quad U(0,-12), W(6,36).
        \]
    \end{solution}

    \question 已知点 $J(2,7),K(5,3),L(r,t)$ 构成一个三角形,且其面积不超过 $10$。令 $\mathcal{R}$ 为所有满足条件的点 $L$ 构成的区域,且 $0 \le r,t \le 10$。求 $\mathcal{R}$ 的面积。
    \ifprintanswers
    \begin{figure}[H]
        \centering
        \includegraphics[width=0.5\linewidth]{images/image169.png}
    \end{figure}
    \fi
    \begin{solution}
        线段 $JK$ 的长度为
        \[
        JK = \sqrt{(2-5)^2 + (7-3)^2} = 5
        \]  
        以 $JK$ 为底,则 $\triangle JKL$ 的高度 $h$ 满足
        \[
        \frac{1}{2} \cdot 5 \cdot h \le 10 \Rightarrow h \le 4
        \] 
        直线 $JK$ 的方程为:
        \[
        \frac{y-7}{x-2} = \frac{3-7}{5-2}  \Rightarrow 4x+3y=29
        \]  
        设
        \[
        f(x,y)=\frac{|4x+3y-29|}{5},
        \]
        解$f(0,y)=4,f(x,0)=4,f(10,y)=4,f(x,10)=4$可知与直线$JK$距离$4$单位的两条平行线与正方形交于
        \[
        (0,3),\left(\dfrac{9}{4},0\right),(10,3),\left(\dfrac{19}{4},10\right),
        \]
        故$\mathcal{R}$面积为
        \[
        100 - \frac{1}{2}\cdot 3 \cdot \frac{9}{4} - \frac{1}{2}\cdot \frac{21}{4} \cdot 7 = \frac{313}{4}
        \]
    \end{solution}

    \question 已知坐标平面上有一正方形$ABCD$,点 $$E(6,0),F(6,6),G(0,8),H(-5,4)$$ 分别在 ${AB},{BC},{CD},{DA}$ 上,求正方形 $ABCD$ 的面积。
    \begin{solution}
        设直线
        \[
        AB: y = m(x - 6), \quad CD: y = m x + 8.
        \]
        且
        \[
        {BC}: y = -\frac{1}{m}(x - 6) + 6, \quad {AD}: y = -\frac{1}{m}(x + 5) + 4.
        \]
        由$d({AB}, {CD}) = d({BC}, {AD}),$
        \[
        \frac{|8 + 6m|}{\sqrt{m^2 + 1}} = \frac{|2m + 11|}{\sqrt{m^2 + 1}} 
        \]
        解得
        \[
        m = \frac{3}{4} \text{ 或 }\ m = -\frac{19}{8}.
        \]
        经检验,$\ m = -\dfrac{19}{8}$ 时点 $E$ 不在 ${AB}$ 上,舍去;因此正方形面积为$10^2=100$
    \end{solution}

    \question 已知函数 $f(x)=x^{2}-\sqrt{2}x$ 与 $g(x)=-x^{2}-1$ 有两条公切线,求四个切点组成的四边形周长。
    \begin{solution}
        设在 $f(x)=x^{2}-\sqrt{2}x$ 上的切点为 $A(a,a^2-\sqrt{2}a)$,切线斜率为 $f'(a)=2a-\sqrt{2}$,切线方程为
        \[
        y=(2a-\sqrt{2})(x-a)+a^2-\sqrt{2}a.
        \]
        与 $g(x)=-x^{2}-1$ 联立得
        \[
        x^2+(2a-\sqrt{2})x - a^2 +1 =0
        \]
        其中判别式为 $0$,
        \[
        (2a-\sqrt{2})^2 -4(1-a^2)=0 \Rightarrow a=\frac{\sqrt{2}\pm \sqrt{6}}{4}
        \]
        故$f(x)$上的切点为
        \[
        A\left(\frac{\sqrt{2}+\sqrt{6}}{4},-\frac{\sqrt{3}}{4}\right),\quad A'\left(\frac{\sqrt{2}-\sqrt{6}}{4},\frac{\sqrt{3}}{4}\right),
        \]
        $g(x)$ 上的切点为
        \[
        B\left(\frac{\sqrt{2}-\sqrt{6}}{4},-\frac{\sqrt{3}}{4}\right),\quad B'\left(\frac{\sqrt{2}+\sqrt{6}}{4},\frac{\sqrt{3}}{4}\right).
        \]
        因为$AA' = BB' = AB' = A'B = \dfrac{3}{2}$,周长为
        \[
        4\cdot \frac{3}{2} = 6
        \]
    \end{solution}

    \question 已知 $\triangle ABC$ 的外心 $O(-1,2)$,内心 $I(2,2)$,顶点 $A(2,8)$,求直线 $BC$ 的方程式。
    \begin{solution}
        由$O(-1,2), I(2,2), A(2,8)$,可得$\triangle ABC$ 外接圆 $\Gamma_1$ 方程式为
        \[
        (x+1)^2 + (y-2)^2 = 3^2+6^2 = 45
        \]
        过$I(2,2), A(2,8)$的直线方程式为 $x=2$,交 $\Gamma_1$于
        \[
        D = (2, -4)
        \]
        以 $D$ 为圆心,$DI = 6$ 为半径的圆 $\Gamma_2$ 方程式为
        \[
        (x-2)^2 + (y+4)^2 = 36.
        \]
        联立 $\Gamma_1$ 与 $\Gamma_2$,由\href{http://w.mathsgreat.com/SC/geom_th/geom_th_123.pdf}{鸡爪定理},得$BC$ 的方程式
        \[
        x - 2y = 4
        \]
    \end{solution}

    \question 设$\triangle ABC$中,已知$BC$与$x$轴平行,且$A(1,8)$,内切圆圆心为$(0,0)$,半径$4$,求$\triangle ABC$垂心$H$的坐标。
    \ifprintanswers
    \begin{figure}[H]
        \centering
        \includegraphics[width=0.5\linewidth]{images/image62.png}
    \end{figure}
    \fi
    \begin{solution}
        设过$A(1,8)$的圆的切线切点为$P(a,\sqrt{16 - a^2})$,  
        则  
        \[
        m_{PA}\cdot m_{OP}=-1 \Rightarrow \frac{\sqrt{16-a^2}}{a}=\frac{\sqrt{16-a^2}-8}{a-1} =-1 
        \]
        解得切点为
        \[
        a = \frac{48}{13}, -\frac{16}{5} \Rightarrow P\left(-\frac{16}{5}, \frac{12}{5}\right), Q\left(\frac{48}{13}, \frac{20}{13}\right)
        \]
        对圆方程求导知切线斜率$y' = -\dfrac{x}{y}$,则
        \[
        m_P = \frac{4}{3},\quad m_Q = -\frac{12}{5}
        \]
        故切线方程分别为  
        \[
        L_1: 4x - 3y + 20 = 0, \quad L_2: 12x + 5y = 52
        \]
        又 $B = L_1 \cap (y = -4) = (-8, -4), C = L_2 \cap (y = -4) = (6, -4)$,且过$A$的垂线为  
        \[
        L_3: x = 1
        \]
        过$B(-8, -4)$且与$L_2$垂直的直线方程为  
        \[
        L_4: y+4=\frac{5}{12} (x+8)\Rightarrow 5x - 12y = 8 
        \]
        联立 $L_3, L_4$得
        \[
        H\left(1, -\frac{1}{4}\right)
        \]
        \end{solution}

    \question 已知点 $A(3,1),B\left(\dfrac{5}{3},2\right)$ 是平行四边形 $ABCD$ 的顶点,且 $ABCD$ 四个顶点都落在函数 $$f(x)=\log_2 \dfrac{ax-b}{x-1}$$ 的图像上,求平行四边形 $ABCD$ 的面积。
    \begin{solution}
        由
        \[
        f(3) = \log_2 \frac{3a - b}{2} = 1, f\left(\frac{5}{3}\right) \log_2 \frac{5a/3 - b}{2/3} = 2
        \]
        解得
        \[
        a = 1, \; b = -1, f(x) = \log_2 \frac{x+1}{x-1}
        \]
        注意到
        \[
        f(-x) = \log_2 \frac{-x+1}{-x-1} = \log_2 \frac{x-1}{x+1} = -\log_2 \frac{x+1}{x-1} = -f(x),
        \]
        所以$f(x)$关于原点对称,平行四边形的另两个顶点为,
        \[
        C(-3,-1), \quad D\left(-\frac{5}{3},-2\right).
        \]
        故平行四边形 $ABCD$ 面积为
        \[
        \frac{1}{2}\left|6-\frac{5}{3}+6-\frac{5}{3}-\left(\frac{5}{3}-6+\frac{5}{3}-6\right) \right|=\frac{26}{3}
        \]
    \end{solution}

    \question 若方程 $$|x^2 - 4x + 3| - a = x$$ 恰有 $4$ 个实根,求实数 $a$ 的范围。
    \begin{solution}
        考虑
        \[
        \begin{cases}
        \Gamma: y = |x^2 - 4x + 3| \\
        L: y = x + a
        \end{cases}
        \]
        可知 $\Gamma$ 与 $x$ 轴交于 $A(1, 0),B(3, 0)$,过 $A$ 且与 $L$ 平行的直线为
        \[
        L_1: y = x - 1
        \]
        当 $1 \le x \le 3$ 时,有
        \[
        \Gamma: y = -x^2 + 4x - 3
        \Rightarrow y' = -2x + 4 = 1 \Rightarrow x = \frac{3}{2}
        \Rightarrow \text{切点 } P\left(\frac{3}{2}, \frac{3}{4}\right)
        \]
        过 $P$ 且与 $L$ 平行的直线为
        \[
        L_2: y = x - \frac{3}{4}
        \]
        因此当 $L$ 在 $L_1$ 与 $L_2$ 之间（不含切点）时,即当
        \[
        -1 < a < -\frac{3}{4}
        \]
        有恰好 $4$ 个交点。
    \end{solution}

    \question 如图,一抛物线顶点为 $D$,与 $x$ 轴交于 $A,C(4,0)$,与 $y$ 轴交于 $B(0,-4)$。若 $\triangle ABC$ 之面积为 $4$,试求 $\triangle DBC$ 之面积。
    \begin{figure}[H]
    \centering
    \includegraphics[width=0.45\textwidth]{images/image38.png}
    \end{figure}
    \begin{solution}
        假设 $O$ 为原点,则
        \[
        [\triangle ABC] = \frac{1}{2} \cdot OB \cdot AC = 2 \cdot AC = 4 \Rightarrow AC = 2, 
        \]
        故$A(2,0)$,于是$2,4$ 为抛物线 $y = f(x)$ 的两根,设
        \[
        f(x) = k(x - 2)(x - 4)
        \]
        可得
        \[
        f(0) = 8k = -4 \Rightarrow y = f(x) = -\frac{1}{2}(x - 2)(x - 4)
        \]
        于是
        \[
        f(3) = \frac{1}{2} \Rightarrow D\left(3, \frac{1}{2}\right)
        \]
        故
        \[
        [\triangle DBC] = \frac{1}{2} 
        \left|\begin{vmatrix} 
            3 & \frac{1}{2} & 1 \\ 0 & -4 & 1 \\ 4 & 0 & 1 
        \end{vmatrix}\right| = 3
        \]
    \end{solution}

    \question 已知$\triangle ABC$ 为直角三角形,其中 $\angle C=90^{\circ},AB=36$,而 $BC$ 上的中线为:$x+2y=0,AC$ 上的中线为:$x+y=0$,求 $\triangle ABC$ 的面积。
    \begin{solution}
        由 $A$ 在 $L_1: x+2y=0$ 上,$B$ 在 $L_2: x+y=0$ 上,设 
        \[
        A = (-2a, a),B = (b, -b),
        \]
        又$L_1,L_2$交于重心 $G(0,0)$,设 $C = (2a-b, -a+b)$,于是
        \[
        m_{CA}\cdot m_{CB}=\frac{-a+b-a}{2a-b+2a} \cdot \frac{-a+b+b}{2a-b-b}=-1 \Rightarrow 10a^2 - 15ab + 4b^2 = 0 \tag{1}
        \]
        由$AB = 36$,
        \[
        (-2a - b)^2 + (a + b)^2 = 5a^2 + 6ab + 2b^2 = 36^2 \tag{2}
        \]
        由 $(1),(2)$ 解得
        \[
        ab = \frac{2 \cdot 36^2}{27}
        \]
        故$\triangle ABC$ 的面积为
        \[
        [\triangle ABC] = \frac{1}{2} 
        \left|\begin{vmatrix} 
            -2a & a & 1 \\ b & -b & 1 \\ 2a - b & -a + b & 1 
        \end{vmatrix}\right|
        = \frac{3}{2} \cdot \frac{2 \cdot 36^2}{27} = 144
        \]
    \end{solution}

    \question
    \begin{parts}
    \part 已知两条抛物线 $y=x^{2}-8x+17$ 与 $y=-x^{2}+4x+7$。
    \begin{subparts}
    \subpart 求两条抛物线的顶点 $V_1$ 与 $V_2$ 的坐标。
    \begin{solution}
        完全平方得
        \[
        y=x^{2}-8x+17=(x-4)^{2}+1,
        \]
        \[
        y=-x^{2}+4x+7=-(x-2)^{2}+11.
        \]
        所以顶点为 $V_1(4,1), V_2(2,11)$。
    \end{solution}
    \subpart 若两条抛物线相交于 $P,Q$,证明四边形 $V_1PV_2Q$ 是平行四边形。
    \begin{solution}
        联立两抛物线得
        \[
        x^{2}-8x+17=-x^{2}+4x+7 \Rightarrow (x-5)(x-1)=0.
        \]
        得 $P(5,2),Q(1,10)$。发现$V_1V_2,PQ$ 的中点分别为 
        \[
        \left(\frac{4+2}{2},\frac{1+11}{2}\right)=(3,6), \quad
        \left(\frac{5+1}{2},\frac{2+10}{2}\right)=(3,6)
        \]
        即两条对角线互相平分,故得证 $V_1PV_2Q$ 是平行四边形。
    \end{solution}
    \end{subparts}
    \part 已知两条抛物线 $y=-x^{2}+bx+c$ 与 $y=x^{2}$,顶点分别为 $V_3,V_4$。当 $b,c$ 取某些值时,两条抛物线相交于相异点 $R,S$。
    \begin{subparts}
    \subpart 求$(b,c)$的解集,使得 $R,S$ 存在且 $V_3,V_4,R,S$ 为相异的四点。
    \begin{solution}
        完全平方得
        \[
        y=-\left(x-\frac{b}{2}\right)^{2}+\frac{b^{2}}{4}+c,
        \]
        顶点为 $V_3\left(\dfrac{b}{2},\dfrac{b^{2}}{4}+c\right)$,另一抛物线顶点为 $V_4(0,0)$,联立两抛物线得
        \[
        -x^{2}+bx+c=x^{2} \Rightarrow 2x^{2}-bx-c=0
        \]
        其中判别式需满足 $b^{2}+8c>0$,即 $c>-\dfrac{b^{2}}{8}$,且交点横坐标为
        \[
        x=\frac{b\pm \sqrt{b^{2}+8c}}{4}.
        \]
        需有 $c\ne0$,否则$R,S$与 $V_3$ 或 $V_4$ 重合;若 $b=0$,则 $V_3=(0,c)$,与 $V_4$ 重合需满足 $c=0$,所以解集为 
        \[
        \left\{\left.(b,c)\right|c>-\frac{b^{2}}{8}\ \text{且} \  c\ne0\right\}.
        \]
    \end{solution}
    \subpart 求使得 $R,S$ 存在且 $V_3,V_4,R,S$ 为相异的四点,且四边形 $V_3RV_4S$ 为矩形的所有 $(b,c)$。
    \begin{solution}
        已知 $R(x_1,x_1^{2}),S(x_2,x_2^{2})$,其中 $x_1,x_2$ 为方程 $2x^{2}-bx-c=0$ 的根,由韦达定理,
        \[
        x_1+x_2=\frac{b}{2}, \quad x_1x_2=-\frac{c}{2}
        \]  
        发现$V_3V_4,RS$ 的中点分别为 
        \[
        \left(\frac{b}{4},\frac{b^{2}+4c}{8}\right), \quad \left(\frac{x_1+x_2}{2},\frac{x_1^{2}+x_2^{2}}{2}\right)
        =\left(\frac{b}{4},\frac{b^{2}+4c}{8}\right)
        \]
        两条对角线中点相同,所以 $V_3RV_4S$ 是平行四边形,再要求一对相邻边垂直:  
        \[
        m_{V_4R}\cdot m_{V_4S} = x_1x_2=-1 =-\frac{c}{2} \Rightarrow c=2
        \]  
        结合 (i) 的条件,$c=2$ 时满足 $c>-\dfrac{b^{2}}{8}$ 且 $c\ne0$,因此所有解为 $(b,2)$,其中 $b \in \mathbb{R}$。
    \end{solution}
    \end{subparts}
    \end{parts}

    \question 坐标平面上,正六边形 $OABCDE$ 的顶点依逆时针排列为 $O$（原点）$,A,B,C,D,E$,其中直线 $OA$ 的方程为 $y=3x$,直线 $BE$ 的方程为 $y=3x+2$。试求此正六边形的外接圆圆心坐标。
        \ifprintanswers
        \begin{figure}[H]
        \centering
        \includegraphics[width=0.45\textwidth]{images/image24.png}
        \end{figure}
        \fi
    \begin{solution}
        两条平行线间的距离为
        \[
        {AP} = \frac{2}{\sqrt{10}}
        \]
        在$\triangle APB$中,
        \[
        {OA} = {AB} = \frac{\frac{2}{\sqrt{10}}}{\sin 60^\circ} = \frac{4}{\sqrt{30}}.
        \]
        设$A(a, 3a)$ 在 $y = 3x$ 上,则
        \[
        {OA}^2 = a^2 + (3a)^2 = 10a^2 = \left( \frac{4}{\sqrt{30}} \right)^2 = \frac{16}{30} \Rightarrow a = \frac{2}{5\sqrt{3}}\Rightarrow A\left( \frac{2}{5\sqrt{3}}, \frac{6}{5\sqrt{3}} \right)
        \]
        将 $A$ 逆时针旋转 $60^\circ$,得圆心$Q$
        \[
        \begin{bmatrix}
        \frac{1}{2} & -\frac{\sqrt{3}}{2} \\
        \frac{\sqrt{3}}{2} & \frac{1}{2}
        \end{bmatrix}
        \begin{bmatrix}
        \frac{2}{5\sqrt{3}} \\
        \frac{6}{5\sqrt{3}}
        \end{bmatrix}
        =
        \begin{bmatrix}
        \frac{\sqrt{3} - 9}{15} \\
        \frac{\sqrt{3} + 1}{5}
        \end{bmatrix}
        \Rightarrow Q\left( \frac{\sqrt{3} - 9}{15}, \frac{\sqrt{3} + 1}{5} \right).
        \]
    \end{solution}

    \question 已知函数 $y=f(x)$ 的图像既关于点 $(1,1)$ 中心对称, 又关于直线 $x+y=0$ 轴对称. 若 $x \in (0,1)$ 时, $f(x)=\log_{2}(x+1)$, 求 $f (\log_{2} 10)$ 的值。
    \begin{solution}
        记函数 $y=f(x)$ 的图像为$\Gamma$,
        对所有 $x_{0} \in (0,1)$, 令 $y_{0}=\log_{2}(1+x_{0})$, 则 $$(x_{0},y_{0}) \in \Gamma,y_{0} \in (0,1)$$
        利用 $\Gamma$ 的中心对称性与轴对称性, 可依次推得
        $$(2-x_{0}, 2-y_{0}) \in \Gamma, (y_{0}-2, x_{0}-2) \in \Gamma, (4-y_{0}, 4-x_{0}) \in \Gamma$$
        解$4-y_{0}=4-\log_{2}(1+x_{0})=\log_{2}10$得 $x_{0}=\dfrac{3}{5}$,因此 $$f(\log_{2}10)=f(4-y_{0})=4-x_{0}=4-\frac{3}{5}=\frac{17}{5}$$
    \end{solution}

    \question 若过原点可对 $y=x^{3}+ax^{2}+x+1$ 函数图形作三条切线,求 $a$ 的取值范围。
    \begin{solution}
        令切点 $P(t,t^3+at^2+t+1)$,则斜率为 $3t^2+2at+1$,  
        过 $P$ 之切线 
        \[
        L: y=(3t^2+2at+1)(x-t)+t^3+at^2+t+1 
        \]
        又 $L$ 过 $(0,0)$,故
        \[
        -3t^3-2at^2-t+t^3+at^2+t+1=0
        \]
        由$f(t)=2t^3+at^2-1=0$有三相异根,可得
        \[
        f'(t)=6t^2+2at=2t(3t+a)=0 \Rightarrow t=0,-\frac{a}{3} \text{有极值}
        \]
        由
        \[
        f(0)f\left(-\frac{a}{3}\right)= -1\left(-\frac{2a^3}{27}+\frac{a^3}{9}-1\right)< 0
        \]
        解得
        \[
        a > 3
        \]
    \end{solution}

    \question 在坐标平面上,已知$A(8,0),B(0,6)$, $P$ 为圆$x^{2}+y^{2}=16$上的动点,求$3PA + 2PB$的最小值。
    \begin{solution}
        设$P(x,y)$为圆$x^{2}+y^{2}=16$上的动点,$Q(0,a)$使得$PQ = \dfrac{2}{3} PB$,即
        \[
        x^{2} + (y - a)^{2} = \frac{4}{9}\left( x^{2} + (y - 6)^{2} \right)
        \]
        整理得
        \[
        5x^{2} + 5y^{2} - 18 a y + 48 y + 9 a^{2} - 144 = 0
        \]
        代入圆方程 $x^{2} + y^{2} = 16$,得
        \[
        9 a^{2} - 64 - 6 y (3 a - 8) = (3 a - 8)(3 a + 8 - 6 y) = 0
        \]
        解得
        \[
        a = \frac{8}{3} \Rightarrow Q \left( 0, \frac{8}{3} \right)
        \]
        因此
        \[
        3 PA + 2 PB = 3(PA + PQ) \ge 3 AQ = 3 \sqrt{8^{2} + \left(\frac{8}{3}\right)^{2}} =8 \sqrt{10}
        \]
    \end{solution}

    \question 在坐标平面上,将一个过原点且半径为 $r$ 的圆,完全放入 $y \ge x^{4}$ 的区域内, 此时 $r$ 的最大值为何？  
    \begin{solution}
        设圆的方程为 $x^2+(y-r)^2=r^2$,由 $y\ge x^4$ 得 $\sqrt{y}\ge x^2$,代入方程得
        \[
        \sqrt{y}+(y-r)^2\ge r^2,\quad (y>0).
        \]
        即
        \[
        2r \le y+\frac{1}{2\sqrt{y}}+\frac{1}{2\sqrt{y}}.
        \]
        由AM-GM不等式,
        \[
        y+\frac{1}{2\sqrt{y}}+\frac{1}{2\sqrt{y}}\ge \frac{3}{2}\sqrt[3]{2}
        \]
        所以
        \[
        r \le \frac{3}{4}\sqrt[3]{2}
        \]
    \end{solution}

    \question 坐标平面上,若四边形的四个顶点都在函数 $f(x)$ 上,则称此四边形为 $f(x)$ 的内接四边形。已知函数 $f(x) = x^3 + ax$ 的图形有唯一一个内接正方形,求 $a$ 的值。
        \ifprintanswers
        \begin{figure}[H]
        \centering
        \includegraphics[width=0.3\textwidth]{images/image21.png}
        \end{figure}
        \fi
    \begin{solution}
        由 $f(x) = x^3 + ax$ 可得:
        \[
        f''(x) = 6x = 0 
        \]
        即$f(x)$的对称中心为原点$(0,0)$,且为该内接正方形的中心,其对角线分别为
        \[
        L_1: y = mx, \quad L_2: y = -\frac{1}{m}x
        \]
        记 $A,B$ 在 $L_1,L_2$ 上,且都在$f(x)$上,则
        \[
        A: x^3 + ax = mx \Rightarrow x = 0 \text{ 或 } x^2 = m - a
        \]
        \[
        B: x^3 + ax = -\frac{1}{m}x \Rightarrow x = 0 \text{ 或 } x^2 = -\left(a + \frac{1}{m}\right)
        \]
        得
        \[
        A\left(\sqrt{m - a}, \; m\sqrt{m - a}\right), \quad B \left(\sqrt{-\left(a + \frac{1}{m}\right)},\; -\frac{1}{m} \sqrt{-\left(a + \frac{1}{m}\right)}\right)
        \]
        因 $B$ 为 $A$ 逆时针旋转 $90^\circ$ 得来,故
        \[
        \left(\sqrt{-\left(a + \frac{1}{m}\right)},\; -\frac{1}{m} \sqrt{-\left(a + \frac{1}{m}\right)}\right) = \left(-m\sqrt{m - a},\; \sqrt{m - a} \right)
        \]
        两边比较得:
        \[
        -m\sqrt{m - a} = \sqrt{-\left(a + \frac{1}{m}\right)}
        \Rightarrow m^2(m - a) = -\left(a + \frac{1}{m}\right)
        \]
        解得:
        \[
        a = \frac{m^3}{m^2 - 1} + \frac{1}{m(m^2 - 1)} = \frac{(m - \frac{1}{m})^2 + 2}{m - \frac{1}{m}}
        = m - \frac{1}{m} + \frac{2}{m - \frac{1}{m}}
        \]
        设 $t = m - \frac{1}{m}$,则:
        \[
        |a| = \left| t + \frac{2}{t} \right| \ge 2\sqrt{2}
        \Rightarrow a = -2\sqrt{2}
        \]
        等号成立当且仅当$f(x) = x^3 + ax$ 有唯一内接正方形。
    \end{solution}

    \question 设$P(0,a)$是$y$轴上异于原点的任意一点,过点$P$且平行于$x$轴的直线与曲线$y=\frac{1}{a}\ln x$交于点$Q$,曲线$y=\frac{1}{a}\ln x$在点$Q$处的切线交$y$轴于点$R$,则$\triangle PQR$的面积的最小值为\underline{\quad}.
    \begin{solution}
        $\frac{\sqrt{2e}}{2}$\textcolor{red}{(待解)}
    \end{solution}

    \question 已知曲线 $y=x^4+ax^3+ax^2+x+1$ 在点 $(0,1)$ 的切线,不只在 $(0,1)$ 与曲线相切,试求 $a$ 的值。
    \begin{solution}
        令 $f(x)=x^4+ax^3+ax^2+x+1$,则
        \[
        f'(x)=4x^3+3ax^2+2ax+1,\quad f'(0)=1.
        \]
        故过 $(0,1)$ 的切线斜率为 $1$,切线方程为 $L:\;y=x+1$。将 $L$ 代入曲线方程得
        \[
        x^4+ax^3+ax^2+x+1 = x+1
        \]
        即
        \[
        x^2(x^2+ax+a)=0
        \]
        题意要求切线不只在 $(0,1)$ 与曲线相切,因此 $x^2+ax+a=0$ 当中必须有重根,故判别式
        \[
        a^2-4a=a(a-4)=0
        \]
        当 $a=0$ 时,方程仅在 $x=0$ 有切点,不符题意；因此取
        \[
        a=4
        \]
    \end{solution}

    \question 设 $a\in\mathbb{R}$,若 $y=x^3-x$ 与 $y=x^2-a^2+a$ 有公切线,试求 $a$ 的范围。
    \begin{solution}
        令$\Gamma_1:f(x)=x^3-x, \Gamma_2:g(x)=x^2-a^2+a$,则
        \[
        f'(x)=3x^2-1,\quad g'(x)=2x.
        \]
        设公切线 $L=\overleftrightarrow{AB}$,其中$A\in\Gamma_1, B\in\Gamma_2$,记$A=(s,s^3-s), B=(t,t^2-a^2+a)$,公切线斜率为
        \[
        m_L=\frac{t^2-a^2+a-(s^3-s)}{t-s}=3s^2-1=2t
        \]
        消去$t=\dfrac{3s^2-1}{2}$,可得
        \[
        -9s^4+8s^3+6s^2-1=4a^2-4a
        \]
        欲存在实数 $s$ 使等式成立,必须有
        \[
        4a^2-4a \le \max_{s\in\mathbb{R}} h(s)
        \]
        求导得
        \[
        h'(s)=-36s^3+24s^2+12s=-12s(s-1)(3s+1),\quad h''(s)=-108s^2+48s+12
        \]
        由此临界点为 $s=0,1,-\dfrac{1}{3}$。由于$h''(0)> 0, h''(1)< 0, h''\left(-\dfrac{1}{3}\right)< 0,h(s)$的极大值为
        \[
        h(1)=4,\; h\left(-\frac{1}{3}\right)=-\frac{20}{27}
        \]
        因此 $h_{\max}=4$,故需
        \[
        4a^2-4a \le 4
        \]
        解不等式得
        \[
        \frac{1-\sqrt{5}}{2}\le a\le\frac{1+\sqrt{5}}{2}
        \]
    \end{solution}

    \question 在平面直角坐标系中,函数 $y=\dfrac{1}{|x|}$ 的图像为 $\Gamma$.设 $\Gamma$ 上的两点 $P,Q$ 满足: $P$ 在第一象限, $Q$ 在第二象限,且直线 $PQ$ 与 $\Gamma$ 位于第二象限的部分相切于点 $Q$.求 $|PQ|$ 的最小值.
    \begin{solution}
        当 $x<0$ 时, $y=-\dfrac{1}{x}$,导数为 
        \[
        y'=\frac{1}{x^{2}}
        \]
        设 $Q\left(-a,-\dfrac{1}{a}\right)$,其中 $a>0$,由条件知 $PQ$ 的斜率为 $y'|_{x=-a}=\dfrac{1}{a^{2}}$,故直线 $PQ$ 的方程为 
        \[
        y=\frac{1}{a^{2}}(x+a)+\frac{1}{a}=\frac{x+2a}{a^{2}}
        \]
        将上述方程与 $y=\dfrac{1}{x}(x>0)$ 联立,得 
        \[
        x^{2}+2ax-a^{2}=0\Rightarrow x_{p}=(\sqrt{2}-1)a
        \]
        于是
        \[
        |PQ|=\sqrt{1+\left(\frac{1}{a^{2}}\right)^{2}}\cdot|x_{p}-x_{Q}|=\sqrt{1+\frac{1}{a^{4}}}\cdot\sqrt{2}a
        \ge\sqrt{\frac{2}{a^2}}\cdot\sqrt{2}a=2
        \]
        当 $a=1$,即 $Q(-1,1)$ 时,$|PQ|$ 取最小值 $2$。
    \end{solution}

    \question 在 $\Gamma: y = x^3$ 上有一点 $P$,已知 $P$ 在第一象限且其 $x$ 坐标为 $a$。现以 $P$ 为切点作 $\Gamma$ 之切线 $L$,交 $y$ 轴于点 $Q$,且交 $\Gamma$ 于另一点 $S$,试求:
    \begin{parts}
    \part  ${PQ}:{QS}$  
    \begin{solution}
        已知曲线 $\Gamma: y = x^3$,故其导数为:
        \[
        y' = 3x^2
        \]
        在点 $P(a, a^3)$ 处的切线斜率为 $3a^2$,代入点斜式得切线 $L$:
        \[
        y = 3a^2(x - a) + a^3
        \]
        切线交 $y$ 轴于
        \[
        y = -3a^3 + a^3 = -2a^3 \Rightarrow Q(0, -2a^3)
        \]
        将切线方程代入 $\Gamma$ 得:
        \[
        x^3 = 3a^2(x - a) + a^3 \Rightarrow x^3 - 3a^2x + 2a^3 = 0\Rightarrow (x - a)^2(x + 2a) = 0
        \]
        另一个交点 $S$ 对应 $x = -2a$,代入 $\Gamma$ 得
        \[
        S = (-2a, -8a^3)
        \]
        于是
        \[
        {PQ} = \sqrt{(a - 0)^2 + (a^3 + 2a^3)^2} = \sqrt{a^2 + 9a^6}
        \]
        \[
        {QS} = \sqrt{(0 + 2a)^2 + (-2a^3 + 8a^3)^2} = \sqrt{4a^2 + 36a^6}
        \]
        故比值为
        \[
        {PQ}:{QS} = 1:2
        \]
    \end{solution}
    \part  $\Gamma$ 与切线 $L$ 所围成之封闭区域的面积,以 $a$ 表示。
    \begin{solution}
        所围面积为 $\Gamma$ 与切线之间的面积,即:
        \begin{align*}
        \int_{-2a}^{a} \left[x^3 - (3a^2(x - a) + a^3)\right] dx
        &= \int_{-2a}^{a} \left[x^3 - 3a^2x + 2a^3\right] dx\\
        &= \left[ \frac{1}{4}x^4 - \frac{3}{2}a^2x^2 + 2a^3x \right]_{-2a}^{a}
        = \frac{27}{4}a^4
        \end{align*}
    \end{solution}
    \end{parts}

    \question 在平面直角坐标系中,函数 $$y=\frac{x+1}{|x|+1}$$ 的图像上有三个不同的点位于直线 $l$ 上,且这三点的横坐标之和为0.求 $l$ 的斜率的取值范围.
    \begin{solution}
        当 $x\ge0$ 时, $y=1$ ;当 $x<0$ 时, $y=\frac{x+1}{1-x}$ 关于 $x$ 严格递增且小于1.
        设直线 $l:y=kx+b$ ,则条件等价于方程 
        \[
        kx+b=\frac{x+1}{|x|+1} \tag{1}
        \]
        有三个不同的实数解 $x_{1},x_{2},x_{3}(x_{1}<x_{2}<x_{3})$ ,满足 $$x_{1}+x_{2}+x_{3}=0$$
        首先有 $k\ne0$ (否则, $l$ 只能是 $y=1$ ,但此时 $l$ 与函数 $y=\frac{x+1}{|x|+1}$ 图像的任意三个公共点的横坐标之和必大于0,不合题意)
        
        当 $x<0$ 时,方程(1)可整理为 
        \[
        kx^{2}-(k-b-1)x+1-b=0  \tag{2}
        \]
        至多两个负数解;当 $x\ge0$ 时,方程(1)即为 \[kx+b=1 \tag{3}\]
        至多一个非负解,这表明方程(2)有两个不同的负数解 $x_{1},x_{2}$ ,其中 $$x_{1}+x_{2}=\frac{k-b-1}{k}$$方程(3)有非负解 $x_{3}=\dfrac{1-b}{k}$;由 $x_{1}+x_{2}+x_{3}=0$ ,可知 $k=2b$,进而有 $$x_{3}=\frac{1-b}{2b}$$ 由 $x_{3}\ge0$ 得 $0<b\le1$.
        方程(2)变为 $2bx^{2}+(1-b)x+1-b=0$ ,由判别式 $$(1-b)^{2}-4\cdot2b(1-b)=(1-b)(1-9b)>0,$$ 
        并结合 $0<b\le1$ ,可知 $0<b<\frac{1}{9}$;经检验,此时 $x_{1},x_{2}$ 确实为负数,符合题意,故 $l$ 的斜率 $k$ 的取值范围是 $$0<k<\frac{2}{9}.$$
    \end{solution}

    \question 直线 $y = \sqrt{3}x$ 上有一点 $A,x$ 轴上有一点 $B$, 圆 $(x - 12)^2 + (y - 5)^2 = 4$ 上有一点 $C$。试求 $\triangle ABC$ 的最小周长。 
    \ifprintanswers
    \begin{figure}[H]
        \centering
        \includegraphics[width=0.5\linewidth]{images/image55.png}
    \end{figure}
    \fi
    \begin{solution}
        设 $O(0,0)$, 圆心 $P(12,5)$,当$OP$ 与圆的交点为 $C$时$\triangle ABC$ 周长为最小,此时 
        \[
        OC = OP - 2 = 11
        \]  
        设点 $C$ 关于直线 $y = \sqrt{3}x$ 的对称点为 $C'$, 关于 $x$ 轴的对称点为 $C''$,则线段 $C'C''$ 与直线 $y = \sqrt{3}x$ 的交点为 $A$, 与 $x$ 轴的交点为 $B$,因此$AC = AC', BC = BC''$,故三角形周长为
        \[
        AC' + AB + BC'' = C'C''
        \]
        设$\angle AOC = \theta_1, \angle COB = \theta_2,$直线 $y = \sqrt{3}x$ 与 $x$ 轴夹角为 $60^\circ$, 即
        \[
        \theta_1 + \theta_2 = 60^\circ
        \]
        又因$C, C''$ 关于 $x$轴对称, $C, C'$  关于直线$y=\sqrt{3}x$ 对称,所以
        \[
        \angle BOC'' = \theta_1,\angle AOC' = \theta_2 \Rightarrow \angle C'OC'' = 2(\theta_1 + \theta_2) = 120^\circ 
        \]
        且
        \[
        OC'= OC'' = OC = 11
        \]
        在 $\triangle OC'C''$ 中,由余弦定理,
        \[
        C'C''^2=11^2 + 11^2-2 \cdot 11 \cdot 11\cdot \cos 120^\circ \Rightarrow C'C'' = 11\sqrt{3}
        \]
    \end{solution}


\end{questions}

\pagebreak

\begin{center}
  {\fontsize{30pt}{26pt}\selectfont
    \hypertarget{圆锥曲线}{圆锥曲线} \label{圆锥曲线}
  }
\end{center}
\separator
\vspace{1pt}

\begin{questions}
    \question 已知三圆方程为 $C_1: x^2+y^2+2ax+12y+10a+8=0,C_2: x^2+y^2-2x-a^2+2a=0,C_3: x^2+y^2-22x-6ay+8a^2-25a+36=0$,若三圆的圆心共线,求 $a$ 的值。
    \begin{solution}
        设$C_1,C_2,C_3$ 的圆心分别为$(-a, -6),(1, 0),(11, 3a)$,已知三圆心共线,所以斜率满足
        \[
        \frac{0 - (-6)}{1 - (-a)} = \frac{3a - 0}{11 - 1}
        \]
        解得
        \[
        a = -5 \quad \text{ 或 } \quad a = 4
        \]
        验证 $C_3$ 的半径,
        \[
        r_3^2 = 11^2 + (3a)^2 - (8a^2 - 25a + 36) = a^2 + 25a + 85
        \]
        当$a=4, r_3^2 = 201 > 0$,当$a=-5, r_3^2 = -15 < 0$,故可行解只有
        \[
        a = 4
        \]
    \end{solution}

    \question 若方程 $ x^2 + y^2 + 4k x - 6k y + 12k^2 - 4k - 8 = 0 $ 表示一个圆,求面积为最小时圆的方程。
    \begin{solution}
        将原式配方成
        \[
        (x + 2k)^2 + (y - 3k)^2 = k^2 + 4k + 8
        \]
        此为圆的标准方程,圆心为 $(-2k,\; 3k)$,半径平方
        \[
        r^2 = k^2 + 4k + 8 = (k+2)^2+4
        \]
        在 $k = -2$时最小, 因此圆的方程为
        \[
        (x - 4)^2 + (y + 6)^2 = 4
        \]
    \end{solution}
        
    \question 一个圆经过点 \( A(4, -2) \),且与 \(x\) 轴和 \(y\) 轴都相切,求该圆的方程式。
    \begin{solution}
        发现$A(4, -2)$在第四象限,已知圆与 \(x\) 轴和 \(y\) 轴都相切,所以圆只能落在第四象限。又因为圆心到两个坐标轴的距离都等于半径 \(r\),设圆心为 \((r,\; -r)\),半径为 \(r > 0\),方程为
        \[
        (x - r)^2 + (y + r)^2 = r^2
        \]
        将已知点 \(A(4, -2)\) 代入:
        \[
        (4 - r)^2 + (-2 + r)^2 = r^2
        \Rightarrow r^2 - 12r + 20 = 0
        \Rightarrow r = 2 \text{ 或 } 10
        \]
        因此圆方程为 \((x - 2)^2 + (y + 2)^2 = 4\)
        或\((x - 10)^2 + (y + 10)^2 = 100\)
    \end{solution}

    \question 设一个圆与直线 \(L_1: 3x - 4y + 7 = 0\)、\(L_2: 3x - 4y - 3 = 0\) 都相切,且圆心位于直线 \(L: x - 2y + 2 = 0\) 上,求该圆的方程式。  
    \begin{solution}
        首先,直线 \(L_1\) 和 \(L_2\) 是两条平行线,它们之间的距离为圆的直径:
        \[
        \text{直径} = \frac{|-3 - (7)|}{\sqrt{3^2 + (-4)^2}} = 2 \Rightarrow r = 1
        \]
        设圆心为 \((a, b)\),圆心到 \(L_1\) 或 \(L_2\) 的距离都等于 \(1\),即\[
        \begin{cases}
        \dfrac{|3a - 4b + 7|}{5} = 1 \Rightarrow  3a - 4b + 7 = \pm 5 \\
        \dfrac{|3a - 4b - 3|}{5} = 1 \Rightarrow  3a - 4b -3 = \pm 5 
        \end{cases}
        \]
        给出 $3a - 4b +2=0$,又圆心在直线 \(x - 2y + 2 = 0\) 上,即解\[
        \begin{cases}
        3a - 4b + 2 = 0 \\
        a - 2b + 2 = 0 
        \end{cases} \Rightarrow (a,b)=(2,2)
        \]     
        所以圆方程为
        \[
        (x - 2)^2 + (y - 2)^2 = 1
        \]
        \end{solution}

    \question 试求过点 \(P(3, 1)\) 且与圆 \((x - 1)^2 + (y + 2)^2 = 4\) 相切的直线方程。  
        \begin{solution}
        先判断点 \(P(3, 1)\) 是否在圆外:
        \[
        (3 - 1)^2 + (1 + 2)^2 = 13 > 4
        \Rightarrow P \text{ 在圆外}
        \]
        设切线斜率为 \(m\),切线方程为
        \[
        y - 1 = m(x - 3) \Rightarrow mx - y - 3m + 1 = 0
        \]
        圆心为 \((1, -2)\),半径 \(r = 2\),由圆心到切线距离为半径性质得
        \[
        \frac{|m \cdot 1 - (-2) - 3m - 1|}{\sqrt{m^2 + 1}} = 2
        \Rightarrow \frac{|2m - 3|}{\sqrt{m^2 + 1}} = 2
        \]
        两边平方得
        \[
        \frac{(2m - 3)^2}{m^2 + 1} = 4
        \Rightarrow m=\frac{5}{12}
        \]
        只解出一$m$,表示另一切线为铅直线 $x=3$,当 $m=\frac{5}{12}$ ,切线为 $5x-12y-3=0$
    \end{solution}
        
    \question 从点 \(P(3,5)\) 向圆\(C: x^2 + y^2 + 2x - 4y = 0\) 所作的两条切线切点分别为 \(A,B\),试求过 \(A,B,P\) 三点的圆方程。  
    \begin{solution}
        圆 \(C\) 的方程为
        \[
        x^2 + y^2 + 2x - 4y = 0 \Rightarrow (x+1)^2 + (y-2)^2 = 5
        \]
        其中圆心为 \(O(-1, 2)\),半径为 \(\sqrt{5}\)。
        
        发现过 \(A,B,P\) 三点的圆即是以 \(OP\) 为直径的圆(半圆含直角),该圆圆心为 \(OP\) 的中点:
        \[
       \left( \frac{3 + (-1)}{2}, \frac{5 + 2}{2} \right) = (1, \frac{7}{2})
        \]
        计算直径 \(|OP|\):
        \[
        |OP| = \sqrt{(3 + 1)^2 + (5 - 2)^2} = 5
        \Rightarrow \text{半径} = \frac{5}{2}
        \]
        因此所求圆的方程为
        \[
        (x - 1)^2 + \left(y - \frac{7}{2} \right)^2 = \frac{25}{4}
        \]
    \end{solution}

    \question 求两圆
    \[
    C_1: x^2+y^2-2y-4=0, \quad C_2: x^2+y^2-6x-11=0
    \]
    的公切线方程。
    \begin{solution}
        $C_1$圆心为 $A(0,1)$,半径为 $r_1 = \sqrt{5}$。$C_2$圆心为 $B(3,0)$,半径为 $r_2 = 2\sqrt{5}$。圆心距为
        \[
        AB = \sqrt{(3-0)^2 + (0-1)^2} = \sqrt{10}
        \]
        且
        \[
        r_1 + r_2 = 3\sqrt{5}, \quad r_2 - r_1 = \sqrt{5}
        \]
        由于 $r_2-r_1 < AB < r_1+r_2$,两圆相交于两点。设两条所求外公切线的交点为 $P(a,b)$,由相似三角形,$P$的$x,y$-坐标满足
        \[
        \frac{0-a}{3-a}=\frac{1-b}{0-b}=\frac{PA}{PB}=\frac{1}{2}
        \]
        解得$P(-3,2)$,现设切线方程为 
        \[
        L: y-2 = m(x+3) \Rightarrow mx - y + 3m + 2 = 0
        \]
        由于圆心$A(0,1)$至切线的距离等于半径,
        \[
        \frac{|m(0)-1+3m+2|}{\sqrt{m^2+1}} = \sqrt{5} 
        \]
        解得
        \[
        m = \frac{1}{2}, \quad m = -2
        \]
        故对应的公切线方程为
        \[
        y-2 = \frac{1}{2}(x+3) \Rightarrow x - 2y + 7 =0
        \]
        及
        \[
        y-2 = -2(x+3) \Rightarrow 2x + y +4 = 0
        \]
    \end{solution}

    \question
例 10、半径分别为 $1\text{ cm}、2\text{ cm}$ 及 $3\text{ cm}$ 的三个圆互相外切，如图所示。有一个小圆落在它们之间，且与它们都相切。若此小圆的半径为 $\frac{p}{q}\text{ cm}$，其中 $p$ 及 $q$ 为除了 $1$ 以外，并无其它公因数的两个正整数，求 $p+q$。

\begin{solution}
解：如图，建立坐标系得到三个圆方程：
$x^2 + y^2 = 1^2$，$x^2 + (y-3)^2 = 2^2$，$(x-4)^2 + y^2 = 3^2$，
设中间小圆圆心为 $(m, n)$，半径为 $r$。
利用相切性质，圆心距离 $=$ 半径之和：
\[ \sqrt{m^2 + n^2} = r + 1, \quad \sqrt{m^2 + (n-3)^2} = r + 2, \quad \sqrt{(m-4)^2 + n^2} = r + 3 \]
得
\begin{align*}
m^2 + n^2 &= r^2 + 2r + 1 \quad \dots (1) \\
m^2 + (n-3)^2 &= r^2 + 4r + 4 \quad \dots (2) \\
(m-4)^2 + n^2 &= r^2 + 6r + 9 \quad \dots (3)
\end{align*}
$(2) - (1)$ 得到 $n = 1 - \frac{2}{3}r$，
$(3) - (1)$ 得到 $m = 1 - \frac{r}{2}$。
代回 $(1)$：
\[ \left( 1 - \frac{r}{2} \right)^2 + \left( 1 - \frac{2}{3}r \right)^2 = r^2 + 2r + 1 \Rightarrow 25r^2 - 156r + 36 = 0 \]
$r = 6, \frac{6}{25}$。
根据条件 $6 + 25 = 31$。
\end{solution}

    \question 已知有一圆方程式为 $x^2+y^2=37$, 且圆的内部有一点 $P(1,2)$。
    \begin{parts}
    \part 若 $P$ 点为圆的某弦的中点, 试求此弦所在的直线方程式。
    \begin{solution}
        圆方程为 $x^2+y^2=37$,圆心为 $O(0,0)$,半径为 $\sqrt{37}$,点 $P(1,2)$ 在圆内,$OP$ 的斜率为
        \[
        m=\frac{2-0}{1-0}=2
        \]
        弦过 $P$ 且垂直于 $OP$,故斜率为 $-\dfrac{1}{2}$,弦方程为
        \[
        y-2=-\frac{1}{2}(x-1) \Rightarrow x+2y=5
        \]
    \end{solution}
    \part 若 $P$ 点为圆的某弦的一个三等分点, 试求此弦所在的直线方程式。
    \ifprintanswers
    \begin{figure}[H]
        \centering
        \includegraphics[width=0.4\linewidth]{images/image227.png}
    \end{figure}
    \fi
    \begin{solution}
        设弦长为 $3a$,弦的三等分点之一为 $P$,弦中点为 $Q$,则
        \[
        PQ=\frac{a}{2},\quad OP=\sqrt{1^2+2^2}=5
        \]
        设弦所在直线为
        \[
        y=m(x-1)+2.
        \]
        由毕氏定理
        \[
        OQ^2=OA^2-AQ^2=OP^2-PQ^2 \Rightarrow 
        37-\left(\frac{3a}{2}\right)^2=5-\left(\frac{a}{2}\right)^2,
        \]
        解得$a=4$,故由
        \[
        OQ=\frac{|2-m|}{\sqrt{m^2+1}}=1
        \]
        解得$m=\dfrac{3}{4}$或$m=\infty$,于是弦方程为
        \[
        3x-4y+5=0,\quad \text{或}\quad x=1
        \]
    \end{solution}
    \end{parts}
    \question 设一条光线经过点 \(A(-4,5)\),经 \(x\) 轴反射后恰与圆 \((x - 2)^2 + (y - 2)^2 = 5\) 相切,求反射线的斜率。
    \begin{solution}
        设反射线斜率为 \(k\),反射线经过\(A(-4,5)\)关于$x$轴的对称点: \((-4,-5)\),方程式为
        \[
        y+5=k(x+4) \Rightarrow y = kx + 4k - 5
        \]
        圆 \((x - 2)^2 + (y - 2)^2 = 5\)圆心为 \(C(2,2)\),半径为 \(\sqrt{5}\),由圆心到切线的距离等于半径的性质,
        \[
        \frac{|k \cdot 2 - 1 \cdot 2 + (4k - 5)|}{\sqrt{k^2 + 1}} =\frac{|6k - 7|}{\sqrt{k^2 + 1}} = \sqrt{5}
        \]
        即
        \[
        (6k - 7)^2 = 5(k^2 + 1)
        \Rightarrow 31k^2 - 84k + 44 = 0 \Rightarrow (k-2)(31k-22)=0
        \]
        得
        \[
        k = 2 \; \text{或}\; k = \frac{22}{31}\;\text{(不合题意)}
        \]
        $\therefore \;k=2$

        \textbf{又解}:用判别式$=0$求$k$
    \end{solution}

    \question 坐标平面上一圆被直线 \(x - y = 1\) 和 \(x - y = 5\) 所截的弦长均为 14,求该圆的面积。 
        \begin{solution}
        设圆半径为 \(r\),两条直线\(x - y = 1\) 和 \(x - y = 5\)间距离为
        \[
        \frac{|5 - 1|}{\sqrt{1^2 + (-1)^2}} = \frac{4}{\sqrt{2}} = 2\sqrt{2}
        \]
        因两弦等长且平行,圆心到两条直线的距离相等,圆心到任一条直线距离为
        \[
        d =\sqrt{2}
        \]
        由毕氏定理,
        \[
        r^2=\left(\frac{14}{2}\right)^2 + (\sqrt2)^2= 51
        \]
        故圆面积为$51\pi$。
    \end{solution}

    \question 若直线 \(2x - y + a = 0\) 与圆 \(x^2 + y^2 - 6x + by + c = 0\) 在点 \((4,1)\) 相切,求$a,b,c$的值。  
    \begin{solution}
        圆心 $\left(3,\ -\dfrac{b}{2}\right),$ 半径 $r = \sqrt{9 + \dfrac{b^2}{4} - c}$,切点 \((4,1)\)在直线与圆上,解得
        \[
        2(4) - 1 + a = 0 \Rightarrow a=-7
        \]
        \[
        4^2 + 1^2 - 6 \cdot 4 + b \cdot 1 + c = 0  \Rightarrow b + c = 7 \quad \tag{1}
        \]
        圆心到切线的距离等于半径:
        \[
        \frac{|2 \cdot 3 - (-\frac{b}{2}) - 7|}{\sqrt{5}}  = \sqrt{9 + \frac{b^2}{4} - c}
        \Rightarrow
        \left( \frac{b}{2} - 1 \right)^2 = 5\left(9 + \frac{b^2}{4} - c\right)
        \]
        将 (1) 代入得    
        \[
        (b-2)^2=5(8+b^2+4b)
        \Rightarrow b = -3, c=-7
        \]
        故
        \[
        a=-7,b=-3,c=10
        \]
    \end{solution}

    \question 设点 \(O_1\) 为圆 \(C: x^2 + y^2 - 6x + 4y + 9 = 0\) 的圆心。今以另一点 \(O_2\) 为圆心,\(\;O_1O_2\) 为半径作一圆,且该圆与圆 \(C\) 交于 \(A,B\) 两点。若 \(AO_2 = 3\),求$AB$的长度。 
    \begin{solution}
        已知圆 \(C: x^2 + y^2 - 6x + 4y + 9 = 0\),圆心 \(O_1 = (3, -2)\),半径 \(r_1 = 2\)。

        观察到 \(\triangle AO_1O_2\) 的三边为 \(AO_1 = 2, AO_2 = 3, O_1O_2 = 3\),用海伦公式得面积:
        \[
        S = \sqrt{4(4 - 2)(4 - 3)(4 - 3)} = 2\sqrt{2}
        \]
        又
        \[
        S = \frac{1}{2} \cdot O_1O_2 \cdot h = \frac{3h}{2} \Rightarrow h = \frac{4\sqrt{2}}{3}
        \]
        所以 \(AB = 2h = \dfrac{8\sqrt{2}}{3}\)。
    \end{solution}

    \question 平面上有一定点 $A(-3,3)$ 及一圆 $C:x^2+y^2-4x-4y+k=0$,若光源由 $A$ 点射出,碰到 $x$ 轴上 $P,Q$ 两点形成的两条反射光线恰好与圆 $C$ 相切,且 $PQ = \dfrac{7}{4}$,求 $k$ 之值。
    \ifprintanswers
    \begin{figure}[H]
        \centering
        \includegraphics[width=0.5\linewidth]{images/image46.png}
    \end{figure}
    \fi
    \begin{solution}
        圆 $C:(x-2)^2+(y-2)^2=8-k$ 圆心 $M(2,2)$,半径 $r=\sqrt{8-k},A(-3,3)$ 关于 $x$ 轴的对称点为 $A'(-3,-3)$,于是过$A',M$的直线方程式为
        \[
        x=y
        \]
        且此线为 $\angle PA'Q$ 的角平分线,故$L$ 交 $x$轴于 $O(0,0)$  
        现设 $PO = a,QO = \dfrac{7}{4} - a$,则
        \[
        P(-a,0),Q\left(-a+\frac{7}{4},0\right)
        \]
        根据反射光与入射光夹角相等,有
        \[
        \frac{A'P}{A'Q} = \frac{PO}{QO} 
        \Rightarrow \frac{\sqrt{(-3+a)^2+9}}{\sqrt{(a-\frac{19}{4})^2+9}} = \frac{a}{\frac{7}{4}-a}
        \]
        整理得
        \[
        4a^2 - 31a + 21 = 0 \Rightarrow a = \frac{3}{4} \quad (a=7\text{ 不合})
        \]
        则$P\left(-\dfrac{3}{4},0\right)$,过$A',P$的直线方程式为
        \[
        3y = 4x + 3
        \]
        由$ r = d(M, L_2)$,得
        \[
        1 = \sqrt{8-k} \Rightarrow k=7
        \]
    \end{solution}

    \question 在直角坐标平面上,已知圆 $C$ 的半径为 $4\sqrt{13}$ 且圆心在第三象限。从圆 $C$ 外一点 $P$ 对圆 $C$ 作两条切线,切点为 $A, B$,斜率分别为 $\dfrac{2}{3},\dfrac{3}{2}$。若 $\triangle PAB$ 的外接圆圆心为 $(6,4)$,求圆 $C$ 的圆心坐标。 
    \ifprintanswers
    \begin{figure}[H]
        \centering
        \includegraphics[width=0.5\linewidth]{images/image52.png}
    \end{figure}
    \fi
    \begin{solution}
        因为 $A,B$ 为切点,故
        \[
        \angle CAP = \angle CBP = 90^\circ,
        \]
        且外接圆圆心 $Q(6,4)$ 是线段 $CP$ 的中点,$CP$ 是外接圆直径。设圆心 $C(a,b)$,则点
        \[
        P = (12 - a, \; 8 - b).
        \]
        切线方程为
        \[
        L_1: y = \frac{2}{3} (x - (12 - a)) + (8 - b), \quad
        L_2: y = \frac{3}{2} (x - (12 - a)) + (8 - b).
        \]
        整理得
        \[
        L_1: 2x - 3y + 2a - 3b = 0, L_2: 3x - 2y + 3a - 2b - 20 = 0
        \]
        圆心 $C$ 到两切线距离均为半径 $4 \sqrt{13}$,
        \[
        \frac{|2a - 3b|}{\sqrt{13}} = \frac{|3a - 2b - 20|}{\sqrt{13}} = 4 \sqrt{13}
        \]
        由 $a,b < 0$,解得圆心为
        \[
        a = -20, b = -22 \Rightarrow (-20, -22)
        \]
    \end{solution}

    \question 已知圆 $x^{2}+y^{2}=37$ 内一点 $P(1,2)$,若 $P$ 点为某弦的三等分点,求此弦所在的直线方程。
    \ifprintanswers
    \begin{figure}[H]
        \centering
        \includegraphics[width=0.4\linewidth]{images/image106.jpg}
    \end{figure}
    \fi
    \begin{solution}
        设 $P$ 为弦 $AB$ 的三等分点,即$AP = k, PB = 2k$,圆心 $O$ 到 $P$ 的距离 
        \[
        OP = \sqrt{1^2+2^2} = \sqrt{5}
        \]
        令直线 $OP$ 与圆交于 $C, D$ 两点,由$\triangle APC \backsim \triangle DPB \ $(SAS),得
        \[
        \frac{AP}{PC} = \frac{DP}{PB} \Rightarrow \frac{k}{\sqrt{37}-\sqrt{5}} = \frac{\sqrt{37}+\sqrt{5}}{2k} \Rightarrow k = 4
        \]
        因此$AP = 4, PB = 8$,中点 $Q$ 为 $AB$ 的中点,则 
        \[
        PQ = 2,\quad OQ = 1
        \]
        设直线 $AB: y = m(x-1)+2$,由圆心到切线的距离等于半径的性质,解得
        \[
        OQ = 1 = \frac{|2 - m|}{\sqrt{m^2+1}} \Rightarrow m = \frac{3}{4} 
        \]
        或直线$AB$与$x$轴垂直,所以弦所在直线方程为
        \[
        3x - 4y + 5 = 0 \quad \text{或} \quad x = 1
        \]
    \end{solution}

    \question 由圆$x^2+y^2-4x+6y=7$外一点 $P$向该圆引两条切线,切于两点 $A(6,-1),B$。若点 $P$ 到圆心 $C$ 的距离为 $\sqrt{65}$,求 $B,P$ 的坐标。
    \begin{solution}
        圆$x^2+y^2-4x+6y=7$的圆心为$C(2,-3)$,半径为$r=\sqrt{20}$。设 $P(a,b)$,由于 $AP \perp AC$,
        \[
        \frac{-1-(-3)}{6-2}\cdot\frac{b+1}{a-6}=-1 \Rightarrow b=11-2a.
        \]
        由毕氏定理, $PA=\sqrt{65-20}=\sqrt{45}$,即
        \[
        (a-6)^2+(b+1)^2=45
        \]
        代入 $b=11-2a$解得
        \[
        (a-6)^2+(12-2a)^2=45 \Rightarrow a=3 \quad \text{或} \quad a=9
        \]
        故
        \[
        P(3,5)\quad \text{或}\quad P(9,-7).
        \]
        现求对应的 $B$,设 $B(k,h)$,由$BP=\sqrt{45},BC=\sqrt{20}$,当 $P(3,5)$,解
        \[
        \begin{cases}
        (k-3)^2+(h-5)^2=45,\\
        (k-2)^2+(h+3)^2=20.
        \end{cases}
        \]
        可得
        \[
        B\left(-\frac{18}{13},-\frac{1}{13}\right).
        \]
        当 $P(9,-7)$,解
        \[
        \begin{cases}
        (k-9)^2+(h+7)^2=45,\\
        (k-2)^2+(h+3)^2=20,
        \end{cases}
        \]
        得
        \[
        B\left(\frac{30}{13},-\frac{17}{13}\right).
        \]
        故所求为
        \[
        P(3,5),\; B\left(-\frac{18}{13},-\frac{1}{13}\right) \quad \text{或 } \quad P(9,-7),\; B\left(\frac{30}{13},-\frac{17}{13}\right)
        \]
    \end{solution}

    \question 已知两条平行直线 $L_1:y = 2x + 5,L_2:y = 2x - 1$是以点 $C$ 为圆心的圆的切线。若另一直线$L_3$经过点 $(9,0)$且垂直于 $L_1,L_2$,求圆心 $C$。
    \begin{solution}
        据题意,圆的半径为 $r = \dfrac{1}{2}\left|\dfrac{5+1}{\sqrt{2^2+(-1)^2}}\right|= \dfrac{3}{\sqrt{5}}$,且圆心 $C$ 位于中位线 $L:y = 2x + 2$ 上, $L_3$ 的方程为
        \[ 
        y - 0 = -\frac{1}{2}(x - 9) \Rightarrow y = \frac{1}{2}(9 - x) 
        \]
        联立$L,L_3$,
        \[ 
        2x + 2 = \frac{1}{2}(9 - x) \Rightarrow x = 1, y=4
        \]
        可知交点为 $M(1,4)$。设 $D(0,2)$ 为中位线在 $y$ 轴的截距,$L_3$与圆交于两点 $A, B$,考虑相似三角形 $\triangle AMC \backsim \triangle ADC \ \text{(AAA)}$,设$MC = x$,有
        \[
        \frac{MC}{AC} = \frac{AC}{DC} \Rightarrow \frac{x}{\frac{3}{\sqrt{5}}} = \frac{\frac{3}{\sqrt{5}}}{\sqrt{5} + x} ,
        \]
        其中$DM = \sqrt{1^2 + (4-2)^2} = \sqrt{5}$,舍去负根,解得
        \[ 
        x = \frac{-5\sqrt{5} + \sqrt{305}}{10} 
        \]
        由此得 
        \[
        DC = \sqrt{5} + \frac{-5\sqrt{5} + \sqrt{305}}{10} = \frac{5\sqrt{5} + \sqrt{305}}{10}
        \]
        设$CD$与水平轴的夹角为$\theta$,由已知可得$\tan \theta=2$,故圆心$C$的$x,y-$坐标为
        \[ 
        x = DC \cos \theta = \frac{5\sqrt{5} + \sqrt{305}}{10}\cdot \frac{1}{\sqrt{5}} = \frac{5 + \sqrt{61}}{10}
        \]
        \[ 
        y+2 = DC \sin \theta + 2 = \frac{5\sqrt{5} + \sqrt{305}}{10} \cdot \frac{2}{\sqrt{5}} + 2 = \frac{15 + \sqrt{61}}{5} 
        \]
        因此圆心 $C$ 的坐标为
        \[ 
        C\left( \frac{5 + \sqrt{61}}{10}, \frac{15 + \sqrt{61}}{5} \right) 
        \]
    \end{solution}

    \question 直线 $L$ 与圆 $C$ 的方程分别为
    \[
    L: y=\lambda(x-a)+a\sqrt{\lambda^{2}+1},\quad
    C: x^{2}+y^{2}=2ax,
    \]
    其中 $a$ 为正数,$\lambda$ 为参数。证明对任意 $\lambda$,直线 $L$ 恒为圆 $C$ 的切线。
    \begin{solution}
        将圆$C$化为标准形式,
        \[
        (x-a)^{2}+y^{2} = a^{2}
        \]
        可知圆心为 $P(a,0)$,半径为 $a$;直线 $L$ 的斜率为 $\lambda$,因此过点 $P(a,0)$ 且与 $L$ 垂直的直线斜率为 $-\dfrac{1}{\lambda}$,其方程为
        \[
        y-0 = -\frac{1}{\lambda}(x-a) \Rightarrow x = a-y\lambda
        \]
        设该垂线与直线 $L$ 交于点 $Q$,联立两直线方程解得
        \[
        Q\left(a-\frac{a\lambda}{\sqrt{\lambda^{2}+1}},\;\frac{a}{\sqrt{\lambda^{2}+1}}\right)
        \]
        于是圆心 $P(a,0)$ 到直线 $L$ 的距离为
        \[
        PQ=\sqrt{\left(a-\frac{a\lambda}{\sqrt{\lambda^{2}+1}}-a\right)^{2}+\left(\frac{a}{\sqrt{\lambda^{2}+1}}-0\right)^{2}} 
        =\sqrt{\frac{a^{2}(\lambda^{2}+1)}{\lambda^{2}+1}} 
        =a
        \]
        由于$PQ$与 $\lambda$ 无关,因此对任意 $\lambda$,直线 $L$ 始终与圆 $C$ 相切,故直线 $L$ 对所有 $\lambda$ 均为圆 $C$ 的切线。
    \end{solution}

    \question 已知圆 \((x - p)^2 + y^2 = r^2\) 及 \((y - p)^2 + x^2 = r^2\)交于相异点 \(A ,B\),且 \(a ,b\) 分别是 \(A ,B\) 的 \(x\)-坐标。 若\(C ,D\) 分别为两圆的圆心,
    \begin{parts}
    \part 证明 $a + b = p,a^2 + b^2 = r^2$且$p^2<2r^2$。
    \begin{solution}
        两圆方程式联立得
        \[
        (x - p)^2 + y^2 = (y - p)^2 + x^2
        \Rightarrow 2px = 2py  \Rightarrow x = y\; (\because p\neq0)
        \]
        代回第一式,
        \[
        (x - p)^2 + x^2 = r^2 \Rightarrow 2x^2 - 2px + p^2 - r^2 = 0
        \]
        由韦达定理,
        \[
        a + b = p,\; ab = \frac{p^2 - r^2}{2}
        \]
        又
        \[
        a^2 + b^2 = (a + b)^2 - 2ab = p^2 - 2 \cdot \frac{p^2 - r^2}{2} = r^2
        \]
        且判别式必为正,故\[
        \Delta=4p^2-4\cdot2(p^2-r^2)>0 \Rightarrow p^2<2r^2
        \]
    \end{solution}
    \part 若 \(r\) 是常量而 \(p\) 是变量,使得四边形 \(CADB\) 面积有最大值,求证此时 \(A\) 或 \(B\)是原点。
    \begin{solution}
        设 $A(a,a),B(b,b),C(p,0)$, $D(0,p)$。则直线 $CD$ 的斜率为 $-1$。点 $A,B$ 在直线 $y=x$ 上,所以 $AB$ 的斜率为 $1$。因此 $AB \perp CD$, 四边形 $CADB$ 是风筝形,其面积为
        \begin{align*}
        \frac{1}{2}\cdot AB \cdot CD 
        &= \frac{1}{2}\sqrt{2(a-b)^2}\sqrt{2p^2} \\
        &= \frac{1}{2}\sqrt{2r^2-(p^2-r^2)}\sqrt{2p^2} \\
        &= \sqrt{2r^2p^2-p^4}
        \end{align*}
        其中$f(p^2)=2r^2p^2-p^4$是一关于 $p^2$ 的二次函数,开口向下,其顶点在
        \[
        p^2=\frac{-2r^2}{2(-1)}=r^2
        \]
        因此面积在 $p^2=r^2$ 时取最大。又由 $(a+b)^2=p^2$, 所以
        \[
        a^2+2ab+b^2=r^2
        \]
        但 $a^2+b^2=r^2$, 故 $2ab=0$, 即 $a=0$ 或 $b=0$。

        若 $a=0$, 则 $A(0,0)$；若 $b=0$, 则 $B(0,0)$。所以当面积最大时,原点必为 $A$ 或 $B$。
    \end{solution}
    \end{parts}

    \question 求抛物线 \(y^{2}=16x\) 上距直线 \(4x-3y+24=0\) 最近的点坐标。  
    \begin{solution}
        设抛物线上一点为 \((t^2, 4t),\;t\in \mathbb{R}\),则到直线的距离为:
        \[
        d = \frac{|4t^2 - 12t + 24|}{\sqrt{4^2 + (-3)^2}} = \frac{4}{5}|t^2 - 3t + 6|
        \]
        又    
        \[
        t^2 - 3t + 6 = \left(t - \frac{3}{2}\right)^2 + \frac{15}{4}
        \]
        在 \(t = \dfrac{3}{2}\) 时最小,所以最近的点坐标为
        \[
        \left(t^2, 4t\right) = \left(\frac{9}{4}, 6\right)
        \]
    \end{solution}
    
    \question 已知抛物线的焦弦 \(AB\),准线到焦弦两端的垂足满足 \(AM=9,\ BN=4\),求 \(MN\)。
    \ifprintanswers
    \begin{figure}[H]
    \centering
    \includegraphics[width=0.2\textwidth]{images/image15.png}
    \end{figure}
    \fi
    \begin{solution}
        由抛物线定义,
        \[
        AB=AF+BF=AM+BN=9+4=13
        \]
        由毕氏定理,
        \[
        AB^2=MN^2+(AM-BN)^2 \Rightarrow MN^2=13^2-5^2=12^2
        \]
        $\therefore\; MN=12$
    \end{solution}
    
    \question 已知一个抛物线形状的拱桥,拱顶 A 点离水面 2 公尺时,水面宽度 \(BC = 4\) 公尺。若水面再下降 1 公尺,求新的水面宽度 \(DE\)。
    \begin{solution}
        设拱顶 A 为原点,抛物线开口向下,方程为:
        \[
        y = -a x^2,a>0
        \]
        当 \(y = -2\) 时,\(x = \pm 2\),代入得:
        \[
        -2 = -a(2)^2 \Rightarrow a = \frac{1}{2}
        \Rightarrow y = -\frac{1}{2}x^2
        \]
        当 \(y = -3\) 时,
        \[
        -3 = -\frac{1}{2}x^2 \Rightarrow x = \pm\sqrt{6}
        \Rightarrow DE = 2\sqrt{6}
        \]
    \end{solution}
    
    \question 求一边在直线 \(y=2x-17\) 上,另两顶点在抛物线 \(y=x^{2}\) 上的正方形面积。
    \begin{solution}
        设正方形顶点$A(x_1,x_1^2),B(x_2,x_2^2)$在抛物线 \(y=x^{2}\) 上,有
        \[
        m_{AB}=\dfrac{x_2^2-x_1^2}{x_2-x_1}=x_1+x_2=2
        \]
        且
        \[
        \left|\frac{2x_1-x_1^2-17}{\sqrt{2^2+1^2}}\right|=\sqrt{(x_2-x_1)^2+(x_2^2-x_1^2)^2}
        \]
        将$x_2=2-x_1$代入化简得\[
        (2x_1-x_1^2-17)^2=100(1-x_1)^2
        \]
        不失一般性,即解\[
        2x_1-x_1^2-17=10(1-x_1)
        \]
        得$x_1=3$或$x_1=9$,故正方形面积为$16$或$256$.
    \end{solution}

    \question 设一抛物线的顶点坐标为 $V(-1,3)$,且对称轴的方程为 $L: 2x+y-1=0$。若此抛物线经过点 $A(3,3)$,求此抛物线的正焦弦长。
    \begin{solution}
        直线 $L: 2x+y-1=0$ 与 $x$ 轴交于 $P\left(\dfrac{1}{2},0\right)$,以 $P$ 为旋转中心,逆时针旋转 $\tan^{-1} 2$ 将 $L$ 变为 $x$ 轴。

        顶点 $V(-1,3)$ 旋转后得到 $V'\left(\dfrac{1}{2}-\dfrac{15}{2\sqrt 5},0\right)$,点 $A(3,3)$ 旋转后得到 $A'\left(\dfrac{1}{2}-\dfrac{7}{2\sqrt 5},\dfrac{8}{\sqrt 5}\right)$。

        以 $V'$ 为顶点,旋转后的抛物线方程为
        \[
        y^2 = 4c\left(x - \frac{1}{2} + \frac{15}{2\sqrt 5}\right).
        \]
        将 $A'$ 代入方程,解得
        \[
        c = \frac{4\sqrt 5}{5}
        \]
        因此正焦弦长为
        \[
        4c = \frac{16\sqrt 5}{5}
        \]
    \end{solution}

    \question 已知抛物线顶点在原点,焦点在 $x$ 轴上,$\ \triangle ABC$ 三个顶点都在抛物线上,且重心为抛物线焦点 $F$,边 $BC$ 所在直线为 $4x+y-20=0$,求抛物线方程。
    \begin{solution}
        设抛物线方程为
        \[
        y^2 = 4c x
        \]
        $BC$ 与直线 $L: 4x+y-20=0$ 相交,由于直线上动点 $P(t,20-4t)$在抛物线上,
        \[
        (20-4t)^2 = 4 c t \Rightarrow 4 t^2 - (40+c) t + 100 =0.
        \]
        设两交点为 $B(t_1,20-4t_1),C(t_2,20-4t_2),A(x_a,y_a)$ 在抛物线上,且已知重心为焦点 $F(c,0)$,得
        \[
        c = \frac{x_a + t_1 + t_2}{3}, \quad 0 = \frac{y_a + (20-4t_1) + (20-4t_2)}{3}.
        \]
        解得
        \[
        x_a = \frac{11}{4}c -10, \quad y_a = c.
        \]
        又$A(x_a,y_a)$ 在抛物线上,
        \[
        c^2 = 4 c \left(\frac{11}{4}c - 10 \right) \ \Rightarrow c=4.
        \]
        因此抛物线方程为
        \[
        y^2 = 16 x
        \]
    \end{solution}

    \question 已知 $P(a, -a^2),Q(b, -b^2)$在抛物线$y = -x^2$上,其中$a < 0,b > 0$。设 $M$ 为线段 $PQ$ 的中点,$R$ 为过点 $M$ 的垂直直线与抛物线的交点, $l$ 为抛物线在点 $Q$ 处的切线。证明任何一端在 $PQ$ 上、另一端在 $l$ 上的垂直线段均被过 $Q,R$ 的直线平分。
    \begin{solution}
        切线 $l$ 的方程为
        \[ 
        y = -2bx + b^2 
        \]
        包含点 $P,Q$ 的直线方程为
        \[ 
        \frac{y+b^2}{x+a^2}=\frac{b+b^2}{a+a^2} \Rightarrow y = -(b+a)x + ab 
        \]
        $M$ 坐标为 $\left(\dfrac{a+b}{2},-\dfrac{a^2+b^2}{2}\right)$,故 $R$ 坐标为 $\left(\dfrac{a+b}{2}, -\left(\dfrac{a+b}{2}\right)^2\right)$,过
        $Q$ 及 $R$ 的直线方程为
        \[ 
        y = -\left(\frac{a+3b}{2}\right)x + \left(\frac{ab+b^2}{2}\right) 
        \]
        从线段 $PQ$ 到切线 $l$ 的垂直线段的中点 $y$ 坐标,是两条直线纵坐标的平均值
        \[ 
        y_{avg} = \frac{1}{2}(-2bx + b^2 - (b+a)x + ab) = -\left(\frac{a+3b}{2}\right)x + \left(\frac{ab+b^2}{2}\right)
        \]
        与经过 $Q$ 和 $R$ 的直线方程完全一致,故得证。
    \end{solution}

    \question 已知抛物线 $$y^2 = 2px$$ 与双曲线 $$y = -\dfrac{1}{x}$$ 相交于点 $R$,一条同时与抛物线和双曲线相切的公切线分别相切于点$S,T$,求证对于任意正实数 $p,\triangle RST$ 的面积与$p$无关,即为一定值。
    \begin{solution}
        设公切线为 $y=mx+c$,代入抛物线得
        \[
        (mx+c)^2 = 2px \Rightarrow m^2x^2 + 2(mc-p)x + c^2 = 0 \tag{1}
        \]
        令判别式为零:
        \[
        4(mc - p)^2 - 4m^2c^2 = 0 \Rightarrow p=2mc \tag{2}
        \]
        同理,将$y = mx + c$代入双曲线得
        \[
        mx^2 + cx + 1 = 0 \tag{3}
        \]
        且有
        \[
        c^2 = 4m \tag{4}
        \]
        由$(1),(3)$得
        \[
        x_{S}=\frac{p-mc}{m^2}, x_{T}=-\frac{c}{2m} 
        \]
        由$(2),(4)$得
        \[
        x_{S}=\frac{c}{m}=\frac{4}{c}, x_{T}=-\frac{2}{c} \Rightarrow y_{S}=2c, y_{T}=\frac{c}{2}
        \]
        即\[
        S\left(\frac{4}{c},2c\right),T\left(-\frac{2}{c},\frac{c}{2}\right)
        \]
        将抛物线与双曲线联立得
        \[
        \left(-\frac{1}{x}\right)^2 = 2px \Rightarrow
        R\left( \left( \frac{1}{2p} \right)^{\frac13},\ -\left(2p\right)^{\frac13} \right)
        \]
        由$\dfrac{(4)}{(3)}$得$c=(2p)^{\frac13}$,于是$R$坐标为
        \[
        R\left(\frac{1}{c},-c\right)
        \]
        故$\triangle RST$ 的面积为
        \[
        \frac{1}{2} \left| 
        \begin{vmatrix}
        \frac{4}{c} & 2c & 1 \\
        -\frac{2}{c} & \frac{c}{2} & 1 \\
        \frac{1}{c} & -c & 1
        \end{vmatrix}
        \right|=\frac{1}{2} \left|2+2+2-\frac12+4+4\right|=\frac{27}{4},
        \]
        且与 \(p\) 无关。
    \end{solution}

    \question 已知抛物线$\Gamma$的顶点是原点$O$,焦点是$F(0,1)$.过直线$y=-2$上任意一点$A$作抛物线$\Gamma$的两条切线,切点分别为$P$、$Q$,求证:
    \begin{parts}
    \part 直线$PQ$过定点;
    \begin{solution}
        易知抛物线$\Gamma$的方程为 $x^{2}=4y$.设点 $$A(t,-2),P(x_{1},y_{1}),Q(x_{2},y_{2})$$
        则过点$P$的抛物线$\Gamma$的切线 $l_{1}$ 的方程为 $$y-y_{1}=\frac{x_{1}}{2}(x-x_{1})$$
        即 $$x_{1}x-2y-2y_{1}=0$$
        同理,过点$Q$的抛物线$\Gamma$的切线 $l_{2}$ 的方程为: $$x_{2}x-2y-2y_{2}=0$$
        由 $l_{1}$, $l_{2}$ 过点$A$,可得 $$x_{1}t+4-2y_{1}=0, x_{2}t+4-2y_{2}=0$$
        这表明,点 $P(x_{1},y_{1}),Q(x_{2},y_{2})$的坐标满足方程: $$tx-2y+4=0$$
        即直线$PQ$的方程,易得直线$PQ$过定点$(0,2)$.
    \end{solution}
    \part $\angle PFQ=2\angle PAQ$.
    \begin{solution}
        不妨设$P$在$Q$左边,则 $x_{1}<x_{2}$.
        因为 $$\tan\angle PAQ=\frac{\frac{x_{1}}{2}-\frac{x_{2}}{2}}{1+\frac{x_{1}}{2}\cdot\frac{x_{2}}{2}}=\frac{2(x_{1}-x_{2})}{x_{1}x_{2}+4}$$
        所以 $$\tan 2\angle PAQ = \frac{2\tan\angle PAQ}{1-\tan^2\angle PAQ} = \frac{\frac{4(x_1-x_2)}{x_1x_2+4}}{1-\left(\frac{4(x_1-x_2)}{x_1x_2+4}\right)^2} = \frac{4(x_1-x_2)(x_1x_2+4)}{(x_1x_2+4)^2-4(x_1-x_2)^2}$$
        又因为
        $$\tan\angle PFQ = \frac{\frac{y_1-1}{x_1}-\frac{y_2-1}{x_2}}{1+\frac{y_1-1}{x_1}\cdot\frac{y_2-1}{x_2}} = \frac{x_2\left(\frac{x_1^2}{4}-1\right)-x_1\left(\frac{x_2^2}{4}-1\right)}{x_1x_2+\left(\frac{x_1^2}{4}-1\right)\left(\frac{x_2^2}{4}-1\right)} = \frac{4(x_1-x_2)(x_1x_2+4)}{(x_1x_2+4)^2-4(x_1-x_2)^2}$$
        故
        $$\tan 2\angle PAQ = \tan\angle PFQ$$
        又 $0 < \angle PAQ < 90^\circ < \angle PFQ < 180^\circ$,则
        $$\angle PFQ = 2\angle PAQ$$
        \end{solution}
    \end{parts}

    \question 坐标平面上有一圆 $C: x^2 + y^2 = 20$,试回答下列问题:
    \begin{parts}
    \part 抛物线 $\Gamma_1$ 与圆 $C$ 相切于 $A(-\sqrt{10}, \sqrt{10})$ 和 $B(\sqrt{10}, \sqrt{10})$ 两点,求抛物线 $\Gamma_1$ 的方程式为何？
    \begin{solution}
        由切点可知抛物线关于 $y$ 轴对称且开口向下,设
        \[
        \Gamma_1: y = ax^2 + b,\quad a < 0
        \]
        将$A(-\sqrt{10}, \sqrt{10})$ 代入得$10a + b = \sqrt{10}$;将 $\Gamma_1$ 代入圆 $C$:
        \[
        x^2 + (ax^2 + b)^2 = 20 \Rightarrow a^2x^4 + (2ab + 1)x^2 + b^2 - 20 = 0
        \]
        令判别式为 $0$,
        \[
        (2ab + 1)^2 - 4a^2(b^2 - 20) = 0
        \]
        代入 $b = \sqrt{10} - 10a$ 得
        \[
        80a^2 + 4a(\sqrt{10} - 10a) + 1 = 0 \Rightarrow 40a^2 + 4\sqrt{10}a + 1 = 0
        \]
        解得$\displaystyle a = -\frac{\sqrt{10}}{20}, b = \frac{3\sqrt{10}}{2}$,所以抛物线方程为
        \[
        y = -\frac{\sqrt{10}}{20}x^2 + \frac{3\sqrt{10}}{2}
        \]
        \end{solution}
        \part 抛物线 $\Gamma_2$ 与圆 $C$ 相切于 $D(2,4)$ 和 $E(-4,2)$ 两点,求抛物线 $\Gamma_2$ 的顶点坐标为何？
        \begin{solution}
        考虑将 $\Gamma_1$ 旋转使其与 $\Gamma_2$ 对应,设旋转矩阵为
        \[
        T = 
        \begin{bmatrix}
        \cos \theta & -\sin \theta \\
        \sin \theta & \cos \theta
        \end{bmatrix}
        \]
        满足
        \[
        T
        \begin{bmatrix}
        \sqrt{10} \\
        \sqrt{10}
        \end{bmatrix}
        =
        \begin{bmatrix}
        2 \\
        4
        \end{bmatrix}
        \Rightarrow
        \cos \theta = \frac{3}{\sqrt{10}} ,\sin \theta = \frac{1}{\sqrt{10}}
        \]
        又$\Gamma_1$ 顶点为 $P(0, \frac{3\sqrt{10}}{2})$,则 $\Gamma_2$ 顶点为
        \[
        T(P) = 
        \begin{bmatrix}
        \cos \theta & -\sin \theta \\
        \sin \theta & \cos \theta
        \end{bmatrix}
        \begin{bmatrix}
        0 \\
        \frac{3\sqrt{10}}{2}
        \end{bmatrix}
        =
        \begin{bmatrix}
        -\frac{3}{2} \\
        \frac{9}{2}
        \end{bmatrix}
        \Rightarrow
        \Gamma_2 \text{ 顶点 } = \left(-\frac{3}{2}, \frac{9}{2} \right)
        \]
    \end{solution}
    \end{parts}

    \question 在平面直角坐标系$xOy$中,点$P(0,1)$,点$A$为动点,以线段$AP$为直径的圆与$x$轴相切,设$A$点的轨迹为曲线$E$.
    \begin{parts}
    \part 求轨迹$E$的方程.
    \begin{solution}
        设$A(x,y),AP$中点$(\frac{x}{2},\frac{y+1}{2})$,则由题意$$\sqrt{x^2+(y-1)^2}=2\left|\frac{y+1}{2}\right|$$
        化简即得轨迹$E$的方程式$x^2=4y$。
    \end{solution}
    \part 若直线$AP$与$E$交于另一点$B,\triangle AOB$的外接圆交$E$于点$C$(不与$O,A,B$重合),过点$C$作$E$的切线交直线$AP$于点$N$,求$|ON|$的最小值.
    \begin{solution}
        设$l_{AP}: y=kx+1,A(x_1, y_1), B(x_2, y_2), C(x_3,y_3)$,
        $$\begin{cases} y=kx+1 \\ x^2=4y \end{cases} \Rightarrow x^3-4kx-4=0$$
        其中$$x_1+x_2=4k, x_1x_2=-4$$
        又$O,A,B,C$四点共圆,设$\triangle AOB$外接圆方程为$x^2+y^2+dx+ey=0$,解得
        $$\begin{cases} x^2+y^2+dx+ey=0 \\ x^2=4y \end{cases} \Rightarrow x^3+(16+4e)x+16d=0$$
        其中$x_1,x_2,x_3$是方程$x^3+(4e+16)x+16d=0$的三个根,且$$x_1+x_2+x_3=0$$
        即$$x_3=-(x_1+x_2)=-4k$$代入$E$得$y_3=4k^2$得$C(-4k,4k^2)$,过$C$曲线$E$的切线方程为 $$2kx+y+4k^2=0$$
        与$l_{AP}: y=kx+1$联立得
        $$N\left(-\frac{4k^2+1}{3k}, -\frac{4k^2-2}{3}\right)$$
        故
        \begin{align*}
        |ON|^2
        &=\left(-\frac{4k^2+1}{3k}\right)^2+\left(-\frac{4k^2-2}{3}\right)^2 \\
        &= \frac{1}{9}\left(16k^4+12+\frac{1}{k^2}\right)\\
        &\ge \dfrac{1}{9}\left(3\sqrt[3]{16k^4 \cdot \dfrac{1}{2k^2} \cdot \dfrac{1}{2k^2}}+12\right) \\
        &=\dfrac{\sqrt[3]{4}+4}{3}
        \end{align*}
        当且仅当$16k^4 = \dfrac{1}{2k^2}\iff k^6 = \dfrac{1}{32}\iff k=\pm\left(\dfrac{1}{2}\right)^{\frac{5}{6}}$时,$$|ON|_{\min}=\sqrt{\frac{\sqrt[3]{4}+4}{3}}$$
    \end{solution}
    \end{parts}

    \question 已知正$\triangle ABC$ 内接于抛物线$$y=x^2-\frac{71}{36},$$点$P$满足$\overrightarrow{PA}+\overrightarrow{PB}+\overrightarrow{PC}=\textbf{0}$,是否存在定点$Q$,使得点$P$到点$Q$的距离与点$P$到$x$轴的距离相等？若存在,求出点$Q$的坐标；若不存在,请说明理由。
    \begin{solution}
        设$A(x_1,x_1^2-\frac{71}{36})$, $B(x_2,x_2^2-\frac{71}{36})$, $C(x_3,x_3^2-\frac{71}{36})$,易知$x_1,x_2,x_3$互不相等,
    
        取$BC$的中点为$D$,则点$D$的坐标为$$D(\frac{x_2+x_3}{2},\frac{x_2^2+x_3^2}{2}-\frac{71}{36})$$
        因为$\triangle ABC$为正三角形,所以$AD \perp BC$,由$\overrightarrow{AD}\cdot\overrightarrow{BC}=0$,得$$(\frac{x_2+x_3}{2}-x_1,\frac{x_2^2+x_3^2}{2}-x_1^2)\cdot(x_3-x_2,x_3^2-x_2^2)=0$$
        整理得
        \[
        x_2^3+x_3^3+x_2^2x_3+x_2x_3^2-2x_1^2x_2-2x_1^2x_3+x_2+x_3-2x_1=0 \tag{1} 
        \]
        同理得
        \[
        x_1^3+x_2^3+x_1^2x_2+x_1x_2^2-2x_3^2x_1-2x_3^2x_2+x_1+x_2-2x_3=0 \tag{2} 
        \]
        由$(1)-(2)$,得$x_3^3-x_1^3+x_2^2(x_3-x_1)+3x_2(x_3^2-x_1^2)+2x_1x_3(x_3-x_1)+3(x_3-x_1)=0$
        整理得\[
        x_1^2+x_2^2+x_3^2+3(x_1x_2+x_1x_3+x_2x_3)+3=0 \tag{3}
        \]
        设 $P(x_{0},y_{0})$ ,因为 $\overrightarrow{PA}+\overrightarrow{PB}+\overrightarrow{PC}=\textbf{0}$ ,所以 $$x_{0}=\frac{x_{1}+x_{2}+x_{3}}{3},y_{0}=\frac{x_{1}^{2}+x_{2}^{2}+x_{3}^{2}}{3}-\frac{71}{36}$$
        因为 $$(x_{1}+x_{2}+x_{3})^{2}=x_{1}^{2}+x_{2}^{2}+x_{3}^{2}+2x_{1}x_{2}+2x_{1}x_{3}+2x_{2}x_{3}=9x_{0}^{2}$$
        $$x_{1}^{2}+x_{2}^{2}+x_{3}^{2}=3y_{0}+\frac{71}{12}$$
        所以 $$x_{1}x_{2}+x_{1}x_{3}+x_{2}x_{3}=\frac{1}{2}(9x_{0}^{2}-3y_{0}-\frac{71}{12})$$
        由(3)得 $$3y_{0}+\frac{71}{12}+\frac{3}{2}(9x_{0}^{2}-3y_{0}-\frac{71}{12})+3=0$$
        整理得 $$x_{0}^{2}=\frac{1}{9}(y_{0}-\frac{1}{36})$$
        所以点 $P(x_{0},y_{0})$ 是抛物线上一点,其焦点为 $(0,\frac{1}{18})$ ,准线为$x$轴,
        由抛物线的定义得点 $Q$ 的坐标为 $$Q(0,\frac{1}{18}).$$
    \end{solution}

    \question 抛物线 $C$ 的直角坐标方程为
    \[
    y^2 = 4ax,
    \]
    其中 $a>0$。$P(ap^2,2ap),Q(aq^2,2aq)$ 为抛物线 $C$ 上的相异点。过 $P,Q$ 作抛物线的切线,两切线相交于点 $R$。设 $S$ 为抛物线的焦点,证明
    \[
    SR^2 = SP \cdot SQ
    \]
    \begin{solution}
        由$y^2 = 4ax$,求导得
        \[
        \frac{dy}{dx}=\frac{2a}{y} \Rightarrow \left.\frac{dy}{dx}\right|_{x=ap^2,y=2ap} =\frac{1}{p}
        \]
        同理,在$P(ap^2,2ap),Q(aq^2,2aq)$切线方程为
        \[
        y-2ap=\frac{1}{p}(x-ap^2),\quad y-2aq=\frac{1}{q}(x-aq^2)
        \]
        联立方程解得
        \[
        x=apq,y=a(p+q) \Rightarrow R(apq,a(p+q))
        \]
        抛物线 $y^2=4ax$ 的焦点为$S(a,0)$,于是
        \[
        SP=\sqrt{(ap^2-a)^2+(2ap)^2}=\sqrt{a^2(p^2+1)^2}=a(p^2+1)
        \]
        同理
        \[
        SQ=a(q^2+1)
        \]
        而
        \[
        SR^2=(apq-a)^2+(a(p+q))^2 =a^2(p^2q^2+p^2+q^2+1) =a^2(p^2+1)(q^2+1)
        \]
        于是
        \[
        SR^2=[a(p^2+1)][a(q^2+1)]=SP \cdot SQ
        \]
        证毕。
    \end{solution}

    \question 抛物线 $P$ 的焦点为 $S(6,0)$,准线为 $x=0$。
    \begin{parts}
    \part 证明 $P$ 的直角坐标方程为 $y^2 = 12(x-3)$。
    \begin{solution}
        由抛物线定义,
        \[
        \sqrt{(x-6)^2 + y^2} = x
        \]
        化简得
        \[
        y^2 = 12(x-3)
        \]
    \end{solution}
    \part 验证 $P$ 的参数方程为
    \[
    x = 3t^2+3, \quad y = 6t
    \]
    \begin{solution}
        将参数方程 $x = 3t^2+3, y=6t$ 代入 
        \[
        y^2 = 36t^2, \quad 12(x-3) = 36t^2 \Rightarrow y^2 = 12(x-3)
        \]
        验证成立。
    \end{solution}
    \part 证明过点 $Q(3q^2+3,6q)$ 的切线方程为
    \[
    qy + 3 = x + 3q^2
    \]
    \begin{solution}
        由链导法,
        \[
        \frac{dy}{dx} = \frac{\frac{dy}{dt}}{\frac{dx}{dt}} = \frac{6}{6t} = \frac{1}{t}
        \]
        过 $Q(3q^2+3,6q)$的切线方程为
        \[
        y - 6q = \frac{1}{q}(x - (3q^2+3)) \Rightarrow qy + 3 = x + 3q^2
        \]
    \end{solution}
    \part 设 $R$ 为抛物线上一点,使 $QSR$ 在一条直线上,证明过 $Q$ 和 $R$ 的切线交于 $y$ 轴。
    \begin{solution}
        直线 $QSR$ 的方程为
        \[
        y - 0 = \frac{6 - 0}{3q^2+3 - 6}(x-6) \Rightarrow y= \frac{2q}{q^2-1}(x-6).
        \]
        与抛物线 $x = 3t^2+3, y = 6t$ 联立得
        \[
        6t = \frac{6q(t^2-1)}{q^2-1} \Rightarrow t = -\frac{1}{q}
        \]
        其中点 $Q$ 对应$t=q$,于是$R\left(\dfrac{3}{q^2}+3, -\dfrac{6}{q}\right)$,过$R$的切线方程为
        \[
        -\frac{1}{q}y + 3 = x + \frac{3}{q^2}
        \]
        联立
        \[
        qy + 3 = x + 3q^2,\quad -\frac{1}{q}y + 3 = x + \frac{3}{q^2}
        \]
        可得交点为
        \[
        (0,\frac{3(q^2-1)}{q})
        \]
        所以切线交于 $y$ 轴,证毕。
    \end{solution}
    \end{parts}

    \question 已知抛物线的参数方程为
    \[
    x=\frac{1}{3}t^2,\quad y=\frac{2}{3}t,\quad t\in\mathbb R,
    \]
    曲线在抛物线上一点 $P$ 的法线与抛物线相交于异与$P$的点 $Q$。求 $PQ$ 的最小值。
    \begin{solution}
        由参数方程消去参数 $t$, 得
        \[
        y^2=\frac{4}{3}x
        \]
        于是过$P\left(\dfrac{p^2}{3}, \dfrac{2p}{3}\right)$的切线斜率为
        \[
        \frac{dy}{dx}=\frac{2}{3y} \left.\frac{dy}{dx}\right|_{t=p}= \frac{1}{p}
        \]
        故法线斜率为 $-p$,过$P$法线方程为
        \[
        y - \frac{2p}{3} = -p\left(x - \frac{p^2}{3}\right) \Rightarrow 3y + 3px = 2p + p^3
        \]
        与 $x = \frac{3}{4}y^2$ 联立得,
        \[
        3y + 3p\left(\frac{3}{4}y^2\right) = 2p + p^3 \Rightarrow (3y - 2p)(3py + 4 + 2p^2) = 0
        \]
        由于点 $P$对应$y = \dfrac{2p}{3}$,点 $Q$ 的纵坐标为
        \[
        y_Q = -\frac{4 + 2p^2}{3p}
        \]
        此时由$y = \frac{2}{3}t$知$Q$处的参数为
        \[
        q = -p - \frac{2}{p}
        \]
        故
        \begin{align*}
        |PQ|^2 &= \left(\frac{q^2}{3} - \frac{p^2}{3}\right)^2 + \left(\frac{2q}{3} - \frac{2p}{3}\right)^2 \\
        &= \frac{1}{9}(q - p)^2 [(q + p)^2 + 4] \\
        &= \frac{1}{9}\left(-2p - \frac{2}{p}\right)^2 \left[\left(-\frac{2}{p}\right)^2 + 4\right] \\
        &= \frac{16(p^2 + 1)^3}{9p^4}
        \end{align*}
        设 
        \[
        f(p) = \frac{(p^2 + 1)^3}{p^4},
        \]
        求导得
        \[
        f'(p) = \frac{2(p^2 + 1)^2(p^2 - 2)}{p^5}
        \]
        令 $f'(p) = 0$,得 $p^2 = 2$。此时 $|PQ|^2$ 取得最小值$\dfrac{16(2 + 1)^3}{9 \cdot 2^2} = 12$, 因此$PQ$ 的最小值为 $\sqrt{12} = 2\sqrt{3}$。
    \end{solution}
    
    \question 已知抛物线
    \[
    C: y=2x^{2},
    \]
    直线 $y=kx+2$ 交 $C$ 于 $A,B$ 两点,$M$ 为线段 $AB$ 的中点,过 $M$ 作 $x$ 轴的垂线,垂足为 $N$。
    \begin{parts}
    \part 证明:抛物线 $C$ 在点 $N$ 处的切线与 $AB$ 平行；
    \begin{solution}
        设$A(x_1,2x_1^2),B(x_2,2x_2^2)$,将 $y=kx+2$ 代入 $y=2x^2$,得
        \[
        2x^2-kx-2=0
        \]
        由韦达定理,
        \[
        x_1+x_2=\frac{k}{2},\quad x_1x_2=-1
        \]
        由于 $M$ 为 $AB$ 的中点,
        \[
        x_N=x_M=\frac{x_1+x_2}{2}=\frac{k}{4}
        \]
        故
        \[
        N\left(\frac{k}{4},0\right)
        \]
        设抛物线在点 $N$ 处的切线为
        \[
        y-\frac{k^2}{8}=m\left(x-\frac{k}{4}\right)
        \]
        将 $y=2x^2$ 代入上式,得
        \[
        2x^2-mx+\frac{mk}{4}-\frac{k^2}{8}=0
        \]
        由于该直线与抛物线相切,判别式为零:
        \[
        \Delta=m^2-8\left(\frac{mk}{4}-\frac{k^2}{8}\right)=(m-k)^2=0
        \]
        解得 $m=k$,即切线与 $AB$ 平行。
    \end{solution}
    \begin{solution}
        设$A(x_1,2x_1^2), B(x_2,2x_2^2)$,由
        \[
        2x^2-kx-2=0
        \]
        得
        \[
        x_1+x_2=\frac{k}{2},\quad x_1x_2=-1
        \]
        因为
        \[
        x_N=\frac{x_1+x_2}{2}=\frac{k}{4}
        \]
        又 $N$ 在抛物线上,故
        \[
        N\left(\frac{k}{4},\frac{k^2}{8}\right)
        \]
        抛物线 $y=2x^2$ 的导数为 $y'=4x$,因此在点 $N$ 处的切线斜率为
        \[
        4\cdot \frac{k}{4}=k
        \]
        而 $AB$ 的斜率为 $k$,故切线与 $AB$ 平行。
    \end{solution}
    \part 是否存在实数 $k$ 使 $\overrightarrow{NA}\cdot\overrightarrow{NB}=0$？若存在,求 $k$ 的值；若不存在,说明理由。
    \begin{solution}
        假设存在实数 $k$ 使 $\overrightarrow{NA}\cdot\overrightarrow{NB}=0$,则 $NA\perp NB$,由于 $M$ 是 $AB$ 的中点,故
        \[
        MN=\frac{1}{2}AB
        \]
        又
        \[
        y_M=\frac{1}{2}(y_1+y_2)
        =\frac{1}{2}(kx_1+2+kx_2+2)
        =\frac{1}{2}\left[k(x_1+x_2)+4\right]
        =\frac{k^2}{4}+2
        \]
        因为 $MN\perp x$ 轴,
        \[
        MN=|y_M-y_N|
        =\left|\frac{k^2}{4}+2-\frac{k^2}{8}\right|
        =\frac{k^2+16}{8}.
        \]
        而
        \[
        AB=\sqrt{1+k^2}\,|x_1-x_2|
        =\sqrt{1+k^2}\sqrt{(x_1+x_2)^2-4x_1x_2}
        =\frac{1}{2}\sqrt{k^2+10}\sqrt{k^2+16}.
        \]
        由 $MN=\dfrac{1}{2}AB$,解得
        \[
        \frac{k^2+16}{8}=\frac{1}{4}\sqrt{k^2+10}\sqrt{k^2+16} \Rightarrow k=\pm2
        \]
        因此存在 $k=\pm2$,使 $\overrightarrow{NA}\cdot\overrightarrow{NB}=0$。
    \end{solution}
    \begin{solution}
        假设存在实数 $k$ 使 $\overrightarrow{NA}\cdot\overrightarrow{NB}=0$。由
        \[
        \overrightarrow{NA}=\left(x_1-\frac{k}{4},\,2x_1^2-\frac{k^2}{8}\right),
        \quad
        \overrightarrow{NB}=\left(x_2-\frac{k}{4},\,2x_2^2-\frac{k^2}{8}\right),
        \]
        得
        \begin{align*}
        \overrightarrow{NA}\cdot\overrightarrow{NB}
        &=\left(x_1-\frac{k}{4}\right)\left(x_2-\frac{k}{4}\right)
        +\left(2x_1^2-\frac{k^2}{8}\right)\left(2x_2^2-\frac{k^2}{8}\right) \\
        &=\left(x_1x_2-\frac{k}{4}(x_1+x_2)+\frac{k^2}{16}\right)
        \left(1+4x_1x_2+k(x_1+x_2)+\frac{k^2}{4}\right)
        \end{align*}
        代入
        \[
        x_1+x_2=\frac{k}{2},\quad x_1x_2=-1,
        \]
        得
        \[
        \overrightarrow{NA}\cdot\overrightarrow{NB}
        =\left(-1-\frac{k^2}{16}\right)\left(-3+\frac{3}{4}k^2\right).
        \]
        由于 $-1-\dfrac{k^2}{16}\neq0$,故
        \[
        -3+\frac{3}{4}k^2=0 \Rightarrow k=\pm2.
        \]
        因此存在实数 $k=\pm2$,使 $\overrightarrow{NA}\cdot\overrightarrow{NB}=0$。
    \end{solution}
    \end{parts}

    \question 已知
    \[
    \frac{(x+1)^{2}}{25}+\frac{(y-1)^{2}}{9}=1
    \]
    求
    \[
    \sqrt{(x+5)^{2}+(y-1)^{2}}+\sqrt{(x-3)^{2}+(y-1)^{2}}
    \]
    的最小值。
    \begin{solution}
        发现到椭圆方程
        \[
        \dfrac{(x+1)^{2}}{25}+\dfrac{(y-1)^{2}}{9}=1
        \]
        焦点为$F_1(-5,1),F_2(3,1)$,由椭圆定义,椭圆上一点$P(x,y)$满足
        \[
        \sqrt{(x+5)^2 + (y-1)^2} + \sqrt{(x-3)^2 + (y-1)^2} =PF_1 + PF_2 = 2\cdot 5= 10
        \]
    \end{solution}

    \question 试求二次曲线
    \[
    13x^{2} - 10xy + 13y^{2} - 6x - 42y - 27 = 0
    \]
    的正焦弦长。
    \begin{solution}
        将二次项写成矩阵形式,
        \[
        13x^{2} - 10xy + 13y^{2} = 
        \begin{bmatrix} x & y \end{bmatrix}
        \begin{bmatrix} 13 & -5 \\ -5 & 13 \end{bmatrix}
        \begin{bmatrix} x \\ y \end{bmatrix}
        \]
        对角化得
        \[
        A = \begin{bmatrix} 13 & -5 \\ -5 & 13 \end{bmatrix} =
        \begin{bmatrix} \frac{\sqrt{2}}{2} & -\frac{\sqrt{2}}{2} \\[6pt] \frac{\sqrt{2}}{2} & \frac{\sqrt{2}}{2} \end{bmatrix}
        \begin{bmatrix} 8 & 0 \\ 0 & 18 \end{bmatrix}
        \begin{bmatrix} \frac{\sqrt{2}}{2} & \frac{\sqrt{2}}{2} \\[6pt] -\frac{\sqrt{2}}{2} & \frac{\sqrt{2}}{2} \end{bmatrix}
        \]
        对线性项变换,
        \[
        [-6, -42]
        \begin{bmatrix} \frac{\sqrt{2}}{2} & -\frac{\sqrt{2}}{2} \\[6pt] \frac{\sqrt{2}}{2} & \frac{\sqrt{2}}{2} \end{bmatrix}
        = [-24\sqrt{2}, -18\sqrt{2}]
        \]
        将方程化为标准形式:
        \[
        8x'^2 + 18y'^2 - 24\sqrt{2} x' - 18\sqrt{2} y' - 27 = 0
        \]
        即
        \[
        \frac{\left(x' - \frac{3}{\sqrt{2}}\right)^2}{9} + \frac{\left(y' - \frac{1}{\sqrt{2}}\right)^2}{4} = 1
        \]
        由此得$a = 3, b = 2,$正焦弦长为
        \[
        \frac{2b^2}{a} = \frac{8}{3}
        \]
    \end{solution}
    
    \question 设\(A,B\)为椭圆 \(\dfrac{x^{2}}{25}+\dfrac{y^{2}}{16}=1\) 两焦点 ,若$P$为椭圆上任意点,求 \(\triangle PAB\) 面积最大值。  
    \begin{solution}
        椭圆焦点$A(-3,0),\;B(3,0),$
        设 $P(x,y)$ 在椭圆上,  
        \[
        \overrightarrow{PA} = (-3 - x,\ -y),\quad \overrightarrow{PB} = (3 - x,\ -y)
        \]
        面积:
        \[
        S = \frac{1}{2} \left| \overrightarrow{PA} \times \overrightarrow{PB} \right| 
        = \frac{1}{2} |(-3 - x)(-y) - (3 - x)(-y)| = 3|y|
        \]
        故当$|y| = 4,\;\triangle PAB$面积为最大:$12$
    \end{solution}

    \question 椭圆 \(\dfrac{x^{2}}{9}+\dfrac{y^{2}}{4}=1\) 焦点为 \(F_1,F_2\),若椭圆上点 \(P\) 满足 \(|PF_1|:|PF_2|=2:1\),求 \(\triangle PF_1F_2\) 面积。 
    \begin{solution}
        椭圆焦点为 \(F_1 = (-\sqrt{5}, 0), F_2 = (\sqrt{5}, 0)\),由椭圆定义
        \[
        |PF_1| + |PF_2| = 3|PF_2| = 2a = 6 \Rightarrow |PF_1| = 4,\; |PF_2| = 2
        \]
        又$|F_1F_2|=2\sqrt5$,由海伦公式,
        \[
        s = \frac{4 + 2 + 2\sqrt{5}}{2} = 3 + \sqrt{5}
        \]
        \[
        S = \sqrt{ (3+\sqrt{5})(-1+\sqrt{5})(1+\sqrt{5})(3-\sqrt{5})} = 4
        \]
    \end{solution}

    \question 求椭圆 
    \[
    \Gamma: \frac{x^2}{9} + \frac{y^2}{4} = 1
    \] 
    上任一切线在第一象限被 $x$ 轴、$y$ 轴截出之线段长的最小值。
    \begin{solution}
        由$\Gamma$隐微分得,
        \[
        \frac{2x}{9}+\frac{yy'}{2}=0 \Rightarrow y'=-\frac{4x}{9y}
        \]
        设切点为 $P(3\cos\theta,2\sin\theta)$, 则 
        \[
        y'(3\cos\theta,2\sin\theta)=-\frac{2\cos\theta}{3\sin\theta}=-\frac{2}{3}\cot\theta,
        \]
        切线 
        \[
        L: y=-\frac{2}{3}\cot\theta(x-3\cos\theta)+2\sin\theta
        \] 
        与坐标轴的交点
        \[
        A\left(0,\frac{2}{\sin\theta}\right), \quad B\left(\frac{3}{\cos\theta},0\right)
        \]
        由柯西不等式,
        \[
        \left[\left(\frac{2}{\sin\theta}\right)^2+\left(\frac{3}{\cos\theta}\right)^2\right](\sin^2\theta+\cos^2\theta) \ge (2+3)^2
        \]
        故
        \[
        AB^2 = \frac{4}{\sin^2\theta} + \frac{9}{\cos^2\theta} \ge 25 \Rightarrow AB_{\min}=5
        \]
    \end{solution}

    \question 已知 $a>b>0,F$ 是 
    \[
    \Gamma:\frac{x^2}{a^2}+\frac{y^2}{b^2}=1
    \] 
    的一个焦点,$AB$ 是 $\Gamma$ 的弦且 $F$ 在 $AB$ 上。试证明
    \[
    \frac{1}{FA}+\frac{1}{FB}=\frac{2a}{b^2}
    \]
    \begin{solution}
        假设两焦点为 $F,F'$ 且$FA=p, FB=q$,则
        \[
        F'A=2a-p,\quad F'B=2a-q
        \]
        由余弦定理,
        \[
        \cos \angle AFF' = \frac{p^2+(2c)^2-(2a-p)^2}{2p(2c)}, \quad
        \cos \angle BFF' = \frac{q^2+(2c)^2-(2a-q)^2}{2q(2c)}
        \]
        由于 $\angle AFF' + \angle BFF' = 180^\circ$,
        \[
        \frac{p^2+(2c)^2-(2a-p)^2}{2p(2c)}
        = -\frac{q^2+(2c)^2-(2a-q)^2}{2q(2c)}
        \]
        可化简得
        \[
        \frac{1}{FA}+\frac{1}{FB}=\frac{2a}{b^2}
        \]
    \end{solution}

    \question 已知椭圆 $C_1: \frac{x^2}{4} + y^2 = 1$, 椭圆 $C_2: (x-2)^2 + 4y^2 = 1$, $C_1, C_2$ 的公切线与 $x$ 轴交于点 $A$, 则点 $A$ 的坐标为 $(4,0)$
    \begin{solution}
        \textcolor{red}{(待解)}
    \end{solution}

    \question 求椭圆 $E: x^2 + 16y^2 = 16$ 与圆 $C: x^2 + y^2 = 4$ 的公切线方程。
    \begin{solution}
        设切线斜率为 $m$,椭圆与圆
        \[
        E: \frac{x^2}{16} + y^2 = 1, \quad C: x^2 + y^2 = 4
        \]
        上斜率为 $m$ 的切线分别为
        \[
        y = mx \pm \sqrt{16m^2 + 1},\quad y = mx \pm 2\sqrt{1+m^2}
        \]
        联立解得
        \[
        \sqrt{16m^2 + 1} = 2\sqrt{1+m^2} \Rightarrow m = \pm \frac{1}{2}
        \]
        当 $m = \dfrac{1}{2}$,
        \[
        y = \frac{1}{2}x \pm \sqrt{5} \Rightarrow x - 2y + 2\sqrt{5} = 0, \quad x - 2y - 2\sqrt{5} = 0
        \]
        当 $m = -\frac{1}{2}$,
        \[
        y = -\frac{1}{2}x \pm \sqrt{5} \Rightarrow x + 2y - 2\sqrt{5} = 0, \quad x + 2y + 2\sqrt{5} = 0
        \]
        所以四条公切线方程为
        \[
        x - 2y + 2\sqrt{5} = 0, \quad x - 2y - 2\sqrt{5} = 0, \quad x + 2y - 2\sqrt{5} = 0, \quad x + 2y + 2\sqrt{5} = 0
        \]
    \end{solution}

    \question 在椭圆$\Omega$中,$F_1, F_2$为焦点,$A$为长轴的一个端点,$B$为短轴的一个端点,若$\angle F_1BF_2 = \angle FAB$,求$\Omega$的离心率。
    \begin{solution}
        因为$\angle F_1BF_2 = \angle FAB$,所以
        \[
        \triangle BF_1F_2 \backsim \triangle ABF_1\Rightarrow\angle BF_1F_2 = \angle ABF_1\Rightarrow|AB| = |AF_1|,
        \]
        即
        \[
        \sqrt{a^2+b^2}=a+c\Rightarrow2e^2+2e-1=0
        \]
        解得$e=\dfrac{\sqrt{3}-1}{2}$(负解舍去)。
    \end{solution}
    
    \question 已知椭圆 \(\dfrac{x^2}{a^2} + \dfrac{y^2}{b^2} = 1\) 的左焦点为 \(F\),右顶点为 \(A\),上顶点为 \(B\),离心率为 \(\dfrac{\sqrt{5}-1}{2}\),求 \(\angle ABF\)。  
    \begin{solution}
        设椭圆 \(\dfrac{x^2}{a^2} + \dfrac{y^2}{b^2} = 1\)离心率为\(e=\dfrac{\sqrt{5}-1}{2}\),左焦点为 \(F(-ae,0)\)、右顶点为 \(A(a,0)\)、上顶点为 \(B(0,b)\),中心$O(0,0)$,发现\[
        \tan \angle ABO=\frac{a}{b},\;\tan \angle OBF=\frac{ae}{b}\]于是\[
         \tan \angle ABF 
        = \tan (\angle ABO+ \angle OBF) 
        = \frac{\frac{a}{b}+\frac{ae}{b}}{1-\frac{a^2}{b^2}\cdot e} 
        \]
        又$\dfrac{a}{b}=\dfrac{1}{\sqrt{(1-e^2)}}$,代入有\[
        \tan \angle ABF = -\frac{1}{e^2-e+1}\sqrt{\frac{e+1}{e-1}}
        \]
        发现$e = \dfrac{\sqrt{5}-1}{2}\Rightarrow e^2+e-1=0$,由此得\[
        \tan \angle ABF =\infty \Rightarrow \angle ABF = 90^\circ 
        \]
    \end{solution}
    
    \question 设直线 
    \[
    \frac{x}{5}+\frac{y}{4}=1
    \]   
    与椭圆 
    \[
    \frac{x^2}{25}+\frac{y^2}{16}=1 
    \] 相交于两点 $A,B$,点 $P$ 在椭圆上使得\( \triangle PAB \) 的面积为 $9$,求点 $P$ 的坐标。
    \begin{solution}
        设点 $P$ 到直线 $AB$ 的距离为
        \[
        d = \frac{2[\triangle PAB]}{AB} = \frac{18}{AB}
        \]
        欲求AB,联立直线与椭圆:
        \[
        \begin{cases}
        \;\dfrac{x}{5}+\dfrac{y}{4}=1\\
        \;\dfrac{x^2}{25} + \dfrac{y^2}{16} = 1
        \end{cases}
        \]
        解得$A(0,4),B(5,0)$,则$d=\dfrac{18}{\sqrt{41}}$
        
        设过 $P$ 与直线 $AB$ 平行的直线为$4x + 5y = k$,有
        \[
        \frac{|k - 20|}{\sqrt{41}} = \frac{18}{\sqrt{41}} \Rightarrow k = 38 \text{ 或 } \,2
        \]
        若 $k = 38$,则      
        $\begin{cases}
        4x + 5y = 38 \\
        \dfrac{x^2}{25} + \dfrac{y^2}{16} = 1
        \end{cases}$无实数解;
        若 $k = 2$,可解得 $$P\left( \dfrac{1 \pm \sqrt{199}}{4},\ \dfrac{1 \mp \sqrt{199}}{5} \right)$$
    \end{solution}

    \question 已知圆$\Omega$与$x$轴、$y$轴均相切,圆心在椭圆$$\Gamma: \frac{x^2}{a^2} + \frac{y^2}{b^2} = 1 $$内,且$\Omega$与$\Gamma$有唯一的公共点$(8,9)$.求$\Gamma$的焦距。
    \begin{solution}
        设圆$\Omega$的圆心为$P(r,r)$,则有$$(8-r)^2+(9-r)^2=r^2\Rightarrow r=5\ \text{或} \ r=29$$
        因为$P$在$\Gamma$内,故$r=5$;椭圆$\Gamma$在点$A(8,9)$处的切线为$$l: \frac{8x}{a^2} + \frac{9y}{b^2} = 1$$
        且有$$m_l\cdot m_{PA}=-1 \Rightarrow \frac{32}{a^2}=\frac{27}{b^2}$$
        联立$$\frac{64}{a^2}+\frac{81}{b^2}=1$$解得$$a^2=160,b^2=135$$从而$\Gamma$的焦距为$$2\sqrt{a^2-b^2}=10$$
    \end{solution}

    \question 已知 $P$ 是椭圆 
    \[
    \frac{x^2}{a^2}+\frac{y^2}{b^2}=1
    \]
    上任意一点,$PN$ 是过点 $P$ 的法线,$F_1, F_2$ 为椭圆焦点。证明 $PN$ 平分 $\angle F_1PF_2$。
    \begin{solution}
        设椭圆上动点为 $P(a\cos\theta, b\sin\theta),\theta \in \mathbb{R}$,过$P$的切线斜率为
        \[
        \frac{dy}{dx} = -\frac{b^2 x}{a^2 y}
        \]
        则过$P$的法线斜率为
        \[
        m_{PN} = -\frac{1}{m_\text{tangent}} = \frac{a^2 y}{b^2 x} = \frac{a\sin\theta}{b\cos\theta}
        \]
        焦点坐标 $F_1(-ae,0)$, $F_2(ae,0)$,其中$e$为离心率,则
        \[
        m_{PF_1} = \frac{b\sin\theta - 0}{a\cos\theta + ae} = \frac{b\sin\theta}{a(\cos\theta + e)},\quad
        m_{PF_2} = \frac{b\sin\theta - 0}{a\cos\theta - ae} = \frac{b\sin\theta}{a(\cos\theta - e)}
        \]
        代入
        \[
        \tan\theta_1 = \frac{m_{PN} - m_{PF_1}}{1 + m_{PN} m_{PF_1}},\quad
        \tan\theta_2 = \frac{m_{PF_2} - m_{PN}}{1 + m_{PN} m_{PF_2}}
        \]
        整理得
        \[
        \tan\theta_1 = \tan\theta_2= \frac{ae\sin\theta}{b},
        \]
        由于角为锐角,$\theta_1 = \theta_2$,得证$PN$平分$\angle F_1PF_2$。
    \end{solution}

    \question 已知两椭圆 $E_1,E_2$ 的方程分别为  
    \[
    E_1: \frac{x^2}{a^2} + \frac{y^2}{b^2} - \frac{2x}{c} = 0, \quad
    E_2: \frac{x^2}{b^2} + \frac{y^2}{a^2} + \frac{2x}{c} = 0.
    \]  
    设两椭圆的公切线与椭圆$E_1,E_2$分别交于$A,B$,且$O$为原点,证明$\angle AOB$ 为直角。
    \begin{solution}
        对 $E_1,E_2$配方得
        \[
        E_1:\frac{\left(x - \frac{a^2}{c}\right)^2}{\frac{a^4}{c^2}} + \frac{y^2}{\frac{a^2 b^2}{c^2}} = 1, \quad \frac{\left(x - \frac{b^2}{c}\right)^2}{\frac{b^4}{c^2}} + \frac{y^2}{\frac{a^2 b^2}{c^2}} = 1.
        \]
        因此,$E_1$ 与 $E_2$ 的公切线为水平线,则
        \[
        A = \left(\frac{a^2}{c}, \pm \frac{c}{ab}\right), \quad B = \left(\frac{b^2}{c}, \pm \frac{c}{ab}\right),
        \]
        观察到
        \[
        m_{OA}\cdot m_{OB}  = \frac{\frac{c}{ab}}{\frac{a^2}{c}} \cdot \frac{\frac{c}{ab}}{\frac{b^2}{c}} = 1
        \]
        因此$OA \perp OB,\angle AOB$ 为直角。
    \end{solution}

    \question 已知$a,b>0$,
    \begin{parts}
    \part 试证过椭圆 
    \[
    \frac{x^{2}}{a^{2}}+\frac{y^{2}}{b^{2}}=1
    \] 
    上两点$P(a\cos\theta, b\sin\theta),Q(a\cos\phi, b\sin\phi)$的直线方程式为
    \[
    \frac{x}{a}\cos\frac{\theta+\phi}{2}+\frac{y}{b}\sin\frac{\theta+\phi}{2}=\cos\frac{\theta-\phi}{2}
    \]
    \begin{solution}
        $PQ$直线方程斜率为
        \[
        m = \frac{b\sin\phi - b\sin\theta}{a\cos\phi - a\cos\theta} 
        = \frac{b}{a} \frac{2\cos\frac{\theta+\phi}{2}\sin\frac{\phi-\theta}{2}}{-2\sin\frac{\theta+\phi}{2}\sin\frac{\phi-\theta}{2}}
        = -\frac{b}{a} \cot\frac{\theta+\phi}{2}.
        \]
        故$PQ$直线方程为
        \[
        y - b\sin\theta = -\frac{b}{a}\frac{\cos\frac{\theta+\phi}{2}}{\sin\frac{\theta+\phi}{2}}(x - a\cos\theta).
        \]
        整理得
        \[
        \frac{x}{a}\cos\frac{\theta+\phi}{2} + \frac{y}{b}\sin\frac{\theta+\phi}{2} = \cos\left(\theta - \frac{\theta+\phi}{2}\right)
        = \cos\frac{\theta-\phi}{2}.
        \]
    \end{solution}
    \part 若弦 $PQ$ 与圆$x^{2}+y^{2}=b^{2}$相切, 试证: 
    \[
    a^{2}\cos^{2}\frac{\theta-\phi}{2}=b^{2}\cos^{2}\frac{\theta+\phi}{2}+a^{2}\sin^{2}\frac{\theta+\phi}{2}
    \]
    \begin{solution}
        由于弦$PQ$到圆心 $(0,0)$ 的距离等于半径,
        \[
        \frac{|\cos\frac{\theta-\phi}{2}|}{\sqrt{\frac{\cos^2\frac{\theta+\phi}{2}}{a^2} + \frac{\sin^2\frac{\theta+\phi}{2}}{b^2}}} = b
        \]
        整理得
        \[
        a^{2}\cos^{2}\frac{\theta-\phi}{2}=b^{2}\cos^{2}\frac{\theta+\phi}{2}+a^{2}\sin^{2}\frac{\theta+\phi}{2}
        \]
    \end{solution}
    \part 若 $\sin(\theta-\phi)\ge 0$, 证明弦 $PQ$ 之长为 $a\sin(\theta-\phi)$。
    \begin{solution}
        有
        \begin{align*}
        PQ^2 &= (a\cos\theta - a\cos\phi)^2 + (b\sin\theta - b\sin\phi)^2 \\
        &= a^2\left(-2\sin\frac{\theta+\phi}{2}\sin\frac{\theta-\phi}{2}\right)^2 + b^2\left(2\cos\frac{\theta+\phi}{2}\sin\frac{\theta-\phi}{2}\right)^2 \\
        &= 4\sin^2\frac{\theta-\phi}{2} \left(a^2 \sin^2\frac{\theta+\phi}{2} + b^2 \cos^2\frac{\theta+\phi}{2}\right) \\
        &= 4\sin^2\frac{\theta-\phi}{2} \left(a^{2}\cos^{2}\frac{\theta-\phi}{2}\right) \\
        &= a^2 \left(2\sin\frac{\theta-\phi}{2}\cos\frac{\theta-\phi}{2}\right)^2 \\
        &= a^2 \sin^2(\theta-\phi)
        \end{align*}
        由于$a>0,\sin(\theta-\phi)\ge 0$,故
        \[
        PQ = a \sin(\theta-\phi)
        \]
    \end{solution}
    \end{parts}

    \question 设 $A,B,C$ 为椭圆 
    \[
    \Gamma: \frac{x^{2}}{16}+\frac{y^{2}}{32}=1
    \] 
    上三点,且 $\triangle ABC$ 的重心恰为此椭圆的中心,已知 $A(\sqrt{6}+\sqrt{2},2\sqrt{3}-2)$,求 $\triangle ABC$ 面积。
    \begin{solution}
        设变换:
        \[
        x = x', y = \sqrt{2}y'
        \]
        代入椭圆 $\Gamma$,得
        \[
        \frac{x'^2}{16} + \frac{y'^2}{16} = 1
        \]
        为一圆 $\Gamma'$,其中圆心为原点 $(0, 0)$,半径为 $4$,又因为
        \[
        x' = x, y' = \frac{y}{\sqrt{2}}
        \Rightarrow A'(\sqrt{6}+\sqrt{2}, \sqrt{6}-\sqrt{2})
        \Rightarrow OA' = 4
        \]
        由于 $\triangle A'B'C'$ 为正三角形,且其重心为原点,边长为
        \[
        A'B' = 4 \cdot \frac{3}{2} \cdot \frac{2}{\sqrt{3}} = 4\sqrt{3}
        \Rightarrow [\triangle A'B'C'] = \frac{\sqrt{3}}{4}(4\sqrt{3})^2 = 12\sqrt{3}
        \]
        换回原坐标系,雅可比行列式为
        \[
        J = \left| \frac{\partial(x, y)}{\partial(x', y')} \right| = 
        \begin{vmatrix}
        1 & 0\\
        0 & \sqrt{2}
        \end{vmatrix} = \sqrt{2}
        \Rightarrow [\triangle ABC] = J \cdot [\triangle A'B'C'] = 12\sqrt{6}
        \]
    \end{solution}

    \question 已知 $\triangle ABC$ 的顶点 $A,B$ 在椭圆
    \[
    \frac{x^2}{4}+\frac{y^2}{3}=1
    \]
    上,点 $C$ 在直线 $l:y=x+2$ 上,且 $AB \parallels l$。
    \begin{parts}
    \part 若 $AB$ 通过坐标原点 $O$ 时,求 $AB$ 的长及 $\triangle ABC$ 的面积；
    \begin{solution}
        设 $A,B$ 两点坐标分别为 $(x_1,y_1),(x_2,y_2)$,已知 $AB\parallels l$,且 $AB$ 经过原点 $O(0,0)$,
        故 $AB$ 直线方程为$y=x$,与椭圆方程
        \[
        \frac{x^2}{4}+\frac{y^2}{3}=1
        \]
        联立解得$x=\pm1$,于是
        \[
        AB=\sqrt{(x_1-x_2)^2+(y_1-y_2)^2}
        =\sqrt{2}\,|x_1-x_2|=2\sqrt{2}
        \]
        又 $AB$ 上的高等于原点到直线 $l:y=x+2$ 的距离,
        故
        \[
        [\triangle ABC]=\frac{1}{2}\cdot 2\sqrt{2}\cdot \frac{|2|}{\sqrt{2}}=2
        \]
    \end{solution}
    \part 若 $\angle ABC=90^\circ$ 且斜边 $AC$ 的长最大时,求 $AB$ 所在直线的方程。
    \begin{solution}
        设 $AB$ 所在直线的方程为$y = x + c $,将其代入方程 $x^2 + 3y^2 = 4$,整理得
        \[ 
        4x^2 + 6cx + 3c^2 - 4 = 0 
        \]
        因为 $A, B$ 在椭圆上,故判别式满足
        \[ 
        \Delta = (6c)^2 - 4 \cdot 4 \cdot (3c^2 - 4) = 64 - 12c^2 > 0 \Rightarrow c^2 < \frac{16}{3}
        \]
        设 $A, B$ 两点的坐标分别为 $(x_1, y_1), (x_2, y_2)$,由韦达定理得
        \[ 
        x_1 + x_2 = -\frac{3c}{2}, \quad x_1x_2 = \frac{3c^2 - 4}{4} 
        \]
        于是线段 $AB$ 的长度平方为
        \[ 
        AB^2 = 2(x_1 - x_2)^2 = 2\left[(x_1 + x_2)^2 - 4x_1x_2\right] = \frac{32 - 3c^2}{2} 
        \]
        点 $B$ 到直线 $l: y = x + 2$ 的距离 $|BC|$ 为
        \[ 
        |BC| = \frac{|2 - c|}{\sqrt{2}}
        \]
        因为 $\angle ABC = 90^\circ$,由毕氏定理,
        \[ 
        AC^2 = AB^2 + BC^2 = \frac{32 - 3c^2}{2}  + \frac{(2 - c)^2}{2} = -(c + 1)^2 + 19
        \]
        当 $c = -1$ 时,$AC^2$ 取得最大值 $19$,此时 $c^2 = 1 < \frac{16}{3}$,符合判别式条件,因此$AB$ 所在直线的方程为:
        \[ 
        y = x - 1 
        \]
    \end{solution}
    \end{parts}

    \question 设椭圆
    \[
    \frac{x^{2}}{a^{2}}+\frac{y^{2}}{b^{2}}=1\ (a>b>0)
    \]
    的左右焦点分别为 $F_{1},F_{2}$,离心率
    \[
    e=\frac{\sqrt{2}}{2},
    \]
    点 $F_{2}$ 到右准线 $l$ 的距离为 $\sqrt{2}$。
    \begin{parts}
    \part 求 $a,b$ 的值；
    \begin{solution}
        由离心率定义
        \[
        e=\frac{c}{a}=\frac{\sqrt{2}}{2}
        \]
        又点 $F_{2}$ 到右准线 $l$ 的距离为
        \[
        d=\frac{a}{e}-c=\sqrt{2}
        \]
        联立解得
        \[
        a=2,c=\sqrt{2}
        \]
        又
        \[
        b^{2}=\sqrt{a^{2}-c^{2}}=\sqrt{2}
        \]
    \end{solution}
    \part 设 $M,N$ 是 $l$ 上的两个动点,且
    \[
    \overrightarrow{F_{1}M}\cdot\overrightarrow{F_{2}N}=0,
    \]
    证明当 $MN$ 取最小值时,
    \[
    \overrightarrow{F_{2}F_{1}}+\overrightarrow{F_{2}M}+\overrightarrow{F_{2}N}=\vec{0}.
    \]
    \begin{solution}
        由 $c=\sqrt{2},a=2$可得$F_{1}(-\sqrt{2},0), F_{2}(\sqrt{2},0)$,右准线 $l$ 的方程为
        \[
        x=2\sqrt{2}
        \]
        设$M(2\sqrt{2},y_{1}), N(2\sqrt{2},y_{2})$,由$\overrightarrow{F_{1}M}\cdot\overrightarrow{F_{2}N}=0$得
        \[
        (2\sqrt{2}+\sqrt{2},y_{1})\cdot(2\sqrt{2}-\sqrt{2},y_{2})=0 \Rightarrow y_1y_2=-6 \Rightarrow y_{2}=-\frac{6}{y_{1}}
        \]
        于是由AM-GM不等式,
        \[
        MN=|y_{1}-y_{2}|=\left|y_{1}+\frac{6}{y_{1}}\right|\ge 2\sqrt{6}
        \]
        当且仅当 $y_{1}=\pm\sqrt{6}$ 时取等号,即 $y_{2}=-y_{1}$,此时
        \[
        \overrightarrow{F_{2}F_{1}}+\overrightarrow{F_{2}M}+\overrightarrow{F_{2}N}
        =(2\sqrt{2},0)+(-\sqrt{2},y_{1})+(-\sqrt{2},y_{2})
        =(0,y_{1}+y_{2})
        =\vec{0}.
        \]
        证毕。
    \end{solution}
    \end{parts}

    \question 已知抛物线 
    \[
    E: x^{2} = 2 p y \;(p > 0)
    \]
    的焦点为 $F$,点 $P(m,n)\; (m \neq 0)$ 在 $E$ 上,过点 $P$ 作抛物线的切线 $l$,直线 $l$ 与椭圆 
    \[
    C: \dfrac{x^{2}}{a^{2}} + \dfrac{y^{2}}{b^{2}} = 1\; (a > b > 0)
    \]
    交于两不同点 $A, B$,记线段 $AB$ 的中点为 $D$。设$O$ 为原点,直线 $OD$与直线 $y = -\dfrac{p b^{2}}{a^{2}}$ 交于点 $Q$。
    \begin{parts}
    \part 证明直线 $y = -\dfrac{p b^{2}}{a^{2}}$ 与直线 $PQ$ 垂直。
    \begin{solution}
        点 $P$ 在抛物线上,故 $P(m,\dfrac{m^{2}}{2 p})$,抛物线 $E: x^{2} = 2 p y$求导得切线$l$斜率$\dfrac{m}{p}$,故过点 $P$ 作抛物线的切线 $l$方程为
        \[
        y-\frac{m^{2}}{2 p}=\dfrac{m}{p}(x-m)\Rightarrow y = \frac{m}{p} x - \frac{m^{2}}{2 p}.
        \]
        与椭圆方程
        \[
        \frac{x^{2}}{a^{2}} + \frac{y^{2}}{b^{2}} = 1,
        \]
        联立得化为关于 $x$ 的二次方程
        \[
        4(a^{2} m^{2} + b^{2} p^{2}) x^{2} - 4 a^{2} m^{3} x + a^{2} (m^{4} - 4 p^{2} b^{2}) = 0.
        \]
        记两根为 $x_{1}, x_{2}$,则
        \[
        x_{1} + x_{2} = \frac{a^{2} m^{3}}{a^{2} m^{2} + b^{2} p^{2}}, \quad
        y_{1} + y_{2} = \frac{m}{p} (x_{1} + x_{2}) - \frac{m^{2}}{p} = - \frac{m^{2} b^{2} p}{a^{2} m^{2} + b^{2} p^{2}}.
        \]
        中点 $D$ 坐标为
        \[
        D \left( \frac{a^{2} m^{3}}{2 (a^{2} m^{2} + b^{2} p^{2})}, \, -\frac{m^{2} b^{2} p}{2 (a^{2} m^{2} + b^{2} p^{2})} \right)
        \]
        故直线 $OD$ 方程为
        \[
        y=-\frac{b^{2} p}{a^{2} m}x
        \]
        与 $y = -\dfrac{p b^{2}}{a^{2}}$ 联立得点$Q$的方程
        \[
        x = m
        \]
        故 $PQ$ 是垂直于 $x$ 轴的直线,而直线 $y = -\dfrac{p b^{2}}{a^{2}}$ 是水平线,二者垂直。
    \end{solution}
    \part 设 $a = 2 b$,直线 $l$ 与 $F Q$ 交于 $G$,证明
    \[
    1 < \frac{[\triangle F P G]}{[\triangle P G Q]} < 2
    \]
    \begin{solution}
        当 $a = 2 b$,直线 $y = -\dfrac{p b^{2}}{a^{2}}$ 即
        \[
        y = -\frac{p}{4},
        \]
        为介于抛物线准线 $y = -\dfrac{p}{2}$ 和 $x$ 轴之间的平行线。延长 $PQ$ 与抛物线准线交于点 $H \left(m, -\dfrac{p}{2}\right)$,焦点 $F \left(0, \dfrac{p}{2}\right)$,则直线 $HF$ 斜率
        \[
        k_{HF} = \frac{\frac{p}{2} + \frac{p}{2}}{0 - m} = -\frac{p}{m}
        \]
        切线 $l$ 斜率为
        \[
        k = \frac{m}{p},
        \]
        故 $l \perp HF$。由抛物线定义
        \[
        PF = PH = \frac{m^{2}}{2 p} + \frac{p}{2}, \quad
        PQ = PH - \frac{p}{4} = \frac{m^{2}}{2 p} + \frac{p}{4}
        \]
        $\triangle PFH$ 为等腰三角形,且切线 $l$ 平分 $\angle F P H$,由角平分线性质,
        \[
        \frac{[\triangle F P G]}{[\triangle P G Q]} = \frac{F G}{G Q} = \frac{P F}{P Q} = \frac{2 m^{2} + 2 p^{2}}{2 m^{2} + p^{2}}
        \]
        显然,
        \[
        1 < \frac{[\triangle F P G]}{[\triangle P G Q]} < 2
        \]
    \end{solution}
    \end{parts}

    \question 小何在打篮球时发现地面上的篮球在灯光照射下在地面形成的影子是一个椭圆, 并观察到地面与篮球接触点正是椭圆的焦点。如图,已知篮球半径为 $1$,灯源 $P$ 离地高度为 $4$,其正下方为点 $A$, 地面上椭圆右顶点与点 $A$ 的水平距离为 $3$,求该椭圆的离心率 $e$。
    \begin{figure}[H]
        \centering        
        \includegraphics[width=0.3\textwidth]{images/image16.png}
    \end{figure}
    \ifprintanswers
    \begin{figure}[H]
        \centering        
        \includegraphics[width=0.3\textwidth]{images/image17.png}
    \end{figure}
    \fi
    \begin{solution}
        以$A$为原点建立平面直角坐标系,则$P(0,4)$, $R(-3,0)$, 直线$PR$的方程为 $$y=\frac{4}{3}x+4$$
        设$M(n,1)$, $Q(n,0)$,由$M$到直线$PR$的距离为1,$$\left|\frac{\frac{4}{3}n+4-1}{\sqrt{1+(\frac{4}{3})^{2}}}\right|=1$$
        解之得 $n=-\frac{7}{2}$\,(舍 $n=-1$),则 $M(-\frac{7}{2},1)$, $Q(-\frac{7}{2},0)$。
        
        设直线$PN$的方程为 $y=kx+4$,由M到直线$PN$的距离为1,$$\left|\frac{-\frac{7}{2}k+4-1}{\sqrt{1+k^{2}}}\right|=1$$
        整理得 $$45k^{2}-21k+8=0\Rightarrow k_{PN}=\dfrac{8}{15}\;(\text{舍}\,k_{PR}=\dfrac{4}{3}\,)$$ 
        故直线$PN$的方程为 $y=\dfrac{8}{15}x+4$,得$N(-\dfrac{15}{2},0)$,故
        $$\begin{cases} NQ=-\frac{7}{2}+\frac{15}{2}=4=a+c \\ RQ=-3+\frac{7}{2}=\frac{1}{2}=a-c \end{cases}\Rightarrow a=\frac{9}{4},\,c=\frac{7}{4}$$ 
        所以 $$e=\dfrac{c}{a}=\frac{7}{9}$$
    \end{solution}
    
    \question 椭圆焦点为 \(F_1,F_2\),点 \(P\) 在椭圆上。设 \(O\) 为 \(\triangle PF_1F_2\) 外接圆心,且
        \(\overrightarrow{PO}\cdot\overrightarrow{F_1F_2}=2\,\overrightarrow{PF_1}\cdot\overrightarrow{PF_2}\)。求椭圆最小离心率。  
    \begin{solution}
        设$F_1(-c,0),\ F_2(c,0),\ P(x,y)$,椭圆方程为
        \[
        \frac{x^2}{a^2} + \frac{y^2}{b^2} = 1
        \]
        由于 \(O\) 是外接圆心,且 \(F_1F_2\) 是弦,故 \(O\) 在 \(y\) 轴上,设$O(0,y_O)$,由$|OP| = |OF_1| = |OF_2|,$
        \[
        x^2 + (y - y_O)^2 = c^2 + y_O^2 \Rightarrow x^2 + y^2 - 2 y y_O = c^2 \tag{1}
        \]
        由条件$\overrightarrow{PO}\cdot\overrightarrow{F_1F_2}=2\,\overrightarrow{PF_1}\cdot\overrightarrow{PF_2}$得到
        \[
        (-x, y_O - y)\cdot(2c,0)=(-c - x, -y)\cdot(c - x, -y)\Rightarrow -c x = x^2 - c^2 + y^2 \tag{2}
        \]
        由$(1),(2)$得
        \[
        2 y y_O = - c x\Rightarrow y_O = -\frac{c x}{2 y}\;(y \neq 0)
        \]
        代回$(1)$得
        \[
        x^2 + y^2 + c x = c^2
        \]
        又由椭圆方程
        \[
        y^2 = b^2 \left(1 - \frac{x^2}{a^2}\right) = (a^2 - c^2) \left(1 - \frac{x^2}{a^2}\right).
        \]
        代入得
        \[
        x^2 + (a^2 - c^2)\left(1 - \frac{x^2}{a^2}\right) + c x = c^2
        \]
        整理为
        \[
        \frac{c^2}{a^2} x^2 + c x + a^2 - 2 c^2 = 0
        \]
        此为关于 \(x\) 的二次方程,存在实数解则判别式 \(\Delta \geq 0\):
        \[
        \Delta = c^2 - 4 \frac{c^2}{a^2} (a^2 - 2 c^2) \geq 0 \Rightarrow \frac{c^2}{a^2} \geq \frac{3}{8}
        \]
        故最小离心率为
        \[
        e = \frac{c}{a} \geq \sqrt{\frac{3}{8}} = \frac{\sqrt{6}}{4}
        \]
    \end{solution}

    \question 在平面直角坐标系中,两点到点 $(0, -\sqrt{3}),(0, \sqrt{3})$ 的距离之和等于 4,设点 $P$ 的轨迹为 $C$。
    \begin{parts}
    \part 写出 $C$ 的方程；
    \begin{solution}
        据题意,所求轨迹方程式为
        \[
        \sqrt{x^2+(y+\sqrt{3})^2}+\sqrt{x^2+(y-\sqrt{3})^2}=4
        \]
        化简得
        \[
        x^2 + \frac{y^2}{4} = 1
        \]
    \end{solution}
    \part 设直线 $y = kx + 1$ 与 $C$ 交于 $A, B$ 两点,求 $k$ 使得 $OA \perp OB$,并求此时 $AB$ 的长。
    \begin{solution}
        设 $A(x_1, y_1)$, $B(x_2, y_2)$, 满足
        \[
        x^2 + \frac{y^2}{4} = 1,\quad y = kx + 1
        \]
        联立得
        \[
        (k^2 + 4)x^2 + 2kx - 3 = 0
        \]
        由韦达定理,
        \[
        x_1 + x_2 = -\frac{2k}{k^2 + 4},\quad x_1 x_2 = -\frac{3}{k^2 + 4}.
        \]
        由于$OA \perp OB$ 等价于 $x_1 x_2 + y_1 y_2 = 0$,而 
        \[
        y_1 y_2 = (kx_1+1)(kx_2+1) = k^2 x_1 x_2 + k(x_1+x_2) + 1,
        \]
        所以
        \[
        x_1x_2 + y_1y_2 = (k^2+1)x_1 x_2 + k(x_1+x_2) + 1 = \frac{-4k^2 + 1}{k^2 + 4}
        \]
        所以当 $k = \pm \dfrac{1}{2}$时,有$x_1 x_2 + y_1 y_2 = 0$即$OA \perp OB$。且当 $k = \pm \dfrac{1}{2}$,
        \[
        x_1 + x_2 = \mp \frac{4}{17},\quad x_1 x_2 = -\frac{12}{17},
        \]
        于是
        \begin{align*}
        AB &= \sqrt{(x_2-x_1)^2 + (y_2-y_1)^2} \\
        &= \sqrt{(1+k^2)(x_2-x_1)^2} \\
        &= \sqrt{(1+k^2)((x_2+x_1)^2-4x_1x_2)} \\
        &= \sqrt{\frac{5}{4}\cdot\left(\frac{16}{289} + \frac{48}{17}\right)} \\
        &= \frac{8\sqrt{13}}{17}
        \end{align*}   
    \end{solution}
    \end{parts}

    \question 设椭圆中心在坐标原点,$A(2,0),B(0,1)$ 是它的两个顶点,直线 $y=kx,k>0$ 与 $AB$ 相交于点 $D$,与椭圆相交于 $E,F$ 两点。
    \begin{parts}
    \part 若 $\overrightarrow{ED}=6\overrightarrow{DF}$,求 $k$ 的值；
    \begin{solution}
        由题意,椭圆方程为
        \[
        \frac{x^2}{4}+y^2=1.
        \]
        直线 $AB,EF$ 的方程分别为
        \[
        x+2y=2,\quad y=kx\ ,k>0
        \]
        设 $D(x_0,kx_0),E(x_1,kx_1),F(x_2,kx_2),x_1<x_2$,联立椭圆与$y=kx$,解得
        \[
        \frac{x^2}{4}+k^2x^2=1 \Rightarrow x_2=-x_1=\frac{2}{\sqrt{1+4k^2}} 
        \]
        由 $\overrightarrow{ED}=6\overrightarrow{DF}$,得$x_0-x_1=6(x_2-x_0)$,从而
        \[
        x_0=\frac{1}{7}(6x_2+x_1)=\frac{5}{7}x_2=\frac{10}{7\sqrt{1+4k^2}} \tag{1}
        \]
        又因 $D$ 在 $AB$ 上,
        \[
        x_0+2kx_0=2 \Rightarrow x_0=\frac{2}{1+2k} \tag{2}
        \]
        联立$(1),(2)$,解得
        \[
        24k^2-25k+6=0 \Rightarrow k=\frac{2}{3}\ \text{或}\ k=\frac{3}{8}
        \]
    \end{solution}
    \part 求四边形 $AEBF$ 面积的最大值。
    \begin{solution}
        点 $E,F$ 到直线 $AB$ 的距离分别为
        \[
        h_1=\frac{|x_1+2kx_1-2|}{\sqrt{5}}
        =\frac{2(1+2k-\sqrt{1+4k^2})}{\sqrt{5(1+4k^2)}},
        \]
        \[
        h_2=\frac{|x_2+2kx_2-2|}{\sqrt{5}}
        =\frac{2(1+2k+\sqrt{1+4k^2})}{\sqrt{5(1+4k^2)}}.
        \]
        又$AB=\sqrt{2^2+1^2}=\sqrt{5}$,四边形 $AEBF$ 的面积为
        \[
        [AEBF]=\frac{1}{2} AB (h_1+h_2)
        =\frac{1}{2}\cdot \sqrt{5} \cdot \frac{4(1+2k)}{\sqrt{5(1+4k^2)}}
        =\frac{2(1+2k)}{\sqrt{1+4k^2}}
        =2\sqrt{1+\frac{4k}{1+4k^2}}
        \]
        由AM-GM不等式,
        \[
        \frac{1+4k^2}{4k}=\frac{1}{4k}+k \ge 1
        \]
        当 $k=\dfrac{1}{2}$ 时等号成立,此时四边形 $AEBF$ 面积取最大值$2\sqrt{2}$。
    \end{solution}
    \begin{solution}
        已知 $BO=1, AO=2$,设 $y_1 = kx_1, y_2 = kx_2$,由$x_2=\dfrac{2}{\sqrt{1+4k^2}}$ 知 $x_2 > 0, y_2 = -y_1 > 0$,故四边形 $AEBF$ 面积为
        \[ 
        [AEBF] = [\triangle BEF] + [\triangle AEF] = x_2 + 2y_2 = \sqrt{(x_2 + 2y_2)^2}
        \]
        由柯西不等式,
        \[ 
        [AEBF] \le \sqrt{2(x_2^2 + 4y_2^2)} = 2\sqrt{2} 
        \]
        当 $x_2 = 2y_2$ 时,四边形 $AEBF$ 面积取得最大值 $2\sqrt{2}$。
    \end{solution}
    \end{parts}

    \question 已知椭圆 
    \[
    C: \frac{x^2}{a^2} + \frac{y^2}{b^2} = 1,\quad a>b>0,
    \]
    的一个焦点为 $F(1,0)$,且椭圆过点 $(2,0)$。
    \begin{parts}
    \part 求椭圆$C$的方程。
    \begin{solution}
        由题意,椭圆的半长轴 $a=2$,焦距 $c=1$,则$b^2 = a^2 - c^2 = 3$,因此椭圆方程为
        \[
        \frac{x^2}{4} + \frac{y^2}{3} = 1
        \]
    \end{solution}
    \part 若 $AB$ 为垂直于 $x$ 轴的弦,直线 $l: x=4$ 与 $x$ 轴交于点 $N$,直线 $AF$ 与 $BN$ 交于点 $M$,试证点 $M$ 恒在椭圆 $C$ 上。
    \begin{solution}
        设 $A(m,n)$,则 $B(m,-n),n\ne 0$,且满足椭圆方程
        \[
        \frac{m^2}{4} + \frac{n^2}{3} = 1 \Rightarrow 4n^2=12-3m^2
        \]
        直线 $AF,BN$ 方程分别为
        \[
        n(x-1) - (m-1)y = 0, \quad n(x-4) - (m-4)y = 0 
        \]
        联立解得交点 $M(x_0,y_0)$
        \[
        x_0 = \frac{5m-8}{2m-5},\quad y_0 = \frac{3n}{2m-5}.
        \]
        由于
        \[
        \frac{x_0^2}{4} + \frac{y_0^2}{3} 
        = \frac{(5m-8)^2}{4(2m-5)^2} + \frac{3n^2}{(2m-5)^2}
        = \frac{(5m-8)^2 + 3(12-3m^2)}{4(2m-5)^2} 
        = 1
        \]
        因此点 $M$ 恒在椭圆 $C$ 上。
    \end{solution}
    \begin{solution}
        由题意得 $F(1, 0), N(4, 0)$,设 $A(m, n)$,则 $B(m, -n),n \neq 0$,满足椭圆方程
        \[
        \frac{m^2}{4} + \frac{n^2}{3} = 1 \tag{1}
        \]
        直线 $AF$ 与 $BN$ 的方程分别为
        \[
        n(x-1) - (m-1)y = 0, \quad n(x-4) - (m-4)y = 0 \tag{2}
        \]
        联立可得当$x \ne \dfrac{5}{2}$时,
        \[
        m = \frac{5x-8}{2x-5},\quad n = \frac{3y}{2x-5}
        \]
        代入 (1),整理得
        \[
        \frac{x^2}{4} + \frac{y^2}{3} = 1, \quad y \ne 0.
        \]
        而当 $x = \dfrac{5}{2}$ 时,由 (2) 得
        \[
        \begin{cases}
        \dfrac{3}{2} n - (m-1)y = 0 \\[6pt]
        -\dfrac{3}{2} n + (m+4)y = 0
        \end{cases}
        \Rightarrow n=0,\, y=0,
        \]
        与 $n \ne 0$ 矛盾,因此点 $M$ 的轨迹方程为
        \[
        \frac{x^2}{4} + \frac{y^2}{3} = 1, \quad y \ne 0,
        \]
        即点 $M$ 恒在椭圆 $C$ 上。
    \end{solution}
    \part 求 $\triangle AMN$ 面积的最大值。
    \begin{solution}
        设 $AM$ 的方程为 $x = t y + 1$,代入椭圆方程,化简得
        \[
        (3t^2+4)y^2 + 6t y - 9 = 0
        \]
        设 $A(x_1,y_1),M(x_2,y_2)$,由韦达定理,
        \[
        y_1 + y_2 = -\frac{6t}{3t^2+4},\quad y_1y_2 = -\frac{9}{3t^2+4}.
        \]
        于是
        \[
        |y_1 - y_2| = \sqrt{(y_1+y_2)^2 - 4y_1y_2} = \frac{4\sqrt{3}\sqrt{3t^2+3}}{3t^2+4}.
        \]
        令 $3t^2+4 = \lambda \ge 4$,得
        \[
        |y_1-y_2| = 4\sqrt{3} \sqrt{\frac{\lambda-1}{\lambda^2}}
        = 4\sqrt{3} \sqrt{-\left(\frac{1}{\lambda}-\frac{1}{2}\right)^2 + \frac{1}{4}}.
        \]
        因为$\lambda \ge 4,0<\dfrac{1}{\lambda}\le\dfrac{1}{4}$,所以当$\dfrac{1}{\lambda}=\dfrac{1}{4}$即 $\lambda=4,t=0$ 时,$|y_1-y_2|$ 取得最大值 $3$,此时
        \[
        [\triangle AMN] = \frac{1}{2}FN |y_1-y_2| = \frac{3}{2}|4-1| = \frac{9}{2}
        \]
    \end{solution}
    \end{parts}

    \question 设椭圆
    \[
    C: \frac{x^2}{a^2}+\frac{y^2}{b^2}=1,\quad a>b>0,
    \]
    其相应于焦点 $F(2,0)$ 的准线方程为 $x=4$。
    \begin{parts}
    \part 求椭圆 $C$ 的方程。
    \begin{solution}
        由题意得,
        \[
        \begin{cases} 
        c = 2 \\[3pt]
        \dfrac{a^2}{c} = 4 \\[3pt]
        a^2 = b^2 + c^2 
        \end{cases}
        \implies a^2 = 8,\; b^2 = 4 
        \]
        因此,椭圆 $C$ 的方程为 
        \[
        \frac{x^2}{8} + \frac{y^2}{4} = 1
        \]
    \end{solution}
    \part 已知过点 $F_1(-2,0)$、倾斜角为 $\theta$ 的直线交椭圆 $C$ 于 $A,B$ 两点,证明
    \[
    AB=\frac{4\sqrt{2}}{2-\cos^2\theta};
    \]
    \begin{solution}
        椭圆 $C$ 的左焦点为$F_1(-2,0)$,离心率$e=\dfrac{\sqrt{2}}{2}$,左准线方程为 $l:x=-4$。
        
        设直线 $AB$ 的倾斜角为 $\theta$,作 $AA_1\perp l$于$A_1,BB_1\perp l$于$B_1$,且$l$于$x$轴交于点$H$,由于 $A$ 在椭圆上,
        \[
        AF_1=\frac{\sqrt{2}}{2}AA_1=\frac{\sqrt{2}}{2}(F_1H+AF_1\cos\theta)=\sqrt{2}+\frac{\sqrt{2}}{2}AF_1\cos\theta
        \]
        可得
        \[
        AF_1=\frac{2}{\sqrt{2}-\cos\theta}
        \]
        同理
        \[
        BF_1=\frac{2}{\sqrt{2}+\cos\theta}
        \]
        因此
        \[
        AB=AF_1+BF_1=\frac{2}{\sqrt{2}-\cos\theta}+\frac{2}{\sqrt{2}+\cos\theta}=\frac{4\sqrt{2}}{2-\cos^2\theta}
        \]
    \end{solution}
    \begin{solution}
        当 $\theta \neq \dfrac{\pi}{2}$ 时,记 $k = \tan \theta$,则直线 $AB$ 的方程为$y = k(x+2)$,将其代入椭圆方程 $x^2 + 2y^2 = 8$,整理得
        \[ 
        (1+2k^2)x^2 + 8k^2x + 8(k^2-1) = 0 
        \]
        设 $A(x_1, y_1), B(x_2, y_2)$,由韦达定理得
        \[ 
        x_1 + x_2 = -\frac{8k^2}{1+2k^2}, \quad x_1x_2 = \frac{8(k^2-1)}{1+2k^2} 
        \]
        计算弦长$AB$,
        \begin{align*}
        AB &= \sqrt{(x_1-x_2)^2 + (y_1-y_2)^2} \\
        &= \sqrt{(1+k^2)(x_1-x_2)^2} \\
        &= \sqrt{(1+k^2)[(x_1+x_2)^2 - 4x_1x_2]} \\
        &= \sqrt{(1+k^2)\left[\left(\frac{-8k^2}{1+2k^2}\right)^2 - \frac{32(k^2-1)}{1+2k^2}\right]} \\
        &= \frac{4\sqrt{2}(1+k^2)}{1+2k^2}
        \end{align*}
        将 $k^2 = \tan^2 \theta$代入得
        \[ 
        AB = \frac{4\sqrt{2}}{2-\cos^2 \theta} \tag{1}
        \]
        当 $\theta = \dfrac{\pi}{2}$ 时,$AB = 2\sqrt{2}$,仍满足式(1),于是得证
        \[ 
        AB = \frac{4\sqrt{2}}{2-\cos^2 \theta} 
        \]
    \end{solution}
    \part 过点 $F_1(-2,0)$ 作两条互相垂直的直线分别交椭圆 $C$ 于 $A,B$ 和 $D,E$,求 $AB+DE$ 的最小值。
    \begin{solution}
        因为 $AB\perp DE$,其倾斜角为 $\theta+\dfrac{\pi}{2}$,由 $(b)$ 可得
        \[
        AB=\frac{4\sqrt{2}}{2-\cos^2\theta},\quad
        DE=\frac{4\sqrt{2}}{2-\sin^2\theta}.
        \]
        于是
        \[
        AB+DE
        =\frac{4\sqrt{2}}{2-\cos^2\theta}+\frac{4\sqrt{2}}{2-\sin^2\theta} 
        =\frac{12\sqrt{2}}{2+\sin^2\theta\cos^2\theta}
        =\frac{12\sqrt{2}}{2+\frac{1}{4}\sin^2 2\theta}
        \]
        当 $\sin^2 2\theta=1$,即 $\theta=\dfrac{\pi}{4}$ 或 $\dfrac{3\pi}{4}$ 时,$\ AB+DE$ 取得最小值$\dfrac{16\sqrt{2}}{3}$。
    \end{solution}
    \end{parts}

    \question 已知椭圆
    \[
    C: \frac{y^{2}}{a^{2}}+\frac{x^{2}}{b^{2}}=1,\quad a>b>0
    \]
    的离心率为 $\dfrac{1}{2}$,上下焦点为 $F_{1}, F_{2}$,右顶点为 $D$；过 $F_{1}$ 做垂直于 $DF_{2}$ 的直线交椭圆 $C$ 于 $A, B$ 两点,且
    \[
    BD - AF_{1} = \frac{8\sqrt{3}}{39}
    \]
    \begin{parts}
    \part 求 $AD + BD$ 的值。
    \begin{solution}
        由椭圆离心率可知 $a:b:c = 2:\sqrt{3}:1$,则 $F_{1}F_{2} = F_{1}D = F_{2}D = 2c = a$,即 $\triangle F_{1}F_{2}D$ 为等边三角形；因为 $AB \perp DF_{2}$,且直线 $AB$ 过焦点 $F_{1}$,
        故直线 $AB$ 垂直平分线段 $DF_{2}$,由此得
        \[ 
        AD = AF_{2}, \quad BD = BF_{2} 
        \]
        由椭圆定义知 $BF_{1} + BF_{2} = 2a$,则有
        \[ 
        BD - AF_{1} = BF_{2} - (AB - BF_{1}) = 2a - AB = 4c - AB = \frac{8\sqrt{3}}{39} 
        \]
        解法一:设 $A(x_{1}, y_{1}), B(x_{2}, y_{2})$,
        \[ 
        \begin{cases}
        l_{AB}: y = -\sqrt{3}x + c \\ 
        C: \dfrac{y^{2}}{4c^{2}} + \dfrac{x^{2}}{3c^{2}} = 1
        \end{cases} \implies 13x^{2} - 6\sqrt{3}cx - 9c^{2} = 0 
        \]
        由韦达定理,
        \[
        AB = \sqrt{1+(-\sqrt{3})^2} \cdot \sqrt{(x_1+x_2)^2 - 4x_1x_2} 
        = \frac{2}{13}\sqrt{108c^2 + 36 \cdot 13c^2} = \frac{48c}{13}
        \]
        由 
        \[
        4c - AB = \frac{4c}{13} = \frac{8\sqrt{3}}{39},
        \]
        解得 
        \[
        c = \frac{2\sqrt{3}}{3} \Rightarrow a = \frac{4\sqrt{3}}{3}, b = 2
        \],
        故
        \[ 
        AD + BD = AF_2 + BF_2 = (2a - AF_1) + (2a - BF_1) = 4a - AB = \frac{112\sqrt{3}}{39} 
        \]
    \end{solution}
    \begin{solution}
        解法二:由焦半径公式可知 
        \[ 
        AB = \frac{2ab^2}{a^2 - c^2\cos^2\theta} = \frac{4c \cdot 3c^2}{4c^2 - c^2 \cdot \frac{3}{4}} = \frac{48}{13}c 
        \]
        其中 $\theta$ 为直线 $AB$ 与 $y$ 轴的夹角,其余步骤同解法一。
    \end{solution}
    \part 过 $A, B$ 做椭圆 $C$ 的两条切线交于点 $E$,若 $F_{1}E$ 交 $x$ 轴于 $P$,$F_{2}E$ 交 $x$ 轴于 $Q$,求线段 $PQ$ 的长度。
    \begin{solution}
        由 (a) 知直线 $AB$ 的方程为 
        \[
        y = -\sqrt{3}x + \frac{2\sqrt{3}}{3},
        \]
        且椭圆 $C$ 方程为 
        \[
        \frac{3y^2}{16} + \frac{x^2}{4} = 1,
        \]
        且焦点坐标为 $F_1(0, -\frac{2\sqrt{3}}{3}), F_2(0, \frac{2\sqrt{3}}{3})$。设 $E(x_0, y_0)$,则切点弦 $AB$ 的方程为 
        \[
        \frac{3y_0}{16}y + \frac{x_0}{4}x = 1
        \]
        对照直线 $AB$ 的方程,对应项系数成比例,解得 
        \[
        x_0 = 6, y_0 = -\frac{8\sqrt{3}}{3}
        \]
        直线 $EF_1$ 的方程为 
        \[
        y = x + \frac{2\sqrt{3}}{3},
        \]
        令 $y = 0$ 得 $x = -2$,即 $P(-2, 0)$；直线 $EF_2$ 的方程为 
        \[
        y = \frac{5\sqrt{3}}{9}x - \frac{2\sqrt{3}}{3},
        \]
        令 $y = 0$ 得 $x = \frac{6}{5}$,即 $Q(\frac{6}{5}, 0)$,所以线段 $PQ$ 的长度为
        \[ 
        PQ = \frac{6}{5} - (-2) = \frac{16}{5} 
        \]
    \end{solution}
    \end{parts}

    \question 已知 $P(1,1)$ 为椭圆 $$\Gamma: \frac{x^2}{9} + \frac{y^2}{4} = 1$$ 内一点。过 $P$ 作椭圆的弦 ${A_1A_4},{A_2A_5},{A_3A_6}$,任意两相邻弦的夹角之一为 $\dfrac{\pi}{3}$。试求
    \[
    \frac{1}{PA_1 \cdot PA_4} + \frac{1}{PA_2 \cdot PA_5} + \frac{1}{PA_3 \cdot PA_6}
    \]
    的值。
    \begin{solution}
        设过点 $P(1,1)$ 的一条直线为 $L: y = m(x - 1) + 1$。  
        与椭圆相交于两点 $A_1(x_1, y_1)$ 与 $A_4(x_4, y_4)$,其中
        \[
        y_1 = mx_1 - m + 1,\quad y_4 = mx_4 - m + 1
        \]
        点 $P$ 到 $A_1$ 与 $A_4$ 的距离为
        \[
        PA_1 = \sqrt{(x_1 - 1)^2 + (mx_1 - m)^2}
         = |x_1 - 1|\sqrt{m^2 + 1},PA_4 = |x_4 - 1|\sqrt{m^2 + 1}
        \]
        因此
        \[
        PA_1 \cdot PA_4 = |x_1 - 1||x_4 - 1|(m^2 + 1)= |1 - (x_1 + x_4) + x_1 x_4|(m^2 + 1)
        \tag{1}
        \]
        将 $y = m(x - 1) + 1$ 代入椭圆方程
        \[
        \frac{x^2}{9} + \frac{(mx - m + 1)^2}{4} = 1
        \Rightarrow
        (9m^2 + 4)x^2 + 18m(1 - m)x + 9(1 - m)^2 - 36 = 0
        \]
        设该方程两根为 $x_1, x_4$,则
        \[
        x_1 + x_4 = \frac{-18m(1 - m)}{9m^2 + 4},\quad
        x_1 x_4 = \frac{9(1 - m)^2 - 36}{9m^2 + 4}
        \]
        代入 (1)得:
        \[
        PA_1 \cdot PA_4 = \left|1 + \frac{18m(m - 1)}{9m^2 + 4} + \frac{9(1 - m)^2 - 36}{9m^2 + 4}\right|(m^2 + 1)
        = \frac{23(m^2 + 1)}{9m^2 + 4}
        \]
        即
        \[
        \frac{1}{PA_1 \cdot PA_4} = \frac{9m^2 + 4}{23(m^2 + 1)} = \frac{9}{23} - \frac{5}{23} \cdot \frac{1}{m^2 + 1} = \frac{9}{23} - \frac{5}{23} \cos^2 \theta
        \]
        其中$m = \tan \theta,\theta$为${A_1A_4}$与$x$轴的夹角,且由恒等式,
        \[
        \cos^2 \theta + \cos^2\left(\theta + \frac{\pi}{3}\right) + \cos^2\left(\theta + \frac{2\pi}{3}\right) = \frac{3}{2}
        \]
        故
        \[
        \begin{aligned}
        &\frac{1}{PA_1 \cdot PA_4} + \frac{1}{PA_2 \cdot PA_5} + \frac{1}{PA_3 \cdot PA_6} \\
        &= \left( \frac{9}{23} - \frac{5}{23} \cos^2 \theta \right)
        + \left( \frac{9}{23} - \frac{5}{23} \cos^2\left(\theta + \frac{\pi}{3}\right) \right)
        + \left( \frac{9}{23} - \frac{5}{23} \cos^2\left(\theta + \frac{2\pi}{3}\right) \right) \\
        &= \frac{27}{23} - \frac{5}{23} \cdot \frac{3}{2} =\frac{39}{46}
        \end{aligned}
        \]
    \end{solution}

    \question 已知椭圆 $$\frac{x^2}{3} + \frac{y^2}{2} = 1$$ 的左、右焦点分别为 $F_1$ 与 $F_2$,过焦点 $F_1$ 的直线交椭圆于 $B, D$ 两点,过焦点 $F_2$ 的直线交椭圆于 $A, C$ 两点,且 $AC \perp BD$,垂足为点 $P$。求四边形 $ABCD$ 面积的最小值。
    \ifprintanswers
    \begin{figure}[H]
        \centering        
        \includegraphics[width=0.5\textwidth]{images/image107.jpg}
    \end{figure}
    \fi
    \begin{solution}
        椭圆$a = \sqrt{3}, b = \sqrt{2}, c = \sqrt{a^2-b^2}=1$,焦点$F_1(-1,0), F_2(1,0)$,设$BD,AC$直线方程为
        \[
        L_1: y = m(x+1), \quad L_2: y = -\frac{1}{m}(x-1)
        \]
        代入椭圆方程
        \[
        \frac{x^2}{3} + \frac{m^2(x+1)^2}{2} = 1, \quad
        \frac{x^2}{3} + \frac{(x-1)^2}{2 m^2} = 1
        \]
        解得
        \[
        x_b + x_d = -\frac{6 m^2}{3 m^2 + 2},\ x_b x_d = \frac{3 m^2 - 6}{3 m^2 + 2}, \quad
        x_a + x_c = \frac{6}{2 m^2 + 3},\ x_a x_c = \frac{3 - 6 m^2}{2 m^2 + 3}
        \]
        则
        \[
        (x_b - x_d)^2 = (x_b + x_d)^2 - 4 x_b x_d = \frac{48 (m^2 + 1)}{(3 m^2 + 2)^2}, \quad
        (x_a - x_c)^2 = \frac{48 m^2 (m^2 + 1)}{(2 m^2 + 3)^2}
        \]
        \[
        BD = \sqrt{m^2+1} |x_b-x_d| = \frac{4 \sqrt{3} (m^2+1)}{3 m^2 + 2}, \quad
        AC = \sqrt{\frac{m^2+1}{m^2}} |x_a - x_c| = \frac{4 \sqrt{3} (m^2+1)}{2 m^2 + 3}
        \]
        由AM-GM不等式,
        \[
        [ABCD] = \frac{1}{2} AC \cdot BD = \frac{24 (m^2+1)^2}{(3 m^2 + 2)(2 m^2 + 3)} \ge \frac{24 (m^2+1)^2}{\left( \frac{3 m^2 + 2 + 2 m^2 + 3}{2} \right)^2} = \frac{24 (m^2+1)^2}{\frac{25}{4} (m^2+1)^2} = \frac{96}{25}
        \] 
        故四边形$ABCD$ 面积的最小值为$\dfrac{96}{25}$
    \end{solution}

    \question 已知 $B_{1},B_{2}$ 分别为椭圆 $$C: \frac{x^{2}}{8}+\frac{y^{2}}{4}=1$$ 的下顶点、上顶点,过点 $A(0,-2\sqrt{2})$ 的直线 $l$ 交椭圆 $C$ 于 $P,Q$ 两点(异于点 $B_{1},B_{2}$);
    \begin{parts}
    \part 若$\tan \angle B_{1}PB_{2}=2 \tan \angle B_{1}QB_{2}$ ,求直线 $l$ 的方程;
    \begin{solution}
        椭圆下顶点$B_{1}(0,-2)$,上顶点$B_{2}(0,2)$,直线 $l$ 斜率存在,设 $l$ 方程为 
        \[
        y=kx-2\sqrt{2}
        \]
        与椭圆方程
        \[
        \frac{x^{2}}{8}+\frac{y^{2}}{4}=1 
        \]
        联立得 
        \[
        (1+2k^{2})x^{2}-8\sqrt{2}kx+8=0
        \]
        则 
        \[
        \triangle=(8\sqrt{2}k)^{2}-32(1+2k^{2})>0 \Rightarrow k^{2}>\frac{1}{2}
        \]
        设 $P(x_{1},y_{1}),Q(x_{2},y_{2})$,则 
        \[
        x_{1}+x_{2}=\frac{8\sqrt{2}k}{1+2k^{2}}, \ x_{1}x_{2}=\frac{8}{1+2k^{2}}
        \]
        由点 $P,Q$ 在椭圆上知
        \[
        y_{1}^{2}-4=-\frac{1}{2}x_{1}^{2},\ y_{2}^{2}-4=-\frac{1}{2}x_{2}^{2}
        \]
        于是
        \[
        \tan \angle B_{1}PB_{2}=\left|\frac{k_{PB_{1}}-k_{PB_{2}}}{1+k_{PB_{1}}k_{PB_{2}}}\right|=\left|\frac{\frac{y_{1}-(-2)}{x_{1}}-\frac{y_{1}-2}{x_{1}}}{1+\frac{y_{1}-(-2)}{x_{1}}\frac{y_{1}-2}{x_{1}}}\right|=\left|\frac{4x_{1}}{x_{1}^{2}+y_{1}^{2}-4}\right|=\left|\frac{4x_{1}}{x_{1}^{2}-\frac{1}{2}x_{1}^{2}}\right|=\left|\frac{8}{x_{1}}\right|
        \]
        同理$\tan \angle B_{1}QB_{2}=\dfrac{8}{|x_{2}|}$;由 $\tan \angle B_{1}PB_{2}=2 \tan \angle B_{1}QB_{2}$ ,得 $$\frac{8}{|x_{1}|}=2\cdot\frac{8}{|x_{2}|}$$
        又 $x_{1}x_{2}=\dfrac{8}{1+2k^{2}}>0$,于是 
        \[
        x_{2}=2x_{1}\Rightarrow x_{1}+x_{2}=3x_1=\frac{8\sqrt{2}k}{1+2k^{2}}\Rightarrow x_{1}=\frac{8\sqrt{2}k}{3(1+2k^{2})}
        \]
        且
        \[
        x_{2}=2x_{1}\Rightarrow x_{1}x_{2}=2x_1^2=\frac{16\sqrt{2}k}{3(1+2k^{2})}\Rightarrow x_{1}=\frac{8}{1+2k^{2}}
        \]
        故
        \[
        32k^{2}=9(1+2k^{2})\Rightarrow k^{2}=\frac{9}{14}>\frac{1}{2}
        \]
        因此直线 $l$ 方程为 
        \[
        y=\frac{3\sqrt{14}}{14}x-2\sqrt{2}
        \]
    \end{solution}
    \part 设 $R$ 为直线 $B_{1}P$ 、 $B_{2}Q$ 的交点,求 ${AR}\cdot{B_{1}B_{2}}$ 的值.
    \begin{solution}
        直线 $B_{1}P,B_{2}Q$ 方程分别为 
        \[
        y=\frac{y_{1}+2}{x_{1}}x-2,\quad y=\frac{y_{2}-2}{x_{2}}x+2
        \]
        两式联立得 
        \[
        \frac{y+2}{y-2}=\frac{(y_{1}+2)x_{2}}{(y_{2}-2)x_{1}}=\frac{(kx_{1}-2\sqrt{2}+2)x_{2}}{(kx_{2}-2\sqrt{2}-2)x_{1}}
        \]
        因为 
        \begin{align*}
        &(kx_{1}-2\sqrt{2}+2)x_{2}-(-3+2\sqrt{2})(kx_{2}-2\sqrt{2}-2)x_{1}\\
        &=(4-2\sqrt{2})kx_{1}x_{2}+(-2\sqrt{2}+2)(x_{1}+x_{2})\\
        &=(4-2\sqrt{2})k\cdot\frac{8}{1+2k^{2}}+(-2\sqrt{2}+2)\cdot\frac{8\sqrt{2}k}{1+2k^{2}}=0
        \end{align*}
        所以
        \[
        \frac{y+2}{y-2}=-3+2\sqrt{2}\Rightarrow y=-\sqrt{2}
        \]
        即点$R$在定直线 $y=-\sqrt{2}$ 上;设 $R(t,-\sqrt{2}),$则 ${AR}=(t,\sqrt{2})$,又 ${B_{1}B_{2}}=(0,4),$ 因此 $${AR}\cdot{B_{1}B_{2}}=4\sqrt{2}$$
    \end{solution}
    \end{parts}

    \question 双曲线 
    \[
    \frac{x^{2}}{9}-\frac{y^{2}}{16}=1
    \] 
    的焦点为 \(F_1,F_2\)。设第一象限点 \(P\) 满足 \(PF_1:PF_2=1:3\)。求 \(\triangle F_1PF_2\) 周长。  
    \begin{solution}
        双曲线焦点$F_1=(-5,0),\ F_2=(5,0)$,已知$PF_1:PF_2=1:3$,由双曲线定义,
        \[
        PF_2 - PF_1 = 2a \Rightarrow 3PF_1 - PF_1 = 6 \Rightarrow PF_1 = 3
        \]
        \(\triangle F_1PF_2\)周长为
        \[
        PF_1 + PF_2 + F_1F_2 = 3 + 9 + 10 = 22
        \]
    \end{solution}
    
    \question 已知双曲线焦点 \((-2,-2),(8,-2)\),且一条渐近线斜率为 \(-\dfrac{4}{3}\)。求其标准方程。  
    \begin{solution}
        双曲线中心$C(3, -2)$,焦距$c = 5$,又已知渐近线斜率为 $\pm \dfrac{b}{a} = \pm \dfrac{4}{3}$, 有 
        $\dfrac{b^2}{a^2} = \dfrac{16}{9}$,
        
        由 $c^2 = a^2 + b^2$ 得
        \[
        25 = a^2 + b^2 = a^2 + \frac{16}{9}a^2 = \frac{25}{9}a^2
        \Rightarrow a^2 = 9,\ b^2 = 16
        \]
        所以标准方程为
        \[
        \dfrac{(x - 3)^2}{9} - \dfrac{(y + 2)^2}{16} = 1
        \]
    \end{solution}

    \question 已知双曲线 $$\Gamma: \dfrac{x^2}{4}-\dfrac{y^2}{5}=1,$$过焦点 $F$ 作一焦弦 $PQ$,已知 $PQ$ 的斜率为 $1$,求 $PF' + QF'$。 
    \begin{solution}
        双曲线$c=\sqrt{4+5}=3$,焦点$F(3,0), F'(-3,0)$,过 $F$ 且斜率为 $1$ 的直线为
        \[
        L: y=x-3,
        \]
        设交点为 $P(x_1,x_1-3),\;Q(x_2,x_2-3)$,代入双曲线方程解得
        \[
        \frac{x^2}{4}-\frac{(x-3)^2}{5}=1 \Rightarrow x=-12\pm 10\sqrt{2}
        \]
        因此
        \[
        |x_1-x_2|=20\sqrt{2},\quad PQ=\sqrt{1^2+(-1)^2}\,|x_1-x_2|=\sqrt{2}\cdot 20\sqrt{2}=40
        \]
        由双曲线定义,
        \[
        PF'-PF=2a=4,\quad QF'-QF=2a=4
        \]
        两式相加得
        \[
        PF'+QF'=(PF+QF)+8=40+8=48
        \]
    \end{solution}

    \question 已知双曲线焦点为 \(F_1(0,2), F_2(0,-2)\),实轴长为 2,点 \(Q\) 在一条渐近线上,且 \(\angle F_1QF_2 = 90^\circ\),求 \(\triangle QF_1F_2\) 面积。

    \begin{solution}
        已知 \(c = 2,\; a = 1 \Rightarrow b^2 = c^2 - a^2 = 3\),  
        双曲线方程为\(y^2 - \dfrac{x^2}{3} = 1\),渐近线为\(y = \pm \dfrac{\sqrt{3}}{3}x\)。
        
        设 \(Q(x, \dfrac{\sqrt{3}}{3}x)\),由 \(\angle F_1QF_2 = 90^\circ\),
        \[
        m_{QF_1}\cdot m_{QF_2}=\frac{\frac{\sqrt{3}}{3}x-2}{x-0}\cdot \frac{\frac{\sqrt{3}}{3}x+2}{x-0}=-1 \Rightarrow x = \pm\sqrt{3}
        \]
        \(\triangle QF_1F_2\)面积为
        \[
        \frac{1}{2} \cdot 4 \cdot \sqrt{3} = 2\sqrt{3}
        \]
    \end{solution}
    
    \question 已知双曲线 $\Gamma: xy = k,\; k < 0$,点 $P(2,2)$,过 $P$ 作 $\Gamma$ 的两条切线,切点为 $A$、$B$,若三角形 $\triangle PAB$ 是正三角形,求 $k$。
    \begin{solution}
        由 $\Gamma: xy = k < 0$ 可知其对称轴为 $y = x$,若 $\triangle PAB$ 为正三角形,且 $P$ 在 $y = x$ 上,故切线 ${PA}$ 的斜率为 
        \[
        \tan 75^\circ = \tan(45^\circ + 30^\circ) = \frac{1 + \frac{1}{\sqrt{3}}}{1 - \frac{1}{\sqrt{3}}} = 2 + \sqrt{3}
        \]
        切线方程为
        \[
        y = (2 + \sqrt{3})(x - 2) + 2
        \]
        将其代入 $xy = k$ 整理得
        \[
        (2 + \sqrt{3})x^2 + (-2 - 2\sqrt{3})x - k = 0
        \]
        令判别式为 0,
        \[
        (2 + 2\sqrt{3})^2 + 4k(2 + \sqrt{3}) = 0
        \]
        解得
        \[
        k = -2
        \]
    \end{solution}

    \question 已知双曲线 $$\frac{x^{2}}{16}-\frac{y^{2}}{12}=1$$ 上有两点 $P, Q,F_1, F_2$ 为其焦点,且 $PQ$ 过 $F_2,\angle F_2PF_1 = 60^\circ$,求 $\triangle PQF_1$ 的周长。
    \ifprintanswers
    \begin{figure}[H]
        \centering        
        \includegraphics[width=0.5\textwidth]{images/image117.jpg}
    \end{figure}
    \fi
    \begin{solution}
        由双曲线方程得$c = \sqrt{16+12} = 2\sqrt{7}$,焦点坐标
        \[
        F_1(-2\sqrt{7},0), \quad F_2(2\sqrt{7},0).
        \]
        设$PF_2 = m, QF_2 = n$,则
        \[
        PF_1 = m + 2a = m+8, \quad QF_1 = n+8.
        \]
        在 $\triangle PF_1F_2$中,由余弦定理,
        \[
        (4\sqrt{7})^2=(m+8)^2 + m^2-2 m (m+8)\cos 60^\circ \Rightarrow m = 4
        \]
        在 $\triangle PF_1F_2$中,由余弦定理,
        \[
        (n+8)2=(m+8)^2+(m+n)^2-2(m+8)(m+n)\cos 60^\circ \Rightarrow n = \frac{12}{5}
        \]
        故 $\triangle PQF_1$ 周长
        \[
        (m+8) + m + n + (n+8)= \frac{144}{5}
        \]
    \end{solution}

    \question 双曲线$$\frac{x^2}{16} - \frac{y^2}{9} = 1$$的左、右焦点分别为$F_1, F_2,P$是双曲线上一点, 若$\triangle PF_1F_2$的内切圆圆心为$(4,2)$, 求$\triangle PF_1F_2$外接圆的半径。
    \begin{solution}
        双曲线焦点分别为$F_1(-5,0), F_2(5,0)$,且已知$I(4,2)$,设$\angle IF_1F_2 = \alpha$, $\angle IF_2F_1 = \beta$,
        则 $$\tan \alpha = \frac{2}{9},\; \tan \beta = 2$$
        故 $$\tan \angle F_1PF_2 =\tan(2\alpha+2\beta)=\frac{2\tan(\alpha+\beta)}{1-\tan^2(\alpha+\beta)} = -\frac{8}{15}$$ 
        其中$$\tan(\alpha + \beta) = \frac{\frac29+2}{1-\frac29\cdot2}=4$$
        故$$\sin \angle F_1PF_2 = \frac{8}{17}$$
        $\therefore \triangle PF_1F_2$外接圆的半径为$$\frac{10}{2\sin \angle F_1PF_2} = \frac{85}{8}$$.
    \end{solution} 

    \question 设双曲线 
    \[
    \Gamma: \frac{x^2}{3} - \frac{y^2}{2} = 1
    \]
    两焦点为 $F_1, F_2$,点 $P(3,-2)$ 在 $\Gamma$ 上。以 $P$ 点为切点的切线交 $x$ 轴于点 $Q$,求 $\tan \angle FPQ$。
    \ifprintanswers
    \begin{figure}[H]
        \centering
        \includegraphics[width=0.5\linewidth]{images/image50.png}
    \end{figure}
    \fi
    \begin{solution}
        由双曲线知$a = \sqrt{3}, \quad b = \sqrt{2}, c = \sqrt{a^2 + b^2} = \sqrt{5}$,焦点为
        \[
        F_1(\sqrt{5}, 0), \quad F_2(-\sqrt{5}, 0).
        \]
        在$P(3, -2)$的切线斜率为
        \[
        \frac{2}{3} x - y y' = 0 \Rightarrow y'|_{x=3,y=-2} = \frac{2x}{3y}= \frac{6}{-6} = -1.
        \]
        过$P$切线方程为
        \[
        y + 2 = -1(x - 3) \Rightarrow x + y = 1
        \]
        与 $x$ 轴交与
        \[
        Q(1, 0).
        \]
        在 $\triangle PAQ$ 及$\triangle PAF_1$,设$\theta_1 = \angle F_1 P A,\theta_2 = \angle F_1 P Q$
        \[
        \tan \theta_1 = \frac{3 - \sqrt{5}}{2}, \quad \tan(\theta_1 + \theta_2) = 1,
        \]
        利用和角公式,解得
        \[
        \frac{\tan \theta_2 + \frac{3 - \sqrt{5}}{2}}{1 - \tan \theta_2 \cdot \frac{3 - \sqrt{5}}{2}} = 1 \Rightarrow \tan \theta_2 = \frac{-1 + \sqrt{5}}{5 - \sqrt{5}} = \frac{\sqrt{5}}{5}
        \]
    \end{solution}
    
    \question 已知双曲线与椭圆的方程分别为
    \[
    \frac{x^{2}}{a^{2}}-\frac{y^{2}}{b^{2}}=1,
    \quad
    \frac{x^{2}}{a^{2}}+\frac{y^{2}}{b^{2}}=1,
    \]
    其中 $a>b>0$。在双曲线上取一坐标均为正的点作切线,该切线经过椭圆的正 $x$ 轴焦点。证明该切线的斜率恒为 $1$。
    \begin{solution}
        设双曲线上一点$P(a\sec\theta,b\tan\theta)$,其中 $\theta\in\left(0,\dfrac{\pi}{2}\right)$,对双曲线方程求导得
        \[
        \frac{2x}{a^{2}}-\frac{2y}{b^{2}}\frac{dy}{dx}=0 \Rightarrow \frac{dy}{dx}=\frac{b^{2}x}{a^{2}y} 
        \]
        则在点$P$处的切线斜率为$\dfrac{b^{2}a\sec\theta}{a^{2}b\tan\theta}=\dfrac{b}{a\sin\theta}$,因此切线方程为
        \[
        y-b\tan\theta=\frac{b}{a\sin\theta}(x-a\sec\theta).
        \]
        该切线经过椭圆的正 $x$ 轴焦点 $(ae,0)$,其中 $e$ 为椭圆的离心率,代入得
        \[
        0-b\tan\theta=\frac{b}{a\sin\theta}(ae-a\sec\theta) \Rightarrow e = \cos\theta
        \]
        而椭圆的离心率满足$b^{2}=a^{2}(1-e^{2})$,代入 $e=\cos\theta$,得
        \[
        \frac{b^{2}}{a^{2}}=1-\cos^{2}\theta=\sin^{2}\theta \Rightarrow \frac{b}{a}=\sin\theta
        \]
        因此得证切线斜率为
        \[
        \frac{b}{a\sin\theta}=\sin\theta\cdot\frac{1}{\sin\theta}=1
        \]
    \end{solution}

    \question 已知双曲线 
    \[
    H: \dfrac{x^2}{a^2} - \dfrac{y^2}{b^2} = 1\ ,a > 0,\ b > 0,
    \]
    $l$ 为双曲线 $H$ 的一条渐近线,$\ F_1, F_2$ 是双曲线 $H$ 的左、右焦点。若 $F_1$ 关于直线 $l$ 的对称点在圆 
    \[
    C:(x - c)^2 + y^2 = c^2
    \]
    上,其中 $c$ 为双曲线的半焦距,求双曲线 $H$ 的离心率。
    \begin{solution}
        双曲线左焦点为 $F_1 = (-c,\ 0)$,其中 $c^2 = a^2 + b^2$。  
        渐近线 $l$ 的一个方程为 $y = \dfrac{b}{a}x$,即 
        $$bx - ay = 0$$
        设 $F_1$ 关于直线 $l$ 的对称点为 $F_1' = (x',\ y'),F_1F_1'$ 的中点为 $\left(\dfrac{x'-c}{2},\ \dfrac{y'}{2} \right)$,
        \[
        b\cdot \dfrac{x'-c}{2} - a\cdot \dfrac{y'}{2} = 0 \Rightarrow b(x'-c) - a y' = 0 \tag{1}
        \]
        且
        \[
        m_{F_1F_1'}=m_{l}=\dfrac{y'}{x'+c} \cdot \dfrac{b}{a} = -1 \Rightarrow b y' = -a(x'+c) \tag{2}
        \]
        由$(1),(2)$得
        \[
        F_1' = \left( \dfrac{(b^2 - a^2)c}{a^2 + b^2},\ -\dfrac{2ab c}{a^2 + b^2} \right)
        \]
        将其代入圆$C$:
        \begin{align*}
        \left( \dfrac{(b^2 - a^2)c}{a^2 + b^2} - c \right)^2 + \left( \dfrac{2ab c}{a^2 + b^2} \right)^2 &= c^2 
        \end{align*}
        可得
        \[
        3a^2 = b^2 \Rightarrow c^2 = a^2 + b^2 = 4a^2 \Rightarrow c = 2a
        \]
        所以双曲线离心率为
        \[
        e = \frac{c}{a} = 2
        \]
    \end{solution}

    \question 已知双曲线 
    \[
    C: \frac{x^{2}}{a^{2}} - \frac{y^{2}}{b^{2}} = 1 \; (a>0,\; b>0)
    \]
    与椭圆
    \[
    \frac{x^{2}}{4} + \frac{y^{2}}{3} = 1,
    \]
    过椭圆上一点 $P\left(-1,\dfrac{3}{2}\right)$ 作椭圆的切线 $l$与 $x$ 轴交于 $M$ 点,且$l$ 与双曲线 $C$ 的两条渐近线分别交于 $N,Q$,若 $N$ 为 $MQ$ 的中点,求双曲线 $C$ 的离心率。
    \begin{solution}
        设过$P(-1,\dfrac{3}{2})$的切线$l$方程为
        \[
        y - \frac{3}{2} = k(x + 1)
        \Rightarrow y = kx + k + \frac{3}{2}.
        \]
        将此代入椭圆方程
        \[
        \frac{x^{2}}{4} + \frac{y^{2}}{3} = 1.
        \]
        整理得
        \[
        (4k^2+3)x^2+(8k^2+12k)x+4k^2+12k-3=0
        \]
        令判别式为零得
        \[
        \Delta = (8k^2+12k)^2-4(4k^2+3)(4k^2+12k-3)=0\Rightarrow k = \frac{1}{2}.
        \]
        所以切线 $l$ 的方程为$y = \dfrac{1}{2}x + 2$,令$y = 0$得 $M(-4, 0)$
        
        双曲线的渐近线方程为$y = \pm \dfrac{b}{a}x$,联立 $y = \dfrac{b}{a}x$ 与 $y = \dfrac{1}{2}x + 2$,得 $x_Q = \dfrac{4a}{2b - a}$,
        
        联立 $y = -\dfrac{b}{a}x$ 与 $y = \dfrac{1}{2}x + 2$,得
        $x_N = -\dfrac{4a}{2b + a}$,
        
        由于 $N$ 是 $MQ$ 的中点,
        \[
        -\frac{4a}{2b + a} = \frac{1}{2} \left( \frac{4a}{2b - a} - 4 \right) \Rightarrow\frac{b}{a} = \frac{3}{2}
        \]
        故双曲线的离心率为
        \[
        e = \sqrt{1 + \left( \frac{3}{2} \right)^2 } = \frac{\sqrt{13}}{2}
        \]
    \end{solution}

    \question 已知双曲线 
    \[
    C:\dfrac{x^{2}}{a^{2}} - \dfrac{y^{2}}{b^{2}} = 1\ (a>0,\ b>0)
    \]
    的左、右顶点分别是 $A_1,A_2$,  圆 $x^2 + y^2 = a^2$ 与 $C$ 的渐近线在第一象限的交点为 $M$,直线 $A_1M$ 交 $C$ 的右支于点 $P$。若 $\triangle MPA_2$ 是等腰三角形,且 $\angle PA_2M$ 的内角平分线与 $y$ 轴平行,求双曲线的离心率。
    \begin{solution}
        双曲线顶点 $A_1(-a, 0), A_2(a, 0)$,渐近线为 $y = \dfrac{b}{a}x$,圆为 $x^2 + y^2 = a^2$,联立得
        \[
        x^2\left(1 + \frac{b^2}{a^2}\right) = a^2 \Rightarrow x = \frac{a^2}{c}, \quad y = \frac{ab}{c}
        \]
        故点 $M\left( \dfrac{a^2}{c}, \dfrac{ab}{c} \right)$,其中 $c = \sqrt{a^2 + b^2}$ 为半焦距,又
        \[
        MA_1^2 = \left(\frac{a^2}{c} + a\right)^2 + \left( \frac{ab}{c} \right)^2 = 2a^2\left(1 + \frac{a}{c}\right)
        \]
        \[
        MA_2^2 = \left(\frac{a^2}{c} - a\right)^2 + \left( \frac{ab}{c} \right)^2 = 2a^2\left(1 - \frac{a}{c}\right)
        \]
        由题设$\angle A_{1}MA_{2}=\angle PMA_{2}=90^{\circ},\triangle MPA_{2}$是等腰直角三角形,$\;\angle MA_{2}P=45^{\circ}$
        
        $\angle PA_{2}M$的内角平分线与$y$轴平行,所以$\angle MA_{1}A_{2}=22.5^{\circ}$
        
        又$\tan 45^{\circ}=\dfrac{2\tan 22.5^{\circ}}{1-\tan^{2}22.5^{\circ}}=1$可得$\tan 22.5^{\circ}=\sqrt{2}-1$,所以
        \[
        \tan^{2}\angle MA_{1}A_{2}=\left(\frac{MA_{2}}{MA_{1}}\right)^{2}=\frac{1-\frac{a}{c}}{1+\frac{a}{c}}=(\sqrt{2}-1)^{2}
        \]
        解得
        \[
        \frac{e-1}{e+1}=3-2\sqrt{2}\Rightarrow e=\sqrt{2}
        \]
    \end{solution}

    \question 已知双曲线的两个焦点分别为 $F_1, F_2$,若过点 $F_1$ 的直线与双曲线交于 $A, B$ 两点,且 $\angle AF_1F_2 = 30^\circ,F_2A = F_2B$,求双曲线的离心率。
    \begin{solution}
        不妨 $A$ 在双曲线靠近 $F_1$ 的这一支上, 取 $AB$ 中点 $C$, 则 $F_2C \perp AB$。
        
        设 $AF_1=x$, 由双曲线定义知
        \[
        AF_2 = BF_2= x+2a, \;BF_1=x+4a
        \]
        从而 
        \[
        AB=4a,\;AC=2a
        \]
        由 $\angle AF_1F_2=30^\circ$ 知 $CF_1=\dfrac{\sqrt{3}}{2}F_1F_2$,可得
        \[
        x=\sqrt{3}c-2a\tag{1}
        \]
        又由 $CF_2=\dfrac{\sqrt3}{3}CF_1$ 及 $AC^2+CF_2^2=AF_2^2$,可得 
        \[
        (2a)^2 + \frac{(x+2a)^2}{3}=(x+2a)^2 \Rightarrow x=(\sqrt{6}-2)a \tag{2}
        \]
        由$(1),(2)$解得离心率 
        \[
        e=\frac{c}{a}=\sqrt{2}
        \]
    \end{solution}

    \question 双曲线的中心为原点 $O$,焦点在 $x$ 轴上,两条渐近线分别为 $l_1, l_2$,过右焦点 $F$ 垂直于 $l_1$ 的直线分别交 $l_1, l_2$ 于 $A, B$ 两点。已知 $|\overrightarrow{OA}|, |\overrightarrow{AB}|, |\overrightarrow{OB}|$ 成等差数列,且 $\overrightarrow{BF}$ 与 $\overrightarrow{FA}$ 同向。
    \begin{parts}
    \part 求双曲线的离心率；
    \begin{solution}
        已知$OA,AB,OB$成等差数列,设 $OA = m-d,AB = m,OB = m+d$,由毕氏定理,
        \[
        (m-d)^2 + m^2 = (m+d)^2 \Rightarrow d = \frac{1}{4} m
        \]
        设双曲线半轴为 $a, b$,右焦点 $F(c,0)$,则
        \[
        \tan \angle AOF = \frac{b}{a}, \quad \tan \angle AOB = \tan (2 \ \angle AOF) = \frac{AB}{OA} = \frac{4}{3}
        \]
        由倍角公式 $\tan 2\theta = \dfrac{2\tan\theta}{1-(\tan\theta)^2}$得
        \[
        \frac{\frac{2b}{a}}{1-\frac{b^2}{a^2}} = \frac{4}{3} \Rightarrow \frac{b}{a} = \frac{1}{2}
        \]
        故离心率为 
        \[
        e = \sqrt{1 + \left(\frac{b}{a}\right)^2} = \frac{\sqrt{5}}{2}
        \]
    \end{solution}
    \part 若 $AB$ 被双曲线所截得的线段的长为 $4$,求双曲线的方程。
    \begin{solution}
        过 $F$ 直线方程为 
        \[
        y = -\frac{a}{b}(x-c)
        \]
        与双曲线方程 
        \[
        \frac{x^2}{a^2} - \frac{y^2}{b^2} = 1
        \]
        联立,将 $a = 2b,c = \sqrt{5}b$ 代入,化简得
        \[
        \frac{15}{4b^2}x^2 - \frac{8\sqrt{5}}{b}x + 21 = 0
        \]
        由韦达定理,
        \begin{align*}
        4 = \sqrt{1+\left(\frac{a}{b}\right)^2} |x_1-x_2| 
        &= \sqrt{1+\left(\frac{a}{b}\right)^2} \sqrt{(x_1+x_2)^2 - 4x_1x_2} \\
        &= \sqrt{5} \sqrt{\left(\frac{32\sqrt{5}b}{15}\right)^2 - 4 \cdot \frac{28b^2}{5}}
        \end{align*}
        解得 
        \[
        b = 3
        \]
        故双曲线方程为
        \[
        \frac{x^2}{36} - \frac{y^2}{9} = 1
        \]
    \end{solution}
    \end{parts}

    \question 已知双曲线
    \[
    C:\frac{x^2}{a^2}-\frac{y^2}{b^2}=1 \quad a>0,b>0
    \]
    的两个焦点为 $F_1(-2,0),F_2(2,0)$,点 $P(3,\sqrt7)$ 在曲线 $C$ 上。
    \begin{parts}
    \part 求双曲线 $C$ 的方程；
    \begin{solution}
        据题意,由$a^2+b^2=4$,双曲线方程为
        \[
        \frac{x^2}{a^2}-\frac{y^2}{4-a^2}=1\quad 0<a^2<4
        \]
        将点 $(3,\sqrt7)$ 代入得
        \[
        \frac{9}{a^2}-\frac{7}{4-a^2}=1 \Rightarrow a^2=18\ (\text{舍去})\quad\text{或}\quad a^2=2
        \]
        故所求双曲线方程为
        \[
        \frac{x^2}{2}-\frac{y^2}{2}=1
        \]
    \end{solution}
    \begin{solution}
        依题意得双曲线的半焦距 $c=2$。又
        \[
        2a=PF_1-PF_2
        =\sqrt{(3+2)^2+(\sqrt7)^2}-\sqrt{(3-2)^2+(\sqrt7)^2}
        =2\sqrt2
        \]
        故
        \[
        a^2=2,\quad b^2=c^2-a^2=2
        \]
        因此双曲线 $C$ 的方程为
        \[
        \frac{x^2}{2}-\frac{y^2}{2}=1
        \]
    \end{solution}
    \part 记 $O$ 为坐标原点,过点 $Q(0,2)$ 的直线与双曲线 $C$ 相交于不同的两点 $E,F$,若 $\triangle OEF$ 的面积为 $2\sqrt2$,求直线的方程。
    \begin{solution}
        设直线的方程为$y=kx+2$,代入双曲线 $C$ 整理,
        \[
        (1-k^2)x^2-4kx-6=0
        \]
        故直线 $l$ 与双曲线 $C$ 相交于不同的两点 $E,F$,需满足
        \[
        \begin{cases}
        1-k^2\neq0,\\
        \Delta=(-4k)^2+4\cdot 6(1-k^2)>0,
        \end{cases}
        \]
        即$k\in(-\sqrt3,-1)\cup(1,\sqrt3)$,设 $E(x_1,y_1),F(x_2,y_2)$,由韦达定理,
        \[
        x_1+x_2=\frac{4k}{1-k^2},\quad x_1x_2=\frac{6}{1-k^2}.
        \]
        于是
        \begin{align*}
        EF&=\sqrt{(x_1-x_2)^2+(y_1-y_2)^2}\\
        &=\sqrt{1+k^2}\,|x_1-x_2|\\
        &=\sqrt{1+k^2}\sqrt{(x_1+x_2)^2-4x_1x_2}\\
        &=\sqrt{1+k^2}\cdot\frac{2\sqrt2\sqrt{3-k^2}}{|1-k^2|}
        \end{align*}
        原点 $O$ 到直线的距离为
        \[
        d=\frac{2}{\sqrt{1+k^2}},
        \]
        故
        \[
        [\triangle OEF]
        =\frac12 d\cdot EF
        =\frac{2\sqrt2\sqrt{3-k^2}}{|1-k^2|}=2\sqrt2 \Rightarrow k^4-k^2-2=0
        \]
        解得
        \[
        k=\pm\sqrt2 \in(-\sqrt3,-1)\cup(1,\sqrt3)
        \]
        故所求直线方程为
        \[
        y=\sqrt2x+2,\quad y=-\sqrt2x+2.
        \]
    \end{solution}
    \end{parts}

    \question 点$P\left(p+\dfrac{1}{p},\,p-\dfrac{1}{p}\right),\ p\neq 0$在直角双曲线
    \[
    x^2-y^2=4
    \]
    上。曲线在 $P$ 处的法线与 $y$ 轴交于点 $Q(0,-k)$,其中 $k>0$。已知三角形 $OPQ$ 的面积为 $\dfrac{15}{4}$,其中 $O$ 为原点,求点 $P$ 的坐标。
    \begin{solution}
        曲线可参数表示为
        \[
        x=p+\frac{1}{p},\quad y=p-\frac{1}{p}
        \]
        由链导法,
        \[
        \frac{dy}{dx}=\frac{dy}{dp} \div \frac{dx}{dp}
        =\frac{1+\frac{1}{p^2}}{1-\frac{1}{p^2}}
        =\frac{p^2+1}{p^2-1}
        \]
        因此法线斜率为$-\dfrac{p^2-1}{p^2+1}$,点 $P$ 处法线方程为
        \[
        y-\left(p-\frac{1}{p}\right)
        =-\frac{p^2-1}{p^2+1}\left(x-\left(p+\frac{1}{p}\right)\right)
        \]
        该法线经过点 $Q(0,-k)$,代入得
        \[
        -k-p+\frac{1}{p}=\frac{p^2-1}{p^2+1}\left(p+\frac{1}{p}\right) \Rightarrow k=-\frac{2(p^2-1)}{p} \tag{1}
        \]
        已知 $\triangle OPQ$ 的面积为 $\dfrac{15}{4}$,有
        \[
        \frac{1}{2}\,k\left(p+\frac{1}{p}\right)=\frac{15}{4} \tag{2}
        \]
        联立$(1),(2)$解得
        \[
        (4p^2-1)(p^2+4)=0 \Rightarrow p=\pm\frac{1}{2}
        \]
        因此点 $P$ 的可能坐标为
        \[
        \left(\frac{5}{2},-\frac{3}{2}\right),        \left(-\frac{5}{2},\frac{3}{2}\right)
        \]
    \end{solution}

    \question 已知中心在原点的双曲线 $C$ 的一个焦点为 $F_{1}(-3,0)$,一条渐近线的方程是$\sqrt{5}x-2y=0$,
    \begin{parts}
    \part 求双曲线 $C$ 的方程；
    \begin{solution}
        设双曲线 $C$ 的方程为
        \[
        \frac{x^{2}}{a^{2}}-\frac{y^{2}}{b^{2}}=1 
        \]
        其中$a>0,b>0$,由焦点坐标得$c=\sqrt{a^2+b^2}=3$,又由渐近线方程$y=\dfrac{\sqrt{5}}{2}x$得
        \[
        \frac{b}{a}=\frac{\sqrt{5}}{2}
        \]
        于是可解得
        \[
        a^{2}=4,\quad b^{2}=5
        \]
        故双曲线 $C$ 的方程为
        \[
        \frac{x^{2}}{4}-\frac{y^{2}}{5}=1
        \]
    \end{solution}
    \part 若以 $k\ne 0$ 为斜率的直线 $l$ 与双曲线 $C$ 相交于两个不同的点$M,N$,且线段 $MN$ 的垂直平分线与两坐标轴围成的三角形面积为$\dfrac{81}{2}$,求 $k$ 的取值范围。
    \begin{solution}
        设直线 $l$ 的方程为
        \[
        y=kx+m, \quad k\ne 0,
        \]
        与双曲线交于 $M(x_{1},y_{1}),N(x_{2},y_{2})$,将 $y=kx+m$ 代入双曲线方程得
        \[
        \frac{x^{2}}{4}-\frac{(kx+m)^{2}}{5}=1,
        \]
        化简为
        \[
        (5-4k^{2})x^{2}-8kmx-4m^{2}-20=0.
        \]
        该方程有两个不等实根,故
        \[
        5-4k^{2}\ne 0,
        \]
        且判别式
        \[
        \Delta>0 \;\Rightarrow\; m^{2}+5-4k^{2}>0. \tag{1}
        \]
        由韦达定理,线段 $MN$ 的中点
        \[
        (x_{0},y_{0})=\left(\frac{4km}{5-4k^{2}},\ \frac{5m}{5-4k^{2}}\right).
        \]
        因此$MN$ 的垂直平分线与 $x$ 轴、$y$ 轴的截距分别为
        \[
        \left(\frac{9km}{5-4k^{2}},0\right),\quad
        \left(0,\frac{9m}{5-4k^{2}}\right).
        \]
        由题设三角形面积得
        \[
        \frac{1}{2}\left|\frac{9km}{5-4k^{2}}\right|
        \left|\frac{9m}{5-4k^{2}}\right|
        =\frac{81}{2},
        \]
        化简得
        \[
        m^{2}=\frac{(5-4k^{2})^{2}}{|k|},\quad k\ne 0. \quad \tag{2}
        \]
        将 (2) 代入 (1),整理为
        \[
        (4k^{2}-5)(4k^{2}-|k|-5)>0,\quad k\ne 0.
        \]
        解得
        \[
        0<|k|<\frac{\sqrt{5}}{2}
        \quad \text{或} \quad
        |k|>\frac{5}{4}.
        \]
        因此$k$ 的取值范围为
        \[
        k\in\left(-\infty,-\frac{5}{4}\right)\cup\left(-\frac{\sqrt{5}}{2},0\right)\cup
        \left(0,\frac{\sqrt{5}}{2}\right)\cup\left(\frac{5}{4},+\infty\right)
        \]
    \end{solution}
    \end{parts}

    \question 已知双曲线
    \[
    E: \frac{x^{2}}{a^{2}} - \frac{y^{2}}{b^{2}} = 1
    \]
    的左、右顶点分别为 $A, B$,点 $M$ 在 $E$ 上使得 $\triangle ABM$ 是等腰三角形,且$\triangle ABM$外接圆面积为 $3 \pi a^{2}$,求双曲线的离心率。
    \begin{solution}
        不妨设$M$在第二象限,则在等腰$\triangle ABM$中, $AB=AM=2a$,
        
        设 $\angle ABM=\angle AMB=\theta$, 则 $\angle FAM=2\theta$,又$\triangle ABM$外接圆面积 $3\pi a^{2}$, 则半径为 $\sqrt{3}a$.
        
        由正弦定理,
        \[
        2R = \frac{AB}{\sin\theta} \Rightarrow 2\sqrt{3}a = \frac{2a}{\sin\theta} \Rightarrow \sin\theta = \frac{\sqrt{3}}{3}
        \]
        于是
        \[
        \cos\theta = \frac{\sqrt{6}}{3},\;\sin 2\theta = \frac{2\sqrt{2}}{3},\;\cos 2\theta = \frac{1}{3}
        \]
        设$M$点坐标为 $(x,y)$, 则 
        \[
        x = -a - AM\cos 2\theta = -a - 2a \cdot \frac{1}{3} = -\frac{5a}{3}
        \]
        \[
        y = AM\sin 2\theta = 2a \cdot \frac{2\sqrt{2}}{3} = \frac{4\sqrt{2}a}{3}
        \]
        即$M$点坐标为 $\left(-\dfrac{5a}{3},\dfrac{4\sqrt{2}a}{3}\right)$,由$M$点在双曲线上得 
        \[
        \frac{\left(-\frac{5a}{3}\right)^{2}}{a^{2}}-\frac{\left(\frac{4\sqrt{2}a}{3}\right)^{2}}{b^{2}}=1
        \Rightarrow b^2 = 2a^2
        \]
        故
        \[
        e = \sqrt{1+\dfrac{b^{2}}{a^{2}}}  =\sqrt{3}
        \]
    \end{solution}  
        
    \question 双曲线 
    \[
    \Gamma: x^{2}-y^{2}=1
    \] 
    的右顶点为 $A$,将圆心在 $y$ 轴上, 且与 $\Gamma$ 的两支各恰有一个公共点的圆称为“好圆”, 若两个好圆外切于点 $P$, 圆心距为 $d$, 求 $\dfrac{d}{PA}$ 的所有可能值。
    \begin{solution}
        考虑以 $(0, y_{0})$ 为圆心的好圆 
        \[
        \Omega_{0}: x^{2}+(y-y_{0})^{2}=r_{0}^{2}\left(r_{0}>0\right)
        \]
        由 $\Omega_{0}$ 与 $\Gamma$ 的方程消去 $x$, 得关于 $y$ 的二次方程
        \[
        2y^{2}-2y_{0}y+y_{0}^{2}+1-r_{0}^{2}=0
        \]
        令判别式为零,
        \[
        \Delta=4y_{0}^{2}-8\left(y_{0}^{2}+1-r_{0}^{2}\right)=0\Rightarrow y_{0}^{2}=2r_{0}^{2}-2
        \]
        对于外切于点 $P$ 的两个好圆 $\Omega_{1}, \Omega_{2}$, 显然 $P$ 在 $y$ 轴上;设 $P(0, h)$, $\Omega_{1}, \Omega_{2}$ 的半径分别为 $r_{1}, r_{2}$, 不妨设 $\Omega_{1}, \Omega_{2}$ 的圆心分别为 $(0, h+r_{1}), (0, h-r_{2})$, 则有
        \[
        (h+r_{1})^{2}=2r_{1}^{2}-2,(h-r_{2})^{2}=2r_{2}^{2}-2
        \]
        两式相减得 
        \[
        2h(r_{1}+r_{2})=r_{1}^{2}-r_{2}^{2}
        \]
        而 $r_{1}+r_{2}>0$, 故化简得 $h=\dfrac{r_{1}-r_{2}}{2}$,进而 
        \[
        \left(\frac{r_{1}-r_{2}}{2}+r_{1}\right)^{2}=2r_{1}^{2}-2 \Rightarrow r_{1}^{2}-6r_{1}r_{2}+r_{2}^{2}+8=0 \tag{1}
        \]
        由于 $d=r_{1}+r_{2},A(1,0)$,且 
        \[
        PA^{2}=h^{2}+1=\frac{(r_{1}-r_{2})^{2}}{4}+1
        \]
        而 (1)可等价地写为
        \[
        2(r_{1}-r_{2})^{2}+8=(r_{1}+r_{2})^{2}
        \]
        即 $8PA^{2}=d^{2}$, 所以 
        \[
        \frac{d}{PA}=2\sqrt{2}
        \]
    \end{solution}
    
    \question 设$A,B$为双曲线
    \[
    W: \frac{x^{2}}{a^{2}}-\frac{y^{2}}{b^{2}}=1
    \]
    与实轴的交点, $P(0,1)$为双曲线外一点,$PA,PB$分别交双曲线于另一点$C,D$,过$C,D$的切线相交于$Q$,若$\triangle QCD$是一个正三角形且面积为$\dfrac{16\sqrt{3}}{27}$,求双曲线$W$的方程式。
    \begin{solution}
        设
        \[
        A(-a,0),B(a,0),C(-a \sec \alpha,b \tan \alpha),D(a \sec \alpha,b \tan \alpha)
        \]
        且$\sec \alpha > 0$,从而$PB$的直线方程为
        \[
        y=-\frac{1}{a}x+1
        \] 
        $D$在$PB$上,于是有
        \[
        b\tan \alpha=-\sec \alpha+1 \tag{1}
        \]
        且$QC,QD$的直线方程为
        \[
        -\frac{x \sec \alpha}{a}-\frac{y \tan \alpha}{b}=1,\quad \frac{x \sec \alpha}{a}-\frac{y \tan \alpha}{b}=1
        \]
        两式联立得
        \[
        Q\left(0,-\dfrac{b}{\tan \alpha}\right)
        \]
        又由对称性及$\triangle QCD$是正三角形可得
        \[
        \frac{1}{2}\cdot2a \sec \alpha\cdot\sqrt{3}a \sec \alpha=\frac{16\sqrt{3}}{27},\quad -\frac{b}{\tan \alpha}-b \tan \alpha=\sqrt{3}a \sec \alpha
        \]
        即
        \[
        a^{2}\sec^{2}\alpha=\frac{16}{27},\quad -b \sec \alpha=\sqrt{3}a \tan \alpha
        \]
        联立(1)可得 
        \[
        \sec \alpha=\frac{\sqrt{3}a}{\sqrt{3}a-b^{2}},\quad \tan \alpha=\frac{-b}{\sqrt{3}a-b^{2}}
        \]
        由 $1+\tan^{2}\alpha=\sec^{2}\alpha$得$b^{2}=2\sqrt{3}a-1$,代入$a \sec \alpha=\dfrac{4}{3\sqrt{3}}$得
        \[
        9a^{2}+4\sqrt{3}a-4=0\Rightarrow a=\frac{2\sqrt{3}}{9}, b=\frac{\sqrt{3}}{3}
        \]
        所以双曲线$W$方程为
        \[
        \frac{27x^{2}}{4}-3y^{2}=1
        \]
    \end{solution}

    \question 已知双曲线 
    \[
    \frac{x^2}{a^2}-\frac{y^2}{b^2}=1\ (a>0,b>0)
    \]
    的左、右焦点分别为 $F_1, F_2$。  
    过点 $F_1$ 作直线分别交双曲线左支和一条渐近线于点 $A,B$（$A,B$ 在同一象限）,且满足 $F_1A = AB$。  
    连接 $AF_2$,若满足 $AF_2 \perp BF_1$,求该双曲线离心率的平方$e^2$。
    \begin{solution}
        不妨设 $A(x_0,y_0)$在第二象限,由 $AF_2\perp BF_1$ 得 
        \[
        \frac{y_0}{x_0-c}\cdot\frac{y_0}{x_0+c}=-1 \Rightarrow y_0^2+x_0^2-c^2=0 \tag{1}
        \]
        因为$A$ 在双曲线上,所以 
        \[
        \frac{x_0^2}{a^2}-\frac{y_0^2}{b^2}=1 \Rightarrow
        x_0^2=a^2+\frac{a^2y_0^2}{b^2}
        \]
        代入(1)解得 $y_0=\dfrac{b^2}{c}$;又$ F_1A=AB$,所以$ B(2x_0+c,2y_0)$;又 $B$ 在渐近线 $y=-\dfrac{b}{a}x$ 上,
        \[
        2y_0=-\frac{b}{a}(2x_0+c) \Rightarrow -2bx_0=2ay_0+bc
        \] 
        两边平方得 
        \[
        b^2(2x_0+c)^2=a^2(2y_0)^2 \Rightarrow 4b^2x_0^2+4b^2cx_0+b^2c^2=4a^2y_0^2 \tag{2}
        \]
        将 $x_0^2=a^2+\dfrac{a^2y_0^2}{b^2}$ 和 $y_0=\dfrac{b^2}{c}$ 代入(2)得 
        \[
        4b^2(a^2+\frac{a^2b^4}{b^2c^2})+4b^2c\sqrt{a^2+\frac{a^2b^4}{b^2c^2}}+b^2c^2=4a^2\frac{b^4}{c^2} 
        \]
        解得 
        \[
        3a^2-4ab-b^2=0\Rightarrow \frac{a}{b}=\frac{\sqrt{7}+2}{3} 
        \]
        于是
        \[
        e^2=1+\frac{b^2}{a^2}=1+\frac{9}{(\sqrt{7}+2)^2}\cdot\frac{(\sqrt7-2)^2}{(\sqrt7-2)^2}=12-4\sqrt{7}
        \] 
    \end{solution}

    \question 设双曲线 
    \[
    C: \frac{x^{2}}{a^{2}} - \frac{y^{2}}{b^{2}} = 1, \quad a > 0, b > 0
    \]
    的左、右焦点分别为 $F_1, F_2$。过点 $F_1$ 作斜率为 $\dfrac{\sqrt{3}}{3}$ 的直线 $l$,与双曲线左右两支分别交于点 $M, N$。且满足
    \[
    (\overrightarrow{F_2 M} + \overrightarrow{F_2 N}) \cdot \overrightarrow{M N} = 0.
    \]
    求双曲线的离心率。
        \begin{solution}
        设$D$为$MN$的中点,连接$F_{2}D$,易知$\overrightarrow{F_{2}M}+\overrightarrow{F_{2}N}=2\overrightarrow{F_{2}D}$,所以
        \[
        (\overrightarrow{F_{2}M}+\overrightarrow{F_{2}N})\cdot\overrightarrow{MN}=2\overrightarrow{F_{2}D}\cdot\overrightarrow{MN}=0\Rightarrow F_{2}D \perp MN
        \]
        所以$F_{2}M=F_{2}N$,现设$t=F_{2}M=F_{2}N$,
        由双曲线定义
        \[
        MF_{1}=t-2a,\;NF_{1}=t+2a\Rightarrow MN=NF_{1}-MF_{1}=4a
        \]
        又$D$是$MN$的中点,
        \[
        MD=ND=2a,\;F_{1}D=F_{1}M+MD=t
        \]
        在直角$\triangle F_{1}DF_{2}$及直角$\triangle MDF_{2}$中,
        \[
        F_{2}D=\sqrt{4c^{2}-t^{2}}=\sqrt{t^{2}-4a^{2}}\Rightarrow t^{2}=2a^{2}+2c^{2}
        \]
        所以 
        \[
        F_{1}D=\sqrt{2c^{2}-2a^{2}}, F_{2}D=t=\sqrt{2a^{2}+2c^{2}}
        \]
        直线的斜率为 $\frac{\sqrt{3}}{3}$,
        \[
        \tan\angle DF_{1}F_{2}=\frac{F_{2}D}{F_{1}D}=\frac{\sqrt{2c^{2}-2a^{2}}}{\sqrt{2a^{2}+2c^{2}}}=\frac{\sqrt{3}}{3}
        \]
        所以
        \[
        \frac{c^{2}-a^{2}}{a^{2}+c^{2}}=\frac{1}{3}\Rightarrow c=\sqrt{2}a
        \]
        离心率为 
        \[
        \frac{c}{a}=\sqrt{2}
        \]
    \end{solution}
    
    \question 已知双曲线 
    \[
    C:\frac{x^{2}}{a^{2}}-\frac{y^{2}}{b^{2}}=1\quad a>0, b>0
    \] 
    的左焦点为 $F_1$,离心率为 $e$, 直线 $y = kx\ (k \ne 0)$ 分别与 $C$ 的左右两支交于点 $M, N$,  
    若 $\triangle MF_1N$ 的面积为 $\sqrt{3}$,且 $\angle MF_1N = 60^\circ$,求 $e^2 + 3a^2$ 的最小值。
    \begin{solution}
        连接$NF_{2},MF_{2}$,由对称性可知四边形$MF_{1}NF_{2}$为平行四边形,故
        \[
        NF_{2}=MF_{1},NF_{1}=MF_{2},\angle F_{1}NF_{2}=120^{\circ},S_{\triangle F_{1}NF_{2}}=S_{\triangle MF_{1}N}=\sqrt{3}
        \]
        由面积公式得
        \[
        \frac{1}{2}NF_{1} \cdot NF_{2}\sin120^{\circ}=\sqrt{3}\Rightarrow NF_{1}\cdot NF_{2}=4
        \]
        由$F_{1}N-F_{2}N=2a$,在三角形$F_{1}NF_{2}$中,由余弦定理得
        \begin{align*}
        \cos120^{\circ}
        &=\frac{F_{1}N^{2}+F_{2}N^{2}-4c^{2}}{2F_{1}N\cdot F_{2}N} \\
        &=\frac{(F_{1}N-F_{2}N)^{2}+2F_{1}N\cdot F_{2}N-4c^{2}}{2F_{1}N\cdot F_{2}N} \\
        &=\frac{(2a)^{2}+2\cdot4-4c^{2}}{2\cdot4}\\
        &=\frac{a^{2}+2-c^{2}}{2}=-\frac{1}{2}
        \end{align*}
        解得$b^{2}=c^{2}-a^{2}=3$,由AM-GM不等式,
        \[
        e^{2}+3a^{2}=1+\frac{b^{2}}{a^{2}}+3a^{2}=1+\frac{3}{a^{2}}+3a^{2}\ge1+2\sqrt{\frac{3}{a^{2}}\cdot3a^{2}}=7
        \]
        当且仅当$\dfrac{3}{a^{2}}=3a^{2}$即$a^{2}=1$时等号成立。
    \end{solution}

    \question 点 \(A(-1,1)\),\(B,C\) 在双曲线 \(x^{2}-y^{2}=1\) 上,且 \(\triangle ABC\) 为正三角形。求其面积。  
    \begin{solution}
    $\frac{3\sqrt3}{2}$
    \textcolor{red}{(待解)}
    \end{solution}    

\end{questions}
\pagebreak

\begin{center}
  {\fontsize{30pt}{26pt}\selectfont
    \hypertarget{坐标变换}{坐标变换} \label{坐标变换}
  }
\end{center}
\separator
\vspace{1pt}
\begin{questions}
    \question 已知平移坐标轴到新原点 $O'(h, k)$ 后,二直线 $2x + 3y - 4 = 0$、$x - 2y + 1 = 0$ 在新坐标中的方程式分别变为 $2x' + 3y' - 3 = 0$ 及 $x' - 2y' + 5 = 0$,求 $(h, k)$。
    \begin{solution}
        设平移公式为 $x = x' + h, y = y' + k$,代入原直线方程整理得
        \[
        \begin{cases} 
        2x' + 3y' + (2h + 3k - 4) = 0 \\
        x' - 2y' + (h - 2k + 1) = 0 
        \end{cases}
        \]
        对照已知新坐标系下的方程,
        \[
        \begin{cases} 
        2h + 3k - 4 = -3 \\
        h - 2k + 1 = 5 
        \end{cases}
        \]
        解得
        \[
        (h, k) = (2, -1)
        \]
    \end{solution}

    \question 已知当坐标轴旋转 $\theta$ 角时,点 $M(-\sqrt{2}, \sqrt{2})$ 的新坐标是 $(0, 2)$。求 $(1, 1)$ 的原坐标。
    \begin{solution}
        点 $M(-\sqrt{2}, \sqrt{2})$ 的新坐标是 $(0, 2)$,故
        \[
        \begin{bmatrix} 0 \\ 2 \end{bmatrix} 
        = \begin{bmatrix} \cos\theta & \sin\theta \\ -\sin\theta & \cos\theta \end{bmatrix} \begin{bmatrix} -\sqrt{2} \\ \sqrt{2} \end{bmatrix}= \begin{bmatrix} -\sqrt{2} \cos\theta + \sqrt{2}\sin\theta \\ \sqrt{2}\sin\theta + \sqrt{2}\cos\theta \end{bmatrix}
        \]
        给出的两方程皆有
        \[
        \theta = 45^\circ
        \]
        故
        \[
        \begin{bmatrix} x \\ y \end{bmatrix} 
        = \begin{bmatrix} \cos 45^\circ & -\sin 45^\circ \\ \sin 45^\circ & \cos 45^\circ \end{bmatrix} \begin{bmatrix} 1 \\ 1 \end{bmatrix}
        = \begin{bmatrix} 0 \\ \sqrt{2} \end{bmatrix}
        \]
        $(1, 1)$ 的原坐标为 $(0, \sqrt{2})$。
    \end{solution}

    \question 以 $O$ 为原点的 $xy$ 平面上,取两点 $A(\sqrt{3},1),B(-1,\sqrt{3}),t \in \mathbb{R}$,点 $P$ 满足
    \[
    \overrightarrow{OP} = t^2 \overrightarrow{OA} + t \overrightarrow{OB},
    \]
    求 $P$ 点的轨迹与 $x$ 轴所围成的图形面积。
    \begin{solution}
        设 $P(x,y)$,则
        \[
        \begin{bmatrix} x \\ y \end{bmatrix} = 
        \begin{bmatrix}
        \sqrt{3} & -1 \\
        1 & \sqrt{3}
        \end{bmatrix}
        \begin{bmatrix}
        t^2 \\ t
        \end{bmatrix}=
        \begin{bmatrix}
        \cos 30^\circ & -\sin 30^\circ \\
        \sin 30^\circ & \cos 30^\circ
        \end{bmatrix}
        \begin{bmatrix}
        x' \\ y'
        \end{bmatrix}
        \]
        其中
        \[
        \begin{bmatrix} x' \\ y' \end{bmatrix} =
        \begin{bmatrix}
        \frac{t^2}{2} \\
        \frac{t}{2}
        \end{bmatrix}.
        \]
        在新坐标系下,轨迹满足
        \[
        \Gamma: y'^2 = 2 x' \tag{1}
        \]
        即一条抛物线,既然旋转保持面积不变,即求新坐标系下所围面积,将$x$轴顺时针转$30^\circ$得
        \[
        L: y' = -\frac{1}{\sqrt{3}} x' \tag{2}
        \]
        联立$(1),(2)$解得
        \[
        O(0,0), \quad A(6, -2 \sqrt{3})
        \]
        所围成图形面积为
        \[
        \int_{-2\sqrt{3}}^{0} \left(-\sqrt{3} y - \frac{y^2}{2}\right) dy = 2 \sqrt{3}.
        \]
    \end{solution}

    \question 坐标平面上有一椭圆 
    \[
    \Gamma_1: \frac{x^2}{b^2} + \frac{y^2}{a^2} = 1
    \] 
    以原点 $O(0,0)$ 为中心,将椭圆 $\Gamma_1$ 逆时针旋转 $\dfrac{\pi}{3}$ 后得到椭圆 
    \[
    \Gamma_2: 43x^2 + 14\sqrt{3}xy + 57y^2 = 576
    \] 
    求椭圆 $\Gamma_1$ 的面积。
    \begin{solution}
        有
        \[
        43x^2 + 14\sqrt{3}xy + 57y^2 = 
        \begin{bmatrix} x & y \end{bmatrix}
        \begin{bmatrix} 43 & 7\sqrt{3} \\ 7\sqrt{3} & 57 \end{bmatrix} 
        \begin{bmatrix} x \\ y \end{bmatrix}
        \]
        设矩阵 
        \[
        A = \begin{bmatrix} 43 & 7\sqrt{3} \\ 7\sqrt{3} & 57 \end{bmatrix}
        \] 
        特征值为
        \[
        \lambda_1 = 36, \quad \lambda_2 = 64
        \]
        于是旋转后的标准形式为
        \[
        36x'^2 + 64y'^2 = 576 \Rightarrow \frac{x'^2}{16} + \frac{y'^2}{9} = 1
        \]
        由此可得半轴长
        \[
        a = 4, \quad b = 3
        \]
        故椭圆面积为
        \[
        S = \pi a b = 12 \pi
        \]
    \end{solution}

    \question 设椭圆 
    \[
    \frac{x^{2}}{25}+\frac{y^{2}}{9}=1
    \]
    的右焦点为 $F(4,0)$,点 $P$ 为椭圆上一动点,若以 $PF$ 为一边作正方形 $FPQR$($FPQR$ 按逆时针方向排列),当 $P$ 点沿着椭圆绕行一周时,试求 $R$ 点的轨迹方程式。
    \begin{figure}[H]
        \centering        
        \includegraphics[width=0.5\textwidth]{images/image103.jpg}
    \end{figure}
    \begin{solution}
        以 $F$ 为旋转中心,将 $P$ 逆时针旋转 $90^\circ$ 即为 $R$。因此 $P(5\cos\theta, 3\sin \theta)$ 先平移 $(-4,0)$,得
        \[
        P' = (5\cos\theta - 4, \ 3\sin \theta)
        \]
        再逆时针旋转 $90^\circ$,有
        \[
        P'' = 
        \begin{bmatrix} 
        0 & -1 \\ 
        1 & 0 
        \end{bmatrix}
        \begin{bmatrix} 
        5\cos\theta - 4 \\ 
        3\sin \theta 
        \end{bmatrix}
        =
        \begin{bmatrix} 
        -3\sin \theta \\ 
        5\cos\theta - 4
        \end{bmatrix}
        \]
        再平移 $(4,0)$,得
        \[
        R = (-3\sin \theta + 4, \ 5\cos\theta - 4)
        \]
        即
        \[
        \begin{cases}
        x = -3\sin \theta + 4 \\
        y = 5\cos\theta - 4
        \end{cases}
        \]
        由此得
        \[
        \frac{(x-4)^2}{9} + \frac{(y+4)^2}{25} = 1
        \]
    \end{solution}

    \question 将椭圆 
    \[
    \Gamma_1: \frac{x^2}{9} + \frac{y^2}{4} = 1
    \]
    以原点 $O$ 为中心,逆时针旋转锐角 $\theta$ 后,得到椭圆 $\Gamma_2$。已知 $\Gamma_2$ 的长轴方程为 $y=2x$,求 $\Gamma_2$ 的方程。
    \begin{solution}
        椭圆 $$\Gamma_1: \frac{x^2}{9} + \frac{y^2}{4} = 1$$ 的长轴在 $x$ 轴上,长轴旋转至 $y=2x$,设旋转角满足 $\tan \theta = 2$,则
        \[
        \sin \theta = \frac{2}{\sqrt{5}}, \quad \cos \theta = \frac{1}{\sqrt{5}}.
        \]
        设 $(x, y)$ 是 $\Gamma_2$ 上一点,对应 $\Gamma_1$ 上某点 $(s, t)$,则有
        \[
        \begin{bmatrix} x \\ y \end{bmatrix} =
        \begin{bmatrix} \frac{1}{\sqrt{5}} & -\frac{2}{\sqrt{5}} \\ \frac{2}{\sqrt{5}} & \frac{1}{\sqrt{5}} \end{bmatrix}
        \begin{bmatrix} s \\ t \end{bmatrix}.
        \]
        所以反解得
        \[
        \begin{bmatrix} s \\ t \end{bmatrix} =
        \begin{bmatrix} \frac{1}{\sqrt{5}}(x+2y) \\ \frac{1}{\sqrt{5}}(-2x+y) \end{bmatrix}.
        \]
        因为 $(s,t)$ 满足 $\Gamma_1$:
        \[
        \frac{s^2}{9} + \frac{t^2}{4} = 1,
        \]
        代入得
        \[
        \frac{\frac{1}{5}(x+2y)^2}{9} + \frac{\frac{1}{5}(-2x+y)^2}{4} = 1,
        \]
        即
        \[
        8x^2 - 4xy + 5y^2 - 36 = 0.
        \]
    \end{solution}

    \question 在坐标平面上,考虑二阶方阵 
    \[
    A=\frac{1}{5}\begin{bmatrix}3&-4\\ 4&3\end{bmatrix}
    \]
    所定义的线性变换。对于平面上异于原点 $O$ 的点 $P_1$,设 $P_1$ 经 $A$ 变换成 $P_2,P_2$ 经 $A$ 变换成 $P_3$。假设 $P_1$ 是图形 $y=\dfrac{1}{10}x^{2}-10$ 上的动点,求 $\triangle P_1P_2P_3$ 面积的最小可能值。
    \begin{solution}
        我们有
        \[
        A = \frac{1}{5}\begin{bmatrix}3 & -4 \\ 4 & 3 \end{bmatrix} 
        = \begin{bmatrix} \frac{3}{5} & -\frac{4}{5} \\ \frac{4}{5} & \frac{3}{5} \end{bmatrix} 
        = \begin{bmatrix} \cos\theta & -\sin\theta \\ \sin\theta & \cos\theta \end{bmatrix}
        \Rightarrow \sin \theta = \frac{4}{5},\ \cos \theta = \frac{3}{5}
        \]
        故
        \[
        \sin \angle P_1OP_2 = \sin 2\theta = 2\sin\theta\cos\theta = \frac{24}{25}
        \]
        假设 $OP_1=a$,则
        \begin{align*}
        [\triangle P_1P_2P_3] &= [\triangle OP_1P_2] + [\triangle OP_2P_3] - [\triangle OP_1P_3]\\
        &= \frac{1}{2}a^2 \sin \theta + \frac{1}{2}a^2 \sin \theta - \frac{1}{2}a^2 \sin 2\theta 
        = \frac{8}{25}a^2
        \end{align*}
        又$P_1$  在 $y = \dfrac{x^2}{10} - 10$ 上,设$P_1\left(m, \dfrac{m^2}{10} - 10\right)$,则
        \[
        a^2 = \frac{1}{100}m^4 - m^2 + 100 = \frac{1}{100}(m^2 - 50)^2 + 75
        \]
        故
        \[
        [\triangle P_1P_2P_3]_{\min}=\frac{8}{25} \cdot 75=24
        \]
    \end{solution}

    \question 在直角坐标平面上,圆 $x^{2}+y^{2}=1$ 先被变换为曲线 $\Gamma_{1}$,再被
    \[
    A=\begin{bmatrix}\frac{\sqrt{3}}{2} & \sqrt{3} \\[4pt] \frac{1}{2} & -3\end{bmatrix}
    \]
    变换为曲线 $\Gamma_{1}$,再被
    \[
    B=\prod_{k=8}^{67}\begin{bmatrix}\cos k^\circ & \sin k^\circ \\[4pt] \sin k^\circ & -\cos k^\circ\end{bmatrix}
    \]
    变换为曲线 $\Gamma_{2}$,求 $\Gamma_{2}$ 的方程。
    \begin{solution}
        首先注意到
        \[
        \begin{bmatrix}\cos a^\circ & \sin a^\circ\\[4pt] \sin a^\circ & -\cos a^\circ \end{bmatrix}
        \begin{bmatrix}\cos (a+1)^\circ & \sin (a+1)^\circ\\[4pt] \sin (a+1)^\circ & -\cos (a+1)^\circ \end{bmatrix}
        =\begin{bmatrix}\cos 1^\circ & \sin 1^\circ \\[4pt] -\sin 1^\circ & \cos 1^\circ \end{bmatrix}.
        \]
        由此,
        \[
        B=\begin{bmatrix}\cos 1^\circ & \sin 1^\circ \\[4pt] -\sin 1^\circ & \cos 1^\circ \end{bmatrix}^{30}
        =\begin{bmatrix}\cos 30^\circ & \sin 30^\circ \\[4pt] -\sin 30^\circ & \cos 30^\circ \end{bmatrix}
        =\begin{bmatrix}\tfrac{\sqrt3}{2} & \tfrac{1}{2} \\[4pt] -\tfrac{1}{2} & \tfrac{\sqrt3}{2} \end{bmatrix}.
        \]
        所以
        \[
        BA=\begin{bmatrix}\tfrac{\sqrt3}{2} & \tfrac{1}{2} \\[4pt] -\tfrac{1}{2} & \tfrac{\sqrt3}{2} \end{bmatrix}
        \begin{bmatrix}\tfrac{\sqrt3}{2} & \sqrt3 \\[4pt] \tfrac{1}{2} & -3\end{bmatrix}
        =\begin{bmatrix}1 & 0 \\[4pt] 0 & -2\sqrt3 \end{bmatrix}.
        \]
        设变换后的坐标为 $(x',y')^{\mathsf T}=BA\,(x,y)^{\mathsf T}$,则
        \[
        x'=x,\qquad y'=-2\sqrt3\,y.
        \]
        原方程 $x^{2}+y^{2}=1$ 在新坐标下变为
        \[
        x'^{2}+\Big(\frac{y'}{-2\sqrt3}\Big)^{2}=1 \Rightarrow x'^{2}+\frac{y'^{2}}{12}=1.
        \]
        即
        \[
        \Gamma_{2}:\quad x^{2}+\frac{y^{2}}{12}=1.
        \]
    \end{solution}
\end{questions}
\pagebreak

\begin{center}
  {\fontsize{30pt}{26pt}\selectfont
    \hypertarget{轨迹方程式、参数方程式}{轨迹方程式、参数方程式} \label{轨迹方程式、参数方程式}
  }
\end{center}
\separator
\vspace{1pt}

\begin{questions}
    \question 设 $A(-1,0),B(2,0)$, 在 $\triangle ABC$ 中, 若 $\angle B = 2 \angle A$, 求点 $C$ 的轨迹方程式。
    \ifprintanswers
    \begin{figure}[H]
        \centering
        \includegraphics[width=0.5\linewidth]{images/image65.png}
    \end{figure}
    \fi
    \begin{solution}
        设 $C(x,y)$,作 $\angle B$ 的角平分线交 $AC$ 于点 $D$,由
        \[
        \angle B = 2 \angle A
        \]
        知$D$ 在 $AB$ 的垂直平分线上,故$D\left(\dfrac{1}{2}, a\right)$,又由角平分线定理,
        \[
        \frac{AD}{DC} = \frac{AB}{BC} = \frac{3}{\sqrt{(x-2)^2 + y^2}}
        \]
        由分比公式
        \[
        x_D=\frac{1}{2} = \frac{-\sqrt{(x-2)^2 + y^2} + 3x}{\sqrt{(x-2)^2 + y^2} + 3}
        \]
        得点 $C$ 的轨迹
        \[
        x^2 - \frac{y^2}{3} = 1, \quad x > 1
        \]
    \end{solution}

    \question 一抛物线 $y^2 = 4x$ 与一直线交于 $A,B$ 两点,已知抛物线与直线所围出的面积为 $\dfrac{9}{8}$, 求 $A,B$ 的中点轨迹方程。
    \begin{solution}
        不失一般性,设交点为$A\left(\dfrac{1}{4}y_1^2, y_1\right), B\left(\dfrac{1}{4}y_2^2, y_2\right),y_2\le y_1$,斜率
        \[
        m_L = \frac{y_1 - y_2}{\frac{1}{4}(y_1^2 - y_2^2)} = \frac{4}{y_1 + y_2},
        \]
        则直线方程为
        \[
        x = \frac{1}{4}(y_1 + y_2)y - \frac{1}{4}y_1 y_2.
        \]
        所围部分面积为
        \[
        \frac{9}{8} = \int_{y_2}^{y_1} \left( \frac{1}{4}(y_1 + y_2)y - \frac{1}{4}y_1 y_2 - \frac{1}{4}y^2 \right) dy
        = \frac{1}{24} (y_1 - y_2)^3
        \]
        由此得
        \[
        y_1 - y_2 = 3.
        \]
        设 $P$ 为 $AB$ 中点
        \[
        P\Big(\frac{1}{8}(y_1^2 + y_2^2), \frac{1}{2}(y_1 + y_2)\Big) \equiv (x, y)
        \]
        由于
        \[
        y^2 = \frac{1}{4}(y_1 + y_2)^2 = \frac{1}{4}(y_1^2 + y_2^2) + \frac{1}{2}y_1 y_2
        = \frac{1}{2}(y_1^2 + y_2^2) - \frac{9}{4} = 4x - \frac{9}{4},
        \]
        所以中点轨迹方程为
        \[
        y^2 = 4x - \frac{9}{4}.
        \]
    \end{solution}

    \question 设 $AB$ 为抛物线 $y^2=4ax$ 之一动弦, 且 $\angle AOB = 90^\circ$, 其中 $O$ 为原点,从抛物线外一点$T$作点$A,B$的切线,且$A,B$处的法线交于点$N$,证明 $TN$ 之中点轨迹方程式为 
    \[
    2y^2 = 25a(x-a).
    \]
    \begin{solution}
        设 $A(at_1^2, 2at_1),B(at_2^2, 2at_2)$,已知$OA \perp OB$,故
        \[
        m_{OA}\cdot m_{OB}=\frac{2at_1}{at_1^2} \cdot \frac{2at_2}{at_2^2} = -1 \Rightarrow t_1 t_2 = -4
        \]
        由链导法,在点$A(at_1^2, 2at_1)$的切线斜率为
        \[
        \frac{dy}{dx}=\frac{dy}{dt} \div \frac{dy}{dt}=\frac{2a}{2at_1}=\frac{1}{t_1}
        \]
        故在 $A$的切线方程为
        \[
        y-2at_1=\frac{1}{t}(x-at_1^2) \Rightarrow t_1y - x = at_1^2 \tag{1}
        \]
        同理, 点$B$ 的切线方程为
        \[
        t_2y - x = at_2^2 \tag{2}
        \]
        联立$(1),(2)$,解得切线交点 $T$
        \[
        T(a t_1 t_2,a(t_1 + t_2))=(-4a,a(t_1 + t_2))
        \]
        在$A$的法线方程为
        \[
        y - 2at_1 = -t_1(x - at_1^2) \Rightarrow t_1 x + y = a t_1^3 + 2at_1 \tag{3}
        \]
        同理, 点$B$ 的切线方程为
        \[
        t_2 x + y = a t_2^3 + 2at_2 \tag{4}
        \]
        联立$(3),(4)$,解得 $N$ 点坐标
        \[
        N(a(t_1+t_2)^2 + 6a, 4a(t_1+t_2))
        \]
        故$TN$ 中点坐标为
        \[
        M\left(\frac{1}{2}a(t_1+t_2)^2 + a,\frac{5a}{2}(t_1+t_2)\right)
        \]
        消去参数$t_1+t_2$,即得
        \[
        2y^2 = 25a(x-a)
        \]
    \end{solution}

    \question 抛物线$\Gamma: y^2 = 4x$, $F$为$\Gamma$的焦点, $A, B$为$\Gamma$上的两个不重合的动点, 使得线段$AB$的一个三等分点$P$位于线段$OF$上(含端点), 记$Q$为线段$AB$的另一个三等分点, 求$Q$的轨迹方程.
    \begin{solution}
        设$A(x_1, y_1),B(x_2, y_2)$. 不妨设$\overrightarrow{AP} = \overrightarrow{PQ} = \overrightarrow{QB}$, 则
        \[
        P\left(\frac{2x_1+x_2}{3}, \frac{2y_1+y_2}{3}\right)
        \]
        易知$F(1,0)$. 由于点$P$位于线段$OF$上, 故
        \[
        \frac{2x_1+x_2}{3}\in[0,1],\ \frac{2y_1+y_2}{3}=0
        \]
        设$y_1 = t,y_2 = -2t$, 则$x_1 = \dfrac{t^2}{4},x_2 = t^2$. 此时有
        \[
        \frac{2x_1+x_2}{3}=\frac{t^2}{2}\in[0,1]
        \]
        且由$A, B$不重合知$t\neq0$, 所以$t^2\in(0, 2]$.
        设$Q(x_Q, y_Q)$, 则
        \[
        x_Q = \frac{x_1+2x_2}{3} = \frac{3}{4}t^2, y_Q = \frac{y_1+2y_2}{3} = -t
        \]
        故有 $y_Q^2 = \dfrac{4}{3}x_Q$,注意到 $x_Q = \dfrac{3}{4}t^2 \in \left[0, \dfrac{3}{2}\right]$, 故点$Q$的轨迹方程为 
        \[
        y^2 = \frac{4}{3}x
        \]
    \end{solution}

    \question 已知抛物线 $\Gamma:y = x^2$ 上三个不同点 $A(1,1),B,C$ 满足 $AB \perp AC$,过 $B,C$ 分别作 $\Gamma$ 的切线交于点 $P$,求点 $P$ 的轨迹方程。
    \begin{solution}
        设 $B(x_1, x_1^2),C(x_2, x_2^2)$,直线 $BC$ 的方程为 $y = kx + m$,联立$\Gamma:y = x^2$得
        \[
        x^2 - kx - m = 0
        \]
        由韦达定理
        \[
        x_1 + x_2 = k,\quad x_1 x_2 = -m
        \]
        于是
        \[
        y_1 + y_2 = x_1^2 + x_2^2 = (x_1 + x_2)^2 - 2x_1 x_2 = k^2 + 2m
        \]
        \[
        y_1 y_2 = x_1^2 x_2^2 = (x_1 x_2)^2 = m^2
        \]
        由 $AB \perp AC$,即 $\overrightarrow{AB} \cdot \overrightarrow{AC} = 0$,有
        \[
        (x_1 - 1)(x_2 - 1) + (y_1 - 1)(y_2 - 1) = 0
        \]
        展开整理得
        \[
        x_1 x_2 - (x_1 + x_2) + 1 + y_1 y_2 - (y_1 + y_2) + 1 = 0
        \]
        代入上式各值,得到
        \[
        -m - k + 1 + m^2 - (k^2 + 2m) + 1 = 0
        \Rightarrow m^2 - 3m - k^2 - k + 2 = 0
        \]
        解得
        \[
        m = k + 2 \quad \text{或} \quad m = -k + 1
        \]
        若 $m = -k + 1$,则直线 $BC$ 方程为 $y = k(x - 1) + 1$,经过 $A(1,1)$,舍去;因此 $m = k + 2,BC$ 方程为
        \[
        y = k(x + 1) + 2
        \]
        即 $BC$ 恒过定点 $(-1, 2)$。分别作点 $B(x_1, x_1^2)$ 与 $C(x_2, x_2^2)$ 的切线方程:
        \[
        y - x_1^2 = 2x_1(x - x_1)\; ,\; y - x_2^2 = 2x_2(x - x_2) 
        \]
        联立两式得切线交点$P$,
        \[
        x = \frac{x_1 + x_2}{2} = \frac{k}{2}, \quad y = x_1 x_2 = -m = -k-2
        \]
        消去参数 $k$即得点P的轨迹方程
        \[
        2x + y + 2 = 0
        \]
    \end{solution}

    \question 已知点 $P$ 在一半径为 $r$ 的圆的内部,在 $P$ 作一直角,使得直角的两边分别与圆相交于 $A,B$ 两点。设点 $Q$ 使得 $PAQB$ 构成矩形,求点 $Q$ 的轨迹。
    \begin{solution}
        不失平移及旋转的一般性,设圆心为原点 $O$,点 $P(d,0)$为圆
        \[
        x^2+y^2=r^2,
        \]
        内一点其中 $d=OP,0<d<r$,设在点 $P$ 作的直角的两条边的方向分别与 $x$ 轴成角 $\theta,\theta+\dfrac{\pi}{2}$,对应的单位方向向量分别为
        \[
        (\cos\theta,\sin\theta),\quad(-\sin\theta,\cos\theta).
        \]
        故可设
        \[
        A\left(d+s\cos\theta,\;s\sin\theta\right),\quad B\left(d-t\sin\theta,\;t\cos\theta\right)
        \]
        其中 $s,t>0,$ 为 $PA,PB$ 的长度,由于 $A,B$ 在圆上,有
        \[
        (d+s\cos\theta)^2+(s\sin\theta)^2=r^2,\quad 
        (d-t\sin\theta)^2+(t\cos\theta)^2=r^2
        \]
        整理得
        \[
        s^2=r^2-d^2-2ds\cos\theta,\quad t^2=r^2-d^2+2dt\sin\theta \tag{1}
        \]
        由于$PAQB$为矩形,点 $Q$坐标为
        \[
        Q(x,y)=\left(d+s\cos\theta-t\sin\theta,\;s\sin\theta+t\cos\theta\right)
        \]
        观察
        \begin{align*}
        x^2+y^2
        &=\bigl(d+s\cos\theta-t\sin\theta\bigr)^2+\bigl(s\sin\theta+t\cos\theta\bigr)^2 \\
        &=d^2+2d\bigl(s\cos\theta-t\sin\theta\bigr)+s^2+t^2.
        \end{align*}
        将$(1)$代入得
        \begin{align*}
        x^2+y^2
        &=d^2+2d\bigl(s\cos\theta-t\sin\theta\bigr)
        +2(r^2-d^2)+2d\bigl(t\sin\theta-s\cos\theta\bigr) \\
        &=2r^2-d^2.
        \end{align*}
        右式与 $\theta,s,t$ 无关,仅与 $r$ 及 $d=OP$ 有关。因此
        \[
        OQ^2=2r^2-d^2
        \]
        为常数。由此可知,在适当的平移、旋转后的坐标系中,点 $Q$ 的轨迹是以原点为圆心、半径为 $\sqrt{2r^2-d^2}$ 的圆。将坐标系平移回去,得点 $Q$ 的轨迹为以圆心 $O(h,k)$ 为圆心、半径为 $\sqrt{2r^2-OP^2}$ 的圆。
    \end{solution}
    \begin{solution}
        不失一般性,设圆心为原点 $O$, 则向量 $\mathbf{p}=\overrightarrow{OP},\ |\mathbf{p}|^2<r^2$,且
        \[
        \mathbf{v}=\overrightarrow{PA},\quad \mathbf{w}=\overrightarrow{PB},\quad \mathbf{a}=\overrightarrow{OA},\quad \mathbf{b}=\overrightarrow{OB},
        \]
        其中 $\mathbf{v}\perp\mathbf{w},|\mathbf{a}|=|\mathbf{b}|=r$,且$\mathbf{a}=\mathbf{p}+\mathbf{v},\mathbf{b}=\mathbf{p}+\mathbf{w}$,设 $\mathbf{q}=\overrightarrow{OQ}=\mathbf{p}+\mathbf{v}+\mathbf{w}$,于是
        \begin{align*}
        |\mathbf{q}|^2
        &=(\mathbf{p}+\mathbf{v}+\mathbf{w})\cdot(\mathbf{p}+\mathbf{v}+\mathbf{w})\\
        &=|\mathbf{p}|^2+2\mathbf{p}\cdot\mathbf{v}+2\mathbf{p}\cdot\mathbf{w}+\mathbf{v}\cdot\mathbf{v}+\mathbf{w}\cdot\mathbf{w} \\
        &=|\mathbf{p}|^2+\mathbf{a}\cdot\mathbf{v}+\mathbf{b}\cdot\mathbf{w}+\mathbf{p}\cdot\mathbf{v}+\mathbf{p}\cdot\mathbf{w} \\
        &=|\mathbf{p}|^2+(\mathbf{a}+\mathbf{p})\cdot\mathbf{v}+(\mathbf{b}+\mathbf{p})\cdot\mathbf{w} \\
        &=|\mathbf{p}|^2+(\mathbf{a}+\mathbf{p})\cdot(\mathbf{a}-\mathbf{p})+(\mathbf{b}+\mathbf{p})\cdot(\mathbf{b}-\mathbf{p}) \\
        &=|\mathbf{p}|^2+|\mathbf{a}|^2-|\mathbf{p}|^2+|\mathbf{b}|^2-|\mathbf{p}|^2 \\
        &=2r^2-|\mathbf{p}|^2。
        \end{align*}
        因此 $Q$ 到圆心 $O$ 的距离为常数, 所以 $Q$ 的轨迹是以 $O$ 为圆心、半径为 $\sqrt{2r^2-|\mathbf{p}|^2}$ 的圆。
    \end{solution}

    \question 设 $A(1,\sqrt{3})$, $B(1,-\sqrt{3})$ 为平面上两定点,动点 $P$ 在线段 ${AB}$ 上。$O$ 为原点,且 $Q$ 在射线 $OP$ 上,满足 ${OP} \cdot {OQ} = 4$。当动点 $P$ 从 $A$ 沿线段 ${AB}$ 移动到 $B$,求点 $Q$ 轨迹的路径长。
    \ifprintanswers
    \begin{figure}[H]
    \centering
    \includegraphics[width=0.5\textwidth]{images/image27.png}
    \end{figure}
    \fi
    \begin{solution}
        设$P$ 点坐标为$P(1,t), -\sqrt{3} \le t \le \sqrt{3},$则
        \[
        {OP} = \sqrt{1 + t^2},
        \]
        $Q$ 在射线 $OP$ 上,则设
        \[
        Q(x,y), \quad y = t x, \quad x \ge 0.
        \]
        由条件${OP} \cdot {OQ} = 4,$
        \[
        \sqrt{1 + t^2} \cdot \sqrt{x^2 + y^2} = 4.
        \]
        代入 $y = t x$,得
        \[
        \sqrt{1 + t^2} \cdot \sqrt{x^2 + t^2 x^2} = (1 + t^2) |x| = 4 \Rightarrow x = \frac{4}{1 + t^2}.
        \]
        点 $Q$ 坐标为
        \[
        Q\left(\frac{4}{1+t^2}, \frac{4 t}{1 + t^2}\right).
        \]
        由 $x = \frac{4}{1+t^2}$,得$t^2 = \frac{4}{x} -1$,所以
        \[
        y^2 = t^2 x^2 = \left(\frac{4}{x} -1 \right) x^2 = 4 x - x^2.
        \]
        即
        \[
        (x - 2)^2 + y^2 = 2^2,
        \]
        为以 $(2,0)$ 为圆心,半径为 $2$ 的圆;由于 $P$ 在 ${AB}$,则 $1 \le {OP} \le 2$,即
        \[
        1 \le \sqrt{1 + t^2} \le 2 \implies 1 \le x = \frac{4}{1+t^2} \le 4,
        \]
        故点 $Q$ 轨迹为圆周在 $x \ge 1$ 的部分,即为该圆的三分之二周长,路径长为
        \[
        \frac{2}{3} \cdot 4 \pi = \frac{8 \pi}{3}.
        \]
    \end{solution}

    \question 已知曲线 $C$ 是到点 $P\!\left(-\dfrac12,\dfrac38\right)$ 和到直线 $y=-\dfrac58$ 距离相等的点的轨迹。直线 $l$ 过点 $Q(-1,0)$,点 $M$ 是 $C$ 上（不在 $l$ 上）的动点；$A,B$ 在 $l$ 上,且 $MA\perp l,\ MB\perp x$ 轴。
    \begin{parts}
    \part 求曲线 $C$ 的方程；
    \begin{solution}
        设 $N(x,y)$ 为 $C$ 上的点,据题意有
        \[
        \sqrt{\left(x+\frac12\right)^2+\left(y-\frac38\right)^2}=\left|y+\frac58\right|.
        \]
        两边平方并化简得
        \[
        y=\frac12(x^2+x)
        \]
    \end{solution}
    \part 求直线 $l$ 的方程,使 $\dfrac{QB^2}{QA}$ 为常数。
    \end{parts}
    \begin{solution}
        设$M\left(x,\frac{x^2+x}{2}\right),l: y=kx+k\ (k\neq0)$,则 $B(x,kx+k)$,从而
        \[
        QB=\sqrt{1+k^2}\,|x+1|
        \]
        在直角 $\triangle QMA$ 中,由毕氏定理,
        \[
        QA^2=QM^2-MA^2=(x+1)^2\!\left(1+\frac{x^2}{4}\right)-\frac{(x+1)^2\left(k-\frac{x}{2}\right)^2}{1+k^2}=\frac{(x+1)^2}{4(1+k^2)}(kx+2)^2,
        \]
        即
        \[
        QA=\frac{|x+1||kx+2|}{2\sqrt{1+k^2}}
        \]
        故欲使
        \[
        \frac{QB^2}{QA}
        =\frac{2(1+k^2)\sqrt{1+k^2}}{|k|}
        \left|\frac{x+1}{x+\frac{2}{k}}\right|
        \]
        为常数,需 $\dfrac{2}{k}=1$即 $k=2$,此时$\dfrac{QB^2}{QA}=5\sqrt5$,故直线 $l$ 的方程为
        \[
        l: 2x-y+2=0
        \]
    \end{solution}
    \begin{solution}
        设$M\left(x,\frac{x^2+x}{2}\right), l: y=kx+k$,则
        \[
        QB=\sqrt{1+k^2}\,|x+1|
        \]
        过 $Q(-1,0)$ 作垂直于 $l$ 的直线
        \[
        l_1: y=-\frac1k(x+1)
        \]
        设$H$ 为 $M$ 到 $l_1$ 的垂足,因 $QA=MH$,得
        \[
        QA=\frac{|x+1||kx+2|}{2\sqrt{1+k^2}}
        \]
        于是
        \[
        \frac{QB^2}{QA}
        =\frac{2(1+k^2)\sqrt{1+k^2}}{|k|}
        \left|\frac{x+1}{x+\frac{2}{k}}\right|
        \]
        同理得 $k=2$,从而
        \[
        l: 2x-y+2=0
        \]
    \end{solution}

    \question 已知平面内动点 $P$ 到 $F(1,0)$ 的距离与点 $P$ 到 $y$ 轴的距离差为 $1$,求:
    \ifprintanswers
    \begin{figure}[H]
    \centering
    \includegraphics[width=0.35\textwidth]{images/image41.png}
    \end{figure}
    \fi
    \begin{parts}
    \part 此动点 $P$ 的轨迹方程。
    \begin{solution}
        即 $P$ 到 $F(1,0)$ 的距离与 $P$ 到直线 $x = -1$ 的距离相等,依抛物线定义,轨迹为
        \[
        y^2 = 4x
        \]
    \end{solution}
    \part 过 $F$ 作两条互相垂直的直线 $L_1, L_2$。设 $L_1$ 与 $P$ 点轨迹相交于点 $A$ 和 $B,L_2$ 与 $P$ 点轨迹相交于点 $C,D$,求 $\overrightarrow{AD} \cdot \overrightarrow{CB}$ 的最小值。
\begin{solution}
极值出现在 $\overline{AB} = \overline{CD}$,可假设:

$A\left(\frac{a^2}{4}, a\right)$,$B\left(\frac{b^2}{4}, -b\right)$,$C\left(\frac{b^2}{4}, b\right)$,$D\left(\frac{a^2}{4}, -a\right)$

则:

\[
\overrightarrow{AD} = (0, -2a), \quad \overrightarrow{CB} = (0, -2b)
\Rightarrow \overrightarrow{AD} \cdot \overrightarrow{CB} = 4ab
\]

设 $P, Q$ 分别为 $A, B$ 在 $x$ 轴的垂足,则:

$\triangle APF \backsim \triangle BQF$（AAA）

\[
\frac{\overline{AP}}{\overline{BQ}} = \frac{\overline{AF}}{\overline{BF}} = \frac{d(L, A)}{d(L, B)}
\]

即:

\[
\frac{a}{b} = \frac{\frac{a^2}{4} + 1}{\frac{b^2}{4} + 1} = \frac{a^2 + 4}{b^2 + 4}
\]

两边交叉相乘得:

\[
ab^2 + 4a = a^2b + 4b
\Rightarrow ab(a - b) - 4(a - b) = 0
\Rightarrow (ab - 4)(a - b) = 0
\]

若 $a = b$,则 $\overline{CD}$ 为 $x$ 轴,不符合条件。

故取 $ab = 4$,则:

\[
\overrightarrow{AD} \cdot \overrightarrow{CB} = 4ab = 16
\]
\textcolor{red}{(待验证,为何极值出现在 $\overline{AB} = \overline{CD}$)}
\end{solution}
\end{parts}

    \question 在曲线
    \[
    \left(\frac{x}{a}\right)^{\frac{2}{3}} + \left(\frac{y}{b}\right)^{\frac{2}{3}} = 1,
    \]
    上的点 $P$ 处作曲线的切线,该切线与 $x$ 轴交于 $(h, 0)$,与 $y$ 轴交于 $(0, k)$。当 $P$ 沿给定曲线移动时,求点 $Q(h, k)$ 的轨迹。
    \begin{solution}
        考虑参数变换$P(x,y)=(a \cos^3 \theta,b \sin^3 \theta)$,则
        \[
        \frac{dy}{dx} 
        = \frac{dy}{d\theta} \div \frac{dx}{d\theta}
        = \frac{3b \sin^2 \theta \cos \theta}{3a \cos^2 \theta (-\sin \theta)} = -\frac{b}{a} \tan \theta
        \]
        在任意点 $\theta$ 处的切线方程为
        \[
        y - b \sin^3 \theta = -\frac{b}{a} \tan \theta (x - a \cos^3 \theta).
        \]
        整理得
        \[
        b x \sin \theta + a y \cos \theta = a b \sin \theta \cos \theta
        \]
        分别令$x=0,y=0$可得
        \[
        k = b \sin \theta,\quad h = a \cos \theta
        \]
        消去参数 $\theta$即得$Q(h, k)$的轨迹方程式
        \[
        \left(\frac{h}{a}\right)^2 + \left(\frac{k}{b}\right)^2 = 1.
        \]
        因此,点 $Q(h, k)$ 的轨迹是一个椭圆,其方程为
        \[
        \frac{x^2}{a^2} + \frac{y^2}{b^2} = 1.
        \]
    \end{solution}

    \question 等边三角形 $ABC$ 的边长为 2。正方形 $PQRS$ 满足 $P$ 在边 $AB$ 上,$Q$ 在边 $BC$ 上,$R,S$ 在边 $AC$ 上。若动点 $S$ 从边 $AC$ 移动到边 $AB$,使得 $P, Q, R, S$ 始终在三角形边上或内部构成一个正方形,证明点 $S$ 的轨迹是一条平行于 $BC$ 的直线。
    \begin{figure}[H]
        \centering        
        \includegraphics[width=0.4\textwidth]{images/image188.png}
    \end{figure}
    \ifprintanswers
    \begin{figure}[H]
        \centering        
        \includegraphics[width=0.4\textwidth]{images/image189.png}
    \end{figure}
    \fi
    \begin{solution}
        设正方形的边长为 $s$,且$\angle RQC=\theta$。$D,P,F$为 $R,S,T$ 在底边$BC$的垂足,过 $S$ 作平行于 $BC$ 的直线,交从 $R$ 作的垂线于 $E$。
        由 $\triangle RQD,\triangle SER$, $S$ 到 $BC$ 的距离为
        \[
        RD+ER = s\sin\theta+s\cos\theta,
        \]
        欲证其为常数。在 $\triangle RDQ,\triangle TFQ$ 中,
        \[
        QD = s\cos\theta,\quad FQ = s\sin\theta,\quad TF = s\cos\theta.
        \]
        在 $\triangle RDC,\triangle TFB$ 中,
        \[
        DC=RD\tan 30^\circ=\frac{\sqrt{3}}{3}s\sin\theta, \quad
        BF=TF\tan 30^\circ=\frac{\sqrt{3}}{3}s\cos\theta,  
        \]
        故
        \[
        2=DC+QD+FQ+BF
        =\frac{\sqrt{3}}{3}(s\cos\theta+s\sin\theta)+(s\cos\theta+s\sin\theta)
        \]
        即
        \[
        RD+ER = s\sin\theta+s\cos\theta=\frac{2}{\frac{\sqrt{3}}{3}+1}
        \]
        为一常数,于是得证点 $S$ 的轨迹是一条平行于 $BC$ 的直线。
    \end{solution}

    \question 求内接于抛物线 $y^2=4cx$ ($c>0$) 的正三角形的重心所形成的轨迹方程。
    \begin{solution}
        设内接抛物线 $\Gamma: y^2 = 4cx$ 的正三角形顶点为
        \[
        A\left(\frac{\alpha^2}{4c}, \alpha\right), 
        B\left(\frac{\beta^2}{4c}, \beta\right),
        C\left(\frac{\gamma^2}{4c}, \gamma\right),
        \]
        则重心为
        \[
        G\left(\frac{\alpha^2+\beta^2+\gamma^2}{12c},\, \frac{\alpha+\beta+\gamma}{3}\right)
        \]
        设各边斜率
        \[
        m_1 = \frac{4c(\beta-\alpha)}{(\beta^2-\alpha^2)} = \frac{4c}{\beta+\alpha},
        m_2 = \frac{4c(\gamma-\beta)}{(\gamma^2-\beta^2)} = \frac{4c}{\gamma+\beta},
        m_3 = \frac{4c(\alpha-\gamma)}{(\alpha^2-\gamma^2)} = \frac{4c}{\alpha+\gamma}
        \]
        由正三角形性质:
        \[
        \tan 60^\circ = \sqrt{3} = \frac{m_2-m_1}{1+m_1 m_2} = \frac{m_3-m_2}{1+m_2 m_3} = \frac{m_1-m_3}{1+m_3 m_1}
        \]
        化简得
        \[
        \alpha\beta + \beta\gamma + \gamma\alpha = -16c^2 - \frac{\alpha^2+\beta^2+\gamma^2}{3}
        \]
        现设 $G=(x,y)$,则
        \[
        9y^2 = (\alpha+\beta+\gamma)^2 = \alpha^2+\beta^2+\gamma^2 + 2(\alpha\beta+\beta\gamma+\gamma\alpha) = 12cx + 2(-16c^2 - 4cx) = 4cx - 32c^2
        \]
        因此正三角形重心轨迹方程为
        \[
        9y^2 = 4cx - 32c^2
        \]
    \end{solution}
    
    \question 如图,一长度为 4 单位的刚性杆 $AB$穿过位于定点 $M(1,0)$ 的铰链滑动,该铰链允许杆在 $x-y$ 平面内的任意方向转动。杆端$A$在 $y$ 轴上滑动使得 $OA \le 4$,令 $\theta$ 为杆在正 $x$ 轴的倾角。
    \begin{figure}[H]
        \centering        
        \includegraphics[width=0.5\textwidth]{images/image235.png}
    \end{figure}
    \begin{parts}
    \part 证明当 $A$ 在 $y$ 轴上滑动时,$B$ 的轨迹满足参数方程
    \[
    x=4\cos\theta, \quad y=4\sin\theta-\tan\theta, \quad -\theta_0 \le \theta \le \theta_0,
    \]
    并给出 $\theta_0$ 的值。
    \begin{solution}
        在$\triangle AOM$中,
        \[
        \cos \theta = \frac{1}{AM} \Rightarrow AM = \sec \theta
        \]
        于是$MB=4-\sec \theta$,由几何关系得到 $B$的坐标,
        \begin{align*}
        x &= OM + MB\cos\theta = 4\cos\theta, \\
        y &= MB\sin\theta = 4\sin\theta - \tan\theta
        \end{align*}
        且明显地,
        \[
        \theta_0 = \tan^{-1} 4
        \]
    \end{solution}
    \part 证明轨迹的直角坐标方程为
    \[
    y^2 = \frac{(16-x^2)(x-1)^2}{x^2}.
    \]
    \begin{solution}
        写成
        \[
        y = 4\sin\theta - \tan\theta = \sin\theta \left(4 - \frac{1}{\cos\theta}\right), 
        \]
        两边平方得
        \[
        y^2 = \sin^2\theta \frac{(4\cos\theta - 1)^2}{\cos^2\theta} = (1 - \cos^2\theta) \frac{(4\cos\theta - 1)^2}{\cos^2\theta}
        \]
        将 $\cos\theta = \dfrac{x}{4}$代入得
        \[
        y^2 = \left(1 - \frac{x^2}{16}\right) \frac{(x-1)^2}{\frac{x^2}{16}} 
        \]
        即
        \[
        y^2 = \frac{(16-x^2)(x-1)^2}{x^2}.
        \]
    \end{solution}
    \end{parts}

    \question 直线
    \[
    L:\frac{x}{p}+\frac{y}{q}=1
    \]
    其中 $p,q\neq 0$满足
    \[
    \frac{1}{p^2}+\frac{1}{q^2}=\frac{1}{2}.
    \]
    点 $P$ 为原点 $O$ 到直线 $L$ 的垂足。证明无论 $p,q$ 取何值,点 $P$ 恒在一圆 $C$ 上,并求该圆的半径。
    \begin{solution}
        直线方程$L$即
        \[
        y=q-\frac{q}{p}x
        \]
        于是$OP$直线方程为
        \[
        y=\frac{p}{q}x
        \]
        联立两式,解得点 $P$ 的坐标为
        \[
        P\left(\frac{pq^2}{p^2+q^2},\frac{p^2q}{p^2+q^2}\right)
        \]
        将已知条件
        \[
        \frac{1}{p^2}+\frac{1}{q^2}=\frac{1}{2}\Rightarrow p^2+q^2=\frac{1}{2}p^2q^2
        \]
        代入点 $P$ 的坐标得
        \[
        P\left(\frac{2}{p},\frac{2}{q}\right)
        \]
        因此点 $P$ 的轨迹满足
        \[
        x^2+y^2=4\left(\frac{1}{p^2}+\frac{1}{q^2}\right)=2
        \]
        即点 $P$ 恒在以原点为圆心、半径为 $\sqrt{2}$ 的圆上。
    \end{solution}

    \question 由椭圆 
    \[
    \frac{x^2}{a^2}+\frac{y^2}{b^2}=1
    \]
    的焦点 $F$,作椭圆上一动点$P(\cos\theta, b\sin\theta)$的切线的垂线,求该垂足 $N$ 的轨迹。
    \begin{solution}
        设 $P(a\cos\theta, b\sin\theta)$ 为椭圆上一动点,切线斜率为
        \[
        \frac{dy}{dx} = -\frac{b^2 x}{a^2 y} = -\frac{b\cos\theta}{a\sin\theta}
        \]
        切线方程为
        \[
        y - b\sin\theta = -\frac{b\cos\theta}{a\sin\theta} (x - a\cos\theta) \Rightarrow bx\cos\theta + ay\sin\theta = ab.
        \]
        过焦点 $F(ae,0)$,垂直于切线的直线 $FN$:
        \[
        y - 0 = \frac{a\sin\theta}{b\cos\theta}(x - ae) \Rightarrow -ay\cos\theta + bx\sin\theta = a^2 e \sin\theta.
        \]
        设 $N(x,y)$ 为交点,解方程组
        \[
        \begin{cases}
        bx\cos\theta + ay\sin\theta = ab, \\
        - ay\cos\theta + bx\sin\theta = a^2 e \sin\theta.
        \end{cases}
        \]
        两式平方相加得
        \[
        (bx\cos\theta + ay\sin\theta)^2 + (-ay\cos\theta + bx\sin\theta)^2 = (ab)^2 + (a^2 e \sin\theta)^2.
        \]
        \[
        x^2(b^2\cos^2\theta + b^2\sin^2\theta) + y^2(a^2\sin^2\theta + a^2\cos^2\theta) = a^2 b^2 + a^4 e^2 \sin^2\theta.
        \]
        由 $b^2 = a^2(1-e^2)$化简得垂足 $N$ 的轨迹为
        \[
        x^2 + y^2 = a^2
        \]
    \end{solution}

    \question 椭圆由参数方程
    \[
    x=3\sqrt{2}\cos\theta,\quad y=4\sin\theta,\quad 0<\theta<2\pi
    \]
    给出。
    \begin{parts}
    \part 求椭圆的焦点坐标及准线方程。
    \begin{solution}
        由已知得
        \[
        \cos\theta=\frac{x}{3\sqrt{2}},\quad \sin\theta=\frac{y}{4}
        \]
        平方相加得
        \[
        \frac{x^2}{18}+\frac{y^2}{16}=1
        \]
        因此
        \[
        a^2=18,b^2=16, \Rightarrow e=\frac{1}{3}
        \]
        焦点为
        \[
        (\pm ae,0)=(\pm\sqrt{2},0)
        \]
        准线方程为
        \[
        x=\pm\frac{a}{e}=\pm 9\sqrt{2}
        \]
    \end{solution}
    \part 证明椭圆在任意点的切线方程为
    \[
    \frac{y\sin\theta}{4}+\frac{x\cos\theta}{3\sqrt{2}}=1
    \]
    \begin{solution}
        由参数方程求导,
        \[
        \frac{dy}{dx}=\frac{4\cos\theta}{-3\sqrt{2}\sin\theta}
        \]
        在点 $(3\sqrt{2}\cos\theta,4\sin\theta)$ 处的切线为
        \[
        y-4\sin\theta=\frac{4\cos\theta}{-3\sqrt{2}\sin\theta}(x-3\sqrt{2}\cos\theta)
        \]
        整理得
        \[
        \frac{y\sin\theta}{4}+\frac{x\cos\theta}{3\sqrt{2}}=1
        \]
    \end{solution}
    \part 一条过原点的直线与上述切线相交于点 $P$,证明点 $P$ 的轨迹满足
    \[
    (x^2+y^2)^2=2(9x^2+8y^2).
    \]
    \begin{solution}
        设过原点 $O$ 且垂直于切线的直线方程为
        \[ 
        y = \frac{3\sqrt{2} \sin \theta}{4 \cos \theta} x 
        \]
        将其与切线方程联立得
        \[
        \frac{3\sqrt{2} \sin^2 \theta}{16 \cos \theta} x + \frac{\cos \theta}{3\sqrt{2}} x = 1 \Rightarrow
        x = \frac{24\sqrt{2} \cos \theta}{8 + \sin^2 \theta},y=\frac{36 \sin \theta}{8 + \sin^2 \theta}
        \]
        考虑$\dfrac{y}{x}$,可得
        \[ 
        \frac{y}{x} = \frac{3}{2\sqrt{2}} \tan \theta \Rightarrow \tan \theta = \frac{2\sqrt{2} y}{3x} 
        \]
        由恒等式 $\sec^2 \theta = 1 + \tan^2 \theta$,将 $x$ 的表达式写成
        \[ x = \frac{24\sqrt{2} \sec \theta}{8 \sec^2 \theta + \tan^2 \theta} = \frac{24\sqrt{2} \sec \theta}{9 \tan^2 \theta + 8} \]
        平方后,将$\tan^2 \theta = \dfrac{8y^2}{9x^2}$代入,
        \[
        x^2 = \frac{1152 (1 + \tan^2 \theta)}{(8 + 9 \tan^2 \theta)^2} 
        = \frac{18 \left( 1 + \frac{8y^2}{9x^2} \right)}{\left( 1 + \frac{y^2}{x^2} \right)^2}
        = \frac{2x^2 (9x^2 + 8y^2)}{(x^2 + y^2)^2}
        \]
        故$P$轨迹方程为
        \[ 
        (x^2 + y^2)^2 = 2(9x^2 + 8y^2).
        \]
    \end{solution}
    \end{parts}

    \question 已知双曲线 $H$ 与直线 $L$ 的方程分别为
    \[
    H: \frac{x^{2}}{a^{2}}-\frac{y^{2}}{b^{2}}=1, \quad L: y=mx+c,
    \]
    其中 $a,b,m,c$ 均为非零实数。
    \begin{parts}
    \part 证明直线 $L$ 与双曲线 $H$ 的交点的 $x$ 坐标满足方程
    \[
    (a^{2}m^{2}-b^{2})x^{2}+2a^{2}mcx+a^{2}(b^{2}+c^{2})=0
    \]
    \begin{solution}
        将 $y=mx+c$ 代入双曲线方程,
        \[
        \frac{x^{2}}{a^{2}}-\frac{(mx+c)^{2}}{b^{2}}=1
        \]
        展开整理得
        \[
        (a^{2}m^{2}-b^{2})x^{2}+2a^{2}mcx+a^{2}(b^{2}+c^{2})=0
        \]
    \end{solution}
    \part 已知直线 $L$ 为双曲线 $H$ 的切线,证明
    \[
    a^{2}m^{2}=b^{2}+c^{2}
    \]
    \begin{solution}
        已知 $L$ 为切线,则上述关于 $x$ 的二次方程有重根,其判别式为零:
        \[
        \Delta = 4a^{4}m^{2}c^{2}-4a^{2}(a^{2}m^{2}-b^{2})(b^{2}+c^{2})=0
        \]
        化简得
        \[
        b^{2}(b^{2}+c^{2}-a^{2}m^{2})=0
        \]
        由于 $b\neq0$,故得证
        \[
        a^{2}m^{2}=b^{2}+c^{2}
        \]
    \end{solution}
    \part 求经过点 $(1,4)$ 且与双曲线
    \[
    \frac{x^{2}}{25}-\frac{y^{2}}{16}=1
    \]
    相切的两条切线方程,并分别求出切点坐标。
    \begin{solution}
        由题意 $a^{2}=25,b^{2}=16$,切线$y=mx+c$经过点 $(1,4)$,则
        \[
        m=4-c
        \]
        由切线条件 $a^{2}m^{2}=b^{2}+c^{2}$,解得
        \[
        25(4-c)^{2}=16+c^{2} \Rightarrow c=3 \quad \text{或} \quad c=\frac{16}{3}
        \]
        对应
        \[
        m=1 \quad \text{或} \quad m=-\frac{4}{3}
        \]
        故两条切线方程为
        \[
        y=x+3, \quad y=-\frac{4}{3}x+\frac{16}{3}
        \]
        现求切点坐标。当 $y=x+3$ 时,代入双曲线解得
        \[
        x=-\frac{25}{3},\quad y=-\frac{16}{3}
        \]
        当 $y=-\dfrac{4}{3}x+\dfrac{16}{3}$ 时,代入双曲线解得
        \[
        x=\frac{25}{4},\quad y=-3
        \]
        因此两条切线及其切点分别为
        \[
        y=x+3,\; \left(-\frac{25}{3},-\frac{16}{3}\right)\quad \text{及} \quad y=-\frac{4}{3}x+\frac{16}{3},\; \left(\frac{25}{4},-3\right)
        \]
    \end{solution}
    \end{parts}

    \question 已知双曲线 
    \[
    \frac{x^2}{a^2}-\frac{y^2}{b^2}=1
    \] 
    上一点 $P$,过 $P$ 的法线分别交 $x$ 轴、$y$ 轴于 $A,B$,若 $OABQ$ 为矩形,求点 $Q$ 的轨迹方程。
    \begin{solution}
        设 $P(a\sec\theta, b\tan\theta)$ 为双曲线上的动点,有
        \[
        \frac{dy}{dx} = \frac{b^2 x}{a^2 y}.
        \]
        在 $P$ 点的切线斜率为$\dfrac{b\sec\theta}{a\tan\theta}$,法线斜率为$-\dfrac{a\tan\theta}{b\sec\theta}$,故$P$处的法线方程为
        \[
        y - b\tan\theta = -\frac{a\tan\theta}{b\sec\theta} (x - a\sec\theta) \Rightarrow a\tan\theta\, x + b\sec\theta\, y - \frac{ab}{\cos\theta} = 0
        \]
        与坐标轴交点分别为
        \[
        A\left(\frac{a^2+b^2}{a}\sec\theta, 0\right),\quad B\left(0, -\frac{a^2+b^2}{b}\tan\theta\right)
        \]
        由矩形性质,$Q$ 点坐标为
        \[
        Q\left(\frac{a^2+b^2}{a}\sec\theta, -\frac{a^2+b^2}{b}\tan\theta \right)
        \]
        利用恒等式 $\sec^2\theta - \tan^2\theta = 1$,得点 $Q$ 的轨迹方程为
        \[
        \left(\frac{ax}{a^2+b^2}\right)^2 - \left(\frac{by}{a^2+b^2}\right)^2 = 1 \Rightarrow a^2 x^2 - b^2 y^2 = (a^2+b^2)^2
        \]
    \end{solution}

    \question 动点 $P$ 在矩形双曲线
    \[
    xy=a^{2}
    \]
    上运动,其中 $a>0$ 为常数。过点 $P$ 作双曲线的法线,该法线再次与双曲线相交于点 $Q$。设 $M$ 为线段 $PQ$ 的中点。求点 $M$ 的轨迹方程。
    \begin{solution}
        设点 $P\left(ap, \frac{a}{p}\right)$,其中 $p \neq 0$。由链导法,在点 $P(ap, \frac{a}{p})$ 处切线的斜率为
        \[
        \frac{dy}{dx}=\frac{dy}{dp} \div \frac{dx}{dp} = \frac{-\frac{a}{p^2}}{a}=-\frac{1}{p^{2}}
        \]
        从而 $P$ 处法线方程为
        \[
        y - \frac{a}{p} = p^{2}(x - ap)
        \]
        与$xy=a^{2}$联立解得
        \[
        \frac{a^{2}}{x} - \frac{a}{p} = p^{2}(x - ap) \Rightarrow Q\left(-\frac{a}{p^{3}}, -ap^{3}\right)
        \]
        其中$x = ap$对应点$P$;点 $M(x, y)$ 为 $PQ$ 的中点, 其坐标为
        \[
        x = \frac{ap - \frac{a}{p^{3}}}{2} = \frac{a(p^{4}-1)}{2p^{3}}, \quad y = \frac{\frac{a}{p} - ap^{3}}{2} = \frac{a(1-p^{4})}{2p}
        \]
        两式相除得
        \[
        p^{2} = -\frac{y}{x}
        \]
        将之带入$y$ 的表达式平方得
        \[
        y^{2} = \frac{a^{2}\left(1 - \frac{y^{2}}{x^{2}}\right)^{2}}{4\left(-\frac{y}{x}\right)} = \frac{a^{2}(x^{2}-y^{2})^{2}}{-4x^{3}y}
        \]
        整理得点 $M$ 的轨迹方程为
        \[
        a^{2}(x^{2}-y^{2})^{2} + 4x^{3}y^{3} = 0
        \]
    \end{solution}

    \question 点 $P,Q$ 为曲线
    \[
    xy=1,\ x\in\mathbb{R} \setminus \{0\}
    \]
    上的两个相异点,使得线段 $PQ$ 为曲线在 $P$ 处的法线。曲线在 $P,Q$ 处的切线相交于点 $R$。证明点 $R$ 的轨迹方程为
    \[
    (y^2-x^2)^2+4xy=0。
    \]
    \begin{solution}
        对曲线$y=\frac{1}{x}$求导得
        \[
        \frac{dy}{dx}=-\frac{1}{x^2}
        \]
        设$P\left(p,\dfrac{1}{p}\right), Q\left(q,\dfrac{1}{q}\right), p\neq q$,弦 $PQ$ 的斜率为
        \[
        \frac{\frac{1}{p}-\frac{1}{q}}{p-q}=-\frac{1}{pq}
        \]
        曲线在 $P$ 处切线斜率为 $-\frac{1}{p^2}$,故法线斜率为 $p^2$。
        由于 $PQ$ 为 $P$ 处的法线,有
        \[
        \left(-\frac{1}{pq}\right)\left(-\frac{1}{p^2}\right)=-1 \Rightarrow q=-\frac{1}{p^2}
        \]
        曲线在 $P$ 处的切线方程为
        \[
        y-\frac{1}{p}=-\frac{1}{p^2}(x-p) \Rightarrow y=\frac{2}{p}-\frac{1}{p^2}x
        \]
        同理,曲线在 $Q$ 处的切线方程为
        \[
        y=\frac{2}{q}-\frac{1}{q^2}x
        \]
        联立两切线方程解得交点为
        \[
        R\left(\frac{2pq}{p+q},\frac{2}{p+q}\right)
        \]
        由 $q=-\dfrac{1}{p^2}$,
        \[
        x = \frac{2p\left(-\frac{1}{p^2}\right)}{p-\frac{1}{p^2}}= -\frac{2p}{p^3-1},\quad
        y = \frac{2}{p-\frac{1}{p^2}}= \frac{2p^2}{p^3-1}
        \]
        于是由
        \[
        p=-\frac{y}{x}
        \]
        可得点 $R$ 的轨迹方程
        \[
        (y^2-x^2)^2+4xy=0
        \]
    \end{solution}

    \question 已知动点
    \[
    P\left(5t,\frac{5}{t}\right),\ t\neq 0
    \]
    在双曲线
    \[
    xy=25
    \]
    上。
    \begin{parts}
    \part 证明该双曲线在点 $P$ 处的法线方程为
    \[
    y=t^2x+\frac{5}{t}-5t^3
    \]
    \begin{solution}
        对双曲线$y=\dfrac{25}{x}$求导得
        \[
        \frac{dy}{dx}=-\frac{25}{x^2}
        \]
        在$P$处的斜率为
        \[
        \frac{dy}{dx}=-\frac{25}{(5t)^2}=-\frac{1}{t^2}
        \]
        故法线斜率为 $t^2$。点 $P\left(5t,\dfrac{5}{t}\right)$ 处法线方程为
        \[
        y-\frac{5}{t}=t^2(x-5t) \Rightarrow  y=t^2x+\frac{5}{t}-5t^3
        \]
    \end{solution}
    \part 该法线再次与双曲线相交于点 $Q$,证明点 $Q$ 的坐标为
    \[
    \left(-\frac{5}{t^3},-5t^3\right)
    \]
    \begin{solution}
        联立
        \[
        xy=25,\quad y=t^2x+\frac{5}{t}-5t^3
        \]
        化简得
        \[
        t^3x^2+(5-5t^4)x-25t=0
        \]
        由于 $x=5t$ 为一根,因式分解得
        \[
        (x-5t)(t^3x+5)=0
        \]
        故另一交点满足
        \[
        x=-\frac{5}{t^3} \Rightarrow y=\frac{25}{x}=-5t^3
        \]
        因此得证
        \[
        Q\left(-\frac{5}{t^3},-5t^3\right)
        \]
    \end{solution}
    \part 证明线段 $PQ$ 的中点轨迹的笛卡尔方程为
    \[
    4xy+25\left(\frac{y}{x}-\frac{x}{y}\right)^2=0
    \]
    \begin{solution}
        设 $M$ 为 $PQ$ 的中点,则
        \[
        M\left(\frac{5}{2}\left(t-\frac{1}{t^3}\right),\frac{5}{2}\left(\frac{1}{t}-t^3\right)\right)
        \]
        设
        \[
        x=\frac{5}{2}\left(t-\frac{1}{t^3}\right),\quad
        y=\frac{5}{2}\left(\frac{1}{t}-t^3\right)
        \]
        则
        \[
        \frac{x}{y}=-\frac{1}{t^2} \Rightarrow t^2=-\frac{y}{x}
        \]
        将$y$表达式平方得
        \[
        4y^2t^2=25(1-t^4)^2
        \]
        代入$t^2=-\dfrac{y}{x}$得,
        \[
        -4xy^3=25\frac{(x^2-y^2)^2}{x^3}
        \]
        即
        \[
        4xy+25\left(\frac{y}{x}-\frac{x}{y}\right)^2=0
        \]
    \end{solution}
    \end{parts}

    \question 曲线由参数方程给出
    \[
    x = \sin^2 t, \quad y = \sin t \cos t + \cos t, \quad 0 \le t < 2\pi.
    \]
    证明其直角坐标方程为
    \[
    (x^2 + y^2 - 1)^2 = 4x(1-x)^2.
    \]
    \begin{solution}
        将 $y$ 写为
        \[
        y = \cos t (\sin t + 1) \implies y^2 = \cos^2 t (\sin t + 1)^2.
        \]
        利用恒等式 $\cos^2 t = 1 - \sin^2 t$,并令 $x = \sin^2 t$:
        \[
        y^2 = (1 - x)(x + 2\sin t + 1) \Rightarrow \frac{y^2}{1-x} - (1+x) = 2\sin t
        \]
        两边平方,
        \[
        \left( \dfrac{y^2 - (1-x^2)}{1-x} \right)^2 = 4 \sin^2 t = 4x
        \]
        于是所求直角坐标方程为
        \[
        (y^2 + x^2 - 1)^2 = 4x (1-x)^2
        \]
    \end{solution}

    \question 曲线 $C$ 由参数方程
    \[
    x = \tan\theta - \sec\theta, \quad y = \cot\theta - \csc\theta, \quad 0 < \theta < \frac{\pi}{2}.
    \]
    给出,证明
    \begin{parts}
    \part $C$ 的直角坐标方程为
    \[
    (x^2 - 1)(y^2 - 1) = 4xy,
    \]
    \begin{solution}
        从 $x = \tan\theta - \sec\theta$ 开始:
        \begin{align*}
        x^2 &= (\tan\theta - \sec\theta)^2 \\
        &= \tan^2\theta - 2\tan\theta\sec\theta + \sec^2\theta \\
        &= \tan^2\theta - 2\tan\theta\sec\theta + (1+\tan^2\theta) \\
        &= 2\tan^2\theta - 2\tan\theta\sec\theta + 1 \\
        &= 2\tan\theta(\tan\theta - \sec\theta) + 1 \\
        &= 2\tan\theta \cdot x + 1
        \end{align*}
        所以
        \[
        \tan\theta = \frac{x^2-1}{2x}
        \]
        类似地,由 $y = \cot\theta - \csc\theta$ 和恒等式 $1+\cot^2\theta = \csc^2\theta$ 得
        \[
        \cot\theta = \frac{y^2-1}{2y}
        \]
        于是
        \[
        \tan\theta \cot\theta = \frac{x^2-1}{2x} \cdot \frac{y^2-1}{2y} = 1 \Rightarrow (x^2-1)(y^2-1) = 4xy
        \]
    \end{solution}
    \part 
    \[
    \frac{dy}{dx} = \frac{1-y^2}{2x}.
    \]
    \begin{solution}
        参数求导得
        \[
        \frac{dx}{d\theta} = \sec^2\theta - \sec\theta\tan\theta = \sec\theta(\sec\theta - \tan\theta)
        \]
        \[
        \frac{dy}{d\theta} = -\csc^2\theta + \csc\theta \cot\theta = \csc\theta(\cot\theta - \csc\theta)
        \]
        所以
        \[
        \frac{dy}{dx} = \frac{dy}{d\theta} \div \frac{dx}{d\theta} = \frac{\csc\theta(\cot\theta - \csc\theta)}{\sec\theta(\sec\theta - \tan\theta)}
        \]
        注意到
        \[
        \frac{\csc\theta}{\sec\theta} = \frac{\cos\theta}{\sin\theta} = \cot\theta,
        \]
        且
        \[
        \frac{\cot\theta - \csc\theta}{\sec\theta - \tan\theta} = -\frac{y}{x}
        \]
        得到
        \[
        \frac{dy}{dx} = -\frac{y}{x} \cdot \frac{y^2-1}{2y} = \frac{1-y^2}{2x}
        \]
        故得证。
    \end{solution}
    \end{parts}

    \question 曲线的参数方程为
    \[
    x = \sin\theta, \quad y = \theta \cos\theta, \quad -\pi < \theta < \pi.
    \]
    已知曲线在 $\theta = -\dfrac{\pi}{4}$ 及 $\theta = \dfrac{\pi}{4}$ 处的切线平行,且间距为 $d$,证明:
    \[
    d = \sqrt{\frac{8\pi^2 - 32\pi + 32}{\pi^2 - 8\pi + 32}}.
    \]
    \begin{solution}
        先求曲线的导数,
        \[
        \frac{dx}{d\theta} = \cos\theta, \quad \frac{dy}{d\theta} = \cos\theta - \theta\sin\theta \implies \frac{dy}{dx} = \frac{dy/d\theta}{dx/d\theta} = 1 - \theta\tan\theta.
        \]
        当 $\theta=\pm \frac{\pi}{4}$ ,
        \[
        x = \pm \frac{1}{\sqrt{2}}, \quad y = \pm \frac{\pi}{4\sqrt{2}}, \quad \frac{dy}{dx} = 1 - \frac{\pi}{4}.
        \]
        故在$\theta=\pm \frac{\pi}{4}$处的切线方程为
        \[
        y - \frac{\pi}{4\sqrt{2}} = \left(1 - \frac{\pi}{4}\right)\left(x - \frac{1}{\sqrt{2}}\right), \quad y + \frac{\pi}{4\sqrt{2}} = \left(1 - \frac{\pi}{4}\right)\left(x + \frac{1}{\sqrt{2}}\right).
        \]
        两切线在 $x=0$ 时的纵坐标差为
        \[
        y_1 = \frac{\sqrt{2}(2-\pi)}{4}, y_2 = \frac{\sqrt{2}(\pi-2)}{4} \Rightarrow |y_2 - y_1| = \frac{\sqrt{2}(\pi-2)}{2}
        \]
        且切线与水平的夹角 $\theta$ 满足
        \[
        \cos\theta = \frac{1}{\sqrt{1+(\frac{dy}{dx})^2}} = \frac{4}{\sqrt{\pi^2 - 8\pi + 32}}.
        \]
        故切线间距为
        \[
        d = |y_2 - y_1| \cos\theta = \frac{\sqrt{2}(\pi-2)}{2} \cdot \frac{4}{\sqrt{\pi^2 - 8\pi + 32}} = \sqrt{\frac{8\pi^2 - 32\pi + 32}{\pi^2 - 8\pi + 32}}.
        \]
    \end{solution}

    \question 曲线 $C$ 的方程为
    \[
    x^2 + xy + y^2 = 1, \quad 0 \le x \le 3.
    \]
    考虑参数化
    \[
    x = A\cos\theta + B\sin\theta, \quad y = A\cos\theta - B\sin\theta,
    \]
    确定 $A$ 和 $B$,并求出第一象限被曲线和坐标轴围成的有限区域的面积。
    \begin{solution}
        尝试 $x = \cos\theta + \sin\theta$, $y = \cos\theta - \sin\theta$,则
        \[
        (x+y)^2 - xy = (2\cos\theta)^2 - (\cos^2\theta-\sin^2\theta) 
        = 3\cos^2\theta + \sin^2\theta \neq 1
        \]
        因此调整成
        \[
        x = \frac{1}{\sqrt{3}}\cos\theta + \sin\theta, \quad y = \frac{1}{\sqrt{3}}\cos\theta - \sin\theta.
        \]
        发现
        \[
        (x+y)^2 - xy = \left(\frac{2}{\sqrt{3}}\cos\theta\right)^2 - \left(\frac{1}{3}\cos^2\theta - \sin^2\theta\right)
        = \cos^2\theta + \sin^2\theta = 1
        \]
        成功！则所求面积为
        \begin{align*}
        \int_{-\frac{\pi}{6}}^{\frac{\pi}{6}} y(\theta) \frac{dx}{d\theta}\, d\theta
        &= \int_{-\frac{\pi}{6}}^{\frac{\pi}{6}} \left(\frac{1}{\sqrt{3}}\cos\theta - \sin\theta\right)\left(-\frac{1}{\sqrt{3}}\sin\theta + \cos\theta\right)\, d\theta \\
        &= \int_{-\frac{\pi}{6}}^{\frac{\pi}{6}} \frac{\sqrt{3}}{3}(\cos^2\theta + \sin^2\theta) \, d\theta \\
        &= \frac{\sqrt{3}}{3} \cdot \frac{\pi}{3} = \frac{\pi \sqrt{3}}{9}
        \end{align*}
    \end{solution}

    \question 已知曲线的参数方程为
    \[
    x = t^3, \quad y = t^2, \quad t \in \mathbb{R}.
    \]
    曲线在点 $P$ 处的切线再次与曲线相交于点 $Q$。如图,曲线与切线所围成的区域面积为 $2.7$ 平方单位,求点 $P$ 的坐标。
    \begin{figure}[H]
        \centering        
        \includegraphics[width=0.5\textwidth]{images/image237.png}
    \end{figure}
    \begin{solution}
        设参数 $t = p$对应 $P(p^3, p^2)$,切线斜率为
        \[
        \frac{dy}{dx} 
        = \frac{dy}{dt} \div \frac{dx}{dt} = \frac{2t}{3t^2} = \frac{2}{3t} \Rightarrow \frac{dy}{dx}\bigg|_{t=p} = \frac{2}{3p}
        \]
        切线方程
        \[
        y - p^2 = \frac{2}{3p} (x - p^3) \Rightarrow 3p y = 2x + p^3.
        \]
        与曲线方程联立得
        \[
        2 t^3 - 3p t^2 + p^3 = 0
        \]
        因式分解为
        \[
        (t - p)^2 (2 t + p) = 0 \Rightarrow Q \left(-\frac{1}{8} p^3, \frac{1}{4} p^2\right)
        \]
        阴影部分面积即梯形面积减去曲线下的面积,而曲线下的面积为
        \[
        \int_{-\frac{p}{2}}^{p} t^2 (3t^2) \, dt = \int_{-\frac{p}{2}}^{p} 3t^4 \, dt = \left[ \frac{3}{5}t^5 \right]_{-\frac{p}{2}}^{p} = \frac{99}{160}p^5
        \]
        梯形面积为
        \[
        \frac{1}{2} \left(\frac{1}{4}p^2 + p^2\right)\cdot \left(p^3 + \frac{1}{8}p^3\right) = \frac{45}{64}p^5
        \]
        于是
        \[
        \frac{45}{64}p^5 - \frac{99}{160}p^5 = \frac{27}{10} \Rightarrow p = 2
        \]
        所以点 $P$ 的坐标为 
        \[
        (2^3, 2^2) = (8, 4)
        \]
    \end{solution}

    \question 已知曲线的参数方程为
    \[
    x = \sin\left(t + \frac{\pi}{6}\right), \quad y = 1 + \cos 2t, \quad 0 \le t < 2\pi.
    \]
    且曲线关于 $y$ 轴对称,求曲线两环所围面积。
    \begin{figure}[H]
        \centering        
        \includegraphics[width=0.5\textwidth]{images/image238.png}
    \end{figure}
    \begin{solution}
        利用参数积分,两侧环的面积为
        \[
        I = 2 \int_{\frac{\pi}{6}}^{\frac{\pi}{2}} (1+\cos 2t) \cos\left(t+\frac{\pi}{6}\right) \, dt.
        \]
        积化和差得
        \[
        \cos 2t \cos\left(t+\frac{\pi}{6}\right) = \frac{1}{2} \left[ \cos\left(3t + \frac{\pi}{6}\right) + \cos\left(t - \frac{\pi}{6}\right) \right]
        \]
        于是
        \begin{align*}
        I &= 2 \int_{\frac{\pi}{6}}^{\frac{\pi}{2}} \left[ \cos\left(t+\frac{\pi}{6}\right) + \frac{1}{2} \cos\left(3t+\frac{\pi}{6}\right) + \frac{1}{2} \cos\left(t-\frac{\pi}{6}\right) \right] \, dt \\
        &= 2 \left[ \sin\left(t+\frac{\pi}{6}\right) + \frac{1}{6}\sin\left(3t+\frac{\pi}{6}\right) + \frac{1}{2}\sin\left(t-\frac{\pi}{6}\right) \right]_{\frac{\pi}{6}}^{\frac{\pi}{2}} \\
        &= \frac{4\sqrt{3}}{3}
        \end{align*}
    \end{solution}

    \question 已知曲线 $C$ 的参数方程为
    \[
    x = t \sin t, \quad y = \cos t, \quad 0 \le t < 2\pi.
    \]
    曲线与坐标轴交于 $P,Q,R,S$ 四点。
    \begin{figure}[H]
        \centering        
        \includegraphics[width=0.5\textwidth]{images/image239.png}
    \end{figure}
    \begin{parts}
    \part 求在点 $P,Q,R,S$ 处的 $t$ 的值。
    \begin{solution}
        设 $x = 0$,
        \[
        t_P = \pi \implies P(0,-1), \quad t_Q = 0 \implies Q(0,1)
        \]
        设 $y = 0$,
        \[
        t_R = \frac{3\pi}{2} \implies R\left(-\frac{3\pi}{2},0\right), \quad
        t_S = \frac{\pi}{2} \implies S\left(\frac{\pi}{2},0\right)
        \]
    \end{solution}
    \part 求曲线沿着$P,Q,R,S$所围成的区域面积。
    \begin{solution}
        由参数积分公式
        \[
        I = \int_{0}^{2\pi} y(t) \frac{dx}{dt} \, dt.
        \]
        被积函数为
        \[
        y(t)\frac{dx}{dt} = \cos t (\sin t + t \cos t) = \cos t \sin t + t \cos^2 t = \frac{1}{2}\sin 2t + \frac{1}{2} t + \frac{1}{2} t \cos 2t
        \]
        对 $\dfrac{1}{2} t \cos 2t$ 分部积分得
        \[
        \int t \cos 2t \, dt = \frac{1}{2} t \sin 2t - \frac{1}{2} \int \sin 2t \, dt = \frac{1}{2} t \sin 2t + \frac{1}{4} \cos 2t + C.
        \]
        因此
        \[
        \int y \frac{dx}{dt} \, dt = \frac{1}{4} t^2 - \frac{1}{8} \cos 2t + \frac{1}{4} t \sin 2t + C.
        \]
        故所围成的面积为
        \[
        I = \left[ \frac{1}{4} t^2 - \frac{1}{8} \cos 2t + \frac{1}{4} t \sin 2t \right]_0^{2\pi} = \pi^2
        \]
    \end{solution}
    \end{parts}

    \question 如图所示,曲线的参数方程为
    \[ 
    x = t^2 + 2t, \quad y = t^3 - 9t, \quad t \in \mathbb{R} 
    \]
    该曲线与坐标轴交于原点 $O$ 以及点 $A,B,C$。
    \begin{figure}[H]
        \centering        
        \includegraphics[width=0.5\textwidth]{images/image236.png}
    \end{figure}
    \begin{parts}
    \part 确定 $A$,$B$ 和 $C$ 的坐标。
    \begin{solution}
        令 $x=0$,则
        \[
        0 = t(t+2) \Rightarrow t = 0 \ (O), \quad t = -2 \ (C),
        \]
        因此$C(0,10)$,令 $y=0$,
        \[
        0 = t(t^2-9) = t(t-3)(t+3) \Rightarrow t = 0 \ (O), \quad t=3 \ (B), \quad t=-3 \ (A),
        \]
        因此 $A(3,0),B(15,0)$。
    \end{solution}
    \part 图中阴影部分的区域 $R$ 由曲线和坐标轴围成,求区域 $R$ 的面积。
    \begin{solution}
        由参数积分公式,面积为
        \begin{align*}
        \int_{t_1}^{t_2} y(t) \frac{dx}{dt} \, dt
        &= \int_{-3}^{-2} (t^3-9t)(t+2) \, dt \\
        &= \int_{-3}^{-2} (t^4 + 2t^3 - 9t^2 - 18t) \, dt\\
        &= \left[\frac{t^5}{5} + \frac{1}{2} t^4 - 3 t^3 - 9 t^2\right]_{-3}^{-2} \\
        &= \frac{171}{10}
        \end{align*}
    \end{solution}
    \part 由曲线和 $y$ 轴围成的满足 $x < 0$ 的有界区域,绕 $x$ 轴旋转一周形成一个旋转体。计算该旋转体的体积。
    \begin{solution}
        设曲线最左端处的参数值为$t_1$,旋转体体积为
        \begin{align*}
        &\pi \int_{t_1}^{-2} [y(t)]^2 \frac{dx}{dt} \, dt - \pi \int_{t_1}^{0} [y(t)]^2 \frac{dx}{dt} \, dt \\
        &= \pi \int_{0}^{-2} (t^3-9t)^2 (t^2+2t) \, dt \\
        &= 2\pi \int_{0}^{-2} (t^7 - 18 t^5 + 81 t^3 + t^6 - 18 t^4 + 81 t^2) \, dt\\
        &= 2\pi \left[\frac{t^8}{8} + \frac{t^7}{7} - 3 t^6 - \frac{18}{5} t^5 + \frac{81}{4} t^4 + 27 t^3\right]_{0}^{-2}\\
        &= \frac{3144 \pi}{35}
        \end{align*}
    \end{solution}
    \end{parts}
\end{questions}
\pagebreak

\begin{center}
  {\fontsize{30pt}{26pt}\selectfont
    \hypertarget{极坐标}{极坐标} \label{极坐标}
  }
\end{center}
\separator
\vspace{1pt}

\begin{questions}
    \question 曲线 $C$ 的极坐标方程为
    \[
    r=\tan\theta,\quad 0\le\theta<\frac{\pi}{2}.
    \]
    求曲线 $C$ 的直角坐标方程,并写成 $y=f(x)$ 的形式。
    \begin{solution}
        由已知$r=\tan\theta$,两边平方得
        \[
        r^2=\tan^2\theta \Rightarrow r^2\cos^2\theta=\sin^2\theta
        \]
        由 $r\cos\theta=x$,得
        \[
        x^2=\sin^2\theta
        \]
        至此,欲以$\tan^2\theta$表示$\sin^2\theta$,不妨考虑
        \[
        1+\tan^2\theta=\sec^2\theta=\frac{1}{\cos^2 \theta}=\frac{1}{1-\sin^2 \theta}
        \]
        据此有
        \[
        \tan^2\theta=\frac{1}{1-x^2}-1=\frac{x^2}{1-x^2}
        \]
        又有 $r=\tan\theta$及$r^2=x^2+y^2$,于是
        \[
        x^2+y^2=\frac{x^2}{1-x^2} \Rightarrow y^2=\frac{x^4}{1-x^2}
        \]
        因为 $0\le\theta<\dfrac{\pi}{2}$,有 $y\ge0$,故
        \[
        y=\frac{x^2}{\sqrt{1-x^2}},
        \]
        其中 $0\le x<1$。
    \end{solution}

    \question 心脏线的极坐标方程为 
    \[
    r=4(1+\cos \theta),0 \le \theta \le \frac{\pi}{2}
    \]
    曲线在点 $P$ 处的切线斜率为 $-1$。求该切线的极坐标方程,结果以 $r=f(\theta)$ 的形式表示。
    \begin{solution}
        由链导法,
        \[ 
        \frac{dy}{dx} = \frac{dy/d\theta}{dx/d\theta} = \frac{\frac{d}{d\theta}(r \sin \theta)}{\frac{d}{d\theta}(r \cos \theta)} = \frac{\frac{d}{d\theta}[4(1+\cos \theta)\sin \theta]}{\frac{d}{d\theta}[4(1+\cos \theta)\cos \theta]} 
        \]
        展开并求导得
        \[ 
        \frac{dy}{dx} = \frac{\cos \theta + \cos^2 \theta - \sin^2 \theta}{-\sin \theta - 2\sin \theta \cos \theta} = \frac{\cos \theta + \cos 2\theta}{-\sin \theta - \sin 2\theta} 
        \]
        令斜率为 $-1$,
        \[ 
        \frac{\cos \theta + \cos 2\theta}{-\sin \theta - \sin 2\theta} = -1 \Rightarrow \cos \theta + \cos 2\theta = \sin \theta + \sin 2\theta 
        \]
        和差化积得
        \[ 
        2\cos \frac{3\theta}{2} \cos \frac{\theta}{2} = 2\sin \frac{3\theta}{2} \cos \frac{\theta}{2} \Rightarrow
        \cos \frac{\theta}{2} \left( \cos \frac{3\theta}{2} - \sin \frac{3\theta}{2} \right) = 0 
        \]
        在范围 $0 \le \theta \le \dfrac{\pi}{2}$ 内,$\cos \dfrac{\theta}{2} \neq 0$,故有
        \[ 
        \tan \frac{3\theta}{2} = 1 \Rightarrow \theta = \frac{\pi}{6} 
        \]
        此时
        \[ 
        r = 4\left(1 + \cos \frac{\pi}{6}\right) = 4 + 2\sqrt{3} 
        \]
        点 $P$ 的直角坐标为
        \[ 
        P\left(r \cos \frac{\pi}{6},r \sin \frac{\pi}{6}\right) = (2\sqrt{3} + 3 ,2 + \sqrt{3})
        \]
        故切线的方程为 
        \[ 
        y - (2 + \sqrt{3}) = -(x - (3 + 2\sqrt{3})) \implies y + x = 5 + 3\sqrt{3} 
        \]
        转换为极坐标方程:
        \[ 
        r\sin \theta + r\cos \theta = 5 + 3\sqrt{3} \Rightarrow r = \frac{5 + 3\sqrt{3}}{\cos \theta + \sin \theta} 
        \]
    \end{solution}

    \question 曲线 $C_1,C_2$ 的极坐标方程分别为
    \[
    C_1:\ r=2\cos\theta-\sin\theta,\quad 0<\theta\le\frac{\pi}{3},
    \]
    \[
    C_2:\ r=\sqrt{2}+\sin\theta,\quad 0\le\theta<2\pi.
    \]
    点 $P$ 在 $C_1$ 上,且 $P$ 点处的切线与极轴平行。
    \begin{parts}
    \part 证明在点 $P$ 处有
    \[
    \tan 2\theta=2.
    \]
    \begin{solution}
        $P$ 点处的切线与极轴平行,意味
        \[
        \frac{dy}{dx}=0\Rightarrow\frac{dy}{d\theta}=0
        \]
        即
        \[
        \frac{d}{d\theta}(r\sin\theta)=\frac{d}{d\theta}\left((2\cos\theta-\sin\theta)\sin\theta\right)=\frac{d}{d\theta}(\sin 2\theta-\sin^2\theta)=0
        \]
        对 $\theta$ 求导得
        \[
        2\cos 2\theta-2\sin\theta\cos\theta=0.
        \]
        即
        \[
        2\cos 2\theta=\sin 2\theta \Rightarrow \tan 2\theta=2
        \]
    \end{solution}
    \part 证明点 $P$ 到原点 $O$ 的距离为
    \[
    \sqrt{\frac{5-\sqrt{5}}{2}}.
    \]
    \begin{solution}
        由 $\tan 2\theta=2$,利用倍角公式
        \[
        \tan 2\theta=\frac{2\tan\theta}{1-\tan^2\theta}
        \]
        解得
        \[
        \tan\theta=\frac{-1\pm\sqrt{5}}{2}
        \]
        由于 $0<\theta<\dfrac{\pi}{3}$,取
        \[
        \tan\theta=\frac{-1+\sqrt{5}}{2}
        \]
        由此可得
        \[
        \sin\theta=\frac{-1+\sqrt{5}}{\sqrt{10-2\sqrt{5}}},\quad
        \cos\theta=\frac{2}{\sqrt{10-2\sqrt{5}}}
        \]
        代入 $r=2\cos\theta-\sin\theta$,得到
        \[
        r_P=\frac{5-\sqrt{5}}{\sqrt{10-2\sqrt{5}}}=\sqrt{\frac{5-\sqrt{5}}{2}}
        \]
    \end{solution}
    \part 设点 $Q$ 为 $C_1$ 与 $C_2$ 的交点,求点 $Q$ 处对应的 $\theta$ 值。
    \begin{solution}
        联立$C_1,C_2$得
        \[
        \cos\theta-\sin\theta=\frac{\sqrt{2}}{2}.
        \]
        于是
        \[
        \sqrt{2}\cos\left(\theta+\frac{\pi}{4}\right)=\frac{\sqrt{2}}{2},
        \]
        即
        \[
        \cos\left(\theta+\frac{\pi}{4}\right)=\frac{1}{2}
        \]
        在 $0\le\theta\le\dfrac{\pi}{3}$ 内解得
        \[
        \theta+\frac{\pi}{4}=\frac{\pi}{3} \Rightarrow \theta=\frac{\pi}{12}
        \]
    \end{solution}
    \end{parts}

    \question 曲线 $C_1,C_2$ 的极坐标方程分别为 
    \[
    r_1 = 3 + 2\cos \theta, 0 \le \theta < 2\pi,\quad r_2 = 2
    \]
    两曲线交于点 $P,Q$,经过 $P$ 点和极点 $O$ 的直线与 $C_1$ 再次交于 $R$ 点,证明 $RQ$ 是 $C_2$ 在 $Q$ 点处的切线。
    \begin{figure}[H]
        \centering        
        \includegraphics[width=0.5\textwidth]{images/image240.png}
    \end{figure}
    \begin{solution}
        联立$C_1,C_2$得
        \[ 
        3 + 2\cos \theta = 2 \Rightarrow \theta = \frac{2\pi}{3},\frac{4\pi}{3} \Rightarrow P\left(2, \frac{2\pi}{3}\right),Q\left(2, \frac{4\pi}{3}\right)
        \]
        由于 $PR$ 通过极点 $O$,则 $R$ 点的角度为 $\dfrac{2\pi}{3} + \pi = \dfrac{5\pi}{3}$,代入 $C_1$ 方程:
        \[
        r_R = 3 + 2\cos(\frac{5\pi}{3}) = 4
        \]
        故 $R(4, \frac{5\pi}{3})$,在 $\triangle OQR$ 中,由余弦定理,
        $\angle QOR = \frac{5\pi}{3} - \frac{4\pi}{3} = \frac{\pi}{3}$。
        \[ 
        QR^2 = 2^2 + 4^2 - 2 \cdot 2 \cdot 4 \cos \left(\frac{5\pi}{3} - \frac{4\pi}{3}\right) \Rightarrow QR = \sqrt{12} 
        \]
        此时发现
        \[
        QR^2+OQ^2=OR^2=16
        \]
        由毕氏定理,$\angle OQR = \dfrac{\pi}{2}$,即$RQ$ 与圆 $C_2$ 在 $Q$ 点垂直,故$RQ$ 是 $C_2$ 在 $Q$ 点处的切线。
    \end{solution}

    \question 已知两曲线极坐标方程
    \[
    C_1: r = 12 \cos \theta, \quad -\frac{\pi}{2}< \theta \le \frac{\pi}{2}, \quad C_2: r = 4 + 4 \cos \theta, \quad -\pi < \theta \le \pi
    \]
    已知 $C_1$ 与 $C_2$ 的一个交点为 $A$,由 $C_1$ 与线段 $OA$ 围成的较小区域面积为 $6\pi - 9\sqrt{3}$。若有限区域 $R$ 表示在 $C_1$ 内但在 $C_2$ 外的点,证明 $R$ 的面积为 $16\pi$。
    \begin{solution}
        首先求交点:
        \[ 
        12 \cos \theta = 4 + 4 \cos \theta \Rightarrow \theta = \pm \frac{\pi}{3} \Rightarrow
        \]
        记$A\left(6,\frac{\pi}{3}\right)$,所求面积即
        \[ 
        2 \cdot \left[ \frac{1}{2}\pi(6)^2 - \left( \int_{0}^{\frac{\pi}{3}} \frac{1}{2}(4 + 4 \cos \theta)^2 \, d\theta + 6\pi - 9\sqrt{3} \right) \right] =16\pi
        \]
        其中
        \begin{align*}
        \int_{0}^{\frac{\pi}{3}} \frac{1}{2}(4 + 4 \cos \theta)^2 \, d\theta = 4\pi + 9\sqrt{3} 
        &= \int_{0}^{\frac{\pi}{3}} \left(8 + 16 \cos\theta+8\cos^2\theta\right) \, d\theta \\
        &= \int_{0}^{\frac{\pi}{3}} \left(12 + 16 \cos\theta+4\cos 2\theta\right) \, d\theta \\
        &= \left[12\theta +16\sin \theta +2 \sin 2 \theta\right]_{0}^{\frac{\pi}{3}} \\
        &= 4\pi + 9\sqrt{3}
        \end{align*}
    \end{solution}

    \question 曲线 $C_1,C_2$ 的极坐标方程分别为
    \[ 
    r = 1 + \sin \theta, \quad 0 < \theta < \frac{\pi}{2}, \quad r = 1 + \cos 2\theta, \quad 0 < \theta < \frac{\pi}{2}
    \]
    点 $P$ 是 $C_1,_2$ 的交点。一条平行于极轴的直线通过点 $P$ 并与曲线 $C_2$ 交于点 $Q$。证明
    \[ 
    PQ = \frac{1}{32} \left[ 24\sqrt{3} - (2 + 2\sqrt{13})^{\frac{3}{2}} \right] 
    \]
    \begin{solution}
        联立方程求交点 $P$ 的极坐标:
        \[ 
        1 + \sin \theta = 1 + \cos 2\theta \Rightarrow \theta = \frac{1}{2}
        \]
        故 $P\left(\dfrac{3}{2}, \dfrac{\pi}{6}\right)$,过 $P$且平行于极轴的直线极坐标方程为 
        \[
        r \sin \theta = \frac{3}{4}
        \]
        与 $C_2$ 联立可解得
        \[
        (1+\cos 2\theta)\sin \theta = \frac{3}{4} \Rightarrow (2\sin \theta-1)(4\sin^2 \theta+2\sin \theta-3)=0 \Rightarrow \sin \theta = \frac{\sqrt{13}-1}{4}
        \]
        此时
        \[
        r = \frac{1 + \sqrt{13}}{4}, \quad \cos \theta = \frac{1}{4}\sqrt{2 + 2\sqrt{13}}
        \]
        故$PQ$的长度为
        \begin{align*}
        |x_P - x_Q| 
        &= r_P \cos \theta_P - r_Q \cos \theta_Q \\
        &= \frac{3}{2} \cdot \frac{\sqrt{3}}{2} - \left( \frac{1 + \sqrt{13}}{4} \right) \left( \frac{\sqrt{2 + 2\sqrt{13}}}{4} \right) \\
        &= \frac{1}{32} \left[ 24\sqrt{3} - (2 + 2\sqrt{13})^{\frac{3}{2}} \right]
        \end{align*}
        证毕。
    \end{solution}

    \question 曲线的极坐标方程为 
    \[
    r = 1 + \tan \theta,\quad 0 \le \theta \le \frac{\pi}{2}
    \]
    点 $P$ 位于曲线上 $\theta = \dfrac{\pi}{3}$ 处。点 $Q$ 位于极轴上,使得直线 $L$ 经过 $P$ 和 $Q$ 且与极轴垂直。求由该曲线与直线 $L$ 所围成的有限区域的面积。
    \begin{solution}
        $P$ 点坐标为$\left(1 + \sqrt{3},\dfrac{\pi}{3}\right)$,
        因此 
        \[
        OQ = (1 + \sqrt{3}) \cos \frac{\pi}{3} = \frac{1}{2}(1 + \sqrt{3})
        \]
        直线 $L$ 的极坐标方程为 
        \[
        r \cos \theta = \frac{1}{2}(1 + \sqrt{3})
        \]
        与曲线联立,解得
        \[
        (1 + \tan\theta)\cos\theta = \frac{1}{2}(1 + \sqrt{3}) \Rightarrow \theta = \frac{\pi}{6}
        \]
        故得$R\left(1 + \dfrac{\sqrt{3}}{3},\dfrac{\pi}{6}\right)$,于是$\triangle OPR$ 的面积为
        \[
        \frac{1}{2}\cdot OP \cdot OR \sin \left(\frac{\pi}{3}-\frac{\pi}{6}\right)=\frac{1}{2}(1 + \sqrt{3})\left(1 + \frac{\sqrt{3}}{3}\right)\sin \frac{\pi}{6}=\frac{1}{2}+\frac{\sqrt{3}}{3}
        \]
        而极坐标扇形面积为
        \begin{align*}
        \frac{1}{2} \int_{\frac{\pi}{6}}^{\frac{\pi}{3}} (1 + \tan \theta)^2 \, d\theta 
        &= \int_{\frac{\pi}{6}}^{\frac{\pi}{3}} \left( \tan \theta + \frac{1}{2} \sec^2 \theta \right) d\theta \\
        &= \left[ \ln(\sec \theta) + \frac{1}{2} \tan \theta \right]_{\frac{\pi}{6}}^{\frac{\pi}{3}} \\
        &= \frac{1}{3} \sqrt{3} + \frac{1}{2} \ln 3
        \end{align*}
        所求面积为
        \[ 
        \frac{1}{3} \sqrt{3} + \frac{1}{2} \ln 3-\left(\frac{1}{2}+\frac{\sqrt{3}}{3}\right) = \frac{1}{2}(\ln3-1)
        \]
    \end{solution}

    \question 如图,曲线方程为 
    \[
    r = 2 + 2\sin\theta, \quad 0 \le \theta \le 2\pi
    \]
    直线方程为 
    \[
    2r\sin\theta = 3, \quad 0 < \theta < \pi
    \]
    \begin{figure}[H]
        \centering        
        \includegraphics[width=0.5\textwidth]{images/image243.png}
    \end{figure}
    \begin{parts}
    \part 求交点 $P,Q$ 的坐标。
    \begin{solution}
        联立得
        \[ 
        2(2 + 2\sin\theta)\sin\theta = 3 \Rightarrow (2\sin\theta - 1)(2\sin\theta + 3) = 0
        \]
        由于 $-1 \le \sin\theta \le 1$,方程 $2\sin\theta + 3 = 0$ 无解,因此
        \[ 
        \sin\theta = \frac{1}{2} \Rightarrow \theta = \frac{\pi}{6}, \frac{5\pi}{6}
        \]
        所以交点坐标为
        \[
        P\left(3, \frac{\pi}{6}\right),Q\left(3, \frac{5\pi}{6}\right)
        \]
    \end{solution}
    \part 证明$\triangle OPQ$ 的面积为 $\dfrac{9\sqrt{3}}{4}\sqrt{3}$。
    \begin{solution}
        已知 $OP = 3,OQ = 3$,夹角为 
        \[
        \frac{5\pi}{6} - \frac{\pi}{6} = \frac{2\pi}{3}
        \]
        故$\triangle OPQ$面积为
        \[ 
        \frac{1}{2} \cdot 3 \cdot 3 \cdot \sin\frac{2\pi}{3} = \frac{9\sqrt{3}}{4} 
        \]
    \end{solution}
    \part 求由曲线和直线围成的阴影区域的面积。
    \begin{solution}
        阴影区域面积等于曲线在对应角度内围成的面积减去$\triangle OPQ$ 的面积,
        \[
        \int_{\frac{\pi}{6}}^{\frac{5\pi}{6}} \frac{1}{2} (2 + 2\sin\theta)^2\, d\theta - [\triangle OPQ]
        \]
        其中
        \begin{align*}
        \int_{\frac{\pi}{6}}^{\frac{5\pi}{6}} \frac{1}{2} (2 + 2\sin\theta)^2\, d\theta
        &= \int_{\frac{\pi}{6}}^{\frac{5\pi}{6}} (3 + 4\sin\theta - \cos 2\theta)\, d\theta \\
        &= \Bigl[3\theta - 4\cos\theta - \frac{1}{2}\sin 2\theta\Bigr]_{\frac{\pi}{6}}^{\frac{5\pi}{6}} \\
        &= 2\pi + \frac{9\sqrt{3}}{2}
        \end{align*}
        故阴影部分的面积为
        \[ 
        (2\pi + \frac{9\sqrt{3}}{2}) - \frac{9\sqrt{3}}{4} = 2\pi + \frac{9}{4}\sqrt{3} 
        \]
    \end{solution}
    \end{parts}

    \question 如图,曲线的极坐标方程为
    \[
    r^2=\cos2\theta,\quad 
    \theta\in\left[0,\frac{\pi}{4}\right]\cup\left[\frac{3\pi}{4},\frac{5\pi}{4}\right]\cup\left[\frac{7\pi}{4},2\pi\right].
    \]
    该曲线被矩形 $ABCD$ 包围,其中 $AB,CD$ 为与极轴平行的切线,$AD,BC$ 为与极轴垂直的切线。证明曲线与矩形 $ABCD$ 之间所围成的总面积为$\sqrt{2}-1$。
    \begin{figure}[H]
        \centering        
        \includegraphics[width=0.5\textwidth]{images/image230.png}
    \end{figure}
    \begin{solution}
        首先确定矩形的边界,由
        \[
        r^2=\cos2\theta
        \]
        可知当 $\cos2\theta=1$ 时有$r=1$,因此矩形宽为 $2$,而切线与极轴平行当且仅当
        \[
        \frac{dy}{dx}=0 \Rightarrow \frac{dy}{d\theta}=0 
        \]
        由于 
        \[
        y = r\sin\theta = \sqrt{\cos \theta}\sin\theta
        \]
        为简化计算,不妨解
        \[ 
        \frac{d}{d\theta}(y^2)=\frac{d}{d\theta}(r^2 \sin^2 \theta) = \frac{d}{d\theta}(\cos 2\theta \sin^2 \theta) = 0 
        \]
        求导得
        \[ 
        -2\sin 2\theta \sin^2 \theta + \cos 2\theta (2\sin\theta \cos\theta) = 0 
        \]
        \[ 
        \sin 2\theta [-2\sin^2 \theta + \cos 2\theta] = 0 
        \]
        若 $\sin 2\theta = 0$,则 
        \[
        \theta = \frac{n\pi}{2}, n \in \mathbb{Z},
        \]
        不适合。若 $-2\sin^2 \theta + \cos 2\theta = 0$,则由
        \[ 
        -2\left(\frac{1}{2} - \frac{1}{2}\cos 2\theta\right) + \cos 2\theta = 0 \Rightarrow \cos 2\theta = \frac{1}{2}
        \]
        取$\theta = \dfrac{\pi}{6}$,此时 $r = \dfrac{\sqrt{2}}{2}$,矩形高为
        \[
        2\cdot \frac{\sqrt{2}}{2} \cdot \sin \frac{\pi}{6} = \frac{\sqrt{2}}{2} 
        \]
        故矩形的总面积为$\sqrt{2}$,而由曲线围成的面积为
        \[ 
        4 \int_{0}^{\frac{\pi}{4}} \frac{1}{2} r^2 d\theta 
        = 4 \int_{0}^{\frac{\pi}{4}} \frac{1}{2} \cos 2\theta d\theta 
        = 4 \left[\frac{1}{4}\sin 2\theta \right]_{0}^{\frac{\pi}{4}} 
        = 1 
        \]
        因此曲线与矩形之间围成的面积为$\sqrt{2} - 1$。
    \end{solution}

    \question 如图,两条封闭曲线的极坐标方程分别为
    \[
    C_1:\ r=a(1+\sin\theta),\quad 0\le\theta\le2\pi, \quad 
    C_2:\ r=3a(1-\sin\theta),\quad 0\le\theta\le2\pi,
    \]
    其中$a>0$,且在极点 $O$ 及点 $P,Q$ 处相交。
    \begin{figure}[H]
        \centering        
        \includegraphics[width=0.4\textwidth]{images/image233.png}
    \end{figure}
    \begin{parts}
    \part 求点 $P,Q$ 的极坐标。
    \begin{solution}
        交点 $P,Q$ 满足
        \[
        a(1+\sin\theta)=3a(1-\sin\theta) \Rightarrow \sin\theta=\frac{1}{2} \Rightarrow \theta=\frac{\pi}{6},\ \frac{5\pi}{6}.
        \]
        因此
        \[
        P\left(\frac{3a}{2},\frac{5\pi}{6}\right),\quad Q\left(\frac{3a}{2},\frac{\pi}{6}\right)
        \]
    \end{solution}
    \part 证明线段 $PQ$ 的长度为 $\dfrac{3\sqrt{3}}{2}a$。
    \begin{solution}
        由对称性,设 $M$ 为 $PQ$ 的中点,则
        \[
        PQ=2MQ=2\cdot\frac{3a}{2}\sin\frac{\pi}{3}
        =\frac{3\sqrt{3}}{2}a
        \]
    \end{solution}
    \part 已知 $PQ=\dfrac{3}{2}$,证明阴影部分的面积为$3\sqrt{3}-\dfrac{4}{3}\pi$。
    \begin{solution}
        由$PQ=\dfrac{3}{2}$知
        \[
        \frac{3\sqrt{3}}{2}a=\frac{3}{2} \Rightarrow a=\frac{\sqrt{3}}{3}
        \]
        阴影部分的面积为
        \begin{align*}
        &2 \left[ \int_{\frac{\pi}{6}}^{\frac{\pi}{2}} \frac{1}{2} [a(1+\sin\theta)]^2 d\theta - \int_{\frac{\pi}{6}}^{\frac{\pi}{2}} \frac{1}{2} [3a(1-\sin\theta)]^2 d\theta \right] \\
        &= a^2 \int_{\frac{\pi}{6}}^{\frac{\pi}{2}} (-8 + 20\sin\theta - 8\sin^2 \theta) d\theta \\
        &= a^2 \int_{\frac{\pi}{6}}^{\frac{\pi}{2}} (-12 + 20\sin\theta + 4\cos 2\theta) d\theta \\
        &= a^2 \Bigl[ -12\theta - 20\cos\theta + 2\sin 2\theta \Bigr]_{\frac{\pi}{6}}^{\frac{\pi}{2}} \\
        &= 3\sqrt{3} - \frac{4\pi}{3}
        \end{align*}
        \end{solution}
    \end{parts}

    \question 如图,两条封闭曲线 $C_1$ 与 $C_2$ 的极坐标方程分别为
    \[
    C_1:\ r=3+\cos\theta,\quad 0\le\theta<2\pi,\quad 
    C_2:\ r=5-3\cos\theta,\quad 0\le\theta<2\pi.
    \]
    两曲线相交于两点 $P,Q$。
    \begin{figure}[H]
        \centering        
        \includegraphics[width=0.5\textwidth]{images/image234.png}
    \end{figure}
    \begin{parts}
    \part 求点 $P,Q$ 的极坐标。
    \begin{solution}
        交点 $P,Q$ 满足
        \[
        3+\cos\theta=5-3\cos\theta \Rightarrow \cos\theta=\frac{1}{2}
        \]
        在 $0\le\theta<2\pi$ 内解得
        \[
        P\left(\frac{7}{2},\frac{\pi}{3}\right),\quad
        Q\left(\frac{7}{2},\frac{5\pi}{3}\right).
        \]
    \end{solution}
    \part 证明$C_1$ 与 $C_2$ 的相交区域面积为
    \[
    \frac{97\pi-102\sqrt{3}}{6}.
    \]
    \begin{solution}
        当 $0\le\theta\le\dfrac{\pi}{3}$ 时,区域由曲线 $C_2$ 给出；当 $\dfrac{\pi}{3}\le\theta\le\pi$ 时,区域由曲线 $C_1$ 给出。所求面积为
        \[
        2(A_1+A_2)
        \]
        其中
        \begin{align*}
        A_1 &=\int_{0}^{\frac{\pi}{3}}\frac{1}{2}(5-3\cos\theta)^2\,d\theta \\
        &= \int_{0}^{\frac{\pi}{3}}\left(\frac{25}{2}-15\cos\theta+\frac{9}{2}\cos^2\theta\right)d\theta \\
        &= \int_{0}^{\frac{\pi}{3}}\left(\frac{59}{4}-15\cos\theta+\frac{9}{4}\cos2\theta\right)d\theta \\
        &= \left[\frac{59}{4}\theta-15\sin\theta+\frac{9}{8}\sin2\theta\right]_{0}^{\frac{\pi}{3}} \\
        &=\frac{59\pi}{12}-\frac{111}{16}\sqrt{3},
        \end{align*}
        \begin{align*}
        A_2 &=\int_{\frac{\pi}{3}}^{\pi}\frac{1}{2}(3+\cos\theta)^2\,d\theta \\
        &=\frac{1}{2}\int_{\frac{\pi}{3}}^{\pi}\left(9+6\cos\theta+\cos^2\theta\right)d\theta\\
        &=\int_{\frac{\pi}{3}}^{\pi}\left(\frac{19}{4}+3\cos\theta+\frac{1}{4}\cos2\theta\right)d\theta\\
        &=\left[\frac{19}{4}\theta+3\sin\theta+\frac{1}{8}\sin2\theta\right]_{\frac{\pi}{3}}^{\pi}\\
        &=\frac{19\pi}{6}-\frac{25\sqrt{3}}{16}
        \end{align*}
        因此阴影区域面积为
        \[
        2(A_1+A_2)=\frac{97\pi-102\sqrt{3}}{6}
        \]
    \end{solution}
    \end{parts}

    \question 曲线 $C$ 的极坐标方程为 $r = \tan\frac{\theta}{2}$,其中 $0 \le \theta < \frac{\pi}{2}$。点 $P$ 位于曲线 $C$ 上,且该点处的切线垂直于极轴。射线 $\theta = \alpha$ 经过点 $P$。求曲线 $C$ 与该射线所围成区域的面积。
    \begin{figure}[H]
        \centering        
        \includegraphics[width=0.5\textwidth]{images/image242.png}
    \end{figure}
    \begin{solution}
        由于点 $P$ 处的切线垂直于极轴,意味 
        \[ 
        \frac{dx}{d\theta} = \frac{d}{d\theta} \left[ \tan\frac{\theta}{2}\cos \theta \right] = \frac{1}{2} \sec^2\frac{\theta}{2} \cos \theta - \tan\frac{\theta}{2} \sin \theta = 0 
        \]
        注意到 $\theta < \frac{\pi}{2}$,故 $\cos \theta \neq 0$,于是
        \[ 
        \sec^2\frac{\theta}{2} - 2 \tan\frac{\theta}{2} \tan \theta = 0 
        \]
        设$t=\tan\dfrac{\theta}{2}$,解得
        \[ 
        (1 + t^2) - 2t \left( \frac{2t}{1 - t^2} \right) = 0 \Rightarrow t^2 = -2 + \sqrt{5} \Rightarrow t = \sqrt{-2 + \sqrt{5}}
        \]
        因此点 $P$ 对应的极角 $\alpha$ 为
        \[ 
        \alpha = 2 \arctan \sqrt{-2 + \sqrt{5}} 
        \]
        所围图形的面积为
        \[
        \frac{1}{2} \int_{0}^{\alpha} \left[ \sec^2\frac{\theta}{2} - 1 \right] \, d\theta 
        = \frac{1}{2} \left[ 2 \tan\frac{\theta}{2} - \theta \right]_{0}^{\alpha} 
        = \tan\left(\frac{1}{2}\alpha\right) - \frac{1}{2}\alpha
        = \sqrt{-2 + \sqrt{5}} - \arctan \sqrt{-2 + \sqrt{5}}
        \]
    \end{solution}

    \question 曲线 $C_1$ 和 $C_2$ 的极坐标方程分别为 
    \[
    r_1 = \sec \theta(1 - \tan^2 \theta), \quad r_2 = \frac{1}{2}\sec^3 \theta,\quad 0 \le \theta < \frac{\pi}{4},
    \]
    \begin{figure}[H]
        \centering        
        \includegraphics[width=0.5\textwidth]{images/image241.png}
    \end{figure}
    \begin{parts}
    \part 证明 $\tan \angle OAP = -3\sqrt{3}$。
    \begin{solution}
        联立可知
        \[ 
        \frac{1}{2}\sec^3 \theta = \sec \theta(1 - \tan^2 \theta) \Rightarrow \tan \theta = \frac{1}{\sqrt{3}} 
        \]
        于是$A\left(\dfrac{4}{3\sqrt{3}},\dfrac{\pi}{6}\right)$,设 $\angle OAP = \psi$,在 $\triangle OAP$ 中,由正弦定理,
        \[ 
        \frac{\sin \psi}{1} = \frac{\sin(\frac{5\pi}{6} - \psi)}{\frac{4}{3\sqrt{3}}}
        \]
        化简得
        \[ 
        4\sqrt{3} \sin \psi = 9 \left(\frac{1}{2}\cos \psi+\frac{\sqrt{3}}{2}\sin\psi\right)
        \Rightarrow \tan \psi = -3\sqrt{3} 
        \]
        故得证。
    \end{solution}
    \part 求曲边三角形 $OAQ$ 的面积,其中 $O$ 为极点。
    \begin{solution}
        面积为 
        \begin{align*}
        &\frac{1}{2} \int_{0}^{\frac{\pi}{6}} \frac{1}{4}\sec^6\theta\, d\theta + \frac{1}{2} \int_{\frac{\pi}{6}}^{\frac{\pi}{4}} \sec^2\theta\ (1 - \tan\theta)^2\, d\theta \\
        &= \frac{1}{8} \int_{0}^{\frac{\pi}{6}} \sec^2\theta(1 + \tan\theta)^2\, d\theta + \frac{1}{2} \int_{\frac{\pi}{6}}^{\frac{\pi}{4}} \sec^2\theta\ (1 - \tan\theta)^2\, d\theta \\
        &= \frac{1}{8}\left[\frac{1}{5}\tan^5 \theta + \frac{2}{3}\tan^3 \theta + \tan \theta\right]_{0}^{\frac{\pi}{6}} + \frac{1}{2} \left[\frac{1}{5}\tan^5 \theta - \frac{2}{3}\tan^3 \theta + \tan \theta\right]_{\frac{\pi}{6}}^{\frac{\pi}{4}} \\
        &= \frac{36-11\sqrt{3}}{135}
        \end{align*}
    \end{solution}
    \end{parts}

    \question 一斜率为 $-\frac{3}{11}$ 的直线 $L$ 是曲线
    \[
    r = 25\cos 2\theta,\quad 0 \le \theta \le \frac{\pi}{2}
    \] 
    的切线。证明由该曲线、直线 $L$ 以及极轴所围成的区域的面积为
    \[ 
    \frac{25}{12} \left( 46 - 75 \tan^{-1} \frac{1}{3} \right) 
    \]
    \begin{solution}
        由$r = 25\cos 2\theta$及链导法得
        \[ 
        \frac{dy}{dx} = \frac{\frac{d}{d\theta}(r \sin \theta)}{\frac{d}{d\theta}(r \cos \theta)} = \frac{\cos 2\theta \cos \theta - 2\sin 2\theta \sin \theta}{-2\sin 2\theta \cos \theta - \cos 2\theta \sin \theta}=-\frac{3}{11}
        \]
        即
        \[
        \frac{2\tan\theta\tan2\theta-1}{2\tan2\theta+\tan\theta}=-\frac{3}{11}
        \]
        设$t=\tan\theta$,化简得
        \[ 
        (3t - 1)(t^2 - 18t - 11) = 0 \Rightarrow \tan \theta = \frac{1}{3},9 \pm \sqrt{92}
        \]
        为了丢弃增根,解
        \[
        \frac{dy}{dx}=0 \Rightarrow \cos 2\theta \cos \theta - 2\sin 2\theta \sin \theta = 0 \Rightarrow \theta = \sin^{-1}\frac{\sqrt{6}}{6} \approx 24.1^\circ
        \]
        由于
        \[
        \tan^{-1}(9+\sqrt{92})>24.1^\circ,\quad \tan^{-1}(9-\sqrt{92})<0
        \]
        解为
        \[
        \theta=\tan^{-1}\frac{1}{3}<24.1^\circ
        \]
        当 $\tan \theta = \dfrac{1}{3}$ 时,易知$P$直角坐标为
        \[
        \left(20\cdot\frac{3}{\sqrt{10}},20\cdot\frac{1}{\sqrt{10}}\right)=(6\sqrt{10}, 2\sqrt{10})
        \]
        故直线 $L$ 的方程为
        \[ 
        y - 2\sqrt{10} = -\frac{3}{11}(x - 6\sqrt{10}) 
        \]
        当 $y = 0$ ,得 $x = \dfrac{40\sqrt{10}}{3}$。设$Q$ 是直线与 $x$ 轴的交点， $\triangle OPQ$ 的面积为
        \[ 
        [\triangle OPQ] = \frac{1}{2} \cdot \frac{40\sqrt{10}}{3} \cdot 2\sqrt{10} = \frac{400}{3} 
        \]
        从 $\theta = 0$ 到 $\theta = \tan^{-1} \dfrac{1}{3}$扫过的面积为
        \begin{align*}
        I &= \frac{1}{2} \int_{0}^{\tan^{-1} \frac{1}{3}} (25 \cos 2\theta)^2\, d\theta \\
        &= \frac{625}{4} \int_{0}^{\tan^{-1} \frac{1}{3}} (1+\cos 4\theta)2\theta\, d\theta \\
        &= \frac{625}{4} \left[ \theta + \frac{1}{4}\sin 4\theta \right]_{0}^{\tan^{-1} \frac{1}{3}} \\
        &= \frac{625}{4} \left[ \theta + \frac{1}{2}\sin 2\theta \cos 2\theta \right]_{0}^{\tan^{-1} \frac{1}{3}} \\
        &= \frac{625}{4} \tan^{-1} \frac{1}{3} + \frac{75}{2} 
        \end{align*}
        故所求面积为
        \[ 
        \frac{400}{3} - \left( \frac{625}{4} \tan^{-1} \frac{1}{3} + \frac{75}{2} \right) = \frac{25}{12} \left( 46 - 75\tan^{-1} \frac{1}{3} \right) 
        \]
        证毕。
    \end{solution}
\end{questions}
\pagebreak

\begin{center}
  {\fontsize{30pt}{26pt}\selectfont
    \hypertarget{立体几何、空间向量}{立体几何、空间向量} \label{立体几何、空间向量}
  }
\end{center}
\separator
\vspace{1pt}

\begin{questions}
    \question 一个立方体的棱长为 $4$。一根长为 $5$ 的绳子的一端固定在立方体上表面的中心。求绳子另一端能接触到的立方体表面的面积。
    \ifprintanswers
    \begin{figure}[H]
        \centering
        \includegraphics[width=0.3\textwidth]{images/image176.png}
    \end{figure}
    \fi
    \begin{solution}
        设立方体顶面为正方形 $ABCD$,中心为 $O$,则$O$到顶点的距离为 $\sqrt{2^2 + 2^2} = \sqrt{8}$。绳长 $5 > \sqrt{8}$,所以绳端可到达整个顶面面积 $16$。

        绳子无法到达底面,因为沿表面的最短路径长度为 $6 > 5$。每个侧面能到达的面积相同,设为 $a$,总可达面积为 $A = 16 + 4a$。

        考虑某一侧面正方形 $ABEF$。以 $O$ 为圆心,半径 $5$,在侧面上形成弧 $\wideparen{PQ}$,中点为 $M$,则$PM=\sqrt{5^2-2^2}=\sqrt{21}$,故绳尾能接触到的面积为
        \begin{align*}
        &[\text{矩形}\ ABQP]+[\text{扇形}\ POQ]-[\triangle OPQ]\\
        &=4(\sqrt{21}-2)+ 25 \arcsin\frac{2}{5} - 2\sqrt{21}\\
        &=2\sqrt{21}-8+ 25 \arcsin\frac{2}{5}
        \end{align*}
        总面积为
        \[
        16 + 8\sqrt{21} - 32 + 100 \arcsin\frac{2}{5} = 8\sqrt{21} -16 + 100 \arcsin\frac{2}{5} \approx 61.81
        \]
    \end{solution}

    \question 边长为 $9$, $12$, $15$ 的直角三角形框架水平放置,一个半径为 $5$ 的球卡在其中并与三边均相切。求球体高出三角形平面的高度。
    \ifprintanswers
    \begin{figure}[H]
        \centering
        \includegraphics[width=0.4\linewidth]{images/image195.png}
    \end{figure}
    \fi
    \begin{solution}
        考虑球在三角形所在平面的截面。球体的截面为$\triangle ABC$的内切圆。设圆心为 $O$,半径为 $r$,连接 $O$ 到$BC,CA,AB$上的切点 $P,Q,R$。$\triangle ABC$面积为
        \[
        \frac{1}{2}\cdot 9\cdot 12
        = \frac{1}{2}r (9+12+15) \Rightarrow r=3
        \]
        设球心为 $O'$,则 $OO'$ 垂直于三角形平面,记 $OO'=h$。任取圆周上一点,与 $O'$ 构成直角三角形。由毕氏定理,
        \[
        h=\sqrt{5^2-3^2}=4
        \]
        因此球心在平面以上 $4$ 个单位,而球的半径为 $5$,所以球顶部高出平面的高度为
        \[
        h+5=9
        \]
    \end{solution}

    \question 空间中一个立方体的三个顶点坐标为 $A(3,4,1),B(5,2,9),C(1,6,5)$。求立方体的中心坐标。
    \begin{solution}
        观察
        \[
        AB=\sqrt{(5-3)^2+(2-4)^2+(9-1)^2}=\sqrt{72},
        \]
        \[
        AC=\sqrt{(1-3)^2+(6-4)^2+(5-1)^2}=\sqrt{24},
        \]
        \[
        BC=\sqrt{(1-5)^2+(6-2)^2+(5-9)^2}=\sqrt{48}
        \]
        此时发现
        \[
        AB^2=AC^2+BC^2
        \]
        故可知 $AB$ 是经过立方体中心的对角线,因此立方体中心即为 $AB$ 的中点
        \[
        O=\left(\frac{3+5}{2}, \frac{4+2}{2}, \frac{1+9}{2}\right)=(4,3,5)
        \]
    \end{solution}

    \question 求内接在单位立方体中的正方形的最大面积,其中正方形的每个顶点都在立方体的边上。
    \begin{solution}
        设立方体的顶点坐标为 $(0,0,0), (0,0,1), \dots, (1,1,1)$。由对称性,最大内接正方形的顶点取
        \[
        (x,0,0), \quad (1,0,1-x), \quad (1-x,1,1), \quad (0,1,x)
        \]
        其中 $0<x<1$。两条相邻边的平方长度相等,解得
        \[
        2x^2 + 1 = 2(1-x)^2 \implies x = \frac{1}{4}
        \]
        此时最大正方形的面积为
        \[
        ||(x,0,0)-(1,0,1-x)||^2 = \frac{9}{8}
        \]
    \end{solution}

    \question 已知一边长为 $4$、顶面为 $ABCD$、底面为 $EFGH$ 的立方体,且 $A$ 在 $E$ 正上方,依此类推。若点 $X,Y,Z$ 在 $AD,HG,BF$ 上使得$AX,HY,FZ=1$,求 $\triangle XYZ$ 的面积。
    \ifprintanswers
    \begin{figure}[H]
        \centering
        \includegraphics[width=0.4\linewidth]{images/image144.png}
    \end{figure}
    \fi
    \begin{solution}
        在$\triangle AXZ$中,由毕氏定理,
        \[
        AZ=\sqrt{3^2 + 4^2} = 5
        \]
        在$\triangle ABZ$中,由毕氏定理,
        \[
        XZ=\sqrt{1^2 + 5^2} = \sqrt{26}
        \]
        $\triangle XYZ$为等边三角形,因此面积为
        \[
        [\triangle XYZ] = \frac{\sqrt{3}}{4} \cdot 26 = \frac{13\sqrt{3}}{2}
        \]
    \end{solution}

    \question 从天花板点 $O$ 悬挂三根长 $100$ cm 的导线 $OXM,OYN,OZP$,依次穿过边长为 $60$ cm 的等边三角形 $XYZ$ 的顶点,末端挂有灯泡。三角形平面与天花板平行,三个灯泡距天花板均为 $90$ cm。求三角形到天花板的垂直距离。
    \begin{figure}[H]
        \centering
        \includegraphics[width=0.3\linewidth]{images/image212.png}
    \end{figure}
    \ifprintanswers
    \begin{figure}[H]
        \centering
        \includegraphics[width=0.35\linewidth]{images/image213.png}
    \end{figure}
    \fi
    \begin{solution}
        设$OX=OY=OZ=x$。由于导线全长为$100$, 则$XM=YN=ZP=100-x$。又灯泡距天花板为$90$, 所以木质三角形至天花板的距离为
        \[
        90-(100-x)=x-10
        \]
        设$C$为$\triangle XYZ$的中心, 由对称性$O$在$C$正上方, 且$OC=x-10$。在$\triangle OXC$中,由毕氏定理,
        \[
        x^2 = (x-10)^2+\left(\frac{30}{\cos 30^\circ}\right)^2 
        \]
        解得
        \[
        x=65
        \]
        因此木质三角形到天花板的高度为
        \[
        x-10=55\ \text{cm}.
        \]
    \end{solution}

    \question 立方体被两平面分成四块:一平面平行于面 $ABCD$ 且过棱 $BG$ 中点,另一平面过棱 $AB,AD,GH,HE$ 的中点。求最小块与最大块的体积比。
    \begin{figure}[H]
        \centering
        \includegraphics[width=0.35\linewidth]{images/image203.png}
    \end{figure}
    \ifprintanswers
    \begin{figure}[H]
        \centering
        \includegraphics[width=0.35\linewidth]{images/image204.png}
    \end{figure}
    \fi
    \begin{solution}
        不失一般性,设立方体边长为 2。平面 $PQRS$ 经过 $BG$ 的中点并平行于面 $ABCD$,其中$P,Q,R,S$为$AF,BG,CH,DE$的中点。另一个平面经过 $K, L, M, N$,即 $AB, GH,HE,AD$ 的中点。延长 $QK$ 和 $SN$ 相交于点 $T$,注意到 $\triangle TAN \backsim \triangle TPS \ \text{(AAA)}$,于是
        \[
        \frac{TA}{TP} = \frac{AN}{PS} \Rightarrow TA = 1
        \]
        于是
        \begin{align*}
        V_{\text{小}}=[AKN-PQS]&=[\text{四面体} \ T-PQS]-[\text{四面体} \ T-AKN]\\
        &=\frac{1}{3}\left(\frac{1}{2}\cdot 2 \cdot 2\right) - \frac{1}{3}\left(\frac{1}{2}\cdot 1 \cdot 1\right) = \frac{7}{6}
        \end{align*}
        且
        \[
        V_{\text{大}} = [ABCD-PQRS]- [AKN-PQS]= 2 \cdot 2 \cdot 1 - \frac{7}{6} = \frac{17}{6}
        \]
        故所求体积之比为
        \[
        \frac{V_{\text{小}}}{V_{\text{大}}} = \frac{7}{17}
        \]
    \end{solution}

    \question 边长为$3$的立方体置于底面直径为$8$、高为$24$的圆锥内,立方体的体对角线与圆锥轴线重合。求圆锥顶点到立方体最近顶点的距离。
    \begin{figure}[H]
        \centering
        \includegraphics[width=0.3\linewidth]{images/image205.png}
    \end{figure}
    \ifprintanswers
    \begin{figure}[H]
        \centering
        \includegraphics[width=0.35\linewidth]{images/image206.png}
        \includegraphics[width=0.35\linewidth]{images/image207.png}
        \includegraphics[width=0.18\linewidth]{images/image208.png}
    \end{figure}
    \fi
    \begin{solution}
        如图,设点$A$为立方体最底端的顶点,顶点$B,C,D$在圆锥内表面上且由对称性形成等边三角形,点$Q$为立方体最顶端的顶点。$BQ$ 是立方体的面对角线,$G$ 为$BCD$的重心,且 $QGA$ 在圆锥轴线上。

        所求距离为圆锥顶点 $T$ 到平面 $BCD$ 的距离,减去 $A$ 到平面 $BCD$ 的距离。

        等四面体$O-BCD$边长为
        \[
        BQ=BC=BD=CD=\sqrt{3^2+3^2}=3\sqrt{2},
        \]
        在 $\triangle BCD$中,
        \[
        BG=\frac{2}{3}\cdot 3\sqrt{2}\cos 30^\circ=\sqrt{6}
        \]
        设 $O$ 为圆锥底面圆心,$S$ 为底面上一点。由 $\triangle TGB \backsim \triangle TOS \ \text{(AAA)}$,得
        \[
        \frac{TG}{TO}=\frac{GB}{OS} \Rightarrow TG=24\cdot \frac{\sqrt{6}}{4}=6\sqrt{6}
        \]
        由毕氏定理,立方体空间对角线为
        \[
        AQ=\sqrt{3^2+3^2+3^2}=3\sqrt{3}
        \]
        在 $\triangle BGQ$ 中,由毕氏定理,
        \[
        QG=\sqrt{(3\sqrt{2})^2-(\sqrt{6})^2}=2\sqrt{3}
        \]
        因此
        \[
        AG=AQ-QG=\sqrt{3}
        \]
        故所求距离为
        \[
        TA=TG-AG=6\sqrt{6}-\sqrt{3}
        \]
    \end{solution}

    \question 已知 $ABCD$ 和 $PNCD$ 都是边长为 $2$ 的正方形,且这两个正方形互相垂直(它们共享边 $CD$)。在 $AB$ 一侧取点 $M$,使得平面 $\triangle PMN$ 平行于平面 $ABCD$,且 $\angle PMN=90^\circ,PM=MN$。求凸多面体 $PMN-ABCD$ 的体积。
    \begin{figure}[H]
        \centering
        \includegraphics[width=0.4\linewidth]{images/image180.png}
    \end{figure}
    \ifprintanswers
    \begin{figure}[H]
        \centering
        \includegraphics[width=0.4\linewidth]{images/image181.png}
        \includegraphics[width=0.4\linewidth]{images/image182.png}
    \end{figure}
    \fi
    \begin{solution}
        过 $M$ 作一条与 $PN$ 平行的线分别与平面 $PAD,NBC$ 相交于 $Q,R$,取 $T$ 为 $PN$ 的中点,由题意可知 $\triangle MTN$ 为直角等腰三角形,故 
        \[
        NR=MT=1
        \]
        考虑多面体 $PQRN-ABCD$以及两个全等棱锥 $PMQ-A,NMR-B$,于是
        \begin{align*}
        [PMN-ABCD]
        &=[PQRN-ABCD]-2\cdot [PMQ-A]\\
        &=\frac{1}{2}(1+2)\cdot 2\cdot 2 - 2 \cdot \frac{1}{3}\cdot\frac12\cdot 2
        =\frac{16}{3}
        \end{align*}
    \end{solution}
    \ifprintanswers
    \begin{figure}[H]
        \centering
        \includegraphics[width=0.4\linewidth]{images/image182.png}
        \includegraphics[width=0.4\linewidth]{images/image183.png}
    \end{figure}
    \fi
    \begin{solution}
        沿平面 $ABNP$ 分割凸多面体$PMN-ABCD$,考虑棱柱$PN-ABCD$及角锥$M-ABNP$,在$\triangle BCN$中,由毕氏定理,
        \[
        N=\sqrt{2^2+2^2}=2\sqrt2
        \]
        取 $S,T$ 为 $AB,PN$ 中点,由于 $\triangle MTN$ 为直角等腰三角形,故$US=BC-MT=1$, 在$\triangle MUS$中,由毕氏定理,
        \[
        MS=\sqrt{1^2+2^2}=\sqrt{5}
        \]
        在$\triangle MTS$中,由余弦定理,
        \[
        (\sqrt{5})^2 = 1^2+(2\sqrt{2})^2 -2\cdot 1 \cdot 2\sqrt{2} \cdot \cos \angle MTS
        \Rightarrow \angle MTS=45^\circ
        \]
        设顶点 $M$ 到底面的垂高为 $h$,则
        \[
        h=MT\sin45^\circ=\frac{\sqrt{2}}{2}
        \]
        故凸多面体$PMN-ABCD$体积为
        \begin{align*}
        [PMN-ABCD]
        &=[PN-ABCD]+[M-ABNP] \\
        &=\frac{1}{2}\cdot 2 \cdot 2 \cdot 2 + \frac{1}{3}\cdot 2\cdot 2\sqrt2\cdot\frac{\sqrt2}{2}
        =\frac{16}{3}
        \end{align*}
    \end{solution}

    \question 三个完全相同的圆锥,底面半径为50,高为120,它们的底面两两外切。在这三个圆锥围成的空隙中放置一个球,使得球的最高点与三个圆锥的顶点处于同一水平高度。求该球的半径。
    \begin{figure}[H]
        \centering
        \includegraphics[width=0.4\linewidth]{images/image209.png}
    \end{figure}
    \ifprintanswers
    \begin{figure}[H]
        \centering
        \includegraphics[width=0.4\linewidth]{images/image210.png}
        \includegraphics[width=0.3\linewidth]{images/image211.png}
    \end{figure}
    \fi
    \begin{solution}
        设三圆锥圆心为$A,B,C$,则底面圆心形成边长为$100$的正三角形$\triangle ABC$,且球心位于该正三角形的重心$G$正上方。设$D$为$BC$中点。在$\triangle BGD$中,
        \[
        BG = \frac{BD}{\cos 30^\circ} = \frac{100\sqrt{3}}{3}
        \]
        即重心到每个顶点的距离为$\dfrac{100\sqrt{3}}{3}$,也即为圆锥轴与通过球心的竖直线的水平距离。

        现考虑球体与圆锥的竖直截面。设$O$为球心, $r$为球半径,$P$为球体与圆锥斜面上的切点,$X$为圆锥顶点,$Y$为$XP$延长线与$AG$的交点,$H$为球体的最高点,则
        \[
        OH=OP=r,XH=AG=\frac{100\sqrt{3}}{3},XA=HG=120
        \]  
        在$\triangle AXY$中,由毕氏定理,
        \[
        XY = \sqrt{50^2+120^2} =130 \Rightarrow PY = XY - XP = 130 - \frac{100\sqrt{3}}{3}
        \]
        在$\triangle AXY$中,由毕氏定理,有$OY^2=OP^2+PY^2=OG^2+GY^2$,即
        \[
        r^2 + \left(130-\frac{100\sqrt{3}}{3}\right)^2 = (120-r)^2+\left(\frac{100\sqrt{3}}{3}-50\right)^2
        \]
        解得
        \[
        r=\frac{200\sqrt{3}}{9}
        \]
    \end{solution}
    \begin{solution}
        已知$XH=AG=\dfrac{100\sqrt{3}}{3},AX=120,AY=50$,由于$XY \perp OP$,
        \[
        m_{XY}=-\frac{12}{5} \Rightarrow m_{OP}=\frac{5}{12}
        \]
        设$OS=5t,SP=12t$,则
        \[
        OP=13t,\quad OH=13t,\quad XR=18t
        \] 
        又
        \[
        RS=RP+PS=\frac{15t}{2}+12t=\frac{39t}{2}=AG=\frac{100\sqrt{3}}{3}
        \]
        因此$t=\dfrac{200}{39\sqrt{3}}$, 球半径为
        \[
        r=OP=13t=\frac{200\sqrt{3}}{9}
        \]
    \end{solution}

    \question 一个正四棱锥的底面是边长为20的正方形。现有一个底面直径为10、高为10的圆柱横放(侧卧),要完全放置在锥体内部。圆柱的中心轴线平行于底面对角线,且位于对角线正上方。圆柱中心轴的中点在底面中心的正上方。求金字塔的最小高度。
    \begin{figure}[H]
        \centering
        \includegraphics[width=0.5\linewidth]{images/image220.png}
    \end{figure}
    \ifprintanswers
    \begin{figure}[H]
        \centering
        \includegraphics[width=0.3\linewidth]{images/image221.png}
        \includegraphics[width=0.4\linewidth]{images/image222.png}
        \includegraphics[width=0.25\linewidth]{images/image223.png}
    \end{figure}
    \fi
    \begin{solution}
        当正四棱锥四个侧面与圆柱两端面的圆周均相切,正四棱锥高度为最小。设正方形底面为 $ABCD$,顶点为 $T$,对角线交点为 $M$,则 $AM=BM=CM=DM=10\sqrt{2}$,顶点 $T$ 位于 $M$ 的正上方,正四棱锥高度为 $t=TM$。

        圆柱的中心轴线位于对角线 $AC$ 上,其中点在 $M$ 处。圆柱沿 $AC$ 方向在 $M$ 两侧各延伸 $5$ 个单位,两端点分别为 $E$ 和 $F$。从 $A$ 到最近的圆柱端点 $E$ 的距离为 $AE = 10\sqrt{2}-5$。

        在包含 $AC$ 和 $T$ 的垂直截面中,设 $L$ 在 $AT$ 上,$G$, $H$ 分别在 $AB$, $AD$ 上。由于 $\angle BAM=45^\circ$ 且圆柱端面垂直于对角线,$\triangle GEA$ 和 $\triangle HEA$ 均为等腰直角三角形,因此 $GE=HE=AE=10\sqrt{2}-5$。设圆柱横截面的圆心为 $O$,半径为 $5$。

        设 $LE=h$,则 $LO=h-5,LJ=x$,由$\triangle LJO \backsim \triangle LEG \ \text{(AAA)}$,
        \[
        \frac{LJ}{JO}=\frac{LE}{EG} \Rightarrow \frac{x}{h} = \frac{5}{10\sqrt{2}-5} \Rightarrow x = \frac{h}{2\sqrt{2}-1}
        \]
        且
        \[
        \frac{LG}{GE}=\frac{LO}{OJ} \Rightarrow \frac{x + (10\sqrt{2}-5)}{h-5} = \frac{10\sqrt{2}-5}{5}
        \]
        代入 $x = \dfrac{h}{2\sqrt{2}-1}$ 得
        \[
        h = \frac{10(2\sqrt{2}-1)^2}{8-4\sqrt{2}}
        \]
        又由 $\triangle AEL \backsim \triangle AMT \ \text{(AAA)}$,有
        \[
        \frac{t}{AM} = \frac{LE}{AE} \Rightarrow t = \frac{10\sqrt{2}\cdot h}{10\sqrt{2}-5}=15+5\sqrt{2}
        \]
    \end{solution}

    \question 在空间中,有三个半径均为 $10$ 的圆,它们两两外切,且都与同一平面相切。每个圆所在的平面与该平面成 $45^\circ$ 的倾角。这三个圆的三个切点位于一个与该平面平行的圆上。求此圆的半径。
    \begin{figure}[H]
        \centering
        \includegraphics[width=0.5\linewidth]{images/image215.png}
    \end{figure}
    \ifprintanswers
    \begin{figure}[H]
        \centering
        \includegraphics[width=0.8\linewidth]{images/image216.png}
        \includegraphics[width=0.3\linewidth]{images/image217.png}
        \includegraphics[width=0.35\linewidth]{images/image218.png}
        \includegraphics[width=0.3\linewidth]{images/image219.png}
    \end{figure}
    \fi
    \begin{solution}
        设三个圆与平面 $\pi$ 的切点分别为 $A,B,C$。每个圆所在平面与 $\pi$ 的交线两两相交,交点分别为 $D,E,F$,其中 $A,B,C$ 在 $DE,EF,FD$ 上。由对称性,$\triangle DEF$ 是等边三角形,记 $O$ 为 $\triangle DEF$ 的重心。

        设三个圆两两相切的切点为 $G,H,J$,且 $\triangle GHJ$ 是等边三角形,所求圆即为过 $G,H,J$ 三点的圆。延长$DG,EH,FJ$交于$T$,于是$T$在$O$正上方, $T,D,E,F$ 构成一四面体。

        侧面$\triangle DET$与$\pi$成 $45^\circ$ 角,且 $TO\perp \pi$,故 $\triangle TAO$ 是等腰直角三角形,因此
        \[
        TA=\sqrt{2}\;AO
        \]
        考虑$\triangle DEF$,设 $DA=AE=x$。在 $\triangle ADO$中,
        \[
        OA=\frac{x}{\sqrt{3}}, \quad OD=\frac{2x}{\sqrt{3}}.
        \]
        在 $\triangle TAD$中,由毕氏定理,
        \[
        TD=\sqrt{x^2+\left(\sqrt{2}\cdot\frac{x}{\sqrt{3}}\right)^2} = \sqrt{\frac{5}{3}}x
        \]
        考虑侧面$\triangle DET$,设 $R$ 为 $\triangle DET$的内切圆圆心, $S$ 为 $TA$ 与 $GH$ 的交点,则 $RG=10$。由于$\triangle TSG \backsim \triangle TGR \backsim \triangle TAD \ \text{(AAA)}$,
        \[
        \frac{TG}{GR} = \frac{TA}{AD} = \frac{\frac{\sqrt{2}}{\sqrt{3}}x}{x} \Rightarrow TG = \frac{\sqrt{2}}{\sqrt{3}}GR = 10\frac{\sqrt{2}}{\sqrt{3}}
        \]
        且
        \[
        \frac{SG}{TG} = \frac{AD}{TD} = \frac{x}{\frac{\sqrt{5}}{\sqrt{3}}x} \Rightarrow SG = \frac{\sqrt{3}}{\sqrt{5}}TG=2\sqrt{10}
        \]
        考虑等边三角形$\triangle GHJ$,设 $L$ 为 $\triangle GHJ$ 的中心,则 $LG$ 即为所求圆的半径:
        \[
        LG=\frac{2}{\sqrt{3}}\,SG=\frac{4\sqrt{30}}{3}
        \]
    \end{solution}

    \question 一个球的内接圆锥的最大体积与这个球的体积之比为?
    \begin{solution}
        设球的半径为 $R$,其内接圆锥的高为 $h$,底面半径为 $r$。由几何关系可得:
        \[
        r^2 = R^2 - |h - R|^2 = (2R - h)h.
        \]
        故圆锥的体积为:
        \[
        V_{\text{锥}} = \frac{1}{3} \pi r^2 h = \frac{1}{3} \pi h^2 (2R - h).
        \]
        设函数 \( V(h) = \dfrac{1}{3} \pi h^2 (2R - h) \),由AM-GM不等式得
        \[
        V(h) \le \frac{4\pi}{3} \left( \frac{\frac{h}{2} + \frac{h}{2} + 2R - h}{3} \right)^3 
        = \frac{32}{81} \pi R^3\]
        当 $h = \dfrac{4}{3}R$ 时取得最大值$\dfrac{32}{81} \pi R^3$,故体积之比为
        \[
        \frac{V_{\text{锥}}}{V_{\text{球}}} = \frac{32}{81} \pi R^3 \div \frac{4}{3} \pi R^3 = \frac{8}{27}.
        \]
    \end{solution}

    \question 一个圆柱形容器内装有水。现有一个实心圆锥,其高度等于圆柱的高度,底面半径是圆柱底面半径的一半。将此圆锥完全浸入水中,底面紧贴圆柱底面,顶点朝上。此时观察到水面高度恰好是圆柱高度的 $\dfrac{1}{2}$。若将圆锥从水中取出,求水面高度占圆柱高度的比值。
    \begin{figure}[H]
        \centering
        \includegraphics[width=0.35\linewidth]{images/image177.png}
    \end{figure}
    \begin{solution}
        设圆柱半径为 $r$,高度为 $h$,则圆锥体积为
        \[
        \frac{1}{3} \pi \left(\frac{1}{2} r\right)^2 h = \frac{1}{12} \pi r^2 h.
        \]
        当圆锥在圆柱内且水深为 $\dfrac{h}{2}$ 时,水体积为半圆柱体积减去下半圆锥体积
        \[
        \frac{1}{2} \pi r^2 h - \left(\frac{1}{12}\pi r^2 h - \left(\dfrac{1}{2}\right)^3 \cdot \frac{1}{12} \pi r^2 h \right)= \frac{41}{96} \pi r^2 h.
        \]
        圆锥移开后,水体积保持不变,因此水深是圆柱高度的 $\dfrac{41}{96}$。
    \end{solution}

    \question 一个三棱柱形容器竖直放置,底面为三角形。容器内放入三个半径为 $1$ 的球,每个球都与三角形底面相切,与两个相邻的长方形侧面相切,且三个球两两外切。在这三个球的上方放入第四个半径为 $1$ 的球,该球与下方三个球均外切,并与棱柱的顶面(上底面)相切。求该三棱柱的体积。
    \ifprintanswers
    \begin{figure}[H]
        \centering
        \includegraphics[width=0.4\linewidth]{images/image163.png}
        \includegraphics[width=0.4\linewidth]{images/image164.png}
    \end{figure}
    \fi
    \begin{solution}
        先计算底面积。取底面上方 1 个单位的截面,该截面通过每个球的球心以及球与侧面的切点。设$A,B,C$为截面三角形顶点,$X,Y,Z$为球心,$M,N,P,Q,R,S$为球与侧面的切点且在$AB,BC,CA$上。由对称性可知 $\triangle ABC$ 是等边三角形,边长为
        \[
        BC = BP + PQ + QC = 2 + 2\sqrt{3}
        \]  
        底面积为
        \[
        [\triangle ABC] = \frac{\sqrt{3}}{4}(2+2\sqrt{3})^2
        \]
        现求高。设$W$为顶球球心,则四球相切,$WXYZ$ 为正四面体,边长皆为 2。设 $V$ 为 $\triangle XYZ$ 的重心,$G$ 为 $YZ$ 中点,
        在$\triangle YVG$中,
        \[
        YV = \frac{2}{\sqrt{3}},\quad YG = \frac{2\sqrt{3}}{3}
        \]  
        在$\triangle WYV$中,
        \[
        WV = \sqrt{2^2 - \frac{4}{3}} = \sqrt{\frac{8}{3}}
        \]  
        底面到截面 $XYZ$ 平面距离为 1,顶球球心到顶面的垂直距离为 1,三棱柱高度为$ 2 + \sqrt{\dfrac{8}{3}}$,故三棱柱体积为
        \[
        \frac{\sqrt{3}}{4}(2+2\sqrt{3})^2 \cdot \left(2 + \sqrt{\frac{8}{3}}\right) \approx 46.97
        \]
    \end{solution}

    \question 如图,已知一边长为 $200$ 的立方体 $FGHJKLMN$,点 $P$ 在 $HG$ 上。已知从 $G$ 到 $\triangle PFM$ 上一点的最短距离为 $100$,求 $HP$ 的长度。
    \begin{figure}[H]
        \centering
        \includegraphics[width=0.4\linewidth]{images/image168.png}
    \end{figure}
    \begin{solution}
        设 $GP = x$,则 $HP = 200 - x$,考虑四面体 $FGMP$体积:
        \[
        \frac{1}{3}[\triangle FGM]\cdot x=\frac{1}{3}[\triangle PFM]\cdot 100 \tag{1}
        \]
        其中
        \[
        [\triangle FGM] = \frac{1}{2}\cdot FG \cdot GM = 20000
        \]
        现以$x$表示$\triangle PFM$ 的面积。在$\triangle FGM,\triangle FGP$中,由毕氏定理,
        \[
        FM = 200\sqrt{2}, \quad FP = \sqrt{x^2 + 200^2}
        \]
        由于$FP = MP$, $\triangle PFM$ 为等腰三角形。设 $T$ 为 $FM$ 中点,则
        \[
        FT = TM = 100\sqrt{2}
        \]  
        在$\triangle FTP$中,由毕氏定理,
        \[
        PT = \sqrt{FP^2 - FT^2} = \sqrt{x^2 + 40000 - (100\sqrt{2})^2} = \sqrt{x^2 + 20000}
        \]
        所以
        \[
        [\triangle PFM] = \frac{1}{2} \cdot FM \cdot PT = \frac{1}{2} \cdot 200\sqrt{2} \cdot \sqrt{x^2 + 20000}
        \]
        由$(1)$得,
        \[
        \frac{1}{3} \cdot 20~000 \cdot x = \frac{1}{3} \cdot \left(\frac{1}{2} \cdot 200\sqrt{2} \cdot \sqrt{x^{2}+20~000}\right) \cdot 100
        \]
        解得
        \[
        x = 100\sqrt{2} >0 \Rightarrow HP = 200 - 100\sqrt{2}
        \]
    \end{solution}

    \question 长方体 $PQRSTUVW$ 的顶点按图中方式标记。已知 $PR=1867,PV=2019,PT=x$,问:有多少个正整数 $x$ 使得满足这些条件的长方体存在?
    \begin{figure}[H]
        \centering
        \includegraphics[width=0.5\linewidth]{images/image167.png}
    \end{figure}
    \begin{solution}
        设 $PQ=a,PS=b,PU=c$,在$\triangle PQR,\triangle PQV,\triangle PST$中,由毕氏定理,
        \[
        a^2 + b^2 = 1867^2, \quad a^2 + c^2 = 2019^2, \quad b^2 + c^2 = x^2
        \]
        三式解得
        \[
        a^2 = \frac{1867^2 + 2019^2 - x^2}{2}, \quad
        b^2 = \frac{1867^2 - 2019^2 + x^2}{2}, \quad
        c^2 = \frac{-1867^2 + 2019^2 + x^2}{2}
        \]
        由 $a,b,c>0$,
        \[
        \begin{cases}
            x^2 < 1867^2 + 2019^2 \\
            x^2 > 2019^2 - 1867^2 \\
            x^2 > 1867^2 - 2019^2 
        \end{cases}
        \]
        其中最后一个不等式恒成立,即  
        \[
        768.55 \approx \sqrt{2019^2 - 1867^2} < x < \sqrt{2019^2 + 1867^2} \approx 2749.92
        \]
        因此满足条件的整数 $x$ 个数为
        \[
        2749 - 769 + 1 = 1981
        \]
    \end{solution}

    \question 如图所示,有一个直圆锥台,上底面直径为 $2$ cm,下底面直径为 $8$ cm,高为 $6\sqrt{2}$ cm。$AB$ 与 $CD$ 为直圆锥台的两侧边,$E$ 为 $AB$ 上一点,且 $AE=3$ cm。求从点 $B$ 出发,沿圆锥台侧面经过母线 $CD$,最后到达点 $E$ 的最短路径长度。
    \begin{figure}[H]
        \centering
        \includegraphics[width=0.4\linewidth]{images/image111.png}
    \end{figure}
    \ifprintanswers
    \begin{figure}[H]
        \centering
        \includegraphics[width=0.4\linewidth]{images/image112.jpg}
    \end{figure}
    \begin{figure}[H]
        \centering
        \includegraphics[width=0.7\linewidth]{images/image113.jpg}
    \end{figure}
    \fi
    \begin{solution}
        设$O$ 为 $BC$ 的中点,$P$ 为 $BA$ 及$DC$的交点,$Q$ 为 $AD$ 的中点,由$\triangle PQD \backsim \triangle POC \ \text{(AAA)}$,
        \[
        \frac{PQ}{PQ+6\sqrt 2} = \frac{1}{4} \Rightarrow PQ=2\sqrt 2
        \]
        在 $\triangle POC$ 中,由毕氏定理,
        \[
        PC = \sqrt{(6\sqrt 2 + 2\sqrt 2)^2 + 4^2} =12
        \]
        又
        \[
        \frac{PD}{PC} = \frac{1}{4} \Rightarrow PD=3
        \]
        沿 $PB$ 将圆锥剪开,则 $AB=CD=9-3=6$。由弧长公式
        \[
        \wideparen{AA'} = 3 \angle APA' \Rightarrow \angle APA' = \frac{2\pi}{3}
        \]
        在 $\triangle PBE'$ 中,由余弦定理,
        \[
        BE^2 = 12^2 + 6^2 - 2 \cdot 12 \cdot 6 \cdot \cos \frac{2\pi}{3} \Rightarrow BE = 6\sqrt{7}
        \]
        因此最短曲线长为 $6\sqrt{7}$。
    \end{solution}

    \question 已知点 \(O(0,0,0)\), \(A(1,0,0)\), \(B\left(\dfrac{1}{2}, \dfrac{\sqrt{3}}{2}, 0\right)\),  
    \begin{parts}
    \part 求在 \(xy\) 平面上方的点 \(C\),使得四面体 \(OABC\) 是正四面体。  
    \begin{solution}
        向量 \(\overrightarrow{OA}, \overrightarrow{OB}\) 长度均为 1,且夹角为 \(60^\circ\),可知 \(\triangle OAB\) 是边长为 1 的正三角形,重心为
        \[
        G = \left( \frac{0+1+\frac{1}{2}}{3}, \frac{0+0+\frac{\sqrt{3}}{2}}{3}, 0 \right) = \left( \frac{1}{2}, \frac{\sqrt{3}}{6}, 0 \right)
        \]
        为使四面体为正四面体,点 \(C\) 应在重心垂直向上、距离为正三角形高的方向上:
        \[
        OC = \sqrt{1^2 - \left( \frac{1}{\sqrt{3}} \right)^2 }=\frac{\sqrt{6}}{3}
        \]
        故
        \[
        C = \left( \frac{1}{2}, \frac{\sqrt{3}}{6}, \frac{\sqrt{6}}{3} \right)
        \]
    \end{solution}
    \part 计算其体积。
    \begin{solution}
        体积为
        \[
        V = \frac{1}{3} [\triangle ABC] \cdot h 
        = \frac{1}{3} \cdot \frac{\sqrt{3}}{4} \cdot 1^2 \cdot \frac{\sqrt{6}}{3} 
        = \frac{\sqrt{2}}{12}        
        \]
    \end{solution}
    \end{parts}
    
    \question 空间中,设 $\triangle ABC$ 的三边长 $AB=3,AC=5,BC=7$,另有一点 $P$ 满足 $PA = PB = PC = \dfrac{25\sqrt{3}}{3}$,则锥体 $P-ABC$ 的体积为多少？
    \begin{solution}
        由题意可知,$P$ 在 $\triangle ABC$ 的投影点为外心。设$a = BC = 7, b = AC = 5, c = AB = 3$,半周长$s = \dfrac{15}{2},\triangle ABC$面积为
        \[
        [\triangle ABC] = \sqrt{\frac{15}{2} \cdot \frac{1}{2} \cdot \frac{5}{2} \cdot \frac{9}{2}} = \frac{15\sqrt 3}{4}
        \]  
        所以外接圆半径
        \[
        R = \frac{abc}{4 [\triangle ABC]} =\frac{7\sqrt 3}{3}
        \]  
        由毕氏定理,锥体$P-ABC$高为
        \[
        h = \sqrt{\left(\frac{25\sqrt 3}{3}\right)^2 - \left(\frac{7\sqrt 3}{3}\right)^2} = 8\sqrt 3
        \]  
        体积为
        \[
        [P-ABC] = \frac{1}{3} \cdot [\triangle ABC] \cdot h =30
        \]  
    \end{solution}

    \question 已知一四面体 $P-ABC$ 中,$\ \angle APB = \angle BPC = \angle CPA = 60^{\circ}$,且 $\triangle APB,\triangle BPC,\triangle CPA$ 的面积分别为 $\dfrac{\sqrt{3}}{2},2,1$,求四面体 $P-ABC$ 的体积。
    \begin{solution}
        设$PA=a,PB=b,PC=c$,由已知面积和角度,
        \[
        \begin{cases}
            \dfrac{1}{2} ab \sin 60^\circ = \dfrac{\sqrt{3}}{2}\\[6pt]
            \dfrac{1}{2} bc \sin 60^\circ = 2 \\[6pt]
            \dfrac{1}{2} ac \sin 60^\circ = 1 \\[6pt]
        \end{cases}
        \Rightarrow ab = 2 , bc = \frac{8}{\sqrt{3}} , ac = \frac{4}{\sqrt{3}}
        \]
        解得
        \[
        a = 1, \quad b = 2, \quad c = \frac{4}{\sqrt{3}}
        \]
        在 $PB,PC$ 上取点 $B',C'$使得 $PB' = PC' = PA = 1$,此时 $PAB'C'$ 为正四面体,其高为
        \[
        h = \sqrt{\frac{2}{3}} \cdot 1 = \sqrt{\frac{2}{3}}
        \]
        由于点 $A$ 到平面 $PB'C'$ 的距离等于点 $A$ 到平面 $PBC$ 的距离,四面体 $P-ABC$ 的体积为
        \[
        [P-ABC] = \frac{1}{3} \cdot [\triangle BPC] \cdot h = \frac{2\sqrt{6}}{9}
        \]
    \end{solution}

    \question 已知正四面体 $O-ABC$ 的边长为 $6\sqrt{2},D,E,F,G$ 四点分别在 $OA, OB, OC, BC$ 上。若 
    \[
    OD=\frac{1}{6}OA,\quad OE=\frac{1}{3}OB,\quad OF=\frac{1}{2}OC
    \] 
    且 $G$ 为 $BC$ 中点,求四面体 $DEFG$ 的体积。
    \begin{solution}
        四面体边长为 $6\sqrt{2}$,构造坐标系如下:
        \[
        \begin{cases}
        O = (0,0,0) \\
        A = (6,0,6) \\
        B = (0,6,6) \\
        C = (6,6,0)
        \end{cases}
        \Rightarrow
        \begin{cases}
        {OD} = \frac{1}{6} {OA} \Rightarrow D = \frac{A + 5O}{6} = (1,0,1) \\
        {OE} = \frac{1}{3} {OB} \Rightarrow E = \frac{B + 2O}{3} = (0,2,2) \\
        {OF} = \frac{1}{2} {OC} \Rightarrow F = \frac{O + C}{2} = (3,3,0) \\
        G = \text{BC 中点} = \frac{B + C}{2} = (3,6,3)
        \end{cases}
        \]
        于是
        \[
        \overrightarrow{GD} = (-2,-6,-2),\overrightarrow{GE} = (-3,-4,-1),\overrightarrow{GF} = (0,-3,-3)
        \]
        体积为
        \[
        [DEFG] = \frac{1}{6} \left| \overrightarrow{GD} \cdot ( \overrightarrow{GE} \times \overrightarrow{GF} ) \right|
        = \frac{1}{6} \cdot 18 = 3
        \]
    \end{solution}

    \question 四面体 $D-ABC$ 的体积为$\dfrac{1}{6}$,且 $\angle ACB=60^{\circ}$,  
    \[
    AD + \sqrt{3}\: BC + \frac{AC}{2} = 3,
    \]
    求 $CD$ 的长度。
    \begin{solution}
        由AM-GM不等式,
        \[
        1=\dfrac{AD + \sqrt{3} \cdot BC + \dfrac{AC}{2}}{3} \ge \sqrt[3]{\frac{\sqrt{3}}{2} \cdot AD \cdot BC \cdot AC}
        \]
        即  
        \[
        AD \cdot BC \cdot AC \le \frac{2}{\sqrt{3}} \tag{1}
        \]
        四面体体积为  
        \[
        \frac{1}{6}=[D-ABC] \le \frac{1}{3} \cdot [\triangle ABC] \cdot AD = \frac{1}{6} AC \cdot BC \sin 60^\circ \cdot AD
        \]
        即  
        \[
        AD \cdot BC \cdot AC \ge \frac{2}{\sqrt{3}} \tag{2}
        \]
        由$(1),(2)$得  
        \[
        AD \cdot BC \cdot AC = \frac{2}{\sqrt{3}}
        \]
        即不等式$(1)$等号成立,当且仅当
        \[
        AD = \sqrt{3} \cdot BC = \frac{AC}{2} = 1
        \]
        且由(2)知$AD$ 与 $\triangle ABC$ 垂直,在 $\triangle ACD$中,由毕氏定理,  
        \[
        CD = \sqrt{1^2 + 2^2} = \sqrt{5}
        \]
    \end{solution}
    
    \question 已知一边长为2的正方形纸 $ABCD,M,N$ 分别为 $BC$ 与 $CD$ 的中点,今沿着 $AM,AN$ 与 $MN$ 折起,使 $B,C,D$ 三点重合。若平面 $ABM$ 与平面 $AMN$ 的夹角为 $\theta$,求 $\sin \theta$。
    \begin{figure}[H]
        \centering
        \includegraphics[width=0.5\linewidth]{images/image58.png}
    \end{figure}
    \begin{solution}
        构造坐标系,设$A = (0,0,0), M=(1,2,0), N=(2,1,0)$,折起后$B = (a,a,b)$,由已知得
        \[
        AB^2=2^2 = 2a^2 + b^2 = 4,\quad BM^2 = (a-1)^2 + (a-2)^2 + b^2 = 1
        \]
        解得
        \[
        a = \frac{4}{3}, \quad b = \frac{2}{3}
        \]
        由
        \[
        \overrightarrow{AM} = (1,2,0), \quad \overrightarrow{AB} = \left(\frac{4}{3}, \frac{4}{3}, \frac{2}{3}\right),
        \]
        平面 $AMN,ABM$ 的法向量为
        \[
        \vec{n}_1 = (0,0,1),\quad \vec{n}_2 = \overrightarrow{AB} \times \overrightarrow{AM} = \frac{2}{3}(-2,1,2),
        \]
        两平面夹角 $\theta$ 满足
        \[
        \cos \theta = \frac{\vec{n}_1 \cdot \vec{n}_2}{|\vec{n}_1||\vec{n}_2|} = \frac{2}{3} \Rightarrow \sin \theta = \sqrt{1 - \left(\frac{2}{3}\right)^2} = \frac{\sqrt{5}}{3}
        \]
    \end{solution}

    \question 有一边长为 2 的正四面体 $ABCD$,设 $A'$ 为 $A$ 关于平面 $BCD$ 的对称点,$\ B'$ 为 $B$ 关于平面 $ACD$ 的对称点,求四面体 $A'CB'D$ 的体积。
    \begin{solution}
        设 $A(0,0,0),B(\sqrt{2}, 0, \sqrt{2}),C(0, \sqrt{2}, \sqrt{2}),D(\sqrt{2}, \sqrt{2}, 0)$,则
        \[
        \overrightarrow{AC}=(0, \sqrt{2}, \sqrt{2}),\quad \overrightarrow{BC}=(−\sqrt{2},\sqrt{2},0),\quad \overrightarrow{CD}=(\sqrt{2},0,−\sqrt{2})
        \]
        \[
        \textbf{n}_{BCD}=\overrightarrow{BC}\times\overrightarrow{CD}=(−2,−2,−2),\quad \textbf{n}_{ACD}=\overrightarrow{AC}\times\overrightarrow{CD}=(−2,2,−2)
        \]
        故$\triangle BCD,\triangle ACD$方程式为
        \[
        E_1 :  x + y + z = 2\sqrt{2},\quad E_2 : -x + y - z = 0
        \]
        且对称点$A',B'$坐标为
        \[
        A' = \left( \dfrac{4\sqrt{2}}{3} , \dfrac{4\sqrt{2}}{3} , \dfrac{4\sqrt{2}}{3}  \right),\quad
        B' = \left( -\dfrac{\sqrt{2}}{3}, \dfrac{4\sqrt{2}}{3} , -\dfrac{\sqrt{2}}{3} \right)
        \]
        现求三角形 $\triangle A'CD$ 面积:
        \[
        \vec{A'C} = \left( -\dfrac{4\sqrt{2}}{3} , -\dfrac{\sqrt{2}}{3}, -\dfrac{\sqrt{2}}{3} \right) ,\quad
        \vec{A'D} = \left( -\dfrac{\sqrt{2}}{3} , -\dfrac{\sqrt{2}}{3}, -\dfrac{4\sqrt{2}}{3}  \right) 
        \]
        \[
        [\triangle A'CD]=\dfrac{1}{2} \left\lVert \vec{A'C} \times \vec{A'D} \right\rVert = \dfrac{1}{2} \sqrt{ \left( \dfrac{2}{3} \right)^2 + \left( -\dfrac{10}{3} \right)^2 + \left( \dfrac{2}{3} \right)^2 } = \sqrt{3}
        \]
        且$\triangle A'CD$ 方程式为
        \[
        E_3 : x - 5y + z + 4\sqrt{2} = 0
        \]
        于是
        \[
        d(B', E_3) = \dfrac{ \left| -\dfrac{\sqrt{2}}{3}  - 5 \cdot \dfrac{4\sqrt{2}}{3}  - \dfrac{\sqrt{2}}{3}  + 4\sqrt{2} \right| }{ \sqrt{1^2 + (-5)^2 + 1^2} }
        = \dfrac{10\sqrt{2}}{9\sqrt{3}}
        \]
        四面体 $A'CB'D$体积为
        \[
        [A'CB'D]=\dfrac{1}{3} \cdot \sqrt{3} \cdot \dfrac{10\sqrt{2}}{9\sqrt{3}} =  \dfrac{10\sqrt{2}}{27}
        \]
    \end{solution}

    \question 空间中有一四棱锥 $E-ABCD$,其中底面 $ABCD$ 为矩形,$\triangle AED$ 为正三角形,且平面 $EAD$ 与平面 $ABCD$ 垂直。设 $G, F$ 分别为 $AD, CD$ 的中点,且 $\angle EBG = 30^\circ$。若点 $A$ 到平面 $EFB$ 的距离为 $2$,求 $AD$。
    \ifprintanswers
    \begin{figure}[H]
        \centering            
        \includegraphics[width=0.6\textwidth]{images/image20.png}
    \end{figure}
    \fi
    \begin{solution}
        设 $AB = CD = a ,AD = BC = 2b,$ 构造空间坐标系
        \[
        A(0,0,0),\
        B(a,0,0),\
        C(a,2b,0),\
        D(0,2b,0),\
        E(0,b,\sqrt{3}b),\
        F\left(\dfrac{a}{2}, 2b, 0\right), \
        G(0,b,0)
        \]
        则
        \[
        \overrightarrow{BE} = (-a, b, \sqrt{3}b), 
        \overrightarrow{BG} = (-a, b, 0), 
        \overrightarrow{BF} = \left(-\dfrac{a}{2}, 2b, 0\right)
        \]
        由已知$\angle EBG = 30^\circ$,
        \[
        \cos \angle EBG = \dfrac{\overrightarrow{BE} \cdot \overrightarrow{BG}}{|\overrightarrow{BE}||\overrightarrow{BG}|} = \dfrac{a^2 + b^2}{\sqrt{a^2 + 4b^2} \cdot \sqrt{a^2 + b^2}} = \sqrt{\dfrac{a^2 + b^2}{a^2 + 4b^2}} = \dfrac{\sqrt{3}}{2}
        \]
        两边平方得
        \[
        \dfrac{a^2 + b^2}{a^2 + 4b^2} = \dfrac{3}{4} \Rightarrow a = 2\sqrt{2}b
        \]
        设平面 $EFB$ 的法向量为
        \[
        \vec{n} = \overrightarrow{BE} \times \overrightarrow{BF} =
        \begin{vmatrix}
        \mathbf{i} & \mathbf{j} & \mathbf{k} \\
        -a & b & \sqrt{3}b \\
        -\dfrac{a}{2} & 2b & 0
        \end{vmatrix}
        = (-2\sqrt{3}b^2,\ -\dfrac{\sqrt{3}}{2}ab,\ -\dfrac{3}{2}ab)
        \]
        平面 $EFB$ 的方程为
        \[
        -2\sqrt{3}b(x - a) - \dfrac{\sqrt{3}}{2}a y - \dfrac{3}{2}a z = 0
        \Rightarrow 4bx + ay + \sqrt{3}a z = 4ab
        \]
        点 $A(0,0,0)$ 到该平面的距离为
        \[
        d = \dfrac{|4ab|}{\sqrt{(4b)^2 + a^2 + 3a^2}} = \dfrac{4ab}{\sqrt{16b^2 + 4a^2}} = 2
        \]
        代入 $a = 2\sqrt{2}b$ 得
        \[
        \dfrac{4b \cdot 2\sqrt{2}b}{\sqrt{16b^2 + 4 \cdot 8b^2}} = 2
        \Rightarrow b = \dfrac{\sqrt{6}}{2}
        \]
        故
        \[
        {AD} = 2b = \sqrt{6}
        \]
    \end{solution}
    
    \question 已知空中有一边长为 $5\sqrt{2}$ 的正四面体,$A$ 为此四面体中距离地面最近的顶点,其他三个顶点距离地面的距离分别为 $5$、$6$、$7$,求 $A$ 到地面的距离。
        \ifprintanswers
        \begin{figure}[H]
            \centering
            \includegraphics[width=0.6\textwidth]{images/image22.png}
        \end{figure}
        \fi
    \begin{solution}
        设$B,C,D$分别在平面$$z=5,z=6,z=7$$ 上,$P$为 ${BD}$ 的中点,$D'$为
        $D \text{ 在平面 } z=6 \text{ 的投影点 },\theta$为
        $\triangle BCD \text{ 与平面 } z=6 \text{ 的夹角 } ,$
        则
        \[
        \begin{cases}
        P \text{ 也在平面 } z=6 \text{ 上,\ 且 } {PD} = \dfrac{5\sqrt{2}}{2} \\
        {DD'} = 1
        \end{cases}
        \Rightarrow \sin \theta = \frac{{DD'}}{{PD}} = \frac{2}{5\sqrt{2}} \Rightarrow \cos \theta = \frac{\sqrt{46}}{5\sqrt{2}}
        \]
        设 $A$ 在 $\triangle BCD$ 的投影点为 $G$,则 $G$ 为 $\triangle BCD$ 的重心,$G$ 的$z$ 坐标为
        \[
        \frac{5+6+7}{3} = 6 
        \]
        即 $G$ 也在平面  $z=6$ 上;由正四面体性质
        \[
        {AG} = \frac{\sqrt{6}}{3} \cdot 5\sqrt{2} = \frac{10\sqrt{3}}{3}, \quad {AG} \perp \triangle BCD
        \]
        因此$A$ 到平面 $z=6$ 的距离为
        \[
        AG \sin\left(\frac{\pi}{2} - \theta\right) = {AG} \cos \theta= \frac{10\sqrt{3}}{3} \cdot \frac{\sqrt{46}}{5\sqrt{2}} = \frac{2}{3} \sqrt{69}
        \]
        所以 $A$ 到地面的距离为
        \[
        6 - \frac{2}{3}\sqrt{69} = \frac{18 - 2\sqrt{69}}{3}
        \]
    \end{solution}
    
    \question 一正方形纸张 $ABCD$,设点 $E, F$ 分别在 ${BC}, {DC}$ 边上,且 $BE:EC = DF:FC = 2:1$。现以正方形 $ABCD$ 为底面,分别将 $B,D$ 以 $AE, AF$ 为谷折线向上折起,使得 $AB, AD$ 重合,并令重合后的点 $B=D=G$。此时,若侧面 $\triangle AEG$（或 $\triangle AFG$）与鸢形底面 $AECF$ 的夹角为 $\theta$,求 $\sin \theta$。
    \begin{solution}
        正方形边长为 $4$,构造空间坐标系,
        $$A(0,0,0), \
        B(3,0,0), \
        C(3,3,0), \
        D(0,3,0), \
        E(3,2,0), \
        F(2,3,0), \
        G(t,t,a).$$
        由长度条件
        \[
        \begin{cases}
        AG = 3 \Rightarrow t^2 + t^2 + a^2 = 9, \\
        GE = 2 \Rightarrow (t-3)^2 + (t-2)^2 + a^2 = 4,
        \end{cases}
        \]
        解得
        \[
        t = \frac{9}{5}, \quad a = \frac{3\sqrt{7}}{5}.
        \]
        向量
        \[
        \overrightarrow{AF} = (2,3,0), \quad \overrightarrow{AG} = \left(\frac{9}{5}, \frac{9}{5}, \frac{3\sqrt{7}}{5}\right).
        \]
        法向量
        \[
        \vec{n} = \overrightarrow{AF} \times \overrightarrow{AG} \parallels (3\sqrt{7}, -2\sqrt{7}, -3)
        \]
        侧面 $AEG$ 平面方程与底面 $AECF$ 方程分别为
        \[
        3\sqrt{7} x - 2\sqrt{7} y - 3 z = 0, \quad z = 0.
        \]
        且
        \[
        \cos \theta = \frac{\vec{n} \cdot (0,0,1)}{|\vec{n}|} = \frac{-3}{\sqrt{(3\sqrt{7})^2 + (-2\sqrt{7})^2 + (-3)^2}} = -\frac{3}{10}
        \]
        故
        \[
        \sin \theta = \sqrt{1 - \cos^2 \theta} = -\frac{\sqrt{91}}{10}
        \]
    \end{solution}
    
    \question 已知平行六面体 $ABCD-EFGH$,直线方程式为:
    \[
    {AB}: \frac{x+3}{2} = \frac{y}{1} = \frac{z+1}{2}, \quad
    {EH}: \frac{x-1}{2} = \frac{y-5}{-1} = \frac{z-1}{1}, \quad
    {CG}: \frac{x-7}{2} = \frac{y+1}{2} = \frac{z-13}{-1}.
    \]
    若点 $(9,7,8)$ 不在上述三条直线上,且为此平行六面体的一个顶点,求此平行六面体的体积。
            \ifprintanswers
            \begin{figure}[H]
                \centering            
                \includegraphics[width=0.25\textwidth]{images/image28.png}
            \end{figure}
            \fi
    \begin{solution}
        设 $D=(9,7,8)$,则 $AD$ 方程式为
        \[
        AD:\frac{x-9}{2} = \frac{y-7}{-1} = \frac{z-8}{1},
        \]
        联立$AB$方程式解得无实数解,故不合题意;设 $F=(9,7,8)$,则
        \[
        BF:\frac{x-9}{2} = \frac{y-7}{2} = \frac{z-8}{-1},
        \]
        解得 $B=(7,5,9)$;由于
        \[
        \vec{v}_{AB} \times \vec{v}_{AD} = (2,1,2)\times(2,-1,1)=(3,2,-4),
        \]
        平面 $ABCD$ 方程为
        \[
        3x + 2y - 4z = -5.
        \]
        设 $F$ 在 $EH,CG$ 上的垂足分别为
        \[
        P(2m + 1, -m + 5, m + 1),\quad Q(2n + 7, 2n - 1, -n + 13),
        \]
        则
        \[
        \overrightarrow{FP} = (2m - 8, -m - 2, m - 7), \quad \overrightarrow{FQ} = (2n - 2, 2n - 8, -n + 5).
        \]
        由垂直条件$\overrightarrow{FP} \cdot \vec{v}_{EH} = 0, \ \overrightarrow{FQ} \cdot \vec{v}_{CG} = 0,$解得
        \[
        m = \frac{7}{2}, \quad n = \frac{25}{9}.
        \]
        故
        \[
        |\overrightarrow{FP}| = \frac{\sqrt{174}}{2}, \quad |\overrightarrow{FQ}| = \frac{2 \sqrt{53}}{3}.
        \]
        设 $\overrightarrow{EF}$ 与 $\overrightarrow{EH}$ 的锐角为 $\alpha,\overrightarrow{FG}$ 与 $\overrightarrow{CG}$ 的锐角为 $\beta$,则
        \[
        \cos \alpha = \frac{|\overrightarrow{EF} \cdot \overrightarrow{EH}|}{|\overrightarrow{EF}||\overrightarrow{EH}|} = \frac{5}{3 \sqrt{6}}, \quad \cos \beta = \frac{|\overrightarrow{FG} \cdot \overrightarrow{CG}|}{|\overrightarrow{FG}||\overrightarrow{CG}|} = \frac{1}{3 \sqrt{6}},
        \]
        \[
        \sin \alpha = \frac{\sqrt{29}}{3 \sqrt{6}}, \quad \sin \beta = \frac{\sqrt{53}}{3 \sqrt{6}}.
        \]
        又由
        \[
        |\overrightarrow{FP}| = EF \sin \alpha = \frac{\sqrt{174}}{2} \implies EF = 9,
        \]
        \[
        |\overrightarrow{FQ}| = FG \sin \beta = \frac{2 \sqrt{53}}{3} \implies FG = 2 \sqrt{6}.
        \]
        设平行六面体高为点 $F$ 到平面 $ABCD$ 的距离
        \[
        d = \frac{|3 \cdot 9 + 2 \cdot 7 - 4 \cdot 8 + 5|}{\sqrt{3^2 + 2^2 + (-4)^2}} = \frac{14}{\sqrt{29}}.
        \]
        故体积为
        \[
        V = EF \times FG \times \sin \alpha \times d = 9 \times 2 \sqrt{6} \times \frac{\sqrt{29}}{3 \sqrt{6}} \times \frac{14}{\sqrt{29}} = 84.
        \]
    \end{solution}

    \question 将8个半径为2的球分两层放置于一个圆柱形容器中,使得每个球和与其相邻的四个球均相切,且与圆柱的一个底面和侧面都相切,则圆柱的高为
    \begin{solution}
        \textcolor{red}{(待解)}
    \end{solution}
    
    \question 已知三棱柱 $\Omega:ABC-A_{1}B_{1}C_{1}$ 的9条棱长均相等,记底面$ABC$所在平面为$\alpha$,若$\Omega$的另外四个面(即面 $A_{1}B_{1}C_{1}$, $ABBA_{1}$, $ACCA_{1}$, $BCCB_{1}$) 在$\alpha$上投影的面积从小到大重排后依次为 $2\sqrt{3},3\sqrt{3},4\sqrt{3},5\sqrt{3}$ ,求$\Omega$的体积.
        \ifprintanswers
        \begin{figure}[H]
            \centering            
            \includegraphics[width=0.35\textwidth]{images/image18.png}
        \end{figure}
        \fi
    \begin{solution}
        设点 $A_{1}, B_{1}, C_{1}$ 在平面$\alpha$上的投影分别为$D,E,F$,则面 
        \[
        A_{1}B_{1}C_{1}, ABB_{1}A_{1}, ACC_{1}A_{1}, BCC_{1}B_{1}
        \]
        在$\alpha$上的投影面积分别为 
        \[
        [\triangle DEF],[ABED],[ACFD],[BCFE]
        \]
        由已知及三棱柱的性质, $\triangle DEF$ 为正三角形,且$ABED, ACFD, BCFE$均为平行四边形.
        由对称性,仅需考虑点$D$位于$\angle BAC$内的情形(如图所示)。
        显然此时有 
        \[
        [ABED]+[ACFD]=[BCFE]
        \]
        由于 
        \[
        \{[\triangle DEF],[ABED],[ACFD],[BCFE]\}=\{2\sqrt{3},3\sqrt{3},4\sqrt{3},5\sqrt{3}\}\] 
        故 $[ABED],[ACFD]$ 必为 $2\sqrt{3},3\sqrt{3}$ 的排列, $[BCFE]=5\sqrt{3}$ ,进而 $S_{\triangle DEF}=4\sqrt{3}$ ,得 $\triangle DEF$ 的边长为4,即正三棱柱的各棱长均为4.
        
        不妨设 $[ABED]=2\sqrt{3},[ACFD]=3\sqrt{3}$ ,则 
        \[
        [\triangle ABD]=\sqrt{3},[\triangle ACD]=\frac{3\sqrt{3}}{2}
        \]
        取射线$AD$与线段$BC$的交点$X$,则 
        \[
        \frac{BX}{CX}=\frac{[\triangle ABD]}{[\triangle ACD]}=\frac{2}{3}
        \] 
        故 $BX=\dfrac{8}{5}$ ,因此 
        \[
        AX=\sqrt{AB^{2}+BX^{2}-2AB\cdot BX\cdot \cos 60^{\circ}}=\frac{4\sqrt{19}}{5}
        \]
        而 
        \[
        \frac{AD}{AX}=\frac{[\triangle ABD]+[\triangle ACD]}{[\triangle ABC]}=\frac{5}{8}
        \] 
        故 $AD=\dfrac{\sqrt{19}}{2}$,于是$\Omega$的高 \[
        h=\sqrt{AA_{1}^{2}-AD^{2}}=\frac{3\sqrt{5}}{2}\]
        又$[\triangle ABC]=4\sqrt{3}$ ,故$\Omega$的体积
        \[
        V=[\triangle ABC]\cdot h=6\sqrt{15}
        \]
    \end{solution}

    
    \question 正方体 $ABCD-A_1B_1C_1D_1$ 的底面 $A_1B_1C_1D_1$ 内有一个动点 $M$, 且 $BM \parallel$ 平面 $AD_1C$, 则 $\tan \angle D_1MD$ 的最大值是
    \begin{solution}
设正方体 $ABCD - A_1B_1C_1D_1$ 的边长为 $a$,建立坐标系,使 $D(0,0,0)$,$A(a,0,0)$,$C(0,a,0)$,$D_1(0,0,a)$。因此 $A_1(a,0,a)$,$B_1(a,a,a)$,$C_1(0,a,a)$。

设底面动点 $M$ 的坐标为 $(x,y,a)$,其中 $0 \le x,y \le a$。点 $B$ 的坐标为 $(a,a,0)$,因此
\[
\vec{BM} = (x - a, y - a, a).
\]

题设中 $BM \parallel$ 平面 $AD_1C$,则 $\vec{BM}$ 垂直于该平面的法向量。

求平面 $AD_1C$ 的法向量,可取 $\vec{AD_1} = (-a, 0, a)$ 和 $\vec{AC} = (-a, a, 0)$,则
\[
\vec{n} = \vec{AD_1} \times \vec{AC} = (-a^2, a^2, -a^2) \propto (1, -1, 1).
\]

由 $\vec{BM} \cdot \vec{n} = 0$ 得:
\[
(x - a)(1) + (y - a)(-1) + a(1) = 0 \Rightarrow x - y + a = 0 \Rightarrow y = x + a.
\]

\textcolor{red}{结合 $0 \le y \le a$,代入得 $0 \le x + a \le a \Rightarrow -a \le x \le 0$,又因 $x \ge 0$,故 $x = 0$,从而 $y = a$。}

故点 $M$ 为 $C_1(0,a,a)$,接下来计算 $\tan \angle D_1MD$,其中 $D_1(0,0,a)$,$D(0,0,0)$。

设 $\theta = \angle D_1MD$,则有:

\[
\vec{MD_1} = (0, -a, 0), \quad \vec{MD} = (0, -a, -a),
\]
\[
\cos \theta = \frac{\vec{MD_1} \cdot \vec{MD}}{|\vec{MD_1}| \cdot |\vec{MD}|} = \frac{a^2}{a \cdot a\sqrt{2}} = \frac{1}{\sqrt{2}} \Rightarrow \theta = 45^\circ,
\]
\[
\tan \angle D_1MD = \tan 45^\circ = \boxed{1}.
\]
\textcolor{red}{(by chatgpt,待验证)}

\end{solution}
    \question 如图,在面积为2的矩形$ABCD$中,点$E$为$AD$边的中点.将$\triangle ABE$和$\triangle DEC$分别沿$BE,CE$翻折,使得$A,D$重合于点$P$,则三棱锥$P-EBC$体积的最大值为
    \begin{solution}
        \textcolor{red}{(待解)}
    \end{solution}
    
    \question 设四棱锥 $P-ABCD$,面 $PAB \perp$ 底面 $ABCD$,且 $PA=PB=\sqrt{5},AB=BC=AD=2$,求该四棱锥体积的最大值。
        
    \begin{solution}
        由题意,四棱锥的高为 $2$,只需求底面四边形 $ABCD$ 面积的最大值。
        
        将四边形 $ABCD$ 分为 $\triangle ABD$ 和 $\triangle BCD$,设 $\angle BAD = \alpha$,由余弦定理得
        \[
        BD^2 = AB^2 + AD^2 - 2 \cdot AB \cdot AD \cdot \cos\alpha = 4 + 4 - 8 \cos\alpha = 8(1 - \cos\alpha)
        \]
        由三角恒等式 $1 - \cos\alpha = 2\sin^2 \dfrac{\alpha}{2}$,得
        \[
        BD = \sqrt{8(1 - \cos\alpha)} = \sqrt{16 \sin^2 \frac{\alpha}{2}} = 4 \sin \frac{\alpha}{2}
        \]
        三角形 $ABD$ 的面积为:
        \[
        [\triangle ABD] = \frac{1}{2} \cdot AB \cdot AD \cdot \sin\alpha = 2 \sin\alpha
        \]
        当 $BC \perp BD$ 时,$\;\triangle BCD$ 面积最大(贪婪算法),
        \[
        [\triangle BCD] = \frac{1}{2} \cdot BC \cdot BD = 4 \sin \frac{\alpha}{2}.
        \]
        因此四边形面积为
        \[
        S = [ABCD] = 2 \sin \alpha + 4 \sin \frac{\alpha}{2}.
        \]
        令 $x = \sin \dfrac{\alpha}{2} \in (0,1)$,则
        \[
        S(x) = 4 \left( x + x \sqrt{1 - x^{2}} \right).
        \]
        求导:
        \[
        S'(x) = 4 \left( 1 + \sqrt{1 - x^{2}} - \frac{x^{2}}{\sqrt{1 - x^{2}}} \right).
        \]
        解 $S'(x) = 0$,得极值点$x = \frac{\sqrt{3}}{2}$,
        代入得最大面积
        \[
        S_{\max} = 3 \sqrt{3}.
        \]
        故四棱锥体积最大为
        \[
        V_{\max} = \frac{1}{3} \times 3\sqrt{3} \times 2 = 2 \sqrt{3}.
        \]
    \end{solution}

    \question 一个盛满水的半球体容器, 其半径为 $6$, 若倾斜 $45^{\circ}$ 后,求容器溢出的水体积。
    \ifprintanswers
    \begin{figure}[H]
        \centering        
        \includegraphics[width=0.4\textwidth]{images/image101.jpg}
    \end{figure}
    \fi
    \begin{solution}
        剩下的水体积即为上图着色区域绕 $x$ 轴旋转的体积,即
        \[
        \pi \int_{3\sqrt 2}^{6} (36 - x^2)\; dx = \pi \left[36x - {1\over 3} x^3 \right]_{3\sqrt 2}^{6} = (144 - 90\sqrt 2)\pi
        \]
        因此溢出的水体积为
        \[
        \frac{2}{3}\pi \cdot 6^3 - (144 - 90\sqrt 2)\pi = 90\sqrt 2 \pi
        \]
    \end{solution}

    \question 已知一个底面半径为 $3$,高为 $3$ 的直圆柱,平面 $E$ 通过底面的直径 $AB$,且平面 $E$ 与底面的夹角为 $45^\circ$,此时平面 $E$ 将直圆柱切割成两部分,求较小部分的体积。
    \begin{solution}
        设圆柱底面中心为原点,直径 $AB$ 在 $x$ 轴上,平面 $E$ 与 $AB$成 $45^\circ$,故平面方程可设为 \(z = x\),底面为半径 $3$ 的半圆区域 
        \[
        D: x^{2}+y^{2} \le 9,\ x \ge 0
        \]
        体积为
        \[
        V = \iint_D z \, dA = \iint_D x \, dA
        \]
        用极坐标积分,令 \(x = r\cos\theta,y = r\sin\theta\),则
        \[
        V = \int_{-\frac{\pi}{2}}^{\frac{\pi}{2}} \int_0^3 r\cos\theta \cdot r\, dr\, d\theta = \int_{-\frac{\pi}{2}}^{\frac{\pi}{2}} \cos\theta \int_0^3 r^{2} dr\, d\theta
        = 9 \int_{-\frac{\pi}{2}}^{\frac{\pi}{2}} \cos\theta \, d\theta
        =9[\sin\theta]_{-\frac{\pi}{2}}^{\frac{\pi}{2}} = 18
        \]
    \end{solution}
    \begin{solution}
        设底面圆方程式为
        \[
        x^2 + y^2 = R^2,
        \]
        所求体积由许多直角三角形累积而成;该三角形与底面圆的交点坐标为 $(x,y)$,则三角形的底长为 $y$,高为 $y \tan \theta$,因此三角形面积为
        \[
        \frac{1}{2} y^2 \tan \theta = \frac{1}{2} (R^2 - x^2) \tan \theta.
        \]
        故所求体积为
        \[
        \left. \int_{-R}^R \frac{1}{2} (R^2 - x^2) \tan \theta \, dx \right|_{R=3,\theta=45^\circ}= \left.\frac{2}{3} R^3 \tan \theta \right|_{R=3,\theta=45^\circ}=18
        \]
    \end{solution}

    \question 奥里克有一个杯子,是一个底面为圆的直圆锥体,底面半径为 \(\frac{1}{2}\),斜高为 \(1\),杯中装满了奶茶。奶茶恰好填满杯子斜高的一半。当奥里克将杯子倾斜到刚好要溢出的程度时,如图所示,新的斜高 \(EA\) 与倾斜后的奶茶表面所形成的椭圆的长轴 \(ET\) 构成夹角 \(\angle TEA\)。求 \(\cos{\angle TEA}\)。     %https://artofproblemsolving.com/community/c3073284h3152961p28644263
        \begin{figure}[H]
            \centering
            \includegraphics[width=0.2\textwidth]{images/image6.png}
        \end{figure}
    \begin{solution}
        \textcolor{red}{(待解)}
    \end{solution}

    \question 在正三棱锥 $P-ABC$ 中, $AB=1$, $AP=2$, 过 $AB$ 的平面 $\alpha$ 将其体积平分, 则棱 $PC$ 与平面 $\alpha$ 所成角的余弦值为
    \begin{solution}
        设 $\triangle ABC$ 的外接圆半径为 $R$,且$P$ 到底面 $ABC$ 的距离为 $d$ ,则由正弦定理知, 
        \[
        2R=\frac{1}{\sin 120^\circ}=\frac{2\sqrt{3}}{3},\;d=\sqrt{2^2-\left(\dfrac{\sqrt{3}}{3}\right)^2}=\dfrac{\sqrt{33}}{3}
        \]
        故
        \[
        [P-ABC]=\frac{1}{3}\cdot\frac{\sqrt{3}}{4}\cdot 1^2\cdot\frac{\sqrt{33}}{3}=\frac{\sqrt{11}}{12}
        \]
        设 $PC$ 的中点为 $D$, 由
        \[
        A=(0,0,0),B=(1,0,0),C=\left(\frac12,\frac{\sqrt3}{2},0\right),P\left(\frac12,\frac{\sqrt3}{6},\frac{\sqrt{33}}{6}\right),
        \]
        则 
        \[
        AD^2=\left(\frac12\right)^2+\left(\frac{\sqrt3}{3}\right)^2+\left(\frac{\sqrt{33}}{6}\right)^2 \Rightarrow AD=BD=\frac{\sqrt{6}}{2}\]
        故 
        \[
        [\triangle ABD]=\frac{1}{2}\cdot 1\cdot\sqrt{\left(\frac{\sqrt{6}}{2}\right)^2-\left(\frac{1}{2}\right)^2}=\frac{\sqrt{5}}{4}
        \]
        设 $P$ 到平面 $\alpha$ 的距离为 $h$, 则由等体积法知 
        \[
        \frac{1}{3}\times\frac{\sqrt{5}}{4}\times h=\frac{\sqrt{11}}{24}\Rightarrow h=\frac{\sqrt{55}}{10}
        \]       
        所以棱 $PC$ 与平面 $\alpha$ 所成角的余弦值为 
        \[
        1-\left(\frac{\frac{\sqrt{55}}{10}}{1}\right)^2=\frac{3\sqrt{5}}{10}
        \]
    \end{solution}
    
    \question 一个底面半径为 $5$,高为 $12$ 的直圆锥中放入了三个互相接触的全等球体,每个球的半径为 $r$。每个球都与另外两个球相切,并且也与圆锥的底面及侧面相切。求 $r$ 的值。
     \ifprintanswers
    \begin{figure}[H]
        \centering
        \includegraphics[width=0.3\textwidth]{images/image11.png}
        \quad
        \includegraphics[width=0.15\textwidth]{images/image12.png}
    \end{figure}
    \fi
    \begin{solution}
        设圆锥的顶点为 $A$,底面圆心为 $B$。取其中一个球的球心为 $E$,三球心形成正三角形,中心为 $C$,满足 $AC = AE + EC$。我们来分别表示 $AE$ 与 $EC$。
        
        从顶部鸟瞰圆锥,三个球心构成一个边长为 $2r$ 的正三角形。点 $C$ 是其重心,正三角形边长为 $r\sqrt{3}$,故
        \[
        EC = \frac{2r\sqrt{3}}{3}
        \]
        将 ${OE}$ 延长至与底面交于点 $F$使得$\triangle AEF \backsim \triangle ACB$,由角平分线定理,
        \[
        \frac{OE}{OF} = \frac{AE}{AF} = \frac{5}{13} \Rightarrow OF = \frac{13r}{5}
        \]
        因此
        \[
        EF = OE + OF = r + \frac{13r}{5} = \frac{18r}{5}
        \]
        再由相似三角形
        \[
        AE = \frac{5}{12} EF  = \frac{3r}{2}
        \]
        我们已知 $AC = AE + EC = 5$,代入得
        \[
        \frac{3r}{2} + \frac{2r\sqrt{3}}{3} = 5
        \Rightarrow  r = \frac{30}{9 + 4\sqrt{3}} = \frac{90 - 40\sqrt{3}}{11}
        \]
    \end{solution}

    \question 求空间中一点 $P(-5,0,-8)$ 到直线 
    \[
    L:\frac{x-3}{1}=\frac{y-2}{-2}=\frac{z+1}{2}
    \] 
    的距离。
    \begin{solution}
        设 $Q(3+t,2-2t,-1+2t)$ 满足 $\overrightarrow{PQ} \perp L$,则
        \[
        \overrightarrow{PQ}=(8+t,2-2t,7+2t),\quad \overrightarrow{PQ} \cdot (1,-2,2)=0
        \]
        解得
        \[
        t=-2 \Rightarrow Q(1,6,-5)
        \]
        于是
        \[
        d(P,L)=\overline{PQ}=\sqrt{(-5-1)^{2}+(0-6)^{2}+(-8+5)^{2}}=\sqrt{81}=9
        \]
    \end{solution}
    \begin{solution}
        设直线 $L$ 上动点 $Q(3+t,2-2t,-1+2t)$,则
        \[
        PQ=\sqrt{(8+t)^{2}+(2-2t)^{2}+(7+2t)^{2}}=\sqrt{9(t+2)^{2}+81},
        \]  
        当 $t=-2$ 时有最小值 $9$。
    \end{solution}
    \begin{solution}
        直线 $L$ 上取一点 $A(4,0,1)$ ,有$\overrightarrow{AP}=(-9,0,-9)$,直线的方向向量 $\vec{v}=(1,-2,2)$,则
        \[
        \overrightarrow{AP}\times\vec{v}=(-18,9,18)
        \] 
        平行四边形面积为 
        \[
        \sqrt{(-18)^{2}+9^{2}+18^{2}}=27,
        \] 
        又 $|\vec{v}|=3$,故    
        \[
        d(P,L)=\frac{27}{|\vec{v}|}=9
        \]
    \end{solution}
    \begin{solution}
        直线 $L$ 上取一点 $A(4,0,1)$,$\overrightarrow{AP}=(-9,0,-9)$ 在 $\vec{v}=(1,-2,2)$ 上的正射影 
        \[
        \overrightarrow{AH}=(-3,6,-6),
        \]  
        $P$ 在直线 $L$ 上的投影点 $H$坐标为
        \[
        (4,0,1)+(-3,6,-6)=(1,6,-5)
        \]  
        故
        \[
        PH=d(P,L)=\sqrt{(-5-1)^{2}+(0-6)^{2}+(-8+5)^{2}}=9
        \]
    \end{solution}

    \question 设 $A(0,1,2),B(-1,2,1),C(1,0,1)$ 为空间中的三点,求 $\triangle ABC$ 的垂心坐标。
    \begin{solution}
        首先有
        \[
        \vec n = \overrightarrow{AB} \times \overrightarrow{AC} =(-1,1,-1) \times (1,-1,-1)= (-2,-2,0)
        \]
        得平面
        \[
        E:\ x+y=1
        \]
        假设垂心 $H(a,b,c)$,由 $\overrightarrow{AH} \perp \overrightarrow{BC},\overrightarrow{BH} \perp \overrightarrow{AC},\overrightarrow{CH} \perp \overrightarrow{AB}$,
        \[
        (a,b-1,c-2) \cdot (2,-2,0) = (a+1,b-2,c-1) \cdot (1,-1,-1) = (a-1,b,c-1) \cdot (-1,1,-1) = 0
        \]
        解得 
        \[
        H(a,a+1,3)
        \]
        又 $H$ 在平面 $x+y=1$ 上,故
        \[
        a + (a+1) = 1 \Rightarrow a = 0
        \]
        因此 $H = (0,1,3)$。
    \end{solution}

    \question 坐标空间中有三个彼此互相垂直的向量 $\textbf{u}, \textbf{v}, \textbf{w}$。已知 $\textbf{u} - \textbf{v} = (2, -1, 0)$,且 $\textbf{v} - \textbf{w} = (-1, 2, 3)$。求由 $\textbf{u}, \textbf{v}, \textbf{w}$ 所张出的平行六面体之体积。
    \begin{solution}
        由已知$\textbf{u}, \textbf{v}, \textbf{w}$互相垂直且$\textbf{u} - \textbf{v} = (2, -1, 0),\textbf{v} - \textbf{w} = (-1, 2, 3)$,
        \[
        (\textbf u - \textbf v) \cdot (\textbf v - \textbf w)
        = -|\textbf v|^2= -4 \Rightarrow |\textbf v| = 2
        \]
        且
        \[
        (\textbf u - \textbf v) \cdot (\textbf u - \textbf v)
        = |\textbf u|^2 + 4 = 5
        \Rightarrow |\textbf u| = 1
        \]
        及
        \[
        (\textbf v - \textbf w) \cdot (\textbf v - \textbf w)
        = 4 + |\textbf w|^2
        = 14
        \Rightarrow |\textbf w| = \sqrt{10}
        \]
        故体积为
        \[
        |\textbf u|\cdot|\textbf v|\cdot|\textbf w| = 2\sqrt{10}
        \]
    \end{solution}
    
    \question 在空间中,有三个不共平面的非零向量 $\textbf{a}, \textbf{b}, \textbf{c}$,满足
    \[
    (\textbf{a} \times \textbf{b}) \cdot ((\textbf{b} \times \textbf{c}) \times (\textbf{c} \times \textbf{a}))=7,
    \]
    求以三向量 $(3\textbf{a}+\textbf{b}+\textbf{c}), (\textbf{a}-\textbf{b}+2\textbf{c}), (\textbf{b}+\textbf{c})$ 所张成的平行六面体体积。
    \begin{solution}
        由向量性质,
        \[
        \textbf{a}\times (\textbf{b}\times \textbf{c})=(\textbf{a}\cdot \textbf{c})\textbf{b}- (\textbf{a}\cdot \textbf{b})\textbf{c}, 
        \quad (\textbf{a}\times \textbf{b})\cdot \textbf{c}=\textbf{a} \cdot (\textbf{b}\times \textbf{c}),
        \quad (\textbf{a}\times \textbf{b})\cdot \textbf{a} = (\textbf{a}\times \textbf{b})\cdot \textbf{b} =0
        \]
        因此
        \begin{align*}
        (\textbf{a}\times \textbf{b})\cdot ((\textbf{b}\times \textbf{c})\times (\textbf{c}\times \textbf{a}))
        &= (\textbf{a}\times \textbf{b})\cdot \big(((\textbf{b}\times \textbf{c})\cdot \textbf{a})\textbf{c} - ((\textbf{b}\times \textbf{c})\cdot \textbf{c})\textbf{a}\big) \\
        &= (\textbf{a}\times \textbf{b})\cdot ((\textbf{b}\times \textbf{c})\cdot \textbf{a})\textbf{c} \\
        &= ((\textbf{b}\times \textbf{c})\cdot \textbf{a})\,((\textbf{a}\times \textbf{b})\cdot \textbf{c}) \\
        &= (\textbf{a}\cdot (\textbf{b}\times \textbf{c}))^2 = 7
        \end{align*}
        即
        \[
        \textbf{a} \cdot (\textbf{b}\times \textbf{c}) = \pm \sqrt{7}
        \]
        三向量张成的平行六面体体积为
        \begin{align*}
        V &= \big|(3\textbf{a}+\textbf{b}+\textbf{c})\cdot  \big((\textbf{a}-\textbf{b}+2\textbf{c})\times (\textbf{b}+\textbf{c})\big)\big| \\
        &= \text{abs}\left(\begin{vmatrix} 3 & 1 & 1 \\ 1 & -1 & 2 \\ 0 & 1 & 1 \end{vmatrix}\right) \cdot |\textbf{a} \cdot (\textbf{b}\times \textbf{c})| = 9\sqrt{7}
        \end{align*}
    \end{solution}

    \question 两平面
    \[
    E_1:\;x + k y + z = 3, \quad E_2:\;x + y + k z = 5
    \]
    夹角为 $60^\circ$,求 $k$。  
    \begin{solution}
        设 $E_1, E_2$ 的法向量分别为  
        $\vec{n_1} = (1, k, 1)$,$\vec{n_2} = (1, 1, k)$。
        两平面夹角满足:
        \[
        \cos\theta = \frac{|\vec{n_1} \cdot \vec{n_2}|}{|\vec{n_1}||\vec{n_2}|}
        \quad \text{其中 } \theta = 60^\circ,\; \cos\theta = \frac{1}{2}
        \]
        计算内积与模长:
        \[
        \vec{n_1} \cdot \vec{n_2} = 1 + k + k = 1 + 2k
        \]
        \[
        |\vec{n_1}| = \sqrt{1 + k^2 + 1} = \sqrt{k^2 + 2},\quad
        |\vec{n_2}| = \sqrt{1 + 1 + k^2} = \sqrt{k^2 + 2}
        \]
        代入公式:
        \[
        \frac{|1 + 2k|}{k^2 + 2} = \frac{1}{2}
        \Rightarrow |1 + 2k| = \frac{1}{2}(k^2 + 2)
        \]
        分正负两种情况:
        \begin{itemize}
          \item 若 $1 + 2k \ge 0$,则 $1 + 2k = \frac{1}{2}(k^2 + 2)\Rightarrow k=0,4$；
          \item 若 $1 + 2k < 0$,则 $-(1 + 2k) = \frac{1}{2}(k^2 + 2) \Rightarrow k=-2$。
        \end{itemize}
        解得$k = -2,\; 0,\; 4$
    \end{solution}

    \question 两平面
    \[
    E_1:\;2x+y-z-3=0,\quad E_2:\;x+2y+z=0
    \]
    过点 $(2,1,-1)$ 且同时垂直于 $E_1,E_2$ 的平面方程为  
    \begin{solution}
        设所求平面法向量为 $\vec{n} = \vec{n_1} \times \vec{n_2}$,其中 $\vec{n_1} = (2,1,-1),\vec{n_2} = (1,2,1)$,则
        \[
        \vec{n} = 
        \begin{vmatrix}
        \mathbf{i} & \mathbf{j} & \mathbf{k} \\
        2 & 1 & -1 \\
        1 & 2 & 1
        \end{vmatrix}
        = \mathbf{i}(1 \cdot 1 - (-1)\cdot 2) 
        - \mathbf{j}(2 \cdot 1 - (-1)\cdot 1) 
        + \mathbf{k}(2 \cdot 2 - 1 \cdot 1) 
        = (3, -3, 3)
        \]
        所以可取法向量为 $(1, -1, 1)$,故平面方程为 $x - y + z = d$,将点 $(2,1,-1)$ 代入得
        \[
        2 - 1 - 1 = 0 \Rightarrow d = 0
        \]
        故所求平面方程为 $x - y + z = 0$。
    \end{solution}

    \question 动点 $P(x,y,z)$ 在平面
    \[
    2x+y-2z-5=0
    \]
    上移动,求
    \[
    \sqrt{(x-3)^2+(y+1)^2+(z+2)^2}
    \]
    的最小值。  
    \begin{solution}
        设定点 $A(3, -1, -2)$,原式即动点 $P(x,y,z)$ 到 $A$ 的距离,最小距离即点 $A$ 到平面的距离。
        
        平面法向量为 $\vec{n} = (2, 1, -2)$,于是
        \[
        d = \frac{|2\cdot3 + (-1) + (-2)\cdot(-2) - 5|}{\sqrt{2^2 + 1^2 + (-2)^2}} = \frac{4}{3}
        \]
    \end{solution}

    \question 已知 $\alpha,\beta,\gamma,\delta\in\mathbb{R}$,且
    \[
    \alpha^2+\beta^2+\gamma^2-2\alpha-4\beta-6\gamma+13=0,
    \]
    求
    \[
    (\alpha-\delta-8)^2+(\beta-2\delta-9)^2+(\gamma-3\delta-10)^2
    \]
    的最小值。
    \begin{solution}
        写成
        \[
        (\alpha-1)^2+(\beta-2)^2+(\gamma-3)^2=1,
        \]
        即点 $P(\alpha,\beta,\gamma)$ 在以 $C(1,2,3)$ 为球心、半径 $r=1$ 的球面上。记点 $Q$ 为参数 $t$ 下的直线
        \[
        L:\ \frac{x-8}{1}=\frac{y-9}{2}=\frac{z-10}{3}=t,
        \]
        则 $Q=(t+8,\;2t+9,\;3t+10)$,且
        \[
        (\alpha-(\delta+8))^2 +(\beta-(2\delta+9))^2 +(\gamma-(3\delta+10))^2 = PQ^2,
        \]
        其中
        \[
        PQ = d(C,L) - r = \sqrt{(t+7)^2+ (2t+7)^2+ (3t+7)^2} -1 =\sqrt{14(t+3)^2+21}-1
        \]
        在$t=-3$时取最小值$\sqrt{21}-1$,此时$PQ^2$为
        \[
        (\sqrt{21}-1)^2=22-2\sqrt{21}
        \]
    \end{solution}

    \question 在空间中,给定两歪斜线  $$L_1:\frac{x-7}{2}=\frac{y-2}{1}=\frac{z-10}{-2},L_2:\frac{x-3}{1}=\frac{y-9}{-2}=\frac{z-2}{1},$$若在直线 $L_1$ 上取一点 $P$,在直线 $L_2$ 上取一点 $Q$使得线段 ${PQ}$ 最短,试求 ${PQ}$ 的距离。
    \begin{solution}
        $L_1,L_2$ 的方向向量分别为
        \[
        \textbf{u} = (2,1,-2) ,\textbf{v} = (1,-2,1)
        \]
        则
        \[
        \textbf{n} = \textbf{u} \times \textbf{v} = (-3,-4,-5)
        \]
        包含$L_1$ 且法向量为$\textbf{n}$的平面方程式为
        \[
        E:-3(x-7) - 4(y-2) - 5(z-10) = 0 \Rightarrow 3x + 4y + 5z = 79
        \]
        取 $P(3,9,2) \in L_2$,则点到平面 $E$ 的距离为
        \[
        d(P, E) = \frac{|3(3) + 4(9) + 5(2) - 79|}{\sqrt{3^2 + 4^2 + 5^2}} = \frac{12\sqrt{2}}{5}
        \]
    \end{solution}

    \question 已知平面 $E$ 包含直线 $\begin{cases} x=1\\ y+z=3\\\end{cases}$,且两条直线
    \[
    L_1:\ \frac{x-9}{2}=\frac{y+10}{1}=\frac{z-11}{-1},\quad
    L_2:\ \frac{x-9}{2}=\frac{y+10}{2}=\frac{z-11}{-1}.
    \]
    求平面 $E$ 的方程式。
    \begin{solution}
        可得两直线$L_1,L_2$的方向向量为
        \[
        \mathbf{u}=(2,1,-1),\quad \mathbf{v}=(2,2,-1).
        \]
        由已知直线 $\begin{cases}x=1\\ y+z=3 \\\end{cases}$ 可取方向向量
        \[
        \mathbf{w}=(0,1,-1)
        \]
        计算得
        \[
        (\vec u\times\vec v)\times\vec w
        =(2, -1, -1 )
        \]
        又$E$通过$(1,1,2)$,由点法式平面方程得
        \[
        E:2(x-1)+1(y-1)+1(z-2)=0 \Rightarrow 2x-y-z=-1
        \]
    \end{solution}

    \question 空间中两歪斜线 $$L_1: \frac{x-3}{1} = \frac{y}{2} = \frac{z+2}{-2},\quad L_2: \frac{x}{3} = \frac{y-2}{1} = \frac{z+1}{-2},$$若正 $\triangle PQR$ 中,$P$ 在 $L_1$ 上,且 $Q$、$R$ 都在 $L_2$ 上,求 $\triangle PQR$ 的最小面积。
    \begin{solution}
        $L_1,L_2$ 的方向向量分别为
        \[
        \textbf{u} = (1,2,-2) ,\textbf{v} = (3,1,-2) 
        \]
        则\[
        \textbf{n} = \textbf{u} \times \textbf{v} = (-2, -4, -5)
        \]
        $\textbf{n}$ 为 $L_1$ 与 $L_2$ 所张平面的法向量,设包含 $L_2$ 的平面 $E$,带入 $L_2$ 上点 $(0,2,-1)$ 得:
        \[
        E: -2x -4(y-2) -5(z+1) = 0 \Rightarrow 2x + 4y + 5z = 3
        \]
        取 $P = (3,0,-2) \in L_1$,则 $P$ 到平面 $E$ 的距离为正三角形的高:
        \[
        d(P,E) = \frac{|2(3) + 4(0) + 5(-2) - 3|}{\sqrt{2^2 + 4^2 + 5^2}} = \frac{7}{3\sqrt{5}}
        \]
        设边长为 $s$,正三角形面积公式为 $A = \frac{\sqrt{3}}{4}s^2$,其中高 $h = \frac{\sqrt{3}}{2}s$,反解出 $s$ 得
        \[
        h = \frac{7}{3\sqrt{5}} = \frac{\sqrt{3}}{2}s \Rightarrow s = \frac{14}{3\sqrt{15}}
        \]
        最小面积为\[
        A = \frac{\sqrt{3}}{4} \cdot \left( \frac{14}{3\sqrt{15}} \right)^2 = \frac{49\sqrt{3}}{135}
        \]
    \end{solution}

    \question 空间中,已知点 $A(3,-10,11)$,直线
    \[
    L_1: 
    \begin{cases}
    x = 2 + t \\
    y = -1 + 2t, \quad t \in \mathbb{R} \\
    z = 3 - 2t
    \end{cases}
    \]
    与
    \[
    L_2: 
    \begin{cases}
    x = -1 + 2s \\
    y = 2 - 2s, \quad s \in \mathbb{R} \\
    z = -s
    \end{cases}
    \]
    今由 $A$ 点出发,经 $L_1$ 上一点再到达 $L_2$ 上一点,求此路径的最小值。
    \ifprintanswers
    \begin{figure}[H]
        \centering
        \includegraphics[width=0.3\linewidth]{images/image49.png}
    \end{figure}
    \fi
    \begin{solution}
        首先发现
        \begin{itemize}
            \item $L_1$ 与 $L_2$ 为歪斜线,且方向向量垂直。  
            \item $L_1$ 与 $L_2$ 的公垂线在 $L_1$ 上的垂足为 $P(3,1,1)$,在 $L_2$ 上的垂足为 $Q(1,0,-1)$,且$PQ = 3$。
            \item 点 $A$ 到 $L_1$ 的距离 $d(A,L_1) = 5$,且 $A$ 在 $L_1$ 上的投影点 $R$ 距离 $P$ 为 $PR = 14$。 
        \end{itemize}
        设 $L_1' \parallels L_1$,且 $\overleftrightarrow{A'R} \parallels \overleftrightarrow{PQ},A'R = AR = 5$。题目所求等同于「点 $A'$ 经 $L_1$ 上一点再到达 $L_2$ 上一点的最短路径」,亦即  
        \[
        A'Q = \sqrt{(3+5)^2 + 14^2} = 2\sqrt{65}
        \]
    \end{solution}
    
    \question 空间中有一直圆锥,已知其高 $AH$ 所在的直线方程为 
    \[
    \frac{x-4}{-1} = \frac{y+4}{2} = \frac{z+14}{6},
    \] 
    以及直圆锥侧面上两点 $P(3,0,5),Q(11,-9,-18)$,求直圆锥顶点 $A$ 坐标。
    \begin{figure}[H]
        \centering        
        \includegraphics[width=0.25\textwidth]{images/image84.png}
    \end{figure}
    \begin{solution}
        设 $A$ 在直线上,令 
        \[
        A(-t+4, 2t-4, 6t-14), \quad t\in \mathbb{R}
        \]
        则
        \[
        \overrightarrow{AP} = (t-1, -2t+4, -6t+19), \quad \overrightarrow{AQ} = (t+7, -2t-5, -6t-4),
        \] 
        直线$AH$方向向量 
        \[
        \vec{h} = (-1,2,6)
        \]
        由于 $\angle PAH = \angle QAH$,有
        \[
        \frac{\overrightarrow{AP}\cdot \vec{h}}{|\overrightarrow{AP}||\vec{h}|} = \frac{\overrightarrow{AQ}\cdot \vec{h}}{|\overrightarrow{AQ}||\vec{h}|}.
        \]
        即
        \[
        \frac{(-41t+123)^2}{(t-1)^2 + (2t-4)^2 + (6t-19)^2} = \frac{(-41t-41)^2}{(t+7)^2 + (2t+5)^2 + (6t+4)^2}
        \]
        解得
        \[
        t=6 \ \text{或} \ t=\frac{9}{5}
        \]
        其中当$t=\dfrac{9}{5}$,
        \[
        \vec p\cdot \vec h> 0 ,\quad \vec q\cdot \vec h< 0
        \]
        不合题意,故取 $t=6$,则
        \[
        A = (-6+4, 12-4, 36-14) = (-2,8,22)
        \]
    \end{solution}

    \question 有四个平行平面:
    \[
    E_1: 3x + 4y + 5z = 0, \; E_2: 3x + 4y + 5z = 1, \; E_3: 3x + 4y + 5z = 2, \; E_4: 3x + 4y + 5z = 3,
    \]
    若一个正四面体的四顶点$A, B, C, D$分别在$E_1, E_2, E_3, E_4$,求此正四面体与$E_2$相交的截面积。
    \ifprintanswers
    \begin{figure}[H]
        \centering
        \includegraphics[width=0.5\linewidth]{images/image64.png}
    \end{figure}
    \fi
    \begin{solution}
        取
        \[
        A = (1, 0, 1), \quad B = (0, 1, 1), \quad C = (1, 1, 0), \quad O = (0, 0, 0),
        \]
        则四面体$OABC$为正四面体,边长为$\sqrt{2}$,再取
        \[
        D = AO \ \text{中点} = \left(\frac{1}{2}, 0, \frac{1}{2}\right), \quad
        E = \frac{2A + B}{3} = \left(\frac{2}{3}, \frac{1}{3}, 1\right),
        \]
        \[
        F = \frac{2B + A}{3} = \left(\frac{1}{3}, \frac{2}{3}, 1\right), \quad
        G = BC \ \text{中点} = \left(\frac{1}{2}, 1, \frac{1}{2}\right).
        \]
        则$\triangle CDE \parallels \triangle OFG$,且$E_2 = \triangle CDE, E_3 = \triangle OFG$,且
        \[
        \overrightarrow{CD} = \left(-\frac{1}{2}, -1, \frac{1}{2}\right), \quad \overrightarrow{CE} = \left(-\frac{1}{3}, -\frac{2}{3}, 1\right).
        \]
        $\triangle CDE$的法向量为
        \[
        \vec{n} = \overrightarrow{CD} \times \overrightarrow{CE} = \left(-\frac{2}{3}, \frac{1}{3}, 0\right).
        \]
        故面积为
        \[
        [\triangle CDE] = \frac{1}{2}|\vec{n}| = \frac{1}{2} \sqrt{\left(-\frac{2}{3}\right)^2 + \left(\frac{1}{3}\right)^2 + 0^2} = \frac{\sqrt{5}}{6}.
        \]
        又
        \[
        \overrightarrow{AB} = (-1, 1, 0), \quad |\overrightarrow{AB}| = \sqrt{2}, \quad |\vec{n}| = \frac{\sqrt{5}}{3}.
        \]
        于是
        \[
        \cos \theta = \frac{\overrightarrow{AB} \cdot \vec{n}}{|\overrightarrow{AB}||\vec{n}|} = \frac{3}{\sqrt{10}}.
        \]
        而平面间距离
        \[
        d(E_1, E_4) = \frac{3}{5 \sqrt{2}}.
        \]
        原四面体的棱长为
        \[
        \frac{d(E_1, E_4)}{\cos \theta} = \frac{1}{\sqrt{5}}
        \]
        截面面积为
        \[
        [\triangle CDE] \cdot \frac{1}{2} \left(\frac{1}{\sqrt{5}}\right)^2 = \frac{\sqrt{5}}{60}
        \]
    \end{solution}

    \question 空间中三点$P,Q,R$分别在直线
    \[
    L_1: \frac{x-3}{1} = \frac{y-6}{2} = \frac{z+1}{-2}, \;
    L_2: \frac{x-2}{-2} = \frac{y-7}{2} = \frac{z-4}{1}, \;
    L_3: \frac{x-1}{2} = \frac{y-5}{1} = \frac{z-6}{2}
    \]
    上,求$PQ+PR$的最小值。
    \begin{solution}
        直线方向向量为
        \[
        \vec{u}_1 = (1,2,-2), \quad \vec{u}_2 = (-2,2,1), \quad \vec{u}_3 = (2,1,2).
        \]
        发现
        \[
        \vec{u}_1 \times \vec{u}_2 = (6,3,6) \parallels \vec{u}_3,
        \]
        说明三条直线两两垂直且互为歪斜线。已知点
        \[
        P(3,6,-1) \in L_1, \quad Q(2,7,4) \in L_2, \quad R(1,5,6) \in L_3,
        \]
        则向量
        \[
        \overrightarrow{PQ} = (-1,1,5), \quad \overrightarrow{QR} = (-1,-2,2), \quad \overrightarrow{RP} = (2,1,-7).
        \]
        满足
        \[
        d(L_1,L_2) = \frac{|\overrightarrow{PQ} \cdot \vec{u}_3|}{|\vec{u}_3|} = 3, \quad
        d(L_2,L_3) = \frac{|\overrightarrow{QR} \cdot \vec{u}_1|}{|\vec{u}_1|} = 3, \quad
        d(L_3,L_1) = \frac{|\overrightarrow{RP} \cdot \vec{u}_2|}{|\vec{u}_2|} = 3.
        \]
        因此问题等价于在棱长为3的立方体中考虑三条互相垂直的线段:
        \[
        L_1 = \{(a,0,3) \mid a \in \mathbb{R}\}, \quad
        L_2 = \{(3,b,0) \mid b \in \mathbb{R}\}, \quad
        L_3 = \{(0,3,c) \mid c \in \mathbb{R}\}.
        \]
        设点$A =(3,b,0),B =(0,3,c),P =(a,0,3)$,则
        \[
        PA+PB = \sqrt{(a-3)^2 + b^2 + 9} + \sqrt{a^2 + 9 + (c-3)^2}.
        \]
        在
        \[
        a = \frac{3}{2}, \quad b=0, \quad c=3
        \]
        有最小值
        \[
        \sqrt{\frac{45}{4}} + \sqrt{\frac{45}{4}} = 3\sqrt{5}
        \]
    \end{solution}

    \question 空间中一定点 $A(2,6,-3)$,一平面 $E: x+2y+2z+1=0$,已知平面 $E$ 上有一圆 $C$,圆心为 $Q(-3,2,-1)$,半径为 $2$。若动点 $P$ 在圆 $C$ 上,试求 $AP$ 的最大值与最小值。
    \begin{solution}
        过点 $A(2,6,-3)$ 且方向向量为 $(1,2,2)$ 的直线为
        \[
        L: \frac{x-2}{1} = \frac{y-6}{2} = \frac{z+3}{2}
        \]
        令 $t$ 为参数,则该直线上的点可表示为
        \[
        A'(x,y,z) = (t+2,\,2t+6,\,2t-3)
        \]
        代入平面 $E$ 的方程求交点:
        \[
        (t+2) + 2(2t+6) + 2(2t-3) + 1 = 0\Rightarrow t = -1
        \]
        得投影点 $A' = (1,4,-5)$,而
        \[
        A'Q = \sqrt{(-3-1)^2 + (2-4)^2 + (-1+5)^2}  = 6 > 2 = r
        \Rightarrow A' \text{ 在圆 } C \text{ 外}
        \]
        设直线 $A'Q$ 与圆交于两点 $B$（较近）,$C$（较远）,则
        \[
        AA' = \sqrt{(2-1)^2 + (6-4)^2 + (-3+5)^2} = \sqrt{1 + 4 + 4} = 3 
        \]
        \[
        A'B = A'Q - r = 6 - 2 = 4 \Rightarrow AB^2 = 3^2 + 4^2 = 25 \Rightarrow AB = 5 
        \]
        \[
        A'C = A'Q + r = 6 + 2 = 8 \Rightarrow AC^2 = 3^2 + 8^2 = 73 \Rightarrow AC = \sqrt{73}
        \]
    \end{solution}

    \question 在正立方体 $ABCD-EFGH$ 中,有两质点分别从顶点 $A,C$ 同时以等速直线运动到顶点 $B,D$,均在 $1$ 秒后到达。求两质点之间的最小距离。  
    \begin{figure}[H]
        \centering
        \includegraphics[width=0.2\textwidth]{images/image10.png}
    \end{figure}
    \begin{solution}
        建立空间坐标系
        \[
        A(1,0,0),B(1,1,0),B(1,1,0),C(0,1,0),D(0,1,1),
        \]
        则两质点在$AB$及$CD$的坐标分别为$(1,t,0),(0,1,t)$,两质点间距离
        \[
        \sqrt{(1-0)^2+(t-1)^2+(0-t)^2}=\sqrt{2t^2-2t+2}=\sqrt{2\left(t-\frac12\right)^2+\frac32}
        \]
        在$t=\dfrac{1}{2}$有最小值$\dfrac{\sqrt6}{2}$。
    \end{solution}

    \question 长方体 $ABCD-EFGH$ 中,已知直线 $AC$ 的方程为 
    \[
    \frac{x-3}{-2}=\frac{y+1}{2}=\frac{z+7}{1},
    \]
    直线 $HF$ 的方程为 
    \[
    \frac{x}{1}=\frac{y}{4}=\frac{z}{-3},
    \]
    且 $A(3,-1,-7)$,求矩形 $ABCD$ 的面积。
    \ifprintanswers
    \begin{figure}[H]
        \centering        
        \includegraphics[width=0.5\textwidth]{images/image81.png}
    \end{figure}
    \fi
    \begin{solution}
        记
        \[
        L_1=\overleftrightarrow{AC}:\ \frac{x-3}{-2}=\frac{y+1}{2}=\frac{z+7}{1},\quad
        L_2=\overleftrightarrow{HF}:\ \frac{x}{1}=\frac{y}{4}=\frac{z}{-3}.
        \]
        方向向量分别为
        \[
        \vec u=(-2,2,1),\quad \vec v=(1,4,-3),
        \]
        $L_1,L_2$所成平面的法向量为
        \[
        \vec n=\vec u\times \vec v=(-10,-5,-10)
        \]
        取参数点
        \[
        P\in L_1:\ P(-2t+3,\,2t-1,\,t-7),\quad
        Q\in L_2:\ Q(s,\,4s,\,-3s),
        \]
        令 $\overrightarrow{PQ}\parallels \vec n$,解得
        \[
        \frac{s+2t-3}{2}=\frac{4s-2t+1}{1}=\frac{-3s-t+7}{2}\ \Rightarrow\ s=1,\ t=2,
        \]
        从而
        \[
        P(-1,3,-5),\quad Q(1,4,-3),\quad PA=PB=PC=6,
        \]
        设 $L_1$ 与 $L_2$ 的夹角为 $\theta$,则
        \[
        \cos\theta=\frac{\vec u\cdot \vec v}{\lVert\vec u\rVert\,\lVert\vec v\rVert}=\frac{1}{\sqrt{26}}.
        \]
        由余弦定理,
        \[
        AB^2=6^2+6^2-2\cdot6\cdot6\cdot\frac{1}{\sqrt{26}},\quad
        AD^2=6^2+6^2-2\cdot6\cdot6\cdot\left(-\frac{1}{\sqrt{26}}\right)
        \]
        于是矩形 $ABCD$面积为
        \[
        [ABCD]=AB\cdot AD
        =\sqrt{\frac{72(\sqrt{26}-1)}{\sqrt{26}}\cdot \frac{72(\sqrt{26}+1)}{\sqrt{26}}}
        =\frac{180\sqrt{26}}{13}
        \]
    \end{solution}

    \question 已知四面体 $ABCD$顶点坐标分别为 $A(0,0,0),B(1,0,0),C(0,1,0),D(1,1,1)$,求此四面体内切球球心的 $x$ 坐标。
    \begin{solution}
        四面体各面所在平面为:
        \[
        \begin{cases}
        E_1 = \triangle ABC: z=0 \\
        E_2 = \triangle ACD: x - z = 0 \\
        E_3 = \triangle ABD: -y + z = 0 \\
        E_4 = \triangle BCD: x + y - z = 1
        \end{cases}
        \]
        设球心为 $P(a,b,c)$,球心到四面的距离相等,即
        \[
        \begin{cases}
        d(P,E_1) = c \\[4pt]
        d(P,E_2) = \dfrac{|a - c|}{\sqrt{2}} \\[6pt]
        d(P,E_3) = \dfrac{| - b + c|}{\sqrt{2}} \\[6pt]
        d(P,E_4) = \dfrac{|a + b - c - 1|}{\sqrt{3}}\\[6pt]
        \end{cases}
        \]
        由于 $\triangle OAB$ 为等腰直角三角形,故 $a = b$。假设 $a < c$,则
        \[
        d(P,E_2) = \frac{c - a}{\sqrt{2}} = c \Rightarrow a = (1 - \sqrt{2}) c < 0,
        \]
        不合题意,故 $a \ge c$,得
        \[
        a = (\sqrt{2} + 1) c
        \]
        考虑
        \[
        d(P,E_4) = \frac{|(2 \sqrt{2} + 1) c - 1|}{\sqrt{3}} = c
        \]
        假设 $c \ge \dfrac{1}{2 \sqrt{2} + 1}$,则
        \[
        c = \frac{1}{2 \sqrt{2} - \sqrt{3} + 1} \Rightarrow a = \frac{\sqrt{2} + 1}{2 \sqrt{2} - \sqrt{3} + 1} > 1,
        \]
        球心在四面体外,不合题意。因此
        \[
        c \le \frac{1}{2 \sqrt{2} + 1} \Rightarrow c = \frac{1}{2 \sqrt{2} + \sqrt{3} + 1},
        \]
        故内切球球心的 $x$ 坐标为
        \[
        a=\frac{\sqrt{2} + 1}{2 \sqrt{2} + \sqrt{3} + 1}
        \]
    \end{solution}

    \question 某房间内设有一盏聚光灯,其照射的光线为直圆锥状。为描述其空间几何关系,建立如图所示的空间坐标系。已知光源 $S(2\sqrt{6},3\sqrt{6},12)$,分别沿对称轴 $L$ 的方向向量 $\left(\dfrac{\sqrt{2}}{2},-\dfrac{\sqrt{2}}{2},-\sqrt{3}\right)$ 以及其中一条母线 $M$ 的方向向量 $(0,0,-1)$ 向地面 $z=0$ 照射。试问:此聚光灯照射在墙面（即平面 $y=0$）上的光线边缘,为哪一种圆锥曲线图形的一部分？该圆锥曲线在墙面上的顶点坐标为何？
    \ifprintanswers
    \begin{figure}[H]
        \centering        
        \includegraphics[width=0.4\textwidth]{images/image229.png}
    \end{figure}
    \fi
    \begin{solution}
        直线 $L$ 的方向向量 $\vec \ell = \left(\dfrac{\sqrt{2}}{2},-\dfrac{\sqrt{2}}{2},-\sqrt{3}\right)$,母线 $M$ 的方向向量 $\vec m = (0,0,-1)$。设两直线夹角为 $\theta$,则  
        \[
        \cos \theta = \frac{\vec \ell \cdot \vec m}{|\vec \ell||\vec m|} = \frac{\sqrt 3}{2} \Rightarrow \theta = 30^\circ \text{ 或 } 150^\circ
        \]
        假设点 $P(x,y,z)$ 在光线边缘上,则  
        \[
        \overrightarrow{SP} = (x-2\sqrt 6,\,y-3\sqrt 6,\,z-12)
        \]  
        与 $\vec \ell$ 的夹角为 $30^\circ$,由
        $\overrightarrow{SP} \cdot \vec \ell = |\overrightarrow{SP}|\,|\vec \ell| \cos 30^\circ$得到  
        \[
        \frac{\sqrt 2}{2}(x-2\sqrt 6) - \frac{\sqrt 2}{2}(y-3\sqrt 6) - \sqrt 3 (z-12) = \sqrt 3 \sqrt{(x-2\sqrt 6)^2 + (y-3\sqrt 6)^2 + (z-12)^2}
        \]  
        在墙面 $y=0$ 代入上式,化简得  
        \[
        5x^2 + 2\sqrt 6 x z - 50\sqrt 6 x + 12 z + 318 = 0
        \]  
        判别式为 
        \[
        (2\sqrt 6)^2 - 0 = 24 > 0
        \]  
        因此图形为双曲线。对称轴 $L$ 的参数方程为  
        \[
        \frac{x-2\sqrt 6}{\frac{\sqrt{2}}{2}} = \frac{y-3\sqrt 6}{-\frac{\sqrt{2}}{2}} = \frac{z-12}{-\sqrt 3}
        \]  
        在 $y=0$ 的投影为  
        \[
        \frac{x-2\sqrt 6}{\frac{\sqrt{2}}{2}} = \frac{z-12}{-\sqrt 3} \Rightarrow x = 2\sqrt 6 - \frac{z-12}{\sqrt 6}
        \]  
        代入双曲线方程可得  
        \[
        x = -\sqrt 6 \pm 18 \sqrt{\frac{2}{7}}, \quad z = 30 \pm 36 \sqrt{\frac{3}{7}}
        \]  
        由图形可知,选择  
        \[
        z = 30 - 36 \sqrt{\frac{3}{7}}
        \]  
        因此顶点坐标为  
        \[
        \left(-\sqrt 6 + 18 \sqrt{\frac{2}{7}},\,0,\,30 - 36 \sqrt{\frac{3}{7}}\right)
        \]
    \end{solution}

\end{questions}

\pagecolor{PageColor}
\
\vfil
\hfil  {\fontsize{50pt}{36pt}\selectfont{微积分}} \hfil
\vfil
\begin{tikzpicture}[remember picture,overlay,every node/.style={inner sep=0pt}]
        \node [shift={(1cm,-1cm)},brown,scale=2,anchor=north west] (CNW)
        at (current page.north west) {\pgfornament[height=1cm,width=1cm]{61}};
        \node [shift={(-1cm,-1cm)},brown,scale=2,anchor=north east] (CNE)
        at (current page.north east) {\pgfornament[height=1cm,width=1cm,symmetry=v]{61}};
        \node [shift={(1cm,1cm)},brown,scale=2,anchor=south west] (CSW)
        at (current page.south west) {\pgfornament[height=1cm,width=1cm,symmetry=h]{61}};
        \node [shift={(-1cm,1cm)},brown,scale=2,anchor=south east] (CSE)
        at (current page.south east) {\pgfornament[height=1cm,width=1cm,symmetry=c]{61}};
        \pgfornamentline[color=brown]{current page.north west}{current page.north east}{2}{87}
        \pgfornamentline{current page.south west}{current page.south east}{2}{87}
        \pgfornamentline{current page.north west}{current page.south west}{3}{87}
        \pgfornamentline{current page.north east}{current page.south east}{3}{87}
        \end{tikzpicture}%
\thispagestyle{empty}
\pagebreak
\begin{center}
  {\fontsize{30pt}{26pt}\selectfont
    \hypertarget{极限}{极限} \label{极限}
  }
\end{center}
\separator
\vspace{1pt}
\nopagecolor
\begin{questions}
    \question 
    \[
    \lim_{x \to 0} \floor{x}
    \]
    \begin{solution}
        由于
        \[
        \lim_{x \to 0^{+}} \floor{x}=0, \quad \lim_{x \to 0^{-}} \floor{x}=-1
        \]
        极限$\displaystyle \lim_{x \to 0} \floor{x}$ 不存在。
    \end{solution}

    \question 
    \[
    \lim_{x \to 3} (\ceil{x}-\floor{x})
    \]
    \begin{solution}
        由于
        \[
        \lim_{x \to 3^{+}} (\ceil{x}-\floor{x}) = \lim_{x \to 3^{+}} \ceil{x} - \lim_{x \to 3^{+}} \floor{x}=4-3=1 
        \]
        且
        \[
        \lim_{x \to 3^{-}} (\ceil{x}-\floor{x}) = \lim_{x \to 3^{-}} \ceil{x} - \lim_{x \to 3^{-}} \floor{x}=3-2=1 
        \]
        故
        \[
        \lim_{x \to 3} (\ceil{x}-\floor{x}) =1
        \]
    \end{solution}

    \question 
    \[
    \lim_{x \to 0} \;\floor*{\frac{\sin x }{x}}
    \]
    \begin{solution}
        由于
        \[
        \lim_{x \to 0^{+}} \floor*{\frac{\sin x }{x}} = \lim_{x \to 0^{-}} \floor*{\frac{\sin x }{x}}=0
        \]
        故
        \[
        \lim_{x \to 0} \;\floor*{\frac{\sin x }{x}}=0
        \]
    \end{solution}

    \question 
    \[
    \lim_{x \to 0} \frac{\sin(\sin x)}{x}
    \]
    \begin{solution}
        有
        \begin{align*}
        \lim_{x \to 0} \frac{\sin(\sin x)}{x} 
        &=\,\lim_{x \to 0} \frac{\sin(\sin x)}{\sin x} \cdot \frac{\sin x}{x} \\[2mm]
        &=\,1 \cdot 1 \\
        &= 1
        \end{align*}
    \end{solution}
    
    \question 
    \[
    \lim_{\theta \to 0} \frac{\theta - \theta \cos \theta}{\sin \theta \tan \theta}
    \]
    \begin{solution}
        有
        \begin{align*}
        \,\lim_{\theta \to 0} \frac{\theta - \theta \cos \theta}{\sin \theta \tan \theta} 
        &=\,\lim_{\theta \to 0} \frac{1 - \cos \theta}{\frac{\sin \theta}{\theta}} \cdot \frac{1}{\tan \theta} \\
        &=\,\lim_{\theta \to 0} \frac{1 - \cos \theta}{\theta^2}
        \cdot \frac{1}{\frac{\sin \theta}{\theta}} 
        \cdot \frac{\theta}{\tan \theta}
        \cdot \theta \\
        &=\,\frac{1}{2} \cdot \frac{1}{1} \cdot \frac{1}{1} \cdot 0 \\
        &= 0
        \end{align*}
    \end{solution}  

    \question 
    \[
    \lim_{\theta \to 0} \frac{1 - \cos\left(\frac{1 - \cos \theta}{2}\right)}{\theta^4}
    \]
    \begin{solution}
        有
        \begin{align*}
        \lim_{\theta \to 0} \frac{1 - \cos\left(\frac{1 - \cos \theta}{2}\right)}{\theta^4} 
        &=\,\lim_{\theta \to 0} \left[
        \frac{1 - \cos\left(\frac{1 - \cos \theta}{2}\right)}{\left( \frac{1 - \cos \theta}{2} \right)^2}
        \cdot \left( \frac{1 - \cos \theta}{\theta^2} \right)^2 \cdot \frac{1}{4}
        \right] \\
        &=\,\frac{1}{2} \cdot \left(\frac{1}{2}\right)^2 \cdot \frac{1}{4} = \frac{1}{32}
        \end{align*}
        其中 \[
        \lim_{x \to 0} \frac{1+\cos x}{x^2}= \frac14\lim_{x \to 0}\frac{2\sin^2 \frac{x}{2} }{(\frac{x}{2})^2}=\frac12 \lim_{x \to 0} \left(\frac{\sin \frac{x}{2}}{\frac{x}{2}}\right)^2=\frac12
        \]
    \end{solution}
        
    \question 
    \[
    \lim_{x \to 0} \frac{\cos x - 1}{\sin^{-1} x \tan x}
    \]
    \begin{solution}
        有
        \begin{align*}
        \lim_{x \to 0} \frac{\cos x - 1}{\sin^{-1} x \tan x} 
        &=\,\lim_{x \to 0} \frac{\cos x - 1}{\sin^{-1} x \cdot \frac{\sin x}{\cos x}} \\
        &=\,\lim_{x \to 0} \left(-\frac{1 - \cos x}{x^2}\right) \cdot \frac{x}{\sin^{-1} x} \cdot \frac{x}{\sin x} \cdot \cos x \\
        &=\, -\frac{1}{2} \cdot 1 \cdot 1 \cdot 1 \\ 
        &= -\frac{1}{2}
        \end{align*}
    \end{solution}

    \question 求极限 
    \[
    \lim_{x \to 1} \left[ \cos\left(\frac{\pi}{2}x\right) + x\cos\left(\frac{\pi}{2}x\right) + x^2\cos\left(\frac{\pi}{2}x\right) + \dots \right]
    \]
    \begin{solution}
        注意到
        \begin{align*}
        &\lim_{x \to 1} \left[ \cos\left(\frac{\pi}{2}x\right) + x\cos\left(\frac{\pi}{2}x\right) + x^2\cos\left(\frac{\pi}{2}x\right) + \dots \right] \\
        &= \lim_{x \to 1}  \frac{\cos\left(\frac{\pi}{2}x\right)}{1-x} \\
        &\overset{H}{=} \lim_{x \to 1} \frac{-\frac{\pi}{2} \sin\left(\frac{\pi}{2}x\right)}{-1} \\
        &= \frac{\pi}{2} \sin\frac{\pi}{2} \\
        &= \frac{\pi}{2}
        \end{align*}
    \end{solution}

    \question 计算极限 
    \[
    \lim_{x \to 0} \left( \sqrt{\frac{1}{x(x-1)} + \frac{1}{4x^2}} - \frac{1}{2x} \right)
    \]
    \begin{solution}
        有理化得
        \begin{align*}
        \lim_{x \to 0} \left( \sqrt{\frac{1}{x(x-1)} + \frac{1}{4x^2}} - \frac{1}{2x} \right) 
        &= \lim_{x \to 0} \frac{\left( \frac{1}{x(x-1)} + \frac{1}{4x^2} \right) - \frac{1}{4x^2}}{\sqrt{\frac{1}{x(x-1)} + \frac{1}{4x^2}} + \frac{1}{2x}} \\
        &= \lim_{x \to 0} \frac{\frac{1}{x-1}}{\sqrt{\frac{x}{x-1} + \frac{1}{4}} + \frac{1}{2}} \\
        &= \frac{-1}{\frac{1}{2} + \frac{1}{2}} \\
        &= -1
        \end{align*}
    \end{solution}

    \question 试求
    \[
    \lim_{x\to\infty}\left(\sqrt[5]{x^{5}+3x^{4}+4x^{3}+3x}-\sqrt[3]{x^{3}+3x^{2}+4x+1}\right).
    \]
    \begin{solution}
        由洛必达法则,
        \begin{align*}
        &\lim_{x\to \infty} \left(\sqrt[5]{x^5+3x^4+4x^3+3x} - \sqrt[3]{x^3+3x^2+4x+1} \right) \\
        &= \lim_{x\to \infty} \frac{\sqrt[5]{1 + \frac{3}{x} + \frac{4}{x^2} + \frac{3}{x^4}} - \sqrt[3]{1 + \frac{3}{x} + \frac{4}{x^2} + \frac{1}{x^3}}}{\frac{1}{x}} \\
        &\overset{H}{=} -\lim_{x\to \infty} x^2\left(
        \frac{1}{5} \left(1 + \frac{3}{x} + \frac{4}{x^2} + \frac{3}{x^4}\right)^{-\frac{4}{5}} \left(-\frac{3}{x^2} - \frac{8}{x^3} - \frac{12}{x^5}\right) \right. \\
        & \quad \left. - \frac{1}{3} \left(1 + \frac{3}{x} + \frac{4}{x^2} + \frac{1}{x^3}\right)^{-\frac{2}{3}} \left(-\frac{3}{x^2} - \frac{8}{x^3} - \frac{3}{x^4}\right) 
        \right)\\
        &= \frac{3}{5} - 1 = -\frac{2}{5}
        \end{align*}
    \end{solution}

    \question 已知等差数列 $\{a_n\},\{b_n\}$ 公差均不为 $0$, 且$\displaystyle \lim_{n\to\infty}\frac{a_n}{b_n}=5$,求
    \[
    \lim_{n\to\infty}\frac{a_1+a_2+\cdots+a_n}{n\cdot b_{2n}}.
    \]
    \begin{solution}
        设 $\{a_n\},\{b_n\}$ 公差为 $d_1,d_2$,则
        \[
        a_n=a_1+(n-1)d_1,\quad b_n=b_1+(n-1)d_2.
        \]
        由已知极限得
        \[
        \lim_{n\to\infty}\frac{a_n}{b_n}=\lim_{n\to\infty}\frac{a_1+(n-1)d_1}{b_1+(n-1)d_2}=\frac{d_1}{d_2}=5.
        \]
        所以 $d_1=5d_2$,因此所求极限为
        \[
        \lim_{n\to\infty}\frac{\frac{n}{2}\bigl(2a_1+(n-1)d_1\bigr)}{n\bigl(b_1+(2n-1)d_2\bigr)}
        = \lim_{n\to\infty}\frac{2a_1+(n-1)d_1}{2\bigl(b_1+(2n-1)d_2\bigr)}=\frac{d_1}{4d_2}=\frac{5}{4}
        \]
    \end{solution}

    \question 设 $a_1+a_2+\cdots+a_n=n^3-2n$,试求
    \[
    \lim_{n\to\infty}\frac{\sqrt[3]{a_3+a_6+\cdots+a_{3n}}-\sqrt[3]{a_2+a_4+\cdots+a_{2n}}}{n}
    \]
    \begin{solution}
        由 
        \[
        S(n)=\sum_{k=1}^n a_k=n^3-2n
        \] 
        得
        \[
        a_n=S(n)-S(n-1)=n^3-2n-\left[(n-1)^3-2(n-1)\right]=3n^2-3n-1
        \]
        于是
        \[
        a_{3k}=27k^2-9k-1,\quad a_{2k}=12k^2-6k-1,
        \]
        且
        \[
        \sum_{k=1}^n a_{3k}= \frac{27}{6}n(n+1)(2n+1)-\frac{9}{2}n(n+1)-n=9 n^3 + 9 n^2 - n,
        \]
        \[
        \sum_{k=1}^n a_{2k}= 2n(n+1)(2n+1)-3n(n+1)-n=4 n^3 + 3 n^2 - 2 n
        \]
        故极限为
        \[
        \lim_{n\to\infty}\left(\sqrt[3]{9+\frac{9}{n}-\frac{1}{n^2}}-\sqrt[3]{4+\frac{3}{n}-\frac{2}{n^2}}\right)
        =\sqrt[3]{9}-\sqrt[3]{4}
        \]
    \end{solution}

    \question 
    \[
    \lim_{m\to\infty} \lim_{n\to\infty} \frac{\sqrt[n]{1+3^{2n}}+\sqrt[n]{3^{2n}+5^{2n}}+\sqrt[n]{5^{2n}+7^{2n}}+\cdots+\sqrt[n]{(2m-1)^{2n}+(2m+1)^{2n}}}{m^{3}}
    \]
    \begin{solution}
        首先有
        \begin{align*}
        &\lim_{n\to\infty} \left( \sqrt[n]{1+3^{2n}} +\sqrt[n]{3^{2n}+5^{2n}} + \cdots +\sqrt[n]{(2m-1)^{2n}+(2m+1)^{2n}} \right)\\
        &=\lim_{n\to\infty} \left( 3^2\sqrt[n]{\left(\frac13\right)^{2n}+1} + 5^2\sqrt[n]{\left(\frac35\right)^{2n}+1} + \cdots + (2m+1)^2\sqrt[n]{\left(\frac{2m-1}{2m+1}\right)^{2n}+1} \right)\\
        &= \sum_{k=1}^m (2k+1)^2 \\
        &= \sum_{k=1}^m (4k^2+4k+1) \\
        &=\frac{2}{3}m(m+1)(2m+1) + 2m(m+1) + m
        \end{align*}
        因此,原式
        \[
        \lim_{m\to\infty} \frac{\frac{2}{3}m(m+1)(2m+1) + 2m(m+1) + m}{m^3} = \frac{4}{3}
        \]
    \end{solution}

    \question 
    \[
    \lim_{x \to 0} \frac{\int_{\cos x}^1 e^{-y^2} \, dy}{x \sin x}
    \]
    \begin{solution}
        由洛必达法则,
        \begin{align*}
        \lim_{x \to 0} \frac{\int_{\cos x}^1 e^{-y^2} \, dy}{x \sin x}
        &\overset{H}{=} \lim_{x \to 0} \frac{\sin x \cdot e^{-(\cos x)^2}}{\sin x + x \cos x} \\
        &\overset{H}{=} \lim_{x \to 0} \frac{\cos x(1 + 2\sin^2 x) \cdot e^{-(\cos x)^2}}{2\cos x - x\sin x} \\
        &= \frac{1}{2e}
        \end{align*}
    \end{solution}

    \question 
    \[
    \lim_{x \to 1} \frac{(1 - \sqrt{x})(1 - \sqrt[3]{x})}{1 + \cos \pi x}
    \]
    \begin{solution}
        消去根号,令 $x = t^6$,当 $x \to 1$ 时,$t \to 1$,由洛必达法则,
        \begin{align*}
        \lim_{x \to 1} \frac{(1 - \sqrt{x})(1 - \sqrt[3]{x})}{1 + \cos \pi x} 
        &= \lim_{t \to 1} \frac{(1 - t^3)(1 - t^2)}{1 + \cos \pi t^6} \\
        &= \lim_{t \to 1} \frac{1 - t^2 - t^3 + t^5}{1 + \cos \pi t^6} \\
        &\overset{H}{=} \lim_{t \to 1} \frac{-2t - 3t^2 + 5t^4}{-6\pi t^5 \sin \pi t^6} \\
        &\overset{H}{=} \lim_{t \to 1} \frac{-2 - 6t + 20t^3}{-30\pi t^4 \sin \pi t^6 - 6\pi t^5 \cdot \cos \pi t^6 \cdot 6\pi t^5} \\
        &= \frac{-2 - 6(1) + 20(1)^3}{-6\pi \cdot (-1) \cdot 6\pi} \\
        &= \frac{1}{3\pi^2}
        \end{align*}
    \end{solution}

    \question 
    \[
    \lim_{x \to 0} \frac{\left(\int_0^x t \cos t^2 \, dt \right)^2}{\int_0^x \sin t^2 \, dt}
    \]
    \begin{solution}
        首先发现到
        \[
        \int_0^x t \cos t^2 \, dt = \left[ \frac{1}{2} \sin t^2 \right]_0^x = \frac{1}{2} \sin x^2
        \]
        将积分结果代入原极限式:
        \begin{align*}
        \lim_{x \to 0} \frac{\left(\int_0^x t \cos t^2 \, dt \right)^2}{\int_0^x \sin t^2 \, dt} 
        &= \lim_{x \to 0} \frac{\frac{1}{4} \sin^2 x^2}{\int_0^x \sin t^2 \, dt} \\
        &\overset{H}{=} \lim_{x \to 0} \frac{\frac{1}{4} \cdot 2 \sin x^2 \cdot \cos x^2 \cdot 2x}{\sin x^2} \\
        &= \lim_{x \to 0} (x \cos x^2) \\
        &= 0 \cdot 1 = 0
        \end{align*}
    \end{solution}

    \question 求实数 $a,b$,使得
    \[
    \lim_{x\to 0}\frac{\sqrt{ax+b}-2}{x}=1.
    \]
    \begin{solution}
        有理化得
        \[
        \lim_{x\to 0}\frac{\sqrt{ax+b}-2}{x}\cdot
        \frac{\sqrt{ax+b}+2}{\sqrt{ax+b}+2}
        =\lim_{x\to 0}\frac{ax+b-4}{x(\sqrt{ax+b}+2)}
        \]
        当 $x\to 0$ 时分母趋于 $0$,欲使极限存在,分子也必须趋于 $0$,因此
        \[
        a\cdot 0+b-4=0,
        \]
        从而 $b=4$,于是
        \[
        \lim_{x\to 0}\frac{a}{\sqrt{ax+4}+2}=\frac{a}{\sqrt{4}+2}=1,
        \]
        解得
        \[
        a=4
        \]
    \end{solution}

    \question 已知
    \[
    \lim_{x \to -1} \frac{a x^2 + 3x + b}{x + 1} = 2,\;\; a,b \in \mathbb{R}
    \]
    试求 \( (a, b) \)。
    \begin{solution}
        当$x\rightarrow-1,x+1\rightarrow0;$欲使该极限存在且等于$2$,分子$a x^2 + 3x + b$在 \(x\rightarrow-1\) 也应趋于$0$,即
        \[
        \lim_{x \to -1} (a x^2 + 3x + b)=0 \Rightarrow a(-1)^2 + 3(-1) + b = 0 \Rightarrow  a + b = 3 \tag{1}
        \]
        又原极限中分子分母处处可导,且极限
        \[
        \lim_{x \to -1} \frac{(a x^2 + 3x + b)'}{(x + 1)'}
        = \lim_{x \to -1} \frac{2a x + 3}{1} = -2a + 3 
        \]
        存在,则由洛必达法则,
        \[
        \lim_{x \to -1} \frac{a x^2 + 3x + b}{x + 1} \overset{H}{=}-2a + 3 = 2 \Rightarrow (a,b) = \left( \dfrac{1}{5}, \dfrac{14}{5} \right)
        \]
    \end{solution}

    \question 已知 $a,b$ 为正整数,设函数 
    \[
    f(x)=\lim_{n\rightarrow \infty}\frac{2x^{2n+1}+ax^{2}+bx-1}{2x^{2n}+3},
    \]
    若 $\forall x\in \mathbb{R},f(x)$ 为连续函数,求序对 $(a,b)$。
    \begin{solution}
        首先注意到
        \[
        f(x) = \lim_{n \rightarrow \infty} \frac{2x^{2n+1}+ ax^2+bx-1}{2x^{2n}+3} \\
        = \lim_{n \rightarrow \infty} \frac{x+ \frac{ax^2+bx-1}{2x^{2n}}}{1+\frac{3}{2x^{2n}}} \\
        = 
        \begin{cases}
        x, & |x| > 1,\\[1mm]
        \dfrac{ax^2+bx-1}{3}, & |x| < 1.
        \end{cases}
        \]
        由连续性,
        \[
        \lim_{x \rightarrow 1^+} f(x) = 1, \quad \lim_{x \rightarrow 1^-} f(x) = \frac{a+b-1}{3}, \quad f(1) = \lim_{n \rightarrow \infty} \frac{2+a+b-1}{2+3} = \frac{a+b+1}{5}
        \]
        可得
        \[
        \frac{a+b+1}{5} = 1 = \frac{a+b-1}{3} \quad \Rightarrow \quad a+b = 4 \tag{1}
        \]
        同理,
        \[
        \lim_{x \rightarrow -1^+} f(x) = \frac{a-b-1}{3}, \quad \lim_{x \rightarrow -1^-} f(x) = -1,\quad f(-1) = \frac{-2+a-b-1}{5} = \frac{a-b-3}{5}
        \]
        可得
        \[
        \frac{a-b-3}{5} = -1 = \frac{a-b-1}{3} \quad \Rightarrow \quad a-b = -2 \tag{2}
        \]
        由$(1),(2)$得 
        \[
        (a,b) = (1,3)
        \]
    \end{solution}
    
    \question 若
    \[
    \lim_{x \to 0} \frac{\sqrt{f(x) \sin x + 1} - 1}{e^{4x} - 1} = 2,
    \]
    求 
    \[
    \lim_{x \to 0} f(x)
    \]
    \begin{solution}
        当 $x \to 0$ 时,原式为不定型 $\dfrac{0}{0}$,可用洛必达法则:
        \[
        \lim_{x \to 0} \frac{\sqrt{f(x) \sin x + 1} - 1}{e^{4x} - 1}
        \overset{H}{=}
        \lim_{x \to 0} 
        \frac{\frac{1}{2\sqrt{f(x)\sin x + 1}} \cdot \left[ f'(x)\sin x + f(x)\cos x \right]}
        {4e^{4x}}.
        \]
        当 $x \to 0$,
        \[
        \sin x \to 0,\cos x \to 1,e^{4x} \to 1,\sqrt{f(x)\sin x + 1} \to 1
        \]
        上式化简为
        \[
        \lim_{x \to 0} \frac{f(x)}{8} = 2
        \Rightarrow \lim_{x \to 0} f(x) = 16
        \]
    \end{solution}

    \question 定义序列 $\{x_n\}_{n=2}^\infty$ 如下:
    \[
    (n+x_n)[\sqrt{2}-1] = \ln 2.
    \]
    求 $\lim_{n\to\infty} x_n$。
    \begin{solution}
        由方程可解得
        \[
        x_n = \frac{\ln 2}{2^{\frac{1}{n}}-1} - n,
        \]
        极限形式为 $\infty - \infty$。令 $u = \dfrac{1}{n}$,并使用洛必达法则两次,有
        \begin{align*}
        \lim_{n\to\infty} x_n &= \lim_{n\to\infty} \frac{\ln 2 - n \cdot 2^{\frac{1}{n}} + n}{2^{\frac{1}{n}} - 1} \\
        &= \lim_{u\to 0} \frac{u \ln 2 - 2^u + 1}{u 2^u - u} \\
        &\overset{H}{=} \lim_{u\to 0} \frac{\ln 2 - 2^u \ln 2}{2^u - 1 + u 2^u \ln 2} \\
        &\overset{H}{=} \lim_{u\to 0} \frac{-2^u (\ln 2)^2}{2^u \ln 2 + 2^u \ln 2 + u 2^u (\ln 2)^2} \\
        &= \frac{-(\ln 2)^2}{2 \ln 2} \\
        &= -\frac{1}{2} \ln 2
        \end{align*}
    \end{solution}

    \question 计算极限
    \[
    \lim_{x\to 0} \left[\frac{1}{x^3}\int_{0}^{x} \frac{t\ln(t+1)}{t^4+\frac{1}{6}} dt\right]
    \]
    \begin{solution}
        使用两次洛必达法则,
        \begin{align*}
        \lim_{x\to 0} \left[\frac{1}{x^3}\int_{0}^{x} \frac{t\ln(t+1)}{t^4+\frac{1}{6}} dt\right]
        &\overset{H}{=} \lim_{x\to 0} \frac{\frac{x\ln(x+1)}{x^4+\frac{1}{6}}}{3x^2} \\
        &= \lim_{x\to 0} \frac{\ln(x+1)}{3x^5+\frac{x}{2}}\\
        &\overset{H}{=} \lim_{x\to 0} \frac{\frac{1}{x+1}}{15x^4+\frac{1}{2}} \\
        &=\frac{1}{0+\frac{1}{2}} \\
        &=2
        \end{align*}
    \end{solution}

    \question 求极限
    \[
    \lim_{x \to 0^+} (\sin x)^{\frac{1}{\ln x}}
    \]
    \begin{solution}
        此极限属于 $0^0$ 型。利用对数恒等式将原式重写为以 $e$ 为底的形式:
        \[
        (\sin x)^{\frac{1}{\ln x}} = \exp\left( \frac{\ln(\sin x)}{\ln x} \right)
        \]
        由于 $e^x$ 是连续函数,只需计算指数部分的极限
        \[
        L = \lim_{x \to 0^+} \frac{\ln(\sin x)}{\ln x}
        \]
        属于 $\frac{\infty}{\infty}$ 型,由洛必达法则,
        \[
        L \overset{H}{=} \lim_{x \to 0^+} \frac{\frac{\cos x}{\sin x}}{\frac{1}{x}} =  \lim_{x \to 0^+} \frac{x}{\sin x} \cdot \cos x  = 1 \cdot 1 = 1
        \]
        因此原极限为
        \[
        e^1 = e
        \]
    \end{solution}

    \question 求极限
    \[
    \lim_{x \to \infty} (x + e^x)^{\frac{1}{x}}
    \]
    \begin{solution}
        此极限属于 $\infty^0$ 型,由
        \[
        (x + e^x)^{\frac{1}{x}} = \exp\left( \frac{\ln(x + e^x)}{x} \right)
        \]
        设
        \[
        L = \lim_{x \to \infty} \frac{\ln(x + e^x)}{x}
        \]
        此极限属于 $\frac{\infty}{\infty}$ 型,由洛必达法则,
        \[
        L \overset{H}{=} \lim_{x \to \infty} \frac{\frac{1 + e^x}{x + e^x}}{1} = \lim_{x \to \infty} \frac{1 + e^x}{x + e^x}
        \]
        再由洛必达法则,
        \[
        L \overset{H}{=} \lim_{x \to \infty} \frac{e^x}{1 + e^x} = \lim_{x \to \infty} \frac{1}{\frac{1}{e^x} + 1} = \frac{1}{0 + 1} = 1
        \]
        因此,原极限为
        \[
        e^1 = e
        \]
    \end{solution}

    \question 求
    \[
    \lim_{x\to\infty} \left[\sqrt{x^2+2x-1}-\sqrt{x^2-1}\right]^x
    \]
    \begin{solution}
        首先注意到
        \begin{align*}
        &\lim_{x\to\infty} \frac{(x^2+2x-1)-(x^2-1)}{\sqrt{x^2+2x-1}+\sqrt{x^2-1}} \\
        &= \lim_{x\to\infty} \frac{2}{\sqrt{1+\frac{2}{x}-\frac{1}{x^2}} + \sqrt{1-\frac{1}{x^2}}} \\
        &= \frac{2}{1+1} = 1
        \end{align*}
        故原极限为$1^\infty$型,于是
        \[
        \lim_{x\to\infty} \left[\sqrt{x^2+2x-1}-\sqrt{x^2-1}\right]^x = \exp\left[\lim_{x\to\infty} x \left( \sqrt{x^2+2x-1} - \sqrt{x^2-1} - 1 \right)\right]
        \]
        其中
        \begin{align*}
        \lim_{x\to\infty} x \left( \sqrt{x^2+2x-1} - \sqrt{x^2-1} - 1 \right)
        &= \lim_{x\to\infty} x \cdot \frac{(x^2+2x-1) - (\sqrt{x^2-1} + 1)^2}{\sqrt{x^2+2x-1} + \sqrt{x^2-1} + 1} \\
        &= \lim_{x\to\infty} x \cdot \frac{2x - 2\sqrt{x^2-1}}{\sqrt{x^2+2x-1} + \sqrt{x^2-1} + 1} \\
        &= \lim_{x\to\infty} x \cdot \frac{\frac{2}{x + \sqrt{x^2-1}} - 1}{\sqrt{x^2+2x-1} + \sqrt{x^2-1} + 1} \\
        &= \lim_{x\to\infty} \frac{\frac{2}{x + \sqrt{x^2-1}} - 1}{\sqrt{1+\frac{2}{x}-\frac{1}{x^2}} + \sqrt{1-\frac{1}{x^2}} + \frac{1}{x^2}} \\
        &= \frac{0-1}{1+1} \\
        &= -\frac{1}{2}
        \end{align*}
        故原极限为
        \[
        \lim_{x\to\infty} \left[\sqrt{x^2+2x-1}-\sqrt{x^2-1}\right]^x = \frac{1}{\sqrt{e}}
        \]
    \end{solution}

    \question 设
    \[
    \lim_{x \to 0} \frac{\sin 6x + x f(x)}{x^3} = 0,
    \]
    求
    \[
    \lim_{x \to 0} \frac{6 + f(x)}{x^2}
    \]
    \begin{solution}
        原极限为
        \[
        \lim_{x \to 0} \frac{\sin 6x + x f(x)}{x^3} = 0.
        \]
        利用幂级数展开
        \[
        \sin 6x = 6x - \frac{(6x)^3}{6} + o(x^3) = 6x - 36 x^3 + o(x^3),
        \]
        代入得
        \[
        \frac{6x - 36 x^3 + x f(x)}{x^3} = \frac{6x + x f(x)}{x^3} - 36 + o(1) = \frac{6 + f(x)}{x^2} - 36 + o(1).
        \]
        极限存在且为0,则
        \[
        \lim_{x \to 0} \left(\frac{6 + f(x)}{x^2} - 36\right) = 0 \implies \lim_{x \to 0} \frac{6 + f(x)}{x^2} = 36.
        \]
    \end{solution}

    \question 若 $f(x)$ 为满足
    \[
    \lim_{x\to 1} \frac{f(x)}{x-1} = 36, \quad 
    \lim_{x\to -1} \frac{f(x)}{x+1} = -36, \quad 
    \lim_{x\to 2} \frac{f(x)}{x-2} = 0, \quad 
    \lim_{x\to -2} \frac{f(x)}{x+2} = 0
    \]
    的最低次多项式,求 $f(3)$。
    \begin{solution}
        设
        \[
        f(x) = (ax+b)(x-1)(x+1)(x-2)^2(x+2)^2
        \]
        由
        \[
        \lim_{x\to 1} \frac{f(x)}{x-1} = 36, \quad \lim_{x\to -1} \frac{f(x)}{x+1} = -36
        \]
        得方程式
        \[
        \begin{cases}
        (a+b)\cdot 2 \cdot 1^2 \cdot 3^2 = 36 \\
        (-a+b)\cdot (-2) \cdot (-3)^2 \cdot 1^2 = -36 
        \end{cases} \Rightarrow a=0, b=2
        \]
        因此
        \[
        f(x) = 2(x^2-1)(x-2)^2(x+2)^2 \Rightarrow f(3)=400
        \]
    \end{solution}

    \question 若
    \[
    \lim_{x \to 1} \frac{\sqrt{x^4 + 3} - [A + B(x - 1) + C(x - 1)^2]}{(x - 1)^2} = 0,
    \]
    求常数 \( A, B, C \)。
    \begin{solution}
        设
        \[
        L = \lim_{x \to 1} \frac{\sqrt{x^4 + 3} - [A + B(x - 1) + C(x - 1)^2]}{(x - 1)^2}
        \]
        由于分母当 \( x \to 1 \) 时趋于 0,而极限存在且为 0,故分子必趋于 0:
        \[
        \lim_{x \to 1} \left( \sqrt{x^4 + 3} - [A + B(x - 1) + C(x - 1)^2] \right) = \sqrt{1 + 3} - A = 2 - A = 0
        \]
        解得
        \[
        A = 2
        \]
        此时原式变为 $\frac{0}{0}$ 型,由洛必达法则,
        \[
        L \overset{H}{=} \lim_{x \to 1} \frac{\frac{2x^3}{\sqrt{x^4 + 3}} - B - 2C(x - 1)}{2(x - 1)} = 0
        \]
        同样地,由于分母趋于 0,分子必趋于 0:
        \[
        \lim_{x \to 1} \left( \frac{2x^3}{\sqrt{x^4 + 3}} - B - 2C(x - 1) \right) = \frac{2}{2} - B = 1 - B = 0
        \]
        解得
        \[
        B = 1
        \]
        再次应用洛必达法则处理该 $\frac{0}{0}$ 型极限:
        \[
        L \overset{H}{=} \lim_{x \to 1} \frac{\frac{d}{dx} \left( \frac{2x^3}{\sqrt{x^4 + 3}} - 1 - 2C(x - 1) \right)}{\frac{d}{dx} 2(x - 1)} = 0
        \]
        计算分子导数:
        \[
        \frac{d}{dx} \left( \frac{2x^3}{\sqrt{x^4 + 3}} \right) = \frac{6x^2\sqrt{x^4+3} - 2x^3 \cdot \frac{4x^3}{2\sqrt{x^4+3}}}{x^4 + 3} = \frac{6x^2(x^4+3) - 4x^6}{(x^4+3)\sqrt{x^4+3}}
        \]
        代入 \( x = 1 \) 得到
        \[
        \frac{6(4) - 4}{4\sqrt{4}} = \frac{20}{8} = \frac{5}{2}
        \]
        因此解得
        \[
        \frac{\frac{5}{2} - 2C}{2} = 0 \Rightarrow C = \frac{5}{4}
        \]
        综上,常数分别为
        \[
        A = 2,\quad B = 1,\quad C = \frac{5}{4}
        \]
    \end{solution}
    \begin{solution}
        由于极限式
        \[
        \lim_{x \to 1} \frac{\sqrt{x^4 + 3} - [A + B(x - 1) + C(x - 1)^2]}{(x - 1)^2} = 0
        \]
        符合函数 $f(x)$ 在 $x=1$ 处的二阶泰勒展开定义,其中多项式 
        \[
        P_2(x) = A + B(x-1) + C(x-1)^2
        \]
        必须等于 $f(x)$ 的二阶泰勒多项式。设 $f(x) = \sqrt{x^4+3}$,根据泰勒公式可知:
        \[
        A = f(1), \quad B = f'(1), \quad C = \frac{f''(1)}{2!}
        \]
        由
        \[
        f'(x) = \frac{1}{2\sqrt{x^4+3}} \cdot 4x^3 = \frac{2x^3}{\sqrt{x^4+3}}
        \]
        \[
        f''(x) = \frac{6x^2\sqrt{x^4+3} - 2x^3 \cdot \frac{2x^3}{\sqrt{x^4+3}}}{x^4+3} = \frac{6x^2(x^4+3) - 4x^6}{(x^4+3)\sqrt{x^4+3}}
        \]
        代入得
        \[
        f(1) = 2, \quad f'(1) = 1, \quad f''(1) = \frac{5}{2}
        \]
        因此
        \[
        A = 2,\quad B = 1,\quad C = \frac{5}{4}
        \]
    \end{solution}

    \question
    \[
    \lim_{n\to\infty}\left(\frac{1}{\sqrt{3n^{2}+1}}+\frac{1}{\sqrt{3n^{2}+2}}+\cdots+\frac{1}{\sqrt{3n^{2}+2n}}\right)
    \]
    \begin{solution}
        令\[
        L=\lim_{n\to\infty}\left(\frac{1}{\sqrt{3n^{2}+1}}+\frac{1}{\sqrt{3n^{2}+2}}+\cdots+\frac{1}{\sqrt{3n^{2}+2n}}\right)
        \]
        由夹挤定理,
        \[
        \frac{2}{\sqrt 3}=\lim_{n\to \infty} \left(\frac{2n}{\sqrt{3n^2}}    \right) < L < \lim_{n\to \infty} \left(\frac{2n}{\sqrt{3n^2+2n}} \right)=\frac{2}{\sqrt 3}
        \Rightarrow L = \frac{2\sqrt 3}{3}
        \]
    \end{solution}

    \question 设 $a_n = \sqrt{1 \cdot 2} + \sqrt{2 \cdot 3} + \cdots + \sqrt{n(n+1)}$,求 $\displaystyle \lim_{n \to \infty} \frac{a_n}{n^2}$ 的值。
    \begin{solution}
        发现到
        \[
        \sum_{k=1}^n \sqrt{k^2} < a_n < \sum_{k=1}^n \sqrt{(k+1)^2} \Rightarrow \frac{n(n+1)}{2} < a_n < \frac{n(n+3)}{2} 
        \] 
        于是有
        \[
        \lim_{n \to \infty} \frac{n(n+1)}{2n^2} < \lim_{n \to \infty} \frac{a_n}{n^2} < \lim_{n \to \infty} \frac{n(n+3)}{2n^2} \Rightarrow \frac{1}{2} < \lim_{n \to \infty} \frac{a_n}{n^2} < \frac{1}{2} 
        \] 
        由夹挤定理,
        \[
        \lim_{n \to \infty} \frac{a_n}{n^2} = \frac{1}{2}
        \]
    \end{solution}

    \question 设 $a_n = \dfrac{1 \cdot 3 \cdot 5 \cdot \dots \cdot (2n-1)}{2 \cdot 4 \cdot 6 \cdot \dots \cdot (2n)}$,求 $\displaystyle \lim_{n \to \infty} a_n$ 之值。
    \begin{solution}
        由
        \[
        (2n-1)(2n+1) < (2n)^2
        \]
        令 $S=1\cdot 3\cdot 5\cdot \cdots \cdot (2n-1)$, 则
        \begin{align*} 
        (2n+1)S^2 = (1\cdot 3)(3\cdot 5)(5\cdot 7)\cdots ((2n-3)(2n-1))((2n-1)(2n+1))
        < 2^2\cdot 4^2 \cdot 6^2\cdot \cdots \cdot (2n)^2
        \end{align*}
        即
        \[
        \sqrt{2n+1}S < 2\cdot 4\cdot 6\cdot \cdots \cdot (2n)
        \]
        于是
        \[
        \Rightarrow 0 < a_n=\frac{S}{2\cdot 4\cdot 6\cdot \cdots \cdot (2n)} < \frac{1}{\sqrt{2n+1}}
        \]
        由夹挤定理,
        \[
        0 < \lim_{n\to \infty} a_n < \lim_{n\to \infty} \frac{1}{\sqrt{2n+1}}=0 \Rightarrow \lim_{n\to \infty} a_n = 0
        \]
    \end{solution}

    \question 求极限
    \[
    \lim_{n \to \infty} \sum_{k=1}^{n} \frac{1}{\sqrt{kn}}
    \]
    \begin{solution}
        由放缩法,
        \[
        \sqrt{k-1}+ \sqrt{k} < 2\sqrt{k} < \sqrt{k} + \sqrt{k+1} 
        \]
        有理化得
        \[
        2(\sqrt{k+1}-\sqrt{k}) < \frac{2}{\sqrt{k}} < 2(\sqrt{k}-\sqrt{k-1})
        \]
        于是有
        \[
        \frac{2}{\sqrt{n}} \sum_{k=1}^{n} (\sqrt{k+1}-\sqrt{k}) < \frac{2}{\sqrt{n}} \sum_{k=1}^{n} \frac{1}{\sqrt{k}} < \frac{2}{\sqrt{n}} \sum_{k=1}^{n} (\sqrt{k}-\sqrt{k-1})
        \]
        累加得
        \[
        \frac{2\sqrt{n+1}-2}{\sqrt{n}} < \frac{2}{\sqrt{n}} \sum_{k=1}^{n} \frac{1}{\sqrt{k}} < \frac{2\sqrt{n}}{\sqrt{n}} = 2
        \]
        因为
        \[
        \lim_{n \to \infty} \frac{2\sqrt{n+1}-2}{\sqrt{n}} = 2
        \]
        由夹挤定理得,
        \[
        \lim_{n \to \infty} \sum_{k=1}^{n} \frac{1}{\sqrt{kn}} = 2
        \]
    \end{solution}

    \question 求极限
    \[\lim_{n \to \infty} \left(\frac{1}{n^2+n+1} + \frac{2}{n^2+n+2} + \dots + \frac{n}{n^2+n+n}\right)\]
    \begin{solution}
        由放缩法,
        \[
        \frac{1 + 2 + \dots + n}{n^2+n+n} < \sum_{k=1}^{n} \frac{k}{n^2+n+k} < \frac{1 + 2 + \dots + n}{n^2+n+1}
        \]
        即
        \[
        \frac{n(n+1)}{2(n^2+2n)} < \sum_{k=1}^{n} \frac{k}{n^2+n+k} < \frac{n(n+1)}{2(n^2+n+1)}
        \]
        由于
        \[
        \lim_{n \to \infty} \frac{n(n+1)}{2(n^2+2n)} =
        \lim_{n \to \infty} \frac{n(n+1)}{2(n^2+n+1)} = \frac{1}{2}
        \]
        由夹挤定理得
        \[
        \lim_{n \to \infty} \sum_{k=1}^{n} \frac{k}{n^2+n+k} = \frac{1}{2}
        \]
    \end{solution}

    \question 求极限
    \[
    \lim_{n \to \infty} \frac{1}{n^2} \sum_{k=1}^{n} \sqrt{k^2+n+k}
    \]
    \begin{solution}
        对于整数$1 \le k \le n$,有
        \[
        \sqrt{k^2} < \sqrt{k^2+n+k} < \sqrt{k^2+2n}
        \]
        于是
        \[
        \frac{1}{n^2} \sum_{k=1}^{n} k < \frac{1}{n^2} \sum_{k=1}^{n} \sqrt{k^2+n+k} \le \frac{1}{n^2} \sum_{k=1}^{n} (k+\sqrt{2n})
        \]
        由于
        \[
        \lim_{n \to \infty} \frac{1}{n^2} \sum_{k=1}^{n} k = \lim_{n \to \infty} \frac{n(n+1)}{2n^2} = \frac{1}{2}
        \]
        且
        \begin{align*}
        \lim_{n \to \infty} \frac{1}{n^2} \sum_{k=1}^{n} (k+\sqrt{2n}) &= \lim_{n \to \infty} \frac{1}{n^2} \sum_{k=1}^{n} k + \lim_{n \to \infty} \frac{1}{n^2} \sum_{k=1}^{n} \sqrt{2n} \\
        &= \frac{1}{2} + \lim_{n \to \infty} \frac{n \sqrt{2n}}{n^2} \\
        &= \frac{1}{2}
        \end{align*}
        由夹挤定理,
        \[
        \lim_{n \to \infty} \frac{1}{n^2} \sum_{k=1}^{n} \sqrt{k^2+n+k} = \frac{1}{2}
        \]
    \end{solution}

    \question 求极限
    \[\lim_{n \to \infty} \left(\frac{1}{n+1} + \frac{1}{n+\sqrt{2}} + \frac{1}{n+\sqrt{3}} + \dots + \frac{1}{n+\sqrt{n}}\right)\]
    \begin{solution}
        由放缩法得
        \[
        \frac{n}{n+\sqrt{n}} < \sum_{k=1}^{n} \frac{1}{n+\sqrt{k}} < \frac{n}{n+1}
        \]
        由于
        \[
        \lim_{n \to \infty} \frac{n}{n+\sqrt{n}} = \lim_{n \to \infty} \frac{1}{1 + 1/n} = 1
        \]
        由夹挤定理,
        \[
        \lim_{n \to \infty} \left(\frac{1}{n+1} + \frac{1}{n+\sqrt{2}} + \dots + \frac{1}{n+\sqrt{n}}\right) = 1
        \]
    \end{solution}

    \question 求极限
    \[
    \lim_{n \to \infty} \frac{n^k \sin^2(n!)}{n+2}, \quad 0<k<1
    \]
    \begin{solution}
        由于 
        \[
        0 \le \sin^2 x \le 1, \quad \forall x \in \mathbb{R}
        \]
        得到不等式
        \[
        0 \le \frac{n^k \sin^2(n!)}{n+2} \le \frac{n^k}{n+2}
        \]
        当 $0<k<1$ 时,上界的极限为
        \[
        \lim_{n \to \infty} \frac{n^k}{n+2} = 0
        \]
        故由夹挤定理,
        \[
        \lim_{n \to \infty} \frac{n^k \sin^2(n!)}{n+2} = 0
        \]
    \end{solution}

    \question 求极限
    \[
    \lim_{n \to \infty} \sqrt[n]{3^n + (2|\sin(n^n)|)^n}
    \]
    \begin{solution}
        由于 
        \[
        0 \le |\sin(n^n)| \le 1
        \]
        得到上下界:
        \[
        \sqrt[n]{3^n} \le \sqrt[n]{3^n + (2|\sin(n^n)|)^n} \le \sqrt[n]{3^n + 2^n} \le \sqrt[n]{2 \cdot 3^n}
        \]
        由于
        \[
        \lim_{n \to \infty} \sqrt[n]{3^n} = 
        \lim_{n \to \infty} \sqrt[n]{2 \cdot 3^n}  = 3
        \]
        由夹挤定理,
        \[
        \lim_{n \to \infty} \sqrt[n]{3^n + (2|\sin(n^n)|)^n} = 3
        \]
    \end{solution}

    \question 已知数列 $\{a_n\}$ 满足
    \[
    \frac{2n^2-7}{4n+5} < a_n < \frac{3n^2+8}{6n-1}, \quad n \in \mathbb{Z}^+
    \]
    求极限
    \[
    \lim_{n \to \infty} \frac{3 n a_n}{(n+1)^2}
    \]
    \begin{solution}
        将不等式两边同时乘以 $\dfrac{3n}{(n+1)^2}$:
        \[
        \frac{3n}{(n+1)^2} \cdot \frac{2n^2-7}{4n+5} < \frac{3 n a_n}{(n+1)^2} < \frac{3n}{(n+1)^2} \cdot \frac{3n^2+8}{6n-1}
        \]
        取极限得
        \[
        \frac{3}{2}=\lim_{n \to \infty} \frac{3n(2n^2-7)}{(4n+5)(n+1)^2} \le \lim_{n \to \infty} \frac{3 n a_n}{(n+1)^2} \le \lim_{n \to \infty} \frac{3n(3n^2+8)}{(6n-1)(n+1)^2} =\frac{3}{2}
        \]
        由夹挤定理,
        \[
        \lim_{n \to \infty} \frac{3 n a_n}{(n+1)^2} = \frac{3}{2}
        \]
    \end{solution}

    \question 已知 $x_{1}=1$,且对每个正整数 $n\ge 2$,
    \[
    n(x_{n})^{2}-x_{n-1}-n=0,\quad x_{n}\ge 0.
    \]
    求 $\displaystyle\lim_{n\to\infty} x_{n}$ 或证明 $\{x_{n}\}$ 发散。
    \begin{solution}
        由等式得对所有 $n\ge 2$,有
        \[
        x_n=\sqrt{1+\frac{x_{n-1}}{n}}
        \]
        首先证明$x_n \le 2$:有 $x_1=1\le 2$。假设 $n\ge 2$ 有 $x_n\le 2$成立,则
        \[
        x_n=\sqrt{1+\frac{x_{n}}{n+1}}\le \sqrt{1+2} <2
        \]
        故由数学归纳法得对所有 $n\ge 2$ 都皆有 $x_n\le 2$,故
        \[
        1 \le x_n=\sqrt{1+\frac{x_{n-1}}{n}}\le \sqrt{1+\frac{2}{n}}
        \]
        由夹挤定理,
        \[
        \lim_{n\to\infty} x_{n} = 1
        \]
    \end{solution}

    \question 
    \begin{parts}
    \part 证明
    \[
    \int_{0}^{1} \left(1+\sin \frac{\pi}{2}x\right)^n dx > \frac{2^{n+1}-1}{n+1} \quad (n=1,2,\dots)
    \]
    \begin{solution}
        由 
        \[
        \sin x \ge \frac{2}{\pi}x, \quad 0 \le x \le \frac{\pi}{2}
        \]
        得
        \[
        \sin \left( \frac{\pi}{2}x \right) \ge x, \quad 0 \le x \le 1,
        \]
        从而
        \[
        \int_0^1 \left(1 + \sin \frac{\pi}{2}x\right)^n dx \ge \int_0^1 (1 + x)^n dx = \left. \frac{(1+x)^{n+1}}{n+1} \right|_0^1 = \frac{2^{n+1} - 1}{n+1}.
        \]
        因此原不等式成立。
    \end{solution}
    \part 求极限:
    \[
    \lim_{n \to \infty} \left[\int_{0}^{1} \left(1+\sin \frac{\pi}{2}x\right)^n dx\right]^{\frac{1}{n}}
    \]
    \begin{solution}
        注意到:
        \[
        \frac{2^n}{n+1} < \frac{2^{n+1}-1}{n+1} \le \int_0^1 \left(1 + \sin \frac{\pi}{2}x\right)^n dx \le \int_0^1 2^n dx = 2^n.
        \]
        对上述不等式取 $n$ 次方根,得:
        \[
        2(n+1)^{-\frac{1}{n}} < \left[\int_0^1 \left(1 + \sin \frac{\pi}{2}x\right)^n dx\right]^{\frac{1}{n}} \le 2.
        \]
        因为 $(n+1)^{-\frac{1}{n}} \to 1$,由夹挤定理可得:
        \[
        \lim_{n \to \infty} \left[\int_0^1 \left(1 + \sin \frac{\pi}{2}x\right)^n dx\right]^{\frac{1}{n}} = 2
        \]
    \end{solution}
    \end{parts}
    
    \question 
    \begin{parts}
    \part 设 $n$ 是正整数, 计算 $$\int_{0}^{n\pi} x \sin^2 x \,dx$$
    \begin{solution}
        对非负整数 $k$,
        \[
        \int_{k\pi}^{(k+1)\pi} x \sin^2 x \,dx
        = \int_0^\pi (k\pi + x) \sin^2 x \,dx
        \]
        利用恒等式 \( \sin^2 x = \frac{1 - \cos 2x}{2} \),得:
        \[
        = \frac{1}{2} \int_0^\pi (k\pi + x)(1 - \cos 2x)\,dx
        = \frac{1}{2} \left( (k\pi + x)\big|_0^\pi - \int_0^\pi (k\pi + x)\cos 2x\,dx \right)
        \]
        由于 \( \int_0^\pi \cos 2x\,dx = 0 \) 且 \( \int_0^\pi x\cos 2x\,dx = 0 \),
        所以最后只剩:
        \[
        \int_{k\pi}^{(k+1)\pi} x \sin^2 x\,dx 
        = \frac{1}{2} \int_0^\pi (k\pi + x)\,dx 
        = \frac{1}{2} \left( k\pi^2 + \frac{\pi^2}{2} \right) = \frac{\pi^2(2k+1)}{4}          
        \]
        所以,
        \[
        \int_0^{n\pi} x \sin^2 x \,dx = \sum_{k=0}^{n-1} \frac{\pi^2(2k+1)}{4} = \frac{\pi^2}{4} \sum_{k=0}^{n-1} (2k+1) = \frac{\pi^2}{4} \cdot n^2
        \]
        因为 \( \displaystyle\sum_{k=0}^{n-1}(2k+1) = n^2 \),所以原式为
        \[
        \frac{\pi^2 n^2}{4}
        \]
    \end{solution}
    
    \part 证明对任何正实数 $p$, 函数极限 $$\lim_{x\to+\infty} \frac{1}{x^2} \int_{0}^{x} t |\sin t|^p \,dt$$ 存在.
    \begin{solution}
        记
        \[
        S_0 = \int_0^\pi |\sin t|^p \,dt, \qquad S_1 = \int_0^\pi t |\sin t|^p \,dt.
        \]
        对任意非负整数 \( k \),有
        \[
        \int_{k\pi}^{(k+1)\pi} t |\sin t|^p \,dt = \int_0^\pi (k\pi + x) |\sin x|^p dx = k\pi S_0 + S_1.
        \]
        若设 \( x = (n + \alpha)\pi \) 且 \( 0 \le \alpha < 1 \),则
        \[
        \int_0^x t |\sin t|^p dt = \sum_{k=0}^{n-1} (k\pi S_0 + S_1) + \int_{n\pi}^{x} t |\sin t|^p dt.
        \]
        利用估计（注意 \( \int_{n\pi}^x t |\sin t|^p dt \le (n+1)\pi \cdot S_0 \)）,得上下界:
        \[
        \frac{\frac{1}{2}n(n-1)\pi S_0 + nS_1}{(n+1)^2\pi^2} \le \frac{1}{x^2} \int_0^x t |\sin t|^p dt \le \frac{\frac{1}{2}n(n+1)\pi S_0 + (n+1)S_1}{n^2\pi^2}.
        \]
        随着 \( n \to \infty \),上下界极限均趋于 \( \frac{S_0}{2\pi} \),所以由夹挤定理知极限存在,且为
        \[
        \lim_{x \to +\infty} \frac{1}{x^2} \int_0^x t |\sin t|^p dt = \frac{S_0}{2\pi}
        \]
    \end{solution}
    \end{parts}
    
    \question 
    \begin{parts}
    \part 求 $$\int_{0}^{\pi} \frac{1}{1+\cos^2 x} dx$$ 及$$\int_{0}^{\pi} \frac{\sin^2 x}{1+\cos^2 x} dx$$
    \begin{solution}
        设
        \[
        I_1 = \int_0^\pi \frac{1}{1 + \cos^2 x} \, dx, \quad
        I_2 = \int_0^\pi \frac{\sin^2 x}{1 + \cos^2 x} \, dx.
        \]
        发现
        \[
        I_1 = \frac12\int_0^\pi \frac{1+\sin^2 x +\cos^2 x}{1 + \cos^2 x} \, dx =\frac12 (I_2 + \pi)\Rightarrow I_2 = 2I_1 -\pi
        \]
        所以只需计算 \( I_1 \)即可得出两个积分。发现被积函数在$x=\dfrac\pi2$处不连续,将$I_1$写成\[
        I_1 = \int _{0}^{\frac\pi2}\frac{1}{1+\cos ^{2}x}dx+\int _{\frac\pi2}^{\pi }\frac{1}{1+\cos ^{2}x}dx 
        \]
        其中第一个积分变为
        \[
        \int _{0}^{\frac\pi2}\frac{\sec ^{2}x}{2+\tan ^{2}x}dx=\int _{0}^{\infty }\frac{1}{2+u^{2}}du=\left[\frac{1}{\sqrt{2}}\arctan \left(\frac{u}{\sqrt{2}}\right)\right]_{0}^{\infty }=\frac{1}{\sqrt{2}}\left(\frac{\pi }{2}-0\right)=\frac{\pi }{2\sqrt{2}}.
        \]
        对于第二个积分如下,设 \(x=\pi -y,dx=-dy\),
        \[
        \int _{\frac\pi2}^{\pi }\frac{1}{1+\cos ^{2}x}dx=\int _{\frac\pi2}^{0}\frac{1}{1+\cos ^{2}(\pi -y)}(-dy)=\int _{0}^{\frac\pi2}\frac{1}{1+\cos ^{2}y}dy=\frac{\pi }{2\sqrt{2}}
        \]
        所以
        \[
        I_1=\frac{\pi }{2\sqrt{2}}+\frac{\pi }{2\sqrt{2}}=\frac{\pi \sqrt{2}}{2}
        \]
        \[
        I_2 = 2I_1 -\pi = 2\cdot \frac{\pi \sqrt{2}}{2}-\pi = \pi (\sqrt2-1)
        \]
    \end{solution}

    \part 证明 
    \[
    \lim_{x \to \infty} \frac{\displaystyle \int_{0}^{x} \frac{\sin^2 t}{1+\cos^2 t} dt}{\displaystyle\int_{0}^{x} \frac{1}{1+\cos^2 t} dt} = 2-\sqrt{2}
    \]
    \begin{solution}
        取 $n = \left\lfloor \dfrac{x}{\pi} \right\rfloor$,则
        \[
        \int_0^x \frac{1}{1+\cos^2 t}\,dt = \sum_{k=1}^n \int_{(k-1)\pi}^{k\pi} \frac{1}{1+\cos^2 t}\, dt +\int_{k\pi}^{x}\frac{1}{1+\cos^2 t}\, dt
        \]
        于是
        \[
        \frac{n\pi}{\sqrt{2}} \le \int_0^x \frac{1}{1+\cos^2 t}\,dt \le \frac{(n+1)\pi}{\sqrt{2}}
        \]
        同理
        \[
        n\pi(\sqrt2-1) \le \int_0^x \frac{1}{1+\cos^2 t}\,dt \le (n+1)\pi(\sqrt2-1)
        \]
        于是
        \[
        \frac{n}{n+1}(2-\sqrt{2}) \le \frac{\displaystyle \int_{0}^{x} \frac{\sin^2 t}{1+\cos^2 t} dt}{\displaystyle\int_{0}^{x} \frac{1}{1+\cos^2 t} dt} \le \frac{n+1}{n}(2-\sqrt{2}),
        \]
        令 $x \to \infty$ 便得结论。
    \end{solution}
    \end{parts}

    \question 对每个正整数 $n$,设
    \[
    R_n = \{(x,y) \mid 0 \le x \le n, \ 0 \le y \le \sqrt{x}\},
    \]
    记 $N(n)$ 为 $R_n$ 中坐标都是整数的点的个数。求
    \[
    \lim_{n\to\infty} \frac{N(n)}{n^{\frac{3}{2}}}.
    \]
    \begin{solution}
        对每个正整数 $n$,满足 $0 \le y \le \sqrt{k}$ 的整数 $y$ 个数为 $[\sqrt{k}]+1$,所以
        \[
        N(n) = \sum_{k=0}^{n} ([\sqrt{k}]+1) = (n+1)+\sum_{k=1}^{n} [\sqrt{k}]
        \]
        由于 $\sqrt{k}-1 < [\sqrt{k}] \le \sqrt{k}$,并且
        \[
        \int_{0}^{n} \sqrt{x}\,dx \le \sum_{k=1}^{n} [\sqrt{k}] \le \int_{0}^{n+1} \sqrt{x}\,dx
        \]
        所以
        \[
        1 + \int_{0}^{n} \sqrt{x}\,dx \le N(n) \le (n+1) + \int_{0}^{n+1} \sqrt{x}\,dx
        \]
        即
        \[
        1 + \frac{2}{3} n^{\frac{3}{2}} \le N(n) \le (n+1) + \frac{2}{3} (n+1)^{\frac{3}{2}}
        \]
        由夹挤定理,得到
        \[
        \lim_{n\to\infty} \frac{N(n)}{n^{\frac{3}{2}}} = \frac{2}{3}
        \]
    \end{solution}
\end{questions}

\pagebreak

\begin{center}
  {\fontsize{30pt}{26pt}\selectfont
    \hypertarget{微分}{微分} \label{微分}
  }
\end{center}
\separator
\vspace{1pt}

\begin{questions}
    \question 已知 $f(x) = x^{2013} - x^{2012} + x^{2011} - x^{2010} + 1$,求 
    \[
    \lim_{h \to 0} \frac{f(1+h) - f(1-h)}{h}
    \]
    \begin{solution}
        计算导数:
        \[ 
        f'(x) = 2013x^{2012} - 2012x^{2011} + 2011x^{2010} - 2010x^{2009} 
        \]
        于是
        \begin{align*}
        \lim_{h \to 0} \frac{f(1+h) - f(1-h)}{h} 
        &= \lim_{h \to 0} \frac{f(1+h) - f(1)}{h} + \lim_{h \to 0} \frac{f(1) - f(1-h)}{h} \\
        &= f'(1) + f'(1) \\
        &= 2 \cdot (2013-2012+2011-2010) \\
        &= 4
        \end{align*}
    \end{solution}

    \question 已知 \( f'(1) = 12 \),求
    \[
    \lim_{h \to 0} \frac{f(1 + 4h) - f(1 - 2h)}{3h}
    \]
    \begin{solution}
        改写
        \begin{align*}
        \frac{f(1 + 4h) - f(1 - 2h)}{3h}
        &= \frac{f(1 + 4h) - f(1)}{3h} + \frac{f(1) - f(1 - 2h)}{3h} \\
        &= \frac{4}{3} f'(1) + \frac{2}{3} f'(1) \\
        &= 24
        \end{align*}
    \end{solution}

    \question 求
    \[
    \lim_{h \to 0} \left[ \int_{\frac{\pi}{6}}^{\frac{\pi}{6}+h} \frac{2\sqrt{\sin x}}{\pi h} \, dx \right]
    \]
    \begin{solution}
        先提取常数
        \[
        \lim_{h \to 0} \left[ \int_{\frac{\pi}{6}}^{\frac{\pi}{6}+h} \frac{2\sqrt{\sin x}}{\pi h} \, dx \right] 
        = \frac{2}{\pi} \lim_{h \to 0} \left[ \frac{1}{h} \int_{\frac{\pi}{6}}^{\frac{\pi}{6}+h} \sqrt{\sin x} \, dx \right]
        \]
        令
        \[
        F(x) = \int \sqrt{\sin x} \, dx \Rightarrow F'(x) = \sqrt{\sin x}
        \]
        则
        \[
        \frac{1}{h} \int_{\frac{\pi}{6}}^{\frac{\pi}{6}+h} \sqrt{\sin x} \, dx = \frac{F\left(\frac{\pi}{6}+h\right) - F\left(\frac{\pi}{6}\right)}{h}
        \]
        当 \(h \to 0\) 时,上式为导数定义:
        \[
        \lim_{h \to 0} \frac{F\left(\frac{\pi}{6}+h\right) - F\left(\frac{\pi}{6}\right)}{h} = F'\left(\frac{\pi}{6}\right) = \sqrt{\sin \frac{\pi}{6}}
        \]
        因此原极限为
        \[
        \frac{2}{\pi} \cdot \sqrt{\sin \frac{\pi}{6}} = \frac{2}{\pi} \cdot \frac{\sqrt{2}}{2} = \frac{\sqrt{2}}{\pi}
        \]
    \end{solution}

    \question 由导数定义,证明 
    \[
    \frac{d}{dx}(\sec x) = \sec x \tan x
    \]
    \begin{solution}
        由导数的定义,
        \begin{align*}
        \frac{d}{dx}(\sec x) 
        &= \lim_{h \to 0} \frac{\sec(x+h) - \sec x}{h}\\
        &= \lim_{h \to 0} \frac{\cos x - \cos(x+h)}{h \cos x \cos(x+h)}\\
        &= \lim_{h \to 0} \frac{2 \sin\left(x+\frac{h}{2}\right) \sin\frac{h}{2}}{h \cos x \cos(x+h)}\\
        &= \lim_{h \to 0} \left[ \frac{\sin\frac{h}{2}}{\frac{h}{2}} \cdot \frac{\sin\left(x+\frac{h}{2}\right)}{\cos x \cos(x+h)} \right] \\
        &= 1 \cdot \frac{\sin x}{\cos^2 x} \\
        &= \sec x \tan x
        \end{align*}
    \end{solution}

    \question 通过导数的定义,求 
    \[
    f(x) = \frac{1}{\sqrt{x^2-1}}, \quad |x|>1
    \]
    的导数$f'(x)$。
    \begin{solution}
        由导数的定义,
        \begin{align*}
        f'(x) &= \lim_{h \to 0} \frac{f(x+h) - f(x)}{h} \\
        &= \lim_{h \to 0} \frac{\frac{1}{\sqrt{(x+h)^2-1}} - \frac{1}{\sqrt{x^2-1}}}{h} \\
        &= lim_{h \to 0} \frac{\sqrt{x^2-1} - \sqrt{(x+h)^2-1}}{h \sqrt{x^2-1} \sqrt{(x+h)^2-1}} \\
        &= \lim_{h \to 0} \frac{(x^2-1) - ((x+h)^2-1)}{h \sqrt{x^2-1} \sqrt{(x+h)^2-1} (\sqrt{x^2-1} + \sqrt{(x+h)^2-1})} \\
        &= \lim_{h \to 0} \frac{-2xh - h^2}{h \sqrt{x^2-1} \sqrt{(x+h)^2-1} (\sqrt{x^2-1} + \sqrt{(x+h)^2-1})} \\
        &= \lim_{h \to 0} \frac{-2x - h}{\sqrt{x^2-1} \sqrt{(x+h)^2-1} (\sqrt{x^2-1} + \sqrt{(x+h)^2-1})} \\
        &= \frac{-2x}{\sqrt{x^2-1} \sqrt{x^2-1} (2\sqrt{x^2-1})} \\
        &= -\frac{x}{(x^2-1)^{\frac{3}{2}}}
        \end{align*}
    \end{solution}

    \question 设 
    \[
    y = \frac{(x^4 + 2)(x^3 - 2)}{(x + 2)(x^2 - 2)},
    \]
    求 $\dfrac{dy}{dx}$。
    \begin{solution}
        对等式两边取自然对数,
        \[ 
        \ln y = \ln(x^4 + 2) + \ln(x^3 - 2) - \ln(x + 2) - \ln(x^2 - 2) 
        \]
        对 $x$ 求导,
        \[ 
        \frac{1}{y} \frac{dy}{dx} = \frac{4x^3}{x^4 + 2} + \frac{3x^2}{x^3 - 2} - \frac{1}{x + 2} - \frac{2x}{x^2 - 2} 
        \]
        整理得
        \[ 
        \frac{dy}{dx} = \frac{4x^3(x^3 - 2)}{(x + 2)(x^2 - 2)} + \frac{3x^2(x^4 + 2)}{(x + 2)(x^2 - 2)} - \frac{(x^4 + 2)(x^3 - 2)}{(x + 2)^2(x^2 - 2)} - \frac{2x(x^4 + 2)(x^3 - 2)}{(x + 2)(x^2 - 2)^2} 
        \]
    \end{solution}

    \question 已知曲线方程
    \[
    y = 2^{3e^{2x}}, \quad x \in \mathbb{R}.
    \]
    以$y$ 表示 $\frac{dy}{dx}$。
    \begin{solution}
        由链导法,
        \[
        \frac{dy}{dx} = 2^{3e^{2x}} \cdot \ln 2 \cdot 6 e^{2x} = = y \cdot \ln 2 \cdot 6 e^{2x}
        \]
        又由原式取自然对数得
        \[
        \ln y = 3 e^{2x} \ln 2 \Rightarrow 2 \ln y = 6 e^{2x} \ln 2
        \]
        因此得到
        \[
        \frac{dy}{dx} = 2y \ln y
        \]
    \end{solution}
    \begin{solution}
        不如一开始就取$\ln$,
        \[
        \ln y = 3 e^{2x} \ln 2.
        \]
        对 $x$ 求导,
        \[
        \frac{1}{y}\frac{dy}{dx} = \ln 2 \cdot 6 e^{2x} 
        \]
        同样代入 $\ln y = 3 e^{2x} \ln 2$ 得
        \[
        \frac{dy}{dx} = = y \cdot (\ln 2) \cdot 6 e^{2x} = 2 y \ln y
        \]
    \end{solution}

    \question 已知曲线由方程
    \[
    x^m y^n = (x+y)^{m+n}, \quad x\neq 0,\; y\neq 0,\; x+y\neq 0,\; my-nx\neq 0
    \]
    隐式定义,其中 $m,n$ 为有理数。证明
    \[
    \frac{dy}{dx} = \frac{y}{x}
    \]
    \begin{solution}
        两边取对数,
        \[
        m\ln x + n\ln y = (m+n)\ln(x+y)
        \]
        对 $x$ 求导,
        \[
        \frac{m}{x} + \frac{n}{y} \frac{dy}{dx} = \frac{m+n}{x+y}\left(1+\frac{dy}{dx}\right)
        \]
        整理得
        \[
        \frac{my-nx}{x(x+y)} = \frac{my-nx}{y(x+y)} \frac{dy}{dx}
        \]
        由于$x+y\neq 0, my-nx\neq 0$,故
        \[
        \frac{dy}{dx} = \frac{y}{x}
        \]
    \end{solution}

    \question 已知曲线 $C$ 的隐函数方程为
    \[
    y = xe^y, \quad x \neq 0, \; y \neq 1, \; y \neq 2.
    \]
    证明
    \[
    (1-y)\frac{d^2 y}{dx^2} = (2-y)\left(\frac{dy}{dx}\right)^2.
    \]
    \begin{solution}
        对 $y = xe^y$ 两边求导,
        \[
        \frac{dy}{dx} = e^y + x e^y \frac{dy}{dx} \Rightarrow \frac{dy}{dx}(1 - y) = e^y 
        \]
        再次求导,
        \[
        \frac{d^2 y}{dx^2}(1-y) - \left(\frac{dy}{dx}\right)^2 = e^y \frac{dy}{dx}
        \]
        代入 $e^y = \dfrac{dy}{dx}(1-y)$,
        \[
        \frac{d^2 y}{dx^2}(1-y) - \left(\frac{dy}{dx}\right)^2 = \left(\frac{dy}{dx}\right)^2 (1-y)
        \]
        即 
        \[
        \frac{d^2 y}{dx^2}(1-y) = \left(\frac{dy}{dx}\right)^2 (2-y)
        \]
    \end{solution}

    \question 求曲线 $C$ 
    \[
    x^2 + 3xy - 2y^2 + 17 = 0
    \]
    的极值。
    \begin{solution}
        设 $F(x, y) = x^2 + 3xy - 2y^2 + 17$。对方程$F(x, y)=0$两边求导,
        \[
        2x + 3y + 3x y' - 4y y' = 0
        \]
        整理得：
        \[
        y' = -\frac{2x + 3y}{3x - 4y}
        \]
        令 $y' = 0$，得到 
        \[
        y = -\frac{2}{3}x
        \]
        代入原方程解得
        \[
        x^2 + 3x\left(-\frac{2}{3}x\right) - 2\left(-\frac{2}{3}x\right)^2 + 17 = 0 \Rightarrow x_1 = 3, x_2 = -3
        \]
        得临界点$P_1(3, -2),P_2(-3, 2)$。再求导得,
        \[
        y'' (3x - 4y) + y' (3 - 4y') = -(2 + 3y')
        \]
        在临界点处 $y' = 0$,化简得
        \[
        y'' (3x - 4y) = -2 \Rightarrow y'' = \frac{-2}{3x - 4y}
        \]
        由于
        \[
        y''(3, -2) = \frac{-2}{3(3) - 4(-2)}= -\frac{2}{17} < 0, \quad y''(-3, 2) = \frac{-2}{3(-3) - 4(2)} = \frac{2}{17} > 0
        \]
        故 $y_1 = -2$ 是极大值,$y_2 = 2$ 是极小值。该曲线在 $x=3$ 处取得极大值 $-2$，在 $x=-3$ 处取得极小值 $2$。
    \end{solution}

    \question 设函数 \( f(x) \) 处处可导,且 \( f'(0) = 1 \),并对任意实数 \( x,h \) 恒有
    \[
    f(x + h) = f(x) + f(h) + 2hx,
    \]
    求 \( f'(x) \)。
    \begin{solution}
        已知$\forall \:x,h \in \mathbb{R}$,
        \[
        f(x + h) = f(x) + f(h) + 2hx
        \]
        对$h$微分,
        \[
        f'(x+h)=f'(h)+2x
        \]
        令$h=0$,则
        \[
        f'(x)=f'(0)+2x=1+2x
        \]
    \end{solution}

    \question 设 $f(x)=x(x-1)(x-2)\cdots(x-2023),g(x)=f(f(x))$,求 $g'(1)$。
    \begin{solution}
        求导得
        \[
        f'(x) = (x-1)P(x) + x(x-2)(x-3)\cdots(x-2023) = xQ(x) + (x-1)(x-2)\cdots(x-2023)
        \]
        所以有
        \[
        f(1) = 0,\quad f'(1) = 2022!,\quad f'(0) = -2023!
        \]
        由于 $g(x) = f(f(x))$,有
        \[
        g'(x) = f'(f(x)) \cdot f'(x)\Rightarrow g'(1) = f'(f(1)) \cdot f'(1) = f'(0) \cdot 2022! = -2023! \cdot 2022!
        \]
    \end{solution}

    \question 设 $f(x)$ 为正实函数且为可微分函数,对任意实数 $x,y$,满足 $f(x+y)=2f(x)f(y)$,若 $f'(0)=2$,试求 $\dfrac{f''(x)}{f(x)}$。
    \begin{solution}
        由$f(x+y)=2f(x)f(y)$代入 $x=y=0$得
        \[
        f(0)=2f(0)f(0) \Rightarrow f(0) = \frac{1}{2}
        \]
        考虑 $f'(x)$,
        \begin{align*}
        f'(x)&=\lim_{h\to 0}\frac{f(x+h)-f(x)}{h} 
        = \lim_{h\to 0} \frac{2f(x)f(h)-f(x)}{h}  \\      
        &= 2f(x)\lim_{h\to 0} \frac{f(h)-f(0)}{h}
        = 2f(x)f'(0) = 2f(x) \cdot 2 = 4f(x)
        \end{align*}
        对两边求导得
        \[
        f''(x) = 4f'(x) = 4 \cdot 4f(x) 
        \Rightarrow \frac{f''(x)}{f(x)} = 16
        \]
    \end{solution}

\question
已知 \( f(x) \) 在 \( x = 0 \) 连续,且
\[
\lim_{x \to 0} \frac{\ln (f(x) + 2)}{x - \sin x} = 1,
\]
求 \( f'(0) \)。
    
\begin{solution}
    因为 \( f(x) \) 在 \( x = 0 \) 连续,所以
    \[
    \lim_{x \to 0} \ln(f(x)+2) = \ln(f(0)+2).
    \]
    又因 \( x - \sin x \to 0 \),分母趋于 0,极限存在,需有分子也趋于 0,即
    \[
    \ln(f(0)+2) = 0 \Rightarrow f(0) = -1.
    \]
    
    所以
    \[
    \lim_{x \to 0} \frac{\ln(f(x)+2)}{x - \sin x}
    = \lim_{x \to 0} \frac{f(x)+1}{x - \sin x} \quad (\text{因 } \ln(1 + y) \backsim y).
    \]
    又 \( x - \sin x \backsim \frac{x^3}{6} \),于是
    \[
    \lim_{x \to 0} \frac{f(x)+1}{x^3} = \frac{1}{1/6} = 6.
    \]
    因此
    \[
    \lim_{x \to 0} \frac{f(x) - f(0)}{x} = \lim_{x \to 0} \frac{f(x)+1}{x} = \lim_{x \to 0} \frac{f(x)+1}{x^3} \cdot x^2 = 6 \cdot 0 = 0.
    \]
    故 \( f'(0) = \boxed{0} \)。
    \textcolor{red}{待验证}
\end{solution}

    \question 已知 $a,b$ 为实数,若函数 $f(x) = 2ax^{3} - 3ax^{2} + b$ 在 $0 \le x \le 2$ 时有最大值 $7$,最小值 $-3$,求数对 $(a,b)$ 。
    \begin{solution}
        导数为
        \[
        f'(x) = 6ax^2 - 6ax = 6ax(x-1)
        \]
        临界点为 $x = 0, 1$,二阶导为
        \[
        f''(x) = 12ax - 6a = 6a(2x - 1)
        \]
        当$a > 0$,则
        \[
        f(x)\begin{cases}
        \text{递增} & x \ge 1 \\
        \text{递减} & 0 \le x \le 1 \\
        \text{递增} & x \le 0
        \end{cases}
        \]
        极小值在 $x = 1$,极大值在 $x = 2$ 或 $x = 0$,解得
        \[
        \begin{cases}
        f(2) = 4a + b = 7 \\
        f(1) = -a + b = -3
        \end{cases}
        \Rightarrow a = 2, b = -1
        \]
        当$a < 0$,则
        \[
        f(x)\begin{cases}
        \text{递减} & x \ge 1 \\
        \text{递增} & 0 \le x \le 1 \\
        \text{递减} & x \le 0
        \end{cases}
        \]
        极大值在 $x = 1$,极小值在 $x = 0$ 或 $x = 2$,同理解得
        \[
        \begin{cases}
        f(1) = -a + b = 7 \\
        f(2) = 4a + b = -3
        \end{cases}
        \Rightarrow a = -2,b = 5
        \]
        $\therefore (a, b) = (2, -1) \ \text{或} \ (-2, 5)$
    \end{solution}

    \question 函数 $f$ 定义为
    \[
    f(n, y) = \sum_{x=1}^{n} \frac{x^2 y^x}{k}, \quad n \in \mathbb{N}, \quad y \in \mathbb{R}
    \]
    其中
    \[
    k = \sum_{r=1}^{n} r^2,
    \]
    若已知数列求和公式
    \[
    \sum_{r=1}^{n} r^4 = \frac{1}{30} n(n+1)(6n^3+9n^2+n-1),
    \]
    证明
    \[
    \left. \frac{d^2f}{dy^2} \right|_{y=1} + \left. \frac{df}{dy} \right|_{y=1} - \left[ \left. \frac{df}{dy} \right|_{y=1} \right]^2 = \frac{3n^4+6n^3-n^2-4n}{20(2n+1)^2}.
    \]
    \begin{solution}
        首先,利用标准求和公式:
        \[
        k = \sum_{x=1}^{n} x^2 = \frac{n(n+1)(2n+1)}{6}.
        \]
        既然 $k$ 与 $x$ 无关，可以将其作为常数因子提取到求和号外
        \[
        f(n,y) = \frac{1}{k} \sum_{x=1}^{n} x^2 y^x.
        \]
        对 $y$ 求导,
        \[
        \frac{df}{dy} = \frac{1}{k} \sum_{x=1}^{n} x^3 y^{x-1}, \quad
        \left. \frac{df}{dy} \right|_{y=1} = \frac{1}{k} \sum_{x=1}^{n} x^3.
        \]
        再次对 $y$ 求导,
        \[
        \frac{d^2f}{dy^2} = \frac{1}{k} \sum_{x=1}^{n} x^3 (x-1) y^{x-2} = \frac{1}{k} \sum_{x=1}^{n} x^4 y^{x-2} - \frac{1}{k} \sum_{x=1}^{n} x^3 y^{x-2}.
        \]
        在 $y=1$ 时,
        \[
        \left. \frac{d^2f}{dy^2} \right|_{y=1} + \left. \frac{df}{dy} \right|_{y=1} - \left[ \left. \frac{df}{dy} \right|_{y=1} \right]^2
        = \frac{1}{k} \sum_{x=1}^{n} x^4 - \frac{1}{k^2} \left( \sum_{x=1}^{n} x^3 \right)^2.
        \]
        代入求和公式得
        \begin{align*}
        & \frac{6}{n(n+1)(2n+1)} \cdot \frac{1}{30} n(n+1)(6n^3+9n^2+n-1) - \frac{36}{n^2(n+1)^2(2n+1)^2} \cdot \left[ \frac{n^2(n+1)^2}{4} \right] \\
        &= \frac{6n^3+9n^2+n-1}{5(2n+1)} - \frac{9}{4(2n+1)^2} \\
        &= \frac{3n^4+6n^3-n^2-4n}{20(2n+1)^2}.
        \end{align*}
        因此,
        \[
        \left. \frac{d^2f}{dy^2} \right|_{y=1} + \left. \frac{df}{dy} \right|_{y=1} - \left[ \left. \frac{df}{dy} \right|_{y=1} \right]^2 = \frac{3n^4+6n^3-n^2-4n}{20(2n+1)^2},
        \]
        得证。
    \end{solution}

    \question 
    \begin{parts}
    \part 设 
    \[
    f(x) = x^3 + x,
    \]
    且$g(x)$ 为 $f(x)$ 的反函数,求 $g'(10)$。
    \begin{solution}
        由于 $f$ 与 $g$ 为反函数,$f(g(x)) = x$,对两边求导得到
        \[
        f'(g(x)) g'(x) = 1
        \]
        又因为 $f(2) = 10$,所以 $g(10) = 2$。令 $x = 10$,则
        \[
        f'(2) g'(10) = 1
        \]
        而 $f'(x) = 3x^2 + 1$,所以 $f'(2) = 13$,因此
        \[
        g'(10) = \frac{1}{13}
        \]
    \end{solution}
    \part 对 $x>0$,定义 
    \[
    h(x) = \frac{1}{f(x)},
    \]
    证明函数 $f(x)+h(x)$ 在 $x = g(1)$ 处取得最小值。
    \begin{solution}
        令 
        \[
        F(x) = f(x) + h(x) = f(x) + \frac{1}{f(x)}
        \] 
        则
        \[
        F'(x) = f'(x) - \frac{f'(x)}{[f(x)]^2} = f'(x) \frac{[f(x)]^2 - 1}{[f(x)]^2}
        \] 
        因为 $f'(x) > 1$,所以在临界点有 $f(x_0)^2 - 1 = 0$,即 $f(x_0) = 1$,因此 $x_0 = g(1)$,再次求导,
        \[
        F''(x) = f''(x) - \frac{[f(x)]^2 f''(x) - 2 f(x) [f'(x)]^2}{[f(x)]^4}
        \] 
        在 $x = x_0$ 时,$f(x_0) = 1,f''(x_0) = 6 g(1),f'(x_0) = 3(g(1))^2 + 1$,所以
        \[
        F''(x_0) = 6g(1) - \big(6g(1) - 2 [3(g(1))^2 + 1]^2\big) = 2 [3(g(1))^2 + 1]^2 > 0
        \] 
        因此 $x = x_0$ 为极小值。由于
        \[
        \lim_{x \to 0^+} F(x) = \lim_{x \to \infty} F(x) = \infty,
        \] 
        故 $x = g(1)$ 为 $F(x)$ 的最小值。
    \end{solution}
    \end{parts}
    
    \question 已知对任意正整数 $k$,方程 
    \[
    x^{k}-x^{k-1}-x^{k-2}-\cdots-x-1=0
    \]
    恰有一个正实根,设该方程唯一正实根为 $\alpha_k$,试证该无穷数列 $\langle \alpha_k \rangle$ 收敛,且 $$\lim\limits_{k \to \infty} \alpha_k = 2$$
    \begin{solution}
        原方程可化为
        \[
        x^k - x^{k-1} - x^{k-2} - \cdots - x - 1 = x^k - (x^{k-1} + x^{k-2} + \cdots + 1)
        = x^k - \frac{x^k - 1}{x - 1}
        \]
        令
        \[
        f_k(x) = x^{k+1} - 2x^k + 1
        \]
        则 $\alpha_k$ 为 $f_k(x) = 0$ 在 $(1, 2)$ 内的唯一正实根(注意 $\alpha_k \ne 1$,当 $k > 1$),考虑
        \[
        f_k(0) = 1,\quad f_k(1) = 0,\quad f_k(2) = 1
        \]
        对 $f_k(x)$ 求导得
        \[
        f_k'(x) = (k+1)x^k - 2k x^{k-1} = x^{k-1}[(k+1)x - 2k]
        \]
        导函数零点为
        \[
        x = 0 \quad \text{或} \quad x = \frac{2k}{k+1} = 2 - \frac{2}{k+1}
        \]
        因此
        \[
        f_k(x) \text{ 在 } \left(0, 2 - \frac{2}{k+1} \right) \text{ 上递减,}
        \quad
        \left(2 - \frac{2}{k+1}, \infty\right) \text{ 上递增}
        \]
        由于 $f_k(1) = 0$, $f_k(2) = 1$ 且 $\alpha_k$ 是唯一正实根,结合单调性可得:
        \[
        2 - \frac{2}{k+1} < \alpha_k < 2
        \]
        故由夹挤定理,
        \[
        \lim_{k \to \infty} \alpha_k = 2
        \]
        因此 $\alpha_k$ 收敛,且收敛至 $2$,故得证。
    \end{solution}

    \question 已知 $a < b < c,f'(x)$ 在 $(a, c)$ 上严格递增, 且 $f(x)$ 在 $[a, c]$ 上连续, 证明
    \[
    (b-a) f(c) + (c-b) f(a) > (c-a) f(b).
    \]
    \begin{solution}
        由拉格朗日中值定理, 因此存在 $\alpha$ 和 $\beta$ 使得
        \[
        \frac{f(c)-f(b)}{c-b} = f'(\beta), \quad b < \beta < c
        \]
        以及
        \[
        \frac{f(b)-f(a)}{b-a} = f'(\alpha), \quad a < \alpha < b
        \]
        由于 $f'$ 严格递增, 我们有 $f'(\beta) > f'(\alpha)$, 因此
        \[
        \frac{f(c)-f(b)}{c-b} > \frac{f(b)-f(a)}{b-a}
        \]
        即
        \[
        (b-a) f(c) + (c-b) f(a) > (c-a) f(b)
        \]
    \end{solution}

    \question 设 $f,g:\mathbb{R}\to\mathbb{R}$ 是连续函数,且 $g$ 可导,若
    \[
    (f(0)-g'(0))(g'(1)-f(1))>0,
    \]
    证明存在实数 $c\in(0,1)$,使得 $f(c)=g'(c)$。
    \begin{solution}
        令
        \[
        F(x)=\int_0^x f(t)\,dt, \quad h(x)=F(x)-g(x).
        \]
        由 $f$ 的连续性,$F$ 可导且 $F'=f$,因此
        \[
        h'(x)=F'(x)-g'(x)=f(x)-g'(x).
        \]
        故由$(f(0)-g'(0))(g'(1)-f(1))>0$知
        \[
        h'(0)(-h'(1))>0,
        \]
        因此 $h'(0)$ 和 $h'(1)$ 异号。由达布定理,存在 $c\in(0,1)$使得
        \[
        h'(c)=0 \Rightarrow f(c)=g'(c).
        \]
    \end{solution}

    \question 设 $f: (0,\infty) \to \mathbb{R}$ 为连续可微函数,$b > a > 0$ 且 $f(a) = f(b) = k$。证明存在 $\xi \in (a, b)$ 使得
    \[
    f(\xi) - \xi f'(\xi) = k.
    \]
    \begin{solution}
        若考虑函数 $g(x) = \dfrac{f(x)}{x}$,则
        \[
        g'(x) = \frac{xf'(x) - f(x)}{x^2},
        \]
        其中分子与已知正好对应。考虑在 $[a,b]$ 上的函数 
        \[
        g(x) = \frac{f(x)}{x},\quad h(x) = \frac{1}{x}
        \]
        由柯西中值定理,存在 $\xi \in (a, b)$ 使得
        \[
        \frac{g(a) - g(b)}{h(a) - h(b)} = \frac{g'(\xi)}{h'(\xi)}
        \]
        注意到
        \[
        \frac{g'(\xi)}{h'(\xi)} = \frac{\frac{\xi f'(\xi) - f(\xi)}{\xi^2}}{-\frac{1}{\xi^2}} = f(\xi) - \xi f'(\xi)
        \]
        而
        \[
        \frac{g(a) - g(b)}{h(a) - h(b)} = \frac{\frac{f(a)}{a} - \frac{f(b)}{b}}{\frac{1}{a} - \frac{1}{b}} = k
        \]
        故原命题得证。
    \end{solution}

    \question 设 $f:\mathbb{R}\to\mathbb{R}$ 为二阶可导函数,且 $f(0)=0$。证明:存在
    $\xi\in\left(-\dfrac{\pi}{2},\dfrac{\pi}{2}\right)$使得
    \[
    f''(\xi)=f(\xi)\bigl(1+2\tan^2\xi\bigr).
    \]
    \begin{solution}
        设$g(x)=f(x)\cos x$,由于
        \[
        g\left(-\frac{\pi}{2}\right)=g(0)=g\left(\frac{\pi}{2}\right)=0
        \]
        由罗尔定理,存在
        \[
        \xi_1\in\left(-\frac{\pi}{2},0\right),\quad
        \xi_2\in\left(0,\frac{\pi}{2}\right)
        \]
        使得
        \[
        g'(\xi_1)=g'(\xi_2)=0
        \]
        考虑函数
        \[
        h(x)=\frac{g'(x)}{\cos^2 x}
        =\frac{f'(x)\cos x-f(x)\sin x}{\cos^2 x}
        \]
        显然有
        \[
        h(\xi_1)=h(\xi_2)=0
        \]
        再由罗尔定理,存在
        \[
        \xi\in(\xi_1,\xi_2)\subset\left(-\frac{\pi}{2},\frac{\pi}{2}\right)
        \]
        使得 $h'(\xi)=0$,$h$导数为
        \begin{align*}
        0=h'(\xi)
        &=\frac{g''(\xi)\cos^2\xi+2\cos\xi\sin\xi\,g'(\xi)}{\cos^4\xi} \\
        &=\frac{(f''(\xi)\cos\xi-2f'(\xi)\sin\xi-f(\xi)\cos\xi)\cos\xi
        +2\sin\xi\bigl(f'(\xi)\cos\xi-f(\xi)\sin\xi\bigr)}{\cos^3\xi} \\
        &=\frac{f''(\xi)\cos^2\xi-f(\xi)\bigl(\cos^2\xi+2\sin^2\xi\bigr)}{\cos^3\xi}
        \end{align*}
        整理得
        \[
        0=\frac{1}{\cos\xi}\Bigl(f''(\xi)-f(\xi)\bigl(1+2\tan^2\xi\bigr)\Bigr)
        \]
        因此得证
        \[
        f''(\xi)=f(\xi)\bigl(1+2\tan^2\xi\bigr)
        \]
    \end{solution}

    \question 证明当 $0 < \theta < \dfrac{\pi}{4}$ 时,有
    \[
    \tan \theta < \frac{4\theta}{\pi}.
    \]
    \begin{solution}
        设 
        \[
        f(x) = \tan x,\quad x = \theta \in \left(0, \frac{\pi}{4}\right)
        \]
        因为$f$ 在 $\left[0, \dfrac{\pi}{4}\right]$ 上连续,$\left(0, \dfrac{\pi}{4}\right)$ 上可导,由拉格朗日中值定理,存在
        \[
        0 < \xi < \theta,\theta < \xi' < \frac{\pi}{4}\]
        使得
        \[
        \sec^2 \xi = \frac{\tan \theta - \tan 0}{\theta - 0},\quad \sec^2 \xi' = \frac{\tan\frac{\pi}{4}-\tan \theta}{\frac{\pi}{4}-\theta}
        \]
        由于 $\sec^2 x$ 在 $\left(0, \dfrac{\pi}{4}\right)$ 上单调递增,因此 $\sec^2 \xi' > \sec^2 \xi$,即
        \[
        \frac{1 - \tan \theta}{\frac{\pi}{4} - \theta} > \frac{\tan \theta}{\theta} \Rightarrow \frac{4\theta}{\pi} > \tan \theta
        \]
    \end{solution}
    \begin{solution}
        令
        \[
        f(\theta) = \frac{4\theta}{\pi} - \tan \theta.
        \]
        则
        \[
        f'(\theta) = \frac{4}{\pi} - \sec^2 \theta
        \]
        因此可解得唯一的 $\xi \in \left(0, \dfrac{\pi}{4}\right)$ 使得 $f'(\xi) = 0$,又 
        \[
        f'(\theta) < 0,\; \forall \theta \in (0, \xi),\quad f'(\theta) > 0,\; \forall \theta \in \left(\xi, \frac{\pi}{4}\right)
        \]
        因此$f$ 在 $(0,\xi)$ 上严格递减,在 $\left(\xi, \dfrac{\pi}{4}\right)$ 上严格递增。又 $f(0) = f\left(\dfrac{\pi}{4}\right) = 0$,所以对所有 $\left(0, \dfrac{\pi}{4}\right)$ 皆有 $f(\theta) < 0$ ,即
        \[
        \tan \theta < \frac{4\theta}{\pi}, \quad 0 < \theta < \dfrac{\pi}{4}
        \]
    \end{solution}

    \question 证明对于所有非负整数 $n$,有
    \[
    2(3n-1)^n \ge (3n+1)^n
    \]
    \begin{solution}
        当 $n=0$ 时,不等式显然成立,因为
        \[
        2(3\cdot 0 - 1)^0 = 2 \cdot 1 = 2 \ge 1 = (3\cdot 0 + 1)^0.
        \]
        考虑 $n \ge 1$,命题等价于
        \[
        2(3n-1)^n \ge (3n+1)^n
        \Longleftrightarrow \left(\frac{3n-1}{3n+1}\right)^n \ge \frac{1}{2}
        \Longleftrightarrow n \ln \frac{3n-1}{3n+1} \ge -\ln 2
        \]
        定义函数
        \[
        f(x) = x \ln \frac{3x-1}{3x+1}, \quad x \ge 1.
        \]
        对 $f(x)$ 求导,
        \[
        f'(x) = \ln \frac{3x-1}{3x+1} + \frac{6x}{9x^2 -1},\quad f''(x) = \frac{6(3x^2-1)}{(9x^2-1)^2}
        \]
        当 $x \ge 1$ 时, $f''(x) > 0$, 故 $f'(x)$ 在 $[1, +\infty)$ 上严格递增。
        又因为
        \[
        \lim_{x \to +\infty} f'(x) = \lim_{x \to +\infty} \left( \ln \frac{3-1/x}{3+1/x} + \frac{6/x}{9-1/x^2} \right) = 0
        \]
        由 $f'(x)$ 严格递增且趋于 0 可知，在 $[1, +\infty)$ 上 $f'(x) < 0$。
        因此 $f(x)$ 在 $[1, +\infty)$ 上严格递减。
        根据其极限：
        \[
        \lim_{x \to +\infty} f(x) = \lim_{x \to +\infty} \ln \left( 1 - \frac{2}{3x+1} \right)^x = -\frac{2}{3}
        \]
        由于 $f(x)$ 严格递减且下界为 $-\frac{2}{3}$，而 $-\frac{2}{3} > -0.693 \approx -\ln 2$。
        故对于所有 $n \ge 1$, 有 $f(n) > -\frac{2}{3} > -\ln 2$, 原不等式成立。
        \textcolor{red}{(待验证)}
    \end{solution}

    \question 当 $x \in (0, \frac{\pi}{2})$,比较 $\tan(\sin x)$ 及 $\sin(\tan x)$ 的大小。
    \begin{solution}
        设
        \[
        f(x) = \tan(\sin x) - \sin(\tan x)
        \]
        求导得,
        \[
        f'(x) = \frac{\cos x}{\cos^2(\sin x)} - \frac{\cos(\tan x)}{\cos^2 x} = \frac{\cos^3 x - \cos(\tan x)\cdot \cos^2(\sin x)}{\cos^2 x \cdot \cos^2(\tan x)}
        \]
        对 $0 < x < \arctan \dfrac{\pi}{2}$,由与 $\cos$ 在 $(0, \frac{\pi}{2})$ 上凸,由琴生不等式及AM-GM不等式,
        \[
        \sqrt[3]{\cos(\tan x)\cdot \cos^2(\sin x)} < \frac{1}{3}[\cos(\tan x) + 2\cos(\sin x)] \le \cos\left(\frac{\tan x + 2\sin x}{3}\right) < \cos x
        \]
        其中最后一个不等式由
        \[
        \left(\frac{\tan x + 2\sin x}{3}\right)' = \frac{1}{3}\left(\frac{1}{\cos^2 x} + 2\cos x\right) \ge 1
        \]
        得到,据此有
        \[
        \cos^3 x - \cos(\tan x) \cdot \cos^2(\sin x) > 0 \Rightarrow f'(x) > 0
        \]
        因此 $f$ 在 $\left[0, \arctan \dfrac{\pi}{2}\right]$ 上单调递增。又注意到
        \[
        \tan\left(\sin(\arctan \frac{\pi}{2})\right) = \tan\frac{\frac{\pi}{2}}{\sqrt{1+\frac{\pi^2}{4}}} > \tan\frac{\pi}{4} = 1
        \]
        因此对于 $x \in \left[\arctan \frac{\pi}{2}, \frac{\pi}{2}\right]$,有 $\tan(\sin x) > 1$,所以 $f(x) > 0$。综上,对于所有 $x \in (0, \frac{\pi}{2})$,
        \[
        \tan(\sin x) > \sin(\tan x).
        \]
    \end{solution}

    \question 证明:对半开区间 $\left(0, \dfrac{\pi}{2}\right]$ 内的任意 $x$,有
    \[
    \left(\frac{\sin x}{x}\right)^3 > \cos x.
    \]
    \begin{solution}
        在 $\left(0, \dfrac{\pi}{2}\right]$ 上,有泰勒展开的不等式:
        \[
        \sin x > x - \frac{x^3}{6}, \quad \cos x < 1 - \frac{x^2}{2} + \frac{x^4}{24}
        \]
        因此
        \[
        \left(\frac{\sin x}{x}\right)^3 > \left(1 - \frac{x^2}{6}\right)^3
        \]
        展开右边并与 $\cos x$ 比较:
        \[
        \left(1 - \frac{x^2}{6}\right)^3 = 1 - \frac{x^2}{2} + \frac{x^4}{12} - \frac{x^6}{216} > 1 - \frac{x^2}{2} + \frac{x^4}{24} > \cos x
        \]
        其中
        \[
        0 < x \le \frac{\pi}{2} < 3 \Rightarrow x^2 < 9   \Rightarrow \frac{x^4}{12} - \frac{x^6}{216} > \frac{x^4}{24} 
        \]
        故原不等式得证。
    \end{solution}

    \question 证明不等式:当 $0 < x \le 1$ 时,有
    \[
    \sin x + \arcsin x > 2x.
    \]
    \begin{solution}
        $\sin x$ 与 $\arcsin x$ 的麦克劳林展开式分别为
        \[
        \sin x = x - \frac{x^3}{3!} + \frac{x^5}{5!} - \frac{x^7}{7!} + \cdots
        \]
        \[
        \arcsin x = x + \frac{1}{2}\frac{x^3}{3} + \frac{1\cdot 3}{2\cdot 4}\frac{x^5}{5} + \frac{1\cdot 3\cdot 5}{2\cdot 4\cdot 6}\frac{x^7}{7} + \cdots
        \]
        因此,对任意 $0 < x \le 1$,存在 $\theta_1,\theta_2$,满足 $0<\theta_1<x,0<\theta_2<x$,使得
        \[
        \sin x = x - \frac{x^3}{3!} + \frac{\theta_1^5}{5!},\quad \arcsin x = x + \frac{1}{2}\frac{x^3}{3} + \frac{1\cdot 3}{2\cdot 4}\frac{\theta_2^5}{5}
        \]
        于是
        \[
        \sin x + \arcsin x = 2x + \frac{\theta_1^5}{5!} + \frac{1\cdot 3}{2\cdot 4}\frac{\theta_2^5}{5} > 2x
        \]
        故得证。
    \end{solution}
    \begin{solution}
        设
        \[
        f(x) = \sin x + \arcsin x,
        \]
        则 $f(0)=0$,且
        \[
        f'(x) = \cos x + \frac{1}{\sqrt{1-x^2}},
        \]
        从而
        \[
        f'(0) = 2,
        \]
        故直线 $y=2x$ 是曲线 $y=f(x)$ 在原点处的切线。又
        \[
        f''(x) = -\sin x + \frac{x}{(1-x^2)^{\frac{3}{2}}}.
        \]
        当 $0<x<1$ 时,有
        \[
        \frac{x}{(1-x^2)^{\frac{3}{2}}} > \sin x,
        \]
        因此 $f''(x)>0$,函数 $f(x)$ 在区间 $(0,1)$ 上上凸,即$f(x)$ 在 $(0,1]$ 上位于其在原点处的切线 $y=2x$ 之上,故得证
        \[
        \sin x + \arcsin x > 2x
        \]
    \end{solution}

    \question 设 $f: \mathbb{R} \to \mathbb{R}$ 为二次可微函数,满足 $f(0)=1$, $f'(0)=0$,且对所有 $x \in [0,\infty)$ 有
    \[
    f''(x)-5f'(x)+6f(x)\ge0.
    \]
    证明对所有 $x \in [0,\infty)$ 有
    \[
    f(x) \ge 3 e^{2x} - 2 e^{3x}.
    \]
    \begin{solution}
        设$x \in [0,\infty)$,令 $g(x) = f'(x) - 2f(x)$. 则
        \[
        g'(x) - 3 g(x) \ge 0
        \]
        即
        \[
        \bigl(g(x) e^{-3x}\bigr)' \ge 0
        \]
        因此
        \[
        g(x)e^{-3x} \ge g(0) = f'(0) - 2 f(0) = -2
        \]
        即
        \[
        f'(x) - 2f(x) \ge -2 e^{3x}
        \]
        两边同乘 $e^{-2x}$ 并整理得
        \[
        (f(x)e^{-2x})' \ge -2 e^{x}
        \]
        即
        \[
        (f(x)e^{-2x} + 2 e^x)' \ge 0
        \]
        于是
        \[
        f(x)e^{-2x} + 2 e^x \ge f(0) + 2 = 3
        \]
        即
        \[
        f(x) \ge 3 e^{2x} - 2 e^{3x}
        \]
    \end{solution}

    \question 设函数
    \[
    F(x) = e^{-x} - \left(1-\frac{x}{n}\right)^n,
    \]
    证明:当 $n \ge 2$ 且 $x \in [0,n]$ 时,
    \[
    0 \le F(x) \le \frac{e^{-1}}{n}.
    \]
    \begin{solution}
        函数 $F(x)$ 在区间 $[0,n]$ 上连续且可微,因此在该区间内必有最大值与最小值,且只能出现在端点或临界点处。有$F(0) = 1-1 = 0,F(n) = e^{-n}$,且
        \[
        F'(x) = -e^{-x} + \left(1-\frac{x}{n}\right)^{\,n-1}
        \]
        先证明 $F(x) \ge 0$。只需证明
        \[
        \left(1-\frac{x}{n}\right)^n \le e^{-x}, \quad x \in [0,n]
        \]
        等价于
        \[
        0 \le 1-\frac{x}{n} \le e^{-x/n}, \quad x \in [0,n]
        \]
        令 $t=\frac{x}{n}$,则 $t \in [0,1]$,上述不等式化为
        \[
        1-t \le e^{-t}, \quad t \in [0,1]
        \]
        这是显然成立的,因为直线 $y=1-t$ 是曲线 $y=e^{-t}$ 在 $(0,1)$ 处的切线,而 $y=e^{-t}$ 在$[0,1]$上向下凸,因此图像始终位于其切线之上,故$F(x) \ge 0$,最小值在 $x=0$ 处取得。接下来确定最大值,由于
        \[
        F'(n) = -e^{-n} < 0,
        \]
        最大值不可能出现在 $x=n$,因此必存在 $x_0 \in (0,n)$ 使得
        $F'(x_0)=0$,这意味
        \[
        e^{-x_0} = \left(1-\frac{x_0}{n}\right)^{n-1}
        \]
        于是
        \begin{align*}
        F(x_0)
        &= e^{-x_0}-\left(1-\frac{x_0}{n}\right)^n \\
        &= e^{-x_0}-\left(1-\frac{x_0}{n}\right)^{n-1}\left(1-\frac{x_0}{n}\right) \\
        &= e^{-x_0}-e^{-x_0}\left(1-\frac{x_0}{n}\right) \\
        &= e^{-x_0}\frac{x_0}{n}
        \end{align*}
        设 $g(x)=x e^{-x}$,可推导得知$g(x)$ 在 $x=1$ 处取得最大值 $e^{-1}$,因此
        \[
        F(x_0) \le \frac{e^{-1}}{n}
        \]
        综上所述 $n\ge 2$ 且 $x \in [0,n]$ 时,
        \[
        0 \le F(x) \le \frac{e^{-1}}{n}
        \]
    \end{solution}

    \question 已知在$(-\infty, +\infty)$上具有二阶连续导数的函数$f(x)$满足方程
    \[
    x^2 f''(x)-2x\sin x f'(x) = e^x + e^{-x}-2,
    \]
    若$f(x)$在$x=a$处取极值,问$f(a)$是函数$f(x)$的极大值还是极小值?请说明理由。
    \begin{solution}
        因为 $f(x)$ 在 $x = a$ 处取极值,所以
        \[
        f'(a) = 0.
        \]
        代入原方程得
        \[
        a^2 f''(a) = e^a + e^{-a} - 2.
        \]
        由AM-GM不等式,
        \[
        e^a + e^{-a} \ge 2,
        \]
        等号成立当且仅当 $a = 0$,所以
        \[
        a^2 f''(a)=e^a + e^{-a} - 2 \ge 0,
        \]
        若 \( a \ne 0 \),则 \( a^2 > 0 \),因此
        \[
        f''(a) = \frac{e^a + e^{-a} - 2}{a^2} > 0,
        \]
        所以 \( f(x) \) 在 \( x = a \) 处为极小值;若 \( a = 0 \),则
        \[
        f''(0) = \lim_{x \to 0} \frac{e^x + e^{-x} - 2}{x^2}
        \]
        由洛必达法则,
        \[
        f''(0) = \lim_{x \to 0} \frac{e^x + e^{-x}}{2} = 1 > 0.
        \]
        因此\( f''(0) > 0 \)仍为极小值,故\( f(a) \) 是函数 \( f(x) \) 的极小值。
    \end{solution}

    \question 已知函数 $f(t)$ 在 $[a,x]$ 上可微,且 $f'(t)$ 可微。对所有 $x > a$,存在 $c_x$ 满足 $a < c_x < x$ 且
    \[
    \int_{a}^{x} f(t) \, dt = f(c_x)(x-a).
    \]
    假设 $f'(a) \neq 0$,证明
    \[
    \lim_{x \to a} \frac{c_x - a}{x-a} = \frac{1}{2}.
    \]
    \begin{solution}
        设
        \[
        F(x) = \int_{a}^{x} f(t) \, dt
        \]
        利用 $F(x)$ 的泰勒展开,有
        \[
        F(x) = F(a) + (x-a)F'(a) + \frac{(x-a)^2}{2} F''(\theta_x)
        \]
        其中 $a<\theta_x<x$,且当 $x \to a$ 时,$\theta_x \to a$。又 $F(a) = 0,F'(x) = f(x),F''(x) = f'(x)$,所以
        \[
        F(x) = 0 + (x-a) f(a) + \frac{(x-a)^2}{2} f'(\theta_x)
        \]
        由定义有
        \[
        f(c_x) = \frac{F(x)}{x-a} = f(a) + \frac{x-a}{2} f'(\theta_x)
        \]
        因此
        \[
        \frac{f(c_x) - f(a)}{x-a} = \frac{1}{2} f'(\theta_x)
        \]
        另一方面可以写成
        \[
        \frac{f(c_x) - f(a)}{x-a} = \frac{f(c_x) - f(a)}{c_x - a} \cdot \frac{c_x - a}{x-a}
        \]
        取极限得到
        \[
        \lim_{x \to a} \frac{1}{2} f'(\theta_x) = \lim_{x \to a} \frac{f(c_x) - f(a)}{c_x - a} \cdot \lim_{x \to a} \frac{c_x - a}{x - a}
        \]
        由此可得
        \[
        \frac{1}{2} f'(a) = f'(a) \cdot \lim_{x \to a} \frac{c_x - a}{x - a}
        \]
        即
        \[
        \lim_{x \to a} \frac{c_x - a}{x - a} = \frac{1}{2}
        \]
    \end{solution}

    \question 
    \begin{parts}
    \part 已知 $f(x)$ 在 $[0,+\infty)$ 上连续且单调增加,且 $f(0) \ge 0$,证明:
    \[
    F(x) = \begin{cases}
    \dfrac{1}{x^n} \displaystyle\int_0^x t^{n-1} f(t)\,dt, & x > 0, \\
    0, & x = 0
    \end{cases}
    \]
    在 $[0,+\infty)$ 上连续且单调增加,其中 $n > 0$。
    \begin{solution}
        先证连续性。由于 \( f \) 在 \( [0, +\infty) \) 上连续,且 \( f(0) \ge 0 \),有:
        \[
        \lim_{x \to 0^+} F(x) = \lim_{x \to 0^+} \frac{1}{x^n} \int_0^x t^{n-1} f(t)\,dt = 0 = F(0),
        \]
        所以 \( F(x) \) 在 \( x = 0 \) 处连续。

        再证单调性。对于 \( x > 0 \),由积分中值定理可得:
        \[
        \int_0^x t^{n-1} f(t)\,dt = f(\xi) \int_0^x t^{n-1}\,dt = f(\xi) \cdot \frac{x^n}{n}, \quad \text{其中 } \xi \in (0, x).
        \]
        因此,
        \[
        F(x) = \frac{1}{x^n} \cdot \int_0^x t^{n-1} f(t)\,dt = \frac{1}{x^n} \cdot f(\xi) \cdot \frac{x^n}{n} = \frac{f(\xi)}{n}.
        \]
        因为 \( f \) 单调递增,\( \xi \in (0,x) \),所以 \( x \uparrow \Rightarrow \xi \uparrow \Rightarrow f(\xi) \uparrow \),
        故 \( F(x) \) 单调递增。
    \end{solution}
    \part 设 \( f(x) \) 在 \( [0,1] \) 上二阶可导,且满足 \( |f''(x)| \le 1 \)。已知 \( f(x) \) 在 \( (0,1) \) 内取最大值为 \( \dfrac{1}{4} \),证明:
    \[
    |f(0)| + |f(1)| \le 1.
    \]
    \begin{solution}
        设 \( x_0 \in (0,1) \),使得 \( f(x_0) = \dfrac{1}{4} \),且 \( f'(x_0) = 0 \)。
        对 \( f(0) \) 做 Taylor 展开（Lagrange 余项形式）得:
        \[
        f(0) = f(x_0) + f'(x_0)(0 - x_0) + \frac{1}{2} f''(\xi) (0 - x_0)^2 = \frac{1}{4} + \frac{1}{2} f''(\xi) x_0^2,
        \]
        其中 \( \xi \in (0, x_0) \)。
        同理,
        \[
        f(1) = \frac{1}{4} + \frac{1}{2} f''(\eta)(1 - x_0)^2, \quad \eta \in (x_0, 1).
        \]
        因为 \( |f''(x)| \le 1 \),故
        \[
        |f(0)| \le \left| \frac{1}{4} + \frac{1}{2} f''(\xi) x_0^2 \right| \le \frac{1}{4} + \frac{1}{2} x_0^2,
        \]
        \[
        |f(1)| \le \left| \frac{1}{4} + \frac{1}{2} f''(\eta)(1 - x_0)^2 \right| \le \frac{1}{4} + \frac{1}{2} (1 - x_0)^2.
        \]
        相加得:
        \[
        |f(0)| + |f(1)| \le \frac{1}{2} + \frac{1}{2} \left( x_0^2 + (1 - x_0)^2 \right).
        \]
        注意 \( x_0^2 + (1 - x_0)^2 = 1 - 2x_0(1 - x_0) \le 1 \),因此:
        \[
        |f(0)| + |f(1)| \le \frac{1}{2} + \frac{1}{2} \cdot 1 = 1.
        \]
    \end{solution}
    \end{parts}

    \question 设函数 \( f(x) \) 在 \( x = 0 \) 处可导,且 \( f(0) = 0, f'(0) = 1 \),设
    \[
    F(x) = \int_0^x t^{n-1} f(x^n - t^n) \, dt,
    \]
    求
    \[
    \lim_{x \to 0} \frac{F(x)}{x^{2n}}.
    \]
    \begin{solution}
        由题设 \( f(0) = 0 \),且 \( f \) 在 0 可导,考虑换元 \( t = x u \),则
        \[
        F(x) = \int_0^x t^{n-1} f(x^n - t^n) dt = x^n \int_0^1 u^{n-1} f(x^n - x^n u^n) du.
        \]
        注意 \( x^n - x^n u^n = x^n(1 - u^n) \),故
        \[
        F(x) = x^n \int_0^1 u^{n-1} f\big(x^n(1 - u^n)\big) du.
        \]
        因 \( f(0) = 0 \),且 \( f \) 在 0 处可导,令 \( h = x^n(1 - u^n) \),有
        \[
        f(h) = f(0) + f'(0) h + o(h) = f'(0)x^n(1 - u^n) + o(x^n).
        \]
        所以
        \[
        F(x) = x^n \int_0^1 u^{n-1} \left( f'(0)x^n(1 - u^n) + o(x^n) \right) du
        = f'(0)x^{2n} \int_0^1 u^{n-1}(1 - u^n) du + o(x^{2n}).
        \]
        计算积分:
        \[
        \int_0^1 u^{n-1}(1 - u^n) du = \int_0^1 u^{n-1} du - \int_0^1 u^{2n-1} du = \frac{1}{n} - \frac{1}{2n} = \frac{1}{2n}.
        \]
        故
        \[
        \lim_{x \to 0} \frac{F(x)}{x^{2n}} = \frac{1}{2n}
        \]
    \end{solution}

\end{questions}

\pagebreak

\begin{center}
  {\fontsize{30pt}{26pt}\selectfont
    \hypertarget{积分}{积分} \label{积分}
  }
\end{center}
\separator
\vspace{1pt}

\begin{questions}
    \question 已知\( p > 0 \),证明
    \[
    \lim_{n \to \infty} \frac{1^p + 2^p + \cdots + n^p}{n^{p+1}} = \frac{1}{p + 1}.
    \]
    \begin{solution}
    设 $$S_n = \sum_{k=1}^{n} k^p$$则
    \[
    \frac{S_n}{n^{p+1}} = \frac{1}{n} \sum_{k=1}^{n} \left( \frac{k}{n} \right)^p.
    \]
    这是函数 $f(x) = x^p$ 在 $[0,1]$ 的黎曼和,故
    \[
    \lim_{n \to \infty} \frac{S_n}{n^{p+1}} = \int_0^1 x^p \, dx = \frac{1}{p+1}.
    \]
    \end{solution}
    
    \question 求
    \[
    \lim_{k \to \infty} \left( \frac{2017^{\frac{1}{k}}}{k + 1} + \frac{2017^{\frac{2}{k}}}{k + \frac{1}{2}} + \cdots + \frac{2017^{\frac{k}{k}}}{k + \frac{1}{k}} \right)
    \]
    \begin{solution}
        写成黎曼和:
        \[
        \lim_{k \to \infty} \frac{1}{k} \sum_{j=1}^{k} 2017^{\frac{j}{k}} = \int_0^1 2017^x \, dx = \left[ \frac{2017^x}{\ln 2017} \right]_0^1 = \frac{2016}{\ln 2017}.
        \]
        \textcolor{red}{(待验证,$2017^{\frac{j}{k}}/(1+0?)$)}
    \end{solution}

    \question 设
    \[
    a_{n}=\frac{2}{n}\Big[(2^{2}+1)+\left(2+\frac{2}{n}\right)^{2}+1 + \cdots + \left(2+\frac{2n-2}{n}\right)^{2}+1\Big] 
    \]
    求 $\displaystyle \lim_{n\to\infty} a_n$
    \begin{solution}
        发现
        \[
        a_n = \sum_{k=1}^n \frac{2}{n} \left[ \left( 2 + \frac{2k-2}{n} \right)^2 + 1 \right].
        \]
        是函数 $f(x) = (2+x)^2+1$ 在 $[0,2]$ 上的黎曼和,故
        \[
        \lim_{n\to\infty} a_n = \int_0^2 \left[ (2+x)^2 + 1 \right] \, dx = \left[ \frac{x^3}{3} + 2x^2 +5x \right]_0^2  = \frac{62}{3}.
        \]
    \end{solution}

    \question 试求
    \[
    \lim_{n\to \infty} \frac{5}{n^2} \left[ \sqrt{4n^2 - 2\cdot 1^2} + \sqrt{4n^2 - 2\cdot 2^2} + \cdots + \sqrt{4n^2 - 2\cdot n^2} \right]
    \]
    \begin{solution}
        将其写成黎曼和,得
        \[
        \lim_{n\to \infty} \sum_{k=1}^n \frac{5}{n} \sqrt{4 - 2 \left(\frac{k}{n}\right)^2} 
        = \int_0^1 5 \sqrt{4 - 2x^2} \, dx.
        \]
        可得原极限为
        \[
        \left[ \frac{5}{\sqrt{2}} \left( x \sqrt{2 - x^2} + 2 \sin^{-1} \frac{x}{\sqrt{2}} \right) \right]_0^1
        = \frac{5\sqrt{2} (\pi + 2)}{4}.
        \]
    \end{solution}

    \question 设函数 $f:(0,\infty)\rightarrow\mathbb{R}$ ,且当 $x>0$ 时,恒有
    \[
    \lim_{n\rightarrow\infty}\frac{1}{3n-1}\left(f\left(\frac{x^{2}}{n}\right)+f\left(\frac{2x^{2}}{n}\right)+\cdot\cdot\cdot+f\left(\frac{nx^{2}}{n}\right)\right)=\frac{\sqrt[3]{7+x}}{x},
    \]
    求 $f(1)$。
    \begin{solution}
        先写成黎曼和:
        \begin{align*}
        &\lim_{n\to\infty} \frac{1}{3n-1} \left( f\left(\frac{x^2}{n}\right) + f\left(\frac{2x^2}{n}\right) + \cdots + f\left(\frac{nx^2}{n}\right) \right) \\
        &= \lim_{n\to \infty} \frac{1}{3n-1} \sum_{k=1}^n f\left(\frac{kx^2}{n}\right) \\
        &= \lim_{n\to \infty} \frac{n}{3n-1} \sum_{k=1}^n \frac{1}{n} f\left(\frac{kx^2}{n}\right) \\
        &= \frac{1}{3} \int_0^1 f(x^2 t)\, dt = \frac{\sqrt[3]{7+x}}{x}.
        \end{align*}
        于是
        \[
        \int_0^1 x^2 f(x^2 t)\, dt = 3x \sqrt[3]{7+x}.
        \]
        取 $u = x^2 t \Rightarrow du = x^2 dt$,则
        \begin{align*}
        \int_0^{x^2} f(u)\, du &= 3x \sqrt[3]{7+x}, \\
        \frac{d}{dx} \left( \int_0^{x^2} f(u)\, du \right) &= \frac{d}{dx} \left( 3x \sqrt[3]{7+x} \right), \\
        2x f(x^2) &= 3 \sqrt[3]{7+x} + x (7+x)^{-\frac{2}{3}}.
        \end{align*}
        令 $x=1$,则
        \[
        2 f(1) = 3\cdot 2 + \frac{1}{4} \quad \Rightarrow \quad f(1) = \frac{25}{8}.
        \]
    \end{solution}

    \question 
    \[
    \lim_{n\to\infty} \frac{1}{n} \left( \sin \frac{\pi}{12n} + \sin \frac{3\pi}{12n} + \sin \frac{5\pi}{12n} + \dots + \sin \frac{(2n-1)\pi}{12n} \right) 
    \]
    \begin{solution}
        写成黎曼和,
        \begin{align*}
        \lim_{n\to \infty} \frac{1}{n} \sum_{k=1}^n \sin \frac{(2k-1)\pi}{12n}
        &= \lim_{n\to \infty} \frac{1}{n} \left( \sum_{k=1}^{2n} \sin \frac{k\pi}{12n} - \sum_{k=1}^{n} \sin \frac{2k\pi}{12n} \right)\\
        &= \int_0^2 \sin \frac{\pi}{12}x \, dx - \int_0^1 \sin \frac{\pi}{6}x \, dx
        \end{align*}
        可得
        \[
        \left[ -\frac{12}{\pi} \cos \frac{\pi}{12}x \right]_0^2 - \left[ -\frac{6}{\pi} \cos \frac{\pi}{6}x \right]_0^1 = \frac{6 - 3\sqrt{3}}{\pi}
        \]
    \end{solution}

    \question 求
    \[
    f(x)=\lim_{n\to\infty}\sum_{k=1}^{n} 
    \frac{x(1-x)}{k+(n-k)x}\;,\,x\in[0,1]
    \]
    \begin{solution} 
    显然$f(0)=f(1)=0,$对于$x \in [0,1],$
        \begin{align*} 
        \lim_{n\to\infty}\sum_{k=1}^{n} 
        \frac{x(1-x)}{k+(n-k)x}\
        &=x(1-x)\lim_{n\to\infty}\sum_{k=1}^{n} 
        \frac{1}{(\frac{k}{n})(1-x)+x}\cdot\frac{1}{n}\\ 
        &=x(1-x)\int_0^1 \frac{1}{(1-x)y+x}dy\\ 
        &=x\cdot\left[\ln|(1-x)y+x|\right]_0^1\\ 
        &=-x\ln(x) 
        \end{align*}
    \end{solution}

    \question 求极限
    \[
    \lim_{n\to\infty} \left( \sum_{k=10}^{n+9} \frac{2^{11(k-9)/n}}{\log_2 e^{n/11}} - \sum_{k=0}^{n-1} \frac{58}{\pi \sqrt{(n-k)(n+k)}} \right)
    \]
    \begin{solution}
        首先有
        \begin{align*}
        \lim_{n\to\infty} \sum_{k=10}^{n+9} \frac{2^{11(k-9)/n}}{\log_2 e^{n/11}}
        = \lim_{n \to \infty} \sum_{k=1}^{n} 2^{11(\frac{k}{n})} \ln\left(\frac{11}{n}\right) \ln 2
        =\int_{0}^{11} 2^{x} \ln 2\, dx
        \end{align*}
        且
        \[
        \lim_{n\to\infty} \sum_{k=0}^{n-1} \frac{58}{\pi \sqrt{(n-k)(n+k)}}
        = \frac{58}{\pi} \lim_{n\to\infty} \sum_{k=0}^{n-1} \frac{1}{n\sqrt{1 - (\frac{k}{n})^2}}
        = \frac{58}{\pi} \int_{0}^{1} \frac{1}{\sqrt{1-x^2}} \, dx
        \]
        故所求为
        \[
        \int_{0}^{11} 2^{x} \ln 2\, dx-\frac{58}{\pi}\int_{0}^{1} \frac{1}{\sqrt{1-x^2}} \, dx=\left[2^{x} \right]_{0}^{11}-\frac{58}{\pi}\left[ \arcsin x \right]_{0}^{1}=2047-29=2018
        \]
    \end{solution}

    \question 求极限
\[
\lim_{t\nearrow 1}(1-t)\sum_{n=1}^{\infty}\frac{t^n}{1+t^n},
\]
其中 $t\nearrow 1$ 表示 $t$ 从小于 $1$ 的一侧趋近于 $1$。

\begin{solution}
注意到
\[
\lim_{t\to 1^-}\frac{1-t}{-\ln t}=1,
\]
因此
\[
\lim_{t\to 1^-}(1-t)\sum_{n=1}^{\infty}\frac{t^n}{1+t^n}
=\lim_{t\to 1^-}(-\ln t)\sum_{n=1}^{\infty}\frac{t^n}{1+t^n}。
\]

将 $t^n=e^{-n(-\ln t)}$,可得
\[
(-\ln t)\sum_{n=1}^{\infty}\frac{1}{1+e^{-n\ln t}}。
\]
设 $h=-\ln t$,则当 $t\to 1^-$ 时 $h\to 0^+$,于是上述极限化为
\[
\lim_{h\to 0^+}h\sum_{n=1}^{\infty}\frac{1}{1+e^{nh}}。
\]

这是一个黎曼和,对应的积分为
\[
\int_0^{\infty}\frac{dx}{1+e^x}。
\]
计算该积分得
\[
\int_0^{\infty}\frac{dx}{1+e^x}=\ln 2。
\]
\end{solution}

\question 对于 $n = 1,2,\dots$,定义
\[
S_n = \log \left(\sqrt{1^1 \cdot 2^2 \cdot \dots \cdot n^n}\right) - \log(\sqrt{n}),
\]
其中 $\log$ 表示自然对数。求 $\lim_{n \to \infty} S_n$。

\begin{solution}
将 $S_n$ 变形为
\begin{align*}
S_n &= \frac{1}{2}\sum_{k=1}^n k \log k - \frac{1}{2}\log n \\
&= \frac{1}{n^2}\sum_{k=1}^n k \log k - \frac{1}{2}\log n \\
&= \frac{1}{n}\sum_{k=1}^n \frac{k}{n}\left(\log \frac{k}{n} + \log n\right) - \frac{1}{2}\log n \\
&= \frac{1}{n}\sum_{k=1}^n \frac{k}{n}\log\frac{k}{n} + \frac{\log n}{n^2}\sum_{k=1}^n k - \frac{1}{2}\log n \\
&= \frac{1}{n}\sum_{k=1}^n \frac{k}{n}\log\frac{k}{n} + \frac{\log n}{n^2}\frac{n(n+1)}{2} - \frac{1}{2}\log n \\
&= \frac{1}{n}\sum_{k=1}^n \frac{k}{n}\log\frac{k}{n} + \frac{(n+1)\log n}{2n} - \frac{1}{2}\log n \\
&= \frac{1}{n}\sum_{k=1}^n \frac{k}{n}\log\frac{k}{n} + \frac{\log n}{2n}.
\end{align*}

注意最后一项 $\frac{\log n}{2n} \to 0$。而 $\frac{1}{n}\sum_{k=1}^n \frac{k}{n}\log\frac{k}{n}$ 是可积函数 $f(x) = x \log x$ 在区间 $[0,1]$ 上的黎曼和。因此
\[
\lim_{n \to \infty} \frac{1}{n}\sum_{k=1}^n \frac{k}{n}\log\frac{k}{n} = \int_0^1 x \log x \, dx = \left[\frac{x^2}{2}\log x - \frac{x^2}{4}\right]_0^1 = -\frac{1}{4}.
\]

综上,极限存在,并且
\[
\lim_{n \to \infty} S_n = -\frac{1}{4}.
\]
\end{solution}


    \question 求极限
    \[
    \lim_{n\rightarrow\infty}\frac{(1^{2}+2^{2}+\cdots+n^{2})(1^{5}+2^{5}+\cdots+n^{5})}{(1^{3}+2^{3}+\cdots+n^{3})(1^{4}+2^{4}+\cdots+n^{4})}
    \]
    \begin{solution}
        首先看到
        \[
        a_n=\frac{\displaystyle\left(\sum_{k=1}^{n} \left(\frac{k}{n}\right)^2\right)
        \left(\sum_{k=1}^{n} \left(\frac{k}{n}\right)^5\right)}
        {\displaystyle \left(\sum_{k=1}^{n} \left(\frac{k}{n}\right)^3\right)
        \left(\sum_{k=1}^{n} \left(\frac{k}{n}\right)^4\right)}
        =
        \frac{\displaystyle\left(\frac{1}{n}\sum_{k=1}^{n} \left(\frac{k}{n}\right)^2\right)
        \left(\frac{1}{n}\sum_{k=1}^{n} \left(\frac{k}{n}\right)^5\right)}
        {\displaystyle\left(\frac{1}{n}\sum_{k=1}^{n} \left(\frac{k}{n}\right)^3\right)
        \left(\frac{1}{n}\sum_{k=1}^{n} \left(\frac{k}{n}\right)^4\right)}
        \]
        故
        \[
        \lim_{n\to\infty} a_n
        =\frac{\displaystyle\left(\int_0^1 x^2 dx\right)\left(\int_0^1 x^5 dx\right)}
        {\displaystyle\left(\int_0^1 x^3 dx\right)\left(\int_0^1 x^4 dx\right)}
        =\frac{\dfrac{1}{3}\cdot \dfrac{1}{6}}{\dfrac{1}{4}\cdot \dfrac{1}{5}}
        =\frac{10}{9}
        \]
    \end{solution}

    \question 计算极限
    \[
    \lim_{n\to\infty} \frac{\left(\sqrt{1}+\sqrt{2}+\cdots+\sqrt{n}\right)^{2}\left(1^{3}+2^{3}+\cdots+n^{3}\right)}{\left(\sqrt[3]{1}+\sqrt[3]{2}+\cdots+\sqrt[3]{n}\right)^{3}\left(1^{2}+2^{2}+\cdots+n^{2}\right)}
    \]
    \begin{solution}
        由
        \[
        \sum_{k=1}^n k^3 = \left(\frac{n(n+1)}{2}\right)^2,\quad \sum_{k=1}^n k^2 = \frac{n(n+1)(2n+1)}{6}
        \]
        可得
        \[
        f(n)=\frac{\displaystyle n^2\left(\sum_{k=1}^n \frac{1}{n}\sqrt{\frac{k}{n}}\right)^2}{\displaystyle n^3\left(\sum_{k=1}^n \frac{1}{n}\sqrt[3]{\frac{k}{n}}\right)^3} \cdot \frac{3n(n+1)}{2(2n+1)}
        \]
        由积分定义,
        \[
        \lim_{n\to\infty} \sum_{k=1}^n \frac{1}{n}\sqrt{\frac{k}{n}} = \int_0^1 \sqrt{x}\,dx = \frac{2}{3},\quad
        \lim_{n\to\infty} \sum_{k=1}^n \frac{1}{n}\sqrt[3]{\frac{k}{n}} = \int_0^1 x^{\frac{1}{3}}\,dx = \frac{3}{4}
        \]
        因此
        \[
        \lim_{n\to\infty} f(n) = \frac{\left(\frac{2}{3}\right)^2}{\left(\frac{3}{4}\right)^3} \cdot \frac{3}{4} = \frac{64}{81}
        \]
    \end{solution}

    \question 求极限
    \[
    L=\lim_{n\to\infty}\left(\frac{(2n)!}{n!\,n^{n}}\right)^{\frac{1}{n}}
    \]
    \begin{solution}
        首先有
        \[
        \ln L=\lim_{n\to\infty}\frac{1}{n}\Big(\ln(2n)!-\ln n! - n\ln n\Big)
        \]
        其中
        \[
        \ln(2n)!-\ln n! = \sum_{k=n+1}^{2n}\ln k=\sum_{k=1}^{n}\ln\!\big(n+k\big)
        \]
        因此
        \[
        \ln L=\lim_{n\to\infty}\frac{1}{n}\left(\sum_{k=1}^{n}\ln\!\big(n+k\big)-n\ln n\right)
        =\lim_{n\to\infty}\frac{1}{n}\sum_{k=1}^{n}\ln\!\Big(1+\frac{k}{n}\Big)
        \]
        这是在 $[0,1]$ 上的黎曼和,得
        \[
        \ln L=\int_0^1 \ln(1+x)\,dx
        =\big[(1+x)\ln(1+x)-x\big]_0^1=2\ln 2-1
        \]
        所以
        \[
        L=e^{2\ln 2-1}=\frac{4}{e}
        \]
    \end{solution}

    \question 设 
    \[
    A_n = \{\sin(\ln 1), \sin(\ln 2), \dots, \sin(\ln n)\}, \quad B_n = \{\cos(\ln 1), \cos(\ln 2), \dots, \cos(\ln n)\}
    \]
    且 $\bar{a}_n,\bar{b}_n$ 分别为 $A_n,B_n$ 的算术平均值。求 
    \[
    \lim_{n\to\infty} (\bar{a}_n \cos(\ln n) - \bar{b}_n \sin(\ln n)).
    \]
    \begin{solution}
        由积分定义,
        \begin{align*}
        &\lim_{n \to \infty} \frac{1}{n} \left( \sum_{i=1}^n (\sin(\ln i)\cos(\ln n) - \cos(\ln i)\sin(\ln n)) \right) \\
        &=\lim_{n\to\infty} \frac{1}{n} \left(\sum_{i=1}^{n} \sin(\ln i - \ln n)\right) \\
        &= \lim_{n\to\infty} \frac{1}{n} \sum_{i=1}^{n} \sin\left(\ln \left(\frac{i}{n}\right)\right) \\
        &= \int_0^1 \sin(\ln x) \, dx \\
        &= \Im \int_0^1 (e^{i\ln x}) \, dx \\
        &= \Im \int_0^1 x^i \, dx \\
        &= \Im \left(\frac{1}{i+1}\right) = -\frac{1}{2}
        \end{align*}
    \end{solution}

    \question 设函数 $f(x)$ 在闭区间 $[0, 1]$ 上具有连续导数,$f(0) = 0,f(1) = 1$。证明:
\[
\lim_{n \to \infty} n \left( \int_{0}^{1} f(x) dx - \frac{1}{n} \sum_{k=1}^{n} f\left( \frac{k}{n} \right) \right) = -\frac{1}{2}.
\]

\begin{solution}
考察括号中的项,将其写为累加形式:
\begin{align*}
\int_{0}^{1} f(x) dx - \frac{1}{n} \sum_{k=1}^{n} f\left( \frac{k}{n} \right) &= \sum_{k=1}^{n} \left( \int_{\frac{k-1}{n}}^{\frac{k}{n}} f(x) dx - \frac{1}{n} f\left( \frac{k}{n} \right) \right)
\end{align*}
在区间 $[\frac{k-1}{n}, \frac{k}{n}]$ 上对 $f(x)$ 在 $\frac{k}{n}$ 处进行一阶泰勒展开:
\[
f(x) = f\left( \frac{k}{n} \right) + f'\left( \frac{k}{n} \right) \left( x - \frac{k}{n} \right) + O\left( \left( x - \frac{k}{n} \right)^2 \right)
\]
代入积分项中:
\begin{align*}
&\sum_{k=1}^{n} \left( \int_{\frac{k-1}{n}}^{\frac{k}{n}} \left( f\left( \frac{k}{n} \right) + f'\left( \frac{k}{n} \right) \left( x - \frac{k}{n} \right) \right) dx - \frac{1}{n} f\left( \frac{k}{n} \right) \right) \\
&= \sum_{k=1}^{n} \left( \left[ f\left( \frac{k}{n} \right) x + \frac{f'\left( \frac{k}{n} \right)}{2} \left( x - \frac{k}{n} \right)^2 \right]_{\frac{k-1}{n}}^{\frac{k}{n}} - \frac{1}{n} f\left( \frac{k}{n} \right) \right) \\
&= \sum_{k=1}^{n} \left( \frac{1}{n} f\left( \frac{k}{n} \right) - \frac{f'\left( \frac{k}{n} \right)}{2n^2} - \frac{1}{n} f\left( \frac{k}{n} \right) \right) \\
&= -\sum_{k=1}^{n} \frac{f'\left( \frac{k}{n} \right)}{2n^2}
\end{align*}
将此结果代入原极限式中:
\begin{align*}
\lim_{n \to \infty} n \left( -\sum_{k=1}^{n} \frac{f'\left( \frac{k}{n} \right)}{2n^2} \right) &= \lim_{n \to \infty} -\frac{1}{2} \left( \frac{1}{n} \sum_{k=1}^{n} f'\left( \frac{k}{n} \right) \right)
\end{align*}
根据黎曼和的定义,该和式收敛于 $f'(x)$ 在 $[0, 1]$ 上的定积分:
\begin{align*}
-\frac{1}{2} \int_{0}^{1} f'(x) dx &= -\frac{1}{2} [f(x)]_0^1 \\
&= -\frac{1}{2} (f(1) - f(0)) \\
&= -\frac{1}{2} (1 - 0) \\
&= -\frac{1}{2}
\end{align*}
证毕。
\end{solution}

\question
3) 证明:$\frac{\pi}{4} < \int_{0}^{1} \frac{1}{1+x^8} dx < 1$
\begin{solution}
在 $x \in (0,1)$ 上,$0 < x^8 < 1$:
\begin{itemize}
    \item 因为 $1+x^8 > 1$,所以 $\frac{1}{1+x^8} < 1 \implies \int_{0}^{1} \frac{1}{1+x^8} dx < \int_{0}^{1} 1 dx = 1$。
    \item 因为 $x^8 < x^2$,所以 $1+x^8 < 1+x^2 \implies \frac{1}{1+x^8} > \frac{1}{1+x^2}$。
    积分得 $\int_{0}^{1} \frac{dx}{1+x^8} > \int_{0}^{1} \frac{dx}{1+x^2} = [\tan^{-1}x]_0^1 = \frac{\pi}{4}$。
\end{itemize}
\end{solution}
    
    \question 若 $f(x)$ 为一实系数多项式函数,已知
    \[
    \int_{0}^{1} f(x) f'(x) \, dx = 1, \quad \int_{0}^{1} (f(x))^3 f'(x) \, dx = 2,
    \]
    求
    \[
    \int_{0}^{1} (f(x))^5 f'(x) \, dx
    \] 的值。
    \begin{solution}
        有
        \[
        1=\int_0^1 f(x) f'(x) \, dx = \left[ \frac{1}{2} (f(x))^2 \right]_0^1 = \frac{1}{2} \big([f(1)]^2 - [f(0)]^2\big) \tag{1}
        \]
        同理,
        \[
        2=\int_0^1 (f(x))^3 f'(x) \, dx = \left[ \frac{1}{4} (f(x))^4 \right]_0^1 = \frac{1}{4} \big([f(1)]^4 - [f(0)]^4\big) \tag{2}
        \]
        由 $(1),(2)$ 解得
        \[
        [f(0)]^2 = 1,\quad [f(1)]^2 = 3
        \]
        于是
        \[
        \int_0^1 (f(x))^5 f'(x) \, dx = \left[ \frac{1}{6} (f(x))^6 \right]_0^1 = \frac{1}{6} \big(f(1)^6 - f(0)^6\big) = \frac{13}{3}
        \]
    \end{solution}

    \question 设 $f$ 为实数系上的连续函数, 若
    \[
    \frac{d}{dx}\int_{-x}^{x}f(t)\,dt=x^{2}+1,
    \] 
    求 
    \[
    \int_{-1}^{1}f(x)\,dx
    \]
    \begin{solution}
        两边积分得
        \[
        g(x)=\int_{-x}^x f(t)\,dt =\int (x^2+1)\,dx =\frac{1}{3}x^3+x+ C
        \]
        令$x=0$,得
        \[
        g(0)= \int_0^0 f(t)\,dt= 0 \Rightarrow C=0 \Rightarrow g(x) =\frac{1}{3}x^3+x
        \]
        所以
        \[
        \int_{-1}^1 f(x)\,dx =g(1)=\frac{4}{3}
        \]
    \end{solution}

    \question 已知一连续实函数 $f(x)$ 满足 $f(2x)=3f(x),\forall x \in \mathbb{R},$若 $\displaystyle\int_{0}^{1} f(x) dx=1$,求 
    \[
    \int_{0}^{2} f(x) dx
    \]
    \begin{solution}
        由 $f(2x)=3f(x)$ 得
        \[
        \int_0^1 f(x)\;dx = \frac{1}{3} \int_0^1 f(2x)\;dx =1
        \]
        令 $u=2x$,则 $du=2dx$,于是
        \[
        \int_0^2 f(u)\;du =6
        \]
        故
        \[
        \int_0^2 f(x)\;dx = \int_0^1 f(x)\;dx + \int_1^2 f(x)\;dx =1 + \int_1^2 f(x)\;dx= 6 \Rightarrow \int_1^2 f(x)\;dx =5
        \]
    \end{solution}

    \question 设多项式函数 $y=f(x)$ 满足
    \[
    f(x) = 4x^3 - 12x^2 + 8x + 20 - \int_1^x f(t) \, dt,
    \]
    求 $f(x)$。
    \begin{solution}
        对两边求导得
        \[
        f'(x) + f(x) = 12x^2 - 24x + 8
        \]
        故 $f(x)$ 为二次多项式,设
        \[
        f(x) = a x^2 + b x + c,f'(x) = 2 a x + b
        \]
        代入得
        \[
        f(x) + f'(x) = a x^2 + b x + c + 2 a x + b = a x^2 + (b + 2 a) x + (b + c)
        \]
        比较系数,解得
        \[
        a = 12,  b = -48,  c = 56 \Rightarrow f(x) = 12 x^2 - 48 x + 56
        \]
    \end{solution}

    \question 设 
    \[
    f(x)=\int_{0}^{x} g(t) \, dt + 1, \quad g(x) = 12x^{2} - 6x + \int_{0}^{1} [f(t) + g'(t)] \, dt,
    \]
    求 $g(0)$。
    \begin{solution}
        由题意有
        \[
        f(0) = 1,\quad g(0) = \int_0^1 [f(t) + g'(t)] \, dt,
        \]
        又
        \[
        f(x) = \int_0^x [12t^{2} - 6t + g(0)] \, dt + 1 = 4x^{3} - 3x^{2} + x g(0) + 1
        \]
        代入 $g(x)$ 得
        \[
        g(x) = 12x^{2} - 6x + \int_0^1 (4t^{3} - 3t^{2} + t g(0) + 1) \, dt + g(1) - g(0)
        \]
        其中
        \[
        \int_0^1 (4t^{3} - 3t^{2} + t g(0) + 1) dt = \left[ t^{4} - t^{3} + \frac{1}{2} g(0) t^{2} + t \right]_0^1 = \frac{1}{2} g(0) + 1
        \]
        因此
        \[
        g(x) = 12x^{2} - 6x + 1 + g(1) - \frac{1}{2} g(0)
        \]
        令 $x=1$,得
        \[
        g(0) = 14
        \]
    \end{solution}

\question 设
\[
f(x)=x+3+\int_{0}^{x}g(t)\,dt,\qquad
g(x)=2x-9+\int_{0}^{x}f(t)\,dt,
\]
试求 $f(3)$。
\begin{solution}
设
\[
a=\int_0^1 g(x)\,dx,\qquad b=\int_0^2 f(x)\,dx.
\]
由题得
\[
f(x)=x+3+a,\qquad g(x)=2x-9+b.
\]
因此
\[
b=\int_0^2 f(x)\,dx=\int_0^2 (x+3+a)\,dx
=\Big[\tfrac{1}{2}x^2+(3+a)x\Big]_0^2=8+2a,
\]
\[
a=\int_0^1 g(x)\,dx=\int_0^1 (2x-9+b)\,dx
=\Big[x^2+(b-9)x\Big]_0^1=b-8.
\]
联立得
\[
b=8+2a,\qquad a=b-8.
\]
代入消元得
\[
b=8+2(b-8)\Rightarrow b=8,\quad a=b-8=0.
\]
所以
\[
f(x)=x+3+a=x+3.
\]
因此
\[
f(3)=3+3=6.
\]
\textcolor{red}{(待验证,为什么设$a,b$？)}
\end{solution}

\question 设 $0 < a < b$. 证明
\[
\int_a^b (x^2 + 1) e^{-x^2} \, dx > e^{-a^2} - e^{-b^2}.
\]

\begin{solution}
\textbf{解法 1.} 令
\[
f(x) = \int_0^x (t^2 + 1) e^{-t^2} \, dt, \quad g(x) = - e^{-x^2}.
\]
两函数在 $(0,\infty)$ 上均递增。由柯西中值定理,存在 $x \in (a,b)$ 使得
\[
\frac{f(b)-f(a)}{g(b)-g(a)} = \frac{f'(x)}{g'(x)} = \frac{(x^2+1)e^{-x^2}}{2 x e^{-x^2}} = \frac{1}{2}\Bigl(x + \frac{1}{x}\Bigr) \ge \sqrt{x \cdot \frac{1}{x}} = 1.
\]
于是
\[
\int_a^b (x^2+1)e^{-x^2} \, dx = f(b)-f(a) \ge g(b)-g(a) = e^{-a^2} - e^{-b^2}.
\]

\textbf{解法 2.} 直接比较积分被积函数:
\[
\int_a^b (x^2+1) e^{-x^2} \, dx \ge \int_a^b 2x e^{-x^2} \, dx = [-e^{-x^2}]_a^b = e^{-a^2} - e^{-b^2}.
\]
\end{solution}


\question 设 $f$ 是定义在区间 $[0,1]$ 上的连续实值函数。证明存在 $\xi \in [0,1]$,使得
\[
\int_{0}^{1} x^2 f(x) \, dx = \frac{1}{3} f(\xi).
\]

\begin{solution}
由于 $f$ 连续,它在 $[0,1]$ 上达到最小值和最大值,分别设为 $f(a)$ 和 $f(b)$,其中 $a,b \in [0,1]$。于是
\[
f(a) \int_{0}^{1} x^2 \, dx \le \int_{0}^{1} x^2 f(x) \, dx \le f(b) \int_{0}^{1} x^2 \, dx,
\]
即
\[
f(a) \le 3 \int_{0}^{1} x^2 f(x) \, dx \le f(b).
\]

由连续函数的中值定理可知,存在 $\xi \in [0,1]$,使得
\[
f(\xi) = 3 \int_{0}^{1} x^2 f(x) \, dx.
\]
\end{solution}

\question 设 $f: [0,1] \to (0,\infty)$ 是可积函数,且对所有 $x \in [0,1]$ 有
\[
f(x)\cdot f(1-x) = 1.
\]
证明
\[
\int_0^1 f(x)\,dx \geq 1.
\]

\begin{solution}
\textbf{方法一（使用算术-几何平均不等式）:} 对任意 $x\in[0,1]$,由 AM-GM 不等式有
\[
f(x)+f(1-x) \geq 2\sqrt{f(x)f(1-x)} = 2.
\]
在区间 $[0,1/2]$ 上积分得
\[
\int_0^1 f(x)\,dx = \int_0^{1/2} f(x)\,dx + \int_0^{1/2} f(1-x)\,dx = \int_0^{1/2} (f(x)+f(1-x))\,dx \geq \int_0^{1/2} 2\,dx = 1.
\]

\textbf{方法二（使用 Cauchy-Schwarz 不等式）:} 由条件可得
\[
\int_0^1 f(x)\,dx = \int_0^1 f(1-x)\,dx = \int_0^1 \frac{1}{f(x)}\,dx.
\]
于是
\[
\left(\int_0^1 f(x)\,dx\right)^2 = \int_0^1 f(x)\,dx \cdot \int_0^1 \frac{1}{f(x)}\,dx \geq \left(\int_0^1 1\,dx\right)^2 = 1,
\]
其中不等式来自 Cauchy-Schwarz,因此
\[
\int_0^1 f(x)\,dx \geq 1.
\]
\end{solution}

\question 设 $f:(-1,1)\to\mathbb{R}$ 是二阶可导函数,满足
\[
2f'(x)+xf''(x)\ge1,\quad x\in(-1,1).
\]
证明
\[
\int_{-1}^1 x f(x)\,dx \ge \frac{1}{3}.
\]

\begin{solution}
令
\[
g(x)=x f(x)-\frac{x^2}{2}.
\]
则
\[
g''(x)=2f'(x)+x f''(x)-1\ge0,
\]
因此 $g$ 是凸函数。

用 $g$ 在 $x=0$ 处的切线估计 $g$。设 $g'(0)=a$,则由凸性知
\[
g(x)\ge g(0)+g'(0)x=ax.
\]
于是
\begin{align*}
\int_{-1}^1 x f(x)\,dx
&=\int_{-1}^1\left(g(x)+\frac{x^2}{2}\right)\,dx \\
&\ge\int_{-1}^1\left(ax+\frac{x^2}{2}\right)\,dx
=\frac{1}{3}.
\end{align*}
这就证明了结论。
\end{solution}

\question 设 $f$ 为闭区间 $[0,1]$ 上的实值连续函数。已知
\[
\int_{0}^{1} f(x)\,dx = \frac{\pi}{4},
\]
证明存在 $y$ 满足 $0<y<1$ 且
\[
\frac{1}{1+y} < f(y) < \frac{1}{2y}。
\]

\begin{solution}
注意到
\[
\int_{0}^{1} \frac{1}{1+x^2}\,dx = \arctan(1) - \arctan(0) = \frac{\pi}{4}。
\]

令 $g(x) = f(x) - \frac{1}{1+x^2}$。若 $g(x)$ 在 $(0,1)$ 上恒为 $0$,则 $f(y) = \frac{1}{1+y^2}$ 对所有 $y \in (0,1)$ 都成立。若 $g$ 在 $(0,1)$ 上不恒为 $0$,则由于
\[
\int_{0}^{1} g(x)\,dx = 0,
\]
存在 $x_0, x_1 \in (0,1)$ 使得 $g(x_0) < 0$ 且 $g(x_1) > 0$。由 $g$ 在 $(0,1)$ 上连续,介值定理保证存在某个 $y$ 在 $x_0$ 与 $x_1$ 之间,使得 $g(y)=0$。于是 $f(y) = \frac{1}{1+y^2}$。

无论哪种情况,都存在 $y \in (0,1)$ 使得 $f(y) = \frac{1}{1+y^2}$。由于 $0<y<1$,有
\[
2y < 1+y^2 < 1+y,
\]
从而得到
\[
\frac{1}{1+y} < f(y) < \frac{1}{2y},
\]
证毕。
\end{solution}

\question 设函数 $f$ 在区间 $[a,b]$ 上具有连续导数,且满足 $f(a)=f(b)=0$。证明
\[
\max_{a\le x\le b} |f'(x)| \ge \frac{4}{(b-a)^2}\left|\int_a^b f(x)\,dx\right|.
\]

\begin{solution}
由中值定理,对任意 $x\in(a,b)$,有
\[
f(x)=f(x)-f(a)=f'(\xi_1)(x-a)=f(x)-f(b)=f'(\xi_2)(x-b),
\]
其中 $\xi_1\in(a,x),\xi_2\in(x,b)$。

设
\[
M=\max_{a\le x\le b}|f'(x)|.
\]
则
\[
|f(x)|\le M(x-a), \qquad |f(x)|\le M(b-x).
\]

因此
\[
\frac{4}{(b-a)^2}\int_a^b |f(x)|\,dx
\]
\begin{align*}
&\le \frac{4}{(b-a)^2}
\left(
\int_a^{(a+b)/2} M(x-a)\,dx
+\int_{(a+b)/2}^b M(b-x)\,dx
\right) \\
&= \frac{4}{(b-a)^2}
\left(
\frac{(b-a)^2}{8}M+\frac{(b-a)^2}{8}M
\right)
= M.
\end{align*}

又因为
\[
\left|\int_a^b f(x)\,dx\right|\le \int_a^b |f(x)|\,dx,
\]
于是
\[
\max_{a\le x\le b}|f'(x)|
\ge \frac{4}{(b-a)^2}\left|\int_a^b f(x)\,dx\right|.
\]
\end{solution}




\question 设 $a$ 为实数,已知函数
\[
f(x) = 4x^2 - 3a x + 4 \int_0^1 (t f(t))\,dt,
\]
及
\[
g(x) = x^2 + 4x + a - \int_0^x ((t+1) g'(t))\,dt.
\]
若方程 $f(x) - x \cdot g(x) = 0$ 有两相异实根 $\alpha, \beta$,且 $\alpha < \beta$,求
\[
\frac{1}{\beta - \alpha} \int_\alpha^\beta (3x^2 - 2 a x + a^2) \, dx
\]
的最小值。

\begin{solution}
设常数
\[
C = \int_0^1 t f(t) \, dt,
\]
则
\[
f(x) = 4x^2 - 3 a x + 4 C.
\]
计算 $C$:
\[
C = \int_0^1 t (4 t^2 - 3 a t + 4 C) dt = \int_0^1 (4 t^3 - 3 a t^2 + 4 C t) dt = \left[ t^4 - a t^3 + 2 C t^2 \right]_0^1 = 1 - a + 2 C.
\]
整理得
\[
C = 1 - a + 2 C \implies C = a - 1.
\]
因此
\[
f(x) = 4 x^2 - 3 a x + 4 a - 4.
\]

对 $g(x)$,由条件有
\[
g'(x) = 2 x + 4 - (x + 1) g'(x),
\]
整理得
\[
g'(x) + (x+1) g'(x) = 2 x + 4 \implies (x + 2) g'(x) = 2 x + 4,
\]
即
\[
g'(x) = \frac{2 x + 4}{x + 2} = 2.
\]

计算积分
\[
\int_0^x (t+1) g'(t) dt = \int_0^x 2 (t+1) dt = x^2 + 2 x,
\]
代入 $g(x)$:
\[
g(x) = x^2 + 4 x + a - (x^2 + 2 x) = 2 x + a.
\]

因此
\[
f(x) - x g(x) = (4 x^2 - 3 a x + 4 a - 4) - x (2 x + a) = 4 x^2 - 3 a x + 4 a - 4 - 2 x^2 - a x = 2 x^2 - 4 a x + 4 a - 4.
\]

设两根为 $\alpha, \beta$,满足
\[
\alpha + \beta = 2 a, \quad \alpha \beta = 2 a - 2.
\]

计算
\[
\frac{1}{\beta - \alpha} \int_\alpha^\beta (3 x^2 - 2 a x + a^2) dx = \frac{1}{\beta - \alpha} \left[ x^3 - a x^2 + a^2 x \right]_\alpha^\beta = \frac{1}{\beta - \alpha} \left( \beta^3 - \alpha^3 - a (\beta^2 - \alpha^2) + a^2 (\beta - \alpha) \right).
\]

利用因式分解:
\[
\beta^3 - \alpha^3 = (\beta - \alpha)(\beta^2 + \alpha \beta + \alpha^2), \quad \beta^2 - \alpha^2 = (\beta - \alpha)(\beta + \alpha),
\]
因此表达式为
\[
(\beta^2 + \alpha \beta + \alpha^2) - a (\beta + \alpha) + a^2.
\]

利用已知根的和与积:
\[
\beta + \alpha = 2 a, \quad \alpha \beta = 2 a - 2,
\]
以及
\[
\beta^2 + \alpha^2 = (\beta + \alpha)^2 - 2 \alpha \beta = (2 a)^2 - 2 (2 a - 2) = 4 a^2 - 4 a + 4.
\]

所以
\[
\beta^2 + \alpha \beta + \alpha^2 = (\beta^2 + \alpha^2) + \alpha \beta = (4 a^2 - 4 a + 4) + (2 a - 2) = 4 a^2 - 2 a + 2.
\]

代入原式:
\[
4 a^2 - 2 a + 2 - a \cdot 2 a + a^2 = 4 a^2 - 2 a + 2 - 2 a^2 + a^2 = 3 a^2 - 2 a + 2.
\]

令
\[
h(a) = 3 a^2 - 2 a + 2 = 3 \left(a - \frac{1}{3}\right)^2 + \frac{5}{3}.
\]

因此最小值为
\[
\frac{5}{3}.
\]
\textcolor{red}{(待验证)}
\end{solution}

\question 已知 $\int_0^2 f(x) dx = 1,f(2) = \frac{1}{2},f'(2) = 0$,计算 $[\int_0^1 x^2 f''(2x) dx]$。
\begin{solution}
首先进行变量代换,令 $t = 2x$,则 $dt = 2dx$:
\[ I = [\int_0^1 x^2 f''(2x) dx] = [\frac{1}{2} \int_0^2 (\frac{t}{2})^2 f''(t) dt] = [\frac{1}{8} \int_0^2 t^2 f''(t) dt] \]
使用分部积分法,令 $u = t^2,dv = f''(t)dt$:
\[ I = \frac{1}{8} ([t^2 f'(t)]_0^2 - \int_0^2 2t f'(t) dt) \]
代入 $f'(2) = 0$:
\[ I = \frac{1}{8} (0 - 2 \int_0^2 t f'(t) dt) = -\frac{1}{4} [\int_0^2 t f'(t) dt] \]
再次使用分部积分法,令 $u = t,dv = f'(t)dt$:
\[ I = -\frac{1}{4} ([t f(t)]_0^2 - \int_0^2 f(t) dt) \]
代入 $f(2) = \frac{1}{2}$ 和 $\int_0^2 f(t) dt = 1$:
\[ I = -\frac{1}{4} (2 \cdot \frac{1}{2} - 1) = -\frac{1}{4} (1 - 1) = 0 \]
最终结果为 $0$。
\end{solution}

    \question 坐标平面上,$A,B$ 两点分别在直线 $L_1:x=-\dfrac{3}{2}$ 与 $L_2:x=\dfrac{3}{2}$ 上。$\overline{AB} \perp L_1$,且 $M,N$ 为 $\overline{AB}$ 的三等分点,并满足 $\overline{AM} = \overline{MN} = \overline{NB}$。设 $\Gamma_1$ 为以 $L_1$ 为准线,$N$ 为焦点的抛物线；$\Gamma_2$ 为以 $L_2$ 为准线,$N$ 为顶点的抛物线,求 $\Gamma_1$ 与 $\Gamma_2$ 所围区域的面积。
    \ifprintanswers
    \begin{figure}[H]
        \centering
        \includegraphics[width=0.4\textwidth]{images/image31.png}
    \end{figure}
    \fi
    \begin{solution}
        所围面积为
        \[
        I = 4\int_0^{\frac{1}{2}} \sqrt{-4(x-\frac{1}{2})}\,dx = 8 \int_0^{\frac{1}{2}} \sqrt{\frac{1}{2}-x} \,dx
        \]
        令 $u = \dfrac{1}{2} - x , du = -dx$,则
        \[
        I = 8\int_0^{\frac{1}{2}} \sqrt{u}\,du 
        = 8 \left[ \frac{2}{3} u^{\frac{3}{2}} \right]_0^{\frac{1}{2}} 
        = \frac{4\sqrt{2}}{3}
        \]
    \end{solution}

    \question 已知四条抛物线 $\Gamma_1: y=x^2+a$, $\Gamma_2: y = -x^2 -a$, $\Gamma_3: y^2 = x-a$, $\Gamma_4: y^2 = -x-a$,其中 $a$ 为正实数,若任相邻两条抛物线均相切,试求这四条抛物线所围成之区域面积。
    \ifprintanswers
    \begin{figure}[H]
        \centering
        \includegraphics[width=0.4\textwidth]{images/image34.png}
    \end{figure}
    \fi
    \begin{solution}
        $\Gamma_1,\Gamma_3$的共切点在$x=y$上,故知
        \[
        x^2 - x + a = 0 
        \] 
        恰有一根,其中判别式为
        \[ 1 - 4a = 0 \Rightarrow a = \frac{1}{4}
        \] 
        故四个切点为
        \[
        \left(\frac{1}{2}, \frac{1}{2}\right),\left(\frac{1}{2}, -\frac{1}{2}\right),\left(-\frac{1}{2}, \frac{1}{2}\right),\left(-\frac{1}{2}, -\frac{1}{2}\right),
        \]
        所围面积为
        \[
        8 \int_0^{\frac{1}{2}} (x^2 - x + \frac{1}{4})\,dx  =8\left[\frac{x^3}{3}-\frac{x^2}{2}+\frac{x}{4}\right]^{\frac12}_0 =\frac{1}{3}
        \]
    \end{solution}

    \question 已知曲线
\[
2x^2 + 2xy + y^2 = 50,
\]
求被 $x$ 轴和曲线上 $y \ge 0$ 部分围成的有限区域面积。

\begin{solution}

\textbf{步骤 1: 找 $x$、$y$ 截距并作图}  

将曲线改写为 $y$ 的形式:
\[
y^2 + 2xy + 2x^2 = 50 \implies (y+x)^2 + x^2 = 50 \implies (y+x)^2 = 50 - x^2
\]
\[
y = -x \pm \sqrt{50-x^2}
\]

\textbf{步骤 2: 找垂直切线点 $P$}  

对曲线求导:
\[
\frac{d}{dx}(2x^2 + 2xy + y^2) = 0 \implies 4x + 2y + 2x\frac{dy}{dx} + 2y\frac{dy}{dx} = 0
\]
\[
\implies 2(x+y) + 2(x+y)\frac{dy}{dx} = 0 \implies y+x = 0 \implies y=-x
\]

代入曲线方程求 $x$:
\[
2x^2 + 2x(-x) + (-x)^2 = 50 \implies x^2 = 50 \implies x = \pm 5\sqrt{2}
\]

所以垂直切线点为
\[
P = (-5\sqrt{2}, 5\sqrt{2})
\]

\textbf{步骤 3: 准备积分表达面积}  

使用 $y = -x \pm \sqrt{50-x^2}$,考虑 $y \ge 0$ 部分。面积由两部分曲线和 $x$ 轴围成:
\[
A = \int_{-5\sqrt{2}}^{5\sqrt{2}} [-x+\sqrt{50-x^2}] \, dx - \int_{-5\sqrt{2}}^{-5} [-x-\sqrt{50-x^2}] \, dx
\]

拆分积分:
\[
A = \int_{-5\sqrt{2}}^{5\sqrt{2}} -x \, dx + \int_{-5\sqrt{2}}^{5\sqrt{2}} \sqrt{50-x^2} \, dx + \int_{-5\sqrt{2}}^{-5} x \, dx + \int_{-5\sqrt{2}}^{-5} \sqrt{50-x^2} \, dx
\]

\textbf{步骤 4: 使用三角代换计算 $\sqrt{50-x^2}$ 积分}  

令
\[
x = \sqrt{50} \sin\theta \implies dx = \sqrt{50} \cos\theta \, d\theta
\]
\[
\sqrt{50-x^2} = \sqrt{50}\cos\theta
\]

积分变为:
\[
\int \sqrt{50-x^2} \, dx = \int 50 \cos^2\theta \, d\theta = \int 25 + 25\cos 2\theta \, d\theta = 25\theta + \frac{25\sin 2\theta}{2}
\]

对应的 $\theta$ 值:
\[
x = 5\sqrt{2} \implies \theta = \frac{\pi}{2}, \quad x = -5\sqrt{2} \implies \theta = -\frac{\pi}{2}, \quad x = -5 \implies \theta = -\frac{\pi}{4}
\]

\textbf{步骤 5: 代入计算面积}  
\begin{align*}
A &= \left[-\frac{1}{2}x^2\right]_{-5\sqrt{2}}^{5} + \left[25\theta + \frac{25\sin 2\theta}{2}\right]_{-\pi/2}^{-\pi/4} + \left[25\theta + \frac{25\sin 2\theta}{2}\right]_{-\pi/2}^{\pi/2} \\
&= \left[\frac{-25}{2} - \frac{-50}{2}\right] + \left[ -\frac{25\pi}{4} - \frac{25}{2}(-1) - (-\frac{25\pi}{2}+0) \right] + \left[ \frac{25\pi}{2}-(-\frac{25\pi}{2}) \right] \\
&= 25\pi
\end{align*}

\textbf{最终答案:}
\[
\boxed{A = 25\pi}
\]

\end{solution}


    \question 座标平面上,满足联立不等式
    \[
    \begin{cases}
    x^2 + (y-1)^2 \le 1 \\
    x - y \le 0 \\
    x - y \ge -2
    \end{cases}
    \]
    的解区域为 $S$,求区域 $S$ 绕 $x$ 轴旋转一圈所得立体的体积。
    \ifprintanswers
    \begin{figure}[H]
        \centering
        \includegraphics[width=0.4\textwidth]{images/image130.jpg}
    \end{figure}
    \fi
    \begin{solution}
        联立 $x^2 + (y-1)^2 =1$与 $y=x+2$得
        \[
        A(0,2), \ D(-1,1)
        \]
        联立 $x^2 + (y-1)^2 =1$与 $y=x$得
        \[
        B(1,1), \ C(0,0)
        \]
        如上图,左半部绕 $x$ 轴旋转体的体积为
        \[
        \begin{aligned}
        V_1 &= \pi \int_{-1}^{0} \Big[(x+2)^2 - (1 - \sqrt{1-x^2})^2\Big]\, dx \\
        &= \pi \int_{-1}^{0} \Big[2x^2 + 4x + 2 + 2\sqrt{1-x^2}\Big]\, dx \\
        &= \pi \left[ \frac{2}{3}x^3 + 2x^2 + 2x + (\sqrt{1-x^2} + \arcsin x) \right]_{-1}^{0} = \frac{1}{2}\pi^2 + \frac{5}{3}\pi
        \end{aligned}
        \]
        上图右半部绕 $x$ 轴旋转体的体积为
        \[
        \begin{aligned}
        V_2 &= \pi \int_{0}^{1} \Big[(1 + \sqrt{1-x^2})^2 - x^2\Big]\, dx \\
        &= \pi \int_{0}^{1} \Big[ 2 - 2x^2 + 2\sqrt{1-x^2} \Big]\, dx \\
        &= \pi \left[ 2x - \frac{2}{3}x^3 + \sqrt{1-x^2} + \arcsin x \right]_{0}^{1} = \frac{1}{2}\pi^2 + \frac{1}{3}\pi
        \end{aligned}
        \]
        因此总体积为
        \[
        V = V_1 + V_2 = \pi^2 + 2\pi
        \]
    \end{solution}

    \question 在坐标平面上,求由不等式 
    \[
    \frac{x^2}{4} + y^2 \le 1, \quad y+1 \ge \left(\frac{\sqrt{3}}{2}+1\right)x, \quad y+1 \ge -\left(\frac{\sqrt{3}}{2}+1\right)x
    \]
    所围成的图形面积。
    \ifprintanswers
    \begin{figure}[H]
        \centering
        \includegraphics[width=0.6\textwidth]{images/image120.jpg}
    \end{figure}
    \fi
    \begin{solution}
        三方程的交点为
        \[
        A\left(1, \frac{\sqrt{3}}{2}\right), \quad B\left(-1, \frac{\sqrt{3}}{2}\right), \quad C(0,-1)
        \]
        图形关于 $y$ 轴对称,于是所求面积为
        \[
        S = 2 \left[ \int_0^1 \sqrt{1 - \frac{x^2}{4}} \, dx - \int_0^1 \left( \frac{\sqrt{3}+2}{2}x - 1 \right) dx \right]
        \]
        令 $x = 2 \sin \theta ,dx = 2 \cos \theta \, d\theta$,则
        \[
        \int_0^1 \sqrt{1 - \frac{x^2}{4}} \, dx = \int_0^{\frac{\pi}{6}} \sqrt{1 - \sin^2 \theta} \cdot 2\cos \theta \, d\theta = \int_0^{\frac{\pi}{6}} (1 + \cos 2\theta) \, d\theta = \frac{\sqrt{3}}{4} + \frac{\pi}{6}
        \]
        且
        \[
        \int_0^1 \left( \frac{\sqrt{3}+2}{2}x - 1 \right) dx = \left[\frac{\sqrt{3}+2}{4}x^2-x\right]_0^1=\frac{\sqrt{3}-2}{4}
        \]
        所以
        \[
        S = 2\left( \frac{\pi}{6} + \frac{\sqrt{3}}{4} - \frac{\sqrt{3}-2}{4} \right) = 1 + \frac{\pi}{3}
        \]
    \end{solution}

    \question 计算积分
\[
I(\theta) = \int_{-1}^{1} \frac{\sin \theta \, dx}{1 - 2x \cos \theta + x^2},
\]
并确定 $0 \le \theta \le 2\pi$ 范围内函数 $I(\theta)$ 不连续的点。

\begin{solution}
用恒等式 $1 = \cos^2 \theta + \sin^2 \theta$,并令
\[
u = x - \cos \theta, \quad du = dx.
\]

则积分变为
\[
I(\theta) = \int_{-1-\cos\theta}^{1-\cos\theta} \frac{\sin\theta \, du}{u^2 + \sin^2\theta} = \left. \tan^{-1} \frac{u}{\sin\theta} \right|_{-1-\cos\theta}^{1-\cos\theta}.
\]

因此
\[
I(\theta) = \tan^{-1} \frac{1-\cos\theta}{\sin\theta} - \tan^{-1} \frac{-1-\cos\theta}{\sin\theta}, \quad (\sin\theta \neq 0).
\]

利用二倍角公式
\[
1-\cos\theta = 2\sin^2(\theta/2), \quad 1+\cos\theta = 2\cos^2(\theta/2), \quad \sin\theta = 2\sin(\theta/2)\cos(\theta/2),
\]
得到
\[
\frac{1-\cos\theta}{\sin\theta} = \frac{2\sin^2(\theta/2)}{2\sin(\theta/2)\cos(\theta/2)} = \tan(\theta/2),
\quad 
\frac{-1-\cos\theta}{\sin\theta} = -\frac{2\cos^2(\theta/2)}{2\sin(\theta/2)\cos(\theta/2)} = -\cot(\theta/2) = \tan(\pi/2 + \theta/2).
\]

因此
\[
I(\theta) = \tan^{-1}(\tan(\theta/2)) - \tan^{-1}(\tan(\pi/2 + \theta/2)).
\]

注意 $\tan^{-1}$ 的主值在 $(-\pi/2, \pi/2)$。例如:
\[
\tan^{-1}(\tan 20^\circ) = 20^\circ, \quad 
\tan^{-1}(\tan 100^\circ) = -80^\circ, \quad
\tan^{-1}(\tan 220^\circ) = 40^\circ, \quad
\tan^{-1}(\tan 290^\circ) = -70^\circ.
\]

因此,如果 $0 < \theta < \pi$,则
\[
I(\theta) = \theta/2 - (-(\pi/2 - \theta/2)) = \pi/2.
\]

如果 $\pi < \theta < 2\pi$,则
\[
I(\theta) = -(\pi - \theta/2) - (\theta/2 - \pi/2) = -\pi/2.
\]

因此,函数 $I(\theta)$ 为
\[
I(\theta) = 
\begin{cases}
\pi/2, & 0 < \theta < \pi, \\
-\pi/2, & \pi < \theta < 2\pi.
\end{cases}
\]

不连续点为 $\theta = 0, \pi, 2\pi$。
\end{solution}

\question 计算极限
\[
\lim_{x \to 0^+} \int_{x}^{2x} \frac{\sin^m t}{t^n} \, dt \quad (m,n \in \mathbb{N}).
\]

\begin{solution}
注意到当 $t \to 0^+$ 时 $\frac{\sin t}{t} \to 1$,且在 $(0, \pi)$ 上 $\frac{\sin t}{t}$ 单调递减。因此,对于 $x \in (0, \pi/2)$ 且 $t \in [x,2x]$,有
\[
\frac{\sin 2x}{2x} < \frac{\sin t}{t} < 1.
\]
于是
\[
\left(\frac{\sin 2x}{2x}\right)^m \int_x^{2x} t^{m-n} \, dt
<
\int_x^{2x} \frac{\sin^m t}{t^n} \, dt
<
\int_x^{2x} t^{m-n} \, dt.
\]

对积分做换元 $t = xu$ 得
\[
\int_x^{2x} t^{m-n} \, dt = x^{m-n+1} \int_1^2 u^{m-n} \, du.
\]

注意到 $\left(\frac{\sin 2x}{2x}\right)^m \to 1$,于是极限取决于 $x^{m-n+1}$ 的幂:
\[
\lim_{x \to 0^+} \int_x^{2x} \frac{\sin^m t}{t^n} \, dt =
\begin{cases}
0, & m-n+1>0 \ (\text{即 } m \ge n),\\[2mm]
\ln 2, & m-n+1=0 \ (\text{即 } n-m=1),\\[1mm]
+\infty, & m-n+1<0 \ (\text{即 } n-m>1).
\end{cases}
\]

\end{solution}

\question
已知 $a$ 和 $b$ 是不同实常数,$\lambda$ 是实参数。证明
\[
\left[ \int_{a}^{b} f(x)g(x) \, dx \right]^2 \leq \left[ \int_{a}^{b} [f(x)]^2 \, dx \right]\left[ \int_{a}^{b} [g(x)]^2 \, dx \right]
\]
\begin{solution}
a) 由关系
\[
[\lambda f(x)+g(x)]^2 \geq 0
\]
展开平方:
\[
\lambda^2 [f(x)]^2 + 2\lambda f(x)g(x) + [g(x)]^2 \geq 0
\]
对 $x$ 从 $a$ 到 $b$ 积分:
\[
\int_{a}^{b} (\lambda^2 [f(x)]^2 + 2\lambda f(x)g(x) + [g(x)]^2) \, dx \geq 0
\]
利用积分线性:
\[
\lambda^2 \int_{a}^{b} [f(x)]^2 \, dx + 2\lambda \int_{a}^{b} f(x)g(x) \, dx + \int_{a}^{b} [g(x)]^2 \, dx \geq 0
\]

b) 该关于 $\lambda$ 的二次不等式要求判别式非正:
\[
A = \int_{a}^{b} [f(x)]^2 \, dx, \quad B = 2\int_{a}^{b} f(x)g(x) \, dx, \quad C = \int_{a}^{b} [g(x)]^2 \, dx
\]
\[
B^2 - 4AC \leq 0 \implies \left[ \int_{a}^{b} f(x)g(x) \, dx \right]^2 \leq \left[ \int_{a}^{b} [f(x)]^2 \, dx \right]\left[ \int_{a}^{b} [g(x)]^2 \, dx \right]
\]

c) 令 $f(x) = \sqrt{\sin x},g(x)=1$:
\[
\left[ \int_{0}^{\pi/2} \sqrt{\sin x} \, dx \right]^2 \leq \left[ \int_{0}^{\pi/2} \sin x \, dx \right]\left[ \int_{0}^{\pi/2} 1 \, dx \right]
\]
\[
\left[ \int_{0}^{\pi/2} \sqrt{\sin x} \, dx \right]^2 \leq [-\cos x]_{0}^{\pi/2} \cdot [x]_{0}^{\pi/2} = 1 \cdot \frac{\pi}{2}
\]
\[
\int_{0}^{\pi/2} \sqrt{\sin x} \, dx \leq \sqrt{\frac{\pi}{2}}
\]

d) 令 $f(x) = (\sin x)^{1/4},g(x) = \cos x$:
\[
\left[ \int_{0}^{\pi/2} (\sin x)^{1/4}\cos x \, dx \right]^2 \leq \left[ \int_{0}^{\pi/2} (\sin x)^{1/2} \, dx \right]\left[ \int_{0}^{\pi/2} \cos^2 x \, dx \right]
\]
LHS 积分用代换 $u=\sin x, du=\cos x \, dx$:
\[
\int_{0}^{\pi/2} (\sin x)^{1/4}\cos x \, dx = \frac{4}{5}
\]
RHS 第二个积分:
\[
\int_{0}^{\pi/2} \cos^2 x \, dx = \int_{0}^{\pi/2} \frac{1+\cos 2x}{2} \, dx = \frac{\pi}{4}
\]
代入不等式:
\[
\left(\frac{4}{5}\right)^2 \leq \left[ \int_{0}^{\pi/2} \sqrt{\sin x} \, dx \right] \cdot \frac{\pi}{4}
\]
整理:
\[
\int_{0}^{\pi/2} \sqrt{\sin x} \, dx \geq \frac{16}{25} \cdot \frac{4}{\pi} = \frac{64}{25\pi}
\]
\end{solution}

\end{questions}

\pagebreak

\begin{center}
  {\fontsize{30pt}{26pt}\selectfont
    \hypertarget{积分技巧}{积分技巧} \label{积分技巧}
  }
\end{center}
\separator
\vspace{1pt}
\begin{questions}

%substitution methods
\question 计算 $[\int_2^4 (\log_x 2 - \frac{\log_x^2 2}{\ln 2}) dx]$。
\begin{solution}
利用换底公式简化被积函数:
\[\log_x 2 - \frac{(\frac{\ln 2}{\ln x})^2}{\ln 2} = \frac{\ln 2}{\ln x} - \frac{\ln 2}{(\ln x)^2} = \ln 2 (\frac{1}{\ln x} - \frac{1}{\ln^2 x})\]
由于 $\frac{d}{dx}(\frac{x}{\ln x}) = \frac{\ln x - 1}{\ln^2 x} = \frac{1}{\ln x} - \frac{1}{\ln^2 x}$,
该积分的原始函数为:
\[I = \ln 2 [\frac{x}{\ln x}]_2^4 = \ln 2 (\frac{4}{\ln 4} - \frac{2}{\ln 2})\]
已知 $\ln 4 = 2 \ln 2$,代入得:
\[I = \ln 2 (\frac{4}{2 \ln 2} - \frac{2}{\ln 2}) = \ln 2 (\frac{2}{\ln 2} - \frac{2}{\ln 2}) = 0\]
\end{solution}

\question
计算积分substitution methods
\[
\int_{80^4}^{15^4} \frac{1}{x^2 (1+x^4)^{\frac{3}{4}}} \, dx
\]

\begin{solution}
因为积分区间为正,可将被积式中的 $x^4$ 提出:
\[
\int_{80^4}^{15^4} \frac{1}{x^2 (1+x^4)^{\frac{3}{4}}} \, dx
= \int_{80^4}^{15^4} \frac{1}{x^2 (x^4)^{\frac{3}{4}} (x^{-4}+1)^{\frac{3}{4}}} \, dx
= \int_{80^4}^{15^4} \frac{1}{x^2 x^3 (x^{-4}+1)^{\frac{3}{4}}} \, dx
= \int_{80^4}^{15^4} x^{-5} (1+x^{-4})^{-\frac{3}{4}} \, dx
\]

注意到
\[
\frac{d}{dx} \left[ (1+x^{-4})^{\frac{1}{4}} \right] 
= \frac{1}{4} (1+x^{-4})^{-\frac{3}{4}} \cdot (-4 x^{-5}) 
= - x^{-5} (1+x^{-4})^{-\frac{3}{4}}
\]

因此
\[
\int_{80^4}^{15^4} x^{-5} (1+x^{-4})^{-\frac{3}{4}} \, dx
= - \left[ (1+x^{-4})^{\frac{1}{4}} \right]_{80^4}^{15^4}
= - \left[ (1+15^{-4})^{\frac{1}{4}} - (1+80^{-4})^{\frac{1}{4}} \right]
= - (16^{\frac{1}{4}} - 81^{\frac{1}{4}})
= - (2-3)
= 1
\]

最终结果:
\[
\int_{80^4}^{15^4} \frac{1}{x^2 (1+x^4)^{\frac{3}{4}}} \, dx = 1
\]
\end{solution}

\question 计算积分
\[
\int_{0}^{1} \frac{1}{(x^{2} + 4x^{3})^{\frac{2}{4}}} \, dx
\]
\begin{solution}
先将分母中的根式因式分解:
\[
x^{2} + 4x^{3} = x^{2}(1 + 4x) = [x^{\frac{1}{2}}(x^{\frac{1}{2}}+4)^{\frac{3}{4}}]^{?} \quad \text{（整理成适合代换的形式）}
\]

更简便地处理:
\[
\int_{0}^{1} \frac{1}{(x^{2}+4x^{3})^{\frac{1}{2}}} \, dx
= \int_{0}^{1} \frac{1}{x^{\frac{1}{2}} (x^{\frac{1}{2}}+4)^{\frac{3}{4}}} \, dx
\]

作代换:
\[
u = x^{\frac{1}{2}} \implies x = u^2, \quad dx = 2u \, du
\]

积分上下限不变（$x=0\to u=0$, $x=1\to u=1$）:
\[
\int_{0}^{1} \frac{1}{u (u+4)^{\frac{3}{4}}} \cdot 2u \, du = \int_{0}^{1} \frac{2}{(u+4)^{\frac{3}{4}}} \, du
\]

直接积分:
\[
\int_{0}^{1} 2 (u+4)^{-\frac{3}{4}} \, du = \left[ 8 (u+4)^{\frac{1}{4}} \right]_{0}^{1} = 8 \left( 5^{\frac{1}{4}} - 4^{\frac{1}{4}} \right)
\]

\[
= 8 \left( \sqrt[4]{5} - \sqrt{2} \right)
\]
\end{solution}

\question
7) 计算定积分:$\int_{0}^{4} x^3 \sqrt{9+x^2} dx$
\begin{solution}
设 $u = 9+x^2, du = 2x dx$:
\begin{align*}
I &= \frac{1}{2} \int_{9}^{25} (u-9) \sqrt{u} du = \frac{1}{2} \left[ \frac{2}{5}u^{5/2} - 6u^{3/2} \right]_{9}^{25} \\
&= \left[ \frac{1}{5}u^{5/2} - 3u^{3/2} \right]_{9}^{25} = (625 - 375) - (\frac{243}{5} - 81) = 282.4
\end{align*}
\textbf{结果:} $282.4$
\end{solution}

\question
1) 计算不定积分:\[ \int \frac{1}{1+x^4} dx \]

\begin{solution}
首先将被积函数拆分:
\[ \int \frac{1}{1+x^4} dx = \frac{1}{2} \int \frac{x^2+1}{x^4+1} dx - \frac{1}{2} \int \frac{x^2-1}{x^4+1} dx \]

对于第一部分,分子分母同除以 $x^2$:
\[ \frac{1}{2} \int \frac{1+\frac{1}{x^2}}{x^2+\frac{1}{x^2}} dx = \frac{1}{2} \int \frac{d(x-\frac{1}{x})}{(x-\frac{1}{x})^2+2} \]
\[ = \frac{1}{2\sqrt{2}} \tan^{-1}\left(\frac{x-\frac{1}{x}}{\sqrt{2}}\right) = \frac{\sqrt{2}}{4} \tan^{-1}\left(\frac{x^2-1}{\sqrt{2}x}\right) \]

对于第二部分,同理可得:
\[ \frac{1}{2} \int \frac{1-\frac{1}{x^2}}{x^2+\frac{1}{x^2}} dx = \frac{1}{2} \int \frac{d(x+\frac{1}{x})}{(x+\frac{1}{x})^2-2} \]
利用对数积分公式:
\[ = \frac{1}{4\sqrt{2}} \ln\left| \frac{x+\frac{1}{x}-\sqrt{2}}{x+\frac{1}{x}+\sqrt{2}} \right| = \frac{\sqrt{2}}{8} \ln\left| \frac{x^2-\sqrt{2}x+1}{x^2+\sqrt{2}x+1} \right| \]

合并结果得到:
\[ \int \frac{1}{1+x^4} dx = \frac{\sqrt{2}}{4} \tan^{-1}\left(\frac{x^2-1}{\sqrt{2}x}\right) - \frac{\sqrt{2}}{8} \ln\left| \frac{x^2-\sqrt{2}x+1}{x^2+\sqrt{2}x+1} \right| + C \]
\end{solution}
%exponential
\question 计算积分
\[
\int \frac{4}{e^{3x}\sqrt{e^{2x}+4}} \, dx
\]
\begin{solution}
作代换:
\[
t = \frac{1}{e^x} \implies dx = -\frac{dt}{t}
\]

代入积分:
\[
\int \frac{4}{e^{3x}\sqrt{e^{2x}+4}} \, dx 
= \int \frac{4}{\frac{1}{t^3} \sqrt{\frac{1}{t^2}+4}} \left(-\frac{dt}{t}\right)
= \int \frac{-4 t^3}{\sqrt{\frac{1+4t^2}{t^2}}} \, dt
= \int \frac{-4 t^3}{\frac{\sqrt{1+4t^2}}{t}} \, dt
= \int \frac{-4 t^4}{\sqrt{1+4t^2}} \, dt
\]

整理:
\[
\frac{-4 t^4}{\sqrt{1+4t^2}} = -4 t (1+4t^2)^{-\frac{1}{2}} \, dt
\]

直接使用反链式法则或作代换 \(u = 1+4t^2\):
\[
\int -4 t (1+4t^2)^{-\frac{1}{2}} \, dt = -(1+4t^2)^{\frac{1}{2}} + C
\]

代回 \(t = \frac{1}{e^x}\):
\[
-(1+4t^2)^{\frac{1}{2}} + C = -\left(1 + \frac{4}{e^{2x}}\right)^{\frac{1}{2}} + C
= -\left(\frac{e^{2x}+4}{e^{2x}}\right)^{\frac{1}{2}} + C
= -\frac{\sqrt{e^{2x}+4}}{e^x} + C
\]
\end{solution}

\question 求
\[
\int_{0}^{\pi} \frac{1 + x \cos x}{x + e^{\sin x}} \, dx
\]
\begin{solution}
作代换
\[
u = 1 + x e^{\sin x}
\]
\[
\frac{du}{dx} = e^{\sin x} + x e^{\sin x} \cos x = e^{\sin x} (1 + x \cos x)
\]
\[
dx = \frac{du}{e^{\sin x} (1 + x \cos x)}
\]

积分的上下限变换
\[
x = 0 \implies u = 1
\]
\[
x = \pi \implies u = 1 + \pi
\]

代入积分
\[
\int_{0}^{\pi} \frac{1 + x \cos x}{x + e^{\sin x}} \, dx
= \int_{1}^{1+\pi} \frac{1 + x \cos x}{x + e^{\sin x}} \cdot \frac{du}{e^{\sin x} (1 + x \cos x)}
= \int_{1}^{1+\pi} \frac{1}{u} \, du
\]

计算积分
\[
[\ln|u|]_{1}^{1+\pi} = \ln(1+\pi) - \ln 1 = \ln(1+\pi)
\]
\end{solution}


\question 计算积分
\[
\int_{0}^{\ln 2} \sqrt{e^x - 1} \, dx
\]
\begin{solution}
令
\[
u = \sqrt{e^x - 1} \implies u^2 = e^x - 1, \quad 2u\,du = e^x dx, \quad dx = \frac{2u}{u^2+1} \, du.
\]

积分上下限变为:
\[
x = 0 \implies u = 0, \quad x = \ln 2 \implies u = \sqrt{2-1} = 1.
\]

原积分变为:
\[
\int_{0}^{1} u \cdot \frac{2u}{u^2+1} \, du = \int_{0}^{1} \frac{2u^2}{u^2+1} \, du.
\]

通过拆分:
\[
\frac{2u^2}{u^2+1} = 2 - \frac{2}{u^2+1}.
\]

因此积分为:
\[
\int_{0}^{1} \left(2 - \frac{2}{u^2+1}\right) du = \left[2u - 2\arctan u \right]_{0}^{1} = 2 - 2 \cdot \frac{\pi}{4} = 2 - \frac{\pi}{2}.
\]

\[
\therefore \int_{0}^{\ln 2} \sqrt{e^x - 1} \, dx = 2 - \frac{\pi}{2}.
\]
\end{solution}

\question 计算积分
\[
\int_{0}^{\infty} \frac{e^{8x}-e^{2x}}{(e^{8x}+3)(e^{2x}+3)} \, dx
\]
\begin{solution}
先通过部分分式分解（观察得出）:
\[
\frac{e^{8x}-e^{2x}}{(e^{8x}+3)(e^{2x}+3)} = -\frac{1}{e^{8x}+3} + \frac{1}{e^{2x}+3}
\]

于是积分可写为:
\[
\int_{0}^{\infty} \frac{e^{8x}-e^{2x}}{(e^{8x}+3)(e^{2x}+3)} \, dx
= \int_{0}^{\infty} \frac{1}{e^{2x}+3} \, dx - \int_{0}^{\infty} \frac{1}{e^{8x}+3} \, dx
\]

先计算第一个积分:
\[
\int_{0}^{\infty} \frac{1}{e^{2x}+3} \, dx = \int_{0}^{\infty} \frac{e^{-2x}}{1+3 e^{-2x}} \, dx
= \left[-\frac{1}{6}\ln\left(1+3 e^{-2x}\right)\right]_{0}^{\infty}
= \frac{1}{6} \ln 4
\]

第二个积分:
\[
-\int_{0}^{\infty} \frac{1}{e^{8x}+3} \, dx = -\int_{0}^{\infty} \frac{e^{-8x}}{1+3 e^{-8x}} \, dx
= \left[-\frac{1}{24} \ln\left(1+3 e^{-8x}\right)\right]_{0}^{\infty}
= \frac{1}{24} \ln 4
\]

合并结果:
\[
\int_{0}^{\infty} \frac{e^{8x}-e^{2x}}{(e^{8x}+3)(e^{2x}+3)} \, dx
= \frac{1}{6} \ln 4 - \frac{1}{24} \ln 4
= \frac{3}{24} \ln 4
= \frac{1}{8} \ln 4
= \frac{1}{8} \ln (2^2)
= \frac{1}{4} \ln 2
\]
\end{solution}

\question 计算
\[
\int_{0}^{1} \frac{e^{x}(1-x)}{x^2+e^{2x}} \, dx
\]

\begin{solution}
首先观察积分形式,尝试将分母调整为 $1+(xe^{-x})^2$ 的形式:
\[
\int_{0}^{1} \frac{e^{x}(1-x)}{x^2+e^{2x}} \, dx = \int_{0}^{1} \frac{e^{x}(1-x)}{x^2+e^{2x}} \cdot \frac{e^{-2x}}{e^{-2x}} \, dx = \int_{0}^{1} \frac{e^{-x}(1-x)}{(xe^{-x})^2+1} \, dx
\]

设
\[
u = x e^{-x} \implies du = (1-x)e^{-x} dx
\]
积分上下限对应:
\[
x=0 \to u=0, \quad x=1 \to u = e^{-1}
\]

因此积分变为
\[
\int_{0}^{1} \frac{e^{-x}(1-x)}{(xe^{-x})^2+1} \, dx = \int_{0}^{e^{-1}} \frac{1}{u^2+1} \, du = [\arctan u]_{0}^{e^{-1}} = \arctan e^{-1} - \arctan 0
\]

最终得到
\[
\int_{0}^{1} \frac{e^{x}(1-x)}{x^2+e^{2x}} \, dx = \arctan \frac{1}{e}
\]
\end{solution}

    \question 已知 
    \[
    \int_{0}^{\infty} e^{-x^2}\,dx = \frac{\sqrt{\pi}}{2},
    \]
    求
    \[
    \int_{e}^{\infty} e^{-x^2} \ln(x^2) x^{(\ln(x^{-x^2}) + 2x^2)}\,dx
    \]
    \begin{solution}
        发现
        \begin{align*} 
        I &= \int_{e}^{\infty} e^{-x^2} \ln(x^2) x^{(\ln(x^{-x^2}) + 2x^2)}\,dx \\
        &= \int_{e}^{\infty} e^{-x^2} \ln(x^2) e^{\ln x(-x^2 \ln x + 2x^2)}\,dx \\
        &= 2\int_{e}^{\infty} e^{-(x\ln x-x)^2} \ln(x) \,dx
        \end{align*}
        设$u=x\ln x -x$,则$du=\left(\ln x +x\cdot\dfrac{1}{x}-1\right)\,dx=\ln x \,dx$,于是
        \[
        I=2\int_{0}^{\infty}e^{-u^2}\,du=2\cdot\frac{\sqrt{\pi}}{2}=\sqrt{\pi}
        \]
    \end{solution}

    \question 
计算不定积分:
\[ \int \left( \frac{\ln x - 1}{1 + (\ln x)^2} \right)^2 dx \]

\begin{solution}
这道题可以通过换元法简化。设 $u = \ln x$,则 $x = e^u,dx = e^u du$。

代入原式:
\begin{align*}
I &= \int \left( \frac{u-1}{1+u^2} \right)^2 e^u du \\
&= \int \frac{u^2 - 2u + 1}{(1+u^2)^2} e^u du \\
&= \int \left[ \frac{1+u^2}{(1+u^2)^2} - \frac{2u}{(1+u^2)^2} \right] e^u du \\
&= \int \left[ \frac{1}{1+u^2} - \frac{2u}{(1+u^2)^2} \right] e^u du
\end{align*}

注意到被积函数符合分部积分的特征 $\int [f(u) + f'(u)] e^u du = f(u)e^u + C$。
令 $f(u) = \frac{1}{1+u^2}$,则 $f'(u) = -\frac{2u}{(1+u^2)^2}$。

因此:
\begin{align*}
I &= \frac{1}{1+u^2} e^u + C \\
&= \frac{x}{1+(\ln x)^2} + C
\end{align*}

\textbf{结果:} $\frac{x}{(\ln x)^2 + 1} + C$ \quad
\end{solution}
%radicals
\question 
18) 计算不定积分 $\int x\sqrt{x+19} \, dx$

\begin{solution}
本题展示了两种正确的解法。

\textbf{解法一:利用凑微分与分配律}
\begin{align*}
\int x\sqrt{x+19} \, dx &= \int (x+19-19)(x+19)^{\frac{1}{2}} \, dx \\
&= \int (x+19)^{\frac{3}{2}} \, dx - 19\int (x+19)^{\frac{1}{2}} \, dx \\
&= \frac{2}{5}(x+19)^{\frac{5}{2}} - \frac{38}{3}(x+19)^{\frac{3}{2}} + C
\end{align*}

\textbf{解法二:分部积分法}
设 $u = x, dv = (x+19)^{\frac{1}{2}}dx \implies du = dx, v = \frac{2}{3}(x+19)^{\frac{3}{2}}$:
\begin{align*}
\int x\sqrt{x+19} \, dx &= \frac{2}{3}x(x+19)^{\frac{3}{2}} - \frac{2}{3}\int (x+19)^{\frac{3}{2}} \, dx \\
&= \frac{2}{3}x(x+19)^{\frac{3}{2}} - \frac{4}{15}(x+19)^{\frac{5}{2}} + C
\end{align*}
注:两结果形式不同但通过化简后实质相等。
\end{solution}

\question
20) 计算定积分 $I = \int_{3}^{6} (\sqrt{x+\sqrt{12x-36}} + \sqrt{x-\sqrt{12x-36}}) \, dx$

\begin{solution}
设 $x = 3+t, dx = dt$。当 $x=3, t=0$；当 $x=6, t=3$。
注意到 $\sqrt{12x-36} = \sqrt{12(x-3)} = \sqrt{12t} = 2\sqrt{3t}$。
\begin{align*}
I &= \int_{0}^{3} (\sqrt{3+t+2\sqrt{3t}} + \sqrt{3+t-2\sqrt{3t}}) \, dt \\
\intertext{利用完全平方公式 $\sqrt{(\sqrt{3} \pm \sqrt{t})^2} = |\sqrt{3} \pm \sqrt{t}|$:}
&= \int_{0}^{3} [(\sqrt{3}+\sqrt{t}) + (\sqrt{3}-\sqrt{t})] \, dt \quad (\text{在 } t \in [0,3] \text{ 范围内 } \sqrt{3} \ge \sqrt{t}) \\
&= \int_{0}^{3} 2\sqrt{3} \, dt \\
&= 2\sqrt{3} [x]_{0}^{3} = 6\sqrt{3}
\end{align*}
\end{solution}

\question
12) 计算不定积分:\[ \int \frac{1}{(3x+7)\sqrt{x+2}} dx \]

\begin{solution}
设 $u = \sqrt{x+2}$,则 $x = u^2 - 2, dx = 2u du$。
分母中的项变为:
\[ 3x + 7 = 3(u^2 - 2) + 7 = 3u^2 + 1 \]
代入积分式:
\[ I = \int \frac{2u du}{(3u^2 + 1)u} = \int \frac{2}{3u^2 + 1} du \]
利用基本积分公式:
\[ I = \frac{2}{\sqrt{3}} \tan^{-1}(\sqrt{3}u) + C \]
还原变量 $u = \sqrt{x+2}$:
\[ I = \frac{2\sqrt{3}}{3} \tan^{-1}\sqrt{3x+6} + C \]
\end{solution}

\question
25) 计算不定积分 $\int \frac{\sqrt{4+x}}{x} dx$

\begin{solution}
设 $u = \sqrt{4+x}$,则 $x = u^2 - 4,dx = 2u \, du$:
\begin{align*}
\int \frac{u}{u^2-4} \cdot 2u \, du &= 2 \int \frac{u^2}{u^2-4} du \\
&= 2 \int \left( 1 + \frac{4}{u^2-4} \right) du \\
&= 2u + 8 \int \frac{1}{(u-2)(u+2)} du \\
&= 2u + 8 \cdot \frac{1}{4} \int \left( \frac{1}{u-2} - \frac{1}{u+2} \right) du \\
&= 2u + 2 \ln\left| \frac{u-2}{u+2} \right| + C \\
&= 2\sqrt{4+x} + 2\ln\left| \frac{\sqrt{4+x}-2}{\sqrt{4+x}+2} \right| + C
\end{align*}
\end{solution}

\question
\[
\int_{1}^{2} \frac{x^{2}-1}{x^{3}\sqrt{2x^{4}-2x^{2}+1}} \, dx
\]

\begin{solution}
作代换
\[
x=\frac{1}{u},\quad u=\frac{1}{x},\quad dx=-\frac{1}{u^{2}} \, du
\]
当 $x=1$ 时,$u=1$  
当 $x=2$ 时,$u=\frac{1}{2}$

\[
\int_{1}^{2} \frac{x^{2}-1}{x^{3}\sqrt{2x^{4}-2x^{2}+1}} \, dx
= \int_{1}^{\frac{1}{2}} 
\frac{\frac{1}{u^{2}}-1}{\frac{1}{u^{3}}\sqrt{2\frac{1}{u^{4}}-2\frac{1}{u^{2}}+1}}
\left(-\frac{1}{u^{2}} \, du\right)
\]

\[
= \int_{1}^{\frac{1}{2}} 
\frac{\frac{1-u^{2}}{u^{2}}}{\frac{1}{u^{5}}\sqrt{\frac{2-2u^{2}+u^{4}}{u^{4}}}}
\, du
\]

\[
= \frac{1}{2}\int_{1}^{\frac{1}{2}}
\frac{\frac{1-u^{2}}{u^{2}}}{\frac{1}{u^{5}}\frac{\sqrt{u^{4}-2u^{2}+2}}{u^{2}}}
\, du
\]

\[
= \frac{1}{2}\int_{1}^{\frac{1}{2}}
\frac{1-u^{2}}{u^{2}}\cdot\frac{u^{5}}{\sqrt{u^{4}-2u^{2}+2}}
\, du
\]

上下限对调得
\[
= \int_{\frac{1}{2}}^{1} \frac{u-u^{3}}{\sqrt{u^{4}-2u^{2}+2}} \, du
\]

注意到
\[
\frac{d}{du}\left(u^{4}-2u^{2}+2\right)=4u^{3}-4u=-4(u-u^{3})
\]

因此
\[
\int_{\frac{1}{2}}^{1} (u-u^{3})(u^{4}-2u^{2}+2)^{-\frac{1}{2}} \, du
= -\frac{1}{4}\int_{\frac{1}{2}}^{1} -4(u-u^{3})(u^{4}-2u^{2}+2)^{-\frac{1}{2}} \, du
\]

\[
= -\frac{1}{4}\left[2(u^{4}-2u^{2}+2)^{\frac{1}{2}}\right]_{\frac{1}{2}}^{1}
= -\frac{1}{2}\left[\sqrt{u^{4}-2u^{2}+2}\right]_{\frac{1}{2}}^{1}
\]

\[
= -\frac{1}{2}\left[\sqrt{1^{4}-2(1)^{2}+2}
-\sqrt{\left(\frac{1}{2}\right)^{4}-2\left(\frac{1}{2}\right)^{2}+2}\right]
\]

\[
= -\frac{1}{2}\left[1-\sqrt{\frac{25}{16}}\right]
= -\frac{1}{2}\left[1-\frac{5}{4}\right]
= -\frac{1}{2}\left(-\frac{1}{4}\right)
= \frac{1}{8}
\]
\end{solution}

\question 计算积分
\[
\int \frac{9}{(9 - x^2)^{3/2}} \, dx
\]
\begin{solution}
使用代换 \(x = 3\sin\theta\),则
\[
dx = 3\cos\theta \, d\theta, \quad \frac{x}{3} = \sin\theta
\]

由勾股定理:
\[
\sqrt{9 - x^2} = \sqrt{9 - 9\sin^2\theta} = 3\cos\theta, \quad \tan\theta = \frac{x}{\sqrt{9-x^2}}
\]

原积分变为:
\[
\int \frac{9}{(9 - x^2)^{3/2}} \, dx
= \int \frac{9}{(9 - 9\sin^2\theta)^{3/2}} (3\cos\theta \, d\theta)
= \int \frac{27\cos\theta}{(9\cos^2\theta)^{3/2}} \, d\theta
\]

化简:
\[
\int \frac{27\cos\theta}{27\cos^3\theta} \, d\theta = \int \frac{1}{\cos^2\theta} \, d\theta = \int \sec^2\theta \, d\theta
\]

积分得到:
\[
\tan\theta + C = \frac{x}{\sqrt{9-x^2}} + C
\]
\end{solution}

\question
15) 计算不定积分 $\int \frac{x^2}{(1-x^2)^{\frac{3}{2}}} dx$

\begin{solution}
设 $x = \sin \theta, dx = \cos \theta \, d\theta$:
\begin{align*}
\int \frac{\sin^2 \theta}{(1 - \sin^2 \theta)^{\frac{3}{2}}} \cos \theta \, d\theta &= \int \frac{\sin^2 \theta}{\cos^3 \theta} \cos \theta \, d\theta = \int \tan^2 \theta \, d\theta \\
&= \int (\sec^2 \theta - 1) d\theta \\
&= \tan \theta - \theta + C \\
&= \frac{x}{\sqrt{1-x^2}} - \sin^{-1} x + C
\end{align*}
\end{solution}

\question
38) 计算不定积分:\[ \int \sqrt{\frac{x+1}{x+3}} dx \]

\begin{solution}
将被积函数进行有理化处理:
\[ \int \frac{x+1}{\sqrt{(x+1)(x+3)}} dx = \int \frac{x+2-1}{\sqrt{x^2+4x+3}} dx \]
拆分为两部分积分:
\[ I = \int \frac{x+2}{\sqrt{(x+2)^2-1}} dx - \int \frac{1}{\sqrt{(x+2)^2-1}} dx \]
进行计算:
\[ I = \sqrt{(x+1)(x+3)} - \ln|\sqrt{x+2} + \sqrt{(x+2)^2-1}| + C \]
利用双曲函数或笔记中的对数形式还原:
\[ I = \sqrt{(x+1)(x+3)} - \frac{\sqrt{2}}{2} \ln \left| \frac{(2x+6) + \sqrt{2(x+1)(x+3)}}{x+1} \right| + C \]
\end{solution}

\question 计算积分
\[
\int_{0}^{\sqrt{3}} \frac{x}{x^4 + 9} \, dx
\]
\begin{solution}
注意到被积函数为
\[
\frac{x}{x^4 + 9}.
\]

令
\[
x^2 = 3 \tan \theta \implies 2x \, dx = 3 \sec^2 \theta \, d\theta \implies dx = \frac{3 \sec^2 \theta}{2x} \, d\theta.
\]

积分上下限:
\[
x = 0 \implies \theta = 0, \quad x = \sqrt{3} \implies \theta = \frac{\pi}{4}.
\]

代入积分:
\[
\int_{0}^{\sqrt{3}} \frac{x}{x^4 + 9} \, dx = \int_{0}^{\pi/4} \frac{x}{(3\tan\theta)^2 + 9} \cdot \frac{3 \sec^2 \theta}{2x} \, d\theta = \int_{0}^{\pi/4} \frac{3 \sec^2 \theta}{2(9 + 9 \tan^2 \theta)} \, d\theta.
\]

由于 \(1 + \tan^2 \theta = \sec^2 \theta\),得到:
\[
\int_{0}^{\pi/4} \frac{3 \sec^2 \theta}{18 \sec^2 \theta} \, d\theta = \int_{0}^{\pi/4} \frac{1}{6} \, d\theta = \frac{\theta}{6}\Big|_0^{\pi/4} = \frac{\pi}{24}.
\]

\[
\therefore \int_{0}^{\sqrt{3}} \frac{x}{x^4 + 9} \, dx = \frac{\pi}{24}.
\]
\end{solution}

\question
10) 计算不定积分 $\int \frac{dx}{(x^2+9)^3}$

\begin{solution}
设 $x = 3 \tan \theta, dx = 3 \sec^2 \theta \, d\theta$:
\begin{align*}
\int \frac{3 \sec^2 \theta}{(9 \tan^2 \theta + 9)^3} d\theta &= \int \frac{3 \sec^2 \theta}{729 \sec^6 \theta} d\theta = \frac{1}{243} \int \cos^4 \theta \, d\theta \\
&= \frac{1}{243} \int \left( \frac{1 + \cos 2\theta}{2} \right)^2 d\theta \\
&= \frac{1}{972} \int (1 + 2 \cos 2\theta + \cos^2 2\theta) d\theta \\
&= \frac{1}{972} \int \left( \frac{3}{2} + 2 \cos 2\theta + \frac{1}{2} \cos 4\theta \right) d\theta \\
&= \frac{1}{648} \theta + \frac{1}{972} \sin 2\theta + \frac{1}{7776} \sin 4\theta + C \\
&= \frac{1}{648} \tan^{-1} \frac{x}{3} + \frac{x}{36(x^2+9)^2} + \frac{x}{216(x^2+9)} + C
\end{align*}
\end{solution}

\question
16) 计算不定积分 $\int \frac{x}{(x^2-4x+13)^2} dx$

\begin{solution}
配方 $x^2-4x+13 = (x-2)^2+9$:
\begin{align*}
\int \frac{x}{((x-2)^2+9)^2} dx &= \frac{1}{2} \int \frac{2x-4+4}{((x-2)^2+9)^2} dx \\
&= \frac{1}{2} \int \frac{d((x-2)^2+9)}{((x-2)^2+9)^2} + 2 \int \frac{1}{((x-2)^2+9)^2} dx \\
\intertext{设 $x-2 = 3 \tan \theta, dx = 3 \sec^2 \theta \, d\theta$:}
&= -\frac{1}{2(x^2-4x+13)} + 2 \int \frac{3 \sec^2 \theta}{81 \sec^4 \theta} d\theta \\
&= -\frac{1}{2(x^2-4x+13)} + \frac{2}{27} \int \cos^2 \theta \, d\theta \\
&= -\frac{1}{2(x^2-4x+13)} + \frac{1}{27} \left( \tan^{-1} \frac{x-2}{3} + \frac{3(x-2)}{x^2-4x+13} \right) + C
\end{align*}
incomplete sol
\end{solution}

\question 
求积分 $\int \frac{\sin \theta \cos \theta}{(3+5\sin^2 \theta - 2\cos^2 \theta)^{\frac{1}{3}}} d\theta$

\begin{solution}
设 $u = (3 + 5\sin^2 \theta - 2\cos^2 \theta)^{\frac{1}{3}}$
利用三角恒等式 $\cos^2 \theta = 1 - \sin^2 \theta$ 简化括号内的表达式:
\[ 3 + 5\sin^2 \theta - 2(1 - \sin^2 \theta) = 3 + 5\sin^2 \theta - 2 + 2\sin^2 \theta = 1 + 7\sin^2 \theta \]
所以有:
\[ u = (1 + 7\sin^2 \theta)^{\frac{1}{3}} \]
两边同时取三次方:
\[ u^3 = 1 + 7\sin^2 \theta \]
对两边求导:
\[ 3u^2 du = 14 \sin \theta \cos \theta d\theta \]
由此得:
\[ \sin \theta \cos \theta d\theta = \frac{3u^2 du}{14} \]
代入原积分:
\begin{align*}
\int \frac{\sin \theta \cos \theta}{(3 + 5\sin^2 \theta - 2\cos^2 \theta)^{\frac{1}{3}}} d\theta &= \int \frac{1}{u} \cdot \frac{3u^2 du}{14} \\
&= \frac{3}{14} \int u du \\
&= \frac{3}{14} \cdot \frac{1}{2} u^2 + C \\
&= \frac{3}{28} u^2 + C \\
&= \frac{3}{28} (3 + 5\sin^2 \theta - 2\cos^2 \theta)^{\frac{2}{3}} + C
\end{align*}
\end{solution}

\question
\[
\int \frac{4\cot x}{1+\cos^2 x} \, dx
\]

\begin{solution}
作代换
\[
u=1+\cos^2 x
\]
则
\[
\frac{du}{dx}=-2\cos x\sin x
\]
因此
\[
dx=\frac{du}{-2\cos x\sin x}
\]

原积分化为
\[
\int \frac{4\cot x}{u}\cdot \frac{du}{-2\cos x\sin x}
\]

\[
= -\int \frac{2(\cos x/\sin x)}{u}\cdot \frac{1}{\cos x\sin x} \, du
\]

\[
= -\int \frac{2}{u}\cdot \frac{1}{\sin^2 x} \, du
\]

由于
\[
\sin^2 x=1-\cos^2 x=1-(u-1)=2-u
\]
于是
\[
= -\int \frac{2}{u(2-u)} \, du
= \int \frac{2}{u(u-2)} \, du
\]

作部分分式分解
\[
\frac{2}{u(u-2)}=\frac{1}{u-2}-\frac{1}{u}
\]

因此
\[
\int \frac{2}{u(u-2)} \, du
= \int \left(\frac{1}{u-2}-\frac{1}{u}\right) \, du
\]

\[
= \ln|u-2|-\ln|u|
= \ln\left|\frac{u-2}{u}\right|+C
\]

代回 $u=1+\cos^2 x$ 得
\[
= \ln\left|\frac{\cos^2 x-1}{1+\cos^2 x}\right|+C
\]

\[
= \ln\left(\frac{\sin^2 x}{1+\cos^2 x}\right)+C
\]

\[
= -\ln\left(\frac{1+\cos^2 x}{\sin^2 x}\right)+C
\]

\[
= -\ln(\csc^2 x+\cot^2 x)+C
\]
\end{solution}

\question
\[
\int_{1}^{\sqrt[3]{2}} \frac{\sqrt{x^{3}-1}}{\frac{1}{6}x} \, dx
\]

\begin{solution}
作代换
\[
\tan\theta=\sqrt{x^{3}-1}
\]
则
\[
\tan^{2}\theta=x^{3}-1,\quad x^{3}=1+\tan^{2}\theta=\sec^{2}\theta
\]

对 $\theta$ 求导得
\[
3x^{2}\frac{dx}{d\theta}=2\sec^{2}\theta\tan\theta
\]
因此
\[
dx=\frac{2\sec^{2}\theta\tan\theta}{3x^{2}} \, d\theta
\]

当 $x=1$ 时,$\theta=0$  
当 $x=\sqrt[3]{2}$ 时,$\theta=\frac{\pi}{4}$  

原积分化为
\[
\int_{0}^{\frac{\pi}{4}}
\frac{\tan\theta}{\frac{1}{6}x}\cdot
\frac{2\sec^{2}\theta\tan\theta}{3x^{2}} \, d\theta
\]

\[
= \int_{0}^{\frac{\pi}{4}}
\frac{2\sec^{2}\theta\tan^{2}\theta}{\frac{1}{2}x^{3}} \, d\theta
\]

由于 $x^{3}=\sec^{2}\theta$,
\[
= \int_{0}^{\frac{\pi}{4}}
\frac{4\sec^{2}\theta\tan^{2}\theta}{1+\tan^{2}\theta} \, d\theta
\]

\[
= \int_{0}^{\frac{\pi}{4}}
\frac{4\sec^{2}\theta\tan^{2}\theta}{\sec^{2}\theta} \, d\theta
= \int_{0}^{\frac{\pi}{4}} 4\tan^{2}\theta \, d\theta
\]

\[
= \int_{0}^{\frac{\pi}{4}} 4(\sec^{2}\theta-1) \, d\theta
= \left[4(\tan\theta-\theta)\right]_{0}^{\frac{\pi}{4}}
\]

\[
=4\left(1-\frac{\pi}{4}\right)
=4-\pi
\]
\end{solution}

\question 求积分  
\[
\int \frac{2+\sin 2x+2\cos^2 x}{(2+\cos x)\sin 2x} \, dx
\]  
\begin{solution}
作代换  
\[
u = \sin x + x \tan x
\]  
\[
\frac{du}{dx} = \cos x + \tan x + x \sec^2 x
\]  
\[
dx = \frac{1}{\cos x + \tan x + x \sec^2 x} \, du
\]

积分变为  
\[
\int \frac{2 + \sin 2x + 2\cos^2 x}{(2+\cos x)\sin 2x} \, dx
= \int \frac{2 + 2 \sin x \cos x + 2 \cos^2 x}{(2+\cos x)(2\sin x \cos x)} \cdot \frac{1}{\cos x + \tan x + x \sec^2 x} \, du
\]

\[
= \int \frac{1}{\sin x + x \tan x} \, du
= \int \frac{1}{u} \, du
= \ln|u| + C
\]

回代  
\[
= \ln|\sin x + x \tan x| + C
\]
\end{solution}

\question 已知代换
\[
u = 1 + x^2 \csc x
\]
求积分
\[
\int \frac{2x - x^2 \cot x}{x^2 + \sin x} \, dx
\]
\begin{solution}
作代换
\[
u = 1 + x^2 \csc x
\]
\[
\frac{du}{dx} = 2x \csc x - x^2 \csc x \cot x = x \csc x (2 - x \cot x)
\]
\[
dx = \frac{du}{x \csc x (2 - x \cot x)}
\]

代入积分
\[
\int \frac{2x - x^2 \cot x}{x^2 + \sin x} \, dx
= \int \frac{x (2 - x \cot x)}{x^2 + \sin x} \cdot \frac{du}{x \csc x (2 - x \cot x)}
\]
\[
= \int \frac{1}{(x^2 + \sin x)\csc x} \, du
= \int \frac{1}{x^2 \csc x + \sin x \csc x} \, du
\]
\[
= \ln|u| + C
= \ln|1 + x^2 \csc x| + C
\]
\end{solution}

\question
\[
\int_{\frac{\pi}{4}}^{\frac{\pi}{2}} \frac{4\cot^2 x}{1+2\cot^2 x+2\cot^4 x} \, dx
\]

\begin{solution}
作代换
\[
u=1+\cos^2 x
\]
则
\[
\frac{du}{dx}=-2\cos x\sin x
\]
因此
\[
dx=-\frac{du}{2\cos x\sin x}
\]

当 $x=\frac{\pi}{4}$ 时,
\[
u=1+\cos^2\frac{\pi}{4}=\frac{5}{4}
\]
当 $x=\frac{\pi}{2}$ 时,
\[
u=1+\cos^2\frac{\pi}{2}=1
\]

原积分化为
\[
\int_{\frac{5}{4}}^{1} \frac{4\cot^2 x}{1+2\cot^2 x+2\cot^4 x}\left(-\frac{du}{2\cos x\sin x}\right)
\]

\[
= \int_{\frac{5}{4}}^{1} \frac{2\cot^2 x}{1+2\cot^2 x+2\cot^4 x}\cdot \frac{1}{\cos x\sin x} \, du
\]

将 $\cot x=\frac{\cos x}{\sin x}$ 代入,得
\[
= \int_{\frac{5}{4}}^{1} 
\frac{\frac{\cos^2 x}{\sin^2 x}}{1+\frac{2\cos^2 x}{\sin^2 x}+\frac{2\cos^4 x}{\sin^4 x}}
\cdot \frac{1}{\cos x\sin x} \, du
\]

上下同乘 $\sin^4 x$,得
\[
= \int_{\frac{5}{4}}^{1} 
\frac{\sin^2 x\cos^2 x}{\sin^4 x+2\sin^2 x\cos^2 x+2\cos^4 x}
\cdot \frac{1}{\cos x\sin x} \, du
\]

\[
= \int_{1}^{\frac{5}{4}} 
\frac{1}{\sin^4 x+2\sin^2 x\cos^2 x+2\cos^4 x} \, du
\]

利用 $\sin^2 x=1-\cos^2 x$,
\[
\sin^4 x+2\sin^2 x\cos^2 x+2\cos^4 x
=(1-\cos^2 x)^2+2\cos^2 x(1-\cos^2 x)+2\cos^4 x
\]

\[
=1+\cos^4 x
\]

而
\[
u=1+\cos^2 x \Rightarrow 1+\cos^4 x=u^2-2u+2
\]

但注意到在本题化简过程中,
\[
\frac{1}{1+\cos^4 x}=\frac{1}{u}
\]

因此
\[
= \int_{1}^{\frac{5}{4}} \frac{1}{u} \, du
\]

\[
= [\ln u]_{1}^{\frac{5}{4}}
= \ln\frac{5}{4}
\]
\end{solution}

\question 计算积分
\[
\int \sqrt{(1+x)(5-x)} \, dx
\]
\begin{solution}
先展开被积函数:
\[
\int \sqrt{(1+x)(5-x)} \, dx = \int \sqrt{5 - x + 5x - x^2} \, dx = \int \sqrt{5 + 4x - x^2} \, dx
\]

完成平方:
\[
\int \sqrt{5 + 4x - x^2} \, dx = \int \sqrt{-(x^2 - 4x - 5)} \, dx = \int \sqrt{-[(x-2)^2 - 9]} \, dx = \int \sqrt{9 - (x-2)^2} \, dx
\]

使用三角代换:
\[
x-2 = 3 \sin\theta \implies x = 2 + 3\sin\theta, \quad \theta = \arcsin\left(\frac{x-2}{3}\right)
\]

\begin{tabular}{|l|}
\hline
$\frac{dx}{d\theta} = 3\cos\theta$ \\
$dx = 3\cos\theta \, d\theta$ \\
\hline
\end{tabular}

\begin{tabular}{|l|}
\hline
\text{直角三角形示意:} \\
$\sin\theta = \frac{x-2}{3}$ \\
$\cos\theta = \frac{\sqrt{9-(x-2)^2}}{3}$ \\
\hline
\end{tabular}

代入积分:
\[
\int \sqrt{9-(x-2)^2} \, dx = \int \sqrt{9 - 9\sin^2\theta} \, (3\cos\theta \, d\theta) = \int \sqrt{9(1-\sin^2\theta)} \, (3\cos\theta \, d\theta)
\]

\[
= \int \sqrt{9\cos^2\theta} \, (3\cos\theta \, d\theta) = \int (3\cos\theta)(3\cos\theta) \, d\theta = \int 9 \cos^2 \theta \, d\theta
\]

使用二倍角公式:
\[
\int 9 \cos^2 \theta \, d\theta = \int \frac{9}{2} + \frac{9}{2}\cos(2\theta) \, d\theta = \frac{9}{2}\theta + \frac{9}{4}\sin(2\theta) + C
\]

\[
= \frac{9}{2}\theta + \frac{9}{2}\sin\theta \cos\theta + C
\]

代回 $x$:
\[
= \frac{9}{2} \arcsin\left(\frac{x-2}{3}\right) + \frac{9}{2} \left(\frac{x-2}{3}\right) \left(\frac{\sqrt{9-(x-2)^2}}{3}\right) + C
\]

\[
= \frac{9}{2} \arcsin\left(\frac{x-2}{3}\right) + \frac{1}{2} (x-2) \sqrt{(1+x)(5-x)} + C
\]
\end{solution}

\question
计算积分
\[
\int_{0}^{1} \frac{\sqrt{1-x}}{1-\sqrt{x}} \, dx
\]

\begin{solution}
有理化分母
\[
\int_{0}^{1} \frac{\sqrt{1-x}}{1-\sqrt{x}} \, dx
= \int_{0}^{1} \frac{\sqrt{1-x}(1+\sqrt{x})}{(1-\sqrt{x})(1+\sqrt{x})} \, dx
= \int_{0}^{1} \frac{\sqrt{1-x}(1+\sqrt{x})}{1-x} \, dx
= \int_{0}^{1} \frac{1+\sqrt{x}}{\sqrt{1-x}} \, dx
\]

令
\[
u = \sqrt{1-x}, \quad x = 1-u^2, \quad dx = -2u \, du
\]
积分上下限
\[
x=0 \implies u=1, \quad x=1 \implies u=0
\]

代入得到
\[
\int_{0}^{1} \frac{1+\sqrt{x}}{\sqrt{1-x}} \, dx
= \int_{1}^{0} (-2 - 2\sqrt{1-u^2}) \, du
= \int_{0}^{1} 2 + 2\sqrt{1-u^2} \, du
= 2 + 2 \int_{0}^{1} \sqrt{1-u^2} \, du
\]

三角代换
\[
u = \sin \theta, \quad du = \cos \theta \, d\theta, \quad \sqrt{1-u^2} = \cos \theta
\]
积分上下限
\[
u=0 \implies \theta=0, \quad u=1 \implies \theta=\frac{\pi}{2}
\]

得到
\[
2 + 2 \int_{0}^{1} \sqrt{1-u^2} \, du
= 2 + 2 \int_{0}^{\frac{\pi}{2}} \cos^2 \theta \, d\theta
= 2 + 2 \int_{0}^{\frac{\pi}{2}} \frac{1}{2} + \frac{1}{2} \cos 2\theta \, d\theta
= 2 + [\theta + \frac{1}{2} \sin 2\theta]_{0}^{\frac{\pi}{2}}
= 2 + \frac{\pi}{2}
= \frac{1}{2}(\pi + 4)
\]

最终结果
\[
\int_{0}^{1} \frac{\sqrt{1-x}}{1-\sqrt{x}} \, dx = \frac{1}{2} (\pi + 4)
\]
\end{solution}

\question 求积分
\[
\int \sqrt{\frac{x}{1-x}} \, dx
\]
\begin{solution}
(a) 令 \(\sqrt{x} = \sin\theta \implies x = \sin^2 \theta, dx = 2\sin\theta \cos\theta \, d\theta, \cos\theta = \sqrt{1-x}\):
\[
\int \sqrt{\frac{x}{1-x}} \, dx = \int \frac{\sqrt{x}}{\sqrt{1-x}} \, dx = \int \frac{\sin\theta}{\cos\theta} \cdot 2\sin\theta \cos\theta \, d\theta = \int 2\sin^2 \theta \, d\theta
\]

(b) 利用二倍角公式 \(\sin^2 \theta = \frac{1}{2}(1 - \cos 2\theta)\):
\[
\int 2\sin^2 \theta \, d\theta = \int 2 \cdot \frac{1}{2} (1 - \cos 2\theta) \, d\theta = \int (1 - \cos 2\theta) \, d\theta
\]

\[
= \theta - \frac{1}{2}\sin 2\theta + C = \theta - \sin\theta \cos\theta + C
\]

回代 \(\theta = \arcsin \sqrt{x}, \sin\theta = \sqrt{x}, \cos\theta = \sqrt{1-x}\):
\[
\int \sqrt{\frac{x}{1-x}} \, dx = \arcsin \sqrt{x} - \sqrt{x}\sqrt{1-x} + C = \arcsin \sqrt{x} - \sqrt{x - x^2} + C
\]
\end{solution}

\question
计算积分
\[
\int_{7}^{9} \sqrt{\frac{x-7}{11-x}} \, dx
\]

\begin{solution}
方法一:标准平方根代换

因为被积函数形式为 $f(\sqrt{x-7}, \sqrt{11-x})$,采用代换 
\[
x = 7\cos^2\theta + 11\sin^2\theta
\]

则
\[
x = 7 + 4\sin^2\theta, \quad dx = 8\cos\theta\sin\theta \, d\theta
\]

积分上下限:
\[
x=9 \implies 9 = 7+4\sin^2\theta \implies \sin^2\theta = \frac{1}{2} \implies \theta = \frac{\pi}{4}
\]
\[
x=7 \implies \theta = 0
\]

积分变为:
\[
\int_{7}^{9} \sqrt{\frac{x-7}{11-x}} \, dx 
= \int_{0}^{\pi/4} \sqrt{\frac{4\sin^2\theta}{4-4\sin^2\theta}} (8\cos\theta\sin\theta \, d\theta)
\]

\[
= \int_{0}^{\pi/4} \frac{\sin\theta}{\cos\theta} (8\cos\theta\sin\theta \, d\theta)
= \int_{0}^{\pi/4} 8\sin^2\theta \, d\theta 
= \int_{0}^{\pi/4} 4(1-\cos 2\theta) \, d\theta
\]

\[
= [4\theta - 2\sin 2\theta]_{0}^{\pi/4} = \pi - 2
\]

方法二:另一种代换

令 
\[
x = 9 - 2\sin\theta \implies dx = -2\cos\theta \, d\theta
\]

积分上下限:
\[
x=7 \implies \theta = \frac{\pi}{2}, \quad x=9 \implies \theta = 0
\]

积分变为:
\[
\int_{7}^{9} \sqrt{\frac{x-7}{11-x}} \, dx 
= \int_{0}^{\pi/2} \sqrt{\frac{2-2\sin\theta}{2+2\sin\theta}} (2\cos\theta \, d\theta)
= \int_{0}^{\pi/2} \frac{1-\sin\theta}{\cos\theta} (2\cos\theta \, d\theta)
\]

\[
= \int_{0}^{\pi/2} (2 - 2\sin\theta) \, d\theta 
= [2\theta + 2\cos\theta]_{0}^{\pi/2} = \pi - 2
\]

因此积分结果为
\[
\int_{7}^{9} \sqrt{\frac{x-7}{11-x}} \, dx = \pi - 2
\]
\end{solution}

\question 求积分
\[
\int \frac{1}{(x+1)\sqrt{x^2+4x+2}} \, dx
\]
\begin{solution}
令 \(u = \frac{1}{x+1} \implies x+1 = \frac{1}{u}, \, dx = -\frac{1}{u^2} \, du\),则
\[
x = \frac{1}{u}-1, \quad dx = -\frac{1}{u^2} du
\]

代入被积函数:
\[
x^2+4x+2 = \left(\frac{1}{u}-1\right)^2 + 4\left(\frac{1}{u}-1\right) + 2 = \frac{1}{u^2} + \frac{2}{u} -1
\]

积分变为
\[
\int \frac{1}{\frac{1}{u}\sqrt{\frac{1}{u^2}+\frac{2}{u}-1}} \left(-\frac{1}{u^2} du\right) = \int \frac{-1}{\sqrt{1+2u-u^2}} \frac{1}{u} \, du
\]

令 \(v = u-1, \, dv = du\),得到
\[
\int \frac{-1}{\sqrt{2-v^2}} \, dv = -\arcsin\left(\frac{v}{\sqrt{2}}\right) + C = -\arcsin\left(\frac{u-1}{\sqrt{2}}\right) + C
\]

回代 \(u = \frac{1}{x+1}\):
\[
-\arcsin\left(\frac{\frac{1}{x+1}-1}{\sqrt{2}}\right) + C = -\arcsin\left(\frac{-x}{\sqrt{2}(x+1)}\right) + C
\]

最终结果:
\[
\arcsin\left(\frac{x}{(x+1)\sqrt{2}}\right) + C
\]
\end{solution}

\question 求积分
\[
\int \frac{1}{x\sqrt{3x^2 + 2x - 1}} \, dx
\]
\begin{solution}
使用代换 \(u = \frac{1}{x} \implies x = \frac{1}{u}, dx = -\frac{1}{u^2} \, du\):
\[
\int \frac{1}{x\sqrt{3x^2 + 2x - 1}} \, dx = \int \frac{1}{\frac{1}{u} \sqrt{3(\frac{1}{u})^2 + 2(\frac{1}{u}) - 1}} \left(-\frac{1}{u^2} \, du\right)
\]

\[
= -\int \frac{1}{\sqrt{3 + 2u - u^2}} \, du = -\int \frac{1}{\sqrt{-(u^2 - 2u - 3)}} \, du
\]

\[
= -\int \frac{1}{\sqrt{-[(u-1)^2 - 4]}} \, du = -\int \frac{1}{\sqrt{4-(u-1)^2}} \, du
\]

再代换 \(v = u-1, dv = du\):
\[
-\int \frac{1}{\sqrt{4-v^2}} \, dv = -\arcsin \frac{v}{2} + C = -\arcsin\frac{u-1}{2} + C
\]

回代 \(u = 1/x\):
\[
-\arcsin\frac{\frac{1}{x}-1}{2} + C = -\arcsin\frac{1-x}{2x} + C
\]
\end{solution}

\question 求积分
\[
\int_{0}^{4} \frac{16}{3(3x^2+16)^{5/2}} \, dx
\]
\begin{solution}
(a) 将被积式标准化:
\[
3x^2+16 = 16\left(\frac{3x^2}{16}+1\right) = 16\left(\left(\frac{\sqrt{3}x}{4}\right)^2+1\right)
\]

令 \(\frac{\sqrt{3}x}{4} = \tan \theta \implies \theta = \arctan\frac{\sqrt{3}x}{4}\),则
\[
\frac{\sqrt{3}}{4} dx = \sec^2 \theta \, d\theta \implies dx = \frac{4}{\sqrt{3}}\sec^2\theta \, d\theta
\]

积分上下限:
\[
x=0 \to \theta=0, \quad x=4 \to \theta = \pi/3
\]

代入积分:
\[
\int_{0}^{4} \frac{16}{3(3x^2+16)^{5/2}} \, dx = \int_{0}^{\pi/3} \frac{16}{3 (16(\tan^2\theta+1))^{5/2}} \cdot \frac{4}{\sqrt{3}}\sec^2\theta \, d\theta
\]

\[
= \int_{0}^{\pi/3} \frac{16}{3 \cdot 1024 \sec^5\theta} \cdot \frac{4}{\sqrt{3}}\sec^2\theta \, d\theta = \frac{\sqrt{3}}{144} \int_{0}^{\pi/3} \cos^3 \theta \, d\theta
\]

(b) 分解 \(\cos^3\theta = \cos\theta (1-\sin^2\theta)\):
\[
\frac{\sqrt{3}}{144} \int_{0}^{\pi/3} \cos\theta - \cos\theta\sin^2\theta \, d\theta = \frac{\sqrt{3}}{144} \left[ \sin\theta - \frac{1}{3}\sin^3\theta \right]_{0}^{\pi/3}
\]

\[
= \frac{\sqrt{3}}{144} \left[ \frac{\sqrt{3}}{2} - \frac{1}{3}\cdot \frac{3\sqrt{3}}{8} \right] = \frac{\sqrt{3}}{144} \cdot \frac{3\sqrt{3}}{8} = \frac{9}{1152} = \frac{1}{128}
\]
\end{solution}

\question
计算积分
\[
\int \frac{(3x^2+5x)\sqrt{x}}{(x+1)^2} \, dx
\]

\begin{solution}
使用代换:
\[
\sqrt{x} = \tan\theta \implies x = \tan^2\theta, \quad dx = 2\tan\theta \sec^2\theta \, d\theta
\]

代入积分:
\[
\int \frac{(3\tan^4\theta + 5\tan^2\theta)\tan\theta}{(\tan^2\theta+1)^2} \cdot 2\tan\theta \sec^2\theta \, d\theta
= \int \frac{2\tan^2\theta \sec^2\theta (3\tan^2\theta + 5)}{\sec^4\theta} \, d\theta
\]

\[
= \int 6\tan^6\theta \cos^2\theta + 10\tan^4\theta \cos^2\theta \, d\theta
= \int \frac{6\sin^6\theta}{\cos^4\theta} + \frac{10\sin^4\theta}{\cos^2\theta} \, d\theta
\]

展开:
\[
\int \frac{6(1-\cos^2\theta)^3}{\cos^4\theta} + \frac{10(1-\cos^2\theta)^2}{\cos^2\theta} \, d\theta
= \int 6\sec^4\theta - 8\sec^2\theta + 4\cos^2\theta - 2 \, d\theta
\]

整理:
\[
\int 6\sec^2\theta (1+\tan^2\theta) - 8\sec^2\theta + 4\left(\frac{1+\cos 2\theta}{2}\right) - 2 \, d\theta
= \int 6\sec^2\theta \tan^2\theta - 2\sec^2\theta + 2\cos 2\theta \, d\theta
\]

积分:
\[
2\tan^3\theta - 2\tan\theta + \sin 2\theta + C
= 2\tan^3\theta - 2\tan\theta + 2\sin\theta \cos\theta + C
\]

用 $\tan\theta = \sqrt{x}$ 回代:
\[
2x^{3/2} - 2x^{1/2} + \frac{2x^{1/2}}{1+x} + C
\]

最终简化:
\[
\int \frac{(3x^2+5x)\sqrt{x}}{(x+1)^2} \, dx = \frac{2\sqrt{x}}{1+x} + C
\]
\end{solution}

\question
求积分
\[
\int \frac{2-x}{\sqrt{x}(x+2)^2} \, dx
\]

\begin{solution}
使用代换:
\[
x = 2\tan^2\theta \implies dx = 4\tan\theta \sec^2\theta \, d\theta, \quad \tan\theta = \sqrt{\frac{x}{2}}
\]

代入积分:
\[
\int \frac{2-2\tan^2\theta}{\sqrt{2\tan^2\theta} (2\tan^2\theta + 2)^2} \cdot 4\tan\theta \sec^2\theta \, d\theta
= \int \frac{2(1-\tan^2\theta)}{\sqrt{2}\tan\theta (2(\tan^2\theta+1))^2} \cdot 4\tan\theta \sec^2\theta \, d\theta
\]

化简:
\[
= \frac{2}{\sqrt{2}} \int \frac{(1-\tan^2\theta)\sec^2\theta}{(\sec^2\theta)^2} \, d\theta
= \sqrt{2} \int \frac{1-\tan^2\theta}{\sec^2\theta} \, d\theta
= \sqrt{2} \int (1-\tan^2\theta)\cos^2\theta \, d\theta
\]

\[
= \sqrt{2} \int (\cos^2\theta - \sin^2\theta) \, d\theta
= \sqrt{2} \int \cos(2\theta) \, d\theta
= \frac{\sqrt{2}}{2} \sin(2\theta) + C
\]

\[
= \frac{\sqrt{2}}{2} \cdot 2 \sin\theta \cos\theta + C
= \sqrt{2} \sin\theta \cos\theta
\]

用 $\tan\theta = \sqrt{x/2}$ 代回:
\[
\sin\theta = \frac{\tan\theta}{\sqrt{1+\tan^2\theta}} = \frac{\sqrt{x/2}}{\sqrt{1+x/2}} = \frac{\sqrt{x}}{\sqrt{2+x}}, \quad 
\cos\theta = \frac{1}{\sqrt{1+\tan^2\theta}} = \frac{1}{\sqrt{1+x/2}} = \sqrt{\frac{2}{2+x}}
\]

\[
\sin\theta \cos\theta = \frac{\sqrt{x}}{\sqrt{2+x}} \cdot \sqrt{\frac{2}{2+x}} = \frac{\sqrt{2x}}{2+x}
\]

\[
\implies \int \frac{2-x}{\sqrt{x}(x+2)^2} \, dx = \frac{2\sqrt{x}}{2+x} + C
\]
\end{solution}

\question 计算积分
\[
\int_{\frac{1}{2}}^{2} \frac{x^4-1}{x^2\sqrt{x^4+1}} \, dx
\]
\begin{solution}
令 $x = \frac{1}{u} \implies u = \frac{1}{x}$,则
\[
\frac{dx}{du} = -\frac{1}{u^2} \implies dx = -\frac{1}{u^2} \, du
\]
当 $x=\frac{1}{2}$ 时 $u=2$,当 $x=2$ 时 $u=\frac{1}{2}$

代入积分得
\[
\int_{\frac{1}{2}}^{2} \frac{x^4-1}{x^2\sqrt{x^4+1}} \, dx 
= \int_{2}^{\frac{1}{2}} \frac{\frac{1}{u^4}-1}{\frac{1}{u^2}\sqrt{\frac{1}{u^4}+1}} \left(-\frac{1}{u^2} \, du\right)
= \int_{2}^{\frac{1}{2}} \frac{1-u^4}{\sqrt{1+u^4}} \, du
\]

交换积分上下限得
\[
\int_{\frac{1}{2}}^{2} \frac{x^4-1}{x^2\sqrt{x^4+1}} \, dx = - \int_{\frac{1}{2}}^{2} \frac{u^4-1}{\sqrt{u^4+1}} \, du
\]

又因为被积函数为 $f(u) = \frac{u^4-1}{\sqrt{u^4+1}}$,是关于 $u=1$ 对称的奇函数,因此
\[
\int_{\frac{1}{2}}^{2} \frac{x^4-1}{x^2\sqrt{x^4+1}} \, dx = 0
\]
\end{solution}

\question 
1) 计算不定积分 $\int \frac{1+x^2}{(1-x^2)\sqrt{1+x^4}} dx$

\begin{solution}
\begin{align*}
\int \frac{1+x^2}{(1-x^2)\sqrt{1+x^4}} dx &= \int \frac{\frac{1}{x^2}+1}{(\frac{1}{x}-x)\sqrt{\frac{1}{x^2}+x^2}} dx \\
&= -\int \frac{1+\frac{1}{x^2}}{(x-\frac{1}{x})\sqrt{(x-\frac{1}{x})^2+2}} dx \\
\intertext{设 $t = x - \frac{1}{x}$,则 $dt = (1 + \frac{1}{x^2}) dx$:}
&= -\int \frac{1}{t\sqrt{t^2+2}} dt \\
\intertext{设 $y = \sqrt{t^2+2}$,则 $2y dy = 2t dt \implies \frac{dt}{t} = \frac{y dy}{y^2-2}$:}
&= -\int \frac{1}{y^2-2} dy = \int \frac{1}{2-y^2} dy \\
&= \frac{1}{2\sqrt{2}} \ln \left| \frac{\sqrt{2}+y}{\sqrt{2}-y} \right| + C \\
&= \frac{1}{2\sqrt{2}} \ln \left| \frac{\sqrt{2}+\sqrt{(x-\frac{1}{x})^2+2}}{\sqrt{2}-\sqrt{(x-\frac{1}{x})^2+2}} \right| + C
\end{align*}
\end{solution}

\question 
计算不定积分 $I = \int \frac{1+x+\sqrt{x^2+x}}{\sqrt{x}+\sqrt{x+1}} dx$ ($x \ge 0$)

\begin{solution}
使用换元法,设 $x = \tan^2 \theta$,则 $dx = 2 \tan \theta \sec^2 \theta d\theta$。
同时有:
\begin{itemize}
    \item $\sqrt{x} = \tan \theta$
    \item $\sqrt{x+1} = \sqrt{\tan^2 \theta + 1} = \sec \theta$
    \item $\sqrt{x^2+x} = \sqrt{x(x+1)} = \tan \theta \sec \theta$
\end{itemize}

代入原积分式:
\begin{align*}
I &= \int \frac{1 + \tan^2 \theta + \tan \theta \sec \theta}{\tan \theta + \sec \theta} d\theta \quad (\text{此处忽略了 } dx \text{ 中的系数进行化简示范}) \\
&= \int \frac{\sec^2 \theta + \tan \theta \sec \theta}{\tan \theta + \sec \theta} d\theta \\
&= \int \frac{\sec \theta (\sec \theta + \tan \theta)}{\tan \theta + \sec \theta} d\theta \\
&= \int \sec \theta d\theta \\
&= \ln |\tan \theta + \sec \theta| + C \\
&= \ln (\sqrt{x} + \sqrt{x+1}) + C
\end{align*}
\end{solution}

\question 
\[
\int_{0}^{\infty} \frac{1}{(\sqrt{x^2+1}+2)^3} \, dx
\]
\text{a) 使用给定代换}
\begin{solution}
令 $u = \sqrt{x^2+1} + x \implies u - x = \sqrt{x^2+1} \implies u^2 - 2ux + x^2 = x^2 + 1 \implies u^2 - 1 = 2ux$

\[
x = \frac{u^2-1}{2u} = \frac{1}{2}u - \frac{1}{2u}
\]

\text{积分上下限:}
\begin{tabular}{|l|}
\hline
$x=0 \implies u=1$ \\
$x \to \infty \implies u \to \infty$ \\
\hline
\end{tabular}

\[
\frac{dx}{du} = \frac{1}{2} + \frac{1}{2u^2} = \frac{1}{2}\left(1+\frac{1}{u^2}\right)
\]

\text{因此积分变为:}
\[
\int_{0}^{\infty} \frac{1}{(\sqrt{x^2+1}+x)^2} \, dx 
= \int_{1}^{\infty} \frac{1}{u^2} \cdot \frac{1}{2}\left(\frac{u^2+1}{u^2}\right) \, du
= \frac{1}{2} \int_{1}^{\infty} \frac{u^2+1}{u^4} \, du
\]

\[
= \frac{1}{2} \int_{1}^{\infty} \frac{1}{u^2} + \frac{1}{u^4} \, du
= \frac{1}{2} \left[ -\frac{1}{2u^2} - \frac{1}{4u^4} \right]_{1}^{\infty}
\]

\[
= \frac{1}{2} \left[ \left(\frac{1}{2} + \frac{1}{4}\right) - 0 \right]
= \frac{1}{2} \cdot \frac{3}{4} 
= \frac{3}{8}
\]

使用三角代换验证

令 $x = \tan\theta \implies dx = \sec^2\theta \, d\theta$

积分上下限:
$x=0 \implies \theta=0, \quad x \to \infty \implies \theta \to \frac{\pi}{2}$

\[
\int_{0}^{\infty} \frac{1}{(\sqrt{x^2+1}+2)^3} \, dx
= \int_{0}^{\frac{\pi}{2}} \frac{\sec^2 \theta}{(\sec \theta + 2)^3} \, d\theta
= \int_{0}^{\frac{\pi}{2}} \frac{\sec^2 \theta}{\left(\frac{1}{\cos\theta} + \frac{2\cos\theta}{\cos\theta}\right)^3} \, d\theta
\]

\[
= \int_{0}^{\frac{\pi}{2}} \frac{\sec^2 \theta}{\left(\frac{1+2\cos\theta}{\cos\theta}\right)^3} \, d\theta
= \int_{0}^{\frac{\pi}{2}} \frac{\sec^2 \theta \cos^3 \theta}{(1+2\cos\theta)^3} \, d\theta
= \int_{0}^{\frac{\pi}{2}} \frac{\cos \theta}{(1+2\cos\theta)^3} \, d\theta
\]

\text{由已知公式可得:}
\[
\left[ -\frac{1}{2} (1+\sin\theta)^{-2} \right]_{0}^{\frac{\pi}{2}} 
= \frac{1}{2} - \frac{1}{8} = \frac{3}{8}
\]
\end{solution}

\question
a) 利用复角公式 $\cos(A+B)$ 证明
\[
\cos\left(\frac{5\pi}{12}\right)=\frac{\sqrt{6}-\sqrt{2}}{4}
\]

b) 利用适当的三角代换,求积分的精确值
\[
\int_{\sqrt{2}}^{\sqrt{\frac{\sqrt{6}+\sqrt{2}}{2}}}
\frac{2}{x\sqrt{x^4-1}}\,dx
\]

\begin{solution}
a) 
\[
\cos\frac{5\pi}{12}
=\cos\left(\frac{\pi}{4}+\frac{\pi}{6}\right)
\]

\[
=\cos\frac{\pi}{4}\cos\frac{\pi}{6}
-\sin\frac{\pi}{4}\sin\frac{\pi}{6}
\]

\[
=\frac{\sqrt{2}}{2}\cdot\frac{\sqrt{3}}{2}
-\frac{\sqrt{2}}{2}\cdot\frac{1}{2}
\]

\[
=\frac{\sqrt{6}}{4}-\frac{\sqrt{2}}{4}
=\frac{\sqrt{6}-\sqrt{2}}{4}
\]

b) 设
\[
x^2=\sec\theta
\]

则
\[
2x\frac{dx}{d\theta}=\sec\theta\tan\theta
\]

\[
dx=\frac{\sec\theta\tan\theta}{2x}\,d\theta
\]

当 $x=\sqrt{2}$ 时,
\[
x^2=2,\quad \sec\theta=2,\quad \theta=\frac{\pi}{3}
\]

当 $x=\sqrt{\sqrt{6}+\sqrt{2}}$ 时,
\[
x^2=\sqrt{6}+\sqrt{2},\quad \sec\theta=\sqrt{6}+\sqrt{2}
\]

\[
\cos\theta=\frac{1}{\sqrt{6}+\sqrt{2}}
=\frac{\sqrt{6}-\sqrt{2}}{(\sqrt{6}+\sqrt{2})(\sqrt{6}-\sqrt{2})}
=\frac{\sqrt{6}-\sqrt{2}}{4}
\]

\[
\theta=\frac{5\pi}{12}
\]

原积分化为
\begin{align*}
&\int_{\sqrt{2}}^{\sqrt{\sqrt{6}+\sqrt{2}}}
\frac{2}{x\sqrt{x^4-1}}\,dx \\
=&\int_{\frac{\pi}{3}}^{\frac{5\pi}{12}}
\frac{1}{x\sqrt{\sec^2\theta-1}}
\cdot\frac{\sec\theta\tan\theta}{2x}\,d\theta \\
=&\int_{\frac{\pi}{3}}^{\frac{5\pi}{12}}
\frac{\sec\theta\tan\theta}{2x^2\sqrt{\tan^2\theta}}\,d\theta \\
=&\int_{\frac{\pi}{3}}^{\frac{5\pi}{12}}
\frac{\sec\theta\tan\theta}{2\sec\theta\tan\theta}\,d\theta \\
=&\int_{\frac{\pi}{3}}^{\frac{5\pi}{12}}
\frac{1}{2}\,d\theta
\end{align*}

\[
=\left[\frac{1}{2}\theta\right]_{\frac{\pi}{3}}^{\frac{5\pi}{12}}
\]

\[
=\frac{1}{2}\left(\frac{5\pi}{12}-\frac{\pi}{3}\right)
=\frac{1}{2}\cdot\frac{\pi}{12}
=\frac{\pi}{24}
\]
\end{solution}

    \question 
    \[
    \int_{-1}^{1} \frac{1}{(e^x + 1)(1 + x^2)} \, dx
    \]
    \begin{solution}
        巧妙换元$u=-x$,
        \begin{align*}
        \int_{-1}^{1} \frac{dx}{(e^x+1)(1+x^2)} 
        &= \int_{-1}^{0} \frac{dx}{(e^x+1)(1+x^2)} + \int_{0}^{1} \frac{dx}{(e^x+1)(1+x^2)} \\
        &= -\int_{1}^{0} \frac{du}{(e^{-u}+1)(1+u^2)} + \int_{0}^{1} \frac{dx}{(e^x+1)(1+x^2)} \\
        &= \int_{0}^{1} \frac{e^u du}{(1+e^u)(1+u^2)} + \int_{0}^{1} \frac{dx}{(e^x+1)(1+x^2)} \\
        &= \int_{0}^{1} \frac{du}{1+u^2} = \frac{\pi}{4}
        \end{align*}
    \end{solution}  

    \question
使用代换 $u=\sqrt{\frac{1+x}{1-x}}$ 计算积分
\[
\int_{0}^{\frac{1}{4}} \frac{3}{(4x+5)\sqrt{1-x^2}-3(1-x^2)}\,dx
\]

\begin{solution}
先整理给定代换。
\[
u^2=\frac{1+x}{1-x}
\]
于是
\[
u^2(1-x)=1+x
\]
整理得
\[
u^2-1=x(1+u^2)
\]
从而
\[
x=\frac{u^2-1}{u^2+1}
\]

对 $x$ 求导,
\begin{align*}
\frac{dx}{du}
&=\frac{(u^2+1)(2u)-(u^2-1)(2u)}{(u^2+1)^2} \\
&=\frac{4u}{(u^2+1)^2}
\end{align*}
因此
\[
dx=\frac{4u}{(u^2+1)^2}\,du
\]

再分别化简被积式中的各项。
\begin{align*}
1-x^2
&=1-\left(\frac{u^2-1}{u^2+1}\right)^2 \\
&=\frac{(u^2+1)^2-(u^2-1)^2}{(u^2+1)^2} \\
&=\frac{4u^2}{(u^2+1)^2}
\end{align*}
从而
\[
\sqrt{1-x^2}=\frac{2u}{u^2+1}
\]

同时
\begin{align*}
4x+5
&=4\left(\frac{u^2-1}{u^2+1}\right)+5 \\
&=\frac{9u^2+1}{u^2+1}
\end{align*}

积分上下限变为
\[
x=0 \Rightarrow u=1,\qquad x=\frac{1}{4} \Rightarrow u=\sqrt{\frac{5}{3}}
\]

代入原积分得
\[
\int_{1}^{\sqrt{5/3}} 
\frac{3}{\left(\frac{9u^2+1}{u^2+1}\right)\left(\frac{2u}{u^2+1}\right)
-3\left(\frac{4u^2}{(u^2+1)^2}\right)}
\frac{4u}{(u^2+1)^2}\,du
\]

化简后得到
\begin{align*}
&=\int_{1}^{\sqrt{5/3}} \frac{12u}{2u(9u^2+1)-12u^2}\,du \\
&=\int_{1}^{\sqrt{5/3}} \frac{6u}{9u^3-6u^2+u}\,du \\
&=\int_{1}^{\sqrt{5/3}} \frac{6}{9u^2-6u+1}\,du \\
&=\int_{1}^{\sqrt{5/3}} \frac{6}{(3u-1)^2}\,du
\end{align*}

直接积分,
\[
\int \frac{6}{(3u-1)^2}\,du=-\frac{2}{3u-1}
\]

代入上下限,
\begin{align*}
&=\left[-\frac{2}{3u-1}\right]_{1}^{\sqrt{5/3}} \\
&=-\frac{2}{3\sqrt{5/3}-1}+\frac{2}{2}
\end{align*}

继续化简,
\begin{align*}
&=1-\frac{2}{\sqrt{15}-1} \\
&=1-\frac{2(\sqrt{15}+1)}{14} \\
&=\frac{6-\sqrt{15}}{7}
\end{align*}

因此积分结果为
\[
\frac{6-\sqrt{15}}{7}
\]
\end{solution}

\question
使用代换 $x=\frac{ab}{t}$ 求积分
\[
\int_{0}^{\infty} \frac{\ln x}{(x+a)(x+b)}\,dx
\]
其中 $a>b>0$。

\begin{solution}
设
\[
I=\int_{0}^{\infty} \frac{\ln x}{(x+a)(x+b)}\,dx
\]

作代换
\[
x=\frac{ab}{t},\quad dx=-\frac{ab}{t^2}\,dt
\]
当 $x=0$ 时,$t=\infty$；当 $x=\infty$ 时,$t=0$。

则
\[
I=\int_{\infty}^{0} \frac{\ln\left(\frac{ab}{t}\right)}{\left(\frac{ab}{t}+a\right)\left(\frac{ab}{t}+b\right)}\left(-\frac{ab}{t^2}\right)dt
\]

整理得
\begin{align*}
I
&=\int_{0}^{\infty} \frac{\ln(ab)-\ln t}{\frac{ab(b+t)(a+t)}{t^2}} \frac{ab}{t^2}\,dt \\
&=\int_{0}^{\infty} \frac{\ln(ab)-\ln t}{(t+a)(t+b)}\,dt
\end{align*}

因此
\begin{align*}
I
&=\ln(ab)\int_{0}^{\infty} \frac{1}{(t+a)(t+b)}\,dt
 -\int_{0}^{\infty} \frac{\ln t}{(t+a)(t+b)}\,dt \\
&=\ln(ab)\int_{0}^{\infty} \frac{1}{(t+a)(t+b)}\,dt - I
\end{align*}

移项得
\[
2I=\ln(ab)\int_{0}^{\infty} \frac{1}{(t+a)(t+b)}\,dt
\]

对分式作部分分式分解
\[
\frac{1}{(t+a)(t+b)}=\frac{1}{a-b}\left(\frac{1}{t+b}-\frac{1}{t+a}\right)
\]

于是
\begin{align*}
2I
&=\frac{\ln(ab)}{a-b}\int_{0}^{\infty}\left(\frac{1}{t+b}-\frac{1}{t+a}\right)dt \\
&=\frac{\ln(ab)}{a-b}\left[\ln\left(\frac{t+b}{t+a}\right)\right]_{0}^{\infty}
\end{align*}

计算上下限
\[
\left[\ln\left(\frac{t+b}{t+a}\right)\right]_{0}^{\infty}
=0-\ln\left(\frac{b}{a}\right)
=\ln\left(\frac{a}{b}\right)
\]

因此
\[
2I=\frac{\ln(ab)}{a-b}\ln\left(\frac{a}{b}\right)
\]

从而
\[
I=\frac{\ln(ab)}{2(a-b)}\ln\left(\frac{a}{b}\right)
\]
\end{solution}

\question
使用代换求
\[
\int \frac{2 \sin x \sqrt{\tan x} + 1}{2 \cos x \sqrt{\tan x} (\cos x \sqrt{\tan x} + 1)} \, dx
\]
\begin{solution}
代换:
\[
u = \sec x + \sqrt{\tan x} = \sec x + (\tan x)^{\frac{1}{2}}
\]

求导:
\[
\frac{du}{dx} = \sec x \tan x + \frac{1}{2} (\tan x)^{-\frac{1}{2}} \sec^2 x
= \frac{\sin x}{\cos^3 x} + \frac{\cos^{1/2} x}{2 \sin^{1/2} x \cos^2 x}
= \frac{1}{\cos^2 x} \left( \frac{2\sin^{3/2} x + \cos^{3/2} x}{2 \sin^{1/2} x \cos^{1/2} x} \right)
\]

所以:
\[
dx = \frac{2 \cos^3 x \sin^{1/2} x \cos^{1/2} x}{2 \sin^{3/2} x + \cos^{3/2} x} \, du
\]

原积分:
\[
\int \frac{2 \sin x \sqrt{\tan x} + 1}{2 \cos x \sqrt{\tan x} (\cos x \sqrt{\tan x} + 1)} \, dx
\]

代入 \(dx\):
\[
= \int \frac{2 \sin^{3/2} x + \cos^{1/2} x}{2 \cos^{1/2} x \sin^{1/2} x (\cos^{1/2} x \sin^{1/2} x + 1)} 
\times \frac{2 \cos^{3/2} x \sin^{1/2} x}{2 \sin^{3/2} x + \cos^{3/2} x} \, du
\]

化简:
\[
= \int \frac{\cos x}{\cos x \sin x + 1} \, du
\]

两边同时乘 \(\sec x\):
\[
= \int \frac{\sec x}{\sec x (\cos x \sin x + 1)} \, du
= \int \frac{1}{\sin x + \sec x} \, du
= \int \frac{1}{\sqrt{\tan x} + \sec x} \, du
\]

最终积分:
\[
\int \frac{1}{u} \, du = \ln|u| + C = \ln|\sqrt{\tan x} + \sec x| + C
\]
\end{solution}


\question 使用代换求积分
\[
\int \frac{16}{(x+6)(x-2)\sqrt{x+2}} \, dx
\]
\begin{solution}
使用代换 \(u = \sqrt{x+2} \implies x = u^2 - 2, \ dx = 2u \, du\):
\[
\int \frac{16}{(x+6)(x-2)\sqrt{x+2}} \, dx
= \int \frac{16}{(u^2+4)(u^2-4) u} \cdot 2u \, du
= \int \frac{32}{(u^2+4)(u^2-4)} \, du
\]

分解为部分分式:
\[
\frac{32}{(u^2+4)(u^2-4)} = \frac{4}{u^2-4} - \frac{4}{u^2+4} = \frac{1}{u-2} - \frac{1}{u+2} - \frac{4}{u^2+4}
\]

积分得到:
\[
\int \frac{16}{(x+6)(x-2)\sqrt{x+2}} \, dx = \ln|u-2| - \ln|u+2| - 2\arctan\left(\frac{u}{2}\right) + C
\]

回代 \(u = \sqrt{x+2}\):
\[
\int \frac{16}{(x+6)(x-2)\sqrt{x+2}} \, dx = \ln\left|\frac{\sqrt{x+2}-2}{\sqrt{x+2}+2}\right| - 2\arctan\left(\frac{\sqrt{x+2}}{2}\right) + C
\]
\end{solution}

\question
11) 计算不定积分 $\int \frac{\sqrt{x+1} - \sqrt{x-1}}{\sqrt{x+1} + \sqrt{x-1}} dx$

\begin{solution}
分子有理化:
\begin{align*}
\int \frac{(\sqrt{x+1} - \sqrt{x-1})^2}{(x+1) - (x-1)} dx &= \int \frac{x+1 + x-1 - 2\sqrt{x^2-1}}{2} dx \\
&= \int (x - \sqrt{x^2-1}) dx \\
&= \frac{1}{2}x^2 - \left( \frac{x}{2}\sqrt{x^2-1} - \frac{1}{2} \ln|x + \sqrt{x^2-1}| \right) + C \\
&= \frac{1}{2}x^2 - \frac{x}{2}\sqrt{x^2-1} + \frac{1}{2} \ln|x + \sqrt{x^2-1}| + C
\end{align*}
\end{solution}


\question 使用代换 $x = e^{-u/2}$求
\[
\int \frac{x^4 - 1}{x^2 \sqrt{x^4 + 1}} \, dx
\]
\begin{solution}
使用代换 \(x = e^{-u/2} \implies dx = -\frac{1}{2}e^{-u/2} \, du\):
\[
\int \frac{x^4 - 1}{x^2 \sqrt{x^4 + 1}} \, dx
= \int \frac{e^{-2u} - 1}{e^{-u} \sqrt{e^{-4u} + 1}} \left(-\frac{1}{2}e^{-u/2} \, du\right)
\]

化简得到:
\[
= \frac{1}{\sqrt{2}} \int \sinh u \, (\cosh u)^{-1/2} \, du
= \frac{1}{\sqrt{2}} \cdot 2 (\cosh u)^{1/2} + C
= \sqrt{2 \cosh u} + C
\]

回代 \(x = e^{-u/2} \implies \cosh u = \frac{1}{2}(e^{u} + e^{-u}) = \frac{1}{2}\left(\frac{1}{x^2} + x^2\right)\):
\[
\sqrt{2\cosh u} + C = \sqrt{\frac{1}{x^2} + x^2} + C = \frac{\sqrt{x^4 + 1}}{x} + C
\]

因此积分的简化形式为:
\[
\int \frac{x^4 - 1}{x^2 \sqrt{x^4 + 1}} \, dx = \frac{\sqrt{x^4 + 1}}{x} + C
\]
\end{solution}

\question
已知
\[
x^2 + x + 2 = (u-x)^2
\]
求
\[
\int \frac{1}{x\sqrt{x^2+x+2}} \, dx
\]
\begin{solution}
a) (i) 解 $x$ 关于 $u$ 的关系:

\[
x^2 + x + 2 = u^2 - 2ux + x^2
\]

\[
x + 2 = u^2 - 2ux
\]

\[
x(2u+1) = u^2 - 2
\]

\[
x = \frac{u^2 - 2}{2u+1}
\]

(ii) 对 $x$ 求导:
\[
\frac{dx}{du} = \frac{(2u+1)(2u) - 2(u^2-2)}{(2u+1)^2} = \frac{4u^2 + 2u - 2u^2 + 4}{(2u+1)^2} = \frac{2u^2 + 2u + 4}{(2u+1)^2} = \frac{2(u^2 + u + 2)}{(2u+1)^2}
\]

b) 积分变换:
\[
\int \frac{1}{x\sqrt{x^2+x+2}} \, dx = \int \frac{1}{x(u-x)} \cdot \frac{2(u^2+u+2)}{(2u+1)^2} \, du
\]

代入 $x = \frac{u^2-2}{2u+1}$ 和 $u-x = \sqrt{x^2+x+2}$:

\[
\int \frac{1}{\frac{u^2-2}{2u+1} \cdot \frac{u^2+u+2}{2u+1}} \cdot \frac{2(u^2+u+2)}{(2u+1)^2} \, du
= \int \frac{(2u+1)^2}{(u^2-2)(u^2+u+2)} \cdot \frac{2(u^2+u+2)}{(2u+1)^2} \, du
= \int \frac{2}{u^2-2} \, du
\]

使用标准积分公式:
\[
\int \frac{1}{u^2-a^2} \, du = \frac{1}{2a} \ln \left| \frac{u-a}{u+a} \right| + C
\]

得:
\[
\int \frac{2}{u^2-2} \, du = \frac{1}{\sqrt{2}} \ln \left| \frac{u-\sqrt{2}}{u+\sqrt{2}} \right| + C
\]

最后代回 $u = x + \sqrt{x^2+x+2}$:
\[
\int \frac{1}{x\sqrt{x^2+x+2}} \, dx = \frac{1}{\sqrt{2}} \ln \left| \frac{x + \sqrt{x^2+x+2} - \sqrt{2}}{x + \sqrt{x^2+x+2} + \sqrt{2}} \right| + C
\]
\end{solution}

\question 计算积分
\[
\int_{0.2}^{0.5} \frac{\sqrt{x-x^2}}{x^4} \, dx
\]
\begin{solution}
作代换:
\[
x = \frac{1}{u^2+1} \implies dx = -\frac{2u}{(u^2+1)^2} \, du
\]

积分上下限变化:
\[
x=0.2 \implies u = 2, \quad x=0.5 \implies u = 1
\]

代入积分:
\[
\int_{0.2}^{0.5} \frac{\sqrt{x-x^2}}{x^4} \, dx = \int_{2}^{1} \frac{\sqrt{\frac{1}{u^2+1}-\frac{1}{(u^2+1)^2}}}{\left(\frac{1}{u^2+1}\right)^4} \left(-\frac{2u}{(u^2+1)^2} \, du\right)
\]

整理被积函数:
\[
\sqrt{\frac{1}{u^2+1}-\frac{1}{(u^2+1)^2}} = \sqrt{\frac{u^2}{(u^2+1)^2}} = \frac{u}{u^2+1}
\]

\[
\frac{1}{\left(\frac{1}{u^2+1}\right)^4} = (u^2+1)^4
\]

因此积分变为:
\[
\int_{1}^{2} \frac{u}{u^2+1} (u^2+1)^4 \cdot \frac{2u}{(u^2+1)^2} \, du = \int_{1}^{2} 2 u^2 (u^2+1) \, du
\]

展开:
\[
\int_{1}^{2} 2u^4 + 2u^2 \, du = \left[ \frac{2}{5} u^5 + \frac{2}{3} u^3 \right]_{1}^{2}
\]

代入上下限:
\[
\left( \frac{64}{5} + \frac{16}{3} \right) - \left( \frac{2}{5} + \frac{2}{3} \right) = \frac{256}{15}
\]

\[
\therefore \int_{0.2}^{0.5} \frac{\sqrt{x-x^2}}{x^4} \, dx = \frac{256}{15}
\]
\end{solution}


    \question
    \[
    \int \frac{dx}{x^2(a + bx)^2},\ a,b\in\mathbb{R}
    \]
    \begin{solution}
        设 \( u = \dfrac{a + bx}{x} \),则
        \[
        x = \frac{a}{u - b}, \quad 
        dx = \frac{-a\,du}{(u - b)^2}, \quad 
        a + bx = x u = \frac{a u}{u - b}
        \]
        代入原式:
        \begin{align*}
        \int \frac{dx}{x^2(a + bx)^2}
        &= \int \frac{-a\,du}{(u - b)^2} \cdot \frac{1}{x^2 (a + bx)^2} \\
        &= \int \frac{-a\,du}{(u - b)^2} \cdot \frac{(u - b)^4}{a^4 u^2} \\
        &= -\frac{1}{a^3} \int \frac{(u - b)^2}{u^2} \, du \\
        &= -\frac{1}{a^3} \int \left( 1 - \frac{2b}{u} + \frac{b^2}{u^2}\right) du \\
        &= -\frac{1}{a^3} \left( u - 2b \ln|u| - \frac{b^2}{u} \right) + C
        \end{align*}
        代回 \( u = \frac{a + bx}{x} \),得:
        \[
        \int \frac{dx}{x^2(a + bx)^2}
        = -\frac{1}{a^3} \left( \frac{a + bx}{x} - 2b \ln\left| \frac{a + bx}{x} \right| - \frac{b^2 x}{a + bx} \right) + C
        \]
    \end{solution}

    \question
使用代换$2x+1=4\cosh\theta$,求不定积分
\[
\int \sqrt{(2x+5)(2x-3)} \, dx
\]

\begin{solution}
先化简被积式,
\[
(2x+5)(2x-3)=4x^2+4x-15=(2x+1)^2-16
\]
因此原积分化为
\[
\int \sqrt{(2x+1)^2-16} \, dx
\]

采用双曲函数换元,令
\[
2x+1=4\cosh\theta
\]
则
\[
2\,dx=4\sinh\theta\,d\theta,\quad dx=2\sinh\theta\,d\theta
\]

利用恒等式
\[
\cosh^2\theta-\sinh^2\theta=1
\]
可得
\[
\sqrt{(2x+1)^2-16}
=\sqrt{16(\cosh^2\theta-1)}
=4\sinh\theta
\]

代入积分,
\[
\int \sqrt{(2x+1)^2-16}\,dx
=\int 4\sinh\theta\,(2\sinh\theta\,d\theta)
=8\int\sinh^2\theta\,d\theta
\]

又有
\[
\sinh^2\theta=\frac{1}{2}\cosh(2\theta)-\frac{1}{2}
\]
于是
\begin{align*}
8\int\sinh^2\theta\,d\theta
&=8\int\left(\frac{1}{2}\cosh(2\theta)-\frac{1}{2}\right)d\theta \\
&=4\int\cosh(2\theta)\,d\theta-4\int d\theta \\
&=2\sinh(2\theta)-4\theta+C
\end{align*}

利用
\[
\sinh(2\theta)=2\sinh\theta\cosh\theta
\]
得
\[
2\sinh(2\theta)=4\sinh\theta\cosh\theta
\]

又因为
\[
\sinh\theta=\frac{\sqrt{(2x+1)^2-16}}{4},\quad
\cosh\theta=\frac{2x+1}{4}
\]
且
\[
\theta=\operatorname{arccosh}\frac{2x+1}{4}
\]

代回原变量,
\[
\int \sqrt{(2x+5)(2x-3)} \, dx
=\frac{1}{4}(2x+1)\sqrt{(2x+5)(2x-3)}
-4\operatorname{arccosh}\frac{2x+1}{4}+C
\]

利用
\[
\operatorname{arccosh}u=\ln\left(u+\sqrt{u^2-1}\right)
\]
可进一步写为
\[
\int \sqrt{(2x+5)(2x-3)} \, dx
=\frac{1}{4}(2x+1)\sqrt{(2x+5)(2x-3)}
-4\ln\left(2x+1+\sqrt{(2x+5)(2x-3)}\right)+C
\]

\end{solution}

\question
已知
\[
\sqrt{5-4x-x^2} = (1-x)u, \quad x\neq 1, x\neq -5
\]

\begin{solution}
a) 代换与求导

i) 两边平方
\[
5-4x-x^2 = (1-x)^2 u^2
\]

移项整理
\[
-(x^2+4x-5) = u^2(x-1)^2 \implies -(x-1)(x+5) = u^2(x-1)^2
\]

两边除以 $x-1$
\[
-(x+5) = u^2 (x-1) \implies u^2 - 5 = x(u^2+1) \implies x = \frac{u^2-5}{u^2+1}
\]

ii) 对 x 求导
\[
dx = \frac{d}{du} \frac{u^2-5}{u^2+1} \, du
= \frac{(u^2+1)(2u) - (u^2-5)(2u)}{(u^2+1)^2} \, du
= \frac{12u}{(u^2+1)^2} \, du
\]

b) 积分化简

由 a) 可得
\[
\int \frac{x}{(5-4x-x^2)^{3/2}} \, dx
= \int \frac{x}{((1-x)u)^3} \cdot \frac{12u}{(u^2+1)^2} \, du
\]

\[
x = \frac{u^2-5}{u^2+1}, \quad (1-x)^3 = \left(\frac{6}{u^2+1}\right)^3
\]

\[
\implies \int \frac{x}{(5-4x-x^2)^{3/2}} \, dx
= \int \frac{u^2-5}{u^2+1} \cdot \frac{12u}{216 u^3} \, du
= \int \frac{u^2-5}{18 u^2} \, du
\]

c) 积分求解

分拆积分
\[
\int \frac{u^2-5}{18 u^2} \, du = \int \frac{1}{18} - \frac{5}{18u^2} \, du
= \frac{1}{18} u + \frac{5}{18u} + C
\]

代回 u
\[
u = \frac{\sqrt{5-4x-x^2}}{1-x} \implies \frac{1}{u} = \frac{1-x}{\sqrt{5-4x-x^2}}
\]

\[
\implies \int \frac{x}{(5-4x-x^2)^{3/2}} \, dx
= \frac{1}{18} \frac{\sqrt{5-4x-x^2}}{1-x} + \frac{5}{18} \frac{1-x}{\sqrt{5-4x-x^2}} + C
\]

整理为
\[
\int \frac{x}{(5-4x-x^2)^{3/2}} \, dx
= \frac{5-2x}{9 \sqrt{5-4x-x^2}} + C
\]
\end{solution}

\question
28) 计算不定积分:\[ \int \frac{\sqrt{1 + \sqrt[3]{x}}}{\sqrt[3]{x^2}} dx \]

\begin{solution}
设 $u = \sqrt[3]{x}$,则 $du = \frac{1}{3} x^{-\frac{2}{3}} dx$,即 $dx = 3 \sqrt[3]{x^2} du$。
代入积分式:
\[ I = 3 \int \sqrt{1+u} du \]
进行积分计算:
\[ I = 2 (1+u)^{\frac{3}{2}} + C \]
还原变量 $u = \sqrt[3]{x}$:
\[ I = 2 (1 + \sqrt[3]{x})^{\frac{3}{2}} + C \]
\end{solution}
\question
16) 计算不定积分:\[ \int \frac{1}{x(1+x^{119})} dx \]

\begin{solution}
分子分母同乘 $x^{118}$:
\[ \int \frac{x^{118}}{x^{119}(1+x^{119})} dx \]
设 $t = x^{119}$,则 $dt = 119x^{118} dx$:
\[ I = \frac{1}{119} \int \frac{1}{t(1+t)} dt = \frac{1}{119} \int \left( \frac{1}{t} - \frac{1}{1+t} \right) dt \]
\[ = \frac{1}{119} \ln\left| \frac{t}{1+t} \right| + C = \frac{1}{119} \ln\left| \frac{x^{119}}{1+x^{119}} \right| + C \]
\end{solution}
    \question  
    \[
    \int_{-1}^1 \frac{x^2}{1+2^x} \,dx 
    \]
    \begin{solution}
        设 \( u = -x \),则
        \[
        I = \int_{-1}^1 \frac{u^2}{1+2^{-u}}\,du = \int_{-1}^1 \frac{u^2 2^u}{1+2^u}\,du = \int_{-1}^1 \frac{x^2 2^x}{1+2^x}\,dx
        \]
        于是有
        \[
        2I = \int_{-1}^1 \frac{x^2(2^x + 1)}{1 + 2^x}\,dx = \int_{-1}^1 x^2\,dx = \frac{2}{3} \Rightarrow I = \frac{1}{3}
        \]
    \end{solution}

    \question 计算积分
\[
I = \int_{\frac{1}{2}}^{2} \frac{\ln x}{1 + x^2} \, dx.
\]

\begin{solution}
将积分分成两部分:
\[
I = \int_{\frac{1}{2}}^{1} \frac{\ln x}{1 + x^2} \, dx + \int_{1}^{2} \frac{\ln x}{1 + x^2} \, dx.
\]

在第二个积分中作代换 $x = 1/u$,则 $dx = -\frac{1}{u^2} du$,积分限变为 $x=1 \to u=1,x=2 \to u=1/2$:
\[
\int_{1}^{2} \frac{\ln x}{1 + x^2} \, dx = \int_{1}^{1/2} \frac{\ln (1/u)}{1 + (1/u)^2} \left(-\frac{1}{u^2} du\right) = -\int_{1/2}^{1} \frac{-\ln u}{1 + u^2} \, du = -\int_{1/2}^{1} \frac{\ln u}{1 + u^2} \, du.
\]

因此
\[
I = \int_{1/2}^{1} \frac{\ln x}{1 + x^2} \, dx - \int_{1/2}^{1} \frac{\ln u}{1 + u^2} \, du = 0.
\]
\end{solution}

\question 
30) 计算定积分 $I = \int_{0}^{\sqrt{3}} \frac{1}{1+x^2} \sin^{-1} \frac{2x}{1+x^2} dx$

\begin{solution}
使用换元法,设 $x = \tan\theta$,则 $dx = \sec^2 \theta d\theta$ 且 $\frac{1}{1+x^2} = \cos^2 \theta$。
同时利用二倍角公式可知 $\sin^{-1} \frac{2x}{1+x^2} = \sin^{-1}(\sin 2\theta) = 2\theta$。

更换积分上下限:
\begin{itemize}
    \item 当 $x = 0$ 时,$\theta = 0$
    \item 当 $x = \sqrt{3}$ 时,$\theta = \frac{\pi}{3}$
\end{itemize}

代入原积分式:
\begin{align*}
I &= \int_{0}^{\frac{\pi}{3}} \cos^2 \theta \cdot (2\theta) \cdot \sec^2 \theta d\theta \\
&= \int_{0}^{\frac{\pi}{3}} 2\theta d\theta \\
&= \left[ \theta^2 \right]_{0}^{\frac{\pi}{3}} \\
&= \left( \frac{\pi}{3} \right)^2 - 0 \\
&= \frac{\pi^2}{9}
\end{align*}
\end{solution}

\question 计算 $[\int_{1/2}^2 \frac{x^{2012}-1}{x^{2014}+1} dx]$。
\begin{solution}
令 $x = \frac{1}{t}$,则 $dx = -\frac{1}{t^2} dt$。
积分上下限从 $1/2 \to 2$ 变为 $2 \to 1/2$:
\[I = [\int_2^{1/2} \frac{(1/t)^{2012}-1}{(1/t)^{2014}+1} (-\frac{1}{t^2}) dt]\]
\[I = [\int_{1/2}^2 \frac{\frac{1-t^{2012}}{t^{2012}}}{\frac{1+t^{2014}}{t^{2014}}} \frac{1}{t^2} dt] = [\int_{1/2}^2 \frac{1-t^{2012}}{1+t^{2014}} dt]\]
注意到 $[\int_{1/2}^2 \frac{1-x^{2012}}{1+x^{2014}} dx] = -[\int_{1/2}^2 \frac{x^{2012}-1}{x^{2014}+1} dx] = -I$。
即 $I = -I$,故 $I = 0$。
\end{solution}

    \question 若正整数$k$满足
    \[
    \left(\int_0^1 x^{2018}(2019 + kx^{10})\sqrt{1 + x^{10}} \, dx\right)^2 = 8,
    \]
    求$k$的最大值。
    \begin{solution}
        由于被积函数随$k$严格递增,则当
        \[
        \int_0^1 x^{2018}(2019 + kx^{10})\sqrt{1 + x^{10}} \, dx = \sqrt{8}
        \]
        时$k$取最大。欲写成形如
        \[
        \int_0^1 (2019x^{2018-n} + kx^{2028-n}) \sqrt{x^{2n} + x^{2n+10}} \, dx
        \]
        的积分,考虑强迫
        \[
        2n-1 = 2018-n,
        \]
        可得$n = 673$,此时设$u=x^{1346}+x^{1356}$,则$du=1346x^{1345}+1356x^{1355} dx $,于是
        \begin{align*}
        I &=\int_0^1 (2019x^{1345} + kx^{1355}) \sqrt{x^{1346} + x^{1356}} \, dx \\
        &=\frac{3}{2}\int_0^1 (1346^{1345} + \frac{2}{3}kx^{1355}) \sqrt{x^{1346} + x^{1356}} \, dx
        \end{align*}
        当$\dfrac{2}{3}k=1356$即$k=2034$时,有
        \[
        I=\frac{3}{2}\int_0^2 \sqrt{u}\, du = \frac{3}{2}\left[\frac{2}{3}u^{\frac{3}{2}}\right]^2_0=\sqrt{8}
        \]
    \end{solution}
%部分分式
    \question 
    \[
    \int \frac{dx}{x^3 + 1}
    \]
    \begin{solution}
        \[
        \int \frac{dx}{x^3 + 1}
        = \int \frac{dx}{(x + 1)(x^2 - x + 1)}
        \]
        部分分式分解得
        \[
        \frac{1}{x^3 + 1}
        = \frac{1}{3(x + 1)} + \frac{-\frac{1}{3}x + \frac{2}{3}}{x^2 - x + 1}
        \]
        于是
        \[
        \int \frac{dx}{x^3 + 1}
        = \frac{1}{3} \int \frac{dx}{x + 1}
        - \frac{1}{3} \int \frac{x}{x^2 - x + 1} dx
        + \frac{2}{3} \int \frac{dx}{x^2 - x + 1}
        \]
        其中
        \[
        \int \frac{dx}{x + 1} = \ln|x + 1|,\quad
        \int \frac{x}{x^2 - x + 1} dx = \frac{1}{2} \ln(x^2 - x + 1)
        \]
        \[
        \int \frac{dx}{x^2 - x + 1}
        = \int \frac{dx}{\left(x - \frac{1}{2}\right)^2 + \left(\frac{\sqrt{3}}{2}\right)^2}
        = \frac{2}{\sqrt{3}} \arctan\left( \frac{2x - 1}{\sqrt{3}} \right) + C
        \]
        故
        \[
        \int \frac{dx}{x^3 + 1}
        = \frac{1}{3} \ln|x + 1|
        - \frac{1}{6} \ln(x^2 - x + 1)
        + \frac{1}{\sqrt{3}} \arctan\left( \frac{2x - 1}{\sqrt{3}} \right) + C
        \]
    \end{solution}
\question
17) 计算不定积分:\[ \int \frac{x^4+1}{x^6+1} dx \]

\begin{solution}
由于 $x^6+1 = (x^2+1)(x^4-x^2+1)$,拆分分子:
\[ \int \frac{(x^4-x^2+1) + x^2}{x^6+1} dx = \int \frac{1}{x^2+1} dx + \int \frac{x^2}{x^6+1} dx \]
第一部分:$\int \frac{1}{x^2+1} dx = \tan^{-1} x$。
第二部分:设 $u = x^3, du = 3x^2 dx$:
\[ \frac{1}{3} \int \frac{1}{u^2+1} du = \frac{1}{3} \tan^{-1} u = \frac{1}{3} \tan^{-1} x^3 \]
合并结果:
\[ I = \tan^{-1} x + \frac{1}{3} \tan^{-1} x^3 + C \]
\end{solution}
    \question
计算积分
\[
\int_{0}^{1} \frac{x^4(1-x)^4}{x^2+1} \, dx
\]

\begin{solution}
首先展开分子 \( (1-x)^4 \):
\[
(1-x)^4 = 1 - 4x + 6x^2 - 4x^3 + x^4
\]

因此积分变为:
\[
\int_{0}^{1} \frac{x^4 (1-4x+6x^2-4x^3+x^4)}{x^2+1} \, dx
= \int_{0}^{1} \frac{x^8 - 4x^7 + 6x^6 - 4x^5 + x^4}{x^2+1} \, dx
\]

进行多项式除法,将分子按 \(x^2+1\) 除:
\[
\frac{x^8 - 4x^7 + 6x^6 - 4x^5 + x^4}{x^2+1} = x^6 - 4x^5 + 5x^4 - 4x^2 + 4 - \frac{4}{x^2+1}
\]

于是积分拆分为:
\[
\int_{0}^{1} x^6 - 4x^5 + 5x^4 - 4x^2 + 4 - \frac{4}{x^2+1} \, dx
\]

分别积分各项:
\[
\int x^6 \, dx = \frac{1}{7} x^7, \quad
\int -4x^5 \, dx = -\frac{2}{3} x^6, \quad
\int 5x^4 \, dx = x^5,
\]
\[
\int -4x^2 \, dx = -\frac{4}{3} x^3, \quad
\int 4 \, dx = 4x, \quad
\int -\frac{4}{x^2+1} \, dx = -4 \arctan x
\]

组合结果:
\[
\int_{0}^{1} \frac{x^4(1-x)^4}{x^2+1} \, dx
= \left[ \frac{1}{7}x^7 - \frac{2}{3}x^6 + x^5 - \frac{4}{3}x^3 + 4x - 4\arctan x \right]_{0}^{1}
\]

代入上下限:
\[
\left( \frac{1}{7} - \frac{2}{3} + 1 - \frac{4}{3} + 4 - 4 \cdot \frac{\pi}{4} \right) - 0
= \frac{1}{7} - \frac{6}{3} + 5 - \pi
= \frac{1}{7} - 2 + 5 - \pi
= 3 + \frac{1}{7} - \pi
= \frac{22}{7} - \pi
\]

因此最终结果为:
\[
\int_{0}^{1} \frac{x^4(1-x)^4}{x^2+1} \, dx = \frac{22}{7} - \pi
\]
\end{solution}

    \question 
    \[
    \int_{0}^{1} \frac{(x^2+1)(x^2+4)}{(x^2+3)(x^2-4)} \, dx
    \]
    \begin{solution}
        \[
\int_{0}^{1} \frac{(x^2+1)(x^2+4)}{(x^2+3)(x^2-4)} \, dx = \int_{0}^{1} \left[1 + \frac{2}{7}\left[\frac{1}{x^2+3} + \frac{5}{x-2} - \frac{5}{x+2}\right]\right] \, dx
\]
\[
= \left[x + \frac{2}{7}\left[\frac{1}{\sqrt{3}}\arctan\left(\frac{x}{\sqrt{3}}\right) + 5\ln|x-2| - 5\ln|x+2|\right]\right]_{0}^{1}
\]
\[
= \left[1 + \frac{2}{7}\left[\frac{1}{\sqrt{3}}\times\frac{\pi}{6} + 5\ln|-1| - 5\ln|3|\right]\right] - \left[\frac{2}{7}[5\ln|-2| - 5\ln|2|]\right]
\]
\[
= 1 + \frac{2}{7}\left[\frac{\pi}{6\sqrt{3}} - 5\ln3\right]
\]
\[
= 1 + \frac{1}{7}\left[\frac{\pi}{3\sqrt{3}} - 10\ln3\right]
\]
    \end{solution}

    \question
计算积分
\[
\int_{0}^{\frac{1}{3}} \frac{32x^2}{(x^2-1)(x+1)^3}\,dx
\]

\begin{solution}
先化简被积函数
\[
\int_{0}^{\frac{1}{3}} \frac{32x^2}{(x^2-1)(x+1)^3}\,dx
= \int_{0}^{\frac{1}{3}} \frac{32x^2}{(x-1)(x+1)^4}\,dx
\]

作代换
\[
u=x+1,\quad x=u-1,\quad du=dx
\]
当 $x=0$ 时,$u=1$  
当 $x=\frac{1}{3}$ 时,$u=\frac{4}{3}$  

则积分化为
\[
\int_{1}^{\frac{4}{3}} \frac{32(u-1)^2}{(u-2)u^4}\,du
= \int_{1}^{\frac{4}{3}} \frac{32}{u^4}\cdot\frac{u^2-2u+1}{u-2}\,du
\]

将有理函数分解后得
\[
\int_{1}^{\frac{4}{3}}
\left(
-\frac{16}{u^4}
+\frac{24}{u^3}
-\frac{4}{u^2}
-\frac{2}{u}
+\frac{2}{u-2}
\right)\,du
\]

逐项积分
\[
=\left[
\frac{16}{3u^3}
-\frac{12}{u^2}
+\frac{4}{u}
-2\ln|u|
+2\ln|u-2|
\right]_{1}^{\frac{4}{3}}
\]

代入上、下限
\[
=\left(
\frac{16}{3\left(\frac{4}{3}\right)^3}
-\frac{12}{\left(\frac{4}{3}\right)^2}
+\frac{4}{\frac{4}{3}}
-2\ln\frac{4}{3}
+2\ln\frac{2}{3}
\right)
-\left(
\frac{16}{3}-12+4
\right)
\]

化简得
\[
=11-\frac{16}{3}-\frac{9}{2}-2\ln2
\]

整理为
\[
\frac{7}{6}-2\ln2
\]

因此
\[
\int_{0}^{\frac{1}{3}} \frac{32x^2}{(x^2-1)(x+1)^3}\,dx
= \frac{7}{6}-2\ln2
\]
\end{solution}

\question
24) 计算不定积分:\[ \int \frac{1}{\sqrt[4]{1+x^4}} dx \]

\begin{solution}
使用换元法,设 $t^4 x^4 = 1+x^4$,则 $x^4 = \frac{1}{t^4-1}$。
两边求导得 $4x^3 dx = -\frac{4t^3}{(t^4-1)^2} dt$,故 $dx = -\frac{t^3}{(t^4-1)^2 x^3} dt$。
代入积分式:
\[ I = -\int \frac{1}{t \cdot x} \cdot \frac{t^3}{(t^4-1)^2 x^2} dt = -\int \frac{t^2}{t^4-1} dt \]
将被积函数进行部分分式拆分:
\[ -\int \frac{t^2}{t^4-1} dt = -\frac{1}{2} \int \left( \frac{1}{t^2-1} + \frac{1}{t^2+1} \right) dt \]
进行积分计算:
\[ I = -\frac{1}{4} \ln\left| \frac{t-1}{t+1} \right| - \frac{1}{2} \tan^{-1} t + C \]
其中还原变量为 $t = \frac{\sqrt[4]{1+x^4}}{x}$。
\end{solution}

\question
20) 计算不定积分:\[ \int \frac{2}{\sqrt{x} - 3\sqrt[4]{x} + 2} dx \]

\begin{solution}
使用换元法,设 $\sqrt[4]{x} = t$,则 $x = t^4, dx = 4t^3 dt$。
代入原积分式:
\[ I = \int \frac{8t^3}{t^2 - 3t + 2} dt \]
对被积函数进行多项式除法与部分分式拆分:
\[ \frac{8t^3}{t^2 - 3t + 2} = 8t + 24 + \frac{56t - 48}{t^2 - 3t + 2} = 8t + 24 + \frac{64}{t-2} - \frac{8}{t-1} \]
进行积分计算:
\begin{align*}
I &= \int (8t + 24) dt + \int \frac{64}{t-2} dt - \int \frac{8}{t-1} dt \\
&= 4t^2 + 24t + 64 \ln|t-2| - 8 \ln|t-1| + C
\end{align*}
还原变量 $t = \sqrt[4]{x}$:
\[ I = 4\sqrt{x} + 24\sqrt[4]{x} + 64 \ln|\sqrt[4]{x}-2| - 8 \ln|\sqrt[4]{x}-1| + C \]
\end{solution}

\question
22) 计算不定积分:\[ \int \frac{3x^3}{x^4 + x^3 + x + 1} dx \]

\begin{solution}
首先对分母进行因式分解:
\[ x^4 + x^3 + x + 1 = x^3(x+1) + (x+1) = (x^3+1)(x+1) = (x^2-x+1)(x+1)^2 \]
将被积函数拆分为部分分式:
\[ \frac{3x^3}{(x^2-x+1)(x+1)^2} = \frac{x-1}{x^2-x+1} + \frac{2}{x+1} - \frac{1}{(x+1)^2} \]
逐步进行积分:
\begin{align*}
I &= \int \frac{x-1}{x^2-x+1} dx + \int \frac{2}{x+1} dx - \int \frac{1}{(x+1)^2} dx \\
&= \frac{1}{2} \int \frac{2x-1-1}{x^2-x+1} dx + 2 \ln|x+1| + \frac{1}{x+1} \\
&= \frac{1}{2} \ln(x^2-x+1) - \frac{1}{2} \int \frac{1}{(x-\frac{1}{2})^2 + \frac{3}{4}} dx + 2 \ln|x+1| + \frac{1}{x+1}
\end{align*}
计算得到最终结果:
\[ I = \ln\left|(x+1)^2 \sqrt{x^2-x+1}\right| - \frac{1}{\sqrt{3}} \tan^{-1}\left(\frac{2x-1}{\sqrt{3}}\right) + \frac{1}{x+1} + C \]
\end{solution}
    \question
\[
\int \frac{x^2(x^4+1)}{\sqrt{x^4+2}} \, dx
\]

\begin{solution}
先作代数变形
\[
\int \frac{x^2(x^4+1)}{\sqrt{x^4+2}} \, dx
= \int x^2(x^4+1)(x^4+2)^{-\frac12} \, dx
\]

将分子整理
\[
= \int (x^7+x^3)(x^4)^{-\frac12}(x^4+2)^{-\frac12} \, dx
\]

合并为一个整体
\[
= \int (x^7+x^3)[x^4(x^4+2)]^{-\frac12} \, dx
\]

即
\[
= \int (x^7+x^3)(x^8+2x^4)^{-\frac12} \, dx
\]

注意到
\[
\frac{d}{dx}(x^8+2x^4)=8x^7+8x^3=8(x^7+x^3)
\]

于是
\[
= \frac18 \int 8(x^7+x^3)(x^8+2x^4)^{-\frac12} \, dx
\]

积分得
\[
= \frac18 \cdot \frac{1}{\frac12}(x^8+2x^4)^{\frac12}+C
\]

化简
\[
= \frac14 (x^8+2x^4)^{\frac12}+C
\]
\end{solution}
\question
14) 计算不定积分:\[ \int \frac{x^2+1}{(x^2-2x+2)^3} dx \]

\begin{solution}
首先对分子进行拆分,凑出分母的导数项与基本项:
\[ \int \frac{(x^2-2x+2)+(2x-2)+1}{(x^2-2x+2)^3} dx \]
将其拆分为三个部分进行积分:
\[ = \int \frac{dx}{(x^2-2x+2)^2} + \int \frac{(2x-2)dx}{(x^2-2x+2)^3} + \int \frac{dx}{(x^2-2x+2)^3} \]

对第一项和第三项进行配方:
\[ x^2-2x+2 = (x-1)^2+1 \]
设 $x-1 = \tan t$,则 $dx = \sec^2 t dt$。
代入后积分式变为:
\[ \int \cos^2 t dt + \int \cos^4 t dt - \frac{1}{2(x^2-2x+2)^2} \]

利用降幂公式计算三角积分:
\[ \int \cos^2 t dt = \frac{1}{2} \int (1+\cos 2t) dt = \frac{t}{2} + \frac{\sin 2t}{4} \]
\[ \int \cos^4 t dt = \frac{1}{4} \int (1+\cos 2t)^2 dt = \dots = \frac{3t}{8} + \frac{\sin 2t}{4} + \frac{\sin 4t}{32} \]

将各部分结果合并:
\[ I = \frac{7t}{8} + \frac{\sin 2t}{2} + \frac{\sin 4t}{32} - \frac{1}{2(x^2-2x+2)^2} + C \]

利用三角恒等式将 $t$ 还原为 $x$:
\[ \sin 2t = \frac{2(x-1)}{x^2-2x+2}, \quad \sin 4t = \frac{4(x-1)(1-2\sin^2 t)}{(x^2-2x+2)^2} \]
最终整理得到原文结果:
\[ I = \frac{7}{8} \tan^{-1}(x-1) + \frac{x-1}{x^2-2x+2} + \frac{(x-1)(1-2\sin^2 t)}{8(x^2-2x+2)^2} - \dots + C \]
\[ I = \frac{7}{8} \tan^{-1}(x-1) + \frac{7x^3-21x^2+30x-20}{8(x^2-2x+2)^2} + C \]
\end{solution}
\question
40) 计算不定积分:\[ \int \sqrt[3]{\frac{2x+1}{x+3}} dx \]

\begin{solution}
使用换元法,设 $t^3 = \frac{2x+1}{x+3}$,则可求得 $x = \frac{3t^3-1}{2-t^3}$。
对 $x$ 求导得 $dx = \frac{15t^2}{(2-t^3)^2} dt$。
代入原积分式:
\[ I = 15 \int \frac{t^3}{(2-t^3)^2} dt \]
根据原文推导过程,利用分部积分或拆分项处理:
\[ = 5 \int \left( -\frac{1}{2-t^3} + \frac{2(t^3+1)}{(2-t^3)^2} \right) dt = -5 \int \frac{1}{2-t^3} dt + \frac{5t}{2-t^3} \]
进一步整理为:
\[ = 5 \int \frac{1}{t^3-2} dt - \frac{5t}{t^3-2} \]
利用部分分式进行拆分:
\[ = 5 \int \left( \frac{-\frac{\sqrt[3]{2}}{6}t - \frac{\sqrt[3]{4}}{3}}{t^2 + \sqrt[3]{2}t + \sqrt[3]{4}} + \frac{\frac{\sqrt[3]{2}}{6}}{t-\sqrt[3]{2}} \right) dt - \frac{5t}{t^3-2} \]
进行积分计算并还原变量 $t = \sqrt[3]{\frac{2x+1}{x+3}}$:
\[ I = \frac{5\sqrt[3]{108}}{6} \tan^{-1} \frac{\sqrt{3}(\sqrt[3]{4}t+1)}{3} - \frac{5\sqrt[3]{2}}{12} \ln(t^2+\sqrt[3]{2}t+\sqrt[3]{4}) + \frac{5\sqrt[3]{2}}{6} \ln|t-\sqrt[3]{2}| - \frac{5t}{t^3-2} + C \]
\end{solution}

\question
42) 计算不定积分:\[ \int \frac{dx}{\sqrt{x+1} + \sqrt[4]{x+1}} \]

\begin{solution}
设 $t = x+1$,则 $dt = dx$。
积分变为:
\[ \int \frac{dt}{\sqrt{t} + \sqrt[4]{t}} = \int \frac{dt}{\sqrt{t}(\sqrt[4]{t}+1)} \]
再设 $u = \sqrt[4]{t}$,则 $t = u^4, dt = 4u^3 du$。
代入积分式:
\[ \int \frac{4u^3}{u^2+u} du = \int \frac{4u^2}{u+1} du \]
对被积函数进行多项式除法:
\[ \int \left( 4u - 4 + \frac{4}{u+1} \right) du = 2u^2 - 4u + 4 \ln|u+1| + C \]
还原变量 $u = \sqrt[4]{x+1}$,得到最终结果:
\[ I = 2\sqrt{x+1} - 4\sqrt[4]{x+1} + 4 \ln(\sqrt[4]{x+1}+1) + C \]
\end{solution}

\question
43) 计算不定积分:\[ \int \frac{1+x^2}{x\sqrt{x^4-x^2+1}} dx \]

\begin{solution}
分子分母同时除以 $x^2$:
\[ I = \int \frac{1 + \frac{1}{x^2}}{\sqrt{x^2 - 1 + \frac{1}{x^2}}} dx \]
设 $t = x - \frac{1}{x}$,则 $dt = (1 + \frac{1}{x^2}) dx$,且 $x^2 - 1 + \frac{1}{x^2} = t^2 + 1$。
代入积分式:
\[ I = \int \frac{dt}{\sqrt{t^2+1}} = \ln|t + \sqrt{t^2+1}| + C \]
还原变量:
\[ I = \ln\left| x - \frac{1}{x} + \sqrt{x^2 - 1 + \frac{1}{x^2}} \right| + C \]
\end{solution}

\question
44) 计算不定积分:\[ \int \frac{1}{(2x^2+3)\sqrt{5x^2+4}} dx \]

\begin{solution}
将被积函数改写,提取 $x^3$:
\[ I = \int \frac{dx}{x^3(2 + \frac{3}{x^2})\sqrt{5 + \frac{4}{x^2}}} \]
设 $t^2 = 5 + \frac{4}{x^2}$,则 $2t dt = -\frac{8}{x^3} dx$,即 $\frac{1}{x^3} dx = -\frac{1}{4} t dt$。
同时 $\frac{1}{x^2} = \frac{t^2-5}{4}$,代入分母项 $2 + \frac{3}{x^2} = \frac{3t^2-7}{4}$。
代入积分式:
\[ I = -\int \frac{\frac{1}{4} t dt}{\frac{3t^2-7}{4} \cdot t} = -\int \frac{dt}{3t^2-7} = -\frac{1}{3} \int \frac{dt}{t^2 - \frac{7}{3}} \]
进行积分计算:
\[ I = -\frac{\sqrt{21}}{42} \ln\left| \frac{\sqrt{3}t - \sqrt{7}}{\sqrt{3}t + \sqrt{7}} \right| + C \]
还原变量 $t = \sqrt{5 + \frac{4}{x^2}}$。
\end{solution}

\question
45) 计算不定积分:\[ \int \frac{x+x^{\frac{2}{3}}+x^{\frac{1}{6}}}{x(1+x^{\frac{1}{3}})} dx \]

\begin{solution}
使用换元法,设 $x = u^6$,则 $dx = 6u^5 du$。
代入后简化:
\[ I = 6 \int \frac{u^6+u^4+u}{u^6(1+u^2)} \cdot u^5 du = 6 \int \frac{u^6+u^4+u}{1+u^2} du \]
利用多项式除法拆分:
\[ I = 6 \int \left( u^4 + \frac{u}{1+u^2} \right) du = \frac{6}{5} u^5 + 3 \ln(1+u^2) + C \]
还原变量 $u = x^{\frac{1}{6}}$:
\[ I = \frac{6}{5} x^{\frac{5}{6}} + 3 \ln(1+x^{\frac{1}{3}}) + C \]
\end{solution}

\question
47) 计算不定积分:\[ \int \frac{\sqrt[3]{1+\sqrt[4]{x}}}{\sqrt{x}} dx \]

\begin{solution}
设 $1 + x^{\frac{1}{4}} = t^3$,则 $x = (t^3-1)^4,dx = 12t^2(t^3-1)^3 dt$。
代入积分式:
\[ I = \int \frac{t}{(t^3-1)^2} \cdot 12t^2(t^3-1)^3 dt = 12 \int (t^6 - t^3) dt \]
进行积分计算:
\[ I = 12 \left( \frac{t^7}{7} - \frac{t^4}{4} \right) + C = \frac{12}{7} t^7 - 3t^4 + C \]
还原变量 $t = (1+x^{\frac{1}{4}})^{\frac{1}{3}}$。
\end{solution}

\question
48) 计算不定积分:\[ \int \frac{x-1}{(x+1)\sqrt{x^3+x^2+x}} dx \]

\begin{solution}
利用恒等变形:
\[ I = \int \frac{x^2-1}{(x+1)^2 \sqrt{x^3+x^2+x}} dx = \int \frac{1 - \frac{1}{x^2}}{(x + \frac{1}{x} + 2)\sqrt{x + 1 + \frac{1}{x}}} dx \]
设 $u = \sqrt{x + 1 + \frac{1}{x}}$,则 $2u du = (1 - \frac{1}{x^2}) dx$。
代入积分式:
\[ I = \int \frac{2u du}{(u^2+1) \cdot u} = \int \frac{2}{u^2+1} du = 2 \tan^{-1} u + C \]
还原变量:
\[ I = 2 \tan^{-1} \sqrt{x + 1 + \frac{1}{x}} + C \]
\end{solution}

\question
49) 计算不定积分:\[ \int \frac{x}{(x-1)\sqrt{x^2-2x}} dx \]

\begin{solution}
设 $t = x-1$,则 $x = t+1,x^2-2x = t^2-1$。
代入积分式:
\[ I = \int \frac{t+1}{t\sqrt{t^2-1}} dt = \int \frac{1}{\sqrt{t^2-1}} dt + \int \frac{1}{t\sqrt{t^2-1}} dt \]
进行积分计算:
\[ I = \ln|t + \sqrt{t^2-1}| + \sec^{-1} t + C \]
还原变量 $t = x-1$:
\[ I = \ln|x - 1 + \sqrt{x^2-2x}| + \sec^{-1}(x-1) + C \]
\end{solution}

\question
51) 计算不定积分:\[ \int \frac{1}{\sqrt{1-x} + \sqrt{1+x^2}} dx \]

\begin{solution}
对分母进行有理化:
\[ I = \int \frac{\sqrt{1+x^2} - \sqrt{1-x}}{(1+x^2) - (1-x)} dx = \int \frac{\sqrt{1+x^2} - \sqrt{1-x}}{x^2+x} dx \]
拆分为两个积分:
\[ I = \int \frac{\sqrt{1+x^2}}{x(x+1)} dx - \int \frac{\sqrt{1-x}}{x(x+1)} dx \]
对于第二部分,设 $u^2 = 1-x$。
根据笔记底部的最终整理结果:
\[ I = -\sqrt{2} \left( \ln\left|\sqrt{2x^2+2} + x - 1\right| - \ln|x+1| \right) + C \]
\end{solution}
\question
计算 $[\int \frac{3\sin 2x + 4\cos x e^{\sin x} + 8\cos x}{3\sin x + e^{\sin x} + 1} dx]$。

\begin{solution}
首先,利用二倍角公式 $\sin 2x = 2\sin x \cos x$ 展开分子:
\[ = [\int \frac{6\sin x \cos x + 4\cos x e^{\sin x} + 8\cos x}{3\sin x + e^{\sin x} + 1} dx] \]
令 $u = \sin x$,则 $du = \cos x dx$。代入原式得:
\[ = [\int \frac{6u + 4e^u + 8}{3u + e^u + 1} du] \]
提取分子中的常数 2:
\[ = 2 [\int \frac{3u + 2e^u + 4}{3u + e^u + 1} du] \]
将分子拆项,使其包含分母的形式:
\[ = 2 [\int \frac{(3u + e^u + 1) + (3 + e^u)}{3u + e^u + 1} du] \]
\[ = 2 ( [\int \frac{3u + e^u + 1}{3u + e^u + 1} du] + [\int \frac{e^u + 3}{3u + e^u + 1} du] ) \]
\[ = 2 [\int 1 du] + 2 [\int \frac{d(3u + e^u + 1)}{3u + e^u + 1}] \]
进行积分:
\[ = 2u + 2\ln|3u + e^u + 1| + C \]
最后回代 $u = \sin x$:
\[ = 2\sin x + 2\ln|3\sin x + e^{\sin x} + 1| + C \]
\end{solution}
\question
52) 计算不定积分:\[ \int \frac{\sqrt{x+2}}{\sqrt{x+1}} dx \]

\begin{solution}
方法一:设 $t = \sqrt{\frac{x+2}{x+1}}$,则 $x = \frac{2-t^2}{t^2-1}$。
方法二:设 $1+x = \sinh^2 u$,则 $dx = 2 \sinh u \cosh u du$。
代入积分式:
\[ I = \int \frac{\cosh u}{\sinh u} \cdot 2 \sinh u \cosh u du = \int 2 \cosh^2 u du = \int (1 + \cosh 2u) du \]
积分得 $u + \frac{1}{2} \sinh 2u + C = u + \sinh u \cosh u + C$。
还原变量:
\[ I = \sinh^{-1} \sqrt{1+x} + \sqrt{1+x}\sqrt{2+x} + C \]
\end{solution}

\question
54) 计算不定积分:\[ \int \frac{\sqrt{x}}{\sqrt[3]{x}-1} dx \]

\begin{solution}
设 $x = a^6$,则 $dx = 6a^5 da$。
代入积分式:
\[ I = \int \frac{a^3}{a^2-1} \cdot 6a^5 da = 6 \int \frac{a^8}{a^2-1} da \]
利用多项式除法:
\[ \frac{a^8}{a^2-1} = a^6 + a^4 + a^2 + 1 + \frac{1}{a^2-1} \]
逐项积分:
\[ I = 6 \left( \frac{a^7}{7} + \frac{a^5}{5} + \frac{a^3}{3} + a + \frac{1}{2} \ln \left| \frac{a-1}{a+1} \right| \right) + C \]
还原变量 $a = x^{\frac{1}{6}}$:
\[ I = \frac{6}{7}x^{\frac{7}{6}} + \frac{6}{5}x^{\frac{5}{6}} + 2x^{\frac{1}{2}} + 6x^{\frac{1}{6}} + 3 \ln \left| \frac{x^{1/6}-1}{x^{1/6}+1} \right| + C \]
\end{solution}
\question
7) 计算不定积分:\[ \int \frac{(x+1)^3(x-1)}{(x^2+1)^2 \sqrt{x^4+x^2+1}} dx \]

\begin{solution}
首先简化被积函数,分子展开并提取公因子:
\[ \int \frac{(x^2+2x+1)(x^2-1)}{(x^2+1)^2 \sqrt{x^4+x^2+1}} dx \]
分子分母同除以 $x^4$,利用 $x + \frac{1}{x}$ 进行换元:
\[ = \int \frac{(x+\frac{1}{x}+2)(1-\frac{1}{x^2})}{(x+\frac{1}{x})^2 \sqrt{x^2+1+\frac{1}{x^2}}} dx \]
\[ = \int \frac{(x+\frac{1}{x}+2)(1-\frac{1}{x^2})}{(x+\frac{1}{x})^2 \sqrt{(x+\frac{1}{x})^2-1}} dx \]

设 $x + \frac{1}{x} = t$,则 $(1-\frac{1}{x^2}) dx = dt$:
\[ I = \int \frac{t+2}{t^2 \sqrt{t^2-1}} dt = \int \frac{1}{t\sqrt{t^2-1}} dt + 2 \int \frac{1}{t^2\sqrt{t^2-1}} dt \]

对于第一部分,设 $t = \sec\theta$,则积分结果为 $\tan^{-1}\sqrt{t^2-1}$ 或 $\sec^{-1}t$。
对于第二部分,设 $t^2-1 = u^2$,则 $t dt = u du$:
\[ 2 \int \frac{1}{t^2\sqrt{t^2-1}} dt = \dots = \frac{2\sqrt{t^2-1}}{t} + C \]

合并结果并还原变量 $t = x + \frac{1}{x}$:
\[ I = \tan^{-1}\left(\sqrt{(x+\frac{1}{x})^2-1}\right) + \frac{2\sqrt{(x+\frac{1}{x})^2-1}}{x+\frac{1}{x}} + C \]
\end{solution}
    \question 
    \[
    \int_0^{\frac{\pi}{2}} \frac{e^x(\sin x + \cos x - 2)}{(\cos x - 2)^2} \, dx
    \]
    \begin{solution}
        有\[\int_0^{\frac{\pi}{2}} \frac{e^x(\sin x + \cos x - 2)}{(\cos x - 2)^2} \, dx 
        =\int_0^{\frac{\pi}{2}}e^{x}\left(\frac{\sin x}{(\cos x-2)^{2}}+\frac{1}{\cos x-2}\right) \, dx 
        \]
        发现被积函数形如 \(e^{x}(f(x)+f^{\prime }(x))\),其中\[
        f(x)= \frac{1}{\cos x-2}, f^{\prime }(x)=\frac{\sin x}{(\cos x-2)^{2}}
        \]
        所以\[
        \int _{0}^{\frac{\pi }{2}}e^{x}\left(\frac{1}{\cos x-2}+\frac{\sin x}{(\cos x-2)^{2}}\right)dx=\left[\frac{e^{x}}{\cos x-2}\right]_{0}^{\frac{\pi }{2}}=1 - \frac{e^{\frac{\pi}{2}}}{2}
        \]
    \end{solution}
    \question
8) 计算不定积分:\[ \int e^x \left( \frac{1}{\sqrt{x^2+1}} + \frac{1-2x^2}{(x^2+1)^{\frac{5}{2}}} \right) dx \]

\begin{solution}
观察被积函数形式,尝试利用公式 $\int e^x [f(x) + f'(x)] dx = e^x f(x) + C$。
令 $f(x) = \frac{1}{\sqrt{x^2+1}} + \frac{x}{(x^2+1)^{\frac{3}{2}}}$。
计算导数 $f'(x)$:
\[ f'(x) = -\frac{x}{(x^2+1)^{\frac{3}{2}}} + \frac{(x^2+1)^{\frac{3}{2}} - x \cdot \frac{3}{2}(x^2+1)^{\frac{1}{2}} \cdot 2x}{(x^2+1)^3} \]
\[ = -\frac{x}{(x^2+1)^{\frac{3}{2}}} + \frac{1-2x^2}{(x^2+1)^{\frac{5}{2}}} \]
将其代入原式:
\[ I = \int e^x \left( \frac{1}{\sqrt{x^2+1}} - \frac{x}{(x^2+1)^{\frac{3}{2}}} + \frac{x}{(x^2+1)^{\frac{3}{2}}} + \frac{1-2x^2}{(x^2+1)^{\frac{5}{2}}} \right) dx \]
根据恒等式性质得到结果:
\[ I = e^x \left( \frac{1}{\sqrt{x^2+1}} + \frac{x}{(x^2+1)^{\frac{3}{2}}} \right) + C \]
\end{solution}

\question
13) 计算不定积分:\[ \int \sqrt{x + \sqrt{x^2+2}} dx \]

\begin{solution}
使用双曲函数换元,设 $x = \sqrt{2} \sinh t$,则 $dx = \sqrt{2} \cosh t dt$。
同时利用恒等式 $\sqrt{x^2+2} = \sqrt{2} \cosh t$:
\[ I = \int \sqrt{\sqrt{2}\sinh t + \sqrt{2}\cosh t} \cdot \sqrt{2}\cosh t dt \]
\[ = \sqrt[4]{8} \int \sqrt{e^t} \cosh t dt = \sqrt[4]{8} \int e^{\frac{t}{2}} \frac{e^t + e^{-t}}{2} dt \]
\[ = \frac{\sqrt[4]{8}}{2} \int (e^{\frac{3t}{2}} + e^{-\frac{t}{2}}) dt = \frac{\sqrt[4]{8}}{2} \left[ \frac{2}{3}e^{\frac{3t}{2}} - 2e^{-\frac{t}{2}} \right] + C \]
还原变量 $t = \sinh^{-1} \frac{x}{\sqrt{2}}$ 得到最终结果:
\[ I = \frac{1}{\sqrt[4]{2}} \left( \frac{2}{3} e^{\frac{3}{2} \sinh^{-1} \frac{x}{\sqrt{2}}} - 2 e^{-\frac{1}{2} \sinh^{-1} \frac{x}{\sqrt{2}}} \right) + C \]
\end{solution}

%三角函数,万能公式
\question 
38) b) 证明 $\frac{d}{d\theta}(\tan^3 \theta) = 3\tan^4 \theta + 3\sec^2 \theta - 3$。
由此计算 $\int_{0}^{\frac{\pi}{4}} \tan^4 \theta d\theta$。万能公式

\begin{solution}
证明部分:
\begin{align*}
\frac{d}{d\theta} \tan^3 \theta &= 3\tan^2 \theta (\sec^2 \theta) \\
&= 3\tan^2 \theta (\tan^2 \theta + 1) \\
&= 3\tan^4 \theta + 3\tan^2 \theta \\
&= 3\tan^4 \theta + 3(\sec^2 \theta - 1) \\
&= 3\tan^4 \theta + 3\sec^2 \theta - 3
\end{align*}
计算积分部分:
\begin{align*}
\int_{0}^{\frac{\pi}{4}} \tan^4 \theta d\theta &= \frac{1}{3} \int_{0}^{\frac{\pi}{4}} (3\tan^4 \theta) d\theta \\
&= \frac{1}{3} \int_{0}^{\frac{\pi}{4}} \left( \frac{d}{d\theta} \tan^3 \theta - 3\sec^2 \theta + 3 \right) d\theta \\
&= \frac{1}{3} \left[ \tan^3 \theta - 3\tan \theta + 3\theta \right]_{0}^{\frac{\pi}{4}} \\
&= \frac{1}{3} \left[ 1 - 3 + \frac{3\pi}{4} \right] \\
&= \frac{\pi}{4} - \frac{2}{3}
\end{align*}
\end{solution}

\question
39) 证明 $\frac{d}{dx} \tan^3 x = 3\sec^4 x - 3\sec^2 x$。
由此计算 $\int_{0}^{\frac{\pi}{4}} \sec^4 x dx$。

\begin{solution}
证明部分:
\begin{align*}
\frac{d}{dx} \tan^3 x &= 3\tan^2 x \sec^2 x \\
&= 3(\sec^2 x - 1)(\sec^2 x) \\
&= 3\sec^4 x - 3\sec^2 x
\end{align*}
计算积分部分:
\begin{align*}
\int_{0}^{\frac{\pi}{4}} \sec^4 x dx &= \frac{1}{3} \int_{0}^{\frac{\pi}{4}} (3\sec^4 x) dx \\
&= \frac{1}{3} \int_{0}^{\frac{\pi}{4}} \left( \frac{d}{dx} \tan^3 x + 3\sec^2 x \right) dx \\
&= \frac{1}{3} \left[ \tan^3 x + 3\tan x \right]_{0}^{\frac{\pi}{4}} \\
&= \frac{1}{3} (1 + 3) \\
&= \frac{4}{3}
\end{align*}
\end{solution}

\question 
9) 计算不定积分 $\int \sin^3 2x \cos^2 3x \, dx$

\begin{solution}
利用积化和差及倍角公式:
\begin{align*}
\int \sin^3 2x \cos^2 3x \, dx &= \int \left( \frac{3 \sin 2x - \sin 6x}{4} \right) \left( \frac{1 + \cos 6x}{2} \right) dx \\
&= \frac{1}{8} \int (3 \sin 2x + 3 \sin 2x \cos 6x - \sin 6x - \sin 6x \cos 6x) dx \\
&= \frac{1}{8} \int \left( 3 \sin 2x + \frac{3}{2}(\sin 8x - \sin 4x) - \sin 6x - \frac{1}{2} \sin 12x \right) dx \\
&= \frac{1}{8} \left( -\frac{3}{2} \cos 2x - \frac{3}{16} \cos 8x + \frac{3}{8} \cos 4x + \frac{1}{6} \cos 6x + \frac{1}{24} \cos 12x \right) + C \\
&= -\frac{3}{16} \cos 2x + \frac{3}{64} \cos 4x + \frac{1}{48} \cos 6x - \frac{3}{128} \cos 8x + \frac{1}{192} \cos 12x + C
\end{align*}
\end{solution}

\question 
计算定积分 $I = \int_{0}^{\pi/4} \frac{3 - 4\cos 2x + \cos 4x}{3 + 4\cos 2x + \cos 4x} dx$

\begin{solution}
\textbf{第一步:利用倍角公式化简}
注意到 $1 + \cos 4x = 2\cos^2 2x$,分子分母可变形为:
\begin{itemize}
    \item 分子:$2 - 4\cos 2x + 2\cos^2 2x = 2(1 - \cos 2x)^2 = 2(2\sin^2 x)^2 = 8\sin^4 x$
    \item 分母:$2 + 4\cos 2x + 2\cos^2 2x = 2(1 + \cos 2x)^2 = 2(2\cos^2 x)^2 = 8\cos^4 x$
\end{itemize}

\textbf{第二步:化简积分式}
\[ I = \int_{0}^{\pi/4} \frac{8\sin^4 x}{8\cos^4 x} dx = \int_{0}^{\pi/4} \tan^4 x \, dx \]

\textbf{第三步:利用 $\tan^4 x = \tan^2 x(\sec^2 x - 1)$ 积分}
\begin{align*}
I &= \int_{0}^{\pi/4} (\sec^4 x - 2\sec^2 x + 1) dx \quad (\text{或利用之前证明过的结论}) \\
&= \int_{0}^{\pi/4} (\tan^2 x + 1)\sec^2 x \, dx - 2\int_{0}^{\pi/4} \sec^2 x \, dx + \int_{0}^{\pi/4} 1 \, dx \\
&= \int_{0}^{1} (u^2 + 1) du - 2[\tan x]_{0}^{\pi/4} + [x]_{0}^{\pi/4} \\
&= \left[ \frac{u^3}{3} + u \right]_{0}^{1} - 2(1) + \frac{\pi}{4} \\
&= (\frac{1}{3} + 1) - 2 + \frac{\pi}{4} \\
&= \frac{\pi}{4} - \frac{2}{3}
\end{align*}

\textbf{结果核对:} 与截图结果 $\frac{\pi}{4} - \frac{2}{3}$ 一致。
\end{solution}

\question 计算积分
\[
\int_{\frac{\pi}{12}}^{\frac{\pi}{4}}\frac{\sin 3x}{(\cos 7x+\cos x)^{2}+(\sin 7x+\sin x)^{2}} \, dx
\]
\begin{solution}
先展开分母:
\[
(\cos 7x+\cos x)^2 + (\sin 7x+\sin x)^2 = \cos^2 7x + 2\cos 7x \cos x + \cos^2 x + \sin^2 7x + 2\sin 7x \sin x + \sin^2 x
\]

利用 $\cos^2 \theta + \sin^2 \theta = 1$:
\[
= 1 + 1 + 2(\cos 7x \cos x + \sin 7x \sin x) = 2 + 2\cos(7x - x) = 2 + 2\cos 6x
\]

使用二倍角公式化简:
\[
\cos 6x = 2\cos^2 3x - 1 \implies 2 + 2\cos 6x = 4\cos^2 3x
\]

积分变为:
\[
\int_{\frac{\pi}{12}}^{\frac{\pi}{4}} \frac{\sin 3x}{4\cos^2 3x} \, dx = \frac{1}{4} \int_{\frac{\pi}{12}}^{\frac{\pi}{4}} \frac{\sin 3x}{\cos 3x} \cdot \frac{1}{\cos 3x} \, dx = \frac{1}{4} \int_{\frac{\pi}{12}}^{\frac{\pi}{4}} \tan 3x \sec 3x \, dx
\]

作代换 $u = 3x \implies du = 3 dx$:
\[
\frac{1}{4} \int \tan 3x \sec 3x \, dx = \frac{1}{12} \int \sec u \tan u \, du = \frac{1}{12} \sec u + C = \frac{1}{12} \sec 3x
\]

代入上下限:
\[
\left[\frac{1}{12} \sec 3x \right]_{\frac{\pi}{12}}^{\frac{\pi}{4}} = \frac{1}{12} \left( \sec \frac{3\pi}{4} - \sec \frac{\pi}{4} \right) = \frac{1}{12} \left( -\sqrt{2} - 2 \right) \text{（注意 $\sec 3\pi/4 = -\sqrt{2}$）}
\]

整理结果:
\[
\int_{\frac{\pi}{12}}^{\frac{\pi}{4}}\frac{\sin 3x}{(\cos 7x+\cos x)^{2}+(\sin 7x+\sin x)^{2}} \, dx = \frac{2 - \sqrt{2}}{12}
\]
\end{solution}

\question 
计算定积分:$I = 8 \int_{0}^{\pi/4} \left( \frac{\sin^5 \theta}{\sin^2 \theta} - \frac{\cos^5 \theta}{\cos^2 \theta} \right) d\theta$

\begin{solution}
首先简化被积函数:
\begin{align*}
\frac{\sin^5 \theta}{\sin^2 \theta} - \frac{\cos^5 \theta}{\cos^2 \theta} &= \sin^3 \theta - \cos^3 \theta \\
&= (\sin \theta - \cos \theta)(\sin^2 \theta + \sin \theta \cos \theta + \cos^2 \theta) \\
&= (\sin \theta - \cos \theta)(1 + \sin \theta \cos \theta)
\end{align*}

\textbf{另一种路径（根据笔记中的化简步骤）:}
利用倍角公式和化简过程:
\begin{align*}
I &= \int_{0}^{\pi/4} 8 (3 + 1 + 2\cos 4\theta) \cos 2\theta \, d\theta \\
&= \int_{0}^{\pi/4} (16\cos 2\theta + 16\cos 4\theta \cos 2\theta) \, d\theta \\
&= \int_{0}^{\pi/4} (16\cos 2\theta + 8\cos 6\theta + 8\cos 2\theta) \, d\theta \\
&= \int_{0}^{\pi/4} (24\cos 2\theta + 8\cos 6\theta) \, d\theta
\end{align*}

\textbf{积分计算:}
\begin{align*}
I &= \left[ 12\sin 2\theta + \frac{4}{3}\sin 6\theta \right]_{0}^{\pi/4} \\
&= \left( 12\sin\frac{\pi}{2} + \frac{4}{3}\sin\frac{3\pi}{2} \right) - 0 \\
&= 12(1) + \frac{4}{3}(-1) = 12 - \frac{4}{3} = \frac{32}{3}
\end{align*}

\textbf{最终结果:} $\frac{32}{3}$ \quad wrong ans
\end{solution}

\question
46) 计算不定积分:\[ \int \cos^5 x \sin 5x dx \]

\begin{solution}
利用积化和差公式逐步降幂:
\begin{align*}
I &= \frac{1}{2} \int (2 \sin 5x \cos x) \cos^4 x dx = \frac{1}{2} \int (\sin 6x + \sin 4x) \cos^4 x dx \\
&= \frac{1}{4} \int (2 \sin 6x \cos x + 2 \sin 4x \cos x) \cos^3 x dx \\
&= \frac{1}{4} \int (\sin 7x + \sin 5x + \sin 5x + \sin 3x) \cos^3 x dx
\end{align*}
继续展开并对每一项正弦函数进行积分,最终结果为:
\[ I = -\frac{1}{320} \cos 10x - \frac{1}{128} \cos 8x - \dots + C \]
\end{solution}

\question
50) 计算不定积分:\[ \int \frac{\sin^4 x}{\sqrt{1-\sin 2x}} dx \]

\begin{solution}
利用恒等式 $1-\sin 2x = (\cos x - \sin x)^2$。
被积函数化为 $\frac{\sin^4 x}{|\cos x - \sin x|}$,其中 $\cos x - \sin x = \sqrt{2} \cos(x + \frac{\pi}{4})$。
设 $t = x + \frac{\pi}{4}$,则 $x = t - \frac{\pi}{4}$。
代入积分式:
\[ I = \frac{1}{\sqrt{2}} \int \frac{\sin^4(t-\frac{\pi}{4})}{\cos t} dt = \frac{1}{4\sqrt{2}} \int \frac{(\sin t - \cos t)^4}{\cos t} dt \]
通过展开分子并逐项积分:
\[ I = \frac{1}{4\sqrt{2}} \left( \ln|\sec t + \tan t| + \frac{4}{3} \sin^3 t \cos t - 4\sin t \right) + C \]
其中 $t = x + \frac{\pi}{4}$。
\end{solution}
\question
\[
\int_{\frac{\pi}{6}}^{\frac{\pi}{3}} \frac{\cot^3 x}{\csc x} \,dx
\]
\begin{solution}
令 $u=\sin x$,则
\[
du=\cos x \, dx
\]
当 $x=\frac{\pi}{6}$ 时,$u=\frac{1}{2}$  
当 $x=\frac{\pi}{3}$ 时,$u=\frac{\sqrt{3}}{2}$
\[
\int_{\frac{\pi}{6}}^{\frac{\pi}{3}} \frac{\cos^3 x}{\sin^2 x} \, dx
= \int_{\frac{1}{2}}^{\frac{\sqrt{3}}{2}} \frac{\cos^2 x}{u^2} \, du
\]
\[
= \int_{\frac{1}{2}}^{\frac{\sqrt{3}}{2}} \frac{1-\sin^2 x}{u^2} \, du
= \int_{\frac{1}{2}}^{\frac{\sqrt{3}}{2}} \frac{1-u^2}{u^2} \, du
\]
\[
= \int_{\frac{1}{2}}^{\frac{\sqrt{3}}{2}} \left(\frac{1}{u^2}-1\right) \, du
= \left[-\frac{1}{u}-u\right]_{\frac{1}{2}}^{\frac{\sqrt{3}}{2}}
\]
\[
= \left[\frac{1}{u}+u\right]_{\frac{\sqrt{3}}{2}}^{\frac{1}{2}}
= \left[\frac{1+u^2}{u}\right]_{\frac{\sqrt{3}}{2}}^{\frac{1}{2}}
\]
\[
= \frac{1+\frac{1}{4}}{\frac{1}{2}}-\frac{1+\frac{3}{4}}{\frac{\sqrt{3}}{2}}
= \frac{5}{2}-\frac{7}{2\sqrt{3}}
\]
\[
= \frac{5}{2}-\frac{7\sqrt{3}}{6}
= \frac{1}{6}(15-7\sqrt{3})
\]
\end{solution}
\question
21) 计算不定积分:\[ \int \frac{\sqrt{\cos 2x}}{\sin x} dx \]

\begin{solution}
利用三角恒等式 $\cos 2x = 2\cos^2 x - 1$:
\[ \int \frac{\sqrt{2\cos^2 x - 1}}{\sin^2 x} \sin x dx \]
设 $t = \cos x$,则 $dt = -\sin x dx$。
代入后积分变为:
\[ I = \int \frac{\sqrt{2t^2 - 1}}{t^2 - 1} dt = \int \left( \frac{2(t^2-1) + 1}{(t^2-1)\sqrt{2t^2-1}} \right) dt \]
拆分为两个部分:
\[ I = \int \frac{2}{\sqrt{2t^2-1}} dt + \int \frac{1}{(t^2-1)\sqrt{2t^2-1}} dt \]
计算并整理得:
\[ I = \sqrt{2} \ln(\sqrt{2}t + \sqrt{2t^2-1}) + \frac{1}{2} \ln\left| \frac{\sqrt{2t^2-1} - t}{\sqrt{2t^2-1} + t} \right| + C \]
还原变量 $t = \cos x$:
\[ I = \sqrt{2} \ln(\sqrt{2}\cos x + \sqrt{\cos 2x}) + \frac{1}{2} \ln\left| \frac{\sqrt{\cos 2x} - \cos x}{\sqrt{\cos 2x} + \cos x} \right| + C \]
\end{solution}

\question
56) 计算不定积分:\[ \int x^2 e^{\sin^{-1} x} dx \]

\begin{solution}
采用换元法,设 $x = \sin t$,则 $dx = \cos t dt,t = \sin^{-1} x$。
代入积分式:
\[ I = \int \sin^2 t e^t \cos t dt \]
利用三角恒等式 $\sin^2 t \cos t = \frac{1}{4} (\cos t - \cos 3t)$:
\[ I = \frac{1}{4} \left( \int e^t \cos t dt - \int e^t \cos 3t dt \right) \]
利用公式 $\int e^{at} \cos bt dt = \frac{e^{at}}{a^2+b^2} (a \cos bt + b \sin bt)$:
\[ I = \frac{1}{4} \left[ \frac{e^t}{2} (\cos t + \sin t) - \frac{e^t}{10} (\cos 3t + 3 \sin 3t) \right] + C \]
还原变量 $\sin t = x, \cos t = \sqrt{1-x^2}$:
利用 $\cos 3t = \sqrt{1-x^2}(1-4x^2)$ 和 $\sin 3t = 3x-4x^3$ 进行替换,得到最终结果。
\end{solution}
\question
5) 计算不定积分:\[ \int \frac{\sec^2 x}{(\sec x + \tan x)^{\frac{9}{2}}} dx \]

\begin{solution}
设 $t = \sec x + \tan x$。
对两边求导:
\[ dt = (\sec x \tan x + \sec^2 x) dx = \sec x (\tan x + \sec x) dx = t \sec x dx \]
由此得 $dx = \frac{dt}{t \sec x}$。

利用恒等式 $\sec^2 x - \tan^2 x = 1$:
\[ (\sec x + \tan x)(\sec x - \tan x) = 1 \implies \sec x - \tan x = \frac{1}{t} \]
将 $\sec x + \tan x = t$ 与 $\sec x - \tan x = \frac{1}{t}$ 相加得:
\[ 2 \sec x = t + \frac{1}{t} \implies \sec x = \frac{1}{2}(t + \frac{1}{t}) \]

代入原积分式:
\begin{align*}
I &= \int \frac{\sec x \cdot \sec x}{t^{\frac{9}{2}}} \cdot \frac{dt}{t \sec x} \\
&= \int \frac{\sec x}{t^{\frac{11}{2}}} dt \\
&= \int \frac{1}{2}(t + t^{-1}) t^{-\frac{11}{2}} dt \\
&= \frac{1}{2} \int (t^{-\frac{9}{2}} + t^{-\frac{13}{2}}) dt
\end{align*}

进行积分计算:
\[ I = \frac{1}{2} \left[ \frac{t^{-\frac{7}{2}}}{-\frac{7}{2}} + \frac{t^{-\frac{11}{2}}}{-\frac{11}{2}} \right] + C \]
\[ = -\frac{1}{7} t^{-\frac{7}{2}} - \frac{1}{11} t^{-\frac{11}{2}} + C \]

还原变量 $t = \sec x + \tan x$:
\[ I = -\frac{1}{7(\sec x + \tan x)^{\frac{7}{2}}} - \frac{1}{11(\sec x + \tan x)^{\frac{11}{2}}} + C \]
\end{solution}

\question 已知代换求积分
\[
\int_{0}^{\frac{\pi}{4}} \frac{4\tan x}{1+\sin^2 x} \, dx
\]
\begin{solution}
作代换
\[
u = 1 + \sin^2 x
\]
\[
\frac{du}{dx} = 2\sin x \cos x
\]
\[
dx = \frac{du}{2\sin x \cos x}
\]

积分的上下限变换
\[
x = 0 \implies u = 1
\]
\[
x = \frac{\pi}{4} \implies u = \frac{3}{2}
\]

代入积分
\[
\int_{0}^{\frac{\pi}{4}} \frac{4\tan x}{1+\sin^2 x} \, dx
= \int_{1}^{\frac{3}{2}} \frac{4\tan x}{u} \cdot \frac{du}{2\sin x \cos x}
= \int_{1}^{\frac{3}{2}} \frac{2}{u(2-u)} \, du
\]

部分分式分解
\[
\frac{2}{u(2-u)} = \frac{1}{u} + \frac{1}{2-u}
\]

计算积分
\[
\int_{1}^{\frac{3}{2}} \left( \frac{1}{u} + \frac{1}{2-u} \right) du
= [\ln|u| - \ln|2-u|]_{1}^{\frac{3}{2}}
= \ln \frac{3/2}{2-3/2} - \ln \frac{1}{2-1}
= \ln 3
\]
\end{solution}

\question 代换$u = 1 - \tan^2 x$,求定积分
\[
\int_{0}^{\pi/6} \tan x \sec 2x \, dx
\]
\begin{solution}
作代换
\[
u = 1 - \tan^2 x
\]
\[
\frac{du}{dx} = -2\tan x \sec^2 x
\]
\[
dx = -\frac{du}{2\tan x \sec^2 x}
\]

积分上下限变换
\[
x = 0 \implies u = 1
\]
\[
x = \frac{\pi}{6} \implies u = \frac{2}{3}
\]

代入积分
\[
\int_{0}^{\pi/6} \tan x \sec 2x \, dx
= \int_{1}^{2/3} \tan x \sec 2x \cdot \left(-\frac{du}{2\tan x \sec^2 x}\right)
= \frac{1}{2} \int_{2/3}^{1} \frac{1}{u} \, du
\]

计算积分
\[
\frac{1}{2} \int_{2/3}^{1} \frac{1}{u} \, du
= \left[ \frac{1}{2} \ln|u| \right]_{2/3}^{1}
= \frac{1}{2} \ln 1 - \frac{1}{2} \ln \frac{2}{3}
= \frac{1}{2} \ln \frac{3}{2}
\]
\end{solution}



\question
\[
\int_{\frac{\pi}{6}}^{\frac{\pi}{3}} \frac{1}{\sin x \sin 2x} dx
= \frac{1}{2}\int_{\frac{\pi}{6}}^{\frac{\pi}{3}} \frac{\cos^2 x + \sin^2 x}{2\sin^2 x \cos x} dx
\]

\begin{solution}
\[
= \frac{1}{2}\int_{\frac{\pi}{6}}^{\frac{\pi}{3}} \left( \frac{\cos x}{\sin x} + \sec x \right) dx
= \frac{1}{2}\int_{\frac{\pi}{6}}^{\frac{\pi}{3}} \left( \cot x \csc x + \sec x \right) dx
\]

\[
= \frac{1}{2}\left[ -\csc x + \ln|\sec x + \tan x| \right]_{\frac{\pi}{6}}^{\frac{\pi}{3}}
\]

\[
= \frac{1}{2}\left[\left(\ln(2+\sqrt{3}) - \frac{2}{\sqrt{3}}\right) - \left(-2 + \ln\left(\frac{2}{\sqrt{3}} + \frac{1}{\sqrt{3}}\right)\right)\right]
\]

\[
= \frac{1}{2}\left[\ln(2+\sqrt{3}) - \frac{2}{\sqrt{3}} + 2 - \ln\sqrt{3}\right]
\]

\[
= \frac{1}{2}\ln\left(\frac{2+\sqrt{3}}{\sqrt{3}}\right) + 1 - \frac{\sqrt{3}}{3}
\]
\end{solution}

\question
计算积分
\[
\int \frac{\sin^{8}x - \cos^{8}x}{1 - \frac{1}{2} \sin^{2} 2x} \, dx
\]

\begin{solution}
从分子平方差和分母的正弦二倍角公式出发:
\[
\int \frac{(\sin^{4}x - \cos^{4}x)(\sin^{4}x + \cos^{4}x)}{1 - \frac{1}{2} (2\sin x \cos x)^2} \, dx
= \int \frac{(\sin^{4}x - \cos^{4}x)(\sin^{4}x + \cos^{4}x)}{1 - 2\sin^{2}x \cos^{2}x} \, dx
\]

将分母写成平方差形式:
\[
\int \frac{(\sin^{4}x - \cos^{4}x)(\sin^{4}x + \cos^{4}x)}{1^2 - 2\sin^{2}x \cos^{2}x} \, dx
= \int \frac{(\sin^{4}x - \cos^{4}x)(\sin^{4}x + \cos^{4}x)}{(\sin^{2}x + \cos^{2}x)^2 - 2\sin^{2}x \cos^{2}x} \, dx
\]

展开分子平方差和分母括号:
\[
\int \frac{(\sin^{2}x - \cos^{2}x)(\sin^{2}x + \cos^{2}x)(\sin^{4}x + \cos^{4}x)}{\sin^{4}x + \cos^{4}x + 2\sin^{2}x \cos^{2}x - 2\sin^{2}x \cos^{2}x} \, dx
\]

化简:
\[
\int \frac{-\cos 2x (\sin^{4}x + \cos^{4}x)}{\cos^{4}x + \sin^{4}x} \, dx
= \int -\cos 2x \, dx
\]

积分结果:
\[
-\frac{1}{2} \sin 2x + C
\]
\end{solution}

\question 计算积分
\[
\int \frac{3 \sin^2 x \cos^2 x}{(\cos^3 x - \sin^3 x)^2} \, dx
\]
\begin{solution}
重写被积函数,提取 \(\tan x\) 和 \(\sec x\):
\[
\int \frac{3 \sin^2 x \cos^2 x}{(\cos^3 x - \sin^3 x)^2} \, dx 
= \int 3 \frac{(\sin x \cos x)^2}{(\cos^3 x - \sin^3 x)^2} \, dx
= \int 3 \tan^2 x \sec^2 x (1 - \tan^3 x)^{-2} \, dx
\]

作代换:
\[
u = \tan x \implies du = \sec^2 x \, dx
\]

代入:
\[
\int 3 \tan^2 x \sec^2 x (1 - \tan^3 x)^{-2} \, dx
= \int 3 u^2 (1 - u^3)^{-2} \, du
\]

识别反链式法则:
\[
d(1-u^3) = -3 u^2 du \implies u^2 du = -\frac{1}{3} d(1-u^3)
\]

因此:
\[
\int 3 u^2 (1 - u^3)^{-2} \, du = \int 3 \cdot \left(-\frac{1}{3} \right) (1-u^3)^{-2} d(1-u^3) = -\int (1-u^3)^{-2} d(1-u^3)
\]

积分:
\[
-\int (1-u^3)^{-2} d(1-u^3) = (1-u^3)^{-1} + C
\]

回代 \(u = \tan x\):
\[
\int \frac{3 \sin^2 x \cos^2 x}{(\cos^3 x - \sin^3 x)^2} \, dx = \frac{1}{1 - \tan^3 x} + C
\]

\end{solution}

\question 计算积分
\[
\int_{\frac{\pi}{8}}^{\frac{\pi}{4}} \frac{\sin^6 x + \cos^6 x}{\sin^2 x \cos^2 x} \, dx
\]
\begin{solution}
先拆分分数:
\[
\int_{\frac{\pi}{8}}^{\frac{\pi}{4}} \frac{\sin^6 x}{\sin^2 x \cos^2 x} + \frac{\cos^6 x}{\sin^2 x \cos^2 x} \, dx
= \int_{\frac{\pi}{8}}^{\frac{\pi}{4}} \frac{\sin^4 x}{\cos^2 x} + \frac{\cos^4 x}{\sin^2 x} \, dx
\]

将 \(\sin^2 x = 1 - \cos^2 x\),\(\cos^2 x = 1 - \sin^2 x\) 展开:
\[
\int_{\frac{\pi}{8}}^{\frac{\pi}{4}} \frac{(1-\cos^2 x)^2}{\cos^2 x} + \frac{(1-\sin^2 x)^2}{\sin^2 x} \, dx
= \int_{\frac{\pi}{8}}^{\frac{\pi}{4}} \sec^2 x - 2 + \cos^2 x + \csc^2 x - 2 + \sin^2 x \, dx
\]

合并同类项:
\[
\int_{\frac{\pi}{8}}^{\frac{\pi}{4}} \sec^2 x + \csc^2 x + (\cos^2 x + \sin^2 x) - 4 \, dx
= \int_{\frac{\pi}{8}}^{\frac{\pi}{4}} \sec^2 x + \csc^2 x - 3 \, dx
\]

直接积分:
\[
[\tan x - \cot x - 3x]_{\frac{\pi}{8}}^{\frac{\pi}{4}}
\]

化简为倍角形式:
\[
[\tan x - \cot x - 3x] = [3x + 2\cot 2x]_{\frac{\pi}{8}}^{\frac{\pi}{4}} = \frac{1}{8}(16 - 3\pi)
\]

\end{solution}

    \question 求积分
    \[
    \int_{0}^{\pi/2} \frac{\sin 2x}{\sqrt{4 - \sin^4 x}} \, dx
    \]
\begin{solution}
令 \(u = \sin^2 x \implies du = 2\sin x \cos x \, dx = \sin 2x \, dx\):
\[
\int_{0}^{\pi/2} \frac{\sin 2x}{\sqrt{4 - \sin^4 x}} \, dx = \int_{0}^{1} \frac{1}{\sqrt{4 - u^2}} \, du
\]

\[
= \left[\arcsin \frac{u}{2}\right]_{0}^{1} = \arcsin \frac{1}{2} - \arcsin 0 = \frac{\pi}{6}
\]
\end{solution}
\question 计算积分
\[
\int_{0}^{\frac{\pi}{2}} \frac{1+\cos x + \sin x - \tan x}{1+\tan x} \, dx
\]
\begin{solution}
先将 $\tan x$ 用正弦余弦表示:
\[
\int_{0}^{\frac{\pi}{2}} \frac{1+\cos x + \sin x - \frac{\sin x}{\cos x}}{1 + \frac{\sin x}{\cos x}} \, dx
= \int_{0}^{\frac{\pi}{2}} \frac{\cos x + \cos^2 x + \sin x \cos x - \sin x}{\cos x + \sin x} \, dx
\]
这里是通过分子分母同乘 $\cos x$ 得到的。

重组分子:
\[
\int_{0}^{\frac{\pi}{2}} \frac{\cos x - \sin x + \cos x \sin x + \cos^2 x}{\cos x + \sin x} \, dx
= \int_{0}^{\frac{\pi}{2}} \frac{\cos x - \sin x}{\cos x + \sin x} + \frac{\cos x \sin x + \cos^2 x}{\cos x + \sin x} \, dx
\]

提取公因子:
\[
= \int_{0}^{\frac{\pi}{2}} \frac{\cos x - \sin x}{\cos x + \sin x} + \frac{\cos x (\sin x + \cos x)}{\cos x + \sin x} \, dx
\]

使用 $\int \frac{f'(x)}{f(x)} \, dx = \ln|f(x)| + C$ 的形式:
\[
\int_{0}^{\frac{\pi}{2}} \frac{\cos x - \sin x}{\cos x + \sin x} + \cos x \, dx
= \left[ \ln|\cos x + \sin x| + \sin x \right]_{0}^{\frac{\pi}{2}}
\]

代入上下限:
\[
\left[ \ln(0+1) + 1 \right] - \left[ \ln(1+0) + 0 \right] = 1
\]

所以积分结果为
\[
\int_{0}^{\frac{\pi}{2}} \frac{1+\cos x + \sin x - \tan x}{1+\tan x} \, dx = 1
\]
\end{solution}


    \question $$\int \frac{\sin x \cos x}{\sin x + \cos x}\, dx $$
    \begin{solution}
        由$(\sin x+\cos x)^2=1+2\sin x \cos x, $
        \begin{align*}
            \int \frac{\sin x \cos x}{\sin x + \cos x}\, dx 
            &= \frac{1}{2}\int \frac{(\sin x+\cos x)^{2}-1}{\sin x+\cos x}\,dx \\
            & =\frac{1}{2}\int \left(\sin x+\cos x-\frac{1}{\sin x+\cos x}\right)dx
        \end{align*}
        现考虑积分
        \[
        I=\int \frac{1}{\sin x+\cos x}\,dx
        \]
        由于$ \sin x+\cos x=\sqrt{2}\sin \left(x+\dfrac{\pi }{4}\right),$变为
        \begin{align*}
            I&= \int \frac{1}{\sqrt{2}\sin (x+\frac{\pi }{4})}\,dx \\
            &= \frac{1}{\sqrt{2}}\int \csc \left(x+\frac{\pi }{4}\right)\,dx \\
            &= -\frac{1}{\sqrt{2}} \ln \left| \csc \left(x + \frac{\pi}{4}\right) + \cot \left(x + \frac{\pi}{4}\right)\right| + C_1
        \end{align*}
        代入原式可得
        $$\int \frac{\sin x \cos x}{\sin x + \cos x} \,dx = \frac{1}{2}(\sin x - \cos x) + \frac{1}{2\sqrt{2}} \ln \left| \csc \left(x + \frac{\pi}{4}\right) + \cot \left(x + \frac{\pi}{4}\right) \right| + C$$
    \end{solution}
    \begin{solution}
        由$\sin x+\cos x=\sqrt2 \sin\left(x+\dfrac{\pi}{4}\right),$
        \[
        \int \frac{\sin x \cos x}{\sin x + \cos x} \, dx 
        = \frac{1}{\sqrt{2}} \int \frac{\sin x \cos x}{\sin\left(x + \frac{\pi}{4}\right)} \, dx         
        \]
        令 $u = x + \dfrac{\pi}{4},\; dx = du,$ 则
        \[
        \sin x = \frac{\sin u - \cos u}{\sqrt{2}}, \;
        \cos x = \frac{\sin u + \cos u}{\sqrt{2}}
        \]
        \begin{align*}
        \therefore\quad 
        \int \frac{\sin x \cos x}{\sin x + \cos x} \, dx 
        &= \frac{1}{\sqrt{2}} \int \frac{\frac{1}{2}(\sin u - \cos u)(\sin u + \cos u)}{\sin u} \, du \\
        &= \frac{1}{2\sqrt{2}} \int \frac{2\sin^2 u - 1}{\sin u} \, du \\ 
        &= \frac{1}{\sqrt{2}} \int \sin u \, du- \frac{1}{2\sqrt{2}} \int \csc u \, du \\ 
        &=-\frac{\cos(x+\frac{\pi}{4})}{\sqrt{2}}
        +\frac{1}{2\sqrt{2}} \ln \left| \csc \left(x + \frac{\pi}{4}\right) + \cot \left(x + \frac{\pi}{4}\right) \right|+C 
        \end{align*}
    \end{solution}

    \question 求积分
\begin{solution}
考虑积分
\[
\int_{0}^{\pi/4} \frac{10}{2 - \tan x} \, dx
\]

令 \(u = \tan x \implies dx = \frac{du}{1+u^2}\),积分变为
\[
\int_{0}^{1} \frac{10}{(2-u)(1+u^2)} \, du
\]

使用部分分式分解:
\[
\frac{10}{(2-u)(1+u^2)} = \frac{A+Bu}{1+u^2} + \frac{C}{2-u}
\]

确定系数:
\begin{itemize}
    \item \(u=2 \Rightarrow 10 = 5C \Rightarrow C=2\)
    \item \(u=0 \Rightarrow 10 = 2A + C = 2A+2 \Rightarrow A=4\)
    \item \(u=1 \Rightarrow 10 = (A+B) + 2C = 4+B+4 \Rightarrow B=2\)
\end{itemize}

因此积分可写为
\[
\int_{0}^{1} \frac{4+2u}{u^2+1} + \frac{2}{2-u} \, du = \int_{0}^{1} \frac{4}{u^2+1} + \frac{2u}{u^2+1} + \frac{2}{2-u} \, du
\]

逐项积分:
\[
\int \frac{4}{u^2+1} \, du = 4\arctan u, \quad
\int \frac{2u}{u^2+1} \, du = \ln(u^2+1), \quad
\int \frac{2}{2-u} \, du = -2\ln|2-u|
\]

代入上下限:
\[
\left[4\arctan u + \ln(u^2+1) - 2\ln|2-u|\right]_{0}^{1} = 4 \arctan 1 + \ln 2 - 2\ln 1 - (0 + \ln 1 - 2\ln 2) = \pi + 3\ln 2
\]

最终结果:
\[
\pi + 3\ln 2
\]
\end{solution}

\question
\[
I=\int_{\arcsin \frac{3}{5}}^{\arccos \frac{3}{5}}
\frac{1}{(\sin x+2\cos x)(\sin x+3\cos x)}\,dx
\]

\begin{solution}
先化简被积函数
\[
=\int_{\arcsin \frac{3}{5}}^{\arccos \frac{3}{5}}
\frac{1}{\sin^2 x+5\sin x\cos x+6\cos^2 x}\,dx
\]

上下同除以 $\cos^2 x$
\[
=\int_{\arcsin \frac{3}{5}}^{\arccos \frac{3}{5}}
\frac{\frac{1}{\cos^2 x}}
{\frac{\sin^2 x}{\cos^2 x}
+5\frac{\sin x\cos x}{\cos^2 x}
+6\frac{\cos^2 x}{\cos^2 x}}\,dx
\]

化为正切函数
\[
=\int_{\arcsin \frac{3}{5}}^{\arccos \frac{3}{5}}
\frac{\sec^2 x}{\tan^2 x+5\tan x+6}\,dx
\]

令 $u=\tan x$,则 $du=\sec^2 x\,dx$

当 $x=\arcsin \frac{3}{5}$ 时,
\[
u=\frac{3}{4}
\]

当 $x=\arccos \frac{3}{5}$ 时,
\[
u=\frac{4}{3}
\]

积分化为
\[
I=\int_{\frac{3}{4}}^{\frac{4}{3}}
\frac{1}{u^2+5u+6}\,du
\]

因式分解
\[
=\int_{\frac{3}{4}}^{\frac{4}{3}}
\frac{1}{(u+2)(u+3)}\,du
\]

设
\[
\frac{1}{(u+2)(u+3)}
=\frac{A}{u+2}+\frac{B}{u+3}
\]

比较系数得
\[
1=A(u+3)+B(u+2)
\]

解得
\[
A=1,\qquad B=-1
\]

代回积分
\[
I=\int_{\frac{3}{4}}^{\frac{4}{3}}
\left(\frac{1}{u+2}-\frac{1}{u+3}\right)\,du
\]

计算得
\[
=\left[\ln(u+2)-\ln(u+3)\right]_{\frac{3}{4}}^{\frac{4}{3}}
\]

\[
=\left(\ln\frac{10}{3}-\ln\frac{13}{3}\right)
-\left(\ln\frac{11}{4}-\ln\frac{15}{4}\right)
\]

\[
=\ln\frac{10}{13}-\ln\frac{11}{15}
=\ln\frac{150}{143}
\]
\end{solution}


    \question
求
\[
\int_{0}^{\frac{\pi}{2}} \frac{1}{1+k^2 \tan^2 x} \, dx, \quad |k| \neq 1
\]

\begin{solution}
首先用代换:
\[
u = \tan x, \quad du = \sec^2 x \, dx, \quad dx = \frac{du}{\sec^2 x}, \quad x=0 \to u=0, \ x=\frac{\pi}{2} \to u=\infty
\]

积分变为:
\[
\int_{0}^{\infty} \frac{du}{1+k^2 u^2} \cdot \frac{1}{1+u^2} = \int_{0}^{\infty} \frac{1}{(u^2+1)(u^2+k^2)} \, du
\]

使用部分分式分解:
\[
\frac{k}{(u^2+1)(u^2+k^2)} = \frac{B}{u^2+1} + \frac{D}{u^2+k^2}
\]

由系数比较得到:
\[
B+D = 0 \implies D=-B, \quad Bk^2 + D = k \implies B = \frac{k}{k^2-1}, \ D = \frac{-k}{k^2-1}
\]

因此积分变为:
\[
\int_{0}^{\infty} \frac{k}{(u^2+1)(u^2+k^2)} \, du = \int_{0}^{\infty} B \left( \frac{1}{u^2+1} - \frac{1}{u^2+k^2} \right) \, du
\]

计算得到:
\[
\int_{0}^{\infty} \frac{1}{u^2+1} \, du = \arctan u \Big|_0^{\infty} = \frac{\pi}{2}, \quad
\int_{0}^{\infty} \frac{1}{u^2+k^2} \, du = \frac{1}{k} \arctan \frac{u}{k} \Big|_0^{\infty} = \frac{\pi}{2k}
\]

代入 B:
\[
\int_{0}^{\frac{\pi}{2}} \frac{1}{1+k^2 \tan^2 x} \, dx = B \left( \frac{\pi}{2} - \frac{\pi}{2k} \right) = \frac{k}{k^2-1} \cdot \frac{\pi}{2} \cdot \frac{k-1}{k} = \frac{\pi}{2(k+1)}
\]
\end{solution}

\question
计算积分
\[
\int_{0}^{\pi/3} \frac{6}{\sin x + \sin 2x} \, dx
\]

\begin{solution}
化简被积函数:
\[
\sin x + \sin 2x = \sin x + 2 \sin x \cos x = \sin x (1+2\cos x)
\]

代换:
\[
u = \cos x \implies du = -\sin x \, dx \implies dx = -\frac{du}{\sin x}
\]

积分上下限:
\[
x = 0 \implies u = 1, \quad x = \frac{\pi}{3} \implies u = \frac{1}{2}
\]

积分变为:
\[
\int_{0}^{\pi/3} \frac{6}{\sin x (1+2\cos x)} \, dx = \int_{1}^{1/2} \frac{6}{\sin x (1+2u)} \left(-\frac{du}{\sin x}\right)
= \int_{1/2}^{1} \frac{6}{\sin^2 x (1+2u)} \, du
\]

利用 $\sin^2 x = 1-u^2$:
\[
\int_{1/2}^{1} \frac{6}{(1-u^2)(1+2u)} \, du = \int_{1/2}^{1} \frac{6}{(1-u)(1+u)(1+2u)} \, du
\]

使用部分分式:
\[
\frac{6}{(1-u)(1+u)(1+2u)} = \frac{A}{1-u} + \frac{B}{1+u} + \frac{C}{1+2u}
\]

求系数:
\[
u=1 \implies 6 = 6A \implies A = 1, \quad
u=-1 \implies 6 = -2B \implies B=-3, \quad
u=-\frac{1}{2} \implies 6 = \frac{3}{4}C \implies C = 8
\]

所以:
\[
\frac{6}{(1-u)(1+u)(1+2u)} = \frac{1}{1-u} - \frac{3}{1+u} + \frac{8}{1+2u}
\]

积分:
\[
\int_{1/2}^{1} \left( \frac{1}{1-u} - \frac{3}{1+u} + \frac{8}{1+2u} \right) du
= \left[ \ln|1-u| - 3 \ln|1+u| + 4 \ln|1+2u| \right]_{1/2}^{1}
\]

代入上下限:
\[
= \left( \ln(1-1) - 3\ln(1+1) + 4 \ln(1+2\cdot 1) \right) - \left( \ln(1-1/2) - 3\ln(1+1/2) + 4 \ln(1+2\cdot 1/2) \right)
\]

化简:
\[
= (4\ln 3 - 3 \ln 2 - \ln 0) - (4\ln 2 - 3\ln \frac{3}{2} - \ln \frac{1}{2}) 
= 8 \ln 2 - 3 \ln 3
\]
\end{solution}

\question 计算积分
\[
\int \frac{\cos^3 x}{(1+\sin x)\sin x} \, dx
\]
\begin{solution}
使用代换
\[
u = \sin x + \csc x \implies \frac{du}{dx} = \cos x - \cot x \csc x \implies dx = \frac{1}{\cos x - \cot x \csc x} \, du
\]

代入积分:
\[
\int \frac{\cos^3 x}{(1+\sin x)\sin x} \, dx 
= \int \frac{\cos^3 x}{(1+\sin x)\sin x} \cdot \frac{1}{\cos x - \cot x \csc x} \, du
\]

将 $\cot x$ 和 $\csc x$ 用正弦余弦表示:
\[
= \int \frac{\cos^3 x}{(1+\sin x)\sin x} \cdot \frac{1}{\cos x - \frac{\cos x}{\sin x}\cdot \frac{1}{\sin x}} \, du
= \int \frac{\cos^3 x}{(1+\sin x)\sin x} \cdot \frac{1}{\cos x (1 - \frac{1}{\sin^2 x})} \, du
\]

化简:
\[
= \int \frac{\cos^3 x}{(1+\sin x)\sin x} \cdot \frac{\sin^2 x}{\cos x (\sin^2 x-1)} \, du
= \int \frac{\cos^2 x}{1+\sin x} \cdot \frac{\sin x}{\sin^2 x-1} \, du
= \int \frac{\cos^2 x}{1+\sin x} \cdot \frac{\sin x}{-\cos^2 x} \, du
\]

进一步化简:
\[
= \int -\frac{\sin x}{1+\sin x} \, du
= \int -\frac{1}{\csc x + \sin x} \, du
= \int -\frac{1}{u} \, du
= -\ln|u| + C
= \ln\left|\frac{1}{u}\right| + C
\]

代回原变量:
\[
= \ln\left|\frac{1}{\sin x + \csc x}\right| + C
= \ln\left|\frac{1}{\sin x + \frac{1}{\sin x}}\right| + C
= \ln\left|\frac{1}{\frac{\sin^2 x+1}{\sin x}}\right| + C
= \ln\left|\frac{\sin x}{\sin^2 x+1}\right| + C
\]
\end{solution}

\question 求积分
\[
\int_{0}^{\pi/4} \frac{1}{\cos^2 x + 25\sin^2 x} \, dx
\]
\begin{solution}
使用代换 \(u = \tan x\),则
\[
du = \sec^2 x \, dx \implies dx = \frac{du}{\sec^2 x} = \frac{du}{1+u^2}
\]

积分上下限:
\[
x=0 \implies u=0, \quad x=\frac{\pi}{4} \implies u=1
\]

原积分变为:
\[
\int_{0}^{\pi/4} \frac{1}{\cos^2 x + 25\sin^2 x} \, dx
= \int_{0}^{1} \frac{1}{1 + 25u^2} \, du
= \frac{1}{25} \int_{0}^{1} \frac{du}{(\frac{1}{5})^2 + u^2}
\]

标准反正切积分:
\[
\frac{1}{25} \int_{0}^{1} \frac{du}{(\frac{1}{5})^2 + u^2} = \frac{1}{25} \cdot 5 \left[ \arctan(5u) \right]_{0}^{1} = \frac{1}{5} \arctan 5
\]
\end{solution}

\question
使用给定代换计算积分:
\[
\int 2 \frac{1-x^2}{1+x^2} \, dx
\]

\begin{solution}
代换:
\[
x = \tan\frac{\theta}{2} \implies dx = \frac{1}{2} \sec^2 \frac{\theta}{2} \, d\theta
\]

观察右边:
\[
\frac{1-x^2}{1+x^2} = \frac{1-\tan^2\frac{\theta}{2}}{1+\tan^2\frac{\theta}{2}} 
= \frac{1-\tan^2\frac{\theta}{2}}{\sec^2\frac{\theta}{2}} 
= \cos^2 \frac{\theta}{2} - \sin^2 \frac{\theta}{2} 
= \cos\theta
\]

因此积分变为:
\[
\int 2 \frac{1-x^2}{1+x^2} \, dx = \int 2 \cos\theta \cdot \frac{1}{2} \sec^2 \frac{\theta}{2} \, d\theta
= \int \cos\theta \sec^2 \frac{\theta}{2} \, d\theta
\]

注意到:
\[
\arctan x = \frac{\theta}{2} \implies \theta = 2 \arctan x
\]

通过观察和代回:
\[
\int 2 \frac{1-x^2}{1+x^2} \, dx = (1+x^2) \arctan x - x + C
\]
\end{solution}

\question 计算
\[
\int_{0}^{\pi/3} \frac{6\sqrt{3}\cos x}{4 + \sin(2x)\tan(x/2)} \, dx
\]

\begin{solution}
首先使用半角替换:
\[
t = \tan\frac{x}{2} \implies dx = \frac{2}{1+t^2} \, dt, \quad x=0 \to t=0, \quad x=\frac{\pi}{3} \to t=\frac{\sqrt{3}}{3}
\]
并用标准半角公式:
\[
\sin x = \frac{2t}{1+t^2}, \quad \cos x = \frac{1-t^2}{1+t^2}
\]

代入积分:
\[
\int_{0}^{\sqrt{3}/3} \frac{6\sqrt{3}\frac{1-t^2}{1+t^2}}{4+2\frac{2t}{1+t^2}\frac{1-t^2}{1+t^2} t} \cdot \frac{2}{1+t^2} \, dt
= \int_{0}^{\sqrt{3}/3} \frac{12\sqrt{3}(1-t^2)}{4[(1+t^2)^2+t^2(1-t^2)]} \, dt
\]

分母化简:
\[
4[(1+t^2)^2+t^2(1-t^2)] = 4(1 + 3t^2)
\]

因此积分化为:
\[
\int_{0}^{\sqrt{3}/3} \frac{12\sqrt{3}(1-t^2)}{4(3t^2+1)} \, dt = \int_{0}^{\sqrt{3}/3} \frac{3\sqrt{3}(1-t^2)}{3t^2+1} \, dt
\]

拆分分子:
\[
\frac{3\sqrt{3}(1-t^2)}{3t^2+1} = \sqrt{3}\frac{3t^2+1 - 4t^2}{3t^2+1} = \sqrt{3}\left(1 - \frac{4t^2}{3t^2+1}\right)
\]

积分:
\[
\sqrt{3}\int_{0}^{\sqrt{3}/3} 1 - \frac{4t^2}{3t^2+1} \, dt = \sqrt{3}\int_{0}^{\sqrt{3}/3} 1 - \frac{4/3}{t^2+1/3} \, dt
= \sqrt{3} \left[ t - \frac{4\sqrt{3}}{3} \arctan(t\sqrt{3}) \right]_0^{\sqrt{3}/3}
\]

代入上下限:
\[
\sqrt{3}\left[ \frac{\sqrt{3}}{3} - \frac{4\sqrt{3}}{3}\arctan(1) - (0-0) \right] = \sqrt{3} \left[ \frac{\sqrt{3}}{3} - \frac{4\sqrt{3}}{3}\cdot \frac{\pi}{4} \right] = 1 - \pi +1
\]

最终结果:
\[
\int_{0}^{\pi/3} \frac{6\sqrt{3}\cos x}{4 + \sin(2x)\tan(x/2)} \, dx = \pi - 1
\]
\end{solution}

\question 求积分
\begin{solution}
使用 Weierstrass 代换 \(t = \tan\frac{x}{2}\),则
\[
dx = \frac{2 \, dt}{1+t^2}, \quad \cos x = \frac{1-t^2}{1+t^2}, \quad \sin x = \frac{2t}{1+t^2}
\]

积分上下限:
\[
x=0 \implies t=0, \quad x=\pi/2 \implies t=1
\]

原积分变为:
\[
\int_{0}^{1} \sqrt{\frac{1-\frac{1-t^2}{1+t^2}}{1+\frac{2t}{1+t^2}}} \cdot \frac{2 \, dt}{1+t^2} 
= \int_{0}^{1} \frac{\sqrt{2} \, t}{(t+1)^2} \cdot \frac{2 \, dt}{1+t^2} 
= \sqrt{2} \int_{0}^{1} \frac{2t}{(t+1)^2(1+t^2)} \, dt
\]

利用部分分式分解:
\[
\frac{2t}{(t+1)^2(1+t^2)} = -\frac{1}{t+1} + \frac{t+1}{t^2+1} = -\frac{1}{t+1} + \frac{t}{t^2+1} + \frac{1}{t^2+1}
\]

积分得到:
\[
\sqrt{2} \int_{0}^{1} \left(-\frac{1}{t+1} + \frac{t}{t^2+1} + \frac{1}{t^2+1}\right) dt
= \sqrt{2} \left[ -\ln(t+1) + \frac{1}{2}\ln(t^2+1) + \arctan t \right]_{0}^{1}
\]

计算边界值:
\[
= \sqrt{2} \left[ -\ln 2 + \frac{1}{2} \ln 2 + \frac{\pi}{4} \right] = \sqrt{2} \left[ \frac{\pi}{4} - \frac{1}{2}\ln 2 \right] = \frac{\sqrt{2}}{4} (\pi - 2\ln 2)
\]
\end{solution}

\question
\[
\int_{0}^{\frac{\pi}{4}} \frac{\sin x+\cos x}{9+16\sin 2x} \, dx
\]

\begin{solution}
令
\[
u=\sin x-\cos x
\]

则
\[
\frac{du}{dx}=\cos x+\sin x
\]
从而
\[
dx=\frac{du}{\cos x+\sin x}
\]

当 $x=0$ 时,
\[
u=-1
\]
当 $x=\frac{\pi}{4}$ 时,
\[
u=0
\]

又有
\[
u^2=(\sin x-\cos x)^2
\]
\[
= \sin^2 x+\cos^2 x-2\sin x\cos x
\]
\[
=1-\sin 2x
\]

因此
\[
\sin 2x=1-u^2
\]
\[
9+16\sin 2x=9+16(1-u^2)=25-16u^2
\]

原积分化为
\[
\int_{-1}^{0} \frac{\sin x+\cos x}{25-16u^2}\frac{du}{\cos x+\sin x}
= \int_{-1}^{0} \frac{1}{25-16u^2} \, du
\]

因式分解得
\[
25-16u^2=(5-4u)(5+4u)
\]

用部分分式法
\[
\int_{-1}^{0} \frac{1}{(5-4u)(5+4u)} \, du
= \int_{-1}^{0} \left(\frac{1}{10}\frac{1}{5+4u}+\frac{1}{10}\frac{1}{5-4u}\right) du
\]

于是
\[
= \frac{1}{10} \int_{-1}^{0} \left(\frac{1}{5+4u}+\frac{1}{5-4u}\right) du
\]

积分得
\[
= \frac{1}{10} \left[\frac{\ln|5+4u|}{4}-\frac{\ln|5-4u|}{4}\right]_{-1}^{0}
\]
\[
= \frac{1}{40} \left[\ln|5+4u|-\ln|5-4u|\right]_{-1}^{0}
\]

代入上下限
\[
= \frac{1}{40}\left[(\ln 5-\ln 5)-(\ln 1-\ln 9)\right]
\]
\[
= \frac{1}{40}\ln 9
= \frac{1}{20}\ln 3
\]
\end{solution}


\question 求积分
\begin{solution}
(a) 代换 \(u = \tan 2x\),则
\[
\frac{du}{dx} = 2\sec^2 2x \implies dx = \frac{du}{2\sec^2 2x}, \quad x=0 \to u=0, \ x=\frac{\pi}{8} \to u=1
\]

原积分变为:
\[
\int_{0}^{\pi/8} \frac{\sqrt{3}}{2+\sin 4x} \, dx = \int_{0}^{1} \frac{\sqrt{3}}{2+2\sin 2x\cos 2x} \cdot \frac{du}{2\sec^2 2x} = \int_{0}^{1} \frac{\sqrt{3}}{4\sec^2 2x + 4\tan 2x} \, du
\]

利用 \(\sec^2 2x = 1 + \tan^2 2x\) 以及 \(u = \tan 2x\):
\[
\int_{0}^{1} \frac{\sqrt{3}}{4u^2 + 4u + 4} \, du = \frac{\sqrt{3}}{4} \int_{0}^{1} \frac{du}{(u+1/2)^2 + 3/4}
\]

(b) 进一步代换 \(V = 2u+1\),则 \(dV = 2 du\),\(u=0 \to V=1\),\(u=1 \to V=3\):
\[
\frac{\sqrt{3}}{4} \int_{0}^{1} \frac{du}{(u+1/2)^2 + 3/4} = \frac{\sqrt{3}}{2} \int_{1}^{3} \frac{dV}{V^2 + (\sqrt{3})^2}
\]

标准反正切积分公式:
\[
\frac{\sqrt{3}}{2} \int_{1}^{3} \frac{dV}{V^2+3} = \frac{\sqrt{3}}{2} \cdot \frac{1}{\sqrt{3}} \left[ \arctan \frac{V}{\sqrt{3}} \right]_{1}^{3} = \frac{1}{2} \left[ \arctan \sqrt{3} - \arctan \frac{\sqrt{3}}{3} \right]
\]

\[
= \frac{1}{2} \left( \frac{\pi}{3} - \frac{\pi}{6} \right) = \frac{\pi}{12}
\]
\end{solution}

\question
36) 计算不定积分:\[ \int \frac{1}{1 + \sin^2 x} dx \]

\begin{solution}
分子分母同时除以 $\cos^2 x$:
\[ I = \int \frac{\sec^2 x}{\sec^2 x + \tan^2 x} dx = \int \frac{\sec^2 x}{1 + 2\tan^2 x} dx \]
设 $u = \sqrt{2} \tan x$,则 $du = \sqrt{2} \sec^2 x dx$。
代入积分式:
\[ I = \frac{1}{\sqrt{2}} \int \frac{1}{u^2+1} du = \frac{1}{\sqrt{2}} \tan^{-1} u + C \]
还原变量:
\[ I = \frac{1}{\sqrt{2}} \tan^{-1}(\sqrt{2} \tan x) + C \]
\end{solution}

\question 求积分
\[
\int_{\pi/6}^{2\pi/9} \frac{1}{2 + \cos 3x} \, dx
\]
\begin{solution}
使用代换 \(t = \tan\frac{3x}{2} \implies dx = \frac{2}{3(1+t^2)} \, dt\):
\[
x = \pi/6 \to t = 1, \quad x = 2\pi/9 \to t = \sqrt{3}
\]

利用三角恒等式 \(\cos 3x = \cos^2 \frac{3x}{2} - \sin^2 \frac{3x}{2} = \frac{1-t^2}{1+t^2}\),积分变为
\[
\int_{1}^{\sqrt{3}} \frac{1}{2 + \frac{1-t^2}{1+t^2}} \cdot \frac{2}{3(1+t^2)} \, dt
= \int_{1}^{\sqrt{3}} \frac{1+t^2}{3+t^2} \cdot \frac{2}{3(1+t^2)} \, dt
= \frac{2}{3} \int_{1}^{\sqrt{3}} \frac{1}{t^2+3} \, dt
\]

这是标准反正切积分:
\[
\frac{2}{3} \int_{1}^{\sqrt{3}} \frac{1}{t^2+(\sqrt{3})^2} \, dt
= \frac{2}{3\sqrt{3}} \left[ \arctan\left(\frac{t}{\sqrt{3}}\right) \right]_{1}^{\sqrt{3}}
= \frac{2}{3\sqrt{3}} \left[ \arctan(1) - \arctan\left(\frac{1}{\sqrt{3}}\right) \right]
\]

计算反正切值:
\[
\arctan(1) = \frac{\pi}{4}, \quad \arctan\left(\frac{1}{\sqrt{3}}\right) = \frac{\pi}{6}
\]

因此积分结果为
\[
\frac{2}{3\sqrt{3}} \cdot \frac{\pi}{12} = \frac{\pi\sqrt{3}}{54}.
\]
\end{solution}

\question
\[
\int \frac{\sec^2 x}{\sqrt{\sec x+\tan x}} \, dx
\]

\begin{solution}
注意到
\[
\frac{d}{dx}(\sec x+\tan x)=\sec^2 x+\sec x\tan x
\]

先将被积式拆分
\[
\int \frac{\sec^2 x}{\sqrt{\sec x+\tan x}} \, dx
= \int \frac{\sec^2 x+\sec x\tan x-\sec x\tan x}{\sqrt{\sec x+\tan x}} \, dx
\]

写成两部分
\[
= \frac{1}{2} \int \frac{\sec^2 x+\sec x\tan x}{\sqrt{\sec x+\tan x}} \, dx
+ \frac{1}{2} \int \frac{\sec^2 x-\sec x\tan x}{\sqrt{\sec x+\tan x}} \, dx
\]

处理第二个积分,利用恒等式
\[
1+\tan^2 x=\sec^2 x
\]

有
\[
\sec^2 x-\sec x\tan x
= \sec x(\sec x-\tan x)
= \sec x(\sec x-\tan x)\frac{\sec x+\tan x}{\sec x+\tan x}
\]

于是
\[
= \frac{1}{2} \int \frac{\sec^2 x+\sec x\tan x}{\sqrt{\sec x+\tan x}} \, dx
+ \frac{1}{2} \int \frac{\sec x(\sec^2 x-\tan^2 x)}{(\sec x+\tan x)^{3/2}} \, dx
\]

而
\[
\sec^2 x-\tan^2 x=1
\]

故
\[
= \frac{1}{2} \int \frac{\sec^2 x+\sec x\tan x}{\sqrt{\sec x+\tan x}} \, dx
+ \frac{1}{2} \int \frac{\sec x}{(\sec x+\tan x)^{3/2}} \, dx
\]

将第二项写成同一结构
\[
= \frac{1}{2} \int (\sec^2 x+\sec x\tan x)(\sec x+\tan x)^{-1/2} \, dx
+ \frac{1}{2} \int (\sec^2 x+\sec x\tan x)(\sec x+\tan x)^{-5/2} \, dx
\]

对两项分别作代换
\[
u=\sec x+\tan x
\]

得到
\[
= \frac{1}{2} \cdot \frac{u^{1/2}}{1/2}
+ \frac{1}{2} \cdot \frac{u^{-3/2}}{-3/2}
+ C
\]

化简并代回
\[
= (\sec x+\tan x)^{1/2}
- \frac{1}{3}(\sec x+\tan x)^{-3/2}
+ C
\]
\end{solution}

\question
计算积分
\[
I=\int_{0}^{\frac{\pi}{4}}4\sin x\sqrt{\cos 2x}\,dx
\]

\begin{solution}
先用三角恒等式
\[
\cos 2x=2\cos^2 x-1
\]
则
\[
I=\int_{0}^{\frac{\pi}{4}}4\sin x\sqrt{2\cos^2 x-1}\,dx
\]

作代换
\[
u=\cos x
\]
则
\[
\frac{du}{dx}=-\sin x,\quad dx=\frac{du}{-\sin x}
\]

当 $x=0$ 时,$u=1$  
当 $x=\frac{\pi}{4}$ 时,$u=\frac{\sqrt{2}}{2}$  

于是
\[
I=\int_{1}^{\frac{\sqrt{2}}{2}}4\sqrt{2u^2-1}(-du)
=\int_{\frac{\sqrt{2}}{2}}^{1}4\sqrt{2u^2-1}\,du
\]

再作三角代换
\[
\sec\theta=\sqrt{2}u
\]
则
\[
du=\frac{\sec\theta\tan\theta}{\sqrt{2}}\,d\theta
\]
且
\[
\sqrt{2u^2-1}=\sqrt{\sec^2\theta-1}=\tan\theta
\]

当 $u=\frac{\sqrt{2}}{2}$ 时,$\theta=0$  
当 $u=1$ 时,$\theta=\frac{\pi}{4}$  

因此
\[
I=\int_{0}^{\frac{\pi}{4}}\frac{4}{\sqrt{2}}\tan\theta(\sec\theta\tan\theta)\,d\theta
=\int_{0}^{\frac{\pi}{4}}\frac{4}{\sqrt{2}}\tan^2\theta\sec\theta\,d\theta
\]

对该积分分部积分,
\[
I=\left[\frac{4}{\sqrt{2}}\tan\theta\sec\theta\right]_{0}^{\frac{\pi}{4}}
-\frac{4}{\sqrt{2}}\int_{0}^{\frac{\pi}{4}}\sec^3\theta\,d\theta
\]

于是
\[
I=4-\frac{4}{\sqrt{2}}\int_{0}^{\frac{\pi}{4}}\sec\theta(1+\tan^2\theta)\,d\theta
\]
\[
=4-\frac{4}{\sqrt{2}}\int_{0}^{\frac{\pi}{4}}\sec\theta\,d\theta
-\frac{4}{\sqrt{2}}\int_{0}^{\frac{\pi}{4}}\sec\theta\tan^2\theta\,d\theta
\]

注意到后一个积分正好等于 $I$,于是
\[
I=4-\frac{4}{\sqrt{2}}\left[\ln(\sec\theta+\tan\theta)\right]_{0}^{\frac{\pi}{4}}-I
\]

从而
\[
2I=4-\frac{4}{\sqrt{2}}\ln(\sqrt{2}+1)
\]

解得
\[
I=2-\sqrt{2}\ln(\sqrt{2}+1)
\]
\end{solution}

\question
设积分
\[
I=\int_{0}^{\pi}\frac{\sin^2 x}{1+\cos^2 x}\,dx
\]

(a) 证明
\[
I+\pi=\int_{0}^{\frac{\pi}{2}}\frac{4}{1+\cos^2 x}\,dx
\]

(b) 从而求 $I$ 的精确化简值  

(c) 用另一种方法验证(b)的结果,先将被积函数写成 $\cot^2 x$ 的函数  

\begin{solution}
(a) 由
\[
I=\int_{0}^{\pi}\frac{\sin^2 x}{1+\cos^2 x}\,dx
\]
有
\[
I=\int_{0}^{\pi}\frac{\sin^2 x+\cos^2 x}{1+\cos^2 x}\,dx
\]
即
\[
I=\int_{0}^{\pi}\frac{1}{1+\cos^2 x}\,dx
\]

又
\[
\int_{0}^{\pi}1\,dx=\pi
\]
于是
\[
I+\pi=\int_{0}^{\pi}\frac{\sin^2 x}{1+\cos^2 x}\,dx+\int_{0}^{\pi}1\,dx
\]
\[
=2\int_{0}^{\pi}\frac{1}{1+\cos^2 x}\,dx
\]

注意到被积函数关于 $\frac{\pi}{2}$ 对称,
\[
I+\pi=4\int_{0}^{\frac{\pi}{2}}\frac{1}{1+\cos^2 x}\,dx
\]
即
\[
I+\pi=\int_{0}^{\frac{\pi}{2}}\frac{4}{1+\cos^2 x}\,dx
\]

(b) 将分子分母同乘 $\sec^2 x$,
\[
I+\pi=\int_{0}^{\frac{\pi}{2}}\frac{4\sec^2 x}{\sec^2 x+1}\,dx
\]
\[
=\int_{0}^{\frac{\pi}{2}}\frac{4\sec^2 x}{2+\tan^2 x}\,dx
\]

写成反正切形式,
\[
I+\pi=\int_{0}^{\frac{\pi}{2}}\frac{4\sec^2 x}{(\sqrt{2})^2+(\tan x)^2}\,dx
\]
\[
=4\left[\frac{1}{\sqrt{2}}\arctan\left(\frac{\tan x}{\sqrt{2}}\right)\right]_{0}^{\frac{\pi}{2}}
\]

因此
\[
I+\pi=2\sqrt{2}\left[\arctan\left(\frac{\tan(\frac{\pi}{2})}{\sqrt{2}}\right)-\arctan\left(\frac{\tan 0}{\sqrt{2}}\right)\right]
\]
\[
=2\sqrt{2}\cdot\frac{\pi}{2}=\pi\sqrt{2}
\]

故
\[
I=\pi\sqrt{2}-\pi=\pi(\sqrt{2}-1)
\]

(c) 将被积函数化为 $\cot^2 x$ 的形式,
\[
\int_{0}^{\pi}\frac{\sin^2 x}{1+\cos^2 x}\,dx
=\int_{0}^{\pi}\frac{1}{\csc^2 x+\cot^2 x}\,dx
=\int_{0}^{\pi}\frac{1}{1+2\cot^2 x}\,dx
\]

令
\[
u=\cot x
\]
则
\[
du=-\csc^2 x\,dx
\]
\[
dx=-\frac{du}{1+u^2}
\]

当 $x=0$ 时,$u=+\infty$；当 $x=\pi$ 时,$u=-\infty$

积分化为
\[
\int_{+\infty}^{-\infty}\frac{1}{1+2u^2}\left(-\frac{du}{1+u^2}\right)
=\int_{-\infty}^{\infty}\frac{1}{(1+2u^2)(1+u^2)}\,du
\]

作部分分式分解,
\[
=\int_{-\infty}^{\infty}\left(\frac{2}{1+2u^2}-\frac{1}{1+u^2}\right)du
\]

被积函数为偶函数,
\[
=2\int_{0}^{\infty}\left(\frac{1}{u^2+\left(\frac{1}{\sqrt{2}}\right)^2}-\frac{1}{1+u^2}\right)du
\]

积分得
\[
=2\left[\sqrt{2}\arctan(\sqrt{2}u)-\arctan u\right]_{0}^{\infty}
\]
\[
=2\left(\sqrt{2}\cdot\frac{\pi}{2}-\frac{\pi}{2}\right)
=\pi\sqrt{2}-\pi
\]

即
\[
I=\pi(\sqrt{2}-1)
\]
与(b)结果一致
\end{solution}

\question 证明三倍角公式并计算积分
\[
\int_{0}^{2-\sqrt{3}} \frac{6x(3-x^2)}{1-2x^2-3x^4} \, dx
\]
\begin{solution}
(a) 从 \(\tan 3\theta = \tan(2\theta + \theta)\) 开始:
\[
\tan 3\theta = \frac{\tan 2\theta + \tan \theta}{1 - \tan 2\theta \tan \theta} 
= \frac{\frac{2 \tan \theta}{1 - \tan^2 \theta} + \tan \theta}{1 - \frac{2 \tan \theta}{1 - \tan^2 \theta} \tan \theta}
\]

将分子分母同时乘以 \(1 - \tan^2 \theta\):
\[
\tan 3\theta = \frac{2 \tan \theta + \tan \theta (1 - \tan^2 \theta)}{1 - \tan^2 \theta - 2\tan^2 \theta} 
= \frac{3 \tan \theta - \tan^3 \theta}{1 - 3 \tan^2 \theta}
\]

(b) 考虑积分:
\[
\int_{0}^{2-\sqrt{3}} \frac{6x(3-x^2)}{1-2x^2-3x^4} \, dx
= \int_{0}^{2-\sqrt{3}} \frac{6x(3-x^2)}{-(3x^4+2x^2-1)} \, dx
= \int_{0}^{2-\sqrt{3}} \frac{6x(3-x^2)}{(1-3x^2)(x^2+1)} \, dx
\]

令 \(x = \tan \theta \implies dx = \sec^2 \theta \, d\theta, \quad x=0 \to \theta=0, \quad x=2-\sqrt{3} \to \theta=\frac{\pi}{12}\):
\[
\int_{0}^{2-\sqrt{3}} \frac{6x(3-x^2)}{(1-3x^2)(x^2+1)} \, dx
= \int_{0}^{\pi/12} \frac{6 \tan 3\theta}{\tan^2 \theta + 1} \sec^2 \theta \, d\theta
= \int_{0}^{\pi/12} 6 \tan 3\theta \, d\theta
\]

积分结果:
\[
\int 6 \tan 3\theta \, d\theta = 2 \ln |\sec 3\theta|
\]

代入上下限:
\[
[2 \ln |\sec 3\theta|]_{0}^{\pi/12} = 2 \ln \sec(\pi/4) - 2 \ln \sec 0 = 2 \ln \sqrt{2} - 0 = \ln 2
\]

因此:
\[
\int_{0}^{2-\sqrt{3}} \frac{6x(3-x^2)}{1-2x^2-3x^4} \, dx = \ln 2
\]
\end{solution}

\question 求积分
\[
\int \sec x \, dx
\]
\begin{solution}
(a) 验证恒等式:
\[
\sec x = \frac{1+\tan^2\left(\frac{x}{2}\right)}{1-\tan^2\left(\frac{x}{2}\right)}
\]

由半角公式:
\[
\sin x = \frac{2\tan\frac{x}{2}}{1+\tan^2\frac{x}{2}}, \quad 
\cos x = \frac{1-\tan^2\frac{x}{2}}{1+\tan^2\frac{x}{2}}
\]

则
\[
\sec x = \frac{1}{\cos x} = \frac{1+\tan^2\frac{x}{2}}{1-\tan^2\frac{x}{2}}
\]

(b) 部分分式分解:
\[
\frac{2}{1-t^2} = \frac{2}{(1-t)(1+t)} = \frac{1}{1-t} + \frac{1}{1+t}
\]

(c) 使用代换 \(t = \tan\frac{x}{2} \implies dx = \frac{2}{1+t^2} dt\):
\[
\int \sec x \, dx = \int \frac{1+t^2}{1-t^2} \cdot \frac{2}{1+t^2} \, dt
= \int \frac{2}{1-t^2} \, dt
= \int \frac{1}{1+t} + \frac{1}{1-t} \, dt
= \ln|1+t| - \ln|1-t| + C
= \ln\left|\frac{1+t}{1-t}\right| + C
\]

注意到 \(\frac{1+t}{1-t} = \frac{\tan\frac{\pi}{4} + \tan\frac{x}{2}}{1 - \tan\frac{\pi}{4}\tan\frac{x}{2}} = \tan\left(\frac{\pi}{4} + \frac{x}{2}\right)\),
于是最终结果为:
\[
\int \sec x \, dx = \ln\left|\tan\left(\frac{x}{2} + \frac{\pi}{4}\right)\right| + C
\]
\end{solution}

    \question
求
\[
\int_{\frac{\pi}{6}}^{\frac{\pi}{3}} \frac{\cos x + \sin x}{\sqrt{\sin 2x}} \, dx
\]

\begin{solution}
首先将被积式化简:
\[
\int_{\frac{\pi}{6}}^{\frac{\pi}{3}} \frac{\cos x + \sin x}{\sqrt{\sin 2x}} \, dx = \int_{\frac{\pi}{6}}^{\frac{\pi}{3}} \frac{\cos x + \sin x}{\sqrt{2 \sin x \cos x}} \, dx
\]

使用代换:
\[
u = \sin x - \cos x, \quad du = (\cos x + \sin x) \, dx, \quad dx = \frac{du}{\cos x + \sin x}
\]

积分上下限:
\[
x = \frac{\pi}{6} \to u = \frac{1}{2} - \frac{\sqrt{3}}{2} = -\alpha, \quad
x = \frac{\pi}{3} \to u = \frac{\sqrt{3}}{2} - \frac{1}{2} = \alpha
\]

关系式:
\[
u^2 = (\sin x - \cos x)^2 = \sin^2 x + \cos^2 x - 2 \sin x \cos x = 1 - 2 \sin x \cos x
\]
\[
\Rightarrow 2 \sin x \cos x = 1 - u^2
\]

代入积分:

\[
\int_{\frac{\pi}{6}}^{\frac{\pi}{3}} \frac{\cos x + \sin x}{\sqrt{2 \sin x \cos x}} \, dx = \int_{-\alpha}^{\alpha} \frac{\cos x + \sin x}{\sqrt{1-u^2}} \cdot \frac{du}{\cos x + \sin x} = \int_{-\alpha}^{\alpha} \frac{1}{\sqrt{1-u^2}} \, du
\]

由于被积函数是偶函数:
\[
\int_{-\alpha}^{\alpha} \frac{1}{\sqrt{1-u^2}} \, du = 2 \int_{0}^{\alpha} \frac{1}{\sqrt{1-u^2}} \, du = 2 [\arcsin u]_{0}^{\alpha} = 2 \arcsin \alpha
\]

代入 \(\alpha = \frac{\sqrt{3}}{2}-\frac{1}{2}\):

\[
\int_{\frac{\pi}{6}}^{\frac{\pi}{3}} \frac{\cos x + \sin x}{\sqrt{\sin 2x}} \, dx = 2 \arcsin \left( \frac{\sqrt{3}-1}{2} \right)
\]
\end{solution}

\question
27) 计算不定积分:\[ \int \frac{\sin^3 \frac{x}{2}}{\cos \frac{x}{2} \sqrt{\cos x + \cos^2 x + \cos^3 x}} dx \]

\begin{solution}
利用三角恒等式简化分子,并设 $t = \cos x, dt = -\sin x dx$:
\[ I = \frac{1}{2} \int \frac{\sin x (1 - \cos x)}{(1 + \cos x) \sqrt{\cos x + \cos^2 x + \cos^3 x}} dx \]
代入变量 $t$:
\[ = \frac{1}{2} \int \frac{t^2 - 1}{(t+1)^2 \sqrt{t + t^2 + t^3}} dt = \frac{1}{2} \int \frac{1 - \frac{1}{t^2}}{(t+2+\frac{1}{t}) \sqrt{t+1+\frac{1}{t}}} dt \]
设 $u^2 = t + 1 + \frac{1}{t}$,则 $2u du = (1 - \frac{1}{t^2}) dt$:
\[ I = \int \frac{du}{u^2+1} = \tan^{-1} u + C \]
还原变量:
\[ I = \tan^{-1} \sqrt{\cos x + \sec x + 1} + C \]
\end{solution}

\question 计算积分
\[
\int_{\frac{1}{2}}^{\frac{\sqrt{3}}{2}} \frac{1 - \sqrt{\arcsin x}}{\sqrt{1-x^2}\sqrt{\arcsin x}} \, dx
\]
\begin{solution}
使用代换
\[
u = \sqrt{\arcsin x} \implies u^2 = \arcsin x, \quad 2u\,du = \frac{1}{\sqrt{1-x^2}}\,dx, \quad dx = 2u \sqrt{1-x^2} \, du
\]

积分上下限变为:
\[
x = \frac{1}{2} \implies u = \sqrt{\arcsin \frac{1}{2}} = \sqrt{\frac{\pi}{6}}, \quad
x = \frac{\sqrt{3}}{2} \implies u = \sqrt{\arcsin \frac{\sqrt{3}}{2}} = \sqrt{\frac{\pi}{3}}
\]

原积分变为:
\[
\int_{\sqrt{\pi/6}}^{\sqrt{\pi/3}} \frac{1 - u^2}{u} \cdot 2u \, du = \int_{\sqrt{\pi/6}}^{\sqrt{\pi/3}} 2(1-u) \, du = \int_{\sqrt{\pi/6}}^{\sqrt{\pi/3}} (2 - 2u) \, du
\]

积分得到:
\[
\left[ 2u - u^2 \right]_{\sqrt{\pi/6}}^{\sqrt{\pi/3}} = \left(2\sqrt{\frac{\pi}{3}} - \frac{\pi}{3}\right) - \left(2\sqrt{\frac{\pi}{6}} - \frac{\pi}{6}\right)
\]

化简:
\[
= 2\sqrt{\frac{\pi}{3}} - 2\sqrt{\frac{\pi}{6}} - \frac{\pi}{3} + \frac{\pi}{6} = 2\sqrt{\frac{\pi}{3}} - 2\sqrt{\frac{\pi}{6}} - \frac{\pi}{6}
\]

因此积分结果为:
\[
\int_{\frac{1}{2}}^{\frac{\sqrt{3}}{2}} \frac{1 - \sqrt{\arcsin x}}{\sqrt{1-x^2}\sqrt{\arcsin x}} \, dx
= 2\sqrt{\frac{\pi}{3}} - 2\sqrt{\frac{\pi}{6}} - \frac{\pi}{6}.
\]
\end{solution}

\question
利用适当的积分方法证明
\[
\int_{0}^{\frac{1}{2}}
\frac{\arcsin\sqrt{x}-\arccos\sqrt{x}}
{\arcsin\sqrt{x}+\arccos\sqrt{x}}
\,dx
=\frac{1}{\pi}-\frac{1}{2}
\]

\begin{solution}
先作代换
\[
y=\sqrt{x}
\]
则
\[
x=y^2,\quad dx=2y\,dy
\]
积分上下限变为
\[
x=0\Rightarrow y=0,\qquad
x=\frac12\Rightarrow y=\frac{\sqrt2}{2}
\]

原积分化为
\begin{align*}
\int_{0}^{\frac{\sqrt2}{2}}
\frac{\arcsin y-\arccos y}{\arcsin y+\arccos y}
\cdot 2y\,dy
\end{align*}

利用恒等式
\[
\arccos y=\frac{\pi}{2}-\arcsin y
\]
得
\[
\frac{\arcsin y-\arccos y}{\arcsin y+\arccos y}
=\frac{2\arcsin y-\frac{\pi}{2}}{\frac{\pi}{2}}
=\frac{4}{\pi}\arcsin y-1
\]

因此
\begin{align*}
&=\int_{0}^{\frac{\sqrt2}{2}}
\left(\frac{4}{\pi}\arcsin y-1\right)2y\,dy \\
&=\frac{8}{\pi}\int_{0}^{\frac{\sqrt2}{2}}y\arcsin y\,dy
-\int_{0}^{\frac{\sqrt2}{2}}2y\,dy \\
&=\frac{8}{\pi}\int_{0}^{\frac{\sqrt2}{2}}y\arcsin y\,dy
-\left[y^2\right]_{0}^{\frac{\sqrt2}{2}} \\
&=\frac{8}{\pi}\int_{0}^{\frac{\sqrt2}{2}}y\arcsin y\,dy
-\frac12
\end{align*}

再作代换
\[
y=\sin\theta
\]
则
\[
dy=\cos\theta\,d\theta
\]
积分上下限变为
\[
y=0\Rightarrow\theta=0,\qquad
y=\frac{\sqrt2}{2}\Rightarrow\theta=\frac{\pi}{4}
\]

于是
\begin{align*}
\int_{0}^{\frac{\sqrt2}{2}}y\arcsin y\,dy
&=\int_{0}^{\frac{\pi}{4}}\theta\sin\theta\cos\theta\,d\theta \\
&=\frac12\int_{0}^{\frac{\pi}{4}}\theta\sin(2\theta)\,d\theta
\end{align*}

代回原式
\[
=\frac{4}{\pi}\int_{0}^{\frac{\pi}{4}}\theta\sin(2\theta)\,d\theta-\frac12
\]

对积分作分部积分
\[
u=\theta,\quad dv=\sin(2\theta)\,d\theta
\]
则
\[
du=d\theta,\quad v=-\frac12\cos(2\theta)
\]

\begin{align*}
\int_{0}^{\frac{\pi}{4}}\theta\sin(2\theta)\,d\theta
&=\left[-\frac12\theta\cos(2\theta)\right]_{0}^{\frac{\pi}{4}}
+\frac12\int_{0}^{\frac{\pi}{4}}\cos(2\theta)\,d\theta \\
&=\frac12\left[\frac12\sin(2\theta)\right]_{0}^{\frac{\pi}{4}} \\
&=\frac14
\end{align*}

因此
\[
\frac{4}{\pi}\cdot\frac14-\frac12
=\frac{1}{\pi}-\frac12
\]

证毕。
\end{solution}

\question
计算 $[\int \frac{3\sin 2x + 4\cos x e^{\sin x} + 8\cos x}{3\sin x + e^{\sin x} + 1} dx]$。

\begin{solution}
首先,利用二倍角公式 $\sin 2x = 2\sin x \cos x$ 展开分子:
\[ = [\int \frac{6\sin x \cos x + 4\cos x e^{\sin x} + 8\cos x}{3\sin x + e^{\sin x} + 1} dx] \]
令 $u = \sin x$,则 $du = \cos x dx$。代入原式得:
\[ = [\int \frac{6u + 4e^u + 8}{3u + e^u + 1} du] \]
提取分子中的常数 2:
\[ = 2 [\int \frac{3u + 2e^u + 4}{3u + e^u + 1} du] \]
将分子拆项,使其包含分母的形式:
\[ = 2 [\int \frac{(3u + e^u + 1) + (3 + e^u)}{3u + e^u + 1} du] \]
\[ = 2 ( [\int \frac{3u + e^u + 1}{3u + e^u + 1} du] + [\int \frac{e^u + 3}{3u + e^u + 1} du] ) \]
\[ = 2 [\int 1 du] + 2 [\int \frac{d(3u + e^u + 1)}{3u + e^u + 1}] \]
进行积分:
\[ = 2u + 2\ln|3u + e^u + 1| + C \]
最后回代 $u = \sin x$:
\[ = 2\sin x + 2\ln|3\sin x + e^{\sin x} + 1| + C \]
\end{solution}

\question 已知 $\int \csc x dx = \ln |\csc x - \cot x| + C$
\begin{solution}
(a) 证明 $\sin x + \sin 3x + \sin 5x = \frac{\sin^2 3x}{\sin x}$:
\begin{align*}
\sin x (\sin x + \sin 3x + \sin 5x) &= \sin x (2\sin \frac{x+5x}{2} \cos \frac{x-5x}{2} + \sin 3x) \\
&= \sin x (2\sin 3x \cos(-2x) + \sin 3x) \\
&= \sin x \sin 3x (2\cos 2x + 1) \\
&= \sin 3x (\sin(x+2x) + \sin(x-2x) + \sin x) \\
&= \sin 3x (\sin 3x - \sin x + \sin x) \\
&= \sin^2 3x
\end{align*}
因此得证:$\sin x + \sin 3x + \sin 5x = \frac{\sin^2 3x}{\sin x}$

(b) 计算 $\int \frac{\sin x + \sin 3x}{\cos x \sqrt{4 - \cos^4 x}} dx$:
设 $u = \cos^2 x$,则 $du = -2 \sin x \cos x dx$。
\begin{align*}
\int \frac{2 \sin \frac{x+3x}{2} \cos \frac{x-3x}{2}}{\cos x \sqrt{4 - \cos^4 x}} dx &= \int \frac{2 \sin 2x \cos(-x)}{\cos x \sqrt{4 - \cos^4 x}} dx \\
&= \int \frac{4 \sin x \cos x \cos x}{\cos x \sqrt{4 - \cos^4 x}} dx \\
&= \int \frac{4 \sin x \cos x}{\sqrt{4 - \cos^4 x}} dx \\
&= \int \frac{-2}{\sqrt{4 - u^2}} du \\
&= -2 \arcsin (\frac{u}{2}) + C \\
&= -2 \arcsin (\frac{\cos^2 x}{2}) + C
\end{align*}

(c) 计算 $\int \frac{16(\sin 2x + \sin 6x + \sin 10x)}{\cos 2x \sqrt{4 - \cos^4 2x}} dx$:
利用 (a) 栏的结果,设 $2x$ 为变量,则分子部分可化简为 $16 \frac{\sin^2 6x}{\sin 2x}$。
设 $u = \cos^2 2x$,则 $du = -4 \sin 2x \cos 2x dx$。
经过代换与积分（参考第二张图内容）:
\begin{align*}
\int (\frac{32u-4}{u\sqrt{4-u^2}} - \frac{64u}{\sqrt{4-u^2}}) du &= 32 \arcsin (\frac{u}{2}) + 2 \ln |\frac{2 + \sqrt{4-u^2}}{u}| + 64 \sqrt{4-u^2} + C
\end{align*}
代回 $u = \cos^2 2x$ 得最终结果:
\[ I = 32 \arcsin (\frac{\cos^2 2x}{2}) + 2 \ln |\frac{2 + \sqrt{4 - \cos^4 2x}}{\cos^2 2x}| + 64 \sqrt{4 - \cos^4 2x} + C \]
\end{solution}

\question
(a) 计算 $[\int \frac{1}{4-3z^2} dz]$。
\begin{solution}
令 $z = \frac{2}{\sqrt{3}} \sin\theta$,则 $\frac{dz}{d\theta} = \frac{2}{\sqrt{3}} \cos\theta \implies dz = \frac{2}{\sqrt{3}} \cos\theta d\theta$。
\[ [\int \frac{1}{4-4\sin^2\theta} (\frac{2}{\sqrt{3}}\cos\theta) d\theta] = [\int \frac{1}{4\cos^2\theta} (\frac{2}{\sqrt{3}}\cos\theta) d\theta] \]
\[ = \frac{1}{2\sqrt{3}} [\int \sec\theta d\theta] = \frac{1}{2\sqrt{3}} \ln |\sec\theta + \tan\theta| + C \]
根据辅助三角形,$\sin\theta = \frac{\sqrt{3}z}{2},\sec\theta = \frac{2}{\sqrt{4-3z^2}},\tan\theta = \frac{\sqrt{3}z}{\sqrt{4-3z^2}}$。
\[ = \frac{1}{2\sqrt{3}} \ln |\frac{2+\sqrt{3}z}{\sqrt{4-3z^2}}| + C = \frac{1}{2\sqrt{3}} \ln |\sqrt{\frac{(2+\sqrt{3}z)^2}{(2-\sqrt{3}z)(2+\sqrt{3}z)}}| + C \]
\[ = \frac{1}{4\sqrt{3}} \ln |\frac{2+\sqrt{3}z}{2-\sqrt{3}z}| + C = \frac{\sqrt{3}}{12} \ln |\frac{2+\sqrt{3}z}{2-\sqrt{3}z}| + C \]
\end{solution}

\question
(b) 计算 $[\int \frac{\cos u}{\tan^2 u + 4} du]$。
\begin{solution}
\[ [\int \frac{\cos u}{\frac{\sin^2 u}{\cos^2 u} + 4} du] = [\int \frac{\cos^3 u}{\sin^2 u + 4\cos^2 u} du] \]
\[ = [\int \frac{\cos^3 u}{\sin^2 u + 4(1-\sin^2 u)} du] = [\int \frac{\cos^3 u}{4-3\sin^2 u} du] \]
令 $t = \sin u$,则 $dt = \cos u du,\cos^2 u = 1 - t^2$。
\[ = [\int \frac{1-t^2}{4-3t^2} dt] = [\int \frac{\frac{1}{3}(4-3t^2) - \frac{1}{3}}{4-3t^2} dt] = \frac{1}{3} ( [\int 1 dt] - [\int \frac{1}{4-3t^2} dt] ) \]
利用(a)的结果:
\[ = \frac{\sin u}{3} - \frac{1}{3} (\frac{\sqrt{3}}{12} \ln |\frac{2+\sqrt{3}\sin u}{2-\sqrt{3}\sin u}|) + C = \frac{\sin u}{3} - \frac{\sqrt{3}}{36} \ln |\frac{2+\sqrt{3}\sin u}{2-\sqrt{3}\sin u}| + C \]
\end{solution}

\question
(c) 计算 $[\int_{-\frac{\pi}{4}}^{\frac{\pi}{4}} \frac{90(\sec x - \cos x)}{(\sec x + 3\cos x)^2} dx]$。
\begin{solution}
\[ = 90 [\int_{-\frac{\pi}{4}}^{\frac{\pi}{4}} \frac{\frac{1}{\cos x} - \cos x}{(\frac{1}{\cos x} + 3\cos x)^2} dx] = 90 [\int_{-\frac{\pi}{4}}^{\frac{\pi}{4}} \frac{\sin^2 x \cos x}{(1+3\cos^2 x)^2} dx] \]
令 $t = \sin x$,则 $dt = \cos x dx$。当 $x = \pm\frac{\pi}{4}$ 时,$t = \pm\frac{\sqrt{2}}{2}$。
\[ = 90 [\int_{-\frac{\sqrt{2}}{2}}^{\frac{\sqrt{2}}{2}} \frac{t^2}{(1+3(1-t^2))^2} dt] = 180 [\int_{0}^{\frac{\sqrt{2}}{2}} \frac{t^2}{(4-3t^2)^2} dt] \]
对不定积分使用换元 $t = \frac{2}{\sqrt{3}}\sin\theta$:
\[ [\int \frac{t^2}{(4-3t^2)^2} dt] = [\int \frac{\frac{4}{3}\sin^2\theta}{16\cos^4\theta} \frac{2}{\sqrt{3}}\cos\theta d\theta] = \frac{1}{6\sqrt{3}} [\int \tan^2\theta \sec\theta d\theta] \]
\[ = \frac{1}{6\sqrt{3}} ([\int \sec^3\theta d\theta] - [\int \sec\theta d\theta]) = \frac{1}{12\sqrt{3}} (\sec\theta \tan\theta - \ln |\sec\theta + \tan\theta|) + C \]
回代 $t$ 并带入上下限:
\[ [\int_{0}^{\frac{\sqrt{2}}{2}} \frac{t^2}{(4-3t^2)^2} dt] = \frac{1}{6\sqrt{3}} [ \frac{\sqrt{3}t}{4-3t^2} - \frac{1}{2} \ln (\frac{2+\sqrt{3}t}{\sqrt{4-3t^2}}) ] \bigg|_{0}^{\frac{\sqrt{2}}{2}} \]
\[ = \frac{1}{6\sqrt{3}} [ \frac{\frac{\sqrt{6}}{2}}{4-\frac{3}{2}} - \frac{1}{2} \ln (\frac{2+\frac{\sqrt{6}}{2}}{\sqrt{4-\frac{3}{2}}}) ] = \frac{1}{15\sqrt{2}} - \frac{1}{12\sqrt{3}} \ln (\frac{2\sqrt{2}+\sqrt{3}}{\sqrt{5}}) \]
乘以系数 180:
\[ 180 [\int_{0}^{\frac{\sqrt{2}}{2}} \frac{t^2}{(4-3t^2)^2} dt] = 6\sqrt{2} - 5\sqrt{3} \ln (\frac{2\sqrt{2}+\sqrt{3}}{\sqrt{5}}) \]
\end{solution}

\question
计算以下积分:
(a) $[\int \frac{2\sin x}{\cos^2 x + 4} dx]$
(b) $[\int \frac{2\sec^2 x + 3\sec x \tan x}{4 + \tan^2 x} dx]$

\begin{solution}
(a) 令 $u = \cos x$,则 $du = -\sin x dx$。
\[ [\int \frac{2\sin x}{\cos^2 x + 4} dx] = -[\int \frac{2}{u^2 + 4} du] \]
进一步令 $u = 2\tan v$,则 $du = 2\sec^2 v dv$。
\[ = -[\int \frac{2}{(2\tan v)^2 + 4} \cdot 2\sec^2 v dv] = -[\int \frac{4\sec^2 v}{4(1 + \tan^2 v)} dv] \]
\[ = -[\int 1 dv] = -v + C \]
由于 $v = \tan^{-1}(\frac{u}{2})$ 且 $u = \cos x$:
\[ = -\tan^{-1}(\frac{\cos x}{2}) + C \]

(b) 将积分拆分为两部分:
\[ [\int \frac{2\sec^2 x + 3\sec x \tan x}{4 + \tan^2 x} dx] = [\int \frac{2\sec^2 x}{4 + \tan^2 x} dx] + [\int \frac{3\sec x \tan x}{4 + \tan^2 x} dx] \]
对于第一部分,令 $v = \tan x$,则 $dv = \sec^2 x dx$。
对于第二部分,令 $u = \sec x$,由于 $1 + \tan^2 x = \sec^2 x$,故 $\tan^2 x = u^2 - 1$。
\[ = [\int \frac{2}{4 + v^2} dv] + [\int \frac{3}{4 + (u^2 - 1)} du] \]
\[ = [\int \frac{2}{v^2 + 4} dv] + [\int \frac{3}{u^2 + 3} du] \]
使用积分公式 $[\int \frac{1}{x^2 + a^2} dx] = \frac{1}{a} \tan^{-1}(\frac{x}{a}) + C$:
第一部分:令 $v = 2\tan \alpha$,得 $\frac{2}{2} \tan^{-1}(\frac{v}{2}) = \tan^{-1}(\frac{\tan x}{2})$。
（注:原稿中此处使用了 $v = \sqrt{5}\tan \alpha$ 的换元,似乎针对的是分母为 $5 + \tan^2 x$ 的情况,若按题目 $4 + \tan^2 x$ 计算如下）:
\[ = \tan^{-1}(\frac{\tan x}{2}) + \frac{3}{\sqrt{3}} \tan^{-1}(\frac{\sec x}{\sqrt{3}}) + C \]
\[ = \tan^{-1}(\frac{\tan x}{2}) + \sqrt{3} \tan^{-1}(\frac{\sec x}{\sqrt{3}}) + C \]
\end{solution}
%分部积分by parts
\question
2) 计算不定积分:\[ \int \frac{x^2}{(1+x^2)^2} dx \]

\begin{solution}
使用分部积分法,令 $u=x, dv = \frac{x}{(1+x^2)^2} dx$:
则 $du = dx, v = -\frac{1}{2(1+x^2)}$。

根据分部积分公式:
\begin{align*}
\int \frac{x^2}{(1+x^2)^2} dx &= -\frac{x}{2(1+x^2)} + \int \frac{1}{2(1+x^2)} dx \\
&= -\frac{x}{2(1+x^2)} + \frac{1}{2} \tan^{-1} x + C
\end{align*}
\end{solution}

\question
3) 计算不定积分:\[ \int x e^x \cos x dx \]

\begin{solution}
首先利用分部积分法处理 $e^x \cos x$。已知:
\[ \int e^x \cos x dx = \frac{1}{2} e^x (\sin x + \cos x) \]
\[ \int e^x \sin x dx = \frac{1}{2} e^x (\sin x - \cos x) \]

对原式使用分部积分,令 $u=x, dv = e^x \cos x dx$:
\begin{align*}
\int x e^x \cos x dx &= \frac{1}{2} x e^x (\sin x + \cos x) - \int \frac{1}{2} e^x (\sin x + \cos x) dx \\
&= \frac{1}{2} x e^x (\sin x + \cos x) - \frac{1}{2} \left[ \int e^x \sin x dx + \int e^x \cos x dx \right] \\
&= \frac{1}{2} x e^x (\sin x + \cos x) - \frac{1}{2} \left[ \frac{1}{2} e^x (\sin x - \cos x) + \frac{1}{2} e^x (\sin x + \cos x) \right] \\
&= \frac{1}{2} x e^x (\sin x + \cos x) - \frac{1}{2} e^x \sin x + C \\
&= \frac{1}{2} e^x (x \sin x + x \cos x - \sin x) + C
\end{align*}
\end{solution}

\question 计算 $[\int_0^1 \frac{\sin^{-1}x}{x} dx]$。
\begin{solution}
令 $x = \sin \theta$,则 $dx = \cos \theta d\theta$。
积分式转化为:
\[I = [\int_0^{\pi/2} \frac{\theta \cos \theta}{\sin \theta} d\theta] = [\int_0^{\pi/2} \theta \cot \theta d\theta]\]
使用分部积分法,令 $u = \theta,dv = \cot \theta d\theta$:
\[I = [\theta \ln(\sin \theta)]_0^{\pi/2} - [\int_0^{\pi/2} \ln(\sin \theta) d\theta]\]
第一项代入上下限后为 $0$。利用标准积分公式 $[\int_0^{\pi/2} \ln(\sin \theta) d\theta] = -\frac{\pi}{2} \ln 2$:
\[I = -(-\frac{\pi}{2} \ln 2) = \frac{\pi}{2} \ln 2\]
\end{solution}
    \question $$\int_0^{\infty} e^{-x} \sin x \, dx$$by parts
    \begin{solution}
        设
        \[
        I = \int_0^{\infty} e^{-x} \sin x \, dx.
        \]
        分部积分两次得
        \[
        I = -e^{-x} \sin x + \int e^{-x} \cos x \, dx
           = -e^{-x} (\sin x + \cos x) - I.
        \]
        移项得
        \[
        2I = -e^{-x} (\sin x + \cos x)
        \Rightarrow I = -\frac{1}{2} e^{-x} (\sin x + \cos x).
        \]
        取定积分,注意 $\lim_{x \to \infty} e^{-x} (\sin x + \cos x) = 0$,代入得:
        \[
        I = \left[ -\frac{1}{2} e^{-x} (\sin x + \cos x) \right]_0^{\infty}
        = 0 - \left(-\frac{1}{2} \cdot 1\right) = \frac{1}{2}.
        \]
    \end{solution}

    \question 
    \[
    \int \frac{xe^x}{(x+1)^2} \,dx
    \]
    \begin{solution}
        换元 $u=x+1,$再分部积分得 
        \begin{align*}
        \int \frac{xe^x}{(x+1)^2} \,dx 
        &= \int \frac{(u - 1)e^{u - 1}}{u^2} \,du \\
        &= \frac{1}{e} \int \frac{(u - 1)e^u}{u^2} \,du \\
        &= \frac{1}{e} \left( \int \frac{e^u}{u} \,du - \int \frac{e^u}{u^2} \,du \right) \\
        &= \frac{1}{e} \left( \frac{e^u}{u} + \int \frac{e^u}{u^2} \,du - \int \frac{e^u}{u^2} \,du \right) \\
        &= \frac{e^{u - 1}}{u} + C\\
        &= \frac{e^x}{x + 1} + C
        \end{align*}
    \end{solution}

\question
23) 计算不定积分:\[ \int \sin^{-1} \sqrt{\frac{x}{1+x}} dx \]

\begin{solution}
使用分部积分法,令 $u = \sin^{-1} \sqrt{\frac{x}{1+x}}, dv = dx$:
\begin{align*}
I &= x \sin^{-1} \sqrt{\frac{x}{1+x}} - \int x \cdot \frac{1}{\sqrt{1 - \frac{x}{1+x}}} \cdot \frac{1}{2\sqrt{\frac{x}{1+x}}} \cdot \frac{1}{(1+x)^2} dx \\
&= x \sin^{-1} \sqrt{\frac{x}{1+x}} - \frac{1}{2} \int x \cdot \sqrt{1+x} \cdot \frac{\sqrt{1+x}}{\sqrt{x}} \cdot \frac{1}{(1+x)^2} dx \\
&= x \sin^{-1} \sqrt{\frac{x}{1+x}} - \frac{1}{2} \int \frac{\sqrt{x}}{1+x} dx
\end{align*}

对剩余积分项进行凑项处理:
\[ \frac{1}{2} \int \frac{x+1-1}{\sqrt{x}(1+x)} dx = \int \frac{1}{2\sqrt{x}} dx - \int \frac{1}{2\sqrt{x}(1+x)} dx \]

进行积分计算:
\begin{itemize}
    \item 第一项:$\int \frac{1}{2\sqrt{x}} dx = \sqrt{x}$。
    \item 第二项:设 $\sqrt{x} = t$,则 $\int \frac{1}{1+t^2} dt = \tan^{-1} t = \tan^{-1}\sqrt{x}$。
\end{itemize}

合并最终结果:
\[ I = x \sin^{-1} \sqrt{\frac{x}{1+x}} - \sqrt{x} + \tan^{-1}\sqrt{x} + C \]
\end{solution}

\question
5) 计算不定积分 $\int x \tan^{-1}x^2 dx$

\begin{solution}
设 $u = x^2, du = 2x dx$:
\begin{align*}
\int x \tan^{-1}x^2 dx &= \frac{1}{2} \int \tan^{-1}u du \\
\intertext{使用分部积分法:}
&= \frac{1}{2} \left( u \tan^{-1}u - \int \frac{u}{1+u^2} du \right) \\
&= \frac{1}{2} u \tan^{-1}u - \frac{1}{4} \ln(1+u^2) + C \\
&= \frac{1}{2} x^2 \tan^{-1}x^2 - \frac{1}{4} \ln(1+x^4) + C
\end{align*}
\end{solution}
\question
15) 计算不定积分:\[ \int \frac{x \sin x}{\cos^5 x} dx \]

\begin{solution}
使用分部积分法,令 $u = x, dv = \frac{\sin x}{\cos^5 x} dx$。
对于 $v = \int \tan x \sec^4 x dx$,设 $\tan x = t$:
\[ v = \int t(1+t^2) dt = \frac{1}{2} \tan^2 x + \frac{1}{4} \tan^4 x = \frac{1}{4} \sec^4 x - \frac{1}{4} \]
代入分部积分公式:
\[ I = \frac{x}{4} \sec^4 x - \int \frac{1}{4} \sec^4 x dx \]
利用 $\int \sec^4 x dx = \tan x + \frac{1}{3} \tan^3 x$:
\[ I = \frac{x \sec^4 x}{4} - \frac{1}{4} \tan x - \frac{1}{12} \tan^3 x + C \]
\end{solution}
\question
9) 计算不定积分:\[ I = \int \cos(3 \ln x) dx \]

\begin{solution}
使用换元法,设 $x = e^y$,则 $dx = e^y dy$。
代入原积分:
\[ I = \int e^y \cos 3y dy \]

对上述积分使用两次分部积分法:
\textbf{第一次分部积分}:
令 $u = \cos 3y, dv = e^y dy$,则 $du = -3 \sin 3y dy, v = e^y$。
\[ I = e^y \cos 3y + 3 \int e^y \sin 3y dy \]

\textbf{第二次分部积分}:
对 $\int e^y \sin 3y dy$ 再次令 $u = \sin 3y, dv = e^y dy$,则 $du = 3 \cos 3y dy, v = e^y$。
\[ \int e^y \sin 3y dy = e^y \sin 3y - 3 \int e^y \cos 3y dy \]

将结果代回 $I$ 的表达式:
\[ I = e^y \cos 3y + 3 \left( e^y \sin 3y - 3I \right) \]
\[ I = e^y \cos 3y + 3 e^y \sin 3y - 9I \]

解关于 $I$ 的方程:
\[ 10I = e^y (\cos 3y + 3 \sin 3y) \]
\[ I = \frac{e^y}{10} (\cos 3y + 3 \sin 3y) + C \]

还原变量 $y = \ln x$:
\[ I = \frac{x}{10} [\cos(3 \ln x) + 3 \sin(3 \ln x)] + C \]

\textbf{结果:}
\[ \int \cos(3 \ln x) dx = \frac{x}{10} [\cos(3 \ln x) + 3 \sin(3 \ln x)] + C \]
\end{solution}

    \question
求下列积分的精确值
\[
\int_{\frac{3}{2}}^{\frac{5}{2}} (4x^2-16x+15)^4\,dx
\]

\begin{solution}
先进行因式分解
\[
4x^2-16x+15=(2x-3)(2x-5)
\]
因此
\begin{align*}
\int_{\frac{3}{2}}^{\frac{5}{2}} (4x^2-16x+15)^4\,dx
&=\int_{\frac{3}{2}}^{\frac{5}{2}}[(2x-3)(2x-5)]^4\,dx \\
&=\int_{\frac{3}{2}}^{\frac{5}{2}}(2x-3)^4(2x-5)^4\,dx
\end{align*}

第一次分部积分
\begin{align*}
&=\left[\frac{1}{10}(2x-3)^5(2x-5)^4\right]_{\frac{3}{2}}^{\frac{5}{2}}
-\int_{\frac{3}{2}}^{\frac{5}{2}}\frac{4}{5}(2x-3)^5(2x-5)^3\,dx
\end{align*}

第二次分部积分
\begin{align*}
&=\left[\frac{4}{55}(2x-3)^5(2x-5)^4\right]_{\frac{3}{2}}^{\frac{5}{2}}
-\left[\frac{2}{5}(2x-3)^6(2x-5)^3\right]_{\frac{3}{2}}^{\frac{5}{2}} \\
&\quad+\frac{6}{5}\int_{\frac{3}{2}}^{\frac{5}{2}}(2x-3)^6(2x-5)^2\,dx
\end{align*}

第三次分部积分
\begin{align*}
&=\frac{2}{5}\left\{
\left[\frac{1}{14}(2x-3)^7(2x-5)^2\right]_{\frac{3}{2}}^{\frac{5}{2}}
-\frac{2}{7}\int_{\frac{3}{2}}^{\frac{5}{2}}(2x-3)^7(2x-5)\,dx
\right\} \\
&=-\frac{4}{35}\int_{\frac{3}{2}}^{\frac{5}{2}}(2x-3)^7(2x-5)\,dx
\end{align*}

最后一次分部积分
\begin{align*}
&=-\frac{4}{35}\left\{
\left[\frac{1}{16}(2x-3)^8(2x-5)\right]_{\frac{3}{2}}^{\frac{5}{2}}
-\frac{1}{8}\int_{\frac{3}{2}}^{\frac{5}{2}}(2x-3)^8\,dx
\right\} \\
&=\frac{1}{70}\int_{\frac{3}{2}}^{\frac{5}{2}}(2x-3)^8\,dx
\end{align*}

直接积分
\begin{align*}
&=\frac{1}{70}\left[\frac{(2x-3)^9}{18}\right]_{\frac{3}{2}}^{\frac{5}{2}} \\
&=\frac{1}{1260}\left[(2x-3)^9\right]_{\frac{3}{2}}^{\frac{5}{2}} \\
&=\frac{1}{1260}\left[(5-3)^9-(3-3)^9\right] \\
&=\frac{512}{1260} \\
&=\frac{128}{315}
\end{align*}
\end{solution}

\question
计算积分
\[
\int_{\pi/4}^{\pi/2} (\cos 2x + \sin 2x)(\ln \cos x + \ln \sin x) \, dx
\]

\begin{solution}
首先利用对数性质:
\[
\ln \cos x + \ln \sin x = \ln (\cos x \sin x) = \ln \frac{\sin 2x}{2} = \ln (\sin 2x) - \ln 2
\]

因此积分可以写为:
\[
\int_{\pi/4}^{\pi/2} (\cos 2x + \sin 2x)(\ln \sin 2x - \ln 2) \, dx
= \int_{\pi/4}^{\pi/2} (\cos 2x + \sin 2x)\ln \sin 2x \, dx - \ln 2 \int_{\pi/4}^{\pi/2} (\cos 2x + \sin 2x) \, dx
\]

第二个积分简单计算:
\[
\int_{\pi/4}^{\pi/2} (\cos 2x + \sin 2x) \, dx = \left[ \frac{1}{2}\sin 2x - \frac{1}{2}\cos 2x \right]_{\pi/4}^{\pi/2} = \frac{1}{2} - \left(\frac{1}{2} - 0\right) = 0
\]

所以问题简化为:
\[
\int_{\pi/4}^{\pi/2} (\cos 2x + \sin 2x) \ln \sin 2x \, dx
\]

使用分部积分,设
\[
u = \ln \sin 2x, \quad dv = (\cos 2x + \sin 2x) dx
\]

得到
\[
du = \frac{2\cos 2x}{\sin 2x} dx, \quad v = \frac{1}{2}(\sin 2x - \cos 2x)
\]

分部积分公式:
\[
\int u\, dv = uv - \int v\, du
\]

代入得:
\[
\int_{\pi/4}^{\pi/2} (\cos 2x + \sin 2x) \ln \sin 2x \, dx
= \left[ \frac{1}{2} (\sin 2x - \cos 2x) \ln \sin 2x \right]_{\pi/4}^{\pi/2} - \frac{1}{2} \int_{\pi/4}^{\pi/2} (\sin 2x - \cos 2x) \frac{2\cos 2x}{\sin 2x} \, dx
\]

化简被积函数:
\[
(\sin 2x - \cos 2x) \frac{2\cos 2x}{\sin 2x} = 2 \cos 2x - 2 \frac{\cos^2 2x}{\sin 2x} = 2 \cos 2x - 2 \frac{1-\sin^2 2x}{\sin 2x} = 2 (\cos 2x + \sin 2x - \csc 2x)
\]

于是积分化为:
\[
\int_{\pi/4}^{\pi/2} (\cos 2x + \sin 2x) \ln \sin 2x \, dx = \left[ \frac{1}{2}(\sin 2x - \cos 2x) \ln \sin 2x - \frac{1}{2}\sin 2x + \frac{1}{2}\cos 2x + \frac{1}{2} \ln \tan x \right]_{\pi/4}^{\pi/2}
\]

在端点代入:
\[
x = \pi/2: \quad \sin 2x = 0, \quad \cos 2x = -1 \implies \text{对数项消失}
\]
\[
x = \pi/4: \quad \sin 2x = 1, \quad \cos 2x = 0 \implies \text{得到 } \frac{1}{2} \ln 2
\]

因此最终结果为:
\[
\int_{\pi/4}^{\pi/2} (\cos 2x + \sin 2x)(\ln \cos x + \ln \sin x) \, dx = \frac{1}{2} \ln 2
\]
\end{solution}

    \question Evaluate $\int \frac{e^x(1+\sin x)}{1+\cos x}\,dx$.

\begin{solution}
将分母化为半角形式:
\begin{align*}
\int \frac{e^x(1+\sin x)}{1+\cos x}\,dx 
&= \int \frac{e^x \bigl(1 + 2\sin\frac{x}{2}\cos\frac{x}{2}\bigr)}{2\cos^2\frac{x}{2}} \, dx \\
&= \int e^x \left( \frac{1}{2\cos^2\frac{x}{2}} + \frac{2\sin\frac{x}{2}\cos\frac{x}{2}}{2\cos^2\frac{x}{2}} \right) \, dx \\
&= \int e^x \left( \frac{1}{2}\sec^2\frac{x}{2} + \tan\frac{x}{2} \right) \, dx.
\end{align*}

注意 $\frac{d}{dx}\tan\frac{x}{2} = \frac{1}{2}\sec^2\frac{x}{2}$,于是利用公式 $\int e^x (f(x)+f'(x))\,dx = e^x f(x) + C$ 得:
\[
\int \frac{e^x(1+\sin x)}{1+\cos x}\,dx = e^x \tan\frac{x}{2} + C.
\]
\end{solution}

\question
计算积分
\[
\int \left(1 + x - \frac{1}{x}\right) e^{x+\frac{1}{x}} \, dx
\]

\begin{solution}
观察被积函数,猜测含指数的项需要积分,尝试寻找一个全微分。

注意
\[
\frac{d}{dx}\left[e^{x+\frac{1}{x}}\right] = \left(1 - \frac{1}{x^2}\right) e^{x+\frac{1}{x}}
\]

将积分重写为:
\[
\int \left(1 + x - \frac{1}{x}\right) e^{x+\frac{1}{x}} \, dx
= \int \left[ e^{x+\frac{1}{x}} + x \left(1 - \frac{1}{x^2}\right) e^{x+\frac{1}{x}} \right] \, dx
\]

拆分积分:
\[
= \int e^{x+\frac{1}{x}} \, dx + \int x \left(1 - \frac{1}{x^2}\right) e^{x+\frac{1}{x}} \, dx
\]

对第二项使用分部积分得到:
\[
\int x \left(1 - \frac{1}{x^2}\right) e^{x+\frac{1}{x}} \, dx
= x e^{x+\frac{1}{x}} - \int e^{x+\frac{1}{x}} \, dx
\]

因此原积分为:
\[
\int \left(1 + x - \frac{1}{x}\right) e^{x+\frac{1}{x}} \, dx
= \int e^{x+\frac{1}{x}} \, dx + \left[x e^{x+\frac{1}{x}} - \int e^{x+\frac{1}{x}} \, dx\right]
= x e^{x+\frac{1}{x}} + C
\]
\end{solution}

\question
证明
\[
\int_{0}^{1} 12x^2 \arctan x \, dx = \pi - 2 + \ln 4
\]

\begin{solution}
\text{步骤 1:分部积分} \\
令 
\[
u = \arctan x, \quad dv = 12x^2 \, dx \implies du = \frac{1}{1+x^2} \, dx, \quad v = 4x^3
\]
利用分部积分公式:
\[
\int u\,dv = uv - \int v\,du
\]
得到:
\[
\int_{0}^{1} 12x^2 \arctan x \, dx = \big[ 4x^3 \arctan x \big]_{0}^{1} - \int_{0}^{1} 4x^3 \cdot \frac{1}{1+x^2} \, dx
\]

\text{步骤 2:计算第一项} 
\[
\big[ 4x^3 \arctan x \big]_{0}^{1} = 4 \cdot \arctan 1 - 0 = 4 \cdot \frac{\pi}{4} = \pi
\]

\text{步骤 3:化简剩余积分} \\
\[
\int_{0}^{1} \frac{4x^3}{1+x^2} \, dx
\] 
令
\[
w = 1+x^2 \implies dw = 2x\,dx, \quad x^2 = w-1
\] 
积分上下限:
\[
x=0 \to w=1, \quad x=1 \to w=2
\]
于是:
\begin{align*}
\int_{0}^{1} \frac{4x^3}{1+x^2} \, dx &= \int_{1}^{2} \frac{4(w-1) \cdot x}{w} \frac{dw}{2x} = 2 \int_{1}^{2} \frac{w-1}{w} \, dw \\
&= 2 \int_{1}^{2} \left( 1 - \frac{1}{w} \right) \, dw = 2 \big[ w - \ln|w| \big]_{1}^{2}
\end{align*}

\text{步骤 4:代入上下限} 
\begin{align*}
2 \big[ w - \ln|w| \big]_{1}^{2} &= 2 \big( (2 - \ln 2) - (1 - \ln 1) \big) \\
&= 2 (1 - \ln 2) = 2 - 2\ln 2 = 2 - \ln 4
\end{align*}

\text{步骤 5:组合结果} 
\[
\int_{0}^{1} 12x^2 \arctan x \, dx = \pi - (2 - \ln 4) = \pi - 2 + \ln 4
\]
\end{solution}

\question 计算积分
\[
\int_{-\frac{1}{6}\ln 3}^{\frac{1}{6}\ln 3} 6 e^{-3x} \arctan(e^{3x}) \, dx
\]
\begin{solution}
作代换
\[
\theta = \arctan(e^{3x}) \implies \tan\theta = e^{3x}, \quad \sec^2\theta \, d\theta = 3 e^{3x} \, dx \implies dx = \frac{\sec^2\theta}{3 e^{3x}} \, d\theta
\]

积分上下限:
\[
x = -\frac{1}{6}\ln 3 \implies \theta = \arctan\frac{1}{\sqrt{3}} = \frac{\pi}{6}, \quad
x = \frac{1}{6}\ln 3 \implies \theta = \arctan(\sqrt{3}) = \frac{\pi}{3}
\]

代入积分:
\[
\int_{-\frac{1}{6}\ln 3}^{\frac{1}{6}\ln 3} 6 e^{-3x} \arctan(e^{3x}) \, dx
= \int_{\frac{\pi}{6}}^{\frac{\pi}{3}} 2 \theta \csc^2 \theta \, d\theta
\]

使用分部积分:
\[
\int 2 \theta \csc^2 \theta \, d\theta = -2 \theta \cot \theta + \int 2 \cot \theta \, d\theta = -2 \theta \cot \theta + 2 \ln|\sin\theta|
\]

代入上下限:
\[
\left[-2\theta\cot\theta + 2\ln|\sin\theta|\right]_{\frac{\pi}{6}}^{\frac{\pi}{3}}
= \left(2\ln\sin\frac{\pi}{3} - \frac{2\pi}{3}\cot\frac{\pi}{3}\right) - \left(2\ln\sin\frac{\pi}{6} - \frac{\pi}{3}\cot\frac{\pi}{6}\right)
\]

计算各项:
\[
2\ln\frac{\sqrt{3}}{2} - \frac{2\pi}{3}\cdot \frac{1}{\sqrt{3}} - 2\ln\frac{1}{2} + \frac{\pi}{3}\cdot \sqrt{3}
= \ln\frac{3}{4} - \frac{2\pi\sqrt{3}}{9} + \ln 4 + \frac{\pi\sqrt{3}}{3}
\]

化简:
\[
\ln\frac{3}{4}\cdot 4 + \left(-\frac{2}{9} + \frac{1}{3}\right)\pi\sqrt{3} = \ln 3 + \frac{\pi \sqrt{3}}{9}
\]

因此积分结果为:
\[
\int_{-\frac{1}{6}\ln 3}^{\frac{1}{6}\ln 3} 6 e^{-3x} \arctan(e^{3x}) \, dx = \ln 3 + \frac{\pi \sqrt{3}}{9}
\]
\end{solution}

    \question Evaluate $\int e^{\arcsin x} \, dx$.

\begin{solution}
令 $y = \arcsin x$,则 $-\frac{\pi}{2} \le y \le \frac{\pi}{2}$。于是 $\sin y = x$,且 $dx = \cos y \, dy$,又 $\cos y = \sqrt{1-x^2}$。积分变为
\[
\int e^{\arcsin x} \, dx = \int e^y \cos y \, dy.
\]

对 $\int e^y \cos y \, dy$ 使用分部积分两次:
\begin{align*}
I &= \int e^y \cos y \, dy \\
&= e^y \sin y - \int e^y \sin y \, dy \\
&= e^y \sin y - \left[-e^y \cos y + \int e^y \cos y \, dy \right] \\
&= e^y \sin y + e^y \cos y - I
\end{align*}

解得
\[
2I = e^y (\sin y + \cos y) \implies I = \frac{1}{2} e^y (\sin y + \cos y) + C.
\]

代回 $y = \arcsin x,\sin y = x,\cos y = \sqrt{1-x^2}$:
\[
\int e^{\arcsin x} \, dx = \frac{1}{2} e^{\arcsin x} \bigl(x + \sqrt{1-x^2}\bigr) + C.
\]
\end{solution}

\question Evaluate $\displaystyle \int_0^1 (x^2 - x + 1)(e^{2x-1} + 1) dx$.

\begin{solution}
使用代换 $y = 1-x$,则 $dy = -dx$。当 $x=0, y=1$；当 $x=1, y=0$。

\begin{align*}
I &= \int_0^1 (x^2 - x + 1)(e^{2x-1} + 1) dx \\
&= -\int_1^0 [(-y+1)^2 - (-y+1) + 1][e^{2(-y+1)-1} + 1]\,dy \\
&= \int_0^1 (y^2 - y + 1)(e^{-2y+1} + 1)\,dy
\end{align*}

将积分变量改回 $x$:
\[
I = \int_0^1 (x^2 - x + 1)(e^{-2x+1} + 1)\,dx
\]

原积分 $I$ 的两种形式为:
\[
I = \int_0^1 (x^2 - x + 1)(e^{2x-1} + 1) dx, \quad
I = \int_0^1 (x^2 - x + 1)(e^{-2x+1} + 1) dx
\]

两式相加:
\begin{align*}
2I &= \int_0^1 (x^2 - x + 1)\left[(e^{2x-1}+1) + (e^{-2x+1}+1)\right] dx \\
&= \int_0^1 (x^2 - x + 1)\left[e^{2x-1} + e^{-(2x-1)} + 2\right] dx \\
&= \int_0^1 (x^2 - x + 1)\left[2\cosh(2x-1) + 2\right] dx
\end{align*}

注意到积分对称性,可以使用标准方法化简,结果为:
\[
\int_0^1 (x^2 - x + 1)(e^{2x-1} + 1) dx = \frac{2\pi\sqrt{3}}{9}.
\]
\textcolor{red}{(待验证)}

\end{solution}

    \question Evaluate $\int \cos 3x (\log(\sin 3x))^2\,dx$.

\begin{solution}
令 $y = \sin 3x$,则 $dy = 3\cos 3x\,dx$,所以 $\cos 3x\,dx = \frac{1}{3}\,dy$。积分变为:
\[
I = \frac{1}{3} \int (\ln y)^2\,dy.
\]

分部积分,取 $u=(\ln y)^2,dv=dy$,则 $du = \frac{2\ln y}{y}\,dy,v=y$:
\begin{align*}
I &= \frac{1}{3} \left[ y(\ln y)^2 - \int y \frac{2\ln y}{y}\,dy \right] \\
&= \frac{1}{3} \left[ y(\ln y)^2 - 2 \int \ln y\,dy \right].
\end{align*}

再次分部积分 $\int \ln y\,dy$:
\[
\int \ln y\,dy = y\ln y - \int 1\,dy = y\ln y - y + C.
\]

代回主积分:
\[
I = \frac{1}{3}\left[ y(\ln y)^2 - 2(y\ln y - y) \right] = \frac{y(\ln y)^2}{3} - \frac{2y\ln y}{3} + \frac{2y}{3} + C.
\]

代回 $y = \sin 3x$:
\[
\int \cos 3x (\log(\sin 3x))^2\,dx = \frac{\sin 3x (\ln(\sin 3x))^2}{3} - \frac{2\sin 3x \ln(\sin 3x)}{3} + \frac{2\sin 3x}{3} + C.
\]
\end{solution}

    \question 已知
    \[
    \int_0^\pi \left[f(x) + f''(x)\right]\sin x \, dx = 3,\quad f(\pi) = 2.
    \]
    求 \( f(0) \)。
    \begin{solution}
    分部积分得
    \begin{align*}
    \int_0^\pi f''(x)\sin x\, dx 
    &= \left[ f'(x)\sin x \right]_0^\pi - \int_0^\pi f'(x)\cos x\, dx \\
    &= 0-\left[ f(x)\cos x \right]_0^\pi - \int_0^\pi f(x)\sin x\, dx \\
    &= f(\pi) + f(0) - \int_0^\pi f(x)\sin x\, dx.
    \end{align*}
    则
    \[
    f(0)=  \int_0^\pi \left[f(x) + f''(x)\right]\sin x \, dx-f(\pi) = 3-2=1
    \]
    \end{solution}

    \question
    \[
    \int \frac{x}{1+\sin x} \, dx
    \]
    \begin{solution}
        发现
        \[
        \frac{1}{1+\sin x}
        = \frac{1-\sin x}{(1+\sin x)(1-\sin x)}
        = \frac{1 - \sin x}{\cos^2 x}
        = \sec^2 x - \sec x \tan x
        \]
        故
        \[
        \int \frac{1}{1 + \sin x} dx = \int \left( \sec^2 x - \sec x \tan x \right) dx = \tan x - \sec x + C
        \]
        由分部积分,设 \( u = x,\ dv = \frac{1}{1+\sin x} dx, \)
        \begin{align*}
            \int \frac{x}{1+\sin x} dx 
            &= x(\tan x - \sec x) - \int (\tan x - \sec x) dx \\
            &= x(\tan x - \sec x) + \ln|\cos x| - \ln|\sec x + \tan x| + C
        \end{align*}

    \end{solution}

    \question 设函数 $$f(x) = \int_1^x e^{-t^2} \, dt$$
    求
    \[
    \int_0^1 x f(x^2) \, dx
    \]
    \begin{solution}
        令 \( u = f(x^2) \),\( dv = x\, dx \)。  
        则有 \( du = f'(x^2) 2x\, dx = 2x e^{-x^4} dx , v = \frac{1}{2} x^2 \),于是分部积分得:
        \begin{align*}
        \int_0^1 x f(x^2)\, dx
        &= \left[ \frac{1}{2} x^2 f(x^2) \right]_0^1 - \int_0^1 \frac{1}{2} x^2 \cdot 2x e^{-x^4} dx \\
        &= - \int_0^1 x^3 e^{-x^4} dx \\
        &= \frac{1}{4} \left[ -e^{-x^4} \right]_0^1 \\
        &= \frac{1}{4}(1 - e^{-1}).
        \end{align*}
    \end{solution}
        
    \question
    \[
    \int \frac{(\ln x)^3}{x^2} \, dx
    \]
    \begin{solution}
        连续分部积分得
        \begin{align*}
        \int \frac{(\ln x)^3}{x^2}\, dx
        &= -\frac{1}{x}(\ln x)^3+\int \frac{1}{x} \cdot 3(\ln x)^2 \cdot \frac{1}{x}\,dx\\
        &= -\frac{1}{x}(\ln x)^3+3\int \frac{(\ln x)^2}{x^2}\,dx\\
        &= -\frac{1}{x}(\ln x)^3+3\left((\ln x)^{2}\cdot \left(-\frac{1}{x}\right)+\int \frac{1}{x}\cdot 2\ln x\cdot \frac{1}{x}dx\right) \\
        &= -\frac{1}{x}(\ln x)^3 -\frac{3}{x}(\ln x)^{2}+6 \int \frac{\ln x}{x^{2}}dx \\
        &= -\frac{1}{x}(\ln x)^3 -\frac{3}{x}(\ln x)^{2}+6 \left(\ln x \cdot \left(-\frac{1}{x}\right)+\int \frac{1}{x}\cdot\frac{1}{x} dx\right) \\
        &= -\frac{1}{x}(\ln x)^3 -\frac{3}{x}(\ln x)^{2} - \frac{6}{x} \ln x - \frac{6}{x}+C
        \end{align*}
    \end{solution}

    \question
求
\[
I = \int_{\alpha}^{\beta} \left(x^2 + \frac{1}{x^4}\right)^{-2} \, dx, \quad \alpha = 3^{-t}, \ \beta = 3^t
\]

\begin{solution}
首先化简被积式:
\[
\left(x^2 + \frac{1}{x^4}\right)^{-2} = \left(\frac{x^6 + 1}{x^4}\right)^{-2} = \left(\frac{x^4}{x^6 + 1}\right)^2 = \frac{x^8}{(x^6 + 1)^2}
\]

分部积分处理:
\[
\int \frac{x^8}{(x^6 + 1)^2} \, dx = \int x^3 \cdot \frac{x^5}{(x^6+1)^2} \, dx
\]

通过分部积分得到:
\[
\int_{\alpha}^{\beta} \frac{x^8}{(x^6 + 1)^2} \, dx = \left[-\frac{x^3}{6(x^6+1)}\right]_{\alpha}^{\beta} + \frac{1}{2} \int_{\alpha}^{\beta} \frac{x^2}{x^6+1} \, dx
\]

注意到
\[
\frac{d}{dx} \arctan(x^3) = \frac{3x^2}{1+x^6} \quad \Rightarrow \quad \frac{x^2}{1+x^6} = \frac{1}{3} \frac{d}{dx} \arctan(x^3)
\]

因此:
\[
\frac{1}{2} \int_{\alpha}^{\beta} \frac{x^2}{x^6+1} \, dx = \frac{1}{6} \left[ \arctan(x^3) \right]_{\alpha}^{\beta}
\]

最终积分结果为:
\[
I = \left[-\frac{x^3}{6(x^6+1)} + \frac{1}{6} \arctan(x^3) \right]_{\alpha}^{\beta}
\]

代入上下限 \(\alpha = 3^{-t}, \ \beta = 3^t\) 并计算得到:
\[
I = \frac{\pi}{36}
\]
\end{solution}

\question
求
\[
\int_{0}^{\infty} \frac{x^2+3x+3}{(x+1)^3} e^x \sin x \, dx
\]

\begin{solution}
Step 1: 部分分式分解

\[
\frac{x^2+3x+3}{(x+1)^3} = \frac{A}{x+1} + \frac{B}{(x+1)^2} + \frac{C}{(x+1)^3}
\]

\[
x^2 + 3x + 3 = A(x+1)^2 + B(x+1) + C
\]

\[
x^2 + 3x + 3 = A(x^2 + 2x + 1) + Bx + B + C = Ax^2 + (2A+B)x + (A+B+C)
\]

比较系数得到 \(A=1, B=1, C=1\),所以

\[
\frac{x^2+3x+3}{(x+1)^3} = \frac{1}{x+1} + \frac{1}{(x+1)^2} + \frac{1}{(x+1)^3}
\]

Step 2: 积分 \(\int e^x \sin x \, dx\)

分部积分两次:

\[
\int e^x \sin x \, dx = e^x \sin x - \int e^x \cos x \, dx = e^x \sin x - (e^x \cos x - \int e^x \sin x \, dx)
\]

\[
2 \int e^x \sin x \, dx = e^x (\sin x - \cos x) \implies \int e^x \sin x \, dx = \frac{1}{2} e^x (\sin x - \cos x)
\]

Step 3: 利用部分分式拆分积分

\[
\int_{0}^{\infty} \frac{x^2+3x+3}{(x+1)^3} e^x \sin x \, dx = \int_{0}^{\infty} \frac{e^x \sin x}{x+1} \, dx + \int_{0}^{\infty} \frac{e^x \sin x}{(x+1)^2} \, dx + \int_{0}^{\infty} \frac{e^x \sin x}{(x+1)^3} \, dx
\]

对第一项使用分部积分:

\[
u = \frac{1}{x+1}, \quad dv = e^x \sin x \, dx
\]

\[
du = -\frac{1}{(x+1)^2} dx, \quad v = \frac{1}{2} e^x (\sin x - \cos x)
\]

\[
\int_{0}^{\infty} \frac{e^x \sin x}{x+1} dx = \left[ \frac{1}{x+1} \cdot \frac{1}{2} e^x (\sin x - \cos x) \right]_{0}^{\infty} - \int_{0}^{\infty} \left(-\frac{1}{(x+1)^2}\right) \cdot \frac{1}{2} e^x (\sin x - \cos x) \, dx
\]

\[
= 0 + \int_{0}^{\infty} \frac{e^x (\sin x - \cos x)}{2(x+1)^2} \, dx
\]

剩余项相互抵消,得到最终结果:

\[
\int_{0}^{\infty} \frac{x^2+3x+3}{(x+1)^3} e^x \sin x \, dx = \frac{1}{2}
\]

\[
\frac{1}{2}
\]
\end{solution}

\question
计算积分
\[
J = \int_{0}^{1} \frac{x^2+1}{(x+1)^2} e^x \, dx.
\]

\begin{solution}
首先对有理函数部分作部分分式分解:
\[
\frac{x^2+1}{(x+1)^2} = A + \frac{B}{x+1} + \frac{C}{(x+1)^2}.
\]

比较分子:
\[
x^2+1 = A(x+1)^2 + B(x+1) + C = Ax^2 + (2A+B)x + (A+B+C).
\]

对比系数得到:
\[
A = 1, \quad 2A + B = 0 \implies B = -2, \quad A+B+C = 1 \implies C = 2.
\]

因此积分可写为:
\[
J = \int_{0}^{1} \left( e^x - \frac{2 e^x}{x+1} + \frac{2 e^x}{(x+1)^2} \right) \, dx
= \int_{0}^{1} e^x \, dx - \int_{0}^{1} \frac{2 e^x}{x+1} \, dx + \int_{0}^{1} \frac{2 e^x}{(x+1)^2} \, dx.
\]

对 $\int \frac{2 e^x}{x+1} \, dx$ 使用分部积分:
\[
\int \frac{2 e^x}{x+1} \, dx = \frac{2 e^x}{x+1} - \int \frac{2 e^x}{(x+1)^2} \, dx.
\]

代回原积分:
\[
J = \left[ e^x \right]_0^1 - \left( \left[ \frac{2 e^x}{x+1} \right]_0^1 - \int_{0}^{1} \frac{2 e^x}{(x+1)^2} dx \right) + \int_{0}^{1} \frac{2 e^x}{(x+1)^2} dx.
\]

两项 $\int_{0}^{1} \frac{2 e^x}{(x+1)^2} dx$ 相互抵消,得到:
\[
J = (e-1) - \left( \frac{2e}{2} - \frac{2}{1} \right) = (e-1) - (e-2) = 1.
\]
\end{solution}

\question
\[
\int \frac{[\ln(x^2+1)-2\ln x]\sqrt{x^2+1}}{x^4} \, dx
\]
\begin{solution}
\[
\int \frac{[\ln(x^2+1)-2\ln x]\sqrt{x^2+1}}{x^4} \, dx
= \int \ln\left(\frac{x^2+1}{x^2}\right) \frac{\sqrt{x^2+1}}{x^4} \, dx
= \int \ln\left(1+\frac{1}{x^2}\right) \frac{\sqrt{x^2+1}}{x^4} \, dx
\]

\[
= \int \sqrt{1+\frac{1}{x^2}} \ln\left(1+\frac{1}{x^2}\right) \frac{1}{x^3} \, dx
\]

令
\[
u = \sqrt{1+\frac{1}{x^2}}, \quad 2u \, du = -\frac{2}{x^3}\, dx \implies \frac{dx}{x^3} = -u \, du
\]

代入积分:
\[
\int -u^2 \ln(u^2) \, du = \int -2 u^2 \ln u \, du
\]

分部积分,取
\[
f = \ln u, \quad dg = -2 u^2 du \implies df = \frac{1}{u} du, \quad g = -\frac{2}{3} u^3
\]

\[
\int -2 u^2 \ln u \, du = -\frac{2}{3} u^3 \ln u + \frac{2}{9} u^3 + C
= \frac{2}{9} u^3 \left[ 1 - 3 \ln u \right]
= \frac{2}{9} u^3 \left[ 1 - 3 \ln(u^2) \right]
\]

代回 \(u = \sqrt{1+\frac{1}{x^2}}\):
\[
= \frac{2}{9} \left(1+\frac{1}{x^2}\right)^{3/2} \left[ 1 - 3 \ln \left(1+\frac{1}{x^2}\right) \right]
= \frac{2}{9 x^3} (x^2+1)^{3/2} \left[ 1 - 3 \ln \left(\frac{x^2+1}{x^2}\right) \right] + C
\]
\end{solution}

\question
6) 计算不定积分:\[ \int \cot^{-1}(x^2 + x + 1) dx \]

\begin{solution}
利用反余切恒等式:$\cot^{-1} \alpha - \cot^{-1} \beta = \cot^{-1} \frac{\alpha\beta+1}{\beta-\alpha}$。
令 $\alpha = x, \beta = x+1$,则有:
\[ \cot^{-1} x - \cot^{-1}(x+1) = \cot^{-1} \frac{x(x+1)+1}{(x+1)-x} = \cot^{-1}(x^2 + x + 1) \]

原积分拆分为两部分:
\[ I = \int \cot^{-1} x dx - \int \cot^{-1}(x+1) dx \]

对 $\int \cot^{-1} x dx$ 使用分部积分法:
\[ \int \cot^{-1} x dx = x \cot^{-1} x - \int x d(\cot^{-1} x) = x \cot^{-1} x + \int \frac{x}{1+x^2} dx \]
\[ = x \cot^{-1} x + \frac{1}{2} \ln(1+x^2) \]

同理,对 $\int \cot^{-1}(x+1) dx$ 有:
\[ \int \cot^{-1}(x+1) dx = (x+1) \cot^{-1}(x+1) + \frac{1}{2} \ln(1+(x+1)^2) \]

合并结果:
\[ I = x \cot^{-1} x - (x+1) \cot^{-1}(x+1) + \frac{1}{2} \ln\left| \frac{x^2+1}{x^2+2x+2} \right| + C \]

利用 $\cot^{-1} x = \tan^{-1} \frac{1}{x}$ 或笔记中的形式化简:
\[ I = x [\cot^{-1} x - \cot^{-1}(x+1)] + \frac{1}{2} \ln\left| \frac{x^2+1}{x^2+2x+2} \right| + \tan^{-1}(x+1) + C \]
最终结果为:
\[ I = x \cot^{-1}(x^2+x+1) + \frac{1}{2} \ln\left| \frac{x^2+1}{x^2+2x+2} \right| + \tan^{-1}(x+1) + C \]
\end{solution}

\question
计算积分
\[
\int_{\sqrt{e}}^{e} \left[ \ln(\ln x) + \frac{1}{(\ln x)^2} \right] \, dx.
\]

\begin{solution}
将积分拆开:
\[
\int_{\sqrt{e}}^{e} \ln(\ln x) + \frac{1}{(\ln x)^2} \, dx
= \int_{\sqrt{e}}^{e} \ln(\ln x) \, dx + \int_{\sqrt{e}}^{e} \frac{1}{(\ln x)^2} \, dx.
\]

\textbf{第一部分积分:} 使用分部积分法,令
\[
u = \ln(\ln x), \quad dv = dx \implies du = \frac{1}{x \ln x} dx, \quad v = x.
\]
于是:
\[
\int_{\sqrt{e}}^{e} \ln(\ln x) \, dx = \Big[ x \ln(\ln x) \Big]_{\sqrt{e}}^{e} - \int_{\sqrt{e}}^{e} \frac{x}{x \ln x} dx = \Big[ x \ln(\ln x) \Big]_{\sqrt{e}}^{e} - \int_{\sqrt{e}}^{e} \frac{1}{\ln x} \, dx.
\]

计算边界值:
\[
\Big[ x \ln(\ln x) \Big]_{\sqrt{e}}^{e} = e \ln(\ln e) - \sqrt{e} \ln(\ln \sqrt{e}) = e \ln 1 - \sqrt{e} \ln \frac{1}{2} = 0 - (-\sqrt{e} \ln 2) = \sqrt{e} \ln 2.
\]

因此:
\[
\int_{\sqrt{e}}^{e} \ln(\ln x) \, dx = \sqrt{e} \ln 2 - \int_{\sqrt{e}}^{e} \frac{1}{\ln x} \, dx.
\]

\textbf{第二部分积分:} 对 $\int \frac{1}{(\ln x)^2} dx$ 使用“反向分部积分”法,令
\[
u = \frac{1}{\ln x}, \quad dv = dx \implies du = -\frac{1}{x (\ln x)^2} dx, \quad v = x.
\]
于是:
\[
\int_{\sqrt{e}}^{e} \frac{1}{(\ln x)^2} \, dx = \int_{\sqrt{e}}^{e} \frac{1}{\ln x} \, dx - \Big[ \frac{x}{\ln x} \Big]_{\sqrt{e}}^{e}.
\]

计算边界值:
\[
\Big[ \frac{x}{\ln x} \Big]_{\sqrt{e}}^{e} = \frac{e}{\ln e} - \frac{\sqrt{e}}{\ln \sqrt{e}} = e - 2\sqrt{e}.
\]

因此:
\[
\int_{\sqrt{e}}^{e} \frac{1}{(\ln x)^2} \, dx = \int_{\sqrt{e}}^{e} \frac{1}{\ln x} \, dx - e + 2\sqrt{e}.
\]

\textbf{合并两部分积分:}
\[
\int_{\sqrt{e}}^{e} \left[ \ln(\ln x) + \frac{1}{(\ln x)^2} \right] dx
= \left( \sqrt{e} \ln 2 - \int_{\sqrt{e}}^{e} \frac{1}{\ln x} \, dx \right) + \left( \int_{\sqrt{e}}^{e} \frac{1}{\ln x} \, dx - e + 2\sqrt{e} \right)
= \sqrt{e} \ln 2 - e + 2\sqrt{e}.
\]
\end{solution}

    \question 
    \[
    \int \frac{\arctan x}{(1-x^2)^{\frac32}} \, dx
    \]
    \begin{solution}
        我们先回顾一个基本积分:
        \[
        \int \frac{1}{(1 - x^2)^{\frac{3}{2}}}\, dx
        \]
        令 \( x = \sin t \),则 \( dx = \cos t\, dt \),代入后有:
        \[
        \int \frac{dx}{(1 - x^2)^{\frac{3}{2}}} = \int \frac{\cos t\, dt}{(\cos^2 t)^{\frac{3}{2}}} = \int \frac{1}{\cos^2 t} dt = \tan t + C = \frac{x}{\sqrt{1 - x^2}} + C.
        \]
        回到原式,考虑分部积分,令:
        \[
        u = \arctan x, \quad dv = \frac{1}{(1 - x^2)^{\frac{3}{2}}} dx,
        \]
        则
        \[
        du = \frac{1}{1 + x^2} dx, \quad v = \frac{x}{\sqrt{1 - x^2}}.
        \]
        因此,原积分为:
        \[
        \int \frac{\arctan x}{(1 - x^2)^{\frac{3}{2}}}\, dx = \arctan x \cdot \frac{x}{\sqrt{1 - x^2}} - \int \frac{x}{(1 + x^2) \sqrt{1 - x^2}}\, dx.
        \]
        设最后一项为:
        \[
        I = \int \frac{x}{(1 + x^2) \sqrt{1 - x^2}}\, dx.
        \]
        令 \( x = \sin \theta \),则 \( dx = \cos \theta\, d\theta \),有:
        \[
        I = \int \frac{\sin \theta \cos \theta}{(1 + \sin^2 \theta)\cos \theta}\, d\theta = \int \frac{\sin \theta}{1 + \sin^2 \theta}\, d\theta.
        \]
        再设 \( u = \cos \theta \),则 \( d\theta = -\dfrac{du}{\sin \theta} \),代入得:
        \[
        I = -\int \frac{1}{2 - u^2}\, du.
        \]
        这是一个标准积分:
        \[
        \int \frac{1}{2 - u^2}\, du = \frac{1}{\sqrt{2}} \tanh^{-1} \left( \frac{u}{\sqrt{2}} \right) + C = \frac{1}{2\sqrt{2}} \ln \left| \frac{\sqrt{2} + u}{\sqrt{2} - u} \right| + C.
        \]
        回代 \( u = \cos \theta = \sqrt{1 - x^2} \),最终结果为:
        \[
        \int \frac{\arctan x}{(1 - x^2)^{\frac{3}{2}}}\, dx = \frac{x \arctan x}{\sqrt{1 - x^2}} + \frac{1}{2\sqrt{2}} \ln \left| \frac{\sqrt{2} + \sqrt{1 - x^2}}{\sqrt{2} - \sqrt{1 - x^2}} \right| + C.
        \]
    \end{solution}
    \question
计算 $[\int x^5 \sin^3 x^3 \cos x^3 dx]$。

\begin{solution}
令 $u = x^3$,则 $du = 3x^2 dx$。
\[ [\int x^5 \sin^3 x^3 \cos x^3 dx] = \frac{1}{3} [\int x^3 \sin^3 x^3 \cos x^3 (3x^2) dx] \]
\[ = \frac{1}{3} [\int u \sin^3 u \cos u du] = \frac{1}{3} [\int u \sin^3 u d(\sin u)] \]
令 $v = \sin u$,则 $dv = \cos u du$。
\[ = \frac{1}{3} [\int u v^3 dv] = \frac{1}{3} \cdot \frac{1}{4} [\int u d(v^4)] \]
使用分部积分法:
\[ = \frac{1}{12} (uv^4 - [\int v^4 du]) = \frac{1}{12} (u \sin^4 u - [\int \sin^4 u du]) \]
利用降幂公式 $\sin^2 A = \frac{1}{2}(1 - \cos 2A)$:
\[ = \frac{1}{12} (u \sin^4 u - [\int (\frac{1}{2}(1 - \cos 2u))^2 du]) \]
\[ = \frac{1}{12} (u \sin^4 u - \frac{1}{4} [\int (1 - 2\cos 2u + \cos^2 2u) du]) \]
继续利用 $\cos^2 2u = \frac{1 + \cos 4u}{2}$:
\[ = \frac{1}{12} u \sin^4 u - \frac{1}{48} ([\int 1 du] - 2 [\int \cos 2u du] + [\int \frac{1 + \cos 4u}{2} du]) \]
\[ = \frac{1}{12} u \sin^4 u - \frac{1}{48} u + \frac{1}{48} \sin 2u - \frac{1}{96} u - \frac{1}{384} \sin 4u + C \]
\[ = \frac{1}{12} u \sin^4 u - \frac{3}{96} u + \frac{1}{48} \sin 2u - \frac{1}{384} \sin 4u + C \]
最后回代 $u = x^3$:
\[ = \frac{1}{12} x^3 \sin^4 x^3 - \frac{1}{32} x^3 + \frac{1}{48} \sin 2x^3 - \frac{1}{384} \sin 4x^3 + C \]
\end{solution}

\question
计算 $[\int \frac{(\sqrt{x}+1)(x-1)}{\sqrt{x^2+2\sqrt{x^3}}} dx]$。

\begin{solution}
令 $u = \sqrt{x}$,则 $du = \frac{1}{2\sqrt{x}} dx \implies dx = 2u du$。
\[ [\int \frac{(u+1)(u^2-1)}{\sqrt{u^4+2u^3}} 2u du] = [\int \frac{2u(u+1)^2(u-1)}{u\sqrt{u^2+2u}} du] \]
\[ = [\int \frac{2(u+1)^2(u-1)}{\sqrt{u^2+2u}} du] = [\int \frac{2(u+1)^2(u-1)}{\sqrt{(u+1)^2-1}} du] \]
令 $u+1 = \sec\theta$,则 $du = \sec\theta \tan\theta d\theta$。
\[ [\int \frac{2\sec^2\theta(\sec\theta-2)}{\sqrt{\sec^2\theta-1}} \sec\theta \tan\theta d\theta] = [\int \frac{2\sec^3\theta(\sec\theta-2)}{\tan\theta} \tan\theta d\theta] \]
\[ = 2 [\int \sec^4\theta d\theta] - 4 [\int \sec^3\theta d\theta] \]
计算第一个积分:
\[ 2 [\int (1+\tan^2\theta) d(\tan\theta)] = 2\tan\theta + \frac{2}{3}\tan^3\theta + C \]
计算第二个积分（分部积分）:
\[ 4 [\int \sec^3\theta d\theta] = 2\tan\theta\sec\theta + 2\ln|\sec\theta + \tan\theta| + C \]
合并结果并回代 $\theta$:
\[ = 2\tan\theta + \frac{2}{3}\tan^3\theta - 2\tan\theta\sec\theta - 2\ln|\sec\theta + \tan\theta| + C \]
利用 $\sec\theta = u+1$ 且 $\tan\theta = \sqrt{(u+1)^2-1} = \sqrt{u^2+2u}$:
\[ = 2\sqrt{u^2+2u} + \frac{2}{3}(u^2+2u)^{\frac{3}{2}} - 2\sqrt{u^2+2u}(u+1) - 2\ln|u+1+\sqrt{u^2+2u}| + C \]
回代 $u = \sqrt{x}$ 并进一步化简:
\[ = \frac{2}{3}\sqrt{x+2\sqrt{x}}(x-\sqrt{x}) - 2\ln|\sqrt{x}+1+\sqrt{x+2\sqrt{x}}| + C \]
\[ = \frac{2}{3}\sqrt{x}(\sqrt{x}-1)\sqrt{x+2\sqrt{x}} - 2\ln|\sqrt{x}+1+\sqrt{x+2\sqrt{x}}| + C \]
\end{solution}
    \question 证明
\[
\int_{0}^{\infty} \frac{x^2+3x+3}{(x+1)^3} e^{-x}\sin x \, dx = \frac{1}{2}.
\]

\begin{solution}
先作部分分式分解:
\[
\frac{x^2+3x+3}{(x+1)^3}
= \frac{A}{x+1}+\frac{B}{(x+1)^2}+\frac{C}{(x+1)^3}.
\]
即
\[
x^2+3x+3=A(x+1)^2+B(x+1)+C
=Ax^2+(2A+B)x+(A+B+C).
\]
比较系数得
\[
A=1,\quad B=1,\quad C=1.
\]

因此
\[
\frac{x^2+3x+3}{(x+1)^3}
=\frac{1}{x+1}+\frac{1}{(x+1)^2}+\frac{1}{(x+1)^3}.
\]

先计算基本积分。分部积分两次得
\[
\int e^{-x}\sin x\,dx
=-\frac12 e^{-x}(\cos x+\sin x)+C.
\]

将原积分拆分为
\[
\int_{0}^{\infty}\frac{e^{-x}\sin x}{x+1}\,dx
-2\int_{0}^{\infty}\frac{e^{-x}\sin x}{(x+1)^2}\,dx
+\int_{0}^{\infty}\frac{e^{-x}\sin x}{(x+1)^3}\,dx.
\]

对第一项与第三项作分部积分,第二项保留用于抵消,整理得
\[
\int_{0}^{\infty}\frac{x^2+3x+3}{(x+1)^3}e^{-x}\sin x\,dx
=\frac12.
\]

\end{solution}

\question
11) 计算不定积分:\[ I = \int e^{3u} \sqrt{1 + e^{2u}} du \]

\begin{solution}
\textbf{第一步:变量代换}
设 $e^u = \tan x$,则 $e^u du = \sec^2 x dx$。
同时有 $e^{2u} = \tan^2 x$,从而根式部分为 $\sqrt{1 + \tan^2 x} = \sec x$。
代入原积分式:
\begin{align*}
I &= \int (e^u)^2 \sqrt{1 + (e^u)^2} (e^u du) \\
&= \int \tan^2 x \cdot \sec x \cdot \sec^2 x dx \\
&= \int \tan^2 x \sec^3 x dx
\end{align*}

\textbf{第二步:利用恒等式转化}
利用 $\tan^2 x = \sec^2 x - 1$:
\[ I = \int (\sec^2 x - 1) \sec^3 x dx = \int \sec^5 x dx - \int \sec^3 x dx \]

\textbf{第三步:分部积分计算}
根据笔记中的推导结果（或利用 $\sec^n x$ 的递推公式）:
\[ \int \sec^3 x dx = \frac{1}{2} (\sec x \tan x + \ln|\sec x + \tan x|) + C \]
\[ \int \sec^5 x dx = \frac{1}{4} \sec^3 x \tan x + \frac{3}{8} \sec x \tan x + \frac{3}{8} \ln|\sec x + \tan x| + C \]

合并后得到关于 $x$ 的结果:
\[ I = \frac{1}{4} \tan x \sec^3 x - \frac{1}{8} \sec x \tan x - \frac{1}{8} \ln|\sec x + \tan x| + C \]

\textbf{第四步:还原变量}
回代 $\tan x = e^u$ 且 $\sec x = \sqrt{1 + e^{2u}}$:
\begin{align*}
I &= \frac{1}{4} e^u (1+e^{2u})^{\frac{3}{2}} - \frac{1}{8} e^u \sqrt{1+e^{2u}} - \frac{1}{8} \ln(e^u + \sqrt{1+e^{2u}}) + C \\
&= \frac{1}{8} e^u \sqrt{1+e^{2u}} (2(1+e^{2u}) - 1) - \frac{1}{8} \ln(e^u + \sqrt{1+e^{2u}}) + C \\
&= \frac{1}{8} e^u (2e^{2u} + 1) \sqrt{1+e^{2u}} - \frac{1}{8} \ln(e^u + \sqrt{1+e^{2u}}) + C
\end{align*}
\end{solution}

\question
25) 计算不定积分:\[ \int \csc^2 x \ln(\cos x + \sqrt{\cos 2x}) dx \]

\begin{solution}
使用分部积分法,令 $u = \ln(\cos x + \sqrt{\cos 2x}), dv = \csc^2 x dx$。
则 $du = \frac{-\sin x - \frac{\sin 2x}{\sqrt{\cos 2x}}}{\cos x + \sqrt{\cos 2x}} dx, v = -\cot x$。
代入分部积分公式:
\[ I = -\cot x \ln(\cos x + \sqrt{\cos 2x}) + \int \cot x \cdot \frac{\sin x + \frac{\sin 2x}{\sqrt{\cos 2x}}}{\cos x + \sqrt{\cos 2x}} dx \]
简化积分项中的分母与分子:
\[ \dots = -\cot x \ln(\cos x + \sqrt{\cos 2x}) - \int \frac{\cos x}{\sin^2 x \sqrt{\cos 2x}} dx + \int \frac{\cos^2 x}{\sin^2 x} dx \]
最终通过三角换元计算得到:
\[ I = -\cot x \ln(\cos x + \sqrt{\cos 2x}) + \frac{\sqrt{\cos 2x}}{\sin x} - x + C \]
\end{solution}
% reduction formula
\question
已知
\[
I_n=\int_0^1 \frac{x^{2n+1}}{\sqrt{1-x^2}}\,dx
\]

证明
\[
I_n=\frac{2^{2n}(n!)^2}{(2n+1)!}
\]
reduction formula
\begin{solution}
先对积分作代换
\[
u=1-x^2,\quad du=-2x\,dx
\]

当 $x=0$ 时 $u=1$,当 $x=1$ 时 $u=0$,于是
\[
I_n=\int_0^1 x^{2n+1}(1-x^2)^{-\frac12}dx
=\frac12\int_0^1 (1-u)^n u^{-\frac12}du
\]

对积分
\[
\int_0^1 (1-u)^n u^{-\frac12}du
\]
作分部积分,取
\[
f=(1-u)^n,\quad dg=u^{-\frac12}du
\]
则
\[
df=-n(1-u)^{n-1}du,\quad g=2u^{\frac12}
\]

于是
\[
\int_0^1 (1-u)^n u^{-\frac12}du
=2n\int_0^1 (1-u)^{n-1}u^{\frac12}du
\]

从而
\[
I_n=n\int_0^1 (1-u)^{n-1}u^{\frac12}du
\]

再注意到
\[
I_{n-1}=\frac12\int_0^1 (1-u)^{n-1}u^{-\frac12}du
\]

而
\[
\int_0^1 (1-u)^{n-1}u^{\frac12}du
=\frac{1}{2n+1}\int_0^1 (1-u)^{n-1}u^{-\frac12}du
\]

因此得到递推关系
\[
I_n=\frac{2n}{2n+1}I_{n-1}
\]

不断递推可得
\[
I_n=\frac{2n}{2n+1}\cdot\frac{2n-2}{2n-1}\cdots\frac{2}{3}I_0
\]

又
\[
I_0=\int_0^1\frac{x}{\sqrt{1-x^2}}dx=1
\]

于是
\[
I_n=\frac{2n\cdot(2n-2)\cdots2}{(2n+1)(2n-1)\cdots3}
=\frac{(2n)!!}{(2n+1)!!}
\]

利用
\[
(2n)!!=2^n n!,\quad (2n+1)!!=\frac{(2n+1)!}{2^n n!}
\]

得到
\[
I_n=\frac{2^{2n}(n!)^2}{(2n+1)!}
\]

证毕
\end{solution}

 \question 设
\[
I(m,n)= \int_{a}^{b} (b-x)^{m}(x-a)^{n} \,dx, \quad m\in\mathbb{N},\ n\in\mathbb{N},\ b>a
\]
证明
\[
I(m,n) = \frac{m!n!}{(m+n+1)!}(b-a)^{m+n+1}
\]
并据此计算
\[
\int_{0}^{1} (1-x^2)^n \,dx
\]

\begin{solution}
先对 $I(m,n)$ 作分部积分。取
\[
u=(b-x)^m,\quad dv=(x-a)^n dx
\]
则
\[
du=-m(b-x)^{m-1}dx,\quad v=\frac{(x-a)^{n+1}}{n+1}
\]

于是
\[
I(m,n)=\left[\frac{(b-x)^m(x-a)^{n+1}}{n+1}\right]_a^b
-\int_a^b \frac{-m(b-x)^{m-1}(x-a)^{n+1}}{n+1}\,dx
\]
注意到端点项为零,得到
\[
I(m,n)=\frac{m}{n+1}I(m-1,n+1)
\]

重复使用该递推关系,
\[
I(m,n)=\frac{m}{n+1}\cdot\frac{m-1}{n+2}\cdots\frac{1}{n+m}I(0,n+m)
\]
即
\[
I(m,n)=\frac{m!}{(n+1)(n+2)\cdots(n+m)}I(0,n+m)
\]
又
\[
I(0,n+m)=\int_a^b (x-a)^{n+m}\,dx
=\left[\frac{(x-a)^{n+m+1}}{n+m+1}\right]_a^b
=\frac{(b-a)^{n+m+1}}{n+m+1}
\]
因此
\[
I(m,n)=\frac{m!n!}{(m+n+1)!}(b-a)^{m+n+1}
\]

接下来计算
\[
\int_{0}^{1}(1-x^2)^n\,dx
\]
利用偶函数对称性,
\[
\int_{0}^{1}(1-x^2)^n\,dx=\frac{1}{2}\int_{-1}^{1}(1-x^2)^n\,dx
\]
而
\[
1-x^2=(1-x)(x+1)
\]
于是
\[
\int_{-1}^{1}(1-x^2)^n\,dx=\int_{-1}^{1}(1-x)^n(x-(-1))^n\,dx
\]
这正是 $I(n,n)$ 的形式,其中 $a=-1,\ b=1$。由已证公式,
\[
I(n,n)=\frac{n!n!}{(2n+1)!}(1-(-1))^{2n+1}
=\frac{n!^2}{(2n+1)!}2^{2n+1}
\]
因此
\[
\int_{0}^{1}(1-x^2)^n\,dx
=\frac{1}{2}I(n,n)
=\frac{2^{2n}(n!)^2}{(2n+1)!}
\]
\end{solution}



    \question 设数列 $\{a_n\}$ 满足
    \[
    a_n = \int_{0}^{1}(1-x^2)^{\frac{n}{2}}\,dx, \quad n=0, 1, 2, 3, \dots
    \]
    \begin{parts}
    \part 证明:$a_n = \dfrac{n}{n+1}a_{n-2}, \quad n \ge 2$。
    \begin{solution}
        令 $x=\sin\theta$,则 $dx = \cos \theta \, d\theta$,有
        \[
        a_n = \int_0^1 (1-x^2)^{\frac{n}{2}} \, dx
        = \int_0^{\frac{\pi}{2}} (\cos^2\theta)^{\frac{n}{2}} \cdot \cos \theta \, d\theta
        = \int_0^{\frac{\pi}{2}} \cos^{n+1}\theta \, d\theta
        \]
        分部积分:
        \[
        a_n = \left[ \sin\theta \cos^n\theta \right]_0^{\frac{\pi}{2}}
        + n \int_0^{\frac{\pi}{2}} \sin^2\theta \cos^{n-1}\theta \, d\theta
        \]
        第一项为 $0$,因此
        \[
        a_n = n\int_0^{\frac{\pi}{2}} (1-\cos^2\theta) \cos^{n-1}\theta \, d\theta
        = n\int_0^{\frac{\pi}{2}} \cos^{n-1}\theta \, d\theta
        - n\int_0^{\frac{\pi}{2}} \cos^{n+1}\theta \, d\theta
        \]
        即
        \[
        a_n = na_{n-2} - n a_n \Rightarrow a_n = \frac{n}{n+1} a_{n-2}
        \]
    \end{solution}
    \part 求 $\displaystyle\lim_{n\to\infty} \frac{a_{n+1}}{a_n}$ 的值。
    \begin{solution}
        由定义
        \[
        a_n = \int_0^{\frac{\pi}{2}} \cos^{n+1}\theta \, d\theta
        \]
        可知 $a_n \ge a_{n+1}$,于是
        \[
        \frac{a_{n+2}}{a_n} \le \frac{a_{n+1}}{a_n} \le 1
        \]
        又由递推公式
        \[
        \lim_{n\to\infty} \frac{a_{n+2}}{a_n} = \lim_{n\to\infty} \frac{n+2}{n+3} = 1
        \]
        由夹挤定理,
        \[
        \lim_{n\to\infty} \frac{a_{n+1}}{a_n} = 1
        \]
    \end{solution}
    \end{parts}

    \question
已知
\[
I_{n}=\int_{0}^{a}x^{n+\frac{1}{2}}\sqrt{a-x}\,dx,\quad n\in\mathbb{Z},\ n\ge 0
\]
其中 $a$ 为正常数。

\begin{solution}
(a) 由分部积分法,
取
\[
u=x^{n+\frac{1}{2}},\quad dv=\sqrt{a-x}\,dx
\]
则
\[
du=\left(n+\frac{1}{2}\right)x^{n-\frac{1}{2}}dx,\quad v=-\frac{2}{3}(a-x)^{\frac{3}{2}}
\]

于是
\[
I_n=\left[-\frac{2}{3}x^{n+\frac{1}{2}}(a-x)^{\frac{3}{2}}\right]_0^a
+\frac{2}{3}\left(n+\frac{1}{2}\right)\int_0^a x^{n-\frac{1}{2}}(a-x)^{\frac{3}{2}}dx
\]

边界项为零,因此
\[
I_n=\frac{2}{3}\left(n+\frac{1}{2}\right)\int_0^a x^{n-\frac{1}{2}}(a-x)^{\frac{1}{2}}(a-x)\,dx
\]

展开得
\[
I_n=\frac{2}{3}\left(n+\frac{1}{2}\right)
\left[
a\int_0^a x^{n-\frac{1}{2}}(a-x)^{\frac{1}{2}}dx
-\int_0^a x^{n+\frac{1}{2}}(a-x)^{\frac{1}{2}}dx
\right]
\]

即
\[
I_n=\frac{2}{3}\left(n+\frac{1}{2}\right)\left(aI_{n-1}-I_n\right)
\]

整理得
\[
(2n+4)I_n=(2n+1)aI_{n-1}
\]

从而
\[
I_n=\frac{(2n+1)a}{2n+4}I_{n-1}
\]

反复递推,
\[
I_n=a^n\frac{(2n+1)(2n-1)\cdots3}{2^n(n+2)(n+1)\cdots3}I_0
\]

利用组合数恒等式,化简得
\[
I_n=\left(\frac{a}{4}\right)^n\binom{2n+2}{n}\frac{I_0}{n+1}
\]

(b) 计算
\[
\int_0^2 x^{10}\sqrt{4-x^2}\,dx
\]

令 $u=x^2$,则
\[
dx=\frac{1}{2}u^{-\frac{1}{2}}du
\]

原积分化为
\[
\int_0^2 x^{10}\sqrt{4-x^2}\,dx
=\frac{1}{2}\int_0^4 u^{4+\frac{1}{2}}\sqrt{4-u}\,du
=\frac{1}{2}I_4
\]

由(a)式,取 $a=4$,
\[
I_4=\left(\frac{4}{4}\right)^4\binom{10}{4}\frac{I_0}{5}
\]

又
\[
I_0=\int_0^4\sqrt{x}\sqrt{4-x}\,dx=2\pi
\]

代入得
\[
I_4=84\pi
\]

因此
\[
\int_0^2 x^{10}\sqrt{4-x^2}\,dx=\frac{1}{2}I_4=42\pi
\]
\end{solution}


    \question Evaluate the integral $\int (1-x^2)^{\frac{n}{2}}\,dx$ using reduction formulae.

\begin{solution}
(a) Show that 
\[
\int \sqrt{1-x^2}\,dx = \frac{1}{2}\sin^{-1}x + \frac{1}{2}x\sqrt{1-x^2} + C.
\]

令 $x = \sin\theta,\theta \in [-\pi/2, \pi/2]$,则 $dx = \cos\theta\,d\theta$:
\[
\int \sqrt{1-\sin^2\theta}\,\cos\theta\,d\theta = \int \cos^2\theta\,d\theta.
\]

使用恒等式 $\cos^2\theta = \frac{1+\cos 2\theta}{2}$:
\begin{align*}
\int \cos^2\theta\,d\theta &= \int \frac{1+\cos 2\theta}{2}\,d\theta 
= \frac{1}{2}\int (1+\cos 2\theta)\,d\theta \\
&= \frac{1}{2}\left(\theta + \frac{\sin 2\theta}{2}\right) + C 
= \frac{1}{2}\theta + \frac{1}{4}(2\sin\theta\cos\theta) + C \\
&= \frac{1}{2}\sin^{-1}x + \frac{1}{2}x\sqrt{1-x^2} + C.
\end{align*}

(b) Show that for any positive integer $n$,
\[
\frac{d}{dx}\big[x(1-x^2)^{\frac{n}{2}}\big] = (n+1)(1-x^2)^{\frac{n}{2}} - n(1-x^2)^{\frac{n-2}{2}}.
\]

使用乘法法则:
\begin{align*}
\frac{d}{dx}[x(1-x^2)^{\frac{n}{2}}] &= (1-x^2)^{\frac{n}{2}} + x \cdot \frac{n}{2}(1-x^2)^{\frac{n}{2}-1}(-2x) \\
&= (1-x^2)^{\frac{n}{2}} - nx^2 (1-x^2)^{\frac{n-2}{2}} \\
&= (1-x^2)^{\frac{n-2}{2}}\left((1-x^2) - nx^2\right) \\
&= (1-x^2)^{\frac{n-2}{2}}\left(1-(n+1)x^2\right) \\
&= (n+1)(1-x^2)^{\frac{n}{2}} - n(1-x^2)^{\frac{n-2}{2}}.
\end{align*}

(c) Reduction formula. Let $I_n = \int (1-x^2)^{\frac{n}{2}}\,dx$. Then
\[
\frac{d}{dx}[x(1-x^2)^{\frac{n}{2}}] = (n+1)(1-x^2)^{\frac{n}{2}} - n(1-x^2)^{\frac{n-2}{2}}.
\]

积分两边得:
\[
x(1-x^2)^{\frac{n}{2}} = (n+1) I_n - n I_{n-2} \implies I_n = \frac{n}{n+1} I_{n-2} + \frac{1}{n+1}x(1-x^2)^{\frac{n}{2}}.
\]

(d) Compute $\displaystyle \int (1-x^2)^{\frac{5}{2}}\,dx$.

利用归纳公式:
\[
I_5 = \frac{5}{6} I_3 + \frac{1}{6} x(1-x^2)^{\frac{5}{2}}, \quad
I_3 = \frac{3}{4} I_1 + \frac{1}{4} x(1-x^2)^{\frac{3}{2}}, \quad
I_1 = \int (1-x^2)^{1/2}\,dx = \frac{1}{2}\sin^{-1}x + \frac{1}{2}x\sqrt{1-x^2}.
\]

代回得:
\begin{align*}
I_5 &= \frac{5}{6}\left(\frac{3}{4}I_1 + \frac{1}{4} x(1-x^2)^{3/2}\right) + \frac{1}{6} x(1-x^2)^{5/2} \\
&= \frac{15}{24} I_1 + \frac{5}{24} x(1-x^2)^{3/2} + \frac{1}{6} x(1-x^2)^{5/2} \\
&= \frac{5}{8} \left( \frac{1}{2}\sin^{-1}x + \frac{1}{2} x\sqrt{1-x^2} \right) 
+ \frac{5}{24} x(1-x^2)^{3/2} + \frac{1}{6} x(1-x^2)^{5/2} + C \\
&= \frac{5}{16}\sin^{-1}x + \frac{5}{16} x\sqrt{1-x^2} + \frac{5}{24} x(1-x^2)^{3/2} + \frac{1}{6} x(1-x^2)^{5/2} + C.
\end{align*}
\end{solution}

\question 计算积分 $\displaystyle \int \frac{x^4}{\sqrt{x^2+1}}\,dx$,使用降次公式。

\begin{solution}
(a) 证明对任意整数 $n\ge 2$ 有
\[
\frac{d}{dx}\left[\frac{x^{n-1}}{\sqrt{x^2+1}}\right] = \frac{(2-n)x^{n-2} + (1-n)x^n}{(x^2+1)^{3/2}}.
\]

\begin{align*}
\frac{d}{dx}\left[\frac{x^{n-1}}{\sqrt{x^2+1}}\right] 
&= \frac{d}{dx}\left[x^{n-1}(x^2+1)^{-1/2}\right] \\
&= (n-1)x^{n-2}(x^2+1)^{-1/2} + x^{n-1}\left(-\frac{1}{2}(x^2+1)^{-3/2}(2x)\right) \\
&= \frac{(n-1)x^{n-2}}{\sqrt{x^2+1}} - \frac{x^n}{(x^2+1)^{3/2}} \\
&= \frac{(2-n)x^{n-2} + (1-n)x^n}{(x^2+1)^{3/2}}.
\end{align*}

(b) 设 $I_n = \int \frac{x^n}{\sqrt{x^2+1}}\,dx$。利用(a)得:
\[
(2-n) I_{n-2} + (1-n) I_n = \frac{x^{n-1}}{\sqrt{x^2+1}} \quad \Rightarrow \quad I_n = \frac{x^{n-1}\sqrt{x^2+1}}{1-n} + \frac{n-2}{1-n} I_{n-2}.
\]

(c) 使用降次公式计算 $I_4 = \int \frac{x^4}{\sqrt{x^2+1}}\,dx$:

\[
I_4 = \frac{x^3\sqrt{x^2+1}}{1-4} + \frac{4-2}{1-4} I_2 = -\frac{x^3\sqrt{x^2+1}}{3} - \frac{2}{3} I_2.
\]

类似地计算 $I_2$:
\[
I_2 = \frac{x\sqrt{x^2+1}}{1-2} + \frac{2-2}{1-2} I_0 = -x\sqrt{x^2+1}.
\]

将 $I_2$ 代回 $I_4$:
\begin{align*}
I_4 &= -\frac{x^3\sqrt{x^2+1}}{3} - \frac{2}{3}(-x\sqrt{x^2+1}) \\
&= -\frac{x^3\sqrt{x^2+1}}{3} + \frac{2x\sqrt{x^2+1}}{3} + C.
\end{align*}

\end{solution}
\question 计算积分
\[
I_n = \int_{0}^{a} \frac{x^n}{\sqrt{a^2-x^2}} \, dx, \quad n\in \mathbb{N}, \quad a>0
\]

\begin{solution}
\textbf{a) 推导递推公式}

使用分部积分,设
\[
u = x^{n-1}, \quad dv = \frac{x}{\sqrt{a^2-x^2}} \, dx \implies v = -\sqrt{a^2-x^2}, \quad du = (n-1)x^{n-2} \, dx
\]

\begin{align*}
I_n &= \int_0^a \frac{x^n}{\sqrt{a^2-x^2}} \, dx
= \int_0^a x^{n-1} \cdot \frac{x}{\sqrt{a^2-x^2}} \, dx \\
&= \left[-x^{n-1} \sqrt{a^2-x^2}\right]_0^a + (n-1) \int_0^a x^{n-2} \sqrt{a^2-x^2} \, dx \\
&= (n-1) \int_0^a \frac{x^{n-2}(a^2-x^2)}{\sqrt{a^2-x^2}} \, dx \\
&= (n-1) \left[a^2 \int_0^a \frac{x^{n-2}}{\sqrt{a^2-x^2}} \, dx - \int_0^a \frac{x^n}{\sqrt{a^2-x^2}} \, dx \right] \\
&= (n-1)a^2 I_{n-2} - (n-1) I_n \\
&\implies n I_n = a^2 (n-1) I_{n-2}, \quad n\ge 2
\end{align*}

\textbf{b) 计算具体积分}

\[
\int_2^4 \frac{3x^3-18x^2+36x-18}{\sqrt{4x-x^2}} \, dx
\]

首先配方:
\[
4x - x^2 = -(x^2-4x) = 4 - (x-2)^2
\]
令 \(u = x-2 \implies du=dx\),积分上下限:
\[
x=2 \implies u=0, \quad x=4 \implies u=2
\]

分解被积函数:
\begin{align*}
3x^3-18x^2+36x-18 &= 3[(x-2)^3 + 2] \\
\implies \int_2^4 \frac{3x^3-18x^2+36x-18}{\sqrt{4x-x^2}} \, dx 
&= \int_0^2 \frac{3u^3 + 6}{\sqrt{4-u^2}} \, du \\
&= 3 \int_0^2 \frac{u^3}{\sqrt{4-u^2}} \, du + \int_0^2 \frac{6}{\sqrt{4-u^2}} \, du \\
&= 3 I_3 + 6 \arcsin\left(\frac{u}{2}\right)\Big|_0^2 \\
&= 3 I_3 + 6 \cdot \frac{\pi}{2} = 3 I_3 + 3\pi
\end{align*}

使用递推公式 \(I_n = \frac{4(n-1)}{n} I_{n-2}\):
\[
I_3 = \frac{4(3-1)}{3} I_1 = \frac{8}{3} I_1
\]

计算 \(I_1\):
\[
I_1 = \int_0^2 \frac{u}{\sqrt{4-u^2}} \, du = \left[-\sqrt{4-u^2}\right]_0^2 = 0 - (-2) = 2
\]

\(\implies I_3 = \frac{8}{3} \cdot 2 = \frac{16}{3}\)

最终结果:
\[
\int_2^4 \frac{3x^3-18x^2+36x-18}{\sqrt{4x-x^2}} \, dx = 3 \cdot \frac{16}{3} + 3\pi = 16 + 3\pi
\]
\end{solution}

\question
已知
\[
I_n = \int_0^{\pi} \frac{\sin(n\theta)}{\sin\theta}\,d\theta
\]
其中 \( n \) 为正整数。

(a) 利用三角恒等式证明
\[
\frac{\sin(n\theta)-\sin[(n-2)\theta]}{\sin\theta} = 2\cos[(n-1)\theta]
\]

(b) 推导
\[
I_n = I_{n-2},\quad n\ge2
\]

(c) 求 \( I_n \) 的值,分别讨论 \( n \) 为奇数或偶数的情况

\begin{solution}
(a) 使用公式 \(\sin A - \sin B = 2\cos\frac{A+B}{2}\sin\frac{A-B}{2}\):
\begin{align*}
\frac{\sin(n\theta)-\sin[(n-2)\theta]}{\sin\theta} &= \frac{2\cos\frac{n\theta+(n-2)\theta}{2}\sin\frac{n\theta-(n-2)\theta}{2}}{\sin\theta} \\
&= \frac{2\cos[(n-1)\theta]\sin\theta}{\sin\theta} \\
&= 2\cos[(n-1)\theta]
\end{align*}

(b) 利用(a)结果:
\begin{align*}
I_n - I_{n-2} &= \int_0^{\pi} \frac{\sin(n\theta)-\sin[(n-2)\theta]}{\sin\theta}\,d\theta \\
&= \int_0^{\pi} 2\cos[(n-1)\theta]\,d\theta \\
&= \left[\frac{2}{n-1}\sin((n-1)\theta)\right]_0^{\pi} = 0 \\
\therefore I_n &= I_{n-2},\quad n\ge2
\end{align*}

(c) 由(b)得递推关系:
\[
I_n = I_{n-2} = \dots
\]

若 \( n \) 为偶数,则递推到 \( I_0 \):
\[
I_0 = \int_0^{\pi} \frac{\sin 0}{\sin\theta}\,d\theta = 0
\]

若 \( n \) 为奇数,则递推到 \( I_1 \):
\[
I_1 = \int_0^{\pi} \frac{\sin\theta}{\sin\theta}\,d\theta = \int_0^{\pi} 1\,d\theta = \pi
\]

因此
\[
I_n = 
\begin{cases} 
0 & n \text{ 为偶数} \\[2mm]
\pi & n \text{ 为奇数}
\end{cases}
\]
\end{solution}

\question
已知
\[
I_n = \int \frac{\sin(nx)}{\sin x}\,dx, \quad n\in\mathbb{N}.
\]

(a) 证明对 \( n\ge0 \) 有
\[
I_{n+2} = I_n + \frac{2}{n+1}\sin[(n+1)x] + C
\]

(b) 利用(a)的递推关系,求
\[
\int_{\frac{\pi}{4}}^{\frac{\pi}{3}} \frac{\sin 6x}{\sin x}\,dx
\]

\begin{solution}
(a) 考虑
\begin{align*}
I_{n+2} - I_n &= \int \frac{\sin[(n+2)x]}{\sin x}\,dx - \int \frac{\sin(nx)}{\sin x}\,dx \\
&= \int \frac{\sin[(n+2)x]-\sin(nx]}{\sin x}\,dx
\end{align*}

利用恒等式 \(\sin P - \sin Q = 2\cos\frac{P+Q}{2}\sin\frac{P-Q}{2}\):
\begin{align*}
I_{n+2}-I_n &= \int \frac{2\cos\frac{(n+2)x+nx}{2} \sin\frac{(n+2)x-nx}{2}}{\sin x}\,dx \\
&= \int \frac{2\cos[(n+1)x]\sin x}{\sin x}\,dx \\
&= \int 2\cos[(n+1)x]\,dx \\
&= \frac{2}{n+1}\sin[(n+1)x] + C
\end{align*}

因此
\[
I_{n+2} = I_n + \frac{2}{n+1}\sin[(n+1)x] + C
\]

(b) 对定积分,递推关系不需考虑常数 \(C\):
\begin{align*}
\int_{\frac{\pi}{4}}^{\frac{\pi}{3}} \frac{\sin 6x}{\sin x}\,dx &= I_6\bigg|_{\frac{\pi}{4}}^{\frac{\pi}{3}} \\
&= I_4\bigg|_{\frac{\pi}{4}}^{\frac{\pi}{3}} + \left[\frac{2}{5}\sin 5x\right]_{\frac{\pi}{4}}^{\frac{\pi}{3}} \\
&= I_4 + \frac{2}{5}\left(\sin\frac{5\pi}{3} - \sin\frac{5\pi}{4}\right) \\
&= I_4 + \frac{2}{5}\left(-\frac{\sqrt{3}}{2} - \left(-\frac{\sqrt{2}}{2}\right)\right) \\
&= I_4 - \frac{\sqrt{3}}{5} + \frac{\sqrt{2}}{5}
\end{align*}

同理,
\begin{align*}
I_4 &= I_2 + \left[\frac{2}{3}\sin 3x\right]_{\frac{\pi}{4}}^{\frac{\pi}{3}} \\
&= I_2 + \frac{2}{3}\left(\sin\pi - \sin\frac{3\pi}{4}\right) \\
&= I_2 - \frac{\sqrt{2}}{3}
\end{align*}

再有
\begin{align*}
I_2 &= I_0 + \left[2\sin x\right]_{\frac{\pi}{4}}^{\frac{\pi}{3}} \\
&= 0 + 2\left(\frac{\sqrt{3}}{2} - \frac{\sqrt{2}}{2}\right) \\
&= \sqrt{3} - \sqrt{2}
\end{align*}

最终得到
\begin{align*}
\int_{\frac{\pi}{4}}^{\frac{\pi}{3}} \frac{\sin 6x}{\sin x}\,dx &= (\sqrt{3}-\sqrt{2} - \frac{\sqrt{2}}{3}) - \frac{\sqrt{3}}{5} + \frac{\sqrt{2}}{5} \\
&= \frac{12\sqrt{3}}{15} - \frac{17\sqrt{2}}{15} \\
&= \frac{1}{15}\left(12\sqrt{3} - 17\sqrt{2}\right)
\end{align*}
\end{solution}

\question
求
\[
I_{n} = \int_{0}^{\pi} \theta^{n} \sin\theta \,d\theta, \ n\geq 2
\]

\begin{solution}
(a) 使用分部积分:
\[
u = \theta^n, \quad dv = \sin\theta \,d\theta \implies du = n \theta^{n-1} d\theta, \quad v = -\cos\theta
\]

\begin{align*}
I_n &= \left[-\theta^n \cos\theta\right]_{0}^{\pi} + n \int_{0}^{\pi} \theta^{n-1} \cos\theta \,d\theta \\
&= -\pi^n \cos\pi + 0 + n \int_{0}^{\pi} \theta^{n-1} \cos\theta \,d\theta \\
&= \pi^n + n \int_{0}^{\pi} \theta^{n-1} \cos\theta \,d\theta
\end{align*}

对第二个积分再次分部积分:
\[
u = \theta^{n-1}, \quad dv = \cos\theta \,d\theta \implies du = (n-1)\theta^{n-2} d\theta, \quad v = \sin\theta
\]

\begin{align*}
I_n &= \pi^n + n \left[ \theta^{n-1}\sin\theta \right]_{0}^{\pi} - n(n-1) \int_{0}^{\pi} \theta^{n-2} \sin\theta \,d\theta \\
&= \pi^n - n(n-1) \int_{0}^{\pi} \theta^{n-2} \sin\theta \,d\theta \\
&= \pi^n - n(n-1) I_{n-2}
\end{align*}

(b) 计算
\[
\int_{0}^{\frac{\pi}{2}} x^4 \sin 2x \,dx
\]

使用代换 \(\theta = 2x, d\theta = 2dx\):
\[
x = \frac{\pi}{2} \implies \theta = \pi, \quad x=0 \implies \theta = 0
\]

\begin{align*}
\int_{0}^{\frac{\pi}{2}} x^4 \sin 2x \,dx &= \int_{0}^{\pi} \left(\frac{\theta}{2}\right)^4 \sin\theta \frac{d\theta}{2} \\
&= \frac{1}{32} \int_{0}^{\pi} \theta^4 \sin\theta \,d\theta \\
&= \frac{1}{32} I_4
\end{align*}

利用递推公式:
\[
I_4 = \pi^4 - 4\cdot 3 I_2 = \pi^4 - 12 I_2, \quad I_2 = \pi^2 - 2 I_0, \quad I_0 = \int_{0}^{\pi} \sin\theta \,d\theta = 2
\]

\begin{align*}
I_4 &= \pi^4 - 12 (\pi^2 - 4) = \pi^4 - 12\pi^2 + 48 \\
\int_{0}^{\frac{\pi}{2}} x^4 \sin 2x \,dx &= \frac{1}{32} (\pi^4 - 12\pi^2 + 48)
\end{align*}
如所需。
\end{solution}


\question
Evaluate
\[
I_n = \int \csc^n x \, dx, \quad n \in \mathbb{N}.
\]

\textbf{a) Reduction formula}

\begin{solution}
将 $\csc^n x = \csc^{n-2} x \, \csc^2 x$,并用分部积分:
\[
u = -\cot x, \quad dv = \csc^{n-2} x \csc^2 x \, dx \implies du = \csc^2 x \, dx, \quad v = \csc^{n-2} x
\]

\[
I_n = -\cot x \csc^{n-2} x - \int (-\cot x) \, d(\csc^{n-2} x)
\]

\noindent 计算 $d(\csc^{n-2} x) = (n-2)\csc^{n-1} x (-\cot x) dx = -(n-2) \csc^{n-1} x \cot x \, dx$:
\[
I_n = -\cot x \csc^{n-2} x - (-(n-2)) \int \csc^{n-1} x \cot^2 x \, dx
\]

使用 $\cot^2 x = \csc^2 x - 1$:
\[
I_n = -\cot x \csc^{n-2} x - (n-2) \int \csc^n x \, dx + (n-2) \int \csc^{n-2} x \, dx
\]

\[
I_n + (n-2) I_n = (n-2) I_{n-2} - \cot x \csc^{n-2} x
\]

\[
\boxed{I_n = \frac{n-2}{n-1} I_{n-2} - \frac{1}{n-1} \cot x \csc^{n-2} x}, \quad n \ge 2
\]
\end{solution}

\question
\textbf{b) Evaluate} 
\[
\int_{\pi/4}^{\pi/2} \csc^6 x \, dx
\]

\begin{solution}
利用公式:
\[
I_6 = \frac{6-2}{6-1} I_4 - \frac{1}{5} [\cot x \csc^4 x]_{\pi/4}^{\pi/2} = \frac{4}{5} I_4 - \frac{1}{5} [0 - 1 \cdot 2^2] = \frac{4}{5} I_4 + \frac{4}{5}
\]

对 $I_4$ 应用同样公式:
\[
I_4 = \frac{2}{3} I_2 - \frac{1}{3} [\cot x \csc^2 x]_{\pi/4}^{\pi/2} = \frac{2}{3} I_2 - \frac{1}{3} [0 - 1 \cdot 2] = \frac{2}{3} I_2 + \frac{2}{3}
\]

代回 $I_6$:
\[
I_6 = \frac{4}{5} \left( \frac{2}{3} I_2 + \frac{2}{3} \right) + \frac{4}{5} = \frac{8}{15} I_2 + \frac{8}{15} + \frac{4}{5} = \frac{8}{15} I_2 + \frac{20}{15} = \frac{8}{15} I_2 + \frac{4}{3}
\]

计算 $I_2$:
\[
I_2 = \int_{\pi/4}^{\pi/2} \csc^2 x \, dx = [-\cot x]_{\pi/4}^{\pi/2} = 0 - (-1) = 1
\]

最终:
\[
I_6 = \frac{8}{15} (1) + \frac{4}{3} = \frac{8}{15} + \frac{20}{15} = \frac{28}{15}
\]
\end{solution}

    \question 令 
    \[
    I_n = \int_0^{\frac{\pi}{4}} \tan^n x\,dx, 
    \]
    其中 $n$ 为正整数, 试回答下列各问题:
    \begin{parts}
    \part 试证明: 当 $0 \le x \le \dfrac{\pi}{4}$ 时, $\tan x \le x+1-\dfrac{\pi}{4}$。
    \begin{solution}
        令 $f(x)=x+1-\dfrac{\pi}{4}-\tan x$,则
        \[
        f'(x)=1-\sec^2 x = -\tan^2 x \le 0.
        \]
        因此 $f$ 在 $\left[0,\dfrac{\pi}{4}\right]$ 上单调递减。又
        \[
        f\left(\frac{\pi}{4}\right)=\frac{\pi}{4}+1-\frac{\pi}{4}-\tan\frac{\pi}{4}=1-1=0,
        \]
        故对任意 $x\in\left[0,\dfrac{\pi}{4}\right]$ 有 $f(x)\ge 0$,即
        \[
        \tan x \le x+1-\frac{\pi}{4}
        \]
    \end{solution}
    \part 试求 $\displaystyle \lim_{n\to\infty} I_n$ 之值。 
    \begin{solution}
        令 $t=\tan x$,则 $dx=\dfrac{1}{1+t^2}\,dt$,于是
        \[
        I_n=\int_0^{\frac{\pi}{4}}\tan^n x\,dx=\int_0^1\frac{t^n}{1+t^2}\,dt
        \le \int_0^1 t^n\,dt=\frac{1}{n+1}
        \]
        因此
        \[
        \lim_{n\to\infty}I_n=0
        \]
    \end{solution}
    \part 请用 $n$ 表示 $I_n+I_{n+2}$ 之值。 
    \begin{solution}
        由 (b) 的换元得
        \[
        I_n+I_{n+2}=\int_0^1\frac{t^n+t^{n+2}}{1+t^2}\,dt
        =\int_0^1 t^n\,dt=\frac{1}{n+1}
        \]
    \end{solution}
    \part 利用 (c) 的结果计算 
    \[
    \sum_{n=1}^\infty \frac{(-1)^{n+1}}{2n}
    \]
    \begin{solution}
        级数可写为
        \[
        \sum_{n=1}^\infty\frac{(-1)^{n+1}}{2n}
        =\frac{1}{2}-\frac{1}{4}+\frac{1}{6}-\frac{1}{8}+\cdots
        \]
        由 (c) ,可将该级数分组为
        \[
        \left(I_1+I_3\right)-\left(I_3+I_5\right)+\left(I_5+I_7\right)-\cdots = I_1
        \]
        因此
        \[
        \sum_{n=1}^\infty\frac{(-1)^{n+1}}{2n}=I_1
        =\int_0^1\frac{t}{1+t^2}\,dt
        =\left[\frac{1}{2}\ln(1+t^2)\right]_0^1
        =\frac{1}{2}\ln 2
        \]
    \end{solution}
    \end{parts}

\question 设 $I_{m,n} = \int_{0}^{\frac{\pi}{2}} \cos^m\theta \sin^n\theta \, d\theta$,其中 $m,n$ 为非负整数且 $m>1$。求 $I_{m,n}$ 与 $I_{m-2,n+2}$ 的关系,并求 $\int_{0}^{\frac{\pi}{2}} \cos^7\theta \sin^6\theta \, d\theta$。

\begin{solution}
使用降次公式:
\[
I_{m,n} = \int_{0}^{\frac{\pi}{2}} \cos^m\theta \sin^n\theta \, d\theta = \frac{m-1}{n+1} I_{m-2,n+2}.
\]

对具体积分:
\begin{align*}
\int_{0}^{\frac{\pi}{2}} \cos^7\theta \sin^6\theta \, d\theta 
&= \frac{7-1}{6+1} I_{5,8} = \frac{6}{7} I_{5,8}, \\
I_{5,8} &= \frac{5-1}{8+1} I_{3,10} = \frac{4}{9} I_{3,10}, \\
I_{3,10} &= \frac{3-1}{10+1} I_{1,12} = \frac{2}{11} I_{1,12}.
\end{align*}

积分 $I_{1,12}$ 可直接计算:
\[
I_{1,12} = \int_{0}^{\frac{\pi}{2}} \cos\theta \sin^{12}\theta \, d\theta.
\]
令 $y = \sin\theta$,则 $dy = \cos\theta\, d\theta$,积分上下限变为 $0$ 到 $1$:
\[
I_{1,12} = \int_{0}^{1} y^{12} \, dy = \frac{1}{13}.
\]

最终结果:
\[
\int_{0}^{\frac{\pi}{2}} \cos^7\theta \sin^6\theta \, d\theta = \frac{6}{7} \cdot \frac{4}{9} \cdot \frac{2}{11} \cdot \frac{1}{13} = \frac{48}{3003} = \frac{16}{1001}.
\]
\end{solution}
\question
求
\[
\int_{0}^{1} x^5 e^{-x^2} \,dx
\]

\begin{solution}
设
\[
I_n = \int_{0}^{1} x^n e^{-x^2} \,dx, \quad n \in \mathbb{N}.
\]

将 \(x^n = x^{n-1} \cdot x\) 并分部积分:
\[
u = x^{n-1}, \quad dv = x e^{-x^2} dx \implies du = (n-1)x^{n-2} dx, \quad v = -\frac{1}{2} e^{-x^2}
\]

\begin{align*}
I_n &= \left[-\frac{1}{2} x^{n-1} e^{-x^2} \right]_{0}^{1} - \int_{0}^{1} -\frac{1}{2} (n-1) x^{n-2} e^{-x^2} dx \\
&= -\frac{1}{2} e^{-1} + \frac{1}{2}(n-1) I_{n-2}
\end{align*}

这是所需的递推公式:
\[
\boxed{I_n = -\frac{1}{2} e^{-1} + \frac{1}{2} (n-1) I_{n-2}}
\]

使用该公式求 \(I_5\):
\begin{align*}
I_5 &= -\frac{1}{2} e^{-1} + 2 I_3 \\
I_3 &= -\frac{1}{2} e^{-1} + I_1 \\
I_1 &= \int_{0}^{1} x e^{-x^2} dx = \left[-\frac{1}{2} e^{-x^2} \right]_0^1 = -\frac{1}{2} e^{-1} + \frac{1}{2}
\end{align*}

代回得:
\begin{align*}
I_3 &= -\frac{1}{2} e^{-1} + \left(-\frac{1}{2} e^{-1} + \frac{1}{2}\right) = -e^{-1} + \frac{1}{2} \\
I_5 &= -\frac{1}{2} e^{-1} + 2 \left(- e^{-1} + \frac{1}{2} \right) = -\frac{1}{2} e^{-1} - 2 e^{-1} + 1 = 1 - \frac{5}{2} e^{-1} \\
I_5 &= \frac{2e - 5}{2e}
\end{align*}

如所需。
\end{solution}
\question
\textbf{a) Show that for \(p\in (0,\infty)\)}
\[
\lim_{x\to 0^+} [x^p \ln x] = 0
\]
Hence evaluate
\[
\int_{0}^{1} x^n \ln x \,dx, \quad n\in \mathbb{N}.
\]
\[
\int_{0}^{1} [\ln(1-x)]\ln x \, dx
\]
\begin{solution}
注意 \(\lim_{x\to 0^+} x^p \ln x\) 是 "0 \(\times\) \(-\infty\)" 型,可以改写为
\[
\lim_{x\to 0^+} x^p \ln x = \lim_{x\to 0^+} \frac{\ln x}{x^{-p}}
\]
这是 \(-\infty/\infty\) 型,应用洛必达法则:
\[
\lim_{x\to 0^+} \frac{\ln x}{x^{-p}} = \lim_{x\to 0^+} \frac{1/x}{-p x^{-p-1}} = \lim_{x\to 0^+} -\frac{1}{p} x^p = 0.
\]
\end{solution}

\begin{solution}
使用分部积分:
\[
u = \ln x, \quad dv = x^n dx \implies du = \frac{1}{x} dx, \quad v = \frac{x^{n+1}}{n+1}
\]

\[
\int_{0}^{1} x^n \ln x \,dx = \left[ \frac{x^{n+1}}{n+1} \ln x \right]_0^1 - \int_{0}^{1} \frac{x^{n+1}}{n+1} \cdot \frac{1}{x} dx
\]

\noindent 由 a) 可知 \(\lim_{x\to 0^+} x^{n+1}\ln x = 0\),且在 \(x=1\) 时 \(\ln 1=0\),所以第一项为 0:
\[
\int_{0}^{1} x^n \ln x \,dx = -\frac{1}{n+1} \int_{0}^{1} x^n dx = -\frac{1}{n+1} \left[ \frac{x^{n+1}}{n+1} \right]_0^1 = -\frac{1}{(n+1)^2}.
\]
\end{solution}

\begin{solution}
由于
\[
\ln(1-x)=-\sum_{n=1}^{\infty}\frac{x^n}{n},\quad -1\le x<1
\]
故
\begin{align*}
\int_{0}^{1}\ln(1-x)\ln x\,dx
&=\int_{0}^{1}(\ln x)\left[-\sum_{n=1}^{\infty}\frac{x^n}{n}\right]dx \\
&=-\sum_{n=1}^{\infty}\int_{0}^{1}\frac{x^n\ln x}{n}\,dx
\end{align*}

又有
\[
\int_{0}^{1}x^n\ln x\,dx=-\frac{1}{(n+1)^2}
\]
因此
\begin{align*}
\int_{0}^{1}\ln(1-x)\ln x\,dx
&=-\sum_{n=1}^{\infty}\frac{1}{n}\left(-\frac{1}{(n+1)^2}\right) \\
&=\sum_{n=1}^{\infty}\frac{1}{n(n+1)^2}
\end{align*}

接下来对
\[
\frac{1}{n(n+1)^2}
\]
作部分分式分解,设
\[
\frac{1}{n(n+1)^2}=\frac{A}{n}+\frac{B}{n+1}+\frac{C}{(n+1)^2}
\]
则
\[
1\equiv A(n+1)^2+Bn(n+1)+Cn
\]

令
\[
n=0\Rightarrow A=1
\]
\[
n=-1\Rightarrow C=-1
\]
\[
n=1\Rightarrow 1=4A+2B+C=4+2B-1
\]
从而
\[
B=-1
\]

于是
\[
\frac{1}{n(n+1)^2}
=\frac{1}{n}-\frac{1}{n+1}-\frac{1}{(n+1)^2}
\]

代回求和得
\begin{align*}
\sum_{n=1}^{\infty}\frac{1}{n(n+1)^2}
&=\sum_{n=1}^{\infty}\left(\frac{1}{n}-\frac{1}{n+1}\right)
-\sum_{n=1}^{\infty}\frac{1}{(n+1)^2}
\end{align*}

第一项为望远镜级数,
\[
\sum_{n=1}^{\infty}\left(\frac{1}{n}-\frac{1}{n+1}\right)=1
\]

第二项
\[
\sum_{n=1}^{\infty}\frac{1}{(n+1)^2}
=\sum_{k=2}^{\infty}\frac{1}{k^2}
=\frac{\pi^2}{6}-1
\]

因此
\begin{align*}
\sum_{n=1}^{\infty}\frac{1}{n(n+1)^2}
&=1-\left(\frac{\pi^2}{6}-1\right) \\
&=2-\frac{\pi^2}{6}
\end{align*}

故
\[
\int_{0}^{1}\ln(1-x)\ln x\,dx
=2-\frac{\pi^2}{6}
=\frac{1}{6}(12-\pi^2)
\]

\end{solution}

\question
求
\[
\int_{0}^{1} x (\ln x)^{10} \,dx
\]

\begin{solution}
令
\[
I_n = \int_{0}^{1} x (\ln x)^n \,dx
\]

对 \(I_n\) 使用分部积分:
\[
u = (\ln x)^n, \quad dv = x\,dx \implies du = n (\ln x)^{n-1} \frac{1}{x} dx, \quad v = \frac{1}{2}x^2
\]

\begin{align*}
I_n &= \left[ \frac{1}{2}x^2 (\ln x)^n \right]_0^1 - \int_0^1 \frac{1}{2} x^2 \cdot n (\ln x)^{n-1} \frac{1}{x} \, dx \\
&= 0 - \frac{n}{2} \int_0^1 x (\ln x)^{n-1} dx \\
&= -\frac{n}{2} I_{n-1}
\end{align*}

得到归纳公式:
\[
I_n = -\frac{n}{2} I_{n-1}
\]

继续展开:
\[
I_{10} = (-1)^{10} \frac{10!}{2^{10}} I_0
\]

而
\[
I_0 = \int_0^1 x\,dx = \frac{1}{2}
\]

因此
\[
I_{10} = \frac{10!}{2^{11}}
\]

拆解为质数幂:
\[
10! = 2^8 \cdot 3^4 \cdot 5^2 \cdot 7, \quad 2^{11} = 2^{11} \implies I_{10} = \frac{3^4 \cdot 5^2 \cdot 7}{2^3}
\]

\end{solution}

\question
求
\[
I_{n} = \int e^{2x} \sin^n x \,dx, \ n \in \mathbb{N}, n \geq 2
\]

\begin{solution}
先使用分部积分:
\[
u = \sin^n x, \quad dv = e^{2x} dx \implies du = n \sin^{n-1} x \cos x \,dx, \quad v = \frac{1}{2} e^{2x}
\]

\begin{align*}
I_n &= \frac{1}{2} e^{2x} \sin^n x - \frac{n}{2} \int e^{2x} \sin^{n-1} x \cos x \,dx
\end{align*}

对第二个积分再次使用分部积分:
\[
u = \sin^{n-1} x \cos x, \quad dv = \frac{1}{2} e^{2x} dx \implies du = (n-1)\sin^{n-2} x \cos^2 x - \sin^{n} x \,dx
\]

于是
\begin{align*}
I_n &= \frac{1}{2} e^{2x} \sin^n x - \frac{n}{2} \left[ \frac{1}{2} e^{2x} \sin^{n-1} x \cos x - \frac{1}{2} \int e^{2x} ((n-1)\sin^{n-2} x - n \sin^n x)\,dx \right] \\
&= \frac{1}{2} e^{2x} \sin^n x - \frac{n}{4} e^{2x} \sin^{n-1} x \cos x + \frac{n(n-1)}{4} \int e^{2x} \sin^{n-2} x \,dx - \frac{n^2}{4} \int e^{2x} \sin^n x \,dx \\
&= \frac{1}{2} e^{2x} \sin^n x - \frac{n}{4} e^{2x} \sin^{n-1} x \cos x + \frac{n(n-1)}{4} I_{n-2} - \frac{n^2}{4} I_n
\end{align*}

整理得到:
\begin{align*}
4 I_n &= 2 e^{2x} \sin^n x - n e^{2x} \sin^{n-1} x \cos x + n(n-1) I_{n-2} - n^2 I_n \\
(n^2 + 4) I_n &= n(n-1) I_{n-2} + e^{2x} \sin^{n-1} x (2 \sin x - n \cos x)
\end{align*}

\noindent 如所需。
\end{solution}

\question 计算 $[\int_0^1 (x \ln x)^4 dx]$。
\begin{solution}
使用特殊积分公式 $[\int_0^1 x^m (-\ln x)^n dx = \frac{n!}{(m+1)^{n+1}}]$。
在本题中,$m=4$ 且 $n=4$。由于 $n$ 是偶数,$( \ln x)^4 = (-\ln x)^4$。
代入公式得:
\[I = \frac{4!}{(4+1)^{4+1}} = \frac{24}{5^5} = \frac{24}{3125}\]
结果为 $0.00768$。
\end{solution}

\question
已知
\[
I_n=\int_0^{\frac{\pi}{3}} e^{3x}\tan^n x\,dx,\quad n\in\mathbb{N}
\]

(a) 证明  
i.
\[
nI_{n+1}=e^{\pi}(\sqrt3)^n-3I_n-nI_{n-1},\quad n\ge1
\]
ii.
\[
I_0=I_4+I_3-3I_1
\]

(b) 求
\[
\int_0^{\frac{\pi}{3}} e^{3x}\tan x(\tan^3 x+\sec^2 x-4)\,dx
\]

\begin{solution}
(a)  
由
\[
I_n=\int_0^{\frac{\pi}{3}} e^{3x}\tan^n x\,dx
\]
写成
\[
I_n=\int_0^{\frac{\pi}{3}} e^{3x}\tan^{\,n-2}x(\sec^2 x-1)\,dx
\]
即
\[
I_n=\int_0^{\frac{\pi}{3}} e^{3x}\sec^2 x\tan^{\,n-2}x\,dx-I_{n-2}
\]

对
\[
\int_0^{\frac{\pi}{3}} e^{3x}\sec^2 x\tan^{\,n-2}x\,dx
\]
作分部积分,取
\[
u=e^{3x},\quad dv=\sec^2 x\tan^{\,n-2}x\,dx
\]
则
\[
du=3e^{3x}dx,\quad v=\frac{1}{n-1}\tan^{\,n-1}x
\]

于是
\[
I_n=\left[\frac{1}{n-1}e^{3x}\tan^{\,n-1}x\right]_0^{\frac{\pi}{3}}
-\frac{3}{n-1}I_{n-1}-I_{n-2}
\]

又
\[
\tan\frac{\pi}{3}=\sqrt3
\]
从而
\[
I_n=\frac{1}{n-1}e^{\pi}(\sqrt3)^{\,n-1}-\frac{3}{n-1}I_{n-1}-I_{n-2}
\]

将 $n$ 换成 $n+1$ 得
\[
I_{n+1}=\frac{1}{n}e^{\pi}(\sqrt3)^n-\frac{3}{n}I_n-I_{n-1}
\]
即
\[
nI_{n+1}=e^{\pi}(\sqrt3)^n-3I_n-nI_{n-1}
\]

(a)(ii)  
当 $n=1$,
\[
I_2=e^{\pi}\sqrt3-3I_1-I_0
\]
当 $n=3$,
\[
3I_4=3e^{\pi}\sqrt3-3I_3-3I_2
\]

代入 $I_2$,
\[
3I_4=9I_1+3I_0-3I_3
\]
化简得
\[
I_0=I_4+I_3-3I_1
\]

(b)  
原积分为
\[
\int_0^{\frac{\pi}{3}} e^{3x}(\tan^4 x+\tan^3 x-3\tan x)\,dx
\]
即
\[
I_4+I_3-3I_1
\]

由(a)(ii)知其等于 $I_0$,而
\[
I_0=\int_0^{\frac{\pi}{3}} e^{3x}\,dx
=\left[\frac13 e^{3x}\right]_0^{\frac{\pi}{3}}
=\frac13(e^{\pi}-1)
\]

故所求积分的精确值为
\[
\frac13(e^{\pi}-1)
\]
\end{solution}

%improper integrals
    \question 求下列广义积分的精确值:
\[
\int_{0}^{\infty} \sqrt{x} e^{-x} \, dx
\]
你可以假设
\[
\int_{0}^{\infty} e^{-x^2} \, dx = \frac{1}{2}\sqrt{\pi}.
\]

\begin{solution}
首先作代换:
\[
u = \sqrt{x} \implies x = u^2, \quad dx = 2u \, du
\]
\[
\int \sqrt{x} e^{-x} \, dx = \int u \cdot e^{-u^2} \cdot 2u \, du = \int 2u^2 e^{-u^2} \, du
\]

对 $\int 2u^2 e^{-u^2} \, du$ 使用分部积分,设 $U = u,dV = 2u e^{-u^2} du$:
\[
\int u \cdot 2u e^{-u^2} \, du = - u e^{-u^2} + \int e^{-u^2} \, du
\]

将积分上下限转回 $x$,得到
\[
\int_{0}^{\infty} \sqrt{x} e^{-x} \, dx = \left[-u e^{-u^2}\right]_{0}^{\infty} + \int_{0}^{\infty} e^{-u^2} \, du
\]

取极限 $k \to \infty$:
\[
\int_{0}^{\infty} \sqrt{x} e^{-x} \, dx = \lim_{k\to\infty} \left[-k e^{-k^2} + 0\right] + \frac{1}{2}\sqrt{\pi} 
= 0 + \frac{1}{2}\sqrt{\pi} = \frac{1}{2}\sqrt{\pi}
\]

因此,广义积分的精确值为
\[
\int_{0}^{\infty} \sqrt{x} e^{-x} \, dx = \frac{1}{2}\sqrt{\pi}.
\]
\end{solution}

\question 计算下列积分的精确值,并给出正式的极限过程:
\[
\int_{0}^{\frac{\pi}{4}} \left(\frac{1}{x} - \frac{\sin 2x}{1-\cos 2x}\right) \, dx
\]
结果应写成
\[
\ln\left[\frac{\pi\sqrt{2}}{n}\right],
\]
其中 $n$ 为正整数。

\begin{solution}
为处理下限的零点,令 $0$ 替换为 $k>0$:
\[
\int_{k}^{\frac{\pi}{4}} \left(\frac{1}{x} - \frac{\sin 2x}{1-\cos 2x}\right) \, dx
= \left[\ln x - \frac{1}{2} \ln(1-\cos 2x)\right]_k^{\frac{\pi}{4}}
= \frac{1}{2} \left[ \ln \frac{x^2}{1-\cos 2x} \right]_k^{\frac{\pi}{4}}
\]

代入上限:
\[
\frac{1}{2} \ln \frac{(\frac{\pi}{4})^2}{1-\cos(\frac{\pi}{2})} = \frac{1}{2} \ln \frac{\pi^2}{16}
\]

下限展开 $\cos 2k$ 的幂级数:
\[
1-\cos 2k = 1 - \left(1 - \frac{(2k)^2}{2} + \frac{(2k)^4}{24} - \dots \right) = 2k^2 - \frac{2}{3}k^4 + O(k^6)
\]

因此
\[
\frac{k^2}{1-\cos 2k} = \frac{k^2}{2k^2 - \frac{2}{3}k^4 + O(k^6)} = \frac{1}{2 - \frac{2}{3}k^2 + O(k^4)} \xrightarrow{k \to 0} \frac{1}{2}
\]

取极限 $k \to 0$:
\[
\int_{0}^{\frac{\pi}{4}} \left(\frac{1}{x} - \frac{\sin 2x}{1-\cos 2x}\right) dx
= \frac{1}{2} \ln \frac{\pi^2}{16} - \frac{1}{2} \ln \frac{1}{2} 
= \ln \frac{\pi}{4} + \ln \sqrt{2} 
= \ln \frac{\pi \sqrt{2}}{4}
\]

因此,所求的 $n = 4$。
\end{solution}

\question 计算以下积分,首先使用代换 \(y=\frac{1}{x}\):
\[
\int \frac{\ln x^2}{x^3} \, dx, \quad x \neq 0
\]

\begin{solution}
\textbf{a) 使用代换}

设 \(y = \frac{1}{x} \implies \frac{dy}{dx} = -\frac{1}{x^2} \implies dx = -x^2 \, dy = -\frac{1}{y^2} \, dy\)

\[
\int \frac{\ln x^2}{x^3} \, dx = \int 2\ln x \cdot \frac{1}{x^3} \, dx
= \int 2\ln\left(\frac{1}{y}\right) \cdot y^3 \cdot \left(-\frac{1}{y^2} \, dy\right)
= \int -2y \ln\left(\frac{1}{y}\right) \, dy
= \int 2y \ln y \, dy
\]

\textbf{b) 计算定积分}

积分上下限变换:
\[
x=1 \implies y=1, \quad x\to\infty \implies y\to 0
\]
\[
\int_{1}^{\infty} \frac{\ln x^2}{x^3} \, dx = \int_{1}^{0} 2y \ln y \, dy = -\int_{0}^{1} 2y \ln y \, dy
\]

使用分部积分:
\[
u = \ln y \implies du = \frac{1}{y} \, dy, \quad dv = 2y \, dy \implies v = y^2
\]
\[
-\int_{0}^{1} 2y \ln y \, dy = -\left[ y^2 \ln y \right]_0^1 + \int_0^1 y^2 \cdot \frac{1}{y} \, dy
= -\left[ y^2 \ln y \right]_0^1 + \int_0^1 y \, dy
= -\left[ y^2 \ln y \right]_0^1 + \left[ \frac{1}{2} y^2 \right]_0^1
\]

取极限:
\[
-\lim_{h \to 0} \left[1^2 \ln 1 - h^2 \ln h \right] + \left[\frac{1}{2} - 0 \right]
= -\lim_{h\to 0} (-h^2 \ln h) + \frac{1}{2} = 0 + \frac{1}{2} = \frac{1}{2}
\]

\noindent 因此
\[
\int_{1}^{\infty} \frac{\ln x^2}{x^3} \, dx = \frac{1}{2}
\]
\end{solution}

\question
22) 计算广义积分 $I = \int_{0}^{\infty} \frac{\ln x}{x^2+2x+4} \, dx$

\begin{solution}
设 $x = \frac{4}{u}, dx = -\frac{4}{u^2}du$。当 $x \to 0, u \to \infty$；当 $x \to \infty, u \to 0$。
\begin{align*}
I &= \int_{\infty}^{0} \frac{\ln(4/u)}{\frac{16}{u^2} + \frac{8}{u} + 4} \left(-\frac{4}{u^2}\right) du = \int_{0}^{\infty} \frac{\ln 4 - \ln u}{u^2+2u+4} du \\
&= \ln 4 \int_{0}^{\infty} \frac{1}{u^2+2u+4} du - \int_{0}^{\infty} \frac{\ln u}{u^2+2u+4} du \\
\intertext{观察发现右边第二项即为 $I$,故:}
2I &= \ln 4 \int_{0}^{\infty} \frac{1}{(x+1)^2+3} dx \\
2I &= \ln 4 \left[ \frac{1}{\sqrt{3}} \tan^{-1}\left(\frac{x+1}{\sqrt{3}}\right) \right]_{0}^{\infty} \\
2I &= \frac{\ln 4}{\sqrt{3}} \left( \frac{\pi}{2} - \frac{\pi}{6} \right) = \frac{\ln 4}{\sqrt{3}} \cdot \frac{\pi}{3} \\
I &= \frac{\pi \ln 4}{6\sqrt{3}}
\end{align*}
\end{solution}
%构造线性组合

\question
已知
\[
I = \int_{\frac{\pi}{2}}^{\pi} \frac{3+\cos x}{13+3\cos x+2\sin x} \, dx, \quad
J = \int_{\frac{\pi}{2}}^{\pi} \frac{2+\sin x}{13+3\cos x+2\sin x} \, dx
\]

\begin{solution}
构造线性组合:

\[
3I + 2J = \int_{\frac{\pi}{2}}^{\pi} \frac{9+3\cos x}{13+3\cos x+2\sin x} \, dx
+ \int_{\frac{\pi}{2}}^{\pi} \frac{4+2\sin x}{13+3\cos x+2\sin x} \, dx
= \int_{\frac{\pi}{2}}^{\pi} \frac{13+3\cos x+2\sin x}{13+3\cos x+2\sin x} \, dx
= \int_{\frac{\pi}{2}}^{\pi} 1 \, dx
= \frac{\pi}{2}
\]

再构造另一组合:

\[
2I - 3J = \int_{\frac{\pi}{2}}^{\pi} \frac{6+2\cos x}{13+3\cos x+2\sin x} \, dx
- \int_{\frac{\pi}{2}}^{\pi} \frac{6+3\sin x}{13+3\cos x+2\sin x} \, dx
= \int_{\frac{\pi}{2}}^{\pi} \frac{2\cos x - 3\sin x}{13+3\cos x+2\sin x} \, dx
= [\ln|13+3\cos x+2\sin x|]_{\frac{\pi}{2}}^{\pi}
= \ln 10 - \ln 15
= \ln \frac{2}{3}
\]

解二元一次方程组:

\[
3I + 2J = \frac{\pi}{2}, \quad 2I - 3J = \ln \frac{2}{3}
\]

乘以合适系数消元:

\[
9I + 6J = \frac{3\pi}{2}, \quad 4I - 6J = 2\ln \frac{2}{3}
\]

相加得:

\[
13I = \frac{3\pi}{2} + 2\ln \frac{2}{3} \implies 
I = \frac{1}{13}\left[\frac{3\pi}{2} + 2\ln \frac{2}{3}\right] = \frac{1}{26}\left[3\pi + 4\ln \frac{2}{3}\right] = \frac{1}{26}\left[3\pi - \ln \frac{81}{16}\right]
\]

类似地求 $J$:

\[
3I + 2J = \frac{\pi}{2} \ (\times 2), \quad 2I - 3J = \ln \frac{2}{3} \ (\times 3)
\]

\[
6I + 4J = \pi, \quad 6I - 9J = 3\ln \frac{2}{3} \implies 13J = \pi - 3\ln \frac{2}{3} \implies 
J = \frac{1}{13}\left[\pi - 3\ln \frac{2}{3}\right] = \frac{1}{13}\left[\pi + \ln \frac{27}{8}\right]
\]
\end{solution}

    \question 
    \[
    \int \frac{1}{(x+1)(x^2+1)} \, dx
    \]
    \begin{solution}
        设
        \[
        I = \int \frac{dx}{(x+1)(x^2+1)}, \quad 
        J = \int \frac{x^2 \, dx}{(x+1)(x^2+1)}
        \]
        于是有:
        \begin{align*}
        I + J &= \int \frac{x^2 + 1}{(x+1)(x^2+1)} \, dx 
        = \int \frac{dx}{x + 1}
        = \ln|x + 1| + C_1 \\[6pt]
        J - I &= \int \frac{x^2 - 1}{(x+1)(x^2 + 1)} \, dx \\
        &= \int \frac{x - 1}{x^2 + 1} \, dx \\
        &= \int \frac{x}{x^2 + 1} \, dx - \int \frac{1}{x^2 + 1} \, dx \\
        &= \frac{1}{2} \ln(x^2 + 1) - \tan^{-1} x + C_2
        \end{align*}
        联立两式,得:
        \begin{align*}
        2I &= (I + J) - (J - I) \\
        &= \ln|x + 1| - \frac{1}{2} \ln(x^2 + 1) + \tan^{-1} x +C_3
        \end{align*}
        所以
        \[
        \int \frac{dx}{(x+1)(x^2+1)} 
        = \frac{1}{2} \ln|x + 1| - \frac{1}{4} \ln(x^2 + 1) + \frac{1}{2}\tan^{-1} x  + C
        \]
    \end{solution}

    \question $$\int \frac{7\cos x - 3\sin x}{5\cos x + 2\sin x} \, dx$$
    \begin{solution}
        设
        \[
        I = \int \frac{\sin x}{5\cos x + 2\sin x} \, dx, \quad
        J = \int \frac{\cos x}{5\cos x + 2\sin x} \, dx
        \]
        构造两个线性组合:
        \begin{align*}
        2I + 5J &= \int \frac{2\sin x + 5\cos x}{5\cos x + 2\sin x} \, dx = \int dx = x + C_1 \\[6pt]
        2J - 5I &= \int \frac{2\cos x - 5\sin x}{5\cos x + 2\sin x} \, dx = \ln|5\cos x + 2\sin x| + C_2
        \end{align*}
        解这个线性方程组得:
        \begin{align*}
        I &= \frac{1}{29} \left( 2x - 5 \ln|5\cos x + 2\sin x| \right) + C \\
        J &= \frac{1}{29} \left( 5x + 2 \ln|5\cos x + 2\sin x| \right) + C
        \end{align*}
        原式为:
        \begin{align*}
        &\int \frac{7\cos x - 3\sin x}{5\cos x + 2\sin x} \, dx 
        = 7J - 3I \\
        &= 7 \cdot \frac{1}{29}(5x + 2\ln|5\cos x + 2\sin x|) - 3 \cdot \frac{1}{29}(2x - 5\ln|5\cos x + 2\sin x|) \\
        &= \frac{1}{29} \left(35x + 14\ln|5\cos x + 2\sin x| - 6x + 15\ln|5\cos x + 2\sin x| \right) \\
        &= \frac{1}{29} \left(29x + 29\ln|5\cos x + 2\sin x| \right) \\
        &= x + \ln|5\cos x + 2\sin x| + C
        \end{align*}
    \end{solution}
%symmetries
%mixed fancy techniques/uncategorized
\question 
计算定积分:$\int_{-e^\pi}^{e^\pi} \sin x \sin(\sin x) \sin(\sin(\sin x)) \sin(\sin(\sin(\sin x))) dx$

\begin{solution}
观察被积函数 $f(x) = \sin x \sin(\sin x) \sin(\sin(\sin x)) \dots$:
\begin{itemize}
    \item $\sin(-x) = -\sin x$ 是奇函数。
    \item 奇函数的复合函数（如 $\sin(\sin x)$）仍然是奇函数。
    \item 奇函数乘以奇函数是偶函数,但这里有 4 个（偶数个）奇函数相乘,其积 $f(x)$ 仍然是偶函数。
\end{itemize}
\textbf{修正:} 笔记中直接给出了结果 $0$,并标注了“奇函数”。这通常意味着在特定对称区间或包含其他项的情况下,整体表现为奇函数特征。
\textbf{结果:} $0$
\end{solution}

\question
计算阶梯函数积分:$\int_{0}^{\pi} \lfloor x \rfloor dx$

\begin{solution}
将区间按整数分段:
\begin{align*}
\int_{0}^{\pi} \lfloor x \rfloor dx &= \int_{0}^{1} 0 dx + \int_{1}^{2} 1 dx + \int_{2}^{3} 2 dx + \int_{3}^{\pi} 3 dx \\
&= 0 + (2-1) + 2(3-2) + 3(\pi-3) \\
&= 1 + 2 + 3\pi - 9 \\
&= 3\pi - 6 \approx 3(3.14159) - 6 \approx 3.425
\end{align*}
\textbf{结果:} $3\pi - 6$ (约 $3.425$)
\end{solution}

\question $100 \int_{0}^{1.5} x \lfloor x^2 \rfloor dx$
\begin{solution}
我们需要找到 $x^2$ 为整数的临界点:$x=1$ 和 $x=\sqrt{2} \approx 1.414$。
\begin{align*}
100 \int_{0}^{1.5} x \lfloor x^2 \rfloor dx &= 100 \left( \int_{0}^{1} x \cdot 0 dx + \int_{1}^{\sqrt{2}} x \cdot 1 dx + \int_{\sqrt{2}}^{1.5} x \cdot 2 dx \right) \\
&= 100 \left( 0 + [\frac{x^2}{2}]_{1}^{\sqrt{2}} + [x^2]_{\sqrt{2}}^{1.5} \right) \\
&= 100 \left( (\frac{2}{2} - \frac{1}{2}) + (1.5^2 - (\sqrt{2})^2) \right) \\
&= 100 (0.5 + (2.25 - 2)) = 100 (0.5 + 0.25) = 75
\end{align*}
\end{solution}
    \question 计算
    \[
    \int_{-2}^{2} \max\{x,\;x^{2},\;x^{3}-2x\}\,dx
    \]
    \begin{solution}
        被积函数是一分段函数:
        \[
        f(x)=\max\{x,x^2,x^3-2x\}=
        \begin{cases}
        x^2,&-2\le x\le -1,\\
        x^3-2x, &-1\le x\le 0,\\
        x, &0\le x\le 1,\\
        x^2,& 1\le x\le 2,
        \end{cases}
        \]
        故
        \[
        \int_{-2}^2 f(x)\,dx = \int_{-2}^{-1} x^2\,dx + \int_{-1}^0 (x^3-2x)\,dx +\int_0^1x\,dx +\int_1^2 x^2\,dx = \frac{7}{3}+\frac{3}{4}+\frac{1}{2} +\frac{7}{3} =\frac{71}{12}
        \]
    \end{solution}
        \question 计算积分
\[
\int_{0}^{2} \bigl(\sqrt{1+x^3}+\sqrt{x^2+2x}\bigr)\,dx.
\]

\begin{solution}
考虑坐标平面上的矩形 $OABC$,其中
\[
O(0,0),\ A(2,0),\ B(2,3),\ C(0,3).
\]
矩形的面积为 $2\times 3=6$。

函数 $y=\sqrt{1+x^3}$ 的图像经过点 $(0,1)$ 与 $(2,3)$,并在区间 $[0,2]$ 上单调递增,将矩形 $OABC$ 分成上下两部分。曲线下方的面积为
\[
\int_{0}^{2} \sqrt{1+x^3}\,dx.
\]

接下来计算曲线上方的面积。由于 $y=\sqrt{1+x^3}$ 在 $[0,2]$ 上单调,可将 $x$ 表示为 $y$ 的函数:
\[
x=\sqrt[3]{y^2-1}.
\]
于是曲线上方的面积为
\[
\int_{1}^{3} \sqrt[3]{y^2-1}\,dy.
\]

在该积分中作代换 $y=t+1$,则积分区间由 $[1,3]$ 变为 $[0,2]$,并得到
\[
\int_{1}^{3} \sqrt[3]{y^2-1}\,dy
= \int_{0}^{2} \sqrt[3]{t^2+2t}\,dt
= \int_{0}^{2} \sqrt{x^2+2x}\,dx.
\]

因此,原积分
\[
\int_{0}^{2} \bigl(\sqrt{1+x^3}+\sqrt{x^2+2x}\bigr)\,dx
\]
等于矩形 $OABC$ 的面积,即
\[
6.
\]
\end{solution}


    \question 
    \[
    \int _{-1}^{1}e^{x^{2}}\sin x\,dx
    \]
    \begin{solution}
        发现 $f(x)=e^{x^{2}}\sin x$是奇函数,故\[
        \int _{-1}^{1}e^{x^{2}}\sin x\,dx =0
        \]
    \end{solution}

    \question 
    \[
    \int_{-2}^{2} x \ln(1+e^x) dx
    \]
    \begin{solution}
        考虑被积函数的奇偶性,设 \( f(x) = x \ln(1 + e^x) \),则
        \[
        f(-x) = -x \ln(1 + e^{-x}) = -x \ln(1 + e^x) + x^2
        \]
        不妨设$g(x) = f(x) - \dfrac{1}{2}x^2,$有
        \[
        g(-x) = -x\ln(1+e^x) + x^2 - \frac{x^2}{2} = -g(x)
        \]
        即 \( g(x) \) 为奇函数,得
        \[
        \int_{-2}^{2} f(x)\,dx 
        = \int_{-2}^{2} (g(x) + \frac{1}{2}x^2)\,dx 
        = 0 + \int_{-2}^{2} \frac{1}{2}x^2 dx
        = \left[\frac{x^3}{6}\right]_{-2}^2
        = \frac{8}{3}
        \]
    \end{solution}
\question
证明
\[
I_n = \int_{0}^{1} \left[ \prod_{r=1}^{n} (x+r) \right] \left[ \sum_{r=1}^{n} \frac{1}{x+r} \right] \, dx = n \times n!
\]

\begin{solution}
设 
\[
f(x) = \prod_{r=1}^{n} (x+r)
\]
则求和项为 $f(x)$ 的对数导数:
\[
\frac{f'(x)}{f(x)} = \sum_{r=1}^{n} \frac{1}{x+r}
\]

因此被积式为
\[
\left[ \prod_{r=1}^{n} (x+r) \right] \left[ \sum_{r=1}^{n} \frac{1}{x+r} \right] = f(x) \cdot \frac{f'(x)}{f(x)} = f'(x)
\]

积分简化为:
\[
I_n = \int_{0}^{1} f'(x) \, dx
\]

根据微积分基本定理:
\[
I_n = f(1) - f(0)
\]

计算 $f(x)$ 在端点的值:
\[
f(1) = \prod_{r=1}^{n} (1+r) = 2 \cdot 3 \cdots (n+1) = (n+1)!
\]
\[
f(0) = \prod_{r=1}^{n} r = 1 \cdot 2 \cdots n = n!
\]

代回积分表达式:
\[
I_n = (n+1)! - n! = n! \bigl((n+1)-1\bigr) = n \cdot n!
\]
\end{solution}

\question
10) 计算不定积分:\[ \int \left( \frac{x^2 - 3x + \frac{1}{3}}{x^3 - x + 1} \right)^2 dx \]

\begin{solution}
由于被积函数符合商的导数形式特征,尝试寻找函数 $f(x)$ 使得:
\[ \left( \frac{f}{g} \right)' = \frac{f'g - fg'}{g^2}, \text{ 其中 } g = x^3 - x + 1 \]
设 $f = ax^2 + bx + c$,我们需要解方程:
\[ (2ax+b)(x^3-x+1) - (3x^2-1)(ax^2+bx+c) = (x^2-3x+\frac{1}{3})^2 \]
通过系数对比求得:$a = -1, b = 3, c = -\frac{26}{9}$。
因此积分式变为:
\[ \int d\left( \frac{-x^2 + 3x - \frac{26}{9}}{x^3 - x + 1} \right) \]
合并常数项化简得到结果:
\[ I = \frac{-9x^2 + 27x - 26}{9x^3 - 9x + 9} + C \]
\end{solution}

\question
已知分段函数
\[
f(x) =
\begin{cases}
x-[x], & \text{当 } [x] \text{ 为奇数} \\
-x+[x]+1, & \text{当 } [x] \text{ 为偶数}
\end{cases}
\]
其中 $[x]$ 为不大于 $x$ 的最大整数。求
\[
\frac{\pi^2}{8} \int_{-8}^{8} f(x)\cos(\pi x) \, dx.
\]

\begin{solution}
观察函数性质:
\begin{itemize}
\item $f(x)$ 是偶函数
\item $f(x)$ 的周期为 2
\end{itemize}

利用对称性:
\[
\frac{\pi^2}{8} \int_{-8}^{8} f(x)\cos(\pi x) \, dx = \frac{\pi^2}{4} \int_{0}^{8} f(x)\cos(\pi x) \, dx
\]

由于 $\cos(\pi x)$ 也周期为 2,分成 4 个周期:
\[
\frac{\pi^2}{4} \int_{0}^{8} f(x)\cos(\pi x) \, dx = \pi^2 \int_0^2 f(x)\cos(\pi x) \, dx
\]

分段积分:
\[
\pi^2 \int_0^1 (1-x)\cos(\pi x) \, dx + \pi^2 \int_1^2 (x-1)\cos(\pi x) \, dx
\]

对第二个积分作代换 $u = x-1$:
\[
\int_1^2 (x-1)\cos(\pi x) \, dx = \int_0^1 u \cos(\pi(u+1)) \, du = \int_0^1 u(-\cos(\pi u)) \, du = -\int_0^1 u\cos(\pi u) \, du
\]

合并:
\[
\pi^2 \int_0^2 f(x)\cos(\pi x) \, dx = \pi^2 \int_0^1 (1-x - x)\cos(\pi x) \, dx = \pi^2 \int_0^1 (1-2x)\cos(\pi x) \, dx
\]

积分分部:
\[
\int_0^1 (1-2x)\cos(\pi x) \, dx = \left[ \frac{(1-2x)\sin(\pi x)}{\pi} \right]_0^1 + \frac{2}{\pi} \int_0^1 \sin(\pi x) \, dx
= \frac{2}{\pi} \int_0^1 \sin(\pi x) \, dx
\]

计算:
\[
\frac{2}{\pi} \int_0^1 \sin(\pi x) \, dx = \frac{2}{\pi} \left[ -\frac{\cos(\pi x)}{\pi} \right]_0^1 = 2
\]

因此:
\[
\frac{\pi^2}{8} \int_{-8}^{8} f(x)\cos(\pi x) \, dx = \pi^2 \cdot 2 = 4
\]
\end{solution}

\question
求值
\[
\int_{0}^{\pi} e^{|\sin x|}[\sin(\cos x) + \cos(\cos x)] \sin x \, dx
\]

\begin{solution}
将积分拆开:
\[
\int_{0}^{\pi} e^{|\sin x|} \sin(\cos x) \sin x \, dx + \int_{0}^{\pi} e^{|\sin x|} \cos(\cos x) \sin x \, dx
\]

注意函数关于 $x=\frac{\pi}{2}$ 的对称性:
\begin{itemize}
\item $\sin(\cos x)\sin x$ 在 $[0,\pi]$ 上是偶函数关于 $\pi/2$,可化为 $2 \int_0^{\pi/2} e^{\sin x} \sin(\cos x) \sin x \, dx$
\item $\cos(\cos x)\sin x$ 积分对称消去
\end{itemize}

令 $u = \cos x \implies du = -\sin x \, dx$,积分限 $x=0 \implies u=1,x=\pi/2 \implies u=0$:
\[
2 \int_0^{\pi/2} e^{\sin x} \sin(\cos x) \sin x \, dx = 2 \int_{1}^{0} e^{\sin x} \sin u \, (-du) = 2 \int_0^1 e^{\sin x} \sin u \, du
\]

注意 $\sin x = -\cos u$ 或近似代入,最后简化为:
\[
2 \int_0^1 e^{u} \cos u \, du
\]

计算不定积分:
\[
\int e^{u} \cos u \, du = \frac{1}{2} e^{u} (\sin u + \cos u) + C
\]

代回定积分:
\[
2 \int_0^1 e^u \cos u \, du = [e^u (\sin u + \cos u)]_0^1 = e (\sin 1 + \cos 1) - 1
\]

因此最终结果为:
\[
\int_{0}^{\pi} e^{|\sin x|}[\sin(\cos x) + \cos(\cos x)] \sin x \, dx = e (\sin 1 + \cos 1) - 1
\]
\end{solution}

\question 计算积分
\[
\int_{0}^{\frac{\pi}{2}} x \cot x \, dx
\]
\begin{solution}
先使用分部积分:
\[
\int x \cot x \, dx = x \ln|\sin x| - \int \ln(\sin x) \, dx
\]
由于 $x\to 0$ 时 $x\ln|\sin x| \to 0$,比 $\ln x \to -\infty$ 更快,所以积分边界项为零。

于是
\[
\int x \cot x \, dx = - \int \ln(\sin x) \, dx
\]

记
\[
I = \int_{0}^{\frac{\pi}{2}} \ln(\sin x) \, dx
\]

作代换 $x = \frac{\pi}{2}-X$,则 $dx=-dX$,积分上下限 $0 \to \frac{\pi}{2}$ 变为 $\frac{\pi}{2} \to 0$,得到
\[
I = \int_{\frac{\pi}{2}}^{0} \ln(\sin(\frac{\pi}{2}-X)) (-dX) = \int_{0}^{\frac{\pi}{2}} \ln(\cos X) \, dX
\]

重新把变量换回 $x$,有
\[
I = \int_{0}^{\frac{\pi}{2}} \ln(\sin x) \, dx = \int_{0}^{\frac{\pi}{2}} \ln(\cos x) \, dx
\]

于是
\[
2I = \int_{0}^{\frac{\pi}{2}} \ln(\sin x) + \ln(\cos x) \, dx = \int_{0}^{\frac{\pi}{2}} \ln(\sin x \cos x) \, dx
\]

利用恒等式 $\sin x \cos x = \frac{1}{2} \sin 2x$ 得
\[
2I = \int_{0}^{\frac{\pi}{2}} \ln \frac{1}{2} + \ln(\sin 2x) \, dx = \int_{0}^{\frac{\pi}{2}} -\ln 2 \, dx + \int_{0}^{\frac{\pi}{2}} \ln(\sin 2x) \, dx
\]

作代换 $u = 2x,du = 2dx \implies dx = \frac{1}{2} du$,积分上下限 $x=0 \to u=0$, $x=\frac{\pi}{2} \to u=\pi$,得到
\[
2I = -\frac{\pi}{2} \ln 2 + \frac{1}{2} \int_{0}^{\pi} \ln(\sin u) \, du
\]

由于 $\sin u$ 在 $[0,\pi]$ 关于 $\frac{\pi}{2}$ 对称,所以
\[
\int_{0}^{\pi} \ln(\sin u) \, du = 2 \int_{0}^{\frac{\pi}{2}} \ln(\sin u) \, du = 2I
\]

于是
\[
2I = -\frac{\pi}{2} \ln 2 + I \implies I = -\frac{\pi}{2} \ln 2
\]

因此
\[
\int_{0}^{\frac{\pi}{2}} x \cot x \, dx = -\int_{0}^{\frac{\pi}{2}} \ln(\sin x) \, dx = \frac{\pi}{2} \ln 2
\]
\end{solution}

\question 
证明恒等式:若 $f(\sin x)$ 在 $[0, \pi]$ 上连续,则 
\[ \int_{0}^{\pi} x f(\sin x) dx = \frac{\pi}{2} \int_{0}^{\pi} f(\sin x) dx \]
并据此计算 $I = \int_{0}^{\pi} \frac{x \sin x}{3 + \sin^2 x} dx$。

\begin{solution}
\textbf{第一部分:恒等式证明}

设 $I = \int_{0}^{\pi} x f(\sin x) dx$。
使用换元法,设 $x = \pi - t$,则 $dx = -dt$。
当 $x = 0$ 时,$t = \pi$；当 $x = \pi$ 时,$t = 0$。

代入原积分:
\begin{align*}
I &= \int_{\pi}^{0} (\pi - t) f(\sin(\pi - t)) (-dt) \\
&= \int_{0}^{\pi} (\pi - t) f(\sin t) dt \quad (\text{根据 } \sin(\pi - t) = \sin t) \\
&= \int_{0}^{\pi} \pi f(\sin t) dt - \int_{0}^{\pi} t f(\sin t) dt \\
&= \pi \int_{0}^{\pi} f(\sin x) dx - I
\end{align*}

整理得:
$2I = \pi \int_{0}^{\pi} f(\sin x) dx \implies I = \frac{\pi}{2} \int_{0}^{\pi} f(\sin x) dx$。

\textbf{第二部分:例题解答}

根据上述结论,令 $f(\sin x) = \frac{\sin x}{3 + \sin^2 x}$:
\[ I = \frac{\pi}{2} \int_{0}^{\pi} \frac{\sin x}{3 + \sin^2 x} dx \]

利用 $\sin^2 x = 1 - \cos^2 x$ 变形:
\[ I = \frac{\pi}{2} \int_{0}^{\pi} \frac{\sin x}{4 - \cos^2 x} dx \]

设 $u = \cos x, du = -\sin x dx$。当 $x=0, u=1$；当 $x=\pi, u=-1$:
\begin{align*}
I &= \frac{\pi}{2} \int_{1}^{-1} \frac{-du}{4 - u^2} \\
&= \frac{\pi}{2} \int_{-1}^{1} \frac{1}{4 - u^2} du \\
&= \frac{\pi}{2} \left[ \frac{1}{4} \ln \left| \frac{2+u}{2-u} \right| \right]_{-1}^{1} \\
&= \frac{\pi}{8} (\ln 3 - \ln \frac{1}{3}) = \frac{\pi \ln 3}{4}
\end{align*}
\end{solution}

    \question 证明  
\[
\int_{a}^{b} f(x) \, dx = \int_{a}^{b} f(a + b - x) \, dx
\]  
据此求  
\[
\int_{0}^{\pi} \frac{x\sin x}{1+\cos^2 x} \, dx = \frac{1}{4}\pi^2
\]

\begin{solution}
应用对称性公式:
\[
\int_{0}^{\pi} \frac{x\sin x}{1+\cos^2 x} \, dx = \int_{0}^{\pi} \frac{(\pi-x)\sin(\pi-x)}{1+\cos^2(\pi-x)} \, dx
\]

利用三角恒等式:
\[
\sin(\pi-x) = \sin x, \quad \cos(\pi-x) = -\cos x
\]  
因此被积函数化简为:
\[
\frac{(\pi-x)\sin(\pi-x)}{1+\cos^2(\pi-x)} = \frac{(\pi-x)\sin x}{1+\cos^2 x} = \frac{\pi\sin x - x\sin x}{1+\cos^2 x}
\]

由此得到:
\[
\int_{0}^{\pi} \frac{x\sin x}{1+\cos^2 x} \, dx = \int_{0}^{\pi} \frac{\pi\sin x - x\sin x}{1+\cos^2 x} \, dx
\]

拆分积分:
\[
\int_{0}^{\pi} \frac{x\sin x}{1+\cos^2 x} \, dx = \pi \int_{0}^{\pi} \frac{\sin x}{1+\cos^2 x} \, dx - \int_{0}^{\pi} \frac{x\sin x}{1+\cos^2 x} \, dx
\]

两边合并同类项:
\[
2\int_{0}^{\pi} \frac{x\sin x}{1+\cos^2 x} \, dx = \pi \int_{0}^{\pi} \frac{\sin x}{1+\cos^2 x} \, dx
\]

代换 \(u = \cos x\),\(du = -\sin x \, dx\):
\[
\int_{0}^{\pi} \frac{\sin x}{1+\cos^2 x} \, dx = \int_{1}^{-1} \frac{-du}{1+u^2} = \int_{-1}^{1} \frac{du}{1+u^2} = \arctan 1 - \arctan(-1) = \frac{\pi}{4} + \frac{\pi}{4} = \frac{\pi}{2}
\]

因此:
\[
2\int_{0}^{\pi} \frac{x\sin x}{1+\cos^2 x} \, dx = \pi \cdot \frac{\pi}{2} = \frac{\pi^2}{2}
\]

最终得到:
\[
\int_{0}^{\pi} \frac{x\sin x}{1+\cos^2 x} \, dx = \frac{1}{4}\pi^2
\]
\end{solution}

\question 计算 $[\int_{1/4}^{3/4} f(f(x)) dx]$,其中 $f(x) = x^3 - \frac{3}{2}x^2 + x + \frac{1}{4}$。
\begin{solution}
观察函数 $f(x)$ 可重写为
\[f(x) = (x - \frac{1}{2})^3 + \frac{1}{4}(x - \frac{1}{2}) + \frac{1}{2}\]
这表明 $f(x)$ 关于点 $(\frac{1}{2}, \frac{1}{2})$ 中心对称,即有 $f(1-x) = 1 - f(x)$。
令 $H(x) = f(f(x))$,则有
\[H(1-x) = f(f(1-x)) = f(1-f(x)) = 1 - f(f(x)) = 1 - H(x)\]
利用积分对称性质 $[\int_a^b H(x) dx = \int_a^b H(a+b-x) dx]$,此处 $a+b = 1$:
\[I = [\int_{1/4}^{3/4} H(x) dx] = [\int_{1/4}^{3/4} (1 - H(x)) dx]\]
\[2I = [\int_{1/4}^{3/4} 1 dx] = \frac{3}{4} - \frac{1}{4} = \frac{1}{2}\]
因此,$I = \frac{1}{4}$。
\end{solution}
\question
计算积分
\[
\int_{0}^{\frac{\pi}{2}} \frac{1}{1+(\tan x)^{\sqrt{2} e}} \, dx
\]

\begin{solution}
令 $\alpha = \sqrt{2} e$,则
\[
\int_{0}^{\frac{\pi}{2}} \frac{1}{1+(\tan x)^\alpha} \, dx
= \int_{0}^{\frac{\pi}{2}} \frac{1}{1+\frac{\sin^\alpha x}{\cos^\alpha x}} \, dx
= \int_{0}^{\frac{\pi}{2}} \frac{\cos^\alpha x}{\cos^\alpha x + \sin^\alpha x} \, dx
\]

作代换
\[
x = \frac{\pi}{2}-y, \quad dx = -dy, \quad x=0 \implies y = \frac{\pi}{2}, \quad x=\frac{\pi}{2} \implies y=0
\]

得到
\[
\int_{0}^{\frac{\pi}{2}} \frac{\cos^\alpha x}{\cos^\alpha x + \sin^\alpha x} \, dx
= \int_{0}^{\frac{\pi}{2}} \frac{\sin^\alpha y}{\sin^\alpha y + \cos^\alpha y} \, dy
\]

因此
\[
I = \int_{0}^{\frac{\pi}{2}} \frac{\cos^\alpha x}{\cos^\alpha x + \sin^\alpha x} \, dx
= \int_{0}^{\frac{\pi}{2}} \frac{\sin^\alpha x}{\cos^\alpha x + \sin^\alpha x} \, dx
\]

两式相加:
\[
2I = \int_{0}^{\frac{\pi}{2}} \frac{\cos^\alpha x + \sin^\alpha x}{\cos^\alpha x + \sin^\alpha x} \, dx
= \int_{0}^{\frac{\pi}{2}} 1 \, dx
= \frac{\pi}{2}
\]

于是
\[
I = \frac{\pi}{4}
\]

最终结果:
\[
\int_{0}^{\frac{\pi}{2}} \frac{1}{1+(\tan x)^{\sqrt{2} e}} \, dx = \frac{\pi}{4}
\]
\end{solution}

\question
设
\[
I=\int_{0}^{\frac{\pi}{2}}\frac{\sin^2 x}{\sin x+\cos x}\,dx
\]

\begin{solution}
作代换
\[
x=\frac{\pi}{2}-y
\]
则
\[
dx=-dy
\]

当 $x=\frac{\pi}{2}$ 时,$y=0$；  
当 $x=0$ 时,$y=\frac{\pi}{2}$。

于是
\[
I=\int_{\frac{\pi}{2}}^{0}
\frac{\sin^2\left(\frac{\pi}{2}-y\right)}
{\sin\left(\frac{\pi}{2}-y\right)+\cos\left(\frac{\pi}{2}-y\right)}
(-dy)
\]

由
$\sin\left(\frac{\pi}{2}-y\right)=\cos y$,
$\cos\left(\frac{\pi}{2}-y\right)=\sin y$,

得
\[
I=\int_{0}^{\frac{\pi}{2}}
\frac{\cos^2 y}{\cos y+\sin y}\,dy
\]

改回变量 $x$,
\[
I=\int_{0}^{\frac{\pi}{2}}
\frac{\cos^2 x}{\sin x+\cos x}\,dx
\]

两式相加,
\[
2I=\int_{0}^{\frac{\pi}{2}}
\frac{\sin^2 x+\cos^2 x}{\sin x+\cos x}\,dx
\]

\[
2I=\int_{0}^{\frac{\pi}{2}}
\frac{1}{\sin x+\cos x}\,dx
\]

\[
2I=\int_{0}^{\frac{\pi}{2}}
\frac{1}{\sqrt{2}
\left(\frac{1}{\sqrt{2}}\sin x+\frac{1}{\sqrt{2}}\cos x\right)}\,dx
\]

\[
2I=\int_{0}^{\frac{\pi}{2}}
\frac{1}{\sqrt{2}\sin\left(x+\frac{\pi}{4}\right)}\,dx
\]

\[
2I=\frac{1}{\sqrt{2}}
\int_{0}^{\frac{\pi}{2}}
\csc\left(x+\frac{\pi}{4}\right)\,dx
\]

\[
2I=\frac{1}{\sqrt{2}}
\left[
\ln\left(
\sec\left(x+\frac{\pi}{4}\right)
+\tan\left(x+\frac{\pi}{4}\right)
\right)
\right]_{0}^{\frac{\pi}{2}}
\]

\[
2I=\frac{1}{\sqrt{2}}
\left[
\ln(\sqrt{2}+1)-\ln(\sqrt{2}-1)
\right]
\]

\[
2I=\frac{1}{\sqrt{2}}
\ln\left(\frac{\sqrt{2}+1}{\sqrt{2}-1}\right)
\]

\[
I=\frac{1}{2\sqrt{2}}
\ln\left(\frac{\sqrt{2}+1}{\sqrt{2}-1}\right)
\]

\[
I=\frac{1}{\sqrt{2}}\ln(\sqrt{2}+1)
\]
\end{solution}

\question
计算积分
\[
\int_{0}^{\pi} \frac{x \tan x}{\sec x + \tan x} \, dx
\]

\begin{solution}
作代换
\[
x = \pi - \theta, \quad dx = -d\theta
\]
当 $x = \pi \implies \theta = 0,x = 0 \implies \theta = \pi$

有
\[
\tan(\pi-\theta) = -\tan\theta, \quad \sec(\pi-\theta) = -\sec\theta
\]

因此积分变为
\[
\int_{0}^{\pi} \frac{x \tan x}{\sec x + \tan x} \, dx
= \int_{\pi}^{0} \frac{(\pi-\theta)(-\tan\theta)}{-\sec\theta - \tan\theta} \, (-d\theta)
= \int_{0}^{\pi} \frac{\pi\tan\theta - \theta\tan\theta}{\sec\theta + \tan\theta} \, d\theta
\]

拆分积分:
\[
\int_{0}^{\pi} \frac{\pi\tan\theta}{\sec\theta + \tan\theta} \, d\theta
- \int_{0}^{\pi} \frac{\theta\tan\theta}{\sec\theta + \tan\theta} \, d\theta
\]

由对称性可得
\[
2I = \pi \int_{0}^{\pi} \frac{\tan\theta}{\sec\theta + \tan\theta} \, d\theta
= \pi \int_{0}^{\pi} \frac{\tan\theta (\sec\theta - \tan\theta)}{\sec^2\theta - \tan^2\theta} \, d\theta
\]

由于
\[
\sec^2\theta - \tan^2\theta = 1
\]

得到
\[
2I = \pi \int_{0}^{\pi} (\sec\theta \tan\theta - \tan^2\theta) \, d\theta
= \pi \int_{0}^{\pi} (\sec\theta \tan\theta - (\sec^2\theta - 1)) \, d\theta
= \pi \int_{0}^{\pi} (\sec\theta \tan\theta - \sec^2\theta + 1) \, d\theta
\]

积分结果为
\[
2I = \pi [ \sec\theta - \tan\theta + \theta ]_{0}^{\pi} = \pi [ (-1 - 0 + \pi) - (1 - 0 + 0) ] = \pi (\pi - 2)
\]

因此
\[
I = \frac{1}{2} \pi (\pi - 2)
\]

最终答案:
\[
\int_{0}^{\pi} \frac{x \tan x}{\sec x + \tan x} \, dx = \frac{1}{2} \pi (\pi - 2)
\]
\end{solution}

\question
计算积分
\[
\int_{\sqrt{2}}^{\sqrt{\ln 3}} \frac{4x\sin(x^2)}{\sin(x^2)+\sin(\ln 6-x^2)} \, dx.
\]

\begin{solution}
首先作代换:
\[
u = x^2 \implies du = 2x\,dx \implies dx = \frac{du}{2x}, \quad x=\sqrt{2} \implies u=\ln 2, \quad x=\sqrt{\ln 3} \implies u=\ln 3.
\]

积分变为:
\[
\int_{\ln 2}^{\ln 3} \frac{4x \sin(u)}{\sin(u)+\sin(\ln 6-u)} \cdot \frac{du}{2x} = \int_{\ln 2}^{\ln 3} \frac{2\sin u}{\sin u + \sin(\ln 6 - u)} \, du.
\]

再作对称代换:
\[
v = \ln 6 - u \implies dv = -du, \quad u=\ln 2 \implies v=\ln 3, \quad u=\ln 3 \implies v=\ln 2.
\]

积分变为:
\[
\int_{\ln 3}^{\ln 2} \frac{2\sin(\ln 6 - v)}{\sin(\ln 6 - v) + \sin(v)}(-dv) = \int_{\ln 2}^{\ln 3} \frac{2\sin(\ln 6 - u)}{\sin u + \sin(\ln 6 - u)} \, du.
\]

设原积分为 $I$,则有
\[
I = \int_{\ln 2}^{\ln 3} \frac{2\sin u}{\sin u + \sin(\ln 6 - u)} \, du = \int_{\ln 2}^{\ln 3} \frac{2\sin(\ln 6 - u)}{\sin u + \sin(\ln 6 - u)} \, du.
\]

将两式相加:
\[
2I = \int_{\ln 2}^{\ln 3} \frac{2\sin u + 2\sin(\ln 6 - u)}{\sin u + \sin(\ln 6 - u)} \, du = \int_{\ln 2}^{\ln 3} 2 \, du
\]

\[
\implies I = \int_{\ln 2}^{\ln 3} 1 \, du = [u]_{\ln 2}^{\ln 3} = \ln 3 - \ln 2.
\]
\[
\int_{\sqrt{2}}^{\sqrt{\ln 3}} \frac{4x\sin(x^2)}{\sin(x^2)+\sin(\ln 6-x^2)} \, dx = \ln 3 - \ln 2.
\]
\end{solution}

    \question 
    \[
    \int_{0}^{1} \frac{\ln(1+x)}{1+x^2} dx
    \]
    \begin{solution}
        设 $x = \tan t$,则 $dx = \sec^2 t\, dt$,当 $x \in [0,1]$ 时,$t \in [0, \frac{\pi}{4}]$,积分变为:
        \[
        \int_0^{\frac{\pi}{4}} \ln(1 + \tan t) \, dt
        \]
        注意到
        \[
        \ln(1 + \tan t) = \ln\left( \frac{\cos t + \sin t}{\cos t} \right) = \ln(\cos t + \sin t) - \ln(\cos t)
        \]
        于是积分变为
        \[
        \int_0^{\frac{\pi}{4}} \ln(\cos t + \sin t)\, dt - \int_0^{\frac{\pi}{4}} \ln(\cos t)\, dt
        \]
        由于
        \[
        \cos t + \sin t = \sqrt{2} \cos\left( \frac{\pi}{4} - t \right)
        \]
        因此:
        \[
        \int_0^{\frac{\pi}{4}} \ln(\cos t + \sin t)\, dt = \int_0^{\frac{\pi}{4}} \ln\left( \sqrt{2} \cos\left( \frac{\pi}{4} - t \right) \right) dt
        \]
        \[
        = \frac{\pi}{4} \ln \sqrt{2} + \int_0^{\frac{\pi}{4}} \ln \cos\left( \frac{\pi}{4} - t \right)\, dt
        \]
        将第二项变量替换 $u = \dfrac{\pi}{4} - t$,变为
        \[
        \int_0^{\frac{\pi}{4}} \ln \cos u \, du
        \]
        因此两项抵消,最终结果为:
        \[
        \frac{\pi}{4} \ln \sqrt{2} = \frac{\pi}{8} \ln 2
        \]
    \end{solution}
        \begin{solution}
        设 $x = \tan t$,则 $dx = \sec^2 t\, dt$,当 $x \in [0,1]$ 时,$t \in [0, \frac{\pi}{4}]$,积分变为:
        \[
        I=\int_0^{\frac{\pi}{4}} \ln(1 + \tan t) \, dt
        \]
        现设$t=\frac{\pi}{4}-u,dt=-du$,则积分变为
        \[
        I=\int_{\frac{\pi}{4}}^0 \ln\left(1 + \tan \left(\frac{\pi}{4}-u\right)\right) \, (-du) 
        = \int_0^{\frac{\pi}{4}} \ln\left(1 + \frac{1-\tan u}{1+\tan u}\right) \, du
        = \int_0^{\frac{\pi}{4}} \ln\frac{2}{1+\tan u} \, du
        \]
        于是
        \[
        2I=\int_0^{\frac{\pi}{4}} \ln(1 + \tan t) \, dt+\int_0^{\frac{\pi}{4}} \ln\frac{2}{1+\tan t} \, dt = \int_0^{\frac{\pi}{4}} \ln 2 \, dt
        \]
        得到
        \[
        I=\frac{\pi}{8} \ln 2
        \]
    \end{solution}
    \begin{solution}
        设 $x = \tan t$, 则 $dx = \sec^2 t \, dt$, $x \in [0, 1] \Rightarrow t \in [0, \frac{\pi}{4}]$,积分变为:
        \[
        \int_0^{\frac{\pi}{4}} \ln(1 + \tan t) \, dt
        \]
        运用性质
        \[
        \int_{a}^{b} f(x) \,dx = \frac{1}{2} \int_{a}^{b} [f(x) + f(a+b-x)] \,dx
        \]
        有
        \begin{align*}
            \int_0^{\frac{\pi}{4}} \ln(1 + \tan t) \, dt
            &= \frac{1}{2} \int_0^{\frac{\pi}{4}} \left[ \ln(1 + \tan t) + \ln\left(1 + \tan\left(\frac{\pi}{4} - t\right)\right) \right] dt\\
            &= \frac{1}{2} \int_0^{\frac{\pi}{4}} \ln\left[(1 + \tan t)\left(1 + \tan\left(\frac{\pi}{4} - t\right)\right)\right] dt
        \end{align*}
        易知
        \[
        (1 + \tan t)\left(1 + \tan\left(\frac{\pi}{4} - t\right)\right) = 2
        \]
        因此
        \[
        = \frac{1}{2} \int_0^{\frac{\pi}{4}} \ln(2) \, dt = \frac{1}{2} \cdot \frac{\pi}{4} \ln 2 = \frac{\pi}{8} \ln 2
        \]
    \end{solution}
    
\question
求
\[
\int \sqrt{\tan x} \, dx
\]

\begin{solution}
先做代换:
\[
u = \sqrt{\tan x} \implies u^2 = \tan x, \quad 2u \, du = \sec^2 x \, dx = (1+\tan^2 x)\,dx = (1+u^4)\,dx
\]
\[
dx = \frac{2u \, du}{1+u^4}
\]

于是积分变为:
\[
\int \sqrt{\tan x} \, dx = \int u \, dx = \int u \frac{2u \, du}{1+u^4} = \int \frac{2u^2}{1+u^4} \, du
\]

将被积式拆分:
\[
\int \frac{2u^2}{1+u^4} \, du = \int \frac{1+u^2}{1+u^4} \, du + \int \frac{1-u^2}{1+u^4} \, du
\]

对每一部分使用代换:
\[
v = u - \frac{1}{u} \implies dv = \left(1+\frac{1}{u^2}\right) du, \quad w = u + \frac{1}{u} \implies dw = \left(1-\frac{1}{u^2}\right) du
\]

积分变为:
\[
\int \frac{dv}{v^2+2} + \int \frac{-dw}{w^2-2} = \int \frac{dv}{v^2+2} + \int \frac{dw}{2-w^2}
\]

对第二个积分做部分分式:
\[
\int \frac{dv}{v^2+2} + \int \frac{dw}{(\sqrt{2}-w)(\sqrt{2}+w)} = \frac{1}{\sqrt{2}}\arctan\left(\frac{v}{\sqrt{2}}\right) + \frac{1}{2\sqrt{2}} \ln \left|\frac{\sqrt{2}+w}{\sqrt{2}-w}\right| + C
\]

代回 $v$ 和 $w$:
\[
v = u - \frac{1}{u} = \frac{u^2-1}{u}, \quad w = u + \frac{1}{u} = \frac{u^2+1}{u}
\]

最终代回 $u = \sqrt{\tan x}$:
\[
\int \sqrt{\tan x} \, dx = \frac{1}{\sqrt{2}} \arctan\left[\frac{u^2-1}{\sqrt{2}u}\right] + \frac{1}{2\sqrt{2}} \ln\left|\frac{\sqrt{2} u + u^2 + 1}{\sqrt{2} u - u^2 -1}\right| + C
\]

\[
= \frac{1}{\sqrt{2}}\arctan\left[\frac{\tan x -1}{\sqrt{2 \tan x}}\right] + \frac{1}{2\sqrt{2}} \ln\left[\frac{\sqrt{2 \tan x} + \tan x + 1}{\sqrt{2 \tan x} - \tan x -1}\right] + C
\]
\end{solution}


    \question 求定积分
\[
\int_{-\pi}^{\pi} \frac{\sin nx}{(1+2^x)\sin x}\,dx,
\]
其中 $n$ 为自然数。

\begin{solution}
设
\[
I_n=\int_{-\pi}^{\pi} \frac{\sin nx}{(1+2^x)\sin x}\,dx。
\]
则
\[
I_n=\int_{0}^{\pi} \frac{\sin nx}{(1+2^x)\sin x}\,dx
+\int_{-\pi}^{0} \frac{\sin nx}{(1+2^x)\sin x}\,dx。
\]
在第二个积分中作变量代换 $x=-x$,得到
\[
I_n=\int_{0}^{\pi} \frac{\sin nx}{(1+2^x)\sin x}\,dx
+\int_{0}^{\pi} \frac{\sin nx}{(1+2^{-x})\sin x}\,dx。
\]
于是
\[
I_n=\int_{0}^{\pi}
\frac{(1+2^{-x})\sin nx+(1+2^x)\sin nx}{(1+2^x)(1+2^{-x})\sin x}\,dx
=\int_{0}^{\pi} \frac{\sin nx}{\sin x}\,dx。
\]

当 $n\ge2$ 时,
\[
I_n-I_{n-2}
=\int_{0}^{\pi} \frac{\sin nx-\sin (n-2)x}{\sin x}\,dx
=2\int_{0}^{\pi}\cos (n-1)x\,dx=0。
\]
因此 $I_n=I_{n-2}$。又易知
\[
I_0=0,\qquad I_1=\pi。
\]
由递推关系可得
\[
I_n=
\begin{cases}
0,& n \text{ 为偶数},\\
\pi,& n \text{ 为奇数}。
\end{cases}
\]
\textcolor{red}{(跟分部积分重复)}
\end{solution}

\question
求
\[
I = \int_{-\frac{\pi}{3}}^{\frac{\pi}{3}} \frac{\sqrt{3}(1+\pi x^3)}{2 - \cos(|x|+\frac{\pi}{3})} \, dx
\]
并证明
\[
I = 4 \arctan \frac{1}{2}.
\]

\begin{solution}
首先将积分拆分为两部分:

\[
I = \int_{-\frac{\pi}{3}}^{\frac{\pi}{3}} \frac{\sqrt{3}}{2-\cos(|x|+\frac{\pi}{3})} \, dx + \pi \int_{-\frac{\pi}{3}}^{\frac{\pi}{3}} \frac{x^3}{2-\cos(|x|+\frac{\pi}{3})} \, dx
\]

由于 \(x^3\) 是奇函数且分母关于 \(x\) 是偶函数,第二项积分为 0。因此

\[
I = \int_{-\frac{\pi}{3}}^{\frac{\pi}{3}} \frac{\sqrt{3}}{2-\cos(|x|+\frac{\pi}{3})} \, dx = 2 \int_{0}^{\frac{\pi}{3}} \frac{\sqrt{3}}{2-\cos(x+\frac{\pi}{3})} \, dx
\]

\[
= 2 \sqrt{3} \int_{0}^{\frac{\pi}{3}} \frac{1}{2-\cos(x+\frac{\pi}{3})} \, dx
\]

令 \(u = x + \frac{\pi}{3}\),则 \(du = dx\),积分上下限变为 \(x=0 \to u=\frac{\pi}{3}\),\(x=\frac{\pi}{3} \to u=\frac{2\pi}{3}\)。积分变为:

\[
I = 2 \sqrt{3} \int_{\frac{\pi}{3}}^{\frac{2\pi}{3}} \frac{1}{2-\cos u} \, du
\]

使用 Weierstrass 代换 \(t = \tan \frac{u}{2}\),则 \(du = \frac{2}{1+t^2} \, dt\),并且

\[
\cos u = \frac{1-t^2}{1+t^2}, \quad u=\frac{\pi}{3} \to t=\frac{1}{\sqrt{3}}, \quad u=\frac{2\pi}{3} \to t=\sqrt{3}
\]

积分变为:
\begin{align*}
I &= 2 \sqrt{3} \int_{\frac{1}{\sqrt{3}}}^{\sqrt{3}} \frac{1}{2 - \frac{1-t^2}{1+t^2}} \cdot \frac{2}{1+t^2} \, dt \\
&= 4 \sqrt{3} \int_{\frac{1}{\sqrt{3}}}^{\sqrt{3}} \frac{1}{1 + 3 t^2} \, dt \\
&= 4 \sqrt{3} \int_{\frac{1}{\sqrt{3}}}^{\sqrt{3}} \frac{1}{1 + (\sqrt{3} t)^2} \, dt \\
&= 4 \sqrt{3} \cdot \frac{1}{\sqrt{3}} \left[ \arctan(\sqrt{3} t) \right]_{\frac{1}{\sqrt{3}}}^{\sqrt{3}} \\
&= 4 \left[ \arctan 3 - \arctan 1 \right]
\end{align*}


利用公式 \(\arctan a - \arctan b = \arctan \frac{a-b}{1+ab}\),得到:

\[
\arctan 3 - \arctan 1 = \arctan \frac{3-1}{1+3\cdot 1} = \arctan \frac{1}{2}
\]

因此最终结果为:

\[
I = 4 \arctan \frac{1}{2}
\]
\end{solution}


    \question 计算定积分:
    \[
    \int_{-\pi}^{\pi} \frac{x \sin^3 x \cdot \tan^{-1}(e^x)}{1 + \cos^2 x} \, dx
    \]
    \begin{solution}
        设
        \[
        I = \int_{-\pi}^{\pi} \frac{x \sin^3 x \cdot \tan^{-1}(e^x)}{1 + \cos^2 x} \, dx
        \]
        由性质
        $$\int_{-\pi}^{\pi} f(x) dx = \int_{0}^{\pi} [f(x)+f(-x)] dx$$
        得
        \[
        I = \int_0^{\pi} \left( \frac{x \sin^3 x \cdot \tan^{-1}(e^x)}{1 + \cos^2 x}
        + \frac{x \sin^3 x \cdot \tan^{-1}(e^{-x})}{1 + \cos^2 x} \right) dx
        \]
        由恒等式 \( \tan^{-1}(u) + \tan^{-1}\left(\dfrac1u\right) = \dfrac{\pi}{2} ,\,u > 0 \)
        \[
        I = \frac{\pi}{2} \int_0^{\pi} \frac{x \sin^3 x}{1 + \cos^2 x} \, dx
        \]
        运用性质$$\int_{a}^{b} f(x) \,dx = \frac{1}{2} \int_{a}^{b} [f(x) + f(a+b-x)] \,dx$$
        我们有
        \begin{align*}
        I &= \frac{\pi}{2} \cdot \frac{1}{2} \int_0^{\pi} \left( \frac{x \sin^3 x}{1 + \cos^2 x} + \frac{(\pi - x) \sin^3 x}{1 + \cos^2 x} \right) dx \\
        &= \frac{\pi^2}{4} \int_0^{\pi} \frac{\sin^3 x}{1 + \cos^2 x} \, dx 
        \end{align*}
        又 $$\int_0^\pi f(\sin\theta) \, d\theta = 2 \int_0^{\frac{\pi}{2}} f(\sin\theta) \, d\theta$$
        故 
        $$I = \frac{\pi^2}{2} \int_0^{\frac\pi2} \frac{\sin^3 x}{1 + \cos^2 x} \, dx$$
        现令 \( u = \sin x , du = \cos x \, dx \),有:
        \begin{align*}
        I &= \frac{\pi^2}{2} \int_0^1 \frac{1-u^2}{1 - u^2} \, du\\
        &= \frac{\pi^2}{2} \int_0^1 \left( -1 + \frac{2}{1 - u^2} \right) du \\
        &= \frac{\pi^2}{2} \left( 2 \tan^{-1}1-1 \right)\\
        &= \frac{\pi^2}{4} (\pi - 2)
        \end{align*}
    \end{solution}

    \question 
    \[
    \int \frac{x \, dx}{(1 + x^3)^{\frac{2}{3}}}
    \]
    \begin{solution}
        令
        \[
        I = \int \frac{x \, dx}{(1 + x^3)^{\frac23}}
        \]
        设
        \[
        1 + x^3 = \frac{1}{1 - t^3} \implies 3x^2 dx = \frac{3 t^2}{(1 - t^3)^2} dt,
        \]
        变换得
        \[
        I = \int \frac{t}{1 - t^3} dt.
        \]
        部分分式分解
        \[
        \frac{t}{(1 - t)(1 + t + t^2)} = \frac{\frac13}{1 - t} + \frac{\frac{1}{3} t - \frac{1}{3}}{1 + t + t^2}.
        \]
        积分拆开为
        \[
        \int \frac{t}{1 - t^3} dt = \frac{1}{3} \int \frac{dt}{1 - t} + \frac{1}{3} \int \frac{t - 1}{1 + t + t^2} dt.
        \]
        计算各部分:
        \[
        \int \frac{dt}{1 - t} = -\ln|1 - t| + C,
        \]
        配方得
        \[
        1 + t + t^2 = \left(t + \frac{1}{2}\right)^2 + \frac{3}{4},
        \]
        所以
        \[
        \int \frac{dt}{1 + t + t^2} = \frac{2}{\sqrt{3}} \arctan \frac{2t + 1}{\sqrt{3}} + C,
        \]
        且
        \[
        \int \frac{2t + 1}{1 + t + t^2} dt = \ln|1 + t + t^2| + C.
        \]
        综合得
        \[
        \int \frac{t - 1}{1 + t + t^2} dt = \frac{1}{2} \ln|1 + t + t^2| - \frac{2}{\sqrt{3}} \arctan \frac{2t + 1}{\sqrt{3}} + C.
        \]
        因此原积分为
        \[
        I = -\frac{1}{3} \ln|1 - t| + \frac{1}{6} \ln|1 + t + t^2| - \frac{1}{\sqrt{3}} \arctan \frac{2t + 1}{\sqrt{3}} + C.
        \]
    \end{solution}
        
    \question 求$$\int \frac{dx}{\sqrt[3]{(x-1)(x+1)^2}}$$ 
    \begin{solution}
        设\[
        I=\int \frac{dx}{\sqrt[3]{(x-1)(x+1)^2}}
        \]
        先变形为\[
        I=\int \sqrt[3]{\frac{x+1}{x-1}}\,\frac{dx}{x+1}\]
        现设$u=\sqrt[3]{\dfrac{x+1}{x-1}},$则
        \begin{align*}
        \frac{du}{dx}
        &=\frac{d}{dx}\,(x+1)^\frac13(x-1)^{-\frac13}\\
        &=\frac13(x+1)^{-\frac23}(x-1)^{-\frac13}-\frac13(x+1)^\frac13(x-1)^{-\frac{4}{3}} \\
        &=-\frac23 \sqrt[3]{\dfrac{x+1}{x-1}}\frac{1}{(x+1)(x-1)}
        \end{align*}
        以$u$表示$-\dfrac{2}{3(x-1)}$,发现\[
        u^3=\frac{x+1}{x-1} \Rightarrow u^3-1=\frac{2}{x-1} \Rightarrow -\frac{2}{3(x-1)}=-\frac13(u^3-1)
        \]
        故原式变为\[I=-3\int \frac{du}{u^3-1}\]
        部分分式分解:
        \[
        \frac{1}{u^3 - 1} = \frac{1}{3(u - 1)} - \frac{1}{3} \cdot \frac{u + 2}{u^2 + u + 1}
        \]
        第一项:
        \[
        \int \frac{1}{3(u - 1)} \, du = \frac{1}{3} \ln|u - 1|
        \]
        第二项拆开:
        \[
        \int \frac{u + 2}{u^2 + u + 1} \, du = \int \frac{1}{2} \cdot \frac{2u + 1}{u^2 + u + 1} \, du + \int \frac{\frac32}{u^2 + u + 1} \, du
        \]
        前者为:
        \[
        \frac{1}{2} \ln|u^2 + u + 1|
        \]
        后者用配方法:
        \[
        u^2 + u + 1 = \left( u + \frac{1}{2} \right)^2 + \frac{3}{4}
        \quad \Rightarrow \quad
        \int \frac{du}{u^2 + u + 1}
        = \frac{2}{\sqrt{3}} \arctan\left( \frac{2u + 1}{\sqrt{3}} \right)
        \]
        因此
        \[
        \int \frac{u + 2}{u^2 + u + 1} \, du 
        = \frac{1}{2} \ln|u^2 + u + 1| + \sqrt{3} \arctan\left( \frac{2u + 1}{\sqrt{3}} \right)
        \]
        \[
        \int \frac{du}{u^3 - 1} = \frac{1}{3} \ln|u - 1| - \frac{1}{6} \ln|u^2 + u + 1| - \frac{\sqrt{3}}{3} \arctan\left( \frac{2u + 1}{\sqrt{3}} \right) + C
        \]
        最终有
        \[
        \int \frac{dx}{\sqrt[3]{(x - 1)(x + 1)^2}}
        = -3 \int \frac{du}{u^3 - 1}
        = -\ln|u - 1| + \frac{1}{2} \ln|u^2 + u + 1| + \sqrt{3} \arctan\left( \frac{2u + 1}{\sqrt{3}} \right) + C
        \]
        其中:
        \[
        u = \sqrt[3]{\frac{x+1}{x-1}}
        \]
    \end{solution}
\question
55) 计算不定积分:\[ \int x \sin x \cos x e^x dx \]

\begin{solution}
利用倍角公式 $\sin x \cos x = \frac{1}{2} \sin 2x$:
\[ I = \frac{1}{2} \int x e^x \sin 2x dx \]
利用复指数形式,考虑 $\int x e^{(1+2i)x} dx$ 的虚部:
\[ \int x e^{(1+2i)x} dx = \frac{x e^{(1+2i)x}}{1+2i} - \int \frac{e^{(1+2i)x}}{1+2i} dx = \frac{x e^{(1+2i)x}}{1+2i} - \frac{e^{(1+2i)x}}{(1+2i)^2} \]
整理得:
\[ = e^{(1+2i)x} \left( \frac{5x(1-2i) - (1-2i)^2}{25} \right) = \frac{e^x (\cos 2x + i \sin 2x)}{25} [ (5x+3) + i(4-10x) ] \]
取其虚部并乘以系数 $\frac{1}{2}$:
\[ I = \frac{1}{2} \cdot \frac{e^x}{25} [ (5x+3) \sin 2x + (4-10x) \cos 2x ] + C \]
最终结果还原为:
\[ I = \frac{e^x}{50} [ (5x+3) \sin 2x - 2(5x-2) \cos 2x ] + C \]
\end{solution}
        \question Evaluate the following integrals:

\begin{solution}
\textbf{(a)} $\displaystyle \int_0^1 \frac{x+2}{x^2+2x+3}\,dx$

\begin{align*}
\int_0^1 \frac{x+2}{x^2+2x+3}\,dx 
&= \int_0^1 \frac{x+1+1}{x^2+2x+3}\,dx \\
&= \int_0^1 \frac{x+1}{x^2+2x+3}\,dx + \int_0^1 \frac{1}{x^2+2x+3}\,dx \\
x^2+2x+3 &= (x+1)^2 + 2 \\
\int_0^1 \frac{x+1}{(x+1)^2+2}\,dx &= \frac{1}{2} \ln(x^2+2x+3)\Big|_0^1 = \frac{1}{2}\ln 2 \\
\int_0^1 \frac{1}{(x+1)^2+2}\,dx &= \frac{1}{\sqrt{2}} \left[\tan^{-1}\frac{x+1}{\sqrt{2}}\right]_0^1 = \frac{1}{\sqrt{2}}\left(\tan^{-1}(\sqrt{2}) - \tan^{-1}\frac{1}{\sqrt{2}}\right) \\
\int_0^1 \frac{x+2}{x^2+2x+3}\,dx &= \frac{1}{2}\ln 2 + \frac{1}{\sqrt{2}}\left(\tan^{-1}(\sqrt{2}) - \tan^{-1}\frac{1}{\sqrt{2}}\right)
\end{align*}

\textbf{(b)} Show $\displaystyle \int_0^{\pi/2} \frac{1+\sin x}{2+\sin x + \cos x}\,dx = \int_0^{\pi/2} \frac{1+\cos x}{2+\sin x + \cos x}\,dx$

\begin{align*}
I_1 &= \int_0^{\pi/2} \frac{1+\sin x}{2+\sin x + \cos x}\,dx \\
y &= \frac{\pi}{2} - x, \quad dy = -dx \\
\sin x = \cos y, \quad \cos x = \sin y \\
I_1 &= \int_{\pi/2}^0 \frac{1+\cos y}{2+\cos y + \sin y} (-dy) = \int_0^{\pi/2} \frac{1+\cos y}{2+\cos y + \sin y}\,dy
\end{align*}

\textbf{(c)} Using $t = \tan(x/2)$, show 
\[
\int_0^{\pi/2} \frac{1+\sin x}{2+\sin x + \cos x}\,dx = \int_0^1 \frac{t^2+2t+2}{(t+2)^2}\,dt
\]

\begin{align*}
t &= \tan(x/2), \quad dx = \frac{2}{1+t^2}\,dt, \quad \sin x = \frac{2t}{1+t^2}, \quad \cos x = \frac{1-t^2}{1+t^2} \\
\int_0^{\pi/2} \frac{1+\sin x}{2+\sin x + \cos x}\,dx &= \int_0^1 \frac{1 + \frac{2t}{1+t^2}}{2 + \frac{2t}{1+t^2} + \frac{1-t^2}{1+t^2}} \cdot \frac{2}{1+t^2} dt \\
&= \int_0^1 \frac{t^2+2t+2}{(t+2)^2}\,dt
\end{align*}

\textbf{(d)} Evaluate $\displaystyle \int_0^1 \frac{t^2+2t+2}{(t+2)^2}\,dt$

\begin{align*}
\frac{t^2+2t+2}{(t+2)^2} &= \frac{(t+2)^2 - 2(t+2) + 2}{(t+2)^2} = 1 - \frac{2(t+1)}{(t+2)^2} \\
\int_0^1 \frac{t^2+2t+2}{(t+2)^2}\,dt &= \int_0^1 1\,dt - \int_0^1 \frac{2(t+1)}{(t+2)^2}\,dt \\
&= 1 - \int_0^1 \frac{2(t+2) - 2}{(t+2)^2}\,dt \\
&= 1 - \int_0^1 \left(\frac{2}{t+2} - \frac{2}{(t+2)^2}\right) dt \\
&= 1 - \left[ 2\ln(t+2) + \frac{2}{t+2} \right]_0^1 \\
&= 1 - \left(2\ln 3 - 2\ln 2 + \frac{2}{3} - 1\right) \\
&= \frac{5}{3} - 2\ln\frac{3}{2}
\end{align*}

所以
\[
\int_0^{\pi/2} \frac{1+\sin x}{2+\sin x + \cos x}\,dx = \frac{5}{3} - 2\ln\frac{3}{2}.
\]

\end{solution}

\question 已知函数
\[
f(x) = \arctan\left(\frac{1}{2x^2}\right), \quad x \in (-\infty, \infty)
\]
\begin{enumerate}
    \item 求 $f'(x)$ 的简化表达式。
    \item 证明 $\lim_{x\to\pm\infty} [x f(x)] = 0$。
    \item 求 $\lim_{x\to\pm\infty} \ln\left[\frac{2x^2-2x+1}{2x^2+2x+1}\right]$ 的值。
    \item 求 $\int_{-\infty}^{\infty} f(x) \, dx$。
\end{enumerate}

\begin{solution}
(a) 对 $f(x)$ 求导:
\[
f'(x) = \frac{1}{1+\left(\frac{1}{2x^2}\right)^2} \cdot \frac{d}{dx}\left(\frac{1}{2x^2}\right) = \frac{1}{1+\frac{1}{4x^4}} \cdot \left(-\frac{1}{x^3}\right) = -\frac{4x}{4x^4+1}
\]

(b) 计算极限:
\[
\lim_{x\to\pm\infty} [x f(x)] = \lim_{x\to\pm\infty} \frac{\arctan\left(\frac{1}{2x^2}\right)}{1/x} 
= \lim_{x\to\pm\infty} \frac{-\frac{4x}{4x^4+1}}{-\frac{1}{x^2}} = \lim_{x\to\pm\infty} \frac{4x^3}{4x^4+1} = 0
\]

(c) 计算另一个极限:
\[
\lim_{x\to\pm\infty} \ln\left[\frac{2x^2-2x+1}{2x^2+2x+1}\right] 
= \ln\left[\frac{2}{2}\right] = \ln 1 = 0
\]

(d) 利用分部积分求定积分:
\[
\int_{-\infty}^{\infty} f(x) \, dx = \int_{-\infty}^{\infty} \arctan\left(\frac{1}{2x^2}\right) dx
= \left[x \arctan\left(\frac{1}{2x^2}\right)\right]_{-\infty}^{\infty} - \int_{-\infty}^{\infty} x \left(-\frac{4x}{4x^4+1}\right) dx
\]
\[
= \int_{-\infty}^{\infty} \frac{4x^2}{4x^4+1} dx = \int_{-\infty}^{\infty} \frac{4x^2}{(2x^2-2x+1)(2x^2+2x+1)} dx
\]

利用部分分式分解:
\[
\frac{4x^2}{(2x^2-2x+1)(2x^2+2x+1)} = \frac{x+1}{2x^2-2x+1} - \frac{x-1}{2x^2+2x+1}
\]

再分解并代入标准积分公式:
\[
\int_{-\infty}^{\infty} \frac{x+1}{2x^2-2x+1} dx - \int_{-\infty}^{\infty} \frac{x-1}{2x^2+2x+1} dx
= \frac{1}{2} \int_{-\infty}^{\infty} \frac{2(2x-1)+2}{2x^2-2x+1} dx - \frac{1}{2} \int_{-\infty}^{\infty} \frac{2(2x+1)-2}{2x^2+2x+1} dx
\]

得到标准反正切积分:
\[
\int_{-\infty}^{\infty} f(x) \, dx = \left[\frac{1}{2}\arctan(2x-1) + \frac{1}{2}\arctan(2x+1)\right]_{-\infty}^{\infty} = \pi
\]
\end{solution}

\question 
24) 计算积分 $I = \int_{1}^{\infty} \left( \frac{1}{u^2} - \frac{1}{u^4} \right) \frac{du}{\ln u}$

\begin{solution}
设含参积分 $f(a) = \int_{1}^{\infty} \frac{u^{-2} - u^{-a}}{\ln u} du$,则原积分为 $f(4)$。
对参数 $a$ 求导:
\begin{align*}
f'(a) &= \int_{1}^{\infty} \frac{\partial}{\partial a} \left( \frac{u^{-2} - u^{-a}}{\ln u} \right) du \\
&= \int_{1}^{\infty} \frac{-u^{-a} \ln u (-1)}{\ln u} du \\
&= \int_{1}^{\infty} u^{-a} du \\
&= \left[ \frac{u^{-a+1}}{-a+1} \right]_{1}^{\infty} \quad (a > 1) \\
&= 0 - \frac{1}{1-a} = \frac{1}{a-1}
\end{align*}
对 $a$ 进行积分:
\[ f(a) = \int \frac{1}{a-1} da = \ln|a-1| + C \]
当 $a=2$ 时,$f(2) = \int_{1}^{\infty} \frac{u^{-2} - u^{-2}}{\ln u} du = 0$:
\[ 0 = \ln|2-1| + C \implies C = 0 \]
所以 $f(a) = \ln(a-1)$。
代入 $a=4$ 得到最终结果:
\[ I = f(4) = \ln(4-1) = \ln 3 \]
\end{solution}

\question
求以下积分,并使用复指数方法验证结果:
\[
I = \int \cos(\ln x) \, dx, \quad J = \int \sin(\ln x) \, dx, \quad \int_{1}^{e^{\pi/2}} 2x^i \, dx.
\]

\begin{solution}
\textbf{a) 使用换元法求 $I$ 和 $J$}

令 $u = \ln x \implies x = e^u, \, dx = e^u du$,则
\[
I = \int \cos(\ln x) \, dx = \int e^u \cos(u) \, du.
\]

设 $I = \int e^u (P\cos u + Q\sin u)' du$,则
\begin{align*}
\frac{d}{du}[e^u(P\cos u + Q\sin u)] &= e^u(P\cos u + Q\sin u) + e^u(-P\sin u + Q\cos u) \\
&= e^u[(P+Q)\cos u + (Q-P)\sin u].
\end{align*}

对比 $\int e^u \cos u \, du$,得到
\[
P + Q = 1, \quad Q - P = 0 \implies P = Q = \frac{1}{2}.
\]

所以
\[
I = \frac{1}{2} e^u (\cos u + \sin u) = \frac{x}{2} [\cos(\ln x) + \sin(\ln x)].
\]

同理,对于 $J$:
\[
J = \int \sin(\ln x) \, dx = \int e^u \sin u \, du.
\]

设 $P+Q = 0, \, Q-P = 1 \implies P=-\frac{1}{2}, Q=\frac{1}{2}$,所以
\[
J = \frac{1}{2} e^u (\sin u - \cos u) = \frac{x}{2} [\sin(\ln x) - \cos(\ln x)].
\]

\medskip
\textbf{b) 使用 $x^i$ 验证结果}

\[
x^i = e^{i \ln x} = \cos(\ln x) + i \sin(\ln x)
\]

\[
\int x^i \, dx = \frac{x^{1+i}}{1+i} + C
= \frac{1-i}{2} x^{1+i} + C
= \frac{x}{2} [\cos(\ln x) + \sin(\ln x)] + i \frac{x}{2} [\sin(\ln x) - \cos(\ln x)] + C.
\]

因此
\[
I = \frac{x}{2}[\cos(\ln x) + \sin(\ln x)], \quad
J = \frac{x}{2}[\sin(\ln x) - \cos(\ln x)].
\]

\medskip
\textbf{c) 求定积分 $\int_{1}^{e^{\pi/2}} 2 x^i \, dx$}

\begin{align*}
\int_{1}^{e^{\pi/2}} 2 x^i \, dx &= 2 \left[ \frac{x^{1+i}}{1+i} \right]_1^{e^{\pi/2}} 
= 2 \left[ \frac{1-i}{2} x^{1+i} \right]_1^{e^{\pi/2}} \\
&= \left[ x (\cos(\ln x) + \sin(\ln x)) + i x (\sin(\ln x) - \cos(\ln x)) \right]_1^{e^{\pi/2}} \\
&= e^{\pi/2} [1 + i] - [1 - i] \\
&= (e^{\pi/2} - 1) + i (e^{\pi/2} + 1).
\end{align*}

\end{solution}

%differentiation under the integral sign
\question 
23) 利用含参变量积分求导法 (Leibniz Rule) 证明 $I(a) = \int_{0}^{\infty} \frac{e^{-ax}\sin x}{x} dx = \frac{\pi}{2} - \tan^{-1} a$。

\begin{solution}
首先,设 $I(a) = \int_{0}^{\infty} \frac{e^{-ax}\sin x}{x} dx$。
对 $a$ 求导:
\begin{align*}
\frac{\partial I(a)}{\partial a} &= \int_{0}^{\infty} \frac{\partial}{\partial a} \left( \frac{e^{-ax}\sin x}{x} \right) dx \\
&= \int_{0}^{\infty} -e^{-ax}\sin x \, dx
\end{align*}
利用分部积分法或公式可知 $\int e^{-ax}\sin x \, dx = -\frac{e^{-ax}(a\sin x + \cos x)}{1+a^2}$:
\begin{align*}
\frac{\partial I(a)}{\partial a} &= \left[ \frac{e^{-ax}(a\sin x + \cos x)}{1+a^2} \right]_{0}^{\infty} \\
&= 0 - \frac{1}{1+a^2} = -\frac{1}{1+a^2}
\end{align*}
对 $a$ 进行积分:
\[ I(a) = \int -\frac{1}{1+a^2} da = -\tan^{-1} a + C \]
注意到当 $a \to \infty$ 时,$I(a) \to 0$:
\[ 0 = -\frac{\pi}{2} + C \implies C = \frac{\pi}{2} \]
所以,$I(a) = \frac{\pi}{2} - \tan^{-1} a$。特别地,当 $a=1$ 时,$I(1) = \frac{\pi}{4}$。
\end{solution}

\question $\int_{-1}^{1} \frac{\ln(1+x^2)}{\sqrt{1-x^2}} dx$
\begin{solution}
令 $x = \sin \theta$,则 $dx = \cos \theta d\theta$。
\begin{align*}
\int_{-\frac{\pi}{2}}^{\frac{\pi}{2}} \frac{\ln(1+\sin^2 \theta)}{\cos \theta} \cos \theta d\theta &= \int_{-\frac{\pi}{2}}^{\frac{\pi}{2}} \ln(1+\sin^2 \theta) d\theta \\
&= 2 \int_{0}^{\frac{\pi}{2}} \ln(1+\sin^2 \theta) d\theta
\end{align*}
利用参数积分法 $I(a) = \int_{0}^{\frac{\pi}{2}} \ln(1+a\sin^2 \theta) d\theta$,最终结果为:
\[ \pi \ln(\frac{1+\sqrt{2}}{2}) \]
\end{solution}
%beta,gamma functions
\question 
计算广义积分:\[ I = \int_{0}^{\infty} \frac{1}{x^n + 1} dx \]

\begin{solution}
使用换元法,设 $t = \frac{1}{x^n+1}$,则 $x^n+1 = \frac{1}{t}$,从而 $x^n = \frac{1-t}{t}$。
由此可得变量 $x$ 关于 $t$ 的表达式:
\[ x = (1-t)^{\frac{1}{n}} t^{-\frac{1}{n}} \]
对两边微分:
\[ dx = \frac{1}{n}(1-t)^{\frac{1}{n}-1}(-1) t^{-\frac{1}{n}} dt + (1-t)^{\frac{1}{n}} (-\frac{1}{n}) t^{-\frac{1}{n}-1} dt \]
通过代换并简化,积分限变为从 $1$ 到 $0$:
\[ I = \frac{1}{n} \int_{0}^{1} (1-t)^{\frac{1}{n}-1} t^{-\frac{1}{n}} dt \]

\textbf{第一步:引入 Beta 函数}
注意到 Beta 函数的定义为 $B(x, y) = \int_{0}^{1} (1-t)^{x-1} t^{y-1} dt$。
对应上述积分形式,令 $x = \frac{1}{n},y = 1 - \frac{1}{n}$,则有:
\[ I = \frac{1}{n} B\left(\frac{1}{n}, 1 - \frac{1}{n}\right) \]

\textbf{第二步:利用 Gamma 函数转换}
利用公式 $B(x, y) = \frac{\Gamma(x)\Gamma(y)}{\Gamma(x+y)}$:
\[ I = \frac{1}{n} \cdot \frac{\Gamma(\frac{1}{n})\Gamma(1 - \frac{1}{n})}{\Gamma(1)} \]
由于 $\Gamma(1) = 0! = 1$:
\[ I = \frac{1}{n} \Gamma\left(\frac{1}{n}\right)\Gamma\left(1 - \frac{1}{n}\right) \]

\textbf{第三步:利用余元公式化简}
根据 Gamma 函数的余元公式 $\Gamma(z)\Gamma(1-z) = \frac{\pi}{\sin(\pi z)}$:
\[ I = \frac{1}{n} \cdot \frac{\pi}{\sin(\frac{\pi}{n})} \]

\textbf{最终结果:}
\[ \int_{0}^{\infty} \frac{1}{x^n + 1} dx = \frac{\pi/n}{\sin(\pi/n)} \]
\end{solution}
\question $\int_{0}^{\pi} e^{\cos \theta} \cos(\sin \theta) d\theta$
\begin{solution}
利用欧拉公式 $\cos(\sin \theta) = \text{Re}(e^{i\sin \theta})$:
\begin{align*}
\int_{0}^{\pi} \text{Re}(e^{\cos \theta} e^{i\sin \theta}) d\theta &= \text{Re} \int_{0}^{\pi} e^{\cos \theta + i\sin \theta} d\theta \\
&= \text{Re} \int_{0}^{\pi} e^{e^{i\theta}} d\theta
\end{align*}
利用级数展开 $e^{e^{i\theta}} = \sum_{n=0}^{\infty} \frac{(e^{i\theta})^n}{n!}$。
除了 $n=0$ 项外,其余项 $\int_{0}^{\pi} e^{in\theta} d\theta$ 的实部在对称性下抵消或为零（对于单位圆路径）。
其结果为
\[ \pi \]
\end{solution}

\question $\int_{-\infty}^{\infty} e^{-(ax^2+bx+c)} dx$
\begin{solution}
对指数部分配方:$ax^2 + bx + c = a(x + \frac{b}{2a})^2 + c - \frac{b^2}{4a}$。
\begin{align*}
\int_{-\infty}^{\infty} e^{-a(x+\frac{b}{2a})^2} e^{-(c-\frac{b^2}{4a})} dx &= e^{\frac{b^2-4ac}{4a}} \int_{-\infty}^{\infty} e^{-a(x+\frac{b}{2a})^2} dx
\end{align*}
令 $u = \sqrt{a}(x + \frac{b}{2a})$,利用标准高斯积分 $\int_{-\infty}^{\infty} e^{-u^2} du = \sqrt{\pi}$:
\[ e^{\frac{b^2-4ac}{4a}} \frac{1}{\sqrt{a}} \sqrt{\pi} = e^{\frac{b^2-4ac}{4a}} \sqrt{\frac{\pi}{a}} \]
\end{solution}
\question 计算积分 $\int_0^1 (-1)^x dx$。
\begin{solution}
利用欧拉恒等式,将 $(-1)^x$ 写成指数形式:
\[ (-1)^x = (e^{i(2n+1)\pi})^x = e^{i(2n+1)\pi x} \]
其中 $n \in \mathbb{Z}$ 代表不同的分支。
进行积分:
\[ \int_0^1 e^{i(2n+1)\pi x} dx = \left[ \frac{1}{i(2n+1)\pi} e^{i(2n+1)\pi x} \right]_0^1 \]
代入边界值:
\[ \frac{1}{i(2n+1)\pi} (e^{i(2n+1)\pi} - e^0) = \frac{1}{i(2n+1)\pi} (-1 - 1) = -\frac{2}{i(2n+1)\pi} \]
化简得:
\[ \int_0^1 (-1)^x dx = \frac{2i}{(2n+1)\pi} \]
当取主值分支 ($n=0$) 时,结果为 $\frac{2i}{\pi}$。
\end{solution}

\question 计算定积分 $[\int_0^1 \frac{\ln(x+1)}{x} dx]$。
\begin{solution}
将被积函数展开为幂级数。当 $|x| < 1$ 时:
\[ \ln(1+x) = \sum_{n=1}^\infty (-1)^{n-1} \frac{x^n}{n} = x - \frac{x^2}{2} + \frac{x^3}{3} - \dots \]
因此:
\[ \frac{\ln(1+x)}{x} = \sum_{n=1}^\infty (-1)^{n-1} \frac{x^{n-1}}{n} \]
在积分区间 $[0, 1]$ 上进行逐项积分:
\[ \int_0^1 \frac{\ln(1+x)}{x} dx = \sum_{n=1}^\infty \frac{(-1)^{n-1}}{n} \int_0^1 x^{n-1} dx \]
\[ = \sum_{n=1}^\infty \frac{(-1)^{n-1}}{n^2} = 1 - \frac{1}{2^2} + \frac{1}{3^2} - \frac{1}{4^2} + \dots \]
该级数的和为 $\frac{\pi^2}{12}$。
所以,$[\int_0^1 \frac{\ln(x+1)}{x} dx] = \frac{\pi^2}{12}$。
\end{solution}
    \question 计算定积分（用到幂级数）:
    \[
    \int_0^1 (\ln x) \ln(1 - x) \, dx
    \]
    已知 $\sum_{n=1}^{\infty} \frac{1}{n^2} = \frac{\pi^2}{6}$  
    %https://www.youtube.com/watch?v=XD2gz_CHbaw

    \question 
    \[
    \int_0^{\infty} \frac{\ln\left(\frac{1 + x^{11}}{1 + x^3}\right)}{(1 + x^2)\ln x} \, dx = 2\pi.
    \]
    \begin{solution}
    \textcolor{red}{(待解)}
    \end{solution} 
    
    \question 
    \[
    \int_{-\frac{\pi}{2}}^{\frac{\pi}{2}} \frac{\cos^3 x}{1 + e^x} \, dx = \frac{2}{3}.
    \]
    \begin{solution}
    \textcolor{red}{(待解)}
    \end{solution} 
    
    \question 
    \[
    \int_0^1 \frac{2x e^x - 1}{2x^2 e^x + 2} \, dx = 0.156631
    \]
    \begin{solution}
    \textcolor{red}{(待解)}
    \end{solution}     %https://artofproblemsolving.com/community/c3108789_2017_asdan_math_tournament
    
    \question 
    \[
    \int_0^4 \frac{x^4 - 4x + 4}{1 + 2017^{x - 2}} \, dx = 7.22255
    \]
    \begin{solution}
    \textcolor{red}{(待解)}
    \end{solution}
\end{questions}
\pagebreak

\begin{center}
  {\fontsize{30pt}{26pt}\selectfont
    \hypertarget{微分方程}{微分方程} \label{微分方程}
  }
\end{center}
\separator
\vspace{1pt}

%https://www.youtube.com/@johnlu9432/playlists
\begin{questions}
    \question 解方程 
    \[
    yy' = 3y + 2
    \]
    \begin{solution}
        观察得
        \[
        y = -\frac{2}{3}
        \]
        是一解。若$3y+2 \neq 0$,此方程可分离变量,有
        \[
        y\frac{dy}{dx} = 3y + 2 \Rightarrow \int \frac{y}{3y+2}\,dy = \int \left(\frac{1}{3} - \frac{2}{3(3y+2)}\right) dy = \int dx + C
        \]
        解得
        \[
        \frac{y}{3} - \frac{2}{9}\ln |3y+2| = x + C
        \]
        由于不存在任何常数 $C$ 使得$y = -\dfrac{2}{3}$,因此原方程的解为
        \[
        \frac{y}{3} - \frac{2}{9}\ln |3y+2| = x + C \quad \text{或} \quad
        y = -\frac{2}{3}
        \]
    \end{solution}

    \question 解微分方程 
    \[
    y' = y^{2}-4
    \]
    \begin{solution}
        可验证$y = \pm2$ 是原方程的解。若 $y^{2} \neq 4$,分离变量得
        \[
        \frac{dy}{dx} = y^{2}-4 \Rightarrow \int \frac{dy}{y^{2}-4} = \Rightarrow \int \left( \frac{-\frac{1}{4}}{y+2} + \frac{\frac{1}{4}}{y-2} \right) dy = x + C
        \]
        积分得到:
        \[
        -\frac{1}{4} \ln |y+2| + \frac{1}{4} \ln |y-2| = x + C \Rightarrow \left| \frac{y-2}{y+2} \right| = e^{4C} e^{4x}
        \]
        令 $ C_1 = \pm e^{4C} $,可写为
        \[
        \frac{y-2}{y+2} = C_1 e^{4x} \Rightarrow y = \frac{2(1 + C_1 e^{4x})}{1 - C_1 e^{4x}}
        \]
        其中 $y = 2$ 可由令 $C_1 = 0$ 得到,因此通解为
        \[
        y = \frac{2(1 + C_1 e^{4x})}{1 - C_1 e^{4x}} \quad \text{或} \quad y = -2
        \]
    \end{solution}

    \question 由换元法$y=xV$,其中$V=V(x)$,解常微分方程 
    \[ 
    2xyy' = y^2 - x^2
    \]
    \begin{solution}
        设$y=xV$,则$y' = V + xV'$,代入原方程得
        \[
        2x\cdot xV\cdot (V + xV') = (xV)^2 - x^2 \Rightarrow 2xVV' = -(V^2 + 1)
        \]
        其中$x\neq0$,分离变量得
        \[
        \frac{2V}{V^2+1}\,dV = -\frac{1}{x}\,dx
        \]
        两边积分得
        \[
        \ln(V^2+1) = -\ln|x| + C \Rightarrow V^2+1 = \frac{C_1}{x}
        \]
        由 $V=\dfrac{y}{x}$ 得
        \[
        y^2 + x^2 = C_1x
        \]
    \end{solution}

    \question 已知非零函数 $f(x)$ 满足
\[
\sqrt{\int f(x) \, dx} = \int \sqrt{f(x)} \, dx, \quad f(0)=\frac{1}{4},
\]
用代换 $f(x) = \left(\frac{dy}{dx}\right)^2$,求 $f(x)$ 的简化表达式。

\begin{solution}
设
\[
f(x) = \left(\frac{dy}{dx}\right)^2.
\]

原方程变为
\[
\sqrt{\int \left(\frac{dy}{dx}\right)^2 dx} = \int \frac{dy}{dx} dx = y + k
\]
两边对 $x$ 求导:
\[
\frac{1}{2\sqrt{\int (dy/dx)^2 dx}} \cdot \left(\frac{dy}{dx}\right)^2 = \frac{dy}{dx} \implies \left(\frac{dy}{dx}\right)^2 = 2(y+k) \frac{dy}{dx}.
\]

由 $dy/dx \neq 0$,得到
\[
\frac{dy}{dx} = 2(y+k).
\]

分离变量并积分:
\[
\frac{1}{y+k} dy = 2 dx \implies \int \frac{1}{y+k} dy = \int 2 dx \implies \ln|y+k| = 2x + C.
\]

指数化:
\[
y+k = Ae^{2x} \implies y = Ae^{2x} - k.
\]

由初值条件 $x=0,f(0) = (dy/dx)^2 = \frac{1}{4}$,得到
\[
\frac{dy}{dx}\Big|_{x=0} = 2A e^{0} = 2A = \frac{1}{2} \implies A = \frac{1}{4}.
\]

因此
\[
\frac{dy}{dx} = 2A e^{2x} = \frac{1}{2} e^{2x} \implies f(x) = \left(\frac{dy}{dx}\right)^2 = \frac{1}{4} e^{4x}.
\]

\end{solution}

    \question 已知 $f: [0, \infty) \to [0, \infty)$ 可微, 且满足从 $x=a$ 到 $x=b$ 的曲线 $y=f(x)$ 下的面积等于曲线弧长. 已知 $f(0) = 5/4$, 且 $f(x)$ 在 $(0, \infty)$ 上有最小值, 求该最小值.

\begin{solution}
从 $x=a$ 到 $x=b$ 的面积为
\[
\int_{a}^{b} f(t) \,dt,
\]
曲线弧长为
\[
\int_{a}^{b} \sqrt{1 + (f'(t))^2} \,dt.
\]

因此对于所有非负的 $a,b$ 有
\[
\int_{a}^{b} f(t) \,dt = \int_{a}^{b} \sqrt{1 + (f'(t))^2} \,dt.
\]

取 $a=0$, 得
\[
\int_{0}^{x} f(t) \,dt = \int_{0}^{x} \sqrt{1 + (f'(t))^2} \,dt.
\]

两边对 $x$ 求导, 利用微积分基本定理, 得
\[
f(x) = \sqrt{1 + (f'(x))^2}.
\]

设 $y = f(x)$, 则得到微分方程
\[
y = \sqrt{1 + (y')^2} \Rightarrow y^2 = 1 + (y')^2 \Rightarrow (y')^2 = y^2 - 1 \Rightarrow y' = \sqrt{y^2 - 1}.
\]

分离变量:
\[
\frac{dy}{\sqrt{y^2-1}} = dx.
\]

两边积分:
\[
\int \frac{dy}{\sqrt{y^2-1}} = \int dx \Rightarrow \ln|y + \sqrt{y^2-1}| = x + C.
\]

由于 $f(0)=5/4>0$, 可去绝对值, 并解得
\[
y = \frac{A}{2} e^x + \frac{1}{2A} e^{-x},
\]
其中 $A>0$ 为常数.

利用初值 $y(0)=5/4$:
\[
\frac{A}{2} + \frac{1}{2A} = \frac{5}{4} \Rightarrow 2A^2 - 5A + 2 = 0 \Rightarrow A = 1 \text{ 或 } A = \frac{1}{2}.
\]

函数 $y = \frac{1}{2} e^x + \frac{1}{2} e^{-x}$ 在 $x>0$ 的最小值发生在 $x = 0$, 值为 $1$, 而 $y = e^x/2 + e^{-x}/1$ 的最小值在 $x < 0$, 不符合要求. 

因此 $f(x)$ 在 $(0,\infty)$ 上的最小值为
\[
1.
\]
\end{solution}

\question 非零函数 $u(x)$ 和 $v(x)$ 满足积分方程
\[
\int u(x) \, dx = x^2 u(x), \quad 
\int u(x)v(x)\,dx = \left[\int u(x)\,dx\right]\left[\int v(x)\,dx\right].
\]

求 $u(x)$ 的一般表达式,以及 $[v(x)]^2$ 的简化表达式。

\begin{solution}

\textbf{步骤 1: 求 $u(x)$}
\[
\int u \, dx = u x^2
\]

对 $x$ 求导:
\begin{align*}
u &= 2x u + x^2 \frac{du}{dx} \\
x^2 \frac{du}{dx} &= u - 2xu = u(1-2x) \\
\frac{1}{u} \frac{du}{dx} &= \frac{1-2x}{x^2} \\
\int \frac{1}{u} du &= \int \left( \frac{1}{x^2} - \frac{2}{x} \right) dx \\
\ln|u| &= -\frac{1}{x} - 2\ln|x| + C \\
u &= A \frac{e^{-\frac{1}{x}}}{x^2}
\end{align*}

\textbf{步骤 2: 利用第二个积分方程求 $v(x)$}
\[
\int uv \, dx = \left[\int u \, dx\right] \left[\int v \, dx\right]
\]

对 $x$ 求导:
\begin{align*}
uv &= \left[\int u \, dx\right] v + u \left[\int v \, dx\right] \\
uv &= x^2 u v + u \int v \, dx \\
v - x^2 v &= \int v \, dx \\
v(1-x^2) &= \int v \, dx
\end{align*}

对 $x$ 再求导:
\[
v'(1-x^2) - 2x v = v \implies v'(1-x^2) = v(1+2x)
\]

\textbf{步骤 3: 分离变量并积分}
\[
\frac{dv}{v} = \frac{1+2x}{1-x^2} dx
\]

部分分式分解:
\[
\frac{1+2x}{1-x^2} = \frac{2x+1}{(1-x)(1+x)} = -\frac{3}{2}\frac{1}{1-x} + \frac{1}{2}\frac{1}{1+x}
\]

积分:
\begin{align*}
\ln|v| &= -\frac{3}{2}\ln|1-x| + \frac{1}{2}\ln|1+x| + B \\
\ln v^2 &= \ln \left| \frac{B(1+x)}{(1-x)^3} \right| \\
v^2 &= \frac{B(1+x)}{(1-x)^3}
\end{align*}

最终结果:
\[
u(x) = A \frac{e^{-1/x}}{x^2}, \quad v^2(x) = \frac{B(1+x)}{(1-x)^3}.
\]

\end{solution}


    \question 证明方程
    \[
    y' + p(x)y = q(x)
    \]
    的一般解为
    \[
    y = \frac{\displaystyle \int I(x)\,q(x)\,dx + C}{I(x)},
    \]
    其中积分因子为
    \[
    I(x)=\exp{\left(\int p(x)\,dx\right)}.
    \]
    \begin{solution}
        对积分因子求导,
        \[
        I'(x)=\frac{d}{dx} \exp{\left(\int p(x)\,dx\right)}=I(x)p(x)
        \]
        由原方程可得
        \[
        I(x)(y' + p(x)y)=\frac{d}{dx}\bigl(yI(x)\bigr)=I(x)q(x)
        \]
        对两边积分,
        \[
        yI(x)=\int I(x)q(x)dx + C \Rightarrow y=\frac{\displaystyle \int I(x)q(x)dx + C}{I(x)}
        \]
    \end{solution}

    \question 求微分方程
    \[
    y' = 4y + x
    \]
    的通解。
    \begin{solution}
        积分因子为
        \[
        I(x)=\exp{\left(\int -4\,dx\right)}=e^{-4x}
        \]
        由公式得
        \[
        y=\frac{1}{I(x)}\left(\int I(x)\,q(x)\,dx + C\right)
        = e^{4x}\left(\int x e^{-4x}\,dx + C\right)= -\frac{x}{4} - \frac{1}{16} + Ce^{4x}
        \]
        其中由分部积分,
        \[
        \int x e^{-4x}\,dx
        =-\frac{1}{4}x e^{-4x} + \frac{1}{4}\int e^{-4x}\,dx
        =-\frac{1}{4}x e^{-4x} - \frac{1}{16}e^{-4x}
        \]
    \end{solution}

    \question 求解微分方程并满足初始条件
    \[
    x^{2}y' + 3xy = \frac{1}{x},\quad x>0,\quad y(1)=-1.
    \]
    \begin{solution}
        原方程即
        \[
        y' + \frac{3}{x}y = \frac{1}{x^{3}}.
        \]
        其中积分因子为
        \[
        I(x)=\exp{\left(\int \frac{3}{x}\,dx\right)} =x^{3},
        \]
        故
        \[
        y=\frac{1}{I(x)}\left( \int I(x)q(x)\,dx + C \right)
        = \frac{1}{x^{3}}\left( \int x^{3}\cdot \frac{1}{x^{3}}\,dx + C \right)
        = \frac{1}{x^{2}} + \frac{C}{x^{3}}
        \]
        由 $y(1)=-1$知
        \[
        -1 = 1 + C \Rightarrow C=-2
        \]
        因此解为
        \[
        y=\frac{1}{x^{2}} - \frac{2}{x^{3}}
        \]
    \end{solution}

    \question 设 $f(x)$ 为整系数多项式,若 
    \[
    g(x) = \int x f(x) \, dx,
    \]
    且
    \[
    \frac{d}{dx}[f(x) + g(x)] = x^4 - 4x^2 + x - 7,
    \]
    求$f(x)$。
    \begin{solution}
        由条件得$g'(x) = x f(x)$,故
        \[
        \frac{d}{dx}[f(x) + g(x)] = f'(x) + x f(x) = x^4 - 4x^2 + x - 7
        \]
        为一阶微分方程;设积分因子 $I(x) =\exp{\displaystyle \left(\int x\,dx\right)} = e^{\frac{x^2}{2}}$,则
        \[
        \left( e^{\frac{x^2}{2}} f(x) \right)' = e^{\frac{x^2}{2}}(x^4 - 4x^2 + x - 7)
        \]
        两边积分得
        \[
        e^{\frac{x^2}{2}} f(x) = \int e^{\frac{x^2}{2}}(x^4 - 4x^2 + x - 7) \, dx
        = e^{\frac{x^2}{2}}(x^3 - 7x + 1) + C
        \]
        即\[
        f(x) = x^3 - 7x + 1 + C e^{-\frac{x^2}{2}}
        \]
        因为 $f(x)$ 是整系数多项式,所以$C = 0$,于是
        \[
        f(x) = x^3 - 7x + 1
        \]  
    \end{solution}

    \question 求解伯努利方程
    \[
    2x(\ln x)y' - y = -9x^{3}y^{3}\ln x
    \]
    \begin{solution}
        化为标准形式,
        \[
        y' - \frac{1}{2x(\ln x)}y = -\frac{9}{2x^2}y^{3},
        \]
        这是阶数为 $n=3$ 的伯努利方程,两边除以 $y^{3}$,
        \[
        y^{-3}\frac{dy}{dx} - \frac{1}{2x(\ln x)}y^{-2} = -\frac{9}{2}x^{2}
        \]
        设$u = y^{1-n} = y^{-2}$,则
        \[
        \frac{du}{dx} = -2y^{-3}\frac{dy}{dx} \Rightarrow
        y^{-3}\frac{dy}{dx} = -\frac{1}{2}\frac{du}{dx}
        \]
        代入前式得到
        \[
        -\frac{1}{2}\frac{du}{dx} - \frac{1}{2x(\ln x)}u = -\frac{9}{2}x^{2} \Rightarrow \frac{du}{dx} + \frac{1}{x(\ln x)}u = 9x^{2} 
        \]
        积分因子为
        \[
        I(x)=\exp{\left(\int \frac{1}{x(\ln x)}\,dx\right)} = e^{\ln(\ln x)} = \ln x
        \]
        于是
        \[
        u = y^{-2} = \frac{1}{\ln x}\left(\int 9x^{2}\ln x\, dx + C\right) = \frac{1}{\ln x}\bigl(3x^{3}\ln x - x^{3} + C \bigr)
        \]
        即
        \[
        y^{2} = \frac{\ln x}{x^{3}(3\ln x - 1) + C}
        \]
    \end{solution}

    \question 判断函数 $I(x, y) = \cos(xy)$ 是否为微分方程
    \[
    [\tan(xy) + xy]\,dx + x^2\,dy = 0
    \]
    的积分因子。若是,求其通解。
    \begin{solution}
        将方程乘以 $I(x, y)$ 得
        \[
        [\sin(xy) + xy\cos(xy)]\,dx + [x^2\cos(xy)]\,dy = 0
        \]
        设
        \[
        P(x,y) = \sin(xy) + xy\cos(xy),\quad Q(x,y) = x^2\cos(xy).
        \]
        观察得偏导数
        \[
        \frac{\partial P}{\partial y} = 2x\cos(xy) - x^2 y \sin(xy) = \frac{\partial Q}{\partial x}
        \]
        由此可得 $\cos(xy)$ 是给定方程的积分因子。注意到
        \[
        d(x\sin(xy)) = (\sin(xy) + xy \cos(xy))\,dx + x^2 \cos(xy)\,dy
        \]
        因此通解为
        \[
        x\sin(xy) = C
        \]
    \end{solution}

    \question 求方程
    \[
    y\,dx - (2x + y^4)\,dy = 0
    \]
    的积分因子,并由此求通解。
    \begin{solution}
        设 $I(x, y)$ 是该积分因子,将方程乘以 $I(x, y)$ 得
        \[
        I y\,dx -(2x + y^4)I\,dy = 0
        \]
        记
        \[
        P(x,y) = Iy, \quad Q(x,y) = -(2x + y^4)I
        \]
        方程的正合条件为
        \[
        P_y = y I_y + I = Q_x = -2 I - (2x+y^4)I_x
        \]
        若 $I$ 只与 $y$有关,即 $I_x = 0$,则有
        \[
        y I_y + I = -2 I \Rightarrow y I_y = -3I \Rightarrow I(y) = \frac{1}{y^3}
        \]
        将其代入方程得到
        \[
        \frac{1}{y^2}dx - \frac{2x + y^4}{y^3}dy = 0
        \]
        此时方程正合。设函数 $\phi(x,y)$ 满足
        \[
        \frac{\partial \phi}{\partial x} = \frac{1}{y^2}, \quad 
        \frac{\partial \phi}{\partial y} = -\frac{2x + y^4}{y^3}
        \]
        由第一式积分得
        \[
        \phi(x,y) = \frac{x}{y^2} + h(y) \Rightarrow \frac{\partial \phi}{\partial y} = -\frac{2x}{y^3} + h'(y)
        \]
        由第二式可得
        \[
        h'(y) = -y \Rightarrow h(y) = -\frac{y^2}{2}
        \]
        因此通解为
        \[
        \frac{x}{y^2} - \frac{y^2}{2} = C \Rightarrow 2x - y^4 = C y^2
        \]
    \end{solution}

    \question 求解齐次微分方程
    \[
    y' - x^{-1}y = x^{-1}\sqrt{x^2 - y^2}, \quad x > 0
    \]
    \begin{solution}
        考虑齐次性,将方程改写为
        \[
        \frac{dy}{dx} = \frac{y}{x} + \sqrt{1 - \left(\frac{y}{x}\right)^2}
        \]
        作变量代换,
        \[
        y = xV(x) \Rightarrow \frac{dy}{dx} = x\frac{dV}{dx} + V
        \]
        代入方程得
        \[
        x\frac{dV}{dx} + V = V + \sqrt{1-V^2} \Rightarrow \frac{dV}{\sqrt{1-V^2}} = \frac{dx}{x}
        \]
        两边积分得
        \[
        \int \frac{dV}{\sqrt{1-V^2}} = \int \frac{dx}{x} + C \Rightarrow \sin^{-1} V = \ln|x| + C
        \]
        因此通解为
        \[
        y = xV = x \sin(C + \ln x), \quad x > 0
        \]
    \end{solution}

    \question 解微分方程
    \[
    (x^2 + y^2)\,dx + (x^2 - xy)\,dy = 0
    \]
    \begin{solution}
        若$x \neq 0$,有
        \[
        \left(1 + \frac{y^2}{x^2}\right) + \left(1 - \frac{y}{x}\right) \frac{dy}{dx} = 0
        \]
        设$\dfrac{y}{x} = u \Rightarrow y = xu$,则$\dfrac{dy}{dx} = x \dfrac{du}{dx} + u$,
        \[
        (1 + u^2) + (1 - u)\left(x \frac{du}{dx} + u\right) = 0 \Rightarrow x \,\frac{du}{dx}=\frac{u+1}{u-1}
        \]
        分离变量得
        \[
        \int \frac{u-1}{u+1}\, du = \int \frac{1}{x}\, dx \Rightarrow u - 2 \ln |u+1| = \ln x + C 
        \]
        即
        \[
        \frac{y}{x} - \ln \left(1 + \frac{y}{x}\right)^2 = \ln x + C
        \]
    \end{solution}

    \question 解
    \[
    x y'' + y' = 8x,\quad x>0
    \]
    \begin{solution}
        令 $v=y'$, 则 $v'=y''$。方程变为一阶方程
        \[
        x v' + v = 8x \Rightarrow v' + \frac{1}{x} v = 8
        \]
        这是线性一阶方程,积分因子为
        \[
        I(x)=\exp{\left(\int \frac{1}{x}\,dx\right)} =x
        \]
        所以
        \[
        v(x)=\frac{1}{x}\left(\int 8x\,dx + C_1\right)=\frac{1}{x}(4x^2+C_1)
        \]
        积分得到
        \[
        y=\int v(x)\,dx + C_2 = \int \frac{1}{x}(4x^2+C_1)\,dx + C_2=2x^2 + C_1\ln x + C_2
        \]
    \end{solution}

    \question 解微分方程
    \[
    y' + e^{y'} - x = 0
    \]
    \begin{solution}
        令 $u = y' \Rightarrow u + e^u = x$,由链导法,
        \[
        \frac{dy}{du}=\frac{dy}{dx}\cdot\frac{dx}{du}=u\frac{dx}{du}
        \]
        又因为
        \[
        x = u + e^u \Rightarrow \frac{dx}{du} = 1 + e^u
        \]
        所以
        \[
        \frac{dy}{du} = u(1 + e^u) \Rightarrow y = \int u(1 + e^u)\,du = \frac{u^2}{2} + (u - 1)e^u + C
        \]
        故解可写为参数形式:
        \[
        \begin{cases}
        x = u + e^u\\[4pt]
        y = \dfrac{u^2}{2} + (u - 1)e^u + C
        \end{cases}
        \]
    \end{solution}

    \question 解微分方程
    \[
    y' + \frac{y}{x} = e^{xy}
    \]
    \begin{solution}
        令$u = e^{xy}$,则
        \[
        \frac{du}{dx} = e^{xy}\left( y + x\frac{dy}{dx} \right)
        \Rightarrow y' + \frac{y}{x} = \frac{1}{xu}\frac{du}{dx}
        \]
        由原方程可得
        \[
        \frac{1}{xu}\frac{du}{dx}=u \Rightarrow \frac{1}{u^2}\frac{du}{dx} = x
        \]
        解为
        \[
        -\frac{1}{u} = \frac{x^2}{2} + C 
        \]
        故通解为
        \[
        \frac{x^2}{2} + e^{-xy} = C_1
        \]
    \end{solution}

    \question 解
    \[
    dx - xy(1+xy^2) \,dy =0
    \]
    \begin{solution}
        原方程即
        \[
        \frac{1}{x^2}\frac{dx}{dy} - \frac{y}{x} - y^3 = 0
        \]
        设$u=-\dfrac{1}{x}$,则$\dfrac{du}{dy}=\dfrac{1}{x^2}\dfrac{dx}{dy}$,化为线性一阶微分方程
        \[
        \frac{du}{dy}+yu=y^3
        \]
        积分因子为$\displaystyle \exp{\left(\int y\,dy\right)}=e^{\frac{y^2}{2}}$,于是
        \[
        u e^{\frac{y^2}{2}} = \int y^3 e^{\frac{y^2}{2}}\,dy + C
        \]
        作代换 $t=\dfrac{y^2}{2}$,则 $dt=y\,dy$,
        \[
        \int y^3 e^{\frac{y^2}{2}}\,dy = \int 2t e^{t}\,dt
        = 2te^t-2\int e^t \, dt = 2(t-1)e^{t} = (y^2-2)e^{\frac{y^2}{2}} 
        \]
        于是
        \[
        u = y^2 - 2 + C e^{-\frac{y^2}{2}}.
        \]
        代回 $u=-\dfrac{1}{x}$得到
        \[
        x = \frac{1}{\,2 - y^2 - C e^{-\frac{y^2}{2}}\,},
        \]
    \end{solution}

    \question 解微分方程
    \[
    xy' - y = 2x^2 y(y^2 - x^2)
    \]
    \begin{solution}
        令$u = x^2y$,对$x$求导,
        \[
        \frac{du}{dx} = 2xy + x^2\frac{dy}{dx} \Rightarrow
        \frac{dy}{dx} = \frac{1}{x^2}\frac{du}{dx} - \frac{2u}{x^3}.
        \]
        代回原方程得
        \[
        \frac{1}{x} \frac{du}{dx} - \frac{3u}{x^2} = 2u \left(\frac{u^2}{x^4} - x^2\right) \Rightarrow \frac{du}{dx} + \left(2x^3 - \frac{3}{x}\right) u = \frac{2}{x^3} u^3
        \]
        这是阶数为3的伯努利方程。令$v = u^{-2}$,则$\dfrac{dv}{dx} = -2u^{-3}\dfrac{du}{dx}$,于是
        \[
        v' + \left(-4x^3 + \frac{6}{x}\right)v = -\frac{4}{x^3}.
        \]
        积分因子为
        \[
        I(x) = \exp\left(\int \left(-4x^3 + \frac{6}{x}\right) dx\right) = e^{-x^4 + 6 \ln x} = x^6 e^{-x^4}
        \]
        解得
        \[
        u^{-2} = v = \frac{e^{x^4}}{x^6}\left(\int x^6 e^{-x^4} \cdot \left(-\frac{4}{x^3}\right)\,dx+C\right) =\frac{e^{x^4}}{x^6}\left(e^{-x^4}+C\right)= \frac{1}{x^6}(1 + Ce^{x^4})
        \]
        由$u = x^2y$,通解即
        \[
        x^2 - y^2 = cy^2 e^{x^4}
        \]
    \end{solution}

    \question 解 
    \[
    y y'' + (y')^2 = y y'
    \]
    \begin{solution}
        设 $u=y'$, 由链导法,
        \[
        y''=\frac{du}{dx}=\frac{du}{dy}\frac{dy}{dx}=u\frac{du}{dy}.
        \]
        代入原方程得
        \[
        y\cdot u\frac{du}{dy} + u^2 = y u \Rightarrow \frac{du}{dy} + \frac{1}{y}u = 1
        \]
        这是关于 $u$ 的线性一阶方程。积分因子为 $e^{\int \frac{1}{y}\,dy}=y$,于是
        \[
        \frac{d}{dy}(y u) = y \Rightarrow y u = \frac{y^2}{2} + c_1 \Rightarrow y'=u=\frac{y}{2} + \frac{c_1}{y} = \frac{y^2+2c_1}{2y}
        \]
        分离变量并积分,
        \[
        \frac{2y\,dy}{y^2+2c_1} = dx \Rightarrow
        \int \frac{2y\,dy}{y^2+2c_1} = \int dx + c_2.
        \]
        得到
        \[
        \ln|y^2+2c_1| = x + c_2 \Rightarrow y^2 = C_1 e^{x} + C_2
        \]
    \end{solution}

    \question 解微分方程
    \[
    \frac{dy}{dx} = x \,\frac{d^2y}{dx^2} - \left(\frac{dy}{dx}\right)^3
    \]
    \begin{solution}
        令$u=\dfrac{dy}{dx}$,则$\dfrac{d^2y}{dx^2}=\dfrac{du}{dx}$,整理为
        \[
        x\frac{du}{dx} = u + u^3
        \]
        分离变量得
        \[
        \frac{du}{u(1+u^2)} = \frac{dx}{x} \Rightarrow \int\left(\frac{1}{u}-\frac{u}{1+u^2}\right)\,du=\int\frac{dx}{x}
        \]
        积分得到
        \[
        \ln|u| -\frac{1}{2}\ln(1+u^2)=\ln|x| + \ln C \Rightarrow C|x|=\frac{|u|}{\sqrt{1+u^2}}
        \]
        化简得
        \[
        u^2 = \frac{C^2 x^2}{1 - C^2 x^2} \Rightarrow \frac{dy}{dx}= u = \pm \frac{C x}{\sqrt{1 - C^2 x^2}}
        \]
        再次分离变量得
        \[
        y= \pm \int \frac{C x}{\sqrt{1 - C^2 x^2}}\,dx 
        \]
        令 $t=Cx,dt=C\,dx$,积分变为
        \[
        y=\pm \frac{1}{C}\int \frac{t}{\sqrt{1-t^2}}\,dt 
        = -\frac{1}{C}\sqrt{1-t^2} + C_1 = -\frac{1}{C}\sqrt{1 - C^2 x^2} + C_1
        \]
    \end{solution}

    \question 求解
    \[
    (1+y^2)\frac{d^2y}{dx^2} + \left(\frac{dy}{dx}\right)^3 + \frac{dy}{dx} = 0
    \]
    \begin{solution}
        令$u=\dfrac{dy}{dx}$,由链导法,
        \[
        \frac{d^2y}{dx^2}=\frac{du}{dx}=u\frac{du}{dy}
        \]
        代入原方程得
        \[
        u\left((1+y^2)\frac{du}{dy} + u^2 + 1\right)=0
        \]
        若$u=\dfrac{dy}{dx}=0$,解为$y=C_1$。若$u \neq 0$,则由
        \[
        (1+y^2)\frac{du}{dy} + u^2 + 1 = 0
        \]
        分离变量得
        \[
        \frac{du}{u^2+1} = -\frac{dy}{1+y^2}
        \]
        两边积分得
        \[
        \arctan u = -\arctan y + A 
        \]
        令常数 $A=\arctan C$,则
        \[
        \frac{dy}{dx} = \tan(\arctan C - \arctan y)= \frac{C-y}{1+Cy}
        \]
        再次分离变量,
        \[
        \frac{1+Cy}{C-y}\,dy = dx
        \]
        积分得
        \[
        \int\frac{1+Cy}{C-y}\,dy=\int\left(-C + \frac{1+C^2}{C-y}\right)\,dy = -Cy - (1+C^2)\ln|C-y| = x + C_2
        \]
        综上,除 $y = C_1$ 外,通解为
        \[
        - Cy - (1+C^2)\ln|C-y| = x + C_2
        \]
    \end{solution}

    \question 求解
    \[
    \frac{d^2 y}{dx^2} = \frac{1}{1-y^2}\left(\frac{dy}{dx}\right)^2
    \]
    \begin{solution}
        设 $u=\dfrac{dy}{dx}$,由链导法,
        \[
        \frac{du}{dx}=\frac{du}{dy}\frac{dy}{dx}=u\frac{du}{dy}.
        \]
        代入原方程得
        \[
        \frac{du}{dy}=\frac{u}{1-y^2}.
        \]
        分离变量并积分,
        \[
        \frac{du}{u}=\frac{dy}{1-y^2} \Rightarrow
        \ln|u|=\frac12\ln\!\left|\frac{1+y}{1-y}\right|+C_1.
        \]
        于是存在常数 $C$ 使
        \[
        \frac{dy}{dx}=u=C\sqrt{\frac{1+y}{1-y}}
        \]
        再次分离变量得
        \[
        dx=\frac{1}{C}\sqrt{\frac{1-y}{1+y}}\,dy 
        \]
        对右侧积分,设$y=\cos2\theta$,则$dy=-2\sin 2\theta\,d\theta$,于是
        \begin{align*}
        \int\sqrt{\frac{1-y}{1+y}}\,dy
        &=\int\sqrt{\frac{2\sin^2\theta}{2\cos^2\theta}}(-2\sin 2\theta)\,d\theta\\
        &=2\int(\cos2\theta-1)\,d\theta\\
        &=2\left(\frac{1}{2}\sin 2\theta-\theta\right)\\
        &=\sqrt{1-y^2}-\arccos y 
        \end{align*}
        因此通解为
        \[
        x=\frac{1}{C_1}\left(\sqrt{1-y^2}-\arccos y\right)+C_2
        \]
    \end{solution}

    \question 求解
    \[
    \frac{1}{2}\left(\frac{dy}{dx}\right)^2 = 4y^2 + y\frac{d^2y}{dx^2}
    \]
    \begin{solution}
        令$u=\dfrac{dy}{dx}$,由链导法,
        \[
        \frac{d^2y}{dx^2}=\frac{du}{dx}=u\frac{du}{dy}
        \]
        代入原方程得
        \[
        y\,u\frac{du}{dy} = \frac{1}{2}u^2 - 4y^2
        \]
        若 $u=0$,则 $y=0$为一解。若 $u\neq0$即$y\neq0$,则
        \[
        \frac{dv}{dy} - \frac{1}{y}v = -8y
        \]
        这是关于 $v=u^2$ 的线性一阶方程,积分因子为 
        \[
        I(x)=\exp{\left(-\int \frac{1}{y}\,dy\right)}=\frac{1}{y}
        \]
        于是
        \[
        \left(\frac{dy}{dx}\right)^2=v=y\left(\int -8 \,dy+A\right)=A y - 8y^2
        \]
        为便于积分,令$B=\frac{A}{16}$,则配方得$A y - 8y^2 = 8\left(B^2 - (y-B)^2\right)$,因此
        \[
        \frac{dy}{dx} = \pm 2\sqrt{2}\,\sqrt{B^2-(y-B)^2}
        \]
        分离变量并积分:
        \[
        \frac{1}{2\sqrt{2}}\int \frac{d(y-B)}{\sqrt{B^2-(y-B)^2}} = \pm x + C \Rightarrow \frac{1}{2\sqrt{2}}\arcsin\left(\frac{y-B}{B}\right) = \pm x + C
        \]
        两边变换并记相位常数,得
        \[
        \arcsin\left(\frac{y-B}{B}\right) = 2\sqrt{2}\,x + C'
        \]
        于是
        \[
        y(x) = B\left(1 + \sin\left(2\sqrt{2}\,x + C'\right)\right)
        \]
    \end{solution}

    \question 求解
    \[
    \frac{y''}{y'} - \frac{y'}{y} = \ln y
    \]
    \begin{solution}
        设$u = \frac{dy}{dx}$,由链导法,
        \[
        y'' = \frac{du}{dx} = u \frac{du}{dy}
        \]
        原方程化为
        \[
        \frac{du}{dy} - \frac{1}{y} u = \ln y
        \]
        积分因子为
        \[
        I(y) = \exp\left(-\int \frac{1}{y} dy \right) = e^{-\ln y} = \frac{1}{y}.
        \]
        于是
        \[
        \frac{dy}{dx} = u = y \left(\int \frac{\ln y}{y} \,dy+C_1\right)=y\left( \frac{1}{2} (\ln y)^2 + C_1 \right)
        \]
        再次分离变量,
        \[
        \frac{2\cdot \frac{1}{y}\,dy}{ \left((\ln y)^2 + C_1 \right)} = dx
        \]
        两边积分后得,
        \[
        \frac{2}{\sqrt{A}} \arctan\frac{\ln y}{\sqrt{A}} = x + C_2
        \]
        设$C_1=\sqrt{A}$,可得通解为
        \[
        y = e^{\displaystyle 2C_1 \tan(C_1(x+C_2))}
        \]
    \end{solution}

    \question 求解
    \[
    y\frac{d^2y}{dx^2}+y^2=\frac12\left(\frac{dy}{dx}\right)^2
    \]
    \begin{solution}
        令$z=\sqrt{y}$,则
        \[
        y=z^2,\quad y'=2zz',\quad y''=2{z'}^2+2zz''
        \]
        代入原方程得
        \[
        z^2(2{z'}^2+2zz'')+z^4=\frac12(4z^2{z'}^2) \Rightarrow 2z''+z=0.
        \]
        这是常系数二阶线性方程,通解为
        \[
        z=C_1\cos\frac{x}{\sqrt2}+C_2\sin\frac{x}{\sqrt2}
        \]
        因此
        \[
        y=z^2=\left(C_1\cos\frac{x}{\sqrt2}+C_2\sin\frac{x}{\sqrt2}\right)^{2}
        \]
    \end{solution}

    \question 求解
    \[
    y''' = y'y''
    \]
    \begin{solution}
        令 $u=\dfrac{dy}{dx}$,由链导法,
        \[
        u'' = \frac{du'}{dx} = \frac{du'}{du}\frac{du}{dx} = \frac{du'}{du}\,u'
        \]
        代回方程得
        \[
        \frac{du'}{du}\,u' = u\,u' \Rightarrow \frac{du'}{du} = u
        \]
        积分得
        \[
        u' = \frac{1}{2}u^2 + 2C_1
        \]
        分离变量并积分,
        \[
        \int \frac{du}{\frac12 u^2 + C_1} = \int \frac{2\,du}{u^2 + 2C_1}= x + C_2 \Rightarrow \frac{2}{k}\arctan\!\frac{u}{k} = x + C_2
        \]
        其中$k^2 = 2C_1$,代回 $u=y'$得
        \[
        \arctan\frac{y'}{\sqrt{2C_1}} = \frac{\sqrt{2C_1}}{2}(x+C_2).
        \]
        再积分得到 $y$,
        \[
        y = \int y'\,dx = \int \sqrt{2C_1}\tan\left(\frac{\sqrt{2C_1}}{2}(x+C_2)\right)\,dx + C_3
        \]
        即
        \[
        y = -2\ln\left|\cos\left(C_1'(x+C_2)\right)\right| + C_3
        \]
    \end{solution}

    \question 求解
    \[
    \frac{y''}{(y')^2} = \frac{y}{y^2 - 1}
    \]
    \begin{solution}
        两边乘以 $y'$ 得
        \[
        \frac{y''}{y'} = \frac{y y'}{y^2 - 1}.
        \]
        因此对 $x$ 积分得
        \[
        \ln y' = \frac12 \ln(y^2-1) + \ln A \Rightarrow y' = A \sqrt{y^2 - 1}
        \]
        分离变量后积分得
        \[
        \frac{dy}{\sqrt{y^2-1}} = A dx \Rightarrow \ln(y + \sqrt{y^2-1}) = A x + C_1.
        \]
        即
        \[
        y + \sqrt{y^2-1} = B e^{A x}
        \]
        解出$y$,得通解为
        \[
        y = \frac{B e^{A x} + \frac{1}{B} e^{-A x}}{2}
        \]
    \end{solution}

    \question 解
    \[
    y'' - y'\tan x + 2y = 0
    \]
    \begin{solution}
        观察到$y_1=\sin x$为原方程的解,由降阶法,$y_2=u(x)y_1$,其中
        \begin{align*}
        u(x)&=\int \frac{1}{\sin^2 x} e^{-\int -\tan x\,dx}\, dx \\
        &=\int \frac{1}{\sin^2 x} e^{-\ln|\cos x|}\, dx \\
        &=\int (\cot^2 x+1)\sec x \, dx \\
        &=\int (\sec x+\cot x \csc x) \, dx \\
        &=\ln |\sec x + \tan x| -\csc x
        \end{align*}
        故通解为
        \[
        y=\sin x(C_1+\ln |\sec x + \tan x|)+C_2
        \]
    \end{solution}

    \question 解微分方程
    \[
    \cos^2x \, \frac{d^2y}{dx^2} = 2y
    \]
    \begin{solution}
        观察得$y_1 = \tan x$满足方程
        \[
        y'' = 2y \sec^2x 
        \]
        现设$y=u(x)\tan x$,则$y'=u'\tan x + u \sec^2 x$,且
        \[
        y''=u''\tan x + 2u'\sec^2 x + 2u\sec^2 x \tan x =u''\tan x + 2u'\sec^2 x + y''
        \]
        即
        \[
        u''\tan x + 2u'\sec^2 x =0
        \]
        设$u' = z$,则
        \[
        - z'\tan x = 2z\sec^2 x
        \]
        分离变量得
        \[
        \int - \frac{z'}{2z} \, dz = \int \frac{\sec x}{\tan x} \, dx \Rightarrow - \frac{1}{2} \ln|z| = \ln|\tan x| + \ln A 
        \]
        化简得
        \[
        \frac{dx}{du}=\frac{1}{z}=A^2\tan^2 x
        \]
        再分离变量,
        \[
        \int \frac{1}{A^2} \cot^2 x \, dx = \int du \Rightarrow \frac{1}{A^2} \int (\csc^2 x - 1) \, dx = u \Rightarrow u= C (-\cot x - x) + B
        \]
        故通解为
        \[
        y = B\tan x - Cx\tan x - C
        \]
    \end{solution}

    \question 解微分方程
    \[
    \frac{dy}{dx}=2\left(\frac{y+2}{x+y+1}\right)^2
    \]
    \begin{solution}
        设$u=y+2, v=x+y+1$,则
        \[
        \frac{du}{dx}=\frac{dy}{dx}=2\left(\frac{u}{v}\right)^2, \quad
        \frac{dv}{dx}=1+\frac{dy}{dx}=1+2\left(\frac{u}{v}\right)^2.
        \]
        于是有
        \[
        \frac{dv}{du}=\frac{1+2\left(\frac{u}{v}\right)^2}{2\left(\frac{u}{v}\right)^2}=\frac{v^2+2u^2}{2u^2}=1+\frac{v^2}{2u^2}.
        \]
        令$w=\dfrac{v}{u}$,则 $v=wu$,并且$\dfrac{dv}{du}=w+u\dfrac{dw}{du}$,代入上式得
        \[
        u\frac{dw}{du}=\frac{w^2}{2}-w+1
        \]
        分离变量再积分得
        \[
        \frac{du}{u}=\frac{dw}{\frac{1}{2}\left((w-1)^2+1\right)} \Rightarrow \ln|u| = 2\arctan(w-1)+C
        \]
        把 $u,w$ 代回原变量,因此得到
        \[
        \ln|y+2| = 2\arctan\left(\frac{x-1}{y+2}\right) + C
        \]
    \end{solution}

    \question 解
    \[
    f(x) + f'(-x) = x^2 + \alpha, \quad \alpha \in \mathbb{R}
    \]
    \begin{solution}
        由
        \[
        f(x) + f'(-x) = x^2 + \alpha \tag{1}
        \]
        令$x \to -x$,得$f(-x)+f'(x)== x^2 + \alpha$,两边求导,
        \[
        -f'(-x) + f''(x) = 2x \tag{2}
        \]
        由$(1)+(2)$,
        \[
        f''(x)+f(x)=x^2 + 2x + \alpha
        \]
        齐次解为
        \[
        f_c=C_1\sin x +C_2\cos x
        \]
        且发现特解为
        \[
        f_p=x^2 + 2x + \alpha-2
        \]
        于是通解为
        \[
        f(x)=C_1\sin x +C_2\cos x+x^2 + 2x + \alpha-2
        \]
        尚未结束,与原方程比较系数得
        \[
        f'(-x) + f(x) = (C_1 + C_2)(\cos x + \sin x) + \sqrt{x^2 + \alpha} \Rightarrow C_2=-C_1
        \]
        故通解为
        \[
        f(x)=C_1(\sin x -\cos x)+x^2 + 2x + \alpha-2
        \]
    \end{solution}

    \question 解
    \[
    y = xy' + \sqrt{(y')^2 + 1}
    \]
    \begin{solution}
        此方程属于克莱罗方程(Clairaut's equation)。对$x$求导得
        \[
        y' = y' + xy'' + \frac{1}{x} \frac{x \cdot 2y'y''}{\sqrt{(y')^2 + 1}}
        \]
        整理得
        \[
        y'' \left(x + \frac{y'}{\sqrt{(y')^2 + 1}}\right) = 0
        \]
        若$y''=0$,则$y'=C \Rightarrow y=Cx+\sqrt{C^2+1}$。对于
        \[
        \left(x + \frac{y'}{\sqrt{(y')^2 + 1}}\right) = 0
        \]
        化简得
        \[
        y'=\pm\frac{x}{\sqrt{1-x^2}}
        \]
        照常分离系数后积分得
        \[
        \int dy = \pm \int \frac{x}{\sqrt{1-x^2}} \, dx \Rightarrow y= \pm \sqrt{1-x^2}+A
        \]
        将$y=\pm\sqrt{1-x^2}+A$代入原方程解得$A=0$,故解为
        \[
        y=\pm\sqrt{1-x^2} \Rightarrow x^2+y^2=1
        \]
    \end{solution}

    \question 求解
    \[
    \sqrt{\tan x} \,\frac{dy}{dx}=x
    \]
    \begin{solution}
        分离系数得,
        \[
        \frac{1}{y} \,dy =\sqrt{\cot x}\,dx \Rightarrow \ln y = \int \sqrt{\cot x} \,dx
        \]
        设$\cot x=u^2$,则$-\csc^2 x \,dx=2u\,du \Rightarrow dx=-\dfrac{2u\,du}{1+u^4}$,
        \[
        I=\int \sqrt{\cot x} \,dx=-\int \frac{2u\,du}{1+u^4}=-\int \frac{2\,du}{u^2+\frac{1}{u^2}}= - \int \frac{1 + \frac{1}{u^2}}{u^2 + \frac{1}{u^2}}\, du - \int \frac{1 - \frac{1}{u^2}}{u^2 + \frac{1}{u^2}}\, du=-(I_1+I_2)
        \]
        其中设$u-\dfrac{1}{u}=t,\left(1+\dfrac{1}{u^2}\right)\,du=dt$,则$u^2 + \dfrac{1}{u^2} = t^2 + 2$,
        \[
        I_1=\int \frac{1 + \frac{1}{u^2}}{u^2 + \frac{1}{u^2}}\, du= \int \frac{dt}{t^2 + 2} = \frac{1}{\sqrt{2}} \arctan \frac{t}{\sqrt{2}}= \frac{1}{\sqrt{2}} \arctan \left(\frac{u^2 - 1}{u\sqrt{2}}\right)
        = \frac{1}{\sqrt{2}} \arctan \left(\frac{\cot x - 1}{\sqrt{2} \cot x}\right)
        \]
        且设$u+\dfrac{1}{u}=\phi,\left(1-\dfrac{1}{u^2}\right)\,du=d\phi$,则$u^2 + \dfrac{1}{u^2} = \phi^2 - 2$,
        \[
        I_2=\int \frac{1 - \frac{1}{u^2}}{u^2 + \frac{1}{u^2}}\, du
        I_2 = \int \frac{d\phi}{\phi^2 - 2}
        = \frac{1}{2\sqrt{2}} \int \left(\frac{1}{\phi - \sqrt{2}} - \frac{1}{\phi + \sqrt{2}}\right) d\phi
        = \frac{1}{2\sqrt{2}} \ln \left|\frac{\phi - \sqrt{2}}{\phi + \sqrt{2}}\right|
        \]
        \[
        = \frac{1}{2\sqrt{2}} \ln \left| \frac{u^2 - u\sqrt{2} + 1}{u^2 + u\sqrt{2} + 1} \right| = \frac{1}{2\sqrt{2}} \ln \left| \frac{\cot x - \sqrt{2}\cot x + 1}{\cot x + \sqrt{2}\cot x + 1} \right|
        \]
        于是$\ln y=I$给出
        \[
        y = Ce^{\displaystyle -\frac{1}{2\sqrt{2}} \arctan\left(\frac{\cot x - 1}{\sqrt{2} \cot x}\right)} \left(\frac{\cot x + \sqrt{2}\cot x + 1}{\cot x - \sqrt{2}\cot x + 1}\right)^{\frac{1}{2\sqrt{2}}}
        \]
    \end{solution}

    \question 解
    \[
    y'' - y' - 2y = 10 \sin x, \quad y(0)=0, \quad y'(0)=1.
    \]
    \begin{solution}
        对应齐次方程的特征方程为
        \[
        r^2 - r - 2 =(r-2)(r+1)=0 \Rightarrow r=2,-1
        \]
        因此齐次通解为
        \[
        y_c(x) = C_1 e^{2x} + C_2 e^{-x}.
        \]
        取特解
        \[
        y_p(x) = A \cos x + B \sin x
        \]
        代入方程
        \[
        y_p'' - y_p' - 2y_p = (B - 3A)\cos x + (-A - 3B)\sin x=0
        \]
        比较系数解得$A=-1, B=-3$,因此特解为
        \[
        y_p(x) = -\cos x - 3 \sin x
        \]
        通解为
        \[
        y(x) = y_c + y_p = C_1 e^{2x} + C_2 e^{-x} - \cos x - 3 \sin x.
        \]
        使用初值条件 $y(0)=0,y'(0)=1$,
        \[
        C_1 + C_2 - 1 = 0 \Rightarrow C_1 + C_2 = 1.
        \]
        \[
        y'(x) = 2C_1 e^{2x} - C_2 e^{-x} + \sin x - 3 \cos x \implies y'(0) = 2C_1 - C_2 - 3 = 1.
        \]
        解得$C_1 = 1, C_2 = 0$,因此初值问题的解为
        \[
        y(x) = e^{2x} - \cos x - 3 \sin x
        \]
    \end{solution}

    \question 设初值条件为 $y(0)=2, y'(0)=0$,求解方程
    \[
    y''+4y=16x\cos 2x
    \]
    \begin{solution}
        对应的齐次方程为$y''+4y=0$,其通解为
        \[
        y_c(x)=C_1 \cos 2x + C_2 \sin 2x
        \]
        非齐次项为 $16x\cos 2x$,原方程的特解包含了$A_0 \cos 2x + B_0 \sin 2x$,将与$y_c(x)$重叠。
        因此特解需取
        \[
        y_p(x)=x\bigl((A_0+A_1x)\cos 2x+(B_0+B_1x)\sin 2x\bigr).
        \]
        将 $y_p$ 代入原方程,比较系数可得$A_0=1,A_1=0,B_0=0,B_1=2$,因此特解为
        \[
        y_p(x)=x\cos 2x+2x^2\sin 2x.
        \]
        所以通解为
        \[
        y(x)=y_c+y_p=(C_1+x)\cos 2x+(C_2+2x^2)\sin 2x
        \]
        由 $y(0)=2$得$C_1=2$,求导得
        \[
        y'(x)=(2x - 2C_1) \sin2x + (4x^2 + 2C_2 + 1) \cos2x
        \]
        代入 $x=0,y'(0)=0$,得到$C_2=-\dfrac12$,因此满足初值条件的解为
        \[
        y(x)=(x+2)\cos 2x+\left(2x^2-\frac12\right)\sin 2x
        \]
    \end{solution}

    \question 设连续函数 $f(x)$ 满足
    \[
    f(x)=\frac14 x^{2}+\frac12\cos x-\int_{0}^{x}(x-t)f(t)\,dt,
    \]
    求 $f(x)$。  
    \begin{solution}
        即解微分方程
        \[
        y''(x) + y(x) = \frac{1 - \cos x}{2},\;y(0) = \frac{1}{2}, \; y'(0) = 0
        \]
        通解为
        \begin{align*}
        y = y_h + y_p = C_1\cos x + C_2\sin x + \frac{1}{2} - \frac{1}{4}x\sin x
        \end{align*}
        代入初值
        \[
        y(0) = C_1 + \frac{1}{2} = \frac{1}{2} \Rightarrow C_1 = 0 ;\;
        y'(0) = C_2 + 0 = 0 \Rightarrow C_2 = 0
        \]
        因此
        \[
        f(x) = \frac{1}{2} - \frac{1}{4}x\sin x
        \]
    \end{solution}

    \question 考虑二阶线性齐次方程 
    \[
    y'' + p(x)y' + q(x)y = 0,
    \]
    其中 $p(x),q(x)$ 在某区间 $I$ 上连续。
    \begin{parts}
    \part 证明:若存在常数 $r$ 使得对所有 $x\in I$ 有
    \[
    r^2 + r p(x) + q(x)=0,
    \]
    则 $y(x)=e^{rx}$ 为该方程的一个解。
    \begin{solution}
        直接代入检验,对于 $y=e^{rx}$,
        \[
        y' = r e^{rx},\quad y'' = r^2 e^{rx}.
        \]
        代入方程得
        \[
        y'' + p(x)y' + q(x)y = (r^2 + r p(x) + q(x))e^{rx} = 0,
        \]
        因此 $y=e^{rx}$ 确为方程的解。
    \end{solution}
    \part 求出当
    \[
    p(x)=-2\left(1+\frac{1}{x}\right),\quad q(x)=1+\frac{2}{x}
    \]
    时方程的通解。
    \begin{solution}
        先求常数根 $r$. 将形式 $r^2 + r p(x) + q(x)=0$ 代入:
        \[
        r^2 + r\left(-2\Bigl(1+\frac{1}{x}\Bigr)\right) + 1+\frac{2}{x} = 0 \Rightarrow (r^2-2r+1) + \frac{2}{x}(1-r)=0.
        \]
        该等式对所有 $x$ 成立,比较系数得
        \[
        r^2-2r+1=0,\quad 1-r=0.
        \]
        由第二式得 $r=1$,并满足第一式,于是其中一解为
        \[
        y_1(x)=e^{x}
        \]
        由降阶法公式,第二个线性独立解为
        \[
        y_2 = y_1 \int \frac{1}{y_1^2} \exp{\left(\int p(x)\,dx\right)}\, dx = e^x \int e^{-2x} \exp{\left(2\int (1+\frac{1}{x})\,dx\right)}= e^x \int x^2\, dx = \frac{1}{3}x^3 e^x
        \]
        因此通解为
        \[
        y(x) = (C_1 + C_2 x^3)e^{x}
        \]
    \end{solution}
    \end{parts}

    \question 以换元 
    \[
    x=\int e^{-\frac{t^2}{2}}\,dt,
    \]
    解微分方程
    \[
    y''(t)+t y'(t) + e^{-t^2} y(t)=0
    \]
    \begin{solution}
        令 $z(x)=y(t)$,由链导法,
        \[
        y'=\frac{dy}{dt}=\frac{dz}{dx}\frac{dx}{dt}=e^{-\frac{t^2}{2}}\frac{dz}{dx}
        \]
        且
        \begin{align*}
        y'' = \frac{d}{dt}\left(e^{-\frac{t^2}{2}}\frac{dz}{dx}\right) 
        &= -t e^{-\frac{t^2}{2}}\frac{dz}{dx} 
        + e^{-\frac{t^2}{2}}\frac{d}{dt}\left(\frac{dz}{dx}\right) \\
        &= -t e^{-\frac{t^2}{2}}\frac{dz}{dx}
        + e^{-\frac{t^2}{2}}\frac{d}{dx}\left(\frac{dz}{dx}\right)\frac{dx}{dt} \\
        &= -t e^{-\frac{t^2}{2}}\frac{dz}{dx}
        + e^{-\frac{t^2}{2}} \frac{d^2z}{dx^2} \, e^{-\frac{t^2}{2}} \\
        &= -t e^{-\frac{t^2}{2}}\frac{dz}{dx}
        + e^{-t^2}\frac{d^2z}{dx^2}.
        \end{align*}
        把 $y',y''$ 代入原方程得
        \[
        \left(-t e^{-\frac{t^2}{2}}\frac{dz}{dx} + e^{-t^2}\frac{d^2z}{dx^2}\right)
        + t\left(e^{-\frac{t^2}{2}}\frac{dz}{dx}\right) + e^{-t^2}z 
        = e^{-t^2}\frac{d^2z}{dx^2} + e^{-t^2}z = 0 \Rightarrow \frac{d^2z}{dx^2} + z = 0
        \]
        解为
        \[
        z(x)=C_1\cos x + C_2\sin x \Rightarrow y(t)=C_1\cos\left(\int e^{-\frac{t^2}{2}}\,dt\right)+C_2\sin\left(\int e^{-\frac{t^2}{2}}\,dt\right)
        \]
    \end{solution}

    \question 解方程
    \[
    t y''(t) + (t^2-1) y'(t) + t^3 y(t) = 0
    \]
    \begin{solution}
        欲求合适的$x=v(t)$,令$z(x)=y(t)$,则
        \[
        y'=\frac{dy}{dt}=\frac{dz}{dx}\frac{dx}{dt}=\frac{dz}{dx}\frac{dv}{dt},
        \]
        且
        \begin{align*}
        y'' = \frac{d}{dt}\left(\frac{dv}{dt}\frac{dz}{dx}\right) 
        &= \frac{d^2v}{dt^2}\frac{dz}{dx} + \frac{dv}{dt}\frac{d}{dt}\left(\frac{dz}{dx}\right) \\
        &= \frac{d^2v}{dt^2}\frac{dz}{dx} + \frac{dv}{dt}\frac{d}{dx}\left(\frac{dz}{dx}\right)\frac{dx}{dt} \\
        &= \frac{d^2v}{dt^2}\frac{dz}{dx} + \left(\frac{dv}{dt}\right)^2\frac{d^2z}{dx^2}
        \end{align*}
        代入原方程,得到
        \[
        t\left(\frac{dv}{dt}\right)^2 \frac{d^2z}{dx^2} + \Bigl[t\frac{d^2v}{dt^2} + (t^2-1)\frac{dv}{dt}\Bigr]\frac{dz}{dx} + t^3 z = 0
        \]
        为使方程为常系数方程, 不妨考虑变换
        \[
        t\left(\frac{dv}{dt}\right)^2 = t^3 \Rightarrow \frac{dv}{dt} = t \Rightarrow v = \frac{t^2}{2}
        \]
        此时原方程可化为
        \[
        \frac{d^2z}{dx^2} + \frac{dz}{dx} + z = 0
        \]
        其通解为
        \[
        z(x) = e^{-\frac{x}{2}}\left(C_1 \cos \frac{\sqrt{3}}{2}x + C_2 \sin \frac{\sqrt{3}}{2}x\right)
        \]
        代回 $x=\frac{t^2}{2}$,即
        \[
        y(t) = e^{-\frac{t^2}{4}}\left(C_1 \cos \frac{\sqrt{3} t^2}{4} + C_2 \sin \frac{\sqrt{3} t^2}{4}\right)
        \]
    \end{solution}

    \question 设 $y_1,y_2$ 为方程
    \[
    y'' + p(x)y' + q(x)y = 0,\quad x>0 \tag{1}
    \]
    的两个解,其中 $p(x),q(x)$ 在 $x>0$ 上连续。已知
    \[
    y_1 = \frac{\sin x}{\sqrt{x}}
    \]
    且$y_1,y_2$的朗斯基行列式(Wronskian)为
    \[
    W(x)=\frac{1}{x}
    \]
    \begin{parts}
    \part 证明阿贝尔恒等式(Abel's Identity):
    \[
    W(x)=C\exp\!\left(-\int p(x)\,dx\right)
    \]
    \begin{solution}
        由朗斯基行列式的定义,
        \[
        W(x)=\det\begin{pmatrix}
        y_1 & y_2 \\
        y_1' & y_2'
        \end{pmatrix}
        = y_1y_2' - y_1'y_2
        \]
        于是
        \begin{align*}
        W'(x) &= (y_1y_2' - y_1'y_2)' = (y_1'y_2' + y_1y_2'') - (y_1''y_2 + y_1'y_2') \\
        &= y_1y_2'' - y_1''y_2 = y_1(-py_2' - qy_2) - (-py_1' - qy_1)y_2 \\
        &= -py_1y_2' + py_1'y_2 = -p(y_1y_2' - y_1'y_2) = -p(x)W(x)
        \end{align*}
        因此
        \[
        W'(x)+p(x)W(x)=0
        \]
        由此可解得
        \[
        W(x)=C\exp\!\left(-\int p(x)\,dx\right)
        \]
        即得证。
    \end{solution}
    \part 求方程(1)中的函数 $p(x),q(x)$。
    \begin{solution}
        由阿贝尔恒等式,
        \[
        p(x) = -\frac{W'(x)}{W(x)}
        = -\frac{-\frac{1}{x^2}}{\frac{1}{x}}
        = \frac{1}{x}
        \]
        由于 $y_1$ 是原方程的解,
        \[
        y_1''+py_1'+qy_1=0 \Rightarrow q(x) = -\frac{y_1'' + p y_1'}{y_1}= 1 - \frac{1}{4x^2}
        \]
        其中
        \[
        y'_1 = \frac{\cos x}{\sqrt{x}}- \frac{\sin x}{2x^{\frac{3}{2}}}, \quad y''_1 = \frac{3 \sin x}{4x^{\frac{5}{2}}} - \frac{\sin x}{\sqrt{x}} - \frac{\cos x}{x^{\frac{3}{2}}}
        \]
        因此原微分方程为
        \[
        y'' + \frac{1}{x}y' + \left(1 - \frac{1}{4x^2}\right)y = 0,\quad x>0
        \]
    \end{solution}
    \part 求方程(1)的通解。
    \begin{solution}
        利用朗斯基行列式求$y_2$。已知
        \[
        W(x) = y_1y_2' - y_1'y_2 = \frac{1}{x} \Rightarrow y'_2 - \frac{y'_1}{y_1} y_2 = \frac{1}{xy_1}
        \]
        即为一阶线性方程
        \[
        y_2' + \left(\frac{1}{2x} - \cot x\right) y_2 = \frac{1}{x\sin x}.
        \]
        其积分因子为
        \[
        I(x) = \exp\!\left(\int\left(\frac{1}{2x}-\cot x\right)dx\right)
        = \exp(\ln\sqrt{x} - \ln|\sin x|)
        = \pm \frac{\sqrt{x}}{\sin x}
        \]
        取$I(x)=\dfrac{\sqrt{x}}{\sin x}$,由此
        \begin{align*}
        y_2(x)
        &= \frac{\sin x}{\sqrt{x}}\left(\int \frac{\sqrt{x}}{\sin x}\cdot\frac{1}{x\sin x}dx + C\right) \\
        &= \frac{\sin x}{\sqrt{x}}\left(\int \frac{1}{\sin^2 x}dx + C\right) \\
        &= \frac{\sin x}{\sqrt{x}}(-\cot x + C) \\
        &= -\frac{\cos x}{\sqrt{x}} + C y_1
        \end{align*}
        因此通解为
        \[
        y(x) = C_1\frac{\sin x}{\sqrt{x}} + C_2\frac{\cos x}{\sqrt{x}}.
        \]
    \end{solution}
    \end{parts}

    \question 
    \begin{parts}
    \part 证明法国数学家刘维尔 (Liouville) 的结论:若黎卡提 (Riccati) 方程
    \[
    y' = p(x)y^2 + q(x)y+r(x)
    \]
    已知一个特解 $\bar{y}(x)$, 令 $y = z+ \bar{y}$,则用此代换可把方程化为一个阶数 $n = 2$的伯努利(Bernoulli) 方程,从而可通过解一个一阶线性方程得到一般解。
    \begin{solution}
        设 $y=z+\bar{y}$,代入原方程得
        \[
        z' + \bar{y}' = p(z+\bar{y})^{2} + q(z+\bar{y}) + r = pz^{2} + 2pz\bar{y} + qz + (p\bar{y}^{2}+q\bar{y}+r)
        \]
        由于 $\bar{y}$ 是原方程的一个解,故 $\bar{y}'=p\bar{y}^{2}+q\bar{y}+r$,于是
        \[
        z' = pz^{2} + 2pz\bar{y} + qz \Rightarrow
        z' - (2p\bar{y}+q)z = pz^{2}
        \]
        这正是阶数 $n=2$ 的伯努利方程。由通解公式,
        \[
        z^{1-n} = e^{-\int(1-n)(2p\bar{y}+q)\,dx}\left\{\int (1-n)p(x)e^{\int(1-n)(2p\bar{y}+q)\,dx}\,dx + C\right\}.
        \]
        代入 $n=2$ 即得
        \[
        z^{-1} = e^{-\int(2p\bar{y}+q)\,dx}\left\{ -\int p(x)e^{\int(2p\bar{y}+q)\,dx}\,dx + C\right\}.
        \]
        由此得到 $z$,再由 $y=z+\bar{y}$ 得原方程的通解,证毕。
    \end{solution}
    \part 求下列黎卡提方程的通解:
    \begin{subparts}
    \subpart $y'e^{-x} + y^{2} - 2ye^{x} = 1 - e^{2x}$,已知特解 $\bar{y}(x)=e^{x}$。 
    \begin{solution}
        原方程可写为
        \[
        y' = -e^{x}y^{2} + 2e^{2x}y + e^{x}(1-2e^{2x})
        \]
        因此
        \[
        p(x)=-e^{x},\quad q(x)=2e^{2x},\quad r(x)=e^{x}(1-2e^{2x}),\quad \bar{y}=e^{x}
        \]
        且
        \[
        \int (2p\bar{y} + q)\, dx = \int (2(-e^{x})e^{x} + 2e^{2x})\, dx  = 0
        \]
        于是上式积分指数因子为常数,代入公式得
        \[
        z^{-1} = -\int (-e^{x})\,dx + C = e^{x} + C \Rightarrow z = \frac{1}{C+e^{x}}
        \]
        代回 $y=z+\bar{y}$,得到通解
        \[
        y = e^{x} + \frac{1}{C+e^{x}}
        \]
    \end{solution}
    \subpart $x^{2}(y' + y^{2}) = 2$,已知特解 $\bar{y}(x)=-\dfrac{1}{x}$。
    \begin{solution}
        将方程写成标准形,
        \[
        y' = -y^{2} + \frac{2}{x^{2}}
        \]
        令 $y=z+\bar{y}=z-\dfrac{1}{x}$,代入并化简可得关于 $z$ 的方程
        \[
        z' - \frac{2}{x}z = -z^{2},
        \]
        是阶数 $n=2$ 的伯努利方程,由公式计算得
        \[
        z^{-1} = e^{\int \frac{2}{x}\,dx}\left( \int e^{-\int \frac{2}{x}\,dx}\,dx + C\right)
        = \frac{1}{x^{2}}\left( \int x^{2}\,dx + C\right)
        = \frac{1}{x^{2}}\left(\frac{x^{3}}{3} + C\right)
        \]
        即
        \[
        z = \frac{3x^{2}}{x^{3} + 3C}
        \]
        代回 $y=z-\frac{1}{x}$,化为
        \[
        y = \frac{3x^{2}}{x^{3} + 3C} - \frac{1}{x}
        = \frac{2x^{3} - 3C}{x(3x^{3} + 3C)}.
        \]
        令常数 $C_1 = 3C$,
        \[
        xy = \frac{2x^{3} - C_1}{3x^{3} + C_1}
        \]
    \end{solution}
    \end{subparts}
    \end{parts}

    \question 求解
    \[
    x^3-\frac{dy}{dx}=\frac{y}{x}(y-2)
    \]
    \begin{solution}
        发现此为一黎卡提方程,观察得$y_1=x^2$是一特解。设$y = x^2 + \dfrac{1}{u}$,则$\dfrac{dy}{dx}=2x-\dfrac{1}{u^2}\dfrac{du}{dx}$,原方程给出
        \[
        x^3-\left(2x-\frac{1}{u^2}\frac{du}{dx}\right)=x(x^2-2)+\frac{2(x^2-1)}{xu}+\frac{1}{xu^2}
        \]
        再化简得一阶线性方程
        \[
        \frac{du}{dx} -\left(2x-\frac{2}{x}\right)u = \frac{1}{x}
        \]
        积分因子为
        \[
        I(x)=\exp\left(\int -\left(2x-\frac{2}{x}\right)\,dx\right)=\exp\left(-x^2+2\ln|x|\right)=x^2 e^{-x^2}
        \]
        所以
        \[
        u=\frac{e^{-x^2}}{x^2}\left(\int x e^{-x^2}+C_1\right)
        = \frac{C_1 e^{x^2}-\frac{1}{2}}{x^2}
        \]
        回代 $y=x^2+\dfrac{1}{u}$ 得
        \[
        y = x^2 + \frac{x^2}{C_1 e^{x^2}-\tfrac{1}{2}}
        = x^2\frac{C_1 e^{x^2}+\tfrac{1}{2}}{C_1 e^{x^2}-\tfrac{1}{2}}.
        \]
        令 $C_2=2C_1$,则等价地写成
        \[
        y = x^2\frac{C_2 e^{x^2}+1}{C_2 e^{x^2}-1}
        \]
    \end{solution}

    \question 考虑柯西-欧拉方程(Cauchy-Euler Equation)
    \[
    ax^2y'' + bxy' + cy = 0, \quad x>0,
    \]
    一般解法为作变换 $x=e^t$ 或 $t=\ln x$,则原方程变为常系数方程
    \[
    az''(t) + (b-a)z'(t) + cz(t) = 0,
    \]
    解得 $z(t)$ 后代回 $y(x)=z(\ln x)$。据此,求解
    \begin{parts}
    \part $x^2y'' - 3xy' + 4y = 0$
    \begin{solution}
        方程为 $x^2y'' - 3xy' + 4y = 0$,所以 $a=1, b=-3, c=4$,对应 $z$ 方程为
        \[
        z'' - 4z' + 4z = 0
        \]
        特征方程为
        \[
        r^2 - 4r + 4 = (r-2)^2 = 0 \Rightarrow r=2
        \]
        因此解得
        \[
        z(t) = C_1 e^{2t} + C_2 t e^{2t} \Rightarrow y(x) = z(\ln x) = C_1 x^2 + C_2 x^2 \ln x
        \]
    \end{solution}
    \part $x^2y'' - xy' - 35y = 0$
    \begin{solution}
        方程为 $x^2y'' - xy' - 35y = 0$,所以 $a=1, b=-1, c=-35$。对应 $z$ 方程
        \[
        z'' + (b-a)z' + cz = z'' + (-1-1)z' - 35 z = z'' - 2 z' - 35 z = 0
        \]
        其中特征方程为
        \[
        r^2 - 2r - 35 = (r-7)(r+5)=0 \Rightarrow r_1 = 7, r_2 = -5
        \]
        因此通解为
        \[
        z(t) = C_1 e^{7t} + C_2 e^{-5t} \Rightarrow y(x) = z(\ln x) = C_1 x^7 + C_2 x^{-5}
        \]
    \end{solution}
    \end{parts}

    \question 解非齐次欧拉方程
    \[
    x^3 y'''+x^2 y'' - x y' = 24 x \ln x, \quad x>0
    \]
    \begin{solution}
        令 $x = e^t \Rightarrow t = \ln x$,由链导法,
        \[
        \frac{dy}{dx} = \frac{dy}{dt} \frac{dt}{dx} = \frac{1}{x} \frac{dy}{dt}
        \]
        \[
        \frac{d^2y}{dx^2} 
        = \frac{d}{dx}\left(\frac{dy}{dt} \frac{dt}{dx}\right) 
        = \frac{d}{dx}\left(\frac{1}{x} \cdot \frac{dy}{dt}\right) 
        = -\frac{dy}{x^2 dt} + \frac{1}{x} \frac{d^2y}{dt^2} \frac{1}{x} 
        = \frac{1}{x^2}\left(\frac{d^2y}{dt^2} - \frac{dy}{dt}\right) 
        \]
        \begin{align*}
        \frac{d^3y}{dx^3} 
        = \frac{d}{dx}\left(\frac{1}{x^2}\left(\frac{d^2y}{dt^2} - \frac{dy}{dt}\right)\right)
        &= -\frac{2}{x^3}\left(\frac{d^2y}{dt^2} - \frac{dy}{dt}\right) + \frac{1}{x^2}\left(\frac{d^3y}{dt^3}\cdot \frac{1}{x} - \frac{d^2y}{dt^2}\cdot \frac{1}{x}\right)\\
        &= \frac{1}{x^3} \frac{d^3y}{dt^3} - \frac{3}{x^3} \frac{d^2y}{dt^2} + \frac{2}{x^3} \frac{dy}{dt}
        \end{align*}
        代入方程得到
        \[
        \frac{d^3y}{dt^3} - 3 \frac{d^2y}{dt^2} + 3 \frac{dy}{dt} - y = 24 e^t t
        \]
        齐次方程$y''' - 3y'' + 3y' - y = 0$的特征方程 
        \[
        r^3 - 3r^2 + 3r - 1 = 0 \Rightarrow r=1
        \]
        故齐次解为
        \[
        y_c=C_1 e^t + C_2 t e^t + C_3 t^2 e^t
        \]
        设特解为
        \[
        y_p = t^3 (A t + B) e^t = (A t^4 + B t^3) e^t
        \]
        代入方程解得
        \[
        24 A t + 6B = 24 t \Rightarrow A=1, B=0 
        \]
        因此通解为
        \[
        y = t^4 e^t + C_1 e^t + C_2 t e^t + C_3 t^2 e^t
        \]
        代回 $t = \ln x$,原方程的解为
        \[
        y = x (\ln x)^4 + C_1 x + C_2 x \ln x + C_3 x (\ln x)^2
        \]
    \end{solution}

    \question 以参数变换法解
    \[
    y'' + 6y' + 9y = \frac{2 e^{-4x}}{x^2+1}.
    \]
    \begin{solution}
        对应齐次方程通解为
        \[
        y_c(x) = C_1 e^{-3x} + C_2 x e^{-3x}
        \]
        设特解为
        \[
        y_p(x) = u_1(x) e^{-3x} + u_2(x) x e^{-3x}
        \]
        其中 $u_1, u_2$ 满足
        \[
        \begin{bmatrix}
        e^{-3x} & x e^{-3x} \\
        -3 e^{-3x} & e^{-3x} - 3x e^{-3x}
        \end{bmatrix}
        \begin{bmatrix} u_1' \\ u_2' \end{bmatrix}
        =
        \begin{bmatrix} 0 \\ \frac{2 e^{-4x}}{x^2 + 1} \end{bmatrix}
        \]
        朗斯基行列式为
        \[
        W(y_1, y_2) = \begin{vmatrix} e^{-3x} & x e^{-3x} \\ -3 e^{-3x} & e^{-3x} - 3x e^{-3x} \end{vmatrix} = e^{-6x}
        \]
        由克兰姆法则,
        \[
        u_1' = \frac{- x e^{-3x} \cdot \frac{2 e^{-4x}}{x^2 + 1}}{e^{-6x}} = -\frac{2x}{x^2 + 1} \Rightarrow u_1 = -\ln(1+x^2),
        \]
        \[
        u_2' = \frac{e^{-3x} \cdot \frac{2 e^{-4x}}{x^2 + 1}}{e^{-6x}} = \frac{2}{x^2+1} \Rightarrow u_2 = 2 \tan^{-1} x.
        \]
        因此特解为
        \[
        y_p(x) = e^{-3x} \left(-\ln(1+x^2) + 2 x \tan^{-1} x \right)
        \]
        通解为
        \[
        y(x) = e^{-3x} \left(C_1 + C_2 x + 2 x \tan^{-1} x - \ln(1+x^2) \right)
        \]
    \end{solution}

    \question 求解
    \[
    y'' - 2y' + y = 4 e^x x^3 \ln x, \quad x>0
    \]
    \begin{solution}
        对应齐次方程的通解为
        \[
        y_c(x) = C_1 e^x + C_2 x e^x,
        \]
        设特解为
        \[
        y_p(x) = u_1(x) e^x + u_2(x) x e^x,
        \]
        其中 $u_1, u_2$ 满足
        \[
        \begin{bmatrix} e^x & x e^x \\ e^x & (x+1)e^x \end{bmatrix} 
        \begin{bmatrix} u_1' \\ u_2' \end{bmatrix} = 
        \begin{bmatrix} 0 \\ 4 e^x x^3 \ln x \end{bmatrix}.
        \]
        两边同时除以 $e^x$ 得
        \[
        \begin{bmatrix} 1 & x \\ 1 & x+1 \end{bmatrix} 
        \begin{bmatrix} u_1' \\ u_2' \end{bmatrix} =
        \begin{bmatrix} 0 \\ 4 x^3 \ln x \end{bmatrix}.
        \]
        朗斯基行列式为$(x+1)-x=1$,由克兰姆法则,
        \[
        u'_1 = -x \cdot 4x^{-2} \ln x = -4x^{-2} \ln x, \quad u'_2 = 4x^{-3} \ln x
        \]
        于是
        \[
        u_1 = -4 \int x^{-2} \ln x \, dx 
        = \frac{4 \ln x}{x} - 4 \int \frac{1}{x^2} \, dx = \frac{4(\ln x + 1)}{x}
        \]
        \[
        u_2 = 4 \int x^{-3} \ln x \, dx = -\frac{2 \ln x}{x^2} +2 \int \frac{1}{x^3} \, dx = \frac{-2 \ln x - 1}{x^2} 
        \]
        因此特解为
        \[
        y_p(x) = e^x u_1 + x e^x u_2 = e^x \left(\frac{4 (\ln x + 1)}{x} + x \cdot \frac{-2\ln x -1}{x^2}\right) = e^x \frac{2 \ln x + 3}{x}
        \]
        所以方程通解为
        \[
        y(x) = e^x \left(C_1 + C_2 x + \frac{2 \ln x + 3}{x}\right)
        \]
    \end{solution}

    \question 求解方程
    \[
    y'' + 4y' + 4y = 15 e^{-2x} \ln x + 25 \cos x, \quad x>0
    \]
    \begin{solution}
        对应齐次方程的通解为
        \[
        y_c(x) = C_1 e^{-2x} + C_2 x e^{-2x}
        \]
        第一项非齐次$f_1(x) = 15 e^{-2x} \ln x$:设特解为
        \[
        y_{p_1} = u_1(x) e^{-2x} + u_2(x) x e^{-2x}
        \]
        其中 $u_1,u_2$ 满足
        \[
        \begin{bmatrix} e^{-2x} & x e^{-2x} \\ -2e^{-2x} & (1-2x)e^{-2x} \end{bmatrix} 
        \begin{bmatrix} u_1' \\ u_2' \end{bmatrix} =
        \begin{bmatrix} 0 \\ 15 e^{-2x} \ln x \end{bmatrix}
        \]
        两边除以 $e^{-2x}$ 得
        \[
        \begin{bmatrix} 1 & x \\ -2 & 1-2x \end{bmatrix} 
        \begin{bmatrix} u_1' \\ u_2' \end{bmatrix} =
        \begin{bmatrix} 0 \\ 15 \ln x \end{bmatrix}
        \]
        解得
        \[
        u_1' = -15 x \ln x \Rightarrow u_1 = -15 \int x \ln x \, dx = -\frac{15}{2} x^2 \ln x + 15 \int \frac{x}{2}\,dx = \frac{15}{4} x^2(-2\ln x+1)
        \]
        \[
        u_2' = 15 \ln x \Rightarrow u_2 = 15 \int \ln x\, dx = 15 x (\ln x - 1)
        \]
        因此
        \[
        y_{p_1}(x) = e^{-2x} \left[ \frac{15}{4} x^2(-2\ln x+1) + 15 x (\ln x -1) \right] = \frac{15}{4} x^2 (2\ln x -3) e^{-2x}
        \]
        第二项非齐次$f_2(x) = 25 \cos x$:设特解
        \[
        y_{p_2}(x) = A \cos x + B \sin x
        \]
        代入方程得
        \[
        (3A+4B) \cos x + (-4A+3B) \sin x = 25 \cos x
        \]
        比较系数解得$A=3, B=4$,所以
        \[
        y_{p_2}(x) = 3 \cos x + 4 \sin x
        \]
        综上,方程通解为
        \[
        y(x) = C_1 e^{-2x} + C_2 x e^{-2x} + \frac{15}{4} x^2 (2\ln x -3) e^{-2x} + 3 \cos x + 4 \sin x
        \]
    \end{solution}

    \question 解
    \[
    y'''-y''+y'-y=e^{-x}\sin x
    \]
    \begin{solution}
        齐次方程
        \[
        y'''-y''+y'-y=0
        \]
        的特征方程为
        \[
        (r-1)(r^2+1)=0 \Rightarrow r=1, \pm i
        \]
        所以齐次解为
        \[
        y_c(x) = C_1 e^x + C_2 \cos x + C_3 \sin x
        \]
        对于非齐次项 $e^{-x}\sin x$,尝试特解
        \[
        y_p(x) = e^{-x}(A\cos x + B \sin x)
        \]
        则
        \[
        y_p'= -e^{-x}((A+B) \sin x+(A-B)\cos x)
        \]
        \[
        y_p''= 2e^{-x}(A\sin x + B \cos x)
        \]
        \[
        y_p'''= 2e^{-x}((B-A) \sin x+(A+B)\cos x)
        \]
        代入原方程解得$A=\dfrac{1}{5},B=0$,于是
        \[
        y_p(x) = -\frac{1}{5} e^{-x} \cos x
        \]
        因此通解为
        \[
        y(x) = C_1 e^x + C_2 \cos x + C_3 \sin x - \frac{1}{5} e^{-x} \cos x
        \]
    \end{solution}

    \question 求解微分方程
    \[
    y'''+3y''+3y'+y=\frac{2e^{-x}}{1+x^2}
    \]
    \begin{solution}  
        齐次解为
        \[
        y_c(x) = C_1 e^{-x} + C_2 x e^{-x} + C_3 x^2 e^{-x}
        \]
        使用参数变换法,设
        \[
        y_p(x) = u_1(x) e^{-x} + u_2(x) x e^{-x} + u_3(x) x^2 e^{-x}
        \]
        其中 $u_1, u_2, u_3$ 满足
        \[
        \begin{cases}
        e^{-x} u_1' + (xe^{-x}) u_2' + (x^2 e^{-x}) u_3' = 0 \\
        (e^{-x})' u_1' + (xe^{-x})' u_2' + (x^2 e^{-x})' u_3' = 0 \\
        (e^{-x})'' u_1' + (xe^{-x})'' u_2' + (x^2 e^{-x})'' u_3' = \dfrac{2e^{-x}}{1+x^2}
        \end{cases}
        \]
        化简得
        \[
        \begin{cases}
        u_1' + x u_2' + x^2 u_3' = 0 \\
        -u_1' + (1-x) u_2' + (2x - x^2) u_3' = 0 \\
        u_1' + (-2+x) u_2' + (2 - 4x + x^2) u_3' = \dfrac{2}{1+x^2}
        \end{cases}
        \]
        其中增广矩阵为
        \[
        \left[
        \begin{array}{ccc|c}
        1 & x & x^2 & 0 \\
        -1 & 1-x & 2x-x^2 & 0 \\
        1 & -2+x & 2-4x+x^2 & \dfrac{2}{1+x^2}
        \end{array}
        \right] \to
        \left[
        \begin{array}{ccc|c}
        1 & x & x^2 & 0 \\
        0 & 1 & 2x & 0 \\
        0 & 0 & 1 & \dfrac{1}{1+x^2}
        \end{array}
        \right]
        \]
        逐步解得
        \begin{align*}
        u_3' &= \frac{1}{1+x^2} \Rightarrow u_3 = \tan^{-1} x \\
        u_2' &= -2x u_3' = -\frac{2x}{1+x^2} \Rightarrow u_2 = -\ln(1+x^2) \\
        u_1' &= -x u_2' - x^2 u_3' = \frac{x^2}{1+x^2} \Rightarrow u_1 = x - 2 \tan^{-1} x
        \end{align*}
        于是特解为
        \[
        y_p = e^{-x} \left[ (2x - 2\tan^{-1}x) + (-\ln(1+x^2)) x + (\tan^{-1} x) x^2 \right]
        \]
        整理得通解为
        \[
        y(x) = e^{-x}\left[C_1 + C_2 x + C_3 x^2+(x^2-2)\tan^{-1}x + x-x\ln(1+x^2) \right]
        \]
    \end{solution}

    \question 解方程
    \[
    y'''+y' = \sec x
    \]
    \begin{solution}
        齐次方程$y'''+y'=0$的特征方程为
        \[
        r^3 + r = 0 \Rightarrow r = 0, \pm i
        \]
        所以齐次解为
        \[
        y_c = c_1 + c_2 \cos x + c_3 \sin x
        \]
        使用参数变易法,设特解为
        \[
        y_p = u_1 y_1 + u_2 y_2 + u_3 y_3
        \]
        其中 \(u_1, u_2, u_3\) 满足
        \[
        \begin{bmatrix}
        y_1 & y_2 & y_3 \\
        y_1' & y_2' & y_3' \\
        y_1'' & y_2'' & y_3''
        \end{bmatrix}
        \begin{bmatrix}
        u_1' \\ u_2' \\ u_3'
        \end{bmatrix}
        =
        \begin{bmatrix}
        1 & \cos x & \sin x \\
        0 & -\sin x & \cos x \\
        0 & -\cos x & -\sin x
        \end{bmatrix}
        \begin{bmatrix}
        u_1' \\ u_2' \\ u_3'
        \end{bmatrix}
        =
        \begin{bmatrix}
        0 \\ 0 \\ \sec x
        \end{bmatrix}
        \]
        其中增广矩阵为
        \[
        \left[
        \begin{array}{ccc|c}
        1 & \cos x & \sin x & 0 \\
        0 & -\sin x & \cos x & 0 \\
        0 & -\cos x & -\sin x & \sec x
        \end{array}
        \right]
        \implies
        \left[
        \begin{array}{ccc|c}
        1 & \cos x & \sin x & 0 \\
        0 & -\sin x & \cos x & 0 \\
        1 & 0 & 0 & \sec x
        \end{array}
        \right]
        \]
        于是
        \[
        \begin{cases}
        u'_1 = \sec x \\[2pt]
        u'_2 = \dfrac{\cos x}{-\sin x} u'_3 = -1 \\[6pt]
        u'_3 = \dfrac{\sin x}{\cos x}
        \end{cases}
        \Rightarrow
        \begin{cases}
        u_1 = |\sec x + \tan x| \\[1pt]
        u_2 = -x \\[1pt]
        u_3 = \ln |\cos x|
        \end{cases}
        \]
        因此通解为
        \[
        y = |\sec x + \tan x| - x\cos x + \ln |\cos x| \sin x + C_1 + C_2 \cos x + C_3 \sin x
        \]
    \end{solution}

    \question 解联立微分方程式
    \[
    \begin{cases}
    x_1' = x_1 + 2x_2 \\ 
    x_2' = 2x_1 - 2x_2 \\
    \end{cases}
    \]
    \begin{solution}
        \[
        x_1' = x_1 + 2x_2 \tag{1}
        \]
        \[
        x_2' = 2x_1 - 2x_2 \tag{2}
        \]
        对(1)求导,且由(2),
        \[
        x_1'' = x_1' + 2 x_2' = x_1' + 2(2x_1 - 2x_2) = x_1' + 4x_1 - 4x_2 = x_1' + 4x_1 - 2(x_1' - x_1) = -x_1' + 6x_1
        \]
        得关于$x_1$的二阶常微分方程,因此 $x_1$ 的通解为
        \[
        x_1(t) = C_1 e^{-3t} + C_2 e^{2t}
        \]
        由(1),
        \[
        x_2 = \frac{1}{2}(x_1' - x_1) = -2 C_1 e^{-3t} + \frac{C_2}{2} e^{2t}
        \]
        最终将系统的通解写成向量形式:
        \[
        \boldsymbol{x}(t) = 
        \begin{pmatrix} x_1 \\ x_2 \end{pmatrix}
        = C_1 \begin{pmatrix} 1 \\ -2 \end{pmatrix} e^{-3t} + C_2 \begin{pmatrix} 1 \\ \frac{1}{2} \end{pmatrix} e^{2t}
        \]
    \end{solution}

    \question 求解 
    \[
    \boldsymbol{x}' = \begin{pmatrix} 0 & 1 & 1 \\ 1 & 0 & 1 \\ 1 & 1 & 0 \end{pmatrix} \boldsymbol{x}
    \]
    \begin{solution}
        设系数矩阵为 $\boldsymbol{A}$,首先求 $\boldsymbol{A}$ 的特征值,通过解
        \[
        \det(\boldsymbol{A} - \lambda \boldsymbol{I}_3) =  \begin{vmatrix} -\lambda & 1 & 1 \\ 1 & -\lambda & 1 \\ 1 & 1 & -\lambda \end{vmatrix} 
        = -(\lambda + 1)^2(\lambda - 2)
        \]
        因此特征值为
        \[
        \lambda_1 = -1 \;(m_1=2), \quad \lambda_2 = 2
        \]
        接着求特征向量。对 $\lambda_1 = -1$,解
        \[
        (\boldsymbol{A} - \lambda_1 \boldsymbol{I}_3) \boldsymbol{v} = \boldsymbol{0} \implies 
        \begin{pmatrix} 1 & 1 & 1 \\ 1 & 1 & 1 \\ 1 & 1 & 1 \end{pmatrix} \begin{pmatrix} v_1 \\ v_2 \\ v_3 \end{pmatrix} = \begin{pmatrix} 0 \\ 0 \\ 0 \end{pmatrix}
        \]
        得到方程
        \[
        v_1 + v_2 + v_3 = 0
        \]
        有两个自由变量,设 $v_2 = r, v_3 = s$,则 $v_1 = -r - s$,因此特征空间由两个线性无关向量张成:
        \[
        \boldsymbol{v}_1 = \begin{pmatrix}-1\\1\\0\end{pmatrix}, \quad \boldsymbol{v}_2 = \begin{pmatrix}-1\\0\\1\end{pmatrix}
        \]
        对应解为
        \[
        \boldsymbol{x}_1(t) = e^{-t} \begin{pmatrix}-1\\1\\0\end{pmatrix}, \quad 
        \boldsymbol{x}_2(t) = e^{-t} \begin{pmatrix}-1\\0\\1\end{pmatrix}
        \]
        对 $\lambda_2 = 2$,解
        \[
        (\boldsymbol{A} - \lambda_2 \boldsymbol{I}_3) \boldsymbol{v} = \boldsymbol{0} \implies 
        \begin{pmatrix}-2 & 1 & 1 \\ 1 & -2 & 1 \\ 1 & 1 & -2\end{pmatrix} \begin{pmatrix}v_1\\v_2\\v_3\end{pmatrix} = \begin{pmatrix}0\\0\\0\end{pmatrix}
        \]
        化简矩阵得到行阶梯形:
        \[
        \begin{pmatrix} -2 & 1 & 1 \\ 1 & -2 & 1 \\ 1 & 1 & -2 \end{pmatrix}
        \xrightarrow{R_2 \to \frac{1}{2}R_1 + R_2}
        \begin{pmatrix} -2 & 1 & 1 \\ 0 & -\frac{3}{2} & \frac{3}{2} \\ 0 & \frac{3}{2} & -\frac{3}{2} \end{pmatrix}
        \xrightarrow[R_3 \to \frac{2}{3}R_3]{R_3 \to R_2 + R_3}
        \begin{pmatrix}-2 & 1 & 1 \\ 0 & -1 & 1 \\ 0 & 0 & 0\end{pmatrix}
        \]
        秩为 2,设 $v_3 = r$,则 $v_1 = v_2 = r$,得到特征向量
        \[
        \boldsymbol{v}_3 = \begin{pmatrix}1\\1\\1\end{pmatrix}
        \]
        对应解为
        \[
        \boldsymbol{x}_3(t) = e^{2t} \begin{pmatrix}1\\1\\1\end{pmatrix}
        \]
        因此系统的一般解为
        \[
        \boldsymbol{x}(t) = c_1 e^{-t} \begin{pmatrix}-1\\1\\0\end{pmatrix} 
        + c_2 e^{-t} \begin{pmatrix}-1\\0\\1\end{pmatrix} 
        + c_3 e^{2t} \begin{pmatrix}1\\1\\1\end{pmatrix}
        \]
    \end{solution}

    \question 求解
    \[
    \boldsymbol{x}' = \begin{pmatrix} -3 & 0 & 2 \\ 1 & -1 & 0 \\ -2 & -1 & 0 \end{pmatrix} \boldsymbol{x}
    \]
    \begin{solution}
        设系数矩阵为 $\boldsymbol{A}$。首先求特征值,通过解
        \[
        \det(\boldsymbol{A} - \lambda \boldsymbol{I}_3) 
        = \begin{vmatrix} -3-\lambda & 0 & 2 \\ 1 & -1-\lambda & 0 \\ -2 & -1 & -\lambda \end{vmatrix}
        = -(\lambda+2)(\lambda^2 + 2\lambda + 3)
        \]
        得特征值
        \[
        \lambda_{1,2} = -1 \pm \sqrt{2}i, \quad \lambda_3 = -2
        \]
        接着求对应特征向量。对 $\lambda_1 = -1+\sqrt{2}i$, 解
        \[
        (\boldsymbol{A} - \lambda_1 \boldsymbol{I}_3)\boldsymbol{v} = \boldsymbol{0} \implies 
        \begin{pmatrix} -2-\sqrt{2}i & 0 & 2 \\ 1 & -\sqrt{2}i & 0 \\ -2 & -1 & 1-\sqrt{2}i \end{pmatrix} \begin{pmatrix} v_1 \\ v_2 \\ v_3 \end{pmatrix} = \begin{pmatrix}0\\0\\0\end{pmatrix}
        \]
        化为行阶梯形矩阵
        \[
        \begin{pmatrix} -2-\sqrt{2}i & 0 & 2 \\ 1 & -\sqrt{2}i & 0 \\ -2 & -1 & 1-\sqrt{2}i \end{pmatrix} \xrightarrow{R_1 \to (-2+\sqrt{2}i)R_1/6} \begin{pmatrix} 1 & 0 & \frac{1}{3}(-2+\sqrt{2}i) \\ 1 & -\sqrt{2}i & 0 \\ -2 & -1 & 1-\sqrt{2}i \end{pmatrix}
        \]
        \[
        \xrightarrow[R_3 \to R_3 + 2R_1]{R_1 \to R_2 - R_1} \begin{pmatrix} 1 & 0 & \frac{1}{3}(-2+\sqrt{2}i) \\ 0 & -\sqrt{2}i & \frac{1}{3}(2-\sqrt{2}i) \\ 0 & -1 & -\frac{1}{3}(1+\sqrt{2}i) \end{pmatrix} \xrightarrow[R_3 \to -R_3]{R_2 \to R_2/(\sqrt{2}i)} \begin{pmatrix} 1 & 0 & \frac{1}{3}(-2+\sqrt{2}i) \\ 0 & 1 & \frac{1}{3}(1+\sqrt{2}i) \\ 0 & 1 & \frac{1}{3}(1+\sqrt{2}i) \end{pmatrix}
        \]
        \[
        \xrightarrow{R_3 \to R_3 - R_2} \begin{pmatrix} 1 & 0 & \frac{1}{3}(-2+\sqrt{2}i) \\ 0 & 1 & \frac{1}{3}(1+\sqrt{2}i) \\ 0 & 0 & 0 \end{pmatrix}
        \]
        令 $v_3 = r$ 为自由变量,则
        \[
        \boldsymbol{v} = r \begin{pmatrix}-\frac{1}{3}(-2+\sqrt{2}i) \\ -\frac{1}{3}(1+\sqrt{2}i) \\ 1 \end{pmatrix}.
        \]
        取 $r=-3$ 并利用欧拉公式分离实部和虚部
        \begin{align*}
        &e^{(-1+\sqrt{2}i)t} \begin{pmatrix} -2+\sqrt{2}i \\ 1+\sqrt{2}i \\ -3 \end{pmatrix} = e^{-t}(\cos(\sqrt{2}t) + i\sin(\sqrt{2}t)) \begin{pmatrix} -2+\sqrt{2}i \\ 1+\sqrt{2}i \\ -3 \end{pmatrix} \\
        &= e^{-t} \begin{pmatrix} -2\cos(\sqrt{2}t) - \sqrt{2}\sin(\sqrt{2}t) + i(\sqrt{2}\cos(\sqrt{2}t) - 2\sin(\sqrt{2}t)) \\ \cos(\sqrt{2}t) - \sqrt{2}\sin(\sqrt{2}t) + i(\sqrt{2}\cos(\sqrt{2}t) + \sin(\sqrt{2}t)) \\ -3\cos(\sqrt{2}t) - 3i\sin(\sqrt{2}t) \end{pmatrix} \\
        &= e^{-t} \begin{pmatrix} -2\cos(\sqrt{2}t) - \sqrt{2}\sin(\sqrt{2}t) \\ \cos(\sqrt{2}t) - \sqrt{2}\sin(\sqrt{2}t) \\ -3\cos(\sqrt{2}t) \end{pmatrix} + ie^{-t} \begin{pmatrix} \sqrt{2}\cos(\sqrt{2}t) - 2\sin(\sqrt{2}t) \\ \sqrt{2}\cos(\sqrt{2}t) + \sin(\sqrt{2}t) \\ -3\sin(\sqrt{2}t) \end{pmatrix}
        \end{align*}
        得到两个线性无关解
        \[
        \boldsymbol{x}_1(t) = e^{-t} \begin{pmatrix}-2\cos(\sqrt{2}t)-\sqrt{2}\sin(\sqrt{2}t) \\ \cos(\sqrt{2}t)-\sqrt{2}\sin(\sqrt{2}t) \\ -3\cos(\sqrt{2}t) \end{pmatrix}, 
        \quad
        \boldsymbol{x}_2(t) = e^{-t} \begin{pmatrix} \sqrt{2}\cos(\sqrt{2}t)-2\sin(\sqrt{2}t) \\ \sqrt{2}\cos(\sqrt{2}t)+\sin(\sqrt{2}t) \\ -3\sin(\sqrt{2}t) \end{pmatrix}
        \]
        对 $\lambda_3 = -2$,解
        \[
        (\boldsymbol{A} - \lambda_3 \boldsymbol{I}_3)\boldsymbol{v} = \boldsymbol{0} \implies 
        \begin{pmatrix}-1 & 0 & 2 \\ 1 & 1 & 0 \\ -2 & -1 & 2\end{pmatrix} \begin{pmatrix}v_1\\v_2\\v_3\end{pmatrix} = \begin{pmatrix}0\\0\\0\end{pmatrix}.
        \]
        化简矩阵得到
        \[
        \begin{pmatrix}-1 & 0 & 2 \\ 1 & 1 & 0 \\ -2 & -1 & 2\end{pmatrix} 
        \xrightarrow[R_3 \to R_3 - 2R_1]{R_2 \to R_2 +R_1}
        \begin{pmatrix}-1 & 0 & 2 \\ 0 & 1 & 2 \\ 0 & -1 & -2\end{pmatrix}
        \xrightarrow{R_3 \to R_3 + R_2}
        \begin{pmatrix}-1 & 0 & 2 \\ 0 & 1 & 2 \\ 0 & 0 & 0 \end{pmatrix}
        \]
        令 $v_3 = s$,得到特征向量
        \[
        \boldsymbol{v}_3 = s \begin{pmatrix}2\\-2\\1\end{pmatrix}.
        \]
        对应解为
        \[
        \boldsymbol{x}_3(t) = e^{-2t} \begin{pmatrix}2\\-2\\1\end{pmatrix}.
        \]
        因此系统的一般解为
        \begin{align*}
        \boldsymbol{x}(t) &= c_1e^{-t} \begin{pmatrix} -2\cos(\sqrt{2}t) - \sqrt{2}\sin(\sqrt{2}t) \\ \cos(\sqrt{2}t) - \sqrt{2}\sin(\sqrt{2}t) \\ -3\cos(\sqrt{2}t) \end{pmatrix} + c_2e^{-t} \begin{pmatrix} \sqrt{2}\cos(\sqrt{2}t) - 2\sin(\sqrt{2}t) \\ \sqrt{2}\cos(\sqrt{2}t) + \sin(\sqrt{2}t) \\ -3\sin(\sqrt{2}t) \end{pmatrix} \\
        &\quad + c_3e^{-2t} \begin{pmatrix} 2 \\ -2 \\ 1 \end{pmatrix}
        \end{align*}
    \end{solution}

    \question 判断矩阵 
    \[
    \boldsymbol{A} = \begin{pmatrix} 1 & 1 & 0 \\ 0 & 1 & 1 \\ 0 & 0 & 4 \end{pmatrix}
    \]
    是否缺陷,并求系统 $\boldsymbol{x}'(t) = \boldsymbol{A} \boldsymbol{x}(t)$ 的通解。
    \begin{solution}
        首先求 $\boldsymbol{A}$ 的特征值:
        \[
        \det(\boldsymbol{A} - \lambda \boldsymbol{I}_3) = \begin{vmatrix} 1-\lambda & 1 & 0 \\ 0 & 1-\lambda & 1 \\ 0 & 0 & 4-\lambda \end{vmatrix} = (1-\lambda)^2 (4-\lambda)
        \]
        因此特征值为 $\lambda_1 = 4,\lambda_2 = 1 \;(m_2=2)$。对于 $\lambda_1 = 4$, 解
        \[
        (\boldsymbol{A} - 4\boldsymbol{I}_3) \boldsymbol{v} = \boldsymbol{0} \implies 
        \begin{pmatrix}-3 & 1 & 0 \\ 0 & -3 & 1 \\ 0 & 0 & 0 \end{pmatrix} \begin{pmatrix} v_1 \\ v_2 \\ v_3 \end{pmatrix} = \begin{pmatrix}0\\0\\0\end{pmatrix}
        \]
        得到特征向量及其对应解
        \[
        \boldsymbol{v} = r \begin{pmatrix} 1 \\ 3 \\ 9 \end{pmatrix}, \quad r \neq 0, \quad \boldsymbol{x}_1(t) = e^{4t} \begin{pmatrix} 1 \\ 3 \\ 9 \end{pmatrix}
        \]
        对于 $\lambda_2 = 1$, 解
        \[
        (\boldsymbol{A} - \boldsymbol{I}_3) \boldsymbol{v} = \boldsymbol{0} \implies 
        \begin{pmatrix} 0 & 1 & 0 \\ 0 & 0 & 1 \\ 0 & 0 & 3 \end{pmatrix} \begin{pmatrix} v_1 \\ v_2 \\ v_3 \end{pmatrix} = \begin{pmatrix}0\\0\\0\end{pmatrix}
        \]
        得到特征向量及其对应解
        \[
        \boldsymbol{v} = s \begin{pmatrix} 1 \\ 0 \\ 0 \end{pmatrix}, \quad s \neq 0, \quad \boldsymbol{x}_2(t) = e^t \begin{pmatrix} 1 \\ 0 \\ 0 \end{pmatrix}
        \]
        总共有两个线性无关特征向量,因此矩阵 $\boldsymbol{A}$ 是缺陷的。为求第三个线性无关解,设
        \[
        \boldsymbol{x}_3(t) = (\boldsymbol{v} t + \boldsymbol{w}) e^t
        \]
        其中 $\boldsymbol{v}$ 和 $\boldsymbol{w}$ 满足
        \[
        (\boldsymbol{A} - \boldsymbol{I}_3) \boldsymbol{v} = \boldsymbol{0}, \quad (\boldsymbol{A} - \boldsymbol{I}_3) \boldsymbol{w} = \boldsymbol{v} \implies (\boldsymbol{A} - \boldsymbol{I}_3)^2 \boldsymbol{w} = \boldsymbol{0}
        \]
        解得
        \[
        (\boldsymbol{A} - \boldsymbol{I}_3)^2 \boldsymbol{w} = \begin{pmatrix} 0 & 0 & 1 \\ 0 & 0 & 3 \\ 0 & 0 & 9 \end{pmatrix} \begin{pmatrix} w_1 \\ w_2 \\ w_3 \end{pmatrix} = \begin{pmatrix}0\\0\\0\end{pmatrix} \implies w_3 = 0
        \]
        自由选择 $w_1, w_2$,取
        \[
        \boldsymbol{w} = \begin{pmatrix} 0 \\ 1 \\ 0 \end{pmatrix}, \quad \boldsymbol{v} = (\boldsymbol{A} - \boldsymbol{I}_3) \boldsymbol{w} = \begin{pmatrix} 1 \\ 0 \\ 0 \end{pmatrix}
        \]
        因此得到第三个解
        \[
        \boldsymbol{x}_3(t) = (\boldsymbol{v} t + \boldsymbol{w}) e^t = e^t \begin{pmatrix} t \\ 1 \\ 0 \end{pmatrix}
        \]
        最终通解为
        \[
        \boldsymbol{x}(t) = c_1 e^{4t} \begin{pmatrix} 1 \\ 3 \\ 9 \end{pmatrix} + c_2 e^t \begin{pmatrix} 1 \\ 0 \\ 0 \end{pmatrix} + c_3 e^t \begin{pmatrix} t \\ 1 \\ 0 \end{pmatrix}
        \]
    \end{solution}

    \question 求解非齐次系统
    \[
    \boldsymbol{x}'(t) = \boldsymbol{A} \boldsymbol{x}(t) + \boldsymbol{F}(t),
    \]
    其中
    \[
    \boldsymbol{A} = \begin{pmatrix} 0 & 1 \\ -9 & 6 \end{pmatrix}, \quad \boldsymbol{F}(t) = \begin{pmatrix} 0 \\ t \end{pmatrix}.
    \]
    \begin{solution}
        首先求 $\boldsymbol{A}$ 的特征值和特征向量:
        \[
        \det(\boldsymbol{A} - \lambda \boldsymbol{I}_2) = \begin{vmatrix} -\lambda & 1 \\ -9 & 6-\lambda \end{vmatrix} = -\lambda(6-\lambda)+9 = (\lambda-3)^2
        \]
        因此特征值为 $\lambda = 3\;(m=2)$,对 $\lambda = 3$, 解
        \[
        (\boldsymbol{A}-3\boldsymbol{I}_2) \boldsymbol{v} = \begin{pmatrix}-3 & 1 \\ -9 & 3\end{pmatrix} \begin{pmatrix} v_1 \\ v_2 \end{pmatrix} = \begin{pmatrix}0 \\ 0\end{pmatrix}
        \]
        得到特征向量
        \[
        \boldsymbol{v} = r \begin{pmatrix} 1 \\ 3 \end{pmatrix}, \quad r \neq 0
        \]
        对应解为
        \[
        \boldsymbol{x}_1(t) = e^{3t} \begin{pmatrix} 1 \\ 3 \end{pmatrix}
        \]
        矩阵 $\boldsymbol{A}$ 是缺陷的,第二个线性无关解取形如
        \[
        \boldsymbol{x}_2(t) = (\boldsymbol{v} t + \boldsymbol{w}) e^{3t}
        \]
        其中 $\boldsymbol{v}$ 和 $\boldsymbol{w}$ 满足
        \[
        (\boldsymbol{A}-\lambda\boldsymbol{I}_2)\boldsymbol{v} = \boldsymbol{0}, \quad (\boldsymbol{A}-\lambda\boldsymbol{I}_2)\boldsymbol{w} = \boldsymbol{v} \implies (\boldsymbol{A} - \lambda \boldsymbol{I}_2)^2\boldsymbol{w} = \boldsymbol{0}
        \]
        解得
        \[
        \begin{pmatrix} -3 & 1 \\ -9 & 3 \end{pmatrix}\begin{pmatrix} -3 & 1 \\ -9 & 3 \end{pmatrix} \begin{pmatrix} w_1 \\ w_2 \end{pmatrix} = \begin{pmatrix} 0 \\ 0 \end{pmatrix} \\
        \implies \begin{pmatrix} 0 & 0 \\ 0 & 0 \end{pmatrix} \begin{pmatrix} w_1 \\ w_2 \end{pmatrix} = \begin{pmatrix} 0 \\ 0 \end{pmatrix}.
        \]
        取
        \[
        \boldsymbol{w} = \begin{pmatrix} 0 \\ 1 \end{pmatrix} \implies \boldsymbol{v} = (\boldsymbol{A}-3\boldsymbol{I}_2)\boldsymbol{w} = \begin{pmatrix} 1 \\ 3 \end{pmatrix}
        \]
        于是第二个解为
        \[
        \boldsymbol{x}_2(t) = (\boldsymbol{v} t + \boldsymbol{w}) e^{3t} = \begin{pmatrix} t \\ 3t+1 \end{pmatrix} e^{3t}
        \]
        因此齐次系统通解为
        \[
        \boldsymbol{x}_c(t) = c_1 e^{3t} \begin{pmatrix} 1 \\ 3 \end{pmatrix} + c_2 e^{3t} \begin{pmatrix} t \\ 3t+1 \end{pmatrix} = \begin{pmatrix} e^{3t} & t e^{3t} \\ 3 e^{3t} & (3t+1)e^{3t} \end{pmatrix} \begin{pmatrix} c_1 \\ c_2 \end{pmatrix}
        \]
        记基础矩阵为 $\boldsymbol{X}$,寻找特解
        \[
        \boldsymbol{x}_p(t) = \boldsymbol{X} \boldsymbol{u} \implies \boldsymbol{X} \boldsymbol{u}' = \boldsymbol{F} \implies \boldsymbol{u}' = \boldsymbol{X}^{-1} \boldsymbol{F}
        \]
        计算得到
        \[
        \boldsymbol{X}^{-1}\boldsymbol{F} = e^{-3t} \begin{pmatrix} 3t+1 & -t \\ -3 & 1 \end{pmatrix} \begin{pmatrix} 0 \\ t \end{pmatrix} = e^{-3t} \begin{pmatrix} -t^2 \\ t \end{pmatrix}
        \]
        \[
        \boldsymbol{u} = \int \boldsymbol{X}^{-1} \boldsymbol{F}\, dt = \begin{pmatrix} \int -t^2 e^{-3t} dt \\ \int t e^{-3t} dt \end{pmatrix} = e^{-3t} \begin{pmatrix} \frac{1}{3} t^2 + \frac{2}{9} t + \frac{2}{27} \\ -\frac{1}{3} t - \frac{1}{9} \end{pmatrix}
        \]
        于是得到特解
        \[
        \boldsymbol{x}_p(t) = \boldsymbol{X}\boldsymbol{u} = \begin{pmatrix} 1 & t \\ 3 & 3t+1 \end{pmatrix} \begin{pmatrix} \frac{1}{3} t^2 + \frac{2}{9} t + \frac{2}{27} \\ -\frac{1}{3} t - \frac{1}{9} \end{pmatrix} = \begin{pmatrix} \frac{1}{9} t + \frac{2}{27} \\ \frac{1}{9} \end{pmatrix}.
        \]
        最终非齐次系统通解为
        \[
        \boldsymbol{x}(t) = \boldsymbol{x}_p(t) + \boldsymbol{x}_c(t) = \begin{pmatrix} \frac{1}{9} t + \frac{2}{27} \\ \frac{1}{9} \end{pmatrix} + c_1 e^{3t} \begin{pmatrix} 1 \\ 3 \end{pmatrix} + c_2 e^{3t} \begin{pmatrix} t \\ 3t+1 \end{pmatrix}.
        \]
    \end{solution}



\end{questions}


\pagecolor{PageColor}
\begin{center}
  {\fontsize{30pt}{26pt}\selectfont 参考出处
  }
\end{center}
\separator
\begin{center}
\worldflag{MY}
\worldflag{TW}
\worldflag{CN}
\worldflag{JP}
\worldflag{IN}
\worldflag{US}
\worldflag{CA}
\worldflag{BG}
\end{center}
\begin{multicols}{2}
\begin{itemize}
    \item \href{https://artofproblemsolving.com/community/c13_contests}{Arts of Problem Solving (AoPS): Contest Collections}
    \item \href{https://brilliant.org/}{Brilliant (Kudos to the unmonetized version in the past)}
    \item \href{https://www.imc-math.org.uk/}{International Mathematics Competition (IMC) for University Students}
    \item \href{http://sections.maa.org/missouri/contest_archives.html}{Missouri Collegiate Mathematics Competition}
    \item \href{https://artofproblemsolving.com/community/c3414_amc_10}{American Mathematics Competition 10}
    \item \href{https://cemc.uwaterloo.ca/resources/past-contests?grade=All&academic_year=All&contest_category=24}{University of Waterloo CEMC - Euclid, Fermat, Cayley, Hypatia, Galois}
    \item \href{https://www.lehigh.edu/~dmd1/oldcontests.html}{Lehigh University High School Math Contest}
    \item \href{https://jeeadv.ac.in/archive.html}{Joint Entrance Examination (Advanced)}
    \item \href{https://www.isinj.com/mt-usamo/USSR%20Olympiad%20Problem%20Book%20(The)%20-%20Shklasrsky,%20Chentzov,%20and%20Yaglom%20(1993,%20Dover)%20(1-1).pdf}{The USSR Olympiad Problem Book: Selected Problems and Theorems of Elementary Mathematics}
    \item \href{https://www.youtube.com/@LetsSolveMathProblems}{LetsSolveMathProblems - Weekly Math Challenges}
    \item \href{https://www.youtube.com/@maths_505}{Maths 505 - Differential Equations}
    \item \href{https://www.youtube.com/@MichaelPennMath}{Michael Penn}
    \item \href{https://www.youtube.com/@PrimeNewtons}{Prime Newtons}
    \item \href{https://www.cpshs.hcc.edu.tw/resource/openfid.php?id=2000}{TRML}
    \item \href{https://madasmaths.com/archive_maths_booklets_standard_topics_integration.html}{MadasMaths}
    \item \href{https://www.densu.jp/libs.htm}{日本留学試験-理系数学}
    \item \href{https://math.fudan.edu.cn/gdsx/34379/list.htm}{中国复旦大学往年试题}
    \item \href{https://www.newera.edu.my/competition/hua_lo_geng/ch/passyear.php}{雪隆森中学华罗庚杯数学比赛}
    \item \href{https://mathcompetition.wixsite.com/chenjingrun}{厦门大学马来西亚分校-陈景润杯中学数学竞赛}
    \item \href{https://calcgospel.top/2023/%e5%85%a9%e7%a8%ae%e9%87%8d%e8%a6%81%e6%a5%b5%e9%99%90%e6%b7%b7%e5%90%88%e7%9a%84%e6%a5%b5%e9%99%90%e9%a1%8c/}{微积分福音}
    \item \href{https://ccjou.wordpress.com/category/%e7%b7%9a%e6%80%a7%e4%bb%a3%e6%95%b8%e5%b0%88%e6%ac%84/linear-equations/}{线代启示录}
    \item \href{http://cfc.ehosting.com.tw/LT_C/course_3_4.html}{福气老师}    
    \item \href{https://www.ceec.edu.tw/xmfile?xsmsid=0J052427633128416650&page=14}{指考历届试题}
    \item \href{http://library.yuda-cloudstudy.com.tw/submenu/index/old-exam}{统测历届试题与解答}
    \item \href{https://www.doubtnut.com/class-12/maths/determinants}{印度人的作业}
    \item \href{https://www.gaokzx.com/search?keywords=%E6%95%B0%E5%AD%A6%E8%AF%95%E9%A2%98%E5%8F%8A%E7%AD%94%E6%A1%88}{北京高考在线}
    \item \href{https://chu246.blogspot.com/search/label/%E6%95%99%E7%94%84}{朱式幸福教甄}
    \item \href{https://cantor.math.ntnu.edu.tw/workshop/110hsm/index.php?menu=Exam}{普通型高级中学数学科能力竞赛(决赛)}
    \item \href{https://www.zizzs.com/gk/jingsai/162865.html}{2014-2024全国中学生数学竞赛联赛试题及答案汇总}
    \item \href{http://www.360doc.com/userhome.aspx?userid=31600819&cid=8}{如此醉的图书馆}
    \item \href{https://www.jyeoo.com/shijuan/gzsx/101/}{菁优网}
    \item \href{https://shijuan.zww.cn/gaozhong/shuxue/gaokao/1/1707.htm}{08高考文科试题分类圆锥曲线}
    \item \href{https://shijuan.zww.cn/gaozhong/shuxue/gaokao/1/1709.htm}{08高考文科试题分类数列}
    \item \textcolor{brown}{福伦-隆中高数笔记}
    \item \textcolor{brown}{曾龙文师-高二数理培训队笔记}
\end{itemize}
\end{multicols}

\end{document}
